% Auto-generated from data/segments.json + layout.json (+ transforms) — do not edit by hand.
% volume: 01Vin01
% mode: printing
% segments: 2548
% transform_rules: 28
% toc_marks: 488

% Volume layout defaults (layout.json ``layout``).
\csromanlayoutapply{1}{1.5}{21.6pt}{8pt}{8pt}{65pt}{2.5em}%

\pitaka{\textbf{วินยปิฏก}}
\csromanmatikaanchor{mtk.0}
\csromantocmarknopage{chapter}{ปาราชิกปาฬิ}{mtk.0}
\bookmark[dest={mtk.0},level=0]{ปาราชิกปาฬิ}
\gambhira{\textbf{ปาราชิกปาฬิ}}
\namakkaram{นโม ตสฺส ภควโต อรหโต สมฺมา\-สมฺพุทฺธสฺส.}
\csromanmatikaanchor{mtk.1}
\csromantocmarknopage{chapter}{เวรญฺชกณฺฑ}{mtk.1}
\bookmark[dest={mtk.1},level=0]{เวรญฺชกณฺฑ}
\chapterhead{\textbf{เวรญฺชกณฺฑ}}
\csromanmatikaanchor{mtk.2}
\csromantocmarkref{section}{ภควโต ปริภวกถา}{mtk.2}
\bookmark[dest={mtk.2},level=1]{ภควโต ปริภวกถา}
\proseitem{1}{\csromanfolio{1}เตน สมเยน พุทฺโธ ภควา เวรญฺชายํ วิหรติ นเฬรุ\-ปุจิมนฺท\-มูเล มหตา ภิกฺขุ\-สํเฆน สทฺธิํ ปญฺจมตฺเตหิ ภิกฺขุสเตหิ. อสฺโสสิ โข เวรญฺโช พฺราหฺมโณ “สมโณ ขลุ โภ โคตโม สกฺยปุตฺโต สกฺยกุลา ปพฺพชิโต เวรญฺชายํ วิหรติ นเฬรุ\-ปุจิมนฺท\-มูเล มหตา ภิกฺขุ\-สํเฆน สทฺธิํ ปญฺจมตฺเตหิ ภิกฺขุสเตหิ, ตํ โข ปน ภวนฺตํ โคตมํ เอวํ กลฺยาโณ กิตฺติสทฺโท อพฺภุคฺคโต ‘อิติปิ โส ภควา อรหํ สมฺมาสมฺพุทฺโธ วิชฺชา\-จรณ\-สมฺปนฺโน สุคโต โลกวิทู อนุตฺตโร ปุริส\-ทมฺม\-สารถิ สตฺถา เทวมนุสฺสานํ พุทฺโธ ภควา’\csromansharedfootnote{665a3d55e62cc309}{ภควาติ (สฺยา), ที 1. 46, 109, 125 ปิฏฺเฐสุ อพฺภุคฺคตากาเรน ปน สเมติ.}, โส อิมํ โลกํ สเทวกํ สมารกํ สพฺรหฺมกํ สสฺสมณ\-พฺราหฺมณิํ ปชํ สเทวมนุสฺสํ สยํ อภิญฺญา สจฺฉิกตฺวา ปเวเทติ, โส ธมฺมํ เทเสติ อาทิกลฺยาณํ มชฺเฌกลฺยาณํ ปริโยสาน\-กลฺยาณํ สาตฺถํ สพฺยญฺชนํ เกวล\-ปริปุณฺณํ ปริสุทฺธํ พฺรหฺมจริยํ ปกาเสติ, สาธุ โข ปน ตถารูปานํ อรหตํ ทสฺสนํ โหตี”ติ.}

\proseitem{2}{\csromansymbolfootnote{*}{อิโต ปรํ ยาว 7 ปิฏฺเฐ ปทกฺขิณํ กตฺวา ปกฺกามีติ ปาโฐ อํ 3. 18 ปิฏฺฐาทีสุปิ อาคโต.}อถ โข เวรญฺโช พฺราหฺมโณ เยน ภควา เตนุปสงฺกมิ, อุปสงฺกมิตฺวา ภควตา สทฺธิํ สมฺโมทิ, สมฺโมทนียํ กถํ สารณียํ วีติสาเรตฺวา เอกมนฺตํ นิสีทิ, เอกมนฺตํ นิสินฺโน โข เวรญฺโช พฺราหฺมโณ ภควนฺตํ เอตทโวจ “สุตํ เมตํ โภ โคตม ‘น สมโณ \csromanfolio{2}โคตโม พฺราหฺมเณ ชิณฺเณ วุฑฺเฒ มหลฺลเก อทฺธคเต วโยอนุปฺปตฺเต อภิวาเทติ วา ปจฺจุฏฺเฐติ วา อาสเนน วา นิมนฺเตตี’ติ, ตยิทํ โภ โคตม ตเถว, น หิ ภวํ โคตโม พฺราหฺมเณ ชิณฺเณ วุฑฺเฒ มหลฺลเก อทฺธคเต วโยอนุปฺปตฺเต อภิวาเทติ วา ปจฺจุฏฺเฐติ วา อาสเนน วา นิมนฺเตติ. ตยิทํ โภ โคตม น สมฺปนฺนเมวา”ติ.}

\prose{นาหํ ตํ พฺราหฺมณ ปสฺสามิ สเทวเก โลเก สมารเก สพฺรหฺมเก สสฺสมณ\-พฺราหฺมณิยา ปชาย สเทว\-มนุสฺสาย ยมหํ อภิวาเทยฺยํ วา ปจฺจุฏฺเฐยฺยํ วา อาสเนน วา นิมนฺเตยฺยํ. ยํ หิ พฺราหฺมณ ตถาคโต อภิวาเทยฺย วา ปจฺจุฏฺเฐยฺย วา อาสเนน วา นิมนฺเตยฺย, มุทฺธาปิ ตสฺส วิปเตยฺยาติ.}

\proseitem{3}{อรสรูโป ภวํ โคตโมติ. อตฺถิ เขฺวส พฺราหฺมณ ปริยาโย, เยน มํ ปริยาเยน สมฺมา วทมาโน วเทยฺย “อรสรูโป สมโณ โคตโม”ติ. เย เต พฺราหฺมณ รูปรสา สทฺทรสา คนฺธรสา รสรสา โผฏฺฐพฺพรสา, เต ตถาคตสฺส ปหีนา อุจฺฉินฺนมูลา ตาลาวตฺถุกตา อนภาวํกตา\csromansharedfootnote{1754982fc472adab}{อนภาวกตา (สี), อนภาวํคตา (สฺยา)} อายติํ อนุปฺปาทธมฺมา. อยํ โข พฺราหฺมณ ปริยาโย, เยน มํ ปริยาเยน สมฺมา วทมาโน วเทยฺย “อรสรูโป สมโณ โคตโม”ติ, โน จ โข ยํ ตฺวํ สนฺธาย วเทสีติ.}

\proseitem{4}{นิพฺโภโค ภวํ โคตโมติ. อตฺถิ เขฺวส พฺราหฺมณ ปริยาโย, เยน มํ ปริยาเยน สมฺมา วทมาโน วเทยฺย “นิพฺโภโค สมโณ โคตโม”ติ. เย เต พฺราหฺมณ รูปโภคา สทฺทโภคา คนฺธโภคา รสโภคา โผฏฺฐพฺพ\-โภคา, เต ตถาคตสฺส ปหีนา อุจฺฉินฺนมูลา ตาลาวตฺถุกตา อนภาวํกตา อายติํ อนุปฺปาทธมฺมา. อยํ โข พฺราหฺมณ ปริยาโย, เยน มํ ปริยาเยน สมฺมา วทมาโน วเทยฺย “นิพฺโภโค สมโณ โคตโม”ติ, โน จ โข ยํ ตฺวํ สนฺธาย วเทสีติ.}

\proseitem{5}{อกิริยวาโท ภวํ โคตโมติ. อตฺถิ เขฺวส พฺราหฺมณ ปริยาโย, เยน มํ ปริยาเยน สมฺมา วทมาโน วเทยฺย “อกิริยวาโท สมโณ โคตโม”ติ. อหํ หิ พฺราหฺมณ อกิริยํ วทามิ \csromanfolio{3}กาย\-ทุจฺจริตสฺส วจีทุจฺจริตสฺส มโน\-ทุจฺจริตสฺส, อเนกวิหิตานํ ปาปกานํ อกุสลานํ ธมฺมานํ อกิริยํ วทามิ. อยํ โข พฺราหฺมณ ปริยาโย, เยน มํ ปริยาเยน สมฺมา วทมาโน วเทยฺย “อกิริยวาโท สมโณ โคตโม”ติ, โน จ โข ยํ ตฺวํ สนฺธาย วเทสีติ.}

\proseitem{6}{อุจฺเฉทวาโท ภวํ โคตโมติ. อตฺถิ เขฺวส พฺราหฺมณ ปริยาโย, เยน มํ ปริยาเยน สมฺมา วทมาโน วเทยฺย “อุจฺเฉทวาโท สมโณ โคตโม”ติ. อหํ หิ พฺราหฺมณ อุจฺเฉทํ วทามิ ราคสฺส โทสสฺส โมหสฺส, อเนกวิหิตานํ ปาปกานํ อกุสลานํ ธมฺมานํ อุจฺเฉทํ วทามิ. อยํ โข พฺราหฺมณ ปริยาโย, เยน มํ ปริยาเยน สมฺมา วทมาโน วเทยฺย “อุจฺเฉทวาโท สมโณ โคตโม”ติ, โน จ โข ยํ ตฺวํ สนฺธาย วเทสีติ.}

\proseitem{7}{เชคุจฺฉี ภวํ โคตโมติ. อตฺถิ เขฺวส พฺราหฺมณ ปริยาโย, เยน มํ ปริยาเยน สมฺมา วทมาโน วเทยฺย “เชคุจฺฉี สมโณ โคตโม”ติ. อหํ หิ พฺราหฺมณ ชิคุจฺฉามิ กาย\-ทุจฺจริเตน วจีทุจฺจริเตน มโน\-ทุจฺจริเตน, อเนกวิหิตานํ ปาปกานํ อกุสลานํ ธมฺมานํ สมาปตฺติยา ชิคุจฺฉามิ. อยํ โข พฺราหฺมณ ปริยาโย, เยน มํ ปริยาเยน สมฺมา วทมาโน วเทยฺย “เชคุจฺฉี สมโณ โคตโม”ติ, โน จ โข ยํ ตฺวํ สนฺธาย วเทสีติ.}

\proseitem{8}{เวนยิโก ภวํ โคตโมติ. อตฺถิ เขฺวส พฺราหฺมณ ปริยาโย, เยน มํ ปริยาเยน สมฺมา วทมาโน วเทยฺย “เวนยิโก สมโณ โคตโม”ติ. อหํ หิ พฺราหฺมณ วินยาย ธมฺมํ เทเสมิ ราคสฺส โทสสฺส โมหสฺส, อเนกวิหิตานํ ปาปกานํ อกุสลานํ ธมฺมานํ วินยาย ธมฺมํ เทเสมิ. อยํ โข พฺราหฺมณ ปริยาโย, เยน มํ ปริยาเยน สมฺมา วทมาโน วเทยฺย “เวนยิโก สมโณ โคตโม”ติ, โน จ โข ยํ ตฺวํ สนฺธาย วเทสีติ.}

\csromanmatikaanchor{mtk.3}
\csromantocmarkref{section}{กุกฺกุฏจฺฉาปกูปมากถา}{mtk.3}
\bookmark[dest={mtk.3},level=1]{กุกฺกุฏจฺฉาปกูปมากถา}
\proseitem{9}{ตปสฺสี ภวํ โคตโมติ. อตฺถิ เขฺวส พฺราหฺมณ ปริยาโย, เยน มํ ปริยาเยน สมฺมา วทมาโน วเทยฺย “ตปสฺสี สมโณ \csromanfolio{4}โคตโม”ติ. ตปนียาหํ พฺราหฺมณ ปาปเก อกุสเล ธมฺเม วทามิ กายทุจฺจริตํ วจีทุจฺจริตํ มโนทุจฺจริตํ. ยสฺส โข พฺราหฺมณ ตปนียา ปาปกา อกุสลา ธมฺมา ปหีนา อุจฺฉินฺนมูลา ตาลาวตฺถุกตา อนภาวํกตา อายติํ อนุปฺปาทธมฺมา, ตมหํ ตปสฺสีติ วทามิ. ตถาคตสฺส โข พฺราหฺมณ ตปนียา ปาปกา อกุสลา ธมฺมา ปหีนา อุจฺฉินฺนมูลา ตาลาวตฺถุกตา อนภาวํกตา อายติํ อนุปฺปาทธมฺมา. อยํ โข พฺราหฺมณ ปริยาโย, เยน มํ ปริยาเยน สมฺมา วทมาโน วเทยฺย “ตปสฺสี สมโณ โคตโม”ติ, โน จ โข ยํ ตฺวํ สนฺธาย วเทสีติ.}

\proseitem{10}{อปคพฺโภ ภวํ โคตโมติ. อตฺถิ เขฺวส พฺราหฺมณ ปริยาโย, เยน มํ ปริยาเยน สมฺมา วทมาโน วเทยฺย “อปคพฺโภ สมโณ โคตโม”ติ. ยสฺส โข พฺราหฺมณ อายติํ คพฺภเสยฺยา ปุนพฺภวา\-ภินิพฺพตฺติ ปหีนา อุจฺฉินฺนมูลา ตาลาวตฺถุกตา อนภาวํกตา อายติํ อนุปฺปาทธมฺมา, ตมหํ อปคพฺโภติ วทามิ. ตถาคตสฺส โข พฺราหฺมณ อายติํ คพฺภเสยฺยา ปุนพฺภวา\-ภินิพฺพตฺติ ปหีนา อุจฺฉินฺนมูลา ตาลาวตฺถุกตา อนภาวํกตา อายติํ อนุปฺปาทธมฺมา. อยํ โข พฺราหฺมณ ปริยาโย, เยน มํ ปริยาเยน สมฺมา วทมาโน วเทยฺย “อปคพฺโภ สมโณ โคตโม”ติ, โน จ โข ยํ ตฺวํ สนฺธาย วเทสิ.}

\proseitem{11}{เสยฺยถาปิ พฺราหฺมณ กุกฺกุฏิยา อณฺฑานิ อฏฺฐ วา ทส วา ทฺวาทส วา, ตานสฺสุ กุกฺกุฏิยา สมฺมา อธิสยิตานิ สมฺมา ปริเสทิตานิ สมฺมา ปริภาวิตานิ. โย นุ โข เตสํ กุกฺกุฏ\-จฺฉาปกานํ ปฐมตรํ ปาท\-นข\-สิขาย วา มุขตุณฺฑเกน วา อณฺฑโกสํ ปทาเลตฺวา โสตฺถินา อภินิพฺภิชฺเชยฺย, กินฺติ สฺวาสฺส วจนีโย เชฏฺโฐ วา กนิฏฺโฐ วาติ. เชฏฺโฐติสฺส โภ โคตม วจนีโย, โส หิ เนสํ เชฏฺโฐ โหตีติ. เอวเมว โข อหํ พฺราหฺมณ อวิชฺชาคตาย ปชาย อณฺฑภูตาย ปริโยนทฺธาย อวิชฺชณฺฑโกสํ ปทาเลตฺวา เอโกว โลเก อนุตฺตรํ สมฺมาสมฺโพธิํ อภิสมฺพุทฺโธ, สฺวาหํ พฺราหฺมณ เชฏฺโฐ เสฏฺโฐ โลกสฺส.}

\prose{อารทฺธํ โข ปน เม พฺราหฺมณ วีริยํ\csromansharedfootnote{29fffb8d80c6a398}{วิริยํ (สี, สฺยา)} อโหสิ อสลฺลีนํ, อุปฏฺฐิตา สติ อสมฺมุฏฺฐา\csromansharedfootnote{fbe59365fda0c2bf}{อปฺปมุฏฺฐา (สี, สฺยา)}, ปสฺสทฺโธ กาโย อสารทฺโธ, สมาหิตํ จิตฺตํ เอกคฺคํ. \csromanfolio{5}โส โข อหํ พฺราหฺมณ วิวิจฺเจว กาเมหิ วิวิจฺจ อกุสเลหิ ธมฺเมหิ สวิตกฺกํ สวิจารํ วิเวกชํ ปีติสุขํ ปฐมํ ฌานํ อุปสมฺปชฺช วิหาสิํ. วิตกฺก\-วิจารานํ วูปสมา อชฺฌตฺตํ สมฺปสาทนํ เจตโส เอโกทิภาวํ อวิตกฺกํ อวิจารํ สมาธิชํ ปีติสุขํ ทุติยํ ฌานํ อุปสมฺปชฺช วิหาสิํ. ปีติยา จ วิราคา อุเปกฺขโก จ วิหาสิํ สโต จ สมฺปชาโน, สุขญฺจ กาเยน ปฏิสํเวเทสิํ, ยํ ตํ อริยา อาจิกฺขนฺติ “อุเปกฺขโก สติมา สุขวิหารี”ติ ตติยํ ฌานํ อุปสมฺปชฺช วิหาสิํ. สุขสฺส จ ปหานา ทุกฺขสฺส จ ปหานา ปุพฺเพว โสมนสฺส\-โทมนสฺสานํ อตฺถงฺคมา อทุกฺขมสุขํ อุเปกฺขา\-สติ\-ปาริสุทฺธิํ จตุตฺถํ ฌานํ อุปสมฺปชฺช วิหาสิํ.}

\proseitem{12}{โส เอวํ สมาหิเต จิตฺเต ปริสุทฺเธ ปริโยทาเต อนงฺคเณ วิคตู\-ปกฺกิเลเส มุทุภูเต กมฺมนิเย ฐิเต อาเนญฺชปฺปตฺเต ปุพฺเพ\-นิวาสา\-นุสฺสติ\-ญาณาย จิตฺตํ อภินินฺนาเมสิํ. โส อเนกวิหิตํ ปุพฺเพนิวาสํ อนุสฺสรามิ, “เสยฺยถิทํ, เอกมฺปิ ชาติํ เทฺวปิ ชาติโย ติสฺโสปิ ชาติโย จตสฺโสปิ ชาติโย ปญฺจปิ ชาติโย ทสปิ ชาติโย วีสมฺปิ ชาติโย ติํสมฺปิ ชาติโย จตฺตาลีสมฺปิ ชาติโย ปญฺญาสมฺปิ ชาติโย ชาติสตมฺปิ ชาติสหสฺสมฺปิ ชาติ\-สต\-สหสฺสมฺปิ อเนเกปิ สํวฏฺฏกปฺเป อเนเกปิ วิวฏฺฏกปฺเป อเนเกปิ สํวฏฺฏ\-วิวฏฺฏ\-กปฺเป อมุตฺราสิํ เอวํนาโม เอวํโคตฺโต เอวํวณฺโณ เอวมาหาโร เอวํ\-สุขทุกฺข\-ปฺปฏิสํเวที เอวมา\-ยุปริยนฺโต, โส ตโต จุโต อมุตฺร อุทปาทิํ, ตตฺราปาสิํ เอวํนาโม เอวํโคตฺโต เอวํวณฺโณ เอวมาหาโร เอวํ\-สุขทุกฺข\-ปฺปฏิสํเวที เอวมา\-ยุปริยนฺโต, โส ตโต จุโต อิธูปปนฺโน”ติ, อิติ สาการํ สอุทฺเทสํ อเนกวิหิตํ ปุพฺเพนิวาสํ อนุสฺสรามิ. อยํ โข เม พฺราหฺมณ รตฺติยา ปฐเม ยาเม ปฐมา วิชฺชา อธิคตา, อวิชฺชา วิหตา, วิชฺชา อุปฺปนฺนา, ตโม วิหโต, อาโลโก อุปฺปนฺโน, ยถา ตํ อปฺปมตฺตสฺส อาตาปิโน ปหิตตฺตสฺส วิหรโต. อยํ โข เม พฺราหฺมณ ปฐมา\-ภินิพฺภิทา อโหสิ กุกฺกุฏจฺฉาปกสฺเสว อณฺฑโกสมฺหา.}

\proseitem{13}{โส เอวํ สมาหิเต จิตฺเต ปริสุทฺเธ ปริโยทาเต อนงฺคเณ วิคตู\-ปกฺกิเลเส มุทุภูเต กมฺมนิเย ฐิเต อาเนญฺชปฺปตฺเต สตฺตานํ จุตูปปาต\-ญาณาย จิตฺตํ อภินินฺนาเมสิํ. โส ทิพฺเพน จกฺขุนา วิสุทฺเธน อติกฺกนฺต\-มานุสเกน\csromansharedfootnote{9698df013f0ee02e}{อติกฺกนฺตมานุสฺสเกน (ก)} สตฺเต ปสฺสามิ จวมาเน อุปปชฺชมาเน หีเน \csromanfolio{6}ปณีเต สุวณฺเณ ทุพฺพณฺเณ สุคเต ทุคฺคเต, ยถากมฺมูปเค สตฺเต ปชานามิ “อิเม วต โภนฺโต สตฺตา กาย\-ทุจฺจริเตน สมนฺนาคตา วจีทุจฺจริเตน สมนฺนาคตา มโน\-ทุจฺจริเตน สมนฺนาคตา อริยานํ อุปวาทกา มิจฺฉาทิฏฺฐิกา มิจฺฉาทิฏฺฐิ\-กมฺม\-สมาทานา, เต กายสฺส เภทา ปรํ มรณา อปายํ ทุคฺคติํ วินิปาตํ นิรยํ อุปปนฺนา. อิเม วา ปน โภนฺโต สตฺตา กายสุจริเตน สมนฺนาคตา วจีสุจริเตน สมนฺนาคตา มโนสุจริเตน สมนฺนาคตา อริยานํ อนุปวาทกา สมฺมาทิฏฺฐิกา สมฺมา\-ทิฏฺฐิ\-กมฺม\-สมาทานา, เต กายสฺส เภทา ปรํ มรณา สุคติํ สคฺคํ โลกํ อุปปนฺนา”ติ. อิติ ทิพฺเพน จกฺขุนา วิสุทฺเธน อติกฺกนฺต\-มานุสเกน สตฺเต ปสฺสามิ จวมาเน อุปปชฺชมาเน หีเน ปณีเต สุวณฺเณ ทุพฺพณฺเณ สุคเต ทุคฺคเต, ยถากมฺมูปเค สตฺเต ปชานามิ. อยํ โข เม พฺราหฺมณ รตฺติยา มชฺฌิเม ยาเม ทุติยา วิชฺชา อธิคตา, อวิชฺชา วิหตา, วิชฺชา อุปฺปนฺนา, ตโม วิหโต, อาโลโก อุปฺปนฺโน, ยถา ตํ อปฺปมตฺตสฺส อาตาปิโน ปหิตตฺตสฺส วิหรโต. อยํ โข เม พฺราหฺมณ ทุติยา\-ภินิพฺภิทา อโหสิ กุกฺกุฏจฺฉาปกสฺเสว อณฺฑโกสมฺหา.}

\proseitem{14}{โส เอวํ สมาหิเต จิตฺเต ปริสุทฺเธ ปริโยทาเต อนงฺคเณ วิคตู\-ปกฺกิเลเส มุทุภูเต กมฺมนิเย ฐิเต อาเนญฺชปฺปตฺเต อาสวานํ ขยญาณาย จิตฺตํ อภินินฺนาเมสิํ. โส อิทํ ทุกฺขนฺติ ยถาภูตํ อพฺภญฺญาสิํ, อยํ ทุกฺข\-สมุทโยติ ยถาภูตํ อพฺภญฺญาสิํ, อยํ ทุกฺข\-นิโรโธติ ยถาภูตํ อพฺภญฺญาสิํ, อยํ ทุกฺขนิโรธ\-คามินี ปฏิปทาติ ยถาภูตํ อพฺภญฺญาสิํ, อิเม อาสวาติ ยถาภูตํ อพฺภญฺญาสิํ, อยํ อาสว\-สมุทโยติ ยถาภูตํ อพฺภญฺญาสิํ, อยํ อาสวนิโรโธติ ยถาภูตํ อพฺภญฺญาสิํ, อยํ อาสว\-นิโรธ\-คามินี ปฏิปทาติ ยถาภูตํ อพฺภญฺญาสิํ. ตสฺส เม เอวํ ชานโต เอวํ ปสฺสโต กามาสวาปิ จิตฺตํ วิมุจฺจิตฺถ, ภวาสวาปิ จิตฺตํ วิมุจฺจิตฺถ, อวิชฺชาสวาปิ จิตฺตํ วิมุจฺจิตฺถ, วิมุตฺตสฺมิํ วิมุตฺตมิติ ญาณํ อโหสิ, ขีณา ชาติ, วุสิตํ พฺรหฺมจริยํ, กตํ กรณียํ, นาปรํ อิตฺถตฺตายาติ อพฺภญฺญาสิํ. อยํ โข เม พฺราหฺมณ รตฺติยา ปจฺฉิเม ยาเม ตติยา วิชฺชา อธิคตา, อวิชฺชา วิหตา, วิชฺชา อุปฺปนฺนา, ตโม วิหโต, อาโลโก อุปฺปนฺโน, ยถา ตํ อปฺปมตฺตสฺส อาตาปิโน ปหิตตฺตสฺส วิหรโต. อยํ โข เม พฺราหฺมณ ตติยา\-ภินิพฺภิทา อโหสิ กุกฺกุฏจฺฉาปกสฺเสว อณฺฑโกสมฺหาติ.}

\csromanmatikaanchor{mtk.4}
\csromantocmarkref{section}{เวรญฺชายํ ทุพฺภิกฺขกถา}{mtk.4}
\bookmark[dest={mtk.4},level=1]{เวรญฺชายํ ทุพฺภิกฺขกถา}
\proseitem{15}{\csromanfolio{7}เอวํ วุตฺเต เวรญฺโช พฺราหฺมโณ ภควนฺตํ เอตทโวจ “เชฏฺโฐ ภวํ โคตโม, เสฏฺโฐ ภวํ โคตโม, อภิกฺกนฺตํ โภ โคตม, อภิกฺกนฺตํ โภ โคตม. เสยฺยถาปิ โภ โคตม นิกฺกุชฺชิตํ วา อุกฺกุชฺเชยฺย, ปฏิจฺฉนฺนํ วา วิวเรยฺย, มูฬฺหสฺส วา มคฺคํ อาจิกฺเขยฺย, อนฺธกาเร วา เตลปชฺโชตํ ธาเรยฺย ‘จกฺขุมนฺโต รูปานิ ทกฺขนฺตี’ติ, เอวเมวํ โภตา โคตเมน อเนก\-ปริยาเยน ธมฺโม ปกาสิโต. เอสาหํ ภวนฺตํ โคตมํ สรณํ คจฺฉามิ ธมฺมญฺจ ภิกฺขุสํฆญฺจ, อุปาสกํ มํ ภวํ โคตโม ธาเรตุ อชฺชตคฺเค ปาณุเปตํ สรณํ คตํ, อธิวาเสตุ จ เม ภวํ โคตโม เวรญฺชายํ วสฺสาวาสํ สทฺธิํ ภิกฺขุสํเฆนา”ติ. อธิวาเสสิ ภควา ตุณฺหีภาเวน. อถ โข เวรญฺโช พฺราหฺมโณ ภควโต อธิวาสนํ วิทิตฺวา อุฏฺฐายาสนา ภควนฺตํ อภิวาเทตฺวา ปทกฺขิณํ กตฺวา ปกฺกามิ.}

\proseitem{16}{เตน โข ปน สมเยน เวรญฺชา ทุพฺภิกฺขา โหติ ทฺวีหิติกา เสตฏฺฐิกา สลากาวุตฺตา, น สุกรา อุญฺเฉน ปคฺคเหน ยาเปตุํ. เตน โข ปน สมเยน อุตฺตราปถกา\csromansharedfootnote{0be49d3a014652a5}{อุตฺตราหกา (สี)} อสฺสวาณิชา\csromansharedfootnote{07d3d4e43adbf080}{อสฺสวณิชา (ก)} ปญฺจมตฺเตหิ อสฺสสเตหิ เวรญฺชํ วสฺสาวาสํ อุปคตา โหนฺติ, เตหิ อสฺสมณฺฑลิกาสุ ภิกฺขูนํ ปตฺถ\-ปตฺถ\-ปุลกํ\csromansharedfootnote{8078a99cfcc8ec52}{ปตฺถปตฺถมูลกํ (ก)} ปญฺญตฺตํ โหติ. ภิกฺขู ปุพฺพณฺหสมยํ นิวาเสตฺวา ปตฺต\-จีวรมา\-ทาย เวรญฺชํ ปิณฺฑาย ปวิสิตฺวา ปิณฺฑํ อลภมานา อสฺสมณฺฑลิกาสุ ปิณฺฑาย จริตฺวา ปตฺถ\-ปตฺถ\-ปุลกํ อารามํ อาหริตฺวา อุทุกฺขเล โกฏฺเฏตฺวา โกฏฺเฏตฺวา ปริภุญฺชนฺติ. อายสฺมา ปนานนฺโท ปตฺถปุลกํ สิลายํ ปิสิตฺวา ภควโต อุปนาเมติ, ตํ ภควา ปริภุญฺชติ.}

\prose{อสฺโสสิ โข ภควา อุทุกฺขล\-สทฺทํ. ชานนฺตาปิ ตถาคตา ปุจฺฉนฺติ, ชานนฺตาปิ น ปุจฺฉนฺติ. กาลํ วิทิตฺวา ปุจฺฉนฺติ, กาลํ วิทิตฺวา น ปุจฺฉนฺติ. อตฺถสํหิตํ ตถาคตา ปุจฺฉนฺติ โน อนตฺถ\-สํหิตํ. อนตฺถสํหิเต เสตุฆาโต ตถาคตานํ. ทฺวีหิ อากาเรหิ พุทฺธา ภควนฺโต ภิกฺขู ปฏิปุจฺฉนฺติ “ธมฺมํ วา เทเสสฺสาม, สาวกานํ วา สิกฺขาปทํ ปญฺญเปสฺสามา”ติ\csromansharedfootnote{03b81508cafe3540}{ปญฺญาเปสฺสามาติ (สี, สฺยา)}. อถ โข ภควา อายสฺมนฺตํ อานนฺทํ อามนฺเตสิ “กิํ นุ โข โส อานนฺท อุทุกฺขลสทฺโท”ติ. อถ โข อายสฺมา อานนฺโท ภควโต เอตมตฺถํ อาโรเจสิ. สาธุ สาธุ อานนฺท, ตุเมฺหหิ อานนฺท สปฺปุริเสหิ วิชิตํ, ปจฺฉิมา ชนตา สาลิมํโสทนํ อติมญฺญิสฺสตีติ.}

\csromanmatikaanchor{mtk.5}
\csromantocmarkref{section}{สิกฺขาปทปญฺญตฺยารมฺภกถา}{mtk.5}
\bookmark[dest={mtk.5},level=1]{สิกฺขาปทปญฺญตฺยารมฺภกถา}
\proseitem{17}{\csromanfolio{8}อถ โข อายสฺมา มหาโมคฺคลฺลาโน\csromansharedfootnote{c938078d1438d43e}{มหาโมคฺคลาโน (ก)} เยน ภควา เตนุปสงฺกมิ, อุปสงฺกมิตฺวา ภควนฺตํ อภิวาเทตฺวา เอกมนฺตํ นิสีทิ, เอกมนฺตํ นิสินฺโน โข อายสฺมา มหาโมคฺคลฺลาโน ภควนฺตํ เอตทโวจ “เอตรหิ ภนฺเต เวรญฺชา ทุพฺภิกฺขา ทฺวีหิติกา เสตฏฺฐิกา สลากาวุตฺตา, น สุกรา อุญฺเฉน ปคฺคเหน ยาเปตุํ. อิมิสฺสา ภนฺเต มหาปถวิยา เหฏฺฐิมตลํ สมฺปนฺนํ, เสยฺยถาปิ ขุทฺทมธุํ อนีลกํ, เอวมสฺสาทํ. สาธาหํ ภนฺเต ปถวิํ ปริวตฺเตยฺยํ, ภิกฺขู ปปฺปฏโกชํ ปริภุญฺชิสฺสนฺตี”ติ. เย ปน เต โมคฺคลฺลาน ปถวินิสฺสิตา ปาณา, เต กถํ กริสฺสสีติ. เอกาหํ ภนฺเต ปาณิํ อภินิมฺมินิสฺสามิ เสยฺยถาปิ มหาปถวี, เย ปถวินิสฺสิตา ปาณา, เต ตตฺถ สงฺกาเมสฺสามิ, เอเกน หตฺเถน ปถวิํ ปริวตฺเตสฺสามีติ. อลํ โมคฺคลฺลาน มา เต รุจฺจิ ปถวิํ ปริวตฺเตตุํ, วิปลฺลาสํปิ สตฺตา ปฏิลเภยฺยุนฺติ. สาธุ ภนฺเต สพฺโพ ภิกฺขุสํโฆ อุตฺตรกุรุํ ปิณฺฑาย คจฺเฉยฺยาติ. อลํ โมคฺคลฺลาน มา เต รุจฺจิ สพฺพสฺส ภิกฺขุ\-สํฆสฺส อุตฺตรกุรุํ ปิณฺฑาย คมนนฺติ.}

\proseitem{18}{อถ โข อายสฺมโต สาริปุตฺตสฺส รโหคตสฺส ปฏิสลฺลีนสฺส เอวํ เจตโส ปริวิตกฺโก อุทปาทิ “กตเมสานํ โข พุทฺธานํ ภควนฺตานํ พฺรหฺมจริยํ น จิรฏฺฐิติกํ อโหสิ, กตเมสานํ พุทฺธานํ ภควนฺตานํ พฺรหฺมจริยํ จิรฏฺฐิติกํ อโหสี”ติ. อถ โข อายสฺมา สาริปุตฺโต สายนฺหสมยํ\csromansharedfootnote{167bf7caa363071b}{สายณฺหสมยํ (สี)} ปฏิสลฺลานา วุฏฺฐิโต เยน ภควา เตนุปสงฺกมิ, อุปสงฺกมิตฺวา ภควนฺตํ อภิวาเทตฺวา เอกมนฺตํ นิสีทิ, เอกมนฺตํ นิสินฺโน โข อายสฺมา สาริปุตฺโต ภควนฺตํ เอตทโวจ “อิธ มยฺหํ ภนฺเต รโหคตสฺส ปฏิสลฺลีนสฺส เอวํ เจตโส ปริวิตกฺโก อุทปาทิ ‘กตเมสานํ โข พุทฺธานํ ภควนฺตานํ พฺรหฺมจริยํ น จิรฏฺฐิติกํ อโหสิ, กตเมสานํ พุทฺธานํ ภควนฺตานํ พฺรหฺมจริยํ จิรฏฺฐิติกํ อโหสี’ติ. กตเมสานํ นุ โข ภนฺเต พุทฺธานํ ภควนฺตานํ พฺรหฺมจริยํ น จิรฏฺฐิติกํ อโหสิ, กตเมสานํ พุทฺธานํ ภควนฺตานํ พฺรหฺมจริยํ จิรฏฺฐิติกํ อโหสี”ติ.}

\prose{ภควโต จ สาริปุตฺต วิปสฺสิสฺส ภควโต จ สิขิสฺส ภควโต จ เวสฺสภุสฺส พฺรหฺมจริยํ น จิรฏฺฐิติกํ อโหสิ. ภควโต จ สาริปุตฺต กกุสนฺธสฺส ภควโต จ โกณาคมนสฺส ภควโต จ กสฺสปสฺส พฺรหฺมจริยํ จิรฏฺฐิติกํ อโหสีติ.}

\proseitem{19}{\csromanfolio{9}โก นุ โข ภนฺเต เหตุ, โก ปจฺจโย, เยน ภควโต จ วิปสฺสิสฺส ภควโต จ สิขิสฺส ภควโต จ เวสฺสภุสฺส พฺรหฺมจรยํ น จิรฏฺฐิติกํ อโหสีติ. ภควา จ สาริปุตฺต วิปสฺสี ภควา จ สิขี ภควา จ เวสฺสภู กิลาสุโน อเหสุํ สาวกานํ วิตฺถาเรน ธมฺมํ เทเสตุํ, อปฺปกญฺจ เนสํ อโหสิ สุตฺตํ เคยฺยํ เวยฺยากรณํ คาถา อุทานํ อิติวุตฺตกํ ชาตกํ อพฺภุตธมฺมํ เวทลฺลํ, อปญฺญตฺตํ สาวกานํ สิกฺขาปทํ, อนุทฺทิฏฺฐํ ปาติโมกฺขํ. เตสํ พุทฺธานํ ภควนฺตานํ อนฺตรธาเนน พุทฺธา\-นุพุทฺธานํ สาวกานํ อนฺตรธาเนน เย เต ปจฺฉิมา สาวกา นานานามา นานาโคตฺตา นานาชจฺจา นานากุลา ปพฺพชิตา, เต ตํ พฺราหฺมจริยํ ขิปฺปญฺเญว อนฺตรธาเปสุํ. เสยฺยถาปิ สาริปุตฺต นานาปุปฺผานิ ผลเก นิกฺขิตฺตานิ สุตฺเตน อสงฺคหิตานิ, ตานิ วาโต วิกิรติ วิธมติ วิทฺธํเสติ, ตํ กิสฺส เหตุ, ยถา ตํ สุตฺเตน อสงฺคหิตตฺตา. เอวเมว โข สาริปุตฺต เตสํ พุทฺธานํ ภควนฺตานํ อนฺตรธาเนน พุทฺธา\-นุพุทฺธานํ สาวกานํ อนฺตรธาเนน เย เต ปจฺฉิมา สาวกา นานานามา นานาโคตฺตา นานาชจฺจา นานากุลา ปพฺพชิตา, เต ตํ พฺรหฺมจริยํ ขิปฺปญฺเญว อนฺตรธาเปสุํ.}

\prose{อกิลาสุโน จ เต ภควนฺโต อเหสุํ สาวเก เจตสา เจโต ปริจฺจ โอวทิตุํ. ภูตปุพฺพํ สาริปุตฺต เวสฺสภู ภควา อรหํ สมฺมาสมฺพุทฺโธ อญฺญตรสฺมิํ ภิํสนเก\csromansharedfootnote{7c1dbdb5104bd797}{ภีสนเก (ก)} วนสณฺเฑ สหสฺสํ ภิกฺขุสํฆํ เจตสา เจโต ปริจฺจ โอวทติ อนุสาสติ “เอวํ วิตกฺเกถ, มา เอวํ วิตกฺกยิตฺถ, เอวํ มนสิกโรถ, มา เอวํ มนสากตฺถ\csromansharedfootnote{61f9c392da9be219}{มนสากริตฺถ (ก)}, อิทํ ปชหถ, อิทํ อุปสมฺปชฺช วิหรถา”ติ. อถ โข สาริปุตฺต ตสฺส ภิกฺขุ\-สหสฺสสฺส เวสฺสภุนา ภควตา อรหตา สมฺมา\-สมฺพุทฺเธน เอวํ โอวทิยมานานํ เอวํ อนุสาสิย\-มานานํ อนุปาทาย อาสเวหิ จิตฺตานิ วิมุจฺจิํสุ. ตตฺร สุทํ สาริปุตฺต ภิํสนกสฺส วนสณฺฑสฺส ภิํสนกตสฺมิํ โหติ, โย โกจิ อวีตราโค ตํ วนสณฺฑํ ปวิสติ, เยภุยฺเยน โลมานิ หํสนฺติ. อยํ โข สาริปุตฺต เหตุ, อยํ ปจฺจโย, เยน ภควโต จ วิปสฺสิสฺส ภควโต จ สิขิสฺส ภควโต จ เวสฺสภุสฺส พฺรหฺมจริยํ น จิรฏฺฐิติกํ อโหสีติ.}

\proseitem{20}{\csromanfolio{10}โก ปน ภนฺเต เหตุ, โก ปจฺจโย, เยน ภควโต จ กกุสนฺธสฺส ภควโต จ โกณาคมนสฺส ภควโต จ กสฺสปสฺส พฺรหฺมจริยํ จิรฏฺฐิติกํ อโหสีติ. ภควา จ สาริปุตฺต กกุสนฺโธ ภควา จ โกณาคมโน ภควา จ กสฺสโป อกิลาสุโน อเหสุํ สาวกานํ วิตฺถาเรน ธมฺมํ เทเสตุํ, พหุญฺจ เนสํ อโหสิ สุตฺตํ เคยฺยํ เวยฺยากรณํ คาถา อุทานํ อิติวุตฺตกํ ชาตกํ อพฺภุตธมฺมํ เวทลฺลํ, ปญฺญตฺตํ สาวกานํ สิกฺขาปทํ, อุทฺทิฏฺฐํ ปาติโมกฺขํ. เตสํ พุทฺธานํ ภควนฺตานํ อนฺตรธาเนน พุทฺธา\-นุพุทฺธานํ สาวกานํ อนฺตรธาเนน เย เต ปจฺฉิมา สาวกา นานานามา นานาโคตฺตา นานาชจฺจา นานากุลา ปพฺพชิตา, เต ตํ พฺรหฺมจริยํ จิรํ ทีฆมทฺธานํ ฐเปสุํ. เสยฺยถาปิ สาริปุตฺต นานาปุปฺผานิ ผลเก นิกฺขิตฺตานิ สุตฺเตน สุสงฺคหิตานิ, ตานิ วาโต น วิกิรติ น วิธมติ น วิทฺธํเสติ, ตํ กิสฺส เหตุ, ยถา ตํ สุตฺเตน สุสงฺคหิตตฺตา. เอวเมว โข สาริปุตฺต เตสํ พุทฺธานํ ภควนฺตานํ อนฺตรธาเนน พุทฺธา\-นุพุทฺธานํ สาวกานํ อนฺตรธาเนน เย เต ปจฺฉิมา สาวกา นานานามา นานาโคตฺตา นานาชจฺจา นานากุลา ปพฺพชิตา, เต ตํ พฺรหฺมจริยํ จิรํ ทีฆมทฺธานํ ฐเปสุํ. อยํ โข สาริปุตฺต เหตุ, อยํ ปจฺจโย, เยน ภควโต จ กกุสนฺธสฺส ภควโต จ โกณาคมนสฺส ภควโต จ กสฺสปสฺส พฺรหฺมจริยํ จิรฏฺฐิติกํ อโหสีติ.}

\proseitem{21}{อถ โข อายสฺมา สาริปุตฺโต อุฏฺฐายาสนา เอกํสํ อุตฺตราสงฺคํ กริตฺวา เยน ภควา เตนญฺชลิํ ปณาเมตฺวา ภควนฺตํ เอตทโวจ “เอตสฺส ภควา กาโล, เอตสฺส สุคต กาโล, ยํ ภควา สาวกานํ สิกฺขาปทํ ปญฺญเปยฺย\csromansharedfootnote{4f0821d0c1622abd}{ปญฺญาเปยฺย (สี, สฺยา)}, อุทฺทิเสยฺย ปาติโมกฺขํ, ยถยิทํ พฺรหฺมจริยํ อทฺธนิยํ อสฺส จิรฏฺฐิติกนฺ”ติ. อาคเมหิ ตฺวํ สาริปุตฺต, อาคเมหิ ตฺวํ สาริปุตฺต, ตถาคโตว ตตฺถ กาลํ ชานิสฺสติ. น ตาว สาริปุตฺต สตฺถา สาวกานํ สิกฺขาปทํ ปญฺญเปติ อุทฺทิสติ\csromansharedfootnote{6c125bcd1c1c8621}{น อุทฺทิสติ (สี)} ปาติโมกฺขํ, ยาว น อิเธกจฺเจ อาสวฏฺฐานียา ธมฺมา สํเฆ ปาตุภวนฺติ. ยโต จ โข สาริปุตฺต อิเธกจฺเจ อาสวฏฺฐานียา ธมฺมา สํเฆ ปาตุภวนฺติ, อถ สตฺถา สาวกานํ สิกฺขาปทํ ปญฺญเปติ อุทฺทิสติ ปาติโมกฺขํ เตสํ เยว\csromansharedfootnote{9a3af51fcf41c94d}{เตสํเยว นิคฺคหีตสนฺธิ \csromandash{} ม.พ.ป.} อาสว\-ฏฺฐานียานํ ธมฺมานํ ปฏิฆาตาย. น ตาว สาริปุตฺต อิเธกจฺเจ \csromanfolio{11}อาสวฏฺฐานียา ธมฺมา สํเฆ ปาตุภวนฺติ, ยาว น สํโฆ รตฺตญฺญุ\-มหตฺตํ ปตฺโต โหติ. ยโต จ โข สาริปุตฺต สํโฆ รตฺตญฺญุ\-มหตฺตํ ปตฺโต โหติ, อถ อิเธกจฺเจ อาสวฏฺฐานียา ธมฺมา สํเฆ ปาตุภวนฺติ, อถ สตฺถา สาวกานํ สิกฺขาปทํ ปญฺญเปติ อุทฺทิสติ ปาติโมกฺขํ เตสํ เยว\csromansharedfootnote{9a3af51fcf41c94d}{เตสํเยว นิคฺคหีตสนฺธิ \csromandash{} ม.พ.ป.} อาสว\-ฏฺฐานียานํ ธมฺมานํ ปฏิฆาตาย. น ตาว สาริปุตฺต อิเธกจฺเจ อาสวฏฺฐานียา ธมฺมา สํเฆ ปาตุภวนฺติ, ยาว น สํโฆ เวปุลฺล\-มหตฺตํ ปตฺโต โหติ. ยโต จ โข สาริปุตฺต สํโฆ เวปุลฺล\-มหตฺตํ ปตฺโต โหติ, อถ อิเธกจฺเจ อาสวฏฺฐานียา ธมฺมา สํเฆ ปาตุภวนฺติ, อถ สตฺถา สาวกานํ สิกฺขาปทํ ปญฺญเปติ อุทฺทิสติ ปาติโมกฺขํ เตสํ เยว\csromansharedfootnote{9a3af51fcf41c94d}{เตสํเยว นิคฺคหีตสนฺธิ \csromandash{} ม.พ.ป.} อาสว\-ฏฺฐานียานํ ธมฺมานํ ปฏิฆาตาย. น ตาว สาริปุตฺต อิเธกจฺเจ อาสวฏฺฐานียา ธมฺมา สํเฆ ปาตุภวนฺติ, ยาว น สํโฆ ลาภคฺค\-มหตฺตํ ปตฺโต โหติ. ยโต จ โข สาริปุตฺต สํโฆ ลาภคฺค\-มหตฺตํ ปตฺโต โหติ, อถ อิเธกจฺเจ อาสวฏฺฐานียา ธมฺมา สํเฆ ปาตุภวนฺติ, อถ สตฺถา สาวกานํ สิกฺขาปทํ ปญฺญเปติ อุทฺทิสติ ปาติโมกฺขํ เตสํ เยว\csromansharedfootnote{9a3af51fcf41c94d}{เตสํเยว นิคฺคหีตสนฺธิ \csromandash{} ม.พ.ป.} อาสว\-ฏฺฐานียานํ ธมฺมานํ ปฏิฆาตาย. น ตาว สาริปุตฺต อิเธกจฺเจ อาสวฏฺฐานียา ธมฺมา สํเฆ ปาตุภวนฺติ, ยาว น สํโฆ พาหุสจฺจ\-มหตฺตํ ปตฺโต โหติ. ยโต จ โข สาริปุตฺต สํโฆ พาหุสจฺจ\-มหตฺตํ ปตฺโต โหติ, อถ อิเธกจฺเจ อาสวฏฺฐานียา ธมฺมา สํเฆ ปาตุภวนฺติ, อถ สตฺถา สาวกานํ สิกฺขาปทํ ปญฺญเปติ อุทฺทิสติ ปาติโมกฺขํ เตสํ เยว\csromansharedfootnote{9a3af51fcf41c94d}{เตสํเยว นิคฺคหีตสนฺธิ \csromandash{} ม.พ.ป.} อาสว\-ฏฺฐานียานํ ธมฺมานํ ปฏิฆาตาย. นิรพฺพุโท หิ สาริปุตฺต ภิกฺขุสํโฆ นิราทีนโว อปคตกาฬโก สุทฺโธ สาเร ปติฏฺฐิโต, อิเมสํ หิ สาริปุตฺต ปญฺจนฺนํ ภิกฺขุสตานํ โย ปจฺฉิมโก ภิกฺขุ, โส โสตาปนฺโน อวินิปาต\-ธมฺโม นิยโต สมฺโพธิ\-ปรายโณติ.}

\proseitem{22}{อถ โข ภควา อายสฺมนฺตํ อานนฺทํ อามนฺเตสิ “อาจิณฺณํ โข ปเนตํ อานนฺท ตถาคตานํ เยหิ นิมนฺติตา วสฺสํ วสนฺติ, น เต อนปโลเกตฺวา ชนปท\-จาริกํ ปกฺกมนฺติ, อายามานนฺท เวรญฺชํ พฺราหฺมณํ อปโลเกสฺสามา”ติ. เอวํ ภนฺเตติ โข อายสฺมา อานนฺโท ภควโต ปจฺจสฺโสสิ. อถ โข ภควา นิวาเสตฺวา ปตฺต\-จีวรมา\-ทาย อายสฺมตา อานนฺเทน ปจฺฉาสมเณน เยน เวรญฺชสฺส พฺราหฺมณสฺส นิเวสนํ เตนุปสงฺกมิ, อุปสงฺกมิตฺวา ปญฺญตฺเต อาสเน นิสีทิ. อถ โข \csromanfolio{12}เวรญฺโช พฺราหฺมโณ เยน ภควา เตนุปสงฺกมิ, อุปสงฺกมิตฺวา ภควนฺตํ อภิวาเทตฺวา เอกมนฺตํ นิสีทิ, เอกมนฺตํ นิสินฺนํ โข \textbf{เวรญฺชํ พฺราหฺมณํ ภควา เอตทโวจ “นิมนฺติตมฺห ตยา พฺราหฺมณ} \textbf{วสฺสํวุฏฺฐา}\csromansharedfootnote{c60f16da0fc44b82}{วสฺสํวุตฺถา (สี, สฺยา, ก)}\textbf{, อปโลเกม ตํ อิจฺฉาม มยํ ชนปท}\-\textbf{จาริกํ} ปกฺกมิตุนฺ”ติ. สจฺจํ โภ โคตม นิมนฺติตตฺถ มยา วสฺสํวุฏฺฐา, อปิ จ โย เทยฺยธมฺโม, โส น ทินฺโน. ตญฺจ โข โน อสนฺตํ, โนปิ อทาตุกมฺยตา, ตํ กุเตตฺถ ลพฺภา พหุกิจฺจา ฆราวาสา พหุกรณียา, อธิวาเสตุ เม ภวํ โคตโม สฺวาตนาย ภตฺตํ สทฺธิํ ภิกฺขุสํเฆนาติ. อธิวาเสสิ ภควา ตุณฺหีภาเวน. อถ โข ภควา เวรญฺชํ พฺราหฺมณํ ธมฺมิยา กถาย สนฺทสฺเสตฺวา สมาทเปตฺวา สมุตฺเตเชตฺวา สมฺปหํเสตฺวา อุฏฺฐายาสนา ปกฺกามิ. อถ โข เวรญฺโช พฺราหฺมโณ ตสฺสา รตฺติยา อจฺจเยน สเก นิเวสเน ปณีตํ ขาทนียํ โภชนียํ ปฏิยาทาเปตฺวา ภควโต กาลํ อาโรจาเปสิ “กาโล โภ โคตม นิฏฺฐิตํ ภตฺตนฺ”ติ.}

\proseitemwithcloser{23}{อถ โข ภควา ปุพฺพณฺหสมยํ นิวาเสตฺวา ปตฺต\-จีวรมา\-ทาย เยน เวรญฺชสฺส พฺราหฺมณสฺส นิเวสนํ เตนุปสงฺกมิ, อุปสงฺกมิตฺวา ปญฺญตฺเต อาสเน นิสีทิ สทฺธิํ ภิกฺขุ\-สํเฆน. อถ โข เวรญฺโช พฺราหฺมโณ พุทฺธ\-ปฺปมุขํ ภิกฺขุสํฆํ ปณีเตน ขาทนีเยน โภชนีเยน สหตฺถา สนฺตปฺเปตฺวา สมฺปวาเรตฺวา ภควนฺตํ ภุตฺตาวิํ โอนีต\-ปตฺต\-ปาณิํ\csromansharedfootnote{13301a4065a0645f}{โอณีตปตฺตปาณิํ (ก)} ติจีวเรน อจฺฉาเทสิ, เอกเมกญฺจ ภิกฺขุํ เอกเมเกน ทุสฺสยุเคน อจฺฉาเทสิ. อถโข\csromansharedfootnote{b8dfa686fd155ee3}{อถ โข เทฺว นิปาตา \csromandash{} ม.พ.ป.} ภควา เวรญฺชํ พฺราหฺมณํ ธมฺมิยา กถาย สนฺทสฺเสตฺวา สมาทเปตฺวา สมุตฺเตเชตฺวา สมฺปหํเสตฺวา อุฏฺฐายาสนา ปกฺกามิ. อถ โข ภควา เวรญฺชายํ ยถาภิรนฺตํ วิหริตฺวา อนุปคมฺม โสเรยฺยํ สงฺกสฺสํ กณฺณกุชฺชํ เยน ปยาค\-ปติฏฺฐานํ เตนุปสงฺกมิ, อุปสงฺกมิตฺวา ปยาค\-ปติฏฺฐาเน คงฺคํ นทิํ อุตฺตริตฺวา เยน พาราณสี ตทวสริ. อถ โข ภควา พาราณสิยํ ยถาภิรนฺตํ วิหริตฺวา เยน เวสาลี เตน จาริกํ ปกฺกามิ. อนุปุพฺเพน จาริกํ จรมาโน เยน เวสาลี ตทวสริ. ตตฺร สุทํ ภควา เวสาลิยํ วิหรติ มหาวเน กูฏาคารสาลา\-ยนฺติ.}{เวรญฺช\-ภาณวาโร นิฏฺฐิโต.}
\csromanmatikaanchor{mtk.6}
\csromantocmarknopage{chapter}{1. ปาราชิกกณฺฑ}{mtk.6}
\bookmark[dest={mtk.6},level=0]{1. ปาราชิกกณฺฑ}
\chapterheadpageread{1. \textbf{ปาราชิกกณฺฑ}}
\csromanmatikaanchor{mtk.7}
\csromanmatikaanchor{mtk.8}
\csromantocmarknopage{section}{1. ปฐมปาราชิก}{mtk.7}
\bookmark[dest={mtk.7},level=1]{1. ปฐมปาราชิก}
\csromantocmarkref{subsection}{สุทินฺนภาณวาร สุทินฺนภิกฺขุวตฺถุ}{mtk.8}
\bookmark[dest={mtk.8},level=2]{สุทินฺนภาณวาร สุทินฺนภิกฺขุวตฺถุ}
\csromanheader{2}{1. \textbf{ปฐม}\-\textbf{ปาราชิก สุทินฺน}\-\textbf{ภาณวาร}}
\proseitem{24}{\csromanfolio{13}เตน โข ปน สมเยน เวสาลิยา อวิทูเร กลนฺทคาโม นาม อตฺถิ\csromansharedfootnote{e5cdaed7dfd88179}{กลนฺทคาโม นาม โหติ (สี), กลนฺทคาโม โหติ (สฺยา)}, ตตฺถ สุทินฺโน นาม กลนฺทปุตฺโต เสฏฺฐิปุตฺโต โหติ. อถ โข สุทินฺโน กลนฺทปุตฺโต สมฺพหุเลหิ\csromansharedfootnote{af23566eab9af4de}{สมฺปหูเลหิ (สี)} สหายเกหิ สทฺธิํ เวสาลิํ อคมาสิ เกนจิเทว กรณีเยน. เตน โข ปน สมเยน ภควา มหติยา ปริสาย ปริวุโต ธมฺมํ เทเสนฺโต นิสินฺโน โหติ. อทฺทส โข สุทินฺโน กลนฺทปุตฺโต ภควนฺตํ มหติยา ปริสาย ปริวุตํ ธมฺมํ เทเสนฺตํ นิสินฺนํ, ทิสฺวานสฺส เอตทโหสิ “ยนฺนูนาหมฺปิ ธมฺมํ สุเณยฺยนฺ”ติ. อถ โข สุทินฺโน กลนฺทปุตฺโต เยน สา ปริสา เตนุปสงฺกมิ, อุปสงฺกมิตฺวา เอกมนฺตํ นิสีทิ, เอกมนฺตํ นิสินฺนสฺส โข สุทินฺนสฺส กลนฺท\-ปุตฺตสฺส เอตทโหสิ “ยถา ยถา โข อหํ ภควตา ธมฺมํ เทสิตํ อาชานามิ, นยิทํ สุกรํ อคารํ อชฺฌาวสตา เอกนฺต\-ปริปุณฺณํ เอกนฺต\-ปริสุทฺธํ สงฺขลิขิตํ พฺรหฺมจริยํ จริตุํ, ยนฺนูนาหํ เกสมสฺสุํ โอหาเรตฺวา กาสายานิ วตฺถานิ อจฺฉาเทตฺวา อคารสฺมา อนคาริยํ ปพฺพเชยฺยนฺ”ติ. อถ โข สา ปริสา ภควตา ธมฺมิยา กถาย สนฺทสฺสิตา สมาทปิตา สมุตฺเตชิตา สมฺปหํสิตา อุฏฺฐายาสนา ภควนฺตํ อภิวาเทตฺวา ปทกฺขิณํ กตฺวา ปกฺกามิ.}

\proseitem{25}{อถ โข สุทินฺโน กลนฺทปุตฺโต อจิร\-วุฏฺฐิตาย ปริสาย เยน ภควา เตนุปสงฺกมิ, อุปสงฺกมิตฺวา ภควนฺตํ อภิวาเทตฺวา เอกมนฺตํ นิสีทิ, เอกมนฺตํ นิสินฺโน โข สุทินฺโน กลนฺทปุตฺโต ภควนฺตํ เอตทโวจ “ยถา ยถาหํ ภนฺเต ภควตา ธมฺมํ เทสิตํ อาชานามิ, นยิทํ สุกรํ อคารํ อชฺฌาวสตา เอกนฺต\-ปริปุณฺณํ เอกนฺต\-ปริสุทฺธํ สงฺขลิขิตํ พฺรหฺมจริยํ จริตุํ, อิจฺฉามหํ ภนฺเต เกสมสฺสุํ โอหาเรตฺวา กาสายานิ วตฺถานิ อจฺฉาเทตฺวา อคารสฺมา อนคาริยํ ปพฺพชิตุํ, ปพฺพาเชตุ มํ ภควา”ติ. อนุญฺญาโตสิ ปน ตฺวํ สุทินฺน มาตาปิตูหิ อคารสฺมา อนคาริยํ ปพฺพชฺชายาติ. น โข อหํ ภนฺเต อนุญฺญาโต มาตาปิตูหิ อคารสฺมา \csromanfolio{14}อนคาริยํ ปพฺพชฺชายาติ. น โข สุทินฺน ตถาคตา อนนุญฺญาตํ มาตาปิตูหิ ปุตฺตํ ปพฺพาเชนฺตีติ. โสหํ ภนฺเต ตถา กริสฺสามิ, ยถา มํ มาตาปิตโร อนุชานิสฺสนฺติ อคารสฺมา อนคาริยํ ปพฺพชฺชายาติ.}

\proseitem{26}{อถ โข สุทินฺโน กลนฺทปุตฺโต เวสาลิยํ ตํ กรณียํ ตีเรตฺวา เยน กลนฺทคาโม เยน มาตาปิตโร เตนุปสงฺกมิ, อุปสงฺกมิตฺวา มาตาปิตโร เอตทโวจ “อมฺมตาตา ยถา ยถาหํ ภควตา ธมฺมํ เทสิตํ อาชานามิ, นยิทํ สุกรํ อคารํ อชฺฌาวสตา เอกนฺต\-ปริปุณฺณํ เอกนฺต\-ปริสุทฺธํ สงฺขลิขิตํ พฺรหฺมจริยํ จริตุํ, อิจฺฉามหํ เกสมสฺสุํ โอหาเรตฺวา กาสายานิ วตฺถานิ อจฺฉาเทตฺวา อคารสฺมา อนคาริยํ ปพฺพชิตุํ, อนุชานาถ มํ อคารสฺมา อนคาริยํ ปพฺพชฺชายา”ติ. เอวํ วุตฺเต สุทินฺนสฺส กลนฺท\-ปุตฺตสฺส มาตาปิตโร สุทินฺนํ กลนฺทปุตฺตํ เอตทโวจุํ “ตฺวํ โขสิ ตาต สุทินฺน อมฺหากํ เอกปุตฺตโก ปิโย มนาโป สุเขธิโต สุขปริหโต, น ตฺวํ ตาต สุทินฺน กิญฺจิ ทุกฺขสฺส ชานาสิ, มรเณนปิ มยํ เต อกามกา วินา ภวิสฺสาม, กิํ ปน มยํ ตํ ชีวนฺตํ อนุชานิสฺสาม อคารสฺมา อนคาริยํ ปพฺพชฺชายา”ติ. ทุติยมฺปิ โข สุทินฺโน กลนฺทปุตฺโต มาตาปิตโร เอตทโวจ “อมฺมตาตา ยถา ยถาหํ ภควตา ธมฺมํ เทสิตํ อาชานามิ, นยิทํ สุกรํ อคารํ อชฺฌาวสตา เอกนฺต\-ปริปุณฺณํ เอกนฺต\-ปริสุทฺธํ สงฺขลิขิตํ พฺรหฺมจริยํ จริตุํ, อิจฺฉามหํ เกสมสฺสุํ โอหาเรตฺวา กาสายานิ วตฺถานิ อจฺฉาเทตฺวา อคารสฺมา อนคาริยํ ปพฺพชิตุํ, อนุชานาถ มํ อคารสฺมา อนคาริยํ ปพฺพชฺชายา”ติ. ทุติยมฺปิ โข สุทินฺนสฺส กลนฺท\-ปุตฺตสฺส มาตาปิตโร สุทินฺนํ กลนฺทปุตฺตํ เอตทโวจุํ “ตฺวํ โขสิ ตาต สุทินฺน อมฺหากํ เอกปุตฺตโก ปิโย มนาโป สุเขธิโต สุขปริหโต, น ตฺวํ ตาต สุทินฺน กิญฺจิ ทุกฺขสฺส ชานาสิ, มรเณนปิ มยํ เต อกามกา วินา ภวิสฺสาม, กิํ ปน มยํ ตํ ชีวนฺตํ อนุชานิสฺสาม อคารสฺมา อนคาริยํ ปพฺพชฺชายา”ติ. ตติยมฺปิ โข สุทินฺโน กลนฺทปุตฺโต มาตาปิตโร เอตทโวจ “อมฺมตาตา ยถา ยถาหํ ภควตา ธมฺมํ เทสิตํ อาชานามิ, นยิทํ สุกรํ อคารํ อชฺฌาวสตา เอกนฺต\-ปริปุณฺณํ เอกนฺต\-ปริสุทฺธํ สงฺขลิขิตํ พฺรหฺมจริยํ จริตุํ, อิจฺฉามหํ เกสมสฺสุํ โอหาเรตฺวา กาสายานิ วตฺถานิ อจฺฉาเทตฺวา อคารสฺมา อนคาริยํ ปพฺพชิตุํ, อนุชานาถ มํ อคารสฺมา อนคาริยํ ปพฺพชฺชายา”ติ. \csromanfolio{15}ตติยมฺปิ โข สุทินฺนสฺส กลนฺท\-ปุตฺตสฺส มาตาปิตโร สุทินฺนํ กลนฺทปุตฺตํ เอตทโวจุํ “ตฺวํ โขสิ ตาต สุทินฺน อมฺหากํ เอกปุตฺตโก ปิโย มนาโป สุเขธิโต สุขปริหโต, น ตฺวํ ตาต สุทินฺน กิญฺจิ ทุกฺขสฺส ชานาสิ, มรเณนปิ มยํ เต อกามกา วินา ภวิสฺสาม, กิํ ปน มยํ ตํ ชีวนฺตํ อนุชานิสฺสาม อคารสฺมา อนคาริยํ ปพฺพชฺชายา”ติ.}

\proseitem{27}{อถ โข สุทินฺโน กลนฺทปุตฺโต “น มํ มาตาปิตโร อนุชานนฺติ อคารสฺมา อนคาริยํ ปพฺพชฺชายา”ติ ตตฺเถว อนนฺตรหิตาย ภูมิยา นิปชฺชิ “อิเธว เม มรณํ ภวิสฺสติ ปพฺพชฺชา วา”ติ. อถ โข สุทินฺโน กลนฺทปุตฺโต เอกมฺปิ ภตฺตํ น ภุญฺชิ, เทฺวปิ ภตฺตานิ น ภุญฺชิ, ตีณิปิ ภตฺตานิ น ภุญฺชิ, จตฺตาริปิ ภตฺตานิ น ภุญฺชิ, ปญฺจปิ ภตฺตานิ น ภุญฺชิ, ฉปิ ภตฺตานิ น ภุญฺชิ, สตฺตปิ ภตฺตานิ น ภุญฺชิ.}

\proseitem{28}{อถ โข สุทินฺนสฺส กลนฺท\-ปุตฺตสฺส มาตาปิตโร สุทินฺนํ กลนฺทปุตฺตํ เอตทโวจุํ “ตฺวํ โขสิ ตาต สุทินฺน อมฺหากํ เอกปุตฺตโก ปิโย มนาโป สุเขธิโต สุขปริหโต, น ตฺวํ ตาต สุทินฺน กิญฺจิ ทุกฺขสฺส ชานาสิ, มรเณนปิ มยํ เต อกามกา วินา ภวิสฺสาม, กิํ ปน มยํ ตํ ชีวนฺตํ อนุชานิสฺสาม อคารสฺมา อนคาริยํ ปพฺพชฺชาย, อุฏฺเฐหิ ตาต สุทินฺน ภุญฺช จ ปิว จ ปริจาเรหิ จ, ภุญฺชนฺโต ปิวนฺโต ปริจาเรนฺโต กาเม ปริภุญฺชนฺโต ปุญฺญานิ กโรนฺโต อภิรมสฺสุ, น ตํ มยํ อนุชานาม อคารสฺมา อนคาริยํ ปพฺพชฺชายา”ติ. เอวํ วุตฺเต สุทินฺโน กลนฺทปุตฺโต ตุณฺหี อโหสิ. ทุติยมฺปิ โข ฯเปฯ. ตติยมฺปิ โข สุทินฺนสฺส กลนฺท\-ปุตฺตสฺส มาตาปิตโร สุทินฺนํ กลนฺทปุตฺตํ เอตทโวจุํ “ตฺวํ โขสิ ตาต สุทินฺน อมฺหากํ เอกปุตฺตโก ปิโย มนาโป สุเขธิโต สุขปริหโต, น ตฺวํ ตาต สุทินฺน กิญฺจิ ทุกฺขสฺส ชานาสิ, มรเณนปิ มยํ เต อกามกา วินา ภวิสฺสาม, กิํ ปน มยํ ตํ ชีวนฺตํ อนุชานิสฺสาม อคารสฺมา อนคาริยํ ปพฺพชฺชาย, อุฏฺเฐหิ ตาต สุทินฺน ภุญฺช จ ปิว จ ปริจาเรหิ จ, ภุญฺชนฺโต ปิวนฺโต ปริจาเรนฺโต กาเม ปริภุญฺชนฺโต ปุญฺญานิ กโรนฺโต อภิรมสฺสุ, น ตํ มยํ อนุชานาม อคารสฺมา อนคาริยํ ปพฺพชฺชายา”ติ. ตติยมฺปิ โข สุทินฺโน กลนฺทปุตฺโต ตุณฺหี อโหสิ.}

\prose{อถ โข สุทินฺนสฺส กลนฺท\-ปุตฺตสฺส สหายกา เยน สุทินฺโน กลนฺทปุตฺโต เตนุ\-ปสงฺกมิํสุ, อุปสงฺกมิตฺวา สุทินฺนํ กลนฺทปุตฺตํ \csromanfolio{16}เอตทโวจุํ “ตฺวํ โขสิ สมฺม สุทินฺน มาตาปิตูนํ เอกปุตฺตโก ปิโย มนาโป สุเขธิโต สุขปริหโต, น ตฺวํ สมฺม สุทินฺน กิญฺจิ ทุกฺขสฺส ชานาสิ, มรเณนปิ เต มาตาปิตโร อกามกา วินา ภวิสฺสนฺติ, กิํ ปน ตํ ชีวนฺตํ อนุชานิสฺสนฺติ อคารสฺมา อนคาริยํ ปพฺพชฺชาย, อุฏฺเฐหิ สมฺม สุทินฺน ภุญฺช จ ปิว จ ปริจาเรหิ จ, ภุญฺชนฺโต ปิวนฺโต ปริจาเรนฺโต กาเม ปริภุญฺชนฺโต ปุญฺญานิ กโรนฺโต อภิรมสฺสุ, น ตํ มาตาปิตโร อนุชานิสฺสนฺติ อคารสฺมา อนคาริยํ ปพฺพชฺชายา”ติ. เอวํ วุตฺเต สุทินฺโน กลนฺทปุตฺโต ตุณฺหี อโหสิ. ทุติยมฺปิ โข ฯเปฯ. ตติยมฺปิ โข สุทินฺนสฺส กลนฺท\-ปุตฺตสฺส สหายกา สุทินฺนํ กลนฺทปุตฺตํ เอตทโวจุํ “ตฺวํ โขสิ สมฺม สุทินฺน ฯเปฯ. ตติยมฺปิ โข สุทินฺโน กลนฺทปุตฺโต ตุณฺหี อโหสิ.}

\proseitem{29}{อถ โข สุทินฺนสฺส กลนฺท\-ปุตฺตสฺส สหายกา เยน สุทินฺนสฺส กลนฺท\-ปุตฺตสฺส มาตาปิตโร เตนุ\-ปสงฺกมิํสุ, อุปสงฺกมิตฺวา สุทินฺนสฺส กลนฺท\-ปุตฺตสฺส มาตาปิตโร เอตทโวจุํ “อมฺมตาตา เอโส สุทินฺโน อนนฺตรหิตาย ภูมิยา นิปนฺโน ‘อิเธว เม มรณํ ภวิสฺสติ ปพฺพชฺชา วา’ติ, สเจ ตุเมฺห สุทินฺนํ นานุชานิสฺสถ อคารสฺมา อนคาริยํ ปพฺพชฺชาย, ตตฺเถว มรณํ อาคมิสฺสติ, สเจ ปน ตุเมฺห สุทินฺนํ อนุชานิสฺสถ อคารสฺมา อนคาริยํ ปพฺพชฺชาย, ปพฺพชิตมฺปิ นํ ทกฺขิสฺสถ, สเจ สุทินฺโน นาภิรมิสฺสติ อคารสฺมา อนคาริยํ ปพฺพชฺชาย, กา ตสฺส อญฺญา คติ ภวิสฺสติ, อิเธว ปจฺจาคมิสฺสติ, อนุชานาถ สุทินฺนํ อคารสฺมา อนคาริยํ ปพฺพชฺชายา”ติ. อนุชานาม ตาตา สุทินฺนํ อคารสฺมา อนคาริยํ ปพฺพชฺชายาติ. อถ โข สุทินฺนสฺส กลนฺท\-ปุตฺตสฺส สหายกา เยน สุทินฺโน กลนฺทปุตฺโต เตนุ\-ปสงฺกมิํสุ, อุปสงฺกมิตฺวา สุทินฺนํ กลนฺทปุตฺตํ เอตทโวจุํ “อุฏฺเฐหิ สมฺม สุทินฺน อนุญฺญาโตสิ มาตาปิตูหิ อคารสฺมา อนคาริยํ ปพฺพชฺชายา”ติ.}

\proseitem{30}{อถ โข สุทินฺโน กลนฺทปุตฺโต “อนุญฺญาโตมฺหิ กิร มาตาปิตูหิ อคารสฺมา อนคาริยํ ปพฺพชฺชายา”ติ หฏฺโฐ อุทคฺโค ปาณินา คตฺตานิ ปริปุญฺฉนฺโต วุฏฺฐาสิ. อถ โข สุทินฺโน กลนฺทปุตฺโต กติปาหํ พลํ คาเหตฺวา เยน ภควา เตนุปสงฺกมิ, อุปสงฺกมิตฺวา ภควนฺตํ อภิวาเทตฺวา เอกมนฺตํ นิสีทิ, เอกมนฺตํ นิสินฺโน โข สุทินฺโน \csromanfolio{17}กลนฺทปุตฺโต ภควนฺตํ เอตทโวจ “อนุญฺญาโต\csromansharedfootnote{9199df826c52f38c}{อนุญฺญาโตมฺหิ (สี, สฺยา)} อหํ ภนฺเต มาตาปิตูหิ อคารสฺมา อนคาริยํ ปพฺพชฺชาย, ปพฺพาเชตุ มํ ภควา”ติ. อลตฺถ โข สุทินฺโน กลนฺทปุตฺโต ภควโต สนฺติเก ปพฺพชฺชํ, อลตฺถ อุปสมฺปทํ. อจิรู\-ปสมฺปนฺโน จ ปนายสฺมา สุทินฺโน เอวรูเป ธุตคุเณ สมาทาย วตฺตติ, อารญฺญิโก โหติ ปิณฺฑปาติโก ปํสุกูลิโก สปทานจาริโก, อญฺญตรํ วชฺชิคามํ อุปนิสฺสาย วิหรติ.}

\prose{เตน โข ปน สมเยน วชฺชี ทุพฺภิกฺขา โหติ ทฺวีหิติกา เสตฏฺฐิกา สลากาวุตฺตา, น สุกรา อุญฺเฉน ปคฺคเหน ยาเปตุํ. อถ โข อายสฺมโต สุทินฺนสฺส เอตทโหสิ “เอตรหิ โข วชฺชี ทุพฺภิกฺขา ทฺวีหิติกา เสตฏฺฐิกา สลากาวุตฺตา, น สุกรา อุญฺเฉน ปคฺคเหน ยาเปตุํ, พหู โข ปน เม เวสาลิยํ ญาตี อฑฺฒามหทฺธนา\csromansharedfootnote{9df40a79f6ac748d}{อฑฺฒา มหทฺธนา เทฺว ปทานิ \csromandash{} ม.พ.ป.} มหาโภคา ปหูต\-ชาตรูป\-รชตา ปหูต\-วิตฺตู\-ปกรณา ปหูต\-ธน\-ธญฺญา, ยํนูนาหํ ญาตี อุปนิสฺสาย วิหเรยฺยํ, ญาตี มํ\csromansharedfootnote{1ed421db6affa681}{ญาตกาปิ มํ (สฺยา)} นิสฺสาย ทานานิ ทสฺสนฺติ, ปุญฺญานิ กริสฺสนฺติ, ภิกฺขู จ ลาภํ ลจฺฉนฺติ, อหญฺจ ปิณฺฑเกน น กิลมิสฺสามี”ติ. อถ โข อายสฺมา สุทินฺโน เสนาสนํ สํสาเมตฺวา ปตฺต\-จีวรมา\-ทาย เยน เวสาลี เตน ปกฺกามิ, อนุปุพฺเพน เยน เวสาลี ตทวสริ. ตตฺร สุทํ อายสฺมา สุทินฺโน เวสาลิยํ วิหรติ มหาวเน กูฏาคาร\-สาลายํ. อสฺโสสุํ โข อายสฺมโต สุทินฺนสฺส ญาตกา “สุทินฺโน กิร กลนฺทปุตฺโต เวสาลิํ อนุปฺปตฺโต”ติ. เต อายสฺมโต สุทินฺนสฺส สฏฺฐิมตฺเต ถาลิปาเก ภตฺตาภิหารํ อภิหริํสุ. อถ โข อายสฺมา สุทินฺโน เต สฏฺฐิมตฺเต ถาลิปาเก ภิกฺขูนํ วิสฺสชฺเชตฺวา ปุพฺพณฺหสมยํ นิวาเสตฺวา ปตฺต\-จีวรมา\-ทาย กลนฺทคามํ ปณฺฑาย ปาวิสิ, กลนฺทคาเม สปทานํ ปิณฺฑาย จรมาโน เยน สกปิตุ นิเวสนํ เตนุปสงฺกมิ.}

\proseitem{31}{เตน โข ปน สมเยน อายสฺมโต สุทินฺนสฺส ญาติทาสี อาภิโทสิกํ กุมฺมาสํ ฉฑฺเฑตุกามา\csromansharedfootnote{3a9d1ee5ae122c84}{ฉฏฺเฏตุกามา (ก)} โหติ. อถ โข อายสฺมา สุทินฺโน ตํ ญาติทาสิํ เอตทโวจ “สเจ ตํ ภคินิ ฉฑฺฑนีย\-ธมฺมํ, อิธ เม ปตฺเต อากิรา”ติ. อถ โข อายสฺมโต สุทินฺนสฺส ญาติทาสี ตํ \csromanfolio{18}อาภิโทสิกํ กุมฺมาสํ อายสฺมโต สุทินฺนสฺส ปตฺเต อากิรนฺตี หตฺถานญฺจ ปาทานญฺจ สรสฺส จ นิมิตฺตํ อคฺคเหสิ. อถ โข อายสฺมโต สุทินฺนสฺส ญาติทาสี เยนายสฺมโต สุทินฺนสฺส มาตา เตนุปสงฺกมิ, อุปสงฺกมิตฺวา อายสฺมโต สุทินฺนสฺส มาตรํ เอตทโวจ “ยคฺเฆยฺเย ชาเนยฺยาสิ อยฺยปุตฺโต สุทินฺโน อนุปฺปตฺโต”ติ. สเจ เช ตฺวํ สจฺจํ ภณสิ, อทาสิํ ตํ กโรมีติ.}

\proseitem{32}{เตน โข ปน สมเยน อายสฺมา สุทินฺโน ตํ อาภิโทสิกํ กุมฺมาสํ อญฺญตรํ กุฏฺฏมูลํ\csromansharedfootnote{b6ce0971fff71f0b}{กุฑฺฑมูลํ (สี, สฺยา)} นิสฺสาย ปริภุญฺชติ. ปิตาปิ โข อายสฺมโต สุทินฺนสฺส กมฺมนฺตา อาคจฺฉนฺโต อทฺทส อายสฺมนฺตํ สุทินฺนํ ตํ อาภิโทสิกํ กุมฺมาสํ อญฺญตรํ กุฏฺฏมูลํ นิสฺสาย ปริภุญฺชนฺตํ, ทิสฺวาน เยนายสฺมา สุทินฺโน เตนุปสงฺกมิ, อุปสงฺกมิตฺวา อายสฺมนฺตํ สุทินฺนํ เอตทโวจ “อตฺถิ นาม ตาต สุทินฺน อาภิโทสิกํ กุมฺมาสํ ปริภุญฺชิสฺสสิ, นนุ นาม ตาต สุทินฺน สกํ เคหํ คนฺตพฺพนฺ”ติ. อคมิมฺห\csromansharedfootnote{9e14c191e176f64a}{อคมมฺหา (ก)} โข เต คหปติ เคหํ, ตตายํ อาภิโทสิโก กุมฺมาโสติ. อถ โข อายสฺมโต สุทินฺนสฺส ปิตา อายสฺมโต สุทินฺนสฺส พาหายํ คเหตฺวา อายสฺมนฺตํ สุทินฺนํ เอตทโวจ “เอหิ ตาต สุทินฺน ฆรํ คมิสฺสามา”ติ. อถ โข อายสฺมา สุทินฺโน เยน สกปิตุ นิเวสนํ เตนุปสงฺกมิ, อุปสงฺกมิตฺวา ปญฺญตฺเต อาสเน นิสีทิ. อถ โข อายสฺมโต สุทินฺนสฺส ปิตา อายสฺมนฺตํ สุทินฺนํ เอตทโวจ “ภุญฺช ตาต สุทินฺนา”ติ. อลํ คหปติ กตํ เม อชฺช ภตฺตกิจฺจนฺติ. อธิวาเสหิ ตาต สุทินฺน สฺวาตนาย ภตฺตนฺติ. อธิวาเสสิ โข อายสฺมา สุทินฺโน ตุณฺหีภาเวน. อถ โข อายสฺมา สุทินฺโน อุฏฺฐายาสนา ปกฺกามิ.}

\proseitem{33}{อถ โข อายสฺมโต สุทินฺนสฺส มาตา ตสฺสา รตฺติยา อจฺจเยน หริเตน โคมเยน ปถวิํ โอปุญฺชาเปตฺวา\csromansharedfootnote{8fc10202fa63333d}{โอปุญฺฉาเปตฺวา (สี, สฺยา)} เทฺว ปุญฺเช การาเปสิ เอกํ หิรญฺญสฺส เอกํ สุวณฺณสฺส, ตาว มหนฺตา ปุญฺชา อเหสุํ โอรโต ฐิโต ปุริโส ปารโต ฐิตํ ปุริสํ น ปสฺสติ, ปารโต ฐิโต ปุริโส โอรโต ฐิตํ ปุริสํ น ปสฺสติ. เต ปุญฺเช กิลญฺเชหิ ปฏิจฺฉาทาเปตฺวา มชฺเฌ อาสนํ ปญฺญาเปตฺวา ติโรกรณียํ ปริกฺขิปิตฺวา \csromanfolio{19}อายสฺมโต สุทินฺนสฺส ปุราณ\-ทุติยิกํ อามนฺเตสิ “เตน หิ วธุ เยน อลงฺกาเรน อลงฺกตา ปุตฺตสฺส เม สุทินฺนสฺส ปิยา อโหสิ มนาปา, เตน อลงฺกาเรน อลงฺกรา”ติ. “เอวํ อยฺเย”ติ โข อายสฺมโต สุทินฺนสฺส ปุราณทุติยิกา อายสฺมโต สุทินฺนสฺส มาตุยา ปจฺจสฺโสสิ.}

\proseitem{34}{อถ โข อายสฺมา สุทินฺโน ปุพฺพณฺหสมยํ นิวาเสตฺวา ปตฺต\-จีวรมา\-ทาย เยน สกปิตุ นิเวสนํ เตนุปสงฺกมิ, อุปสงฺกมิตฺวา ปญฺญตฺเต อาสเน นิสีทิ. อถ โข อายสฺมโต สุทินฺนสฺส ปิตา เยนายสฺมา สุทินฺโน เตนุปสงฺกมิ, อุปสงฺกมิตฺวา เต ปุญฺเช วิวราเปตฺวา อายสฺมนฺตํ สุทินฺนํ เอตทโวจ “อิทํ เต ตาต สุทินฺน มาตุ มตฺติกํ, อิตฺถิกาย อิตฺถิธนํ, อญฺญํ เปตฺติกํ, อญฺญํ ปิตามหํ. ลพฺภา ตาต สุทินฺน หีนายาวตฺติตฺวา โภคา จ ภุญฺชิตุํ ปุญฺญานิ จ กาตุํ, เอหิ ตฺวํ ตาต สุทินฺน หีนายาวตฺติตฺวา โภเค จ ภุญฺชสฺสุ ปุญฺญานิ จ กโรหี”ติ. ตาต น อุสฺสหามิ น วิสหามิ, อภิรโต อหํ พฺรหฺมจริยํ จรามีติ. ทุติยมฺปิ โข ฯเปฯ. ตติยมฺปิ โข อายสฺมโต สุทินฺนสฺส ปิตา อายสฺมนฺตํ สุทินฺนํ เอตทโวจ “อิทํ เต ตาต สุทินฺน มาตุ มตฺติกํ, อิตฺถิกาย อิตฺถิธนํ, อญฺญํ เปตฺติกํ, อญฺญํ ปิตามหํ. ลพฺภา ตาต สุทินฺน หีนายาวตฺติตฺวา โภคา จ ภุญฺชิตุํ ปุญฺญานิ จ กาตุํ, เอหิ ตฺวํ ตาต สุทินฺน หีนายาวตฺติตฺวา โภเค จ ภุญฺชสฺสุ ปุญฺญานิ จ กโรหี”ติ. วเทยฺยาม โข ตํ คหปติ สเจ ตฺวํ นาติกฑฺเฒยฺยาสีติ. วเทหิ ตาต สุทินฺนาติ. เตน หิ ตฺวํ คหปติ มหนฺเต มหนฺเต สาณิปสิพฺพเก การาเปตฺวา หิรญฺญ\-สุวณฺณสฺส ปูราเปตฺวา สกเฏหิ นิพฺพาหาเปตฺวา มชฺเฌ คงฺคาย โสเต โอปาเตหิ\csromansharedfootnote{2308c487b8e7657c}{โอสาเทหิ (สี, สฺยา)}. ตํ กิสฺส เหตุ, ยํ หิ เต คหปติ ภวิสฺสติ ตโตนิทานํ ภยํ วา ฉมฺภิตตฺตํ วา โลมหํโส วา อารกฺโข วา, โส เต น ภวิสฺสตีติ. เอวํ วุตฺเต อายสฺมโต สุทินฺนสฺส ปิตา อนตฺตมโน อโหสิ “กถํ หิ นาม ปุตฺโต สุทินฺโน เอวํ วกฺขตี”ติ.}

\proseitem{35}{อถ โข อายสฺมโต สุทินฺนสฺส ปิตา อายสฺมโต สุทินฺนสฺส ปุราณ\-ทุติยิกํ อามนฺเตสิ “เตน หิ วธุ ตฺวํ ปิยา จ มนาปา จ\csromansharedfootnote{327e6c16243ad5b7}{ตฺวมฺปิ ยาจ (สี)}, \csromanfolio{20}อปฺเปว นาม ปุตฺโต สุทินฺโน ตุยฺหมฺปิ วจนํ กเรยฺยา”ติ. อถ โข อายสฺมโต สุทินฺนสฺส ปุราณทุติยิกา อายสฺมโต สุทินฺนสฺส ปาเทสุ คเหตฺวา อายสฺมนฺตํ สุทินฺนํ เอตทโวจ “กีทิสา นาม ตา อยฺยปุตฺต อจฺฉราโย, ยาสํ ตฺวํ เหตุ พฺรหฺมจริยํ จรสี”ติ. น โข อหํ ภคินิ อจฺฉรานํ เหตุ พฺรหฺมจริยํ จรามีติ. อถ โข อายสฺมโต สุทินฺนสฺส ปุราณทุติยิกา “อชฺชตคฺเค มํ อยฺยปุตฺโต สุทินฺโน ภคินิวาเทน สมุทาจรตี”ติ ตตฺเถว มุจฺฉิตา ปปตา.}

\prose{อถ โข อายสฺมา สุทินฺโน ปิตรํ เอตทโวจ “สเจ คหปติ โภชนํ ทาตพฺพํ เทถ, มา โน วิเหฐยิตฺถา”ติ. ภุญฺช ตาต สุทินฺนาติ. อถ โข อายสฺมโต สุทินฺนสฺส มาตา จ ปิตา จ อายสฺมนฺตํ สุทินฺนํ ปณีเตน ขาทนีเยน โภชนีเยน สหตฺถา สนฺตปฺเปสุํ สมฺปวาเรสุํ. อถ โข อายสฺมโต สุทินฺนสฺส มาตา อายสฺมนฺตํ สุทินฺนํ ภุตฺตาวิํ โอนีต\-ปตฺต\-ปาณิํ เอตทโวจ “อิทํ ตาต สุทินฺน กุลํ อฑฺฒํ มหทฺธนํ มหาโภคํ ปหูต\-ชาตรูป\-รชตํ ปหูต\-วิตฺตู\-ปกรณํ ปหูต\-ธน\-ธญฺญํ, ลพฺภา ตาต สุทินฺน หีนายาวตฺติตฺวา โภคา จ ภุญฺชิตุํ ปุญฺญานิ จ กาตุํ, เอหิ ตฺวํ ตาต สุทินฺน หีนายาวตฺติตฺวา โภเค จ ภุญฺชสฺสุ ปุญฺญานิ จ กโรหี”ติ. อมฺม น อุสฺสหามิ น วิสหามิ, อภิรโต อหํ พฺรหฺมจริยํ จรามีติ. ทุติยมฺปิ โข ฯเปฯ. ตติยมฺปิ โข อายสฺมโต สุทินฺนสฺส มาตา อายสฺมนฺตํ สุทินฺนํ เอตทโวจ “อิทํ ตาต สุทินฺน กุลํ อฑฺฒํ มหทฺธนํ มหาโภคํ ปหูต\-ชาตรูป\-รชตํ ปหูต\-วิตฺตู\-ปกรณํ ปหูต\-ธน\-ธญฺญํ\csromansharedfootnote{fd294af43851c355}{ปหูตธนธญฺญํ (ฯเปฯ จรามีติ) อิติปาโฐ สพฺพตฺถ นตฺถิ, อูโน มญฺเญ.}, เตน หิ ตาต สุทินฺน พีชกมฺปิ เทหิ, มา โน อปุตฺตกํ สาปเตยฺยํ ลิจฺฉวโย อติหราเปสุนฺ”ติ. เอตํ โข เม อมฺม สกฺกา กาตุนฺติ. กหํ ปน ตาต สุทินฺน เอตรหิ วิหรสีติ. มหาวเน อมฺมาติ. อถ โข อายสฺมา สุทินฺโน อุฏฺฐายาสนา ปกฺกามิ.}

\proseitem{36}{อถ โข อายสฺมโต สุทินฺนสฺส มาตา อายสฺมโต สุทินฺนสฺส ปุราณ\-ทุติยิกํ อามนฺเตสิ “เตน หิ วธุ ยทา อุตุนี โหสิ, ปุปฺผํ เต อุปฺปนฺนํ โหติ, อถ เม อาโรเจยฺยาสี”ติ. “เอวํ อยฺเย”ติ โข อายสฺมโต สุทินฺนสฺส ปุราณทุติยิกา อายสฺมโต สุทินฺนสฺส มาตุยา ปจฺจสฺโสสิ. อถ โข อายสฺมโต สุทินฺนสฺส ปุราณทุติยิกา \csromanfolio{21}นจิรสฺเสว อุตุนี อโหสิ, ปุปฺผํสา อุปฺปชฺชิ. อถ โข อายสฺมโต สุทินฺนสฺส ปุราณทุติยิกา อายสฺมโต สุทินฺนสฺส มาตรํ เอตทโวจ “อุตุนีมฺหิ อยฺเย ปุปฺผํ เม อุปฺปนฺนนฺ”ติ. เตน หิ วธุ เยน อลงฺกาเรน อลงฺกตา ปุตฺตสฺส สุทินฺนสฺส ปิยา อโหสิ มนาปา, เตน อลงฺกาเรน อลงฺกราติ. “เอวํ อยฺเย”ติ โข อายสฺมโต สุทินฺนสฺส ปุราณทุติยิกา อายสฺมโต สุทินฺนสฺส มาตุยา ปจฺจสฺโสสิ. อถ โข อายสฺมโต สุทินฺนสฺส มาตา อายสฺมโต สุทินฺนสฺส ปุราณ\-ทุติยิกํ อาทาย เยน มหาวนํ เยนายสฺมา สุทินฺโน เตนุปสงฺกมิ, อุปสงฺกมิตฺวา อายสฺมนฺตํ สุทินฺนํ เอตทโวจ “อิทํ ตาต สุทินฺน กุลํ อฑฺฒํ มหทฺธนํ มหาโภคํ ปหูต\-ชาตรูป\-รชตํ ปหูต\-วิตฺตู\-ปกรณํ ปหูต\-ธน\-ธญฺญํ, ลพฺภา ตาต สุทินฺน หีนายาวตฺติตฺวา โภคา จ ภุญฺชิตุํ ปุญฺญานิ จ กาตุํ, เอหิ ตฺวํ ตาต สุทินฺน หีนายาวตฺติตฺวา โภเค จ ภุญฺชสฺสุ ปุญฺญานิ จ กโรหี”ติ. อมฺม น อุสฺสหามิ น วิสหามิ, อภิรโต อหํ พฺรหฺมจริยํ จรามีติ. ทุติยมฺปิ โข ฯเปฯ. ตติยมฺปิ โข อายสฺมโต สุทินฺนสฺส มาตา อายสฺมนฺตํ สุทินฺนํ เอตทโวจ “อิทํ ตาต สุทินฺน กุลํ อฑฺฒํ มหทฺธนํ มหาโภคํ ปหูต\-ชาตรูป\-รชตํ ปหูต\-วิตฺตู\-ปกรณํ ปหูต\-ธน\-ธญฺญํ, เตน หิ ตาต สุทินฺน พีชกํปิ เทหิ, มา โน อปุตฺตกํ สาปเตยฺยํ ลิจฺฉวโย อติหราเปสุนฺ”ติ. เอตํ โข เม อมฺม สกฺกา กาตุนฺติ ปุราณ\-ทุติยิกาย พาหายํ คเหตฺวา มหาวนํ อชฺโฌคาเหตฺวา อปญฺญตฺเต สิกฺขาปเท อนาทีนวทสฺโส ปุราณ\-ทุติยิกาย ติกฺขตฺตุํ เมถุนํ ธมฺมํ อภิวิญฺญาเปสิ. สา เตน คพฺภํ คณฺหิ. ภุมฺมา เทวา สทฺทม\-นุสฺสาเวสุํ “นิรพฺพุโท วต โภ ภิกฺขุสํโฆ นิราทีนโว, สุทินฺเนน กลนฺทปุตฺเตน อพฺพุทํ อุปฺปาทิตํ, อาทีนโว อุปฺปาทิโต”ติ. ภุมฺมานํ เทวานํ สทฺทํ สุตฺวา จาตุมหาราชิกา\csromansharedfootnote{68e067275c9c308d}{จาตุมฺมหาราชิกา (สี, สฺยา)} เทวา สทฺทม\-นุสฺสาเวสุํ ฯเปฯ. ตาวติํสา เทวา. ยามา เทวา. ตุสิตา เทวา. นิมฺมานรตี เทวา. ปรนิมฺมิต\-วส\-วตฺตี เทวา. พฺรหฺมกายิกา เทวา สทฺทม\-นุสฺสาเวสุํ “นิรพฺพุโท วต โภ ภิกฺขุสํโฆ นิราทีนโว, สุทินฺเนน กลนฺทปุตฺเตน อพฺพุทํ อุปฺปาทิตํ, อาทีนโว อุปฺปาทิโต”ติ. อิติห เตน ขเณน เตน มุหุตฺเตน ยาว พฺรหฺมโลกา สทฺโท อพฺภุคฺคจฺฉิ.}

\prose{\csromanfolio{22}อถ โข อายสฺมโต สุทินฺนสฺส ปุราณทุติยิกา ตสฺส คพฺภสฺส ปริปากม\-นฺวาย ปุตฺตํ วิชายิ. อถ โข อายสฺมโต สุทินฺนสฺส สหายกา ตสฺส ทารกสฺส พีชโกติ นามํ อกํสุ, อายสฺมโต สุทินฺนสฺส ปุราณ\-ทุติยิกาย พีชกมาตาติ นามํ อกํสุ, อายสฺมโต สุทินฺนสฺส พีชกปิตาติ นามํ อกํสุ. เต อปเรน สมเยน อุโภ อคารสฺมา อนคาริยํ ปพฺพชิตฺวา อรหตฺตํ สจฺฉากํสุ.}

\proseitem{37}{อถ โข อายสฺมโต สุทินฺนสฺส อหุเทว กุกฺกุจฺจํ อหุ วิปฺปฏิสาโร “อลาภา วต เม, น วต เม ลาภา, ทุลฺลทฺธํ วต เม, น วต เม สุลทฺธํ, โยหํ เอวํ สฺวากฺขาเต ธมฺมวินเย ปพฺพชิตฺวา นาสกฺขิํ ยาวชีวํ ปริปุณฺณํ ปริสุทฺธํ พฺรหฺมจริยํ จริตุนฺ”ติ. โส เตเนว กุกฺกุจฺเจน เตน วิปฺปฏิสาเรน กิโส อโหสิ ลูโข ทุพฺพณฺโณ อุปฺปณฺฑุ\-ปฺปณฺฑุก\-ชาโต ธมนิสนฺถต\-คตฺโต, อนฺโตมโน ลีนมโน ทุกฺขี ทุมฺมโน วิปฺปฏิสารี ปชฺฌายิ.}

\proseitem{38}{อถ โข อายสฺมโต สุทินฺนสฺส สหายกา ภิกฺขู อายสฺมนฺตํ สุทินฺนํ เอตทโวจุํ “ปุพฺเพ โข ตฺวํ อาวุโส สุทินฺน วณฺณวา อโหสิ ปีณินฺทฺริโย ปสนฺน\-มุขวณฺโณ วิปฺปสนฺน\-ฉวิวณฺโณ, โส ทานิ ตฺวํ เอตรหิ กิโส ลูโข ทุพฺพณฺโณ อุปฺปณฺฑุ\-ปฺปณฺฑุก\-ชาโต ธมนิสนฺถต\-คตฺโต อนฺโตมโน ลีนมโน ทุกฺขี ทุมฺมโน วิปฺปฏิสารี ปชฺฌายสิ, กจฺจิ โน ตฺวํ อาวุโส สุทินฺน อนภิรโต พฺรหฺมจริยํ จรสี”ติ. น โข อหํ อาวุโส อนภิรโต พฺรหฺมจริยํ จรามิ, อตฺถิ เม ปาปกมฺมํ กตํ, ปุราณ\-ทุติยิกาย เมถุโน ธมฺโม ปฏิเสวิโต, ตสฺส มยฺหํ อาวุโส อหุเทว กุกฺกุจฺจํ, อหุ วิปฺปฏิสาโร “อลาภา วต เม, น วต เม ลาภา, ทุลฺลทฺธํ วต เม, น วต เม สุลทฺธํ, โยหํ เอวํ สฺวากฺขาเต ธมฺมวินเย ปพฺพชิตฺวา นาสกฺขิํ ยาวชีวํ ปริปุณฺณํ ปริสุทฺธํ พฺรหฺมจริยํ จริตุนฺ”ติ. อลํ หิ เต อาวุโส สุทินฺน กุกฺกุจฺจาย อลํ วิปฺปฏิสาราย, ยํ ตฺวํ เอวํ สฺวากฺขาเต ธมฺมวินเย ปพฺพชิตฺวา น สกฺขิสฺสสิ ยาวชีวํ ปริปุณฺณํ ปริสุทฺธํ พฺรหฺมจริยํ จริตุํ. นนุ อาวุโส ภควตา อเนก\-ปริยาเยน วิราคาย ธมฺโม เทสิโต โน สราคาย, วิสํโยคาย ธมฺโม เทสิโต โน สํโยคาย, อนุปาทานาย ธมฺโม เทสิโต โน สฺเอาปาทานาย. ตตฺถ นาม ตฺวํ อาวุโส \csromanfolio{23}ภควตา วิราคาย ธมฺเม เทสิเต สราคาย เจเตสฺสสิ, วิสํโยคาย ธมฺเม เทสิเต สํโยคาย เจเตสฺสสิ, อนุปาทานาย ธมฺเม เทสิเต สฺเอาปาทานาย เจเตสฺสสิ. นนุ อาวุโส ภควตา อเนก\-ปริยาเยน ราควิราคาย ธมฺโม เทสิโต, มทนิมฺมทนาย ปิปาสวินยาย อาลย\-สมุคฺฆาตาย วฏฺฏุ\-ปจฺเฉทาย ตณฺหากฺขยาย วิราคาย นิโรธาย นิพฺพานาย ธมฺโม เทสิโต, นนุ อาวุโส ภควตา อเนก\-ปริยาเยน กามานํ ปหานํ อกฺขาตํ, กามสญฺญานํ ปริญฺญา อกฺขาตา, กามปิปาสานํ ปฏิวินโย อกฺขาโต, กามวิตกฺกานํ สมุคฺฆาโต อกฺขาโต, กาม\-ปริฬาหานํ วูปสโม อกฺขาโต. เนตํ อาวุโส อปฺปสนฺนานํ วา ปสาทาย ปสนฺนานํ วา ภิยฺโยภาวาย, อถ เขฺวตํ อาวุโส อปฺปสนฺนาน\-ญฺเจว อปฺปสาทาย ปสนฺนานญฺจ เอกจฺจานํ อญฺญถตฺตายาติ.}

\proseitem{39}{อถ โข เต ภิกฺขู อายสฺมนฺตํ สุทินฺนํ อเนก\-ปริยาเยน วิครหิตฺวา ภควโต เอตมตฺถํ อาโรเจสุํ. อถ โข ภควา เอตสฺมิํ นิทาเน เอตสฺมิํ ปกรเณ ภิกฺขุสํฆํ สนฺนิปาตาเปตฺวา อายสฺมนฺตํ สุทินฺนํ ปฏิปุจฺฉิ “สจฺจํ กิร ตฺวํ สุทินฺน ปุราณ\-ทุติยิกาย เมถุนํ ธมฺมํ ปฏิเสวี”ติ. สจฺจํ ภควาติ. วิครหิ พุทฺโธ ภควา “อนนุจฺฉวิกํ\csromansharedfootnote{95846e3080cc32bb}{อนนุจฺฉวิยํ (สี)} โมฆปุริส อนนุโลมิกํ อปฺปติรูปํ อสฺสามณกํ อกปฺปิยํ อกรณียํ, กถํ หิ นาม ตฺวํ โมฆปุริส เอวํ สฺวากฺขาเต ธมฺมวินเย ปพฺพชิตฺวา น สกฺขิสฺสสิ ยาวชีวํ ปริปุณฺณํ ปริสุทฺธํ พฺรหฺมจริยํ จริตุํ. นนุ มยา โมฆปุริส อเนก\-ปริยาเยน วิราคาย ธมฺโม เทสิโต โน สราคาย, วิสํโยคาย ธมฺโม เทสิโต โน สํโยคาย, อนุปาทานาย ธมฺโม เทสิโต โน สฺเอาปาทานาย. ตตฺถ นาม ตฺวํ โมฆปุริส มยา วิราคาย ธมฺเม เทสิเต สราคาย เจเตสฺสสิ, วิสํโยคาย ธมฺเม เทสิเต สํโยคาย เจเตสฺสสิ, อนุปาทานาย ธมฺเม เทสิเต สฺเอาปาทานาย เจเตสฺสสิ. นนุ มยา โมฆปุริส อเนก\-ปริยาเยน ราควิราคาย ธมฺโม เทสิโต, มทนิมฺมทนาย ปิปาสวินยาย อาลย\-สมุคฺฆาตาย วฏฺฏุ\-ปจฺเฉทาย ตณฺหากฺขายาย วิราคาย นิโรธาย นิพฺพานาย ธมฺโม เทสิโต, นนุ มยา โมฆปุริส อเนก\-ปริยาเยน \csromanfolio{24}\textbf{กามานํ ปหานํ อกฺขาตํ, กามสญฺญานํ ปริญฺญา อกฺขาตา,} \textbf{กามปิปาสานํ ปฏิวินโย อกฺขาโต, กามวิตกฺกานํ สมุคฺฆาโต อกฺขาโต,} \textbf{กาม}\-\textbf{ปริฬาหานํ วูปสโม อกฺขาโต. วรํ เต โมฆปุริส อาสิวิสสฺส}\csromansharedfootnote{8e7009d1cf91e47d}{อาสีวิสสฺส (สี, สฺยา)} โฆรวิสสฺส มุเข องฺคชาตํ ปกฺขิตฺตํ, น เตฺวว มาตุคามสฺส องฺคชาเต องฺคชาตํ ปกฺขิตฺตํ. วรํ เต โมฆปุริส กณฺหสปฺปสฺส มุเข องฺคชาตํ ปกฺขิตฺตํ, น เตฺวว มาตุคามสฺส องฺคชาเต องฺคชาตํ ปกฺขิตฺตํ. วรํ เต โมฆปุริส องฺคารกาสุยา อาทิตฺตาย สมฺปชฺชลิตาย สโชติภูตาย องฺคชาตํ ปกฺขิตฺตํ, น เตฺวว มาตุคามสฺส องฺคชาเต องฺคชาตํ ปกฺขิตฺตํ. ตํ กิสฺส เหตุ, ตโตนิทานํ หิ โมฆปุริส มรณํ วา นิคจฺเฉยฺย มรณมตฺตํ วา ทุกฺขํ, น เตฺวว ตปฺปจฺจยา กายสฺส เภทา ปรํ มรณา อปายํ ทุคฺคติํ วินิปาตํ นิรยํ อุปปชฺเชยฺย, อิโตนิทานญฺจ โข โมฆปุริส กายสฺส เภทา ปรํ มรณา อปายํ ทุคฺคติํ วินิปาตํ นิรยํ อุปปชฺเชยฺย. ตตฺถ นาม ตฺวํ โมฆปุริส ยํ ตฺวํ อสทฺธมฺมํ คามธมฺมํ วสลธมฺมํ ทุฏฺฐุลฺลํ โอทกนฺติกํ รหสฺสํ ทฺวยํ\-ทฺวย\-สมาปตฺติํ สมาปชฺชิสฺสสิ. พหูนํ โข ตฺวํ โมฆปุริส อกุสลานํ ธมฺมานํ อาทิกตฺตา ปุพฺพงฺคโม. เนตํ โมฆปุริส อปฺปสนฺนานํ วา ปสาทาย ปสนฺนานํ วา ภิยฺโยภาวาย, อถ เขฺวตํ โมฆปุริส อปฺปสนฺนาน\-ญฺเจว อปฺปสาทาย ปสนฺนานญฺจ เอกจฺจานํ อญฺญถตฺตายา”ติ.}

\prose{อถ โข ภควา อายสฺมนฺตํ สุทินฺนํ อเนก\-ปริยาเยน วิครหิตฺวา ทุพฺภรตาย ทุปฺโปสตาย มหิจฺฉตาย อสนฺตุฏฺฐิตาย\csromansharedfootnote{32b6fcac3b6425c8}{อสนฺตุฏฺฐตาย (สฺยา)} สงฺคณิกาย โกสชฺชสฺส อวณฺณํ ภาสิตฺวา อเนก\-ปริยาเยน สุภรตาย สุโปสตาย อปฺปิจฺฉสฺส สนฺตุฏฺฐสฺส สลฺเลขสฺส ธุตสฺส ปาสาทิกสฺส อปจยสฺส วีริยา\-รมฺภสฺส\csromansharedfootnote{82452c5bb53baf2b}{วีริยารพฺภสฺส (ก)} วณฺณํ ภาสิตฺวา ภิกฺขูนํ ตทนุจฺฉวิกํ ตทนุโลมิกํ ธมฺมิํ กถํ กตฺวา ภิกฺขู อามนฺเตสิ\csromandash{} เตน หิ ภิกฺขเว ภิกฺขุนํ สิกฺขาปทํ ปญฺญเปสฺสามิ\csromansharedfootnote{057f462d8f730225}{ปญฺญาเปสฺสามิ (สี, สฺยา)} ทส อตฺถวเส ปฏิจฺจ สํฆสุฏฺฐุตาย สํฆผาสุตาย ทุมฺมงฺกูนํ ปุคฺคลานํ นิคฺคหาย เปสลานํ ภิกฺขูนํ ผาสุวิหาราย ทิฏฺฐ\-ธมฺมิกานํ อาสวานํ สํวราย สมฺปรายิกานํ อาสวานํ ปฏิฆาตาย อปฺปสนฺนานํ ปสาทาย ปสนฺนานํ ภิยฺโยภาวาย สทฺธมฺม\-ฏฺฐิติยา วินยา\-นุคฺคหาย. เอวญฺจ ปน ภิกฺขเว อิมํ สิกฺขาปทํ อุทฺทิเสยฺยาถ\csromandash{}}

\csromanmatikaanchor{mtk.9}
\csromantocmarkref{subsection}{ปฐมปญฺญตฺติ}{mtk.9}
\bookmark[dest={mtk.9},level=2]{ปฐมปญฺญตฺติ}
\prose{\csromanfolio{25}\textbf{“โย ปน ภิกฺขุ เมถุนํ ธมฺมํ ปฏิเสเวยฺย, ปาราชิโก โหติ อสํวาโส”}ติ.}

\prosewithcloserruled{เอวญฺจิทํ ภควตา ภิกฺขูนํ สิกฺขาปทํ ปญฺญตฺตํ โหติ.}{สุทินฺน\-ภาณวาโร นิฏฺฐิโต.}
\csromanmatikaanchor{mtk.10}
\csromantocmarkref{subsection}{มกฺกฏีวตฺถุ, ปฐมานุปญฺญตฺติ}{mtk.10}
\bookmark[dest={mtk.10},level=2]{มกฺกฏีวตฺถุ, ปฐมานุปญฺญตฺติ}
\csromanheader{2}{\textbf{มกฺกฏีวตฺถุ}}
\proseitem{40}{เตน โข ปน สมเยน อญฺญตโร ภิกฺขุ เวสาลิยํ มหาวเน มกฺกฏิํ อามิเสน อุปลาเปตฺวา ตสฺสา เมถุนํ ธมฺมํ ปฏิเสวติ. อถ โข โส ภิกฺขุ ปุพฺพณฺหสมยํ นิวาเสตฺวา ปตฺตจีวรํ อาทาย เวสาลิํ ปิณฺฑาย ปาวิสิ. เตน โข ปน สมเยน สมฺพหุลา ภิกฺขู เสนาสนจาริกํ อาหิณฺฑนฺตา เยน ตสฺส ภิกฺขุโน วิหาโร เตนุ\-ปสงฺกมิํสุ, อทฺทส โข สา มกฺกฏี เต ภิกฺขู ทูรโตว อาคจฺฉนฺเต, ทิสฺวาน เยน เต ภิกฺขู เตนุปสงฺกมิ, อุปสงฺกมิตฺวา เตสํ ภิกฺขูนํ ปุรโต กฏิมฺปิ จาเลสิ, เฉปฺปมฺปิ จาเลสิ, กฏิมฺปิ โอฑฺฑิ, นิมิตฺตมฺปิ อกาสิ. อถ โข เตสํ ภิกฺขูนํ เอตทโหสิ “นิสฺสํสยํ โข โส ภิกฺขุ อิมิสฺสา มกฺกฏิยา เมถุนํ ธมฺมํ ปฏิเสวตี”ติ เอกมนฺตํ นิลียิํสุ. อถ โข โส ภิกฺขุ เวสาลิยํ ปิณฺฑาย จริตฺวา ปิณฺฑปาตํ อาทาย ปฏิกฺกมิ.}

\proseitem{41}{อถ โข สา มกฺกฏี เยน โส ภิกฺขุ เตนุปสงฺกมิ. อถ โข โส ภิกฺขุ ตํ ปิณฺฑปาตํ เอกเทสํ ภุญฺชิตฺวา เอกเทสํ ตสฺสา มกฺกฏิยา อทาสิ. อถ โข สา มกฺกฏี ตํ ปิณฺฑปาตํ ภุญฺชิตฺวา ตสฺส ภิกฺขุโน กฏิํ โอฑฺฑิ. อถ โข โส ภิกฺขุ ตสฺสา มกฺกฏิยา เมถุนํ ธมฺมํ ปฏิเสวติ. อถ โข เต ภิกฺขู ตํ ภิกฺขุํ เอตทโวจุํ “นนุ อาวุโส ภควตา สิกฺขาปทํ ปญฺญตฺตํ, กิสฺส ตฺวํ อาวุโส มกฺกฏิยา เมถุนํ ธมฺมํ ปฏิเสวสี”ติ. สจฺจํ อาวุโส ภควตา สิกฺขาปทํ ปญฺญตฺตํ, ตญฺจ โข \csromanfolio{26}\textbf{มนุสฺสิตฺถิยา โน ติรจฺฉาน}\-\textbf{คตา}\-\textbf{ยาติ. นนุ อาวุโส ตเถว ตํ โหติ,} \textbf{อนนุจฺฉวิกํ อาวุโส อนนุโลมิกํ อปฺปติรูปํ อสฺสามณกํ} \textbf{อกปฺปิยํ อกรณียํ, กถํ หิ นาม ตฺวํ อาวุโส เอวํ สฺวากฺขาเต} ธมฺมวินเย ปพฺพชิตฺวา น สกฺขิสฺสสิ ยาวชีวํ ปริปุณฺณํ ปริสุทฺธํ พฺรหฺมจริยํ จริตุํ. นนุ อาวุโส ภควตา อเนก\-ปริยาเยน วิราคาย ธมฺโม เทสิโต โน สราคาย ฯเปฯ กาม\-ปริฬาหานํ วูปสโม อกฺขาโต. เนตํ อาวุโส อปฺปสนฺนานํ วา ปสาทาย ปสนฺนานํ วา ภิยฺโยภาวาย, อถเขฺวตํ\csromansharedfootnote{318eb6d047e97acd}{อถ เขฺวตํ เทฺว ปทานิ \csromandash{} ม.พ.ป.} อาวุโส อปฺปสนฺนาน\-ญฺเจว อปฺปสาทาย ปสนฺนานญฺจ เอกจฺจานํ อญฺญถตฺตายาติ. อถ โข เต ภิกฺขู ตํ ภิกฺขุํ อเนก\-ปริยาเยน วิครหิตฺวา ภควโต เอตมตฺถํ อาโรเจสุํ.}

\proseitem{42}{อถ โข ภควา เอตสฺมิํ นิทาเน เอตสฺมิํ ปกรเณ ภิกฺขุสํฆํ สนฺนิปาตาเปตฺวา ตํ ภิกฺขุํ ปฏิปุจฺฉิ “สจฺจํ กิร ตฺวํ ภิกฺขุ มกฺกฏิยา เมถุนํ ธมฺมํ ปฏิเสวี”ติ. สจฺจํ ภควาติ. วิครหิ พุทฺโธ ภควา “อนนุจฺฉวิกํ โมฆปุริส อนนุโลมิกํ อปฺปติรูปํ อสฺสามณกํ อกปฺปิยํ อกรณียํ, กถํ หิ นาม ตฺวํ โมฆปุริส เอวํ สฺวากฺขาเต ธมฺมวินเย ปพฺพชิตฺวา น สกฺขิสฺสสิ ยาวชีวํ ปริปุณฺณํ ปริสุทฺธํ พฺรหฺมจริยํ จริตุํ. นนุ มยา โมฆปุริส อเนก\-ปริยาเยน วิราคาย ธมฺโม เทสิโต โน สราคาย ฯเปฯ กาม\-ปริฬาหานํ วูปสโม อกฺขาโต. วรํ เต โมฆปุริส อาสิวิสสฺส โฆรวิสสฺส มุเข องฺคชาตํ ปกฺขิตฺตํ, น เตฺวว มกฺกฏิยา องฺคชาเต องฺคชาตํ ปกฺขิตฺตํ. วรํ เต โมฆปุริส กณฺหสปฺปสฺส มุเข องฺคชาตํ ปกฺขิตฺตํ, น เตฺวว มกฺกฏิยา องฺคชาเต องฺคชาตํ ปกฺขิตฺตํ. วรํ เต โมฆปุริส องฺคารกาสุยา อาทิตฺตาย สมฺปชฺชลิตาย สโชติภูตาย องฺคชาตํ ปกฺขิตฺตํ, น เตฺวว มกฺกฏิยา องฺคชาเต องฺคชาตํ ปกฺขิตฺตํ. ตํ กิสฺส เหตุ, ตโตนิทานํ หิ โมฆปุริส มรณํ วา นิคจฺเฉยฺย มรณมตฺตํ วา ทุกฺขํ, น เตฺวว ตปฺปจฺจยา กายสฺส เภทา ปรํ มรณา อปายํ ทุคฺคติํ วินิปาตํ นิรยํ อุปปชฺเชยฺย, อิโตนิทานญฺจ โข โมฆปุริส กายสฺส เภทา ปรํ มรณา อปายํ ทุคฺคติํ วินิปาตํ นิรยํ อุปปชฺเชยฺย. ตตฺถ นาม ตฺวํ โมฆปุริส ยํ ตฺวํ อสทฺธมฺมํ คามธมฺมํ วสลธมฺมํ ทุฏฺฐุลฺลํ โอทกนฺติกํ รหสฺสํ ทฺวยํ\-ทฺวย\-สมาปตฺติํ สมาปชฺชิสฺสสิ. เนตํ โมฆปุริส อปฺปสนฺนานํ วา ปสาทาย ฯเปฯ. เอวญฺจ ปน ภิกฺขเว อิมํ สิกฺขาปทํ อุทฺทิเสยฺยาถ\csromandash{}}

\prose{\csromanfolio{27}\textbf{“โย ปน ภิกฺขุ เมถุนํ ธมฺมํ ปฏิเสเวยฺย, อนฺตมโส} \textbf{ติรจฺฉาน}\-\textbf{คตายปิ, ปาราชิโก โหติ อสํวาโส”}ติ.}

\prosewithcloserruled{เอวญฺจิทํ ภควตา ภิกฺขูนํ สิกฺขาปทํ ปญฺญตฺตํ โหติ.}{มกฺกฏีวตฺถุ นิฏฺฐิตํ.}
\csromanmatikaanchor{mtk.11}
\csromantocmarkref{subsection}{สนฺถตภาณวาร, ทุติยานุปญฺญตฺติ}{mtk.11}
\bookmark[dest={mtk.11},level=2]{สนฺถตภาณวาร, ทุติยานุปญฺญตฺติ}
\csromanheader{2}{\textbf{สนฺถต}\-\textbf{ภาณวาร}}
\proseitem{43}{เตน โข ปน สมเยน สมฺพหุลา เวสาลิกา วชฺชิปุตฺตกา ภิกฺขู ยาวทตฺถํ ภุญฺชิํสุ, ยาวทตฺถํ สุปิํสุ, ยาวทตฺถํ นฺหายิํสุ, ยาวทตฺถํ ภุญฺชิตฺวา ยาวทตฺถํ สุปิตฺวา ยาวทตฺถํ นฺหายิตฺวา อโยนิโส มนสิ กริตฺวา สิกฺขํ อปจฺจกฺขาย ทุพฺพลฺยํ อนาวิกตฺวา เมถุนํ ธมฺมํ ปฏิเสวิํสุ. เต อปเรน สมเยน ญาติพฺยสเนนปิ ผุฏฺฐา โภคพฺยสเนนปิ ผุฏฺฐา โรคพฺยสเนนปิ ผุฏฺฐา อายสฺมนฺตํ อานนฺทํ อุปสงฺกมิตฺวา เอวํ วทนฺติ “น มยํ ภนฺเต อานนฺท พุทฺธครหิโน, น ธมฺมครหิโน, น สํฆครหิโน, อตฺตครหิโน มยํ ภนฺเต อานนฺท อนญฺญครหิโน, มยเมวมฺหา อลกฺขิกา มยํ อปฺปปุญฺญา, เย มยํ เอวํ สฺวากฺขาเต ธมฺมวินเย ปพฺพชิตฺวา นาสกฺขิมฺหา ยาวชีวํ ปริปุณฺณํ ปริสุทฺธํ พฺรหฺมจริยํ จริตุํ, อิทานิ เจปิ\csromansharedfootnote{e7f80ee95cd8c727}{อิทานิปิ เจ (สฺยา)} มยํ ภนฺเต อานนฺท ลเภยฺยาม ภควโต สนฺติเก ปพฺพชฺชํ ลเภยฺยาม อุปสมฺปทํ, อิทานิปิ มยํ วิปสฺสกา กุสลานํ ธมฺมานํ ปุพฺพรตฺตาปรรตฺตํ โพธิปกฺขิกานํ ธมฺมานํ ภาวนา\-นุโยคม\-นุยุตฺตา วิหเรยฺยาม, สาธุ ภนฺเต อานนฺท ภควโต เอตมตฺถํ อาโรเจหี”ติ. เอวมาวุโสติ โข อายสฺมา อานนฺโท เวสาลิกานํ วชฺชิ\-ปุตฺต\-กานํ ปฏิสฺสุณิตฺวา เยน ภควา เตนุปสงฺกมิ, อุปสงฺกมิตฺวา ภควโต เอตมตฺถํ อาโรเจสิ.}

\prose{อฏฺฐานเมตํ อานนฺท อนวกาโส, ยํ ตถาคโต วชฺชีนํ วา วชฺชิ\-ปุตฺต\-กานํ วา การณา สาวกานํ ปาราชิกํ สิกฺขาปทํ ปญฺญตฺตํ สมูหเนยฺยาติ.}

\csromanmatikaanchor{mtk.12}
\csromantocmarkref{subsection}{สิกฺขาปทวิภงฺค, ปทภาชนีย}{mtk.12}
\bookmark[dest={mtk.12},level=2]{สิกฺขาปทวิภงฺค, ปทภาชนีย}
\prose{\csromanfolio{28}อถ โข ภควา เอตสฺมิํ นิทาเน เอตสฺมิํ ปกรเณ ธมฺมิํ กถํ กตฺวา ภิกฺขู อามนฺเตสิ โย ภิกฺขเว\csromansharedfootnote{6dbf2e489a144776}{โย ปน ภิกฺขเว ภิกฺขุ (สี) โย โข ภิกฺขเว ภิกฺขุ (สฺยา)} สิกฺขํ อปจฺจกฺขาย ทุพฺพลฺยํ อนาวิกตฺวา เมถุนํ ธมฺมํ ปฏิเสวติ, โส อาคโต น อุปสมฺ\-ปาเท\-ตพฺโพ. โย จ โข ภิกฺขเว\csromansharedfootnote{be112fcb1a59c69d}{โย จ โข ภิกฺขเว ภิกฺขุ (สี, สฺยา)} สิกฺขํ ปจฺจกฺขาย ทุพฺพลฺยํ อาวิกตฺวา เมถุนํ ธมฺมํ ปฏิเสวติ, โส อาคโต อุปสมฺ\-ปาเท\-ตพฺโพ. เอวญฺจ ปน ภิกฺขเว อิมํ สิกฺขาปทํ อุทฺทิเสยฺยาถ\csromandash{}}

\proseitem{44}{“\textbf{โย ปน ภิกฺขุ ภิกฺขูนํ สิกฺขา}\-\textbf{สาชีว}\-\textbf{สมาปนฺโน สิกฺขํ อปจฺจกฺขาย ทุพฺพลฺยํ อนาวิกตฺวา เมถุนํ ธมฺมํ ปฏิเสเวยฺย, อนฺตมโส ติรจฺฉาน}\-\textbf{คตายปิ, ปาราชิโก โหติ อสํวาโส”}ติ.}

\proseitem{45}{\textbf{โยปนา}ติ\csromansharedfootnote{423cb82fe87b6ac9}{โย ปนาติ เทฺว ปทานิ \csromandash{} ม.พ.ป.} โย ยาทิโส ยถายุตฺโต ยถาชจฺโจ ยถานาโม ยถาโคตฺโต ยถาสีโล ยถาวิหารี ยถาโคจโร เถโร วา นโว วา มชฺฌิโม วา, เอโส วุจฺจติ “โย ปนา”ติ.}

\prose{\csromansymbolfootnote{*}{อภิ 2. 254 ปิฏฺเฐ ฌานวิภงฺเคปิ.}\textbf{ภิกฺขู}ติ ภิกฺขโกติ ภิกฺขุ, ภิกฺขาจริยํ อชฺฌุปคโตติ ภิกฺขุ, ภินฺน\-ปฏ\-ธโรติ ภิกฺขุ, สมญฺญาย ภิกฺขุ, ปฏิญฺญาย ภิกฺขุ, เอหิ ภิกฺขูติ ภิกฺขุ, ตีหิ สรณคมเนหิ อุปสมฺปนฺโนติ ภิกฺขุ, ภโทฺร ภิกฺขุ, สาโร ภิกฺขุ, เสโข ภิกฺขุ, อเสโข ภิกฺขุ, สมคฺเคน สํเฆน ญตฺติ\-จตุตฺเถน กมฺเมน อกุปฺเปน ฐานารเหน อุปสมฺปนฺโนติ ภิกฺขุ. ตตฺร ยฺวายํ ภิกฺขุ สมคฺเคน สํเฆน ญตฺติ\-จตุตฺเถน กมฺเมน อกุปฺเปน ฐานารเหน อุปสมฺปนฺโน, อยํ อิมสฺมิํ อตฺเถ อธิปฺเปโต “ภิกฺขู”ติ.}

\prose{\csromansymbolfootnote{+}{ที 3. 183 ปิฏฺเฐปิ.}\textbf{สิกฺขา}ติ ติสฺโส สิกฺขา อธิสีลสิกฺขา อธิจิตฺต\-สิกฺขา อธิปญฺญา\-สิกฺขา. ตตฺร ยายํ อธิสีลสิกฺขา, อยํ อิมสฺมิํ อตฺเถ อธิปฺเปตา “สิกฺขา”ติ.}

\prose{\textbf{สาชีวํ} นาม ยํ ภควตา ปญฺญตฺตํ สิกฺขาปทํ, เอตํ สาชีวํ นาม. ตสฺมิํ สิกฺขติ, เตน วุจฺจติ “สาชีว\-สมาปนฺโน”ติ.}

\prose{\csromanfolio{29}\textbf{สิกฺขํ อปจฺจกฺขาย ทุพฺพลฺยํ อนาวิกตฺวา}ติ อตฺถิ ภิกฺขเว ทุพฺพลฺยาวิกมฺม\-ญฺเจว โหติ สิกฺขา จ อปจฺจกฺขาตา, อตฺถิ ภิกฺขเว ทุพฺพลฺยาวิกมฺม\-ญฺเจว โหติ สิกฺขา จ ปจฺจกฺขาตา.}

\prose{กถญฺจ ภิกฺขเว ทุพฺพลฺยาวิกมฺม\-ญฺเจว โหติ สิกฺขา จ อปจฺจกฺขาตา. อิธ ภิกฺขเว ภิกฺขุ อุกฺกณฺฐิโต อนภิรโต สามญฺญา จวิตุกาโม ภิกฺขุภาวํ อฏฺฏียมาโน หรายมาโน ชิคุจฺฉมาโน คิหิภาวํ ปตฺถยมาโน อุปาสกภาวํ ปตฺถยมาโน อารามิกภาวํ ปตฺถยมาโน สามเณรภาวํ ปตฺถยมาโน ติตฺถิยภาวํ ปตฺถยมาโน ติตฺถิยสาวก\-ภาวํ ปตฺถยมาโน อสฺสมณภาวํ ปตฺถยมาโน อสกฺยปุตฺติย\-ภาวํ ปตฺถยมาโน ยํนูนาหํ พุทฺธํ ปจฺจกฺเขยฺยนฺติ วทติ วิญฺญาเปติ. เอวํปิ ภิกฺขเว ทุพฺพลฺยาวิกมฺม\-ญฺเจว โหติ สิกฺขา จ อปจฺจกฺขาตา.}

\prose{อถ วา ปน อุกฺกณฺฐิโต อนภิรโต สามญฺญา จวิตุกาโม ภิกฺขุภาวํ อฏฺฏียมาโน หรายมาโน ชิคุจฺฉมาโน คิหิภาวํ ปตฺถยมาโน ฯเปฯ อสกฺยปุตฺติย\-ภาวํ ปตฺถยมาโน ยนฺนูนาหํ ธมฺมํ ปจฺจกฺเขยฺยนฺติ วทติ วิญฺญาเปติ ฯเปฯ ยนฺนูนาหํ สํฆํ. ยนฺนูนาหํ สิกฺขํ. ยนฺนูนาหํ วินยํ. ยนฺนูนาหํ ปาติโมกฺขํ. ยนฺนูนาหํ อุทฺเทสํ. ยนฺนูนาหํ อุปชฺฌายํ. ยนฺนูนาหํ อาจริยํ. ยนฺนูนาหํ สทฺธิ\-วิหาริกํ. ยนฺนูนาหํ อนฺเตวาสิกํ. ยนฺนูนาหํ สมานุ\-ปชฺฌายกํ. ยนฺนูนาหํ สมานา\-จริยกํ. ยนฺนูนาหํ สพฺรหฺมจาริํ ปจฺจกฺเขยฺยนฺติ วทติ วิญฺญาเปติ. ยนฺนูนาหํ คิหี อสฺสนฺติ วทติ วิญฺญาเปติ. ยนฺนูนาหํ อุปาสโก อสฺสนฺติ. ยนฺนูนาหํ อารามิโก อสฺสนฺติ. ยนฺนูนาหํ สามเณโร อสฺสนฺติ. ยนฺนูนาหํ ติตฺถิโย อสฺสนฺติ. ยนฺนูนาหํ ติตฺถิยสาวโก อสฺสนฺติ. ยนฺนูนาหํ อสฺสมโณ อสฺสนฺติ. ยนฺนูนาหํ อสกฺยปุตฺติโย อสฺสนฺติ วทติ วิญฺญาเปติ. เอวํปิ ภิกฺขเว ทุพฺพลฺยาวิกมฺม\-ญฺเจว โหติ สิกฺขา จ อปจฺจกฺขาตา.}

\proseitem{46}{อถ วา ปน อุกฺกณฺฐิโต อนภิรโต สามญฺญา จวิตุกาโม ภิกฺขุภาวํ อฏฺฏียมาโน หรายมาโน ชิคุจฺฉมาโน คิหิภาวํ ปตฺถยมาโน ฯเปฯ อสกฺยปุตฺติย\-ภาวํ ปตฺถยมาโน ยทิ ปนาหํ พุทฺธํ ปจฺจกฺเขยฺยนฺติ \csromanfolio{30}วทติ วิญฺญาเปติ ฯเปฯ ยทิ ปนาหํ อสกฺยปุตฺติโย อสฺสนฺติ วทติ วิญฺญาเปติ ฯเปฯ อปาหํ พุทฺธํ ปจฺจกฺเขยฺยนฺติ วทติ วิญฺญาเปติ ฯเปฯ อปาหํ อสกฺยปุตฺติโย อสฺสนฺติ วทติ วิญฺญาเปติ ฯเปฯ หนฺทาหํ พุทฺธํ ปจฺจกฺเขยฺยนฺติ วทติ วิญฺญาเปติ ฯเปฯ หนฺทาหํ อสกฺยปุตฺติโย อสฺสนฺติ วทติ วิญฺญาเปติ ฯเปฯ โหติ เม พุทฺธํ ปจฺจกฺเขยฺยนฺติ วทติ วิญฺญาเปติ ฯเปฯ โหติ เม อสกฺยปุตฺติโย อสฺสนฺติ วทติ วิญฺญาเปติ. เอวํปิ ภิกฺขเว ทุพฺพลฺยาวิกมฺม\-ญฺเจว โหติ สิกฺขา จ อปจฺจกฺขาตา.}

\proseitem{47}{อถ วา ปน อุกฺกณฺฐิโต อนภิรโต สามญฺญา จวิตุกาโม ภิกฺขุภาวํ อฏฺฏียมาโน หรายมาโน ชิคุจฺฉมาโน คิหิภาวํ ปตฺถยมาโน ฯเปฯ อสกฺยปุตฺติย\-ภาวํ ปตฺถยมาโน มาตรํ สรามีติ วทติ วิญฺญาเปติ. ปิตรํ สรามีติ วทติ วิญฺญาเปติ. ภาตรํ สรามีติ วทติ วิญฺญาเปติ. ภคินิํ สรามีติ วทติ วิญฺญาเปติ. ปุตฺตํ สรามีติ วทติ วิญฺญาเปติ. ธีตรํ สรามีติ วทติ วิญฺญาเปติ. ปชาปติํ สรมีติ วทติ วิญฺญาเปติ. ญาตเก สรมีติ วทติ วิญฺญาเปติ. มิตฺเต สรามีติ วทติ วิญฺญาเปติ. คามํ สรามีติ วทติ วิญฺญาเปติ. นิคมํ สรามีติ วทติ วิญฺญาเปติ. เขตฺตํ สรามีติ วทติ วิญฺญาเปติ. วตฺถุํ สรามีติ วทติ วิญฺญาเปติ. หิรญฺญํ สรามีติ วทติ วิญฺญาเปติ. สุวณฺณํ สรามีติ วทติ วิญฺญาเปติ. สิปฺปํ สรามีติ วทติ วิญฺญาเปติ. ปุพฺเพ หสิตํ ลปิตํ กีฬิตํ สมนุสฺสรามีติ วทติ วิญฺญาเปติ. เอวํปิ ภิกฺขเว ทุพฺพลฺยาวิกมฺม\-ญฺเจว โหติ สิกฺขา จ อปจฺจกฺขาตา.}

\proseitem{48}{อถ วา ปน อุกฺกณฺฐิโต อนภิรโต สามญฺญา จวิตุกาโม ภิกฺขุภาวํ อฏฺฏียมาโน หรายมาโน ชิคุจฺฉมาโน คิหิภาวํ ปตฺถยมาโน ฯเปฯ อสกฺยปุตฺติย\-ภาวํ ปตฺถยมาโน มาตา เม อตฺถิ, สา มยา โปเสตพฺพาติ วทติ วิญฺญาเปติ. ปิตา เม อตฺถิ, โส มยา โปเสตพฺโพติ วทติ วิญฺญาเปติ. ภาตา เม อตฺถิ, โส มยา โปเสตพฺโพติ วทติ วิญฺญาเปติ. ภคินี เม อตฺถิ, สา มยา โปเสตพฺพาติ วทติ วิญฺญาเปติ. ปุตฺโต เม อตฺถิ, โส มยา โปเสตพฺโพติ วทติ วิญฺญาเปติ. ธีตา เม อตฺถิ, สา มยา โปเสตพฺพาติ วทติ วิญฺญาเปติ. ปชาปติ เม อตฺถิ, สา มยา โปเสตพฺพาติ วทติ \csromanfolio{31}วิญฺญาเปติ. ญาตกา เม อตฺถิ, เต มยา โปเสตพฺพาติ วทติ วิญฺญาเปติ. มิตฺตา เม อตฺถิ, เต มยา โปเสตพฺพาติ วทติ วิญฺญาเปติ. เอวํปิ ภิกฺขเว ทุพฺพลฺยาวิกมฺม\-ญฺเจว โหติ สิกฺขา จ อปจฺจกฺขาตา.}

\proseitem{49}{อถ วา ปน อุกฺกณฺฐิโต อนภิรโต สามญฺญา จวิตุกาโม ภิกฺขุภาวํ อฏฺฏียมาโน หรายมาโน ชิคุจฺฉมาโน คิหิภาวํ ปตฺถยมาโน ฯเปฯ อสกฺยปุตฺติย\-ภาวํ ปตฺถยมาโน มาตา เม อตฺถิ, สา มํ โปเสสฺสตีติ วทติ วิญฺญาเปติ. ปิตา เม อตฺถิ, โส มํ โปเสสฺสตีติ วทติ วิญฺญาเปติ. ภาตา เม อตฺถิ, โส มํ โปเสสฺสตีติ วทติ วิญฺญาเปติ. ภคินี เม อตฺถิ, สา มํ โปเสสฺสตีติ วทติ วิญฺญาเปติ. ปุตฺโต เม อตฺถิ, โส มํ โปเสสฺสตีติ วทติ วิญฺญาเปติ. ธีตา เม อตฺถิ, สา มํ โปเสสฺสตีติ วทติ วิญฺญาเปติ. ปชาปติ เม อตฺถิ, สา มํ โปเสสฺสตีติ วทติ วิญฺญาเปติ. ญาตกา เม อตฺถิ, เต มํ โปเสสฺสนฺตีติ วทติ วิญฺญาเปติ. มิตฺตา เม อตฺถิ, เต มํ โปเสสฺสนฺตีติ วทติ วิญฺญาเปติ. คาโม เม อตฺถิ, เตนาหํ ชีวิสฺสามีติ วทติ วิญฺญาเปติ. นิคโม เม อตฺถิ, เตนาหํ ชีวิสฺสามีติ วทติ วิญฺญาเปติ. เขตฺตํ เม อตฺถิ, เตนาหํ ชีวิสฺสามีติ วทติ วิญฺญาเปติ. วตฺถุ เม อตฺถิ, เตนาหํ ชีวิสฺสามีติ วทติ วิญฺญาเปติ. หิรญฺญํ เม อตฺถิ, เตนาหํ ชีวิสฺสามีติ วทติ วิญฺญาเปติ. สุวณฺณํ เม อตฺถิ, เตนาหํ ชีวิสฺสามีติ วทติ วิญฺญาเปติ. สิปฺปํ เม อตฺถิ, เตนาหํ ชีวิสฺสามีติ วทติ วิญฺญาเปติ. เอวํปิ ภิกฺขเว ทุพฺพลฺยาวิกมฺม\-ญฺเจว โหติ สิกฺขา จ อปจฺจกฺขาตา.}

\proseitem{50}{อถ วา ปน อุกฺกณฺฐิโต อนภิรโต สามญฺญา จวิตุกาโม ภิกฺขุภาวํ อฏฺฏียมาโน หรายมาโน ชิคุจฺฉมาโน คิหิภาวํ ปตฺถยมาโน ฯเปฯ อสกฺยปุตฺติย\-ภาวํ ปตฺถยมาโน ทุกฺกรนฺติ วทติ วิญฺญาเปติ. น สุกรนฺติ วทติ วิญฺญาเปติ. ทุจฺจรนฺติ วทติ วิญฺญาเปติ. น สุจรนฺติ วทติ วิญฺญาเปติ. น อุสฺสหามีติ วทติ วิญฺญาเปติ. น วิสหามีติ วทติ วิญฺญาเปติ. น รมามีติ วทติ วิญฺญาเปติ. นาภิรมามีติ วทติ วิญฺญาเปติ. เอวํปิ โข ภิกฺขเว ทุพฺพลฺยาวิกมฺม\-ญฺเจว โหติ สิกฺขา จ อปจฺจกฺขาตา.}

\proseitem{51}{\csromanfolio{32}กถญฺจ ภิกฺขเว ทุพฺพลฺยาวิกมฺม\-ญฺเจว โหติ สิกฺขา จ ปจฺจกฺขาตา. อิธ ภิกฺขเว ภิกฺขุ อุกฺกณฺฐิโต อนภิรโต สามญฺญา จวิตุกาโม ภิกฺขุภาวํ อฏฺฏียมาโน หรายมาโน ชิคุจฺฉมาโน คิหิภาวํ ปตฺถยมาโน ฯเปฯ อสกฺยปุตฺติย\-ภาวํ ปตฺถยมาโน พุทฺธํ ปจฺจกฺขามีติ วทติ วิญฺญาเปติ. เอวํปิ ภิกฺขเว ทุพฺพลฺยาวิกมฺม\-ญฺเจว โหติ สิกฺขา จ ปจฺจกฺขาตา.}

\prose{อถ วา ปน อุกฺกณฺฐิโต อนภิรโต สามญฺญา จวิตุกาโม ภิกฺขุภาวํ อฏฺฏียมาโน หรายมาโน ชิคุจฺฉมาโน คิหิภาวํ ปตฺถยมาโน ฯเปฯ อสกฺยปุตฺติย\-ภาวํ ปตฺถยมาโน ธมฺมํ ปจฺจกฺขามีติ วทติ วิญฺญาเปติ. สํฆํ ปจฺจกฺขามีติ วทติ วิญฺญาเปติ. สิกฺขํ ปจฺจกฺขามีติ วทติ วิญฺญาเปติ. วินยํ ปจฺจกฺขามีติ วทติ วิญฺญาเปติ. ปาติโมกฺขํ ปจฺจกฺขามีติ วทติ วิญฺญาเปติ. อุทฺเทสํ ปจฺจกฺขามีติ วทติ วิญฺญาเปติ. อุปชฺฌายํ ปจฺจกฺขามีติ วทติ วิญฺญาเปติ. อาจริยํ ปจฺจกฺขามีติ วทติ วิญฺญาเปติ. สทฺธิ\-วิหาริกํ ปจฺจกฺขามีติ วทติ วิญฺญาเปติ. อนฺเตวาสิกํ ปจฺจกฺขามีติ วทติ วิญฺญาเปติ. สมานุ\-ปชฺฌายกํ ปจฺจกฺขามีติ วทติ วิญฺญาเปติ. สมานา\-จริยกํ ปจฺจกฺขามีติ วทติ วิญฺญาเปติ. สพฺรหฺมจาริํ ปจฺจกฺขามีติ วทติ วิญฺญาเปติ. คิหีติ มํ ธาเรหีติ วทติ วิญฺญาเปติ. อุปาสโกติ มํ ธาเรหีติ วทติ วิญฺญาเปติ. อารามิโกติ มํ ธาเรหีติ วทติ วิญฺญาเปติ. สามเณโรติ มํ ธาเรหีติ วทติ วิญฺญาเปติ. ติตฺถิโยติ มํ ธาเรหีติ วทติ วิญฺญาเปติ. ติตฺถิย\-สาวโกติ มํ ธาเรหีติ วทติ วิญฺญาเปติ. อสฺสมโณติ มํ ธาเรหีติ วทติ วิญฺญาเปติ. อสกฺยปุตฺติโยติ มํ ธาเรหีติ วทติ วิญฺญาเปติ. เอวํปิ ภิกฺขเว ทุพฺพลฺยาวิกมฺม\-ญฺเจว โหติ สิกฺขา จ ปจฺจกฺขาตา.}

\proseitem{52}{อถ วา ปน อุกฺกณฺฐิโต อนภิรโต สามญฺญา จวิตุกาโม ภิกฺขุภาวํ อฏฺฏียมาโน หรายมาโน ชิคุจฺฉมาโน คิหิภาวํ ปตฺถยมาโน ฯเปฯ อสกฺยปุตฺติย\-ภาวํ ปตฺถยมาโน อลํ เม พุทฺเธนาติ วทติ วิญฺญาเปติ ฯเปฯ อลํ เม สพฺรหฺมจารีหีติ วทติ วิญฺญาเปติ. เอวํปิ ฯเปฯ. อถ วา ปน ฯเปฯ. กิํ นุ เม พุทฺเธนาติ วทติ วิญฺญาเปติ ฯเปฯ. กิํ นุ เม สพฺรหฺมจารีหีติ วทติ วิญฺญาเปติ. น มมตฺโถ พุทฺเธนาติ วทติ วิญฺญาเปติ ฯเปฯ. น มมตฺโถ สพฺรหฺมจารีหีติ วทติ วิญฺญาเปติ. สุมุตฺตาหํ \csromanfolio{33}พุทฺเธนาติ วทติ วิญฺญาเปติ ฯเปฯ. สุมุตฺตาหํ สพฺรหฺมจารีหีติ วทติ วิญฺญาเปติ. เอวํปิ ภิกฺขเว ทุพฺพลฺยาวิกมฺม\-ญฺเจว โหติ สิกฺขา จ ปจฺจกฺขาตา.}

\proseitem{53}{ยานิ วา ปนญฺญานิปิ อตฺถิ พุทฺธ\-เววจนานิ วา ธมฺม\-เววจนานิ วา สํฆเววจนานิ วา สิกฺขา\-เววจนานิ วา วินย\-เววจนานิ วา ปาติโมกฺข\-เววจนานิ วา อุทฺเทส\-เววจนานิ วา อุปชฺฌาย\-เววจนานิ วา อาจริย\-เววจนานิ วา สทฺธิวิหาริก\-เววจนานิ วา อนฺเตวาสิก\-เววจนานิ วา สมานุปชฺฌายก\-เววจนานิ วา สมานาจริยก\-เววจนานิ วา สพฺรหฺมจาริ\-เววจนานิ วา คิหิเววจนานิ วา อุปาสก\-เววจนานิ วา อารามิก\-เววจนานิ วา สามเณร\-เววจนานิ วา ติตฺถิย\-เววจนานิ วา ติตฺถิยสาวก\-เววจนานิ วา อสฺสมณ\-เววจนานิ วา อสกฺยปุตฺติย\-เววจนานิ วา, เตหิ อากาเรหิ เตหิ ลิงฺเคหิ เตหิ นิมิตฺเตหิ วทติ วิญฺญาเปติ. เอวํ โข ภิกฺขเว ทุพฺพลฺยาวิกมฺม\-ญฺเจว โหติ สิกฺขา จ ปจฺจกฺขาตา.}

\proseitem{54}{กถญฺจ ภิกฺขเว อปจฺจกฺขาตา โหติ สิกฺขา. อิธ ภิกฺขเว เยหิ อากาเรหิ เยหิ ลิงฺเคหิ เยหิ นิมิตฺเตหิ สิกฺขา ปจฺจกฺขาตา โหติ, เตหิ อากาเรหิ เตหิ ลิงฺเคหิ เตหิ นิมิตฺเตหิ อุมฺมตฺตโก สิกฺขํ ปจฺจกฺขาติ, อปจฺจกฺขาตา โหติ สิกฺขา. อุมฺมตฺตกสฺส สนฺติเก สิกฺขํ ปจฺจกฺขาติ, อปจฺจกฺขาตา โหติ สิกฺขา. ขิตฺตจิตฺโต สิกฺขํ ปจฺจกฺขาติ, อปจฺจกฺขาตา โหติ สิกฺขา. ขิตฺตจิตฺตสฺส สนฺติเก สิกฺขํ ปจฺจกฺขาติ, อปจฺจกฺขาตา โหติ สิกฺขา. เวทนาฏฺโฏ สิกฺขํ ปจฺจกฺขาติ, อปจฺจกฺขาตา โหติ สิกฺขา. เวทนาฏฺฏสฺส สนฺติเก สิกฺขํ ปจฺจกฺขาติ, อปจฺจกฺขาตา โหติ สิกฺขา. เทวตาย สนฺติเก สิกฺขํ ปจฺจกฺขาติ, อปจฺจกฺขาตา โหติ สิกฺขา. ติรจฺฉาน\-คตสฺส สนฺติเก สิกฺขํ ปจฺจกฺขาติ, อปจฺจกฺขาตา โหติ สิกฺขา. อริยเกน มิลกฺขสฺส\csromansharedfootnote{fe1bdaef38e7bfe0}{มิลกฺขกสฺส (สี, สฺยา) มิลกฺขุสฺส (ก)} สนฺติเก สิกฺขํ ปจฺจกฺขาติ, โส จ น ปฏิวิชานาติ, อปจฺจกฺขาตา โหติ สิกฺขา. มิลกฺขเกน อริยสฺส สนฺติเก สิกฺขํ ปจฺจกฺขาติ, โส จ น ปฏิวิชานาติ, อปจฺจกฺขาตา โหติ สิกฺขา. อริยเกน อริยสฺส สนฺติเก สิกฺขํ ปจฺจกฺขาติ, โส จ น \csromanfolio{34}ปฏิวิชานาติ, อปจฺจกฺขาตา โหติ สิกฺขา. มิลกฺขเกน มิลกฺขสฺส สนฺติเก สิกฺขํ ปจฺจกฺขาติ, โส จ น ปฏิวิชานาติ, อปจฺจกฺขาตา โหติ สิกฺขา. ทวาย สิกฺขํ ปจฺจกฺขาติ, อปจฺจกฺขาตา โหติ สิกฺขา. รวาย สิกฺขํ ปจฺจกฺขาติ, อปจฺจกฺขาตา โหติ สิกฺขา. อสาเวตุกาโม สาเวติ, อปจฺจกฺขาตา โหติ สิกฺขา. สาเวตุกาโม น สาเวติ, อปจฺจกฺขาตา โหติ สิกฺขา. อวิญฺญุสฺส สาเวติ, อปจฺจกฺขาตา โหติ สิกฺขา. วิญฺญุสฺส น สาเวติ, อปจฺจกฺขาตา โหติ สิกฺขา. สพฺพโส วา ปน น สาเวติ, อปจฺจกฺขาตา โหติ สิกฺขา. เอวํ โข ภิกฺขเว อปจฺจกฺขาตา โหติ สิกฺขา.}

\proseitem{55}{\csromansymbolfootnote{*}{ขุ 7. 107, 110, 112 ปิฏฺเฐ มหานิทฺเทเสปิ.}\textbf{เมถุนธมฺโม} นาม โย โส อสทฺธมฺโม คามธมฺโม วสลธมฺโม ทุฏฺฐุลฺลํ โอทกนฺติกํ รหสฺสํ ทฺวยํ\-ทฺวย\-สมาปตฺติ, เอโส เมถุนธมฺโม นาม.}

\prose{\textbf{ปฏิเสวติ} นาม โย นิมิตฺเตน นิมิตฺตํ องฺคชาเตน องฺคชาตํ อนฺตมโส ติลผล\-มตฺตํปิ ปเวเสติ, เอโส ปฏิเสวติ นาม.}

\prose{\textbf{อนฺตมโส ติรจฺฉาน}\-\textbf{คตา}\-\textbf{ยปี}ติ ติรจฺฉาน\-คติ\-ตฺถิยาปิ เมถุนํ ธมฺมํ ปฏิเสวิตฺวา อสฺสมโณ โหติ อสกฺยปุตฺติโย, ปเคว มนุสฺสิตฺถิยา, เตน วุจฺจติ “อนฺตมโส ติรจฺฉานคตายปี”ติ.}

\prose{\textbf{ปาราชิโก โหตี}ติ เสยฺยถาปิ นาม ปุริโส สีสจฺฉินฺโน อภพฺโพ เตน สรีร\-พนฺธเนน ชีวิตุํ, เอวเมว ภิกฺขุ เมถุนํ ธมฺมํ ปฏิเสวิตฺวา อสฺสมโณ โหติ อสกฺยปุตฺติโย, เตน วุจฺจติ “ปาราชิโก โหตี”ติ.}

\prose{\textbf{อสํวาโส}ติ สํวาโส นาม เอกกมฺมํ เอกุทฺเทโส สมสิกฺขตา, เอโส สํวาโส นาม. โส เตน สทฺธิํ นตฺถิ, เตน วุจฺจติ “อสํวาโส”ติ.}

\proseitem{56}{ติสฺโส อิตฺถิโย มนุสฺสิตฺถี อมนุสฺสิตฺถี ติรจฺฉาน\-คติตฺถี. ตโย อุภโต\-พฺยญฺชนกา มนุสฺสุ\-ภโตพฺยญฺชนโก อมนุสฺสุ\-ภโตพฺยญฺชนโก ติรจฺฉานคตุ\-ภโตพฺยญฺชนโก. ตโย ปณฺฑกา มนุสฺสปณฺฑโก \csromanfolio{35}อมนุสฺส\-ปณฺฑโก ติรจฺฉานคต\-ปณฺฑโก. ตโย ปุริสา มนุสฺสปุริโส อมนุสฺสปุริโส ติรจฺฉานคต\-ปุริโส.}

\prose{มนุสฺสิตฺถิยา ตโย มคฺเค เมถุนํ ธมฺมํ ปฏิเสวนฺตสฺส อาปตฺติ ปาราชิกสฺส วจฺจมคฺเค ปสฺสาวมคฺเค มุเข. อมนุสฺสิตฺถิยา ฯเปฯ. ติรจฺฉาน\-คติ\-ตฺถิยา ตโย มคฺเค เมถุนํ ธมฺมํ ปฏิเสวนฺตสฺส อาปตฺติ ปาราชิกสฺส วจฺจมคฺเค ปสฺสาวมคฺเค มุเข. มนุสฺสุ\-ภโต\-พฺยญฺชนกสฺส. อมนุสฺสุ\-ภโตพฺยญฺชนกสฺส. ติรจฺฉานคตุ\-ภโตพฺยญฺชนกสฺส ตโย มคฺเค เมถุนํ ธมฺมํ ปฏิเสวนฺตสฺส อาปตฺติ ปาราชิกสฺส วจฺจมคฺเค ปสฺสาวมคฺเค มุเข. มนุสฺส\-ปณฺฑกสฺส เทฺว มคฺเค เมถุนํ ธมฺมํ ปฏิเสวนฺตสฺส อาปตฺติ ปาราชิกสฺส วจฺจมคฺเค มุเข. อมนุสฺส\-ปณฺฑกสฺส. ติรจฺฉานคต\-ปณฺฑกสฺส. มนุสฺส\-ปุริสสฺส. อมนุสฺส\-ปุริสสฺส. ติรจฺฉานคต\-ปุริสสฺส เทฺว มคฺเค เมถุนํ ธมฺมํ ปฏิเสวนฺตสฺส อาปตฺติ ปาราชิกสฺส วจฺจมคฺเค มุเข.}

\proseitem{57}{ภิกฺขุสฺส เสวนจิตฺตํ อุปฏฺฐิเต มนุสฺสิตฺถิยา วจฺจมคฺคํ องฺคชาตํ ปเวเสนฺตสฺส อาปตฺติ ปาราชิกสฺส. ภิกฺขุสฺส เสวนจิตฺตํ อุปฏฺฐิเต มนุสฺสิตฺถิยา ปสฺสาวมคฺคํ. มุขํ องฺคชาตํ ปเวเสนฺตสฺส อาปตฺติ ปาราชิกสฺส. ภิกฺขุสฺส เสวนจิตฺตํ อุปฏฺฐิเต อมนุสฺสิตฺถิยา. ติรจฺฉาน\-คติ\-ตฺถิยา. มนุสฺสุ\-ภโต\-พฺยญฺชนกสฺส. อมนุสฺสุ\-ภโตพฺยญฺชนกสฺส. ติรจฺฉานคตุ\-ภโตพฺยญฺชนกสฺส วจฺจมคฺคํ. ปสฺสาวมคฺคํ. มุขํ องฺคชาตํ ปเวเสนฺตสฺส อาปตฺติ ปาราชิกสฺส. ภิกฺขุสฺส เสวนจิตฺตํ อุปฏฺฐิเต มนุสฺส\-ปณฺฑกสฺส วจฺจมคฺคํ. มุขํ องฺคชาตํ ปเวเสนฺตสฺส อาปตฺติ ปาราชิกสฺส. ภิกฺขุสฺส เสวนจิตฺตํ อุปฏฺฐิเต อมนุสฺส\-ปณฺฑกสฺส. ติรจฺฉานคต\-ปณฺฑกสฺส. มนุสฺส\-ปุริสสฺส. อมนุสฺส\-ปุริสสฺส. ติรจฺฉานคต\-ปุริสสฺส. วจฺจมคฺคํ. มุขํ องฺคชาตํ ปเวเสนฺตสฺส อาปตฺติ ปาราชิกสฺส.}

\proseitem{58}{ภิกฺขุ\-ปจฺจตฺถิกา มนุสฺสิตฺถิํ ภิกฺขุสฺส สนฺติเก อาเนตฺวา วจฺจมคฺเคน องฺคชาตํ อภินิสีเทนฺติ, โส เจ ปเวสนํ สาทิยติ\csromansharedfootnote{657c357bbf86b4d1}{ปวิสนํ สาทยติ (ก)}, ปวิฏฺฐํ สาทิยติ, ฐิตํ สาทิยติ, อุทฺธรณํ สาทิยติ, อาปตฺติ ปาราชิกสฺส. ภิกฺขุ\-ปจฺจตฺถิกา มนุสฺสิตฺถิํ ภิกฺขุสฺส สนฺติเก อาเนตฺวา วจฺจมคฺเคน องฺคชาตํ อภินิสีเทนฺติ, โส เจ ปเวสนํ น สาทิยติ, ปวิฏฺฐํ สาทิยติ, ฐิตํ สาทิยติ, อุทฺธรณํ สาทิยติ, อาปตฺติ ปาราชิกสฺส. ภิกฺขุ\-ปจฺจตฺถิกา มนุสฺสิตฺถิํ ภิกฺขุสฺส \csromanfolio{36}สนฺติเก อาเนตฺวา วจฺจมคฺเคน องฺคชาตํ อภินิสีเทนฺติ, โส เจ ปเวสนํ น สาทิยติ, ปวิฏฺฐํ น สาทิยติ, ฐิตํ สาทิยติ, อุทฺธรณํ สาทิยติ, อาปตฺติ ปาราชิกสฺส. ภิกฺขุ\-ปจฺจตฺถิกา มนุสฺสิตฺถิํ ภิกฺขุสฺส สนฺติเก อาเนตฺวา วจฺจมคฺเคน องฺคชาตํ อภินิสีเทนฺติ, โส เจ ปเวสนํ น สาทิยติ, ปวิฏฺฐํ น สาทิยติ, ฐิตํ น สาทิยติ, อุทฺธรณํ สาทิยติ, อาปตฺติ ปาราชิกสฺส. ภิกฺขุ\-ปจฺจตฺถิกา มนุสฺสิตฺถิํ ภิกฺขุสฺส สนฺติเก อาเนตฺวา วจฺจมคฺเคน องฺคชาตํ อภินิสีเทนฺติ, โส เจ ปเวสนํ น สาทิยติ, ปวิฏฺฐํ น สาทิยติ, ฐิตํ น สาทิยติ, อุทฺธรณํ น สาทิยติ, อนาปตฺติ.}

\prose{ภิกฺขุ\-ปจฺจตฺถิกา มนุสฺสิตฺถิํ ภิกฺขุสฺส สนฺติเก อาเนตฺวา ปสฺสาวมคฺเคน. มุเขน องฺคชาตํ อภินิสีเทนฺติ, โส เจ ปเวสนํ สาทิยติ, ปวิฏฺฐํ สาทิยติ, ฐิตํ สาทิยติ, อุทฺธรณํ สาทิยติ, อาปตฺติ ปาราชิกสฺส ฯเปฯ น สาทิยติ, อนาปตฺติ.}

\proseitem{59}{ภิกฺขุ\-ปจฺจตฺถิกา มนุสฺสิตฺถิํ ชาครนฺติํ. สุตฺตํ. มตฺตํ. อุมฺมตฺตํ. ปมตฺตํ. มตํ อกฺขายิตํ. มตํ เยภุยฺเยน อกฺขายิตํ ฯเปฯ อาปตฺติ ปาราชิกสฺส. มตํ เยภุยฺเยน ขายิตํ ภิกฺขุสฺส สนฺติเก อาเนตฺวา วจฺจมคฺเคน. ปสฺสาวมคฺเคน. มุเขน องฺคชาตํ อภินิสีเทนฺติ, โส เจ ปเวสนํ สาทิยติ, ปวิฏฺฐํ สาทิยติ, ฐิตํ สาทิยติ, อุทฺธรณํ สาทิยติ, อาปตฺติ ถุลฺลจฺจยสฺส ฯเปฯ น สาทิยติ, อนาปตฺติ.}

\prose{ภิกฺขุ\-ปจฺจตฺถิกา อมนุสฺสิตฺถิํ. ติรจฺฉาน\-คติตฺถิํ. มนุสฺสุ\-ภโตพฺยญฺชนกํ. อมนุสฺสุ\-ภโตพฺยญฺชนกํ. ติรจฺฉานคตุ\-ภโตพฺยญฺชนกํ ภิกฺขุสฺส สนฺติเก อาเนตฺวา วจฺจมคฺเคน. ปสฺสาวมคฺเคน. มุเขน องฺคชาตํ อภินิสีเทนฺติ, โส เจ ปเวสนํ สาทิยติ, ปวิฏฺฐํ สาทิยติ, ฐิตํ สาทิยติ, อุทฺธรณํ สาทิยติ, อาปตฺติ ปาราชิกสฺส ฯเปฯ น สาทิยติ, อนาปตฺติ.}

\prose{ภิกฺขุ\-ปจฺจตฺถิกา ติรจฺฉานคตุ\-ภโตพฺยญฺชนกํ ชาครนฺตํ. สุตฺตํ. มตฺตํ. อุมฺมตฺตํ. ปมตฺตํ. มตํ อกฺขายิตํ. มตํ เยภุยฺเยน อกฺขายิตํ ฯเปฯ อาปตฺติ ปาราชิกสฺส. มตํ เยภุยฺเยน ขายิตํ ภิกฺขุสฺส สนฺติเก อาเนตฺวา วจฺจมคฺเคน. ปสฺสาวมคฺเคน. มุเขน องฺคชาตํ อภินิสีเทนฺติ, โส เจ ปเวสนํ \csromanfolio{37}สาทิยติ, ปวิฏฺฐํ สาทิยติ, ฐิตํ สาทิยติ, อุทฺธรณํ สาทิยติ, อาปตฺติ ถุลฺลจฺจยสฺส ฯเปฯ น สาทิยติ, อนาปตฺติ.}

\prose{ภิกฺขุ\-ปจฺจตฺถิกา มนุสฺส\-ปณฺฑกํ. อมนุสฺส\-ปณฺฑกํ. ติรจฺฉานคต\-ปณฺฑกํ ภิกฺขุสฺส สนฺติเก อาเนตฺวา วจฺจมคฺเคน. มุเขน องฺคชาตํ อภินิสีเทนฺติ, โส เจ ปเวสนํ สาทิยติ, ปวิฏฺฐํ สาทิยติ, ฐิตํ สาทิยติ, อุทฺธรณํ สาทิยติ, อาปตฺติ ปาราชิกสฺส ฯเปฯ น สาทิยติ, อนาปตฺติ.}

\prose{ภิกฺขุ\-ปจฺจตฺถิกา ติรจฺฉานคต\-ปณฺฑกํ ชาครนฺตํ. สุตฺตํ. มตฺตํ. อุมฺมตฺตํ. ปมตฺตํ. มตํ อกฺขายิตํ. มตํ เยภุยฺเยน อกฺขายิตํ ฯเปฯ อาปตฺติ ปาราชิกสฺส. มตํ เยภุยฺเยน ขายิตํ ภิกฺขุสฺส สนฺติเก อาเนตฺวา วจฺจมคฺเคน. มุเขน องฺคชาตํ อภินิสีเทนฺติ, โส เจ ปเวสนํ สาทิยติ, ปวิฏฺฐํ สาทิยติ, ฐิตํ สาทิยติ, อุทฺธรณํ สาทิยติ, อาปตฺติ ถุลฺลจฺจยสฺส ฯเปฯ น สาทิยติ, อนาปตฺติ.}

\proseitem{60}{ภิกฺขุ\-ปจฺจตฺถิกา มนุสฺสปุริสํ. อมนุสฺส\-ปุริสํ. ติรจฺฉานคต\-ปุริสํ ภิกฺขุสฺส สนฺติเก อาเนตฺวา วจฺจมคฺเคน. มุเขน องฺคชาตํ อภินิสีเทนฺติ, โส เจ ปเวสนํ สาทิยติ, ปวิฏฺฐํ สาทิยติ, ฐิตํ สาทิยติ, อุทฺธรณํ สาทิยติ, อาปตฺติ ปาราชิกสฺส ฯเปฯ น สาทิยติ, อนาปตฺติ.}

\prose{ภิกฺขุ\-ปจฺจตฺถิกา ติรจฺฉานคต\-ปุริสํ ชาครนฺตํ. สุตฺตํ. มตฺตํ. อุมฺมตฺตํ. ปมตฺตํ. มตํ อกฺขายิตํ. มตํ เยภุยฺเยน อกฺขายิตํ ฯเปฯ อาปตฺติ ปาราชิกสฺส. มตํ เยภุยฺเยน ขายิตํ ภิกฺขุสฺส สนฺติเก อาเนตฺวา วจฺจมคฺเคน. มุเขน องฺคชาตํ อภินิสีเทนฺติ, โส เจ ปเวสนํ สาทิยติ, ปวิฏฺฐํ สาทิยติ, ฐิตํ สาทิยติ, อุทฺธรณํ สาทิยติ, อาปตฺติ ถุลฺลจฺจยสฺส ฯเปฯ น สาทิยติ, อนาปตฺติ.}

\proseitem{61}{ภิกฺขุ\-ปจฺจตฺถิกา มนุสฺสิตฺถิํ ภิกฺขุสฺส สนฺติเก อาเนตฺวา วจฺจมคฺเคน. ปสฺสาวมคฺเคน. มุเขน องฺคชาตํ อภินิสีเทนฺติ สนฺถตาย อสนฺถตสฺส. อสนฺถตาย สนฺถตสฺส. สนฺถตาย สนฺถตสฺส. อสนฺถตาย อสนฺถตสฺส. โส เจ ปเวสนํ สาทิยติ, ปวิฏฺฐํ สาทิยติ, ฐิตํ สาทิยติ, อุทฺธรณํ สาทิยติ, อาปตฺติ ปาราชิกสฺส ฯเปฯ น สาทิยติ, อนาปตฺติ.}

\prose{\csromanfolio{38}ภิกฺขุ\-ปจฺจตฺถิกา มนุสฺสิตฺถิํ ชาครนฺติํ. สุตฺตํ. มตฺตํ. อุมฺมตฺตํ. ปมตฺตํ. มตํ อกฺขายิตํ. มตํ เยภุยฺเยน อกฺขายิตํ ฯเปฯ อาปตฺติ ปาราชิกสฺส. มตํ เยภุยฺเยน ขายิตํ ภิกฺขุสฺส สนฺติเก อาเนตฺวา วจฺจมคฺเคน. ปสฺสาวมคฺเคน. มุเขน องฺคชาตํ อภินิสีเทนฺติ สนฺถตาย อสนฺถตสฺส. อสนฺถตาย สนฺถตสฺส. สนฺถตาย สนฺถตสฺส. อสนฺถตาย อสนฺถตสฺส. โส เจ ปเวสนํ สาทิยติ, ปวิฏฺฐํ สาทิยติ, ฐิตํ สาทิยติ, อุทฺธรณํ สาทิยติ, อาปตฺติ ถุลฺลจฺจยสฺส ฯเปฯ น สาทิยติ, อนาปตฺติ.}

\prose{ภิกฺขุ\-ปจฺจตฺถิกา อมนุสฺสิตฺถิํ. ติรจฺฉาน\-คติตฺถิํ. มนุสฺสุ\-ภโตพฺยญฺชนกํ. อมนุสฺสุ\-ภโตพฺยญฺชนกํ. ติรจฺฉานคตุ\-ภโตพฺยญฺชนกํ ภิกฺขุสฺส สนฺติเก อาเนตฺวา วจฺจมคฺเคน. ปสฺสาวมคฺเคน. มุเขน องฺคชาตํ อภินิสีเทนฺติ สนฺถตสฺส อสนฺถตสฺส. อสนฺถตสฺส สนฺถตสฺส. สนฺถตสฺส สนฺถตสฺส. อสนฺถตสฺส อสนฺถตสฺส. โส เจ ปเวสนํ สาทิยติ, ปวิฏฺฐํ สาทิยติ, ฐิตํ สาทิยติ, อุทฺธรณํ สาทิยติ, อาปตฺติ ปาราชิกสฺส ฯเปฯ น สาทิยติ, อนาปตฺติ.}

\prose{ภิกฺขุ\-ปจฺจตฺถิกา ติรจฺฉานคตุ\-ภโตพฺยญฺชนกํ ชาครนฺตํ. สุตฺตํ. มตฺตํ. อุมฺมตฺตํ. ปมตฺตํ. มตํ อกฺขายิตํ. มตํ เยภุยฺเยน อกฺขายิตํ ฯเปฯ อาปตฺติ ปาราชิกสฺส. มตํ เยภุยฺเยน ขายิตํ ภิกฺขุสฺส สนฺติเก อาเนตฺวา วจฺจมคฺเคน. ปสฺสาวมคฺเคน. มุเขน องฺคชาตํ อภินิสีเทนฺติ สนฺถตสฺส อสนฺถตสฺส. อสนฺถตสฺส สนฺถตสฺส. สนฺถตสฺส สนฺถตสฺส. อสนฺถตสฺส อสนฺถตสฺส. โส เจ ปเวสนํ สาทิยติ, ปวิฏฺฐํ สาทิยติ, ฐิตํ สาทิยติ, อุทฺธรณํ สาทิยติ, อาปตฺติ ถุลฺลจฺจยสฺส ฯเปฯ น สาทิยติ, อนาปตฺติ.}

\proseitem{62}{ภิกฺขุ\-ปจฺจตฺถิกา มนุสฺส\-ปณฺฑกํ. อมนุสฺส\-ปณฺฑกํ. ติรจฺฉานคต\-ปณฺฑกํ. มนุสฺสปุริสํ. อมนุสฺส\-ปุริสํ. ติรจฺฉานคต\-ปุริสํ ภิกฺขุสฺส สนฺติเก อาเนตฺวา วจฺจมคฺเคน, มุเขน องฺคชาตํ อภินิสีเทนฺติ สนฺถตสฺส อสนฺถตสฺส. อสนฺถตสฺส สนฺถตสฺส. สนฺถตสฺส สนฺถตสฺส. อสนฺถตสฺส อสนฺถตสฺส. โส เจ ปเวสนํ สาทิยติ, ปวิฏฺฐํ สาทิยติ, ฐิตํ สาทิยติ, อุทฺธรณํ สาทิยติ, อาปตฺติ ปาราชิกสฺส ฯเปฯ น สาทิยติ, อนาปตฺติ.}

\prose{\csromanfolio{39}ภิกฺขุ\-ปจฺจตฺถิกา ติรจฺฉานคต\-ปุริสํ ชาครนฺตํ. สุตฺตํ. มตฺตํ. อุมฺมตฺตํ. ปมตฺตํ. มตํ อกฺขายิตํ. มตํ เยภุยฺเยน อกฺขายิตํ ฯเปฯ อาปตฺติ ปาราชิกสฺส. มตํ เยภุยฺเยน ขายิตํ ภิกฺขุสฺส สนฺติเก อาเนตฺวา วจฺจมคฺเคน, มุเขน องฺคชาตํ อภินิสีเทนฺติ สนฺถตสฺส อสนฺถตสฺส. อสนฺถตสฺส สนฺถตสฺส. สนฺถตสฺส สนฺถตสฺส. อสนฺถตสฺส อสนฺถตสฺส. โส เจ ปเวสนํ สาทิยติ, ปวิฏฺฐํ สาทิยติ, ฐิตํ สาทิยติ, อุทฺธรณํ สาทิยติ, อาปตฺติ ถุลฺลจฺจยสฺส ฯเปฯ น สาทิยติ, อนาปตฺติ.}

\proseitem{63}{ภิกฺขุ\-ปจฺจตฺถิกา ภิกฺขุํ มนุสฺสิตฺถิยา สนฺติเก อาเนตฺวา องฺคชาเตน วจฺจมคฺคํ. ปสฺสาวมคฺคํ. มุขํ อภินิสิเทนฺติ. โส เจ ปเวสนํ สาทิยติ, ปวิฏฺฐํ สาทิยติ, ฐิตํ สาทิยติ, อุทฺธรณํ สาทิยติ, อาปตฺติ ปาราชิกสฺส ฯเปฯ น สาทิยติ, อนาปตฺติ.}

\prose{ภิกฺขุ\-ปจฺจตฺถิกา ภิกฺขุํ มนุสฺสิตฺถิยา ชาครนฺติยา. สุตฺตาย. มตฺตาย. อุมฺมตฺตาย. ปมตฺตาย. มตาย อกฺขายิตาย. มตาย เยภุยฺเยน อกฺขายิตาย ฯเปฯ อาปตฺติ ปาราชิกสฺส. มตาย เยภุยฺเยน ขายิตาย สนฺติเก อาเนตฺวา องฺคชาเตน วจฺจมคฺคํ. ปสฺสาวมคฺคํ. มุขํ อภินิสีเทนฺติ. โส เจ ปเวสนํ สาทิยติ, ปวิฏฺฐํ สาทิยติ, ฐิตํ สาทิยติ, อุทฺธรณํ สาทิยติ, อาปตฺติ ถุลฺลจฺจยสฺส ฯเปฯ น สาทิยติ, อนาปตฺติ.}

\prose{ภิกฺขุ\-ปจฺจตฺถิกา ภิกฺขุํ อมนุสฺสิตฺถิยา. ติรจฺฉาน\-คติ\-ตฺถิยา. มนุสฺสุ\-ภโต\-พฺยญฺชนกสฺส. อมนุสฺสุ\-ภโตพฺยญฺชนกสฺส. ติรจฺฉานคตุ\-ภโตพฺยญฺชนกสฺส. มนุสฺส\-ปณฺฑกสฺส. อมนุสฺส\-ปณฺฑกสฺส. ติรจฺฉานคต\-ปณฺฑกสฺส. มนุสฺส\-ปุริสสฺส. อมนุสฺส\-ปุริสสฺส. ติรจฺฉานคต\-ปุริสสฺส สนฺติเก อาเนตฺวา องฺคชาเตน วจฺจมคฺคํ. มุขํ อภินิสีเทนฺติ. โส เจ ปเวสนํ สาทิยติ, ปวิฏฺฐํ สาทิยติ, ฐิตํ สาทิยติ, อุทฺธรณํ สาทิยติ, อาปตฺติ ปาราชิกสฺส ฯเปฯ น สาทิยติ, อนาปตฺติ.}

\prose{ภิกฺขุ\-ปจฺจตฺถิกา ภิกฺขุํ ติรจฺฉานคต\-ปุริสสฺส ชาครนฺตสฺส. สุตฺตสฺส. มตฺตสฺส. อุมฺมตฺตสฺส. ปมตฺตสฺส. มตสฺส อกฺขายิตสฺส. มตสฺส เยภุยฺเยน อกฺขายิตสฺส ฯเปฯ อาปตฺติ ปาราชิกสฺส. มตสฺส เยภุยฺเยน ขายิตสฺส สนฺติเก อาเนตฺวา องฺคชาเตน วจฺจมคฺคํ. มุขํ อภินิสีเทนฺติ. \csromanfolio{40}โส เจ ปเวสนํ สาทิยติ, ปวิฏฺฐํ สาทิยติ, ฐิตํ สาทิยติ, อุทฺธรณํ สาทิยติ, อาปตฺติ ถุลฺลจฺจยสฺส ฯเปฯ น สาทิยติ, อนาปตฺติ.}

\proseitem{64}{ภิกฺขุ\-ปจฺจตฺถิกา ภิกฺขุํ มนุสฺสิตฺถิยา สนฺติเก อาเนตฺวา องฺคชาเตน วจฺจมคฺคํ. ปสฺสาวมคฺคํ. มุขํ อภินิสีเทนฺติ สนฺถตสฺส อสนฺถตาย. อสนฺถตสฺส สนฺถตาย. สนฺถตสฺส สนฺถตาย. อสนฺถตสฺส อสนฺถตาย. โส เจ ปเวสนํ สาทิยติ, ปวิฏฺฐํ สาทิยติ, ฐิตํ สาทิยติ, อุทฺธรณํ สาทิยติ, อาปตฺติ ปาราชิกสฺส ฯเปฯ น สาทิยติ, อนาปตฺติ.}

\prose{ภิกฺขุ\-ปจฺจตฺถิกา ภิกฺขุํ มนุสฺสิตฺถิยา ชาครนฺติยา. สุตฺตาย. มตฺตาย. อุมฺมตฺตาย. ปมตฺตาย. มตาย อกฺขายิตาย. มตาย เยภุยฺเยน อกฺขายิตาย ฯเปฯ อาปตฺติ ปาราชิกสฺส. มตาย เยภุยฺเยน ขายิตาย สนฺติเก อาเนตฺวา องฺคชาเตน วจฺจมคฺคํ. ปสฺสาวมคฺคํ. มุขํ อภินิสีเทนฺติ สนฺถตสฺส อสนฺถตาย. อสนฺถตสฺส สนฺถตาย. สนฺถตสฺส สนฺถตาย. อสนฺถตสฺส อสนฺถตาย. โส เจ ปเวสนํ สาทิยติ, ปวิฏฺฐํ สาทิยติ, ฐิตํ สาทิยติ, อุทฺธรณํ สาทิยติ, อาปตฺติ ถุลฺลจฺจยสฺส ฯเปฯ น สาทิยติ, อนาปตฺติ.}

\prose{ภิกฺขุ\-ปจฺจตฺถิกา ภิกฺขุํ อมนุสฺสิตฺถิยา. ติรจฺฉานคติตฺติยา. มนุสฺสุ\-ภโต\-พฺยญฺชนกสฺส. อมนุสฺสุ\-ภโตพฺยญฺชนกสฺส. ติรจฺฉานนคตุภโต\-พฺยญฺชนกสฺส. มนุสฺส\-ปณฺฑกสฺส. อมนุสฺส\-ปณฺฑกสฺส. ติรจฺฉานคต\-ปณฺฑกสฺส. มนุสฺส\-ปุริสสฺส. อมนุสฺส\-ปุริสสฺส. ติรจฺฉานคต\-ปุริสสฺส สนฺติเก อาเนตฺวา องฺคชาเตน วจฺจมคฺคํ. มุขํ อภินิสีเทนฺติ สนฺถตสฺส อสนฺถตสฺส. อสนฺถตสฺส สนฺถตสฺส. สนฺถตสฺส สนฺถตสฺส. อสนฺถตสฺส อสนฺถตสฺส. โส เจ ปเวสนํ สาทิยติ, ปวิฏฺฐํ สาทิยติ, ฐิตํ สาทิยติ, อุทฺธรณํ สาทิยติ, อาปตฺติ ปาราชิกสฺส ฯเปฯ น สาทิยติ, อนาปตฺติ.}

\proseitem{65}{ภิกฺขุ\-ปจฺจตฺถิกา ภิกฺขุํ ติรจฺฉานคต\-ปุริสสฺส ชาครนฺตสฺส. สุตฺตสฺส. มตฺตสฺส. อุมฺมตฺตสฺส. ปมตฺตสฺส. มตสฺส อกฺขายิตสฺส. มตสฺส เยภุยฺเยน อกฺขายิตสฺส ฯเปฯ อาปตฺติ ปาราชิกสฺส. มตสฺส เยภุยฺเยน ขายิตสฺส สนฺติเก อาเนตฺวา องฺคชาเตน วจฺจมคฺคํ. มุขํ อภินิสีเทนฺติ สนฺถตสฺส อสนฺถตสฺส. อสนฺถตสฺส สนฺถตสฺส. สนฺถตสฺส \csromanfolio{41}สนฺถตสฺส. อสนฺถตสฺส อสนฺถตสฺส. โส เจ ปเวสนํ สาทิยติ, ปวิฏฺฐํ สาทิยติ, ฐิตํ สาทิยติ, อุทฺธรณํ สาทิยติ, อาปตฺติ ถุลฺลจฺจยสฺส ฯเปฯ น สาทิยติ อนาปตฺติ.}

\prose{ยถา ภิกฺขุ\-ปจฺจตฺถิกา วิตฺถาริตา, เอวํ วิตฺถาเรตพฺพา.}

\prose{ราช\-ปจฺจตฺถิกา. โจรปจฺจตฺติกา. ธุตฺต\-ปจฺจตฺถิกา. อุปฺปฬ\-คนฺธ\-ปจฺจตฺถิกา. สํขิตฺตํ.}

\proseitem{66}{มคฺเคน มคฺคํ ปเวเสติ, อาปตฺติ ปาราชิกสฺส. มคฺเคน อมคฺคํ ปเวเสติ, อาปตฺติ ปาราชิกสฺส. อมคฺเคน มคฺคํ ปเวเสติ, อาปตฺติ ปาราชิกสฺส. อมคฺเคน อมคฺคํ ปเวเสติ, อาปตฺติ ถุลฺลจฺจยสฺส.}

\prose{ภิกฺขุ สุตฺต\-ภิกฺขุมฺหิ วิปฺปฏิปชฺชติ, ปฏิพุทฺโธ สาทิยติ, อุโภ นาเสตพฺพา. ปฏิพุทฺโธ น สาทิยติ, ทูสโก นาเสตพฺโพ. ภิกฺขุ สุตฺต\-สามเณรมฺหิ วิปฺปฏิปชฺชติ, ปฏิพุทฺโธ สาทิยติ, อุโภ นาเสตพฺพา. ปฏิพุทฺโธ น สาทิยติ, ทูสโก นาเสตพฺโพ. สามเณโร สุตฺต\-ภิกฺขุมฺหิ วิปฺปฏิปชฺชติ, ปฏิพุทฺโธ สาทิยติ, อุโภ นาเสตพฺพา. ปฏิพุทฺโธ น สาทิยติ, ทูสโก นาเสตพฺโพ. สามเณโร สุตฺต\-สามเณรมฺหิ วิปฺปฏิปชฺชติ, ปฏิพุทฺโธ สาทิยติ, อุโภ นาเสตพฺพา. ปฏิพุทฺโธ น สาทิยติ, ทูสโก นาเสตพฺโพ.}

\prosewithcloserruled{อนาปตฺติ อชานนฺตสฺส อสาทิยนฺตสฺส อุมฺมตฺตกสฺส ขิตฺตจิตฺตสฺส เวทนาฏฺฏสฺส อาทิกมฺมิกสฺสาติ.}{สนฺถต\-ภาณวาโร นิฏฺฐิโต.}
\csromanmatikaanchor{mtk.13}
\csromantocmarkref{subsection}{วินีตวตฺถุอุทฺทานคาถา}{mtk.13}
\bookmark[dest={mtk.13},level=2]{วินีตวตฺถุอุทฺทานคาถา}
\csromanheader{2}{วินีตวตฺถุ\-อุทฺทานคาถา}
\setlength{\csromangathapagewidth}{0pt}
\setlength{\csromangathawakwidth}{0pt}
\csromangathameasuregroup{มกฺกฏี วชฺชิปุตฺตา จ, \\ ทาริกุปฺปลวณฺณา จ, \\ มาตา ธีตา ภคินี จ, \\ เทฺว วณา เลปจิตฺตญฺจ, \\ สุนฺทเรน สห ปญฺจ, \\ นาคี ยกฺขี จ เปตี จ, \\ ภทฺทิเย อรหํ สุตฺโต, \\ เวสาลิยา ตโย มาลา, \\ สุปพฺพา สทฺธา ภิกฺขุนี, \\ เวสิยา ปณฺฑโก คิหี,}{\csromangathabat{มกฺกฏี วชฺชิปุตฺตา จ,}{คิหี นคฺโค จ ติตฺถิยา.} \\ \csromangathabat{ทาริกุปฺปลวณฺณา จ,}{พฺยญฺชเนหิปเร ทุเว.} \\ \csromangathabat{มาตา ธีตา ภคินี จ,}{ชายา จ มุทุ ลมฺพินา.} \\ \csromangathabat{เทฺว วณา เลปจิตฺตญฺจ,}{ทารุธีตลิกาย จ.} \\ \csromangathabat{สุนฺทเรน สห ปญฺจ,}{ปญฺจ สิวถิกฏฺฐิกา.} \\ \csromangathabat{นาคี ยกฺขี จ เปตี จ,}{ปณฺฑโกปหโต ฉุเป.} \\ \csromangathabat{ภทฺทิเย อรหํ สุตฺโต,}{สาวตฺถิยา จตุโร ปเร.} \\ \csromangathabat{เวสาลิยา ตโย มาลา,}{สุปิเน ภารุกจฺฉโก.} \\ \csromangathabat{สุปพฺพา สทฺธา ภิกฺขุนี,}{สิกฺขมานา สามเณรี จ.} \\ \csromangathabat{เวสิยา ปณฺฑโก คิหี,}{อญฺญมญฺญํ วุฑฺฒปพฺพชิโต มิโคติ.}}
\csromangathasetleft{มกฺกฏี วชฺชิปุตฺตา จ, \\ ทาริกุปฺปลวณฺณา จ, \\ มาตา ธีตา ภคินี จ, \\ เทฺว วณา เลปจิตฺตญฺจ, \\ สุนฺทเรน สห ปญฺจ, \\ นาคี ยกฺขี จ เปตี จ, \\ ภทฺทิเย อรหํ สุตฺโต, \\ เวสาลิยา ตโย มาลา, \\ สุปพฺพา สทฺธา ภิกฺขุนี, \\ เวสิยา ปณฺฑโก คิหี,}
\csromangathagroup{\csromangathastanza{\csromangathabat{มกฺกฏี วชฺชิปุตฺตา จ,}{คิหี นคฺโค จ ติตฺถิยา.} \\ \csromangathabat{ทาริกุปฺปลวณฺณา จ,}{พฺยญฺชเนหิปเร ทุเว.}}\csromangathastanza{\csromanfolio{42}\csromangathabat{มาตา ธีตา ภคินี จ,}{ชายา จ มุทุ ลมฺพินา.} \\ \csromangathabat{เทฺว วณา เลปจิตฺตญฺจ,}{ทารุ\-ธีตลิกาย จ.}}\csromangathastanza{\csromangathabat{สุนฺทเรน สห ปญฺจ,}{ปญฺจ สิวถิกฏฺฐิกา.} \\ \csromangathabat{นาคี ยกฺขี จ เปตี จ,}{ปณฺฑโกปหโต ฉุเป.}}\csromangathastanza{\csromangathabat{ภทฺทิเย อรหํ สุตฺโต,}{สาวตฺถิยา จตุโร ปเร.} \\ \csromangathabat{เวสาลิยา ตโย มาลา,}{สุปิเน ภารุกจฺฉโก.}}\csromangathastanza{\csromangathabat{สุปพฺพา สทฺธา ภิกฺขุนี,}{สิกฺขมานา สามเณรี จ.} \\ \csromangathabat{เวสิยา ปณฺฑโก คิหี,}{อญฺญมญฺญํ วุฑฺฒ\-ปพฺพชิโต มิโคติ.}}}

\csromanmatikaanchor{mtk.14}
\csromantocmarkref{subsection}{วินีตวตฺถุ}{mtk.14}
\bookmark[dest={mtk.14},level=2]{วินีตวตฺถุ}
\csromanheader{2}{\textbf{วินีตวตฺถุ}}
\proseitem{67}{เตน โข ปน สมเยน อญฺญตโร ภิกฺขุ มกฺกฏิยา เมถุนํ ธมฺมํ ปฏิเสวิ. ตสฺส กุกฺกุจฺจํ อโหสิ “ภควตา สิกฺขาปทํ ปญฺญตฺตํ, กจฺจิ นุ โข อหํ ปาราชิกํ อาปตฺติํ อาปนฺโน”ติ. ภควโต เอตมตฺถํ อาโรเจสิ. อาปตฺติํ ตฺวํ ภิกฺขุ อาปนฺโน ปาราชิกนฺติ. (1)}

\prose{เตน โข ปน สมเยน สมฺพหุลา เวสาลิกา วชฺชิปุตฺตกา ภิกฺขู สิกฺขํ อปจฺจกฺขาย ทุพฺพลฺยํ อนาวิกตฺวา เมถุนํ ธมฺมํ ปฏิเสวิํสุ. เตสํ กุกฺกุจฺจํ อโหสิ “ภควตา สิกฺขาปทํ ปญฺญตฺตํ, กจฺจิ นุ โข มยํ ปาราชิกํ อาปตฺติํ อาปนฺนา”ติ. ภควโต เอตมตฺถํ อาโรเจสุํ. อาปตฺติํ ตุเมฺห ภิกฺขเว อาปนฺนา ปาราชิกนฺติ. (2)}

\prose{เตน โข ปน สมเยน อญฺญตโร ภิกฺขุ เอวํ เม อนาปตฺติ ภวิสฺสตีติ คิหิลิงฺเคน เมถุนํ ธมฺมํ ปฏิเสวิ. ตสฺส กุกฺกุจฺจํ อโหสิ “ภควตา สิกฺขาปทํ ปญฺญตฺตํ, กจฺจิ นุ โข อหํ ปาราชิกํ อาปตฺติํ อาปนฺโน”ติ. ภควโต เอตมตฺถํ อาโรเจสิ. อาปตฺติํ ตฺวํ ภิกฺขุ อาปนฺโน ปาราชิกนฺติ. (3)}

\prose{เตน โข ปน สมเยน อญฺญตโร ภิกฺขุ เอวํ เม อนาปตฺติ ภวิสฺสตีติ นคฺโค หุตฺวา เมถุนํ ธมฺมํ ปฏิเสวิ. ตสฺส กุกฺกุจฺจํ อโหสิ ฯเปฯ. อาปตฺติํ ตฺวํ ภิกฺขุ อาปนฺโน ปาราชิกนฺติ. (4)}

\prose{\csromanfolio{43}เตน โข ปน สมเยน อญฺญตโร ภิกฺขุ เอวํ เม อนาปตฺติ ภวิสฺสตีติ กุสจีรํ นิวาเสตฺวา. วากจีรํ นิวาเสตฺวา. ผลกจีรํ นิวาเสตฺวา. เกสกมฺพลํ นิวาเสตฺวา. วาลกมฺพลํ นิวาเสตฺวา. อุลูกปกฺขิกํ นิวาเสตฺวา. อชินกฺขิปํ นิวาเสตฺวา เมถุนํ ธมฺมํ ปฏิเสวิ. ตสฺส กุกฺกุจฺจํ อโหสิ ฯเปฯ. อาปตฺติํ ตฺวํ ภิกฺขุ อาปนฺโน ปาราชิกนฺติ. \mbox{(5-11)}}

\prose{เตน โข ปน สมเยน อญฺญตโร ปิณฺฑจาริโก ภิกฺขุ ปีฐเก นิปนฺนํ ทาริกํ ปสฺสิตฺวา สารตฺโต องฺคุฏฺฐํ องฺคชาตํ ปเวเสสิ, สา กาลมกาสิ. ตสฺส กุกฺกุจฺจํ อโหสิ ฯเปฯ. อนาปตฺติ ภิกฺขุ ปาราชิกสฺส, อาปตฺติ สํฆาทิเสสสฺสาติ. (12)}

\proseitem{68}{เตน โข ปน สมเยน อญฺญตโร มาณวโก อุปฺปลวณฺณาย ภิกฺขุนิยา ปฏิพทฺธ\-จิตฺโต โหติ, อถ โข โส มาณวโก อุปฺปลวณฺณาย ภิกฺขุนิยา คามํ ปิณฺฑาย ปวิฏฺฐาย กุฏิกํ ปวิสิตฺวา นิลีโน อจฺฉิ. อุปฺปลวณฺณา ภิกฺขุนี ปจฺฉาภตฺตํ ปิณฺฑปาต\-ปฏิกฺกนฺตา ปาเท ปกฺขาเลตฺวา กุฏิกํ ปวิสิตฺวา มญฺจเก นิสีทิ. อถ โข โส มาณวโก อุปฺปลวณฺณํ ภิกฺขุนิํ อุคฺคเหตฺวา ทูเสสิ. อุปฺปลวณฺณา ภิกฺขุนี ภิกฺขุนีนํ เอตมตฺถํ อาโรเจสิ. ภิกฺขุนิโย ภิกฺขูนํ เอตมตฺถํ อาโรเจสุํ. ภิกฺขู ภควโต เอตมตฺถํ อาโรเจสุํ. อนาปตฺติ ภิกฺขเว อสาทิยนฺติยาติ. (13)}

\proseitem{69}{เตน โข ปน สมเยน อญฺญตรสฺส ภิกฺขุโน อิตฺถิลิงฺคํ ปาตุภูตํ โหติ. ภควโต เอตมตฺถํ อาโรเจสุํ. อนุชานามิ ภิกฺขเว ตํ เยว\csromansharedfootnote{cf7dc7e6d0b8f44d}{ตํเยว นิคฺคหีตสนฺธิ \csromandash{} ม.พ.ป.} อุปชฺฌํ ตเมว อุปสมฺปทํ ตานิเยว\csromansharedfootnote{13839f432a97e126}{ตานิ (สี, สฺยา)} วสฺสานิ ภิกฺขุนีหิ สงฺคมิตุํ\csromansharedfootnote{bdd8dd8e6c509882}{สงฺกมิตุํ (สี, สฺยา)}. ยา อาปตฺติโย ภิกฺขูนํ ภิกฺขุนีหิ สาธารณา, ตา อาปตฺติโย ภิกฺขุนีนํ สนฺติเก วุฏฺฐาตุํ. ยา อาปตฺติโย ภิกฺขูนํ ภิกฺขุนีหิ อสาธารณา, ตาหิ อาปตฺตีหิ อนาปตฺตีติ. (14)}

\prose{เตน โข ปน สมเยน อญฺญตริสฺสา ภิกฺขุนิยา ปุริสลิงฺคํ ปาตุภูตํ โหติ. ภควโต เอตมตฺถํ อาโรเจสุํ. อนุชานามิ ภิกฺขเว ตํ เยว\csromansharedfootnote{cf7dc7e6d0b8f44d}{ตํเยว นิคฺคหีตสนฺธิ \csromandash{} ม.พ.ป.} อุปชฺฌํ ตเมว อุปสมฺปทํ ตานิเยว\csromansharedfootnote{13839f432a97e126}{ตานิ (สี, สฺยา)} วสฺสานิ ภิกฺขูหิ สงฺคมิตุํ\csromansharedfootnote{bdd8dd8e6c509882}{สงฺกมิตุํ (สี, สฺยา)}. ยา \csromanfolio{44}อาปตฺติโย ภิกฺขุนีนํ ภิกฺขูหิ สาธารณา, ตา อาปตฺติโย ภิกฺขูนํ สนฺติเก วุฏฺฐาตุํ. ยา อาปตฺติโย ภิกฺขุนีนํ ภิกฺขูหิ อสาธารณา, ตาหิ อาปตฺตีหิ อนาปตฺตีติ. (15)}

\proseitem{70}{เตน โข ปน สมเยน อญฺญตโร ภิกฺขุ เอวํ เม อนาปตฺติ ภวิสฺสตีติ มาตุยา เมถุนํ ธมฺมํ ปฏิเสวิ. ธีตุยา เมถุนํ ธมฺมํ ปฏิเสวิ. ภคินิยา เมถุนํ ธมฺมํ ปฏิเสวิ. ตสฺส กุกฺกุจฺจํ อโหสิ ฯเปฯ. อาปตฺติํ ตฺวํ ภิกฺขุ อาปนฺโน ปาราชิกนฺติ. \mbox{(16-18)}}

\prose{เตน โข ปน สมเยน อญฺญตโร ภิกฺขุ ปุราณ\-ทุติยิกาย เมถุนํ ธมฺมํ ปฏิเสวิ, ตสฺส กุกฺกุจฺจํ อโหสิ ฯเปฯ. อาปตฺติํ ตฺวํ ภิกฺขุ อาปนฺโน ปาราชิกนฺติ. (19)}

\proseitem{71}{เตน โข ปน สมเยน อญฺญตโร ภิกฺขุ มุทุปิฏฺฐิโก โหติ. โส อนภิรติยา ปีฬิโต อตฺตโน องฺคชาตํ มุเขน อคฺคเหสิ. ตสฺสกุกฺกุจฺจํ\csromansharedfootnote{8e07e006f6a09b57}{ตสฺส กุกฺกุจฺจํ เทฺว ปทานิ \csromandash{} ม.พ.ป.} อโหสิ ฯเปฯ. อาปตฺติํ ตฺวํ ภิกฺขุ อาปนฺโน ปาราชิกนฺติ. (20)}

\prose{เตน โข ปน สมเยน อญฺญตโร ภิกฺขุ ลมฺพี โหติ. โส อนภิรติยา ปีฬิโต อตฺตโน องฺคชาตํ อตฺตโน วจฺจมคฺคํ ปเวเสสิ. ตสฺส กุกฺกุจฺจํ อโหสิ ฯเปฯ. อาปตฺติํ ตฺวํ ภิกฺขุ อาปนฺโน ปาราชิกนฺติ. (21)}

\prose{เตน โข ปน สมเยน อญฺญตโร ภิกฺขุ มตสรีรํ ปสฺสิ. ตสฺมิํ จ สรีเร องฺคชาต\-สามนฺตา วโณ โหติ. โส เอวํ เม อนาปตฺติ ภวิสฺสตีติ องฺคชาเต องฺคชาตํ ปเวเสตฺวา วเณน นีหริ. ตสฺส กุกฺกุจฺจํ อโหสิ ฯเปฯ. อาปตฺติํ ตฺวํ ภิกฺขุ อาปนฺโน ปาราชิกนฺติ. (22)}

\prose{เตน โข ปน สมเยน อญฺญตโร ภิกฺขุ มตสรีรํ ปสฺสิ. ตสฺมิํ จ สรีเร องฺคชาต\-สามนฺตา วโณ โหติ. โส เอวํ เม อนาปตฺติ ภวิสฺสตีติ วเณ องฺคชาตํ ปเวเสตฺวา องฺคชาเตน นีหริ. ตสฺส กุกฺกุจฺจํ อโหสิ ฯเปฯ. อาปตฺติํ ตฺวํ ภิกฺขุ อาปนฺโน ปาราชิกนฺติ. (23)}

\prose{เตน โข ปน สมเยน อญฺญตโร ภิกฺขุ สารตฺโต เลปจิตฺตสฺส นิมิตฺตํ องฺคชาเตน ฉุปิ. ตสฺส กุกฺกุจฺจํ อโหสิ ฯเปฯ. อนาปตฺติ ภิกฺขุ ปาราชิกสฺส, อาปตฺติ ทุกฺกฏสฺสาติ. (24)}

\prose{\csromanfolio{45}เตน โข ปน สมเยน อญฺญตโร ภิกฺขุ สารตฺโต ทารุ\-ธีตลิกาย นิมิตฺตํ องฺคชาเตน ฉุปิ. ตสฺส กุกฺกุจฺจํ อโหสิ ฯเปฯ. อนาปตฺติ ภิกฺขุ ปาราชิกสฺส, อาปตฺติ ทุกฺกฏสฺสาติ. (25)}

\proseitem{72}{เตน โข ปน สมเยน สุนฺทโร นาม ภิกฺขุ ราชคหา ปพฺพชิโต รถิกาย\csromansharedfootnote{7608ea908ad7b28f}{รถิยาย (ก)} คจฺฉติ. อญฺญตรา อิตฺถี “มุหุตฺตํ\csromansharedfootnote{627493397976a09e}{อิตฺถี ตํ ปสฺสิตฺวา เอตทโวจ มุหุตฺตํ (สฺยา)} ภนฺเต อาคเมหิ วนฺทิสฺสามี”ติ สา วนฺทนฺตี อนฺตรวาสกํ อุกฺขิปิตฺวา มุเขน องฺคชาตํ อคฺคเหสิ. ตสฺส กุกฺกุจฺจํ อโหสิ ฯเปฯ. สาทิยิ ตฺวํ ภิกฺขูติ. นาหํ ภควา สาทิยินฺติ. อนาปตฺติ ภิกฺขุ อสาทิยนฺตสฺสาติ. (26)}

\prose{เตน โข ปน สมเยน อญฺญตรา อิตฺถี ภิกฺขุํ ปสฺสิตฺวา เอตทโวจ “เอหิ ภนฺเต เมถุนํ ธมฺมํ ปฏิเสวา”ติ. อลํ ภคินิ เนตํ กปฺปตีติ. เอหิ ภนฺเต อหํ วายมิสฺสามิ, ตฺวํ มา วายมิ, เอวํ เต อนาปตฺติ ภวิสฺสตีติ. โส ภิกฺขุ ตถา อกาสิ. ตสฺส กุกฺกุจฺจํ อโหสิ ฯเปฯ. อาปตฺติํ ตฺวํ ภิกฺขุ อาปนฺโน ปาราชิกนฺติ. (27)}

\prose{เตน โข ปน สมเยน อญฺญตรา อิตฺถี ภิกฺขุํ ปสฺสิตฺวา เอตทโวจ “เอหิ ภนฺเต เมถุนํ ธมฺมํ ปฏิเสวา”ติ. อลํ ภคินิ เนตํ กปฺปตีติ. เอหิ ภนฺเต ตฺวํ วายม, อหํ น วายมิสฺสามิ, เอวํ เต อนาปตฺติ ภวิสฺสตีติ. โส ภิกฺขุ ตถา อกาสิ. ตสฺส กุกฺกุจฺจํ อโหสิ ฯเปฯ. อาปตฺติํ ตฺวํ ภิกฺขุ อาปนฺโน ปาราชิกนฺติ. (28)}

\prose{เตน โข ปน สมเยน อญฺญตรา อิตฺถี ภิกฺขุํ ปสฺสิตฺวา เอตทโวจ “เอหิ ภนฺเต เมถุนํ ธมฺมํ ปฏิเสวา”ติ. อลํ ภคินิ เนตํ กปฺปตีติ. เอหิ ภนฺเต อพฺภนฺตรํ ฆฏฺเฏตฺวา พหิ โมเจหิ ฯเปฯ พหิ ฆฏฺเฏตฺวา อพฺภนฺตรํ โมเจหิ, เอวํ เต อนาปตฺติ ภวิสฺสตีติ. โส ภิกฺขุ ตถา อกาสิ. ตสฺส กุกฺกุจฺจํ อโหสิ ฯเปฯ. อาปตฺติํ ตฺวํ ภิกฺขุ อาปนฺโน ปาราชิกนฺติ. \mbox{(29-30)}}

\proseitem{73}{เตน โข ปน สมเยน อญฺญตโร ภิกฺขุ สิวถิกํ คนฺตฺวา อกฺขายิตํ สรีรํ ปสฺสิตฺวา ตสฺมิํ เมถุนํ ธมฺมํ ปฏิเสวิ. ตสฺส กุกฺกุจฺจํ อโหสิ ฯเปฯ. อาปตฺติํ ตฺวํ ภิกฺขุ อาปนฺโน ปาราชิกนฺติ. (31)}

\prose{\csromanfolio{46}เตน โข ปน สมเยน อญฺญตโร ภิกฺขุ สิวถิกํ คนฺตฺวา เยภุยฺเยน อกฺขายิตํ สรีรํ ปสฺสิตฺวา ตสฺมิํ เมถุนํ ธมฺมํ ปฏิเสวิ. ตสฺส กุกฺกุจฺจํ อโหสิ ฯเปฯ. อาปตฺติํ ตฺวํ ภิกฺขุ อาปนฺโน ปาราชิกนฺติ. (32)}

\prose{เตน โข ปน สมเยน อญฺญตโร ภิกฺขุ สิวถิกํ คนฺตฺวา เยภุยฺเยน ขายิตํ สรีรํ ปสฺสิตฺวา ตสฺมิํ เมถุนํ ธมฺมํ ปฏิเสวิ. ตสฺส กุกฺกุจฺจํ อโหสิ ฯเปฯ. อนาปตฺติ ภิกฺขุ ปาราชิกสฺส, อาปตฺติ ถุลฺลจฺจยสฺสาติ. (33)}

\prose{เตน โข ปน สมเยน อญฺญตโร ภิกฺขุ สิวถิกํ คนฺตฺวา ฉินฺนสีสํ ปสฺสิตฺวา วฏฺฏกเต มุเข ฉุปนฺตํ องฺคชาตํ ปเวเสสิ. ตสฺส กุกฺกุจฺจํ อโหสิ ฯเปฯ. อาปตฺติํ ตฺวํ ภิกฺขุ อาปนฺโน ปาราชิกนฺติ. (34)}

\prose{เตน โข ปน สมเยน อญฺญตโร ภิกฺขุ สิวถิกํ คนฺตฺวา ฉินฺนสีสํ ปสฺสิตฺวา วฏฺฏกเต มุเข อจฺฉุปนฺตํ องฺคชาตํ ปเวเสสิ. ตสฺส กุกฺกุจฺจํ อโหสิ ฯเปฯ. อนาปตฺติ ภิกฺขุ ปาราชิกสฺส, อาปตฺติ ทุกฺกฏสฺสาติ. (35)}

\prose{เตน โข ปน สมเยน อญฺญตโร ภิกฺขุ อญฺญตริสฺสา อิตฺถิยา ปฏิพทฺธ\-จิตฺโต โหติ. สา กาลงฺกตา\csromansharedfootnote{eec123c088830038}{กาลกตา (สี, สฺยา)} สุสาเน ฉฑฺฑิตา อฏฺฐิกานิ วิปฺปกิณฺณานิ โหนฺติ. อถ โข โส ภิกฺขุ สิวถิกํ คนฺตฺวา อฏฺฐิกานิ สงฺกฑฺฒิตฺวา นิมิตฺเต องฺคชาตํ ปฏิปาเทสิ. ตสฺส กุกฺกุจฺจํ อโหสิ ฯเปฯ. อนาปตฺติ ภิกฺขุ ปาราชิกสฺส, อาปตฺติ ทุกฺกฏสฺสาติ. (36)}

\prose{เตน โข ปน สมเยน อญฺญตโร ภิกฺขุ นาคิยา เมถุนํ ธมฺมํ ปฏิเสวิ. ยกฺขินิยา เมถุนํ ธมฺมํ ปฏิเสวิ. เปติยา เมถุนํ ธมฺมํ ปฏิเสวิ. ปณฺฑกสฺส เมถุนํ ธมฺมํ ปฏิเสวิ. ตสฺส กุกฺกุจฺจํ อโหสิ ฯเปฯ. อาปตฺติํ ตฺวํ ภิกฺขุ อาปนฺโน ปาราชิกนฺติ. \mbox{(37-40)}}

\prose{เตน โข ปน สมเยน อญฺญตโร ภิกฺขุ อุปหตินฺทฺริโย โหติ. โส “นาหํ เวทิยามิ\csromansharedfootnote{e61cfdbc2a037366}{เวทยามิ (ก)} สุขํ วา ทุกฺขํ วา, อนาปตฺติ เม ภวิสฺสตี”ติ เมถุนํ ธมฺมํ ปฏิเสวิ ฯเปฯ. ภควโต เอตมตฺถํ อาโรเจสุํ. เวทิยิ วา โส ภิกฺขเว โมฆปุริโส น วา เวทิยิ, อาปตฺติ ปาราชิกสฺสาติ. (41)}

\prose{\csromanfolio{47}เตน โข ปน สมเยน อญฺญตโร ภิกฺขุ อิตฺถิยา เมถุนํ ธมฺมํ ปฏิเสวิสฺสามีติ ฉุปิตมตฺเต วิปฺปฏิสารี อโหสิ. ตสฺส กุกฺกุจฺจํ อโหสิ ฯเปฯ. อนาปตฺติ ภิกฺขุ ปาราชิกสฺส, อาปตฺติ สํฆาทิเสสสฺสาติ. (42)}

\proseitem{74}{เตน โข ปน สมเยน อญฺญตโร ภิกฺขุ ภทฺทิเย ชาติยาวเน ทิวาวิหารคโต นิปนฺโน โหติ. ตสฺส องฺคมงฺคานิ วาตู\-ปตฺถทฺธานิ โหนฺติ. อญฺญตรา อิตฺถี ปสฺสิตฺวา องฺคชาเต อภินิสีทิตฺวา ยาวทตฺถํ กตฺวา ปกฺกามิ. ภิกฺขู กิลินฺนํ ปสฺสิตฺวา ภควโต เอตมตฺถํ อาโรเจสุํ. ปญฺจหิ ภิกฺขเว อากาเรหิ องฺคชาตํ กมฺมนิยํ โหติ ราเคน วจฺเจน ปสฺสาเวน วาเตน อุจฺจาลิงฺคปาณก\-ทฏฺเฐน, อิเมหิ โข ภิกฺขเว ปญฺจหากาเรหิ องฺคชาตํ กมฺมนิยํ โหติ. อฏฺฐานเมตํ ภิกฺขเว อนวกาโส, ยํ ตสฺส ภิกฺขุโน ราเคน องฺคชาตํ กมฺมนิยํ อสฺส, อรหํ โส ภิกฺขเว ภิกฺขุ, อนาปตฺติ ภิกฺขเว ตสฺส ภิกฺขุโนติ. (43)}

\prose{เตน โข ปน สมเยน อญฺญตโร ภิกฺขุ สาวตฺถิยา อนฺธวเน ทิวาวิหารคโต นิปนฺโน โหติ. อญฺญตรา โคปาลิกา ปสฺสิตฺวา องฺคชาเต อภินิสีทิ. โส ภิกฺขุ ปเวสนํ สาทิยิ, ปวิฏฺฐํ สาทิยิ, ฐิตํ สาทิยิ, อุทฺธรณํ สาทิยิ. ตสฺส กุกฺกุจฺจํ อโหสิ ฯเปฯ. อาปตฺติํ ตฺวํ ภิกฺขุ อาปนฺโน ปาราชิกนฺติ. (44)}

\prose{เตน โข ปน สมเยน อญฺญตโร ภิกฺขุ สาวตฺถิยา อนฺธวเน ทิวาวิหารคโต นิปนฺโน โหติ. อญฺญตรา อชปาลิกา ปสฺสิตฺวา. อญฺญตรา กฏฺฐหาริกา ปสฺสิตฺวา. อญฺญตรา โคมยหาริกา ปสฺสิตฺวา องฺคชาเต อภินิสีทิ. โส ภิกฺขุ ปเวสนํ สาทิยิ, ปวิฏฺฐํ สาทิยิ, ฐิตํ สาทิยิ, อุทฺธรณํ สาทิยิ. ตสฺส กุกฺกุจฺจํ อโหสิ ฯเปฯ. อาปตฺติํ ตฺวํ ภิกฺขุ อาปนฺโน ปาราชิกนฺติ. \mbox{(45-47)}}

\proseitem{75}{เตน โข ปน สมเยน อญฺญตโร ภิกฺขุ เวสาลิยํ มหาวเน ทิวาวิหารคโต นิปนฺโน โหติ. อญฺญตรา อิตฺถี ปสฺสิตฺวา องฺคชาเต อภินิสีทิตฺวา ยาวทตฺถํ กตฺวา สามนฺตา หสมานา ฐิตา \csromanfolio{48}โหติ. โส ภิกฺขุ ปฏิพุชฺฌิตฺวา ตํ อิตฺถิํ เอตทโวจ “ตุยฺหิทํ กมฺมนฺ”ติ. อาม มยฺหํ กมฺมนฺติ. ตสฺส กุกฺกุจฺจํ อโหสิ ฯเปฯ. สาทิยิ ตฺวํ ภิกฺขูติ. นาหํ ภควา ชานามีติ. อนาปตฺติ ภิกฺขุ อชานนฺตสฺสาติ. (48)}

\proseitem{76}{เตน โข ปน สมเยน อญฺญตโร ภิกฺขุ เวสาลิยํ มหาวเน ทิวาวิหารคโต รุกฺขํ อปสฺสาย นิปนฺโน โหติ. อญฺญตรา อิตฺถี ปสฺสิตฺวา องฺคชาเต อภินิสีทิ. โส ภิกฺขุ สหสา วุฏฺฐาสิ, ตสฺส กุกฺกุจฺจํ อโหสิ ฯเปฯ. สาทิยิ ตฺวํ ภิกฺขูติ. นาหํ ภควา สาทิยินฺติ. อนาปตฺติ ภิกฺขุ อสาทิยนฺตสฺสาติ. (49)}

\prose{เตน โข ปน สมเยน อญฺญตโร ภิกฺขุ เวสาลิยํ มหาวเน ทิวาวิหารคโต รุกฺขํ อปสฺสาย นิปนฺโน โหติ. อญฺญตรา อิตฺถี ปสฺสิตฺวา องฺคชาเต อภินิสีทิ. โส ภิกฺขุ อกฺกมิตฺวา ปวตฺเตสิ\csromansharedfootnote{3d4552d890bd3813}{ปวฏฺเฏสิ (สี, สฺยา)}. ตสฺส กุกฺกุจฺจํ อโหสิ ฯเปฯ. สาทิยิ ตฺวํ ภิกฺขูติ. นาหํ ภควา สาทิยินฺติ. อนาปตฺติ ภิกฺขุ อสาทิยนฺตสฺสาติ. (50)}

\proseitem{77}{เตน โข ปน สมเยน อญฺญตโร ภิกฺขุ เวสาลิยํ มหาวเน กูฏาคาร\-สาลายํ ทิวาวิหารคโต ทฺวารํ วิวริตฺวา นิปนฺโน โหติ. ตสฺส องฺคมงฺคานิ วาตู\-ปตฺถทฺธานิ โหนฺติ. เตน โข ปน สมเยน สมฺพหุลา อิตฺถิโย คนฺธญฺจ มาลญฺจ อาทาย อารามํ อาคมํสุ วิหารเปกฺขิกาโย. อถ โข ตา อิตฺถิโย ตํ ภิกฺขุํ ปสฺสิตฺวา องฺคชาเต อภินิสีทิตฺวา ยาวทตฺถํ กตฺวา “ปุริสูสโภ วตายนฺ”ติ วตฺวา คนฺธญฺจ มาลญฺจ อาโรเปตฺวา ปกฺกมิํสุ. ภิกฺขู กิลินฺนํ ปสฺสิตฺวา ภควโต เอตมตฺถํ อาโรเจสุํ. ปญฺจหิ ภิกฺขเว อากาเรหิ องฺคชาตํ กมฺมนิยํ โหติ ราเคน วจฺเจน ปสฺสาเวน วาเตน อุจฺจาลิงฺคปาณก\-ทฏฺเฐน, อิเมหิ โข ภิกฺขเว ปญฺจหากาเรหิ องฺคชาตํ กมฺมนิยํ โหติ. อฏฺฐานเมตํ ภิกฺขเว อนวกาโส, ยํ ตสฺส ภิกฺขุโน ราเคน องฺคชาตํ กมฺมนิยํ อสฺส, อรหํ โส ภิกฺขเว ภิกฺขุ, อนาปตฺติ ภิกฺขเว ตสฺส ภิกฺขุโน. อนุชานามิ ภิกฺขเว ทิวา ปฏิสลฺลียนฺเตน ทฺวารํ สํวริตฺวา ปฏิสลฺลียิตุนฺติ. (51)}

\proseitem{78}{\csromanfolio{49}เตน โข ปน สมเยน อญฺญตโร ภารุกจฺฉโก ภิกฺขุ สุปินนฺเต\csromansharedfootnote{78563ce828708c85}{สุปินนฺเตน (สี, สฺยา)} ปุราณ\-ทุติยิกาย เมถุนํ ธมฺมํ ปฏิเสวิตฺวา “อสฺสมโณ อหํ วิพฺภมิสฺสามี”ติ ภารุกจฺฉํ คจฺฉนฺโต อนฺตรามคฺเค อายสฺมนฺตํ อุปาลิํ ปสฺสิตฺวา เอตมตฺถํ อาโรเจสิ. อายสฺมา อุปาลิ เอวมาห “อนาปตฺติ อาวุโส สุปินนฺเตนา”ติ. (52)}

\prose{เตน โข ปน สมเยน ราชคเห สุปพฺพา นาม อุปาสิกา มุธปฺปสนฺนา\csromansharedfootnote{d28eb23bb7399eb6}{มุทฺธปฺปสนฺนา (สี)} โหติ. สา เอวํทิฏฺฐิกา โหติ “ยา เมถุนํ ธมฺมํ เทติ, สา อคฺคทานํ เทตี”ติ. สา ภิกฺขุํ ปสฺสิตฺวา เอตทโวจ “เอหิ ภนฺเต เมถุนํ ธมฺมํ ปฏิเสวา”ติ. อลํ ภคินิ เนตํ กปฺปตีติ. เอหิ ภนฺเต อูรุนฺตริกาย\csromansharedfootnote{9b4227a4f7e87a9b}{อูรนฺตริกาย (สี)} ฆฏฺเฏหิ. เอวํ เต อนาปตฺติ ภวิสฺสตีติ ฯเปฯ. เอหิ ภนฺเต นาภิยํ ฆฏฺเฏหิ. เอหิ ภนฺเต อุทรวฏฺฏิยํ ฆฏฺเฏหิ. เอหิ ภนฺเต อุปกจฺฉเก ฆฏฺเฏหิ. เอหิ ภนฺเต คีวายํ ฆฏฺเฏหิ. เอหิ ภนฺเต กณฺณจฺฉิทฺเท ฆฏฺเฏหิ. เอหิ ภนฺเต เกสวฏฺฏิยํ ฆฏฺเฏหิ. เอหิ ภนฺเต องฺคุล\-นฺตริกาย ฆฏฺเฏหิ. เอหิ ภนฺเต หตฺเถน อุปกฺกมิตฺวา โมเจสฺสามิ, เอวํ เต อนาปตฺติ ภวิสฺสตีติ. โส ภิกฺขุ ตถา อกาสิ. ตสฺส กุกฺกุจฺจํ อโหสิ ฯเปฯ. อนาปตฺติ ภิกฺขุ ปาราชิกสฺส, อาปตฺถิ สํฆาทิเสสสฺสาติ. \mbox{(53-61)}}

\proseitem{79}{เตน โข ปน สมเยน สาวตฺถิยํ สทฺธา นาม อุปาสิกา มุธปฺปสนฺนา โหติ. สา เอวํทิฏฺฐิกา โหติ “ยา เมถุนํ ธมฺมํ เทติ, สา อคฺคทานํ เทตี”ติ. สา ภิกฺขุํ ปสฺสิตฺวา เอตทโวจ “เอหิ ภนฺเต เมถุนํ ธมฺมํ ปฏิเสวา”ติ. อลํ ภคินิ เนตํ กปฺปตีติ. เอหิ ภนฺเต อูรุนฺตริกาย ฆฏฺเฏหิ ฯเปฯ. เอหิ ภนฺเต หตฺเถน อุปกฺกมิตฺวา โมเจสฺสามิ, เอวํ เต อนาปตฺติ ภวิสฺสตีติ. โส ภิกฺขุ ตถา อกาสิ. ตสฺส กุกฺกุจฺจํ อโหสิ ฯเปฯ. อนาปตฺติ ภิกฺขุ ปาราชิกสฺส, อาปตฺติ สํฆาทิเสสสฺสาติ. \mbox{(62-70)}}

\proseitem{80}{เตน โข ปน สมเยน เวสาลิยํ ลิจฺฉวิ\-กุมารกา ภิกฺขุํ คเหตฺวา ภิกฺขุนิยา วิปฺปฏิปาเทสุํ. สิกฺขมานาย วิปฺปฏิปาเทสุํ. สามเณริยา วิปฺปฏิปาเทสุํ. อุโภ สาทิยิํสุ, อุโภ นาเสตพฺพา. อุโภ น สาทิยิํสุ, อุภินฺนํ อนาปตฺติ. \mbox{(71-76)}}

\csromanmatikaanchor{mtk.15}
\csromanmatikaanchor{mtk.16}
\csromantocmarknopage{section}{2. ทุติยปาราชิก}{mtk.15}
\bookmark[dest={mtk.15},level=1]{2. ทุติยปาราชิก}
\csromantocmarkref{subsection}{ธนิยภิกฺขุวตฺถุ}{mtk.16}
\bookmark[dest={mtk.16},level=2]{ธนิยภิกฺขุวตฺถุ}
\proseitem{81}{\csromanfolio{50}เตน โข ปน สมเยน เวสาลิยํ ลิจฺจวิกุมารกา ภิกฺขุํ คเหตฺวา เวสิยา วิปฺปฏิปาเทสุํ. ปณฺฑเก วิปฺปฏิปาเทสุํ. คิหินิยา วิปฺปฏิปาเทสุํ. ภิกฺขุ สาทิยิ, ภิกฺขุ นาเสตพฺโพ. ภิกฺขุ น สาทิยิ, ภิกฺขุสฺส อนาปตฺติ. \mbox{(77-82)}}

\prose{เตน โข ปน สมเยน เวสาลิยํ ลิจฺฉวิ\-กุมารกา ภิกฺขู คเหตฺวา อญฺญมญฺญํ วิปฺปฏิปาเทสุํ. อุโภ สาทิยิํสุ, อุโภ นาเสตพฺพา. อุโภ น สาทิยิํสุ, อุภินฺนํ อนาปตฺติ. \mbox{(83-84)}}

\proseitem{82}{เตน โข ปน สมเยน อญฺญตโร วุฑฺฒ\-ปพฺพชิโต ภิกฺขุ ปุราณ\-ทุติยิกาย ทสฺสนํ อคมาสิ. สา “เอหิ ภนฺเต วิพฺภมา”ติ อคฺคเหสิ. โส ภิกฺขุ ปฏิกฺกมนฺโต อุตฺตาโน ปริปติ. สา อุพฺภชิตฺวา\csromansharedfootnote{1192f0aabc11ca8b}{อุพฺภุชิตฺวา (สี, สฺยา)} องฺคชาเต\csromansharedfootnote{f6a28f308941ecb5}{องฺคชาเตน (สี)} อภินิสีทิ. ตสฺส กุกฺกุจฺจํ อโหสิ ฯเปฯ. สาทิยิ ตฺวํ ภิกฺขูติ. นาหํ ภควา สาทิยินฺติ. อนาปตฺติ ภิกฺขุ อสาทิยนฺตสฺสาติ. (85)}

\proseitemwithcloserruled{83}{เตน โข ปน สมเยน อญฺญตโร ภิกฺขุ อรญฺเญ วิหรติ. มิคโปตโก ตสฺส ปสฺสาวฏฺฐานํ อาคนฺตฺวา ปสฺสาวํ ปิวนฺโต มุเขน องฺคชาตํ อคฺคเหสิ. โส ภิกฺขุ สาทิยิ. ตสฺส กุกฺกุจฺจํ อโหสิ ฯเปฯ. อาปตฺติํ ตฺวํ ภิกฺขุ อาปนฺโน ปาราชิกนฺติ. (86)}{ปฐม\-ปาราชิกํ สมตฺตํ.}
\csromanheader{1}{2. \textbf{ทุติยปาราชิก}}
\proseitem{84}{เตน สมเยน พุทฺโธ ภควา ราชคเห วิหรติ คิชฺฌกูเฏ ปพฺพเต. เตน โข ปน สมเยน สมฺพหุลา สนฺทิฏฺฐา สมฺภตฺตา ภิกฺขู อิสิคิลิปสฺเส ติณกุฏิโย กริตฺวา วสฺสํ อุปคจฺฉิํสุ. อายสฺมาปิ ธนิโย กุมฺภการ\-ปุตฺโต ติณกุฏิกํ กริตฺวา วสฺสํ อุปคจฺฉิ. อถ โข \csromanfolio{51}เต ภิกฺขู วสฺสํวุฏฺฐา เตมาสจฺจเยน ติณกุฏิโย ภินฺทิตฺวา ติณญฺจ กฏฺฐญฺจ ปฏิสาเมตฺวา ชนปท\-จาริกํ ปกฺกมิํสุ. อายสฺมา ปน ธนิโย กุมฺภการ\-ปุตฺโต ตตฺเถว วสฺสํ วสิ, ตตฺถ เหมนฺตํ, ตตฺถ คิมฺหํ. อถ โข อายสฺมโต ธนิยสฺส กุมฺภการ\-ปุตฺตสฺส คามํ ปิณฺฑาย ปวิฏฺฐสฺส ติณหาริโย กฏฺฐหาริโย ติณกุฏิกํ ภินฺทิตฺวา ติณญฺจ กฏฺฐญฺจ อาทาย อคมํสุ. ทุติยมฺปิ โข อายสฺมา ธนิโย กุมฺภการ\-ปุตฺโต ติณญฺจ กฏฺฐญฺจ สํกฑฺฒิตฺวา ติณกุฏิกํ อกาสิ. ทุติยมฺปิ โข อายสฺมโต ธนิยสฺส กุมฺภการ\-ปุตฺตสฺส คามํ ปิณฺฑาย ปวิฏฺฐสฺส ติณหาริโย กฏฺฐหาริโย ติณกุฏิกํ ภินฺทิตฺวา ติณญฺจ กฏฺฐญฺจ อาทาย อคมํสุ. ตติยมฺปิ โข อายสฺมา ธนิโย กุมฺภการ\-ปุตฺโต ติณญฺจ กฏฺฐญฺจ สํกฑฺฒิตฺวา ติณกุฏิกํ อกาสิ. ตติยมฺปิ โข อายสฺมโต ธนิยสฺส กุมฺภการ\-ปุตฺตสฺส คามํ ปิณฺฑาย ปวิฏฺฐสฺส ติณหาริโย กฏฺฐหาริโย ติณกุฏิกํ ภินฺทิตฺวา ติณญฺจ กฏฺฐญฺจ อาทาย อคมํสุ.}

\prose{อถ โข อายสฺมโต ธนิยสฺส กุมฺภการ\-ปุตฺตสฺส เอตทโหสิ “ยาวตติยกํ โข เม คามํ ปิณฺฑาย ปวิฏฺฐสฺส ติณหาริโย กฏฺฐหาริโย ติณกุฏิกํ ภินฺทิตฺวา ติณญฺจ กฏฺฐญฺจ อาทาย อคมํสุ, อหํ โข ปน สุสิกฺขิโต อนวโย สเก อาจริยเก กุมฺภการ\-กมฺเม ปริโยทาต\-สิปฺโป, ยํนูนาหํ สามํ จิกฺขลฺลํ มทฺทิตฺวา สพฺพ\-มตฺติกามยํ กุฏิกํ กเรยฺยนฺ”ติ. อถ โข อายสฺมา ธนิโย กุมฺภการ\-ปุตฺโต สามํ จิกฺขลฺลํ มทฺทิตฺวา สพฺพ\-มตฺติกามยํ กุฏิกํ กริตฺวา ติณญฺจ กฏฺฐญฺจ โคมยญฺจ สํกฑฺฒิตฺวา ตํ กุฏิกํ ปจิ, สา อโหสิ กุฏิกา อภิรูปา ทสฺสนียา ปาสาทิกา โลหิติกา\csromansharedfootnote{d74d0c5328486f8d}{โลหิตกา (สฺยา)}, เสยฺยถาปิ อินฺทโคปโก. เสยฺยถาปิ นาม กิงฺกณิกสทฺโท\csromansharedfootnote{c2ca6fe1b682b877}{กิงฺกิณิกสทฺโท (สี, สฺยา)}, เอวเมวํ ตสฺสา กุฏิกาย สทฺโท อโหสิ.}

\proseitem{85}{อถ โข ภควา สมฺพหุเลหิ ภิกฺขูหิ สทฺธิํ คิชฺฌกูฏา ปพฺพตา โอโรหนฺโต อทฺทส ตํ กุฏิกํ อภิรูปํ ทสฺสนียํ ปาสาทิกํ โลหิติกํ, ทิสฺวาน ภิกฺขู อามนฺเตสิ “กิํ เอตํ ภิกฺขเว อภิรูปํ ทสฺสนียํ ปาสาทิกํ โลหิติกํ, เสยฺยถาปิ อินฺทโคปโก”ติ. อถ โข เต ภิกฺขู ภควโต เอตมตฺถํ อาโรเจสุํ. วิครหิ พุทฺโธ ภควา “อนนุจฺฉวิกํ ภิกฺขเว ตสฺส โมฆปุริสสฺส อนนุโลมิกํ อปฺปติรูปํ \csromanfolio{52}อสฺสามณกํ อกปฺปิยํ อกรณียํ, กถํ หิ นาม โส ภิกฺขเว โมฆปุริโส สพฺพ\-มตฺติกามยํ กุฏิกํ กริสฺสติ, น หิ นาม ภิกฺขเว ตสฺส โมฆปุริสสฺส ปาเณสุ อนุทฺทยา อนุกมฺปา อวิเหสา ภวิสฺสติ. คจฺฉเถตํ ภิกฺขเว กุฏิกํ ภินฺทถ, มา ปจฺฉิมา ชนตา ปาเณสุ ปาตพฺยตํ อาปชฺชิ. น จ ภิกฺขเว สพฺพ\-มตฺติกามยา กุฏิกา กาตพฺพา, โย กเรยฺย, อาปตฺติ ทุกฺกฏสฺสา”ติ. เอวํ ภนฺเตติ โข เต ภิกฺขู ภควโต ปฏิสฺสุณิตฺวา เยน สา กุฏิกา เตนุ\-ปสงฺกมิํสุ, อุปสงฺกมิตฺวา ตํ กุฏิกํ ภินฺทิํสุ. อถ โข อายสฺมา ธนิโย กุมฺภการ\-ปุตฺโต เต ภิกฺขู เอตทโวจ “กิสฺส เม ตุเมฺห อาวุโส กุฏิกํ ภินฺทถา”ติ. ภควา อาวุโส เภทาเปตีติ. ภินฺทถาวุโส สเจ ธมฺมสฺสามี เภทาเปตีติ.}

\proseitem{86}{อถ โข อายสฺมโต ธนิยสฺส กุมฺภการ\-ปุตฺตสฺส เอตทโหสิ “ยาวตติยกํ โข เม คามํ ปิณฺฑาย ปวิฏฺฐสฺส ติณหาริโย กฏฺฐหาริโย ติณกุฏิกํ ภินฺทิตฺวา ติณญฺจ กฏฺฐญฺจ อาทาย อคมํสุ, ยาปิ มยา สพฺพ\-มตฺติกามยา กุฏิกา กตา, สาปิ ภควตา เภทาปิตา, อตฺถิ จ เม ทารุคเห คณโก สนฺทิฏฺโฐ, ยํนูนาหํ ทารุคเห คณกํ ทารูนิ ยาจิตฺวา ทารุกุฏิกํ กเรยฺยนฺ”ติ. อถ โข อายสฺมา ธนิโย กุมฺภการ\-ปุตฺโต เยน ทารุคเห คณโก เตนุปสงฺกมิ, อุปสงฺกมิตฺวา ทารุคเห คณกํ เอตทโวจ “ยาวตติยกํ โข เม อาวุโส คามํ ปิณฺฑาย ปวิฏฺฐสฺส ติณหาริโย กฏฺฐหาริโย ติณกุฏิกํ ภินฺทิตฺวา ติณญฺจ กฏฺฐญฺจ อาทาย อคมํสุ, ยาปิ มยา สพฺพ\-มตฺติกามยา กุฏิกา กตา, สาปิ ภควตา เภทาปิตา, เทหิ เม อาวุโส ทารูนิ อิจฺฉามิ ทารุกุฏิกํ\csromansharedfootnote{71a242f5f77fd376}{ทารุกุฑฺฑิกํ กุฏิกํ (สี)} กาตุนฺ”ติ. นตฺถิ ภนฺเต ตาทิสานิ ทารูนิ, ยานาหํ อยฺยสฺส ทเทยฺยํ, อตฺถิ ภนฺเต เทวคหทารูนิ นคร\-ปฏิสงฺขาริกานิ อาปทตฺถาย นิกฺขิตฺตานิ, สเจ ตานิ ทารูนิ ราชา ทาเปติ หราเปถ ภนฺเตติ. ทินฺนานิ อาวุโส รญฺญาติ. อถ โข ทารุคเห คณกสฺส เอตทโหสิ “อิเม โข สมณา สกฺยปุตฺติยา ธมฺมจาริโน สมจาริโน\csromansharedfootnote{94476af58838d477}{สมฺมจาริโน (ก)} พฺรหฺมจริโน สจฺจวาทิโน สีลวนฺโต กลฺยาณธมฺมา, ราชาปิเมสํ อภิปฺปสนฺโน, นารหติ อทินฺนํ ทินฺนนฺติ วตฺตุนฺ”ติ. อถ โข \csromanfolio{53}ทารุคเห คณโก อายสฺมนฺตํ ธนิยํ กุมฺภการ\-ปุตฺตํ เอตทโวจ “หราเปถ ภนฺเต”ติ. อถ โข อายสฺมา ธนิโย กุมฺภการ\-ปุตฺโต ตานิ ทารูนิ ขณฺฑาขณฺฑิกํ เฉทาเปตฺวา สกเฏหิ นิพฺพาหาเปตฺวา ทารุกุฏิกํ อกาสิ.}

\proseitem{87}{อถ โข วสฺสกาโร พฺราหฺมโณ มคธ\-มหามตฺโต ราชคเห กมฺมนฺเต อนุสญฺญายมาโน เยน ทารุคเห คณโก เตนุปสงฺกมิ, อุปสงฺกมิตฺวา ทารุคเห คณกํ เอตทโวจ “ยานิ ตานิ ภเณ เทวคหทารูนิ นคร\-ปฏิสงฺขาริกานิ อาปทตฺถาย นิกฺขิตฺตานิ, กหํ ตานิ ทารูนี”ติ. ตานิ สามิ ทารูนิ เทเวน อยฺยสฺส ธนิยสฺส กุมฺภการ\-ปุตฺตสฺส ทินฺนานีติ. อถโข\csromansharedfootnote{b8dfa686fd155ee3}{อถ โข เทฺว นิปาตา \csromandash{} ม.พ.ป.} วสฺสกาโร พฺราหฺมโณ มคธ\-มหามตฺโต อนตฺตมโน อโหสิ “กถํ หิ นาม เทโว เทวคหทารูนิ นคร\-ปฏิสงฺขาริกานิ อาปทตฺถาย นิกฺขิตฺตานิ ธนิยสฺส กุมฺภการ\-ปุตฺตสฺส ทสฺสตี”ติ. อถโข\csromansharedfootnote{b8dfa686fd155ee3}{อถ โข เทฺว นิปาตา \csromandash{} ม.พ.ป.} วสฺสกาโร พฺราหฺมโณ มคธ\-มหามตฺโต เยน ราชา มาคโธ เสนิโย พิมฺพิสาโร เตนุปสงฺกมิ, อุปสงฺกมิตฺวา ราชานํ มาคธํ เสนิยํ พิมฺพิสารํ เอตทโวจ “สจฺจํ กิร เทเวน\csromansharedfootnote{7cf73981753dc88c}{สจฺจํ กิร เทว เทเวน (สี)} เทวคหทารูนิ นคร\-ปฏิสงฺขาริกานิ อาปทตฺถาย นิกฺขิตฺตานิ ธนิยสฺส กุมฺภการ\-ปุตฺตสฺส ทินฺนานี”ติ. โก เอวมาหาติ. ทารุคเห คณโก เทวาติ. เตน หิ พฺราหฺมณ ทารุคเห คณกํ อาณาเปหีติ. อถ โข วสฺสกาโร พฺราหฺมโณ มคธ\-มหามตฺโต ทารุคเห คณกํ พนฺธํ\csromansharedfootnote{b0b14e5888d7aa9e}{พทฺธํ (สี)} อาณาเปสิ. อทฺทส โข อายสฺมา ธนิโย กุมฺภการ\-ปุตฺโต ทารุคเห คณกํ พนฺธํ นิยฺยมานํ, ทิสฺวาน ทารุคเห คณกํ เอตทโวจ “กิสฺส ตฺวํ อาวุโส พนฺโธ นิยฺยาสี”ติ. เตสํ ภนฺเต ทารูนํ กิจฺจาติ. คจฺฉาวุโส อหํปิ อาคจฺฉามีติ. เอยฺยาสิ ภนฺเต ปุราหํ หญฺญามีติ.}

\proseitem{88}{อถ โข อายสฺมา ธนิโย กุมฺภการ\-ปุตฺโต เยน รญฺโญ มาคธสฺส เสนิยสฺส พิมฺพิสารสฺส นิเวสนํ เตนุปสงฺกมิ, อุปสงฺกมิตฺวา ปญฺญตฺเต อาสเน นิสีทิ. อถ โข ราชา มาคโธ เสนิโย พิมฺพิสาโร เยนายสฺมา ธนิโย กุมฺภการ\-ปุตฺโต เตนุปสงฺกมิ, อุปสงฺกมิตฺวา อายสฺมนฺตํ ธนิยํ กุมฺภการ\-ปุตฺตํ อภิวาเทตฺวา เอกมนฺตํ นิสีทิ, เอกมนฺตํ นิสินฺโน โข ราชา มาคโธ เสนิโย พิมฺพิสาโร อายสฺมนฺตํ ธนิยํ \csromanfolio{54}กุมฺภการ\-ปุตฺตํ เอตทโวจ “สจฺจํ กิร มยา ภนฺเต เทวคหทารูนิ นคร\-ปฏิสงฺขาริกานิ อาปทตฺถาย นิกฺขิตฺตานิ อยฺยสฺส ทินฺนานี”ติ. เอวํ มหาราชาติ. มยํ โข ภนฺเต ราชาโน นาม พหุกิจฺจา พหุกรณียา ทตฺวาปิ น สเรยฺยาม, อิงฺฆ ภนฺเต สราเปหีติ. สรสิ ตฺวํ มหาราช ปฐมา\-ภิสิตฺโต เอวรูปิํ วาจํ ภาสิตา “ทินฺนญฺเญว สมณ\-พฺราหฺมณานํ ติณกฏฺโฐทกํ ปริภุญฺชนฺตู”ติ. สรามหํ ภนฺเต, สนฺติ ภนฺเต สมณพฺราหฺมณา ลชฺชิโน กุกฺกุจฺจกา สิกฺขากามา, เตสํ อปฺปมตฺตเกปิ กุกฺกุจฺจํ อุปฺปชฺชติ, เตสํ มยา สนฺธาย ภาสิตํ, ตญฺจ โข อรญฺเญ อปริคฺคหิตํ, โส ตฺวํ ภนฺเต เตน เลเสน ทารูนิ อทินฺนํ หริตุํ มญฺญสิ, \textbf{กถํ หิ นาม มาทิโส สมณํ วา พฺราหฺมณํ วา วิชิเต วสนฺตํ} \textbf{หเนยฺย วา พนฺเธยฺย วา ปพฺพาเชยฺย วา, คจฺฉ ภนฺเต โลเมน ตฺวํ} \textbf{มุตฺโตสิ, มาสฺสุ ปุนปิ เอวรูปํ อกาสีติ. มนุสฺสา อุชฺฌายนฺติ ขิยฺยนฺติ} \textbf{วิปาเจนฺติ “อลชฺชิโน อิเม สมณา สกฺยปุตฺติยา ทุสฺสีลา มุสาวาทิโน, อิเม หิ} \textbf{นาม ธมฺมจาริโน สมจาริโน พฺราหฺมจาริโน สจฺจวาทิโน สีลวนฺโต} กลฺยาณธมฺมา ปฏิชานิสฺสนฺติ. นตฺถิ อิเมสํ สามญฺญํ, นตฺถิ อิเมสํ พฺรหฺมญฺญํ, นฏฺฐํ อิเมสํ สามญฺญํ, นฏฺฐํ อิเมสํ พฺรหฺมญฺญํ, กุโต อิเมสํ สามญฺญํ, กุโต อิเมสํ พฺรหฺมญฺญํ, อปคตา อิเม สามญฺญา, อปคตา อิเม พฺรหฺมญฺญา, ราชานํปิ อิเม วญฺเจนฺติ, กิํ ปนญฺเญ มนุสฺเส”ติ. อสฺโสสุํ โข ภิกฺขู เตสํ มนุสฺสานํ อุชฺฌายนฺตานํ ขิยฺยนฺตานํ วิปาเจนฺตานํ. เย เต ภิกฺขู อปฺปิจฺฉา สนฺตุฏฺฐา ลชฺชิโน กุกฺกุจฺจกา สิกฺขากามา, เต อุชฺฌายนฺติ ขิยฺยนฺติ วิปาเจนฺติ “กถํ หิ นาม อายสฺมา ธนิโย กุมฺภการ\-ปุตฺโต รญฺโญ ทารูนิ อทินฺนํ อาทิยิสฺสตี”ติ. อถ โข เต ภิกฺขู อายสฺมนฺตํ ธนิยํ กุมฺภการ\-ปุตฺตํ อเนก\-ปริยาเยน วิครหิตฺวา ภควโต เอตมตฺถํ อาโรเจสุํ. อถ โข ภควา เอตสฺมิํ นิทาเน เอตสฺมิํ ปกรเณ ภิกฺขุสํฆํ สนฺนิปาตาเปตฺวา อายสฺมนฺตํ ธนิยํ กุมฺภการ\-ปุตฺตํ ปฏิปุจฺฉิ “สจฺจํ กิร ตฺวํ ธนิย รญฺโญ ทารูนิ อทินฺนํ อาทิยี”ติ. สจฺจํ ภควาติ. วิครหิ พุทฺโธ ภควา “อนนุจฺฉวิกํ โมฆปุริส อนนุโลมิกํ อปฺปติรูปํ อสฺสามณกํ อกปฺปิยํ อกรณียํ, กถํ หิ นาม ตฺวํ โมฆปุริส รญฺโญ ทารูนิ อทินฺนํ อาทิยิสฺสสิ, เนตํ โมฆปุริส อปฺปสนฺนานํ วา ปสาทาย ปสนฺนานํ วา ภิยฺโยภาวาย, อถเขฺวตํ\csromansharedfootnote{318eb6d047e97acd}{อถ เขฺวตํ เทฺว ปทานิ \csromandash{} ม.พ.ป.} โมฆปุริส อปฺปสนฺนาน\-ญฺเจว อปฺปสาทาย ปสนฺนานฺนญฺจ เอกจฺจานํ อญฺญถตฺตายา”ติ.}

\csromanmatikaanchor{mtk.17}
\csromantocmarkref{subsection}{ปฐมปญฺญตฺติ}{mtk.17}
\bookmark[dest={mtk.17},level=2]{ปฐมปญฺญตฺติ}
\prose{\csromanfolio{55}เตน โข ปน สมเยน อญฺญตโร ปุราณโวหาริโก มหามตฺโต ภิกฺขูสุ ปพฺพชิโต ภควโต อวิทูเร นิสินฺโน โหติ. อถ โข ภควา ตํ ภิกฺขุํ เอตทโวจ “กิตฺตเกน โข ภิกฺขุ ราชา มาคโธ เสนิโย พิมฺพิสาโร โจรํ คเหตฺวา หนติ วา พนฺธติ วา ปพฺพาเชติ วา”ติ. ปาเทน วา ภควา ปาทารเหน วาติ\csromansharedfootnote{68dc52f5150898f9}{ปาทารเหนวา อติเรกปาเทนวาติ (สฺยา)}. เตน โข ปน สมเยน ราชคเห ปญฺจมาสโก ปาโท โหติ. อถ โข ภควา อายสฺมนฺตํ ธนิยํ กุมฺภการ\-ปุตฺตํ อเนก\-ปริยาเยน วิครหิตฺวา ทุพฺภรตาย ฯเปฯ. เอวญฺจปน ภิกฺขเว อิมํ สิกฺขาปทํ อุทฺทิเสยฺยาถ\csromandash{}}

\proseitem{89}{\textbf{“โย ปน ภิกฺขุ อทินฺนํ เถยฺยสงฺขาตํ อาทิเยยฺย, ยถารูเป อทินฺนาทาเน ราชาโน โจรํ คเหตฺวา หเนยฺยุํ วา พนฺเธยฺยุํ วา ปพฺพาเชยฺยุํ วา โจโรสิ พาโลสิ มูโฬฺหสิ เถโนสีติ, ตถารูปํ ภิกฺขุ อทินฺนํ อาทิยมาโน อยมฺปิ ปาราชิโก โหติ อสํวาโส”ติ}.}

\prose{เอวญฺจิทํ ภควตา ภิกฺขูนํ สิกฺขาปทํ ปญฺญตฺตํ โหติ.}

\csromanmatikaanchor{mtk.18}
\csromanmatikaanchor{mtk.19}
\csromantocmarkref{subsection}{อนุปญฺญตฺติ}{mtk.18}
\bookmark[dest={mtk.18},level=2]{อนุปญฺญตฺติ}
\csromantocmarkref{subsection}{สิกฺขาปทวิภงฺค, ปทภาชนีย}{mtk.19}
\bookmark[dest={mtk.19},level=2]{สิกฺขาปทวิภงฺค, ปทภาชนีย}
\proseitem{90}{เตน โข ปน สมเยน ฉพฺพคฺคิยา ภิกฺขู รชก\-ตฺถรณํ คนฺตฺวา รชก\-ภณฺฑิกํ อวหริตฺวา อารามํ หริตฺวา ภาเชสุํ. ภิกฺขู เอวมาหํสุ “มหาปุญฺญตฺถ ตุเมฺห อาวุโส, พหุํ ตุมฺหากํ จีวรํ อุปฺปนฺนนฺ”ติ. กุโต อาวุโส อมฺหากํ ปุญฺญํ, อิทานิ มยํ รชก\-ตฺถรณํ คนฺตฺวา รชก\-ภณฺฑิกํ อวหริมฺหาติ. นนุ อาวุโส ภควตา สิกฺขาปทํ ปญฺญตฺตํ, กิสฺส ตุเมฺห อาวุโส รชก\-ภณฺฑิกํ อวหริตฺถาติ. สจฺจํ อาวุโส ภควตา สิกฺขาปทํ ปญฺญตฺตํ, ตญฺจ โข คาเม, โน อรญฺเญติ. นนุ อาวุโส ตเถเวตํ โหติ, อนนุจฺฉวิกํ อาวุโส อนนุโลมิกํ อปฺปติรูปํ อสฺสามณกํ อกปฺปิยํ อกรณียํ, กถํ หิ นาม ตุเมฺห อาวุโส รชก\-ภณฺฑิกํ อวหริสฺสถ, เนตํ อาวุโส อปฺปสนฺนานํ วา ปสาทาย ปสนฺนานํ วา ภิยฺโยภาวาย. อถเขฺวตํ\csromansharedfootnote{318eb6d047e97acd}{อถ เขฺวตํ เทฺว ปทานิ \csromandash{} ม.พ.ป.} อาวุโส อปฺปสนฺนาน\-ญฺเจว อปฺปสาทาย ปสนฺนานญฺจ เอกจฺจานํ อญฺญถตฺตายาติ. อถ โข เต ภิกฺขู ฉพฺพคฺคิเย ภิกฺขู อเนก\-ปริยาเยน วิครหิตฺวา ภควโต เอตมตฺถํ อาโรเจสุํ. \csromanfolio{56}อถ โข ภควา เอตสฺมิํ นิทาเน เอตสฺมิํ ปกรเณ ภิกฺขุสํฆํ สนฺนิปาตาเปตฺวา ฉพฺพคฺคิเย ภิกฺขู ปฏิปุจฺฉิ “สจฺจํ กิร ตุเมฺห ภิกฺขเว รชก\-ตฺถรณํ คนฺตฺวา รชก\-ภณฺฑิกํ อวหริตฺถา”ติ. สจฺจํ ภควาติ. วิครหิ พุทฺโธ ภควา “อนนุจฺฉวิกํ โมฆปุริสา อนนุโลมิกํ อปฺปติรูปํ อสฺสามณกํ อกปฺปิยํ อกรณียํ, กถํ หิ นาม ตุเมฺห โมฆปุริสา รชก\-ภณฺฑิกํ อวหริสฺสถ, เนตํ โมฆปุริสา อปฺปสนฺนานํ วา ปสาทาย ปสนฺนานํ วา ภิยฺโยภาวาย, อถ เขฺวตํ โมฆปุริสา อปฺปสนฺนาน\-ญฺเจว อปฺปสาทาย ปสนฺนานญฺจ เอกจฺจานํ อญฺญถตฺตายา”ติ. อถ โข ภควา ฉพฺพคฺคิเย ภิกฺขู อเนก\-ปริยาเยน วิครหิตฺวา ทุพฺภรตาย ฯเปฯ วีริยารพฺภสฺส วณฺณํ ภาสิตฺวา ภิกฺขูนํ ตทนุจฺฉวิกํ ตทนุโลมิกํ ธมฺมิํ กถํ กตฺวา ภิกฺขู อามนฺเตสิ ฯเปฯ. เอวญฺจ ปน ภิกฺขเว อิมํ สิกฺขาปทํ อุทฺทิเสยฺยาถ\csromandash{}}

\proseitem{91}{\textbf{“โย ปน ภิกฺขุ คามา วา อรญฺญา วา อทินฺนํ เถยฺยสงฺขาตํ อาทิเยยฺย, ยถารูเป อทินฺนาทาเน ราชาโน โจรํ คเหตฺวา หเนยฺยุํ วา พนฺเธยฺยุํ วา ปพฺพาเชยฺยุํ วา โจโรสิ พาโลสิ มูโฬฺหสิ เถโนสีติ, ตถารูปํ ภิกฺขุ อทินฺนํ อาทิยมาโน อยมฺปิ ปาราชิโก โหติ อสํวาโส”}ติ.}

\proseitem{92}{\textbf{โยปนา}ติ\csromansharedfootnote{423cb82fe87b6ac9}{โย ปนาติ เทฺว ปทานิ \csromandash{} ม.พ.ป.} โย ยาทิโส ฯเปฯ. \textbf{ภิกฺขู}ติ ฯเปฯ อยํ อิมสฺมิํ อตฺเถ อธิปฺเปโต ภิกฺขูติ.}

\prose{\textbf{คาโม} นาม เอกกุฏิโกปิ คาโม, ทฺวิกุฏิโกปิ คาโม, ติกุฏิโกปิ คาโม, จตุกุฏิโกปิ คาโม, สมนุสฺโสปิ คาโม, อมนุสฺโสปิ คาโม, ปริกฺขิตฺโตปิ คาโม, อปริกฺขิตฺโตปิ คาโม, โคนิสาทิ\-นิวิฏฺโฐปิ คาโม, โยปิ สตฺโถ อติเรก\-จตุมาส\-นิวิฏฺโฐ, โสปิ วุจฺจติ คาโม.}

\prose{\textbf{คามูปจาโร} นาม ปริกฺขิตฺตสฺส คามสฺส อินฺทขีเล\csromansharedfootnote{9f3e0858805c2f0d}{อินฺทขิเล (ก)} ฐิตสฺส มชฺฌิมสฺส ปุริสสฺส เลฑฺฑุปาโต, อปริกฺขิตฺตสฺส คามสฺส ฆรูปจาเร ฐิตสฺส มชฺฌิมสฺส ปุริสสฺส เลฑฺฒุปาโต.}

\prose{\csromanfolio{57}\textbf{อรญฺญํ} นาม ฐเปตฺวา คามญฺจ คามูปจารญฺจ อวเสสํ อรญฺญํ นาม.}

\prose{\textbf{อทินฺนํ} นาม ยํ อทินฺนํ อนิสฺสฏฺฐํ อปริจฺจตฺตํ รกฺขิตํ โคปิตํ มมายิตํ ปรปริคฺคหิตํ, เอตํ อทินฺนํ นาม.}

\prose{\textbf{เถยฺย}\-\textbf{สงฺขาต}นฺติ เถยฺยจิตฺโต อวหรณจิตฺโต.}

\prose{\textbf{อาทิเยยฺยา}ติ อาทิเยยฺย, หเรยฺย, อวหเรยฺย, อิริยาปถํ วิโกเปยฺย, ฐานา จาเวยฺย, สงฺเกตํ วีตินาเมยฺย.}

\prose{\textbf{ยถารูปํ} นาม ปาทํ วา ปาทารหํ วา อติเรกปาทํ วา.}

\prose{\textbf{ราชาโน} นาม ปถพฺยาราชา ปเทสราชา มณฺฑลิกา อนฺตรโภคิกา อกฺขทสฺสา มหามตฺตา, เย วา ปน เฉชฺชเภชฺชํ กโรนฺตา อนุสาสนฺติ, เอเต ราชาโน นาม.}

\prose{\textbf{โจโร} นาม โย ปญฺจมาสกํ วา อติเรก\-ปญฺจ\-มาสกํ วา อคฺฆนกํ อทินฺนํ เถยฺยสงฺขาตํ อาทิยติ, เอโส โจโร นาม.}

\prose{\textbf{หเนยฺยุํ วา}ติ หตฺเถน วา ปาเทน วา กสาย วา เวตฺเตน วา อฑฺฒทณฺฑเกน วา เฉชฺชาย วา หเนยฺยุํ.}

\prose{\textbf{พนฺเธยฺยุํ วา}ติ รชฺชุ\-พนฺธเนน วา อนฺทุพนฺธเนน วา สงฺขลิก\-พนฺธเนน วา ฆรพนฺธเนน วา นคร\-พนฺธเนน วา คามพนฺธเนน วา นิคม\-พนฺธเนน วา พนฺเธยฺยุํ, ปุริสคุตฺติํ วา กเรยฺยุํ.}

\prose{\textbf{ปพฺพาเชยฺยุํ วา}ติ คามา วา นิคมา วา นครา วา ชนปทา วา ชนปทปเทสา วา ปพฺพาเชยฺยุํ.}

\prose{\textbf{โจโรสิ พาโลสิ มูโฬฺหสิ เถโนสี}ติ ปริภาโส เอโส.}

\prose{\textbf{ตถารูปํ} นาม ปาทํ วา ปาทารหํ วา อติเรกปาทํ วา.}

\prose{\textbf{อาทิยมาโน}ติ อาทิยมาโน, หรมาโน, อวหรมาโน, อิริยาปถํ วิโกปยมาโน, ฐานา จาวยมาโน, สงฺเกตํ วีตินามยมาโน.}

\prose{\csromanfolio{58}\textbf{อยมฺปี}ติ ปุริมํ อุปาทาย วุจฺจติ.}

\prose{\textbf{ปาราชิโก โหตี}ติ เสยฺยถาปิ นาม ปณฺฑุปลาโส พนฺธนา ปวุตฺโต\csromansharedfootnote{b7c9a2e5614d85d7}{หริตตฺตาย (สี, สฺยา)} อภพฺโพ หริตตฺถาย\csromansharedfootnote{b7c9a2e5614d85d7}{หริตตฺตาย (สี, สฺยา)}, เอวเมว ภิกฺขุ ปาทํ วา ปาทารหํ วา อติเรกปาทํ วา อทินฺนํ เถยฺยสงฺขาตํ อาทิยิตฺวา อสฺสมโณ โหติ อสกฺยปุตฺติโย, เตน วุจฺจติ “ปาราชิโก โหตี”ติ.}

\prose{\textbf{อสํวาโส}ติ สํวาโส นาม เอกกมฺมํ เอกุทฺเทโส สมสิกฺขตา, เอโส สํวาโส นามา. โส เตน สทฺธิํ นตฺถิ, เตน วุจฺจติ “อสํวาโส”ติ.}

\proseitem{93}{ภูมฏฺฐํ ถลฏฺฐํ อากาสฏฺฐํ เวหาสฏฺฐํ อุทกฏฺฐํ นาวฏฺฐํ ยานฏฺฐํ ภารฏฺฐํ อารามฏฺฐํ วิหารฏฺฐํ เขตฺตฏฺฐํ วตฺถุฏฺฐํ คามฏฺฐํ อรญฺญฏฺฐํ อุทกํ ทนฺตโปณํ\csromansharedfootnote{feb297c07d149a36}{ทนฺตโปนํ (สี, ก)} วนปฺปติ หรณกํ อุปนิธิ สุงฺกฆาตํ ปาโณ อปทํ ทฺวิปทํ จตุปฺปทํ พหุปฺปทํ โอจรโก โอณิรกฺโข สํวิทาวหาโร สงฺเกตกมฺมํ นิมิตกมฺมนฺติ.}

\proseitem{94}{\textbf{ภูมฏฺฐํ} นาม ภณฺฑํ ภูมิยํ นิกฺขิตฺตํ โหติ นิขาตํ ปฏิจฺฉนฺนํ. “ภูมฏฺฐํ ภณฺฑํ อวหริสฺสามี”ติ เถยฺยจิตฺโต ทุติยํ วา ปริเยสติ กุทาลํ วา ปิฏกํ วา ปริเยสติ คจฺฉติ วา, อาปตฺติ ทุกฺกฏสฺส. ตตฺถชาตกํ กฏฺฐํ วา ลตํ วา ฉินฺทติ, อาปตฺติ ทุกฺกฏสฺส. ตตฺถ ปํสุํ ขณติ วา พฺยูหติ\csromansharedfootnote{1ff636b3db6b49b7}{วิยูหติ (สฺยา)} วา อุทฺธรติ วา, อาปตฺติ ทุกฺกฏสฺส. กุมฺภิํ อามสติ, อาปตฺติ ทุกฺกฏสฺส. ผนฺทาเปติ, อาปตฺติ ถุลฺลจฺจยสฺส. ฐานา จาเวติ, อาปตฺติ ปาราชิกสฺส. อตฺตโน ภาชนํ ปเวเสตฺวา ปญฺจมาสกํ วา อติเรก\-ปญฺจ\-มาสกํ วา อคฺฆนกํ เถยฺยจิตฺโต อามสติ, อาปตฺติ ทุกฺกฏสฺส. ผนฺทาเปติ, อาปตฺติ ถุลฺลจฺจยสฺส. อตฺตโน ภาชนคตํ วา กโรติ, มุฏฺฐิํ วา ฉินฺทติ, อาปตฺติ ปาราชิกสฺส. สุตฺตารุฬฺหํ ภณฺฑํ ปามงฺคํ วา กณฺฐสุตฺตกํ วา กฏิสุตฺตกํ วา สาฏกํ วา เวฐนํ วา เถยฺยจิตฺโต อามสติ, อาปตฺติ ทุกฺกฏสฺส. ผนฺทาเปติ, อาปตฺติ ถุลฺลจฺจยสฺส. โกฏิยํ คเหตฺวา อุจฺจาเรติ, อาปตฺติ ถุลฺลจฺจยสฺส. ฆํสนฺโต นีหรติ, \csromanfolio{59}อาปตฺติ ถุลฺลจฺจยสฺส. อนฺตมโส เกสคฺค\-มตฺตํปิ กุมฺภิมุขา โมเจติ, อาปตฺติ ปาราชิกสฺส. สปฺปิํ วา เตลํ วา เตลํ วา มธุํ วา ผาณิตํ วา ปญฺจมาสกํ วา อติเรก\-ปญฺจ\-มาสกํ วา อคฺฆนกํ เถยฺยจิตฺโต เอเกน ปโยเคน ปิวติ, อาปตฺติ ปาราชิกสฺส. ตตฺเถว ภินฺทติ ฉฑฺเฑติ วา ฌาเปติ วา อปริโภคํ วา กโรติ, อาปตฺติ ทุกฺกฏสฺส.}

\proseitem{95}{\textbf{ถลฏฺฐํ} นาม ภณฺฑํ ถเล นิกฺขิตฺตํ โหติ. “ถลฏฺฐํ ภณฺฑํ อวหริสฺสามี”ติ เถยฺยจิตฺโต ทุติยํ วา ปริเยสติ คจฺฉติ วา, อาปตฺติ ทุกฺกฏสฺส. อามสติ, อาปตฺติ ทุกฺกฏสฺส. ผนฺทาเปติ, อาปตฺติ ถุลฺลจฺจยสฺส. ฐานา จาเวติ, อาปตฺติ ปาราชิกสฺส.}

\proseitem{96}{\textbf{อากาสฏฺฐํ} นาม ภณฺฑํ อากาสคตํ โหติ, โมโร วา กปิญฺชโร วา ติตฺติโร วา วฏฺฏโก วา สาฏกํ วา เวฐนํ วา หิรญฺญํ วา สุวณฺณํ วา ฉิชฺชมานํ ปตติ. “อากาสฏฺฐํ ภณฺฑํ อวหริสฺสามี”ติ เถยฺยจิตฺโต ทุติยํ วา ปริเยสติ คจฺฉติ วา, อาปตฺติ ทุกฺกฏสฺส. คมนํ อุปจฺฉินฺทติ, อาปตฺติ ทุกฺกฏสฺส. อามสติ, อาปตฺติ ทุกฺกฏสฺส. ผนฺทาเปติ, อาปตฺติ ถุลฺลจฺจยสฺส. ฐานา จาเวติ, อาปตฺติ ปาราชิกสฺส.}

\proseitem{97}{\textbf{เวหาสฏฺฐํ} นาม ภณฺฑํ เวหาสคตํ โหติ, มญฺเจ วา ปีเฐ วา จีวรวํเส วา จีวรรชฺชุยา วา ภิตฺติขิเล วา นาคทนฺเต วา รุกฺเข วา ลคฺคิตํ โหติ อนฺตมโส ปตฺตาธารเกปิ. “เวหาสฏฺฐํ ภณฺฑํ อวหริสฺสามี”ติ เถยฺยจิตฺโต ทุติยํ วา ปริเยสติ คจฺฉติ วา, อาปตฺติ ทุกฺกฏสฺส. อามสติ, อาปตฺติ ทุกฺกฏสฺส. ผนฺทาเปติ, อาปตฺติ ถุลฺลจฺจยสฺส. ฐานา จาเวติ, อาปตฺติ ปาราชิกสฺส.}

\proseitem{98}{\textbf{อุทกฏฺฐํ} นาม ภณฺฑํ อุทเก นิกฺขิตฺตํ โหติ. “อุทกฏฺฐํ ภณฺฑํ อวหริสฺสามี”ติ เถยฺยจิตฺโต ทุติยํ วา ปริเยสติ คจฺฉติ วา, อาปตฺติ ทุกฺกฏสฺส. นิมุชฺชติ วา อุมฺมุชฺชติ วา, อาปตฺติ ทุกฺกฏสฺส. อามสติ, อาปตฺติ ทุกฺกฏสฺส. ผนฺทาเปติ, อาปตฺติ ถุลฺลจฺจยสฺส. ฐานา จาเวติ, อาปตฺติ ปาราชิกสฺส. ตตฺถ ชาตกํ อุปฺปลํ วา ปทุมํ วา ปุณฺฑรีกํ วา ภิสํ วา มจฺฉํ วา กจฺฉปํ วา ปญฺจมาสกํ วา อติเรก\-ปญฺจ\-มาสกํ วา อคฺฆนกํ เถยฺยจิตฺโต อามสติ, อาปตฺติ ทุกฺกฏสฺส. ผนฺทาเปติ, อาปตฺติ ถุลฺลจฺจยสฺส. ฐานา จาเวติ, อาปตฺติ ปาราชิกสฺส.}

\proseitem{99}{\csromanfolio{60}\textbf{นาวา} นาม ยาย ตรติ. \textbf{นาวฏฺฐํ} นาม ภณฺฑํ นาวาย นิกฺขิตฺตํ โหติ. “นาวฏฺฐํ ภณฺฑํ อวหริสฺสามี”ติ เถยฺยจิตฺโต ทุติยํ วา ปริเยสติ คจฺฉติ วา, อาปตฺติ ทุกฺกฏสฺส. อามสติ, อาปตฺติ ทุกฺกฏสฺส. ผนฺทาเปติ, อาปตฺติ ถุลฺลจฺจยสฺส. ฐานา จาเวติ, อาปตฺติ ปาราชิกสฺส. “นาวํ อวหริสฺสามี”ติ เถยฺยจิตฺโต ทุติยํ วา ปริเยสติ คจฺฉติ วา, อาปตฺติ ทุกฺกฏสฺส. อามสติ, อาปตฺติ ทุกฺกฏสฺส. ผนฺทาเปติ, อาปตฺติ ถุลฺลจฺจยสฺส. พนฺธนํ โมเจติ, อาปตฺติ ทุกฺกฏสฺส. พนฺธนํ โมเจตฺวา อามสติ, อาปตฺติ ทุกฺกฏสฺส. ผนฺทาเปติ, อาปตฺติ ถุลฺลจฺจยสฺส. อุทฺธํ วา อโธ วา ติริยํ วา อนฺตมโส เกสคฺค\-มตฺตํปิ สงฺกาเมติ, อาปตฺติ ปาราชิกสฺส.}

\proseitem{100}{\textbf{ยานํ} นาม วยฺหํ รโถ สกฏํ สนฺทมานิกา. \textbf{ยานฏฺฐํ} นาม ภณฺฑํ ยาเน นิกฺขิตฺตํ โหติ. “ยานฏฺฐํ ภณฺฑํ อวหริสฺสามี”ติ เถยฺยจิตฺโต ทุติยํ วา ปริเยสติ คจฺฉติ วา, อาปตฺติ ทุกฺกฏสฺส. อามสติ, อาปตฺติ ทุกฺกฏสฺส. ผนฺทาเปติ, อาปตฺติ ถุลฺลจฺจยสฺส. ฐานา จาเวติ, อาปตฺติ ปาราชิกสฺส. “ยานํ อวหริสฺสามี”ติ เถยฺยจิตฺโต ทุติยํ วา ปริเยสติ คจฺฉติ วา, อาปตฺติ ทุกฺกฏสฺส. อามสติ, อาปตฺติ ทุกฺกฏสฺส. ผนฺทาเปติ, อาปตฺติ ถุลฺลจฺจยสฺส. ฐานา จาเวติ, อาปตฺติ ปาราชิกสฺส.}

\proseitem{101}{\textbf{ภาโร} นาม สีสภาโร ขนฺธภาโร กฏิภาโร โอลมฺพโก. สีเส ภารํ เถยฺยจิตฺโต อามสติ, อาปตฺติ ทุกฺกฏสฺส. ผนฺทาเปติ, อาปตฺติ ถุลฺลจฺจยสฺส. ขนฺธํ โอโรเปติ, อาปตฺติ ปาราชิกสฺส. ขนฺเธ ภารํ เถยฺยจิตฺโต อามสติ, อาปตฺติ ทุกฺกฏสฺส. ผนฺทาเปติ, อาปตฺติ ถุลฺลจฺจยสฺส. กฏิํ โอโรเปติ, อาปตฺติ ปาราชิกสฺส. กฏิยา ภารํ เถยฺยจิตฺโต อามสติ, อาปตฺติ ทุกฺกฏสฺส. ผนฺทาเปติ, อาปตฺติ ถุลฺลจฺจยสฺส. หตฺเถน คณฺหาติ, อาปตฺติ ปาราชิกสฺส. หตฺเถ ภารํ เถยฺยจิตฺโต ภูมิยํ นิกฺขิปติ, อาปตฺติ ปาราชิกสฺส. เถยฺยจิตฺโต ภูมิโต คณฺหาติ, อาปตฺติ ปาราชิกสฺส.}

\proseitem{102}{\textbf{อาราโม} นาม ปุปฺผาราโม ผลาราโม. \textbf{อารามฏฺฐํ} นาม ภณฺฑํ อาราเม จตูหิ ฐาเนหิ นิกฺขิตฺตํ โหติ ภูมฏฺฐํ ถลฏฺฐํ อากาสฏฺฐํ เวหาสฏฺฐํ. “อารามฏฺฐํ ภณฺฑํ อวหริสฺสามี”ติ เถยฺยจิตฺโต ทุติยํ วา \csromanfolio{61}ปริเยสติ คจฺฉติ วา, อาปตฺติ ทุกฺกฏสฺส. อามสติ, อาปตฺติ ทุกฺกฏสฺส. ผนฺทาเปติ, อาปตฺติ ถุลฺลจฺจยสฺส. ฐานา จาเวติ, อาปตฺติ ปาราชิกสฺส. ตตฺถ ชาตกํ มูลํ วา ตจํ วา ปตฺตํ วา ปุปฺผํ วา ผลํ วา ปญฺจมาสกํ วา อติเรก\-ปญฺจ\-มาสกํ วา อคฺฆนกํ เถยฺยจิตฺโต อามสติ, อาปตฺติ ทุกฺกฏสฺส. ผนฺทาเปติ, อาปตฺติ ถุลฺลจฺจยสฺส. ฐานา จาเวติ, อาปตฺติ ปาราชิกสฺส. อารามํ อภิยุญฺชติ, อาปตฺติ ทุกฺกฏสฺส. สามิกสฺส วิมติํ อุปฺปาเทติ, อาปตฺติ ถุลฺลจฺจยสฺส. สามิโก “น มยฺหํ ภวิสฺสตี”ติ ธุรํ นิกฺขิปติ, อาปตฺติ ปาราชิกสฺส. ธมฺมํ จรนฺโต สามิกํ ปราเชติ, อาปตฺติ ปาราชิกสฺส. ธมฺมํ จรนฺโต ปรชฺชติ, อาปตฺติ ถุลฺลจฺจยสฺส.}

\proseitem{103}{\textbf{วิหารฏฺฐํ} นาม ภณฺฑํ วิหาเร จตูหิ ฐาเนหิ นิกฺขิตฺตํ โหติ ภูมฏฺฐํ ถลฏฺฐํ อากาสฏฺฐํ เวหาสฏฺฐํ. “วิหารฏฺฐํ ภณฺฑํ อวหริสฺสามี”ติ เถยฺยจิตฺโต ทุติยํ วา ปริเยสติ คจฺฉติ วา, อาปตฺติ ทุกฺกฏสฺส. อามสติ, อาปตฺติ ทุกฺกฏสฺส. ผนฺทาเปติ, อาปตฺติ ถุลฺลจฺจยสฺส. ฐานา จาเวติ, อาปตฺติ ปาราชิกสฺส. วิหารํ อภิยุญฺชติ, อาปตฺติ ทุกฺกฏสฺส. สามิกสฺส วิมติํ อุปฺปาเทติ, อาปตฺติ ถุลฺลจฺจยสฺส. สามิโก “น มยฺหํ ภวิสฺสตี”ติ ธุรํ นิกฺขิปติ, อาปตฺติ ปาราชิกสฺส. ธมฺมํ จรนฺโต สามิกํ ปราเชติ, อาปตฺติ ปาราชิกสฺส. ธมฺมํ จรนฺโต ปรชฺชติ, อาปตฺติ ถุลฺลจฺจยสฺส.}

\proseitem{104}{\textbf{เขตฺตํ} นาม ยตฺถ ปุพฺพณฺณํ วา อปรณฺณํ วา ชายติ. เขตฺตฏฺฐํ นาม ภณฺฑํ เขตฺเต จตูหิ ฐาเนหิ นิกฺขิตฺตํ โหติ ภูมฏฺฐํ ถลฏฺฐํ อากาสฏฺฐํ เวหาสฏฺฐํ. “\textbf{เขตฺตฏฺฐํ} ภณฺฑํ อวหริสฺสามี”ติ เถยฺยจิตฺโต ทุติยํ วา ปริเยสติ คจฺฉติ วา, อาปตฺติ ทุกฺกฏสฺส. อามสติ, อาปตฺติ ทุกฺกฏสฺส. ผนฺทาเปติ, อาปตฺติ ถุลฺลจฺจยสฺส. ฐานา จาเวติ, อาปตฺติ ปาราชิกสฺส. ตตฺถ ชาตกํ ปุพฺพณฺณํ วา อปรณฺณํ วา ปญฺจมาสกํ วา อติเรก\-ปญฺจ\-มาสกํ วา อคฺฆนกํ เถยฺยจิตฺโต อามสติ, อาปตฺติ ทุกฺกฏสฺส. ผนฺทาเปติ, อาปตฺติ ถุลฺลจฺจยสฺส. ฐานา จาเวติ, อาปตฺติ ปาราชิกสฺส. เขตฺตํ อภิยุญฺชติ, อาปตฺติ ทุกฺกฏสฺส. สามิกสฺส วิมติํ อุปฺปาเทติ, อาปตฺติ ถุลฺลจฺจยสฺส. สามิโก “น มยฺหํ ภวิสฺสตี”ติ ธุรํ นิกฺขิปติ, อาปตฺติ ปาราชิกสฺส. ธมฺมํ จรนฺโต สามิกํ ปราเชติ, อาปตฺติ ปาราชิกสฺส. ธมฺมํ จรนฺโต ปรชฺชติ, อาปตฺติ ถุลฺลจฺจยสฺส. ขีลํ วา \csromanfolio{62}รชฺชุํ วา วติํ วา มริยาทํ วา สงฺกาเมติ, อาปตฺติ ทุกฺกฏสฺส. เอกํ ปโยคํ อนาคเต อาปตฺติ ถุลฺลจฺจยสฺส. ตสฺมิํ ปโยเค อาคเต อาปตฺติ ปาราชิกสฺส.}

\proseitem{105}{\textbf{วตฺถุ} นาม อารามวตฺถุ วิหารวตฺถุ. \textbf{วตฺถุฏฺฐํ} นาม ภณฺฑํ วตฺถุสฺมิํ จตูหิ ฐาเนหิ นิกฺขิตฺตํ โหติ ภูมฏฺฐํ ถลฏฺฐํ อากาสฏฺฐํ เวหาสฏฺฐํ. “วตฺถุฏฺฐํ ภณฺฑํ อวหริสฺสามี”ติ เถยฺยจิตฺโต ทุติยํ วา ปริเยสติ คจฺฉติ วา, อาปตฺติ ทุกฺกฏสฺส. อามสติ, อาปตฺติ ทุกฺกฏสฺส. ผนฺทาเปติ, อาปตฺติ ถุลฺลจฺจยสฺส. ฐานา จาเวติ, อาปตฺติ ปาราชิกสฺส. วตฺถุํ อภิยุญฺชติ, อาปตฺติ ทุกฺกฏสฺส. สามิกสฺส วิมติํ อุปฺปาเทติ, อาปตฺติ ถุลฺลจฺจยสฺส. สามิโก “น มยฺหํ ภวิสฺสตี”ติ ธุรํ นิกฺขิปติ, อาปตฺติ ปาราชิกสฺส. ธมฺมํ จรนฺโต สามิกํ ปราเชติ, อาปตฺติ ปาราชิกสฺส. ธมฺมํ จรนฺโต ปรชฺชติ, อาปตฺติ ถุลฺลจฺจยสฺส. ขีลํ วา รชฺชุํ วา วติํ วา ปาการํ วา สงฺกาเมติ, อาปตฺติ ทุกฺกฏสฺส. เอกํ ปโยคํ อนาคเต อาปตฺติ ถุลฺลจฺจยสฺส. ตสฺมิํ ปโยเค อาคเต อาปตฺติ ปาราชิกสฺส.}

\proseitem{106}{\textbf{คามฏฺฐํ} นาม ภณฺฑํ คาเม จตูหิ ฐาเนหิ นิกฺขิตฺตํ โหติ ภูมฏฺฐํ ถลฏฺฐํ อากาสฏฺฐํ เวหาสฏฺฐํ. “คามฏฺฐํ ภณฺฑํ อวหริสฺสามี”ติ เถยฺยจิตฺโต ทุติยํ วา ปริเยสติ คจฺฉติ วา, อาปตฺติ ทุกฺกฏสฺส. อามสติ, อาปตฺติ ทุกฺกฏสฺส. ผนฺทาเปติ, อาปตฺติ ถุลฺลจฺจยสฺส. ฐานา จาเวติ, อาปตฺติ ปาราชิกสฺส.}

\proseitem{107}{\textbf{อรญฺญํ} นาม ยํ มนุสฺสานํ ปริคฺคหิตํ โหติ, ตํ อรญฺญํ. อรญฺญฏฺฐํ นาม ภณฺฑํ อรญฺเญ จตูหิ ฐาเนหิ นิกฺขิตฺตํ โหติ ภูมฏฺฐํ ถลฏฺฐํ อากาสฏฺฐํ เวหาสฏฺฐํ. \textbf{อรญฺญฏฺฐํ} ภณฺฑํ อวหริสฺสามีติ เถยฺยจิตฺโต ทุติยํ วา ปริเยสติ คจฺฉติ วา, อาปตฺติ ทุกฺกฏสฺส. อามสติ, อาปตฺติ ทุกฺกฏสฺส. ผนฺทาเปติ, อาปตฺติ ถุลฺลจฺจยสฺส. ฐานา จาเวติ, อาปตฺติ ปาราชิกสฺส. ตตฺถชาตกํ กฏฺฐํ วา ลตํ วา ติณํ วา ปญฺจมาสกํ วา อติเรก\-ปญฺจ\-มาสกํ วา อคฺฆนกํ เถยฺยจิตฺโต อามสติ, อาปตฺติ ทุกฺกฏสฺส. ผนฺทาเปติ, อาปตฺติ ถุลฺลจฺจยสฺส. ฐานา จาเวติ, อาปตฺติ ปาราชิกสฺส.}

\proseitem{108}{\csromanfolio{63}\textbf{อุทกํ} นาม ภาชนคตํ วา โหติ โปกฺขรณิยา วา ตฬาเก วา. เถยฺยจิตฺโต อามสติ, อาปตฺติ ทุกฺกฏสฺส. ผนฺทาเปติ, อาปตฺติ ถุลฺลจฺจยสฺส. ฐานา จาเวติ, อาปตฺติ ปาราชิกสฺส. อตฺตโน ภาชนํ ปเวเสตฺวา ปญฺจมาสกํ วา อติเรก\-ปญฺจ\-มาสกํ วา อคฺฆนกํ อุทกํ เถยฺยจิตฺโต อามสติ, อาปตฺติ ทุกฺกฏสฺส. ผนฺทาเปติ, อาปตฺติ ถุลฺลจฺจยสฺส. อตฺตโน ภาชนคตํ กโรติ, อาปตฺติ ปาราชิกสฺส. มริยาทํ ภินฺทติ, อาปตฺติ ทุกฺกฏสฺส. มริยาทํ ภินฺทิตฺวา ปญฺจมาสกํ วา อติเรก\-ปญฺจ\-มาสกํ วา อคฺฆนกํ อุทกํ นิกฺขาเมติ, อาปตฺติ ปาราชิกสฺส. อติเรกมาสกํ วา อูนปญฺจมาสกํ วา อคฺฆนกํ อุทกํ นิกฺขาเมติ, อาปตฺติ ถุลฺลจฺจยสฺส, มาสกํ วา อูนมาสกํ วา อคฺฆนกํ อุทกํ นิกฺขาเมติ, อาปตฺติ ทุกฺกฏสฺส.}

\proseitem{109}{\textbf{ทนฺตโปณํ} นาม ฉินฺนํ วา อจฺฉินฺนํ วา. ปญฺจมาสกํ วา อติเรก\-ปญฺจ\-มาสกํ วา อคฺฆนกํ เถยฺยจิตฺโต อามสติ, อาปตฺติ ทุกฺกฏสฺส. ผนฺทาเปติ, อาปตฺติ ถุลฺลจฺจยสฺส. ฐานา จาเวติ, อาปตฺติ ปาราชิกสฺส.}

\proseitem{110}{\textbf{วนปฺปติ} นาม โย มนุสฺสานํ ปริคฺคหิโต โหติ รุกฺโข ปริโภโค. เถยฺยจิตฺโต ฉินฺทติ, ปหาเร ปหาเร อาปตฺติ ทุกฺกฏสฺส. เอกํ ปหารํ อนาคเต อาปตฺติ ถุลฺลจฺจยสฺส. ตสฺมิํ ปหาเร อาคเต อาปตฺติ ปาราชิกสฺส.}

\proseitem{111}{\textbf{หรณกํ} นาม อญฺญสฺส หรณกํ ภณฺฑํ. เถยฺยจิตฺโต อามสติ, อาปตฺติ ทุกฺกฏสฺส. ผนฺทาเปติ, อาปตฺติ ถุลฺลจฺจยสฺส. ฐานา จาเวติ, อาปตฺติ ปาราชิกสฺส. “สหภณฺฑหารกํ ปทสา เนสฺสามี”ติ ปฐมํ ปาทํ สงฺกาเมติ, อาปตฺติ ถุลฺลจฺจยสฺส. ทุติยํ ปาทํ สงฺกาเมติ, อาปตฺติ ปาราชิกสฺส. “ปติตํ ภณฺฑํ คเหสฺสามี”ติ ปาตาเปติ, อาปตฺติ ทุกฺกฏสฺส. ปติตํ ภณฺฑํ ปญฺจมาสกํ วา อติเรก\-ปญฺจ\-มาสกํ วา อคฺฆนกํ เถยฺยจิตฺโต อามสติ, อาปตฺติ ทุกฺกฏสฺส. ผนฺทาเปติ, อาปตฺติ ถุลฺลจฺจยสฺส. ฐานา จาเวติ, อาปตฺติ ปาราชิกสฺส.}

\proseitem{112}{\textbf{อุปนิธิ} นาม อุปนิกฺขิตฺตํ ภณฺฑํ. “เทหิ เม ภณฺฑนฺ”ติ วุจฺจมาโน “นาหํ คณฺหามี”ติ ภณติ, อาปตฺติ ทุกฺกฏสฺส. สามิกสฺส วิมติํ \csromanfolio{64}อุปฺปาเทติ, อาปตฺติ ถุลฺลจฺจยสฺส. สามิโก “น มยฺหํ ทสฺสตี”ติ ธุรํ นิกฺขิปติ, อาปตฺติ ปาราชิกสฺส. ธมฺมํ จรนฺโต สามิกํ ปราเชติ, อาปตฺติ ปาราชิกสฺส. ธมฺมํ จรนฺโต ปรชฺชติ, อาปตฺติ ถุลฺลจฺจยสฺส.}

\proseitem{113}{\textbf{สุงฺกฆาตํ} นาม รญฺญา ฐปิตํ โหติ ปพฺพตขณฺเฑ วา นทีติตฺเถ วา คามทฺวาเร วา “อตฺร ปวิฏฺฐสฺส สุงฺกํ คณฺหนฺตู”ติ. ตตฺร ปวิสิตฺวา รชคฺคํ\csromansharedfootnote{731ab76cac3c229f}{ราชคฺฆํ (สี, สฺยา)} ภณฺฑํ ปญฺจมาสกํ วา อติเรก\-ปญฺจ\-มาสกํ วา อคฺฆนกํ เถยฺยจิตฺโต อามสติ, อาปตฺติ ทุกฺกฏสฺส. ผนฺทาเปติ, อาปตฺติ ถุลฺลจฺจยสฺส. ปฐมํ ปาทํ สุงฺกฆาตํ อติกฺกาเมติ, อาปตฺติ ถุลฺลจฺจยสฺส. ทุติยํ ปาทํ อติกฺกาเมติ, อาปตฺติ ปาราชิกสฺส. อนฺโตสุงฺกฆาเต ฐิโต พหิสุงฺกฆาตํ ปาเตติ, อาปตฺติ ปาราชิกสฺส. สุงฺกํ ปริหรติ, อาปตฺติ ทุกฺกฏสฺส.}

\proseitem{114}{\textbf{ปาโณ} นาม มนุสฺสปาโณ วุจฺจติ. เถยฺยจิตฺโต อามสติ, อาปตฺติ ทุกฺกฏสฺส. ผนฺทาเปติ, อาปตฺติ ถุลฺลจฺจยสฺส. ฐานา จาเวติ, อาปตฺติ ปาราชิกสฺส. “ปทสา เนสฺสามี”ติ ปฐมํ ปาทํ สงฺกาเมติ, อาปตฺติ ถุลฺลจฺจยสฺส. ทุติยํ ปาทํ สงฺกาเมติ, อาปตฺติ ปาราชิกสฺส.}

\prose{\textbf{อปทํ} นาม อหิ มจฺฉา. ปญฺจมาสกํ วา อติเรก\-ปญฺจ\-มาสกํ วา อคฺฆนกํ เถยฺยจิตฺโต อามสติ, อาปตฺติ ทุกฺกฏสฺส. ผนฺทาเปติ, อาปตฺติ ถุลฺลจฺจยสฺส. ฐานา จาเวติ, อาปตฺติ ปาราชิกสฺส.}

\proseitem{115}{\textbf{ทฺวิปทํ} นาม มนุสฺสา ปกฺขชาตา. เถยฺยจิตฺโต อามสติ, อาปตฺติ ทุกฺกฏสฺส. ผนฺทาเปติ, อาปตฺติ ถุลฺลจฺจยสฺส. ฐานา จาเวติ, อาปตฺติ ปาราชิกสฺส. “ปทสา เนสฺสามี”ติ ปฐมํ ปาทํ สงฺกาเมติ, อาปตฺติ ถุลฺลจฺจยสฺส. ทุติยํ ปาทํ สงฺกาเมติ, อาปตฺติ ปาราชิกสฺส.}

\proseitem{116}{\textbf{จตุปฺปทํ} นาม หตฺถี อสฺสา โอฏฺฐา โคณา คทฺรภา ปสุกา. เถยฺยจิตฺโต อามสติ, อาปตฺติ ทุกฺกฏสฺส. ผนฺทาเปติ, อาปตฺติ ถุลฺลจฺจยสฺส. ฐานา จาเวติ, อาปตฺติ ปาราชิกสฺส. “ปทสา เนสฺสามี”ติ ปฐมํ ปาทํ สงฺกาเมติ, อาปตฺติ ถุลฺลจฺจยสฺส. ทุติยํ ปาทํ สงฺกาเมติ, \csromanfolio{65}อาปตฺติ ถุลฺลจฺจยสฺส. ตติยํ ปาทํ สงฺกาเมติ, อาปตฺติ ถุลฺลจฺจยสฺส. จตุตฺถํ ปาทํ สงฺกาเมติ, อาปตฺติ ปาราชิกสฺส.}

\proseitem{117}{\textbf{พหุปฺปทํ} นาม วิจฺฉิกา สตปที อุจฺจาลิงฺค\-ปาณกา. ปญฺจมาสกํ วา อติเรก\-ปญฺจ\-มาสกํ วา อคฺฆนกํ เถยฺยจิตฺโต อามสติ, อาปตฺติ ทุกฺกฏสฺส. ผนฺทาเปติ, อาปตฺติ ถุลฺลจฺจยสฺส. ฐานา จาเวติ, อาปตฺติ ปาราชิกสฺส. “ปทสา เนสฺสามี”ติ สงฺกาเมติ, ปเท ปเท อาปตฺติ ถุลฺลจฺจยสฺส. ปจฺฉิมํ ปาทํ สงฺกาเมติ, อาปตฺติ ปาราชิกสฺส.}

\proseitem{118}{\textbf{โอจรโก} นาม ภณฺฑํ โอจริตฺวา อาจิกฺขติ “อิตฺถนฺนามํ ภณฺฑํ อวหรา”ติ, อาปตฺติ ทุกฺกฏสฺส. โส ตํ ภณฺฑํ อวหรติ, อาปตฺติ อุภินฺนํ ปาราชิกสฺส.}

\prose{\textbf{โอณิรกฺโข} นาม อาหฏํ ภณฺฑํ โคเปนฺโต ปญฺจมาสกํ วา อติเรก\-ปญฺจ\-มาสกํ วา อคฺฆนกํ เถยฺยจิตฺโต อามสติ, อาปตฺติ ทุกฺกฏสฺส. ผนฺทาเปติ, อาปตฺติ ถุลฺลจฺจยสฺส. ฐานา จาเวติ, อาปตฺติ ปาราชิกสฺส.}

\prose{\textbf{สํวิทาวหาโร} นาม สมฺพหุลา สํวิทหิตฺวา เอโก ภณฺฑํ อวหรติ, อาปตฺติ สพฺเพสํ ปาราชิกสฺส.}

\proseitem{119}{\textbf{สงฺเกตกมฺมํ} นาม สงฺเกตํ กโรติ “ปุเรภตฺตํ วา ปจฺฉาภตฺตํ วา รตฺติํ วา ทิวา วา เตน สงฺเกเตน ตํ ภณฺฑํ อวหรา”ติ, อาปตฺติ ทุกฺกฏสฺส. เตน สงฺเกเตน ตํ ภณฺฑํ อวหรติ, อาปตฺติ อุภินฺนํ ปาราชิกสฺส. ตํ สงฺเกตํ ปุเร วา ปจฺฉา วา ตํ ภณฺฑํ อวหรติ, มูลฏฺฐสฺส อนาปตฺติ. อวหารกสฺส อาปตฺติ ปาราชิกสฺส.}

\proseitem{120}{\textbf{นิมิตฺตกมฺมํ} นาม นิมิตฺตํ กโรติ “อกฺขิํ วา นิขณิสฺสามิ, ภมุกํ วา อุกฺขิปิสฺสามิ, สีสํ วา อุกฺขิปิสฺสามิ, เตน นิมิตฺเตน ตํ ภณฺฑํ อวหรา”ติ, อาปตฺติ ทุกฺกฏสฺส. เตน นิมิตฺเตน ตํ ภณฺฑํ อวหรติ, อาปตฺติ อุภินฺนํ ปาราชิกสฺส. ตํ นิมิตฺตํ ปุเร วา ปจฺฉา วา ตํ ภณฺฑํ อวหรติ, มูลฏฺฐสฺส อนาปตฺติ. อวหารกสฺส อาปตฺติ ปาราชิกสฺส.}

\proseitem{121}{\csromanfolio{66}ภิกฺขุ ภิกฺขุํ อาณาเปติ “อิตฺถนฺนามํ ภณฺฑํ อวหรา”ติ, อาปตฺติ ทุกฺกฏสฺส. โส ตํ มญฺญมาโน ตํ อวหรติ, อาปตฺติ อุภินฺนํ ปาราชิกสฺส.}

\prose{ภิกฺขุ ภิกฺขุํ อาณาเปติ “อิตฺถนฺนามํ ภณฺฑํ อวหรา”ติ, อาปตฺติ ทุกฺกฏสฺส. โส ตํ มญฺญมาโน อญฺญํ อวหรติ, มูลฏฺฐสฺส อนาปตฺติ. อวหารกสฺส อาปตฺติ ปาราชิกสฺส.}

\prose{ภิกฺขุ ภิกฺขุํ อาณาเปติ “อิตฺถนฺนามํ ภณฺฑํ อวหรา”ติ, อาปตฺติ ทุกฺกฏสฺส. โส อญฺญํ มญฺญมาโน ตํ อวหรติ, อาปตฺติ อุภินฺนํ ปาราชิกสฺส.}

\prose{ภิกฺขุ ภิกฺขุํ อาณาเปติ “อิตฺถนฺนามํ ภณฺฑํ อวหรา”ติ, อาปตฺติ ทุกฺกฏสฺส. โส อญฺญํ มญฺญมาโน อญฺญํ อวหรติ, มูลฏฺฐสฺส อนาปตฺติ. อวหารกส อาปตฺติ ปาราชิกสฺส.}

\prose{ภิกฺขุ ภิกฺขุํ อาณาเปติ “อิตฺถนฺนามสฺส ปาวท, อิตฺถนฺนาโม อิตฺถนฺนามสฺส ปาวทตุ อิตฺถนฺนาโม อิตฺถนฺนามํ ภณฺฑํ อวหรตู”ติ, อาปตฺติ ทุกฺกฏสฺส. โส อิตรสฺส อาโรเจติ, อาปตฺติ ทุกฺกฏสฺส. อวหารโก ปฏิคฺคณฺหาติ, มูลฏฺฐสฺส อาปตฺติ ถุลฺลจฺจยสฺส. โส ตํ ภณฺฑํ อวหรติ, อาปตฺติ สพฺเพสํ ปาราชิกสฺส.}

\prose{ภิกฺขุ ภิกฺขุํ อาณาเปติ “อิตฺถนฺนามสฺส ปาวท, อิตฺถนฺนาโม อิตฺถนฺนามสฺส ปาวทตุ อิตฺถนฺนาโม อิตฺถนฺนามํ ภณฺฑํ อวหรตู”ติ, อาปตฺติ ทุกฺกฏสฺส. โส อญฺญํ อาณาเปติ, อาปตฺติ ทุกฺกฏสฺส. อวหารโก ปฏิคฺคณฺหาติ, อาปตฺติ ทุกฺกฏสฺส. โส ตํ ภณฺฑํ อวหรติ, มูลฏฺฐสฺส อนาปตฺติ. อาณาปกสฺส จ อวหารกสฺส จ อาปตฺติ ปาราชิกสฺส.}

\prose{ภิกฺขุ ภิกฺขุํ อาณาเปติ “อิตฺถนฺนามํ ภณฺฑํ อวหรา”ติ, อาปตฺติ ทุกฺกฏสฺส. โส คนฺตฺวา ปุน ปจฺจาคจฺฉติ “นาหํ สกฺโกมิ ตํ ภณฺฑํ อวหริตุนฺ”ติ. โส ปุน อาณาเปติ “ยทา สกฺโกสิ, ตทา ตํ ภณฺฑํ อวหรา”ติ, อาปตฺติ ทุกฺกฏสฺส. โส ตํ ภณฺฑํ อวหรติ, อาปตฺติ อุภินฺนํ ปาราชิกสฺส.}

\prose{\csromanfolio{67}ภิกฺขุ ภิกฺขุํ อาณาเปติ “อิตฺถนฺนามํ ภณฺฑํ อวหรา”ติ, อาปตฺติ ทุกฺกฏสฺส. โส อาณาเปตฺวา วิปฺปฏิสารี น สาเวติ “มา อวหรี”ติ. โส ตํ ภณฺฑํ อวหรติ, อาปตฺติ อุภินฺนํ ปาราชิกสฺส.}

\prose{ภิกฺขุ ภิกฺขุํ อาณาเปติ “อิตฺถนฺนามํ ภณฺฑํ อวหรา”ติ, อาปตฺติ ทุกฺกฏสฺส. โส อาณาเปตฺวา วิปฺปฏิสารี สาเวติ “มา อวหรี”ติ. โส “อาณตฺโต อหํ ตยา”ติ ตํ ภณฺฑํ อวหรติ, มูลฏฺฐสฺส อนาปตฺติ. อวหารกสฺส อาปตฺติ ปาราชิกสฺส.}

\prose{ภิกฺขุ ภิกฺขุํ อาณาเปติ “อิตฺถนฺนามํ ภณฺฑํ อวหรา”ติ, อาปตฺติ ทุกฺกฏสฺส. โส อาณาเปตฺวา วิปฺปฏิสารี สาเวติ “มา อวหรี”ติ. โส “สาธู”ติ\csromansharedfootnote{5881e4c76a2ee19c}{สุฏฺฐูติ (ก)} โอรมติ, อุภินฺนํ อนาปตฺติ.}

\proseitem{122}{ปญฺจหิ อากาเรหิ อทินฺนํ อาทิยนฺตสฺส อาปตฺติ ปาราชิกสฺส. ปรปริคฺคหิตญฺจ โหติ, ปรปริคฺคหิต\-สญฺญี จ, ครุโก จ โหติ ปริกฺขาโร ปญฺจมาสโก วา อติเรก\-ปญฺจ\-มาสโก วา, เถยฺยจิตฺตญฺจ ปจฺจุปฏฺฐิตํ โหติ, อามสติ, อาปตฺติ ทุกฺกฏสฺส. ผนฺทาเปติ, อาปตฺติ ถุลฺลจฺจยสฺส. ฐานา จาเวติ, อาปตฺติ ปาราชิกสฺส.}

\proseitem{123}{ปญฺจหิ อากาเรหิ อทินฺนํ อาทิยนฺตสฺส อาปตฺติ ถุลฺลจฺจยสฺส. ปรปริคฺคหิตญฺจ โหติ, ปรปริคฺคหิต\-สญฺญี จ, ลหุโก จ โหติ ปริกฺขาโร อติเรกมาสโก วา อูนปญฺจมาสโก วา, เถยฺยจิตฺตญฺจ ปจฺจุปฏฺฐิตํ โหติ, อามสติ, อาปตฺติ ทุกฺกฏสฺส. ผนฺทาเปติ, อาปตฺติ ทุกฺกฏสฺส. ฐานา จาเวติ, อาปตฺติ ถุลฺลจฺจยสฺส.}

\proseitem{124}{ปญฺจหิ อากาเรหิ อทินฺนํ อาทิยนฺตสฺส อาปตฺติ ทุกฺกฏสฺส. ปรปริคฺคหิตญฺจ โหติ, ปรปริคฺคหิต\-สญฺญี จ, ลหุโก จ โหติ ปริกฺขาโร มาสโก วา อูนมาสโก วา, เถยฺยจิตฺตญฺจ ปจฺจุปฏฺฐิตํ โหติ, อามสติ, อาปตฺติ ทุกฺกฏสฺส. ผนฺทาเปติ, อาปตฺติ ทุกฺกฏสฺส. ฐานา จาเวติ, อาปตฺติ ทุกฺกฏสฺส.}

\proseitem{125}{ฉหิ อากาเรหิ อทินฺนํ อาทิยนฺตสฺส อาปตฺติ ปาราชิกสฺส. น จ สกสญฺญี, น จ วิสฺสาสคฺคาหี, น จ ตาวกาลิกํ, ครุโก จ \csromanfolio{68}โหติ ปริกฺขาโร ปญฺจมาสโก วา อติเรก\-ปญฺจ\-มาสโก วา, เถยฺยจิตฺตญฺจ ปจฺจุปฏฺฐิตํ โหติ, อามสติ, อาปตฺติ ทุกฺกฏสฺส. ผนฺทาเปติ, อาปตฺติ ถุลฺลจฺจยสฺส. ฐานา จาเวติ, อาปตฺติ ปาราชิกสฺส.}

\proseitem{126}{ฉหิ อากาเรหิ อทินฺนํ อาทิยนฺตสฺส อาปตฺติ ถุลฺลจฺจยสฺส. น จ สกสญฺญี, น จ วิสฺสาสคฺคาหี, น จ ตาวกาลิกํ, ลหุโก จ โหติ ปริกฺขาโร อติเรกมาสโก วา อูนปญฺจมาสโก วา, เถยฺยจิตฺตญฺจ ปจฺจุปฏฺฐิตํ โหติ, อามสติ, อาปตฺติ ทุกฺกฏสฺส. ผนฺทาเปติ, อาปตฺติ ทุกฺกฏสฺส. ฐานา จาเวติ, อาปตฺติ ถุลฺลจฺจยสฺส.}

\proseitem{127}{ฉหิ อากาเรหิ อทินฺนํ อาทิยนฺตสฺส อาปตฺติ ทุกฺกฏสฺส. น จ สกสญฺญี, น จ วิสฺสาสคฺคาหี, น จ ตาวกาลิกํ, ลหุโก จ โหติ ปริกฺขาโร มาสโก วา อูนมาสโก วา, เถยฺยจิตฺตญฺจ ปจฺจุปฏฺฐิตํ โหติ, อามสติ, อาปตฺติ ทุกฺกฏสฺส. ผนฺทาเปติ, อาปตฺติ ทุกฺกฏสฺส. ฐานา จาเวติ, อาปตฺติ ทุกฺกฏสฺส.}

\proseitem{128}{ปญฺจหิ อากาเรหิ อทินฺนํ อาทิยนฺตสฺส อาปตฺติ ทุกฺกฏสฺส. น จ ปรปริคฺคหิตํ โหติ, ปรปริคฺคหิต\-สญฺญี จ, ครุโก จ โหติ ปริกฺขาโร ปญฺจมาสโก วา อติเรก\-ปญฺจ\-มาสโก วา, เถยฺยจิตฺตญฺจ ปจฺจุปฏฺฐิตํ โหติ, อามสติ, อาปตฺติ ทุกฺกฏสฺส. ผนฺทาเปติ, อาปตฺติ ทุกฺกฏสฺส. ฐานา จาเวติ, อาปตฺติ ทุกฺกฏสฺส.}

\proseitem{129}{ปญฺจหิ อากาเรหิ อทินฺนํ อาทิยนฺตสฺส อาปตฺติ ทุกฺกฏสฺส. น จ ปรปริคฺคหิตํ โหติ, ปรปริคฺคหิต\-สญฺญี จ, ลหุโก จ โหติ ปริกฺขาโร อติเรกมาสโก วา อูนปญฺจมาสโก วา, เถยฺยจิตฺตญฺจ ปจฺจุปฏฺฐิตํ โหติ, อามสติ, อาปตฺติ ทุกฺกฏสฺส. ผนฺทาเปติ, อาปตฺติ ทุกฺกฏสฺส. ฐานา จาเวติ, อาปตฺติ ทุกฺกฏสฺส.}

\proseitem{130}{ปญฺจหิ อากาเรหิ อทินฺนํ อาทิยนฺตสฺส อาปตฺติ ทุกฺกฏสฺส. น จ ปรปริคฺคหิตํ โหติ, ปรปริคฺคหิต\-สญฺญี จ, ลหุโก จ โหติ ปริกฺขาโร มาสโก วา อูนมาสโก วา, เถยฺยจิตฺตญฺจ ปจฺจุปฏฺฐิตํ โหติ, อามสติ, อาปตฺติ ทุกฺกฏสฺส. ผนฺทาเปติ, อาปตฺติ ทุกฺกฏสฺส. ฐานา จาเวติ, อาปตฺติ ทุกฺกฏสฺส.}

\proseitemwithcloserruled{131}{\csromanfolio{69}อนาปตฺติ สสญฺญิสฺส วิสฺสาสคฺคาเห ตาวกาลิเก เปตปริคฺคเห ติรจฺฉานคต\-ปริคฺคเห ปํสุกูล\-สญฺญิสฺส อุมฺมตฺตกสฺส (ขิตฺตจิตฺตสฺส เวทนาฏฺฏสฺส)\csromansharedfootnote{39593217a53c162e}{(ขิตฺตจิตฺตสฺส เวทนาฏฺฏสฺส) กตฺถจิ นตฺถิ.} อาทิกมฺมิกสฺสาติ.}{อทินฺนาทานมฺหิ ปฐม\-ภาณวาโร นิฏฺฐิโต.}
\csromanmatikaanchor{mtk.20}
\csromantocmarkref{subsection}{วินีตวตฺถุอุทฺทานคาถา}{mtk.20}
\bookmark[dest={mtk.20},level=2]{วินีตวตฺถุอุทฺทานคาถา}
\csromanheader{2}{วินีตวตฺถุ\-อุทฺทานคาถา}
\setlength{\csromangathapagewidth}{0pt}
\setlength{\csromangathawakwidth}{0pt}
\csromangathameasuregroup{รชเกหิ ปญฺจ อกฺขาตา, \\ อนฺธกาเรน เว ปญฺจ, \\ นิรุตฺติยา ปญฺจ อกฺขาตา, \\ อสมฺภินฺเน กุสาปาโต, \\ วิฆาเสหิ ปญฺจ อกฺขาตา, \\ ทุพฺภิกฺเข กุรมํสญฺจ\textsuperscript{1}, \\ ฉปริกฺขารถวิกา, \\ ขาทนียญฺจ วิสฺสาสํ, \\ สตฺต นาวหรามาติ, \\ สํฆสฺส อวหรุํ สตฺต, \\ ตโย จ วุตฺตวาทิโน, \\ สูกรา จ มิคา มจฺฉา, \\ ทุเว เปสี ทุเว ทารู, \\ อนุปุพฺพวิธาเนน, \\ สาวตฺถิยา จตุโร มุฏฺฐี, \\ สํฆสฺส ภาชยุํ สตฺต, \\ จมฺปา ราชคเห เจว, \\ พาราณสี จ โกสมฺพี,}{\csromangathabat{รชเกหิ ปญฺจ อกฺขาตา,}{จตุโร อตฺถรเณหิ จ.} \\ \csromangathabat{อนฺธกาเรน เว ปญฺจ,}{ปญฺจ หารณเกน จ.} \\ \csromangathabat{นิรุตฺติยา ปญฺจ อกฺขาตา,}{วาเตหิ อปเร ทุเว.} \\ \csromangathabat{อสมฺภินฺเน กุสาปาโต,}{ชนฺตคฺเคน\textsuperscript{1} สหา ทส.} \\ \csromangathabat{วิฆาเสหิ ปญฺจ อกฺขาตา,}{ปญฺจ เจว อมูลกา.} \\ \csromangathabat{ทุพฺภิกฺเข กุรมํสญฺจ\textsuperscript{1},}{ปูวสกฺขลิโมทกา.} \\ \csromangathabat{ฉปริกฺขารถวิกา,}{ภิสิวํสา น นิกฺขเม.} \\ \csromangathabat{ขาทนียญฺจ วิสฺสาสํ,}{สสญฺญา ยปเร ทุเว.} \\ \csromangathabat{สตฺต นาวหรามาติ,}{สตฺต เจว อวาหรุํ.} \\ \csromangathabat{สํฆสฺส อวหรุํ สตฺต,}{ปุปฺเผหิ อปเร ทุเว.} \\ \csromangathabat{ตโย จ วุตฺตวาทิโน,}{มณิ ตีณิ อติกฺกเม.} \\ \csromangathabat{สูกรา จ มิคา มจฺฉา,}{ยานญฺจาปิ ปวตฺตยิ.} \\ \csromangathabat{ทุเว เปสี ทุเว ทารู,}{ปํสุกูลํ ทุเว ทกา.} \\ \csromangathabat{อนุปุพฺพวิธาเนน,}{ตทญฺโญ น ปริปูรยิ.} \\ \csromangathabat{สาวตฺถิยา จตุโร มุฏฺฐี,}{เทฺว วิฆาสา ทุเว ติณา.} \\ \csromangathabat{สํฆสฺส ภาชยุํ สตฺต,}{สตฺต เจว อสฺสามิกา.} \\ ทารุทกา มตฺติกา เทฺว ติณานิ, \\ สํฆสฺส สตฺต อวหาสิ เสยฺยํ. \\ สสฺสามิกํ น จาปิ นีหเรยฺย, \\ หเรยฺย สสฺสามิกํ ตาวกาลิกํ. \\ \csromangathabat{จมฺปา ราชคเห เจว,}{เวสาลิยา จ อชฺชุโก.} \\ \csromangathabat{พาราณสี จ โกสมฺพี,}{สาคลา ทฬฺหิเกน จาติ.}}
\csromangathameasurewak{ทารุทกา มตฺติกา เทฺว ติณานิ, \\ สํฆสฺส สตฺต อวหาสิ เสยฺยํ. \\ สสฺสามิกํ น จาปิ นีหเรยฺย, \\ หเรยฺย สสฺสามิกํ ตาวกาลิกํ.}
\csromangathasetleft{รชเกหิ ปญฺจ อกฺขาตา, \\ อนฺธกาเรน เว ปญฺจ, \\ นิรุตฺติยา ปญฺจ อกฺขาตา, \\ อสมฺภินฺเน กุสาปาโต, \\ วิฆาเสหิ ปญฺจ อกฺขาตา, \\ ทุพฺภิกฺเข กุรมํสญฺจ\textsuperscript{1}, \\ ฉปริกฺขารถวิกา, \\ ขาทนียญฺจ วิสฺสาสํ, \\ สตฺต นาวหรามาติ, \\ สํฆสฺส อวหรุํ สตฺต, \\ ตโย จ วุตฺตวาทิโน, \\ สูกรา จ มิคา มจฺฉา, \\ ทุเว เปสี ทุเว ทารู, \\ อนุปุพฺพวิธาเนน, \\ สาวตฺถิยา จตุโร มุฏฺฐี, \\ สํฆสฺส ภาชยุํ สตฺต, \\ จมฺปา ราชคเห เจว, \\ พาราณสี จ โกสมฺพี,}
\csromangathagroup{\csromangathastanza{\csromangathabat{รชเกหิ ปญฺจ อกฺขาตา,}{จตุโร อตฺถรเณหิ จ.} \\ \csromangathabat{อนฺธกาเรน เว ปญฺจ,}{ปญฺจ หารณเกน จ.}}\csromangathastanza{\csromangathabat{นิรุตฺติยา ปญฺจ อกฺขาตา,}{วาเตหิ อปเร ทุเว.} \\ \csromangathabat{อสมฺภินฺเน กุสาปาโต,}{ชนฺตคฺเคน\csromansharedfootnote{5063acc3f9acbbe3}{ชนฺตาฆเรน (สฺยา)} สหา ทส.}}\csromangathastanza{\csromangathabat{วิฆาเสหิ ปญฺจ อกฺขาตา,}{ปญฺจ เจว อมูลกา.} \\ \csromangathabat{ทุพฺภิกฺเข กุรมํสญฺจ\csromansharedfootnote{c88d4b420eb13a03}{กูรมํสญฺจ (สฺยา)},}{ปูวสกฺขลิโมทกา.}}\csromangathastanza{\csromangathabat{ฉปริกฺขารถวิกา,}{ภิสิวํสา น นิกฺขเม.} \\ \csromangathabat{ขาทนียญฺจ วิสฺสาสํ,}{สสญฺญา ยปเร ทุเว.}}\csromangathastanza{\csromangathabat{สตฺต นาวหรามาติ,}{สตฺต เจว อวาหรุํ.} \\ \csromangathabat{สํฆสฺส อวหรุํ สตฺต,}{ปุปฺเผหิ อปเร ทุเว.}}\csromangathastanza{\csromangathabat{ตโย จ วุตฺตวาทิโน,}{มณิ ตีณิ อติกฺกเม.} \\ \csromangathabat{สูกรา จ มิคา มจฺฉา,}{ยานญฺจาปิ ปวตฺตยิ.}}\csromangathastanza{\csromangathabat{ทุเว เปสี ทุเว ทารู,}{ปํสุกูลํ ทุเว ทกา.} \\ \csromangathabat{อนุปุพฺพวิธาเนน,}{ตทญฺโญ น ปริปูรยิ.}}\csromangathastanza{\csromanfolio{70}\csromangathabat{สาวตฺถิยา จตุโร มุฏฺฐี,}{เทฺว วิฆาสา ทุเว ติณา.} \\ \csromangathabat{สํฆสฺส ภาชยุํ สตฺต,}{สตฺต เจว อสฺสามิกา.}}\csromangathastanza{\csromangathawak{ทารุทกา มตฺติกา เทฺว ติณานิ, \\ สํฆสฺส สตฺต อวหาสิ เสยฺยํ. \\ สสฺสามิกํ น จาปิ นีหเรยฺย, \\ หเรยฺย สสฺสามิกํ ตาวกาลิกํ.}}\csromangathastanza{\csromangathabat{จมฺปา ราชคเห เจว,}{เวสาลิยา จ อชฺชุโก.} \\ \csromangathabat{พาราณสี จ โกสมฺพี,}{สาคลา ทฬฺหิเกน จาติ.}}}

\csromanmatikaanchor{mtk.21}
\csromantocmarkref{subsection}{วินีตวตฺถุ}{mtk.21}
\bookmark[dest={mtk.21},level=2]{วินีตวตฺถุ}
\csromanheader{2}{\textbf{วินีตวตฺถุ}}
\proseitem{132}{เตน โข ปน สมเยน ฉพฺพคฺคิยา ภิกฺขู รชก\-ตฺถรณํ คนฺตฺวา รชก\-ภณฺฑิกํ อวหริํสุ. เตสํ กุกฺกุจฺจํ อโหสิ “ภควตา สิกฺขาปทํ ปญฺญตฺตํ, กจฺจิ นุ โข มยํ ปาราชิกํ อาปตฺติํ อาปนฺนา”ติ. ภควโต เอตมตฺถํ อาโรเจสุํ. อาปตฺติํ ตุเมฺห ภิกฺขเว อาปนฺนา ปาราชิกนฺติ. (1)}

\prose{เตน โข ปน สมเยน อญฺญตโร ภิกฺขุ รชก\-ตฺถรณํ คนฺตฺวา มหคฺฆํ ทุสฺสํ ปสฺสิตฺวา เถยฺยจิตฺตํ อุปฺปาเทสิ. ตสฺส กุกฺกุจฺจํ อโหสิ “ภควตา สิกฺขาปทํ ปญฺญตฺตํ, กจฺจิ นุ โข อหํ ปาราชิกํ อาปตฺติํ อาปนฺโน”ติ. ภควโต เอตมตฺถํ อาโรเจสิ. อนาปตฺติ ภิกฺขุ จิตฺตุปฺปาเทติ. (2)}

\prose{เตน โข ปน สมเยน อญฺญตโร ภิกฺขุ รชก\-ตฺถรณํ คนฺตฺวา มหคฺฆํ ทุสฺสํ ปสฺสิตฺวา เถยฺยจิตฺโต อามสิ. ตสฺส กุกฺกุจฺจํ อโหสิ ฯเปฯ. อนาปตฺติ ภิกฺขุ ปาราชิกสฺส, อาปตฺติ ทุกฺกฏสฺสาติ. (3)}

\prose{เตน โข ปน สมเยน อญฺญตโร ภิกฺขุ รชก\-ตฺถรณํ คนฺตฺวา มหคฺฆํ ทุสฺสํ ปสฺสิตฺวา เถยฺยจิตฺโต ผนฺทาเปสิ. ตสฺส กุกฺกุจฺจํ อโหสิ ฯเปฯ. อนาปตฺติ ภิกฺขุ ปาราชิกสฺส, อาปตฺติ ถุลฺลจฺจยสฺสาติ. (4)}

\prose{\csromanfolio{71}เตน โข ปน สมเยน อญฺญตโร ภิกฺขุ รชก\-ตฺถรณํ คนฺตฺวา มหคฺฆํ ทุสฺสํ ปสฺสิตฺวา เถยฺยจิตฺตํ ฐานา จาเวสิ. ตสฺส กุกฺกุจฺจํ อโหสิ ฯเปฯ. อาปตฺติํ ตฺวํ ภิกฺขุ อาปนฺโน ปาราชิกนฺติ. (5)}

\proseitem{133}{เตน โข ปน สมเยน อญฺญตโร ปิณฺฑจาริโก ภิกฺขุ มหคฺฆํ อุตฺตร\-ตฺถรณํ ปสฺสิตฺวา เถยฺยจิตฺตํ อุปฺปาเทสิ ฯเปฯ เถยฺยจิตฺโต อามสิ ฯเปฯ เถยฺยจิตฺโต ผนฺทาเปสิ ฯเปฯ เถยฺยจิตฺโต ฐานา จาเวสิ. ตสฺส กุกฺกุจฺจํ อโหสิ ฯเปฯ. อาปตฺติํ ตฺวํ ภิกฺขุ อาปนฺโน ปาราชิกนฺติ. \mbox{(6-9)}}

\proseitem{134}{เตน โข ปน สมเยน อญฺญตโร ภิกฺขุ ทิวา ภณฺฑํ ปสฺสิตฺวา นิมิตฺตํ อกาสิ “รตฺติํ อวหริสฺสามี”ติ. โส ตํ มญฺญมาโน ตํ อวหริ ฯเปฯ ตํ มญฺญมาโน อญฺญํ อวหริ ฯเปฯ อญฺญํ มญฺญมาโน ตํ อวหริ ฯเปฯ อญฺญํ มญฺญมาโน อญฺญํ อวหริ. ตสฺส กุกฺกุจฺจํ อโหสิ ฯเปฯ. อาปตฺติํ ตฺวํ ภิกฺขุ อาปนฺโน ปาราชิกนฺติ. \mbox{(10-13)}}

\prose{เตน โข ปน สมเยน อญฺญตโร ภิกฺขุ ทิวา ภณฺฑํ ปสฺสิตฺวา นิมิตฺตํ อกาสิ “รตฺติํ อวหริสฺสามี”ติ. โส ตํ มญฺญมาโน อตฺตโน ภณฺฑํ อวหริ. ตสฺส กุกฺกุจฺจํ อโหสิ ฯเปฯ. อนาปตฺติ ภิกฺขุ ปาราชิกสฺส, อาปตฺติ ทุกฺกฏสฺสาติ. (14)}

\prose{เตน โข ปน สมเยน อญฺญตโร ภิกฺขุ อญฺญสฺส ภณฺฑํ หรนฺโต สีเส ภารํ เถยฺยจิตฺโต อามสิ ฯเปฯ เถยฺยจิตฺโต ผนฺทาเปสิ ฯเปฯ เถยฺยจิตฺโต ขนฺธํ โอโรเปสิ ฯเปฯ ขนฺเธ ภารํ เถยฺยจิตฺโต อามสิ ฯเปฯ เถยฺยจิตฺโต ผนฺทาเปสิ ฯเปฯ เถยฺยจิตฺโต กฏิํ โอโรเปสิ ฯเปฯ กฏิยา ภารํ เถยฺยจิตฺโต อามสิ ฯเปฯ เถยฺยจิตฺโต ผนฺทาเปสิ ฯเปฯ เถยฺยจิตฺโต หตฺเถน อคฺคเหสิ ฯเปฯ หตฺเถ ภารํ เถยฺยจิตฺโต ภูมิยํ นิกฺขิปิ ฯเปฯ เถยฺยจิตฺโต ภูมิโต อคฺคเหสิ. ตสฺส กุกฺกุจฺจํ อโหสิ ฯเปฯ. อาปตฺติํ ตฺวํ ภิกฺขุ อาปนฺโน ปาราชิกนฺติ. \mbox{(15-25)}}

\proseitem{135}{เตน โข ปน สมเยน อญฺญตโร ภิกฺขุ อชฺโฌกาเส จีวรํ ปตฺถริตฺวา วิหารํ ปาวิสิ. อญฺญตโร ภิกฺขุ “มายิทํ จีวรํ นสฺสี”ติ ปฏิสาเมสิ. โส นิกฺขมิตฺวา ตํ ภิกฺขุํ ปุจฺฉิ “อาวุโส มยฺหํ จีวรํ เกน \csromanfolio{72}อวหฏนฺ”ติ. โส เอวมาห “มยา อวหฏนฺ”ติ. โส ตํ อาทิยิ “อสฺสมโณสิ ตฺวนฺ”ติ. ตสฺส กุกฺกุจฺจํ อโหสิ ฯเปฯ. ภควโต เอตมตฺถํ อาโรเจสิ. กิํ จิตฺโต ตฺวํ ภิกฺขูติ. นิรุตฺติปโถ อหํ ภควาติ. อนาปตฺติ ภิกฺขุ นิรุตฺติปเถติ. (26)}

\prose{เตน โข ปน สมเยน อญฺญตโร ภิกฺขุ ปีเฐ จีวรํ นิกฺขิปิตฺวา. ปีเฐ นิสีทนํ นิกฺขิปิตฺวา. เหฏฺฐาปีเฐ ปตฺตํ นิกฺขิปิตฺวา วิหารํ ปาวิสิ. อญฺญตโร ภิกฺขุ “มายํ ปตฺโต นสฺสี”ติ ปฏิสาเมสิ. โส นิกฺขมิตฺวา ตํ ภิกฺขุํ ปุจฺฉิ “อาวุโส มยฺหํ ปตฺโต เกน อวหโฏ”ติ. โส เอวมาห “มยา อวหโฏ”ติ. โส ตํ อาทิยิ “อสฺสมโณสิ ตฺวนฺ”ติ. ตสฺส กุกฺกุจฺจํ อโหสิ ฯเปฯ. อนาปตฺติ ภิกฺขุ นิรุตฺติปเถติ. \mbox{(27-29)}}

\prose{เตน โข ปน สมเยน อญฺญตรา ภิกฺขุนี วติยา จีวรํ ปตฺถริตฺวา วิหารํ ปาวิสิ. อญฺญตรา ภิกฺขุนี “มายิทํ จีวรํ นสฺสี”ติ ปฏิสาเมสิ. สา นิกฺขมิตฺวา ตํ ภิกฺขุนิํ ปุจฺฉิ “อยฺเย มยฺหํ จีวรํ เกน อวหฏนฺ”ติ. สา เอวมาห “มยา อวหฏนฺ”ติ. สา ตํ อาทิยิ “อสฺสมณีสิ ตฺวนฺ”ติ. ตสฺสา กุกฺกุจฺจํ อโหสิ ฯเปฯ. อถ โข สา ภิกฺขุนี ภิกฺขุนีนํ เอตมตฺถํ อาโรเจสิ. ภิกฺขุนิโย ภิกฺขูนํ เอตมตฺถํ อาโรเจสุํ. ภิกฺขู ภควโต เอตมตฺถํ อาโรเจสุํ. อนาปตฺติ ภิกฺขเว นิรุตฺติปเถติ. (30)}

\proseitem{136}{เตน โข ปน สมเยน อญฺญตโร ภิกฺขุ วาต\-มณฺฑลิกาย อุกฺขิตฺตํ สาฏกํ ปสฺสิตฺวา “สามิกานํ ทสฺสามี”ติ อคฺคเหสิ. สามิกา ตํ ภิกฺขุํ โจเทสุํ “อสฺสมโณสิ ตฺวนฺ”ติ. ตสฺส กุกฺกุจฺจํ อโหสิ ฯเปฯ. กิํจิตฺโต ตฺวํ ภิกฺขูติ. อเถยฺยจิตฺโต อหํ ภควาติ. อนาปตฺติ ภิกฺขุ อเถยฺยจิตฺตสฺสาติ. (31)}

\prose{เตน โข ปน สมเยน อญฺญตโร ภิกฺขุ วาต\-มณฺฑลิกาย อุกฺขิตฺตํ เวฐนํ ปสฺสิตฺวา “ปุเร สามิกา ปสฺสนฺตี”ติ เถยฺยจิตฺโต อคฺคเหสิ. สามิกา ตํ ภิกฺขุํ โจเทสุํ “อสฺสมโณสิ ตฺวนฺ”ติ. ตสฺส กุกฺกุจฺจํ อโหสิ ฯเปฯ. อาปตฺติํ ตฺวํ ภิกฺขุ อาปนฺโน ปาราชิกนฺติ. (32)}

\proseitem{137}{เตน โข ปน สมเยน อญฺญตโร ภิกฺขุ สุสานํ คนฺตฺวา อภินฺเน สรีเร ปํสุกูลํ อคฺคเหสิ. ตสฺมิํ จ สรีเร เปโต อธิวตฺโถ \csromanfolio{73}โหติ. อถ โข โส เปโต ตํ ภิกฺขุํ เอตทโวจ “มา ภนฺเต มยฺหํ สาฏกํ อคฺคเหสี”ติ. โส ภิกฺขุ อนาทิยนฺโต อคมาสิ. อถ โข ตํ สรีรํ อุฏฺฐหิตฺวา ตสฺส ภิกฺขุโน ปิฏฺฐิโต ปิฏฺฐิโต อนุพนฺธิ. อถ โข โส ภิกฺขุ วิหารํ ปวิสิตฺวา ทฺวารํ ถเกสิ. อถ โข ตํ สรีรํ ตตฺเถว ปริปติ. ตสฺส กุกฺกุจฺจํ อโหสิ ฯเปฯ. อนาปตฺติ ภิกฺขุ ปาราชิกสฺส, น จ ภิกฺขเว อภินฺเน สรีเร ปํสุกูลํ คเหตพฺพํ, โย คเณฺหยฺย, อาปตฺติ ทุกฺกฏสฺสาติ. (33)}

\proseitem{138}{เตน โข ปน สมเยน อญฺญตโร ภิกฺขุ สํฆสฺส จีวเร ภาชียมาเน เถยฺยจิตฺโต กุสํ สงฺกาเมตฺวา จีวรํ อคฺคเหสิ. ตสฺส กุกฺกุจฺจํ อโหสิ ฯเปฯ. อาปตฺติํ ตฺวํ ภิกฺขุ อาปนฺโน ปาราชิกนฺติ. (34)}

\proseitem{139}{เตน โข ปน สมเยน อายสฺมา อานนฺโท ชนฺตาฆเร อญฺญตรสฺส ภิกฺขุโน อนฺตรวาสกํ อตฺตโน มญฺญมาโน นิวาเสสิ. อถ โข โส ภิกฺขุ อายสฺมนฺตํ อานนฺทํ เอตทโวจ “กิสฺส เม ตฺวํ อาวุโส อานนฺท อนฺตรวาสกํ นิวาเสสี”ติ. สกสญฺญี อหํ อาวุโสติ. ภควโต เอตมตฺถํ อาโรเจสุํ. อนาปตฺติ ภิกฺขเว สกสญฺญิสฺสาติ. (35)}

\proseitem{140}{เตน โข ปน สมเยน สมฺพหุลา ภิกฺขู คิชฺฌกูฏา ปพฺพตา โอโรหนฺตา สีควิฆาสํ ปสฺสิตฺวา ปจาเปตฺวา ปริภุญฺชิํสุ. เตสํ กุกฺกุจฺจํ อโหสิ ฯเปฯ. อนาปตฺติ ภิกฺขเว สีหวิฆาเสติ. (36)}

\prose{เตน โข ปน สมเยน สมฺพหุลา ภิกฺขู คิชฺฌกูฏา ปพฺพตา โอโรหนฺตา พฺยคฺฆ\-วิฆาสํ ปสฺสิตฺวา. ทีปิวิฆาสํ ปสฺสิตฺวา. ตรจฺฉ\-วิฆาสํ ปสฺสิตฺวา. โกกวิฆาสํ ปสฺสิตฺวา ปจาเปตฺวา ปริภุญฺชิํสุ. เตสํ กุกฺกุจฺจํ อโหสิ ฯเปฯ. อนาปตฺติ ภิกฺขเว ติรจฺฉานคต\-ปริคฺคเหติ. \mbox{(37-40)}}

\proseitem{141}{เตน โข ปน สมเยน อญฺญตโร ภิกฺขุ สํฆสฺส โอทเน ภาชียมาเน “อปรสฺส ภาคํ เทหี”ติ อมูลกํ อคฺคเหสิ. ตสฺส กุกฺกุจฺจํ อโหสิ ฯเปฯ. อนาปตฺติ ภิกฺขุ ปาราชิกสฺส, อาปตฺติ สมฺปชาน\-มุสาวาเท ปาจิตฺติยสฺสาติ. (41)}

\prose{\csromanfolio{74}เตน โข ปน สมเยน อญฺญตโร ภิกฺขุ สํฆสฺส ขาทนีเย ภาชียมาเน. สํฆสฺส ปูเว ภาชียมาเน. สํฆสฺส อุจฺฉุมฺหิ ภาชียมาเน. สํฆสฺส ติมฺพรูสเก ภาชียมาเน “อปรสฺส ภาคํ เทหี”ติ อมูลกํ อคฺคเหสิ. ตสฺส กุกฺกุจฺจํ อโหสิ ฯเปฯ. อนาปตฺติ ภิกฺขุ ปาราชิกสฺส, อาปตฺติ สมฺปชาน\-มุสาวาเท ปาจิตฺติยสฺสาติ. \mbox{(42-45)}}

\proseitem{142}{เตน โข ปน สมเยน อญฺญตโร ภิกฺขุ ทุพฺภิกฺเข โอทนียฆรํ ปวิสิตฺวา ปตฺตปูรํ โอทนํ เถยฺยจิตฺโต อวหริ. ตสฺส กุกฺกุจฺจํ อโหสิ ฯเปฯ. อาปตฺติํ ตฺวํ ภิกฺขุ อาปนฺโน ปาราชิกนฺติ. (46)}

\prose{เตน โข ปน สมเยน อญฺญตโร ภิกฺขุ ทุพฺภิกฺเข สูนฆรํ\csromansharedfootnote{fb214f2e13df4ccf}{สูนาฆรํ (สี, สฺยา)} ปวิสิตฺวา ปตฺตปูรํ มํสํ เถยฺยจิตฺโต อวหริ. ตสฺส กุกฺกุจฺจํ อโหสิ ฯเปฯ. อาปตฺติํ ตฺวํ ภิกฺขุ อาปนฺโน ปาราชิกนฺติ. (47)}

\prose{เตน โข ปน สมเยน อญฺญตโร ภิกฺขุ ทุพฺภิกฺเข ปูวฆรํ ปวิสิตฺวา ปตฺตปูรํ ปูวํ เถยฺยจิตฺโต อวหริ ฯเปฯ ปตฺตปูรา สกฺขลิโย เถยฺยจิตฺโต อวหริ ฯเปฯ ปตฺตปูเร โมทเก เถยฺยจิตฺโต อวหริ. ตสฺส กุกฺกุจฺจํ อโหสิ ฯเปฯ. อาปตฺติํ ตฺวํ ภิกฺขุ อาปนฺโน ปาราชิกนฺติ. \mbox{(48-50)}}

\proseitem{143}{เตน โข ปน สมเยน อญฺญตโร ภิกฺขุ ทิวา ปริกฺขารํ ปสฺสิตฺวา นิมิตฺตํ อกาสิ “รตฺติํ อวหริสฺสามี”ติ. โส ตํ มญฺญมาโน ตํ อวหริ ฯเปฯ ตํ มญฺญมาโน อญฺญํ อวหริ ฯเปฯ อญฺญํ อญฺญมาโน ตํ อวหริ ฯเปฯ อญฺญํ มญฺญมาโน อญฺญํ อวหริ. ตสฺส กุกฺกุจฺจํ อโหสิ ฯเปฯ. อาปตฺติํ ตฺวํ ภิกฺขุ อาปนฺโน ปาราชิกนฺติ. \mbox{(51-54)}}

\prose{เตน โข ปน สมเยน อญฺญตโร ภิกฺขุ ทิวา ปริกฺขารํ ปสฺสิตฺวา นิมิตฺตํ อกาสิ “รตฺติํ อวหริสฺสามี”ติ. โส ตํ มญฺญมาโน อตฺตโน ปริกฺขารํ อวหริ. ตสฺส กุกฺกุจฺจํ อโหสิ ฯเปฯ. อนาปตฺติ ภิกฺขุ ปาราชิกสฺส, อาปตฺติ ทุกฺกฏสฺสาติ. (55)}

\proseitem{144}{เตน โข ปน สมเยน อญฺญตโร ภิกฺขุ ปีเฐ ถวิกํ ปสฺสิตฺวา “อิโต คณฺหนฺโต ปาราชิโก ภวิสฺสามี”ติ สห ปีฐเกน \csromanfolio{75}สงฺกาเมตฺวา อคฺคเหสิ. ตสฺส กุกฺกุจฺจํ อโหสิ ฯเปฯ. อาปตฺติํ ตฺวํ ภิกฺขุ อาปนฺโน ปาราชิกนฺติ. (56)}

\prose{เตน โข ปน สมเยน อญฺญตโร ภิกฺขุ สํฆสฺส ภิสิํ เถยฺยจิตฺโต อวหริ. ตสฺส กุกฺกุจฺจํ อโหสิ ฯเปฯ. อาปตฺติํ ตฺวํ ภิกฺขุ อาปนฺโน ปาราชิกนฺติ. (57)}

\proseitem{145}{เตน โข ปน สมเยน อญฺญตโร ภิกฺขุ จีวรวํเส จีวรํ เถยฺยจิตฺโต อวหริ. ตสฺส กุกฺกุจฺจํ อโหสิ ฯเปฯ. อาปตฺติํ ตฺวํ ภิกฺขุ อาปนฺโน ปาราชิกนฺติ. (58)}

\prose{เตน โข ปน สมเยน อญฺญตโร ภิกฺขุ วิหาเร จีวรํ อวหริตฺวา “อิโต นิกฺขมนฺโต ปาราชิโก ภวิสฺสามี”ติ วิหารา น นิกฺขมิ ฯเปฯ ภควโต เอตมตฺถํ อาโรเจสุํ. นิกฺขมิ\csromansharedfootnote{a6e671119c87b461}{นิกฺขเมยฺย (สี, สฺยา)} วา โส ภิกฺขเว โมฆปุริโส น วา นิกฺขมิ\csromansharedfootnote{a6e671119c87b461}{นิกฺขเมยฺย (สี, สฺยา)}, อาปตฺติ ปาราชิกสฺสาติ. (59)}

\proseitem{146}{เตน โข ปน สมเยน เทฺว ภิกฺขู สหายกา โหนฺติ. เอโก ภิกฺขุ คามํ ปิณฺฑาย ปาวิสิ. ทุติโย ภิกฺขุ สํฆสฺส ขาทนีเย ภาชียมาเน สหายกสฺส ภาคํ คเหตฺวา ตสฺส วิสฺสสนฺโต ปริภุญฺชิ. โส ชานิตฺวา ตํ โจเทสิ “อสฺสมโณสิ ตฺวนฺ”ติ. ตสฺส กุกฺกุจฺจํ อโหสิ ฯเปฯ. กิํ จิตฺโต ตฺวํ ภิกฺขูติ. วิสฺสาสคฺคาโห อหํ ภควาติ. อนาปตฺติ ภิกฺขุ วิสฺสาสคฺคาเหติ. (60)}

\proseitem{147}{เตน โข ปน สมเยน สมฺพหุลา ภิกฺขู จีวรกมฺมํ กโรนฺติ. สํฆสฺส ขาทนีเย ภาชียมาเน สพฺเพสํ ปฏิวิสา อาหริตฺวา อุปนิกฺขิตฺตา โหนฺติ. อญฺญตโร ภิกฺขุ อญฺญตรสฺส ภิกฺขุโน ปฏิวิสํ อตฺตโน มญฺญมาโน ปริภุญฺชิ. โส ชานิตฺวา ตํ โจเทสิ “อสฺสมโณสิ ตฺวนฺ”ติ. ตสฺส กุกฺกุจฺจํ อโหสิ ฯเปฯ. กิํ จิตฺโต ตฺวํ ภิกฺขูติ. สกสญฺญี อหํ ภควาติ. อนาปตฺติ ภิกฺขุ สกสญฺญิสฺสาติ. (61)}

\prose{เตน โข ปน สมเยน สมฺพหุลา ภิกฺขู จีวรกมฺมํ กโรนฺติ. สํฆสฺส ขาทนีเย ภาชียมาเน อญฺญตรสฺส ภิกฺขุโน ปตฺเตน อญฺญตรสฺส \csromanfolio{76}ภิกฺขุโน ปฏิวิโส อาหริตฺวา อุปนิกฺขิตฺโต โหติ. ปตฺตสามิโก ภิกฺขุ อตฺตโน มญฺญมาโน ปริภุญฺชิ. โส ชานิตฺวา ตํ โจเทสิ “อสฺสมโณสิ ตฺวนฺ”ติ. ตสฺส กุกฺกุจฺจํ อโหสิ ฯเปฯ. อนาปตฺติ ภิกฺขุ สกสญฺญิสฺสาติ. (62)}

\proseitem{148}{เตน โข ปน สมเยน อมฺพโจรกา อมฺพํ ปาเตตฺวา ภณฺฑิกํ อาทาย อคมํสุ. สามิกา เต โจรเก อนุพนฺธิํสุ. โจรกา สามิเก ปสฺสิตฺวา ภณฺฑิกํ ปาเตตฺวา ปลายิํสุ. ภิกฺขู ปํสุกูล\-สญฺญิโน ปฏิคฺคหาเปตฺวา ปริภุญฺชิํสุ. สามิกา เต ภิกฺขู โจเทสุํ “อสฺสมณาตฺถ ตุเมฺห”ติ. เตสํ กุกฺกุจฺจํ อโหสิ ฯเปฯ. ภควโต เอตมตฺถํ อาโรเจสุํ. กิํจิตฺตา ตุเมฺห ภิกฺขเวติ. ปํสุกูล\-สญฺญิโน มยํ ภควาติ. อนาปตฺติ ภิกฺขเว ปํสุกูลสญฺญิสฺสาติ. (63)}

\prose{เตน โข ปน สมเยน ชมฺพุโจรกา. ลพุชโจรกา. ปนส โจรกา\csromansharedfootnote{f4d4b483bf7c95a1}{ปนสโจรกา สมาโส \csromandash{} ม.พ.ป.}. ตาลปกฺก\-โจรกา. อุจฺฉุโจรกา. ติมฺพรูสก\-โจรกา ติมฺพรูสเก อุจฺจินิตฺวา ภณฺฑิกํ อาทาย อคมํสุ. สามิกา เต โจรเก อนุพนฺธิํสุ. โจรกา สามิเก ปสฺสิตฺวา ภณฺฑิกํ ปาเตตฺวา ปลายิํสุ. ภิกฺขู ปํสุกูล\-สญฺญิโน ปฏิคฺคหาเปตฺวา ปริภุญฺชิํสุ. สามิกา เต ภิกฺขู โจเทสุํ “อสฺสมณาตฺถ ตุเมฺห”ติ. เตสํ กุกฺกุจฺจํ อโหสิ ฯเปฯ. อนาปตฺติ ภิกฺขเว ปํสุกูลสญฺญิสฺสาติ. \mbox{(64-69)}}

\prose{เตน โข ปน สมเยน อมฺพโจรกา อมฺพํ ปาเตตฺวา ภณฺฑิกํ อาทาย อคมํสุ. สามิกา เต โจรเก อนุพนฺธิํสุ. โจรกา สามิเก ปสฺสิตฺวา ภณฺฑิกํ ปาเตตฺวา ปลายิํสุ. ภิกฺขู “ปุเร สามิกา ปสฺสนฺตี”ติ เถยฺยจิตฺตา ปริภุญฺชิํสุ. สามิกา เต ภิกฺขู โจเทสุํ “อสฺสมณาตฺถ ตุเมฺห”ติ. เตสํ กุกฺกุจฺจํ อโหสิ ฯเปฯ. อาปตฺติํ ตุเมฺห ภิกฺขเว อาปนฺนา ปาราชิกนฺติ. (70)}

\prose{เตน โข ปน สมเยน ชมฺพุโจรกา. ลพุชโจรกา. ปนสโจรกา. ตาลปกฺก\-โจรกา. อุจฺฉุโจรกา. ติมฺพรูสก\-โจรกา ติมฺพรูสเก อุจฺจินิตฺวา ภณฺฑิกํ อาทาย อคมํสุ. สามิกา เต โจรเก อนุพนฺธิํสุ. โจรกา สามิเก ปสฺสิตฺวา ภณฺฑิกํ \csromanfolio{77}ปาเตตฺวา ปลายิํสุ. ภิกฺขู “ปุเร สามิกา ปสฺสนฺตี”ติ เตยฺยจิตฺตา ปริภุญฺชิํสุ. สามิกา เต ภิกฺขู โจเทสุํ “อสฺสมณาตฺถ ตุเมฺห”ติ. เตสํ กุกฺกุจฺจํ อโหสิ ฯเปฯ. อาปตฺติํ ตุเมฺห ภิกฺขเว อาปนฺนา ปาราชิกนฺติ. \mbox{(71-76)}}

\prose{เตน โข ปน สมเยน อญฺญตโร ภิกฺขุ สํฆสฺส อมฺพํ เถยฺยจิตฺโต อวหริ. สํฆสฺส ชมฺพุํ. สํฆสฺส ลพุชํ. สํฆสฺส ปนสํ. สํฆสฺส ตาลปกฺกํ. สํฆสฺส อุจฺฉุํ. สํฆสฺส ติมฺพรูสกํ เถยฺยจิตฺโต อวหริ. ตสฺส กุกฺกุจฺจํ อโหสิ ฯเปฯ. อาปตฺติํ ตฺวํ ภิกฺขุ อาปนฺโน ปาราชิกนฺติ. \mbox{(77-83)}}

\proseitem{149}{เตน โข ปน สมเยน อญฺญตโร ภิกฺขุ ปุปฺผารามํ คนฺตฺวา โอจิตํ ปุปฺผํ ปญฺจมาสคฺฆนกํ เถยฺยจิตฺโต อวหริ. ตสฺส กุกฺกุจฺจํ อโหสิ ฯเปฯ. อาปตฺติํ ตฺวํ ภิกฺขุ อาปนฺโน ปาราชิกนฺติ. (84)}

\prose{เตน โข ปน สมเยน อญฺญตโร ภิกฺขุ ปุปฺผารามํ คนฺตฺวา ปุปฺผํ โอจินิตฺวา ปญฺจมาสคฺฆนกํ เถยฺยจิตฺโต อวหริ. ตสฺส กุกฺกุจฺจํ อโหสิ ฯเปฯ. อาปตฺติํ ตฺวํ ภิกฺขุ อาปนฺโน ปาราชิกนฺติ. (85)}

\proseitem{150}{เตน โข ปน สมเยน อญฺญตโร ภิกฺขุ คามกํ คจฺฉนฺโต อญฺญตรํ ภิกฺขุํ เอตทโวจ “อาวุโส ตุยฺหํ อุปฏฺฐากกุลํ วุตฺโต วชฺเชมี”ติ. โส คนฺตฺวา เอกํ สาฏกํ อาหราเปตฺวา อตฺตนา ปริภุญฺชิ. โส ชานิตฺวา ตํ โจเทสิ “อสฺสมโณสิ ตฺวนฺ”ติ. ตสฺส กุกฺกุจฺจํ อโหสิ ฯเปฯ. อนาปตฺติ ภิกฺขุ ปาราชิกสฺส, น จ ภิกฺขเว “วุตฺโต วชฺเชมี”ติ วตฺตพฺโพ, โย วเทยฺย, อาปตฺติ ทุกฺกฏสฺสาติ. (86)}

\prose{เตน โข ปน สมเยน อญฺญตโร ภิกฺขุ คามกํ คจฺฉติ. อญฺญตโร ภิกฺขุ ตํ ภิกฺขุํ เอตทโวจ “อาวุโส มยฺหํ อุปฏฺฐากกุลํ วุตฺโต วชฺเชหี”ติ. โส คนฺตฺวา ยุคสาฏกํ อาหราเปตฺวา เอกํ อตฺตนา ปริภุญฺชิ. เอกํ ตสฺส ภิกฺขุโน อทาสิ. โส ชานิตฺวา ตํ โจเทสิ “อสฺสมโณสิ ตฺวนฺ”ติ. ตสฺส กุกฺกุจฺจํ อโหสิ ฯเปฯ. อนาปตฺติ ภิกฺขุ ปาราชิกสฺส, น จ ภิกฺขเว “วุตฺโต วชฺเชหี”ติ วตฺตพฺโพ, โย วเทยฺย, อาปตฺติ ทุกฺกฏสฺสาติ. (87)}

\prose{\csromanfolio{78}เตน โข ปน สมเยน อญฺญตโร ภิกฺขุ คามกํ คจฺฉนฺโต อญฺญตรํ ภิกฺขุํ เอตทโวจ “อาวุโส ตุยฺหํ อุปฏฺฐากกุลํ วุตฺโต วชฺเชมี”ติ. โสปิ เอวมาห “วุตฺโต วชฺเชหี”ติ. โส คนฺตฺวา อาฬฺหกํ สปฺปิํ ตุลํ คุฬํ โทณํ ตณฺฑุลํ อาหราเปตฺวา อตฺตนา ปริภุญฺชิ. โส ชานิตฺวา ตํ โจเทสิ “อสฺสมโณสิ ตฺวนฺ”ติ. ตสฺส กุกฺกุจฺจํ อโหสิ ฯเปฯ. อนาปตฺติ ภิกฺขุ ปาราชิกสฺส, น จ ภิกฺขเว “วุตฺโต วชฺเชมี”ติ วตฺตพฺโพ, น จ “วุตฺโต วชฺเชหี”ติ วตฺตพฺโพ, โย วเทยฺย, อาปตฺติ ทุกฺกฏสฺสาติ. (88)}

\proseitem{151}{เตน โข ปน สมเยน อญฺญตโร ปุริโส มหคฺฆํ มณิํ อาทาย อญฺญตเรน ภิกฺขุนา สทฺธิํ อทฺธาน\-มคฺค\-ปฺปฏิปนฺโน โหติ. อถ โข โส ปุริโส สุงฺกฏฺฐานํ ปสฺสิตฺวา ตสฺส ภิกฺขุโน อชานนฺตสฺส ถวิกาย มณิํ ปกฺขิปิตฺวา สุงฺกฏฺฐานํ อติกฺกมิตฺวา อคฺคเหสิ. ตสฺส กุกฺกุจฺจํ อโหสิ ฯเปฯ. กิํจิตฺโต ตฺวํ ภิกฺขูติ. นาหํ ภควา ชานามีติ. อนาปตฺติ ภิกฺขุ อชานนฺตสฺสาติ. (89)}

\prose{เตน โข ปน สมเยน อญฺญตโร ปุริโส มหคฺฆํ มณิํ อาทาย อญฺญตเรน ภิกฺขุนา สทฺธิํ อทฺธาน\-มคฺค\-ปฺปฏิปนฺโน โหติ. อถ โข โส ปุริโส สุงฺกฏฺฐานํ ปสฺสิตฺวา คิลานาลยํ กริตฺวา อตฺตโน ภณฺฑิกํ ตสฺส ภิกฺขุโน อทาสิ. อถ โข โส ปุริโส สุงฺกฏฺฐานํ อติกฺกมิตฺวา ตํ ภิกฺขุํ เอตทโวจ “อาหร เม ภนฺเต ภณฺฑิกํ นาหํ อกลฺลโก”ติ. กิสฺส ปน ตฺวํ อาวุโส เอวรูปํ อกาสีติ. อถ โข โส ปุริโส ตสฺส ภิกฺขุโน เอตมตฺถํ อาโรเจสิ. ตสฺส กุกฺกุจฺจํ อโหสิ ฯเปฯ. กิํจิตฺโต ตฺวํ ภิกฺขูติ. นาหํ ภควา ชานามีติ. อนาปตฺติ ภิกฺขุ อชานนฺตสฺสาติ. (90)}

\proseitem{152}{เตน โข ปน สมเยน อญฺญตโร ภิกฺขุ สตฺเถน สทฺธิํ อทฺธาน\-มคฺค\-ปฺปฏิปนฺโน โหติ. อญฺญตโร ปุริโส ตํ ภิกฺขุํ อามิเสน อุปลาเปตฺวา สุงฺกฏฺฐานํ ปสฺสิตฺวา มหคฺฆํ มณิํ ตสฺส ภิกฺขุโน อทาสิ “อิมํ ภนฺเต มณิํ สุงฺกฏฺฐานํ อติกฺกาเมหี”ติ. อถ โข โส ภิกฺขุ ตํ มณิํ สุงฺกฏฺฐานํ อติกฺกาเมสิ. ตสฺส กุกฺกุจฺจํ อโหสิ ฯเปฯ. อาปตฺติํ ตฺวํ ภิกฺขุ อาปนฺโน ปาราชิกนฺติ. (91)}

\proseitem{153}{\csromanfolio{79}เตน โข ปน สมเยน อญฺญตโร ภิกฺขุ ปาเส พนฺธํ สูกรํ การุญฺเญน มุญฺจิ. ตสฺส กุกฺกุจฺจํ อโหสิ ฯเปฯ. กิํจิตฺโต ตฺวํ ภิกฺขูติ. การุญฺญา\-ธิปฺปาโย อหํ ภควาติ. อนาปตฺติ ภิกฺขุ การุญฺญาธิปฺปายสฺสาติ. (92)}

\prose{เตน โข ปน สมเยน อญฺญตโร ภิกฺขุ ปาเส พนฺธํ สูกรํ “ปุเร สามิกา ปสฺสนฺตี”ติ เถยฺยจิตฺโต มุญฺจิ. ตสฺส กุกฺกุจฺจํ อโหสิ ฯเปฯ. อาปตฺติํ ตฺวํ ภิกฺขุ อาปนฺโน ปาราชิกนฺติ. (93)}

\prose{เตน โข ปน สมเยน อญฺญตโร ภิกฺขุ ปาเส พนฺธํ มิคํ การุญฺเญน มุญฺจิ. ปาเส พนฺธํ มิคํ “ปุเร สามิกา ปสฺสนฺตี”ติ เถยฺยจิตฺโต มุญฺจิ. กุมิเน พนฺเธ มจฺเฉ การุญฺเญน มุญฺจิ. กุมิเน พนฺเธ มจฺเฉ “ปุเร สามิกา ปสฺสนฺตี”ติ เถยฺยจิตฺโต มุญฺจิ. ตสฺส กุกฺกุจฺจํ อโหสิ ฯเปฯ. อาปตฺติํ ตฺวํ ภิกฺขุ อาปนฺโน ปาราชิกนฺติ. \mbox{(94-97)}}

\prose{เตน โข ปน สมเยน อญฺญตโร ภิกฺขุ ยาเน ภณฺฑํ ปสฺสิตฺวา “อิโต คณฺหนฺโต ปาราชิโก ภวิสฺสามี”ติ อติกฺกมิตฺวา ปวฏฺเฏตฺวา\csromansharedfootnote{99660804f4972140}{ปวตฺเตตฺวา (ก)} อคฺคเหสิ. ตสฺส กุกฺกุจฺจํ อโหสิ ฯเปฯ. อาปตฺติํ ตฺวํ ภิกฺขุ อาปนฺโน ปาราชิกนฺติ. (98)}

\prose{เตน โข ปน สมเยน อญฺญตโร ภิกฺขุ กุลเลน อุกฺขิตฺตํ มํสเปสิํ “สามิกานํ ทสฺสามี”ติ อคฺคเหสิ. สามิกา ตํ ภิกฺขุํ โจเทสุํ “อสฺสมโณสิ ตฺวนฺ”ติ. ตสฺส กุกฺกุจฺจํ อโหสิ ฯเปฯ. อนาปตฺติ ภิกฺขุ อเถยฺยจิตฺตสฺสาติ. (99)}

\prose{เตน โข ปน สมเยน อญฺญตโร ภิกฺขุ กุลเลน อุกฺขิตฺตํ มํสเปสิํ “ปุเร สามิกา ปสฺสนฺตี”ติ เถยฺยจิตฺโต อคฺคเหสิ. สามิกา ตํ ภิกฺขุํ โจเทสุํ “อสฺสมโณสิ ตฺวนฺ”ติ. ตสฺส กุกฺกุจฺจํ อโหสิ ฯเปฯ. อาปตฺติํ ตฺวํ ภิกฺขุ อาปนฺโน ปาราชิกนฺติ. (100)}

\proseitem{154}{เตน โข ปน สมเยน มนุสฺสา อุฬุมฺปํ พนฺธิตฺวา อจิรวติยา นทิยา โอสาเรนฺติ. พนฺธเน ฉินฺเน กฏฺฐานิ วิปฺปกิณฺณานิ อคมํสุ. ภิกฺขู ปํสุกูล\-สญฺญิโน อุตฺตาเรสุํ. สามิกา เต ภิกฺขู โจเทสุํ \csromanfolio{80}“อสฺสมณาตฺถ ตุเมฺห”ติ. เตสํ กุกฺกุจฺจํ อโหสิ ฯเปฯ. อนาปตฺติ ภิกฺขเว ปํสุกูลสญฺญิสฺสาติ. (101)}

\prose{เตน โข ปน สมเยน มนุสฺสา อุฬุมฺปํ พนฺธิตฺวา อจิรวติยา นทิยา โอสาเรนฺติ. พนฺธเน ฉินฺเน กฏฺฐานิ วิปฺปกิณฺณานิ อคมํสุ. ภิกฺขู “ปุเร สามิกา ปสฺสนฺตี”ติ เถยฺยจิตฺตา อุตฺตาเรสุํ. สามิกา เต ภิกฺขู โจเทสุํ “อสฺสมณาตฺถ ตุเมฺห”ติ. เตสํ กุกฺกุจฺจํ อโหสิ ฯเปฯ. อาปตฺติํ ตุเมฺห ภิกฺขเว อาปนฺนา ปาราชิกนฺติ. (102)}

\prose{เตน โข ปน สมเยน อญฺญตโร โคปาลโก รุกฺเข สาฏกํ อาลคฺเคตฺวา อุจฺจารํ อคมาสิ. อญฺญรโต ภิกฺขุ ปํสุกูลสญฺญี อคฺคเหสิ. อถ โข โส โคปาลโก ตํ ภิกฺขุํ โจเทสิ “อสฺสมโณสิ ตฺวนฺ”ติ. ตสฺส กุกฺกุจฺจํ อโหสิ ฯเปฯ. อนาปตฺติ ภิกฺขุ ปํ สุกูลสญฺญิสฺสาติ. (103)}

\prose{เตน โข ปน สมเยน อญฺญตโร ภิกฺขุโน นทิํ ตรนฺตสฺส รชกานํ หตฺถโต มุตฺตํ สาฏกํ ปาเท ลคฺคํ โหติ. โส ภิกฺขุ สามิกานํ ทสฺสามีติ อคฺคเหสิ. สามิกา ตํ ภิกฺขุํ โจเทสุํ “อสฺสมโณสิ ตฺวนฺ”ติ. ตสฺส กุกฺกุจฺจํ อโหสิ ฯเปฯ. อนาปตฺติ ภิกฺขุ อเถยฺยจิตฺตสฺสาติ. (104)}

\prose{เตน โข ปน สมเยน อญฺญตรสฺส ภิกฺขุโน นทิํ ตรนฺตสฺส รชกานํ หตฺถโต มุตฺตํ สาฏกํ ปาเท ลคฺคํ โหติ. โส ภิกฺขุ “ปุเร สามิกา ปสฺสนฺตี”ติ เถยฺยจิตฺโต อคฺคเหสิ. สามิกา ตํ ภิกฺขุํ โจเทสุํ “อสฺสมโณสิ ตฺวนฺ”ติ. ตสฺส กุกฺกุจฺจํ อโหสิ ฯเปฯ. อาปตฺติํ ตฺวํ ภิกฺขุ อาปนฺโน ปาราชิกนฺติ. (105)}

\proseitem{155}{เตน โข ปน สมเยน อญฺญตโร ภิกฺขุ สปฺปิกุมฺภิํ ปสฺสิตฺวา โถกํ โถกํ ปริภุญฺชิ, ตสฺส กุกฺกุจฺจํ อโหสิ ฯเปฯ. อนาปตฺติ ภิกฺขุ ปาราชิกสฺส, อาปตฺติ ทุกฺกฏสฺสาติ. (106)}

\prose{เตน โข ปน สมเยน สมฺพหุลา ภิกฺขู สํวิทหิตฺวา อคมํสุ “ภณฺฑํ อวหริสฺสามา”ติ. เอโก ภณฺฑํ อวหริ. เต เอวมาหํสุ \csromanfolio{81}“น มยํ ปาราชิกา, โย อวหโฏ, โส ปาราชิโก”ติ. ภควโต เอตมตฺตํ อาโรเจสุํ. อาปตฺติํ ตุเมฺห ภิกฺขเว อาปนฺนา ปาราชิกนฺติ. (107)}

\prose{เตน โข ปน สมเยน สมฺพหุลา ภิกฺขู สํวิทหิตฺวา ภณฺฑํ อวหริตฺวา ภาเชสุํ. เตหิ ภาชียมาเน เอกเมกสฺส ปฏิวิโส น ปญฺจมาสโก ปูริ. เต เอว มาหํสุ “น มยํ ปาราชิกา”ติ. ภควโต เอตมตฺถํ อาโรเจสุํ. อาปตฺติํ ตุเมฺห ภิกฺขเว อาปนฺนา ปาราชิกนฺติ. (108)}

\prose{เตน โข ปน สมเยน อญฺญตโร ภิกฺขุ สาวตฺถิยํ ทุพฺภิกฺเข อาปณิกสฺส ตณฺฑุลมุฏฺฐิํ เถยฺยจิตฺโต อวหริ. ตสฺส กุกฺกุจฺจํ อโหสิ ฯเปฯ. อาปตฺติํ ตฺวํ ภิกฺขุ อาปนฺโน ปาราชิกนฺติ. (109)}

\prose{เตน โข ปน สมเยน อญฺญตโร ภิกฺขุ สาวตฺถิยํ ทุพฺภิกฺเข อาปณิกสฺส มุคฺคมุฏฺฐิํ. มาสมุฏฺฐิํ. ติลมุฏฺฐิํ เถยฺยจิตฺโต อวหริ. ตสฺส กุกฺกุจฺจํ อโหสิ ฯเปฯ. อาปตฺติํ ตฺวํ ภิกฺขุ อาปนฺโน ปาราชิกนฺติ. \mbox{(110-112)}}

\prose{เตน โข ปน สมเยน สาวตฺถิยํ อนฺธวเน โจรกา คาวิํ หนฺตฺวา มํสํ ขาทิตฺวา เสสกํ ปฏิสาเมตฺวา อคมํสุ. ภิกฺขู ปํสุกูล\-สญฺญิโน ปฏิคฺคหาเปตฺวา ปริภุญฺชิํสุ. โจรกา เต ภิกฺขู โจเทสุํ “อสฺสมณาตฺถ ตุเมฺห”ติ. เตสํ กุกฺกุจฺจํ อโหสิ ฯเปฯ. อนาปตฺติ ภิกฺขเว ปํสุกูลสญฺญิสฺสาติ. (113)}

\prose{เตน โข ปน สมเยน สาวตฺถิยํ อนฺธวเน โจรกา สูกรํ หนฺตฺวา มํสํ ขาทิตฺวา เสสกํ ปฏิสาเมตฺวา อคมํสุ. ภิกฺขู ปํสุกูล\-สญฺญิโน ปฏิคฺคหาเปตฺวา ปริภุญฺชิํสุ. โจรกา เต ภิกฺขู โจเทสุํ “อสฺสมณาตฺถ ตุเมฺห”ติ. เตสํ กุกฺกุจฺจํ อโหสิ ฯเปฯ. อนาปตฺติ ภิกฺขเว ปํสุกูลสญฺญิสฺสาติ. (114)}

\prose{เตน โข ปน สมเยน อญฺญตโร ภิกฺขุ ติณกฺเขตฺตํ คนฺตฺวา ลูตํ ติณํ ปญฺจมาสคฺฆนกํ เถยฺยจิตฺโต อวหริ. ตสฺส กุกฺกุจฺจํ อโหสิ ฯเปฯ. อาปตฺติํ ตฺวํ ภิกฺขุ อาปนฺโน ปาราชิกนฺติ. (115)}

\prose{\csromanfolio{82}เตน โข ปน สมเยน อญฺญตโร ภิกฺขุ ติณกฺเขตฺตํ คนฺตฺวา ติณํ ลายิตฺวา ปญฺจมาสคฺฆนกํ เถยฺยจิตฺโต อวหริ. ตสฺส กุกฺกุจฺจํ อโหสิ ฯเปฯ. อาปตฺติํ ตฺวํ ภิกฺขุ อาปนฺโน ปาราชิกนฺติ. (116)}

\proseitem{156}{เตน โข ปน สมเยน อาคนฺตุกา ภิกฺขู สํฆสฺส อมฺพํ ภาชาเปตฺวา ปริภุญฺชิํสุ. อาวาสิกา ภิกฺขู เต ภิกฺขู โจเทสุํ “อสฺสมณาตฺถ ตุเมฺห”ติ. เตสํ กุกฺกุจฺจํ อโหสิ ฯเปฯ. ภควโต เอตมตฺถํ อาโรเจสุํ. กิํจิตฺตา ตุเมฺห ภิกฺขเวติ. ปริโภค\-ตฺถาย\csromansharedfootnote{8d65449cf3a100a9}{ปริโภคตฺถา (สี)} มยํ ภควาติ. อนาปตฺติ ภิกฺขเว ปริโภค\-ตฺถา\-ยาติ. (117)}

\prose{เตน โข ปน สมเยน อาคนฺตุกา ภิกฺขู สํฆสฺส ชมฺพุํ. สํฆสฺส ลพุชํ. สํฆสฺส ปนสํ. สํฆสฺส ตาลปกฺกํ. สํฆสฺส อุจฺฉุํ. สํฆสฺส ติมฺพรูสกํ ภาชาเปตฺวา ปริภุญฺชิํสุ. อาวาสิกา ภิกฺขู เต ภิกฺขู โจเทสุํ “อสฺสมณาตฺถ ตุเมฺห”ติ. เตสํ กุกฺกุจฺจํ อโหสิ ฯเปฯ. อนาปตฺติ ภิกฺขเว ปริโภค\-ตฺถา\-ยาติ. \mbox{(118-123)}}

\prose{เตน โข ปน สมเยน อมฺพปาลกา ภิกฺขูนํ อมฺพผลํ เทนฺติ. ภิกฺขู “โคเปตุํ อิเม อิสฺสรา, น ยิเม ทาตุนฺ”ติ กุกฺกุจฺจายนฺตา น ปฏิคฺคณฺหนฺติ. ภควโต เอตมตฺถํ อาโรเจสุํ. อนาปตฺติ ภิกฺขเว โคปกสฺส ทาเนติ. (124)}

\prose{เตน โข ปน สมเยน ชมฺพุปาลกา. ลพุชปาลกา. ปนสปาลกา. ตาลปกฺก\-ปาลกา. อุจฺฉุปาลกา. ติมฺพรูสก\-ปาลกา ภิกฺขูนํ ติมฺพรูสกํ เทนฺติ. ภิกฺขู “โคเปตุํ อิเม อิสฺสรา, น ยิเม ทาตุนฺ”ติ กุกฺกุจฺจายนฺตา น ปฏิคฺคณฺหนฺติ. ภควโต เอตมตฺถํ อาโรเจสุํ. อนาปตฺติ ภิกฺขเว โคปกสฺส ทาเนติ. \mbox{(125-130)}}

\prose{เตน โข ปน สมเยน อญฺญตโร ภิกฺขุ สํฆสฺส ทารุํ ตาวกาลิกํ หริตฺวา อตฺตโน วิหารสฺส กุฏฺฏํอุปตฺถมฺเภสิ. ภิกฺขู ตํ ภิกฺขุํ โจเทสุํ “อสฺสมโณสิ ตฺวนฺ”ติ. ตสฺส กุกฺกุจฺจํ อโหสิ ฯเปฯ. ภควโต เอตมตฺถํ อาโรเจสิ. กิํจิตฺโต ตฺวํ ภิกฺขูติ. ตาวกาลิโก อหํ ภควาติ. อนาปตฺติ ภิกฺขุ ตาวกาลิเกติ. (131)}

\prose{\csromanfolio{83}เตน โข ปน สมเยน อญฺญตโร ภิกฺขุ สํฆสฺส อุทกํ เถยฺยจิตฺโต อวหริ. สํฆสฺส มตฺติกํ เถยฺยจิตฺโต อวหริ. สํฆสฺส ปุญฺชกิตํ ติณํ เถยฺยจิตฺโต อวหริ. ตสฺส กุกฺกุจฺจํ อโหสิ ฯเปฯ. อาปตฺติํ ตฺวํ ภิกฺขุ อาปนฺโน ปาราชิกนฺติ. \mbox{(132-134)}}

\prose{เตน โข ปน สมเยน อญฺญตโร ภิกฺขุ สํฆสฺส ปุญฺชกิตํ ติณํ เถยฺยจิตฺโต ฌาเปสิ. ตสฺส กุกฺกุจฺจํ อโหสิ ฯเปฯ. อนาปตฺติ ภิกฺขุ ปาราชิกสฺส, อาปตฺติ ทุกฺกฏสฺสาติ. (135)}

\prose{เตน โข ปน สมเยน อญฺญตโร ภิกฺขุ สํฆสฺส มญฺจํ เถยฺยจิตฺโต อวหริ. ตสฺส กุกฺกุจฺจํ อโหสิ ฯเปฯ. อาปตฺติํ ตฺวํ ภิกฺขุ อาปนฺโน ปาราชิกนฺติ. (136)}

\prose{เตน โข ปน สมเยน อญฺญตโร ภิกฺขุ สํฆสฺส ปีฐํ. สํฆสฺส ภิสิํ. สํฆสฺส พิพฺโพหนํ\csromansharedfootnote{cb8b705f1f808b97}{พิมฺโพหนํ (สี, สฺยา)}. สํฆสฺส กวาฏํ. สํฆสฺส อาโลกสนฺธิํ. สํฆสฺส โคปานสิํ เถยฺยจิตฺโต อวหริ. ตสฺส กุกฺกุจฺจํ อโหสิ ฯเปฯ. อาปตฺติํ ตฺวํ ภิกฺขุ อาปนฺโน ปาราชิกนฺติ. \mbox{(137-142)}}

\proseitem{157}{\csromansymbolfootnote{*}{วิ 4. 327 ปิฏฺเฐปิ.}เตน โข ปน สมเยน ภิกฺขู อญฺญตรสฺส อุปาสกสฺส วิหาร\-ปริโภคํ เสนาสนํ อญฺญตฺร ปริภุญฺชนฺติ. อถ โข โส อุปาสโก อุชฺฌายติ ขิยฺยติ วิปาเจติ “กถํ หิ นาม ภทนฺตา อญฺญตฺร ปริโภคํ อญฺญตฺร ปริภุญฺชิสฺสนฺตี”ติ. ภควโต เอตมตฺถํ อาโรเจสุํ. น ภิกฺขเว อญฺญตฺร ปริโภโค อญฺญตฺร ปริภุญฺชิ\-ตพฺโพ, โย ปริภุญฺเชยฺย, อาปตฺติ ทุกฺกฏสฺสาติ. (143)}

\prose{\csromansymbolfootnote{+}{วิ 4. 328 ปิฏฺเฐปิ.}เตน โข ปน สมเยน ภิกฺขู อุโปสถคฺคํปิ สนฺนิสชฺชํปิ หริตุํ กุกฺกุจฺจายนฺตา ฉมายํ นิสีทนฺติ, คตฺตานิปิ จีวรานิปิ ปํสุกิตานิ โหนฺติ. ภควโต เอตมตฺถํ อาโรเจสุํ. อนุชานามิ ภิกฺขเว ตาวกาลิกํ หริตุนฺติ. (144)}

\prose{เตน โข ปน สมเยน จมฺปายํ ถุลฺลนนฺทาย ภิกฺขุนิยา อนฺเตวาสินี ภิกฺขุนี ถุลฺลนนฺทาย ภิกฺขุนิยา อุปฏฺฐากกุลํ คนฺตฺวา “อยฺยา \csromanfolio{84}อิจฺฉติ เตกฏุลยาคุํ ปาตุนฺ”ติ ปจาเปตฺวา หริตฺวา อตฺตนา ปริภุญฺชิ. สา ชานิตฺวา ตํ โจเทสิ “อสฺสมณีสิ ตฺวนฺ”ติ. ตสฺสา กุกฺกุจฺจํ อโหสิ ฯเปฯ. อถ โข สา ภิกฺขุนี ภิกฺขุนีนํ เอตมตฺถํ อาโรเจสิ. ภิกฺขุนิโย ภิกฺขูนํ เอตมตฺถํ อาโรเจสุํ. ภิกฺขู ภควโต เอตมตฺถํ อาโรเจสุํ. อนาปตฺติ ภิกฺขเว ปาราชิกสฺส, อาปตฺติ สมฺปชาน\-มุสาวาเท ปาจิตฺติยสฺสาติ. (145)}

\prose{เตน โข ปน สมเยน ราชคเห ถุลฺลนนฺทาย ภิกฺขุนิยา อนฺเตวาสินี ภิกฺขุนี ถุลฺลนนฺทาย ภิกฺขุนิยา อุปฏฺฐากกุลํ คนฺตฺวา “อยฺยา อิจฺฉติ มธุโคฬกํ ขาทิตุนฺ”ติ ปจาเปตฺวา หริตฺวา อตฺตนา ปริภุญฺชิ. สา ชานิตฺวา ตํ โจเทสิ “อสฺสมณีสิ ตฺวนฺ”ติ. ตสฺส กุกฺกุจฺจํ อโหสิ ฯเปฯ. อนาปตฺติ ภิกฺขเว ปาราชิกสฺส, อาปตฺติ สมฺปชาน\-มุสาวาเท ปาจิตฺติยสฺสาติ. (146)}

\proseitem{158}{เตน โข ปน สมเยน เวสาลิยํ อายสฺมโต อชฺชุกสฺส อุปฏฺฐากสฺส คหปติโน เทฺว ทารกา โหนฺติ ปุตฺโต จ ภาคิเนยฺโย จ. อถ โข โส คหปติ อายสฺมนฺตํ อชฺชุกํ เอตทโวจ “อิมํ ภนฺเต โอกาสํ โย อิเมสํ ทฺวินฺนํ ทารกานํ สทฺโธ โหติ ปสนฺโน, ตสฺส อาจิกฺเขยฺยาสี”ติ\csromansharedfootnote{75a630ac8f24b7da}{อาจิกฺเขยฺยาสีติ โส กาลมกาสิ (สฺยา)}. เตน โข ปน สมเยน ตสฺส คหปติโน ภาคิเนยฺโย สทฺโธ โหติ ปสนฺโน. อถ โข อายสฺมา อชฺชุโก ตํ โอกาสํ ตสฺส ทารกสฺส อาจิกฺขิ. โส เตน สาปเตยฺเยน กุฏุมฺพญฺจ สณฺฐเปสิ, ทานญฺจ ปฏฺฐเปสิ. อถ โข ตสฺส คหปติโน ปุตฺโต อายสฺมนฺตํ อานนฺทํ เอตทโวจ “โก นุ โข ภนฺเต อานนฺท ปิตุโน ทายชฺโช ปุตฺโต วา ภาคิเนยฺโย วา”ติ. ปุตฺโต โข อาวุโส ปิตุโน ทายชฺโชติ. อยํ ภนฺเต อยฺโย อชฺชุโก อมฺหากํ สาปเตยฺยํ อมฺหากํ เมถุนกสฺส อาจิกฺขีติ. อสฺสมโณ อาวุโส อายสฺมา อชฺชุโกติ. อถ โข อายสฺมา อชฺชุโก อายสฺมนฺตํ อานนฺทํ เอตทโวจ “เทหิ เม อาวุโส อานนฺท วินิจฺฉยนฺ”ติ. เตน โข ปน สมเยน อายสฺมา อุปาลิ อายสฺมโต อชฺชุกสฺส ปกฺโข โหติ. อถ โข อายสฺมา อุปาลิ อายสฺมนฺตํ อานนฺทํ เอตทโวจ \csromanfolio{85}“โย นุโข อาวุโส อานนฺท สามิเกน ‘อิมํ โอกาสํ อิตฺถนฺนามสฺส อาจิกฺเขยฺยาสี’ติ วุตฺโต ตสฺส อาจิกฺขติ, กิํ โส อาปชฺชตี”ติ. น ภนฺเต กิญฺจิ อาปชฺชติ อนฺตมโส ทุกฺกฏมตฺตํปีติ. อยํ อาวุโส อายสฺมา อชฺชุโก สามิเกน อิมํ โอกาสํ อิตฺถนฺนามสฺส อาจิกฺขาติ วุตฺโต ตสฺส อาจิกฺขติ, อนาปตฺติ อาวุโส อายสฺมโต อชฺชุกสฺสาติ. (147)}

\proseitem{159}{เตน โข ปน สมเยน พาราณสิยํ อายสฺมโต ปิลินฺท\-วจฺฉสฺส อุปฏฺฐากกุลํ โจเรหิ อุปทฺทุตํ โหติ. เทฺว จ ทารกา นีตา โหนฺติ. อถ โข อายสฺมา ปิลินฺทวจฺโฉ เต ทารเก อิทฺธิยา อาเนตฺวา ปาสาเท ฐเปสิ. มนุสฺสา เต ทารเก ปสฺสิตฺวา “อยฺยสฺสายํ ปิลินฺท\-วจฺฉสฺส อิทฺธานุภาโว”ติ อายสฺมนฺเต ปิลินฺทวจฺเฉ อภิปฺปสีทิํสุ. ภิกฺขู อุชฺฌายนฺติ ขิยฺยนฺติ วิปาเจนฺติ “กถํ หิ นาม อายสฺมา ปิลินฺทวจฺโฉ โจเรหิ นีเต ทารเก อาเนสฺสตี”ติ. ภควโต เอตมตฺถํ อาโรเจสุํ. อนาปตฺติ ภิกฺขเว อิทฺธิมสฺส\csromansharedfootnote{9633f322ce591603}{อิทฺธิมโต (สี), อิทฺธิมนฺตสฺส (สฺยา)} อิทฺธิวิสเยติ. (148)}

\proseitem{160}{เตน โข ปน สมเยน เทฺว ภิกฺขู สหายกา โหนฺติ ปณฺฑุโก จ กปิโล จ. เอโก คามเก วิหรติ, เอโก โกสมฺพิยํ. อถ โข ตสฺส ภิกฺขุโน คามกา โกสมฺพิํ คจฺฉนฺตสฺส อนฺตรามคฺเค นทิํ ตรนฺตสฺส สูกริกานํ หตฺถโก มุตฺตา เมทวฏฺฏิ ปาเท ลคฺคา โหติ. โส ภิกฺขุ “สามิกานํ ทสฺสามี”ติ อคฺคเหสิ. สามิกา ตํ ภิกฺขุํ โจเทสุํ “อสฺสมโณสิ ตฺวนฺ”ติ. ตํ อุตฺติณฺณํ โคปาลิกา\csromansharedfootnote{4c0696da6e360242}{อญฺญตรา โคปาลิกา (สี, สฺยา)} ปสฺสิตฺวา เอตทโวจ “เอหิ ภนฺเต เมถุนํ ธมฺมํ ปฏิเสวา”ติ. โส “ปกติยาปาหํ อสฺสมโณ”ติ ตสฺสา เมถุนํ ธมฺมํ ปฏิเสวิตฺวา โกสมฺพิํ คนฺตฺวา ภิกฺขูนํ เอตมตฺถํ อาโรเจสิ. ภิกฺขู ภควโต เอตมตฺตํ อาโรเจสุํ. อนาปตฺติ ภิกฺขเว อทินฺนาทาเน ปาราชิกสฺส, อาปตฺติ เมถุนธมฺม\-สมาโยเค ปาราชิกสฺสติ. (149)}

\csromanmatikaanchor{mtk.22}
\csromanmatikaanchor{mtk.23}
\csromantocmarknopage{section}{3. ตติยปาราชิก}{mtk.22}
\bookmark[dest={mtk.22},level=1]{3. ตติยปาราชิก}
\csromantocmarkref{subsection}{สมฺพหุลภิกฺขุวตฺถุ}{mtk.23}
\bookmark[dest={mtk.23},level=2]{สมฺพหุลภิกฺขุวตฺถุ}
\proseitemwithcloserruled{161}{เตน โข ปน สมเยน สาคลายํ อายสฺมโต ทฬฺหิกสฺส สทฺธิวิหาริโก ภิกฺขุ อนภิรติยา ปีฬิโต อาปณิกสฺส เวฐนํ \csromanfolio{86}อวหริตฺวา อายสฺมนฺตํ ทฬฺหิกํ เอตทโวจ “อสฺสมโณ อหํ ภนฺเต ปิพฺภมิสฺสามี”ติ. กิํตยา อาวุโส กตนฺติ, โส ตมตฺถํ อาโรเจสิ, อาหราเปตฺวา อคฺฆาเปสิ. ตํ อคฺฆาเปนฺตํ น ปญฺจมาสเก อคฺฆติ. “อนาปตฺติ อาวุโส ปาราชิกสฺสา”ติ ธมฺมกถํ อกาสิ. โส ภิกฺขุ อภิรมตีติ\csromansharedfootnote{b1c341e3b16f85c2}{อภิรมีติ (สี, สฺยา)}. (150)}{ทุติย\-ปาราชิกํ สมตฺตํ.}
\csromanheader{1}{3. \textbf{ตติยปาราชิก}}
\proseitem{162}{\csromansymbolfootnote{*}{อิทํ วตฺถุ สํ 3. 278 ปิฏฺเฐปิ อาคตํ.}เตน สมเยน พุทฺโธ ภควา เวสาลิยํ วิหรติ มหาวเน กูฏาคาร\-สาลายํ. เตน โข ปน สมเยน ภควา ภิกฺขูนํ อเนก\-ปริยาเยน อสุภกถํ กเถติ, อสุภาย วณฺณํ ภาสติ, อสุภ\-ภาวนาย วณฺณํ ภาสติ, อาทิสฺส อาทิสฺส อสุภ\-สมาปตฺติยา วณฺณํ ภาสติ. อถ โข ภควา ภิกฺขู อามนฺเตสิ “อิจฺฉามหํ ภิกฺขเว อทฺธมาสํ ปฏิสลฺลียิตุํ, นมฺหิ เกนจิ อุปสงฺกมิตพฺโพ อญฺญตฺร เอเกน ปิณฺฑปาต\-นีหารเกนา”ติ. เอวํ ภนฺเตติ โข เต ภิกฺขู ภควโต ปฏิสฺสุณิตฺวา นาสฺสุธ โกจิ ภควนฺตํ อุปสงฺกมติ อญฺญตฺร เอเกน ปิณฺฑปาต\-นีหารเกน. ภิกฺขู “ภควา โข อเนก\-ปริยาเยน อสุภกถํ กเถติ, อสุภาย วณฺณํ ภาสติ, อสุภ\-ภาวนาย วณฺณํ ภาสติ, อาทิสฺส อาทิสฺส อสุภ\-สมาปตฺติยา วณฺณํ ภาสตี”ติ (เต)\csromansharedfootnote{a0db7c0123cc8398}{( ) (?) เอวมุปริปิ อีทิเสสุ ฐาเนสุ.} อเนกา\-การ\-โวการํ อสุภ\-ภาวนา\-นุโยคม\-นุยุตฺตา วิหรนฺติ. เต สเกน กาเยน อฏฺฏียนฺติ หรายนฺติ ชิคุจฺฉนฺติ, เสยฺยถาปิ นาม อิตฺถี วา ปุริโส วา ทหโร ยุวา มณฺฑนก\-ชาติโก สีสํนฺหาโต อหิกุณเปน วา กุกฺกุร\-กุณเปน วา มนุสฺส\-กุณเปน วา กณฺเฐ อาสตฺเตน อฏฺฏีเยยฺย หราเยยฺย ชิคุจฺเฉยฺย, เอวเมว เต ภิกฺขู สเกน กาเยน อฏฺฏียนฺตา หรายนฺตา ชิคุจฺฉนฺตา อตฺตนาปิ อตฺตานํ ชีวิตา โวโรเปนฺติ, อญฺญมญฺญํปิ ชีวิตา โวโรเปนฺติ, มิคลณฺฑิกํปิ สมณกุตฺตกํ อุปสงฺกมิตฺวา เอวํ วทนฺติ “สาธุ โน อาวุโส ชีวิตา โวโรเปหิ \csromanfolio{87}, อิทํ เต ปตฺตจีวรํ ภวิสฺสตี”ติ. อถ โข มิคลณฺฑิโก สมณกุตฺตโก ปตฺตจีวเรหิ ภโฏ สมฺพหุเล ภิกฺขู ชีวิตา โวโรเปตฺวา โลหิตกํ\csromansharedfootnote{71357fdf22b6a0ab}{โลหิตคตํ (ก)} อสิํ อาทาย เยน วคฺคุมุทา นที เตนุปสงฺกมิ.}

\proseitem{163}{อถ โข มิคลณฺฑิกสฺส สมณ\-กุตฺตกสฺส โลหิตกํ ตํ อสิํ โธวนฺตสฺส อหุเทว กุกฺกุจฺจํ อหุ วิปฺปฏิสาโร “อลาภา วต เม, น วต เม ลาภา, ทุลฺลทฺธํ วต เม, น วต เม สุลทฺธํ, พหุํ วต มยา อปุญฺญํ ปสุตํ, โยหํ ภิกฺขู สีลวนฺเต กลฺยาณธมฺเม ชีวิตา โวโรเปสินฺ”ติ. อถ โข อญฺญตรา มารกายิกา เทวตา อภิชฺชมาเน อุทเก อาคนฺตฺวา มิคลณฺฑิกํ สมณกุตฺตกํ เอตทโวจ “สาธุ สาธุ สปฺปุริส, ลาภา เต สปฺปุริส, สุลทฺธํ เต สปฺปุริส, พหุํ ตยา สปฺปุริส ปุญฺญํ ปสุตํ, ยํ ตฺวํ อติณฺเณ ตาเรสี”ติ. อถ โข มิคลณฺฑิโก สมณกุตฺตโก “ลาภา กิร เม, สุลทฺธํ กิร เม, พหุํ กิร มยา ปุญฺญํ ปสุตํ, อติณฺเณ กิราหํ ตาเรมี”ติ ติณฺหํ อสิํ อาทาย วิหาเรน วิหารํ ปริเวเณน ปริเวณํ อุปสงฺกมิตฺวา เอวํ วเทติ “โก อติณฺโณ, กํ ตาเรมี”ติ. ตตฺถ เย เต ภิกฺขู อวีตราคา, เตสํ ตสฺมิํ สมเย โหติ เยว ภยํ, โหติ ฉมฺภิ ตตฺตํ, โหติ โลมหํโส. เย ปน เต ภิกฺขู วีตราคา, เตสํ ตสฺมิํ สมเย น โหติ ภยํ, น โหติ ฉมฺภิ ตตฺตํ, น โหติ โลมหํโส. อถ โข มิคลณฺฑิโก สมณกุตฺตโก เอกํปิ ภิกฺขุํ เอกาเหน ชีวิตา โวโรเปสิ, เทฺวปิ ภิกฺขู เอกาเหน ชีวิตา โวราเปสิ, ตโยปิ ภิกฺขู เอกาเหน ชีวิตา โวโรเปสิ, จตฺตาโรปิ ภิกฺขู เอกาเหน ชีวิตา โวโรเปสิ, ปญฺจปิ ภิกฺขู เอกาเหน ชีวิตา โวโรเปสิ, ทสปิ ภิกฺขู เอกาเหน ชีวิตา โวโรเปสิ, วีสมฺปิ ภิกฺขู เอกาเหน ชีวิตา โวโรเปสิ, ติํสํปิ ภิกฺขู เอกาเหน ชีวิตา โวโรเปสิ, จตฺตาลีสํปิ ภิกฺขู เอกาเหน ชีวิตา โวโรเปสิ, ปญฺญาสํปิ ภิกฺขู เอกาเหน ชีวิตา โวโรเปสิ, สฏฺฐิํปิ ภิกฺขู เอกาเหน ชีวิตา โวโรเปสิ.}

\proseitem{164}{อถ โข ภควา ตสฺส อทฺธมาสสฺส อจฺจเยน ปฏิสลฺลานา วุฏฺฐิโต อายสฺมนฺตํ อานนฺทํ อามนฺเตสิ “กิํ นุ โข อานนฺท ตนุภูโต วิย ภิกฺขุสํโฆ”ติ. ตถา หิ ปน ภนฺเต ภควา ภิกฺขูนํ อเนก\-ปริยาเยน \csromanfolio{88}อสุภกถํ กเถติ, อสุภาย วณฺณํ ภาสติ, อสุภ\-ภาวนาย วณฺณํ ภาสติ, อาทิสฺส อาทิสฺส อสุภ\-สมาปตฺติยา วณฺณํ ภาสติ. เต จ ภนฺเต ภิกฺขู “ภควา โข อเนก\-ปริยาเยน อสุภกถํ กเถติ, อสุภาย วณฺณํ ภาสติ, อสุภ\-ภาวนาย วณฺณํ ภาสติ, อาทิสฺส อาทิสฺส อสุภ\-สมาปตฺติยา วณฺณํ ภาสตี”ติ เต อเนกา\-การ\-โวการํ อสุภ\-ภาวนา\-นุโยคม\-นุยุตฺตา วิหรนฺติ. เต สเกน กาเยน อฏฺฏียนฺติ หรายนฺติ ชิคุจฺฉนฺติ, เสยฺยถาปิ นาม อิตฺถี วา ปุริโส วา ทหโร ยุวา มณฺฑนก\-ชาติโก สีสํนฺหาโต อหิกุณเปน วา กุกฺกุร\-กุณเปน วา มนุสฺส\-กุณเปน วา กณฺเฐ อาสตฺเตน อฏฺฏีเยยฺย หราเยยฺย ชิคุจฺเฉยฺย, เอวเมว เต ภิกฺขู สเกน กาเยน อฏฺฏียนฺตา หรายนฺตา ชิคุจฺฉนฺตา อตฺตนาปิ อตฺตานํ ชีวิตา โวโรเปนฺติ, อญฺญมญฺญํปิ ชีวิตา โวโรเปนฺติ, มิคลณฺฑิกํปิ สมณกุตฺตกํ อุปสงฺกมิตฺวา เอวํ วทนฺติ “สาธุ โน อาวุโส ชีวิตา โวโรเปหิ, อิทํ เต ปตฺตจีวรํ ภวิสฺสตี”ติ. อถ โข ภนฺเต มิคลณฺฑิโก สมณกุตฺตโก ปตฺตจีวเรหิ ภโฏ เอกํปิ ภิกฺขุํ เอกาเหน ชีวิตา โวโรเปสิ ฯเปฯ สฏฺฐิํปิ ภิกฺขู เอกาเหน ชีวิตา โวโรเปสิ. สาธุ ภนฺเต ภควา อญฺญํ ปริยายํ อาจิกฺขตุ, ยถายํ ภิกฺขุสํโฆ อญฺญาย สณฺฐเหยฺยาติ. เตนหานนฺท ยาวติกา ภิกฺขู เวสาลิํ อุปนิสฺสาย วิหรนฺติ, เต สพฺเพ อุปฏฺฐาน\-สาลายํ สนฺนิปาเตหีติ. เอวํ ภนฺเตติ โข อายสฺมา อานนฺโท ภควโต ปฏิสฺสุณิตฺวา ยาวติกา ภิกฺขู เวสาลิํ อุปนิสฺสาย วิหรนฺติ, เต สพฺเพ อุปฏฺฐาน\-สาลายํ สนฺนิปาเตตฺวา เยน ภควา เตนุปสงฺกมิ, อุปสงฺกมิตฺวา ภควนฺตํ เอตทโวจ “สนฺนิปติโต ภนฺเต ภิกฺขุสํโฆ, ยสฺส ทานิ ภนฺเต ภควา กาลํ มญฺญตี”ติ.}

\proseitem{165}{อถ โข ภควา เยน อุปฏฺฐานสาลา เตนุปสงฺกมิ, อุปสงฺกมิตฺวา ปญฺญตฺเต อาสเน นิสีทิ, นิสชฺช โข ภควา ภิกฺขู อามนฺเตสิ “อยํ ปิ โข ภิกฺขเว อานาปาน\-สฺสติ\-สมาธิ ภาวิโต พหุลีกโต สนฺโต เจว ปณีโต จ อเสจนโก จ สุโข จ วิหาโร, อุปฺปนฺนุปฺปนฺเน จ ปาปเก อกุสเล ธมฺเม ฐานโส อนฺตรธาเปติ วูปสเมติ, เสยฺยถาปิ ภิกฺขเว คิมฺหานํ ปจฺฉิเม มาเส อุหตํ\csromansharedfootnote{536da9ef52ac6009}{อูหตํ (ก)} รโชชลฺลํ, \csromanfolio{89}ตเมนํ มหา อกาลเมโฆ ฐานโส อนฺตรธาเปติ วูปสเมติ, เอวเมว โข ภิกฺขเว อานาปาน\-สฺสติ\-สมาธิ ภาวิโต พหุลีกโต สนฺโต เจว ปณีโต จ อเสจนโก จ สุโข จ วิหาโร, อุปฺปนฺนุปฺปนฺเน จ ปาปเก อกุสเล ธมฺเม ฐานโส อนฺตรธาเปติ วูปสเมติ. กถํ ภาวิโต จ ภิกฺขเว อานาปาน\-สฺสติ\-สมาธิ กถํ พหุลีกโต สนฺโต เจว ปณีโต จ อเสจนโก จ สุโข จ วิหาโร, อุปฺปนฺนุปฺปนฺเน จ ปาปเก อกุสเล ธมฺเม ฐานโส อนฺตรธาเปติ วูปสเมติ. อิธ ภิกฺขเว ภิกฺขุ อรญฺญคโต วา รุกฺขมูลคโต วา สุญฺญาคารคโต วา นิสีทติ ปลฺลงฺกํ อาภุชิตฺวา อุชุํ กายํ ปณิธาย ปริมุขํ สติํ อุปฏฺฐเปตฺวา. โส สโตว อสฺสสติ, สโต ปสฺสสติ. ทีฆํ วา อสฺสสนฺโต ทีฆํ อสฺสสามีติ ปชานาติ, ทีฆํ วา \textbf{ปสฺสสนฺโต ทีฆํ ปสฺสสามีติ ปชานาติ. รสฺสํ วา อสฺสสนฺโต รสฺสํ} \textbf{อสฺสสามีติ ปชานาติ, รสฺสํ วา ปสฺสสนฺโต รสฺสํ ปสฺสสามีติ ปชานาติ.} \textbf{สพฺพกาย}\-\textbf{ปฺปฏิสํเวที อสฺสสิสฺสามีติ สิกฺขติ, สพฺพกาย}\-\textbf{ปฺปฏิสํเวที} \textbf{ปสฺสสิสฺสามีติ สิกฺขติ. ปสฺสมฺภยํ กายสงฺขารํ อสฺสสิสฺสามีติ} สิกฺขติ, ปสฺสมฺภยํ กายสงฺขารํ ปสฺสสิสฺสามีติ สิกฺขติ. ปีติปฺปฏิสํเวที อสฺสสิสฺสามีติ สิกฺขติ, ปีติปฺปฏิสํเวที ปสฺสสิสฺสามีติ สิกฺขติ. สุขปฺปฏิสํเวที อสฺสสิสฺสามีติ สิกฺขติ, สุขปฺปฏิสํเวที ปสฺสสิสฺสามีติ สิกฺขติ. จิตฺตสงฺขาร\-ปฺปฏิสํเวที อสฺสสิสฺสามีติ สิกฺขติ, จิตฺตสงฺขาร\-ปฺปฏิสํเวที ปสฺสสิสฺสามีติ สิกฺขติ. ปสฺสมฺภยํ จิตฺตสงฺขารํ อสฺสสิสฺสามีติ สิกฺขติ, ปสฺสมฺภยํ จิตฺตสงฺขารํ ปสฺสสิสฺสามีติ สิกฺขติ. จิตฺต\-ปฺปฏิสํเวที อสฺสสิสฺสามีติ สิกฺขติ, จิตฺต\-ปฺปฏิสํเวที ปสฺสสิสฺสามีติ สิกฺขติ. อภิปฺปโมทยํ จิตฺตํ ฯเปฯ. สมาทหํ จิตฺตํ ฯเปฯ. วิโมจยํ จิตฺตํ ฯเปฯ. อนิจฺจานุปสฺสี ฯเปฯ. วิราคานุปสฺสี ฯเปฯ. นิโรธานุปสฺสี ฯเปฯ. ปฏินิสฺสคฺคา\-นุปสฺสี อสฺสสิสฺสามีติ สิกฺขติ, ปฏินิสฺสคฺคา\-นุปสฺสี ปสฺสาสิสฺสามีติ สิกฺขติ. เอวํ ภวิโต โข ภิกฺขเว อานาปาน\-สฺสติ\-สมาธิ เอวํ พหุลีกโต สนฺโต เจว ปณีโต จ อเสจนโก จ สุโข จ วิหาโร อุปฺปนฺนุปฺปนฺเน จ ปาปเก อกุสเล ธมฺเม ฐานโส อนฺตรธาเปติ วูปสเมตี”ติ.}

\csromanmatikaanchor{mtk.24}
\csromantocmarkref{subsection}{ปฐมปญฺญตฺติ}{mtk.24}
\bookmark[dest={mtk.24},level=2]{ปฐมปญฺญตฺติ}
\proseitem{166}{อถ โข ภควา เอตสฺมิํ นิทาเน เอตสฺมิํ ปกรเณ ภิกฺขุสํฆํ สนฺนิปาตาเปตฺวา ภิกฺขู ปฏิปุจฺฉิ “สจฺจํ กิร ภิกฺขเว ภิกฺขู อตฺตนาปิ อตฺตานํ ชีวิตา โวราเปนฺติ, อญฺญมญฺญํปิ ชีวิตา โวโรเปนฺติ, \csromanfolio{90}มิคลณฺฑิกํปิ สมณกุตฺตกํ อุปสงฺกมิตฺวา เอวํ วทนฺติ ‘สาธุ โน อาวุโส ชีวิตา โวโรเปหิ, อิทํ เต ปตฺตจีวรํ ภวิสฺสตี”ติ. สจฺจํ ภควาติ. วิครหิ พุทฺโธ ภควา “อนนุจฺฉวิกํ ภิกฺขเว เตสํ ภิกฺขูนํ อนนุโลมิกํ อปฺปติรูปํ อสฺสามณกํ อกปฺปิยํ อกรณียํ, กถํ หิ นาม เต ภิกฺขเว ภิกฺขู อตฺตนาปิ อตฺตานํ ชีวิตา โวโรเปสฺสนฺติ, อญฺญมญฺญํปิ ชีวิตา โวโรเปสฺสนฺติ, มิคลณฺฑิกํปิ สมณกุตฺตกํ อุปสงฺกมิตฺวา เอวํ วกฺขนฺติ ‘สาธุ โน อาวุโส ชีวิตา โวโรเปหิ, อิทํ เต ปตฺตจีวรํ ภวิสฺสตี’ติ, เนตํ ภิกฺขเว อปฺปสนฺนานํ วา ปสาทาย ฯเปฯ. เอวญฺจ ปน ภิกฺขเว อิมํ สิกฺขาปทํ อุทฺทิเสยฺยาถ\csromandash{}}

\proseitem{167}{\textbf{“โย ปน ภิกฺขุ สญฺจิจฺจ มนุสฺส}\-\textbf{วิคฺคหํ ชีวิตา โวโรเปยฺย, สตฺถหารกํ วาสฺส ปริเยเสยฺย, อยมฺปิ ปาราชิโก โหติ อสํวาโส”ติ}.}

\prose{เอวญฺจิทํ ภควตา ภิกฺขูนํ สิกฺขาปทํ ปญฺญตฺตํ โหติ.}

\proseitem{168}{เตน โข ปน สมเยน อญฺญตโร อุปาสโก คิลาโน โหติ, ตสฺส ปชาปติ อภิรูปา โหติ ทสฺสนียา ปาสาทิกา. ฉพฺพคฺคิยา ภิกฺขู ตสฺสา อิตฺถิยา ปฏิพทฺธ\-จิตฺตา โหนฺติ. อถ โข ฉพฺพคฺคิยานํ ภิกฺขูนํ เอตทโหสิ “สเจ โข โส อาวุโส อุปาสโก ชีวิสฺสติ, น มยํ ตํ อิตฺถิํ ลภิสฺสาม, หนฺท มยํ อาวุโส ตสฺส อุปาสกสฺส มรณวณฺณํ สํวณฺเณมา”ติ. อถ โข ฉพฺพคฺคิยา ภิกฺขู เยน โส อุปาสโก เตนุ\-ปสงฺกมิํสุ, อุปสงฺกมิตฺวา ตํ อุปาสกํ เอตทโวจุํ “ตฺวํ โขสิ อุปาสก กตกลฺยาโณ กตกุสโล กตภีรุตฺตาโณ, อกตปาโป อกตลุทฺโท อกตกิพฺพิโส, กตํ ตยา กลฺยาณํ, อกตํ ตยา ปาปํ, กิํ ตุยฺหิมินา ปาปเกน ทุชฺชีวิเตน, มตํ เต ชีวิตา เสยฺโย, อิโต ตฺวํ กาลงฺกโต กายสฺส เภทา ปรํ มรณา สุคติํ สคฺคํ โลกํ อุปปชฺชิสฺสสิ, ตตฺถ ทิพฺเพหิ ปญฺจหิ กามคุเณหิ สมปฺปิโต สมงฺคีภูโต ปริจาเรสฺสสี”ติ.}

\proseitem{169}{อถ โข โส อุปาสโก “สจฺจํ โข อยฺยา อาหํสุ, อหญฺหิ กตกลฺยาโณ กตกุสโล กตภีรุตฺตาโณ, อกตปาโป อกตลุทฺโท อกตกิพฺพิโส, กตํ มยา กลฺยาณํ, อกตํ มยา \csromanfolio{91}ปาปํ, กิํ มยฺหิมินา ปาปเกน ทุชฺชีวิเตน, มตํ เม ชีวิตา เสยฺโย, อิโต อหํ กาลงฺกโต กายสฺส เภทา ปรํ มรณา สุคติํ สคฺคํ โลกํ อุปปชฺชิสฺสามิ, ตตฺถ ทิพฺเพหิ ปญฺจหิ กามคุเณหิ สมปฺปิโต สมงฺคีภูโต ปริจาเรสฺสามี”ติ โส อสปฺปายานิ เจว โภชนานิ ภุญฺชิ, อสปฺปายานิ จ \textbf{ขาทนียานิ ขาทิ, อสปฺปายานิ จ สายนียานิ สายิ, อสปฺปายานิ จ ปานานิ} \textbf{ปิวิ. ตสฺส อสปฺปายานิ เจว โภชนานิ ภุญฺชโต อสปฺปายานิ จ ขาทนียานิ} \textbf{ขาทโต อสปฺปายานิ จ สายนียานิ สายโต อสปฺปายานิ จ ปานานิ ปิวโต ขโร} \textbf{อาพาโธ อุปฺปชฺชิ, โส เตเนว อาพาเธน กาลมกาสิ. ตสฺส ปชาปติ อุชฺฌายติ} \textbf{ขิยฺยติ วิปาเจติ “อลชฺชิโน อิเม สมณา สกฺยปุตฺติยา ทุสฺสีลา มุสาวาทิโน,} \textbf{อิเม หิ นาม ธมฺมจาริโน สมจาริโน พฺรหฺมจาริโน สจฺจวาทิโน สีลวนฺโต} \textbf{กลฺยาณธมฺมา ปฏิชานิสฺสนฺติ. นตฺถิ อิเมสํ สามญฺญํ, นตฺถิ อิเมสํ} \textbf{พฺรหฺมญฺญํ, นฏฺฐํ อิเมสํ สามญฺญํ, นฏฺฐํ อิเมสํ พฺรหฺมญฺญํ,} กุโต อิเมสํ สามญฺญํ, กุโต อิเมสํ พฺรหฺมญฺญํ, อปคตา อิเม สามญฺญา, อปคตา อิเม พฺรหฺมญฺญา, อิเม เม สามิกสฺส มรณวณฺณํ สํวณฺเณสุํ, อิเมหิ เม สามิโก มาริโต”ติ. อญฺเญปิ มนุสฺสา อุชฺฌายนฺติ ขิยฺยนฺติ วิปาเจนฺติ “อลชฺชิโน อิเม สมณา สกฺยปุตฺติยา ทุสฺสีลา มุสาวาทิโน, อิเม หิ นาม ธมฺมจาริโน สมจาริโน พฺรหฺมจาริโน สจฺจวาทิโน สีลวนฺโต กลฺยาณธมฺมา ปฏิชานิสฺสนฺติ. นตฺถิ อิเมสํ สามญฺญํ, นตฺถิ อิเมสํ พฺรหฺมญฺญํ, นฏฺฐํ อิเมสํ สามญฺญํ, นฏฺฐํ อิเมสํ พฺรหฺมญฺญํ, กุโต อิเมสํ สามญฺญํ, กุโต อิเมสํ พฺรหฺมญฺญํ, อปคตา อิเม สามญฺญา, อปคตา อิเม พฺรหฺมญฺญา, อิเม อุปาสกสฺส มรณวณฺณํ สํวณฺเณสุํ อิเมหิ อุปาสโก มาริโต”ติ. อสฺโสสุํ โข ภิกฺขู เตสํ มนุสฺสานํ อุชฺฌายนฺตานํ ขิยฺยนฺตานํ วิปาเจนฺตานํ, เย เต ภิกฺขู อปฺปิจฺฉา ฯเปฯ เต อุชฺฌายนฺติ ขิยฺยนฺติ วิปาเจนฺติ “กถํ หิ นาม ฉพฺพคฺคิยา ภิกฺขู อุปาสกสฺส มรณวณฺณํ สํวณฺณิสฺสนฺตี”ติ.}

\csromanmatikaanchor{mtk.25}
\csromanmatikaanchor{mtk.26}
\csromantocmarkref{subsection}{อนุปญฺญตฺติ}{mtk.25}
\bookmark[dest={mtk.25},level=2]{อนุปญฺญตฺติ}
\csromantocmarkref{subsection}{สิกฺขาปทวิภงฺค, ปทภาชนีย}{mtk.26}
\bookmark[dest={mtk.26},level=2]{สิกฺขาปทวิภงฺค, ปทภาชนีย}
\proseitem{170}{อถ โข เต ภิกฺขู ฉพฺพคฺคิเย ภิกฺขู อเนก\-ปริยาเยน วิครหิตฺวา ภควโต เอตมตฺถํ อาโรเจสุํ ฯเปฯ “สจฺจํ กิร ตุเมฺห ภิกฺขเว อุปาสกสฺส มรณวณฺณํ สํวณฺเณถา”ติ. สจฺจํ ภควาติ. วิครหิ พุทฺโธ ภควา “อนนุจฺฉวิกํ โมฆปุริสา อนนุโลมิกํ อปฺปติรูปํ อสฺสามณกํ อกปฺปิยํ อกรณียํ, กถํ หิ นาม ตุเมฺห โมฆปุริสา \csromanfolio{92}อุปาสกสฺส มรณวณฺณํ สํวณฺณิสฺสถ, เนตํ โมฆปุริสา \textbf{อปฺปสนฺนานํ วา ปสาทาย ฯเปฯ. เอวญฺจ ปน ภิกฺขเว อิมํ} \textbf{สิกฺขาปทํ อุทฺทิเสยฺยาถ\csromandash{}}}

\proseitem{171}{\textbf{“โย ปน ภิกฺขุ สญฺจิจฺจ มนุสฺส}\-\textbf{วิคฺคหํ ชีวิตา โวโรเปยฺย, สตฺถหารกํ วาสฺส ปริเยเสยฺย, มรณวณฺณํ วา สํวณฺเณยฺย, มรณาย วา สมาทเปยฺย ‘อมฺโภ ปุริส กิํ ตุยฺหิมินา ปาปเกน ทุชฺชีวิเตน, มตํ เต ชีวิตา เสยฺโย’ติ, อิติ จิตฺตมโน จิตฺตสงฺกปฺโป อเนก}\-\textbf{ปริยาเยน มรณวณฺณํ วา สํวณฺเณยฺย มรณาย วา สมาทเปยฺย, อยํปิ ปาราชิโก โหติ อสํวาโส”}ติ.}

\proseitem{172}{\textbf{โยปนา}ติ\csromansharedfootnote{423cb82fe87b6ac9}{โย ปนาติ เทฺว ปทานิ \csromandash{} ม.พ.ป.} โย ยาทิโส ฯเปฯ. \textbf{ภิกฺขู}ติ ฯเปฯ. อยํ อิมสฺมิํ อตฺเถ อธิปฺเปโต ภิกฺขูติ.}

\prose{\textbf{สญฺจิจฺจา}ติ ชานนฺโต สญฺชานนฺโต เจจฺจ อภิวิตริตฺวา วีติกฺกโม.}

\prose{\textbf{มนุสฺสวิคฺคโห} นาม ยํ มาตุกุจฺฉิสฺมิํ ปฐมํ จิตฺตํ อุปฺปนฺนํ ปฐมํ วิญฺญาณํ ปาตุภูตํ, ยาว มรณกาลา เอตฺถนฺตเร เอโส มนุสฺสวิคฺคโห นาม.}

\prose{\textbf{ชีวิตา โวโรเปยฺยา}ติ ชีวิตินฺทฺริยํ อุปจฺฉินฺทติ, อุปโรเธติ, สนฺตติํ วิโกเปติ.}

\prose{\textbf{สตฺถหารกํ วาสฺส ปริเยเสยฺยา}ติ อสิํ วา สตฺติํ วา เภณฺฑิํ วา\csromansharedfootnote{e87fbea3afe8ae4c}{เภนฺทิํ วา (ก)} ลคุฬํ วา\csromansharedfootnote{6346c4c19f2b6bba}{สูลํวา ลคุฬํวา (สฺยา)} ปาสาณํ วาสตฺถํ วา วิสํ วา รชฺชุํ วา.}

\prose{\textbf{มรณวณฺณํ วา สํวณฺเณยฺยา}ติ ชีวิเต อาทีนวํ ทสฺเสติ, มรเณ วณฺณํ ภณติ.}

\prose{\textbf{มรณาย วา สมาทเปยฺยา}ติ สตฺถํ วา อาหร, วิสํ วา ขาท, รชฺชุยา วา อุพฺพนฺธิตฺวา กาลงฺกโรหีติ.}

\prose{\csromanfolio{93}\textbf{อมฺโภ ปุริสา}ติ อาลปนาธิวจนเมตํ.}

\prose{\textbf{กิํ ตุยฺหิมินา ปาปเกน ทุชฺชีวิเตนา}ติ ปาปกํ นาม ชีวิตํ อฑฺฒานํ ชีวิตํ อุปาทาย ทลิทฺทานํ ชีวิตํ ปาปกํ ลามกํ, สธนานํ ชีวิตํ อุปาทาย อธนานํ ชีวิตํ ปาปกํ, เทวานํ ชีวิตํ อุปาทาย มนุสฺสานํ ชีวิตํ ปาปกํ.}

\prose{\textbf{ทุชฺชีวิตํ} นาม หตฺถ\-จฺฉินฺนสฺส ปาทจฺฉินฺนสฺส หตฺถ\-ปาท\-จฺฉินฺนสฺส กณฺณ\-จฺฉินฺนสฺส นาสจฺฉินฺนสฺส กณฺณ\-นาส\-จฺฉินฺนสฺส, อิมินา จ ปาปเกน อิมินา ทุชฺชีวิเตน มตํ เต ชิวิตา เสยฺโยติ.}

\prose{\textbf{อิติ จิตฺตมโน}ติ ยํ จิตฺตํ ตํ มโน, ยํ มโน ตํ จิตฺตํ.}

\prose{\textbf{จิตฺต}\-\textbf{สงฺกปฺโป}ติ มรณสญฺญี มรณเจตโน มรณธิปฺปาโย.}

\prose{\textbf{อเนก}\-\textbf{ปริยาเย}\-\textbf{นา}ติ อุจฺจาวเจหิ อากาเรหิ.}

\prose{\textbf{มรณวณฺณํ วา สํวณฺเณยฺยา}ติ ชีวิเต อาทีนวํ ทสฺเสติ, มรเณ วณฺณํ ภณติ “อิโต ตฺวํ กาลงฺกโต กายสฺส เภทา ปรํ มรณา สุคติํ สคฺคํ โลกํ อุปปชฺชิสฺสสิ, ตตฺถ ทิพฺเพหิ ปญฺจหิ กามคุเณหิ สมปฺปิโต สมงฺคีภูโต ปริจาเรสฺสสี”ติ.}

\prose{\textbf{มรณาย วา สมาทเปยฺยา}ติ สตฺถํ วา อาหร, วิสํ วา ขาท, รชฺชุยา วา อุพฺพนฺธิตฺวา กาลงฺกโรหิ, โสพฺเภ วา นรเก วา ปปาเต วา ปปตาติ.}

\prose{\textbf{อยมฺปี}ติ ปุริเม อุปาทาย วุจฺจติ.}

\prose{\textbf{ปาราชิโก โหตี}ติ เสยฺยถาปิ นาม ปุถุสิลา ทฺวิธา ภินฺนา\csromansharedfootnote{b528b422ca434791}{เทฺวธา ภินฺนา (สฺยา)} อปฺปฏิสนฺธิกา โหติ, เอวเมว ภิกฺขุ สญฺจิจฺจ มนุสฺส\-วิคฺคหํ ชีวิตา โวโรเปตฺวา อสฺสมโณ โหติ อสกฺยปุตฺติโย, เตน วุจฺจติ “ปาราชิโก โหตี”ติ.}

\prose{\csromanfolio{94}\textbf{อสํวาโส}ติ สํวาโส นาม เอกกมฺมํ เอกุทฺเทโส สมสิกฺขตา, เอโส สํวาโส นาม. โส เตน สทฺธิํ นตฺถิ, เตน วุจฺจติ “อสํวาโส”ติ.}

\proseitem{173}{สามํ, อธิฏฺฐาย, ทูเตน, ทูตปรํปราย, วิสกฺกิเยน ทูเตน, คตปจฺจาเตน ทูเตน, อรโห รโหสญฺญี, รโห อรโหสญฺญี, อรโห อรโหสญฺญี, รโห รโหสญฺญี, กาเยน สํวณฺเณติ, วาจาย สํวณฺเณติ, กาเยน วาจาย สํวณฺเณติ, ทูเตน สํวณฺเณติ, เลขาย สํวณฺเณติ, โอปาตํ, อปสฺเสนํ, อุปนิกฺขิปนํ, เภสชฺชํ, รูปูปหาโร, สทฺทูปหาโร, คนฺธูปหาโร, รสูปหาโร, โผฏฺฐพฺพู\-ปหาโร, ธมฺมูปหาโร, อาจิกฺขนา, อนุสาสนี, สงฺเกตกมฺมํ, นิมิตฺต\-กมฺมนฺติ.}

\proseitem{174}{\textbf{สามนฺ}ติ สยํ หนติ กาเยน วา กาย\-ปฏิพทฺเธน วา นิสฺสคฺคิเยน วา.}

\prose{\textbf{อธิฏฺฐายา}ติ อธิฏฺฐหิตฺวา อาณาเปติ “เอวํ วิชฺฌ, เอวํ ปหร, เอวํ ฆาเตหี”ติ.}

\prose{ภิกฺขุ ภิกฺขุํ อาณาเปติ “อิตฺถนฺนามํ ชีวิตา โวโรเปหี”ติ, อาปตฺติ ทุกฺกฏสฺส. โส ตํ มญฺญมาโน ตํ ชีวิตา โวโรเปติ, อาปตฺติ อุภินฺนํ ปาราชิกสฺส.}

\prose{ภิกฺขุ ภิกฺขุํ อาณาเปติ “อิตฺถนฺนามํ ชีวิตา โวโรเปหี”ติ, อาปตฺติ ทุกฺกฏสฺส. โส ตํ มญฺญมาโน อญฺญํ ชีวิตา โวโรเปติ, มูลฏฺฐสฺส อนาปตฺติ, วธกสฺส อาปตฺติ ปาราชิกสฺส.}

\prose{ภิกฺขุ ภิกฺขุํ อาณาเปติ “อิตฺถนฺนามํ ชีวิตา โวโรเปหี”ติ, อาปตฺติ ทุกฺกฏสฺส. โส อญฺญํ มญฺญมาโน ตํ ชีวิตา โวโรเปติ, อาปตฺติ อุภินฺนํ ปาราชิกสฺส.}

\prose{ภิกฺขุ ภิกฺขุํ อาณาเปติ “อิตฺถนฺนามํ ชีวิตา โวโรเปหี”ติ, อาปตฺติ ทุกฺกฏสฺส. โส อญฺญํ มญฺญมาโน อญฺญํ ชีวิตา โวโรเปติ, มูลฏฺฐสฺส อนาปตฺติ, วธกสฺส อาปตฺติ ปาราชิกสฺส.}

\prose{\csromanfolio{95}ภิกฺขุ ภิกฺขุํ อาณาเปติ “อิตฺถนฺนามสฺส ปาวท, ‘อิตฺถนฺนาโม อิตฺถนฺนามสฺส ปาวทตุ, ‘อิตฺถนฺนาโม อิตฺถนฺนามํ ชีวิตา โวโรเปตู”ติ. อาปตฺติ ทุกฺกฏสฺส. โส อิตรสฺส อาโรเจติ, อาปตฺติ ทุกฺกฏสฺส. วธโก ปฏิคฺคณฺหาติ, มูลฏฺฐสฺส อาปตฺติ ถุลฺลจฺจยสฺส. โส ตํ ชีวิตา โวโรเปติ. อาปตฺติ สพฺเพสํ ปาราชิกสฺส.}

\prose{ภิกฺขุ ภิกฺขุํ อาณาเปติ “อิตฺถนฺนามสฺส ปาวท, ‘อิตฺถนฺนาโม อิตฺถนฺนามสฺส ปาวทตุ, ‘อิตฺถนฺนาโม อิตฺถนฺนามํ ชีวิตา โวโรเปตู”ติ. อาปตฺติ ทุกฺกฏสฺส. โส อญฺญํ อาณาเปติ, อาปตฺติ ทุกฺกฏสฺส. วธโก ปฏิคฺคณฺหาติ อาปตฺติ ทุกฺกฏสฺส. โส ตํ ชีวิตา โวโรเปติ, มูลฏฺฐสฺส อนาปตฺติ. อาณาปกสฺส จ วธกสฺส จ อาปตฺติ ปาราชิกสฺส.}

\prose{ภิกฺขุ ภิกฺขุํ อาณาเปติ “อิตฺถนฺนามํ ชีวิตา โวโรเปหี”ติ, อาปตฺติ ทุกฺกฏสฺส. โส คนฺตฺวา ปุน ปจฺจาคจฺฉติ “นาหํ สกฺโกมิ ตํ ชีวิตา โวโรเปตุนฺ”ติ. โส ปุน อาณาเปติ “ยทา สกฺโกสิ, ตทา ตํ ชีวิตา โวโรเปหี”ติ, อาปตฺติ ทุกฺกฏสฺส. โส ตํ ชีวิตา โวโรเปติ, อาปตฺติ อุภินฺนํ ปาราชิกสฺส.}

\prose{ภิกฺขุ ภิกฺขุํ อาณาเปติ “อิตฺถนฺนามํ ชีวิตา โวโรเปหี”ติ, อาปตฺติ ทุกฺกฏสฺส. โส อาณาเปตฺวา วิปฺปฏิสารี น สาเวติ “มา ฆาเตหี”ติ. โส ตํ ชีวิตา โวโรเปติ, อาปตฺติ อุภินฺนํ ปาราชิกสฺส.}

\prose{ภิกฺขุ ภิกฺขุํ อาณาเปติ “อิตฺถนฺนามํ ชีวิตา โวโรเปหี”ติ, อาปตฺติ ทุกฺกฏสฺส. โส อาณาเปตฺวา วิปฺปฏสารี สาเวติ “มา ฆาเตหี”ติ. โส “อาณตฺโต อหํ ตยา”ติ ตํ ชีวิตา โวโรเปติ, มูลฏฺฐสฺส อนาปตฺติ. วธกสฺส อาปตฺติ ปาราชิกสฺส.}

\prose{ภิกฺขุ ภิกฺขุํ อาณาเปติ “อิตฺถนฺนามํ ชีวิตา โวโรเปหี”ติ, อาปตฺติ ทุกฺกฏสฺส. โส อาณาเปตฺวา วิปฺปฏิสารี สาเวติ “มา ฆาเตหี”ติ. โส “สาธู”ติ โอรมติ, อุภินฺนํ อนาปตฺติ.}

\proseitem{175}{อรโห รโหสญฺญี อุลฺลปติ “อโห อิตฺถนฺนาโม หโต อสฺสา”ติ, อาปตฺติ ทุกฺกฏสฺส. รโห อรโหสญฺญี อุลฺลปติ “อโห \csromanfolio{96}อิตฺถนฺนาโม หโต อสฺสา”ติ, อาปตฺติ ทุกฺกฏสฺส. อรโห อรโหสญฺญี อุลฺลปติ “อโห อิตฺถนฺนาโม หโต อสฺสา”ติ, อาปตฺติ ทุกฺกฏสฺส. รโห รโหสญฺญี อุลฺลปติ “อโห อิตฺถนาโม หโต อสฺสา”ติ, อาปตฺติ ทุกฺกฏสฺส.}

\prose{\textbf{กาเยน สํวณฺเณติ} นาม กาเยน วิการํ กโรติ “โย เอวํ มรติ, โส ธนํ วา ลภติ, ยสํ วา ลภติ, สคฺคํ วา คจฺฉตี”ติ, อาปตฺติ ทุกฺกฏสฺส. ตาย สํวณฺณนาย มริสฺสามีติ ทุกฺขํ เวทนํ อุปฺปาเทติ, อาปตฺติ ถุลฺลจฺจยสฺส. มรติ, อาปตฺติ ปาราชิกสฺส.}

\prose{\textbf{วาจาย สํวณฺเณติ} นาม วาจาย ภณติ “โย เอวํ มรติ, โส ธนํ วา ลภติ, ยสํ วา ลภติ, สคฺคํ วา คจฺฉตี”ติ, อาปตฺติ ทุกฺกฏสฺส. ตาย สํวณฺณนาย มริสฺสามีติ ทุกฺขํ เวทนํ อุปฺปาเทติ, อาปตฺติ ถุลฺลจฺจยสฺส. มรติ, อาปตฺติ ปาราชิกสฺส.}

\prose{\textbf{กาเยน วาจาย สํวณฺเณติ} นาม กาเยน จ วิการํ กโรติ, วาจาย จ ภณติ “โย เอวํ มรติ, โส ธนํ วา ลภติ, ยสํ วา ลภติ, สคฺคํ วา คจฺฉตี”ติ, อาปตฺติ ทุกฺกฏสฺส. ตาย สํวณฺณนาย มริสฺสามีติ ทุกฺขํ เวทนํ อุปฺปาเทติ, อาปตฺติ ถุลฺลจฺจยสฺส. มรติ, อาปตฺติ ปาราชิกสฺส.}

\prose{\textbf{ทูเตน สํวณฺเณติ} นาม ทูตสฺส สาสนํ อาโรเจติ “โย เอวํ มรติ, โส ธนํ วา ลภติ, ยสํ วา ลภติ, สคฺคํ วา คจฺฉตี”ติ, อาปตฺติ ทุกฺกฏสฺส. ทูตสฺส สาสนํ สุตฺวา มริสฺสามีติ ทุกฺขํ เวทนํ อุปฺปาเทติ, อาปตฺติ ถุลฺลจฺจยสฺส. มรติ, อาปตฺติ ปาราชิกสฺส.}

\proseitem{176}{\textbf{เลขาย สํวณฺเณติ} นาม เลขํ ฉินฺทติ “โย เอวํ มรติ, โส ธนํ วา ลภติ, ยสํ วา ลภติ, สคฺคํ วา คจฺฉตี”ติ, อกฺขร\-กฺขราย อาปตฺติ ทุกฺกฏสฺส. เลขํ ปสฺสิตฺวา มริสฺสามีติ ทุกฺขํ เวทนํ อุปฺปาเทติ, อาปตฺติ ถุลฺลจฺจยสฺส. มรติ, อาปตฺติ ปาราชิกสฺส.}

\prose{\textbf{โอปาตํ} นาม มนุสฺสํ อุทฺทิสฺส โอปาตํ ขนติ “ปปติตฺวา มริสฺสตี”ติ, อาปตฺติ ทุกฺกฏสฺส. ปปติเต ทุกฺขา เวทนา อุปฺปชฺชติ, อาปตฺติ ถุลฺลจฺจยสฺส. มรติ, อาปตฺติ ปาราชิกสฺส. อโนทิสฺส โอปาตํ ขนติ “โย โกจิ \csromanfolio{97}ปปติตฺวา มริสฺสตี”ติ, อาปตฺติ ทุกฺกฏสฺส. มนุสฺโส ตสฺมิํ ปปตติ, อาปตฺติ ทุกฺกฏสฺส. ปปติเต ทุกฺขา เวทนา อุปฺปชฺชติ, อาปตฺติ ถุลฺลจฺจยสฺส. มรติ, อาปตฺติ ปาราชิกสฺส. ยกฺโข วา เปโต วา ติรจฺฉานคต\-มนุสฺสวิคฺคโห วา ตสฺมิํ ปปตติ, อาปตฺติ ทุกฺกฏสฺส. ปปติเต ทุกฺขา เวทนา อุปฺปชฺชติ, อาปตฺติ ทุกฺกฏสฺส. มรติ, อาปตฺติ ถุลฺลจฺจยสฺส. ติรจฺฉานคโต ตสฺมิํ ปปตติ, อาปตฺติ ทุกฺกฏสฺส. ปปติเต ทุกฺขา เวทนา อุปฺปชฺชติ, อาปตฺติ ทุกฺกฏสฺส. มรติ, อาปตฺติ ปาจิตฺติยสฺส.}

\proseitem{177}{\textbf{อปสฺเสนํ} นาม อปสฺเสเน สตฺถํ วา ฐเปติ, วิเสน วา มกฺเขติ, ทุพฺพลํ วา กโรติ, โสพฺเภ วา นรเก วา ปปาเต วา ฐเปติ “ปปติตฺวา มริสฺสตี”ติ, อาปตฺติ ทุกฺกฏสฺส. สตฺเถน วา วิเสน วา ปปติเตน วา ทุกฺขา เวทนา อุปฺปชฺชติ, อาปตฺติ ถุลฺลจฺจยสฺส. มรติ, อาปตฺติ ปาราชิกสฺส.}

\prose{\textbf{อุปนิกฺขิปนํ} นาม อสิํ วา สตฺติํ วา เภณฺฑิํ วา ลคุฬํ วา ปาสาณํ วา สตฺถํ วา วิสํ วา รชฺชุํ วา อุปนิกฺขิปติ “อิมินา มริสฺสตี”ติ, อาปตฺติ ทุกฺกฏสฺส. เตน มริสฺสามีติ ทุกฺขํ เวทนํ อุปฺปาเทติ, อาปตฺติ ถุลฺลจฺจยสฺส. มรติ, อาปตฺติ ปาราชิกสฺส.}

\prose{\textbf{เภสชฺชํ} นาม สปฺปิํ วา นวนีตํ วา เตลํ วา มธุํ วา ผาณิตํ วา เทติ “อิมํ สายิตฺวา มริสฺสตี”ติ, อาปตฺติ ทุกฺกฏสฺส. ตํ สายิเต ทุกฺขา เวทนา อุปฺปชฺชติ, อาปตฺติ ถุลฺลจฺจยสฺส. มรติ, อาปตฺติ ปาราชิกสฺส.}

\proseitem{178}{\textbf{รูปูปหาโร} นาม อมนาปิกํ รูปํ อุปสํหรติ ภยานกํ เภรวํ “อิมํ ปสฺสิตฺวา อุตฺตสิตฺวา มริสฺสตี”ติ, อาปตฺติ ทุกฺกฏสฺส. ตํ ปสฺสิตฺวา อุตฺตสติ, อาปตฺติ ถุลฺลจฺจยสฺส. มรติ, อาปตฺติ ปาราชิกสฺส. มนาปิกํ รูปํ อุปสํหรติ\csromansharedfootnote{3c2cf7576aebc3af}{อุปสํหรติ เปมนียํ หทยงฺคมํ (สฺยา)} “อิมํ ปสฺสิตฺวา อลาภเกน สุสฺสิตฺวา มริสฺสตี”ติ, อาปตฺติ ทุกฺกฏสฺส. ตํ ปสฺสิตฺวา อลาภเกน สุสฺสติ, อาปตฺติ ถุลฺลจฺจยสฺส. มรติ, อาปตฺติ ปาราชิกสฺส.}

\prose{\textbf{สทฺทูปหาโร} นาม อมนาปิกํ สทฺทํ อุปสํหรติ ภยานกํ เภรวํ “อิมํ สุตฺวา อุตฺตสิตฺวา มริสฺสตี”ติ, อาปตฺติ ทุกฺกฏสฺส. ตํ สุตฺวา \csromanfolio{98}อุตฺตสติ, อาปตฺติ ถุลฺลจฺจยสฺส. มรติ, อาปตฺติ ปาราชิกสฺส. มนาปิกํ สทฺทํ อุปสํหรติ เปมนียํ หทยงฺคมํ “อิมํ สุตฺวา อลาภเกน สุสฺสิตฺวา มริสฺสตี”ติ, อาปตฺติ ทุกฺกฏสฺส. ตํ สุตฺวา อลาภเกน สุสฺสติ, อาปตฺติ ถุลฺลจฺจยสฺส. มรติ, อาปตฺติ ปาราชิกสฺส.}

\prose{\textbf{คนฺธูปหาโร} นาม อมนาปิกํ คนฺธํ อุปสํหรติ เชคุจฺฉํ ปาฏิกุลฺยํ\csromansharedfootnote{ec9dab5c519d6eaf}{ปฏิกูลํ (?)} “อิมํ ฆายิตฺวา เชคุจฺฉตา ปาฏิกุลฺยตา มริสฺสตี”ติ, อาปตฺติ ทุกฺกฏสฺส. ตํ ฆายิเต เชคุจฺฉตา ปาฏิกุลฺยตา ทุกฺขา เวทนา อุปฺปชฺชติ, อาปตฺติ ถุลฺลจฺจยสฺส. มรติ, อาปตฺติ ปาราชิกสฺส. มนาปิกํ คนฺธํ อุปสํหรติ “อิมํ ฆายิตฺวา อลาภเกน สุสฺสิตฺวา มริสฺสตี”ติ, อาปตฺติ ทุกฺกฏสฺส. ตํ ฆายิตฺวา อลาภเกน สุสฺสติ, อาปตฺติ ถุลฺลจฺจยสฺส. มรติ, อาปตฺติ ปาราชิกสฺส.}

\prose{\textbf{รสูปหาโร} นาม อมนาปิกํ รสํ อุปสํหรติ เชคุจฺฉํ ปาฏิกุลฺยํ\csromansharedfootnote{ec9dab5c519d6eaf}{ปฏิกูลํ (?)} “อิมํ สายิตฺวา เชคุจฺฉตา ปฏิกุลฺยตา มริสฺสตี”ติ, อาปตฺติ ทุกฺกฏสฺส. ตํ สายิเต เชคุจฺฉตา ปาฏิกุลฺยตา ทุกฺขา เวทนา อุปฺปชฺชติ, อาปตฺติ ถุลฺลจฺจยสฺส. มรติ, อาปตฺติ ปาราชิกสฺส. มนาปิกํ รสํ อุปสํหรติ “อิมํ สายิตฺวา อลาภเกน สุสฺสิตฺวา มริสฺสตี”ติ, อาปตฺติ ทุกฺกฏสฺส. ตํ สายิตฺวา อลาภเกน สุสฺสติ, อาปตฺติ ถุลฺลจฺจยสฺส. มรติ, อาปตฺติ ปาราชิกสฺส.}

\prose{\textbf{โผฏฺฐพฺพู}\-\textbf{ปหาโร} นาม อมนาปิกํ โผฏฺฐพฺพํ อุปสํหรติ ทุกฺข\-สมฺผสฺสํ ขรมฺผสฺสํ “อิมินา ผุฏฺโฐ มริสฺสตี”ติ, อาปตฺติ ทุกฺกฏสฺส. เตน ผุฏฺฐสฺส ทุกฺขา เวทนา อุปฺปชฺชติ, อาปตฺติ ถุลฺลจฺจยสฺส. มรติ, อาปตฺติ ปาราชิกสฺส. มนาปิกํ โผฏฺฐพฺพํ อุปสํหรติ สุขสมฺผสฺสํ มุทุสมฺผสฺสํ “อิมินา ผุฏฺโฐ อลาภเกน สุสฺสิตฺวา มริสฺสตี”ติ, อาปตฺติ ทุกฺกฏสฺส. เตน ผุฏฺโฐ อลาภเกน สุสฺสติ, อาปตฺติ ถุลฺลจฺจยสฺส. มรติ, อาปตฺติ ปาราชิกสฺส.}

\prose{\textbf{ธมฺมูปหาโร} นาม เนรยิกสฺส นิรยกถํ กเถติ “อิมํ สุตฺวา อุตฺตสิตฺวา มริสฺสตี”ติ, อาปตฺติ ทุกฺกฏสฺส. ตํ สุตฺวา อุตฺตสติ, อาปตฺติ ถุลฺลจฺจยสฺส. มรติ, อาปตฺติ ปาราชิกสฺส. กลฺยาณ\-กมฺมสฺส สคฺคกถํ กเถติ “อิมํ สุตฺวา อธิมุตฺโต มริสฺสตี”ติ, อาปตฺติ ทุกฺกฏสฺส. ตํ \csromanfolio{99}สุตฺวา อธิมุตฺโต มริสฺสามีติ ทุกฺขํ เวทนํ อุปฺปาเทติ, อาปตฺติ ถุลฺลจฺจยสฺส. มรติ, อาปตฺติ ปาราชิกสฺส.}

\proseitem{179}{\textbf{อาจิกฺขนา} นาม ปุฏฺโฐ ภณติ, “เอวํ มรสฺสุ, โย เอวํ มรติ, โส ธนํ วา ลภติ, ยสํ วา ลภติ, สคฺคํ วา คจฺฉตี”ติ, อาปตฺติ ทุกฺกฏสฺส. ตาย อาจิกฺขนาย มริสฺสามีติ ทุกฺขํ เวทนํ อุปฺปาเทติ, อาปตฺติ ถุลฺลจฺจยสฺส. มรติ, อาปตฺติ ปาราชิกสฺส.}

\prose{\textbf{อนุสาสนี} นาม อปุฏฺโฐ ภณติ “เอวํ มรสฺสุ, โย เอวํ มรติ, โส ธนํ วา ลภติ, ยสํ วา ลภติ, สคฺคํ วา คจฺฉตี”ติ, อาปตฺติ ทุกฺกฏสฺส. ตาย อนุสาสนิยา มริสฺสามีติ ทุกฺขํ เวทนํ อุปฺปาเทติ, อาปตฺติ ถุลฺลจฺจยสฺส. มรติ, อาปตฺติ ปาราชิกสฺส.}

\prose{\textbf{สงฺเกตกมฺมํ} นาม สงฺเกตํ กโรติ ปุเรภตฺตํ วา ปจฺฉาภตฺตํ วา รตฺติํ วา ทิวา วา “เตน สงฺเกเตน ตํ ชีวิตา โวโรเปหี”ติ, อาปตฺติ ทุกฺกฏสฺส. เตน สงฺเกเตน ตํ ชีวิตา โวโรเปติ, อาปตฺติ อุภินฺนํ ปาราชิกสฺส. ตํ สงฺเกตํ ปุเร วา ปจฺฉา วา ตํ ชีวิตา โวโรเปติ, มูลฏฺฐสฺส อนาปตฺติ. วธกสฺส อาปตฺติ ปาราชิกสฺส.}

\prose{\textbf{นิมิตฺตกมฺมํ} นาม นิมิตฺตํ กโรติ “อกฺขิํ วา นิขณิสฺสามิ, ภมุกํ วา อุกฺขิปิสฺสามิ, สีสํ วา อุกฺขิปิสฺสามิ, เตน นิมิตฺเตน ตํ ชีวิตา โวโรเปหี”ติ, อาปตฺติ ทุกฺกฏสฺส. เตน นิมิตฺเตน ตํ ชีวิตา โวโรเปติ, อาปตฺติ อุภินฺนํ ปาราชิกสฺส. ตํ นิมิตฺตํ ปุเร วา ปจฺฉา วา ตํ ชีวิตา โวโรเปติ, มูลฏฺฐสฺส อนาปตฺติ. วธกสฺส อาปตฺติ ปาราชิกสฺส.}

\prosewithcloser{อนาปตฺติ อสญฺจิจฺจ อชานนฺตสฺส นมรณา\-ธิปฺปายสฺส อุมฺมตฺตกสฺส\csromansharedfootnote{ea970a7cbbf62c4b}{อุมฺมตฺตกสฺส ขิตฺตจิกฺกสฺส เวทนาวฏฺฏสฺส (สฺยา)} อาทิกมฺมิกสฺสาติ.}{มนุสฺสวิคฺคหปาชิกมฺหิ ปฐม\-ภาณวาโร นิฏฺฐิโต.}
\csromanmatikaanchor{mtk.27}
\csromantocmarkref{subsection}{วินีตวตฺถุอุทฺทานคาถา}{mtk.27}
\bookmark[dest={mtk.27},level=2]{วินีตวตฺถุอุทฺทานคาถา}
\csromanheader{2}{\textbf{วินีตวตฺถุ}\-อุทฺทานคาถา}
\setlength{\csromangathapagewidth}{0pt}
\setlength{\csromangathawakwidth}{0pt}
\csromangathameasuregroup{สํวณฺณนา นิสีทนฺโต, \\ วุฑฺฒปพฺพชิตาภิสนฺโน, \\ ตโย จ วตฺถุกมฺเมหิ, \\ วาสี โคปานสี เจว, \\ เสทํ นตฺถุญฺจ สมฺพาโห, \\ อุฏฺฐาเปนฺโต นิปาเตนฺโต, \\ ชารคพฺโภ สปตฺตี จ, \\ อุโภ น มิยฺยเร มทฺทา, \\ ปโตทํ นิคฺคเห ยกฺโข, \\ ตํ มญฺญมาโน ปหริ, \\ อาฬวิยา ตโย รุกฺขา, \\ มา กิลเมสิ น ตุยฺหํ,}{\csromangathabat{สํวณฺณนา นิสีทนฺโต,}{มุสโลทุกฺขเลน จ.} \\ \csromangathabat{วุฑฺฒปพฺพชิตาภิสนฺโน,}{อคฺควีมํสนาวิสํ.} \\ \csromangathabat{ตโย จ วตฺถุกมฺเมหิ,}{อิฏฺฐกาหิปเร ตโย.} \\ \csromangathabat{วาสี โคปานสี เจว,}{อฏฺฏโกตรณํ ปติ.} \\ \csromangathabat{เสทํ นตฺถุญฺจ สมฺพาโห,}{นฺหาปนพฺภญฺชเนน จ.} \\ \csromangathabat{อุฏฺฐาเปนฺโต นิปาเตนฺโต,}{อนฺนปาเนน มารณํ.} \\ \csromangathabat{ชารคพฺโภ สปตฺตี จ,}{มาตา ปุตฺตํ อุโภ วธิ.} \\ \csromangathabat{อุโภ น มิยฺยเร มทฺทา,}{ตาปํ วญฺฌา วิชายินี.} \\ \csromangathabat{ปโตทํ นิคฺคเห ยกฺโข,}{วาฬยกฺขญฺจ ปาหิณิ.} \\ \csromangathabat{ตํ มญฺญมาโน ปหริ,}{สคฺคญฺจ นิรยํ ภเณ.} \\ \csromangathabat{อาฬวิยา ตโย รุกฺขา,}{ทาเยหิ อปเร ตโย.} \\ \csromangathabat{มา กิลเมสิ น ตุยฺหํ,}{ตกฺกํ โสวีรเกน จาติ.}}
\csromangathasetleft{สํวณฺณนา นิสีทนฺโต, \\ วุฑฺฒปพฺพชิตาภิสนฺโน, \\ ตโย จ วตฺถุกมฺเมหิ, \\ วาสี โคปานสี เจว, \\ เสทํ นตฺถุญฺจ สมฺพาโห, \\ อุฏฺฐาเปนฺโต นิปาเตนฺโต, \\ ชารคพฺโภ สปตฺตี จ, \\ อุโภ น มิยฺยเร มทฺทา, \\ ปโตทํ นิคฺคเห ยกฺโข, \\ ตํ มญฺญมาโน ปหริ, \\ อาฬวิยา ตโย รุกฺขา, \\ มา กิลเมสิ น ตุยฺหํ,}
\csromangathagroup{\csromangathastanza{\csromanfolio{100}\csromangathabat{สํวณฺณนา นิสีทนฺโต,}{มุสโลทุกฺขเลน จ.} \\ \csromangathabat{วุฑฺฒปพฺพชิตาภิสนฺโน,}{อคฺควีมํสนาวิสํ.}}\csromangathastanza{\csromangathabat{ตโย จ วตฺถุกมฺเมหิ,}{อิฏฺฐกาหิปเร ตโย.} \\ \csromangathabat{วาสี โคปานสี เจว,}{อฏฺฏโกตรณํ ปติ.}}\csromangathastanza{\csromangathabat{เสทํ นตฺถุญฺจ สมฺพาโห,}{นฺหาปนพฺภญฺชเนน จ.} \\ \csromangathabat{อุฏฺฐาเปนฺโต นิปาเตนฺโต,}{อนฺนปาเนน มารณํ.}}\csromangathastanza{\csromangathabat{ชารคพฺโภ สปตฺตี จ,}{มาตา ปุตฺตํ อุโภ วธิ.} \\ \csromangathabat{อุโภ น มิยฺยเร มทฺทา,}{ตาปํ วญฺฌา วิชายินี.}}\csromangathastanza{\csromangathabat{ปโตทํ นิคฺคเห ยกฺโข,}{วาฬยกฺขญฺจ ปาหิณิ.} \\ \csromangathabat{ตํ มญฺญมาโน ปหริ,}{สคฺคญฺจ นิรยํ ภเณ.}}\csromangathastanza{\csromangathabat{อาฬวิยา ตโย รุกฺขา,}{ทาเยหิ อปเร ตโย.} \\ \csromangathabat{มา กิลเมสิ น ตุยฺหํ,}{ตกฺกํ โสวีรเกน จาติ.}}}

\csromanmatikaanchor{mtk.28}
\csromantocmarkref{subsection}{วินีตวตฺถุ}{mtk.28}
\bookmark[dest={mtk.28},level=2]{วินีตวตฺถุ}
\csromanheader{2}{\textbf{วินีตวตฺถุ}}
\proseitem{180}{เตน โข ปน สมเยน อญฺญตโร ภิกฺขุ คิลาโน โหติ. ตสฺส ภิกฺขู การุญฺเญน มรณวณฺณํ สํวณฺเณสุํ. โส ภิกฺขุ กาลมกาสิ. เตสํ กุกฺกุจฺจํ อโหสิ “ภควตา สิกฺขาปทํ ปญฺญตฺตํ, กจฺจิ นุ โข มยํ ปาราชิกํ อาปตฺติํ อาปนฺนา”ติ. ภควโต เอตมตฺถํ อาโรเจสุํ. อาปตฺติํ ตุเมฺห ภิกฺขเว อาปนฺนา ปาราชิกนฺติ. (1)}

\prose{เตน โข ปน สมเยน อญฺญตโร ปิณฺฑจาริโก ภิกฺขุ ปีฐเก ปิโลติกาย ปฏิจฺฉนฺนํ นิสีทนฺโต โอตฺถริตฺวา มาเรสิ. ตสฺส กุกฺกุจฺจํ อโหสิ “ภควตา สิกฺขาปทํ ปญฺญตฺตํ, กจฺจิ นุ โข อหํ ปาราชิกํ อาปตฺติํ อาปนฺโน”ติ. ภควโต เอตมตฺถํ อาโรเจสิ. อนาปตฺติ ภิกฺขุ ปาราชิกสฺส. น จ ภิกฺขเว อปฺปฏิเวกฺขิตฺวา อาสเน นิสีทิตพฺพํ, โย นิสีเทยฺย, อาปตฺติ ทุกฺกฏสฺสาติ. (2)}

\prose{เตน โข ปน สมเยน อญฺญตโร ภิกฺขุ ภตฺตคฺเค อนฺตรฆเร อาสนํ ปญฺญเปนฺโต มุสเล อุสฺสิเต เอกํ มุสลํ อคฺคเหสิ. ทุติโย มุลโส ปริปติตฺวา อญฺญตรสฺส ทารกสฺส มตฺถเก อวตฺถาสิ. \csromanfolio{101}โส กาลมกาสิ. ตสฺส กุกฺกุจฺจํ อโหสิ ฯเปฯ. กิํจิตฺโต ตฺวํ ภิกฺขูติ. อสญฺจิจฺจ อหํ ภควาติ. อนาปตฺติ ภิกฺขุ อสญฺจิจฺจาติ. (3)}

\prose{เตน โข ปน สมเยน อญฺญตโร ภิกฺขุ ภตฺตคฺเค อนฺตรฆเร อาสนํ ปญฺญเปนฺโต อุทุกฺขล\-ภณฺฑิกํ อกฺกมิตฺวา ปวฏฺเฏสิ. อญฺญตรํ ทารกํ โอตฺถริตฺวา มาเรสิ. ตสฺส กุกฺกุจฺจํ อโหสิ ฯเปฯ. อนาปตฺติ ภิกฺขุ อสญฺจิจฺจาติ. (4)}

\prose{เตน โข ปน สมเยน ปิตาปุตฺตา ภิกฺขูสุ ปพฺพชิตา โหนฺติ. กาเล อาโรจิเต ปุตฺโต ปิตรํ เอตทโวจ “คจฺฉ ภนฺเต สํโฆ ตํ ปติมาเนตี”ติ ปิฏฺฐิยํ คเหตฺวา ปณาเมสิ. โส ปปติตฺวา กาลมกาสิ. ตสฺส กุกฺกุจฺจํ อโหสิ ฯเปฯ. กิํจิตฺโต ตฺวํ ภิกฺขูติ. นาหํ ภควา มรณา\-ธิปฺปาโยติ. อนาปตฺติ ภิกฺขุ นมรณา\-ธิปฺปายสฺสาติ. (5)}

\prose{เตน โข ปน สมเยน ปิตาปุตฺตา ภิกฺขูสุ ปพฺพชิตา โหนฺติ. กาเล อาโรจิเต ปุตฺโต ปิตรํ เอตทโวจ “คจฺฉ ภนฺเต สํโฆ ตํ ปติมาเนตี”ติ มรณาธิปฺปาโย ปิฏฺฐิยํ คเหตฺวา ปณาเมสิ. โส ปปติตฺวา กาลมกาสิ. ตสฺส กุกฺกุจฺจํ อโหสิ ฯเปฯ. อาปตฺติํ ตฺวํ ภิกฺขุ อาปนฺโน ปาราชิกนฺติ. (6)}

\prose{เตน โข ปน สมเยน ปิตาปุตฺตา ภิกฺขูสุ ปพฺพชิตา โหนฺติ. กาเล อาโรจิเต ปุตฺโต ปิตรํ เอตทโวจ “คจฺฉ ภนฺเต สํโฆ ตํ ปติมาเนตี”ติ มรณาธิปฺปาโย ปิฏฺฐิยํ คเหตฺวา ปณาเมสิ. โส ปปติตฺวา กาลมกาสิ. ตสฺส กุกฺกุจฺจํ อโหสิ ฯเปฯ. อนาปตฺติ ภิกฺขุ ปาราชิกสฺส, อาปตฺติ ถุลฺลจฺจยสฺสาติ. (7)}

\proseitem{181}{เตน โข ปน สมเยน อญฺญตรสฺส ภิกฺขุโน ภุญฺชนฺตสฺส มํสํ กณฺเฐ วิลคฺคํ โหติ. อญฺญตโร ภิกฺขุ ตสฺส ภิกฺขุโน คีวายํ ปหารํ อทาสิ. สโลหิตํ มํสํ ปติ. โส ภิกฺขุ กาลมกาสิ. ตสฺส กุกฺกุจฺจํ อโหสิ ฯเปฯ. อนาปตฺติ ภิกฺขุ นมรณา\-ธิปฺปายสฺสาติ. (8)}

\prose{เตน โข ปน สมเยน อญฺญตรสฺส ภิกฺขุโน ภุญฺชนฺตสฺส มํสํ กณฺเฐ วิลคฺคํ โหติ. อญฺญตโร ภิกฺขุ มรณาธิปฺปาโย ตสฺส ภิกฺขุโน คีวายํ ปหารํ อทาสิ. สโลหิตํ มํสํ ปติ. โส ภิกฺขุ \csromanfolio{102}กาลมกาสิ. ตสฺส กุกฺกุจฺจํ อโหสิ ฯเปฯ. อาปตฺติํ ตฺวํ ภิกฺขุ อาปนฺโน ปาราชิกนฺติ. (9)}

\prose{เตน โข ปน สมเยน อญฺญตรสฺส ภิกฺขุโน ภุญฺชนฺตสฺส มํสํ กณฺเฐ วิลคฺคํ โหติ. อญฺญตโร ภิกฺขุ มรณาธิปฺปาโย ตสฺส ภิกฺขุโน คีวายํ ปหารํ อทาสิ. สโลหิตํ มํสํ ปติ. โส ภิกฺขุ กาลมกาสิ. ตสฺส กุกฺกุจฺจํ อโหสิ ฯเปฯ. อนาปตฺติ ภิกฺขุ ปาราชิกสฺส, อาปตฺติ ถุลฺลจฺจยสฺสาติ. (10)}

\prose{เตน โข ปน สมเยน อญฺญตโร ปิณฺฑจาริโก ภิกฺขุ วิสคตํ ปิณฺฑปาตํ ลภิตฺวา ปฏิกฺกมนํ หริตฺวา ภิกฺขูนํ อคฺคการิกํ อทาสิ. เต ภิกฺขู กาลมกํสุ. ตสฺส กุกฺกุจฺจํ อโหสิ ฯเปฯ. กิํจิตฺโต ตฺวํ ภิกฺขูติ. นาหํ ภควา ชานามีติ. อนาปตฺติ ภิกฺขุ อชานนฺตสฺสาติ. (11)}

\prose{เตน โข ปน สมเยน อญฺญตโร ภิกฺขุ วีมํสา\-ธิปฺปาโย อญฺญตรสฺส ภิกฺขุโน วิสํ อทาสิ. โส ภิกฺขุ กาลมกาสิ. ตสฺส กุกฺกุจฺจํ อโหสิ ฯเปฯ. กิํจิตฺโต ตฺวํ ภิกฺขูติ. วิมํสาธิปฺปาโย อหํ ภควาติ. อนาปตฺติ ภิกฺขุ ปาราชิกสฺส, อาปตฺติ ถุลฺลจฺจยสฺสาติ. (12)}

\proseitem{182}{เตน โข ปน สมเยน อาฬวกา\csromansharedfootnote{280dfb87d1c4a0c6}{อาฬวิกา (สฺยา)} ภิกฺขู วิหารวตฺถุํ กโรนฺติ. อญฺญตโร ภิกฺขุ เหฏฺฐา หุตฺวา สิลํ อุจฺจาเรสิ. อุปริเมน ภิกฺขุนา ทุคฺคหิตา สิลา เหฏฺฐิมสฺส ภิกฺขุโน มตฺถเก อวตฺถาสิ. โส ภิกฺขุ กาลมกาสิ. ตสฺส กุกฺกุจฺจํ อโหสิ ฯเปฯ. อนาปตฺติ ภิกฺขุ อสญฺจิจฺจาติ. (13)}

\prose{เตน โข ปน สมเยน อาฬวกา ภิกฺขุ วิหารวตฺถุํ กโรนฺติ. อญฺญตโร ภิกฺขุ เหฏฺฐา หุตฺวา สิลํ อุจฺจาเรสิ. อุปริโม ภิกฺขุ มรณาธิปฺปาโย เหฏฺฐิมสฺส ภิกฺขุโน มตฺถเก สิลํ มุญฺจิ. โส ภิกฺขุ กาลมกาสิ ฯเปฯ. โส ภิกฺขุ น กาลมกาสิ. ตสฺส กุกฺกุจฺจํ อโหสิ ฯเปฯ. อนาปตฺติ ภิกฺขุ ปาราชิกสฺส, อาปตฺติ ถุลฺลจฺจยสฺสติ. \mbox{(14-15)}}

\prose{\csromanfolio{103}เตน โข ปน สมเยน อาฬวกา ภิกฺขู วิหารสฺส กุฏฺฏํ อุฏฺฐาเปนฺติ. อญฺญตโร ภิกฺขุ เหฏฺฐา หุตฺวา อิฏฺฐกํ อุจฺจาเรสิ. อุปริเมน ภิกฺขุนา ทุคฺคหิตา อิฏฺฐกา เหฏฺฐิมสฺส ภิกฺขุโน มตฺถเก อวตฺถาสิ. โส ภิกฺขุ กาลมกาสิ. ตสฺส กุกฺกุจฺจํ อโหสิ ฯเปฯ. อนาปตฺติ ภิกฺขุ อสญฺจิจฺจาติ. (16)}

\prose{เตน โข ปน สมเยน อาฬวกา ภิกฺขู วิหารสฺส กุฏฺฏํ อุฏฺฐาเปนฺติ. อญฺญตโร ภิกฺขุ เหฏฺฐา หุตฺวา อิฏฺฐกํ อุจฺจาเรสิ. อุปริโม ภิกฺขุ มรณาธิปฺปาโย เหฏฺฐิมสฺส ภิกฺขุโน มตฺถเก อิฏฺฐกํ มุญฺจิ. โส ภิกฺขุ กาลมกาสิ ฯเปฯ. โส ภิกฺขุ น กาลมกาสิ. ตสฺส กุกฺกุจฺจํ อโหสิ ฯเปฯ. อนาปตฺติ ภิกฺขุ ปาราชิกสฺส, อาปตฺติ ถุลฺลจฺจยสฺสาติ. \mbox{(17-18)}}

\proseitem{183}{เตน โข ปน สมเยน อาฬวกา ภิกฺขู นวกมฺมํ กโรนฺติ. อญฺญตโร ภิกฺขุ เหฏฺฐา หุตฺวา วาสิํ อุจฺจาเรสิ. อุปริเมน ภิกฺขุนา ทุคฺคหิตา วาสี เหฏฺฐิมสฺส ภิกฺขุโน มตฺถเก อวตฺถาสิ. โส ภิกฺขุ กาลมกาสิ. ตสฺส กุกฺกุจฺจํ อโหสิ ฯเปฯ. อนาปตฺติ ภิกฺขุ อสญฺจิจฺจาติ. (19)}

\prose{เตน โข ปน สมเยน อาฬวกา ภิกฺขู นวกมฺมํ กโรนฺติ. อญฺญตโร ภิกฺขุ เหฏฺฐา หุตฺวา วาสิํ อุจฺจาเรสิ. อุปริโม ภิกฺขุ มรณาธิปฺปาโย เหฏฺฐิมสฺส ภิกฺขุโน มตฺถเก วาสิํ มุญฺจิ. โส ภิกฺขุ กาลมกาสิ ฯเปฯ. โส ภิกฺขุ น กาลมกาสิ. ตสฺส กุกฺกุจฺจํ อโหสิ ฯเปฯ. อนาปตฺติ ภิกฺขุ ปาราชิกสฺส, อาปตฺติ ถุลฺลจฺจยสฺสาติ. \mbox{(20-21)}}

\prose{เตน โข ปน สมเยน อาฬวกา ภิกฺขู นวกมฺมํ กโรนฺติ. อญฺญตโร ภิกฺขุ เหฏฺฐา หุตฺวา โคปานสิํ อุจฺจาเรสิ. อุปริเมน ภิกฺขุนา ทุคฺคหิตา โคปานสี เหฏฺฐิมสฺส ภิกฺขุโน มตฺถเก อวตฺถาสิ. โส ภิกฺขุ กาลมกาสิ. ตสฺส กุกฺกุจฺจํ อโหสิ ฯเปฯ. อนาปตฺติ ภิกฺขุ อสญฺจิจฺจาติ. (22)}

\prose{เตน โข ปน สมเยน อาฬวกา ภิกฺขู นวกมฺมํ กโรนฺติ. อญฺญตโร ภิกฺขุ เหฏฺฐา หุตฺวา โคปานสิํ อุจฺจาเรสิ. อุปริโม ภิกฺขุ มรณาธิปฺปาโย เหฏฺฐิมสฺส ภิกฺขุโน มตฺถเก โคปานสิํ มุญฺจิ. \csromanfolio{104}โส ภิกฺขุ กาลมกาสิ ฯเปฯ. โส ภิกฺขุ น กาลมกาสิ. ตสฺส กุกฺกุจฺจํ อโหสิ ฯเปฯ. อนาปตฺติ ภิกฺขุ ปาราชิกสฺส, อาปตฺติ ถุลฺลจฺจยสฺสาติ. \mbox{(23-24)}}

\prose{เตน โข ปน สมเยน อาฬวกา ภิกฺขู นวกมฺมํ กโรนฺตา อฏฺฏกํ พนฺธนฺติ. อญฺญตโร ภิกฺขุ อญฺญตรํ ภิกฺขุํ เอตทโวจ “อาวุโส อตฺรฏฺฐิโต พนฺธาหี”ติ. โส ตตฺรฏฺฐิโต พนฺธนฺโต ปริปติตฺวา กาลมกาสิ. ตสฺส กุกฺกุจฺจํ อโหสิ ฯเปฯ. กิํจิตฺโต ตฺวํ ภิกฺขูติ. นาหํ ภควา มรณธิปฺปาโยติ. อนาปตฺติ ภิกฺขุ นมรณา\-ธิปฺปายสฺสาติ. (25)}

\prose{เตน โข ปน สมเยน อาฬวกา ภิกฺขู นวกมฺมํ กโรนฺตา อฏฺฏกํ พนฺธนฺติ. อญฺญตโร ภิกฺขุ มรณาธิปฺปาโย อญฺญตรํ ภิกฺขุํ เอตทโวจ “อาวุโส อตฺรฏฺฐิโต พนฺธาหี”ติ. โส ตตฺรฏฺฐิโต พนฺธนฺโต ปริปติตฺวา กาลมกาสิ. ตสฺส กุกฺกุจฺจํ อโหสิ ฯเปฯ ปริปติตฺวา น กาลมกาสิ. ตสฺส กุกฺกุจฺจํ อโหสิ ฯเปฯ. อนาปตฺติ ภิกฺขุ ปาราชิกสฺส, อาปตฺติ ถุลฺลจฺจยสฺสาติ. \mbox{(26-27)}}

\prose{เตน โข ปน สมเยน อญฺญตโร ภิกฺขุ วิหารํ ฉาเทตฺวา โอตรติ. อญฺญตโร ภิกฺขุ ตํ ภิกฺขุํ เอตทโวจ “อาวุโส อิโต โอตราหี”ติ. โส เตน โอตรนฺโต ปริปติตฺวา กาลมกาสิ. ตสฺส กุกฺกุจฺจํ อโหสิ ฯเปฯ. อนาปตฺติ ภิกฺขุ นมรณา\-ธิปฺปายสฺสาติ. (28)}

\prose{เตน โข ปน สมเยน อญฺญตโร ภิกฺขุ วิหารํ ฉาเทตฺวา โอตรติ. อญฺญตโร ภิกฺขุ มรณาธิปฺปาโย ตํ ภิกฺขุํ เอตทโวจ “อาวุโส อิโต โอตราหี”ติ. โส เตน โอตรนฺโต ปริปติตฺวา กาลมกาสิ ฯเปฯ ปริปติตฺวา น กาลมกาสิ. ตสฺส กุกฺกุจฺจํ อโหสิ ฯเปฯ. อนาปตฺติ ภิกฺขุ ปาราชิกสฺส, อาปตฺติ ถุลฺลจฺจยสฺสาติ. \mbox{(29-30)}}

\prose{เตน โข ปน สมเยน อญฺญตโร ภิกฺขุ อนภิรติยา ปีฬิโต คิชฺฌกูฏํ ปพฺพตํ อภิรุหิตฺวา ปปาเต ปปตนฺโต อญฺญตรํ วิลีวการํ โอตฺถริตฺวา มาเรสิ. ตสฺส กุกฺกุจฺจํ อโหสิ ฯเปฯ. อนาปตฺติ ภิกฺขุ ปาราชิกสฺส, น จ ภิกฺขเว อตฺตานํ ปาเตตพฺพํ, โย ปาเตยฺย, อาปตฺติ ทุกฺกฏสฺสาติ. (31)}

\prose{\csromanfolio{105}เตน โข ปน สมเยน ฉพฺพคฺคิยา ภิกฺขู คิชฺฌกูฏํ ปพฺพตํ อภิรุหิตฺวา ทวาย สิลํ ปวิชฺฌิํสุ, สา อญฺญตรํ โคปาลกํ โอตฺถริตฺวา มาเรสิ. เตสํ กุกฺกุจฺจํ อโหสิ ฯเปฯ. อนาปตฺติ ภิกฺขเว ปาราชิกสฺส, น จ ภิกฺขเว ทวาย สิลา ปวิชฺฌิตพฺพา, โย ปวิชฺเฌยฺย, อาปตฺติ ทุกฺกฏสฺสาติ. (32)}

\proseitem{184}{เตน โข ปน สมเยน อญฺญตโร ภิกฺขุ คิลาโน โหติ. ตํ ภิกฺขู เสเทสุํ. โส ภิกฺขุ กาลมกาสิ. เตสํ กุกฺกุจฺจํ อโหสิ ฯเปฯ. อนาปตฺติ ภิกฺขเว นมรณา\-ธิปฺปายสฺสาติ. (33)}

\prose{เตน โข ปน สมเยน อญฺญตโร ภิกฺขุ คิลาโน โหติ. ตํ ภิกฺขู มรณาธิปฺปาย เสเทสุํ. โส ภิกฺขุ กาลมกาสิ ฯเปฯ. โส ภิกฺขุ น กาลมกาสิ. เตสํ กุกฺกุจฺจํ อโหสิ ฯเปฯ. อนาปตฺติ ภิกฺขเว ปาราชิกสฺส, อาปตฺติ ถุลฺลจฺจยสฺสาติ. \mbox{(34-35)}}

\prose{เตน โข ปน สมเยน อญฺญตรสฺส ภิกฺขุโน สีสาภิตาโป โหติ. ตสฺส ภิกฺขู นตฺถุํ อทํสุ. โส ภิกฺขุ กาลมกาสิ. เตสํ กุกฺกุจฺจํ อโหสิ ฯเปฯ. อนาปตฺติ ภิกฺขเว นมรณา\-ธิปฺปายสฺสาติ. (36)}

\prose{เตน โข ปน สมเยน อญฺญตรสฺส ภิกฺขุโน สีสาภิตาโป โหติ. ตสฺส ภิกฺขู มรณาธิปฺปายา นตฺถุํ อทํสุ. โส ภิกฺขุ กาลมกาสิ ฯเปฯ. โส ภิกฺขุ น กาลมกาสิ. เตสํ กุกฺกุจฺจํ อโหสิ ฯเปฯ. อนาปตฺติ ภิกฺขเว ปาราชิกสฺส, อาปตฺติ ถุลฺลจฺจยสฺสาติ. \mbox{(37-38)}}

\prose{เตน โข ปน สมเยน อญฺญตโร ภิกฺขุ คิลาโน โหติ. ตํ ภิกฺขู สมฺพาเหสุํ. โส ภิกฺขุ กาลมกาสิ. เตสํ กุกฺกุจฺจํ อโหสิ ฯเปฯ. อนาปตฺติ ภิกฺขเว นมรณา\-ธิปฺปายสฺสาติ. (39)}

\prose{เตน โข ปน สมเยน อญฺญตโร ภิกฺขุ คิลาโน โหติ. ตํ ภิกฺขู มรณาธิปฺปายา สมฺพาเหสุํ. โส ภิกฺขุ กาลมกาสิ ฯเปฯ. โส ภิกฺขุ น กาลมกาสิ. เตสํ กุกฺกุจฺจํ อโหสิ ฯเปฯ. อนาปตฺติ ภิกฺขเว ปาราชิกสฺส, อาปตฺติ ถุลฺลจฺจยสฺสาติ. \mbox{(40-41)}}

\prose{เตน โข ปน สมเยน อญฺญตโร ภิกฺขุ คิลาโน โหติ. ตํ ภิกฺขู นฺหาเปสุํ. โส ภิกฺขุ กาลมกาสิ. เตสํ กุกฺกุจฺจํ อโหสิ ฯเปฯ. อนาปตฺติ ภิกฺขเว นมรณา\-ธิปฺปายสฺสาติ. (42)}

\prose{\csromanfolio{106}เตน โข ปน สมเยน อญฺญตโร ภิกฺขุ คิลาโน โหติ. ตํ ภิกฺขู มรณาธิปฺปายา นฺหาเปสุํ. โส ภิกฺขุ กาลมกาสิ ฯเปฯ. โส ภิกฺขุ น กาลมกาสิ. เตสํ กุกฺกุจฺจํ อโหสิ ฯเปฯ. อนาปตฺติ ภิกฺขเว ปาราชิกสฺส, อาปตฺติ ถุลฺลจฺจยสฺสาติ. \mbox{(43-44)}}

\prose{เตน โข ปน สมเยน อญฺญตโร ภิกฺขุ คิลาโน โหติ. ตํ ภิกฺขู เตเลน อพฺภญฺชิํสุ. โส ภิกฺขุ กาลมกาสิ. เตสํ กุกฺกุจฺจํ อโหสิ ฯเปฯ. อนาปตฺติ ภิกฺขเว นมรณา\-ธิปฺปายสฺสาติ. (45)}

\prose{เตน โข ปน สมเยน อญฺญตโร ภิกฺขุ คิลาโน โหติ. ตํ ภิกฺขู มรณาธิปฺปายา เตเลน อพฺภิญฺชิํสุ. โส ภิกฺขุ กาลมกาสิ ฯเปฯ. โส ภิกฺขุ น กาลมกาสิ. เตสํ กุกฺกุจฺจํ อโหสิ ฯเปฯ. อนาปตฺติ ภิกฺขเว ปาราชิกสฺส, อาปตฺติ ถุลฺลจฺจยสฺสาติ. \mbox{(46-47)}}

\proseitem{185}{เตน โข ปน สมเยน อญฺญตโร ภิกฺขุ คิลาโน โหติ. ตํ ภิกฺขู อุฏฺฐาเปสุํ. โส ภิกฺขุ กาลมกาสิ. เตสํ กุกฺกุจฺจํ อโหสิ ฯเปฯ. อนาปตฺติ ภิกฺขเว นมรณา\-ธิปฺปายสฺสาติ. (48)}

\prose{เตน โข ปน สมเยน อญฺญตโร ภิกฺขุ คิลาโน โหติ. ตํ ภิกฺขู มรณาธิปฺปายา อุฏฺฐาเปสุํ. โส ภิกฺขุ กาลมกาสิ ฯเปฯ. โส ภิกฺขุ น กาลมลาสิ. เตสํ กุกฺกุจฺจํ อโหสิ ฯเปฯ. อนาปตฺติ ภิกฺขเว ปาราชิกสฺส, อาปตฺติ ถุลฺลจฺจยสฺสาติ. \mbox{(49-50)}}

\prose{เตน โข ปน สมเยน อญฺญตโร ภิกฺขุ คิลาโน โหติ. ตํ ภิกฺขู นิปาเตสุํ. โส ภิกฺขุ กาลมกาสิ. เตสํ กุกฺกุจฺจํ อโหสิ ฯเปฯ. อนาปตฺติ ภิกฺขเว นมรณา\-ธิปฺปายสฺสาติ. (51)}

\prose{เตน โข ปน สมเยน อญฺญตโร ภิกฺขุ คิลาโน โหติ. ตํ ภิกฺขู มรณาธิปฺปายา นิปาเตสุํ. โส ภิกฺขุ กาลมกาสิ ฯเปฯ. โส ภิกฺขุ น กาลมกาสิ. เตสํ กุกฺกุจฺจํ อโหสิ ฯเปฯ. อนาปตฺติ ภิกฺขเว ปาราชิกสฺส, อาปตฺติ ถุลฺลจฺจยสฺสาติ. \mbox{(52-53)}}

\prose{เตน โข ปน สมเยน อญฺญตโร ภิกฺขุ คิลาโน โหติ. ตสฺส ภิกฺขู อนฺนํ อทํสุ. โส ภิกฺขุ กาลมกาสิ. เตสํ กุกฺกุจฺจํ อโหสิ ฯเปฯ. อนาปตฺติ ภิกฺขเว นมรณา\-ธิปฺปายสฺสาติ. (54)}

\prose{\csromanfolio{107}เตน โข ปน สมเยน อญฺญตโร ภิกฺขุ คิลาโน โหติ. ตสฺส ภิกฺขู มรณาธิปฺปายา อนฺนํ อทํสุ. โส ภิกฺขุ กาลมกาสิ ฯเปฯ. โส ภิกฺขุ น กาลมกาสิ. เตสํ กุกฺกุจฺจํ อโหสิ ฯเปฯ. อนาปตฺติ ภิกฺขเว ปาราชิกสฺส, อาปตฺติ ถุลฺลจฺจยสฺสาติ. \mbox{(55-56)}}

\prose{เตน โข ปน สมเยน อญฺญตโร ภิกฺขุ คิลาโน โหติ. ตสฺส ภิกฺขู ปานํ อทํสุ. โส ภิกฺขุ กาลมกาสิ. เตสํ กุกฺกุจฺจํ อโหสิ ฯเปฯ. อนาปตฺติ ภิกฺขเว นมรณา\-ธิปฺปายสฺสาติ. (57)}

\prose{เตน โข ปน สมเยน อญฺญตโร ภิกฺขุ คิลาโน โหติ. ตสฺส ภิกฺขู มรณาธิปฺปายา ปานํ อทํสุ. โส ภิกฺขุ กาลมกาสิ ฯเปฯ. โส ภิกฺขุ น กาลมกาสิ. เตสํ กุกฺกุจฺจํ อโหสิ ฯเปฯ. อนาปตฺติ ภิกฺขเว ปาราชิกสฺส, อาปตฺติ ถุลฺลจฺจยสฺสาติ. \mbox{(58-59)}}

\proseitem{186}{เตน โข ปน สมเยน อญฺญตรา อิตฺถี ปวุตฺถปติกา ชาเรน คพฺภินี โหติ. สา กุลูปกํ ภิกฺขุํ เอตทโวจ “อิงฺฆายฺย คพฺภปาตนํ ชานาหี”ติ. “สุฏฺฐุ ภคินี”ติ ตสฺสา คพฺภปาตนํ อทาสิ. ทารโก กาลมกาสิ. ตสฺส กุกฺกุจฺจํ อโหสิ ฯเปฯ. อาปตฺติํ ตฺวํ ภิกฺขุ อาปนฺโน ปาราชิกนฺติ. (60)}

\prose{เตน โข ปน สมเยน อญฺญตรสฺส ปุริสสฺส เทฺว ปชาปติโย โหนฺติ เอกา วญฺฌา เอกา วิชายินี. วญฺฌา อิตฺถี กุลูปกํ ภิกฺขุํ เอตทโวจ “สเจ สา ภนฺเต วิชายิสฺสติ, สพฺพสฺส กุฏุมฺพสฺส อิสฺสรา ภวิสฺสติ, อิงฺฆายฺย ตสฺสา คพฺภปาตนํ ชานาหี”ติ. “สุฏฺฐุ ภคินี”ติ ตสฺสา คพฺภปาตนํ อทาสิ. ทารโก กาลมกาสิ. มาตา น กาลมกาสิ. ตสฺส กุกฺกุจฺจํ อโหสิ ฯเปฯ. อาปตฺติํ ตฺวํ ภิกฺขุ อาปนฺโน ปาราชิกนฺติ. (61)}

\prose{เตน โข ปน สมเยน อญฺญตรสฺส ปุริสสฺส เทฺว ปชาปติโย โหนฺติ เอกา วญฺฌา เอกา วิชายินี. วญฺฌา อิตฺถี กุลูปกํ ภิกฺขุํ เอตทโวจ “สเจ สา ภนฺเต วิชายิสฺสติ, สพฺพสฺส กุฏุมฺพสฺส อิสฺสรา ภวิสฺสติ, อิงฺฆายฺย ตสฺสา คพฺภปาตนํ ชานาหี”ติ. “สุฏฺฐุ ภคินี”ติ ตสฺสา คพฺภปาตนํ อาทาสิ. มาตา กาลมกาสิ. ทารโก น กาลมกาสิ. ตสฺส กุกฺกุจฺจํ อโหสิ ฯเปฯ. อนาปตฺติ ภิกฺขุ ปาราชิกสฺส, อาปตฺติ ถุลฺลจฺจยสฺสาติ. (62)}

\prose{\csromanfolio{108}เตน โข ปน สมเยน อญฺญตรสฺส ปุริสสฺส เทฺว ปชาปติโย โหนฺติ เอกา วญฺฌา เอกา วิชายินี. วญฺฌา อิตฺถี กุลูปกํ ภิกฺขุํ เอตทโวจ “สเจ สา ภนฺเต วิชายิสฺสติ, สพฺพสฺส กุฏุมฺพสฺส อิสฺสรา ภวิสฺสติ, อิงฺฆายฺย ตสฺสา คพฺภปาตนํ ชานาหี”ติ. “สุฏฺฐุ ภคินี”ติ ตสฺสา คพฺภปาตนํ อทาสิ. อุโภ กาลมกํสุ ฯเปฯ. อุโภ น กาลมกํสุ. ตสฺส กุกฺกุจฺจํ อโหสิ ฯเปฯ. อนาปตฺติ ภิกฺขุ ปาราชิกสฺส, อาปตฺติ ถุลฺลจฺจยสฺสาติ. \mbox{(63-64)}}

\proseitem{187}{เตน โข ปน สมเยน อญฺญตรา คพฺภินี อิตฺถี กุลูปกํ ภิกฺขุํ เอตทโวจ “อิงฺฆายฺย คพฺภปาตนํ ชานาหี”ติ. เตน หิ ภคินิ มทฺทสฺสูติ. สา มทฺทาเปตฺวา คพฺภํ ปาเตสิ. ตสฺส กุกฺกุจฺจํ อโหสิ ฯเปฯ. อาปตฺติํ ตฺวํ ภิกฺขุ อาปนฺโน ปาราชิกนฺติ. (65)}

\prose{เตน โข ปน สมเยน อญฺญตรา คพฺภินี อิตฺถี กุลูปกํ ภิกฺขุํ เอตทโวจ “อิงฺฆายฺย คพฺภปาตนํ ชานาหี”ติ. เตน หิ ภคินิ ตาเปหีติ. สา ตาเปตฺวา คพฺภํ ปาเตสิ. ตสฺส กุกฺกุจฺจํ อโหสิ ฯเปฯ. อาปตฺติํ ตฺวํ ภิกฺขุ อาปนฺโน ปาราชิกนฺติ. (66)}

\prose{เตน โข ปน สมเยน อญฺญตรา วญฺฌา อิตฺถี กุลูปกํ ภิกฺขุํ เอตทโวจ “อิงฺฆายฺย เภสชฺชํ ชานาหิ, เยนาหํ วิชาเยยฺยนฺ”ติ. “สุฏฺฐุ ภคินี”ติ ตสฺสา เภสชฺชํ อทาสิ. สา กาลมกาสิ. ตสฺส กุกฺกุจฺจํ อโหสิ ฯเปฯ. อนาปตฺติ ภิกฺขุ ปาราชิกสฺส, อาปตฺติ ทุกฺกฏสฺสาติ. (67)}

\prose{เตน โข ปน สมเยน อญฺญตรา วิชายินี อิตฺถี กุลูปกํ ภิกฺขุํ เอตทโวจ “อิงฺฆายฺย เภสชฺชํ ชานาหิ, เยนาหํ วิชาเยยฺยนฺ”ติ. “สุฏฺฐุ ภคินี”ติ ตสฺสา เภสชฺชํ อทาสิ. สา กาลมกาสิ. ตสฺส กุกฺกุจฺจํ อโหสิ ฯเปฯ. อนาปตฺติ ภิกฺขุ ปาราชิกสฺส, อาปตฺติ ทุกฺกฏสฺสาติ. (68)}

\prose{เตน โข ปน สมเยน ฉพฺพคฺคิยา ภิกฺขู สตฺตรส\-วคฺคิยํ ภิกฺขุํ องฺคุลิ\-ปโตทเกน หาเสสุํ. โส ภิกฺขุ อุตฺตนฺโต อนสฺสาสโก กาลมกาสิ. เตสํ กุกฺกุจฺจํ อโหสิ ฯเปฯ. อนาปตฺติ ภิกฺขเว ปาราชิกสฺสาติ\csromansharedfootnote{f0ab2658c5a86ed5}{ปาราชิกสฺส, อาปตฺติ ปาจิตฺติยสฺสาติ (สฺยา)}. (69)}

\prose{\csromanfolio{109}เตน โข ปน สมเยน สตฺตรส\-วคฺคิยา ภิกฺขู ฉพฺพคฺคิยํ ภิกฺขุํ กมฺมํ กริสฺสามาติ โอตฺถริตฺวา มาเรสุํ. เตสํ กุกฺกุจฺจํ อโหสิ ฯเปฯ. อนาปตฺติ ภิกฺขเว ปาราชิกสฺสาติ. (70)}

\prose{เตน โข ปน สมเยน อญฺญตโร ภูตเวชฺชโก ภิกฺขุ ยกฺขํ ชีวิตา โวโรเปสิ. ตสฺส กุกฺกุจฺจํ อโหสิ ฯเปฯ. อนาปตฺติ ภิกฺขุ ปาราชิกสฺส, อาปตฺติ ถุลฺลจฺจยสฺสาติ. (71)}

\prose{เตน โข ปน สมเยน อญฺญตโร ภิกฺขุ อญฺญตรํ ภิกฺขุํ วาฬยกฺข\-วิหารํ ปาเหสิ. ตํ ยกฺขา ชีวิตา โวโรเปสุํ. ตสฺส กุกฺกุจฺจํ อโหสิ ฯเปฯ. อนาปตฺติ ภิกฺขุ นมรณา\-ธิปฺปายสฺสาติ. (72)}

\prose{เตน โข ปน สมเยน อญฺญตโร ภิกฺขุ มรณาธิปฺปาโย อญฺญตรํ ภิกฺขุํ วาฬยกฺข\-วิหารํ ปาเหสิ. ตํ ยกฺขา ชีวิตา โวโรเปสุํ ฯเปฯ. ตํ ยกฺขา ชีวิตา น โวโรเปสุํ. ตสฺส กุกฺกุจฺจํ อโหสิ ฯเปฯ. อนาปตฺติ ภิกฺขุ ปาราชิกสฺส, อาปตฺติ ถุลฺลจฺจยสฺสาติ. \mbox{(73-74)}}

\prose{เตน โข ปน สมเยน อญฺญตโร ภิกฺขุ อญฺญตรํ ภิกฺขุํ วาฬกนฺตารํ ปาเหสิ. ตํ วาฬา ชีวิตา โวโรเปสุํ. ตสฺส กุกฺกุจฺจํ อโหสิ ฯเปฯ. อนาปตฺติ ภิกฺขุ นมรณา\-ธิปฺปายสฺสาติ. (75)}

\prose{เตน โข ปน สมเยน อญฺญตโร ภิกฺขุ มรณาธิปฺปาโย อญฺญตรํ ภิกฺขุํ วาฬกนฺตารํ ปาเหสิ. ตํ วาฬา ชีวิตา โวโรเปสุํ ฯเปฯ. ตํ วาฬา ชีวิตา น โวโรเปสุํ. ตสฺส กุกฺกุจฺจํ อโหสิ ฯเปฯ. อนาปตฺติ ภิกฺขุ ปาราชิกสฺส, อาปตฺติ ถุลฺลจฺจยสฺสาติ. \mbox{(76-77)}}

\prose{เตน โข ปน สมเยน อญฺญตโร ภิกฺขุ อญฺญตรํ ภิกฺขุํ โจรกนฺตารํ ปาเหสิ. ตํ โจรา ชีวิตา โวโรเปสุํ. ตสฺส กุกฺกุจฺจํ อโหสิ ฯเปฯ. อนาปตฺติ ภิกฺขุ นมรณา\-ธิปฺปายสฺสาติ. (78)}

\prose{เตน โข ปน สมเยน อญฺญตโร ภิกฺขุ มรณาธิปฺปาโย อญฺญตรํ ภิกฺขุํ โจรกนฺตารํ ปาเหสิ. ตํ โจรา ชีวิตา โวโรเปสุํ ฯเปฯ. ตํ โจรา ชีวิตา น โวโรเปสุํ. ตสฺส กุกฺกุจฺจํ อโหสิ ฯเปฯ. อนาปตฺติ ภิกฺขุ ปาราชิกสฺส, อาปตฺติ ถุลฺลจฺจยสฺสาติ. \mbox{(79-80)}}

\proseitem{188}{เตน โข ปน สมเยน อญฺญตโร ภิกฺขุ ตํ มญฺญมาโน ตํ ชีวิตา โวโรเปสิ ฯเปฯ ตํ มญฺญมาโน อญฺญํ ชีวิตา โวโรเปสิ ฯเปฯ \csromanfolio{110}อญฺญํ มญฺญมาโน ตํ ชีวิตา โวโรเปสิ ฯเปฯ อญฺญํ มญฺญมาโน อญฺญํ ชีวิตา โวโรเปสิ. ตสฺส กุกฺกุจฺจํ อโหสิ ฯเปฯ. อาปตฺติํ ตฺวํ ภิกฺขุ อาปนฺโน ปาราชิกนฺติ. \mbox{(81-84)}}

\prose{เตน โข ปน สมเยน อญฺญตโร ภิกฺขุ อมนุสฺเสน คหิโต โหติ. อญฺญตโร ภิกฺขุ ตสฺส ภิกฺขุโน ปหารํ อทาสิ. โส ภิกฺขุ กาลมกาสิ. ตสฺส กุกฺกุจฺจํ อโหสิ ฯเปฯ. อนาปตฺติ ภิกฺขุ นมรณา\-ธิปฺปายสฺสาติ. (85)}

\prose{เตน โข ปน สมเยน อญฺญตโร ภิกฺขุ อมนุสฺเสน คหิโต โหติ. อญฺญตโร ภิกฺขุ มรณาธิปฺปาโย ตสฺส ภิกฺขุโน ปหารํ อทาสิ. โส ภิกฺขุ น กาลมกาสิ ฯเปฯ. โส ภิกฺขุ กาลมกาสิ ฯเปฯ. ตสฺส กุกฺกุจฺจํ อโหสิ ฯเปฯ. อนาปตฺติ ภิกฺขุ ปาราชิกสฺส, อาปตฺติ ถุลฺลจฺจยสฺสาติ. \mbox{(86-87)}}

\prose{เตน โข ปน สมเยน อญฺญตโร ภิกฺขุ กลฺยาณ\-กมฺมสฺส สคฺคกถํ กเถสิ. โส อธิมุตฺโต กาลมกาสิ. ตสฺส กุกฺกุจฺจํ อโหสิ ฯเปฯ. อนาปตฺติ ภิกฺขุ นมรณา\-ธิปฺปายสฺสาติ. (88)}

\prose{เตน โข ปน สมเยน อญฺญตโร ภิกฺขุ มรณาธิปฺปาโย กลฺยาณ\-กมฺมสฺส สคฺคกถํ กเถสิ. โส อธิมุตฺโต กาลมกาสิ ฯเปฯ. โส อธิมุตฺโต น กาลมกาสิ. ตสฺส กุกฺกุจฺจํ อโหสิ ฯเปฯ. อนาปตฺติ ภิกฺขุ ปาราชิกสฺส, อาปตฺติ ถุลฺลจฺจยสฺสาติ. \mbox{(89-90)}}

\prose{เตน โข ปน สมเยน อญฺญตโร ภิกฺขุ เนรยิกสฺส นิรยกถํ กเถสิ. โส อุตฺตสิตฺวา กาลมกาสิ. ตสฺส กุกฺกุจฺจํ อโหสิ ฯเปฯ. อนาปตฺติ ภิกฺขุ นมรณา\-ธิปฺปายสฺสาติ. (91)}

\prose{เตน โข ปน สมเยน อญฺญตโร ภิกฺขุ มรณาธิปฺปาโย เนรยิกสฺส นิรยกถํ กเถสิ. โส อุตฺตสิตฺวา กาลมกาสิ ฯเปฯ. โส อุตฺตสิตฺวา น กาลมกาสิ. ตสฺส กุกฺกุจฺจํ อโหสิ ฯเปฯ. อนาปตฺติ ภิกฺขุ ปาราชิกสฺส, อาปตฺติ ถุลฺลจฺจยสฺสาติ. \mbox{(92-93)}}

\proseitem{189}{เตน โข ปน สมเยน อาฬวกา ภิกฺขู นวกมฺมํ กโรนฺตา รุกฺขํ ฉินฺทนฺติ. อญฺญตโร ภิกฺขุ อญฺญตรํ ภิกฺขุํ เอตทโวจ “อาวุโส อตฺรฏฺฐิโต ฉินฺทาหี”ติ. ตํ ตตฺรฏฺฐิตํ ฉินฺทนฺตํ รุกฺโข โอตฺถริตฺวา}

\prose{\csromanfolio{111}เตน โข ปน สมเยน อาฬวกา ภิกฺขู นวกมฺมํ กโรนฺตา รุกฺขํ ฉินฺทนฺติ. อญฺญตโร ภิกฺขุ มรณาธิปฺปาโย อญฺญตรํ ภิกฺขุํ เอตทโวจ “อาวุโส อตฺรฏฺฐิโต ฉินฺทาหี”ติ. ตํ ตตฺรฏฺฐิตํ ฉินฺทนฺตํ รุกฺโข โอตฺถริตฺวา มาเรสิ ฯเปฯ รุกฺโข โอตฺถริตฺวา น มาเรสิ. ตสฺส กุกฺกุจฺจํ อโหสิ ฯเปฯ. อนาปตฺติ ภิกฺขุ ปาราชิกสฺส, อาปตฺติ ถุลฺลจฺจยสฺสาติ. \mbox{(95-96)}}

\proseitem{190}{เตน โข ปน สมเยน ฉพฺพคฺคิยา ภิกฺขู ทายํ อาลิมฺเปสุํ\csromansharedfootnote{2eb55cb2e7d9b494}{อาฬิมฺเปสุํ (สฺยา, ก)}. มนุสฺสา ทฑฺฒา กาลมกํสุ. เตสํ กุกฺกุจฺจํ อโหสิ ฯเปฯ. อนาปตฺติ ภิกฺขเว นมรณา\-ธิปฺปายสฺสาติ. (97)}

\prose{เตน โข ปน สมเยน ฉพฺพคฺคิยา ภิกฺขู มรณาธิปฺปายา ทายํ อาลิมฺเปสุํ. มนุสฺสา ทฑฺฒา กาลมกํสุ ฯเปฯ. มนุสฺสา ทฑฺฒา น กาลมกํสุ. เตสํ กุกฺกุจฺจํ อโหสิ ฯเปฯ. อนาปตฺติ ภิกฺขเว ปาราชิกสฺส, อาปตฺติ ถุลฺลจฺจยสฺสาติ. \mbox{(98-99)}}

\proseitem{191}{เตน โข ปน สมเยน อญฺญตโร ภิกฺขุ อาฆาตนํ คนฺตฺวา โจรฆาตํ เอตทโวจ “อาวุโส มายิมํ กิลเมสิ, เอเกน ปหาเรน ชีวิตา โวโรเปหี”ติ. “สุฏฺฐุ ภนฺเต”ติ เอเกน ปหาเรน ชีวิตา โวโรเปสิ. ตสฺส กุกฺกุจฺจํ อโหสิ ฯเปฯ. อาปตฺติํ ตฺวํ ภิกฺขุ อาปนฺโน ปาราชิกนฺติ. (100)}

\prose{เตน โข ปน สมเยน อญฺญตโร ภิกฺขุ อาฆาตนํ คนฺตฺวา โจรฆาตํ เอตทโวจ “อาวุโส มายิมํ กิลเมสิ, เอเกน ปหาเรน ชีวิตา โวโรเปหี”ติ. โส “นาหํ ตุยฺหํ วจนํ กริสฺสามี”ติ ตํ ชีวิตา โวโรเปสิ. ตสฺส กุกฺกุจฺจํ อโหสิ ฯเปฯ. อนาปตฺติ ภิกฺขุ ปาราชิกสฺส, อาปตฺติ ทุกฺกฏสฺสาติ. (101)}

\csromanmatikaanchor{mtk.29}
\csromanmatikaanchor{mtk.30}
\csromantocmarknopage{section}{4. จตุตฺถปาราชิก}{mtk.29}
\bookmark[dest={mtk.29},level=1]{4. จตุตฺถปาราชิก}
\csromantocmarkref{subsection}{วคฺคุมุทาตีริยภิกฺขุวตฺถุ}{mtk.30}
\bookmark[dest={mtk.30},level=2]{วคฺคุมุทาตีริยภิกฺขุวตฺถุ}
\proseitem{192}{เตน โข ปน สมเยน อญฺญตโร ปุริโส ญาติฆเร หตฺถ\-ปาท\-จฺฉินฺโน ญาตเกหิ สํปริกิณฺโณ โหติ. อญฺญตโร ภิกฺขุ เต มนุสฺเส เอตทโวจ “อาวุโส อิจฺฉถ อิมสฺส มรณนฺ”ติ. อาม ภนฺเต อิจฺฉามาติ. เตน หิ ตกฺกํ ปาเยถาติ. เต ตํ ตกฺกํ ปาเยสุํ. \csromanfolio{112}โส กาลมกาสิ. ตสฺส กุกฺกุจฺจํ อโหสิ ฯเปฯ. อาปตฺติํ ตฺวํ ภิกฺขุ อาปนฺโน ปาราชิกนฺติ. (102)}

\prosewithcloserruled{เตน โข ปน สมเยน อญฺญตโร ปุริโส กุลฆเร หตฺถ\-ปาท\-จฺฉินฺโน ญาตเกหิ สํปริกิณฺโณ โหติ. อญฺญตรา ภิกฺขุนี เต มนุสฺเส เอตทโวจ “อาวุโส อิจฺฉถ อิมสฺส มรณนฺ”ติ. อามายฺเย อิจฺฉามาติ. เตน หิ โลณโสวีรกํ ปาเยถาติ. เต ตํ โลณโสวีรกํ ปาเยสุํ. โส กาลมกาสิ. ตสฺสา กุกฺกุจฺจํ อโหสิ ฯเปฯ. อถ โข สา ภิกฺขุนี ภิกฺขุนีนํ เอตมตฺถํ อาโรเจสิ. ภิกฺขุนิโย ภิกฺขูนํ เอตมตฺถํ อาโรเจสุํ. ภิกฺขู ภควโต เอตมตฺถํ อาโรเจสุํ. อาปตฺติํ สา ภิกฺขเว ภิกฺขุนี อาปนฺนา ปาราชิกนฺติ. (103)}{ตติย\-ปาราชิกํ สมตฺตํ.}
\csromanheader{1}{4. \textbf{จตุตฺถ}\-\textbf{ปาราชิก}}
\proseitem{193}{\csromansymbolfootnote{*}{อิทํ วตฺถุ วิ 2. 35 ปิฏฺเฐปิ อาคตํ.}เตน สมเยน พุทฺโธ ภควา เวสาลิยํ วิหรติ มหาวเน กูฏาคาร\-สาลายํ. เตน โข ปน สมเยน สมฺพหุลา สนฺทิฏฺฐา สมฺภตฺตา ภิกฺขู วคฺคุมุทาย นทิยา ตีเร วสฺสํ อุปคจฺฉิํสุ. เตน โข ปน สมเยน วชฺชี ทุพฺภิกฺขา โหติ ทฺวีหิติกา เสตฏฺฐิกา สลากาวุตฺตา, น สุกรา อุญฺเฉน ปคฺคเหน ยาเปตุํ. อถ โข เตสํ ภิกฺขูนํ เอตทโหสิ “เอตรหิ โข วชฺชี ทุพฺภิกฺขา ทฺวีหิติกา เสตฏฺฐิกา สลากาวุตฺตา, น สุกรา อุญฺเฉน ปคฺคเหน ยาเปตุํ, เกน นุ โข มยํ อุปาเยน สมคฺคา สมฺโมทมานา อวิวทมานา ผาสุกํ วสฺสํ วเสยฺยาม, น จ ปิณฺฑเกน กิลเมยฺยามา”ติ. เอกจฺเจ เอวมาหํสุ “หนฺท มยํ อาวุโส คิหีนํ กมฺมนฺตํ อธิฏฺเฐม, เอวํ เต อมฺหากํ ทาตุํ มญฺญิสฺสนฺติ, เอวํ มยํ สมคฺคา สมฺโมทมานา อวิวทมานา ผาสุกํ วสฺสํ วสิสฺสาม, น จ ปิณฺฑเกน กิลมิสฺสามา”ติ. เอกจฺเจ เอวมาหํสุ “อลํ อาวุโส กิํ คิหีนํ กมฺมนฺตํ อธิฏฺฐิเตน, หนฺท มยํ อาวุโส คิหีนํ ทูเตยฺยํ หราม, เอวํ เต อมฺหากํ ทาตุํ มญฺญิสฺสนฺติ, เอวํ มยํ สมคฺคา สมฺโมทมานา อวิวทมานา ผาสุกํ วสฺสํ วสิสฺสาม, น จ ปิณฺฑเกน กิลมิสฺสามา”ติ. เอกจฺเจ เอวมาหํสุ “อลํ อาวุโส กิํ \csromanfolio{113}คิหีนํ กมฺมนฺตํ อธิฏฺฐิเตน, กิํ คิหีนํ ทูเตยฺยํ หเฏน, หนฺท มยํ อาวุโส คิหีนํ อญฺญมญฺญสฺส อุตฺตริ\-มนุสฺสธมฺมสฺส วณฺณํ ภาสิสฺสาม ‘อสุโก ภิกฺขุ ปฐมสฺส ฌานสฺส ลาภี, อสุโก ภิกฺขุ ทุติยสฺส ฌานสฺส ลาภี, อสุโก ภิกฺขุ ตติยสฺส ฌานสฺส ลาภี, อสุโก ภิกฺขุ จตุตฺถสฺส ฌานสฺส ลาภี, อสุโก ภิกฺขุ โสตาปนฺโน, อสุโก ภิกฺขุ สกทาคามี, อสุโก ภิกฺขุ อนาคามี, อสุโก ภิกฺขุ อรหา, อสุโก ภิกฺขุ เตวิชฺโช, อสุโก ภิกฺขุ ฉฬภิญฺโญ’ติ, เอวํ เต อมฺหากํ ทาตุํ มญฺญิสฺสนฺติ, เอวํ มยํ สมคฺคา สมฺโมทมานา อวิวทมานา ผาสุกํ วสฺสํ วสิสฺสาม, น จ ปิณฺฑเกน กิลมิสฺสามา”ติ. เอโส เยว โข อาวุโส เสยฺโย, โย อมฺหากํ คิหีนํ อญฺญมญฺญสฺส อุตฺตริ\-มนุสฺสธมฺมสฺส วณฺโณ ภาสิโตติ.}

\proseitem{194}{อถ โข เต ภิกฺขู คิหีนํ อญฺญมญฺญสฺส อุตฺตริ\-มนุสฺสธมฺมสฺส วณฺณํ ภาสิํสุ “อสุโก ภิกฺขุ ปฐมสฺส ฌานสฺส ลาภี ฯเปฯ อสุโก ภิกฺขุ ฉฬภิญฺโญ”ติ. อถ โข เตน มนุสฺสา “ลาภา วต โน, สุลทฺธํ วต โน, เยสํ โน เอวรูปา ภิกฺขู วสฺสํ อุปคตา, น วต โน อิโต ปุพฺเพ เอวรูปา ภิกฺขู วสฺสํ อุปคตา, ยถยิเม ภิกฺขู สีลวนฺโต กลฺยาณธมฺมา”ติ, เต น ตาทิสานิ โภชนานิ อตฺตานา ปริภุญฺชนฺติ มาตาปิตูนํ เทนฺติ ปุตฺตทารสฺส เทนฺติ ทาส\-กมฺมกร\-โปริสสฺส เทนฺติ มิตฺตามจฺจานํ เทนฺติ ญาติ\-สาโลหิตานํ เทนฺติ, ยาทิสานิ ภิกฺขูนํ เทนฺติ. เต น ตาทิสานิ ขาทนียานิ สายนียานิ ปานานิ อตฺตนา ขาทนฺติ สายนฺติ ปิวนฺติ มาตาปิตูนํ เทนฺติ ปุตฺตทารสฺส เทนฺติ ทาส\-กมฺมกร\-โปริสสฺส เทนฺติ มิตามจฺจานํ เทนฺติ ญาติ\-สาโลหิตานํ เทนฺติ, ยาทิสานิ ภิกฺขูนํ เทนฺติ. อถ โข เต ภิกฺขู วณฺณวา อเหสุํ ปีณินฺทฺริยา ปสนฺน\-มุขวณฺณา วิปฺปสนฺน\-ฉวิวณฺณา.}

\prose{อาจิณฺณํ โข ปเนตํ วสฺสํ\-วุฏฺฐานํ ภิกฺขูนํ ภควนฺตํ ทสฺสนาย อุปสงฺกมิตุํ. อถ โข เต ภิกฺขู วสฺสํวุฏฺฐา เตมาสจฺจเยน เสนาสนํ สํสาเมตฺวา ปตฺตจีวรํ อาทาย เยน เวสาลี เตน ปกฺกมิํสุ, อนุปุพฺเพน เยน เวสาลี มหาวนํ กูฏคารสาลา เยน ภควา เตนุ\-ปสงฺกมิํสุ, อุปสงฺกมิตฺวา ภควนฺตํ อภิวาเทตฺวา เอกมนฺตํ นิสีทิํสุ.}

\prose{\csromanfolio{114}เตน โข ปน สมเยน ทิสาสุ วสฺสํวุฏฺฐา ภิกฺขู กิสา โหนฺติ ลูขา ทุพฺพณฺณา อุปฺปณฺฑุ\-ปฺปณฺฑุก\-ชาตา ธมนิสนฺถต\-คตฺตา. วคฺคุมุทา\-ตีริยา ปน ภิกฺขู วณฺณวา โหนฺติ ปีณินฺทฺริยา ปสนฺน\-มุขวณฺณา วิปฺปสนฺน\-ฉวิวณฺณา. อาจิณฺณํ โข ปเนตํ พุทฺธานํ ภควนฺตานํ อาคนฺตุเกหิ ภิกฺขูหิ สทฺธิํ ปฏิสมฺโมทิตุํ. อถ โข ภควา วคฺคุมุทา\-ตีริเย ภิกฺขู เอตทโวจ “กจฺจิ ภิกฺขเว ขมนียํ, กจฺจิ ยาปนียํ, กจฺจิ สมคฺคา สมฺโมทมานา อวิวทมานา ผาสุกํ วสฺสํ วสิตฺถ, น จ ปิณฺฑเกน กิลมิตฺถา”ติ. ขมนียํ ภควา, ยาปนียํ ภควา, สมคฺคา จ มยํ ภนฺเต สมฺโมทมานา อวิวทมานา ผาสุกํ วสฺสํ วสิมฺหา, น จ ปิณฺฑเกน กิลมิมฺหาติ. ชานนฺตาปิ ตถาคตา ปุจฺฉนฺติ, ชานนฺตาปิ น ปุจฺฉนฺติ ฯเปฯ. ทฺวีหากาเรหิ พุทฺธา ภควนฺโต ภิกฺขู ปฏิปุจฺฉนฺติ “ธมฺมํ วา เทเสสฺสาม, สาวกานํ วา สิกฺขาปทํ ปญฺญเปสฺสามา”ติ. อถ โข ภควา วคฺคุมุทา\-ตีริเย ภิกฺขู เอตทโวจ “ยถา กถํ ปน ตุเมฺห ภิกฺขเว สมคฺคา สมฺโมทมานา อวิวทมานา ผาสุกํ วสฺสํ วสิตฺถ, น จ ปิณฺฑเกน กิลมิตฺถา”ติ. อถ โข เต ภิกฺขู ภควโต เอตมตฺถํ อาโรเจสุํ. กจฺจิ ปน โว ภิกฺขเว ภูตนฺติ. อภูตํ ภควาติ. วิครหิ พุทฺโธ ภควา “อนนุจฺฉวิกํ โมฆปุริสา อนนุโลมิกํ อปฺปติรูปํ อสฺสามณกํ อกปฺปิยํ อกรณียํ, กถํ หิ นาม ตุเมฺห โมฆปุริสา อุทรสฺส การณา คิหีนํ อญฺญมญฺญสฺส อุตฺตริ\-มนุสฺสธมฺมสฺส วณฺณํ ภาสิสฺสถ. วรํ ตุเมฺหหิ โมฆปุริสา ติเณฺหน โควิกนฺตเนน\csromansharedfootnote{d4efb0c6038f156c}{โควิกตฺตเนน (สี, ก)} กุจฺฉิ ปริกนฺโต, น เตฺวว อุทรสฺส การณา คิหีนํ อญฺญมญฺญสฺส อุตฺตริ\-มนุสฺสธมฺมสฺส วณฺโณ ภาสิโต. ตํ กิสฺส เหตุ, ตโต นิทานํ หิ โมฆปุริสา มรณํ วา นิคจฺเฉยฺย มรณมตฺตํ วา ทุกฺขํ, น เตฺวว ตปฺปจฺจยา กายสฺส เภทา ปรํ มรณา อปายํ ทุคฺคติํ วินิปาตํ นิรยํ อุปปชฺเชยฺย, อิโต นิทานญฺจ โข โมฆปุริสา กายสฺส เภทา ปรํ มรณา อปายํ ทุคฺคติํ วินิปาตํ นิรยํ อุปปชฺเชยฺย. เนตํ โมฆปุริสา อปฺปสนฺนานํ วา ปสาทาย ฯเปฯ. วิครหิตฺวา ธมฺมิํ กถํ กตฺวา ภิกฺขู อามนฺเตสิ\csromandash{}}

\proseitem{195}{ปญฺจิเม ภิกฺขเว มหาโจรา สนฺโต สํวิชฺชมาโน โลกสฺมิํ, กตเม ปญฺจ, อิธ ภิกฺขเว เอกจฺจสฺส มหาโจรสฺส เอวํ โหติ “กุทาสฺสุ นามาหํ สเตน วา สหสฺเสน วา ปริวุโต คามนิคม\-ราชธานีสุ \csromanfolio{115}อาหิณฺฑิสฺสามิ หนนฺโต ฆาเตนฺโต ฉินฺทนฺโต เฉทาเปนฺโต ปจนฺโต ปาเจนฺโต”ติ, โส อปเรน สมเยน สเตน วา สหสฺเสน วา ปริวุโต คามนิคม\-ราชธานีสุ อาหิณฺฑติ หนนฺโต ฆาเตนฺโต, ฉินฺทนฺโต เฉทาเปนฺโต ปจนฺโต ปาเจนฺโต, เอวเมว โข ภิกฺขเว อิเธกจฺจสฺส ปาปภิกฺขุโน เอวํ โหติ “กุทาสฺสุ นามาหํ สเตน วา สหสฺเสน วา ปริวุโต คามนิคม\-ราชธานีสุ จาริกํ จริสฺสามิ สกฺกโต ครุกโต มานิโต ปูชิโต อปจิโต คหฏฺฐาน\-ญฺเจว ปพฺพชิตานญฺจ ลาภี จีวร\-ปิณฺฑปาต\-เสนาสน\-คิลานปฺปจฺจย\-เภสชฺช\-ปริกฺขารานนฺ”ติ, โส อปเรน สมเยน สเตน วา สหสฺเสน วา ปริวุโต คามนิคม\-ราชธานีสุ จาริกํ จรติ สกฺกโต ครุกโต มานิโต ปูชิโต อปจิโต คหฏฺฐาน\-ญฺเจว ปพฺพชิตานญฺจ ลาภี \textbf{จีวร}\-\textbf{ปิณฺฑปาต}\-\textbf{เสนาสน}\-\textbf{คิลานปฺปจฺจย}\-\textbf{เภสชฺช}\-\textbf{ปริกฺขารานํ. อยํ} \textbf{ภิกฺขเว ปฐโม มหาโจโร สนฺโต สํวิชฺชมาโน โลกสฺมิํ.}}

\prose{ปุน จ ปรํ ภิกฺขเว อิเธกจฺโจ ปาปภิกฺขุ ตถาคต\-ปฺปเวทิตํ ธมฺมวินยํ ปริยาปุณิตฺวา อตฺตโน ทหติ. อยํ ภิกฺขเว ทุติโย มหาโจโร สนฺโต สํวิชฺชมาโน โลกสฺมิํ.}

\prose{ปุน จ ปรํ ภิกฺขเว อิเธกจฺโจ ปาปภิกฺขุ สุทฺธํ พฺรหฺมจาริํ ปริสุทฺธํ พฺรหฺมจริยํ จรนฺตํ อมูลเกน อพฺรหฺมจริเยน อนุทฺธํเสติ. อยํ ภิกฺขเว ตติโย มหาโจโร สนฺโต สํวิชฺชมาโน โลกสฺมิํ.}

\prose{ปุน จ ปรํ ภิกฺขเว อิเธกจฺโจ ปาปภิกฺขุ ยานิ ตานิ สํฆสฺส ครุภณฺฑานิ ครุปริกฺขารานิ, เสยฺยถิทํ, อาราโม อารามวตฺถุ วิหาโร วิหารวตฺถุ มญฺโจ ปีฐํ ภิสิ พิพฺโพหนํ\csromansharedfootnote{cb8b705f1f808b97}{พิมฺโพหนํ (สี, สฺยา)} โลหกุมฺภี โลหภาณกํ โลหวารโก โลหกฏหํ วาสิ ปรสุ\csromansharedfootnote{930289dec4135d4a}{ผรสุ (สี, สฺยา)} กุฐารี กุทาโล นิขาทนํ วลฺลิ เวฬุ มุญฺชํ ปพฺพชํ ติณํ มตฺติกา ทารุภณฺฑํ มตฺติกาภณฺฑํ, เตหิ คิหิํ สงฺคณฺหาติ อุปลาเปติ. อยํ ภิกฺขเว จตุตฺโถ มหาโจโร สนฺโต สํวิชฺชมาโน โลกสฺมิํ.}

\prose{สเทวเก ภิกฺขเว โลเก สมารเก สพฺรหฺมเก สสฺสมณ\-พฺราหฺมณิยา ปชาย สเทว\-มนุสฺสาย อยํ อคฺโค มหาโจโร, โย อสนฺตํ อภูตํ อุตฺตริ\-มนุสฺสธมฺมํ อุลฺลปติ, ตํ กิสฺส เหตุ, เถยฺยาย โว ภิกฺขเว รฏฺฐปิณฺโฑ ภุตฺโตติ.}

\csromanmatikaanchor{mtk.31}
\csromantocmarkref{subsection}{ปฐมปญฺญตฺติ}{mtk.31}
\bookmark[dest={mtk.31},level=2]{ปฐมปญฺญตฺติ}
\setlength{\csromangathapagewidth}{0pt}
\setlength{\csromangathawakwidth}{0pt}
\csromangathameasuregroup{\textsuperscript{*}“อญฺญาถา สนฺตมตฺตานํ, \\ นิกจฺจ กิตวสฺเสว, \\ \textsuperscript{+}กาสาวกณฺฐา พหโว, \\ ปาปา ปาเปหิ กมฺเมหิ, \\ \textsuperscript{+}เสยฺโย อโยคุโฬ ภุตฺโต, \\ ยญฺเจ ภุญฺเชยฺย ทุสฺสีโล,}{\csromangathabat{\textsuperscript{*}“อญฺญาถา สนฺตมตฺตานํ,}{อญฺญถา โย ปเวทเย.} \\ \csromangathabat{นิกจฺจ กิตวสฺเสว,}{ภุตฺตํ เถยฺเยน ตสฺส ตํ.} \\ \csromangathabat{\textsuperscript{+}กาสาวกณฺฐา พหโว,}{ปาปธมฺมา อสญฺญตา.} \\ \csromangathabat{ปาปา ปาเปหิ กมฺเมหิ,}{นิรยํ เต อุปปชฺชเร.} \\ \csromangathabat{\textsuperscript{+}เสยฺโย อโยคุโฬ ภุตฺโต,}{ตตฺโต อคฺคิสิขูปโม.} \\ \csromangathabat{ยญฺเจ ภุญฺเชยฺย ทุสฺสีโล,}{รฏฺฐปิณฺฑํ อสญฺญโต”ติ.}}
\csromangathasetleft{\textsuperscript{*}“อญฺญาถา สนฺตมตฺตานํ, \\ นิกจฺจ กิตวสฺเสว, \\ \textsuperscript{+}กาสาวกณฺฐา พหโว, \\ ปาปา ปาเปหิ กมฺเมหิ, \\ \textsuperscript{+}เสยฺโย อโยคุโฬ ภุตฺโต, \\ ยญฺเจ ภุญฺเชยฺย ทุสฺสีโล,}
\csromangathagroup{\csromangathastanza{\csromanfolio{116}\csromangathabat{\csromansymbolfootnote{*}{สํ 1. 22 ปิฏฺเฐปิ.}“อญฺญาถา สนฺตมตฺตานํ,}{อญฺญถา โย ปเวทเย.} \\ \csromangathabat{นิกจฺจ กิตวสฺเสว,}{ภุตฺตํ เถยฺเยน ตสฺส ตํ.}}\csromangathastanza{\csromangathabat{\csromansymbolfootnote{+}{ขุ 1. 57 ปิฏฺเฐ ธมฺมปเทปิ.}กาสาวกณฺฐา พหโว,}{ปาปธมฺมา อสญฺญตา.} \\ \csromangathabat{ปาปา ปาเปหิ กมฺเมหิ,}{นิรยํ เต อุปปชฺชเร.}}\csromangathastanza{\csromangathabat{\csromansymbolmark{+}เสยฺโย อโยคุโฬ ภุตฺโต,}{ตตฺโต อคฺคิสิขูปโม.} \\ \csromangathabat{ยญฺเจ ภุญฺเชยฺย ทุสฺสีโล,}{รฏฺฐปิณฺฑํ อสญฺญโต”ติ.}}}

\prose{อถ โข ภควา วคฺคุมุทา\-ตีริเย ภิกฺขู อเนก\-ปริยาเยน วิครหิตฺวา ทุพฺภรตาย ทุปฺโปสตาย ฯเปฯ. เอวญฺจ ปน ภิกฺขเว อิมํ สิกฺขาปทํ อุทฺทิเสยฺยาถ\csromandash{}}

\prose{\textbf{“โย ปน ภิกฺขุ อนภิชานํ อุตฺตริ}\-\textbf{มนุสฺสธมฺมํ อตฺตุปนายิกํ อลมริย}\-\textbf{ญาณทสฺสนํ สมุทาจเรยฺย ‘อิติ ชานามิ อิติ ปสฺสามี’ติ, ตโต อปเรนสมเยน สมนุ}\-\textbf{คฺคาหีย}\-\textbf{มาโน วา อสมนุ}\-\textbf{คฺคาหีย}\-\textbf{มาโน วา อาปนฺโน วิสุทฺธาเปกฺโข เอวํ วเทยฺย ‘อชานเมวํ อาวุโส อวจํ ชานามิ อปสฺสํ ปสฺสามิ ตุจฺฉํ มุสา วิลปินฺ’ติ, อยมฺปิ ปาราชิโก โหติ อสํวาโส”ติ}.}

\prose{เอวญฺจิทํ ภควตา ภิกฺขูนํ สิกฺขาปทํ ปญฺญตฺตํ โหติ.}

\proseitem{196}{เตน โข ปน สมเยน สมฺพหุลา ภิกฺขู อทิฏฺเฐ ทิฏฺฐสญฺญิโน อปตฺเต ปตฺตสญฺญิโน อนธิคเต อธิคต\-สญฺญิโน อสจฺฉิกเต สจฺฉิกต\-สญฺญิโน อธิมาเนน อญฺญํ พฺยากริํสุ. เตสํ อปเรน สมเยน ราคายปิ จิตฺตํ นมติ, โทสายปิ จิตฺตํ นมติ, โมหายปิ จิตฺตํ นมติ. เตสํ กุกฺกุจฺจํ อโหสิ “ภควตา สิกฺขาปทํ ปญฺญตฺตํ, มยญฺจมฺห อทิฏฺเฐ ทิฏฺฐสญฺญิโน อปตฺเต ปตฺตสญฺญิโน อนธิคเต อธิคต\-สญฺญิโน อสจฺฉิกเต สจฺฉิกต\-สญฺญิโน อธิมาเนน อญฺญํ พฺยากริมฺหา, กจฺจิ นุ โข มยํ ปาราชิกํ อาปตฺติํ อาปนฺนา”ติ. เต อายสฺมโต อานนฺทสฺส เอตมตฺถํ อาโรเจสุํ. อายสฺมา อานนฺโท ภควโต เอตมตฺถํ อาโรเจสิ. โหนฺติ เย เต อานนฺท\csromansharedfootnote{3a8e7c861c4e0d46}{โหนฺติ เยวานนฺท (สฺยา), โหนฺติ เต อานนฺท (สี)} ภิกฺขู อทิฏฺเฐ ทิฏฺฐสญฺญิโน อปตฺเต ปตฺตสญฺญิโน อนธิคเต อธิคต\-สญฺญิโน อสจฺฉิกเต สจฺฉิกต\-สญฺญิโน อธิมาเนน อญฺญํ พฺยากโรนฺติ, ตญฺจ โข เอตํ อพฺโพหาริกนฺติ. เอวญฺจ ปน ภิกฺขเว อิมํ สิกฺขาปทํ อุทฺทิเสยฺยาถ\csromandash{}}

\csromanmatikaanchor{mtk.32}
\csromanmatikaanchor{mtk.33}
\csromantocmarkref{subsection}{อนุปญฺญตฺติ}{mtk.32}
\bookmark[dest={mtk.32},level=2]{อนุปญฺญตฺติ}
\csromantocmarkref{subsection}{สิกฺขาปทวิภงฺค, ปทภาชนีย}{mtk.33}
\bookmark[dest={mtk.33},level=2]{สิกฺขาปทวิภงฺค, ปทภาชนีย}
\proseitem{197}{\csromanfolio{117}\textbf{“โย ปน ภิกฺขุ อนภิชานํ อุตฺตริ}\-\textbf{มนุสฺสธมฺมํ อตฺตุปนายิกํ อลมริย}\-\textbf{ญาณทสฺสนํ สมุทาจเรยฺย ‘อิติ ชานามิ อิติ ปสฺสามี’ติ, ตโต อปเรน สมเยน สมนุ}\-\textbf{คฺคาหีย}\-\textbf{มาโน วา อสมนุ}\-\textbf{คฺคาหีย}\-\textbf{มาโน วา อาปนฺโน วิสุทฺธาเปกฺโข เอวํ วเทยฺย ‘อชานเมวํ อาวุโส อวจํ ชานามิ อปสฺสํ ปสฺสามิ ตุจฺฉํ มุสา วิลปินฺ’ติ, อญฺญตฺร อธิมานา, อยมฺปิ ปาราชิโก โหติ อสํวาโส”ติ}.}

\proseitem{198}{\textbf{โยปนา}ติ\csromansharedfootnote{423cb82fe87b6ac9}{โย ปนาติ เทฺว ปทานิ \csromandash{} ม.พ.ป.} โย ยาทิโส ฯเปฯ. ภิกฺขูติ ฯเปฯ. อยํ อิมสฺมิํ อตฺเถ อธิปฺเปโต “ภิกฺขู”ติ.}

\prose{\textbf{อนภิชานนฺ}ติ อสนฺตํ อภูตํ อสํวิชฺชมานํ อชานนฺโต อปสฺสนฺโต อตฺตนิ กุสลํ ธมฺมํ “อตฺถิ เม กุสโล ธมฺโม”ติ.}

\prose{\csromansymbolfootnote{*}{วิ 2. 37 ปิฏฺเฐปิ.}\textbf{อุตฺตริ}\-\textbf{มนุสฺสธมฺโม} นาม ฌานํ วิโมกฺโข\csromansharedfootnote{91182df06eb2c4a9}{วิโมกฺขํ (สี, สฺยา)} สมาธิ สมาปตฺติ ญาณทสฺสนํ มคฺคภาวนา ผลสจฺฉิกิริยา กิเลสปฺปหานํ วินีวรณตา จิตฺตสฺส สุญฺญาคาเร อภิรติ.}

\prose{\textbf{อตฺตุ}\-\textbf{ปนายิกนฺ}ติ เต วา กุสเล ธมฺเม อตฺตนิ อุปเนติ, อตฺตานํ วา เตสุ กุสเลสุ ธมฺเมสุ อุปเนติ.}

\prose{\textbf{ญาณนฺ}ติ ติสฺโส วิชฺชา. \textbf{ทสฺสนนฺ}ติ ยํ ญาณํ ตํ ทสฺสนํ, ยํ ทสฺสนํ ตํ ญาณํ.}

\prose{\textbf{สมุทาจเรยฺยา}ติ อาโรเจยฺย อิตฺถิยา วา ปุริสสฺส วา คหฏฺฐสฺส วา ปพฺพชิตสฺส วา.}

\prose{\textbf{อิติ ชานามิ อิติ ปสฺสามี}ติ “ชานามยํ เอเต ธมฺเม, ปสฺสามหํ เอเต ธมฺเม, อตฺถิ จ เอเต ธมฺมา มยิ, อหญฺจ เอเตสุ ธมฺเมสุ สนฺทิสฺสามี”ติ.}

\prose{\textbf{ตโต อปเรน สมเยนา}ติ ยสฺมิํ ขเณ สมุทาจิณฺณํ โหติ, ตํ ขณํ ตํ ลยํ ตํ มุหุตฺตํ วีติวตฺเต.}

\prose{\textbf{สมนุ}\-\textbf{คฺคาหีย}\-\textbf{มาโน}ติ ยํ วตฺถุ ปฏิญฺญาตํ โหติ, ตสฺมิํ วตฺถุสฺมิํ สมนุ\-คฺคาหีย\-มาโน “กิํ เต อธิคตํ, กินฺติ เต อธิคตํ, กทา เต อธิคตํ, กตฺเถ เต อธิคตํ, กตเม เต กิเลสา ปหีนา, กเตเมสํ ตฺวํ ธมฺมานํ ลาภี”ติ.}

\prose{\csromanfolio{118}\textbf{อสมนุคฺคาหียมาโน}ติ น เกนจิ วุจฺจมาโน.}

\prose{\textbf{อาปนฺโน}ติ ปาปิจฺโฉ อิจฺฉาปกโต อสนฺตํ อภูตํ อุตฺตริ\-มนุสฺสธมฺมํ อุลฺลปิตฺวา ปาราชิกํ อาปตฺติํ อาปนฺโน โหติ.}

\prose{\textbf{วิสุทฺธา}\-\textbf{เปกฺโข}ติ คิหี วา โหตุกาโม อุปาสโก วา โหตุกาโม อารามิโก วา โหตุกาโม สามเณโร วา โหตุกาโม.}

\prose{\textbf{อชานเมวํ อาวุโส อวจํ ชานามิ อปสฺสํ ปสฺสามี}ติ “นาหํ เอเต ธมฺเม ชานามิ, นาหํ เอเต ธมฺเม ปสฺสามิ, นตฺถิ จ เอเต ธมฺมา มยิ, น จาหํ เอเตสุ ธมฺเมสุ สนฺทิสฺสามี”ติ.}

\prose{\textbf{ตุจฺฉํ มุสา วิลปินฺ}ติ ตุจฺฉกํ มยา ภณิตํ, มุสา มยา ภณิตํ, อภูตํ มยา ภณิตํ, อชานนฺเตน มยา ภณิตํ.}

\prose{\textbf{อญฺญตฺร อธิมานา}ติ ฐเปตฺวา อธิมานํ.}

\prose{\textbf{อยมฺปี}ติ ปุริเม อุปาทาย วุจฺจติ.}

\prose{\textbf{ปาราชิโก โหตี}ติ เสยฺยถาปิ นาม ตาโล มตฺถก\-จฺฉินฺโน อภพฺโพ ปุน วิรูฬฺหิยา, เอวเมว ภิกฺขุ ปาปิจฺโฉ อิจฺฉาปกโต อสนฺตํ อภูตํ อุตฺตริ\-มนุสฺสธมฺมํ อุลฺลปิตฺวา อสฺสมโณ โหติ อสกฺยปุตฺติโย, เตน วุจฺจติ “ปาราชิโก โหตี”ติ.}

\prose{\textbf{อสํวาโส}ติ สํวาโส นาม เอกกมฺมํ เอกุทฺเทโส สมสิกฺขตา, เอโส สํวาโส นาม, โส เตน สทฺธิํ นตฺถิ, เตน วุจฺจติ “อสํวาโส”ติ.}

\proseitem{199}{\textbf{อุตฺตริ}\-\textbf{มนุสฺสธมฺโม} นาม ฌานํ วิโมกฺโข สมาธิ สมาปตฺติ ญาณทสฺสนํ มคฺคภาวนา ผลสจฺฉิกิริยา กิเลสปฺปหานํ วินีวรณตา จิตฺตสฺส สุญฺญาคาเร อภิรติ.}

\prose{\csromansymbolfootnote{*}{วิ 2. 38 ปิฏฺเฐปิ.}\textbf{ฌานนฺ}ติ ปฐมํ ฌานํ ทุติยํ ฌานํ ตติยํ ฌานํ จตุตฺถํ ฌานํ.}

\prose{\csromansymbolmark{*}\textbf{วิโมกฺโข}ติ สุญฺญโต วิโมกฺโข อนิมิตฺโต วิโมกฺโข อปฺปณิหิโต วิโมกฺโข.}

\prose{\csromansymbolmark{*}\textbf{สมาธิ}ติ สุญฺญโต สมาธิ อนิมิตฺโต สมาธิ อปฺปณิหิโต สมาธิ.}

\prose{\csromanfolio{119}\csromansymbolfootnote{*}{วิ 2. 38 ปิฏฺเฐปิ.}\textbf{สมาปตฺตี}ติ สุญฺญตา สมาปตฺติ อนิมิตฺตา สมาปตฺติ อปฺปณิหิตา สมาปตฺติ.}

\prose{\csromansymbolmark{*}\textbf{ญาณทสฺสนนฺ}ติ ติสฺโส วิชฺชา.}

\prose{\csromansymbolmark{*}\textbf{มคฺคภาวนา}ติ จตฺตาโร สติปฏฺฐานา จตฺตาโร สมฺมปฺปธานา จตฺตาโร อิทฺธิปาทา ปญฺจินฺทฺริยานิ ปญฺจ พลานิ สตฺต โพชฺฌงฺคา อริโย อฏฺฐงฺคิโก มคฺโค.}

\prose{\csromansymbolmark{*}\textbf{ผลสจฺฉิกิริยา}ติ โสตาปตฺติ\-ผลสฺส สจฺฉิกิริยา สกทาคามิ\-ผลสฺส สจฺฉิกิริยา อนาคามิ\-ผลสฺส สจฺฉิกิริยา อรหตฺตสฺส\csromansharedfootnote{50a831860e03267f}{อรหตฺตผลสฺส (สฺยา)} สจฺฉิกิริยา.}

\prose{\csromansymbolmark{*}\textbf{กิเลส}\-\textbf{ปฺปหานนฺ}ติ ราคสฺส ปหานํ โทสสฺส ปหานํ โมหสฺส ปหานํ.}

\prose{\csromansymbolmark{*}\textbf{วินีวรณตา จิตฺตสฺสา}ติ ราคา จิตฺตํ วินีวรณตา โทสา จิตฺตํ วินีวรณตา โมหา จิตฺตํ วินีวรณตา.}

\prose{\csromansymbolmark{*}\textbf{สุญฺญาคาเร อภิรตี}ติ ปฐเมน ฌาเนน สุญฺญาคาเร อภิรติ, ทุติเยน ฌาเนน สุญฺญาคาเร อภิรติ, ตติเยน ฌาเนน สุญฺญาคาเร อภิรติ, จตุตฺเถน ฌาเนน สุญฺญาคาเร อภิรติ.}

\proseitem{200}{ตีหากาเรหิ ปฐมํ ฌานํ สมาปชฺชินฺติ สมฺปชานมุสา ภณนฺตสฺส อาปตฺติ ปาราชิกสฺส, ปุพฺเพวสฺส โหติ “มุสา ภณิสฺสนฺ”ติ, ภณนฺตสฺส โหติ “มุสา ภณามี”ติ, ภณิตสฺส โหติ “มุสา มยา ภณิตนฺ”ติ.}

\prose{จตูหากาเรหิ ปฐมํ ฌานํ สมาปชฺชินฺติ สมฺปชานมุสา ภณนฺตสฺส อาปตฺติ ปาราชิกสฺส, ปุพฺเพวสฺส โหติ “มุสา ภณิสฺสนฺ”ติ, ภณนฺตสฺส โหติ “มุสา ภณามี”ติ, ภณิตสฺส โหติ “มุสา มยา ภณิตนฺ”ติ, วินิธาย ทิฏฺฐิํ.}

\prose{ปญฺจหากาเรหิ ปฐมํ ฌานํ สมาปชฺชินฺติ สมฺปชานมุสา ภณนฺตสฺส อาปตฺติ ปาราชิกสฺส, ปุพฺเพวสฺส โหติ “มุสา ภณิสฺสนฺ”ติ, ภณนฺตสฺส โหติ “มุสา ภณามี”ติ, ภณิตสฺส โหติ “มุสา มยา ภณิตนฺ”ติ, วินิธายทิฏฺฐิํ, วินิธาย ขนฺติํ.}

\prose{\csromanfolio{120}ฉหากาเรหิ ปฐมํ ฌานํ สมาปชฺชินฺติ สมฺปชานมุสา ภณนฺตสฺส อาปตฺติ ปาราชิกสฺส, ปุพฺเพวสฺส โหติ “มุสา ภณิสฺสนฺ”ติ, ภณนฺตสฺส โหติ “มุสา ภณามี”ติ, ภณิตสฺส โหติ “มุสา มยา ภณิตนฺ”ติ, วินิธาย ทิฏฺฐิํ, วินิธาย ขนฺติํ, วินิธาย รุจิํ.}

\prose{สตฺตหากาเรหิ ปฐมํ ฌานํ สมาปชฺชินฺติ สมฺปชานมุสา ภณนฺตสฺส อาปตฺติ ปาราชิกสฺส, ปุพฺเพวสฺส โหติ “มุสา ภณิสฺสนฺ”ติ, ภณนฺตสฺส โหติ “มุสา ภณามี”ติ, ภณิตสฺส โหติ “มุสา มยา ภณิตนฺ”ติ, วินิธาย ทิฏฺฐิํ, วินิธาย ขนฺติํ, วินิธาย รุจิํ, วินิธาย ภาวํ.}

\proseitem{201}{ตีหากาเรหิ ปฐมํ ฌานํ สมาปชฺชามีติ สมฺปชานมุสา ภณนฺตสฺส อาปตฺติ ปาราชิกสฺส, ปุพฺเพวสฺส โหติ “มุสา ภณิสฺสนฺ”ติ, ภณนฺตสฺส โหติ “มุสา ภณามี”ติ, ภณิตสฺส โหติ “มุสา มยา ภณิตนฺ”ติ.}

\prose{จตูหากาเรหิ ปฐมํ ฌานํ สมาปชฺชามีติ สมฺปชานมุสา ภณนฺตสฺส อาปตฺติ ปาราชิกสฺส, ปุพฺเพวสฺส โหติ “มุสา ภณิสฺสนฺ”ติ, ภณนฺตสฺส โหติ “มุสา ภณามี”ติ, ภณิตสฺส โหติ “มุสา มยา ภณิตนฺ”ติ, วินิธาย ทิฏฺฐิํ.}

\prose{ปญฺจหากาเรหิ ปฐมํ ฌานํ สมาปชฺชามีติ สมฺปชานมุสา ภณนฺตสฺส อาปตฺติ ปาราชิกสฺส, ปุพฺเพวสฺส โหติ “มุสา ภณิสฺสนฺ”ติ, ภณนฺตสฺส โหติ “มุสา ภณามี”ติ, ภณิตสฺส โหติ “มุสา มยา ภณิตนฺ”ติ, วินิธาย ทิฏฺฐิํ, วินิธาย ขนฺติํ.}

\prose{ฉหากาเรหิ ปฐมํ ฌานํ สมาปชฺชามีติ สมฺปชานมุสา ภณนฺตสฺส อาปตฺติ ปาราชิกสฺส, ปุพฺเพวสฺส โหติ “มุสา ภณิสฺสนฺ”ติ, ภณนฺตสฺส โหติ “มุสา ภณามี”ติ, ภณิตสฺส โหติ “มุสา มยา ภณิตนฺ”ติ, วินิธาย ทิฏฺฐิํ, วินิธาย ขนฺติํ, วินิธาย รุจิํ.}

\prose{สตฺตหากาเรหิ ปฐมํ ฌานํ สมาปชฺชามีติ สมฺปชานมุสา ภณนฺตสฺส อาปตฺติ ปาราชิกสฺส, ปุพฺเพวสฺส โหติ “มุสา ภณิสฺสนฺ”ติ, ภณนฺตสฺส โหติ “มุสา ภณามี”ติ, ภณิตสฺส โหติ “มุสา มยา ภณิตนฺ”ติ, วินิธาย ทิฏฺฐิํ, วินิธาย ขนฺติํ, วินิธาย รุจิํ, วินิธาย ภาวํ.}

\proseitem{202}{\csromanfolio{121}ตีหากาเรหิ ปฐมํ ฌานํ สมาปนฺโนติ สมฺปชานมุสา ภณนฺตสฺส อาปตฺติ ปาราชิกสฺส, ปุพฺเพวสฺส โหติ “มุสา ภณิสฺสนฺ”ติ, ภณนฺตสฺส โหติ “มุสา ภณามี”ติ, ภณิตสฺส โหติ “มุสา มยา ภณิตนฺ”ติ.}

\prose{จตูหากาเรหิ ปฐมํ ฌานํ สมาปนฺโนติ สมฺปชานมุสา ภณนฺตสฺส อาปตฺติ ปาราชิกสฺส, ปุพฺเพวสฺส โหติ “มุสา ภณิสฺสนฺ”ติ, ภณนฺตสฺส โหติ “มุสา ภณามี”ติ, ภณิตสฺส โหติ “มุสา มยา ภณิตนฺ”ติ, วินิธาย ทิฏฺฐิํ.}

\prose{ปญฺจหากาเรหิ ปฐมํ ฌานํ สมาปนฺโนติ สมฺปชานมุสา ภณนฺตสฺส อาปตฺติ ปาราชิกสฺส, ปุพฺเพวสฺส โหติ “มุสา ภณิสฺสนฺ”ติ, ภณนฺตสฺส โหติ “มุสา ภณามี”ติ, ภณิตสฺส โหติ “มุสา มยา ภณิตนฺ”ติ, วินิธาย ทิฏฺฐิํ, วินิธาย ขนฺติํ.}

\prose{ฉหากาเรหิ ปฐมํ ฌานํ สมาปนฺโนติ สมฺปชานมุสา ภณนฺตสฺส อาปตฺติ ปาราชิกสฺส, ปุพฺเพวสฺส โหติ “มุสา ภณิสฺสนฺ”ติ, ภณนฺตสฺส โหติ “มุสา ภณามี”ติ, ภณิตสฺส โหติ “มุสา มยา ภณิตนฺ”ติ, วินิธาย ทิฏฺฐิํ, วินิธาย ขนฺติํ, วินิธาย รุจิํ.}

\prose{สตฺตหากาเรหิ ปฐมํ ฌานํ สมาปนฺโนติ สมฺปชานมุสา ภณนฺตสฺส อาปตฺติ ปาราชิกสฺส, ปุพฺเพวสฺส โหติ “มุสา ภณิสฺสนฺ”ติ, ภณนฺตสฺส โหติ “มุสา ภณามี”ติ, ภณิตสฺส โหติ “มุสา มยา ภณิตนฺ”ติ, วินิธาย ทิฏฺฐิํ, วินิธาย ขนฺติํ, วินิธาย รุจิํ, วินิธาย ภาวํ.}

\proseitem{203}{ตีหากาเรหิ ปฐมสฺส ฌานสฺส ลาภิมฺหีติ สมฺปชานมุสา ภณนฺตสฺส อาปตฺติ ปาราชิกสฺส, ปุพฺเพวสฺส โหติ “มุสา ภณิสฺสนฺ”ติ, ภณนฺตสฺส โหติ “มุสา ภณามี”ติ, ภณิตสฺส โหติ “มุสา มยา ภณิตนฺ”ติ.}

\prose{จตูหากาเรหิ ปฐมสฺส ฌานสฺส ลาภิมฺหีติ สมฺปชานมุสา ภณนฺตสฺส อาปตฺติ ปาราชิกสฺส, ปุพฺเพวสฺส โหติ “มุสา ภณิสฺสนฺ”ติ, ภณนฺตสฺส โหติ “มุสา ภณามี”ติ, ภณิตสฺส โหติ “มุสา มยา ภณิตนฺ”ติ, วินิธาย ทิฏฺฐิํ.}

\prose{\csromanfolio{122}ปญฺจหากาเรหิ ปฐมสฺส ฌานสฺส ลาภิมฺหิติ สมฺปชานมุสา ภณนฺตสฺส อาปตฺติ ปาราชิกสฺส, ปุพฺเพวสฺส โหติ “มุสา ภณิสฺสนฺ”ติ, ภณนฺตสฺส โหติ “มุสา ภณามี”ติ, ภณิตสฺส โหติ “มุสา มยา ภณิตนฺ”ติ, วินิธาย ทิฏฺฐิํ, วินิธาย ขนฺติํ.}

\prose{ฉหากาเรหิ ปฐมสฺส ฌานสฺส ลาภิมฺหีติ สมฺปชานมุสา ภณนฺตสฺส อาปตฺติ ปาราชิกสฺส, ปุพฺเพวสฺส โหติ “มุสา ภณิสฺสนฺ”ติ, ภณนฺตสฺส โหติ “มุสา ภณามี”ติ, ภณิตสฺส โหติ “มุสา มยา ภณิตนฺ”ติ, วินิธาย ทิฏฺฐิํ, วินิธาย ขนฺติํ, วินิธาย รุจิํ.}

\prose{สตฺตหากาเรหิ ปฐมสฺส ฌานสฺส ลาภิมฺหีติ สมฺปชานมุสา ภณนฺตสฺส อาปตฺติ ปาราชิกสฺส, ปุพฺเพวสฺส โหติ “มุสา ภณิสฺสนฺ”ติ, ภณนฺตสฺส โหติ “มุสา ภณามี”ติ, ภณิตสฺส โหติ “มุสา มยา ภณิตนฺ”ติ, วินิธาย ทิฏฺฐิํ, วินิธาย ขนฺติํ, วินิธาย รุจิํ, วินิธาย ภาวํ.}

\proseitem{204}{ตีหากาเรหิ ปฐมสฺส ฌานสฺส วสิมฺหีติ สมฺปชานมุสา ภณนฺตสฺส อาปตฺติ ปาราชิกสฺส, ปุพฺเพวสฺส โหติ “มุสา ภณิสฺสนฺ”ติ, ภณนฺตสฺส โหติ “มุสา ภณามี”ติ, ภณิตสฺส โหติ “มุสา มยา ภณิตนฺ”ติ.}

\prose{จตูหากาเรหิ ปฐมสฺส ฌานสฺส วสิมฺหีติ สมฺปชานมุสา ภณนฺตสฺส อาปตฺติ ปาราชิกสฺส, ปุพฺเพวสฺส โหติ “มุสา ภณิสฺสนฺ”ติ, ภณนฺตสฺส โหติ “มุสา ภณามี”ติ, ภณิตสฺส โหติ “มุสา มยา ภณิตนฺ”ติ, วินิธาย ทิฏฺฐิํ.}

\prose{ปญฺจหากาเรหิ ปฐมสฺส ฌานสฺส วสิมฺหีติ สมฺปชานมุสา ภณนฺตสฺส อาปตฺติ ปาราชิกสฺส, ปุพฺเพวสฺส โหติ “มุสา ภณิสฺสนฺ”ติ, ภณนฺตสฺส โหติ “มุสา ภณามี”ติ, ภณิตสฺส โหติ “มุสา มยา ภณิตนฺ”ติ, วินิธาย ทิฏฺฐิํ, วินิธาย ขนฺติํ.}

\prose{ฉหากาเรหิ ปฐมสฺส ฌานสฺส วสิมฺหีติ สมฺปชานมุสา ภณนฺตสฺส อาปตฺติ ปาราชิกสฺส, ปุพฺเพวสฺส โหติ “มุสา ภณิสฺสนฺ”ติ, ภณนฺตสฺส โหติ “มุสา ภณามี”ติ, ภณิตสฺส โหติ “มุสา มยา ภณิตนฺ”ติ, วินิธาย ทิฏฺฐิํ, วินิธาย ขนฺติํ, วินิธาย รุจิํ.}

\prose{\csromanfolio{123}สตฺตหากาเรหิ ปฐมสฺส ฌานสฺส วสิมฺหีติ สมฺปชานมุสา ภณนฺตสฺส อาปตฺติ ปาราชิกสฺส, ปุพฺเพวสฺส โหติ “มุสา ภณิสฺสนฺ”ติ, ภณนฺตสฺส โหติ “มุสา ภณามี”ติ, ภณิตสฺส โหติ “มุสา มยา ภณิตนฺ”ติ, วินิธาย ทิฏฺฐิํ, วินิธาย ขนฺติํ, วินิธาย รุจิํ, วินิธาย ภาวํ.}

\proseitem{205}{ตีหากาเรหิ ปฐมํ ฌานํ สจฺฉิกตํ มยาติ สมฺปชานมุสา ภณนฺตสฺส อาปตฺติ ปาราชิกสฺส, ปุพฺเพวสฺส โหติ “มุสา ภณิสฺสนฺ”ติ, ภณนฺตสฺส โหติ “มุสา ภณามี”ติ, ภณิตสฺส โหติ “มุสา มยา ภณิตนฺ”ติ.}

\prose{จตูหากาเรหิ ปฐมํ ฌานํ สจฺฉิกตํ มยาติ สมฺปชานมุสา ภณนฺตสฺส อาปตฺติ ปาราชิกสฺส, ปุพฺเพวสฺส โหติ “มุสา ภณิสฺสนฺ”ติ, ภณนฺตสฺส โหติ “มุสา ภณามี”ติ, ภณิตสฺส โหติ “มุสา มยา ภณิตนฺ”ติ, วินิธาย ทิฏฺฐิํ.}

\prose{ปญฺจหากาเรหิ ปฐมํ ฌานํ สจฺฉิกตํ มยาติ สมฺปชานมุสา ภณนฺตสฺส อาปตฺติ ปาราชิกสฺส, ปุพฺเพวสฺส โหติ “มุสา ภณิสฺสนฺ”ติ, ภณนฺตสฺส โหติ “มุสา ภณามี”ติ, ภณิตสฺส โหติ “มุสา มยา ภณิตนฺ”ติ, วินิธาย ทิฏฺฐิํ, วินิธาย ขนฺติํ.}

\prose{ฉหากาเรหิ ปฐมํ ฌานํ สจฺฉิกตํ มยาติ สมฺปชานมุสา ภณนฺตสฺส อาปตฺติ ปาราชิกสฺส, ปุพฺเพวสฺส โหติ “มุสา ภณิสฺสนฺ”ติ, ภณนฺตสฺส โหติ “มุสา ภณามี”ติ, ภณิตสฺส โหติ “มุสา มยา ภณิตนฺ”ติ, วินิธาย ทิฏฺฐิํ, วินิธาย ขนฺติํ, วินิธาย รุจิํ.}

\prose{สตฺตหากาเรหิ ปฐมํ ฌานํ สจฺฉิกตํ มยาติ สมฺปชานมุสา ภณนฺตสฺส อาปตฺติ ปาราชิกสฺส, ปุพฺเพวสฺส โหติ “มุสา ภณิสฺสนฺ”ติ, ภณนฺตสฺส โหติ “มุสา ภณามี”ติ, ภณิตสฺส โหติ “มุสา มยา ภณิตนฺ”ติ, วินิธาย ทิฏฺฐิํ, วินิธาย ขนฺติํ, วินิธาย รุจิํ, วินิธาย ภาวํ.}

\prose{ยถา อิทํ ปฐมํ วิตฺถาริตํ, เอวํ สพฺพํปิ วิตฺถาเรตพฺพํ.}

\proseitem{206}{ตีหากาเรหิ ฯเปฯ. สตฺตหากาเรหิ ทุติยํ ฌานํ ฯเปฯ ตติยํ ฌานํ ฯเปฯ จตุตฺถํ ฌานํ สมาปชฺชิํ. สมาปชฺชามิ. สมาปนฺโน. \csromanfolio{124}จตุตฺถสฺส ฌานสฺส ลาภีมฺหิ. วสีมฺหิ. จตุตฺถํ ฌานํ สจฺฉิกตํ มยาติ สมฺปชานมุสา ภณนฺตสฺส อาปตฺติ ปาราชิกสฺส, ปุพฺเพวสฺส โหติ “มุสา ภณิสฺสนฺ”ติ, ภณนฺตสฺส โหติ “มุสา มยา ภณิตนฺ”ติ, วินิธาย ทิฏฺฐิํ, วินิธาย ขนฺติํ, วินิธาย รุจิํ, วินิธาย ภาวํ.}

\proseitem{207}{ตีหากาเรหิ สุญฺญตํ วิโมกฺขํ. อนิมิตฺตํ วิโมกฺขํ. อปฺปณิหิตํ วิโมกฺขํ สมาปชฺชิํ. สมาปชฺชามิ. สมาปนฺโน. อปฺปณิหิตสฺส วิโมกฺขสฺส ลาภิมฺหิ. วสิมฺหิ. อปฺปณิหิโต วิโมกฺโข สจฺฉิกโต มยาติ สมฺปชานมุสา ภณนฺตสฺส อาปตฺติ ปาราชิกสฺส ฯเปฯ.}

\prose{ตีหากเรหิ สุญฺญตํ สมาธิํ. อนิมิตฺตํ สมาธิํ. อปฺปณิหิตํ สมาธิํ สมาปชฺชิํ. สมาปชฺชามิ. สมาปนฺโน. อปฺปณิหิตสฺส สมาธิสฺส ลาภีมฺหิ. วสีมฺหิ. อปฺปณิหิโต สมาธิ สจฺฉิกโต มยาติ สมฺปชานมุสา ภณนฺตสฺส อาปตฺติ ปาราชิกสฺส.}

\prose{ตีหากาเรหิ สุญฺญตํ สมาปตฺติํ. อนิมิตฺตํ สมาปตฺติํ. อปฺปณิหิตํ สมาปตฺติํ สมาปชฺชิํ. สมาปชฺชามิ. สมาปนฺโน. อปฺปณิหิตาย สมาปตฺติยา ลาภิมฺหิ. วสิมฺหิ. อปฺปณิหิตา สมาปตฺติ สจฺฉิกตา มยาติ สมฺปชานมุสา ภณนฺตสฺส อาปตฺติ ปาราชิกสฺส.}

\prose{ตีหากาเรหิ ติสฺโส วิชฺชา สมาปชฺชิํ. สมาปชฺชามิ. สมาปนฺโน. ติสฺสนฺนํ วิชฺชานํ ลาภิมฺหิ. วสิมฺหิ. ติสฺโส วิชฺชา สจฺฉิกตา มยาติ สมฺปชานมุสา ภณนฺตสฺส อาปตฺติ ปาราชิกสฺส.}

\prose{ตีหากาเรหิ จตฺตาโร สติปฏฺฐาเน. จตฺตาโร สมฺมปฺปธาเน. จตฺตาโร อิทฺธิปาเท สมาปชฺชิํ. สมาปชฺชามิ. สมาปนฺโน. จตุนฺนํ อิทฺธิปาทานํ ลาภิมฺหิ. วสิมฺหิ. จตฺตาโร อิทฺธิปาทา สจฺฉิกตา มยาติ สมฺปชานมุสา ภณนฺตสฺส อาปตฺติ ปาราชิกสฺส.}

\prose{ตีหากาเรหิ ปญฺจินฺทฺริยานิ. ปญฺจ พลานิ สมาปชฺชิํ. สมาปชฺชามิ. สมาปนฺโน. ปญฺจนฺนํ พลานํ ลาภิมฺหิ. วสิมฺหิ. ปญฺจพลานิ สจฺฉิกตานิ มยาติ สมฺปชานมุสา ภณนฺตสฺส อาปตฺติ ปาราชิกสฺส.}

\prose{ตีหากาเรหิ สตฺต โพชฺฌงฺเค สมาปชฺชิํ. สมาปชฺชามิ. สมาปนฺโน. สตฺตนฺนํ โพชฺฌงฺคานํ ลาภิมฺหิ. วสิมฺหิ. สตฺต โพชฺฌงฺคา สจฺฉิกตา มยาติ สมฺปชานมุสา ภณนฺตสฺส อาปตฺติ ปาราชิกสฺส.}

\prose{\csromanfolio{125}ตีหากาเรหิ อริยํ อฏฺฐงฺคิกํ มคฺคํ สมาปชฺชิํ. สมาปชฺชามิ. สมาปนฺโน. อริยสฺส อฏฺฐงฺคิกสฺส มคฺคสฺส ลาภิมฺหิ. วสิมฺหิ. อริโย อฏฺฐงฺคิโก มคฺโค สจฺฉิกตา มยาติ สมฺปชานมุสา ภณนฺตสฺส อาปตฺติ ปาราชิกสฺส.}

\prose{ตีหากาเรหิ โสตาปตฺติผลํ. สกทาคามิผลํ. อนาคามิผลํ. อรหตฺตํ สเมปชฺชิํ. สมาปชฺชามิ. สมาปนฺโน. อรหตฺตสฺส ลาภิมฺหิ. วสิมฺหิ. อรหตฺตํ สจฺฉิกตํ มยาติ สมฺปชานมุสา ภณนฺตสฺส อาปตฺติ ปาราชิกสฺส.}

\prose{ตีหากาเรหิ ราโค เม จตฺโต วนฺโต มุตฺโต ปหีโน ปฏินิสฺสฏฺโฐ อุกฺเขฏิโต สมุกฺเขฏิโตติ สมฺปชานมุสา ภณนฺตสฺส อาปตฺติ ปาราชิกสฺส.}

\prose{ตีหากาเรหิ โทโส เม จตฺโต วนฺโต มุตฺโต ปหีโน ปฏินิสฺสฏฺโฐ อุกฺเขฏิโต สมุกฺเขฏิโตติ สมฺปชานมุสา ภณนฺตสฺส อาปตฺติ ปาราชิกสฺส.}

\prose{ตีหากาเรหิ โมโห เม จตฺโต วนฺโต มุตฺโต ปหีโน ปฏินิสฺสฏฺโฐ อุกฺเขฏิโต สมุกฺเขฏิโตติ สมฺปชานมุสา ภณนฺตสฺส อาปตฺติ ปาราชิกสฺส.}

\prose{ตีหากาเรหิ ราคา เม จิตฺตํ วินีวรณนฺติ สมฺปชานมุสา ภณนฺตสฺส อาปตฺติ ปาราชิกสฺส.}

\prose{ตีหากาเรหิ โทสา เม จิตฺตํ วินีวรณนฺติ สมฺปชานมุสา ภณนฺตสฺส อาปตฺติ ปาราชิกสฺส.}

\prosewithcloser{ตีหากาเรหิ ฯเปฯ. สตฺตหากาเรหิ โมหา เม จิตฺตํ วินีวรณนฺติ สมฺปชานมุสา ภณนฺตสฺส อาปตฺติ ปาราชิกสฺส, ปุพฺเพวสฺส โหติ “มุสา ภณิสฺสนฺ”ติ, ภณนฺตสฺส โหติ “มุสา ภณามี”ติ, ภณิตสฺส โหติ “มุสา มยา ภณิตนฺ”ติ, วินิธาย ทิฏฺฐิํ, วินิธาย ขนฺติํ, วินิธาย รุจิํ, วินิธาย ภาวํ.}{สุทฺธิกํ นิฏฺฐิตํ.}
\proseitem{208}{\csromanfolio{126}ตีหากาเรหิ ปฐมญฺจ ฌานํ ทุติยญฺจ ฌานํ สมาปชฺชิํ. สมาปชฺชามิ. สมาปนฺโน. ปฐมสฺส จ ฌานสฺส ทุติยสฺส จ ฌานสฺส ลาภิมฺหิ. วสิมฺหิ. ปฐมญฺจ ฌานํ ทุติยญฺจ ฌานํ สจฺฉิกตํ มยาติ สมฺปชานมุสา ภณนฺตสฺส อาปตฺติ ปาราชิกสฺส ฯเปฯ.}

\prose{ตีหากาเรหิ ปฐมญฺจ ฌานํ ตติยญฺจ ฌานํ สมาปชฺชิํ. สมาปชฺชามิ. สมาปนฺโน. ปฐมสฺส จ ฌานสฺส ตติยสฺส จ ฌานสฺส ลาภิมฺหิ. วสิมฺหิ. ปฐมญฺจ ฌานํ ตติยญฺจ ฌานํ สจฺฉิกตํ มยาติ สมฺปชานมุสา ภณนฺตสฺส อาปตฺติ ปาราชิกสฺส.}

\prose{ตีหากาเรหิ ปฐมญฺจ ฌานํ จตุตฺถญฺจ ฌานํ สมาปชฺชิํ. สมาปชฺชามิ. สมาปนฺโน. ปฐมสฺส จ ฌานสฺส จตุตฺถสฺส จ ฌานสฺส ลาภิมฺหิ. วสิมฺหิ. ปฐมญฺจ ฌานํ จตุตฺถญฺจ ฌานํ สจฺฉิกตํ มยาติ สมฺปชานมุสา ภณนฺตสฺส อาปตฺติ ปาราชิกสฺส.}

\prose{ตีหากาเรหิ ปฐมญฺจ ฌานํ สุญฺญตญฺจ วิโมกฺขํ. ปฐมญฺจ ฌานํ อนิมิตฺตญฺจ วิโมกฺขํ. ปฐมญฺจ ฌานํ อปฺปณิหิตญฺจ วิโมกฺขํ สมาปชฺชิํ. สมาปชฺชามิ. สมาปนฺโน. ปฐมสฺส จ ฌานสฺส อปฺปณิหิตสฺส จ วิโมกฺขสฺส ลาภิมฺหิ. วสิมฺหิ. ปฐมญฺจ ฌานํ อปฺปณิหิโต จ วิโมกฺโข สจฺฉิกโต มยาติ สมฺปชานมุสา ภณนฺตสฺส อาปตฺติ ปาราชิกสฺส.}

\prose{ตีหากาเรหิ ปฐมญฺจ ฌานํ สุญฺญตญฺจ สมาธิํ. ปฐมญฺจ ฌานํ อนิมิตฺตญฺจ สมาธิํ. ปฐมญฺจ ฌานํ อปฺปณิหิตญฺจ สมาธิํ สมาปชฺชิํ. สมาปชฺชามิ. สมาปนฺโน. ปฐมสฺส จ ฌานสฺส อปฺปณิหิตสฺส จสมาธิสฺส ลาภิมฺหิ. วสิมฺหิ. ปฐมญฺจ ฌานํ อปฺปณิหิโต จ สมาธิ สจฺฉิกโต มยาติ สมฺปชานมุสา ภณนฺตสฺส อาปตฺติ ปาราชิกสฺส.}

\prose{ตีหากาเรหิ ปฐมญฺจ ฌานํ สุญฺญตญฺจ สมาปตฺติํ. ปฐมญฺจ ฌานํ อนิมิตฺตญฺจ สมาปตฺติํ. ปฐมญฺจ ฌานํ อปฺปณิหิตญฺจ สมาปตฺติํ สมาปชฺชิํ. สมาปชฺชามิ. สมาปนฺโน. ปฐมสฺส จ ฌานสฺส อปฺปณิหิตาย จ สมาปตฺติยา ลาภิมฺหิ. วสิมฺหิ. ปฐมญฺจ ฌานํ อปฺปณิหิตา จ สมาปตฺติ สจฺฉิกตา มยาติ สมฺปชานมุสา ภณนฺตสฺส อาปตฺติ ปาราชิกสฺส.}

\prose{ตีหากาเรหิ ปฐมญฺจ ฌานํ ติสฺโส จ วิชฺชา สมาปชฺชิํ. สมาปชฺชามิ. สมาปนฺโน. ปฐมสฺส จ ฌานสฺส ติสฺสนฺนญฺจ วิชฺชานํ ลาภิมฺหิ. วสิมฺหิ. ปฐมญฺจ ฌานํ ติสฺโส จ วิชฺชา สจฺฉิกตา มยาติ สมฺปชานมุสา ภณนฺตสฺส อาปตฺติ ปาราชิกสฺส.}

\prose{\csromanfolio{127}ตีหากาเรหิ ปฐมญฺจ ฌานํ จตฺตาโร จ สติปฏฺฐาเน. ปฐมญฺจ ฌานํ จตฺตาโร จ สมฺมปฺปธาเน. ปฐมญฺจ ฌานํ จตฺตาโร จ อิทฺธิปาเท สมาปชฺชิํ. สมาปชฺชามิ. สมาปนฺโน. ปฐมสฺส จ ฌานสฺส จตุนฺนญฺจ อิทฺธิปาทานํ ลาภิมฺหิ. วสิมฺหิ. ปฐมญฺจ ฌานํ จตฺตาโร จ อิทฺธิปาทา สจฺฉิกตา มยาติ สมฺปชานมุสา ภณนฺตสฺส อาปตฺติ ปาราชิกสฺส.}

\prose{ตีหากาเรหิ ปฐมญฺจ ฌานํ ปญฺจ จ อินฺทฺริยานิ. ปฐมญฺจ ฌานํ ปญฺจ จ พลานิ สมาปชฺชิํ. สมาปชฺชามิ. สมาปนฺโน. ปฐมสฺส จ ฌานสฺส ปญฺจนฺนญฺจ พลานํ ลาภิมฺหิ. วสิมฺหิ. ปฐมญฺจ ฌานํ ปญฺจ จ พลานิ สจฺฉิกตานิ มยาติ สมฺปชานมุสา ภณนฺตสฺส อาปตฺติ ปาราชิกสฺส.}

\proseitem{209}{ตีหากาเรหิ ปฐมญฺจ ฌานํ สตฺต จ โพชฺฌงฺเค สมาปชฺชิํ. สมาปชฺชามิ. สมาปนฺโน. ปฐมสฺส จ ฌานสฺส สตฺตนฺนญฺจ โพชฺฌงฺคานํ ลาภิมฺหิ. วสิมฺหิ. ปฐมญฺจ ฌานํ สตฺต จ โพชฺฌงฺคา สจฺฉิกตา มยาติ สมฺปชานมุสา ภณนฺตสฺส อาปตฺติ ปาราชิกสฺส.}

\prose{ตีหากาเรหิ ปฐมญฺจ ฌานํ อริยญฺจ อฏฺฐงฺคิกํ มคฺคํ สมาปชฺชิํ. สมาปชฺชามิ. สมาปนฺโน. ปฐมสฺส จ ฌานสฺส อริยสฺส จ อฏฺฐงฺคิกสฺส มคฺคสฺส ลาภิมฺหิ. วสิมฺหิ. ปฐมญฺจ ฌานํ อริโย จ อฏฺฐงฺคิโก มคฺโค สจฺฉิกโต มยาติ สมฺปชานมุสา ภณนฺตสฺส อาปตฺติ ปาราชิกสฺส.}

\prose{ตีหากาเรหิ ปฐมญฺจ ฌานํ โสตาปตฺติผลญฺจ. ปฐมญฺจ ฌานํ สกทาคามิ\-ผลญฺจ. ปฐมญฺจ ฌานํ อนาคามิผลญฺจ. ปฐมญฺจ ฌานํ อรหตฺตญฺจ สเมปชฺชิํ. สมาปชฺชามิ. สมาปนฺโน. ปฐมสฺส จ ฌานสฺส อรหตฺตสฺส จ ลาภิมฺหิ. วสิมฺหิ. ปฐมญฺจ ฌานํ อรหตฺตญฺจ สจฺฉิกตํ มยาติ สมฺปชานมุสา ภณนฺตสฺส อาปตฺติ ปาราชิกสฺส.}

\prose{ตีหากาเรหิ ปฐมญฺจ ฌานํ สมาปชฺชิํ. สมาปชฺชามิ. สมาปนฺโน. ปฐมสฺส จ ฌานสฺส ลาภิมฺหิ. วสิมฺหิ. ปฐมญฺจ ฌานํ สจฺฉิกตํ มยาติ ราโค จ เม จตฺโต. โทโส จ เม จตฺโต. โมโห จ เม จตฺโต วนฺโต มุตฺโต ปหีโน ปฏินิสฺสฏฺโฐ อุกฺเขฏิโต สมุกฺเขฏิโตติ สมฺปชานมุสา ภณนฺตสฺส อาปตฺติ ปาราชิกสฺส.}

\prosewithcloser{ตีหากาเรหิ ฯเปฯ. สตฺตหากาเรหิ ปฐมญฺจ ฌานํ สมาปชฺชิํ. สมาปชฺชามิ. สมาปนฺโน. ปฐมสฺส จ ฌานสฺส ลาภิมฺหิ. วสิมฺหิ. ปฐมญฺจ ฌานํ สจฺฉิกตํ มยา ราคา จ เม จิตฺตํ วินีวรณํ. โทสา จ เม จิตฺตํ วินีวรณํ. \csromanfolio{128}โมหา จ เม จิตฺตํ วินีวรณนฺติ สมฺปชานมุสา ภณนฺตสฺส อาปตฺติ ปาราชิกสฺส, ปุพฺเพวสฺส โหติ “มุสา ภณิสฺสนฺ”ติ, ภณนฺตสฺส โหติ “มุสา ภณามี”ติ, ภณิตสฺส โหติ “มุสา มยา ภณิตนฺ”ติ, วินิธาย ทิฏฺฐิํ, วินิธาย ขนฺติํ, วินิธาย รุจิํ, วินิธาย ภาวํ.}{ขณฺฑจกฺกํ นิฏฺฐิตํ.}
\proseitem{210}{ตีหากาเรหิ ทุติยญฺจ ฌานํ ตติยญฺจ ฌานํ สมาปชฺชิํ. สมาปชฺชามิ. สมาปนฺโน. ทุติยสฺส จ ฌานสฺส ตติยสฺส จ ฌานสฺส ลาภิมฺหิ. วสิมฺหิ. ทุติยญฺจ ฌานํ ตติยญฺจ ฌานํ สจฺฉิกตํ มยาติ สมฺปชานมุสา ภณนฺตสฺส อาปตฺติ ปาราชิกสฺส.}

\prose{ตีหากาเรหิ ทุติยญฺจ ฌานํ จตุตฺถญฺจ ฌานํ สมาปชฺชิํ. สมาปชฺชามิ. สมาปนฺโน. ทุติยสฺส จ ฌานสฺส จตุตฺถสฺส จ ฌานสฺส ลาภิมฺหิ. วสิมฺหิ. ทุติยญฺจ ฌานํ จตุตฺถญฺจ ฌานํ สจฺฉิกตํ มยาติ สมฺปชานมุสา ภณนฺตสฺส อาปตฺติ ปาราชิกสฺส.}

\prose{ตีหากาเรหิ ทุติยญฺจ ฌานํ สุญฺญตญฺจ วิโมกฺขํ. อนิมิตฺตญฺจ วิโมกฺขํ. อปฺปณิหิตญฺจ วิโมกฺขํ. สุญฺญตญฺจ สมาธิํ. อนิมิตฺตญฺจ สมาธิํ. อปฺปณิหิตญฺจ สมาธิํ. สุญฺญตญฺจ สมาปตฺติํ. อนิมิตฺตญฺจ สมาปตฺติํ. อปฺปณิหิตญฺจ สมาปตฺติํ. ติสฺโส จ วิชฺชา. จตฺตาโร จ สติปฏฺฐาเน. จตฺตาโร จ สมฺมปฺปธาเน. จตฺตาโร จ อิทฺธิปาเท. ปญฺจ จ อินฺทฺริยานิ. ปญฺจ จ พลานิ. สตฺต จ โพชฺฌงฺเค. อริยญฺจ อฏฺฐงฺคิกํ มคฺคํ. โสตาปตฺติผลญฺจ. สกทาคามิ\-ผลญฺจ. อนาคามิผลญฺจ. อรหตฺตญฺจ สมาปชฺชิํ. สมาปชฺชามิ. สมาปนฺโน. ทุติยสฺส จ ฌานสฺส อรหตฺตสฺส จ ลาภิมฺหิ. วสิมฺหิ. ทุติยญฺจ ฌานํ อรหตฺตญฺจ สจฺฉิกตํ มยาติ สมฺปชานมุสา ภณนฺตสฺส อาปตฺติ ปาราชิกสฺส.}

\prose{ตีหากาเรหิ ทุติยญฺจ ฌานํ สมาปชฺชิํ. สมาปชฺชามิ. สมาปนฺโน. ทุติยสฺส จ ฌานสฺส ลาภิมฺหิ. วสิมฺหิ. ทุติยญฺจ ฌานํ สจฺฉิกตํ มยา ราโค จ เม จตฺโต. โทโส จ เม จตฺโต. โมโห จ เม จตฺโต วนฺโต มุตฺโต ปหีโน ปฏินิสฺสฏฺโฐ อุกฺเขฏิโต สมุกฺเขฏิโต. ราคา จ เม จิตฺตํ วินีวรณํ. โทสา จ เม จิตฺตํ วินีวรณํ. โมหา จ เม จิตฺตํ วินีวรณนฺติ สมฺปชานมุสา ภณนฺตสฺส อาปตฺติ ปาราชิกสฺส.}

\prose{\csromanfolio{129}ตีหากาเรหิ ฯเปฯ. สตฺตหากาเรหิ ทุติยญฺจ ฌานํ ปฐมญฺจ ฌานํ สมาปชฺชิํ. สมาปชฺชามิ. สมาปนฺโน. ทุติยสฺส จ ฌานสฺส ปฐมสฺส จ ฌานสฺส ลาภิมฺหิ. วสีมฺหิ. ทุติยญฺจ ฌานํ ปฐมญฺจ ฌานํ สจฺฉิกตํ มยาติ สมฺปชานมุสา ภณนฺตสฺส อาปตฺติ ปาราชิกสฺส ฯเปฯ วินิธาย ภาวํ.}

\vspace{\tipitakaparskip}
\csromancenter{พทฺธจกฺกํ.}
\csromancenter{เอวํ เอเกกํ มูลํ กาตุน พทฺธจกฺกํ ปริวตฺตกํ กตฺตพฺพํ.}
\csromancenter{อิทํ สํขิตฺตํ.}
\proseitem{211}{ตีหากาเรหิ ตติยญฺจ ฌานํ จตุตฺถญฺจ ฌานํ ฯเปฯ ตติยญฺจ ฌานํ อรหตฺตญฺจ สมาปชฺชิํ. สมาปชฺชามิ. สมาปนฺโน. ตติยสฺส จ ฌานสฺส อรหตฺตสฺส จ ลาภีมฺหิ. วสีมฺหิ. ตติยญฺจ ฌานํ อรหตฺตญฺจ สจฺฉิกตํ มยาติ สมฺปชานมุสา ภณนฺตสฺส อาปตฺติ ปาราชิกสฺส.}

\prose{ตีหากาเรหิ ตติยญฺจ ฌานํ สมาปชฺชิํ. สมาปชฺชามิ. สมาปนฺโน. ตติยสฺส จ ฌานสฺส ลาภีมฺหิ. วสีมฺหิ. ตติยญฺจ ฌานํ สจฺฉิกตํ มยา ราโค จ เม จตฺโต. โทโส จ เม จตฺโต. โมโห จ เม จตฺโต วนฺโต มุตฺโต ปหีโน ปฏินิสฺสฏฺโฐ อุกฺเขฏิโต สมุกฺเขฏิโต. ราคา จ เม จิตฺตํ วินีวรณํ. โทสา จ เม จิตฺตํ วินีวรณํ. โมหา จ เม จิตฺตํ วินีวรณนฺติ สมฺปชานมุสา ภณนฺตสฺส อาปตฺติ ปาราชิกสฺส.}

\prose{ตีหากาเรหิ ตติยญฺจ ฌานํ ปฐมญฺจ ฌานํ. ตติยญฺจ ฌานํ ทุติยญฺจ ฌานํ สมาปชฺชิํ. สมาปชฺชามิ. สมาปนฺโน. ตติยสฺส จ ฌานสฺส ทุติยสฺส จ ฌานสฺส ลาภีมฺหิ. วสีมฺหิ. ตติยญฺจ ฌานํ ทุติยญฺจ ฌานํ สจฺฉิกตํ มยาติ สมฺปชานมุสา ภณนฺตสฺส อาปตฺติ ปาราชิกสฺส.}

\prose{ตีหากาเรหิ โมหา จ เม จิตฺตํ วินีวรณํ ปฐมญฺจ ฌานํ ฯเปฯ ทุติยญฺจ ฌานํ. ตติยญฺจ ฌานํ. จตุตฺถญฺจ ฌานํ สมาปชฺชิํ. สมาปชฺชามิ. สมาปนฺโน. โมหา จ เม จิตฺตํ วินีวรณํ จตุตฺถสฺส จ ฌานสฺส ลาภิมฺหิ. วสิมฺหิ. โมหา จ เม จิตฺตํ วินีวรณํ จตุตฺถญฺจ ฌานํ สจฺฉิกตํ มยาติ สมฺปชนมุสา ภณนฺตสฺส อาปตฺติ ปาราชิกสฺส.}

\proseitem{212}{ตีหากาเรหิ โมหา จ เม จิตฺตํ วินีวรณํ สุญฺญตญฺจ วิโมกฺขํ. อนิมิตฺตญฺจ วิโมกฺขํ. อปฺปณิหิตญฺจ วิโมกฺขํ สมาปชฺชิํ. สมาปชฺชามิ. สมาปนฺโน. โมหา จ เม จิตฺตํ วินีวรณํ อปฺปณิหิตสฺส จ \csromanfolio{130}วิโมกฺขสฺส ลาภิมฺหิ. วสิมฺหิ. โมหา จ เม จิตฺตํ วินีวรณํ อปฺปณิหิโต จ วิโมกฺโข สจฺฉิกโต มยาติ สมฺปชานมุสา ภณนฺตสฺส อาปตฺติ ปาราชิกสฺส.}

\prose{ตีหากาเรหิ โมหา จ เม จิตฺตํ วินีวรณํ สุญฺญตญฺจ สมาธิํ. อนิมิตฺตญฺจ สมาธิํ. อปฺปณิหิตญฺจ สมาธิํ. สมาปชฺชิํ. สมาปชฺชามิ. สมาปนฺโน. โมหา จ เม จิตฺตํ วินีวรณํ อปฺปณิหิตสฺส จ สมาธิสฺส ลาภิมฺหิ. วสิมฺหิ. โมหา จ เม จิตฺตํ วินีวรณํ อปฺปณิหิโต จ สมาธิ สจฺฉิกโต มยาติ สมฺปชานมุสา ภณนฺตสฺส อาปตฺติ ปาราชิกสฺส.}

\prose{ตีหากาเรหิ โมหา จ เม จิตฺตํ วินีวรณํสุญฺญตญฺจ สมาปตฺติํ. อนิมิตฺตญฺจ สมาปตฺติํ. อปฺปณิหิตญฺจ สมาปตฺติํ สมาปชฺชิํ. สมาปชฺชามิ. สมาปนฺโน. โมหา จ เม จิตฺตํ วินีวรณํ อปฺปณิหิตาย จ สมาปตฺติยา ลาภิมฺหิ. วสิมฺหิ. โมหา จ เม จิตฺตํ วินีวรณํ อปฺปณิหิตา จ สมาปตฺติ สจฺฉิกตา มยาติ สมฺปชานมุสา ภณนฺตสฺส อาปตฺติ ปาราชิกสฺส.}

\prose{ตีหากาเรหิ โมหา จ เม จิตฺตํ วินีวรณํ ติสฺโส จ วิชฺชา สมาปชฺชิํ. สมาปชฺชามิ. สมาปนฺโน. โมหา จ เม จิตฺตํ วินีวรณํ ติสฺสนฺนญฺจ วิชฺชานํ ลาภิมฺหิ. วสิมฺหิ. โมหา จ เม จิตฺตํ วินีวรณํ ติสฺโส จ วิชฺชา สจฺฉิกตา มยาติ สมฺปชานมุสา ภณนฺตสฺส อาปตฺติ ปาราชิกสฺส.}

\prose{ตีหากาเรหิ โมหา จ เม จิตฺตํ วินีวรณํ จตฺตาโร จ สติปฏฺฐาเน. จตฺตาโร จ สมฺมปฺปธาเน. จตฺตาโร จ อิทฺธิปาเท สมาปชฺชิํ. สมาปชฺชามิ. สมาปนฺโน. โมหา จ เม จิตฺตํ วินีวรณํ จตุนฺนญฺจ อิทฺธิปาทานํ ลาภิมฺหิ. วสิมฺหิ. โมหา จ เม จิตฺตํ วินีวรณํ จตฺตาโร จ อิทฺธิปาทา สจฺฉิกตา มยาติ สมฺปชานมุสา ภณนฺตสฺส อาปตฺติ ปาราชิกสฺส.}

\proseitem{213}{ตีหากาเรหิ โมหา จ เม จิตฺตํ วินีวรณํ ปญฺจ จ อินฺทฺริยานิ. ปญฺจ จ พลานิ สมาปชฺชิํ. สมาปชฺชามิ. สมาปนฺโน. โมหา จ เม จิตฺตํ วินีวรณํ ปญฺจนฺนญฺจ พลานํ ลาภิมฺหิ. วสิมฺหิ. โมหา จ เม จิตฺตํ วินีวรณํ ปญฺจ จ พลานิ สจฺฉิกตานิ มยาติ สมฺปชานมุสา ภณนฺตสฺส อาปตฺติ ปาราชิกสฺส.}

\prose{\csromanfolio{131}ตีหากาเรหิ โมหา จ เม จิตฺตํ วินีวรณํ สตฺต จ โพชฺฌงฺเค สมาปชฺชิํ. สมาปชฺชามิ. สมาปนฺโน. โมหา จ เม จิตฺตํ วินีวรณํ สตฺตนฺนญฺจ โพชฺฌงฺคานํ ลาภิมฺหิ. วสิมฺหิ. โมหา จ เม จิตฺตํ วินีวรณํ สตฺต จ โพชฺฌงฺคา สจฺฉิกตา มยาติ สมฺปชานมุสา ภณนฺตสฺส อาปตฺติ ปาราชิกสฺส.}

\prose{ตีหากาเรหิ โมหา จ เม จิตฺตํ วินีวรณํ อริยญฺจ อฏฺฐงฺคิกํ มคฺคํ สมาปชฺชิํ. สมาปชฺชามิ. สมาปนฺโน. โมหา จ เม จิตฺตํ วินีวรณํ อริยสฺส จ อฏฺฐงฺคิกสฺส มคฺคสฺส ลาภิมฺหิ. วสิมฺหิ. โมหา จ เม จิตฺตํ วินีวรณํ อริโย จ อฏฺฐงฺคิโก มคฺโค สจฺฉิกโต มยาติ สมฺปชานมุสา ภณนฺตสฺส อาปตฺติ ปาราชิกสฺส.}

\prose{ตีหากาเรหิ โมหา จ เม จิตฺตํ วินีวรณํ โสตาปตฺติผลญฺจ. สกทาคามิ\-ผลญฺจ. อนาคามิผลญฺจ. อรหตฺตญฺจ สมาปชฺชิํ. สมาปชฺชามิ. สมาปนฺโน. โมหา จ เม จิตฺตํ วินีวรณํ อรหตฺตสฺส จ ลาภิมฺหิ. วสิมฺหิ. โมหา จ เม จิตฺตํ วินีวรณํ อรหตฺตญฺจ สจฺฉิกตํ มยาติ สมฺปชานมุสา ภณนฺตสฺส อาปตฺติ ปาราชิกสฺส.}

\prose{ตีหากาเรหิ โมหา จ เม จิตฺตํ วินีวรณํ ราโค จ เม จตฺโต. โทโส จ เม จตฺโต. โมโห จ เม จตฺโต วนฺโต มุตฺโต ปหีโน ปฏินิสฺสฏฺโฐ อุกฺเขฏิโต สมุกฺเขฏิโตติ สมฺปชานมุสา ภณนฺตสฺส อาปตฺติ ปาราชิกสฺส.}

\prosewithcloser{ตีหากาเรหิ ฯเปฯ. สตฺตหากาเรหิ โมหา จ เม จิตฺตํ วินีวรณํ ราคา จ เม จิตฺตํ วินีวรณํ. โทสา จ เม จิตฺตํ วินีวรณนฺติ สมฺปชานมุสา ภณนฺตสฺส อาปตฺติ ปาราชิกสฺส, ปุพฺเพวสฺส โหติ “มุสา ภณิสฺสนฺ”ติ, ภณนฺตสฺส โหติ “มุสา ภณามี”ติ, ภณิตสฺส โหติ “มุสา มยา ภณิตนฺ”ติ, วินิธาย ทิฏฺฐิํ, วินิธาย ขนฺติํ, วินิธาย รุจิํ, วินิธาย ภาวํ.}{เอกมูลกํ นิฏฺฐิตํ\csromansharedfootnote{65650f5ca0bf0374}{เอกมูลกํ สงฺขิตฺตํ นิฏฺฐิตํ (สฺยา)}.}
\csromancenter{ยถา เอกมูลกํ วิตฺถาริตํ, เอวเมว ทุมูลกาทิปิ วิตฺถาเรตพฺพํ.}
\csromancenter{อิทํ สพฺพมูลกํ}
\proseitemwithcloser{214}{\csromanfolio{132}ตีหากาเรหิ ฯเปฯ. สตฺตหากาเรหิ ปฐมญฺจ ฌานํ ทุติยญฺจ ฌานํ ตติยญฺจ ฌานํ จตุตฺถญฺจ ฌานํ สุญฺญตญฺจ วิโมกฺขํ อนิมิตฺตญฺจ วิโมกฺขํ อปฺปณิหิตญฺจ วิโมกฺขํ สุญฺญตญฺจ สมาธิํ อนิมิตฺตญฺจ สมาธิํ อปฺปณิหิตญฺจ สมาธิํ สุญฺญตญฺจ สมาปตฺติํ อนิมิตฺตญฺจ สมาปตฺติํ อปฺปณิหิตญฺจ สมาปตฺติํ ติสฺโส จ วิชฺชา จตฺตาโร จ สติปฏฺฐาเน จตฺตาโร จ สมฺมปฺปธาเน จตฺตาโร จ อิทฺธิปาเท ปญฺจ จ อินฺทฺริยานิ ปญฺจ จ พลานิ สตฺต จ โพชฺฌงฺเค อริยญฺจ อฏฺฐงฺคิกํ มคฺคํ โสตาปตฺติผลญฺจ สกทาคามิ\-ผลญฺจ อนาคามิผลญฺจ อรหตฺตญฺจ สมาปชฺชิํ, สมาปชฺชามิ, สมาปนฺโน ฯเปฯ ราโค จ เม จตฺโต, โทโส จ เม จตฺโต, โมโห จ เม จตฺโต วนฺโต มุตฺโต ปหีโน ปฏินิสฺสฏฺโฐ อุกฺเขฏิโต สมุกฺเขฏิโต ราคา จ เม จิตฺตํ วินีวรณํ โทสา จ เม จิตฺตํ วินีวรณํ โมหา จ เม จิตฺตํ วินีวรณนฺติ สมฺปชานมุสา ภณนฺตสฺส อาปตฺติ ปาราชิกสฺส, ปุพฺเพวสฺส โหติ “มุสา ภณิสฺสนฺ”ติ, ภณนฺตสฺส โหติ “มุสา ภณามี”ติ, ภณิตสฺส โหติ “มุสา มยา ภณิตนฺ”ติ, วินิธาย ทิฏฺฐิํ, วินิธาย ขนฺติํ, วินิธาย รุจิํ, วินิธาย ภาวํ.}{สพฺพมูลกํ นิฏฺฐิตํ.}
\nitthitam{สุทฺธิก\-วาร\-กถา นิฏฺฐิตา.}
\proseitem{215}{ตีหากาเรหิ ปฐมํ ฌานํ สมาปชฺชินฺติ วตฺตุกาโม ทุติยํ ฌานํ สมาปชฺชินฺติ สมฺปชานมุสา ภณนฺตสฺส ปฏิวิชานนฺตสฺส อาปตฺติ ปาราชิกสฺส, น ปฏิวิชานนฺตสฺส อาปตฺติ ถุลฺลจฺจยสฺส.}

\prose{ตีหากาเรหิ ปฐมํ ฌานํ สมาปชฺชินฺติ วตฺตุกาโม ตติยํ ฌานํ สมาปชฺชินฺติ สมฺปชานมุสา ภณนฺตสฺส ปฏิวิชานนฺตสฺส อาปตฺติ ปาราชิกสฺส, น ปฏิวิชานนฺตสฺส อาปตฺติ ถุลฺลจฺจยสฺส.}

\prose{ตีหากาเรหิ ปฐมํ ฌานํ สมาปชฺชินฺติ วตฺตุกาโม จตุตฺถํ ฌานํ สมาปชฺชินฺติ สมฺปชานมุสา ภณนฺตสฺส ปฏิวิชานนฺตสฺส อาปตฺติ ปาราชิกสฺส, น ปฏิวิชานนฺตสฺส อาปตฺติ ถุลฺลจฺจยสฺส.}

\prosewithcloser{\csromanfolio{133}ตีหากาเรหิ ฯเปฯ. สตฺตหากาเรหิ ปฐมํ ฌานํ สมาปชฺชินฺติ วตฺตุกาโม สุญฺญตํ วิโมกฺขํ. อนิมิตฺตํ วิโมกฺขํ. อปฺปณิหิตํ วิโมกฺขํ. สุญฺญตํ สมาธิํ. อนิมิตฺตํ สมาธิํ. อปฺปณิหิตํ สมาธิํ. สุญฺญตํ สมาปตฺติํ. อนิมิตฺตํ สมาปตฺติํ. อปฺปณิหิตํ สมาปตฺติํ. ติสฺโส วิชฺชา. จตฺตาโร สติปฏฺฐาเน. จตฺตาโร สมฺมปฺปธาเน. จตฺตาโร อิทฺธิปาเท. ปญฺจินฺทฺริยานิ. ปญฺจ พลานิ. สตฺต โพชฺฌงฺเค. อริยํ อฏฺฐงฺคิกํ มคฺคํ. โสตาปตฺติผลํ. สกทาคามิผลํ. อนาคามิผลํ. อรหตฺตํ สมาปชฺชิํ ฯเปฯ ราโค เม จตฺโต. โทโส เม จตฺโต. โมโห เม จตฺโต วนฺโต มุตฺโต ปหีโน ปฏินิสฺสฏฺโฐ อุกฺเขฏิโต สมุกฺเขฏิโต. ราคา เม จิตฺตํ วินีวรณํ. โทสา เม จิตฺตํ วินีวรณํ. โมหา เม จิตฺตํ วินีวรณนฺติ สมฺปชานมุสา ภณนฺตสฺส ปฏิวิชานนฺตสฺส อาปตฺติ ปาราชิกสฺส, น ปฏิวิชานนฺตสฺส อาปตฺติ ถุลฺลจฺจยสฺส, ปุพฺเพวสฺส โหติ “มุสา ภณิสฺสนฺ”ติ, ภณนฺตสฺส โหติ “มุสา ภณามี”ติ, ภณิตสฺส โหติ “มุสา มยา ภณิตนฺ”ติ, วินิธาย ทิฏฺฐิํ, วินิธาย ขนฺติํ, วินิธาย รุจิํ, วินิธาย ภาวํ.}{วตฺถุ\-วิสารกสฺส เอกมูลกสฺส ขณฺฑจกฺกํ นิฏฺฐิตํ.}
\proseitem{216}{ตีหากาเรหิ ทุติยํ ฌานํ สมาปชฺชินฺติ วตฺตุกาโม ตติยํ ฌานํ สมาปชฺชินฺติ สมฺปชานมุสา ภณนฺตสฺส ปฏิวิชานนฺตสฺส อาปตฺติ ปาราชิกสฺส, น ปฏิวิชานนฺตสฺส อาปตฺติ ถุลฺลจฺจยสฺส.}

\prose{ตีหากาเรหิ ทุติยํ ฌานํ สมาปชฺชินฺติ วตฺตุกาโม จตุตฺถํ ฌานํ สมาปชฺชินฺติ ฯเปฯ โมหา เม จิตฺตํ วินีวรณนฺติ สมฺปชานมุสา ภณนฺตสฺส ปฏิวิชานนฺตสฺส อาปตฺติ ปาราชิกสฺส, น ปติวิชนนฺตสฺส อาปตฺติ ถุลฺลจฺจยสฺส.}

\prose{ตีหากาเรหิ ฯเปฯ. สตฺตหากาเรหิ ทุติยํ ฌานํ สมาปชฺชินฺติ วตฺตุกาโม ปฐมํ ฌานํ สมาปชฺชินฺติ สมฺปชานมุสา ภณนฺตสฺส ปฏิวิชานนฺตสฺส อาปตฺติ ปาราชิกสฺส, น ปฏิวิชานนฺตสฺส อาปตฺติ ถุลฺลจฺจยสฺส ฯเปฯ วินิธาย ภาวํ.}

\vspace{\tipitakaparskip}
\csromancenter{วตฺถุ\-วิสารกสฺส เอกมูลกสฺส พทฺธจกฺกํ.}
\csromancenter{มูลํ สํขิตฺตํ.}
\proseitem{217}{\csromanfolio{134}ตีหากาเรหิ โมหา เม จิตฺตํ วินีวรณนฺติ วตฺตุกาโม ปฐมํ ฌานํ สเมปชฺชินฺติ สมฺปชานมุสา ภณนฺตสฺส ปฏิวิชานนฺตสฺส อาปตฺติ ปาราชิกสฺส, น ปฏิวิชานนฺตสฺส อาปตฺติ ถุลฺลจฺจยสฺส.}

\prosewithcloser{ตีหากาเรหิ ฯเปฯ. สตฺตหากาเรหิ โมหา เม จิตฺตํ วินีวรณนฺติ วตฺตุกาโม โทสา เม จิตฺตํ วินีวรณนฺติ สมฺปชานมุสา ภณนฺตสฺส ปฏิวิชานนฺตสฺส อาปตฺติ ปาราชิกสฺส, น ปฏิวิชานนฺตสฺส อาปตฺติ ถุลฺลจฺจยสฺส ฯเปฯ วินิธาย ภาวํ.}{วตฺถุ\-วิสารกสฺส เอกมูลกํ นิฏฺฐิตํ.}
\csromancenter{ยถา เอกมูลกํ วิตฺถาริตํ, เอวเมว ทุมูลกาทิปิ วิตฺถาเรตพฺพํ.}
\csromancenter{อิทํ สพฺพมูลกํ}
\proseitem{218}{ตีหากาเรหิ ฯเปฯ. สตฺตหากาเรหิ ปฐมญฺจ ฌานํ ทุติยญฺจ ฌานํ ตติยญฺจ ฌานํ จตุตฺถญฺจ ฌานํ สุญฺญตญฺจ วิโมกฺขํ อนิมิตฺตญฺจ วิโมกฺขํ อปฺปณิหิตญฺจ วิโมกฺขํ สุญฺญตญฺจ สมาธิํ อนิมิตฺตญฺจ สมาธิํ อปฺปณิหิตญฺจ สมาธิํ สุญฺญตญฺจ สมาปตฺติํ อนิมิตฺตญฺจ สมาปตฺติํ อปฺปณิหิตญฺจ สมาปตฺติํ ติสฺโส จ วิชฺชา จตฺตาโร จ สติปฏฺฐาเน จตฺตาโร จ สมฺมปฺปธาเน จตฺตาโร จ อิทฺธิปาเท ปญฺจ จ อินฺทฺริยานิ ปญฺจ จ พลานิ สตฺต จ โพชฺฌงฺเค อริยญฺจ อฏฺฐงฺคิกํ มคฺคํ โสตาปตฺติผลญฺจ สกทาคามิ\-ผลญฺจ อนาคามิผลญฺจ อรหตฺตญฺจ สมาปชฺชิํ ฯเปฯ ราโค จ เม จตฺโต, โทสา จ เม จตฺโต, โมโห จ เม จตฺโต วนฺโต มุตฺโต ปหีโน ปฏินิสฺสฏฺโฐ อุกฺเขฏิโต สมุกฺเขฏิโต ราโค จ เม จิตฺตํ วินีวรณํ โทสา จ เม จิตฺตํ วินีวรณนฺติ วตฺตุกาโม โมหา เม จิตฺตํ วินีวรณนฺติ สมฺปชานมุสา ภณนฺตสฺส ปฏิวิชานนฺตสฺส อาปตฺติ ปาราชิกสฺส, น ปฏิวิชานนฺตสฺส อาปตฺติ ถุลฺลจฺจยสฺส.}

\prose{[\csromansymbolfoottext{[ ]}{เอตฺถนฺตเร ปาฐา สฺยามโปตฺถเก นตฺถิ.} 219. ตีหากาเรหิ ทุติยญฺจ ฌานํ ตติยญฺจ ฌานํ จตุตฺถญฺจ ฌานํ สุญฺญตญฺจ วิโมกฺขํ อนิมิตฺตญฺจ วิโมกฺขํ อปฺปณิหิตญฺจ วิโมกฺขํ สุญฺญตญฺจ สมาธิํ อนิมิตฺตญฺจ สมาธิํ อปฺปณิหิตญฺจ สมาธิํ สุญฺญตญฺจ สมาปตฺติํ อนิมิตฺตญฺจ สมาปตฺติํ อปฺปณิหิตญฺจ สมาปตฺติํ ติสฺโส จ วิชฺชา จตฺตาโร จ สติปฏฺฐาเน จตฺตาโร จ สมฺมปฺปธาเน จตฺตาโร จ อิทฺธิปาเท \csromanfolio{135}ปญฺจ จ อินฺทฺริยานิ ปญฺจ จ พลานิ สตฺต จ โพชฺฌงฺเค อริยญฺจ อฏฺฐงฺคิกํ มคฺคํ โสตาปตฺติผลญฺจ สกทาคามิ\-ผลญฺจ อนาคามิผลญฺจ อรหตฺตญฺจ สมาปชฺชิํ ฯเปฯ ราโค จ เม จตฺโต, โทโส จ เม จตฺโต, โมโห จ เม จตฺโต วนฺโต มุตฺโต ปหีโน ปฏินิสฺสฏฺโฐ อุกฺเขฏิโต สมุกฺเขฏิโต ราคา จ เม จิตฺตํ วินีวรณํ โทสา จ เม จิตฺตํ วินีวรณํ โมหา จ เม จิตฺตํ วินีวรณนฺติ วตฺตุกาโม ปฐมํ ฌานํ สมาปชฺชินฺติ สมฺปชานมุสา ภณนฺตสฺส ปฏิวิชานนฺตสฺส อาปตฺติ ปาราชิกสฺส, น ปฏิวิชานนฺตสฺส อาปตฺติ ถุลฺลจฺจยสฺส.}

\prose{ตีหากาเรหิ ตติยญฺจ ฌานํ จตุตฺถญฺจ ฌานํ ฯเปฯ โมหา จ เม จิตฺตํ วินีวรณํ ปฐมญฺจ ฌานํ สมาปชฺชินฺติ วตฺตุกาโม ทุติยํ ฌานํ สมาปชฺชินฺติ สมฺปชานมุสา ภณนฺตสฺส ปฏิวิชานนฺตสฺส อาปตฺติ ปาราชิกสฺส, น ปฏิวิชานนฺตสฺส อาปตฺติ ถุลฺลจฺจยสฺส.}

\prosewithcloser{ตีหากาเรหิ ฯเปฯ. สตฺตหากาเรหิ โมหา จ เม จิตฺตํ วินีวรณํ ปฐมญฺจ ฌานํ ทุติยญฺจ ฌานํ ตติยญฺจ ฌานํ จตุตฺถญฺจ ฌานํ ฯเปฯ ราคา จ เม จิตฺตํ วินีวรณนฺติ วตฺตุกาโม โทสา เม จิตฺตํ วินีวรณนฺติ สมฺปชานมุสา ภณนฺตสฺส ปฏิวิชานนฺตสฺส อาปตฺติ ปาราชิกสฺส, น ปฏิวิชานนฺตสฺส อาปตฺติ ถุลฺลจฺจยสฺส], . ปุพฺเพวสฺส โหติ “มุสา ภณิสฺสนฺ”ติ, ภณนฺตสฺส โหติ “มุสา ภณามี”ติ, ภณิตสฺส โหติ “มุสา มยา ภณิตนฺ”ติ, วินิธาย ทิฏฺฐิํ, วินิธาย ขนฺติํ, วินิธาย รุจิํ, วินิธาย ภาวํ.}{วตฺถุ\-วิสารกสฺส สพฺพมูลกํ นิฏฺฐิตํ.}
\nitthitam{วตฺถุ\-วิสารกสฺส จกฺกเปยฺยาลํ นิฏฺฐิตํ.}
\nitthitam{วตฺถุกาม\-วาร\-กถา นิฏฺฐิตา.}
\proseitem{220}{ตีหากาเรหิ โย เต วิหาเร วสิ, โส ภิกฺขุ ปฐมํ ฌานํ สมาปชฺชิ. สมาปชฺชติ. สมาปนฺโน. โส ภิกฺขุ ปฐมสฺส ฌานสฺส ลาภี. วสี. เตน ภิกฺขุนา ปฐมํ ฌานํ สจฺฉิกตนฺติ สมฺปชานมุสา ภณนฺตสฺส ปฏิวิชานนฺตสฺส อาปตฺติ ถุลฺลจฺจยสฺส, น ปฏิวิชานนฺตสฺส อาปตฺติ ทุกฺกฏสฺส.}

\prose{\csromanfolio{136}จตูหากาเรหิ. ปญฺจหากาเรหิ. ฉหากาเรหิ. สตฺตหากาเรหิ โย เต วิหาเร วสิ, โส ภิกฺขุ ปฐมํ ฌานํ สมาปชฺชิ. สมาปชฺชติ. สมาปนฺโน. โส ภิกฺขุ ปฐมสฺส ฌานสฺส ลาภี. วสี. เตน ภิกฺขุนา ปฐมํ ฌานํ สจฺฉิกตนฺติ สมฺปชานมุสา ภณนฺตสฺส ปฏิวิชานนฺตสฺส อาปตฺติ ถุลฺลจฺจยสฺส, น ปฏิวิชานนฺตสฺส อาปตฺติ ทุกฺกฏสฺส, ปุพฺเพวสฺส โหติ “มุสา ภณิสฺสนฺ”ติ, ภณนฺตสฺส โหติ “มุสา ภณามี”ติ, ภณิตสฺส โหติ “มุสา มยา ภณิตนฺ”ติ, วินิธาย ทิฏฺฐิํ, วินิธาย ขนฺติํ, วินิธาย รุจิํ, วินิธาย ภาวํ.}

\prose{ตีหากาเรหิ โย เต วิหาเร วสิ, โส ภิกฺขุ ทุติยํ ฌานํ. ตติยํ ฌานํ. จตุตฺถํ ฌานํ. สุญฺญตํ วิโมกฺขํ. อนิมิตฺตํ วิโมกฺขํ. อปฺปณิหิต วิโมกฺขํ. สุญฺญตํ สมาธิํ. อนิมิตฺตํ สมาธิํ. อปฺปณิหิตํ สมาธิํ. สุญฺญตํ สมาปตฺติํ. อนิมิตฺตํ สมาปตฺติํ. อปฺปณิหิตํ สมาปตฺติํ. ติสฺโส วิชฺชา. จตฺตาโร สติปฏฺฐาเน. จตฺตาโร สมฺมปฺปธาเน. จตฺตาโร อิทฺธิปาเท. ปญฺจ อินฺทฺริยานิ. ปญฺจ พลานิ. สตฺต โพชฺฌงฺเค. อริยํ อฏฺฐงฺคิกํ มคฺคํ. โสตาปตฺติผลํ. สกทาคามิผลํ. อนาคามิผลํ. อรหตฺตํ สมาปชฺชิ. สมาปชฺชติ. สมาปนฺโน. โส ภิกฺขุ อรหตฺตสฺส ลาภี. วสี. เตน ภิกฺขุนา อรหตฺตํ สจฺฉิกตนฺติ สมฺปชานมุสา ภณนฺตสฺส ปฏิวิชานนฺตสฺส อาปตฺติ ถุลฺลจฺจยสฺส, น ปฏิวิชานนฺตสฺส อาปตฺติ ทุกฺกฏสฺส.}

\prose{ตีหากาเรหิ โย เต วิหาเร วสิ, ตสฺส ภิกฺขุโน ราโค จตฺโต. โทโส จตฺโต. โมโห จตฺโต วนฺโต มุตฺโต ปหีโน ปฏินิสฺสฏฺโฐ อุกฺเขฏิโต สมุกฺเขฏิโตติ สมฺปชานมุสา ภณนฺตสฺส ปฏิวิชานนฺตสฺส อาปตฺติ ถุลฺลจฺจยสฺส, น ปฏิวิชานนฺตสฺส อาปตฺติ ทุกฺกฏสฺส.}

\prose{ตีหากาเรหิ ฯเปฯ. สตฺตหากาเรหิ โย เต วิหาเร วสิ, ตสฺส ภิกฺขุโน ราคา จิตฺตํ วินีวรณํ. โทสา จิตฺตํ วินีวรณํ. โมหา จิตฺตํ วินีวรณนฺติ สมฺปชานมุสา ภณนฺตสฺส ปฏิวิชานนฺตสฺส อาปตฺติ ถุลฺลจฺจยสฺส, น ปฏิวิชานนฺตสฺส อาปตฺติ ทุกฺกฏสฺส, ปุพฺเพวสฺส โหติ “มุสา ภณิสฺสนฺ”ติ, ภณนฺตสฺส โหติ “มุสา ภณามี”ติ, ภณิตสฺส โหติ “มุสา มยา ภณิตนฺ”ติ, วินิธาย ทิฏฺฐิํ, วินิธาย ขนฺติํ, วินิธาย รุจิํ, วินิธาย ภาวํ.}

\prose{ตีหากาเรหิ ฯเปฯ. สตฺตหากาเรหิ โย เต วิหาเร วสิ, โส ภิกฺขุ สุญฺญาคาเร ปฐมํ ฌานํ. ทุติยํ ฌานํ. ตติยํ ฌานํ. จตุตฺถํ ฌานํ \csromanfolio{137}สมาปชฺชิ. สมาปชฺชติ. สมาปนฺโน. โส ภิกฺขุ สุญฺญาคาเร จตุตฺถสฺส ฌานสฺส ลาภี. วสี. เตน ภิกฺขุนา สุญฺญาคาเร จตุตฺถํ ฌานํ สจฺฉิกตนฺติ สมฺปชานมุสา ภณนฺตสฺส ปฏิวิชานนฺตสฺส อาปตฺติ ถุลฺลจฺจยสฺส, น ปฏิวิชานนฺตสฺส อาปตฺติ ทุกฺกฏสฺส, ปุพฺเพวสฺส โหติ “มุสา ภณิสฺสนฺ”ติ, ภณนฺตสฺส โหติ “มุสา ภณามี”ติ, ภณิตสฺส โหติ “มุสา มยา ภณิตนฺ”ติ, วินิธาย ทิฏฺฐิํ, วินิธาย ขนฺติํ, วินิธาย รุจิํ, วินิธาย ภาวํ.}

\prose{ยถา อิทํ วิตฺถาริตํ, เอวเมว เสสานิปิ วิตฺถาเร\-ตพฺพานิ.}

\proseitem{221}{ตีหากาเรหิ ฯเปฯ. สตฺตหากาเรหิ โย เต จีวรํ ปริภุญฺชิ. โย เต ปิณฺฑปาตํ ปริภุญฺชิ. โย เต เสนาสนํ ปริภุญฺชิ. โย เต คิลานปฺปจฺจย\-เภสชฺช\-ปริกฺขารํ ปริภุญฺชิ, โส ภิกฺขุ สุญฺญาคาเร จตุตฺถํ ฌานํ สมาปชฺชิ. สมาปชฺชติ. สมาปนฺโน. โส ภิกฺขุ สุญฺญาคาเร จตุตฺถสฺส ฌานสฺส ลาภี. วสี. เตน ภิกฺขุนา สุญฺญาคาเร จตุตฺถํ ฌานํ สจฺฉิกตนฺติ สมฺปชานมุสา ภณนฺตสฺส ปฏิวิชานนฺตสฺส อาปตฺติ ถุลฺลจฺจยสฺส, น ปฏิวิชานนฺตสฺส อาปตฺติ ทุกฺกฏสฺส ฯเปฯ วินิธาย ภาวํ.}

\prose{ตีหากาเรหิ ฯเปฯ. สตฺตหากาเรหิ เยน เต วิหาโร ปริภุตฺโต. เยน เต จีวรํ ปริภุตฺตํ. เยน เต ปิณฺฑปาโต ปริภุตฺโต. เยน เต เสนาสนํ ปริภุตฺตํ. เยน เต คเลนปฺปจฺจยเภสชฺชปริกฺขาโร ปริภุตฺโต, โส ภิกฺขุ สุญฺญาคาเร จตุตฺถํ ฌานํ สมาปชฺชิ. สมาปชฺชติ. สมาปนฺโน. โส ภิกฺขุ สุญฺญาคาเร จตุตฺถสฺส ฌานสฺส ลาภี. วสี. เตน ภิกฺขุนา สุญฺญาคาเร จตุตฺถํ ฌานํ สจฺฉิกตนฺติ สมฺปชานมุสา ภณนฺตสฺส ปฏิวิชานนฺตสฺส อาปตฺติ ถุลฺลจฺจยสฺส, น ปฏิวิชานนฺตสฺส อาปตฺติ ทุกฺกฏสฺส ฯเปฯ วินิธาย ภาวํ.}

\prosewithcloser{ตีหากาเรหิ ฯเปฯ. สตฺตหากาเรหิ ยํ ตฺวํ อาคมฺม วิหารํ อทาสิ. จีวรํ อทาสิ. ปิณฺฑปาตํ อทาสิ. เสนาสนํ อทาสิ. คิลานปฺปจฺจย\-เภสชฺช\-ปริกฺขารํ อทาสิ, โส ภิกฺขุ สุญฺญาคาเร จตุตฺถํ ฌานํ สมาปชฺชิ. สมาปชฺชติ. สมาปนฺโน. โส ภิกฺขุ สุญฺญาคาเร จตุตฺถสฺส ฌานสฺส ลาภี. วสี. เตน ภิกฺขุนา สุญฺญาคาเร จตุตฺถํ ฌานํ สจฺฉิกตนฺติ สมฺปชานมุสา ภณนฺตสฺส ปฏิวิชานนฺตสฺส อาปตฺติ ถุลฺลจฺจยสฺส, น ปฏิวิชานนฺตสฺส อาปตฺติ ทุกฺกฏสฺส, ปุพฺเพวสฺส โหติ “มุสา \csromanfolio{138}ภณิสฺสนฺ”ติ, ภณนฺตสฺส โหติ “มุสา ภณามี”ติ, ภณิตสฺส โหติ “มุสา มยา ภณิตนฺ”ติ, วินิธาย ทิฏฺฐิํ, วินิธาย ขนฺติํ, วินิธาย รุจิํ, วินิธาย ภาวํ.}{เปยฺยาล\-ปนฺนรสกํ นิฏฺฐิตํ.}
\nitthitam{ปจฺจย\-ปฺปฏิสํยุตฺต\-วาร\-กถา นิฏฺฐิตา.}
\nitthitam{อุตฺตริมนุสฺสธมฺม\-จกฺกเปยฺยาลํ นิฏฺฐิตํ.}
\proseitem{222}{อนาปตฺติ อธิมาเนน อนุลฺลปนา\-ธิปฺปายสฺส อุมฺมตฺตกสฺส ขิตฺตจิตฺตสฺส เวทนาฏฺฏสฺส อาทิกมฺมิกสฺสาติ.}
\csromansectionrule

\csromanmatikaanchor{mtk.34}
\csromantocmarkref{subsection}{วินีตวตฺถุอุทฺทานคาถา}{mtk.34}
\bookmark[dest={mtk.34},level=2]{วินีตวตฺถุอุทฺทานคาถา}
\csromanheader{2}{วินีตวตฺถุ\-อุทฺทานคาถา}
\setlength{\csromangathapagewidth}{0pt}
\setlength{\csromangathawakwidth}{0pt}
\csromangathameasuregroup{อธิมาเน\textsuperscript{1} อรญฺญมฺหิ, \\ สญฺโญชนา รโหธมฺมา, \\ พฺราหฺมเณ ปญฺจ วตฺถูนิ, \\ อคาราวรณา กามา, \\ ตโปทา ราชคเห ยุทฺธํ, \\ โสภิโต อรหํ ภิกฺขุ,}{\csromangathabat{อธิมาเน\textsuperscript{1} อรญฺญมฺหิ,}{ปิณฺโฑปชฺฌาริยาปโถ.} \\ \csromangathabat{สญฺโญชนา รโหธมฺมา,}{วิหาโร ปจฺจุปฏฺฐิโต.} \\ น ทุกฺกรํ วิริยมโถปิ มจฺจุโน, \\ ภายาวุโส วิปฺปฏิสาริ สมฺมา. \\ วิริเยน โยเคน อาราธนาย, \\ อถ เวทนาย อธิวาสนา ทุเว. \\ \csromangathabat{พฺราหฺมเณ ปญฺจ วตฺถูนิ,}{อญฺญํ พฺยากรณา ตโย.} \\ \csromangathabat{อคาราวรณา กามา,}{รติ จาปิ อปกฺกมิ.} \\ อฏฺฐิ เปสิ อุโภ คาวฆาตกา, \\ ปิณฺโฑ สากุณิโก นิจฺฉวิ โอรพฺภิ. \\ อสิ จ สูกริโก สตฺติ มาควิ, \\ อุสุ จ การณิโก สูจิ สารถิ. \\ โย จ สิพฺพียติ สูจโก หิ โส, \\ อณฺฑภาริ อหุ คามกูฏโก. \\ กูเป นิมุคฺโค หิ โส ปารทาริโก, \\ คูถขาที อหุ ทุฏฺฐพฺราหฺมโณ. \\ นิจฺฉวิตฺถี อติจารินี อหุ, \\ มงฺคุลิตฺถี อหุ อิกฺขณิตฺถิกา. \\ โอกิลินี หิ สปตฺตงฺคาโรกิริ, \\ สีสจฺฉินฺโน อหุ โจรฆาตโก. \\ ภิกฺขุ ภิกฺขุนี สิกฺขมานา, \\ สามเณโร อถ สามเณริกา. \\ กสฺสปสฺส วินยสฺมิํ ปพฺพชํ, \\ ปาปกมฺมมกริํสุ ตาวเท. \\ \csromangathabat{ตโปทา ราชคเห ยุทฺธํ,}{นาคาโนคาหเนน จ.} \\ \csromangathabat{โสภิโต อรหํ ภิกฺขุ,}{ปญฺจกปฺปสตํ สเรติ.}}
\csromangathameasurewak{น ทุกฺกรํ วิริยมโถปิ มจฺจุโน, \\ ภายาวุโส วิปฺปฏิสาริ สมฺมา. \\ วิริเยน โยเคน อาราธนาย, \\ อถ เวทนาย อธิวาสนา ทุเว. \\ อฏฺฐิ เปสิ อุโภ คาวฆาตกา, \\ ปิณฺโฑ สากุณิโก นิจฺฉวิ โอรพฺภิ. \\ อสิ จ สูกริโก สตฺติ มาควิ, \\ อุสุ จ การณิโก สูจิ สารถิ. \\ โย จ สิพฺพียติ สูจโก หิ โส, \\ อณฺฑภาริ อหุ คามกูฏโก. \\ กูเป นิมุคฺโค หิ โส ปารทาริโก, \\ คูถขาที อหุ ทุฏฺฐพฺราหฺมโณ. \\ นิจฺฉวิตฺถี อติจารินี อหุ, \\ มงฺคุลิตฺถี อหุ อิกฺขณิตฺถิกา. \\ โอกิลินี หิ สปตฺตงฺคาโรกิริ, \\ สีสจฺฉินฺโน อหุ โจรฆาตโก. \\ ภิกฺขุ ภิกฺขุนี สิกฺขมานา, \\ สามเณโร อถ สามเณริกา. \\ กสฺสปสฺส วินยสฺมิํ ปพฺพชํ, \\ ปาปกมฺมมกริํสุ ตาวเท.}
\csromangathasetleft{อธิมาเน\textsuperscript{1} อรญฺญมฺหิ, \\ สญฺโญชนา รโหธมฺมา, \\ พฺราหฺมเณ ปญฺจ วตฺถูนิ, \\ อคาราวรณา กามา, \\ ตโปทา ราชคเห ยุทฺธํ, \\ โสภิโต อรหํ ภิกฺขุ,}
\csromangathagroup{\csromangathastanza{\csromangathabat{อธิมาเน\csromansharedfootnote{59708a5e7f1ef11d}{อธิมาเนน (อิ)} อรญฺญมฺหิ,}{ปิณฺโฑปชฺฌาริยาปโถ.} \\ \csromangathabat{สญฺโญชนา รโหธมฺมา,}{วิหาโร ปจฺจุปฏฺฐิโต.}}\csromangathastanza{\csromangathawak{น ทุกฺกรํ วิริยมโถปิ มจฺจุโน, \\ ภายาวุโส วิปฺปฏิสาริ สมฺมา. \\ วิริเยน โยเคน อาราธนาย, \\ อถ เวทนาย อธิวาสนา ทุเว.}}\csromangathastanza{\csromangathabat{พฺราหฺมเณ ปญฺจ วตฺถูนิ,}{อญฺญํ พฺยากรณา ตโย.} \\ \csromangathabat{อคาราวรณา กามา,}{รติ จาปิ อปกฺกมิ.}}\csromangathastanza{\csromangathawak{อฏฺฐิ เปสิ อุโภ คาวฆาตกา, \\ ปิณฺโฑ สากุณิโก นิจฺฉวิ โอรพฺภิ. \\ อสิ จ สูกริโก สตฺติ มาควิ, \\ อุสุ จ การณิโก สูจิ สารถิ.}}\csromangathastanza{\csromangathawak{โย จ สิพฺพียติ สูจโก หิ โส, \\ อณฺฑภาริ อหุ คามกูฏโก. \\ กูเป นิมุคฺโค หิ โส ปารทาริโก, \\ คูถขาที อหุ ทุฏฺฐพฺราหฺมโณ.}}\csromangathastanza{\csromangathawak{\csromanfolio{139}นิจฺฉวิตฺถี อติจารินี อหุ, \\ มงฺคุลิตฺถี อหุ อิกฺขณิตฺถิกา. \\ โอกิลินี หิ สปตฺตงฺคาโรกิริ, \\ สีสจฺฉินฺโน อหุ โจรฆาตโก.}}\csromangathastanza{\csromangathawak{ภิกฺขุ ภิกฺขุนี สิกฺขมานา, \\ สามเณโร อถ สามเณริกา. \\ กสฺสปสฺส วินยสฺมิํ ปพฺพชํ, \\ ปาปกมฺมมกริํสุ ตาวเท.}}\csromangathastanza{\csromangathabat{ตโปทา ราชคเห ยุทฺธํ,}{นาคาโนคาหเนน จ.} \\ \csromangathabat{โสภิโต อรหํ ภิกฺขุ,}{ปญฺจกปฺปสตํ สเรติ.}}}

\csromanmatikaanchor{mtk.35}
\csromantocmarkref{subsection}{วินีตวตฺถุ}{mtk.35}
\bookmark[dest={mtk.35},level=2]{วินีตวตฺถุ}
\csromanheader{2}{\textbf{วินีตวตฺถุ}}
\proseitem{223}{เตน โข ปน สมเยน อญฺญตโร ภิกฺขุ อธิมาเนน อญฺญํ พฺยากาสิ. ตสฺส กุกฺกุจฺจํ อโหสิ “ภควตา สิกฺขาปทํ ปญฺญตฺตํ, กจฺจิ นุ โข อหํ ปาราชิกํ อาปตฺติํ อาปนฺโน”ติ. ภควโต เอตมตฺถํ อาโรเจสิ. อนาปตฺติ ภิกฺขุ อธิมาเนนาติ. (1)}

\prose{เตน โข ปน สมเยน อญฺญตโร ภิกฺขุ ปณิธาย อรญฺเญ วิหรติ “เอวํ มํ ชโน สมฺภาเวสฺสตี”ติ. ตํ ชโน สมฺภาเวสิ. ตสฺส กุกฺกุจฺจํ อโหสิ ฯเปฯ. อนาปตฺติ ภิกฺขุ ปาราชิกสฺส, น จ ภิกฺขเว ปณิธาย อรญฺเญ วตฺถพฺพํ, โย วเสยฺย, อาปตฺติ ทุกฺกฏสฺสาติ. (2)}

\prose{เตน โข ปน สมเยน อญฺญตโร ภิกฺขุ ปณิธาย ปิณฺฑาย จรติ “เอวํ มํ ชโน สมฺภาเวสฺสตี”ติ. ตํ ชโน สมฺภาเวสิ. ตสฺส กุกฺกุจฺจํ อโหสิ ฯเปฯ. อนาปตฺติ ภิกฺขุ ปาราชิกสฺส, น จ ภิกฺขเว ปณิธาย ปิณฺฑาย จริตพฺพํ, โย จเรยฺย, อาปตฺติ ทุกฺกฏสฺสาติ. (3)}

\prose{เตน โข ปน สมเยน อญฺญตโร ภิกฺขุ อญฺญตรํ ภิกฺขุํ เอตทโวจ “เย อาวุโส อมฺหากํ อุปชฺฌายสฺส สทฺธิวิหาริกา สพฺเพว อรหตฺโต”ติ. ตสฺส กุกฺกุจฺจํ อโหสิ ฯเปฯ. กิํจิตฺโต ตฺวํ ภิกฺขูติ. \csromanfolio{140}อุลฺลปนา\-ธิปฺปาโย อหํ ภควาติ. อนาปตฺติ ภิกฺขุ ปาราชิกสฺส, อาปตฺติ ถุลฺลจฺจยสฺสาติ. (4)}

\prose{เตน โข ปน สมเยน อญฺญตโร ภิกฺขุ อญฺญตรํ ภิกฺขุํ เอตทโวจ “เย อาวุโส อมฺหากํ อุปชฺฌายสฺส อนฺเตวาสิกา สพฺเพว มหิทฺธิกา มหานุภาวา”ติ. ตสฺส กุกฺกุจฺจํ อโหสิ ฯเปฯ. กิํจิตฺโต ตฺวํ ภิกฺขูติ. อุลฺลปนา\-ธิปฺปาโย อหํ ภควาติ. อนาปตฺติ ภิกฺขุ ปาราชิกสฺส, อาปตฺติ ถุลฺลจฺจยสฺสาติ. (5)}

\prose{เตน โข ปน สมเยน อญฺญตโร ภิกฺขุ ปณิธาย จงฺกมติ. ปณิธาย ติฏฺฐติ. ปณิธาย นิสีทติ. ปณิธาย เสยฺยํ กปฺเปติ “เอวํ มํ ชโน สมฺภาเวสฺสตี”ติ. ตํ ชโน สมฺภาเวสิ. ตสฺส กุกฺกุจฺจํ อโหสิ ฯเปฯ. อนาปตฺติ ภิกฺขุ ปาราชิกสฺส, น จ ภิกฺขเว ปณิธาย เสยฺยา กปฺเปตพฺพา, โย กปฺเปยฺย, อาปตฺติ ทุกฺกฏสฺสาติ. \mbox{(6-9)}}

\prose{เตน โข ปน สมเยน อญฺญตโร ภิกฺขุ อญฺญตรสฺส ภิกฺขุโน อุตฺตริ\-มนุสฺสธมฺมํ อุลฺลปติ. โสปิ เอวมาห “มยฺหํปิ อาวุโส สญฺโญชนา ปหีนา”ติ. ตสฺส กุกฺกุจฺจํ อโหสิ ฯเปฯ. อาปตฺติํ ตฺวํ ภิกฺขุ อาปนฺโน ปาราชิกนฺติ. (10)}

\proseitem{224}{เตน โข ปน สมเยน อญฺญตโร ภิกฺขุ รโหคโต อุตฺตริ\-มนุสฺสธมฺมํ อุลฺลปติ. ปรจิตฺตวิทู ภิกฺขุ ตํ ภิกฺขุํ อปสาเทสิ “มา อาวุโส เอวรูปํ อภณิ, นตฺเถโส ตุยฺหนฺ”ติ. ตสฺส กุกฺกุจฺจํ อโหสิ ฯเปฯ. อนาปตฺติ ภิกฺขุ ปาราชิกสฺส, อาปตฺติ ทุกฺกฏสฺสาติ. (11)}

\prose{เตน โข ปน สมเยน อญฺญตโร ภิกฺขุ รโหคโต อุตฺตริ\-มนุสฺสธมฺมํ อุลฺลปติ. เทวตา ตํ ภิกฺขุํ อปสาเทสิ “มา ภนฺเต เอวรูปํ อภณิ, นตฺเถโส ตุยฺหนฺ”ติ. ตสฺส กุกฺกุจฺจํ อโหสิ ฯเปฯ. อนาปตฺติ ภิกฺขุ ปาราชิกสฺส, อาปตฺติ ทุกฺกฏสฺสาติ. (12)}

\prose{เตน โข ปน สมเยน อญฺญตโร ภิกฺขุ อญฺญตรํ อุปาสกํ เอตทโวจ “โย อาวุโส ตุยฺหํ วิหาเร วสติ, โส ภิกฺขุ อรหา”ติ. โส จ ตสฺส วิหาเร วสติ, ตสฺส กุกฺกุจฺจํ อโหสิ ฯเปฯ. กิํจิตฺโต ตฺวํ ภิกฺขูติ. อุลฺลปนา\-ธิปฺปาโย อหํ ภควาติ. อนาปตฺติ ภิกฺขุ ปาราชิกสฺส, อาปตฺติ ถุลฺลจฺจยสฺสาติ. (13)}

\prose{\csromanfolio{141}เตน โข ปน สมเยน อญฺญตโร ภิกฺขุ อญฺญตรํ อุปาสกํ เอตทโวจ “ยํ ตฺวํ อาวุโส อุปฏฺเฐสิ จีวร\-ปิณฺฑปาต\-เสนาสน\-คิลานปฺปจฺจย\-เภสชฺช\-ปริกฺขาเรน, โส ภิกฺขุ อรหา”ติ. โส จ ตํ อุปฏฺเฐติ จีวร\-ปิณฺฑปาต\-เสนาสน\-คิลานปฺปจฺจย\-เภสชฺช\-ปริกฺขาเตน, ตสฺส กุกฺกุจฺจํ อโหสิ ฯเปฯ. กิํจิตฺโต ตฺวํ ภิกฺขูติ. อุลฺลปนา\-ธิปฺปาโย อหํ ภควาติ. อนาปตฺติ ภิกฺขุ ปาราชิกสฺส, อาปตฺติ ถุลฺลจฺจยสฺสาติ. (14)}

\proseitem{225}{เตน โข ปน สมเยน อญฺญตโร ภิกฺขุ คิลาโน โหติ. ตํ ภิกฺขู เอตทโวจุํ “อตฺถายสฺมโต อุตฺตริ\-มนุสฺสธมฺโม”ติ. นาวุโส ทุกฺกรํ อญฺญํ พฺยากาตุนฺติ. ตสฺส กุกฺกุจฺจํ อโหสิ “เย โข เต ภควโต สาวกา, เต เอวํ วเทยฺยุํ, อหญฺจมฺหิ น ภควโต สาวโก, กจฺจิ นุ โข อหํ ปาราชิกํ อาปตฺติํ อาปนฺโน”ติ. ภควโต เอตมตฺถํ อาโรเจสิ. กิํจิตฺโต ตฺวํ ภิกฺขูติ. อนุลฺลปนา\-ธิปฺปาโย อหํ ภควาติ. อนาปตฺติ ภิกฺขุ อนุลฺลปนาธิปฺปายสฺสาติ. (15)}

\prose{เตน โข ปน สมเยน อญฺญตโร ภิกฺขุ คิลาโน โหติ. ตํ ภิกฺขู เอตทโวจุํ “อตฺถายสฺมโต อุตฺตริ\-มนุสฺสธมฺโม”ติ. อาราธนีโย โข อาวุโส ธมฺโม อารทฺธวีริเยนาติ. ตสฺส กุกฺกุจฺจํ อโหสิ ฯเปฯ. อนาปตฺติ ภิกฺขุ อนุลฺลปนาธิปฺปายสฺสาติ. (16)}

\prose{เตน โข ปน สมเยน อญฺญตโร ภิกฺขุ คิลาโน โหติ. ตํ ภิกฺขู เอตทโวจุํ “มา โข อาวุโส ภายี”ติ. นาหํ อาวุโส มจฺจุโน ภายามีติ. ตสฺส กุกฺกุจฺจํ อโหสิ ฯเปฯ. อนาปตฺติ ภิกฺขุ อนุลฺลปนาธิปฺปายสฺสาติ. (17)}

\prose{เตน โข ปน สมเยน อญฺญตโร ภิกฺขุ คิลาโน โหติ. ตํ ภิกฺขู เอตทโวจุํ “มา โข อาวุโส ภายี”ติ. โย นูนาวุโส วิปฺปฏิสารี อสฺส, โส ภาเยยฺยาติ. ตสฺส กุกฺกุจฺจํ อโหสิ ฯเปฯ. อนาปตฺติ ภิกฺขุ อนุลฺลปนาธิปฺปายสฺสาติ. (18)}

\prose{เตน โข ปน สมเยน อญฺญตโร ภิกฺขุ คิลาโน โหติ. ตํ ภิกฺขู เอตทโวจุํ “อตฺถายสฺมโต อุตฺตริ\-มนุสฺสธมฺโม”ติ. อาราธนีโย โข อาวุโส ธมฺโม สมฺมาปยุตฺเตนาติ. ตสฺส กุกฺกุจฺจํ อโหสิ ฯเปฯ. อนาปตฺติ ภิกฺขุ อนุลฺลปนาธิปฺปายสฺสาติ. (19)}

\prose{\csromanfolio{142}เตน โข ปน สมเยน อญฺญตโร ภิกฺขุ คิลาโน โหติ. ตํ ภิกฺขู เอตทโวจุํ “อตฺถายสฺมโต อุตฺตริ\-มนุสฺสธมฺโม”ติ. อาราธนีโย โข อาวุโส ธมฺโม อารทฺธวีริเยนาติ. ตสฺส กุกฺกุจฺจํ อโหสิ ฯเปฯ. อนาปตฺติ ภิกฺขุ อนุลฺลปนาธิปฺปายสฺสาติ. (20)}

\prose{เตน โข ปน สมเยน อญฺญตโร ภิกฺขุ คิลาโน โหติ. ตํ ภิกฺขู เอตทโวจุํ “อตฺถายสฺมโต อุตฺตริ\-มนุสฺสธมฺโม”ติ. อาราธนีโย โข อาวุโส ธมฺโม ยุตฺตโยเคนาติ. ตสฺส กุกฺกุจฺจํ อโหสิ ฯเปฯ. อนาปตฺติ ภิกฺขุ อนุลฺลปนาธิปฺปายสฺสาติ. (21)}

\prose{เตน โข ปน สมเยน อญฺญตโร ภิกฺขุ คิลาโน โหติ. ตํ ภิกฺขู เอตทโวจุํ “กจฺจาวุโส ขมนียํ, กจฺจิ ยาปนียนฺ”ติ. นาวุโส สกฺกา เยน วา เตน วา อธิวาเสตุนฺติ. ตสฺส กุกฺกุจฺจํ อโหสิ ฯเปฯ. อนาปตฺติ ภิกฺขุ อนุลฺลปนาธิปฺปายสฺสาติ. (22)}

\prose{เตน โข ปน สมเยน อญฺญตโร ภิกฺขุ คิลาโน โหติ. ตํ ภิกฺขู เอตทโวจุํ “กจฺจาวุโส ขมนียํ, กจฺจิ ยาปนียนฺ”ติ. นาวุโส สกฺกา ปุถุชฺชเนน อธิวาเสตุนฺติ. ตสฺส กุกฺกุจฺจํ อโหสิ ฯเปฯ. กิํจิตฺโต ตฺวํ ภิกฺขูติ. อุลฺลปนา\-ธิปฺปาโย อหํ ภควาติ. อนาปตฺติ ภิกฺขุ ปาราชิกสฺส, อาปตฺติ ถุลฺลจฺจยสฺสาติ. (23)}

\proseitem{226}{เตน โข ปน สมเยน อญฺญตโร พฺราหฺมโณ ภิกฺขู นิมนฺเตตฺวา เอตทโวจ “อายนฺตุ โภนฺโต อรหตฺโต”ติ. เตสํ กุกฺกุจฺจํ อโหสิ “มยญฺจมฺห น อรหตฺโต\csromansharedfootnote{f9bdf7221e9e795f}{อนรหตฺโต (สี)}, อยญฺจ พฺราหฺมโณ อเมฺห อรหนฺตวาเทน สมุทาจรติ, กถํ นุ โข อเมฺหหิ ปฏิปชฺชิตพฺพนฺ”ติ. ภควโต เอตมตฺตํ อาโรเจสุํ. อนาปตฺติ ภิกฺขเว ปสาทภญฺเญติ. (24)}

\prose{เตน โข ปน สมเยน อญฺญตโร พฺราหฺมโณ ภิกฺขู นิมนฺเตตฺวา เอตทโวจ “นิสีทนฺตุ โภนฺโต อรหตฺโต”ติ. “ภุญฺชนฺตุ โภนฺโต อรหตฺโต”ติ. “ตปฺเปนฺตุ โภนฺโต อรหตฺโต”ติ. “คจฺฉนฺตุ โภนฺโต อรหตฺโต”ติ. เตสํ กุกฺกุจฺจํ อโหสิ “มยญฺจมฺห น \csromanfolio{143}อรหตฺโต, อยญฺจ พฺราหฺมโณ อเมฺห อรหนฺตวาเทน สมุทาจรติ, กถํ นุ โข อเมฺหหิ ปฏิปชฺชิตพฺพนฺ”ติ. ภควโต เอตมตฺตํ อาโรเจสุํ. อนาปตฺติ ภิกฺขเว ปสาทภญฺเญติ. \mbox{(25-28)}}

\prose{เตน โข ปน สมเยน อญฺญตโร ภิกฺขุ อญฺญตรสฺส ภิกฺขุโน อุตฺตริ\-มนุสฺสธมฺมํ อุลฺลปติ. โสปิ เอวมาห “มยฺหํปิ อาวุโส อาสวา ปหีนา”ติ. ตสฺส กุกฺกุจฺจํ อโหสิ ฯเปฯ. อาปตฺติํ ตฺวํ ภิกฺขุ อาปนฺโน ปาราชิกนฺติ. (29)}

\prose{เตน โข ปน สมเยน อญฺญตโร ภิกฺขุ อญฺญตรสฺส ภิกฺขุโน อุตฺตริ\-มนุสฺสธมฺมํ อุลฺลปติ. โสปิ เอวมาห “มยฺหํปิ อาวุโส เอเต ธมฺมา สํวิชฺชนฺตี”ติ. ตสฺส กุกฺกุจฺจํ อโหสิ ฯเปฯ. อาปตฺติํ ตฺวํ ภิกฺขุ อาปนฺโน ปาราชิกนฺติ. (30)}

\prose{เตน โข ปน สมเยน อญฺญตโร ภิกฺขุ อญฺญตรสฺส ภิกฺขุโน อุตฺตริ\-มนุสฺสธมฺมํ อุลฺลปติ. โสปิ เอวมาห “อหมฺปาวุโส เตสุ ธมฺเมสุ สนฺทิสฺสามี”ติ. ตสฺส กุกฺกุจฺจํ อโหสิ ฯเปฯ. อาปตฺติํ ตฺวํ ภิกฺขุ อาปนฺโน ปาราชิกนฺติ. (31)}

\prose{เตน โข ปน สมเยน อญฺญตรํ ภิกฺขุํ ญาตกา เอตทโวจุํ “เอหิ ภนฺเต อคารํ อชฺฌาวสา”ติ. อภพฺโพ โข อาวุโส มาทิโส อคารํ อชฺฌาวสิตุนฺติ. ตสฺส กุกฺกุจฺจํ อโหสิ ฯเปฯ. อนาปตฺติ ภิกฺขุ อนุลฺลปนาธิปฺปายสฺสาติ. (32)}

\proseitem{227}{เตน โข ปน สมเยน อญฺญตรํ ภิกฺขุํ ญาตกา เอตทโวจุํ “เอหิ ภนฺเต กาเม ปริภุญฺชา”ติ. อาวฏา เม อาวุโส กามาติ. ตสฺส กุกฺกุจฺจํ อโหสิ ฯเปฯ. อนาปตฺติ ภิกฺขุ อนุลฺลปนาธิปฺปายสฺสาติ. (33)}

\prose{เตน โข ปน สมเยน อญฺญตรํ ภิกฺขุํ ญาตกา เอตทโวจุํ “อภิรมสิ ภนฺเต”ติ. อภิรโต อหํ อาวุโส ปรมาย อภิรติยาติ. ตสฺส กุกฺกุจฺจํ อโหสิ “เย โข เต ภควโต สาวกา, เต เอวํ วเทยฺยุํ, อหญฺจมฺหิ น ภควโต สาวโก, กจฺจิ นุ โข อหํ ปาราชิกํ อาปตฺติํ อาปนฺโน”ติ. ภควโต เอตมตฺถํ อาโรเจสิ. กิํจิตฺโต ตฺวํ ภิกฺขูติ. อนุลฺลปนา\-ธิปฺปาโย อหํ ภควาติ. อนาปตฺติ ภิกฺขุ อนุลฺลปนาธิปฺปายสฺสาติ. (34)}

\prose{\csromanfolio{144}เตน โข ปน สมเยน สมฺพหุลา ภิกฺขู กติกํ กตฺวา อญฺญตรสฺมิํ อาวาเส วสฺสํ อุปคจฺฉิํสุ “โย อิมมฺหา อาวาสา ปฐมํ ปกฺกมิสฺสติ, ตํ มยํ อรหาติ ชานิสฺสามา”ติ. อญฺญตโร ภิกฺขุ “มํ อรหาติ ชานนฺตู”ติ ตมฺหา อาวาสา ปฐมํ ปกฺกามิ. ตสฺส กุกฺกุจฺจํ อโหสิ ฯเปฯ. อาปตฺติํ ตฺวํ ภิกฺขุ อาปนฺโน ปาราชิกนฺติ. (35)}

\proseitem{228}{\csromansymbolfootnote{*}{อิมานิ วตฺถูนิ สํ 1. 446 ปิฏฺฐาทีสุปิ อาคตานิ.}เตน สมเยน พุทฺโธ ภควา ราชคเห วิหรติ เวฬุวเน กลนฺทก\-นิวาเป. เตน โข ปน สมเยน อายสฺมา จ ลกฺขโณ อายสฺมา จ มหาโมคฺคลฺลาโน คิชฺฌกูเฏ ปพฺพเต วิหรนฺติ. อถ โข อายสฺมา มหาโมคฺคลฺลาโน ปุพฺพณฺหสมยํ นิวาเสตฺวา ปตฺตจีวรํ อาทาย เยนายสฺมา ลกฺขโณ เตนุปสงฺกมิ, อุปสงฺกมิตฺวา อายสฺมนฺตํ ลกฺขณํ เอตทโวจ “อายามาวุโส ลกฺขณ ราชคหํ ปิณฺฑาย ปวิสิสฺสามา”ติ. “เอวมาวุโส”ติ โข อายสฺมา ลกฺขโณ อายสฺมโต มหา\-โมคฺคลฺลานสฺส ปจฺจสฺโสสิ. อถ โข อายสฺมา มหาโมคฺคลฺลาโน คิชฺฌกูฏา ปพฺพตา โอโรหนฺโต อญฺญตรสฺมิํ ปเทเส สิตํ ปาตฺวากาสิ. อถ โข อายสฺมา ลกฺขโณ อายสฺมนฺตํ มหา\-โมคฺคลฺลานํ เอตทโวจ “โก นุ โข อาวุโส โมคฺคลฺลาน เหตุ, โก ปจฺจโย สิตสฺส ปาตุกมฺมายา”ติ. อกาโล โข อาวุโส ลกฺขณ เอตสฺส ปญฺหสฺส\csromansharedfootnote{7f7be9600f0a92e9}{ปญฺหสฺส พฺยากรณาย (สฺยา)}, ภควโต มํ สนฺติเก เอตํ ปญฺหํ ปุจฺฉาติ. อถ โข อายสฺมา จ ลกฺขโณ อายสฺมา จ มหาโมคฺคลฺลาโน ราชคเห ปิณฺฑาย จริตฺวา ปจฺฉาภตฺตํ ปิณฺฑปาต\-ปฏิกฺกนฺตา เยน ภควา เตนุ\-ปสงฺกมิํสุ, อุปสงฺกมิตฺวา ภควนฺตํ อภิวาเทตฺวา เอกมนฺตํ นิสีทิํสุ, เอกมนฺตํ นิสินฺโน โข อายสฺมา ลกฺขโณ อายสฺมนฺตํ มหา\-โมคฺคลฺลานํ เอตทโวจ “อิธายสฺมา มหาโมคฺคลฺลาโน คิชฺฌกูฏา ปพฺพตา โอโรหนฺโต อญฺญตรสฺมิํ ปเทเส สิตํ ปาตฺวากาสิ, โก นุ โข อาวุโส โมคฺคลฺลาน เหตุ, โก ปจฺจโย สิตสฺส ปาตุกมฺมายา”ติ. อิธาหํ อาวุโส คิชฺฌกูฏา ปพฺพตา โอโรหนฺโต อทฺทสํ อฏฺฐิก\-สงฺขลิกํ เวหาสํ คจฺฉนฺตํ, ตเมนํ คิชฺฌาปิ กากาปิ กุลลาปิ อนุปติตฺวา อนุปติตฺวา ผาสุฬ\-นฺตริกาหิ วิตุเฑนฺติ\csromansharedfootnote{410bda430326baef}{วิตุเทนฺติ วิตจฺเฉนฺติ วิราเชนฺติ (สฺยา)}, สา สุทํ อฏฺฏสฺสรํ กโรติ, ตสฺส มยฺหํ อาวุโส เอตทโหสิ “อจฺฉริยํ วต โภ อพฺภุตํ วต โภ, เอวรูโปปิ นาม สตฺโต ภวิสฺสติ, เอวรูโปปิ \csromanfolio{145}นาม ยกฺโข ภวิสฺสติ, เอวรูโปปิ นาม อตฺตภาว\-ปฺปฏิลาโภ ภวิสฺสตี”ติ. ภิกฺขู อุชฺฌายนฺติ ขิยฺยนฺติ วิปาเจนฺติ “อุตฺตริ\-มนุสฺสธมฺมํ อายสฺมา มหาโมคฺคลฺลาโน อุลฺลปตี”ติ. อถ โข ภควา ภิกฺขู อามนฺเตสิ “จกฺขุภูตา วต ภิกฺขเว สาวกา วิหรนฺติ, ญาณภูตา วต ภิกฺขเว สาวกา วิหรนฺติ, ยตฺร หิ นาม สาวโก เอวรูปํ ญสฺสติ วา ทกฺขติ วา สกฺขิํ วา กริสฺสติ, ปุพฺเพว เม โส ภิกฺขเว สตฺโต ทิฏฺโฐ อโหสิ, อปิ จาหํ น พฺยากาสิํ, อหญฺเจตํ พฺยากเรยฺยํ, ปเร จ เม น สทฺทเหยฺยุํ, เย เม น สทฺทเหยฺยุํ, เตสํ ตํ อสฺส ทีฆรตฺตํ อหิตาย ทุกฺขาย. เอโส ภิกฺขเว สตฺโต อิมสฺมิํ เยว ราชคเห โคฆาตโก อโหสิ, โส ตสฺส กมฺมสฺส วิปาเกน พหูนิ วสฺสานิ พหูนิ วสฺสสตานิ พหูนิ วสฺสสหสฺสานิ พหูนิ วสฺส\-สตสหสฺสานิ นิรเย ปจฺจิตฺวา ตสฺเสว กมฺมสฺส วิปากาวเสเสน เอวรูปํ อตฺตภาว\-ปฺปฏิลาภํ ปฏิสํเวเทติ. สจฺจํ ภิกฺขเว โมคฺคลฺลาโน อาห, อนาปตฺติ ภิกฺขเว โมคลฺลานสฺสา”ติ. (36)}

\proseitem{229}{อิธาหํ อาวุโส คิชฺฌกูฏา ปพฺพตา โอโรหนฺโต อทฺทสํ มํสเปสิํ เวหาสํ คจฺฉนฺติํ, ตเมนํ คิชฺฌาปิ กากาปิ กุลลาปิ อนุปติตฺวา อนุปติตฺวา วิตจฺเฉนฺติ วิภชฺเชนฺติ\csromansharedfootnote{410bda430326baef}{วิตุเทนฺติ วิตจฺเฉนฺติ วิราเชนฺติ (สฺยา)}. สา สุทํ อฏฺฏสฺสรํ กโรติ ฯเปฯ. เอโส ภิกฺขเว สตฺโต อิมสฺมิํ, เยว ราชคเห โคฆาตโก อโหสิ ฯเปฯ. (37)}

\prose{อิธาหํ อาวุโส คิชฺฌกูฏา ปพฺพตา โอโรหนฺโต อทฺทสํ มํสปิณฺฑํ เวหาสํ คจฺฉนฺติํ. ตเมนํ คิชฺฌาปิ กากาปิ กุลลาปิ อนุปติตฺวา อนุปติตฺวา วิตจฺเฉนฺติ วิภชฺเชนฺติ, โส สุทํ อฏฺฏสฺสรํ กโรติ ฯเปฯ. เอโส ภิกฺขเว สตฺโต อิมสฺมิํ, เยว ราชคเห สากุณิโก อโหสิ ฯเปฯ. (38)}

\prose{อิธาหํ อาวุโส คิชฺฌกูฏา ปพฺพตา โอโรหนฺโต อทฺทสํ นิจฺฉวิํ ปุริสํ เวหาสํ คจฺฉนฺติํ. ตเมนํ คิชฺฌาปิ กากาปิ กุลลาปิ อนุปติตฺวา อนุปติตฺวา วิตจฺเฉนฺติ วิภชฺเชนฺติ. โส สุทํ อฏฺฏสฺสรํ กโรติ}

\prose{\csromanfolio{146}อิธาหํ อาวุโส คิชฺฌกูฏา ปพฺพตา โอโรหนฺโต อทฺทสํ อสิโลมํ ปุริสํ เวหาสํ คจฺฉนฺตํ. ตสฺส เต อสี อุปฺปติตฺวา อุปฺปติตฺวา ตสฺเสว กาเย นิปตนฺติ. โส สุทํ อฏฺฏสฺสรํ กโรติ ฯเปฯ. เอโส ภิกฺขเว สตฺโต อิมสฺมิํ เยว ราชคเห สูกริโก อโหสิ ฯเปฯ. (40)}

\prose{อิธาหํ อาวุโส คิชฺฌกูฏา ปพฺพตา โอโรหนฺโต อทฺทสํ สตฺติโลมํ ปุริสํ เวหาสํ คจฺฉนฺตํ. ตสฺส ตา สตฺติโย อุปฺปติตฺวา อุปฺปติตฺวา ตสฺเสว กาเย นิปตนฺติ. โส สุทํ อฏฺฏสฺสรํ กโรติ ฯเปฯ. เอโส ภิกฺขเว สตฺโต อิมสฺมิํ เยว ราชคเห มาควิโก อโหสิ ฯเปฯ. (41)}

\prose{อิธาหํ อาวุโส คิชฺฌกูฏา ปพฺพตา โอโรหนฺโต อทฺทสํ อุสุโลมํ ปุริสํ เวหาสํ คจฺฉนฺตํ. ตสฺส เต อุสู อุปฺปติตฺวา อุปฺปติตฺวา ตสฺเสว กาเย นิปตนฺติ. โส สุทํ อฏฺฏสฺสรํ กโรติ ฯเปฯ. เอโส ภิกฺขเว สตฺโต อิมสฺมิํ เยว ราชคเห การณิโก อโหสิ ฯเปฯ. (42)}

\prose{อิธาหํ อาวุโส คิชฺฌกูฏา ปพฺพตา โอโรหนฺโต อทฺทสํ สูจิโลมํ ปุริสํ เวหาสํ คจฺฉนฺตํ. ตสฺส ตา สูจิโย อุปฺปติตฺวา อุปฺปติตฺวา ตสฺเสว กาเย นิปตนฺติ. โส สุทํ อฏฺฏสฺสรํ กโรติ ฯเปฯ. เอโส ภิกฺขเว สตฺโต อิมสฺมิํ เยว ราชคเห สารถิโก อโหสิ ฯเปฯ. (43)}

\prose{อิธาหํ อาวุโส คิชฺฌกูฏา ปพฺพตา โอโรหนฺโต อทฺทสํ สูจิโลมํ ปุริสํ เวหาสํ คจฺฉนฺตํ. ตสฺส ตา สูจิโย สีเส ปวิสิตฺวา มุขโต นิกฺขมนฺติ, มุเข ปวิสิตฺวา อุรโต นิกฺขมนฺติ, อุเร ปวิสิตฺวา อุทรโต นิกฺขมนฺติ, อุทเร ปวิสิตฺวา อูรูหิ นิกฺขมนฺติ, อูรูสุ ปวิสิตฺวา ชงฺฆาหิ นิกฺขมนฺติ, ชงฺฆาสุ ปวิสิตฺวา ปาเทหิ นิกฺขมนฺติ. โส สุทํ อฏฺฏสฺสรํ กโรติ ฯเปฯ. เอโส ภิกฺขเว สตฺโต อิมสฺมิํ เยว ราชคเห สูจโก อโหสิ ฯเปฯ. (44)}

\prose{อิธาหํ อาวุโส คิชฺฌกูฏา ปพฺพตา โอโรหนฺโต อทฺทสํ กุมฺภณฺฑํ ปุริสํ เวหาสํ คจฺฉนฺตํ. โส คจฺฉนฺโตปิ เตว อณฺเฑ ขนฺเธ อาโรเปตฺวา คจฺฉติ. นิสีทนฺโตปิ เตเสฺวว อณฺเฑสุ นิสีทติ. ตเมนํ คิชฺฌาปิ \csromanfolio{147}กากาปิ กุลลาปิ อนุปติตฺวา อนุปติตฺวา วิตจฺเฉนฺติ วิภชฺเชนฺติ. โส สุทํ อฏฺฏสฺสรํ กโรติ ฯเปฯ. เอโส ภิกฺขเว สตฺโต อิมสฺมิํ เยว ราชคเห คามกูโฏ อโหสิ ฯเปฯ. (45)}

\prose{อิธาหํ อาวุโส คิชฺฌกูฏา ปพฺพตา โอโรหนฺโต อทฺทสํ ปุริสํ คูถกูเป สสีสกํ นิมุคฺคํ ฯเปฯ. เอโส ภิกฺขเว สตฺโต อิมสฺมิํ เยว ราชคเห ปารทาริโก อโหสิ ฯเปฯ. (46)}

\prose{อิธาหํ อาวุโส คิชฺฌกูฏา ปพฺพตา โอโรหนฺโต อทฺทสํ ปุริสํ คูถกูเป สสีสกํ นิมุคฺคํ อุโภหิ หตฺเถหิ คูถํ ขาทนฺตํ ฯเปฯ. เอโส ภิกฺขเว สตฺโต อิมสฺมิํ เยว ราชคเห ทุฏฺฐพฺราหฺมโณ อโหสิ. โส กสฺสปสฺส สมฺมา\-สมฺพุทฺธสฺส ปาวจเน ภิกฺขุสํฆํ ภตฺเตน นิมนฺเตตฺวา โทณิโย\csromansharedfootnote{56e83cd8445d301b}{โทณิยา (อิติปิ)} คูถสฺส ปูราเปตฺวา กาลํ อาโรจาเปตฺวา เอตทโวจ “อโต\csromansharedfootnote{7a27c9c34c85dad1}{อิโต (สฺยา) อโห (อิติปิ)} โภนฺโต ยาวทตฺถํ ภุญฺชนฺตุ เจว หรนฺตุ จา”ติ ฯเปฯ. (47)}

\proseitem{230}{อิธาหํ อาวุโส คิชฺฌกูฏา ปพฺพตา โอโรหนฺโต อทฺทสํ นิจฺฉวิํ อิตฺถิํ เวหาสํ คจฺฉนฺติํ. ตเมนํ คิชฺฌาปิ กากาปิ กุลลาปิ อนุปติตฺวา อนุปติตฺวา วิตจฺเฉนฺติ วิภชฺเชนฺติ. สา สุทํ อฏฺฏสฺสรํ กโรติ ฯเปฯ. เอสา ภิกฺขเว อิตฺถี อิมสฺมิํ เยว ราชคเห อติจารินี อโหสิ ฯเปฯ. (48)}

\prose{อิธาหํ อาวุโส คิชฺฌกูฏา ปพฺพตา โอโรหนฺโต อทฺทสํ อิตฺถิํ ทุคฺคนฺธํ มงฺคุลิํ เวหาสํ คจฺฉนฺติํ. ตเมนํ คิชฺฌาปิ กากาปิ กุลลาปิ อนุปติตฺวา อนุปติตฺวา วิตจฺเฉนฺติ วิภชฺเชนฺติ. สา สุทํ อฏฺฏสฺสรํ กโรติ ฯเปฯ. เอสา ภิกฺขเว อิตฺถี อิมสฺมิํ เยว ราชคเห อิกฺขณิกา อโหสิ ฯเปฯ. (49)}

\prose{อิธาหํ อาวุโส คิชฺฌกูฏา ปพฺพตา โอโรหนฺโต อทฺทสํ อิตฺถิํ อุปฺปกฺกํ โอกิลินิํ โอกิรินิํ เวหาสํ คจฺฉนฺติํ. สา สุทํ อฏฺฏสฺสรํ กโรติ ฯเปฯ. เอสา ภิกฺขเว อิตฺถี กาลิงฺคสฺส รญฺโญ อคฺคมเหสี อโหสิ. สา อิสฺสาปกตา สปตฺติํ องฺคารกฏาเหน โอกิริ ฯเปฯ. (50)}

\prose{อิธาหํ อาวุโส คิชฺฌกูฏา ปพฺพตา โอโรหนฺโต อทฺทสํ อสีสกํ กพนฺธํ เวหาสํ คจฺฉนฺตํ. ตสฺส อุเร อกฺขีนิ เจว โหนฺติ มุขํ จ. ตเมนํ คิชฺฌาปิ กากาปิ กุลลาปิ อนุปติตฺวา อนุปติตฺวา วิตจฺเฉนฺติ วิภชฺเชนฺติ. \csromanfolio{148}โส สุทํ อฏฺฏสฺสรํ กโรติ ฯเปฯ. เอโส ภิกฺขเว สตฺโต อิมสฺมิํ เยว ราชคเห หาริโก นาม โจรฆาตโก อโหสิ ฯเปฯ. (51)}

\prose{อิธาหํ อาวุโส คิชฺฌกูฏา ปพฺพตา โอโรหนฺโต อทฺทสํ ภิกฺขุํ เวหาสํ คจฺฉนฺติํ. ตสฺส สํฆาฏิปิ อาทิตฺตา สมฺปชฺชลิตา สโชติภูตา, ปตฺโตปิ อาทิตฺโต สมฺปชฺชลิโต สโชติภูโต, กาย\-พนฺธนํปิ อาทิตฺตํ สมฺปชฺชลิตํ สโชติภูตํ, กาโยปิ อาทิตฺโต สมฺปชฺชลิโต สโชติภูโต. โส สุทํ อฏฺฏสฺสรํ กโรติ ฯเปฯ. เอโส ภิกฺขเว ภิกฺขุ กสฺสปสฺส สมฺมา\-สมฺพุทฺธสฺส ปาวจเน ปาปภิกฺขุ อโหสิ ฯเปฯ. (52)}

\prose{อิธาหํ อาวุโส คิชฺฌกูฏา ปพฺพตา โอโรหนฺโต อทฺทสํ ภิกฺขุนิํ. อทฺทสํ สิกฺขมานํ. อทฺทสํ สามเณรํ. อทฺทสํ สามเณริํ เวหาสํ คจฺฉนฺติํ. ตสฺสา สํฆาฏิปิ อาทิตฺตา สมฺปชฺชลิตา สโชติภูตา, ปตฺโตปิ อาทิตฺโต สมฺปชฺชลิโต สโชติภูโต, กาย\-พนฺธนํปิ อาทิตฺตํ สมฺปชฺชลิตํ สโชติภูตํ, กาโยปิ อาทิตฺโต สมฺปชฺชลิโต สโชติภูโต. สา สุทํ อฏฺฏสฺสรํ กโรติ. ตสฺส มยฺหํ อาวุโส เอตทโหสิ “อจฺฉริยํ วต โภ อพฺภุตํ วต โภ, เอวรูโปปิ นาม สตฺโต ภวิสฺสติ, เอวรูโปปิ นาม ยกฺโข ภวิสฺสติ, เอวรูโปปิ นาม อตฺตภาว\-ปฺปฏิลาโภ ภวิสฺสตี”ติ. ภิกฺขู อุชฺฌายนฺติ ขิยฺยนฺติ วิปาเจนฺติ “อุตฺตริ\-มนุสฺสธมฺมํ อายสฺมา มหาโมคฺคลฺลาโน อุลฺลปตี”ติ.}

\prose{อถ โข ภควา ภิกฺขู อามนฺเตสิ “จกฺขุภูตา วต ภิกฺขเว สาวกา วิหรนฺติ, ญาณภูตา วต ภิกฺขเว สาวกา วิหรนฺติ, ยตฺร หิ นาม สาวโก เอวรูปํ ญสฺสติ วา ทกฺขติ วา สกฺขิํ วา กริสฺสติ, ปุพฺเพว เม สา ภิกฺขเว สามเณรี ทิฏฺฐา อโหสิ, อปิ จาหํ น พฺยากาสิํ, อหญฺเจตํ พฺยากเรยฺยํ, ปเร จ เม น สทฺทเหยฺยุํ, เย เม น สทฺทเหยฺยุํ, เตสํ ตํ อสฺส ทีฆรตฺตํ อหิตาย ทุกฺขาย. เอสา ภิกฺขเว สามเณรี กสฺสปสฺส สมฺมา\-สมฺพุทฺธสฺส ปาวจเน ปาปสามเณรี อโหสิ, สา ตสฺส กมฺมสฺส วิปาเกน พหูนิ วสฺสานิ พหูนิ วสฺสสตานิ พหูนิ วสฺสสหสฺสานิ พหูนิ วสฺส\-สตสหสฺสานิ นิรเย ปจฺจิตฺวา ตสฺเสว กมฺมสฺส วิปากาวเสเสน เอวรูปํ อตฺตภาว\-ปฺปฏิลาภํ ปฏิสํเวเทติ. สจฺจํ ภิกฺขเว โมคฺคลฺลาโน อาห, อนาปตฺติ ภิกฺขเว โมคฺคลฺลานสฺสา”ติ. \mbox{(53-56)}}

\proseitem{231}{\csromanfolio{149}อถ โข อายสฺมา มหาโมคฺคลฺลาโน ภิกฺขู อามนฺเตสิ “ยตายํ อาวุโส ตโปทา สนฺทติ, โส ทโห อจฺโฉทโก สีโตทโก สาโตทโก เสตโก สุปฺปติตฺโถ รมณีโย ปหูต\-มจฺฉ\-กจฺฉโป, จกฺกมตฺตานิ จ ปทุมานิ ปุปฺผนฺตี”ติ. ภิกฺขู อุชฺฌายนฺติ ขิยฺยนฺติ วิปาเจนฺติ “กถํ หิ นาม อายสฺมา มหาโมคฺคลฺลาโน เอวํ วกฺขติ ‘ยตายํ อาวุโส ตโปทา สนฺทติ, โส ทโห อจฺโฉทโก สีโตทโก สาโตทโก เสตโก สุปฺปติตฺโถ รมณิโย ปหูต\-มจฺฉ\-กจฺฉโป, จกฺกมตฺตานิ จ ปทุมานิ ปุปฺผนฺตี’ติ, อถ จ ปนาหํ ตโปทา กุถิตา สนฺทติ, อุตฺตริ\-มนุสฺสธมฺมํ อายสฺมา มหาโมคฺคลฺลาโน อุลฺลปตี”ติ. ภควโต เอตมตฺถํ อาโรเจสุํ. ยตายํ ภิกฺขเว ตโปทา สนฺทติ, โส ทโห อจฺโฉทโก สีโตทโก สาโตทโก เสตโก สุปฺปติตฺโถ รมณีโย ปหูต\-มจฺฉ\-กจฺฉโป, จกฺกมตฺตานิ จ ปทุมานิ ปุปฺผนฺติ. อปิ จายํ ภิกฺขเว ตโปทา ทฺวินฺนํ มหานิรยานํ อนฺตริกาย อาคจฺฉติ, เตนายํ ตโปทา กุถิตา สนฺทติ. สจฺจํ ภิกฺขเว โมคฺคลฺลาโน อาห, อนาปตฺติ ภิกฺขเว โมคฺคลฺลานสฺสาติ. (57)}

\prose{เตน โข ปน สมเยน ราชา มาคโธ เสนิโย พิมฺพิสาโร ลิจฺฉวีหิ สทฺธิํ สงฺคาเมนฺโต ปภคฺโค อโหสิ. อถ ราชา ปจฺฉา เสนํ สงฺกฑฺฒิตฺวา ลิจฺฉวโย ปราเชสิ. สงฺคาเม จ นนฺที จรติ “รญฺญา ลิจฺฉวี ปภคฺคา”ติ. อถ โข อายสฺมา มหาโมคฺคลฺลาโน ภิกฺขู อามนฺเตสิ “ราชา อาวุโส ลิจฺฉวีหิ ปภคฺโค”ติ. ภิกฺขู อุชฺฌายนฺติ ขิยฺยนฺติ วิปาเจนฺติ “กถํ หิ นาม อายสฺมา มหาโมคฺคลฺลาโน เอวํ วกฺขติ ‘ราชา อาวุโส ลิจฺฉวีหิ ปภคฺโค’ติ, สงฺคาเม จ นนฺทิํ จรติ ‘รญฺญา ลิจฺฉวี ปภคฺคา’ติ, อุตฺตริ\-มนุสฺสธมฺมํ อายสฺมา มหาโมคฺคลฺลาโน อุลฺลปตี”ติ. ภควโต เอตมตฺถํ อาโรเจสุํ. ปฐมํ ภิกฺขเว ราชา ลิจฺฉวีหิ ปภคฺโค, อถ ราชา ปจฺฉา เสนํ สงฺกฑฺฒิตฺวา ลิจฺฉวโย ปราเชสิ. สจฺจํ ภิกฺขเว โมคฺคลฺลาโน อาห, อนาปตฺติ ภิกฺขเว โมคฺคลฺลานสฺสาติ. (58)}

\proseitem{232}{อถ โข อายสฺมา มหาโมคฺคลฺลาโน ภิกฺขู อามนฺเตสิ “อิธาหํ อาวุโส สปฺปินิกาย นทิยา ตีเร อาเนญฺชํ สมาธิํ สมาปนฺโน นาคานํ โอคยฺห อุตฺตรนฺตานํ โกญฺจํ กโรนฺตานํ สทฺทํ อสฺโสสินฺ”ติ. \csromanfolio{150}ภิกฺขู อุชฺฌายนฺติ ขิยฺยนฺติ วิปาเจนฺติ “กถํ หิ นาม อายสฺมา มหาโมคฺคลฺลาโน อาเนญฺชํ สมาธิํ สมาปนฺโน สทฺทํ โสสฺสติ, อุตฺตริ\-มนุสฺสธมฺมํ อายสฺมา มหาโมคฺคลฺลาโน อุลฺลปตี”ติ. ภควโต เอตมตฺถํ อาโรเจสุํ. อตฺเถโส ภิกฺขเว สมาธิ, โส จ โข อปริสุทฺโธ. สจฺจํ ภิกฺขเว โมคฺคลฺลาโน อาห, อนาปตฺติ ภิกฺขเว โมคฺคลฺลานสฺสาติ. (59)}

\prosewithcloser{อถ โข อายสฺมา โสภิโต ภิกฺขู อามนฺเตสิ “อหํ อาวุโส ปญฺจ กปฺปสตานิ อนุสฺสรามี”ติ. ภิกฺขู อุชฺฌายนฺติ ขิยฺยนฺติ วิปาเจนฺติ “กถํ หิ นาม อายสฺมา โสภิโต เอวํ วกฺขติ ‘อหํ อาวุโส ปญฺจ กปฺปสตานิ อนุสฺสรามี’ติ, อุตฺตริ\-มนุสฺสธมฺมํ อายสฺมา โสภิโต อุลฺลปตี”ติ. ภควโต เอตมตฺถํ อาโรเจสุํ. อตฺเถสา ภิกฺขเว โสภิตสฺส, สา จ โข เอกา เยว ชาติ. สจฺจํ ภิกฺขเว โสภิโต อาห, อนาปตฺติ ภิกฺขเว โสภิตสฺสาติ. (60)}{จตุตฺถ\-ปาราชิกํ สมตฺตํ.}
\proseitemwithcloser{233}{อุทฺทิฏฺฐา โข อายสฺมนฺโต จตฺตาโร ปาราชิกา ธมฺมา, เยสํ ภิกฺขุ อญฺญตรํ วา อญฺญตรํ วา อาปชฺชิตฺวา น ลภติ ภิกฺขูหิ สทฺธิํ สํวาสํ ยถา ปุเร ตถา ปจฺฉา, ปาราชิโก โหติ อสํวาโส. ตตฺถายสฺมนฺเต ปุจฺฉามิ กจฺจิตฺถ ปริสุทฺธา, ทุติยมฺปิ ปุจฺฉามิ กจฺจิตฺถ ปริสุทฺธา, ตติยมฺปิ ปุจฺฉามิ กจฺจิตฺถ ปริสุทฺธา, ปริสุทฺเธ\-ตฺถา\-ยสฺมนฺโต, ตสฺมา ตุณฺหี, เอวเมตํ ธารยามีติ.}{ปาราชิกํ นิฏฺฐิตํ.}
\csromancenter{ตสฺสุทฺทานํ}
\setlength{\csromangathapagewidth}{0pt}
\setlength{\csromangathawakwidth}{0pt}
\csromangathameasuregroup{เมถุนาทินฺนาทานญฺจ, \\ ปาราชิกานิ จตฺตาริ,}{\csromangathabat{เมถุนาทินฺนาทานญฺจ,}{มนุสฺสวิคฺคหุตฺตริ.} \\ \csromangathabat{ปาราชิกานิ จตฺตาริ,}{เฉชฺชวตฺถู อสํสยาติ.}}
\csromangathasetleft{เมถุนาทินฺนาทานญฺจ, \\ ปาราชิกานิ จตฺตาริ,}
\csromangathagroupwithclosermajor{\csromangathastanza{\csromangathabat{เมถุนาทินฺนาทานญฺจ,}{มนุสฺสวิคฺคหุตฺตริ.} \\ \csromangathabat{ปาราชิกานิ จตฺตาริ,}{เฉชฺชวตฺถู อสํสยาติ.}}}{\textbf{ปาราชิกกณฺฑํ นิฏฺฐิตํ.}}
\csromanmatikaanchor{mtk.36}
\csromantocmarknopage{chapter}{2. สํฆาทิเสสกณฺฑ}{mtk.36}
\bookmark[dest={mtk.36},level=0]{2. สํฆาทิเสสกณฺฑ}
\chapterheadpageread{2. \textbf{สํฆาทิเสส}\-\textbf{กณฺฑ}}
\csromanmatikaanchor{mtk.37}
\csromantocmarknopage{section}{1. สุกฺกวิสฺสฏฺฐิสิกฺขาปท}{mtk.37}
\bookmark[dest={mtk.37},level=1]{1. สุกฺกวิสฺสฏฺฐิสิกฺขาปท}
\csromanheader{1}{1. \textbf{สุกฺกวิสฺสฏฺฐิ}\-\textbf{สิกฺขาปท}}
\vspace{\tipitakaparskip}
\csromancenter{อิเม โข ปนายสฺมนฺโต เตรส สํฆาทิเสสา \\ ธมฺมา อุทฺเทสํ อาคจฺฉนฺติ.}
\csromanmatikaanchor{mtk.38}
\csromantocmarkref{subsection}{เสยฺยสกภิกฺขุวตฺถุ}{mtk.38}
\bookmark[dest={mtk.38},level=2]{เสยฺยสกภิกฺขุวตฺถุ}
\proseitem{234}{\csromanfolio{151}เตน สมเยน พุทฺโธ ภควา สาวตฺถิยํ วิหรติ เชตวเน อนาถ\-ปิณฺฑิกสฺส อาราเม. เตน โข ปน สมเยน อายสฺมา เสยฺยสโก อนภิรโต พฺรหฺมจริยํ จรติ, โส เตน กิโส โหติ ลูโข ทุพฺพณฺโณ อุปฺปณฺฑุ\-ปฺปณฺฑุก\-ชาโต ธมนิสนฺถต\-คตฺโต. อทฺทส โข อายสฺมา อุทายี อายสฺมนฺตํ เสยฺยสกํ กิสํ ลูขํ ทุพฺพณฺณํ อุปฺปณฺฑุ\-ปฺปณฺฑุก\-ชาตํ ธมนิสนฺถต\-คตฺตํ, ทิสฺวาน อายสฺมนฺตํ เสยฺยสกํ เอตทโวจ “กิสฺส ตฺวํ อาวุโส เสยฺยสก กิโส ลูโข ทุพฺพณฺโณ อุปฺปณฺฑุ\-ปฺปณฺฑุก\-ชาโต ธมนิสนฺถต\-คตฺโต, กจฺจิ โน ตฺวํ อาวุโส เสยฺยสก อนภิรโต พฺรหฺมจริยํ จรสี”ติ, เอวมาวุโสติ. “เตน หิ ตฺวํ อาวุโส เสยฺยสก ยาวทตฺถํ ภุญฺช ยาวทตฺถํ สุป ยาวทตฺถํ นฺหาย, ยาวทตฺถํ ภุญฺชิตฺวา ยาวทตฺถํ สุปิตฺวา ยาวทตฺถํ นฺหายิตฺวา ยทา เต อนภิรติ อุปฺปชฺชติ, ราโค จิตฺตํ อนุทฺธํเสติ, ตทา หตฺเถน อุปกฺกมิตฺวา อสุจิํ โมเจหี”ติ. กิํ นุ โข อาวุโส กปฺปติ เอวรูปํ กาตุนฺติ. อาม อาวุโส อหมฺปิ เอวํ กโรมีติ.}

\prose{อถ โข อายสฺมา เสยฺยสโก ยาวทตฺถํ ภุญฺชิ ยาวทตฺถํ สุปิ ยาวทตฺถํ นฺหายิ, ยาวทตฺถํ ภุญฺชิตฺวา ยาวทตฺถํ สุปิตฺวา ยาวทตฺถํ นฺหายิตฺวา ยทา อนภิรติ อุปฺปชฺชติ, ราโค จิตฺตํ อนุทฺธํเสติ, ตทา หตฺเถน อุปกฺกมิตฺวา อสุจิํ โมเจสิ. อถ โข อายสฺมา เสยฺยสโก อปเรน สมเยน วณฺณวา อโหสิ ปีณินฺทฺริโย ปสนฺน\-มุขวณฺโณ วิปฺปสนฺน\-ฉวิวณฺโณ. อถ โข อายสฺมโต เสยฺยสกสฺส สหายกา ภิกฺขู อายสฺมนฺตํ เสยฺยสกํ เอตทโวจุํ “ปุพฺเพ โข ตฺวํ อาวุโส เสยฺยสก กิโส อโหสิ ลูโข ทุพฺพณฺโณ อุปฺปณฺฑุ\-ปฺปณฺฑุก\-ชาโต ธมนิสนฺถต\-คตฺโต, โส ทานิ ตฺวํ เอตรหิ วณฺณวา ปีณินฺทฺริโย ปสนฺน\-มุขวณฺโณ วิปฺปสนฺน\-ฉวิวณฺโณ, กิํ นุ โข ตฺวํ อาวุโส เสยฺยสก เภสชฺชํ กโรสี”ติ. น โข อหํ อาวุโส เภสชฺชํ \csromanfolio{152}กโรมิ, อปิจาหํ ยาวทตฺถํ ภุญฺชามิ ยาวทตฺถํ สุปามิ ยาวทตฺถํ นฺหายามิ, ยาวทตฺถํ ภุญฺชิตฺวา ยาวทตฺถํ สุปิตฺวา ยาวทตฺถํ นฺหายิตฺวา ยทา เม อนภิรติ อุปฺปชฺชติ, ราโค จิตฺตํ อนุทฺธํเสติ, ตทา หตฺเถน อุปกฺกมิตฺวา อสุจิํ โมเจมีติ. กิํ ปน ตฺวํ อาวุโส เสยฺยสก เยเนว หตฺเถน สทฺธาเทยฺยํ ภุญฺชสิ, เตเนว หตฺเถน อุปกฺกมิตฺวา อสุจิํ โมเจสีติ. เอวมาวุโสติ. เย เต ภิกฺขู อปฺปิจฺฉา ฯเปฯ เต อุชฺฌายนฺติ ขิยฺยนฺติ วิปาเจนฺติ “กถํ หิ นาม อายสฺมา เสยฺยสโก หตฺเถน อุปกฺกมิตฺวา อสุจิํ โมเจสฺสตี”ติ.}

\csromanmatikaanchor{mtk.39}
\csromanmatikaanchor{mtk.40}
\csromantocmarkref{subsection}{ปฐมปญฺญตฺติ, อนุปญฺญตฺติ}{mtk.39}
\bookmark[dest={mtk.39},level=2]{ปฐมปญฺญตฺติ, อนุปญฺญตฺติ}
\csromantocmarkref{subsection}{สิกฺขาปทวิภงฺค, ปทภาชนีย}{mtk.40}
\bookmark[dest={mtk.40},level=2]{สิกฺขาปทวิภงฺค, ปทภาชนีย}
\prose{อถ โข เต ภิกฺขู อายสฺมนฺตํ เสยฺยสกํ อเนก\-ปริยาเยน วิครหิตฺวา ภควโต เอตมตฺถํ อาโรเจสุํ. อถ โข ภควา เอตสฺมิํ นิทาเน เอตสฺมิํ ปกรเณ ภิกฺขุสํฆํ สนฺนิปาตาเปตฺวา อายสฺมนฺตํ เสยฺยสกํ ปฏิปุจฺฉิ “สจฺจํ กิร ตฺวํ เสยฺยสก หตฺเถน อุปกฺกมิตฺวา อสุจิํ โมเจสี”ติ. สจฺจํ ภควาติ. วิครหิ พุทฺโธ ภควา “อนนุจฺฉวิกํ โมฆปุริส อนนุโลมิกํ อปฺปติรูปํ อสฺสามณกํ อกปฺปิยํ อกรณียํ, กถํ หิ นาม ตฺวํ โมฆปุริส หตฺเถน อุปกฺกมิตฺวา อสุจิํ โมเจสฺสสิ. นนุ มยา โมฆปุริส อเนก\-ปริยาเยน วิราคาย ธมฺโม เทสิโต, โน สราคาย, วิสญฺโญคาย ธมฺโม เทสิโต, โน สญฺโญคาย, อนุปาทานาย ธมฺโม เทสิโต, โน สอุปาทานาย. ตตฺถ นาม ตฺวํ โมฆปุริส มยา วิราคาย ธมฺเม เทสิเต สราคาย เจเตสฺสสิ, วิสญฺโญคาย ธมฺเม เทสิเต สญฺโญคาย เจเตสฺสสิ, อนุปาทานาย ธมฺเม เทสิเต สอุปาทานาย เจเตสฺสสิ. นนุ มยา โมฆปุริส อเนก\-ปริยาเยน ราควิราคาย ธมฺโม เทสิโต, มทนิมฺมทนาย ปิปาสวินยาย อาลย\-สมุคฺฆาตาย วฏฺฏุ\-ปจฺเฉทาย ตณฺหากฺขยาย วิราคาย นิโรธาย นิพฺพานาย ธมฺโม เทสิโต, นนุ มยา โมฆปุริส อเนก\-ปริยาเยน กามานํ ปหานํ อกฺขาตํ, กามสญฺญานํ ปริญฺญา อกฺขาตา, กามปิปาสานํ ปฏิวินโย อกฺขาโต, กามวิตกฺกานํ สมุคฺฆาโต อกฺขาโต, กาม\-ปริฬาหานํ วูปสโม อกฺขาโต. เนตํ โมฆปุริส อปฺปสนฺนานํ วา ปสาทาย ปสนฺนานํ วา ภิยฺโยภาวาย, อถ เขฺวตํ โมฆปุริส อปฺปสนฺนาน\-ญฺเจว อปฺปสาทาย ปสนฺนานญฺจ เอกจฺจานํ อญฺญถตฺตายา”ติ. อถ โข ภควา อายสฺมนฺตํ เสยฺยสกํ \csromanfolio{153}อเนก\-ปริยาเยน วิครหิตฺวา ทุพฺภรตาย ฯเปฯ. เอวญฺจ ปน ภิกฺขเว อิมํ สิกฺขาปทํ อุทฺทิเสยฺยาถ\csromandash{}}

\prose{\textbf{“สญฺเจตนิกา สุกฺกวิสฺสฏฺฐิ}\csromansharedfootnote{60ed41d8d8586701}{วิสฏฺฐิ (สี, สฺยา)} \textbf{สํฆาทิเสโส”}ติ.}

\prose{เอวญฺจิทํ ภควตา ภิกฺขูนํ สิกฺขาปทํ ปญฺญตฺตํ โหติ.}

\proseitem{235}{เตน โข ปน สมเยน ภิกฺขู ปณีต\-โภชนานิ ภุญฺชิตฺวา มุฏฺฐสฺสตี อสมฺปชานา นิทฺทํ โอกฺกมนฺติ, เตสํ มุฏฺฐสฺสตีนํ อสมฺปชานานํ นิทฺทํ โอกฺกมนฺตานํ สุปินนฺเตน อสุจิ มุจฺจติ. เตสํ กุกฺกุจฺจํ อโหสิ “ภควตา สิกฺขาปทํ ปญฺญตฺตํ ‘สญฺเจตนิกา สุกฺกวิสฺสฏฺฐิ สํฆาทิเสโส’ติ, อมฺหากญฺจ สุปินนฺเตน อสุจิ มุจฺจติ, อตฺถิ เจตฺถ เจตนา ลพฺภติ, กจฺจิ นุ โข มยํ สํฆาทิเสสํ อาปตฺติํ อาปนฺนา”ติ. ภควโต เอตมตฺถํ อาโรเจสุํ. อตฺเถสา ภิกฺขเว เจตนา, สา จ โข อพฺโพหาริกาติ. เอวญฺจ ปน ภิกฺขเว อิมํ สิกฺขาปทํ อุทฺทิเสยฺยาถ\csromandash{}}

\proseitem{236}{\textbf{“สญฺเจตนิกา สุกฺกวิสฺสฏฺฐิ อญฺญตฺร สุปินนฺตา สํฆาทิเสโส”}ติ.}

\proseitem{237}{\textbf{สญฺเจตนิกา}ติ ชานนฺโต สญฺชานนฺโต เจจฺจ อภิวิตริตฺวา วีตกฺกโม.}

\prose{\textbf{สุกฺกนฺ}ติ ทส สุกฺกานิ นีลํ ปีตกํ โลหิตกํ โอทาตํ ตกฺกวณฺณํ ทกวณฺณํ เตลวณฺณํ ขีรวณฺณํ ทธิวณฺณํ สปฺปิวณฺณํ.}

\prose{\textbf{วิสฺสฏฺฐี}ติ ฐานโต จาวนา วุจฺจติ วิสฺสฏฺฐีติ.}

\prose{\textbf{อญฺญตฺร สุปินนฺตา}ติ ฐเปตฺวา สุปินนฺตํ.}

\prose{\textbf{สํฆาทิเสโส}ติ สํโฆว ตสฺสา อาปตฺติยา ปริวาสํ เทติ มูลาย ปฏิกสฺสติ มานตฺตํ เทติ อพฺเภติ, น สมฺพหุลา น เอกปุคฺคโล, เตน วุจฺจติ “สํฆาทิเสโส”ติ. ตสฺเสว อาปตฺติ\-นิกายสฺส นามกมฺมํ อธิวจนํ, เตนปิ วุจฺจติ “สํฆาทิเสโส”ติ.}

\prose{อชฺฌตฺตรูเป โมเจติ, พหิทฺธารูเป โมเจติ, อชฺฌตฺต\-พหิทฺธา\-รูเป โมเจติ, อากาเส กฏิํ กมฺเปนฺโต โมเจติ, ราคูปตฺถมฺเภ โมเจติ, วจฺจู\-ปตฺถมฺเภ โมเจติ, ปสฺสาวู\-ปตฺถมฺเภ โมเจติ, วาตูปตฺถมฺเภ โมเจติ, \csromanfolio{154}อุจฺจาลิงฺคปาณกทฏฺฐู\-ปตฺถมฺเภ โมเจติ, อาโรคฺยตฺถาย โมเจติ, สุขตฺถาย โมเจติ, เภสชฺชตฺถาย โมเจติ, ทานตฺถาย โมเจติ, ปุญฺญตฺถาย โมเจติ, ยญฺญตฺถาย โมเจติ, สคฺคตฺถาย โมเจติ, พีชตฺถาย โมเจติ, วีมํสตฺถาย โมเจติ, ทวตฺถาย โมเจติ. นีลํ โมเจติ, ปีตกํ โมเจติ, โลหิตกํ โมเจติ, โอทาตํ โมเจติ, ตกฺกวณฺณํ โมเจติ, ทกวณฺณํ โมเจติ, เตลวณฺณํ โมเจติ, ขีรวณฺณํ โมเจติ, ทธิวณฺณํ โมเจติ, สปฺปิวณฺณํ โมเจติ.}

\proseitem{238}{\textbf{อชฺฌตฺตรูเป}ติ อชฺฌตฺตํ อุปาทินฺเน\csromansharedfootnote{644519e87df886ef}{อุปาทิณฺเณ (สี, ก)} รูเป.}

\prose{\textbf{พหิทฺธารูเป}ติ พหิทฺธา อุปาทินฺเน วา อนุปาทินฺเน วา.}

\prose{\textbf{อชฺฌตฺต}\-\textbf{พหิทฺธา}\-\textbf{รูเป}ติ ตทุภเย.}

\prose{\textbf{อากาเส กฏิํ กมฺเปนฺโต}ติ อากาเส วายมนฺตสฺส องฺคชาตํ กมฺมนิยํ โหติ.}

\prose{\textbf{ราคูปตฺเถมฺเภ}ติ ราเคน ปีฬิตสฺส องฺคชาตํ กมฺมนิยํ โหติ.}

\prose{\textbf{วจฺจู}\-\textbf{ปตฺถมฺเภ}ติ วจฺเจน ปีฬิตสฺส องฺคชาตํ กมฺมนิยํ โหติ.}

\prose{\textbf{ปสฺสาวู}\-\textbf{ปตฺถมฺเภ}ติ ปสฺสาเวน ปีฬิตสฺส องฺคชาตํ กมฺมนิยํ โหติ.}

\prose{\textbf{วาตู}\-\textbf{ปตฺถมฺเภ}ติ วาเตน ปีฬิตสฺส องฺคชาตํ กมฺมนิยํ โหติ.}

\prose{\textbf{อุจฺจาลิงฺคปาณกทฏฺฐู}\-\textbf{ปตฺถมฺเภ}ติ อุจฺจาลิงฺคปาณก\-ทฏฺเฐน องฺคชาตํ กมฺมนิยํ โหติ.}

\proseitem{239}{\textbf{อาโรคฺยตฺถายา}ติ อโรโค ภวิสฺสามิ.}

\prose{\textbf{สุขตฺถายา}ติ สุขํ เวทนํ อุปฺปาเทสฺสามิ.}

\prose{\textbf{เภสชฺช}\-\textbf{ตฺถา}\-\textbf{ยา}ติ เภสชฺชํ ภวิสฺสติ.}

\prose{\textbf{ทานตฺถายา}ติ ทานํ ทสฺสามิ.}

\prose{\textbf{ปุญฺญตฺถายา}ติ ปุญฺญํ ภวิสฺสติ.}

\prose{\textbf{ยญฺญตฺถายา}ติ ยญฺญํ ยชิสฺสามิ.}

\prose{\textbf{สคฺคตฺถายา}ติ สคฺคํ คมิสฺสามิ.}

\prose{\textbf{พีชตฺถายา}ติ พีชํ ภวิสฺสติ.}

\prose{\csromanfolio{155}\textbf{วีมํสตฺถายา}ติ นีลํ ภวิสฺสติ, ปีตกํ ภวิสฺสติ, โลหิตกํ ภวิสฺสติ, โอทาตํ ภวิสฺสติ, ตกฺกวณฺณํ ภวิสฺสติ, ทกวณฺณํ ภวิสฺสติ, เตลวณฺณํ ภวิสฺสติ, ขีรวณฺณํ ภวิสฺสติ, ทธิวณฺณํ ภวิสฺสติ, สปฺปิวณฺณํ ภวิสฺสตีติ.}

\prose{\textbf{ทวตฺถายา}ติ ขิฑฺฑาธิปฺปาโย.}

\proseitem{240}{อชฺฌตฺตรูเป เจเตติ อุปกฺกมติ มุจฺจติ, อาปตฺติ สํฆาทิเสสสฺส.}

\prose{พหิทฺธารูเป เจเตติ อุปกฺกมติ มุจฺจติ, อาปตฺติ สํฆาทิเสสสฺส.}

\prose{อชฺฌตฺต\-พหิทฺธา\-รูเป เจเตติ อุปกฺกมติ มุจฺจติ, อาปตฺติ สํฆาทิเสสสฺส.}

\prose{อากาเส กฏิํ กมฺเปนฺโต เจเตติ อุปกฺกมติ มุจฺจติ, อาปตฺติ สํฆาทิเสสสฺส.}

\prose{ราคูปตฺถมฺเภ เจเตติ อุปกฺกมติ มุจฺจติ, อาปตฺติ สํฆาทิเสสสฺส.}

\prose{วจฺจู\-ปตฺถมฺเภ เจเตติ อุปกฺกมติ มุจฺจติ, อาปตฺติ สํฆาทิเสสสฺส.}

\prose{ปสฺสาวู\-ปตฺถมฺเภ เจเตติ อุปกฺกมติ มุจฺจติ, อาปตฺติ สํฆาทิเสสสฺส.}

\prose{วาตูปตฺถมฺเภ เจเตติ อุปกฺกมติ มุจฺจติ, อาปตฺติ สํฆาทิเสสสฺส.}

\prose{อุจฺจาลิงฺคปาณกทฏฺฐู\-ปตฺถมฺเภ เจเตติ อุปกฺกมติ มุจฺจติ, อาปตฺติ สํฆาทิเสสสฺส.}

\prose{อาโรคฺยตฺถาย เจเตติ อุปกฺกมติ มุจฺจติ, อาปตฺติ สํฆาทิเสสสฺส.}

\prose{สุขตฺถาย ฯเปฯ. เภสชฺชตฺถาย. ทานตฺถาย. ปุญฺญตฺถาย. ยญฺญตฺถาย. สคฺคตฺถาย. พีชตฺถาย. วีมํสตฺถาย. ทวตฺถาย เจเตติ อุปกฺกมติ มุจฺจติ, อาปตฺติ สํฆาทิเสสสฺส.}

\prose{นีลํ เจเตติ อุปกฺกมติ มุจฺจติ, อาปตฺติ สํฆาทิเสสสฺส.}

\prosewithcloser{ปีตกํ. โลหิตกํ. โอทาตํ. ตกฺกวณฺณํ. ทกวณฺณํ. เตลวณฺณํ. ขีรวณฺณํ. ทธิวณฺณํ. สปฺปิวณฺณํ เจเตติ อุปกฺกมติ มุจฺจติ, อาปตฺติ สํฆาทิเสสสฺส.}{สุทฺธิกํ นิฏฺฐิตํ.}
\prose{\csromanfolio{156}อาโรคฺยตฺถญฺจ สุขตฺถญฺจ เจเตติ อุปกฺกมติ มุจฺจติ, อาปตฺติ สํฆาทิเสสสฺส.}

\prosewithcloser{อาโรคฺยตฺถญฺจ เภสชฺช\-ตฺถญฺจ ฯเปฯ. อาโรคฺยตฺถญฺจ ทานตฺถญฺจ. อาโรคฺยตฺถญฺจ ปุญฺญตฺถญฺจ. อาโรคฺยตฺถญฺจ ยญฺญตฺถญฺจ. อาโรคฺยตฺถญฺจ สคฺคตฺถญฺจ. อาโรคฺยตฺถญฺจ พีชตฺถญฺจ. อาโรคฺยตฺถญฺจ วีมํสตฺถญฺจ. อาโรคฺยตฺถญฺจ ทวตฺถญฺจ เจเตติ อุปกฺกมติ มุจฺจติ, อาปตฺติ สํฆาทิเสสสฺส.}{เอกมูลกสฺส ขณฺฑจกฺกํ นิฏฺฐิตํ.}
\proseitem{241}{สุขตฺถญฺจ เภสชฺช\-ตฺถญฺจ เจเตติ อุปกฺกมติ มุจฺจติ, อาปตฺติ สํฆาทิเสสสฺส.}

\prose{สุขตฺถญฺจ ทานตฺถญฺจ ฯเปฯ. สุขตฺถญฺจ ปุญฺญตฺถญฺจ. สุขตฺถญฺจ ยญฺญตฺถญฺจ. สุขตฺถญฺจ สคฺคตฺถญฺจ. สุขตฺถญฺจ พีชตฺถญฺจ. สุขตฺถญฺจ วีมํสตฺถญฺจ. สุขตฺถญฺจ ทวตฺถญฺจ เจเตติ อุปกฺกมติ มุจฺจติ, อาปตฺติ สํฆาทิเสสสฺส.}

\prose{สุขตฺถญฺจ อาโรคฺยตฺถญฺจ เจเตติ อุปกฺกมติ มุจฺจติ, อาปตฺติ สํฆาทิเสสสฺส.}

\proseitem{242}{เภสชฺช\-ตฺถญฺจ ทานตฺถญฺจ ฯเปฯ. เภสชฺช\-ตฺถญฺจ ปุญฺญตฺถญฺจ. เภสชฺช\-ตฺถญฺจ ยญฺญตฺถญฺจ. เภสชฺช\-ตฺถญฺจ สคฺคตฺถญฺจ. เภสชฺช\-ตฺถญฺจ พีชตฺถญฺจ. เภสชฺช\-ตฺถญฺจ วีมํสตฺถญฺจ. เภสชฺช\-ตฺถญฺจ ทวตฺถญฺจ เจเตติ อุปกฺกมติ มุจฺจติ, อาปตฺติ สํฆาทิเสสสฺส.}

\prose{เภสชฺช\-ตฺถญฺจ อาโรคฺยตฺถญฺจ ฯเปฯ. เภสชฺช\-ตฺถญฺจ สุขตฺถญฺจ เจเตติ อุปกฺกมติ มุจฺจติ, อาปตฺติ สํฆาทิเสสสฺส.}

\prose{ทานตฺถญฺจ ปุญฺญตฺถญฺจ ฯเปฯ. ทานตฺถญฺจ ยญฺญตฺถญฺจ. ทานตฺถญฺจ สคฺคตฺถญฺจ. ทานตฺถญฺจ พีชตฺถญฺจ. ทานตฺถญฺจ วีมํสตฺถญฺจ. ทานตฺถญฺจ ทวตฺถญฺจ เจเตติ อุปกฺกมติ มุจฺจติ, อาปตฺติ สํฆาทิเสสสฺส.}

\prose{ทานตฺถญฺจ อาโรคฺยตฺถญฺจ ฯเปฯ. ทานตฺถญฺจ สุขตฺถญฺจ. ทานตฺถญฺจ เภสชฺช\-ตฺถญฺจ เจเตติ อุปกฺกมติ มุจฺจติ, อาปตฺติ สํฆาทิเสสสฺส.}

\prose{ปุญฺญตฺถญฺจ ยญฺญตฺถญฺจ ฯเปฯ. ปุญฺญตฺถญฺจ สคฺคตฺถญฺจ. ปุญฺญตฺถญฺจ พีชตฺถญฺจ. ปุญฺญตฺถญฺจ วีมํสตฺถญฺจ. ปุญฺญตฺถญฺจ ทวตฺถญฺจ เจเตติ อุปกฺกมติ มุจฺจติ, อาปตฺติ สํฆาทิเสสสฺส.}

\prose{\csromanfolio{157}ปุญฺญตฺถญฺจ อาโรคฺยตฺถญฺจ ฯเปฯ. ปุญฺญตฺถญฺจ สุขตฺถญฺจ. ปุญฺญตฺถญฺจ เภสชฺช\-ตฺถญฺจ. ปุญฺญตฺถญฺจ ทานตฺถญฺจ เจเตติ อุปกฺกมติ มุจฺจติ, อาปตฺติ สํฆาทิเสสสฺส.}

\prose{ยญฺญตฺถญฺจ สคฺคตฺถญฺจ ฯเปฯ. ยญฺญตฺถญฺจ พีชตฺถญฺจ. ยญฺญตฺถญฺจ วีมํสตฺถญฺจ. ยญฺญตฺถญฺจ ทวตฺถญฺจ เจเตติ อุปกฺกมติ มุจฺจติ, อาปตฺติ สํฆาทิเสสสฺส.}

\prose{ยญฺญตฺถญฺจ อาโรคฺยตฺถญฺจ ฯเปฯ. ยญฺญตฺถญฺจ สุขตฺถญฺจ. ยญฺญตฺถญฺจ เภสชฺช\-ตฺถญฺจ. ยญฺญตฺถญฺจ ทานตฺถญฺจ. ยญฺญตฺถญฺจ ปุญฺญตฺถญฺจ เจเตติ อุปกฺกมติ มุจฺจติ, อาปตฺติ สํฆาทิเสสสฺส.}

\prose{สคฺคตฺถญฺจ พีชตฺถญฺจ ฯเปฯ. สคฺคตฺถญฺจ วีมํสตฺถญฺจ. สคฺคตฺถญฺจ ทวตฺถญฺจ เจเตติ อุปกฺกมติ มุจฺจติ, อาปตฺติ สํฆาทิเสสสฺส.}

\prose{สคฺคตฺถญฺจ อาโรคฺยตฺถญฺจ ฯเปฯ. สคฺคตฺถญฺจ สุขตฺถญฺจ. สคฺคตฺถญฺจ เภสชฺช\-ตฺถญฺจ. สคฺคตฺถญฺจ ทานตฺถญฺจ. สคฺคตฺถญฺจ ปุญฺญตฺถญฺจ. สคฺคตฺถญฺจ ยญฺญตฺถญฺจ เจเตติ อุปกฺกมติ มุจฺจติ, อาปตฺติ สํฆาทิเสสสฺส.}

\prose{พีชตฺถญฺจ วีมํสตฺถญฺจ ฯเปฯ. พีชตฺถญฺจ ทวตฺถญฺจ เจเตติ อุปกฺกมติ มุจฺจติ, อาปตฺติ สํฆาทิเสสสฺส.}

\prose{พีชตฺถญฺจ อาโรคฺยตฺถญฺจ ฯเปฯ. พีชตฺถญฺจ สุขตฺถญฺจ. พีชตฺถญฺจ เภสชฺช\-ตฺถญฺจ. พีชตฺถญฺจ ทานตฺถญฺจ. พีชตฺถญฺจ ปุญฺญตฺถญฺจ. พีชตฺถญฺจ ยญฺญตฺถญฺจ. พีชตฺถญฺจ สคฺคตฺถญฺจ เจเตติ อุปกฺกมติ มุจฺจติ, อาปตฺติ สํฆาทิเสสสฺส.}

\prose{วีมํสตฺถญฺจ ทวตฺถญฺจ เจเตติ อุปกฺกมติ มุจฺจติ, อาปตฺติ สํฆาทิเสสสฺส.}

\prose{วีมํสตฺถญฺจ อาโรคฺยตฺถญฺจ ฯเปฯ. วีมํสตฺถญฺจ สุขตฺถญฺจ. วีมํสตฺถญฺจ เภสชฺช\-ตฺถญฺจ. วีมํสตฺถญฺจ ทานตฺถญฺจ. วีมํสตฺถญฺจ ปุญฺญตฺถญฺจ. วีมํสตฺถญฺจ ยญฺญตฺถญฺจ. วีมํสตฺถญฺจ สคฺคตฺถญฺจ. วีมํสตฺถญฺจ พีชตฺถญฺจ เจเตติ อุปกฺกมติ มุจฺจติ, อาปตฺติ สํฆาทิเสสสฺส.}

\prosewithcloser{ทวตฺถญฺจ อาโรคฺยตฺถญฺจ ฯเปฯ. ทวตฺถญฺจ สุขตฺถญฺจ. ทวตฺถญฺจ เภสชฺช\-ตฺถญฺจ. ทวตฺถญฺจ ทานตฺถญฺจ. ทวตฺถญฺจ ปุญฺญตฺถญฺจ. ทวตฺถญฺจ ยญฺญตฺถญฺจ. ทวตฺถญฺจ สคฺคตฺถญฺจ. ทวตฺถญฺจ พีชตฺถญฺจ. ทวตฺถญฺจ วีมํสตฺถญฺจ เจเตติ อุปกฺกมติ มุจฺจติ, อาปตฺติ สํฆาทิเสสสฺส.}{เอกมูลกสฺส พทฺธจกฺกํ นิฏฺฐิตํ.}
\prose{\csromanfolio{158}อาโรคฺยตฺถญฺจ สุขตฺถญฺจ เภสชฺช\-ตฺถญฺจ เจเตติ อุปกฺกมติ มุจฺจติ, อาปตฺติ สํฆาทิเสสสฺส ฯเปฯ. อาโรคฺยตฺถญฺจ สุขตฺถญฺจ ทวตฺถญฺจ เจเตติ อุปกฺกมติ มุจฺจติ, อาปตฺติ สํฆาทิเสสสฺส.}

\vspace{\tipitakaparskip}
\csromancenter{ทุมูลกสฺส ขณฺฑจกฺกํ.}
\prose{สุขตฺถญฺจ เภสชฺช\-ตฺถญฺจ ทานตฺถญฺจ เจเตติ อุปกฺกมติ มุจฺจติ, อาปตฺติ สํฆาทิเสสสฺส ฯเปฯ. สุขตฺถญฺจ เภสชฺช\-ตฺถญฺจ ทวตฺถญฺจ ฯเปฯ. สุขตฺถญฺจ เภสชฺช\-ตฺถญฺจ อาโรคฺยตฺถญฺจ เจเตติ อุปกฺกมติ มุจฺจติ, อาปตฺติ สํฆาทิเสสสฺส.}

\setlength{\csromangathapagewidth}{0pt}
\setlength{\csromangathawakwidth}{0pt}
\csromangathameasuregroup{ทุมูลกสฺส พทฺธจกฺกํ,}{\csromangathabat{ทุมูลกสฺส พทฺธจกฺกํ,}{สํขิตฺตํ.}}
\csromangathasetleft{ทุมูลกสฺส พทฺธจกฺกํ,}
\csromangathagroup{\csromangathastanza{\csromangathabat{ทุมูลกสฺส พทฺธจกฺกํ,}{สํขิตฺตํ.}}}

\prosewithcloser{วีมํสตฺถญฺจ ทวตฺถญฺจ อาโรคฺยตฺถญฺจ เจเตติ อุปกฺกมติ มุจฺจติ, อาปตฺติ สํฆาทิเสสสฺส ฯเปฯ. วีมํสตฺถญฺจ ทวตฺถญฺจ พีชตฺถญฺจ เจเตติ อุปกฺกมติ มุจฺจติ, อาปตฺติ สํฆาทิเสสสฺส.}{ทุมูลกํ นิฏฺฐิตํ.}
\prose{ติมูลกมฺปิ จตุมูลกมฺปิ ปญฺจมูลกมฺปิ ฉมูลกมฺปิ สตฺตมูลกมฺปิ อฏฺฐมูลกมฺปิ นวมูลกมฺปิ เอวเมว วิตฺถาเรตพฺพํ.}

\vspace{\tipitakaparskip}
\csromancenter{อิทํ สพฺพมูลกํ}
\proseitemwithcloser{243}{อาโรคฺยตฺถญฺจ สุขตฺถญฺจ เภสชฺช\-ตฺถญฺจ ทานตฺถญฺจ ปุญฺญตฺถญฺจ ยญฺญตฺถญฺจ สคฺคตฺถญฺจ พีชตฺถญฺจ วีมํสตฺถญฺจ ทวตฺถญฺจ เจเตติ อุปกฺกมติ มุจฺจติ, อาปตฺติ สํฆาทิเสสสฺส.}{สพฺพมูลกํ นิฏฺฐิตํ.}
\proseitem{244}{นีลญฺจ ปีตกญฺจ เจเตติ อุปกฺกมติ มุจฺจติ, อาปตฺติ สํฆาทิเสสสฺส.}

\prosewithcloser{นีลญฺจ โลหิตกญฺจ ฯเปฯ. นีลญฺจ โอทาตญฺจ. นีลญฺจ ตกฺกวณฺณญฺจ. นีลญฺจ ทกวณฺณญฺจ. นีลญฺจ เตลวณฺณญฺจ. นีลญฺจ ขีรวณฺณญฺจ. นีลญฺจ ทธิวณฺณญฺจ. นีลญฺจ สปฺปิวณฺณญฺจ เจเตติ อุปกฺกมติ มุจฺจติ, อาปตฺติ สํฆาทิเสสสฺส.}{เอกมูลกสฺส ขณฺฑจกฺกํ นิฏฺฐิตํ.}
\proseitem{245}{\csromanfolio{159}ปีตกญฺจ โลหิตกญฺจ เจเตติ อุปกฺกมติ มุจฺจติ, อาปตฺติ สํฆาทิเสสสฺส.}

\prose{ปีตกญฺจ โอทาตญฺจ ฯเปฯ. ปีตกญฺจ ตกฺกวณฺณญฺจ. ปีตกญฺจ ทกวณฺณญฺจ. ปีตกญฺจ เตลวณฺณญฺจ. ปีตกญฺจ ขีรวณฺณญฺจ. ปีตกญฺจ ทธิวณฺณญฺจ. ปีตกญฺจ สปฺปิวณฺณญฺจ เจเตติ อุปกฺกมติ มุจฺจติ, อาปตฺติ สํฆาทิเสสสฺส.}

\prose{ปีตกญฺจ นีลญฺจ เจเตติ อุปกฺกมติ มุจฺจติ, อาปตฺติ สํฆาทิเสสสฺส.}

\vspace{\tipitakaparskip}
\csromancenter{เอกมูลกสฺส พทฺธจกฺกํ.}
\proseitem{246}{โลหิตกญฺจ โอทาตญฺจ เจเตติ อุปกฺกมติ มุจฺจติ, อาปตฺติ สํฆาทิเสสสฺส.}

\prose{โลหิตกญฺจ ตกฺกวณฺณญฺจ ฯเปฯ. โลหิตกญฺจ ทกวณฺณญฺจ. โลหิตกญฺจ เตลวณฺณญฺจ. โลหิตกญฺจ ขีรวณฺณญฺจ. โลหิตกญฺจ ทธิวณฺณญฺจ. โลหิตกญฺจ สปฺปิวณฺณญฺจ เจเตติ อุปกฺกมติ มุจฺจติ, อาปตฺติ สํฆาทิเสสสฺส.}

\prose{โลหิตกญฺจ นีลญฺจ ฯเปฯ. โลหิตกญฺจ ปีตกญฺจ เจเตติ อุปกฺกมติ มุจฺจติ, อาปตฺติ สํฆาทิเสสสฺส.}

\prose{โอทาตญฺจ ตกฺกวณฺณญฺจ ฯเปฯ. โอทาตญฺจ ทกวณฺณญฺจ. โอทาตญฺจ เตลวณฺณญฺจ. โอทาตญฺจ ขีรวณฺณญฺจ. โอทาตญฺจ ทธิวณฺณญฺจ. โอทาตญฺจ สปฺปิวณฺณญฺจ เจเตติ อุปกฺกมติ มุจฺจติ, อาปตฺติ สํฆาทิเสสสฺส.}

\prose{โอทาตญฺจ นีลญฺจ ฯเปฯ. โอทาตญฺจ ปีตกญฺจ. โอทาตญฺจ โลหิตกญฺจ เจเตติ อุปกฺกมติ มุจฺจติ, อาปตฺติ สํฆาทิเสสสฺส.}

\prose{ตกฺกวณฺณญฺจ ทกวณฺณญฺจ ฯเปฯ. ตกฺกวณฺณญฺจ เตลวณฺณญฺจ. ตกฺกวณฺณญฺจ ขีรวณฺณญฺจ. ตกฺกวณฺณญฺจ ทธิวณฺณญฺจ. ตกฺกวณฺณญฺจ สปฺปิวณฺณญฺจ เจเตติ อุปกฺกมติ มุจฺจติ, อาปตฺติ สํฆาทิเสสสฺส.}

\prose{ตกฺกวณฺณญฺจ นีลญฺจ ฯเปฯ. ตกฺกวณฺณญฺจ ปีตกญฺจ. ตกฺกวณฺณญฺจ โลหิตกญฺจ. ตกฺกวณฺณญฺจ โอทาตญฺจ เจเตติ อุปกฺกมติ มุจฺจติ, อาปตฺติ สํฆาทิเสสสฺส.}

\prose{\csromanfolio{160}ทกวณฺณญฺจ เตลวณฺณญฺจ ฯเปฯ. ทกวณฺณญฺจ ขีรวณฺณญฺจ. ทกวณฺณญฺจ ทธิวณฺณญฺจ. ทกวณฺณญฺจ สปฺปิวณฺณญฺจ เจเตติ อุปกฺกมติ มุจฺจติ, อาปตฺติ สํฆาทิเสสสฺส.}

\prose{ทกวณฺณญฺจ นีลญฺจ ฯเปฯ. ทกวณฺณญฺจ ปีตกญฺจ. ทกวณฺณญฺจ โลหิตกญฺจ. ทกวณฺณญฺจ โอทาตญฺจ. ทกวณฺณญฺจ ตกฺกวณฺณญฺจ เจเตติ อุปกฺกมติ มุจฺจติ, อาปตฺติ สํฆาทิเสสสฺส.}

\prose{เตลวณฺณญฺจ ขีรวณฺณญฺจ ฯเปฯ. เตลวณฺณญฺจ ทธิวณฺณญฺจ. เตลวณฺณญฺจ สปฺปิวณฺณญฺจ เจเตติ อุปกฺกมติ มุจฺจติ, อาปตฺติ สํฆาทิเสสสฺส.}

\prose{เตลวณฺณญฺจ นีลญฺจ ฯเปฯ. เตลวณฺณญฺจ ปีตกญฺจ. เตลวณฺณญฺจ โลหิตกญฺจ. เตลวณฺณญฺจ โอทาตญฺจ. เตลวณฺณญฺจ ตกฺกวณฺณญฺจ. เตลวณฺณญฺจ ทกวณฺณญฺจ เจเตติ อุปกฺกมติ มุจฺจติ, อาปตฺติ สํฆาทิเสสสฺส.}

\prose{ขีรวณฺณญฺจ ทธิวณฺณญฺจ ฯเปฯ. ขีรวณฺณญฺจ สปฺปิวณฺณญฺจ เจเตติ อุปกฺกมติ มุจฺจติ, อาปตฺติ สํฆาทิเสสสฺส.}

\prose{ขีรวณฺณญฺจ นีลญฺจ ฯเปฯ. ขีรวณฺณญฺจ ปีตกญฺจ. ขีรวณฺณญฺจ โลหิตกญฺจ. ขีรวณฺณญฺจ โอทาตญฺจ. ขีรวณฺณญฺจ ตกฺกวณฺณญฺจ. ขีรวณฺณญฺจ ทกวณฺณญฺจ. ขีรวณฺณญฺจ เตลวณฺณญฺจ เจเตติ อุปกฺกมติ มุจฺจติ, อาปตฺติ สํฆาทิเสสสฺส.}

\prose{ทธิวณฺณญฺจ สปฺปิวณฺณญฺจ เจเตติ อุปกฺกมติ มุจฺจติ, อาปตฺติ สํฆาทิเสสสฺส.}

\prose{ทธิวณฺณญฺจ นีลญฺจ ฯเปฯ. ทธิวณฺณญฺจ ปีตกญฺจ. ทธิวณฺณญฺจ โลหิตกญฺจ. ทธิวณฺณญฺจ โอทาตญฺจ. ทธิวณฺณญฺจ ตกฺกวณฺณญฺจ. ทธิวณฺณญฺจ ทกวณฺณญฺจ. ทธิวณฺณญฺจ เตลวณฺณญฺจ. ทธิวณฺณญฺจ ขีรวณฺณญฺจ เจเตติ อุปกฺกมติ มุจฺจติ, อาปตฺติ สํฆาทิเสสสฺส.}

\prose{สปฺปิวณฺณญฺจ นีลญฺจ เจเตติ อุปกฺกมติ มุจฺจติ, อาปตฺติ สํฆาทิเสสสฺส.}

\prosewithcloser{สปฺปิวณฺณญฺจ ปีตกญฺจ ฯเปฯ. สปฺปิวณฺณญฺจ โลหิตกญฺจ. สปฺปิวณฺณญฺจ โอทาตญฺจ. สปฺปิวณฺณญฺจ ตกฺกวณฺณญฺจ. สปฺปิวณฺณญฺจ ทกวณฺณญฺจ. สปฺปิวณฺณญฺจ เตลวณฺณญฺจ. สปฺปิวณฺณญฺจ ขีรวณฺณญฺจ. สปฺปิวณฺณญฺจ ทธิวณฺณญฺจ เจเตติ อุปกฺกมติ มุจฺจติ, อาปตฺติ สํฆาทิเสสสฺส.}{เอกมูลกสฺส พทฺธจกฺกํ นิฏฺฐิตํ.}
\prose{\csromanfolio{161}นีลญฺจ ปีตกญฺจ โลหิตกญฺจ เจเตติ อุปกฺกมติ มุจฺจติ, อาปตฺติ สํฆาทิเสสสฺส ฯเปฯ. นีลญฺจ ปีตกญฺจ สปฺปิวณฺณญฺจ เจเตติ อุปกฺกมติ มุจฺจติ, อาปตฺติ สํฆาทิเสสสฺส.}

\vspace{\tipitakaparskip}
\csromancenter{ทุมูลกสฺส ขณฺฑจกฺกํ.}
\prose{ปีตกญฺจ โลหิตกญฺจ โอทาตญฺจ เจเตติ อุปกฺกมติ มุจฺจติ, อาปตฺติ สํฆาทิเสสสฺส ฯเปฯ. ปีตกญฺจ โลหิตกญฺจ สปฺปิวณฺณญฺจ ฯเปฯ. ปีตกญฺจ โลหิตกญฺจ นีลญฺจ เจเตติ อุปกฺกมติ มุจฺจติ, อาปตฺติ สํฆาทิเสสสฺส.}

\setlength{\csromangathapagewidth}{0pt}
\setlength{\csromangathawakwidth}{0pt}
\csromangathameasuregroup{ทุมูลกสฺส พทฺธจกฺกํ,}{\csromangathabat{ทุมูลกสฺส พทฺธจกฺกํ,}{สํขิตฺตํ.}}
\csromangathasetleft{ทุมูลกสฺส พทฺธจกฺกํ,}
\csromangathagroup{\csromangathastanza{\csromangathabat{ทุมูลกสฺส พทฺธจกฺกํ,}{สํขิตฺตํ.}}}

\prosewithcloser{ทธิวณฺณญฺจ สปฺปิวณฺณญฺจ นีลญฺจ เจเตติ อุปกฺกมติ มุจฺจติ, อาปตฺติ สํฆาทิเสสสฺส ฯเปฯ. ทธิวณฺณญฺจ สปฺปิวณฺณญฺจ ขีรวณฺณญฺจ เจเตติ อุปกฺกมติ มุจฺจติ, อาปตฺติ สํฆาทิเสสสฺส.}{ทุมูลกํ นิฏฺฐิตํ.}
\prose{ติมูลกมฺปิ จตุมูลกมฺปิ ปญฺจมูลกมฺปิ ฉมูลกมฺปิ สตฺตมูลกมฺปิ อฏฺฐมูลกมฺปิ นวมูลกมฺปิ เอวเมว วิตฺถาเรตพฺพํ.}

\vspace{\tipitakaparskip}
\csromancenter{อิทํ สพฺพมูลกํ}
\proseitemwithcloser{247}{นีลญฺจ ปีตกญฺจ โลหิตกญฺจ โอทาตญฺจ ตกฺกวณฺณญฺจ ทกวณฺณญฺจ เตลวณฺณญฺจ ขีรวณฺณญฺจ ทธิวณฺณญฺจ สปฺปิวณฺณญฺจ เจเตติ อุปกฺกมติ มุจฺจติ, อาปตฺติ สํฆาทิเสสสฺส.}{สพฺพมูลกํ นิฏฺฐิตํ.}
\proseitem{248}{อาโรคฺยตฺถญฺจ นีลญฺจ\csromansharedfootnote{9b43b6504e9d91b9}{อาโรคฺยตฺถํ นีลํ (?)} เจเตติ อุปกฺกมติ มุจฺจติ, อาปตฺติ สํฆาทิเสสสฺส.}

\prose{อาโรคฺยตฺถญฺจ สุขตฺถญฺจ นีลญฺจ ปีตกญฺจ เจเตติ อุปกฺกมติ มุจฺจติ, อาปตฺติ สํฆาทิเสสสฺส.}

\prose{\csromanfolio{162}อาโรคฺยตฺถญฺจ สุขตฺถญฺจ เภสชฺช\-ตฺถญฺจ นีลญฺจ ปีตกญฺจ โลหิตกญฺจ เจเตติ อุปกฺกมติ มุจฺจติ, อาปตฺติ สํฆาทิเสสสฺส.}

\vspace{\tipitakaparskip}
\csromancenter{อุภโต วฑฺฒกํ เอวเมว วฑฺเฒตพฺพํ.}
\prosewithcloser{อาโรคฺยตฺถญฺจ สุขตฺถญฺจ เภสชฺช\-ตฺถญฺจ ทานตฺถญฺจ ปุญฺญตฺถญฺจ ยญฺญตฺถญฺจ สคฺคตฺถญฺจ พีชตฺถญฺจ วีมํสตฺถญฺจ ทวตฺถญฺจ นีลญฺจ ปีตกญฺจ โลหิตกญฺจ โอทาตญฺจ ตกฺกวณฺณญฺจ ทกวณฺณญฺจ เตลวณฺณญฺจ ขีรวณฺณญฺจ ทธิวณฺณญฺจ สปฺปิวณฺณญฺจ เจเตติ อุปกฺกมติ มุจฺจติ, อาปตฺติ สํฆาทิเสสสฺส.}{มิสฺสกจกฺกํ นิฏฺฐิตํ.}
\proseitem{249}{นีลํ โมเจสฺสามีติ เจเตติ อุปกฺกมติ, ปีตกํ มุจฺจติ, อาปตฺติ สํฆาทิเสสสฺส.}

\prosewithcloser{นีลํ โมเจสฺสามีติ เจเตติ อุปกฺกมติ, โลหิตกํ ฯเปฯ. โอทาตํ. ตกฺกวณฺณํ. ทกวณฺณํ. เตลวณฺณํ. ขีรวณฺณํ. ทธิวณฺณํ. สปฺปิวณฺณํ มุจฺจติ, อาปตฺติ สํฆาทิเสสสฺส.}{ขณฺฑจกฺกํ นิฏฺฐิตํ.}
\proseitem{250}{ปีตกํ โมเจสฺสามีติ เจเตติ อุปกฺกมติ, โลหิตกํ มุจฺจติ, อาปตฺติ สํฆาทิเสสสฺส.}

\prose{ปีตกํ โมเจสฺสามีติ เจเตติ อุปกฺกมติ, โอทาตํ ฯเปฯ. ตกฺกวณฺณํ. ทกวณฺณํ. เตลวณฺณํ. ขีรวณฺณํ. ทธิวณฺณํ. สปฺปิวณฺณํ. นีลํ มุจฺจติ, อาปตฺติ สํฆาทิเสสสฺส.}

\vspace{\tipitakaparskip}
\csromancenter{พทฺธจกฺกมูลํ สํขิตฺตํ.}
\proseitem{251}{สปฺปิวณฺณํ โมเจสฺสามีติ เจเตติ อุปกฺกมติ, นีลํ มุจฺจติ, อาปตฺติ สํฆาทิเสสสฺส.}

\prosewithcloser{สปฺปิวณฺณํ โมเจสฺสามีติ เจเตติ อุปกฺกมติ, ปีตกํ ฯเปฯ. โลหิตกํ. โอทาตํ. ตกฺกวณฺณํ. ทกวณฺณํ. เตลวณฺณํ. ขีรวณฺณํ. ทธิวณฺณํ มุจฺจติ, อาปตฺติ สํฆาทิเสสสฺส.}{กุจฺฉิจกฺกํ นิฏฺฐิตํ.}
\proseitem{252}{\csromanfolio{163}ปีตกํ โมเจสฺสามีติ เจเตติ อุปกฺกมติ, นีลํ มุจฺจติ, อาปตฺติ สํฆาทิเสสสฺส.}

\prosewithcloser{โลหิตกํ ฯเปฯ. โอทาตํ. ตกฺกวณฺณํ. ทกวณฺณํ. เตลวณฺณํ. ขีรวณฺณํ. ทธิวณฺณํ. สปฺปิวณฺณํ โมเจสฺสามีติ เจเตติ อุปกฺกมติ, นีลํ มุจฺจติ, อาปตฺติ สํฆาทิเสสสฺส.}{ปิฏฺฐิจกฺกสฺส ปฐมํ คมนํ นิฏฺฐิตํ.}
\proseitem{253}{โลหิตกํ โมเจสฺสามีติ เจเตติ อุปกฺกมติ, ปีตกํ มุจฺจติ, อาปตฺติ สํฆาทิเสสสฺส.}

\prosewithcloser{โอทาตํ ฯเปฯ. ตกฺกวณฺณํ. ทกวณฺณํ. เตลวณฺณํ. ขีรวณฺณํ. ทธิวณฺณํ. สปฺปิวณฺณํ. นีลํ โมเจสฺสามีติ เจเตติ อุปกฺกมติ, ปีตกํ มุจฺจติ, อาปตฺติ สํฆาทิเสสสฺส.}{ปิฏฺฐิจกฺกสฺส ทุติยํ คมนํ นิฏฺฐิตํ.}
\proseitem{254}{โอทาตํ โมเจสฺสามีติ เจเตติ อุปกฺกมติ, โลหิตกํ มุจฺจติ, อาปตฺติ สํฆาทิเสสสฺส.}

\prosewithcloser{ตกฺกวณฺณํ ฯเปฯ. ทกวณฺณํ. เตลวณฺณํ. ขีรวณฺณํ. ทธิวณฺณํ. สปฺปิวณฺณํ. นีลํ. ปีตกํ โมเจสฺสามีติ เจเตติ อุปกฺกมติ, โลหิตกํ มุจฺจติ, อาปตฺติ สํฆาทิเสสสฺส.}{ปิฏฺฐิจกฺกสฺส ตติยํ คมนํ นิฏฺฐิตํ.}
\proseitem{255}{ตกฺกวณฺณํ โมเจสฺสามีติ เจเตติ อุปกฺกมติ, โอทาตํ มุจฺจติ, อาปตฺติ สํฆาทิเสสสฺส.}

\prosewithcloser{ทกวณฺณํ ฯเปฯ. เตลวณฺณํ. ขีรวณฺณํ. ทธิวณฺณํ. สปฺปิวณฺณํ. นีลํ. ปีตกํ. โลหิตกํ โมเจสฺสามีติ เจเตติ อุปกฺกมติ, โอทาตํ มุจฺจติ, อาปตฺติ สํฆาทิเสสสฺส.}{ปิฏฺฐิจกฺกสฺส จตุตฺถํ คมนํ นิฏฺฐิตํ.}
\proseitem{256}{ทกวณฺณํ โมเจสฺสามีติ เจเตติ อุปกฺกมติ, ตกฺกวณฺณํ มุจฺจติ, อาปตฺติ สํฆาทิเสสสฺส.}

\prosewithcloser{\csromanfolio{164}เตลวณฺณํ ฯเปฯ. ขีรวณฺณํ. ทธิวณฺณํ. สปฺปิวณฺณํ. นีลํ. ปีตกํ. โลหิตกํ. โอทาตํ โมเจสฺสามีติ เจเตติ อุปกฺกมติ, ตกฺกวณฺณํ มุจฺจติ, อาปตฺติ สํฆาทิเสสสฺส.}{ปิฏฺฐิจกฺกสฺส ปญฺจมํ คมนํ นิฏฺฐิตํ.}
\proseitem{257}{เตลวณฺณํ โมเจสฺสามีติ เจเตติ อุปกฺกมติ, ทกวณฺณํ มุจฺจติ, อาปตฺติ สํฆาทิเสสสฺส.}

\prosewithcloser{ขีรวณฺณํ ฯเปฯ. ทธิวณฺณํ. สปฺปิวณฺณํ. นีลํ. ปีตกํ. โลหิตกํ. โอทาตํ. ตกฺกวณฺณํ โมเจสฺสามีติ เจเตติ อุปกฺกมติ, ทกวณฺณํ มุจฺจติ, อาปตฺติ สํฆาทิเสสสฺส.}{ปิฏฺฐิจกฺกสฺส ฉฏฺฐํ คมนํ นิฏฺฐิตํ.}
\proseitem{258}{ขีรวณฺณํ โมเจสฺสามีติ เจเตติ อุปกฺกมติ, เตลวณฺณํ มุจฺจติ, อาปตฺติ สํฆาทิเสสสฺส.}

\prosewithcloser{ทธิวณฺณํ ฯเปฯ. สปฺปิวณฺณํ. นิลํ. ปีตกํ. โลหิตกํ. โอทาตํ. ตกฺกวณฺณํ. ทกวณฺณํ โมเจสฺสามีติ เจเตติ อุปกฺกมติ, เตลวณฺณํ มุจฺจติ, อาปตฺติ สํฆาทิเสสสฺส.}{ปิฏฺฐิจกฺกสฺส สตฺตมํ คมนํ นิฏฺฐิตํ.}
\proseitem{259}{ทธิวณฺณํ โมเจสฺสามีติ เจเตติ อุปกฺกมติ, ขีรวณฺณํ มุจฺจติ, อาปตฺติ สํฆาทิเสสสฺส.}

\prosewithcloser{สปฺปิวณฺณํ ฯเปฯ. นีลํ. ปีตกํ. โลหิตกํ. โอทาตํ. ตกฺกวณฺณํ. ทกวณฺณํ. เตลวณฺณํ โมเจสฺสามีติ เจเตติ อุปกฺกมติ, ขีรวณฺณํ มุจฺจติ, อาปตฺติ สํฆาทิเสสสฺส.}{ปิฏฺฐิจกฺกสฺส อฏฺฐมํ คมนํ นิฏฺฐิตํ.}
\proseitem{260}{สปฺปิวณฺณํ โมเจสฺสามีติ เจเตติ อุปกฺกมติ, ทธิวณฺณํ มุจฺจติ, อาปตฺติ สํฆาทิเสสสฺส.}

\prosewithcloser{นีลํ ฯเปฯ. ปีตกํ. โลหิตกํ. โอทาตํ. ตกฺกวณฺณํ. ทกวณฺณํ. เตลวณฺณํ. ขีรวณฺณํ โมเจสฺสามีติ เจเตติ อุปกฺกมติ, ทธิวณฺณํ มุจฺจติ, อาปตฺติ สํฆาทิเสสสฺส.}{ปิฏฺฐิจกฺกสฺส นวมํ คมนํ นิฏฺฐิตํ.}
\proseitem{261}{\csromanfolio{165}นีลํ โมเจสฺสามีติ เจเตติ อุปกฺกมติ, สปฺปิวณฺณํ มุจฺจติ, อาปตฺติ สํฆาทิเสสสฺส.}

\prosewithcloser{ปีตกํ ฯเปฯ. โลหิตกํ. โอทาตํ. ตกฺกวณฺณํ. ทกวณฺณํ. เตลวณฺณํ. ขีรวณฺณํ. ทธิวณฺณํ โมเจสฺสามีติ เจเตติ อุปกฺกมติ, สปฺปิวณฺณํ มุจฺจติ, อาปตฺติ สํฆาทิเสสสฺส.}{ปิฏฺฐิจกฺกสฺส ทสมํ คมนํ นิฏฺฐิตํ.}
\nitthitam{ปิฏฺฐิจกฺกํ นิฏฺฐิตํ.}
\setlength{\csromangathapagewidth}{0pt}
\setlength{\csromangathawakwidth}{0pt}
\csromangathameasuregroup{เจเตติ อุปกฺกมติ มุจฺจติ, \\ เจเตติ อุปกฺกมติ น มุจฺจติ, \\ เจเตติ น อุปกฺกมติ มุจฺจติ, \\ เจเตติ น อุปกฺกมติ น มุจฺจติ, \\ น เจเตติ อุปกฺกมติ มุจฺจติ, \\ น เจเตติ อุปกฺกมติ น มุจฺจติ, \\ น เจเตติ น อุปกฺกมติ มุจฺจติ, \\ น เจเตติ น อุปกฺกมติ น มุจฺจติ,}{\csromangathabat{เจเตติ อุปกฺกมติ มุจฺจติ,}{อาปตฺติ สํฆาทิเสสสฺส.} \\ \csromangathabat{เจเตติ อุปกฺกมติ น มุจฺจติ,}{อาปตฺติ ถุลฺลจฺจยสฺส.} \\ \csromangathabat{เจเตติ น อุปกฺกมติ มุจฺจติ,}{อนาปตฺติ.} \\ \csromangathabat{เจเตติ น อุปกฺกมติ น มุจฺจติ,}{อนาปตฺติ.} \\ \csromangathabat{น เจเตติ อุปกฺกมติ มุจฺจติ,}{อนาปตฺติ.} \\ \csromangathabat{น เจเตติ อุปกฺกมติ น มุจฺจติ,}{อนาปตฺติ.} \\ \csromangathabat{น เจเตติ น อุปกฺกมติ มุจฺจติ,}{อนาปตฺติ.} \\ \csromangathabat{น เจเตติ น อุปกฺกมติ น มุจฺจติ,}{อนาปตฺติ.}}
\csromangathasetleft{เจเตติ อุปกฺกมติ มุจฺจติ, \\ เจเตติ อุปกฺกมติ น มุจฺจติ, \\ เจเตติ น อุปกฺกมติ มุจฺจติ, \\ เจเตติ น อุปกฺกมติ น มุจฺจติ, \\ น เจเตติ อุปกฺกมติ มุจฺจติ, \\ น เจเตติ อุปกฺกมติ น มุจฺจติ, \\ น เจเตติ น อุปกฺกมติ มุจฺจติ, \\ น เจเตติ น อุปกฺกมติ น มุจฺจติ,}
\csromangathagroup{\csromangathastanza{\csromangathaitembat{262}{เจเตติ อุปกฺกมติ มุจฺจติ,}{อาปตฺติ สํฆาทิเสสสฺส.} \\ \csromangathacontbat{เจเตติ อุปกฺกมติ น มุจฺจติ,}{อาปตฺติ ถุลฺลจฺจยสฺส.}}\csromangathastanza{\csromangathabat{เจเตติ น อุปกฺกมติ มุจฺจติ,}{อนาปตฺติ.} \\ \csromangathabat{เจเตติ น อุปกฺกมติ น มุจฺจติ,}{อนาปตฺติ.}}\csromangathastanza{\csromangathabat{น เจเตติ อุปกฺกมติ มุจฺจติ,}{อนาปตฺติ.} \\ \csromangathabat{น เจเตติ อุปกฺกมติ น มุจฺจติ,}{อนาปตฺติ.}}\csromangathastanza{\csromangathabat{น เจเตติ น อุปกฺกมติ มุจฺจติ,}{อนาปตฺติ.} \\ \csromangathabat{น เจเตติ น อุปกฺกมติ น มุจฺจติ,}{อนาปตฺติ.}}}

\prose{อนาปตฺติ สุปินนฺเตน นโมจนา\-ธิปฺปายสฺส อุมฺมตฺตกสฺส ขิตฺตจิตฺตสฺส เวทนาฏฺฏสฺส อาทิกมฺมิกสฺสาติ.}
\csromansectionrule

\csromanmatikaanchor{mtk.41}
\csromantocmarkref{subsection}{วินีตวตฺถุอุทฺทานคาถา}{mtk.41}
\bookmark[dest={mtk.41},level=2]{วินีตวตฺถุอุทฺทานคาถา}
\csromanheader{2}{วินีตวตฺถุ\-อุทฺทานคาถา}
\setlength{\csromangathapagewidth}{0pt}
\setlength{\csromangathawakwidth}{0pt}
\csromangathameasuregroup{สุปิโนจฺจารปสฺสาโว, \\ เภสชฺชํ กณฺฑุวํ มคฺโค, \\ สามเณโร จ สุตฺโต จ, \\ อากาเส ถมฺภํ นิชฺฌายิ, \\ โสโต อุทญฺชลํ ธาวํ, \\ วาลิกา กทฺทมุสฺเสโก\textsuperscript{1},}{\csromangathabat{สุปิโนจฺจารปสฺสาโว,}{วิตกฺกุโณฺหทเกน จ.} \\ \csromangathabat{เภสชฺชํ กณฺฑุวํ มคฺโค,}{วตฺถิ ชนฺตาฆรุปกฺกโม\textsuperscript{1}.} \\ \csromangathabat{สามเณโร จ สุตฺโต จ,}{อูรุ มุฏฺฐินา ปีฬยิ.} \\ \csromangathabat{อากาเส ถมฺภํ นิชฺฌายิ,}{ฉิทฺทํ กฏฺเฐน ฆฏฺฏยิ.} \\ \csromangathabat{โสโต อุทญฺชลํ ธาวํ,}{ปุปฺผาวลิยํ โปกฺขรํ.} \\ \csromangathabat{วาลิกา กทฺทมุสฺเสโก\textsuperscript{1},}{สยนงฺคุฏฺฐเกน จาติ.}}
\csromangathasetleft{สุปิโนจฺจารปสฺสาโว, \\ เภสชฺชํ กณฺฑุวํ มคฺโค, \\ สามเณโร จ สุตฺโต จ, \\ อากาเส ถมฺภํ นิชฺฌายิ, \\ โสโต อุทญฺชลํ ธาวํ, \\ วาลิกา กทฺทมุสฺเสโก\textsuperscript{1},}
\csromangathagroup{\csromangathastanza{\csromangathabat{สุปิโนจฺจารปสฺสาโว,}{วิตกฺกุโณฺหทเกน จ.} \\ \csromangathabat{เภสชฺชํ กณฺฑุวํ มคฺโค,}{วตฺถิ ชนฺตาฆรุปกฺกโม\csromansharedfootnote{6282ef23fa852832}{ชนฺตคฺคุปกฺกโม (สี), ชนฺตาฆรํ อูรุ (สฺยา)}.}}\csromangathastanza{\csromangathabat{สามเณโร จ สุตฺโต จ,}{อูรุ มุฏฺฐินา ปีฬยิ.} \\ \csromangathabat{อากาเส ถมฺภํ นิชฺฌายิ,}{ฉิทฺทํ กฏฺเฐน ฆฏฺฏยิ.}}\csromangathastanza{\csromangathabat{โสโต อุทญฺชลํ ธาวํ,}{ปุปฺผาวลิยํ โปกฺขรํ.} \\ \csromangathabat{วาลิกา กทฺทมุสฺเสโก\csromansharedfootnote{5bc70fe4bcdcaf8c}{กทฺทโมเสโก (?)},}{สยนงฺคุฏฺฐเกน จาติ.}}}

\csromanmatikaanchor{mtk.42}
\csromantocmarkref{subsection}{วินีตวตฺถุ}{mtk.42}
\bookmark[dest={mtk.42},level=2]{วินีตวตฺถุ}
\csromanheader{2}{\textbf{วินีตวตฺถุ}}
\proseitem{263}{\csromanfolio{166}เตน โข ปน สมเยน อญฺญตรสฺส ภิกฺขุโน สุปินนฺเตน อสุจิ มุจฺจิ. ตสฺส กุกฺกุจฺจํ อโหสิ “ภควตา สิกฺขาปทํ ปญฺญตฺตํ, กจฺจิ นุ โข อหํ สํฆาทิเสสํ อาปตฺติํ อาปนฺโน”ติ. ภควโต เอตมตฺถํ อาโรเจสิ. อนาปตฺติ ภิกฺขุ สุปินนฺเตนาติ. (1)}

\prose{เตน โข ปน สมเยน อญฺญตรสฺส ภิกฺขุโน อุจฺจารํ กโรนฺตสฺส อสุจิ มุจฺจิ. ตสฺส กุกฺกุจฺจํ อโหสิ ฯเปฯ. กิํจิตฺโต ตฺวํ ภิกฺขูติ. นาหํ ภควา โมจนา\-ธิปฺปาโยติ. อนาปตฺติ ภิกฺขุ นโมจนา\-ธิปฺปายสฺสาติ. (2)}

\prose{เตน โข ปน สมเยน อญฺญตรสฺส ภิกฺขุโน ปสฺสาวํ กโรนฺตสฺส อสุจิ มุจฺจิ. ตสฺส กุกฺกุจฺจํ อโหสิ ฯเปฯ. อนาปตฺติ ภิกฺขุ นโมจนา\-ธิปฺปายสฺสาติ. (3)}

\prose{เตน โข ปน สมเยน อญฺญตรสฺส ภิกฺขุโน กามวิตกฺกํ วิตกฺเกนฺตสฺส อสุจิ มุจฺจิ. ตสฺส กุกฺกุจฺจํ อโหสิ ฯเปฯ. อนาปตฺติ ภิกฺขุ วิตกฺเก\-นฺตสฺสาติ. (4)}

\prose{เตน โข ปน สมเยน อญฺญตรสฺส ภิกฺขุโน อุโณฺหทเกน นฺหายนฺตสฺส อสุจิ มุจฺจิ. ตสฺส กุกฺกุจฺจํ อโหสิ ฯเปฯ. กิํจิตฺโต ตฺวํ ภิกฺขูติ. นาหํ ภควา โมจนา\-ธิปฺปาโยติ. อนาปตฺติ ภิกฺขุ นโมจนา\-ธิปฺปายสฺสาติ. (5)}

\prose{เตน โข ปน สมเยน อญฺญตรสฺส ภิกฺขุโน โมจนา\-ธิปฺปายสฺส อุโณฺหทเกน นฺหายนฺตสฺส อสุจิ มุจฺจิ. ตสฺส กุกฺกุจฺจํ อโหสิ ฯเปฯ. อาปตฺติํ ตฺวํ ภิกฺขุ อาปนฺโน สํฆาทิเสสนฺติ. (6)}

\prose{เตน โข ปน สมเยน อญฺญตรสฺส ภิกฺขุโน โมจนา\-ธิปฺปายสฺส อุโณฺหทเกน นฺหายนฺตสฺส อสุจิ น มุจฺจิ. ตสฺส กุกฺกุจฺจํ อโหสิ ฯเปฯ. อนาปตฺติ ภิกฺขุ สํฆาทิเสสสฺส, อาปตฺติ ถุลฺลจฺจยสฺสาติ. (7)}

\prose{เตน โข ปน สมเยน อญฺญตรสฺส ภิกฺขุโน องฺคชาเต วโณ โหติ. เภสชฺเชน\csromansharedfootnote{795fef713506b262}{ตสฺส เภสชฺเชน (?)} อาลิมฺเปนฺตสฺส อสุจิ มุจฺจิ. ตสฺส กุกฺกุจฺจํ อโหสิ ฯเปฯ. อนาปตฺติ ภิกฺขุ นโมจนา\-ธิปฺปายสฺสาติ. (8)}

\prose{\csromanfolio{167}เตน โข ปน สมเยน อญฺญตรสฺส ภิกฺขุโน องฺคชาเต วโณ โหติ. โมจนา\-ธิปฺปายสฺส\csromansharedfootnote{4cdd54e170dbdec3}{ตสฺส โมจนาธิปฺปายสฺส (สฺยา)} เภสชฺเชน อาลิมฺเปนฺตสฺส อสุจิ มุจฺจิ ฯเปฯ อสุจิ น มุจฺจิ. ตสฺส กุกฺกุจฺจํ อโหสิ ฯเปฯ. อนาปตฺติ ภิกฺขุ สํฆาทิเสสสฺส, อาปตฺติ ถุลฺลจฺจยสฺสาติ. \mbox{(9-10)}}

\prose{เตน โข ปน สมเยน อญฺญตรสฺส ภิกฺขุโน อณฺฑํ กณฺฑุวนฺตสฺส อสุจิ มุจฺจิ. ตสฺส กุกฺกุจฺจํ อโหสิ ฯเปฯ. อนาปตฺติ ภิกฺขุ นโมจนา\-ธิปฺปายสฺสาติ. (11)}

\prose{เตน โข ปน สมเยน อญฺญตรสฺส ภิกฺขุโน โมจนา\-ธิปฺปายสฺส อณฺฑํ กณฺฑุวนฺตสฺส\csromansharedfootnote{50f41d65468f4f97}{กณฺฑูวนฺตสฺส (สี)} อสุจิ มุจฺจิ ฯเปฯ อสุจิ น มุจฺจิ. ตสฺส กุกฺกุจฺจํ อโหสิ ฯเปฯ. อนาปตฺติ ภิกฺขุ สํฆาทิเสสสฺส, อาปตฺติ ถุลฺลจฺจยสฺสาติ. \mbox{(12-13)}}

\proseitem{264}{เตน โข ปน สมเยน อญฺญตรสฺส ภิกฺขุโน มคฺคํ คจฺฉนฺตสฺส อสุจิ มุจฺจิ. ตสฺส กุกฺกุจฺจํ อโหสิ ฯเปฯ. อนาปตฺติ ภิกฺขุ นโมจนา\-ธิปฺปายสฺสาติ. (14)}

\prose{เตน โข ปน สมเยน อญฺญตรสฺส ภิกฺขุโน โมจนา\-ธิปฺปายสฺส มคฺคํ คจฺฉนฺตสฺส อสุจิ มุจฺจิ ฯเปฯ อสุจิ น มุจฺจิ. ตสฺส กุกฺกุจฺจํ อโหสิ ฯเปฯ. อนาปตฺติ ภิกฺขุ สํฆาทิเสสสฺสา, อาปตฺติ ถุลฺลจฺจยสฺสาติ. \mbox{(15-16)}}

\prose{เตน โข ปน สมเยน อญฺญตรสฺส ภิกฺขุโน วตฺถิํ คเหตฺวา ปสฺสาวํ กโรนฺตสฺส อสุจิ มุจฺจิ. ตสฺส กุกฺกุจฺจํ อโหสิ ฯเปฯ. อนาปตฺติ ภิกฺขุ นโมจนา\-ธิปฺปายสฺสาติ. (17)}

\prose{เตน โข ปน สมเยน อญฺญตรสฺส ภิกฺขุโน โมจนา\-ธิปฺปายสฺส วตฺถิํ คเหตฺวา ปสฺสาวํ กโรนฺตสฺส อสุจิ มุจฺจิ ฯเปฯ อสุจิ น มุจฺจิ. ตสฺส กุกฺกุจฺจํ อโหสิ ฯเปฯ. อนาปตฺติ ภิกฺขุ สํฆาทิเสสสฺส, อาปตฺติ ถุลฺลจฺจยสฺสาติ. \mbox{(18-19)}}

\prose{เตน โข ปน สมเยน อญฺญตรสฺส ภิกฺขุโน ชนฺตาฆเร อุทรวฏฺฏิํ ตาเปนฺตสฺส อสุจิ มุจฺจิ. ตสฺส กุกฺกุจฺจํ อโหสิ ฯเปฯ. อนาปตฺติ ภิกฺขุ นโมจนา\-ธิปฺปายสฺสาติ. (20)}

\prose{\csromanfolio{168}เตน โข ปน สมเยน อญฺญตรสฺส ภิกฺขุโน โมจนา\-ธิปฺปายสฺส ชนฺตาฆเร อุทรวฏฺฏิํ ตาเปนฺตสฺส อสุจิ มุจฺจิ ฯเปฯ อสุจิ น มุจฺจิ. ตสฺส กุกฺกุจฺจํ อโหสิ ฯเปฯ. อนาปตฺติ ภิกฺขุ สํฆาทิเสสสฺส, อาปตฺติ ถุลฺลจฺจยสฺสาติ. \mbox{(21-22)}}

\prose{เตน โข ปน สมเยน อญฺญตรสฺส ภิกฺขุโน ชนฺตาฆเร อุปชฺฌายสฺส ปิฏฺฐิ\-ปริกมฺมํ กโรนฺตสฺส อสุจิ มุจฺจิ. ตสฺส กุกฺกุจฺจํ อโหสิ ฯเปฯ. อนาปตฺติ ภิกฺขุ นโมจนา\-ธิปฺปายสฺสาติ. (23)}

\prose{เตน โข ปน สมเยน อญฺญตรสฺส ภิกฺขุโน โมจนา\-ธิปฺปายสฺส ชนฺตาฆเร อุปชฺฌายสฺส ปิฏฺฐิ\-ปริกมฺมํ กโรนฺตสฺส อสุจิ มุจฺจิ ฯเปฯ อสุจิ น มุจฺจิ. ตสฺส กุกฺกุจฺจํ อโหสิ ฯเปฯ. อนาปตฺติ ภิกฺขุ สํฆาทิเสสสฺส, อาปตฺติ ถุลฺลจฺจยสฺสาติ. \mbox{(24-25)}}

\proseitem{265}{เตน โข ปน สมเยน อญฺญตรสฺส ภิกฺขุโน อูรุํ ฆฏฺฏาเปนฺตสฺส อสุจิ มุจฺจิ. ตสฺส กุกฺกุจฺจํ อโหสิ ฯเปฯ. อนาปตฺติ ภิกฺขุ นโมจนา\-ธิปฺปายสฺสาติ. (26)}

\prose{เตน โข ปน สมเยน อญฺญตรสฺส ภิกฺขุโน โมจนา\-ธิปฺปายสฺส อูรุํ ฆฏฺฏาเปนฺตสฺส อสุจิ มุจฺจิ ฯเปฯ อสุจิ น มุจฺจิ. ตสฺส กุกฺกุจฺจํ อโหสิ ฯเปฯ. อนาปตฺติ ภิกฺขุ สํฆาทิเสสสฺส, อาปตฺติ ถุลฺลจฺจยสฺสาติ. \mbox{(27-28)}}

\prose{เตน โข ปน สมเยน อญฺญตโร ภิกฺขุ โมจนาธิปฺปาโย อญฺญตรํ สามเณรํ เอตทโวจ “เอหิ เม ตฺวํ อาวุโส สมเณร องฺคชาตํ คณฺหาหี”ติ. โส ตสฺส องฺคชาตํ อคฺคเหสิ. ตสฺเสว อสุจิ มุจฺจิ. ตสฺส กุกฺกุจฺจํ อโหสิ ฯเปฯ. อาปตฺติํ ตฺวํ ภิกฺขุ อาปนฺโน สํฆาทิเสสนฺติ. (29)}

\prose{เตน โข ปน สมเยน อญฺญตโร ภิกฺขุ สุตฺตสฺส สามเณรสฺส องฺคชาตํ อคฺคเหสิ. ตสฺเสว อสุจิ มุจฺจิ. ตสฺส กุกฺกุจฺจํ อโหสิ ฯเปฯ. อนาปตฺติ ภิกฺขุ สํฆาทิเสสสฺส, อาปตฺติ ทุกฺกฏสฺสาติ. (30)}

\proseitem{266}{เตน โข ปน สมเยน อญฺญตรสฺส ภิกฺขุโน โมจนา\-ธิปฺปายสฺส อูรูหิ องฺคชาตํ ปีเฬนฺตสฺส อสุจิ มุจฺจิ ฯเปฯ อสุจิ น มุจฺจิ. ตสฺส กุกฺกุจฺจํ อโหสิ ฯเปฯ. อนาปตฺติ ภิกฺขุ สํฆาทิเสสสฺส, อาปตฺติ ถุลฺลจฺจยสฺสาติ. \mbox{(31-32)}}

\prose{\csromanfolio{169}เตน โข ปน สมเยน อญฺญตรสฺส ภิกฺขุโน โมจนา\-ธิปฺปายสฺส มุฏฺฐินา องฺคชาตํ ปีเฬนฺตสฺส อสุจิ มุจฺจิ ฯเปฯ อสุจิ น มุจฺจิ. ตสฺส กุกฺกุจฺจํ อโหสิ ฯเปฯ. อนาปตฺติ ภิกฺขุ สํฆาทิเสสสฺส, อาปตฺติ ถุลฺลจฺจยสฺสาติ. \mbox{(33-34)}}

\prose{เตน โข ปน สมเยน อญฺญตรสฺส ภิกฺขุโน โมจนา\-ธิปฺปายสฺส อากาเส กฏิํ กมฺเปนฺตสฺส อสุจิ มุจฺจิ ฯเปฯ อสุจิ น มุจฺจิ. ตสฺส กุกฺกุจฺจํ อโหสิ ฯเปฯ. อนาปตฺติ ภิกฺขุ สํฆาทิเสสสฺส, อาปตฺติ ถุลฺลจฺจยสฺสาติ. \mbox{(35-36)}}

\prose{เตน โข ปน สมเยน อญฺญตรสฺส ภิกฺขุโน กายํ ถมฺเภนฺตสฺส อสุจิ มุจฺจิ. ตสฺส กุกฺกุจฺจํ อโหสิ ฯเปฯ. อนาปตฺติ ภิกฺขุ นโมจนา\-ธิปฺปายสฺสาติ. (37)}

\prose{เตน โข ปน สมเยน อญฺญตรสฺส ภิกฺขุโน โมจนา\-ธิปฺปายสฺส กายํ ถมฺเภนฺตสฺส อสุจิ มุจฺจิ ฯเปฯ อสุจิ น มุจฺจิ. ตสฺส กุกฺกุจฺจํ อโหสิ ฯเปฯ. อนาปตฺติ ภิกฺขุ สํฆาทิเสสสฺส, อาปตฺติ ถุลฺลจฺจยสฺสาติ. \mbox{(38-39)}}

\prose{เตน โข ปน สมเยน อญฺญตโร ภิกฺขุ สารตฺโต มาตุคามสฺส องฺคชาตํ อุปนิชฺฌายิ. ตสฺส อสุจิ มุจฺจิ. ตสฺส กุกฺกุจฺจํ อโหสิ ฯเปฯ. อนาปตฺติ ภิกฺขุ สํฆาทิเสสสฺส. น จ ภิกฺขเว สารตฺเตน มาตุคามสฺส องฺคชาตํ อุปนิชฺฌายิตพฺพํ, โย อุปนิชฺฌาเยยฺย, อาปตฺติ ทุกฺกฏสฺสาติ. (40)}

\proseitem{267}{เตน โข ปน สมเยน อญฺญตรสฺส ภิกฺขุโน โมจนา\-ธิปฺปายสฺส ตาฬจฺฉิทฺทํ องฺคชาตํ ปเวเสนฺตสฺส อสุจิ มุจฺจิ ฯเปฯ อสุจิ น มุจฺจิ. ตสฺส กุกฺกุจฺจํ อโหสิ ฯเปฯ. อนาปตฺติ ภิกฺขุ สํฆาทิเสสสฺส, อาปตฺติ ถุลฺลจฺจยสฺสาติ. \mbox{(41-42)}}

\prose{เตน โข ปน สมเยน อญฺญตรสฺส ภิกฺขุโน โมจนา\-ธิปฺปายสฺส กฏฺเฐน องฺคชาตํ ฆฏฺเฏนฺตสฺส อสุจิ มุจฺจิ ฯเปฯ อสุจิ น มุจฺจิ. ตสฺส กุกฺกุจฺจํ อโหสิ ฯเปฯ. อนาปตฺติ ภิกฺขุ สํฆาทิเสสสฺส, อาปตฺติ ถุลฺลจฺจยสฺสาติ. \mbox{(43-44)}}

\prose{\csromanfolio{170}เตน โข ปน สมเยน อญฺญตรสฺส ภิกฺขุโน ปฏิโสเต นฺหายนฺตสฺส อสุจิ มุจฺจิ. ตสฺส กุกฺกุจฺจํ อโหสิ ฯเปฯ. อนาปตฺติ ภิกฺขุ น โมจนาธิปฺปายสฺสาติ. (45)}

\prose{เตน โข ปน สมเยน อญฺญตรสฺส ภิกฺขุโน โมจนา\-ธิปฺปายสฺส ปฏิโสเต นฺหายนฺตสฺส อสุจิ มุจฺจิ ฯเปฯ อสุจิ น มุจฺจิ. ตสฺส กุกฺกุจฺจํ อโหสิ ฯเปฯ. อนาปตฺติ ภิกฺขุ สํฆาทิเสสสฺส, อาปตฺติ ถุลฺลจฺจยสฺสาติ. \mbox{(46-47)}}

\prose{เตน โข ปน สมเยน อญฺญตรสฺส ภิกฺขุโน อุทญฺชลํ กีฬนฺตสฺส อสุจิ มุจฺจิ. ตสฺส กุกฺกุจฺจํ อโหสิ ฯเปฯ. อนาปตฺติ ภิกฺขุ นโมจนา\-ธิปฺปายสฺสาติ. (48)}

\prose{เตน โข ปน สมเยน อญฺญตรสฺส ภิกฺขุโน โมจนา\-ธิปฺปายสฺส อุทญฺชลํ กีฬนฺตสฺส อสุจิ มุจฺจิ ฯเปฯ อสุจิ น มุจฺจิ. ตสฺส กุกฺกุจฺจํ อโหสิ ฯเปฯ. อนาปตฺติ ภิกฺขุ สํฆาทิเสสสฺส, อาปตฺติ ถุลฺลจฺจยสฺสาติ. \mbox{(49-50)}}

\prose{เตน โข ปน สมเยน อญฺญตรสฺส ภิกฺขุโน อุทเก ธาวนฺตสฺส อสุจิ มุจฺจิ. ตสฺส กุกฺกุจฺจํ อโหสิ ฯเปฯ. อนาปตฺติ ภิกฺขุ นโมจนา\-ธิปฺปายสฺสาติ. (51)}

\prose{เตน โข ปน สมเยน อญฺญตรสฺส ภิกฺขุโน โมจนา\-ธิปฺปายสฺส อุทเก ธาวนฺตสฺส อสุจิ มุจฺจิ ฯเปฯ อสุจิ น มุจฺจิ. ตสฺส กุกฺกุจฺจํ อโหสิ ฯเปฯ. อนาปตฺติ ภิกฺขุ สํฆาทิเสสสฺส, อาปตฺติ ถุลฺลจฺจยสฺสาติ. \mbox{(52-53)}}

\prose{เตน โข ปน สมเยน อญฺญตรสฺส ภิกฺขุโน ปุปฺผาวลิกํ\csromansharedfootnote{043c513eaa0caff8}{ปุปฺผาวฬิยํ (สฺยา, ก)} กีฬนฺตสฺส อสุจิ มุจฺจิ. ตสฺส กุกฺกุจฺจํ อโหสิ ฯเปฯ. อนาปตฺติ ภิกฺขุ นโมจนา\-ธิปฺปายสฺสาติ. (54)}

\prose{เตน โข ปน สมเยน อญฺญตรสฺส ภิกฺขุโน โมจนา\-ธิปฺปายสฺส ปุปฺผาวลิยํ กีฬนฺตสฺส อสุจิ มุจฺจิ ฯเปฯ อสุจิ น มุจฺจิ. ตสฺส กุกฺกุจฺจํ อโหสิ ฯเปฯ. อนาปตฺติ ภิกฺขุ สํฆาทิเสสสฺส, อาปตฺติ ถุลฺลจฺจยสฺสาติ. \mbox{(55-56)}}

\proseitem{268}{\csromanfolio{171}เตน โข ปน สมเยน อญฺญตรสฺส ภิกฺขุโน โปกฺขรวเน ธาวนฺตสฺส อสุจิ มุจฺจิ. ตสฺส กุกฺกุจฺจํ อโหสิ ฯเปฯ. อนาปตฺติ ภิกฺขุ นโมจนา\-ธิปฺปายสฺสาติ. (57)}

\prose{เตน โข ปน สมเยน อญฺญตรสฺส ภิกฺขุโน โมจนา\-ธิปฺปายสฺส โปกฺขรวเน ธาวนฺตสฺส อสุจิ มุจฺจิ ฯเปฯ อสุจิ น มุจฺจิ. ตสฺส กุกฺกุจฺจํ อโหสิ ฯเปฯ. อนาปตฺติ ภิกฺขุ สํฆาทิเสสสฺส, อาปตฺติ ถุลฺลจฺจยสฺสาติ. \mbox{(58-59)}}

\prose{เตน โข ปน สมเยน อญฺญตรสฺส ภิกฺขุโน โมจนา\-ธิปฺปายสฺส วาลิกํ องฺคชาตํ ปเวเสนฺตสฺส อสุจิ มุจฺจิ ฯเปฯ อสุจิ น มุจฺจิ. ตสฺส กุกฺกุจฺจํ อโหสิ ฯเปฯ. อนาปตฺติ ภิกฺขุ สํฆาทิเสสสฺส, อาปตฺติ ถุลฺลจฺจยสฺสาติ. \mbox{(60-61)}}

\prose{เตน โข ปน สมเยน อญฺญตรสฺส ภิกฺขุโน โมจนา\-ธิปฺปายสฺส กทฺทมํ องฺคชาตํ ปเวเสนฺตสฺส อสุจิ มุจฺจิ ฯเปฯ อสุจิ น มุจฺจิ. ตสฺส กุกฺกุจฺจํ อโหสิ ฯเปฯ. อนาปตฺติ ภิกฺขุ สํฆาทิเสสสฺส, อาปตฺติ ถุลฺลจฺจยสฺสาติ. \mbox{(62-63)}}

\prose{เตน โข ปน สมเยน อญฺญตรสฺส ภิกฺขุโน อุทเกน องฺคชาตํ โอสิญฺจนฺตสฺส อสุจิ มุจฺจิ. ตสฺส กุกฺกุจฺจํ อโหสิ ฯเปฯ. อนาปตฺติ ภิกฺขุ นโมจนา\-ธิปฺปายสฺสาติ. (64)}

\prose{เตน โข ปน สมเยน อญฺญตรสฺส ภิกฺขุโน โมจนา\-ธิปฺปายสฺส อุทเกน องฺคชาตํ โอสิญฺจนฺตสฺส อสุจิ มุจฺจิ ฯเปฯ อสุจิ น มุจฺจิ. ตสฺส กุกฺกุจฺจํ อโหสิ ฯเปฯ. อนาปตฺติ ภิกฺขุ สํฆาทิเสสสฺส, อาปตฺติ ถุลฺลจฺจยสฺสาติ. \mbox{(65-66)}}

\prose{เตน โข ปน สมเยน อญฺญตรสฺส ภิกฺขุโน โมจนา\-ธิปฺปายสฺส สยเน องฺคชาตํ ฆฏฺเฏนฺตสฺส อสุจิ มุจฺจิ ฯเปฯ อสุจิ น มุจฺจิ. ตสฺส กุกฺกุจฺจํ อโหสิ ฯเปฯ. อนาปตฺติ ภิกฺขุ สํฆาทิเสสสฺส, อาปตฺติ ถุลฺลจฺจยสฺสาติ. \mbox{(67-68)}}

\csromanmatikaanchor{mtk.43}
\csromanmatikaanchor{mtk.44}
\csromantocmarknopage{section}{2. กายสํสคฺคสิกฺขาปท}{mtk.43}
\bookmark[dest={mtk.43},level=1]{2. กายสํสคฺคสิกฺขาปท}
\csromantocmarkref{subsection}{อุทายิภิกฺขุวตฺถุ}{mtk.44}
\bookmark[dest={mtk.44},level=2]{อุทายิภิกฺขุวตฺถุ}
\prosewithcloserruled{เตน โข ปน สมเยน อญฺญตรสฺส ภิกฺขุโน โมจนา\-ธิปฺปายสฺส องฺคุฏฺเฐน องฺคชาตํ ฆฏฺเฏนฺตสฺส อสุจิ มุจฺจิ ฯเปฯ อสุจิ น มุจฺจิ. ตสฺส กุกฺกุจฺจํ \csromanfolio{172}อโหสิ “ภควตา สิกฺขาปทํ ปญฺญตฺตํ, กจฺจิ นุ โข อหํ สํฆาทิเสสํ อาปตฺติํ อาปนฺโน”ติ. ภควโต เอตมตฺตํ อาโรเจสิ. อนาปตฺติ ภิกฺขุ สํฆาทิเสสสฺส, อาปตฺติ ถุลฺลจฺจยสฺสาติ. \mbox{(69-70)}}{สุกฺกวิสฺสฏฺฐิ\-สิกฺขาปทํ นิฏฺฐิตํ ปฐมํ.}
\csromanheader{1}{2. \textbf{กายสํสคฺค}\-\textbf{สิกฺขาปท}}
\proseitem{269}{เตน สมเยน พุทฺโธ ภควา สาวตฺถิยํ วิหรติ เชตวเน อนาถ\-ปิณฺฑิกสฺส อาราเม. เตน โข ปน สมเยน อายสฺมา อุทายี อรญฺเญ วิหรติ, ตสฺสายสฺมโต วิหาโร อภิรูโป โหติ ทสฺสนีโย ปาสาทิโก มชฺเฌคพฺโภ สมนฺตา\-ปริยาคาโร, สุปญฺญตฺตํ มญฺจปีฐํ ภิสิพิพฺโพหนํ, ปานียํ ปริโภชนียํ สูปฏฺฐิตํ, ปริเวณํ สุสมฺมฏฺฐํ. พหู มนุสฺสา อายสฺมโต อุทายิสฺส วิหารเปกฺขกา อาคจฺฉนฺติ. อญฺญตโรปิ พฺราหฺมโณ สปชาปติโก เยนายสฺมา อุทายี เตนุปสงฺกมิ, อุปสงฺกมิตฺวา อายสฺมนฺตํ อุทายิํ เอตทโวจ “อิจฺฉาม มยํ โภโต อุทายิสฺส วิหารํ เปกฺขิตุนฺ”ติ. เตน หิ พฺราหฺมณ เปกฺขสฺสูติ อวาปุรณํ\csromansharedfootnote{f8c093e53ab6fc20}{อปาปุรณํ (สฺยา)} อาทาย ฆฏิกํ อุคฺฆาเฏตฺวา กวาฏํ ปณาเมตฺวา วิหารํ ปาวิสิ, โสปิ โข พฺราหฺมโณ อายสฺมโต อุทายิสฺส ปิฏฺฐิโต ปาวิสิ, สาปิ โข พฺราหฺมณี ตสฺส พฺราหฺมณสฺส ปิฏฺฐิโต ปาวิสิ. อถ โข อายสฺมา อุทายี เอกจฺเจ วาตปาเน วิวรนฺโต เอกจฺเจ วาตปาเน ถเกนฺโต คพฺภํ อนุปริคนฺตฺวา ปิฏฺฐิโต อาคนฺตฺวา ตสฺสา พฺราหฺมณิยา องฺคมงฺคานิ ปรามสิ. อถ โข โส พฺราหฺมโณ อายสฺมตา อุทายินา สทฺธิํ ปฏิสมฺโมทิตฺวา อคมาสิ. อถ โข โส พฺราหฺมโณ อตฺตมโน อตฺตมนวาจํ นิจฺฉาเรสิ “อุฬารา อิเม สมณา สกฺยปุตฺติยา, เย อิเม เอวรูเป อรญฺเญ วิหรนฺติ, ภวํปิ อุทายี อุฬาโร, โย เอวรูเป อรญฺเญ วิหรตี”ติ.}

\csromanmatikaanchor{mtk.45}
\csromanmatikaanchor{mtk.46}
\csromantocmarkref{subsection}{ปญฺญตฺติ}{mtk.45}
\bookmark[dest={mtk.45},level=2]{ปญฺญตฺติ}
\csromantocmarkref{subsection}{สิกฺขาปทวิภงฺค, ปทภาชนีย}{mtk.46}
\bookmark[dest={mtk.46},level=2]{สิกฺขาปทวิภงฺค, ปทภาชนีย}
\prose{เอวํ วุตฺเต สา พฺราหฺมณี ตํ พฺราหฺมณํ เอตทโวจ “กุโต ตสฺส อุฬารตฺตตา, ยเถว เม ตฺวํ องฺคมงฺคานิ ปรามสิ, เอวเมว เม สมโณ อุทายี องฺคมงฺคานิ ปรามสี”ติ. อถ โข โส พฺราหฺมโณ อุชฺฌายติ ขิยฺยติ วิปาเจติ “อลชฺชิโน อิเม สมณา สกฺยปุตฺติยา ทุสฺสีลา มุสาวาทิโน, อิเม หิ นาม ธมฺมจาริโน สมจาริโน พฺรหฺมจาริโน \csromanfolio{173}สจฺจวาทิโน สีลวนฺโต กลฺยาณธมฺมา ปฏิชานิสฺสนฺติ. นตฺถิ อิเมสํ สามญฺญํ, นตฺถิ อิเมสํ พฺรหฺมญฺญํ, นฏฺฐํ อิเมสํ สามญฺญํ, นฏฺฐํ อิเมสํ พฺรหฺมญฺญํ, กุโต อิเมสํ สามญฺญํ, กุโต อิเมสํ พฺรหฺมญฺญํ, อปคตา อิเม สามญฺญา, อปคตา อิเม พฺรหฺมญฺญา. กถํ หิ นาม สมโณ อุทายี มม ภริยาย องฺคมงฺคานิ ปรามสิสฺสติ, น หิ สกฺกา กุลิตฺถีหิ กุลธีตาหิ กุลกุมารีหิ กุลสุณฺหาหิ กุลทาสีหิ อารามํ วา วิหารํ วา คนฺตุํ, สเจ\csromansharedfootnote{5d56274247fa0e72}{สเจ หิ (สฺยา)} กุลิตฺถิโย กุลธีตโร\csromansharedfootnote{bb0886a8927dd2d0}{กุลธีตาโย (สี, สฺยา)} กุลกุมาริโย กุลสุณฺหาโย กุลทาสิโย อารามํ วา วิหารํ วา คจฺเฉยฺยุํ, ตาปิ สมณา \textbf{สกฺยปุตฺติยา ทูเสยฺยุนฺ”ติ.}}

\prose{อสฺโสสุํ โข ภิกฺขู ตสฺส พฺราหฺมณสฺส อุชฺฌยนฺตสฺส ขิยฺยนฺตสฺส วิปาเจนฺตสฺส, เย เต ภิกฺขู อปฺปิจฺฉา ฯเปฯ เต อุชฺฌายนฺติ ขิยฺยนฺติ วิปาเจนฺติ “กถํ หิ นาม อายสฺมา อุทายี มาตุคาเมน สทฺธิํ กายสํสคฺคํ สมาปชฺชิสฺสตี”ติ. อถ โข เต ภิกฺขู อายสฺมนฺตํ อุทายิํ อเนก\-ปริยาเยน วิครหิตฺวา ภควโต เอตมตฺถํ อาโรเจสุํ. อถ โข ภควา เอตสฺมิํ นิทาเน เอตสฺมิํ ปกรเณ ภิกฺขุสํฆํ สนฺนิปาตาเปตฺวา อายสฺมนฺตํ อุทายิํ ปฏิปุจฺฉิ “สจฺจํ กิร ตฺวํ อุทายิ มาตุคาเมน สทฺธิํ กายสํสคฺคํ สมาปชฺชสี”ติ. สจฺจํ ภควาติ. วิครหิ พุทฺโธ ภควา “อนนุจฺฉวิกํ โมฆปุริส อนนุโลมิกํ อปฺปติรูปํ อสฺสามณกํ อกปฺปิยํ อกรณียํ, กถํ หิ นาม ตฺวํ โมฆปุริส มาตุคาเมน สทฺธิํ กายสํสคฺคํ สมาปชฺชิสฺสสิ. นนุ มยา โมฆปุริส อเนก\-ปริยาเยน วิราคาย ธมฺโม เทสิโต, โน สราคาย ฯเปฯ กาม\-ปริฬาหานํ วูปสโม อกฺขาโต. เนตํ โมฆปุริส อปฺปสนฺนานํ วา ปสาทาย ฯเปฯ. เอวญฺจ ปน ภิกฺขเว อิมํ สิกฺขาปทํ อุทฺทิเสยฺยาถ\csromandash{}}

\proseitem{270}{\textbf{“โย ปน ภิกฺขุ โอติณฺโณ วิปริณเตน จิตฺเตน มาตุคาเมน สทฺธิํ กายสํสคฺคํ สมาปชฺเชยฺย หตฺถคฺคาหํ วา เวณิคฺคาหํ วา อญฺญตรสฺส วา อญฺญตรสฺส วา องฺคสฺส ปรามสนํ, สํฆาทิเสโส”ติ}.}

\proseitem{271}{\textbf{โย ปนา}ติ โย ยาทิโส ฯเปฯ. \textbf{ภิกฺขู}ติ ฯเปฯ. อยํ อิมสฺมิํ อตฺเถ อธิปฺเปโต “ภิกฺขู”ติ.}

\prose{\csromanfolio{174}\textbf{โอติณฺโณ} นาม สารตฺโต อเปกฺขวา ปฏิพทฺธ\-จิตฺโต.}

\prose{\textbf{วิปริณตนฺ}ติ รตฺตํปิ จิตฺตํ วิปริณตํ, ทุฏฺฐํปิ จิตฺตํ วิปริณตํ, มูฬฺหํปิ จิตฺตํ วิปริณตํ. อปิ จ รตฺตํ จิตฺตํ อิมสฺมิํ อตฺเถ อธิปฺเปตํ “วิปริณตนฺ”ติ.}

\prose{\textbf{มาตุคาโม} นาม มนุสฺสิตฺถี, น ยกฺขี, น เปตี, น ติรจฺฉานคตา, อนฺตมโส ตทหุชาตาปิ ทาริกา, ปเคว มหตฺตรี.}

\prose{\textbf{สทฺธินฺ}ติ เอกโต.}

\prose{\textbf{กายสํสคฺคํ สมาปชฺเชยฺยา}ติ อชฺฌาจาโร วุจฺจติ.}

\prose{\textbf{หตฺโถ} นาม กปฺปรํ อุปาทาย ยาว อคฺคนขา.}

\prose{\textbf{เวณี} นาม สุทฺธเกสา วา สุตฺตมิสฺสา วา มาลามิสฺสา วา หิรญฺญมิสฺสา วา สุวณฺณมิสฺสา วา มุตฺตามิสฺสา วา มณิมิสฺสา วา.}

\prose{\textbf{องฺคํ} นาม หตฺถญฺจ เวณิญฺจ ฐเปตฺวา อวเสสํ องฺคํ นาม.}

\proseitem{272}{อามสนา, ปรามสนา, โอมสนา, อุมฺมสนา, โอลงฺฆนา, อุลฺลงฺฆนา, อากฑฺฒนา, ปติกฑฺฒนา, อภินิคฺคณฺหนา, อภินิปฺปีฬนา, คหณํ, ฉุปนํ.}

\prose{\textbf{อามสนา} นาม อามฏฺฐมตฺตา.}

\prose{\textbf{ปรามสนา} นาม อิโตจิโต จ สํโจปนา.}

\prose{\textbf{โอมสนา} นาม เหฏฺฐา โอโรปนา.}

\prose{\textbf{อุมฺมสนา} นาม อุทฺธํ อุจฺจารณา.}

\prose{\textbf{โอลงฺฆนา} นาม เหฏฺฐา โอนมนา.}

\prose{\textbf{อุลฺลงฺฆนา} นาม อุทฺธํ อุจฺจารณา.}

\prose{\textbf{อากฑฺฒนา} นาม อาวิญฺฉนา\csromansharedfootnote{55f80e3faeb122fe}{อาวิญฺชนา (สี, สฺยา)}.}

\prose{\textbf{ปติกฑฺฒนา} นาม ปติปฺปณามนา.}

\prose{\textbf{อภินิคฺคณฺหนา} นาม องฺคํ คเหตฺวา นิปฺปีฬนา.}

\prose{\textbf{อภินิปฺปีฬนา} นาม เกนจิ สห นิปฺปีฬนา.}

\prose{\textbf{คหณํ} นาม คหิตมตฺตํ.}

\prose{\textbf{ฉุปนํ} นาม ผุฏฺฐมตฺตํ.}

\prose{\textbf{สํฆาทิเสโส}ติ ฯเปฯ เตนปิ วุจฺจติ “สํฆาทิเสโส”ติ.}

\proseitem{273}{\csromanfolio{175}อิตฺถี จ โหติ อิตฺถิสญฺญี สารตฺโต จ, ภิกฺขุ จ นํ อิตฺถิยา กาเยน กายํ อามสติ ปรามสติ โอมสติ อุมฺมสติ โอลงฺเฆติ อุลฺลงฺเฆติ อากฑฺฒติ ปติกฑฺฒติ อภินิคฺคณฺหาติ อภินิปฺปีเฬติ คณฺหาติ ฉุปติ, อาปตฺติ สํฆาทิเสสสฺส.}

\prose{อิตฺถี จ โหติ เวมติโก สารตฺโต จ, ภิกฺขุ จ นํ อิตฺถิยา กาเยน กายํ อามสติ ปรามสติ ฯเปฯ คณฺหาติ ฉุปติ, อาปตฺติ ถุลฺลจฺจยสฺส.}

\prose{อิตฺถี จ โหติ ปณฺฑกสญฺญี สารตฺโต จ, ภิกฺขุ จ นํ อิตฺถิยา กาเยน กายํ อามสติ ปรามสติ ฯเปฯ คณฺหาติ ฉุปติ, อาปตฺติ ถุลฺลจฺจยสฺส.}

\prose{อิตฺถี จ โหติ ปุริสสญฺญี สารตฺโต จ, ภิกฺขุ จ นํ อิตฺถิยา กาเยน กายํ อามสติ ปรามสติ ฯเปฯ คณฺหาติ ฉุปติ, อาปตฺติ ถุลฺลจฺจยสฺส.}

\prose{อิตฺถี จ โหติ ติรจฺฉานคต\-สญฺญี สารตฺโต จ, ภิกฺขุ จ นํ อิตฺถิยา กาเยน กายํ อามสติ ปรามสติ ฯเปฯ คณฺหาติ ฉุปติ, อาปตฺติ ถุลฺลจฺจยสฺส.}

\prose{ปณฺฑโก จ โหติ ปณฺฑกสญฺญี สารตฺโต จ, ภิกฺขุ จ นํ ปณฺฑกสฺส กาเยน กายํ อามสติ ปรามสติ ฯเปฯ คณฺหาติ ฉุปติ, อาปตฺติ ถุลฺลจฺจยสฺส.}

\prose{ปณฺฑโก จ โหติ เวมติโก สารตฺโต จ, ภิกฺขุ จ นํ ปณฺฑกสฺส กาเยน กายํ อามสติ ปรามสติ ฯเปฯ คณฺหาติ ฉุปติ, อาปตฺติ ทุกฺกฏสฺส.}

\prose{ปณฺฑโก จ โหติ ปุริสสญฺญี สารตฺโต จ, ภิกฺขุ จ นํ ปณฺฑกสฺส กาเยน กายํ อามสติ ปรามสติ ฯเปฯ คณฺหาติ ฉุปติ, อาปตฺติ ทุกฺกฏสฺส.}

\prose{ปณฺฑโก จ โหติ ติรจฺฉานคต\-สญฺญี สารตฺโต จ, ภิกฺขุ จ นํ ปณฺฑกสฺส กาเยน กายํ อามสติ ปรามสติ ฯเปฯ คณฺหาติ ฉุปติ, อาปตฺติ ทุกฺกฏสฺส.}

\prose{\csromanfolio{176}ปณฺฑโก จ โหติ อิตฺถิสญฺญี สารตฺโต จ, ภิกฺขุ จ นํ ปณฺฑกสฺส กาเยน กายํ อามสติ ปรามสติ ฯเปฯ คณฺหาติ ฉุปติ, อาปตฺติ ทุกฺกฏสฺส.}

\prose{ปุริโส จ โหติ ปุริสสญฺญี สารตฺโต จ, ภิกฺขุ จ นํ ปุริสสฺส กาเยน กายํ อามสติ ปรามสติ ฯเปฯ คณฺหาติ ฉุปติ, อาปตฺติ ทุกฺกฏสฺส.}

\prose{ปุริโส จ โหติ เวมติโก ฯเปฯ. ปุริโส จ โหติ ติรจฺฉานคต\-สญฺญี. ปุริโส จ โหติ อิตฺถิสญฺญี. ปุริโส จ โหติ ปณฺฑกสญฺญี สารตฺโต จ, ภิกฺขุ จ นํ ปุริสสฺส กาเยน กายํ อามสติ ปรามสติ ฯเปฯ คณฺหาติ ฉุปติ, อาปตฺติ ทุกฺกฏสฺส.}

\prose{ติรจฺฉานคโต จ โหติ ติรจฺฉานคต\-สญฺญี สารตฺโต จ, ภิกฺขุ จ นํ ติรจฺฉาน\-คตสฺส กาเยน กายํ อามสติ ปรามสติ ฯเปฯ คณฺหาติ ฉุปติ, อาปตฺติ ทุกฺกฏสฺส.}

\prosewithcloser{ติรจฺฉานคโต จ โหติ เวมติโก ฯเปฯ. ติรจฺฉานคโต จ โหติ อิตฺถิสญฺญี. ติรจฺฉานคโต จ โหติ ปณฺฑกสญฺญี. ติรจฺฉานคโต จ โหติ ปุริสสญฺญี สารตฺโต จ, ภิกฺขุ จ นํ ติรจฺฉาน\-คตสฺส กาเยน กายํ อามสติ ปรามสติ ฯเปฯ คณฺหาติ ฉุปติ, อาปตฺติ ทุกฺกฏสฺส.}{เอกมูลกํ นิฏฺฐิตํ.}
\proseitem{274}{เทฺว อิตฺถิโย ทฺวินฺนํ อิตฺถีนํ อิตฺถิสญฺญี สารตฺโต จ, ภิกฺขุ จ นํ ทฺวินฺนํ อิตฺถีนํ กาเยน กายํ อามสติ ปรามสติ ฯเปฯ คณฺหาติ ฉุปติ, อาปตฺติ ทฺวินฺนํ สํฆา\-ทิ\-เสสานํ.}

\prose{เทฺว อิตฺถิโย ทฺวินฺนํ อิตฺถีนํ เวมติโก สารตฺโต จ, ภิกฺขุ จ นํ ทฺวินฺนํ อิตฺถีนํ กาเยน กายํ อามสติ ปรามสติ ฯเปฯ คณฺหาติ ฉุปติ, อาปตฺติ ทฺวินฺนํ ถุลฺลจฺจยานํ.}

\prose{เทฺว อิตฺถิโย ทฺวินฺนํ อิตฺถีนํ ปณฺฑกสญฺญี ฯเปฯ ปุริสสญฺญี. ติรจฺฉานคต\-สญฺญี สารตฺโต จ, ภิกฺขุ จ นํ ทฺวินฺนํ อิตฺถีนํ กาเยน กายํ อามสติ ปรามสติ ฯเปฯ คณฺหาติ ฉุปติ, อาปตฺติ ทฺวินฺนํ ถุลฺลจฺจยานํ.}

\prose{\csromanfolio{177}เทฺว ปณฺฑกา ทฺวินฺนํ ปณฺฑกานํ ปณฺฑกสญฺญิ สารตฺโต จ, ภิกฺขุ จ นํ ทฺวินฺนํ ปณฺฑกานํ กาเยน กายํ อามสติ ปรามสติ ฯเปฯ คณฺหาติ ฉุปติ, อาปตฺติ ทฺวินฺนํ ถุลฺลจฺจยานํ.}

\prose{เทฺว ปณฺฑกา ทฺวินฺนํ ปณฺฑกานํ เวมติโก ฯเปฯ ปุริสสญฺญี. ติรจฺฉานคต\-สญฺญี. อิตฺถิสญฺญี สารตฺโต จ, ภิกฺขุ จ นํ ทฺวินฺนํ ปณฺฑกานํ กาเยน กายํ อามสติ ปรามสติ ฯเปฯ คณฺหาติ ฉุปติ, อาปตฺติ ทฺวินฺนํ ทุกฺกฏานํ.}

\prose{เทฺว ปุริสา ทฺวินฺนํ ปุริสานํ ปุริสสญฺญี สารตฺโต จ, ภิกฺขุ จ นํ ทฺวินฺนํ ปุริสานํ กาเยน กายํ อามสติ ปรามสติ ฯเปฯ คณฺหาติ ฉุปติ, อาปตฺติ ทฺวินฺนํ ทุกฺกฏานํ.}

\prose{เทฺว ปุริสา ทฺวินฺนํ ปุริสานํ เวมติโก ฯเปฯ ติรจฺฉานคต\-สญฺญี. อิตฺถิสญฺญี. ปณฺฑกสญฺญี สารตฺโต จ, ภิกฺขุ จ นํ ทฺวินฺนํ ปุริสานํ กาเยน กายํ อามสติ ปรามสติ ฯเปฯ คณฺหาติ ฉุปติ, อาปตฺติ ทฺวินฺนํ ทุกฺกฏานํ.}

\prose{เทฺว ติรจฺฉานคตา ทฺวินฺนํ ติรจฺฉาน\-คตานํ ติรจฺฉานคต\-สญฺญี สารตฺโต จ, ภิกฺขุ จ นํ ทฺวินฺนํ ติรจฺฉาน\-คตานํ กาเยน กายํ อามสติ ปรามสติ ฯเปฯ คณฺหาติ ฉุปติ, อาปตฺติ ทฺวินฺนํ ทุกฺกฏานํ.}

\prose{เทฺว ติรจฺฉานคตา ทฺวินฺนํ ติรจฺฉาน\-คตานํ เวมติโก ฯเปฯ อิตฺถิสญฺญี. ปณฺฑกสญฺญี. ปุริสสญฺญี สารตฺโต จ, ภิกฺขุ จ นํ ทฺวินฺนํ ติรจฺฉาน\-คตานํ กาเยน กายํ อามสติ ปรามสติ ฯเปฯ คณฺหาติ ฉุปติ, อาปตฺติ ทฺวินฺนํ ทุกฺกฏานํ.}

\proseitem{275}{อิตฺถี จ ปณฺฑโก จ อุภินฺนํ อิตฺถิสญฺญี สารตฺโต จ, ภิกฺขุ จ นํ อุภินฺนํ กาเยน กายํ อามสติ ปรามสติ ฯเปฯ คณฺหาติ ฉุปติ, อาปตฺติ สํฆาทิเสเสน ทุกฺกฏสฺส.}

\prose{อิตฺถี จ ปณฺฑโก จ อุภินฺนํ เวมติโก สารตฺโต จ, ภิกฺขุ จ นํ อุภินฺนํ กาเยน กายํ อามสติ ปรามสติ ฯเปฯ คณฺหาติ ฉุปติ, อาปตฺติ ถุลฺลจฺจเยน ทุกฺกฏสฺส.}

\prose{อิตฺถี จ ปณฺฑโก จ อุภินฺนํ ปณฺฑกสญฺญี สารตฺโต จ, ภิกฺขุ จ นํ อุภินฺนํ กาเยน กายํ อามสติ ปรามสติ ฯเปฯ คณฺหาติ ฉุปติ, อาปตฺติ ทฺวินฺนํ ถุลฺลจฺจยานํ.}

\prose{\csromanfolio{178}อิตฺถี จ ปณฺฑโก จ อุภินฺนํ ปุริสสญฺญี สารตฺโต จ, ภิกฺขุ จ นํ อุภินฺนํ กาเยน กายํ อามสติ ปรามสติ ฯเปฯ คณฺหาติ ฉุปติ, อาปตฺติ ถุลฺลจฺจเยน ทุกฺกฏสฺส.}

\prose{อิตฺถี จ ปณฺฑโก จ อุภินฺนํ ติรจฺฉานคต\-สญฺญี สารตฺโต จ, ภิกฺขุ จ นํ อุภินฺนํ กาเยน กายํ อามสติ ปรามสติ ฯเปฯ คณฺหาติ ฉุปติ, อาปตฺติ ถุลฺลจฺจเยน ทุกฺกฏสฺส.}

\prose{อิตฺถี จ ปุริโส จ อุภินฺนํ อิตฺถิสญฺญี สารตฺโต จ, ภิกฺขุ จ นํ อุภินฺนํ กาเยน กายํ อามสติ ปรามสติ ฯเปฯ คณฺหาติ ฉุปติ, อาปตฺติ สํฆาทิเสเสน ทุกฺกฏสฺส.}

\prose{อิตฺถี จ ปุริโส จ อุภินฺนํ เวมติโก ฯเปฯ ปณฺฑกสญฺญี. ปุริสสญฺญี. ติรจฺฉานคต\-สญฺญี สารตฺโต จ, ภิกฺขุ จ นํ อุภินฺนํ กาเยน กายํ อามสติ ปรามสติ ฯเปฯ คณฺหาติ ฉุปติ, อาปตฺติ ถุลฺลจฺจเยน ทุกฺกฏสฺส.}

\prose{อิตฺถี จ ติรจฺฉานคโต จ อุภินฺนํ อิตฺถิสญฺญี สารตฺโต จ, ภิกฺขุ จ นํ อุภินฺนํ กาเยน กายํ อามสติ ปรามสติ ฯเปฯ คณฺหาติ ฉุปติ, อาปตฺติ สํฆาทิเสเสน ทุกฺกฏสฺส.}

\prose{อิตฺถี จ ติรจฺฉานคโต จ อุภินฺนํ เวมติโก ฯเปฯ ปณฺฑกสญฺญี. ปุริสสญฺญี. ติรจฺฉานคต\-สญฺญี สารตฺโต จ, ภิกฺขุ จ นํ อุภินฺนํ กาเยน กายํ อามสติ ปรามสติ ฯเปฯ คณฺหาติ ฉุปติ, อาปตฺติ ถุลฺลจฺจเยน ทุกฺกฏสฺส.}

\prose{ปณฺฑโก จ ปุริโส จ อุภินฺนํ ปณฺฑกสญฺญี สารตฺโต จ, ภิกฺขุ จ นํ อุภินฺนํ กาเยน กายํ อามสติ ปรามสติ ฯเปฯ คณฺหาติ ฉุปติ, อาปตฺติ ถุลฺลจฺจเยน ทุกฺกฏสฺส.}

\prose{ปณฺฑโก จ ปุริโส จ อุภินฺนํ เวมติโก ฯเปฯ ปุริสสญฺญี. ติรจฺฉานคต\-สญฺญี. อิตฺถิสญฺญี สารตฺโต จ, ภิกฺขุ จ นํ อุภินฺนํ กาเยน กายํ อามสติ ปรามสติ ฯเปฯ คณฺหาติ ฉุปติ, อาปตฺติ ทฺวินฺนํ ทุกฺกฏานํ.}

\prose{ปณฺฑโก จ ติรจฺฉานคโต จ อุภินฺนํ ปณฺฑกสญฺญี สารตฺโต จ, ภิกฺขุ จ นํ อุภินฺนํ กาเยน กายํ อามสติ ปรามสติ ฯเปฯ คณฺหาติ ฉุปติ, อาปตฺติ ถุลฺลจฺจเยน ทุกฺกฏสฺส.}

\prose{\csromanfolio{179}ปณฺฑโก จ ติรจฺฉานคโต จ อุภินฺนํ เวมติโก ฯเปฯ ปุริสสญฺญี. ติรจฺฉานคต\-สญฺญี. อิตฺถิสญฺญี สารตฺโต จ, ภิกฺขุ จ นํ อุภินฺนํ กาเยน กายํ อามสติ ปรามสติ ฯเปฯ คณฺหาติ ฉุปติ, อาปตฺติ ทฺวินฺนํ ทุกฺกฏานํ.}

\prose{ปุริโส จ ติรจฺฉานคโต จ อุภินฺนํ ปุริสสญฺญี สารตฺโต จ, ภิกฺขุ จ นํ อุภินฺนํ กาเยน กายํ อามสติ ปรามสติ ฯเปฯ คณฺหาติ ฉุปติ, อาปตฺติ ทฺวินฺนํ ทุกฺกฏานํ.}

\prosewithcloser{ปุริโส จ ติรจฺฉานคโต จ อุภินฺนํ เวมติโก ฯเปฯ ติรจฺฉานคต\-สญฺญี. อิตฺถิสญฺญี. ปณฺฑกสญฺญี สารตฺโต จ, ภิกฺขุ จ นํ อุภินฺนํ กาเยน กายํ อามสติ ปรามสติ ฯเปฯ คณฺหาติ ฉุปติ, อาปตฺติ ทฺวินฺนํ ทุกฺกฏานํ.}{ทุมูลกํ นิฏฺฐิตํ.}
\proseitem{276}{อิตฺถี จ โหติ อิตฺถิสญฺญี สารตฺโต จ, ภิกฺขุ จ นํ อิตฺถิยา กาเยน กาย\-ปฏิพทฺธํ อามสติ ปรามสติ ฯเปฯ คณฺหาติ ฉุปติ, อาปตฺติ ถุลฺลจฺจยสฺส ฯเปฯ.}

\prose{เทฺว อิตฺถิโย ทฺวินฺนํ อิตฺถีนํ อิตฺถิสญฺญี สารตฺโต จ, ภิกฺขุ จ นํ ทฺวินฺนํ อิตฺถีนํ กาเยน กาย\-ปฏิพทฺธํ อามสติ ปรามสติ ฯเปฯ คณฺหาติ ฉุปติ, อาปตฺติ ทฺวินฺนํ ถุลฺลจฺจยานํ ฯเปฯ.}

\prose{อิตฺถี จ ปณฺฑโก จ อุภินฺนํ อิตฺถิสญฺญี สารตฺโต จ, ภิกฺขุ จ นํ อุภินฺนํ กาเยน กาย\-ปฏิพทฺธํ อามสติ ปรามสติ ฯเปฯ คณฺหาติ ฉุปติ, อาปตฺติ ถุลฺลจฺจเยน ทุกฺกฏสฺส ฯเปฯ.}

\prose{อิตฺถี จ โหติ อิตฺถิสญฺญี สารตฺโต จ, ภิกฺขุ จ นํ อิตฺถิยา กาย\-ปฏิพทฺเธน กายํ อามสติ ปรามสติ ฯเปฯ คณฺหาติ ฉุปติ, อาปตฺติ ถุลฺลจฺจยสฺส ฯเปฯ.}

\prose{เทฺว อิตฺถิโย ทฺวินฺนํ อิตฺถีนํ อิตฺถิสญฺญี สารตฺโต จ, ภิกฺขุ จ นํ ทฺวินฺนํ อิตฺถีนํ กาย\-ปฏิพทฺเธน กายํ อามสติ ปรามสติ ฯเปฯ คณฺหาติ ฉุปติ, อาปตฺติ ทฺวินฺนํ ถุลฺลจฺจยานํ ฯเปฯ.}

\prose{อิตฺถี จ ปณฺฑโก จ อุภินฺนํ อิตฺถิสญฺญี สารตฺโต จ, ภิกฺขุ จ นํ อุภินฺนํ กาย\-ปฏิพทฺเธน กายํ อามสติ ปรามสติ ฯเปฯ คณฺหาติ ฉุปติ, อาปตฺติ ถุลฺลจฺจเยน ทุกฺกฏสฺส ฯเปฯ.}

\prose{\csromanfolio{180}อิตฺถี จ โหติ อิตฺถิสญฺญี สารตฺโต จ, ภิกฺขุ จ นํ อิตฺถิยา กาย\-ปฏิพทฺเธน กาย\-ปฏิพทฺธํ อามสติ ปรามสติ ฯเปฯ คณฺหาติ ฉุปติ, อาปตฺติ ทุกฺกฏสฺส ฯเปฯ.}

\prose{เทฺว อิตฺถิโย ทฺวินฺนํ อิตฺถีนํ อิตฺถิสญฺญี สารตฺโต จ, ภิกฺขุ จ นํ ทฺวินฺนํ อิตฺถีนํ กาย\-ปฏิพทฺเธน กาย\-ปฏิพทฺธํ อามสติ ปรามสติ ฯเปฯ คณฺหาติ ฉุปติ, อาปตฺติ ทฺวินฺนํ ทุกฺกฏานํ ฯเปฯ.}

\prose{อิตฺถี จ ปณฺฑโก จ อุภินฺนํ อิตฺถิสญฺญี สารตฺโต จ, ภิกฺขุ จ นํ อุภินฺนํ กาย\-ปฏิพทฺเธน กาย\-ปฏิพทฺธํ อามสติ ปรามสติ ฯเปฯ คณฺหาติ ฉุปติ, อาปตฺติ ทฺวินฺนํ ทุกฺกฏานํ ฯเปฯ.}

\prose{อิตฺถี จ โหติ อิตฺถิสญฺญี สารตฺโต จ, ภิกฺขุ จ นํ อิตฺถิยา นิสฺสคฺคิเยน กายํ อามสติ, อาปตฺติ ทุกฺกฏสฺส ฯเปฯ.}

\prose{เทฺว อิตฺถิโย ทฺวินฺนํ อิตฺถีนํ อิตฺถิสญฺญี สารตฺโต จ, ภิกฺขุ จ นํ ทฺวินฺนํ อิตฺถีนํ นิสฺสคฺคิเยน กายํ อามสติ, อาปตฺติ ทฺวินฺนํ ทุกฺกฏานํ ฯเปฯ.}

\prose{อิตฺถี จ ปณฺฑโก จ อุภินฺนํ อิตฺถิสญฺญี สารตฺโต จ, ภิกฺขุ จ นํ อุภินฺนํ นิสฺสคฺคิเยน กายํ อามสติ, อาปตฺติ ทฺวินฺนํ ทุกฺกฏานํ ฯเปฯ.}

\prose{อิตฺถี จ โหติ อิตฺถิสญฺญี สารตฺโต จ, ภิกฺขุ จ นํ อิตฺถิยา นิสฺสคฺคิเยน กาย\-ปฏิพทฺธํ อามสติ, อาปตฺติ ทุกฺกฏสฺส ฯเปฯ.}

\prose{เทฺว อิตฺถิโย ทฺวินฺนํ อิตฺถีนํ อิตฺถิสญฺญี สารตฺโต จ, ภิกฺขุ จ นํ ทฺวินฺนํ อิตฺถีนํ นิสฺสคฺคิเยน กาย\-ปฏิพทฺธํ อามสติ, อาปตฺติ ทฺวินฺนํ ทุกฺกฏานํ ฯเปฯ.}

\prose{อิตฺถี จ ปณฺฑโก จ อุภินฺนํ อิตฺถิสญฺญี สารตฺโต จ, ภิกฺขุ จ นํ อุภินฺนํ นิสฺสคฺคิเยน กาย\-ปฏิพทฺธํ อามสติ, อาปตฺติ ทฺวินฺนํ ทุกฺกฏานํ ฯเปฯ.}

\prose{อิตฺถี จ โหติ อิตฺถิสญฺญี สารตฺโต จ, ภิกฺขุ จ นํ อิตฺถิยา นิสฺสคฺคิเยน นิสฺสคฺคิยํ อามสติ, อาปตฺติ ทุกฺกฏสฺส ฯเปฯ.}

\prose{เทฺว อิตฺถิโย ทฺวินฺนํ อิตฺถีนํ อิตฺถิสญฺญี สารตฺโต จ, ภิกฺขุ จ นํ ทฺวินฺนํ อิตฺถีนํ นิสฺสคฺคิเยน นิสฺสคฺคิยํ อามสติ, อาปตฺติ ทฺวินฺนํ ทุกฺกฏานํ ฯเปฯ.}

\prosewithcloser{อิตฺถี จ ปณฺฑโก จ อุภินฺนํ อิตฺถิสญฺญี สารตฺโต จ, ภิกฺขุ จ นํ อุภินฺนํ นิสฺสคฺคิเยน นิสฺสคฺคิยํ อามสติ, อาปตฺติ ทฺวินฺนํ ทุกฺกฏานํ ฯเปฯ.}{ภิกฺขุเปยฺยาโล นิฏฺฐิโต.}
\proseitem{277}{\csromanfolio{181}อิตฺถี จ โหติ อิตฺถิสญฺญี สารตฺโต จ, อิตฺถี จ นํ ภิกฺขุสฺส กาเยน กายํ อามสติ ปรามสติ โอมสติ อุมฺมสติ โอลงฺเฆติ อุลฺลงฺเฆติ อากฑฺฒติ ปติกฑฺฒติ อภินิคฺคณฺหาติ อภินิปฺปีเฬติ คณฺหาติ ฉุปติ, เสวนาธิปฺปาโย กาเยน วายมติ ผสฺสํ ปฏิวิชานาติ, อาปตฺติ สํฆาทิเสสสฺส ฯเปฯ.}

\prose{เทฺว อิตฺถิโย ทฺวินฺนํ อิตฺถีนํ อิตฺถิสญฺญี สารตฺโต จ, อิตฺถิโย จ นํ ภิกฺขุสฺส กาเยน กายํ อามสนฺติ ปรามสนฺติ โอมสนฺติ อุมฺมสนฺติ โอลงฺเฆนฺติ อุลฺลงฺเฆนฺติ อากฑฺฒนฺติ ปติกฑฺฒนฺติ อภินิ\-คฺคณฺหนฺติ อภินิปฺปีเฬนฺติ คณฺหนฺติ ฉุปนฺติ, เสวนาธิปฺปาโย กาเยน วายมติ ผสฺสํ ปฏิวิชานาติ, อาปตฺติ ทฺวินฺนํ สํฆา\-ทิ\-เสสานํ ฯเปฯ.}

\prose{อิตฺถี จ ปณฺฑโก จ อุภินฺนํ อิตฺถิสญฺญี สารตฺโต จ, อุโภ จ นํ ภิกฺขุสฺส กาเยน กายํ อามสนฺติ ปรามสนฺติ ฯเปฯ คณฺหนฺติ ฉุปนฺติ, เสวนาธิปฺปาโย กาเยน วายมติ ผสฺสํ ปฏิวิชานาติ, อาปตฺติ สํฆาทิเสเสน ทุกฺกฏสฺส ฯเปฯ.}

\prose{อิตฺถี จ โหติ อิตฺถิสญฺญี สารตฺโต จ, อิตฺถี จ นํ ภิกฺขุสฺส กาเยน กาย\-ปฏิพทฺธํ อามสติ ปรามสติ ฯเปฯ คณฺหาติ ฉุปติ, เสวนาธิปฺปาโย กาเยน วายมติ ผสฺสํ ปฏิวิชานาติ, อาปตฺติ ถุลฺลจฺจยสฺส ฯเปฯ.}

\prose{เทฺว อิตฺถิโย ทฺวินฺนํ อิตฺถีนํ อิตฺถิสญฺญี สารตฺโต จ, อิตฺถิโย จ นํ ภิกฺขุสฺส กาเยน กาย\-ปฏิพทฺธํ อามสนฺติ ปรามสนฺติ ฯเปฯ คณฺหนฺติ ฉุปนฺติ, เสวนาธิปฺปาโย กาเยน วายมติ ผสฺสํ ปฏิวิชานาติ, อาปตฺติ ทฺวินฺนํ ถุลฺลจฺจยานํ ฯเปฯ.}

\prose{อิตฺถี จ ปณฺฑโก จ อุภินฺนํ อิตฺถิสญฺญี สารตฺโต จ, อุโภ จ นํ ภิกฺขุสฺส กาเยน กาย\-ปฏิพทฺธํ อามสนฺติ ปรามสนฺติ ฯเปฯ คณฺหนฺติ ฉุปนฺติ, เสวนาธิปฺปาโย กาเยน วายมติ ผสฺสํ ปฏิวิชานาติ, อาปตฺติ ถุลฺลจฺจเยน ทุกฺกฏสฺส ฯเปฯ.}

\prose{อิตฺถี จ โหติ อิตฺถิสญฺญี สารตฺโต จ, อิตฺถี จ นํ ภิกฺขุสฺส กาย\-ปฏิพทฺเธน กายํ อามสติ ปรามสติ ฯเปฯ. คณฺหาติ ฉุปติ, เสวนาธิปฺปาโย กาเยน วายมติ ผสฺสํ ปฏิวิชานาติ, อาปตฺติ ถุลฺลจฺจยสฺส ฯเปฯ.}

\prose{\csromanfolio{182}เทฺว อิตฺถิโย ทฺวินฺนํ อิตฺถีนํ อิตฺถิสญฺญี สารตฺโต จ, อิตฺถิโย จ นํ ภิกฺขุสฺส กาย\-ปฏิพทฺเธน กายํ อามสนฺติ ปรามสนฺติ ฯเปฯ คณฺหนฺติ ฉุปนฺติ, เสวนาธิปฺปาโย กาเยน วายมติ ผสฺสํ ปฏิวิชานาติ, อาปตฺติ ทฺวินฺนํ ถุลฺลจฺจยานํ ฯเปฯ.}

\prose{อิตฺถี จ ปณฺฑโก จ อุภินฺนํ อิตฺถิสญฺญี สารตฺโต จ, อุโภ จ นํ ภิกฺขุสฺส กาย\-ปฏิพทฺเธน กายํ อามสนฺติ ปรามสนฺติ ฯเปฯ คณฺหนฺติ ฉุปนฺติ, เสวนาธิปฺปาโย กาเยน วายมติ ผสฺสํ ปฏิวิชานาติ, อาปตฺติ ถุลฺลจฺจเยน ทุกฺกฏสฺส ฯเปฯ.}

\proseitem{278}{อิตฺถี จ โหติ อิตฺถิสญฺญี สารตฺโต จ, อิตฺถี จ นํ ภิกฺขุสฺส กาย\-ปฏิพทฺเธน กาย\-ปฏิพทฺธํ อามสติ ปรามสติ ฯเปฯ คณฺหาติ ฉุปติ, เสวนาธิปฺปาโย กาเยน วายมติ ผสฺสํ ปฏิวิชานาติ, อาปตฺติ ทุกฺกฏสฺส ฯเปฯ.}

\prose{เทฺว อิตฺถิโย ทฺวินฺนํ อิตฺถีนํ อิตฺถิสญฺญี สารตฺโต จ, อิตฺถิโย จ นํ ภิกฺขุสฺส กาย\-ปฏิพทฺเธน กาย\-ปฏิพทฺธํ อามสนฺติ ปรามสนฺติ ฯเปฯ คณฺหนฺติ ฉุปนฺติ, เสวนาธิปฺปาโย กาเยน วายมติ ผสฺสํ ปฏิวิชานาติ, อาปตฺติ ทฺวินฺนํ ทุกฺกฏานํ ฯเปฯ.}

\prose{อิตฺถี จ ปณฺฑโก จ อุภินฺนํ อิตฺถิสญฺญี สารตฺโต จ, อุโภ จ นํ ภิกฺขุสฺส กาย\-ปฏิพทฺเธน กาย\-ปฏิพทฺธํ อามสนฺติ ปรามสนฺติ ฯเปฯ คณฺหนฺติ ฉุปนฺติ, เสวนาธิปฺปาโย กาเยน วายมติ ผสฺสํ ปฏิวิชานาติ, อาปตฺติ ทฺวินฺนํ ทุกฺกฏานํ ฯเปฯ.}

\prose{อิตฺถี จ โหติ อิตฺถิสญฺญี สารตฺโต จ, อิตฺถี จ นํ ภิกฺขุสฺส นิสฺสคฺคิเยน กายํ อามสติ เสวนาธิปฺปาโย กาเยน วายมติ ผสฺสํ ปฏิวิชานาติ, อาปตฺติ ทุกฺกฏสฺส ฯเปฯ.}

\prose{เทฺว อิตฺถิโย ทฺวินฺนํ อิตฺถีนํ อิตฺถิสญฺญี สารตฺโต จ, อิตฺถิโย จ นํ ภิกฺขุสฺส นิสฺสคฺคิเยน กายํ อามสนฺติ, เสวนาธิปฺปาโย กาเยน วายมติ ผสฺสํ ปฏิวิชานาติ, อาปตฺติ ทฺวินฺนํ ทุกฺกฏานํ ฯเปฯ.}

\prose{อิตฺถี จ ปณฺฑโก จ อุภินฺนํ อิตฺถิสญฺญี สารตฺโต จ, อุโภ จ นํ ภิกฺขุสฺส นิสฺสคฺคิเยน กายํ อามสนฺติ, เสวนาธิปฺปาโย กาเยน วายมติ ผสฺสํ ปฏิวิชานาติ, อาปตฺติ ทฺวินฺนํ ทุกฺกฏานํ ฯเปฯ.}

\prose{\csromanfolio{183}อิตฺถี จ โหติ อิตฺถิสญฺญี สารตฺโต จ, อิตฺถี จ นํ ภิกฺขุสฺส นิสฺสคฺคิเยน กาย\-ปฏิพทฺธํ อามสติ, เสวนาธิปฺปาโย กาเยน วายมติ ผสฺสํ ปฏิวิชานาติ, อาปตฺติ ทุกฺกฏสฺส ฯเปฯ.}

\prose{เทฺว อิตฺถิโย ทฺวินฺนํ อิตฺถีนํ อิตฺถิสญฺญี สารตฺโต จ, อิตฺถิโย จ นํ ภิกฺขุสฺส นิสฺสคฺคิเยน กาย\-ปฏิพทฺธํ อามสนฺติ, เสวนาธิปฺปาโย กาเยน วายมติ ผสฺสํ ปฏิวิชานาติ, อาปตฺติ ทฺวินฺนํ ทุกฺกฏานํ ฯเปฯ.}

\prose{อิตฺถี จ ปณฺฑโก จ อุภินฺนํ อิตฺถิสญฺญี สารตฺโต จ, อุโภ จ นํ ภิกฺขุสฺส นิสฺสคฺคิเยน กาย\-ปฏิพทฺธํ อามสนฺติ, เสวนาธิปฺปาโย กาเยน วายมติ ผสฺสํ ปฏิวิชานาติ, อาปตฺติ ทฺวินฺนํ ทุกฺกฏานํ ฯเปฯ.}

\prose{อิตฺถี จ โหติ อิตฺถิสญฺญี สารตฺโต จ, อิตฺถี จ นํ ภิกฺขุสฺส นิสฺสคฺคิเยน นิสฺสคฺคิยํ อามสติ, เสวนาธิปฺปาโย กาเยน วายมติ, น จ ผสฺสํ ปฏิวิชานาติ, อาปตฺติ ทุกฺกฏสฺส ฯเปฯ.}

\prose{เทฺว อิตฺถิโย ทฺวินฺนํ อิตฺถีนํ อิตฺถิสญฺญี สารตฺโต จ, อิตฺถิโย จ นํ ภิกฺขุสฺส นิสฺสคฺคิเยน นิสฺสคฺคิยํ อามสนฺติ, เสวนาธิปฺปาโย กาเยน วายมติ, น จ ผสฺสํ ปฏิวิชานาติ, อาปตฺติ ทฺวินฺนํ ทุกฺกฏานํ ฯเปฯ.}

\prose{อิตฺถี จ ปณฺฑโก จ อุภินฺนํ อิตฺถิสญฺญี สารตฺโต จ, อุโภ จ นํ ภิกฺขุสฺส นิสฺสคฺคิเยน นิสฺสคฺคิยํ อามสนฺติ, เสวนาธิปฺปาโย กาเยน วายมติ ผสฺสํ ปฏิวิชานาติ, อาปตฺติ ทฺวินฺนํ ทุกฺกฏานํ ฯเปฯ.}

\proseitem{279}{เสวนาธิปฺปาโย กาเยน วายมติ ผสฺสํ ปฏิวิชานาติ, อาปตฺติ สํฆาทิเสสสฺส.}

\prose{เสวนาธิปฺปาโย กาเยน วายมติ, น จ ผสฺสํ ปฏิวิชานาติ, อาปตฺติ ทุกฺกฏสฺส.}

\prose{เสวนาธิปฺปาโย น จ กาเยน วายมติ, ผสฺสํ ปฏิวิชานาติ, อนาปตฺติ.}

\prose{เสวนาธิปฺปาโย น จ กาเยน วายมติ, น จ ผสฺสํ ปฏิวิชานาติ, อนาปตฺติ.}

\prose{โมกฺขาธิปฺปาโย กาเยน วายมติ, ผสฺสํ ปฏิวิชานาติ, อนาปตฺติ.}

\prose{\csromanfolio{184}โมกฺขาธิปฺปาโย กาเยน วายมติ, น จ ผสฺสํ ปฏิวิชานาติ, อนาปตฺติ.}

\prose{โมกฺขาธิปฺปาโย น จ กาเยน วายมติ, ผสฺสํ ปฏิวิชานาติ, อนาปตฺติ.}

\prose{โมกฺขาธิปฺปาโย น จ กาเยน วายมติ, น จ ผสฺสํ ปฏิวิชานาติ, อนาปตฺติ.}

\proseitem{280}{อนาปตฺติ อสญฺจิจฺจ อสติยา อชานนฺตสฺส อสาทิยนฺตสฺส อุมฺมตฺตกสฺส ขิตฺตจิตฺตสฺส เวทนาฏฺฏสฺส อาทิกมฺมิกสฺสาติ.}
\csromansectionrule

\csromanmatikaanchor{mtk.47}
\csromantocmarkref{subsection}{วินีตวตฺถุอุทฺทานคาถา}{mtk.47}
\bookmark[dest={mtk.47},level=2]{วินีตวตฺถุอุทฺทานคาถา}
\csromanheader{2}{\textbf{วินีตวตฺถุ}\-อุทฺทานคาถา}
\setlength{\csromangathapagewidth}{0pt}
\setlength{\csromangathawakwidth}{0pt}
\csromangathameasuregroup{มาตา ธีตา ภคินี จ, \\ สุตฺตา มตา ติรจฺฉานา, \\ สมฺปีเฬ สงฺกโม มคฺโค, \\ ทณฺโฑ ปตฺตํ ปณาเมสิ,}{\csromangathabat{มาตา ธีตา ภคินี จ,}{ชายา ยกฺขี จ ปณฺฑโก.} \\ \csromangathabat{สุตฺตา มตา ติรจฺฉานา,}{ทารุธีตลิกาย จ.} \\ \csromangathabat{สมฺปีเฬ สงฺกโม มคฺโค,}{รุกฺโข นาวา จ รชฺชุ จ.} \\ \csromangathabat{ทณฺโฑ ปตฺตํ ปณาเมสิ,}{วนฺเท วายมิ นจฺฉุเปติ.}}
\csromangathasetleft{มาตา ธีตา ภคินี จ, \\ สุตฺตา มตา ติรจฺฉานา, \\ สมฺปีเฬ สงฺกโม มคฺโค, \\ ทณฺโฑ ปตฺตํ ปณาเมสิ,}
\csromangathagroup{\csromangathastanza{\csromangathabat{มาตา ธีตา ภคินี จ,}{ชายา ยกฺขี จ ปณฺฑโก.} \\ \csromangathabat{สุตฺตา มตา ติรจฺฉานา,}{ทารุ\-ธีตลิกาย จ.}}\csromangathastanza{\csromangathabat{สมฺปีเฬ สงฺกโม มคฺโค,}{รุกฺโข นาวา จ รชฺชุ จ.} \\ \csromangathabat{ทณฺโฑ ปตฺตํ ปณาเมสิ,}{วนฺเท วายมิ นจฺฉุเปติ.}}}

\csromanmatikaanchor{mtk.48}
\csromantocmarkref{subsection}{วินีตวตฺถุ}{mtk.48}
\bookmark[dest={mtk.48},level=2]{วินีตวตฺถุ}
\csromanheader{2}{\textbf{วินีตวตฺถุ}}
\proseitem{281}{เตน โข ปน สมเยน อญฺญตโร ภิกฺขุ มาตุยา มาตุเปเมน อามสิ. ตสฺส กุกฺกุจฺจํ อโหสิ “ภควตา สิกฺขาปทํ ปญฺญตฺตํ, กจฺจิ นุ โข อหํ สํฆาทิเสสํ อาปตฺติํ อาปนฺโน”ติ. ภควโต เอตมตฺถํ อาโรเจสิ. อนาปตฺติ ภิกฺขุ สํฆาทิเสสสฺส, อาปตฺติ ทุกฺกฏสฺสาติ. (1)}

\prose{เตน โข ปน สมเยน อญฺญตโร ภิกฺขุ ธีตุยา ธีตุเปเมน อามสิ ฯเปฯ ภคินิยา ภคินิเปเมน อามสิ. ตสฺส กุกฺกุจฺจํ อโหสิ ฯเปฯ. อนาปตฺติ ภิกฺขุ สํฆาทิเสสสฺส, อาปตฺติ ทุกฺกฏสฺสาติ. \mbox{(2-3)}}

\prose{เตน โข ปน สมเยน อญฺญตโร ภิกฺขุ ปุราณ\-ทุติยิกาย กายสํสคฺคํ สมาปชฺชิ. ตสฺส กุกฺกุจฺจํ อโหสิ ฯเปฯ. อาปตฺติํ ตฺวํ ภิกฺขุ อาปนฺโน สํฆาทิเสสนฺติ. (4)}

\prose{เตน โข ปน สมเยน อญฺญตโร ภิกฺขุ ยกฺขินิยา กายสํสคฺคํ สมาปชฺชิ. ตสฺส กุกฺกุจฺจํ อโหสิ ฯเปฯ. อนาปตฺติ ภิกฺขุ สํฆาทิเสสสฺส, อาปตฺติ ถุลฺลจฺจยสฺสาติ. (5)}

\prose{\csromanfolio{185}เตน โข ปน สมเยน อญฺญตโร ภิกฺขุ ปณฺฑกสฺส กายสํสคฺคํ สมาปชฺชิ. ตสฺส กุกฺกุจฺจํ อโหสิ ฯเปฯ. อนาปตฺติ สํฆาทิเสสสฺส, อาปตฺติ ถุลฺลจฺจยสฺสาติ. (6)}

\prose{เตน โข ปน สมเยน อญฺญตโร ภิกฺขุ สุตฺติตฺถิยา กายสํสคฺคํ สมาปชฺชิ. ตสฺส กุกฺกุจฺจํ อโหสิ ฯเปฯ. อาปตฺติํ ตฺวํ ภิกฺขุ อาปนฺโน สํฆาทิเสสนฺติ. (7)}

\prose{เตน โข ปน สมเยน อญฺญตโร ภิกฺขุ มติตฺถิยา กายสํสคฺคํ สมาปชฺชิ. ตสฺส กุกฺกุจฺจํ อโหสิ ฯเปฯ. อนาปตฺติ ภิกฺขุ สํฆาทิเสสสฺส, อาปตฺติ ถุลฺลจฺจยสฺสาติ. (8)}

\prose{เตน โข ปน สมเยน อญฺญตโร ภิกฺขุ ติรจฺฉาน\-คติ\-ตฺถิยา กายสํสคฺคํ สมาปชฺชิ. ตสฺส กุกฺกุจฺจํ อโหสิ ฯเปฯ. อนาปตฺติ ภิกฺขุ สํฆาทิเสสสฺส, อาปตฺติ ทุกฺกฏสฺสาติ. (9)}

\prose{เตน โข ปน สมเยน อญฺญตโร ภิกฺขุ ทารุ\-ธีตลิกาย กายสํสคฺคํ สเมปชฺชิ. ตสฺส กุกฺกุจฺจํ อโหสิ ฯเปฯ. อนาปตฺติ ภิกฺขุ สํฆาทิเสสสฺส, อาปตฺติ ทุกฺกฏสฺสาติ. (10)}

\proseitem{282}{เตน โข ปน สมเยน สมฺพหุลา อิตฺถิโย อญฺญตรํ ภิกฺขุํ สมฺปีเฬตฺวา พาหา\-ปรมฺปราย อาเนสุํ. ตสฺส กุกฺกุจฺจํ อโหสิ ฯเปฯ. สาทิยิ ตฺวํ ภิกฺขูติ. นาหํ ภควา สาทิยินฺติ. อนาปตฺติ ภิกฺขุ อสาทิยนฺตสฺสาติ. (11)}

\prose{เตน โข ปน สมเยน อญฺญตโร ภิกฺขุ อิตฺถิยา อภิรูฬฺหํ สงฺกมํ สารตฺโต สญฺจาเลสิ. ตสฺส กุกฺกุจฺจํ อโหสิ ฯเปฯ. อนาปตฺติ ภิกฺขุ สํฆาทิเสสสฺส, อาปตฺติ ทุกฺกฏสฺสาติ. (12)}

\prose{เตน โข ปน สมเยน อญฺญตโร ภิกฺขุ อิตฺถิํ ปฏิปเถ ปสฺสิตฺวา สารตฺโต อํสกูเฏน ปหารํ อทาสิ. ตสฺส กุกฺกุจฺจํ อโหสิ ฯเปฯ. อาปตฺติํ ตฺวํ ภิกฺขุ อาปนฺโน สํฆาทิเสสนฺติ. (13)}

\prose{เตน โข ปน สมเยน อญฺญตโร ภิกฺขุ อิตฺถิยา อภิรูฬฺหํ รุกฺขํ สารตฺโต สญฺจาเลสิ. ตสฺส กุกฺกุจฺจํ อโหสิ ฯเปฯ. อนาปตฺติ ภิกฺขุ สํฆาทิเสสสฺส, อาปตฺติ ทุกฺกฏสฺสาติ. (14)}

\csromanmatikaanchor{mtk.49}
\csromanmatikaanchor{mtk.50}
\csromantocmarknopage{section}{3. ทุฏฺฐุลฺลวาจาสิกฺขาปท}{mtk.49}
\bookmark[dest={mtk.49},level=1]{3. ทุฏฺฐุลฺลวาจาสิกฺขาปท}
\csromantocmarkref{subsection}{อุทายิภิกฺขุวตฺถุ}{mtk.50}
\bookmark[dest={mtk.50},level=2]{อุทายิภิกฺขุวตฺถุ}
\prose{\csromanfolio{186}เตน โข ปน สมเยน อญฺญตโร ภิกฺขุ อิตฺถิยา อภิรูฬฺหํ นาวํ สารตฺโต สญฺจาเลสิ. ตสฺส กุกฺกุจฺจํ อโหสิ ฯเปฯ. อนาปตฺติ ภิกฺขุ สํฆาทิเสสสฺส, อาปตฺติ ทุกฺกฏสฺสาติ. (15)}

\prose{เตน โข ปน สมเยน อญฺญตโร ภิกฺขุ อิตฺถิยา คหิตํ รชฺชุํ สารตฺโต อาวิญฺฉิ\csromansharedfootnote{d9378e62aedb2b83}{อาวิญฺชิ (สี, สฺยา)}. ตสฺส กุกฺกุจฺจํ อโหสิ ฯเปฯ. อนาปตฺติ ภิกฺขุ สํฆาทิเสสสฺส, อาปตฺติ ถุลฺลจฺจยสฺสาติ. (16)}

\prose{เตน โข ปน สมเยน อญฺญตโร ภิกฺขุ อิตฺถิยา คหิตํ ทณฺฑํ สารตฺโต อาวิญฺฉิ. ตสฺส กุกฺกุจฺจํ อโหสิ ฯเปฯ. อนาปตฺติ ภิกฺขุ สํฆาทิเสสสฺส, อาปตฺติ ถุลฺลจฺจยสฺสาติ. (17)}

\prose{เตน โข ปน สมเยน อญฺญตโร ภิกฺขุ สารตฺโต อิตฺถิํ ปตฺเตน ปณาเมสิ. ตสฺส กุกฺกุจฺจํ อโหสิ ฯเปฯ. อนาปตฺติ ภิกฺขุ สํฆาทิเสสสฺส, อาปตฺติ ถุลฺลจฺจยสฺสาติ. (18)}

\prose{เตน โข ปน สมเยน อญฺญตโร ภิกฺขุ อิตฺถิยา วนฺทนฺติยา สารตฺโต ปาทํ อุจฺจาเรสิ. ตสฺส กุกฺกุจฺจํ อโหสิ ฯเปฯ. อาปตฺติํ ตฺวํ ภิกฺขุ อาปนฺโน สํฆาทิเสสนฺติ. (19)}

\prosewithcloserruled{เตน โข ปน สมเยน อญฺญตโร ภิกฺขุ อิตฺถิํ คเหสฺสามีติ วายมิตฺวา น ฉุปิ. ตสฺส กุกฺกุจฺจํ อโหสิ ฯเปฯ. อนาปตฺติ ภิกฺขุ สํฆาทิเสสสฺส, อาปตฺติ ทุกฺกฏสฺสาติ. (20)}{กายสํสคฺค\-สิกฺขาปทํ นิฏฺฐิตํ ทุติยํ.}
\csromanheader{1}{3. \textbf{ทุฏฺฐุลฺลวาจา}\-\textbf{สิกฺขาปท}}
\csromanmatikaanchor{mtk.51}
\csromanmatikaanchor{mtk.52}
\csromantocmarkref{subsection}{ปญฺญตฺติ}{mtk.51}
\bookmark[dest={mtk.51},level=2]{ปญฺญตฺติ}
\csromantocmarkref{subsection}{สิกฺขาปทวิภงฺค, ปทภาชนีย}{mtk.52}
\bookmark[dest={mtk.52},level=2]{สิกฺขาปทวิภงฺค, ปทภาชนีย}
\proseitem{238}{เตน สมเยน พุทฺโธ ภควา สาวตฺถิยํ วิหรติ เชตวเน อนาถ\-ปิณฺฑิกสฺส อาราเม. เตน โข ปน สมเยน อายสฺมา อุทายี อรญฺเญ วิหรติ. ตสฺสายสฺมโต วิหาโร อภิรูโป โหติ ทสฺสนีโย ปาสาทิโก. เตน โข ปน สมเยน สมฺพหุลา อิตฺถิโย อารามํ อาคมํสุ วิหารเปกฺขิกาโย. อถ โข ตา \csromanfolio{187}อิตฺถิโย เยนายสฺมา อุทายี เตนุ\-ปสงฺกมิํสุ, อุปสงฺกมิตฺวา อายสฺมนฺตํ อุทายิํ เอตทโวจุํ “อิจฺฉาม มยํ ภนฺเต อยฺยสฺส วิหารํ เปกฺขิตุนฺ”ติ. อถ โข อายสฺมา อุทายี ตา อิตฺถิโย วิหารํ เปกฺขาเปตฺวา ตาสํ อิตฺถีนํ วจฺจมคฺคํ ปสฺสาวมคฺคํ อาทิสฺส วณฺณํปิ ภณติ อวณฺณํปิ ภณติ ยาจติปิ อายาจติปิ ปุจฺฉติปิ ปฏิปุจฺฉติปิ อาจิกฺขติปิ อนุสาสติปิ อกฺโกสติปิ. ยา ตา อิตฺถิโย ฉินฺนิกา ธุตฺติกา อหิริกาโย, ตา อายสฺมตา อุทายินา สทฺธิํ อุหสนฺติปิ อุลฺลปนฺติปิ อุชฺชคฺฆนฺติปิ อุปฺปณฺเฑนฺติปิ. ยา ปน ตา อิตฺถิโย หิริมนา, ตา นิกฺขมิตฺวา ภิกฺขู อุชฺฌาเปนฺติ “อิทํ ภนฺเต นจฺฉินฺนํ นปฺปติรูปํ สามิเกนปิ มยํ เอวํ วุตฺตา น อิจฺเฉยฺยาม, กิํ ปนายฺเยน อุทายินา”ติ. เย เต ภิกฺขู อปฺปิจฺฉา ฯเปฯ เต อุชฺฌายนฺติ ขิยฺยนฺติ วิปาเจนฺติ “กถํ หิ นาม อายสฺมา อุทายี มาตุคามํ ทุฏฺฐุลฺลาหิ วาจาหิ โอภาสิสฺสตี”ติ. อถ โข เต ภิกฺขู อายสฺมนฺตํ อุทายิํ อเนก\-ปริยาเยน วิครหิตฺวา ภควโต เอตมตฺถํ อาโรเจสุํ. อถ โข ภควา เอตสฺมิํ นิทาเน เอตสฺมิํ ปกรเณ ภิกฺขุสํฆํ สนฺนิปาตาเปตฺวา อายสฺมนฺตํ อุทายิํ ปฏิปุจฺฉิ “สจฺจํ กิร ตฺวํ อุทายิ มาตุคามํ ทุฏฺฐุลฺลาหิ วาจาหิ โอภาสสี”ติ. สจฺจํ ภควาติ. วิครหิ พุทฺโธ ภควา “อนนุจฺฉวิกํ โมฆปุริส อนนุโลมิกํ อปฺปติรูปํ อสฺสามณกํ อกปฺปิยํ อกรณียํ, กถํ หิ นาม โมฆปุริส มาตุคามํ ทุฏฺฐุลฺลาหิ วาจาหิ โอภาสิสฺสสิ. นนุ มยา โมฆปุริส อเนก\-ปริยาเยน วิราคาย ธมฺโม เทสิโต, โน สราคาย ฯเปฯ กาม\-ปริฬาหานํ วูปสโม อกฺขาโต, เนตํ โมฆปุริส อปฺปสนฺนานํ วา ปสาทาย ฯเปฯ. เอวญฺจ ปน ภิกฺขเว อิมํ สิกฺขาปทํ อุทฺทิเสยฺยาถ.}

\proseitem{284}{\textbf{“โย ปน ภิกฺขุ โอติณฺโณ วิปริณเตน จิตฺเตน มาตุคามํ ทุฏฺฐุลฺลาหิ วาจาหิ โอภาเสยฺย ยถา ตํ ยุวา ยุวติํ เมถุนุ}\-\textbf{ปสํหิตาหิ, สํฆาทิเสโส”}ติ.}

\proseitem{285}{\textbf{โย ปนา}ติ โย ยาทิโส ฯเปฯ. \textbf{ภิกฺขู}ติ ฯเปฯ. อยํ อิมสฺมิํ อตฺเถ อธิปฺเปโต “ภิกฺขู”ติ.}

\prose{\textbf{โอติณฺโณ} นาม สารตฺโต อเปกฺขวา ปฏิพทฺธ\-จิตฺโต.}

\prose{\csromanfolio{188}\textbf{วิปริณตนฺ}ติ รตฺตํปิ จิตฺตํ วิปริณตํ, ทุฏฺฐํปิ จิตฺตํ วิปริณตํ, มูฬฺหํปิ จิตฺตํ วิปริณตํ. อปิ จ รตฺตํ จิตฺตํ อิมสฺมิํ อตฺเถ อธิปฺเปตํ “วิปริณตนฺ”ติ.}

\prose{\textbf{มาตุคาโม} นาม มนุสฺสิตฺถี, น ยกฺขี, น เปตี, น ติรจฺฉานคตา, วิญฺญู ปฏิพลา สุภาสิต\-ทุพฺภาสิตํ ทุฏฺฐุลฺลา\-ทุฏฺฐุลฺลํ อาชานิตุํ.}

\prose{\textbf{ทุฏฺฐุลฺลา} นาม วาจา วจฺจมคฺค\-ปสฺสาวมคฺค\-เมถุนธมฺม\-ปฺปฏิสํยุตฺตา \textbf{วาจา}.}

\prose{\textbf{โอภาเสยฺยา}ติ อชฺฌาจาโร วุจฺจติ.}

\prose{\textbf{ยถา ตํ ยุวา ยุวตินฺ}ติ ทหโร ทหริํ ตรุโณ ตรุณิํ กามโภคี กามโภคินิํ.}

\prose{\textbf{เมถุนุปสํหิตาหี}ติ เมถุนธมฺม\-ปฺปฏิสํยุตฺตาหิ.}

\prose{\textbf{สํฆาทิเสโส}ติ ฯเปฯ เตนปิ วุจฺจติ “สํฆาทิเสโส”ติ.}

\prose{เทฺว มคฺเค อาทิสฺส วณฺณํปิ ภณติ อวณฺณํปิ ภณติ ยาจติปิ อายาจติปิ ปุจฺฉติปิ ปฏิปุจฺฉติปิ อาจิกฺขติปิ อนุสาสติปิ อกฺโกสติปิ.}

\prose{\textbf{วณฺณํ ภณติ} นาม เทฺว มคฺเค โถเมติ วณฺเณติ ปสํสติ.}

\prose{\textbf{อวณฺณํ ภณติ} นาม เทฺว มคฺเค ขุํเสติ วมฺเภติ ครหติ.}

\prose{\textbf{ยาจติ} นาม เทหิ เม อรหสิ เม ทาตุนฺติ.}

\prose{\textbf{อายาจติ} นาม กทา เต มาตา ปสีทิสฺสติ, กทา เต ปิตา ปสีทิสฺสติ, กทา เต เทวตาโย ปสีทิสฺสนฺติ, กทา\csromansharedfootnote{3f8a8c2be0eb0087}{กทา เต (สฺยา)} สุขโณ สุลโย สุมุหุตฺโต ภวิสฺสติ, กทา เต เมถุนํ ธมฺมํ ลภิสฺสามีติ.}

\prose{\textbf{ปุจฺฉติ} นาม กถํ ตฺวํ สามิกสฺส เทสิ, กถํ ชารสฺส เทสีติ.}

\prose{\textbf{ปฏิปุจฺฉติ} นาม เอวํ กิร ตฺวํ สามิกสฺส เทสิ เอวํ ชารสฺส เทสีติ.}

\prose{\textbf{อาจิกฺขติ} นาม ปุฏฺโฐ ภณติ “เอวํ เทหิ, เอวํ เทนฺตา สามิกสฺส ปิยา ภวิสฺสติ มนาปา จา”ติ.}

\prose{\textbf{อนุสาสติ} นาม อปุฏฺโฐ ภณติ “เอวํ เทหิ, เอวํ เทนฺตา สามิกสฺส ปิยา ภวิสฺสติ มนาปา จา”ติ.}

\prose{\csromanfolio{189}\textbf{อกฺโกสติ} นาม อนิมิตฺตาสิ นิมิตฺตมตฺตาสิ อโลหิตาสิ ธุวโลหิตาสิ ธุวโจฬาสิ ปคฺฆรนฺตีสิ สิขรณีสิ อิตฺถิปณฺฑกาสิ เวปุริสิกาสิ สมฺภินฺนาสิ อุภโตพฺยญฺชนาสีติ.}

\proseitem{286}{อิตฺถี จ โหติ อิตฺถิสญฺญี สารตฺโต จ, ภิกฺขุ จ นํ อิตฺถิยา วจฺจมคฺคํ ปสฺสาวมคฺคํ อาทิสฺส วณฺณํปิ ภณติ อวณฺณํปิ ภณติ ยาจติปิ อายาจติปิ ปุจฺฉติปิ ปฏิปุจฺฉติปิ อาจิกฺขติปิ อนุสาสติปิ อกฺโกสติปิ, อาปตฺติ สํฆาทิเสสสฺส ฯเปฯ.}

\prose{เทฺว อิตฺถิโย ทฺวินฺนํ อิตฺถีนํ อิตฺถิสญฺญี สารตฺโต จ, ภิกฺขุ จ นํ ทฺวินฺนํ อิตฺถีนํ วจฺจมคฺคํ ปสฺสาวมคฺคํ อาทิสฺส วณฺณํปิ ภณติ อวณฺณํปิ ภณติ ฯเปฯ อกฺกาสติปิ, อาปตฺติ ทฺวินฺนํ สํฆา\-ทิ\-เสสานํ ฯเปฯ.}

\prose{อิตฺถี จ ปณฺฑโก จ อุภินฺนํ อิตฺถิสญฺญี สารตฺโต จ, ภิกฺขุ จ นํ อุภินฺนํ วจฺจมคฺคํ ปสฺสาวมคฺคํ อาทิสฺส วณฺณํปิ ภณติ อวณฺณํปิ ภณติ ฯเปฯ อกฺโกสติปิ, อาปตฺติ สํฆาทิเสเสน ทุกฺกฏสฺส ฯเปฯ.}

\prose{อิตฺถี จ โหติ อิตฺถิสญฺญี สารตฺโต จ, ภิกฺขุ จ นํ อิตฺถิยา วจฺจมคฺคํ ปสฺสาวมคฺคํ ฐเปตฺวา อธกฺเขกํ อุพฺภ\-ชาณุมณฺฑลํ อาทิสฺส วณฺณํปิ ภณติ อวณฺณํปิ ภณติ ฯเปฯ อกฺโกสติปิ, อาปตฺติ ถุลฺลจฺจยสฺส ฯเปฯ.}

\prose{เทฺว อิตฺถิโย ทฺวินฺนํ อิตฺถีนํ อิตฺถิสญฺญี สารตฺโต จ, ภิกฺขุ จ นํ ทฺวินฺนํ อิตฺถีนํ วจฺจมคฺคํ ปสฺสาวมคฺคํ ฐาเปตฺวา อธกฺขกํ อุพฺภ\-ชาณุมณฺฑลํ อาทิสฺส วณฺณํปิ ภณติ อวณฺณํปิ ภณติ ฯเปฯ อกฺโกสติปิ, อาปตฺติ ทฺวินฺนํ ถุลฺลจฺจยานํ ฯเปฯ.}

\prose{อิตฺถี จ ปณฺฑโก จ อุภินฺนํ อิตฺถิสญฺญี สารตฺโต จ, ภิกฺขุ จ นํ อุภินฺนํ วจฺจมคฺคํ ปสฺสาวมคฺคํ ฐเปตฺวา อธกฺขกํ อุพฺภ\-ชาณุมณฺฑลํ อาทิสฺส วณฺณํปิ ภณติ อวณฺณํปิ ภณติ ฯเปฯ อกฺโกสติปิ, อาปตฺติ ถุลฺลจฺจเยน ทุกฺกฏสฺส ฯเปฯ.}

\prose{อิตฺถี จ โหติ อิตฺถิสญฺญี สารตฺโต จ, ภิกฺขุ จ นํ อิตฺถิยา อุพฺภกฺขกํ อโธ\-ชาณุมณฺฑลํ อาทิสฺส วณฺณํปิ ภณติ อวณฺณํปิ ภณติ ฯเปฯ อกฺโกสติปิ, อาปตฺติ ทุกฺกฏสฺส ฯเปฯ.}

\prose{เทฺว อิตฺถิโย ทฺวินฺนํ อิตฺถีนํ อิตฺถิสญฺญี สารตฺโต จ, ภิกฺขุ จ นํ ทฺวินฺนํ อิตฺถีนํ อุพฺภกฺขกํ อโธ\-ชาณุมณฺฑลํ อาทิสฺส วณฺณํปิ ภณติ อวณฺณํปิ ภณติ ฯเปฯ อกฺโกสติปิ, อาปตฺติ ทฺวินฺนํ ทุกฺกฏานํ ฯเปฯ.}

\prose{\csromanfolio{190}อิตฺถี จ ปณฺฑโก จ อุภินฺนํ อิตฺถิสญฺญี สารตฺโต จ, ภิกฺขุ จ นํ อุภินฺนํ อุพฺภกฺขกํ อโธ\-ชาณุมณฺฑลํ อาทิสฺส วณฺณํปิ ภณติ อวณฺณํปิ ภณติ ฯเปฯ อกฺโกสติปิ, อาปตฺติ ทฺวินฺนํ ทุกฺกฏานํ ฯเปฯ.}

\prose{อิตฺถี จ โหติ อิตฺถิสญฺญี สารตฺโต จ, ภิกฺขุ จ นํ อิตฺถิยา กาย\-ปฏิพทฺธํ อาทิสฺส วณฺณํปิ ภณติ อวณฺณํปิ ภณติ ฯเปฯ อกฺโกสติปิ, อาปตฺติ ทุกฺกฏสฺส ฯเปฯ.}

\prose{เทฺว อิตฺถิโย ทฺวินฺนํ อิตฺถีนํ อิตฺถิสญฺญี สารตฺโต จ, ภิกฺขุ จ นํ ทฺวินฺนํ อิตฺถีนํ กาย\-ปฏิพทฺธํ อาทิสฺส วณฺณํปิ ภณติ อวณฺณํปิ ภณติ ฯเปฯ อกฺโกสติปิ, อาปตฺติ ทฺวินฺนํ ทุกฺกฏานํ ฯเปฯ.}

\prose{อิตฺถี จ ปณฺฑโก จ อุภินฺนํ อิตฺถิสญฺญี สารตฺโต จ, ภิกฺขุ จ นํ อุภินฺนํ กาย\-ปฏิพทฺธํ อาทิสฺส วณฺณํปิ ภณติ อวณฺณํปิ ภณติ ฯเปฯ อกฺโกสติปิ, อาปตฺติ ทฺวินฺนํ ทุกฺกฏานํ ฯเปฯ.}

\proseitem{287}{อนาปตฺติ อตฺถ\-ปุเรกฺขารสฺส ธมฺม\-ปุเรกฺขารสฺส อนุสาสนิ\-ปุเรกฺขารสฺส อุมฺมตฺตกสฺส อาทิกมฺมิกสฺสาติ.}
\csromansectionrule

\csromanmatikaanchor{mtk.53}
\csromantocmarkref{subsection}{วินีตวตฺถุอุทฺทานคาถา}{mtk.53}
\bookmark[dest={mtk.53},level=2]{วินีตวตฺถุอุทฺทานคาถา}
\csromanheader{2}{\textbf{วินีตวตฺถุ}\-อุทฺทานคาถา}
\setlength{\csromangathapagewidth}{0pt}
\setlength{\csromangathawakwidth}{0pt}
\csromangathameasuregroup{โลหิตํ กกฺกสากิณฺณํ, \\ กจฺจิ สํสีทติ มคฺโค,}{\csromangathabat{โลหิตํ กกฺกสากิณฺณํ,}{ขรํ ทีฆญฺจ วาปิตํ.} \\ \csromangathabat{กจฺจิ สํสีทติ มคฺโค,}{สทฺธา ทาเนน กมฺมุนาติ.}}
\csromangathasetleft{โลหิตํ กกฺกสากิณฺณํ, \\ กจฺจิ สํสีทติ มคฺโค,}
\csromangathagroup{\csromangathastanza{\csromangathabat{โลหิตํ กกฺกสากิณฺณํ,}{ขรํ ทีฆญฺจ วาปิตํ.} \\ \csromangathabat{กจฺจิ สํสีทติ มคฺโค,}{สทฺธา ทาเนน กมฺมุนาติ.}}}

\csromanmatikaanchor{mtk.54}
\csromantocmarkref{subsection}{วินีตวตฺถุ}{mtk.54}
\bookmark[dest={mtk.54},level=2]{วินีตวตฺถุ}
\csromanheader{2}{\textbf{วินีตวตฺถุ}}
\proseitem{288}{เตน โข ปน สมเยน อญฺญตรา อิตฺถี นวรตฺตํ กมฺพลํ ปารุตา โหติ. อญฺญตโร ภิกฺขุ สารตฺโต ตํ อิตฺถิํ เอตทโวจ “โลหิตํ โข เต ภคินี”ติ. สา น ปฏิวิชานิ, อามายฺย นวรตฺโต กมฺพโลติ. ตสฺส กุกฺกุจฺจํ อโหสิ “ภควตา สิกฺขาปทํ ปญฺญตฺตํ, กจฺจิ นุ โข อหํ สํฆาทิเสสํ อาปตฺติํ อาปนฺโน”ติ. ภควโต เอตมตฺถํ อาโรเจสิ. อนาปตฺติ ภิกฺขุ สํฆาทิเสสสฺส, อาปตฺติ ทุกฺกฏสฺสาติ. (1)}

\prose{เตน โข ปน สมเยน อญฺญตรา อิตฺถี ขรกมฺพลํ ปารุตา โหติ. อญฺญตโร ภิกฺขุ สารตฺโต ตํ อิตฺถิํ เอตทโวจ “กกฺกสโลมํ โข เต ภคินี”ติ. สา น ปฏิวิชานิ, อามายฺย ขรกมฺพโลติ. ตสฺส}

\prose{\csromanfolio{191}เตน โข ปน สมเยน อญฺญตรา อิตฺถี นวาวุตํ กมฺพลํ ปารุตา โหติ. อญฺญตโร ภิกฺขุ สารตฺโต ตํ อิตฺถิํ เอตทโวจ “อากิณฺณโลมํ โข เต ภคินี”ติ. สา น ปฏิวิชานิ, อามายฺย นวาวุโต กมฺพโลติ. ตสฺส กุกฺกุจฺจํ อโหสิ ฯเปฯ. อนาปตฺติ ภิกฺขุ สํฆาทิเสสสฺส, อาปตฺติ ทุกฺกฏสฺสาติ. (3)}

\prose{เตน โข ปน สมเยน อญฺญตรา อิตฺถี ขรกมฺพลกํ ปารุตา โหติ. อญฺญตโร ภิกฺขุ สารตฺโต ตํ อิตฺถิํ เอตทโวจ “ขรโลมํ โข เต ภคินี”ติ. สา น ปฏิวิชานิ, อามายฺย ขรกมฺพล\-โกติ. ตสฺส กุกฺกุจฺจํ อโหสิ ฯเปฯ. อนาปตฺติ ภิกฺขุ สํฆาทิเสสสฺส, อาปตฺติ ทุกฺกฏสฺสาติ. (4)}

\prose{เตน โข ปน สมเยน อญฺญตรา อิตฺถี ปาวารํ\csromansharedfootnote{c79972b951f939e6}{ทีฆปาวารํ (สฺยา)} ปารุตา โหติ. อญฺญตโร ภิกฺขุ สารตฺโต ตํ อิตฺถิํ เอตทโวจ “ทีฆโลมํ โข เต ภคินี”ติ. สา น ปฏิวิชานิ, อามายฺย ปาวาโร\csromansharedfootnote{c10714ea2da6d98a}{ทีฆปาวาโร (สฺยา)} ติ. ตสฺส กุกฺกุจฺจํ อโหสิ ฯเปฯ. อนาปตฺติ ภิกฺขุ สํฆาทิเสสสฺส, อาปตฺติ ทุกฺกฏสฺสาติ. (5)}

\proseitem{289}{เตน โข ปน สมเยน อญฺญตรา อิตฺถี เขตฺตํ วปาเปตฺวา อาคจฺฉติ. อญฺญตโร ภิกฺขุ สารตฺโต ตํ อิตฺถิํ เอตทโวจ “วาปิตํ โข เต ภคินี”ติ. สา น ปฏิวิชานิ, อามายฺย โน จ โข ปฏิวุตฺตนฺติ. ตสฺส กุกฺกุจฺจํ อโหสิ ฯเปฯ. อนาปตฺติ ภิกฺขุ สํฆาทิเสสสฺส, อาปตฺติ ทุกฺกฏสฺสาติ. (6)}

\prose{เตน โข ปน สมเยน อญฺญตโร ภิกฺขุ ปริพฺพาชิกํ ปฏิปเถ ปสฺสิตฺวา สารตฺโต ตํ ปริพฺพาชิกํ เอตทโวจ “กจฺจิ ภคินิ มคฺโค สํสีทตี”ติ. สา น ปฏิวิชานิ, อาม ภิกฺขุ ปฏิปชฺชิสฺสสีติ. ตสฺส กุกฺกุจฺจํ อโหสิ ฯเปฯ. อนาปตฺติ ภิกฺขุ สํฆาทิเสสสฺส, อาปตฺติ ถุลฺลจฺจยสฺสาติ. (7)}

\csromanmatikaanchor{mtk.55}
\csromanmatikaanchor{mtk.56}
\csromantocmarknopage{section}{4. อตฺตกามปาริจริยสิกฺขาปท}{mtk.55}
\bookmark[dest={mtk.55},level=1]{4. อตฺตกามปาริจริยสิกฺขาปท}
\csromantocmarkref{subsection}{อุทายิภิกฺขุวตฺถุ}{mtk.56}
\bookmark[dest={mtk.56},level=2]{อุทายิภิกฺขุวตฺถุ}
\prose{เตน โข ปน สมเยน อญฺญตโร ภิกฺขุ สารตฺโต อญฺญตรํ อิตฺถิํ เอตทโวจ “สทฺธาสิ ตฺวํ ภคินิ, อปิ จ ยํ สามิกสฺส เทสิ ตํ นามฺหากํ \csromanfolio{192}เทสี”ติ. กิํ ภนฺเตติ. เมถุน\-ธมฺมนฺติ. ตสฺส กุกฺกุจฺจํ อโหสิ ฯเปฯ. อาปตฺติํ ตฺวํ ภิกฺขุ อาปนฺโน สํฆาทิเสสนฺติ. (8)}

\prose{เตน โข ปน สมเยน อญฺญตโร ภิกฺขุ สารตฺโต อญฺญตรํ อิตฺถิํ เอตทโวจ “สทฺธาสิ ตฺวํ ภคินิ, อปิ จ ยํ อคฺคทานํ ตํ นามฺหากํ เทสี”ติ. กิํ ภนฺเต อคฺคทานนฺติ. เมถุน\-ธมฺมนฺติ. ตสฺส กุกฺกุจฺจํ อโหสิ ฯเปฯ. อาปตฺติํ ตฺวํ ภิกฺขุ อาปนฺโน สํฆาทิเสสนฺติ. (9)}

\prosewithcloserruled{เตน โข ปน สมเยน อญฺญตรา อิตฺถี กมฺมํ กโรติ. อญฺญตโร ภิกฺขุ สารตฺโต ตํ อิตฺถิํ เอตทโวจ “ติฏฺฐ ภคินิ อหํ กริสฺสามี”ติ ฯเปฯ “นิสีท ภคินิ อหํ กริสฺสามี”ติ ฯเปฯ “นิปชฺช ภคินิ อหํ กริสฺสามี”ติ. สา น ปฏิวิชานิ. ตสฺส กุกฺกุจฺจํ อโหสิ ฯเปฯ. อนาปตฺติ ภิกฺขุ สํฆาทิเสสสฺส, อาปตฺติ ทุกฺกฏสฺสาติ. \mbox{(10-12)}}{ทุฏฺฐุลฺลวาจา\-สิกฺขาปทํ นิฏฺฐิตํ ตติยํ.}
\csromanheader{1}{4. \textbf{อตฺตกามปาริจริย}\-\textbf{สิกฺขาปท}}
\proseitem{290}{เตน สมเยน พุทฺโธ ภควา สาวตฺถิยํ วิหรติ เชตวเน อนาถ\-ปิณฺฑิกสฺส อาราเม. เตน โข ปน สมเยน อายสฺมา อุทายี สาวตฺถิยํ กุลูปโก โหติ, พหุกานิ กุลานิ อุปสงฺกมติ. เตน โข ปน สมเยน อญฺญตรา อิตฺถี มตปติกา อภิรูปา โหติ ทสฺสนียา ปาสาทิกา. อถ โข อายสฺมา อุทายี ปุพฺพณฺหสมยํ นิวาเสตฺวา ปตฺต\-จีวรมา\-ทาย เยน ตสฺสา อิตฺถิยา นิเวสนํ เตนุปสงฺกมิ, อุปสงฺกมิตฺวา ปญฺญตฺเต อาสเน นิสีทิ. อถ โข สา อิตฺถี เยนายสฺมา อุทายี เตนุปสงฺกมิ, อุปสงฺกมิตฺวา อายสฺมนฺตํ อุทายิํ อภิวาเทตฺวา เอกมนฺตํ นิสีทิ, เอกมนฺตํ นิสินฺนํ โข ตํ อิตฺถิํ อายสฺมา อุทายี ธมฺมิยา กถาย สนฺทสฺเสสิ สมาทเปสิ สมุตฺเตเชสิ สํปหํเสสิ. อถ โข สา อิตฺถี อายสฺมตา อุทายินา ธมฺมิยา กถาย สนฺทสฺสิตา สมาทปิตา สมุตฺเตชิตา สํปหํสิตา อายสฺมนฺตํ อุทายิํ เอตทโวจ “วเทยฺยาถ ภนฺเต เยน อตฺโถ, ปฏิพลา มยํ อยฺยสฺส ทาตุํ, ยทิทํ จีวร\-ปิณฺฑปาต\-เสนาสน\-คิลานปฺปจฺจย\-เภสชฺช\-ปริกฺขารนฺ”ติ.}

\prose{\csromanfolio{193}น โข เต ภคินิ อมฺหากํ ทุลฺลภา ยทิทํ จีวร\-ปิณฺฑปาต\-เสนาสน\-คิลานปฺปจฺจย\-เภสชฺช\-ปริกฺขารา, อปิ จ โย อมฺหากํ ทุลฺลโภ, ตํ เทหีติ. กิํ ภนฺเตติ. เมถุน\-ธมฺมนฺติ. อตฺโถ ภนฺเตติ. อตฺโถ ภคินีติ. เอหิ ภนฺเตติ โอวรกํ ปวิสิตฺวา สาฏกํ นิกฺขิปิตฺวา มญฺจเก อุตฺตานา นิปชฺชิ. อถ โข อายสฺมา อุทายี เยน สา อิตฺถี เตนุปสงฺกมิ, อุปสงฺกมิตฺวา “โก อิมํ วสลํ ทุคฺคนฺธํ อามสิสฺสตี”ติ นิฏฺฐุภิตฺวา ปกฺกามิ. อถ โข สา อิตฺถี อุชฺฌายติ ขิยฺยติ วิปาเจติ “อลชฺชิโน อิเม สมณา สกฺยปุตฺติยา ทุสฺสีลา มุสาวาทิโน, อิเม หิ นาม ธมฺมจาริโน สมจาริโน พฺรหฺมจาริโน สจฺจวาทิโน สีลวนฺโต กลฺยาณธมฺมา ปฏิชานิสฺสนฺติ. นตฺถิ อิเมสํ สามญฺญํ, นตฺถิ อิเมสํ พฺรหฺมญฺญํ, นฏฺฐํ อิเมสํ สามญฺญํ, นฏฺฐํ อิเมสํ พฺรหฺมญฺญํ, กุโต อิเมสํ สามญฺญํ, กุโต อิเมสํ พฺรหฺมญฺญํ, อปคตา อิเม สามญฺญา, อปคตา อิเม พฺรหฺมญฺญา. กถํ หิ นาม สมโณ อุทายี มํ สามํ เมถุนธมฺมํ ยาจิตฺวา ‘โก อิมํ วสลํ ทุคฺคนฺธํ อามสิสฺสตี’ติ นิฏฺฐุภิตฺวา ปกฺกมิสฺสติ, กิํ เม ปาปกํ, กิํ เม ทุคฺคนฺธํ, กสฺสาหํ เกน หายามี”ติ. อญฺญาปิ อิตฺถิโย อุชฺฌายนฺติ ขิยฺยนฺติ วิปาเจนฺติ “อลชฺชิโน อิเม สมณา สกฺยปุตฺติยา ทุสฺสีลา มุสาวาทิโน ฯเปฯ. กถํ หิ นาม สมโณ อุทายี อิมิสฺสา สามํ เมถุนธมฺมํ ยาจิตฺวา ‘โก อิมํ วสลํ ทุคฺคนฺธํ อามสิสฺสตี’ติ นิฏฺฐุภิตฺวา ปกฺกมิสฺสติ, กิํ อิมิสฺสา ปาปกํ, กิํ อิมิสฺสา ทุคฺคนฺธํ, กสฺสายํ เกน หายตี”ติ. อสฺโสสุํ โข ภิกฺขู ตาสํ อิตฺถีนํ อุชฺฌายนฺตีนํ ขิยฺยนฺตีนํ วิปาเจนฺตีนํ. เย เต ภิกฺขู อปฺปิจฺฉา ฯเปฯ เต อุชฺฌายนฺติ ขิยฺยนฺติ วิปาเจนฺติ “กถํ หิ นาม อายสฺมา อุทายี มาตุคามสฺส สนฺติเก อตฺต\-กาม\-ปาริจริยาย วณฺณํ ภาสิสฺสตี”ติ.}

\csromanmatikaanchor{mtk.57}
\csromanmatikaanchor{mtk.58}
\csromantocmarkref{subsection}{ปญฺญตฺติ}{mtk.57}
\bookmark[dest={mtk.57},level=2]{ปญฺญตฺติ}
\csromantocmarkref{subsection}{สิกฺขาปทวิภงฺค, ปทภาชนีย}{mtk.58}
\bookmark[dest={mtk.58},level=2]{สิกฺขาปทวิภงฺค, ปทภาชนีย}
\prose{อถ โข เต ภิกฺขู อายสฺมนฺตํ อุทายิํ อเนก\-ปริยาเยน วิครหิตฺวา ภควโต เอตมตฺถํ อาโรเจสุํ. อถ โข ภควา เอตสฺมิํ นิทาเน เอตสฺมิํ ปกรเณ ภิกฺขุสํฆํ สนฺนิปาตาเปตฺวา อายสฺมนฺตํ อุทายิํ ปฏิปุจฺฉิ “สจฺจํ กิร ตฺวํ อุทายิ มาตุคามสฺส สนฺติเก อตฺต\-กาม\-ปาริจริยาย วณฺณํ ภาสสี’ติ. สจฺจํ ภควาติ. วิครหิ พุทฺโธ ภควา “อนนุจฺฉวิกํ โมฆปุริส อนนุโลมิกํ อปฺปติรูปํ อสฺสามณกํ อกปฺปิยํ อกรณียํ, กถํ หิ นาม ตฺวํ โมฆปุริส มาตุคามสฺส สนฺติเก อตฺต\-กาม\-ปาริจริยาย วณฺณํ ภาสิสฺสสิ. นนุ \csromanfolio{194}มยา โมฆปุริส อเนก\-ปริยาเยน วิราคาย ธมฺโม เทสิโต, โน สราคาย ฯเปฯ กาม\-ปริฬาหานํ วูปสโม อกฺขาโต. เนตํ โมฆปุริส อปฺปสนฺนานํ วา ปสาทาย ฯเปฯ. เอวญฺจ ปน ภิกฺขเว อิมํ \textbf{สิกฺขาปทํ อุทฺทิเสยฺยาถ\csromandash{}}}

\proseitem{291}{\textbf{“โย ปน ภิกฺขุ โอติณฺโณ วิปริณเตน จิตฺเตน มาตุคามสฺส สนฺติเก อตฺต}\-\textbf{กาม}\-\textbf{ปาริจริยาย วณฺณํ ภาเสยฺย ‘เอตทคฺคํ ภคินิ ปาริจริยานํ ยา มาทิสํ สีลวนฺตํ กลฺยาณธมฺมํ พฺรหฺมจาริํ เอเตน ธมฺเมน ปริจเรยฺยา’ติ เมถุนุ}\-\textbf{ปสํหิเตน, สํฆาทิเสโส”}ติ.}

\proseitem{292}{\textbf{โย ปนา}ติ โย ยาทิโส ฯเปฯ. \textbf{ภิกฺขู}ติ ฯเปฯ อยํ อิมสฺมิํ อตฺเถ อธิปฺเปโต “ภิกฺขู”ติ.}

\prose{\textbf{โอติณฺโณ} นาม สารตฺโต อเปกฺขวา ปฏิพทฺธ\-จิตฺโต.}

\prose{\textbf{วิปริณตนฺ}ติ รตฺตํปิ จิตฺตํ วิปริณตํ, ทุฏฺฐํปิ จิตฺตํ วิปริณตํ, มูฬฺหํปิ จิตฺตํ วิปริณตํ. อปิ จ รตฺตํ จิตฺตํ อิมสฺมิํ อตฺเถ อธิปฺเปตํ “วิปริณตนฺ”ติ.}

\prose{\textbf{มาตุคาโม} นาม มนุสฺสิตฺถี, น ยกฺขี, น เปตี, น ติรจฺฉานคตา, วิญฺญู ปฏิพลา สุภาสิต\-ทุพฺภาสิตํ ทุฏฺฐุลฺลา\-ทุฏฺฐุลฺลํ อาชานิตุํ.}

\prose{\textbf{มาตุคามสฺส สนฺติเก}ติ มาตุคามสฺส สามนฺตา มาตุคามสฺส อวิทูเร.}

\prose{\textbf{อตฺตกามนฺ}ติ อตฺตโน กามํ อตฺตโน เหตุํ อตฺตโน อธิปฺปายํ อตฺตโน ปาริจริยํ.}

\prose{\textbf{เอตทคฺคนฺ}ติ เอตํ อคฺคํ เอตํ เสฏฺฐํ เอตํ โมกฺขํ เอตํ อุตฺตมํ เอตํ ปวรํ.}

\prose{\textbf{ยา}ติ ขตฺติยี\csromansharedfootnote{e4fd77aeb30134d7}{ขตฺติยา (สฺยา)} วา พฺรหฺมณี วา เวสฺสี วา สุทฺที วา.}

\prose{\textbf{มาทิสนฺ}ติ ขตฺติยํ วา พฺราหฺมณํ วา เวสฺสํ วา สุทฺทํ วา.}

\prose{\textbf{สีลวนฺตนฺ}ติ ปาณาติปาตา ปฏิวิรตํ อทินฺนาทานา ปฏิวิรตํ มุสาวาทา ปฏิวิรตํ.}

\prose{\textbf{พฺรหฺมจารินฺ}ติ เมถุนธมฺมา ปฏิวิรตํ.}

\prose{\textbf{กลฺยาณธมฺโม} นาม เตน จ สีเลน เตน จ พฺรหฺมจริเยน กลฺยาณธมฺโม โหติ.}

\prose{\csromanfolio{195}\textbf{เอเตน ธมฺเมนา}ติ เมถุน\-ธมฺเมน.}

\prose{\textbf{ปริจเรยฺยา}ติ อภิรเมยฺย.}

\prose{\textbf{เมถุนุปสํหิเตนา}ติ เมถุนธมฺม\-ปฺปฏิสํยุตฺเตน.}

\prose{\textbf{สํฆาทิเสโส}ติ ฯเปฯ เตนปิ วุจฺจติ “สํฆาทิเสโส”ติ.}

\proseitem{293}{อิตฺถี จ โหติ อิตฺถิสญฺญี สารตฺโต จ, ภิกฺขุ จ นํ อิตฺถิยา สนฺติเก อตฺต\-กาม\-ปาริจริยาย วณฺณํ ภาสติ, อาปตฺติ สํฆาทิเสสสฺส.}

\prose{อิตฺถี จ โหติ เวมติโก ฯเปฯ ปณฺฑกสญฺญี. ปุริสสญฺญี. ติรจฺฉานคต\-สญฺญี สารตฺโต จ, ภิกฺขุ จ นํ อิตฺถิยา สนฺติเก อตฺต\-กาม\-ปาริจริยาย วณฺณํ ภาสติ, อาปตฺติ ถุลฺลจฺจยสฺส.}

\prose{ปณฺฑโก จ โหติ ปณฺฑกสญฺญี สารตฺโต จ, ภิกฺขุ จ นํ ปณฺฑกสฺส สนฺติเก อตฺต\-กาม\-ปาริจริยาย วณฺณํ ภาสติ, อาปตฺติ ถุลฺลจฺจยสฺส.}

\prose{ปณฺฑโก จ โหติ เวมติโก ฯเปฯ ปุริสสญฺญี. ติรจฺฉานคต\-สญฺญี. อิตฺถิสญฺญี สารตฺโต จ, ภิกฺขุ จ นํ ปณฺฑกสฺส สนฺติเก อตฺต\-กาม\-ปาริจริยาย วณฺณํ ภาสติ, อาปตฺติ ทุกฺกฏสฺส.}

\prose{ปุริโส จ โหติ ฯเปฯ. ติรจฺฉานคโต จ โหติ ติรจฺฉานคต\-สญฺญี ฯเปฯ. เวมติโก. อิตฺถิสญฺญี. ปณฺฑกสญฺญี. ปุริสสญฺญี สารตฺโต จ, ภิกฺขุ จ นํ ติรจฺฉาน\-คตสฺส สนฺติเก อตฺต\-กาม\-ปาริจริยาย วณฺณํ ภาสติ, อาปตฺติ ทุกฺกฏสฺส.}

\prose{เทฺว อิตฺถิโย ทฺวินฺนํ อิตฺถีนํ อิตฺถิสญฺญี สารตฺโต จ, ภิกฺขุ จ นํ ทฺวินฺนํ อิตฺถีนํ สนฺติเก อตฺต\-กาม\-ปาริจริยาย วณฺณํ ภาสติ, อาปตฺติ ทฺวินฺนํ สํฆา\-ทิ\-เสสานํ ฯเปฯ.}

\prose{อิตฺถี จ ปณฺฑโก จ อุภินฺนํ อิตฺถิสญฺญี สารตฺโต จ, ภิกฺขุ จ นํ อุภินฺนํ สนฺติเก อตฺต\-กาม\-ปาริจริยาย วณฺณํ ภาสติ, อาปตฺติ สํฆาทิเสเสน ทุกฺกฏสฺส ฯเปฯ.}

\proseitem{294}{อนาปตฺติ “จีวร ปิณฺฑปาต เสนาสน คิลานปฺปจฺจย เภสชฺช\-ปริกฺขาเรน อุปฏฺฐหา”ติ ภณติ, อุมฺมตฺตกสฺส อาทิกมฺมิกสฺสาติ.}

\csromanmatikaanchor{mtk.59}
\csromantocmarkref{subsection}{วินีตวตฺถุอุทฺทานคาถา}{mtk.59}
\bookmark[dest={mtk.59},level=2]{วินีตวตฺถุอุทฺทานคาถา}
\csromanheader{2}{\textbf{วินีตวตฺถุ}\-อุทฺทานคาถา}
\setlength{\csromangathapagewidth}{0pt}
\setlength{\csromangathawakwidth}{0pt}
\csromangathameasuregroup{กถํ วญฺฌา ลเภ ปุตฺตํ, \\ กิํ ทชฺชํ เกนุปฏฺเฐยฺยํ,}{\csromangathabat{กถํ วญฺฌา ลเภ ปุตฺตํ,}{ปิยา จ สุภคา สิยํ.} \\ \csromangathabat{กิํ ทชฺชํ เกนุปฏฺเฐยฺยํ,}{กถํ คจฺเฉยฺยํ สุคฺคตินฺติ.}}
\csromangathasetleft{กถํ วญฺฌา ลเภ ปุตฺตํ, \\ กิํ ทชฺชํ เกนุปฏฺเฐยฺยํ,}
\csromangathagroup{\csromangathastanza{\csromanfolio{196}\csromangathabat{กถํ วญฺฌา ลเภ ปุตฺตํ,}{ปิยา จ สุภคา สิยํ.} \\ \csromangathabat{กิํ ทชฺชํ เกนุปฏฺเฐยฺยํ,}{กถํ คจฺเฉยฺยํ สุคฺคตินฺติ.}}}

\csromanmatikaanchor{mtk.60}
\csromantocmarkref{subsection}{วินีตวตฺถุ}{mtk.60}
\bookmark[dest={mtk.60},level=2]{วินีตวตฺถุ}
\csromanheader{2}{\textbf{วินีตวตฺถุ}}
\proseitem{295}{เตน โข ปน สมเยน อญฺญตรา วญฺฌา อิตฺถี กุลูปกํ ภิกฺขุํ เอตทโวจ “กถาหํ ภนฺเต วิชาเยยฺยนฺ”ติ. เตนหิ ภคินิ อคฺคทานํ เทหีติ. กิํ ภนฺเต อคฺคทานนฺติ. เมถุน\-ธมฺมนฺติ. ตสฺส กุกฺกุจฺจํ อโหสิ ฯเปฯ. อาปตฺติํ ตฺวํ ภิกฺขุ อาปนฺโน สํฆาทิเสสนฺติ. (1)}

\prose{เตน โข ปน สมเยน อญฺญตรา วิชายินี อิตฺถี กุลูปกํ ภิกฺขุํ เอตทโวจ “กถาหํ ภนฺเต ปุตฺตํ ลเภยฺยนฺ”ติ. เตนหิ ภคินิ อคฺคทานํ เทหีติ. กิํ ภนฺเต อคฺคทานนฺติ. เมถุน\-ธมฺมนฺติ. ตสฺส กุกฺกุจฺจํ อโหสิ ฯเปฯ. อาปตฺติํ ตฺวํ ภิกฺขุ อาปนฺโน สํฆาทิเสสนฺติ. (2)}

\prose{เตน โข ปน สมเยน อญฺญตรา อิตฺถี กุลูปกํ ภิกฺขุํ เอตทโวจ “กถาหํ ภนฺเต สามิกสฺส ปิยา อสฺสนฺ”ติ. เตนหิ ภคินิ อคฺคทานํ เทหีติ. กิํ ภนฺเต อคฺคทานนฺติ. เมถุน\-ธมฺมนฺติ. ตสฺส กุกฺกุจฺจํ อโหสิ ฯเปฯ. อาปตฺติํ ตฺวํ ภิกฺขุ อาปนฺโน สํฆาทิเสสนฺติ. (3)}

\prose{เตน โข ปน สมเยน อญฺญตรา อิตฺถี กุลูปกํ ภิกฺขุํ เอตทโวจ “กถาหํ ภนฺเต สุภคา อสฺสนฺ”ติ. เตนหิ ภคินิ อคฺคทานํ เทหีติ. กิํ ภนฺเต อคฺคทานนฺติ. เมถุน\-ธมฺมนฺติ. ตสฺส กุกฺกุจฺจํ อโหสิ ฯเปฯ. อาปตฺติํ ตฺวํ ภิกฺขุ อาปนฺโน สํฆาทิเสสนฺติ. (4)}

\prose{เตน โข ปน สมเยน อญฺญตรา อิตฺถี กุลูปกํ ภิกฺขุํ เอตทโวจ “กฺยาหํ ภนฺเต อยฺยสฺส ทชฺชามี”ติ. อคฺคทานํ ภคินีติ. กิํ ภนฺเต อคฺคทานนฺติ. เมถุน\-ธมฺมนฺติ. ตสฺส กุกฺกุจฺจํ อโหสิ ฯเปฯ. อาปตฺติํ ตฺวํ ภิกฺขุ อาปนฺโน สํฆาทิเสสนฺติ. (5)}

\prose{เตน โข ปน สมเยน อญฺญตรา อิตฺถี กุลูปกํ ภิกฺขุํ เอตทโวจ “เกนาหํ ภนฺเต อยฺยํ อุปฏฺเฐมี”ติ. อคฺคาทาเนน ภคินีติ. กิํ ภนฺเต อคฺคทานนฺติ. เมถุน\-ธมฺมนฺติ. ตสฺส กุกฺกุจฺจํ อโหสิ ฯเปฯ. อาปตฺติํ ตฺวํ ภิกฺขุ อาปนฺโน สํฆาทิเสสนฺติ. (6)}

\csromanmatikaanchor{mtk.61}
\csromanmatikaanchor{mtk.62}
\csromantocmarknopage{section}{5. สญฺจริตฺตสิกฺขาปท}{mtk.61}
\bookmark[dest={mtk.61},level=1]{5. สญฺจริตฺตสิกฺขาปท}
\csromantocmarkref{subsection}{อุทายิภิกฺขุวตฺถุ}{mtk.62}
\bookmark[dest={mtk.62},level=2]{อุทายิภิกฺขุวตฺถุ}
\prosewithcloserruled{\csromanfolio{197}เตน โข ปน สมเยน อญฺญตรา อิตฺถี กุลูปกํ ภิกฺขุํ เอตทโวจ “กถาหํ ภนฺเต สุคติํ คจฺเฉยฺยนฺ”ติ. เตน หิ ภคินิ อคฺคทานํ เทหีติ. กิํ ภนฺเต อคฺคทานนฺติ. เมถุน\-ธมฺมนฺติ. ตสฺส กุกฺกุจฺจํ อโหสิ ฯเปฯ. อาปตฺติํ ตฺวํ ภิกฺขุ อาปนฺโน สํฆาทิเสสนฺติ. (7)}{อตฺตกามปาริจริย\-สิกฺขาปทํ นิฏฺฐิตํ จตุตฺถํ.}
\csromanheader{1}{5. \textbf{สญฺจริตฺต}\-\textbf{สิกฺขาปท}}
\proseitem{296}{เตน สมเยน พุทฺโธ ภควา สาวตฺถิยํ วิหรติ เชตวเน อนาถ\-ปิณฺฑิกสฺส อาราเม. เตน โข ปน สมเยน อายสฺมา อุทายี สาวตฺถิยํ กุลูปโก โหติ, พหุกานิ กุลานิ อุปสงฺกมติ. ยตฺถ ปสฺสติ กุมารกํ วา อปชาปติกํ กุมาริกํ วา อปติกํ, กุมารกสฺส มาตาปิตูนํ สนฺติเก กุมาริกาย วณฺณํ ภณติ “อมุกสฺส กุลสฺส กุมาริกา อภิรูปา ทสฺสนียา ปาสาทิกา ปณฺฑิตา พฺยตฺตา เมธาวินี ทกฺขา อนลสา, ฉนฺนา สา กุมาริกา อิมสฺส กุมารกสฺสา”ติ. เต เอวํ วทนฺติ “เอเต โข ภนฺเต อเมฺห น ชานนฺติ ‘เก วา อิเม กสฺส วา’ติ, สเจ ภนฺเต อยฺโย ทาเปยฺย อาเนยฺยาม มยํ ตํ กุมาริกํ อิมสฺส กุมารกสฺสา”ติ. กุมาริกาย มาตาปิตูนํ สนฺติเก กุมารกสฺส วณฺณํ ภณติ “อมุกสฺส กุลสฺส กุมารโก อภิรูโป ทสฺสนีโย ปาสาทิโก ปณฺฑิโต พฺยตฺโต เมธาวี ทกฺโข อนลโส, ฉนฺนายํ กุมาริกา ตสฺส กุมารกสฺสา”ติ\csromansharedfootnote{314069fc0a83c1cc}{ฉนฺโน โส กุมารโก อิมิสฺสา กุมาริกายาติ (สฺยา)}. เต เอวํ วทนฺติ “เอเต โข ภนฺเต อเมฺห น ชานนฺติ ‘เก วา อิเม กสฺส วา’ติ, กสฺมิํ วิย กุมาริกาย วตฺถุํ, สเจ ภนฺเต อยฺโย ยาจาเปยฺย, ทชฺเชยฺยาม มยํ อิมํ กุมาริกํ ตสฺส กุมารกสฺสา”ติ. เอเตเนว อุปาเยน อาวาหานิปิ การาเปติ, วิวาหานิปิ การาเปติ, วาเรยฺยานิปิ การาเปติ.}

\proseitem{297}{เตน โข ปน สมเยน อญฺญตริสฺสา ปุราณคณกิยา ธีตา อภิรูปา โหติ ทสฺสนียา ปาสาทิกา. ติโรคามกา อาชีวกสาวกา อาคนฺตฺวา ตํ คณกิํ เอตทโวจุํ “เทหายฺเย อิมํ กุมาริกํ อมฺหากํ กุมารกสฺสา”ติ. สา เอวมาห “อหํ ขฺวยฺโย\csromansharedfootnote{de5fdd7fab7397cb}{ขฺวยฺยา (สฺยา) ขฺวายฺโย (ก)} \csromanfolio{198}ตุเมฺห น ชานามิ ‘เก วา อิเม กสฺส วา’ติ, อยญฺจ เม เอกธีติกา, ติโรคาโม จ คนฺตพฺโพ, นาหํ ทสฺสามี”ติ. มนุสฺสา เต อาชีวกสาวเก เอตทโวจุํ “กิสฺส ตุเมฺห อยฺโย อาคตตฺถา”ติ. อิธ มยํ อยฺโย อมุกํ นาม คณกิํ ธีตรํ ยาจิมฺหา อมฺหากํ กุมารกสฺส, สา เอว มาห “อหํ ขฺวยฺโย ตุเมฺห น ชานามิ ‘เก วา อิเม กสฺส วา’ติ, อยญฺจ เม เอกธีติกา, ติโรคาโม จ คนฺตพฺโพ, นาหํ ทสฺสามี”ติ. กิสฺส ตุเมฺห อยฺโย ตํ คณกิํ ธีตรํ ยาจิตฺถ, นนุ อยฺโย อุทายี วตฺตพฺโพ, อยฺโย อุทายี ทาเปสฺสตีติ.}

\prose{อถ โข เต อาชีวกสาวกา เยนายสฺมา อุทายี เตนุ\-ปสงฺกมิํสุ, อุปสงฺกมิตฺวา อายสฺมนฺตํ อุทายิํ เอตทโวจุํ “อิธ มยํ ภนฺเต อมุกํ นาม คณกิํ ธีตรํ ยาจิมฺหา อมฺหากํ กุมารกสฺส, สา เอว มาห “อหํ ขฺวยฺโย ตุเมฺห น ชานามิ ‘เก วา อิเม กสฺส วา’ติ, อยญฺจ เม เอกธีติกา, ติโรคาโม จ คนฺตพฺโพ, นาหํ ทสฺสามี”ติ, สาธุ ภนฺเต, อยฺโย ตํ คณกิํ ธีตรํ ทาเปตุ อมฺหากํ กุมารกสฺสา”ติ. อถ โข อายสฺมา อุทายี เยน สา คณกี เตนุปสงฺกมิ, อุปสงฺกมิตฺวา ตํ คณกิํ เอตทโวจ “กิสฺสิเมสํ ธีตรํ น เทสี”ติ. อหํ ขฺวยฺย อิเม น ชานามิ “เก วา อิเม กสฺส วา”ติ, อยญฺจ เม เอกธีติกา, ติโรคาโม จ คนฺตพฺโพ, นาหํ ทสฺสามีติ. เทหิเมสํ อหํ อิเม ชานามีติ. สเจ ภนฺเต อยฺโย ชานาติ ทสฺสามีติ. อถ โข สา คณกี เตสํ อาชีวก\-สาวกานํ ธีตรํ อทาสิ. อถ โข เต อาชีวกสาวกา ตํ กุมาริกํ เนตฺวา มาสํ เยว สุณิสโภเคน ภุญฺชิํสุ. ตโต อปเรน ทาสิโภเคน ภุญฺชนฺติ.}

\prose{อถ โข สา กุมาริกา มาตุยา สนฺติเก ทูตํ ปาเหสิ “อหมฺหิ ทุคฺคตา ทุกฺขิตา น สุขํ ลภามิ, มาสํ เยว มํ สุณิสโภเคน ภุญฺชิํสุ, ตโต อปเรน ทาสิโภเคน ภุญฺชนฺติ, อาคจฺฉตุ เม มาตา มํ เนสฺสตู”ติ. อถ โข สา คณกี เยน เต อาชีวกสาวกา เตนุปสงฺกมิ, อุปสงฺกมิตฺวา เต อาชีวกสาวเก เอตทโวจ “มายฺโย อิมํ กุมาริกํ ทาสิโภเคน ภุญฺชิตฺถ, สุณิสโภเคน อิมํ กุมาริกํ ภุญฺชถา”ติ. เต เอวมาหํสุ “นตฺถมฺหากํ ตยา สทฺธิํ อาหารูปหาโร, สมเณน สทฺธิํ อมฺหากํ อาหารูป \csromanfolio{199}หาโร, คจฺฉ ตฺวํ, น มยํ ตํ ชานามา”ติ. อถ โข สา คณกี เตหิ อาชีวก\-สาวเกหิ อปสาทิตา ปุนเทว สาวตฺถิํ ปจฺจาคญฺฉิ. ทุติยมฺปิ โข สา กุมาริกา มาตุยา สนฺติเก ทูตํ ปาเหสิ “อหมฺหิ ทุคฺคตา ทุกฺขิตา น สุขํ ลภามิ, มาสํ เยว มํ สุณิสโภเคน ภุญฺชิํสุ, ตโต อปเรน ทาสิโภเคน ภุญฺชนฺติ, อาคจฺฉตุ เม มาตา มํ เนสฺสตู”ติ. อถ โข สา คณกี เยนายสฺมา อุทายี เตนุปสงฺกมิ, อุปสงฺกมิตฺวา อายสฺมนฺตํ อุทายิํ เอตทโวจ “สา กิร ภนฺเต กุมาริกา ทุคฺคตา ทุกฺขิตา น สุขํ ลภติ, มาสํ เยว นํ สุณิสโภเคน ภุญฺชิํสุ, ตโต อปเรน ทาสิโภเคน ภุญฺชนฺติ, วเทยฺยาถ ภนฺเต มายฺโย อิมํ กุมาริกํ ทาสิโภเคน ภุญฺชิตฺถ, สุณิสโภเคน อิมํ กุมาริกํ ภุญฺชถา”ติ.}

\prose{อถ โข อายสฺมา อุทายี เยน เต อาชีวกสาวกา เตนุปสงฺกมิ, อุปสงฺกมิตฺวา เต อาชีวกสาวเก เอตทโวจ “มายฺโย อิมํ กุมาริกํ ทาสิโภเคน ภุญฺชิตฺถ, สุณิสโภเคน อิมํ กุมาริกํ ภุญฺชถา”ติ. เต เอวมาหํสุ “นตฺถมฺหากํ ตยา สทฺธิํ อาหารูป หาโร, คณกิยา สทฺธิํ อมฺหากํ อาหารูปหาโร, สมเณน ภวิตพฺพํ อพฺยาวเฏน, สมโณ อสฺส สุสมโณ, คจฺฉ ตฺวํ, น มยํ ตํ ชานามา”ติ. อถ โข อายสฺมา อุทายี เตหิ อาชีวก\-สาวเกหิ อปสาทิโต ปุนเทว สาวตฺถิํ ปจฺจาคญฺฉิ. ตติยมฺปิ โข สา กุมาริกา มาตุยา สนฺติเก ทูตํ ปาเหสิ “อหมฺหิ ทุคฺคตา ทุกฺขิตา น สุขํ ลภามิ, มาสํ เยว มํ สุณิสโภเคน ภุญฺชิํสุ, ตโต อปเรน ทาสิโภเคน ภุญฺชินฺติ, อาคจฺฉตุ เม มาตา มํ เนสฺสตู”ติ. ทุติยมฺปิ โข สา คณกี เยนายสฺมา อุทายี เตนุปสงฺกมิ, อุปสงฺกมิตฺวา อายสฺมนฺตํ อุทายิํ เอตทโวจ “สา กิร ภนฺเต กุมาริกา ทุคฺคตา ทุกฺขิตา น สุขํ ลภติ, มาสํ เยว นํ สุณิสโภเคน ภุญฺชิํสุ, ตโต อปเรน ทาสิโภเคน ภุญฺชนฺติ, วเทยฺยาถ ภนฺเต มายฺโย อิมํ กุมาริกํ ทาสิโภเคน ภุญฺชิตฺถ, สุณิสโภเคน อิมํ กุมาริกํ ภุญฺชิถา”ติ. ปฐมํปาหํ เตหิ อาชีวก\-สาวเกหิ อปสาทิโต, คจฺฉ ตฺวํ, นาหํ คมิสฺสามีติ.}

\csromanmatikaanchor{mtk.63}
\csromantocmarkref{subsection}{ปฐมปญฺญตฺติ}{mtk.63}
\bookmark[dest={mtk.63},level=2]{ปฐมปญฺญตฺติ}
\proseitem{298}{อถ โข สา คณกี อุชฺฌายติ ขิยฺยติ วิปาเจติ “เอวํ ทุคฺคโต โหตุ อยฺโย อุทายี, เอวํ ทุกฺขิโต โหตุ อยฺโย อุทายี, \csromanfolio{200}เอวํ มา สุขํ ลภตุ อยฺโย อุทายี, ยถา เม กุมาริกา ทุคฺคตา ทุกฺขิตา น สุขํ ลภติ ปาปิกาย สสฺสุยา ปาปเกน สสุเรน ปาปเกน สามิเกนา”ติ. สาปิ โข กุมาริกา อุชฺฌายติ ขิยฺยติ วิปาเจติ “เอวํ ทุคฺคโต โหตุ อยฺโย อุทายี, เอวํ ทุกฺขิโต โหตุ อยฺโย อุทายี, เอวํ มา สุขํ ลภตุ อยฺโย อุทายี, ยถาหํ ทุคฺคตา ทุกฺขิตา น สุขํ ลภามิ ปาปิกาย สสฺสุยา ปาปเกน สสุเรน ปาปเกน สามิเกนา”ติ. อญฺญาปิ อิตฺถิโย อสนฺตุฏฺฐา สสฺสูหิ วา สสุเรหิ วา สามิเกหิ วา, ตา เอวํ โอยาจนฺติ “เอวํ ทุคฺคโต โหตุ อยฺโย อุทายี, เอวํ ทุกฺขิโต โหตุ อยฺโย อุทายี, เอวํ มา สุขํ ลภตุ อยฺโย อุทายี, ยถา มยํ ทุคฺคตา ทุกฺขิตา น สุขํ ลภาม ปาปิกาหิ สสฺสูหิ ปาปเกหิ สสุเรหิ ปาปเกหิ สามิเกหี”ติ. ยา ปน ตา อิตฺถิโย สนฺตุฏฺฐา สสฺสูหิ วา สสุเรหิ วา สามิเกหิ วา, ตา เอวํ อายาจนฺติ “เอวํ สุขิโต โหตุ อยฺโย อุทายี, เอวํ สชฺชิโต โหตุ อยฺโย อุทายี, เอวํ สุขเมโธ\csromansharedfootnote{26cd8158840bda39}{สุขเมธิโต (สี, ก)} โหตุ อยฺโย อุทายี, ยถา มยํ สุขิตา สชฺชิตา สุขเมธา ภทฺทิกาหิ สสฺสูหิ ภทฺทเกหิ สสุเรหิ ภทฺทเกหิ สามิเกหี”ติ.}

\prose{อสฺโสสุํ โข ภิกฺขู เอกจฺจานํ อิตฺถีนํ โอยาจนฺตีนํ เอกจฺจานํ อิตฺถีนํ อายาจนฺตีนํ. เย เต ภิกฺขู อปฺปิจฺฉา ฯเปฯ เต อุชฺฌายนฺติ ขิยฺยนฺติ วิปาเจนฺติ “กถํ หิ นาม อายสฺมา อุทายี สญฺจริตฺตํ สมาปชฺชิสฺสตี”ติ. อถ โข เต ภิกฺขู อายสฺมนฺตํ อุทายิํ อเนก\-ปริยาเยน วิครหิตฺวา ภควโต เอตมตฺถํ อาโรเจสุํ. อถ โข ภควา เอตสฺมิํ นิทาเน เอตสฺมิํ ปกรเณ ภิกฺขุสํฆํ สนฺนิปาตาเปตฺวา อายสฺมนฺตํ อุทายิํ ปฏิปุจฺฉิ “สจฺจํ กิร ตฺวํ อุทายิ สญฺจริตฺตํ สมาปชฺชสี”ติ. สจฺจํ ภควาติ. วิครหิ พุทฺโธ ภควา ฯเปฯ. กถํ หิ นาม ตฺวํ โมฆปุริส สญฺจริตฺตํ สมาปชฺชิสฺสสิ, เนตํ โมฆปุริส อปฺปสนฺนานํ วา ปสาทาย ฯเปฯ. เอวญฺจ ปน ภิกฺขเว อิมํ สิกฺขาปทํ อุทฺทิเสยฺยาถ\textbf{\csromandash{}}}

\proseitem{299}{\textbf{“โย ปน ภิกฺขุ สญฺจริตฺตํ สมาปชฺเชยฺย อิตฺถิยา วา ปุริสมติํ ปุริสสฺส วา อิตฺถิมติํ ชายตฺตเน วา ชารตฺตเน วา, สํฆาทิเสโส”}ติ.}

\prose{\csromanfolio{201}เอวญฺจิทํ ภควโต ภิกฺขูนํ สิกฺขาปทํ ปญฺญตฺตํ โหติ.}

\proseitem{300}{เตน โข ปน สมเยน สมฺพหุลา ธุตฺตา อุยฺยาเน ปริจาเรนฺตา อญฺญตริสฺสา เวสิยา สนฺติเก ทูตํ ปาเหสุํ “อาคจฺฉตุ อุยฺยาเน ปริจาเรสฺสามา”ติ. สา เอวมาห “อหํ ขฺวยฺโย ตุเมฺห น ชานามิ ‘เก วา อิเม กสฺส วา’ติ, อหญฺจมฺหิ พหุภณฺฑา พหุปริกฺขารา, พหินครญฺจ คนฺตพฺพํ, นาหํ คมิสฺสามี”ติ. อถ โข โส ทูโต เตสํ ธุตฺตานํ เอตมตฺถํ อาโรเจสิ. เอวํ วุตฺเต อญฺญตโร ปุริโส เต ธุตฺเต เอตทโวจ “กิสฺส ตุเมฺห อยฺโย เอตํ เวสิํ ยาจิตฺถ, นนุ อยฺโย อุทายี วตฺตพฺโพ, อยฺโย อุทายี อุยฺโยเชสฺสตี”ติ. เอวํ วุตฺเต อญฺญตโร อุปาสโก ตํ ปุริสํ เอตทโวจ “มายฺโย เอวํ อวจ, น กปฺปติ สมณานํ สกฺย\-ปุตฺติ\-ยานํ เอวรูปํ กาตุํ, นายฺโย อุทายี เอวํ กริสฺสตี”ติ. เอวํวุตฺเต “กริสฺสติ น กริสฺสตี”ติ อพฺภุตมกํสุ. อถ โข เต ธุตฺตา เยนายสฺมา อุทายี เตนุ\-ปสงฺกมิํสุ, อุปสงฺกมิตฺวา อายสฺมนฺตํ อุทายิํ เอตทโวจุํ “อิธ มยํ ภนฺเต อุยฺยาเน ปริจาเรนฺตา อสุกาย นาม เวสิยา สนฺติเก ทูตํ ปหิณิมฺหา ‘อาคจฺฉตุ อุยฺยาเน ปริจาเรสฺสามา’ติ, สา เอวมาห ‘อหํ ขฺวยฺโย ตุเมฺห น ชานามิ เก วา อิเม กสฺส วาติ, อหญฺจมฺหิ พหุภณฺฑา พหุปริกฺขารา, พหินครญฺจ คนฺตพฺพํ, นาหํ คมิสฺสามี’ติ, สาธุ ภนฺเต อยฺโย ตํ เวสิํ อุยฺโยเชตู”ติ.}

\csromanmatikaanchor{mtk.64}
\csromanmatikaanchor{mtk.65}
\csromantocmarkref{subsection}{อนุปญฺญตฺติ}{mtk.64}
\bookmark[dest={mtk.64},level=2]{อนุปญฺญตฺติ}
\csromantocmarkref{subsection}{สิกฺขาปทวิภงฺค, ปทภาชนีย}{mtk.65}
\bookmark[dest={mtk.65},level=2]{สิกฺขาปทวิภงฺค, ปทภาชนีย}
\prose{อถ โข อายสฺมา อุทายี เยน สา เวสี เตนุปสงฺกมิ, อุปสงฺกมิตฺวา ตํ เวสิํ เอตทโวจ “กิสฺสิเมสํ น คจฺฉสี”ติ. “อหํ ขฺวยฺย อิเม น ชานามิ ‘เก วา อิเม กสฺส วา’ติ, อหญฺจมฺหิ พหุภณฺฑา พหุปริกฺขารา, พหินครญฺจ คนฺตพฺพํ, นาหํ คมิสฺสามี”ติ. คจฺฉิเมสํ, อหํ อิเม ชานามีติ. สเจ ภนฺเต อยฺโย ชานาติ อหํ คมิสฺสามีติ. อถ โข เต ธุตฺตา ตํ เวสิํ อาทาย อุยฺยานํ อคมํสุ. อถ โข โส อุปาสโก อุชฺฌายติ ขิยฺยติ วิปาเจติ “กถํ หิ นาม อยฺโย อุทายี ตงฺขณิกํ สญฺจริตฺตํ สมาปชฺชิสฺสตี”ติ. อสฺโสสุํ โข ภิกฺขู ตสฺส อุปาสกสฺส อุชฺฌายนฺตสฺส ขิยฺยนฺตสฺส วิปาเจนฺตสฺส. เย เต ภิกฺขู อปฺปิจฺฉา ฯเปฯ เต อุชฺฌายนฺติ ขิยฺยนฺติ วิปาเจนฺติ “กถํ หิ นาม อายสฺมา อุทายี ตงฺขณิกํ สญฺจริตฺตํ สมาปชฺชิสฺสตี”ติ. อถ โข เต ภิกฺขู \csromanfolio{202}อายสฺมนฺตํ อุทายิํ อเนก\-ปริยาเยน วิครหิตฺวา ภควโต เอตมตฺถํ อาโรเจสุํ ฯเปฯ “สจฺจํ กิร ตฺวํ อุทายิ ตงฺขณิกํ สญฺจริตฺตํ สมาปชฺชสี”ติ. สจฺจํ ภควาติ. วิครหิ พุทฺโธ ภควา ฯเปฯ. กถํ หิ นาม ตฺวํ โมฆปุริส ตงฺขณิกํ สญฺจริตฺตํ สมาปชฺชิสฺสสิ, เนตํ โมฆปุริส อปฺปสนฺนานํ วา ปสาทาย ฯเปฯ. เอวญฺจ ปน ภิกฺขเว อิมํ สิกฺขาปทํ อุทฺทิเสยฺยาถ\csromandash{}}

\proseitem{301}{\textbf{“โย ปน ภิกฺขุ สญฺจริตฺตํ สมาปชฺเชยฺย อิตฺถิยา วา ปุริสมติํ ปุริสสฺส วา อิตฺถิมติํ ชายตฺตเน วา ชารตฺตเน วา อนฺตมโส ตงฺขณิกายปิ, สํฆาทิเสโส”ติ}.}

\proseitem{302}{\textbf{โย ปนา}ติ โย ยาทิโส ฯเปฯ. \textbf{ภิกฺขู}ติ ฯเปฯ อยํ อิมสฺมิํ อตฺเถ อธิปฺเปโต “ภิกฺขู”ติ.}

\prose{\textbf{สญฺจริตฺตํ สมาปชฺเชยฺยา}ติ อิตฺถิยา วา ปหิโต ปุริสสฺส สนฺติเก คจฺฉติ, ปุริเสน วา ปหิโต อิตฺถิยา สนฺติเก คจฺฉติ.}

\prose{\textbf{อิตฺถิยา วา ปุริสมตินฺ}ติ ปุริสสฺส มติํ อิตฺถิยา อาโรเจติ.}

\prose{\textbf{ปุริสสฺส วา อิตฺถิมตินฺ}ติ อิตฺถิยา มติํ ปุริสสฺส อาโรเจติ.}

\prose{\textbf{ชายตฺตเน วา}ติ ชายา ภวิสฺสสิ.}

\prose{\textbf{ชารตฺตเน วา}ติ ชารี ภวิสฺสสิ.}

\prose{\textbf{อนฺตมโส ตงฺขณิกายปี}ติ มุหุตฺติกา ภวิสฺสสิ.}

\prose{\textbf{สํฆาทิเสโส}ติ ฯเปฯ เตนปิ วุจฺจติ “สํฆาทิเสโส”ติ.}

\proseitem{303}{ทส อิตฺถิโย มาตุรกฺขิตา ปิตุรกฺขิตา มาตาปิตุ\-รกฺขิตา ภาตุรกฺขิตา ภคินิ\-รกฺขิตา ญาติรกฺขิตา โคตฺตรกฺขิตา ธมฺมรกฺขิตา สารกฺขา สปริทณฺฑา.}

\prose{ทส ภริยาโย ธนกฺกีตา ฉนฺทวาสินี โภควาสินี ปฏวาสินี โอทปตฺตกินี โอภฏจุมฺพฏา, ทาสี จ ภริยา จ, กมฺมการี จ ภริยา จ, ธชาหฏา มุหุตฺติกา.}

\proseitem{304}{\textbf{มาตุรกฺขิตา} นาม มาตา รกฺขติ โคเปติ อิสฺสริยํ กาเรติ วสํ วตฺเตติ.}

\prose{\csromanfolio{203}\textbf{ปิตุรกฺขิตา} นาม ปิตา รกฺขติ โคเปหิ อิสฺสริยํ กาเรติ วสํ วตฺเตติ.}

\prose{\textbf{มาตาปิตุ}\-\textbf{รกฺขิตา} นาม มาตาปิตโร รกฺขนฺติ โคเปนฺติ อิสฺสริยํ กาเรนฺติ วสํ วตฺเตนฺติ.}

\prose{\textbf{ภาตุรกฺขิตา} นาม ภาตา รกฺขติ โคเปติ อิสฺสริยํ กาเรติ วสํ วตฺเตติ.}

\prose{\textbf{ภคินิ}\-\textbf{รกฺขิตา} นาม ภคินี รกฺขติ โคเปติ อิสฺสริยํ กาเรติ วสํ วตฺเตติ.}

\prose{\textbf{ญาติรกฺขิตา} นาม ญาตกา รกฺขนฺติ โคเปนฺติ อิสฺสริยํ กาเรนฺติ วสํ วตฺเตนฺติ.}

\prose{\textbf{โคตฺตรกฺขิตา} นาม สโคตฺตา รกฺขนฺติ โคเปนฺติ อิสฺสริยํ กาเรนฺติ วสํ วตฺเตนฺติ.}

\prose{\textbf{ธมฺมรกฺขิตา} นาม สหธมฺมิกา รกฺขนฺติ โคเปนฺติ อิสฺสริยํ กาเรนฺติ วสํ วตฺเตนฺติ.}

\prose{\textbf{สารกฺขา} นาม คพฺเภปิ ปริคฺคหิตา โหติ “มยฺหํ เอสา”ติ อนฺตมโส มาลาคุฬปริกฺขิตฺตาปิ.}

\prose{\textbf{สปริทณฺฑา} นาม เกหิจิ ทณฺโฑ ฐปิโต โหติ “โย อิตฺถนฺนามํ อิตฺถิํ คจฺฉติ เอตฺตโก ทณฺโฑ”ติ.}

\prose{\textbf{ธนกฺกีตา} นาม ธเนน กิณิตฺวา วาเสติ.}

\prose{\textbf{ฉนฺทวาสินี} นาม ปิโย ปิยํ วาเสติ.}

\prose{\textbf{โภควาสินี} นาม โภคํ ทตฺวา วาเสติ.}

\prose{\textbf{ปฏวาสินี} นาม ปฏํ ทตฺวา วาเสติ.}

\prose{\textbf{โอทปตฺตกินี} นาม อุทกปตฺตํ อามสิตฺวา วาเสติ.}

\prose{\textbf{โอภฏจุมฺพฏา} นาม จุมฺพฏํ โอโรเปตฺวา วาเสติ.}

\prose{\textbf{ทาสี} นาม ทาสี เจว โหติ ภริยา จ.}

\prose{\textbf{กมฺมการี} นาม กมฺมการี เจว โหติ ภริยา จ.}

\prose{\textbf{ธชาหฏา} นาม กรมรานีตา วุจฺจติ.}

\prose{\textbf{มุหุตฺติกา} นาม ตงฺขณิกา วุจฺจติ.}

\proseitem{305}{\csromanfolio{204}ปุริโส ภิกฺขุํ ปหิณติ “คจฺฉ ภนฺเต อิตฺถนฺนามํ มาตุรกฺขิตํ พฺรูหิ โหหิ กิร อิตฺถนฺนามสฺส ภริยา ธนกฺกีตา”ติ, ปฏิคฺคณฺหาติ วีมํสกิ ปจฺจาหรติ, อาปตฺติ สํฆาทิเสสสฺส.}

\prose{ปุริโส ภิกฺขุํ ปหิณติ “คจฺฉ ภนฺเต อิตฺถนฺนามํ ปิตุรกฺขิตํ พฺรูหิ ฯเปฯ มาตาปิตุ\-รกฺขิตํ พฺรูหิ. ภาตุรกฺขิตํ พฺรูหิ. ภคินิ\-รกฺขิตํ พฺรูหิ. ญาติรกฺขิตํ พฺรูหิ. โคตฺตรกฺขิตํ พฺรูหิ. ธมฺม\-รกฺขิตํ พฺรูหิ. สารกฺขํ พฺรูหิ. สปริทณฺฑํ พฺรูหิ โหหิ กิร อิตฺถนฺนามสฺส ภริยา ธนกฺกีตา”ติ, ปฏิคฺคณฺหาติ วีมํสติ ปจฺจาหรติ, อาปตฺติ สํฆาทิเสสสฺส.}

\vspace{\tipitakaparskip}
\csromancenter{นิกฺเขปปทานิ.}
\proseitem{306}{ปุริโส ภิกฺขุํ ปหิณติ “คจฺฉ ภนฺเต อิตฺถนฺนามํ มาตุรกฺขิตญฺจ ปิตุรกฺขิตญฺจ พฺรูหิ โหถ กิร อิตฺถนฺนามสฺส ภริยาโย ธนกฺกีตา”ติ, ปฏิคฺคณฺหาติ วีมํสติ ปจฺจาหรติ, อาปตฺติ สํฆาทิเสสสฺส.}

\prose{ปุริโส ภิกฺขุํ ปหิณติ “คจฺฉ ภนฺเต อิตฺถนฺนามํ มาตุรกฺขิตญฺจ มาตาปิตุรกฺขิตญฺจ ฯเปฯ มาตุรกฺขิตญฺจ ภาตุรกฺขิตญฺจ. มาตุรกฺขิตญฺจ ภคินิรกฺขิตญฺจ. มาตุรกฺขิตญฺจ ญาติรกฺขิตญฺจ. มาตุรกฺขิตญฺจ โคตฺตรกฺขิตญฺจ. มาตุรกฺขิตญฺจ ธมฺมรกฺขิตญฺจ. มาตุรกฺขิตญฺจ สารกฺขญฺจ. มาตุรกฺขิตญฺจ สปริทณฺฑญฺจ พฺรูหิ โหถ กิร อิตฺถนฺนามสฺส ภริยาโย ธนกฺกีตา”ติ, ปฏิคฺคณฺหาติ วีมํสติ ปจฺจาหรติ, อาปตฺติ สํฆาทิเสสสฺส.}

\vspace{\tipitakaparskip}
\csromancenter{ขณฺฑจกฺกํ.}
\proseitem{307}{ปุริโส ภิกฺขุํ ปหิณติ “คจฺฉ ภนฺเต อิตฺถนฺนามํ ปิตุรกฺขิตญฺจ มาตาปิตุรกฺขิตญฺจ พฺรูหิ โหถ กิร อิตฺถนฺนามสฺส ภริยาโย ธนกฺกีตา”ติ, ปฏิคฺคณฺหาติ วีมํสติ ปจฺจาหรติ, อาปตฺติ สํฆาทิเสสสฺส.}

\prose{ปุริโส ภิกฺขุํ ปหิณติ “คจฺฉ ภนฺเต อิตฺถนฺนามํ ปิตุรกฺขิตญฺจ ภาตุรกฺขิตญฺจ ฯเปฯ ปิตุรกฺขิตญฺจ ภคินิรกฺขิตญฺจ. ปิตุรกฺขิตญฺจ ญาติรกฺขิตญฺจ. ปิตุรกฺขิตญฺจ โคตฺตรกฺขิตญฺจ. ปิตุรกฺขิตญฺจ ธมฺมรกฺขิตญฺจ. ปิตุรกฺขิตญฺจ สารกฺขญฺจ. ปิตุรกฺขิตญฺจ สปริทณฺฑญฺจ พฺรูหิ โหถ กิร อิตฺถนฺนามสฺส ภริยาโย ธนกฺกีตา”ติ, ปฏิคฺคณฺหาติ วีมํสติ ปจฺจาหรติ, อาปตฺติ สํฆาทิเสสสฺส.}

\prose{\csromanfolio{205}ปุริโส ภิกฺขุํ ปหิณติ “คจฺฉ ภนฺเต อิตฺถนฺนามํ ปิตรกฺขิตญฺจ มาตุรกฺขิตญฺจ พฺรูหิ โหถ กิร อิตฺถนฺนามสฺส ภริยาโย ธนกฺกีตา”ติ, ปฏิคฺคณฺหาติ วีมํสติ ปจฺจาหรติ, อาปตฺติ สํฆาทิเสสสฺส.}

\vspace{\tipitakaparskip}
\csromancenter{พทฺธจกฺกํ, มูลํ สํขิตํ.}
\proseitem{308}{ปุริโส ภิกฺขุํ ปหิณติ “คจฺฉ ภนฺเต อิตฺถนฺนามํ สปริทณฺฑญฺจ มาตุรกฺขิตญฺจ พฺรูหิ โหถ กิร อิตฺถนฺนามสฺส ภริยาโย ธนกฺกีตา”ติ, ปฏิคฺคณฺหาติ วีมํสติ ปจฺจาหรติ, อาปตฺติ สํฆาทิเสสสฺส.}

\prosewithcloser{ปุริโส ภิกฺขุํ ปหิณติ “คจฺฉ ภนฺเต อิตฺถนฺนามํ สปริทณฺฑญฺจ ปิตุรกฺขิตญฺจ ฯเปฯ สปริทณฺฑญฺจ มาตาปิตุรกฺขิตญฺจ. สปริทณฺฑญฺจ ภาตุรกฺขิตญฺจ. สปริทณฺฑญฺจ ภคินิรกฺขิตญฺจ. สปริทณฺฑญฺจ ญาติรกฺขิตญฺจ. สปริทณฺฑญฺจ โคตฺตรกฺขิตญฺจ. สปริทณฺฑญฺจ ธมฺมรกฺขิตญฺจ. สปริทณฺฑญฺจ สารกฺขญฺจ พฺรูหิ โหถ กิร อิตฺถนฺนามสฺส ภริยาโย ธนกฺกีตา”ติ, ปฏิคฺคณฺหาติ วีมํสติ ปจฺจาหรติ, อาปตฺติ สํฆาทิเสสสฺส.}{เอกมูลกํ นิฏฺฐิตํ.}
\prose{เอวํ ทุมูลกํปิ ติมูลกํปิ ยาว นวมูลกํ กาตพฺพํ.}

\proseitemwithcloser{309}{ปุริโส ภิกฺขุํ ปหิณติ “คจฺฉ ภนฺเต อิตฺถนฺนามํ มาตุรกฺขิตญฺจ ปิตุรกฺขิตญฺจ มาตาปิตุรกฺขิตญฺจ ภาตุรกฺขิตญฺจ ภคินิรกฺขิตญฺจ ญาติรกฺขิตญฺจ โคตฺตรกฺขิตญฺจ ธมฺมรกฺขิตญฺจ สารกฺขญฺจ สปริทณฺฑญฺจ พฺรูหิ โหถ กิร อิตฺถนฺนามสฺส ภริยาโย ธนกฺกีตา”ติ, ปฏิคฺคณฺหาติ วีมํสติ ปจฺจาหรติ, อาปตฺติ สํฆาทิเสสสฺส.}{ธนกฺกีตา\-จกฺกํ นิฏฺฐิตํ.}
\proseitem{310}{ปุริโส ภิกฺขุํ ปหิณติ “คจฺฉ ภนฺเต อิตฺถนฺนามํ มาตุรกฺขิตํ พฺรูหิ โหหิ กิร อิตฺถนฺนามสฺส ภริยา ฉนฺทวาสินี ฯเปฯ โภควาสินี. ปฏวาสินี. โอทปตฺตกินี. โอภฏจุมฺพฏา. ทาสี จ ภริยา จ. กมฺมการี จ ภริยา จ. ธชาหฏา. มุหุตฺติกา”ติ, ปฏิคฺคณฺหาติ วีมํสติ ปจฺจาหรติ, อาปตฺติ สํฆาทิเสสสฺส.}

\prose{\csromanfolio{206}ปุริโส ภิกฺขุํ ปหิณติ “คจฺฉ ภนฺเต อิตฺถนฺนามํ ปิตุรกฺขิตํ พฺรูหิ ฯเปฯ มาตาปิตุ\-รกฺขิตํ พฺรูหิ. ภาตุรกฺขิตํ พฺรูหิ. ภคินิ\-รกฺขิตํ พฺรูหิ. ญาติรกฺขิตํ พฺรูหิ. โคตฺตรกฺขิตํ พฺรูหิ. ธมฺม\-รกฺขิตํ พฺรูหิ. สารกฺขํ พฺรูหิ. สปริทณฺฑํ พฺรูหิ โหหิ กิร อิตฺถนฺนามสฺส ภริยา มุหุตฺติกา”ติ, ปฏิคฺคณฺหาติ วีมํสติ ปจฺจาหรติ, อาปตฺติ สํฆาทิเสสสฺส.}

\vspace{\tipitakaparskip}
\csromancenter{นิกฺเขปปทานิ.}
\proseitem{311}{ปุริโส ภิกฺขุํ ปหิณติ “คจฺฉ ภนฺเต อิตฺถนฺนามํ มาตุรกฺขิตญฺจ ปิตุรกฺขิตญฺจ พฺรูหิ โหถ กิร อิตฺถนฺนามสฺส ภริยาโย มุหุตฺติกา”ติ, ปฏิคฺคณฺหาติ วีมํสติ ปจฺจาหรติ, อาปตฺติ สํฆาเทเสสสฺส.}

\prose{ปุริโส ภิกฺขุํ ปหิณติ “คจฺฉ ภนฺเต อิตฺถนฺนามํ มาตุรกฺขิตญฺจ มาตาปิตุรกฺขิตญฺจ ฯเปฯ มาตุรกฺขิตญฺจ สปริทณฺฑญฺจ พฺรูหิ โหถ กิร อิตฺถนฺนามสฺส ภริยาโย มุหุตฺติกา”ติ, ปฏิคฺคณฺหาติ วีมํสติ ปจฺจาหรติ, อาปตฺติ สํฆาทิเสสสฺส.}

\vspace{\tipitakaparskip}
\csromancenter{ขณฺฑจกฺกํ.}
\proseitem{312}{ปุริโส ภิกฺขุํ ปหิณติ “คจฺฉ ภนฺเต อิตฺถนฺนามํ ปิตุรกฺขิตญฺจ มาตาปิตุรกฺขิตญฺจ พฺรูหิ โหถ กิร อิตฺถนฺนามสฺส ภริยาโย มุหุตฺติกา”ติ, ปฏิคฺคณฺหาติ วีมํสติ ปจฺจาหรติ, อาปตฺติ สํฆาทิเสสสฺส.}

\prose{ปุริโส ภิกฺขุํ ปหิณติ “คจฺฉ ภนฺเต อิตฺถนฺนามํ ปิตุรกฺขิตญฺจ ภาตุรกฺขิตญฺจ ฯเปฯ ปิตุรกฺขิตญฺจ สปริทณฺฑญฺจ พฺรูหิ โหถ กิร อิถนฺนามสฺส ภริยาโย มุหุตฺติกา”ติ, ปฏิคฺคณฺหาติ วีมํสติ ปจฺจาหรติ, อาปตฺติ สํฆาทิเสสสฺส.}

\prose{ปุริโส ภิกฺขุํ ปหิณติ “คจฺฉ ภนฺเต อิตฺถนฺนามํ ปิตุรกฺขิตญฺจ มาตุรกฺขิตญฺจ พฺรูหิ โหถ กิร อิตฺถนฺนามสฺส ภริยาโย มุหุตฺติกา”ติ, ปฏิคฺคณฺหาติ วีมํสติ ปจฺจาหรติ, อาปตฺติ สํฆาทิเสสสฺส.}

\vspace{\tipitakaparskip}
\csromancenter{พทฺธจกฺกํ, มูลํ สํขิตฺตํ.}
\proseitem{313}{ปุริโส ภิกฺขุํ ปหิณติ “คจฺฉ ภนฺเต อิตฺถนฺนามํ สปริทณฺฑญฺจ มาตุรกฺขิตญฺจ พฺรูหิ โหถ กิร อิตฺถนฺนามสฺส ภริยาโย \csromanfolio{207}มุหุตฺติกา”ติ, ปฏิคฺคณฺหาติ วีมํสติ ปจฺจาหรติ, อาปตฺติ สํฆาทิเสสสฺส.}

\prosewithcloser{ปุริโส ภิกฺขุํ ปหิณติ “คจฺฉ ภนฺเต อิตฺถนฺนามํ สปริทณฺฑญฺจ ปิตุรกฺขิตญฺจ ฯเปฯ สปริทณฺฑญฺจ สารกฺขญฺจ พฺรูหิ โหถ กิร อิตฺถนฺนามสฺส ภริยาโย มุหุตฺติกา”ติ, ปฏิคฺคณฺหาติ วีมํสติ ปจฺจาหรติ, อาปตฺติ สํฆาทิเสสสฺส.}{เอกมูลกํ นิฏฺฐิตํ.}
\csromancenter{ทุมูลกาทีนิปิ เอวเมว กาตพฺพานิ.}
\proseitemwithcloser{314}{ปุริโส ภิกฺขุํ ปหิณติ “คจฺฉ ภนฺเต อิตฺถนฺนามํ มาตุรกฺขิตญฺจ ปิตุรกฺขิตญฺจ มาตาปิตุรกฺขิตญฺจ ภาตุรกฺขิตญฺจ ภคินิรกฺขิตญฺจ ญาติรกฺขิตญฺจ โคตฺตรกฺขิตญฺจ ธมฺมรกฺขิตญฺจ สารกฺขญฺจ สปริทณฺฑญฺจ พฺรูหิ โหถ กิร อิตฺถนฺนามสฺส ภริยาโย มุหุตฺติกา”ติ, ปฏิคฺคณฺหาติ วีมํสติ ปจฺจาหรติ, อาปตฺติ สํฆาทิเสสสฺส.}{มุหุตฺติกา\-จกฺกํ นิฏฺฐิตํ.}
\proseitem{315}{ปุริโส ภิกฺขุํ ปหิณติ “คจฺฉ ภนฺเต อิตฺถนฺนามํ มาตุรกฺขิตํ พฺรูหิ โหหิ กิร อิตฺถนฺนามสฺส ภริยา ธนกฺกีตา”ติ, ปฏิคฺคณฺหาติ วีมํสติ ปจฺจาหรติ, อาปตฺติ สํฆาทิเสสสฺส.}

\prose{ปุริโส ภิกฺขุํ ปหิณติ “คจฺฉ ภนฺเต อิตฺถนฺนามํ มาตุรกฺขิตํ พฺรูหิ โหหิ กิร อิตฺถนฺนามสฺส ภริยา ฉนฺทวาสินี ฯเปฯ โภควาสินี. ปฏวาสินี. โอทปตฺตกินี. โอภฏจุมฺพฏา. ทาสี จ ภริยา จ. กมฺมการี จ ภริยา จ. ธชาหฏา. มุหุตฺติกา”ติ, ปฏิคฺคณฺหาติ วีมํสติ ปจฺจาหรติ, อาปตฺติ สํฆาทิเสสสฺส.}

\vspace{\tipitakaparskip}
\csromancenter{นิกฺเขปปทานิ.}
\proseitem{316}{ปุริโส ภิกฺขุํ ปหิณติ “คจฺฉ ภนฺเต อิตฺถนฺนามํ มาตุรกฺขิตํ พฺรูหิ โหหิ กิร อิตฺถนฺนามสฺส ภริยา ธนกฺกีตา จ ฉนฺทวาสินี จา”ติ, ปฏิคฺคณฺหาติ วีมํสติ ปจฺจาหรติ, อาปตฺติ สํฆาทิเสสสฺส.}

\prose{\csromanfolio{208}ปุริโส ภิกฺขุํ ปหิณติ “คจฺฉ ภนฺเต อิตฺถนฺนามํ มาตุรกฺขิตํ พฺรูหิ โหหิ กิร อิตฺถนฺนามสฺส ภริยา ธนกฺกีตา จ โภควาสินี จ ฯเปฯ ธนกฺกีตา จ ปฏิวาสินี จ. ธนกฺกีตา จ โอทปตฺตกินี จ. ธนกฺกีตา จ โอภฏจุมฺพฏา จ. ธนกฺกีตา จ ทาสี จ ภริยา จ. ธนกฺกีตา จ กมฺมการี จ ภริยา จ. ธนกฺกีตา จ ธชาหฏา จ. ธนกฺกีตา จ มุหุตฺติกา จา”ติ, ปฏิคฺคณฺหาติ วีมํสติ ปจฺจาหรติ, อาปตฺติ สํฆาทิเสสสฺส.}

\vspace{\tipitakaparskip}
\csromancenter{ขณฺฑจกฺกํ.}
\proseitem{317}{ปุริโส ภิกฺขุํ ปหิณติ “คจฺฉ ภนฺเต อิตฺถนฺนามํ มาตุรกฺขิตํ พฺรูหิ โหหิ กิร อิตฺถนฺนามสฺส ภริยา ฉนฺทวาสินี จ โภควาสินี จ ฯเปฯ ฉนฺทวาสินี จ มุหุตฺติกา จ. ฉนฺทวาสินี จ ธนกฺกีตา จา”ติ, ปฏิคฺคณฺหาติ วีมํสติ ปจฺจาหรติ, อาปตฺติ สํฆาทิเสสสฺส.}

\vspace{\tipitakaparskip}
\csromancenter{พทฺธจกฺกํ, มูลํ สํขิตฺตํ.}
\proseitemwithcloser{318}{ปุริโส ภิกฺขุํ ปหิณติ “คจฺฉ ภนฺเต อิตฺถนฺนามํ มาตุรกฺขิตํ พฺรูหิ โหหิ กิร อิตฺถนฺนามสฺส ภริยา มุหุตฺติกา จ ธนกฺกีตา จ ฯเปฯ มุหุตฺติกา จ ฉนฺทวาสินี จ ฯเปฯ มุหุตฺติกา จ ธชาหฏา จา”ติ, ปฏิคฺคณฺหาติ วีมํสติ ปจฺจาหรติ, อาปตฺติ สํฆาทิเสสสฺส.}{เอกมูลกํ นิฏฺฐิตํ.}
\csromancenter{ทุมูลกาทีนิปิ เอวเมว กาตพฺพานิ.}
\proseitemwithcloser{319}{ปุริโส ภิกฺขุํ ปหิณติ “คจฺฉ ภนฺเต อิตฺถนฺนามํ มาตุรกฺขิตํ พฺรูหิ โหหิ กิร อิตฺถนฺนามสฺส ภริยา ธนกฺกีตา จ ฉนฺทวาสินี จ โภควาสินี จ ปฏิวาสินี จ โอทปตฺตกินี จ โอภฏจุมฺพฏา จ ทาสี จ ภริยา จ กมฺมการี จ ภริยา จ ธชาหฏา จ มุหุตฺติกา จา”ติ, ปฏิคฺคณฺหาติ วีมํสติ ปจฺจาหรติ, อาปตฺติ สํฆาทิเสสสฺส.}{มาตุรกฺขิตา\-จกฺกํ นิฏฺฐิตํ.}
\prose{ปุริโส ภิกฺขุํ ปหิณติ “คจฺฉ ภนฺเต อิตฺถนฺนามํ ปิตุรกฺขิตํ ฯเปฯ มาตาปิตุ\-รกฺขิตํ. ภาตุรกฺขิตํ. ภคินิ\-รกฺขิตํ. ญาติรกฺขิตํ. โคตฺตรกฺขิตํ.}

\prose{\csromanfolio{209}ธมฺม\-รกฺขิตํ. สารกฺขํ. สปริทณฺฑํ พฺรูหิ โหหิ กิร อิตฺถนฺนามสฺส ภริยา ธนกฺกีตา”ติ, ปฏิคฺคณฺหาติ วีมํสติ ปจฺจาหรติ, อาปตฺติ สํฆาทิเสสสฺส.}

\prose{ปุริโส ภิกฺขุํ ปหิณติ “คจฺฉ ภนฺเต อิตฺถนฺนามํ สปริทณฺฑํ พฺรูหิ โหหิ กิร อิตฺถนฺนามสฺส ภริยา ฉนฺทวาสินี ฯเปฯ โภควาสินี. ปฏวาสินี. โอทปตฺตกินี. โอภฏจุมฺพฏา. ทาสี จ ภริยา จ. กมฺมการี จ ภริยา จ. ธชาหฏา. มุหุตฺติกา”ติ, ปฏิคฺคณฺหาติ วีมํสติ ปจฺจาหรติ, อาปตฺติ สํฆาทิเสสสฺส.}

\vspace{\tipitakaparskip}
\csromancenter{นิกฺเขปปทานิ.}
\proseitem{320}{ปุริโส ภิกฺขุํ ปหิณติ “คจฺฉ ภนฺเต อิตฺถนฺนามํ สปริทณฺฑํ พฺรูหิ โหหิ กิร อิตฺถนฺนามสฺส ภริยา ธนกฺกีตา จ ฉนฺทวาสินี จา”ติ, ปฏิคฺคณฺหาติ วีมํสติ ปจฺจาหรติ, อาปตฺติ สํฆาทิเสสสฺส.}

\prose{ปุริโส ภิกฺขุํ ปหิณติ “คจฺฉ ภนฺเต อิตฺถนฺนามํ สปริทณฺฑํ พฺรูหิ โหหิ กิร อิตฺถนฺนามสฺส ภริยา ธนกฺกีตา จ โภควาสินี จ ฯเปฯ ธนกฺกีตา จ มุหุตฺติกา จา”ติ, ปฏิคฺคณฺหาติ วีมํสติ ปจฺจาหรติ, อาปตฺติ สํฆาทิเสสสฺส.}

\vspace{\tipitakaparskip}
\csromancenter{ขณฺฑจกฺกํ.}
\prose{ปุริโส ภิกฺขุํ ปหิณติ “คจฺฉ ภนฺเต อิตฺถนฺนามํ สปริทณฺฑํ พฺรูหิ โหหิ กิร อิตฺถนฺนามสฺส ภริยา ฉนฺทวาสินี จ โภควาสินี จ ฯเปฯ ฉนฺทวาสินี จ มุหุตฺติกา จ. ฉนฺทวาสินี จ ธนกฺกีตาจา”ติ, ปฏิคฺคณฺหาติ วีมํสติ ปจฺจาหรติ, อาปตฺติ สํฆาทิเสสสฺส.}

\vspace{\tipitakaparskip}
\csromancenter{พทฺธจกฺกํ, มูลํ สขิตฺตํ.}
\prosewithcloser{ปุริโส ภิกฺขุํ ปหิณติ “คจฺฉ ภนฺเต อิตฺถนฺนามํ สปริทณฺฑํ พฺรูหิ โหหิ กิร อิตฺถนฺนามสฺส ภริยา มุหุตฺติกา จ ธนกฺกีตา จ ฯเปฯ มุหุตฺติกา จ ฉนฺทวาสินี จ ฯเปฯ มุหุตฺติกา จ ธชาหฏา จา”ติ, ปฏิคฺคณฺหาติ วีมํสติ ปจฺจาหรติ, อาปตฺติ สํฆาทิเสสสฺส.}{เอกมูลกํ นิฏฺฐิตํ.}
\prose{ทุมูลกํปิ ติมูลกํปิ ยาว นวมูลกํ เอวเมว กาตพฺพํ.}

\proseitemwithcloser{321}{\csromanfolio{210}ปุริโส ภิกฺขุํ ปหิณติ “คจฺฉ ภนฺเต อิตฺถนฺนามํ สปริทณฺฑํ พฺรูหิ โหหิ กิร อิตฺถนฺนามสฺส ภริยา ธนกฺกีตา จ ฉนฺทวาสินี จ โภควาสินี จ ปฏวาสินี จ โอทปตฺตกินี จ โอภฏจุมฺพฏา จ ทาสี จ ภริยา จ กมฺมการี จ ภริยา จ ธชาหฏา จ มุหุตฺติกา จา”ติ, ปฏิคฺคณฺหาติ วีมํสติ ปจฺจาหรติ, อาปตฺติ สํฆาทิเสสสฺส.}{สปริทณฺฑา\-จกฺกํ นิฏฺฐิตํ.}
\prose{ปุริโส ภิกฺขุํ ปหิณติ “คจฺฉ ภนฺเต อิตฺถนฺนามํ มาตุรกฺขิตํ พฺรูหิ โหหิ กิร อิตฺถนฺนามสฺส ภริยา ธนกฺกีตา”ติ, ปฏิคฺคณฺหาติ วีมํสติ ปจฺจาหรติ, อาปตฺติ สํฆาทิเสสสฺส.}

\prose{ปุริโส ภิกฺขุํ ปหิณติ “คจฺฉ ภนฺเต อิตฺถนฺนามํ มาตุรกฺขิตญฺจ ปิตุรกฺขิตญฺจ พฺรูหิ โหถ กิร อิตฺถนฺนามสฺส ภริยาโย ธนกฺกีตา จ ฉนฺทวาสินี จา”ติ, ปฏิคฺคณฺหาติ วีมํสติ ปจฺจาหรติ, อาปตฺติ สํฆาทิเสสสฺส.}

\prose{ปุริโส ภิกฺขุํ ปหิณติ “คจฺฉ ภนฺเต อิตฺถนฺนามํ มาตุรกฺขิตญฺจ ปิตุรกฺขิตญฺจ มาตาปิตุรกฺขิตญฺจ พฺรูหิ โหถ กิร อิตฺถนฺนามสฺส ภริยาโย ธนกฺกีตา จ ฉนฺทวาสินี จ โภควาสินี จา”ติ, ปฏิคฺคณฺหาติ วีมํสติ ปจฺจาหรติ, อาปตฺติ สํฆาทิเสสสฺส.}

\vspace{\tipitakaparskip}
\csromancenter{เอวํ อุภโต\-วฑฺฒกํ กาตพฺพํ.}
\prosewithcloser{ปุริโส ภิกฺขุํ ปหิณติ “คจฺฉ ภนฺเต อิตฺถนฺนามํ มาตุรกฺขิตญฺจ ปิตุรกฺขิตญฺจ มาตาปิตุรกฺขิตญฺจ ภาตุรกฺขิตญฺจ ภคินิรกฺขิตญฺจ ญาติรกฺขิตญฺจ โคตฺตรกฺขิตญฺจ ธมฺมรกฺขิตญฺจ สารกฺขญฺจ สปริทณฺฑญฺจ พฺรูหิ โหถ กิร อิตฺถนฺนามสฺส ภริยาโย ธนกฺกีตา จ ฉนฺทวาสินี จ โภควาสินี จ ปฏวาสินี จ โอทปตฺตกินี จ โอภฏจุมฺพฏา จ ทาสี จ ภริยา จ กมฺมการี จ ภริยา จ ธชาหฏา จ มุหุตฺติกา จา”ติ, ปฏิคฺคณฺหาติ วีมํสติ ปจฺจาหรติ, อาปตฺติ สํฆาทิเสสสฺส.}{อุภโต\-วฑฺฒกํ นิฏฺฐิตํ.}
\prose{ปุริสสฺส มาตา ภิกฺขุํ ปหิณติ ฯเปฯ. ปุริสสฺส ปิตา ภิกฺขุํ ปหิณติ ฯเปฯ. ปุริสสฺส มาตาปิตโร ภิกฺขุํ ปหิณนฺติ ฯเปฯ. ปุริสสฺส ภาตา ภิกฺขุํ ปหิณติ ฯเปฯ. ปุริสสฺส ภคินี ภิกฺขุํ ปหิณติ ฯเปฯ. ปุริสสฺส}

\prose{\csromanfolio{211}ญาตกา ภิกฺขุํ ปหิณนฺติ ฯเปฯ. ปุริสสฺส โคตฺตา ภิกฺขุํ ปหิณนฺติ ฯเปฯ. ปุริสสฺส สหธมฺมิกา ภิกฺขุํ ปหิณนฺติ ฯเปฯ.}

\prose{ปุริสสฺส เปยฺยาโล วิตฺถาเรตพฺโพ.}

\prose{อุภโต\-วฑฺฒกํ ยถา ปุริมนโย ตเถว วิตฺถาเรตพฺพํ.}

\proseitem{322}{มาตุรกฺขิตาย มาตา ภิกฺขุํ ปหิณติ “คจฺฉ ภนฺเต อิตฺถนฺนามํ พฺรูหิ โหตุ อิตฺถนฺนามสฺส ภริยา ธนกฺกีตา”ติ, ปฏิคฺคณฺหาติ วีมํสติ ปจฺจาหรติ, อาปตฺติ สํฆาทิเสสสฺส.}

\prose{มาตุรกฺขิตาย มาตา ภิกฺขุํ ปหิณติ “คจฺฉ ภนฺเต อิตฺถนฺนามํ พฺรูหิ โหตุ อิตฺถนฺนามสฺส ภริยา ฉนฺทวาสินี ฯเปฯ โภควาสินี. ปฏวาสินี. โอทปตฺตกินี. โอภฏจุมฺพฏา. ทาสี จ ภริยา จ. กมฺมการี จ ภริยา จ. ธชาหฏา. มุหุตฺติกา”ติ, ปฏิคฺคณฺหาติ วีมํสติ ปจฺจาหรติ, อาปตฺติ สํฆาทิเสสสฺส.}

\vspace{\tipitakaparskip}
\csromancenter{นิกฺเขปปทานิ.}
\proseitem{323}{มาตุรกฺขิตาย มาตา ภิกฺขุํ ปหิณติ “คจฺฉ ภนฺเต อิตฺถนฺนามํ พฺรูหิ โหตุ อิตฺถนฺนามสฺส ภริยา ธนกฺกีตา จ ฉนฺทวาสินี จ ฯเปฯ ธนกฺกีตา จ โภควาสินี จ ฯเปฯ ธนกฺกีตา จ มุหุตฺติกา จา”ติ, ปฏิคฺคณฺหาติ วีมํสติ ปจฺจาหรติ, อาปตฺติ สํฆาทิเสสสฺส.}

\vspace{\tipitakaparskip}
\csromancenter{ขณฺฑจกฺกํ.}
\proseitem{324}{มาตุรกฺขิตาย มาตา ภิกฺขุํ ปหิณติ “คจฺฉ ภนฺเต อิตฺถนฺนามํ พฺรูหิ โหตุ อิตฺถนฺนามสฺส ภริยา ฉนฺทวาสินี จ โภควาสินี จ ฯเปฯ ฉนฺทวาสินี จ มุหุตฺติกา จ. ฉนฺทวาสินี จ ธนกฺกีตา จา”ติ, ปฏิคฺคณฺหาติ วีมํสติ ปจฺจาหรติ, อาปตฺติ สํฆาทิเสสสฺส.}

\vspace{\tipitakaparskip}
\csromancenter{พทฺธจกฺกํ, มูลํ สํขิตฺตํ.}
\proseitemwithcloser{325}{มาตุรกฺขิตาย มาตา ภิกฺขุํ ปหิณติ “คจฺฉ ภนฺเต อิตฺถนฺนามํ พฺรูหิ โหตุ อิตฺถนฺนามสฺส ภริยา มุหุตฺติกา จ ธนกฺกีตา จ ฯเปฯ มุหุตฺติกา จ ฉนฺทวาสินี จ ฯเปฯ มุหุตฺติกา จ ธชาหฏา จา”ติ, ปฏิคฺคณฺหาติ วีมํสติ ปจฺจาหรติ, อาปตฺติ สํฆาทิเสสสฺส.}{เอกมูลกํ นิฏฺฐิตํ.}
\prose{ทุมูลกํปิ ติมูลกํปิ ยาว นวมูลกํ เอวเมว กาตพฺพํ.}

\proseitemwithcloser{326}{\csromanfolio{212}มาตุรกฺขิตาย มาตา ภิกฺขุํ ปหิณติ “คจฺฉ ภนฺเต อิตฺถนฺนามํ พฺรูหิ โหตุ อิตฺถนฺนามสฺส ภริยา ธนกฺกีตา จ ฉนฺทวาสินี จ โภควาสินี จ ปฏวาสินี จ โอทปตฺตกินี จ โอภฏจุมฺพฏา จ ทาสี จ ภริยา จ กมฺมการี จ ภริยา จ ธชาหฏา จ มุหุตฺติกา จา”ติ, ปฏิคฺคณฺหาติ วีมํสติ ปจฺจาหรติ, อาปตฺติ สํฆาทิเสสสฺส.}{มาตุจกฺกํ นิฏฺฐิตํ.}
\prose{ปิตุรกฺขิตาย ปิตา ภิกฺขุํ ปหิณติ ฯเปฯ มาตาปิตุ\-รกฺขิตาย มาตาปิตโร ภิกฺขุํ ปหิณนฺติ. ภาตุ\-รกฺขิตาย ภาตา ภิกฺขุํ ปหิณติ. ภคินิ\-รกฺขิตาย ภคินี ภิกฺขุํ ปหิณติ. ญาติรกฺขิตาย ญาตกา ภิกฺขุํ ปหิณนฺติ. โคตฺต\-รกฺขิตาย โคตฺตา\csromansharedfootnote{e9a63e0d7f794cca}{สโคตฺตา (?)} ภิกฺขุํ ปหิณนฺติ. ธมฺม\-รกฺขิตาย สหธมฺมิกา ภิกฺขุํ ปหิณนฺติ. สารกฺขาย เยน ปริคฺคหิตา โหติ, โส ภิกฺขุํ ปหิณติ. สปริทณฺฑาย เยน ทณฺโฑ ฐปิโต โหติ, โส ภิกฺขุํ ปหิณติ “คจฺฉ ภนฺเต อิตฺถนฺนามํ พฺรูหิ โหตุ อิตฺถนฺนามสฺส ภริยา ธนกฺกีตา”ติ, ปฏิคฺคณฺหาติ วีมํสติ ปจฺจาหรติ, อาปตฺติ สํฆาทิเสสสฺส.}

\prose{สปริทณฺฑาย เยน ทณฺโฑ ฐปิโต โหติ, โส ภิกฺขุํ ปหิณติ “คจฺฉ ภนฺเต อิตฺถนฺนามํ พฺรูหิ โหตุ อิตฺถนฺนามสฺส ภริยา ฉนฺทวาสินี ฯเปฯ โภควาสินี. ปฏวาสินี. โอทปตฺตกินี. โอภฏจุมฺพฏา. ทาสี จ ภริยา จ. กมฺมการี จ ภริยา จ. ธชาหฏา. มุหุตฺติกา”ติ, ปฏิคฺคณฺหาติ วีมํสติ ปจฺจาหรติ, อาปตฺติ สํฆาทิเสสสฺส.}

\vspace{\tipitakaparskip}
\csromancenter{นิกฺเขปปทานิ.}
\proseitem{327}{สปริทณฺฑาย เยน ทณฺโฑ ฐปิโต โหติ, โส ภิกฺขุํ ปหิณติ “คจฺฉ ภนฺเต อิตฺถนฺนามํ พฺรูหิ โหตุ อิตฺถนฺนามสฺส ภริยา ธนกฺกีตา จ ฉนฺทวาสินี จ ฯเปฯ ธนกฺกีตา จ โภควาสินี จ ฯเปฯ ธนกฺกีตา จ มุหุตฺติกา จา”ติ, ปฏิคฺคณฺหาติ วีมํสติ ปจฺจาหรติ, อาปตฺติ สํฆาทิเสสสฺส. ขณฺฑจกฺกํ.}

\proseitem{328}{\csromanfolio{213}สปริทณฺฑาย เยน ทณฺโฑ ฐปิโต โหติ, โส ภิกฺขุํ ปหิณติ “คจฺฉ ภนฺเต อิตฺถนฺนามํ พฺรูหิ โหตุ อิตฺถนฺนามสฺส ภริยา ฉนฺทวาสินี จ โภควาสินี จ ฯเปฯ ฉนฺทวาสินี จ มุหุตฺติกา จ. ฉนฺทวาสินี จ ธนกฺกีตา จา”ติ, ปฏิคฺคณฺหาติ วีมํสติ ปจฺจาหรติ, อาปตฺติ สํฆาทิเสสสฺส.}

\vspace{\tipitakaparskip}
\csromancenter{พทฺธจกฺกํ, มูลํ สํขิตฺตํ.}
\proseitemwithcloser{329}{สปริทณฺฑาย เยน ทณฺโฑ ฐปิโต โหติ, โส ภิกฺขุํ ปหิณติ “คจฺฉ ภนฺเต อิตฺถนฺนามํ พฺรูหิ โหตุ อิตฺถนฺนามสฺส ภริยา มุหุตฺติกา จ ธนกฺกีตา จ ฯเปฯ มุหุตฺติกา จ ฉนฺทวาสินี จ ฯเปฯ มุหุตฺติกา จ ธชาหฏา จา”ติ, ปฏิคฺคณฺหาติ วีมํสติ ปจฺจาหรติ, อาปตฺติ สํฆาทิเสสสฺส.}{เอกมูลกํ นิฏฺฐิตํ.}
\prose{ทุมูลกํปิ ติมูลกํปิ ยาว นวมูลกํ เอวเมว กาตพฺพํ.}

\proseitemwithcloser{330}{สปริทณฺฑาย เยน ทณฺโฑ ฐปิโต โหติ, โส ภิกฺขุํ ปหิณติ “คจฺฉ ภนฺเต อิตฺถนฺนามํ พฺรูหิ โหตุ อิตฺถนฺนามสฺส ภริยา ธนกฺกีตา จ ฉนฺทวาสินี จ โภควาสินี จ ปฏวาสินี จ โอทปตฺตกินี จ โอภฏจุมฺพฏา จ ทาสี จ ภริยา จ กมฺมการี จ ภริยา จ ธชาหฏา จ มุหุตฺติกา จา”ติ, ปฏิคฺคณฺหาติ วีมํสติ ปจฺจาหรติ, อาปตฺติ สํฆาทิเสสสฺส.}{ทณฺฑ\-ฐปิต\-จกฺกํ นิฏฺฐิตํ.}
\prose{มาตุรกฺขิตา ภิกฺขุํ ปหิณติ “คจฺฉ ภนฺเต อิตฺถนฺนามํ พฺรูหิ โหมิ อิตฺถนฺนามสฺส ภริยา ธนกฺกีตา”ติ, ปฏิคฺคณฺหาติ วีมํสติ ปจฺจาหรติ, อาปตฺติ สํฆาทิเสสสฺส.}

\prose{มาตุรกฺขิตา ภิกฺขุํ ปหิณติ “คจฺฉ ภนฺเต อิตฺถนฺนามํ พฺรูหิ โหมิ อิตฺถนฺนามสฺส ภริยา ฉนฺทวาสินี ฯเปฯ โภควาสินี. ปฏวาสินี. โอทปตฺตกินี. โอภฏจุมฺพฏา. ทาสี จ ภริยา จ. กมฺมการี จ ภริยา จ. ธชาหฏา. มุหิตฺติกา”ติ, ปฏิคฺคณฺหาติ วีมํสติ ปจฺจาหรติ, อาปตฺติ สํฆาทิเสสสฺส.}

\vspace{\tipitakaparskip}
\csromancenter{นิกฺเขปปทานิ.}
\proseitem{331}{\csromanfolio{214}มาตุรกฺขิตา ภิกฺขุํ ปหิณติ “คจฺฉ ภนฺเต อิตฺถนฺนามํ พฺรูหิ โหมิ อิตฺถนฺนามสฺส ภริยา ธนกฺกีตา จ ฉนฺทวาสินี จา”ติ, ปฏิคฺคณฺหาติ วีมํสติ ปจฺจาหรติ, อาปตฺติ สํฆาทิเสสสฺส.}

\prose{มาตุรกฺขิตา ภิกฺขุํ ปหิณติ “คจฺฉ ภนฺเต อิตฺถนฺนามํ พฺรูหิ โหมิ อิตฺถนฺนามสฺส ภริยา ธนกฺกีตา จ โภควาสินี จ ฯเปฯ ธนกฺกีตา จ ปฏวาสินี จ ฯเปฯ ธนกฺกีตา จ มุหุตฺติกา จา”ติ, ปฏิคฺคณฺหาติ วีมํสติ ปจฺจาหรติ, อาปตฺติ สํฆาทิเสสสฺส.}

\vspace{\tipitakaparskip}
\csromancenter{ขณฺฑจกฺกํ.}
\proseitem{332}{มาตุรกฺขิตา ภิกฺขุํ ปหิณติ “คจฺฉ ภนฺเต อิตฺถนฺนามํ พฺรูหิ โหมิ อิตฺถนฺนามสฺส ภริยา ฉนฺทวาสินิ จ โภควาสินี จ ฯเปฯ ฉนฺทวาสินี จ มุหุตฺติกา จ. ฉนฺทวาสินี จ ธนกฺกีตา จา”ติ, ปฏิคฺคณฺหาติ วีมํสติ ปจฺจาหรติ, อาปตฺติ สํฆาทิเสสสฺส.}

\vspace{\tipitakaparskip}
\csromancenter{พทฺธจกฺกํ, มูลํ สํขิตฺตํ.}
\prosewithcloser{มาตุรกฺขิตา ภิกฺขุํ ปหิณติ “คจฺฉ ภนฺเต อิตฺถนฺนามํ พฺรูหิ โหมิ อิตฺถนฺนามสฺส ภริยา มุหิตฺติกา จ ธนกฺกีตา จ ฯเปฯ มุหุตฺติกา จ ฉนฺทวาสินี จ ฯเปฯ มุหุตฺติกา จ ธชาหฏา จา”ติ, ปฏิคฺคณฺหาติ วีมํสติ ปจฺจาหรติ, อาปตฺติ สํฆาทิเสสสฺส.}{เอกมูลกํ นิฏฺฐิตํ.}
\csromancenter{ทุมูลกาทีนิปิ เอวเมว กาตพฺพานิ.}
\proseitemwithcloser{333}{มาตุรกฺขิตา ภิกฺขุํ ปหิณติ “คจฺฉ ภนฺเต อิตฺถนฺนามํ พฺรูหิ โหมิ อิตฺถนฺนามสฺส ภริยา ธนกฺกีตา จ ฉนฺทวาสินี จ โภควาสินี จ ปฏวาสินี จ โอทปตฺตกินี จ โอภฏจุมฺพฏา จ ทาสี จ ภริยา จ กมฺมการี จ ภริยา จ ธชาหฏา จ มุหุตฺติกา จา”ติ, ปฏิคฺคณฺหาติ วีมํสติ ปจฺจาหรติ, อาปตฺติ สํฆาทิเสสสฺส.}{อปรํ มาตุรกฺขิตา\-จกฺกํ นิฏฺฐิตํ.}
\prose{ปิตุรกฺขิตา ภิกฺขุํ ปหิณติ ฯเปฯ. มาตาปิตุ\-รกฺขิตา ภิกฺขุํ ปหิณติ. ภาตุรกฺขิตา ภิกฺขุํ ปหิณติ. ภคินิ\-รกฺขิตา ภิกฺขุํ ปหิณติ.}

\prose{\csromanfolio{215}ญาติรกฺขิตา ภิกฺขุํ ปหิณติ. โคตฺตรกฺขิตา ภิกฺขุํ ปหิณติ. ธมฺมรกฺขิตา ภิกฺขุํ ปหิณติ. สารกฺขา ภิกฺขุํ ปหิณติ. สปริทณฺฑา ภิกฺขุํ ปหิณติ “คจฺฉ ภนฺเต อิตฺถนฺนามํ พฺรูหิ โหมิ อิตฺถนฺนามสฺส ภริยา ธนกฺกีตา”ติ, ปฏิคฺคณฺหาติ วีมํสติ ปจฺจาหรติ, อาปตฺติ สํฆาทิเสสสฺส.}

\prose{สปริทณฺฑา ภิกฺขุํ ปหิณติ “คจฺฉ ภนฺเต อิตฺถนฺนามํ พฺรูหิ โหมิ อิตฺถนฺนามสฺส ภริยา ฉนฺทวาสินี ฯเปฯ โภควาสินี. ปฏวาสินี. โอทปตฺตกินี. โอภฏจุมฺพฏา. ทาสี จ ภริยา จ. กมฺมการี จ ภริยา จ. ธชาหฏา. มุหุตฺติกา”ติ, ปฏิคฺคณฺหาติ วีมํสติ ปจฺจาหรติ, อาปตฺติ สํฆาทิเสสสฺส.}

\vspace{\tipitakaparskip}
\csromancenter{นิกฺเขปปทานิ.}
\proseitem{334}{สปริทณฺฑา ภิกฺขุํ ปหิณติ “คจฺฉ ภนฺเต อิตฺถนฺนามํ พฺรูหิ โหมิ อิตฺถนฺนามสฺส ภริยา ธนกฺกีตา จ ฉนฺทวาสินี จ ฯเปฯ ธนกฺกีตา จ มุหุตฺติกา จา”ติ, ปฏิคฺคณฺหาติ วีมํสติ ปจฺจาหรติ, อาปตฺติ สํฆาทิเสสสฺส.}

\vspace{\tipitakaparskip}
\csromancenter{ขณฺฑจกฺกํ.}
\proseitem{335}{สปริทณฺฑา ภิกฺขุํ ปหิณติ “คจฺฉ ภนฺเต อิตฺถนฺนามํ พฺรูหิ โหมิ อิตฺถนฺนามสฺส ภริยา ฉนฺทวาสินี จ โภควาสินี จ ฯเปฯ ฉนฺทวาสินี จ มุหุตฺติกา จ. ฉนฺทวาสินี จ ธนกฺกีตา จา”ติ, ปฏิคฺคณฺหาติ วีมํสติ ปจฺจาหรติ, อาปตฺติ สํฆาทิเสสสฺส.}

\vspace{\tipitakaparskip}
\csromancenter{พทฺธจกฺกํ, มูลํ สํขิตฺตํ.}
\proseitemwithcloser{336}{สปริทณฺฑา ภิกฺขุํ ปหิณติ “คจฺฉ ภนฺเต อิตฺถนฺนามํ พฺรูหิ โหมิ อิตฺถนฺนามสฺส ภริยา มุหุตฺติกา จ ธนกฺกีตา จ ฯเปฯ มุหุตฺติกา จ ฉนฺทวาสินี จ ฯเปฯ มุหุตฺติกา จ ธชาหฏา จา”ติ, ปฏิคฺคณฺหาติ วีมํสติ ปจฺจาหรติ, อาปตฺติ สํฆาทิเสสสฺส.}{เอกมูลกํ นิฏฺฐิตํ.}
\csromancenter{ทุมูลกาทีนิปิ เอวเมว กาตพฺพานิ.}
\proseitemwithcloser{337}{\csromanfolio{216}สปริทณฺฑา ภิกฺขุํ ปหิณติ “คจฺฉ ภนฺเต อิตฺถนฺนามํ พฺรูหิ โหมิ อิตฺถนฺนามสฺส ภริยา ธนกฺกีตา จ ฉนฺทวาสินี จ โภควาสินี จ ปฏวาสินี จ โอทปตฺตกินี จ โอภฏจุมฺพฏา จ ทาสี จ ภริยา จ กมฺมการี จ ภริยา จ ธชาหฏา จ มุหุตฺติกา จา”ติ, ปฏิคฺคณฺหาติ วีมํสติ ปจฺจาหรติ, อาปตฺติ สํฆาทิเสสสฺส.}{อปรํ สปริทณฺฑจกฺกํ นิฏฺฐิตํ.}
\nitthitam{สพฺพํ จกฺกเปยฺยาลํ นิฏฺฐิตํ.}
\setlength{\csromangathapagewidth}{0pt}
\setlength{\csromangathawakwidth}{0pt}
\csromangathameasuregroup{ปฏิคฺคณฺหาติ วีมํสติ ปจฺจาหรติ, \\ ปฏิคฺคณฺหาติ วีมํสติ น ปจฺจาหรติ, \\ ปฏิคฺคณฺหาติ น วีมํสติ ปจฺจาหรติ, \\ ปฏิคฺคณฺหาติ น วีมํสติ น ปจฺจาหรติ, \\ น ปฏิคฺคณฺหาติ วีมํสติ ปจฺจาหรติ, \\ น ปฏิคฺคณฺหาติ วีมํสติ น ปจฺจาหรติ, \\ น ปฏิคฺคณฺหาติ น วีมํสติ ปจฺจาหรติ, \\ น ปฏิคฺคณฺหาติ น วีมํสติ น ปจฺจาหรติ,}{\csromangathabat{ปฏิคฺคณฺหาติ วีมํสติ ปจฺจาหรติ,}{อาปตฺติ สํฆาทิเสสสฺส.} \\ \csromangathabat{ปฏิคฺคณฺหาติ วีมํสติ น ปจฺจาหรติ,}{อาปตฺติ ถุลฺลจฺจยสฺส.} \\ \csromangathabat{ปฏิคฺคณฺหาติ น วีมํสติ ปจฺจาหรติ,}{อาปตฺติ ถุลฺลจฺจยสฺส.} \\ \csromangathabat{ปฏิคฺคณฺหาติ น วีมํสติ น ปจฺจาหรติ,}{อาปตฺติ ทุกฺกฏสฺส.} \\ \csromangathabat{น ปฏิคฺคณฺหาติ วีมํสติ ปจฺจาหรติ,}{อาปตฺติ ถุลฺลจฺจยสฺส.} \\ \csromangathabat{น ปฏิคฺคณฺหาติ วีมํสติ น ปจฺจาหรติ,}{อาปตฺติ ทุกฺกฏสฺส.} \\ \csromangathabat{น ปฏิคฺคณฺหาติ น วีมํสติ ปจฺจาหรติ,}{อาปตฺติ ทุกฺกฏสฺส.} \\ \csromangathabat{น ปฏิคฺคณฺหาติ น วีมํสติ น ปจฺจาหรติ,}{อนาปตฺติ.}}
\csromangathasetleft{ปฏิคฺคณฺหาติ วีมํสติ ปจฺจาหรติ, \\ ปฏิคฺคณฺหาติ วีมํสติ น ปจฺจาหรติ, \\ ปฏิคฺคณฺหาติ น วีมํสติ ปจฺจาหรติ, \\ ปฏิคฺคณฺหาติ น วีมํสติ น ปจฺจาหรติ, \\ น ปฏิคฺคณฺหาติ วีมํสติ ปจฺจาหรติ, \\ น ปฏิคฺคณฺหาติ วีมํสติ น ปจฺจาหรติ, \\ น ปฏิคฺคณฺหาติ น วีมํสติ ปจฺจาหรติ, \\ น ปฏิคฺคณฺหาติ น วีมํสติ น ปจฺจาหรติ,}
\csromangathagroup{\csromangathastanza{\csromangathaitembat{338}{ปฏิคฺคณฺหาติ วีมํสติ ปจฺจาหรติ,}{อาปตฺติ สํฆาทิเสสสฺส.} \\ \csromangathacontbat{ปฏิคฺคณฺหาติ วีมํสติ น ปจฺจาหรติ,}{อาปตฺติ ถุลฺลจฺจยสฺส.}}\csromangathastanza{\csromangathabat{ปฏิคฺคณฺหาติ น วีมํสติ ปจฺจาหรติ,}{อาปตฺติ ถุลฺลจฺจยสฺส.} \\ \csromangathabat{ปฏิคฺคณฺหาติ น วีมํสติ น ปจฺจาหรติ,}{อาปตฺติ ทุกฺกฏสฺส.}}\csromangathastanza{\csromangathabat{น ปฏิคฺคณฺหาติ วีมํสติ ปจฺจาหรติ,}{อาปตฺติ ถุลฺลจฺจยสฺส.} \\ \csromangathabat{น ปฏิคฺคณฺหาติ วีมํสติ น ปจฺจาหรติ,}{อาปตฺติ ทุกฺกฏสฺส.}}\csromangathastanza{\csromangathabat{น ปฏิคฺคณฺหาติ น วีมํสติ ปจฺจาหรติ,}{อาปตฺติ ทุกฺกฏสฺส.} \\ \csromangathabat{น ปฏิคฺคณฺหาติ น วีมํสติ น ปจฺจาหรติ,}{อนาปตฺติ.}}}

\prose{ปุริโส สมฺพหุเล ภิกฺขู อาณาเปติ “คจฺฉถ ภนฺเต อิตฺถนฺนามํ อิตฺถิํ วีมํสถา”ติ, สพฺเพ ปฏิคฺคณฺหนฺติ สพฺเพ วีมํสนฺติ สพฺเพ ปจฺจาหรนฺติ, อาปตฺติ สพฺเพสํ สํฆาทิเสสสฺส.}

\prose{ปุริโส สมฺพหุเล ภิกฺขู อาณาเปติ “คจฺฉถ ภนฺเต อิตฺถนฺนามํ อิตฺถิํ วีมํสถา”ติ, สพฺเพ ปฏิคฺคณฺหนฺติ สพฺเพ วีมํสนฺติ เอกํ ปจฺจาหราเปนฺติ, อาปตฺติ สพฺเพสํ สํฆาทิเสสสฺส.}

\prose{ปุริโส สมฺพหุเล ภิกฺขู อาณาเปติ “คจฺฉถ ภนฺเต อิตฺถนฺนามํ อิตฺถิํ วีมํสถา”ติ, สพฺเพ ปฏิคฺคณฺหนฺติ เอกํ วีมํสาเปตฺวา สพฺเพ ปจฺจาหรนฺติ, อาปตฺติ สพฺเพสํ สํฆาทิเสสสฺส.}

\prose{ปุริโส สมฺพหุเล ภิกฺขู อาณาเปติ “คจฺฉถ ภนฺเต อิตฺถนฺนามํ อิตฺถิํ วีมํสถา”ติ, สพฺเพ ปฏิคฺคณฺหนฺติ เอกํ วีมํสาเปตฺวา เอกํ ปจฺจาหราเปนฺติ, อาปตฺติ สพฺเพสํ สํฆาทิเสสสฺส.}

\prose{\csromanfolio{217}ปุริโส ภิกฺขุํ อาณาเปติ “คจฺฉ ภนฺเต อิตฺถนฺนามํ อิตฺถิํ วีมํสา”ติ, ปฏิคฺคณฺหาติ วีมํสติ ปจฺจาหรติ, อาปตฺติ สํฆาทิเสสสฺส.}

\prose{ปุริโส ภิกฺขุํ อาณาเปติ “คจฺฉ ภนฺเต อิตฺถนฺนามํ อิตฺถิํ วีมํสา”ติ, ปฏิคฺคณฺหาติ วีมํสติ อนฺเตวาสิํ ปจฺจาหราเปติ, อาปตฺติ สํฆาทิเสสสฺส.}

\prose{ปุริโส ภิกฺขุํ อาณาเปติ “คจฺฉ ภนฺเต อิตฺถนฺนามํ อิตฺถิํ วีมํสา”ติ, ปฏิคฺคณฺหาติ อนฺเตวาสิํ วีมํสาเปตฺวา อตฺตนา ปจฺจาหรติ, อาปตฺติ สํฆาทิเสสสฺส.}

\prose{ปุริโส ภิกฺขุํ อาณาเปติ “คจฺฉ ภนฺเต อิตฺถนฺนามํ อิตฺถิํ วีมํสา”ติ, ปฏิคฺคณฺหาติ อนฺเตวาสิํ วีมํสาเปติ, อนฺเตวาสี วีมํสิตฺวา พหิทฺธา ปจฺจาหรติ, อาปตฺติ อุภินฺนํ ถุลฺลจฺจยสฺส.}

\proseitem{339}{คจฺฉนฺโต สมฺปาเทติ, อาคจฺฉนฺโต วิสํวาเทติ, อาปตฺติ ถุลฺลจฺจยสฺส.}

\prose{คจฺฉนฺโต วิสํวาเทติ, อาคจฺฉนฺโต สมฺปาเทติ, อาปตฺติ ถุลฺลจฺจยสฺส.}

\prose{คจฺฉนฺโต สมฺปาเทติ, อาคจฺฉนฺโต สมฺปาเทติ, อาปตฺติ สํฆาทิเสสสฺส.}

\prose{คจฺฉนฺโต วิสํวาเทติ, อาคจฺฉนฺโต วิสํวาเทติ, อนาปตฺติ.}

\proseitem{340}{อนาปตฺติ สํฆสฺส วา เจติยสฺส วา คิลานสฺส วา กรณีเยน คจฺฉติ, อุมฺมตฺตกสฺส อาทิกมฺมิกสฺสาติ.}
\csromansectionrule

\csromanmatikaanchor{mtk.66}
\csromantocmarkref{subsection}{วินีตวตฺถุอุทฺทานคาถา}{mtk.66}
\bookmark[dest={mtk.66},level=2]{วินีตวตฺถุอุทฺทานคาถา}
\csromanheader{2}{\textbf{วินีตวตฺถุ}\-อุทฺทานคาถา}
\setlength{\csromangathapagewidth}{0pt}
\setlength{\csromangathawakwidth}{0pt}
\csromangathameasuregroup{สุตฺตา มตา จ นิกฺขนฺตา, \\ กลหํ กตฺวาน สมฺโพทิ,}{\csromangathabat{สุตฺตา มตา จ นิกฺขนฺตา,}{อนิตฺถี อิตฺถิปณฺฑกา.} \\ \csromangathabat{กลหํ กตฺวาน สมฺโพทิ,}{สญฺจริตฺตญฺจ ปณฺฑเกติ.}}
\csromangathasetleft{สุตฺตา มตา จ นิกฺขนฺตา, \\ กลหํ กตฺวาน สมฺโพทิ,}
\csromangathagroup{\csromangathastanza{\csromangathabat{สุตฺตา มตา จ นิกฺขนฺตา,}{อนิตฺถี อิตฺถิปณฺฑกา.} \\ \csromangathabat{กลหํ กตฺวาน สมฺโพทิ,}{สญฺจริตฺตญฺจ ปณฺฑเกติ.}}}

\csromanmatikaanchor{mtk.67}
\csromantocmarkref{subsection}{วินีตวตฺถุ}{mtk.67}
\bookmark[dest={mtk.67},level=2]{วินีตวตฺถุ}
\csromanheader{2}{\textbf{วินีตวตฺถุ}}
\csromanmatikaanchor{mtk.68}
\csromanmatikaanchor{mtk.69}
\csromantocmarknopage{section}{6. กุฏิการสิกฺขาปท}{mtk.68}
\bookmark[dest={mtk.68},level=1]{6. กุฏิการสิกฺขาปท}
\csromantocmarkref{subsection}{อาฬวกภิกฺขุวตฺถุ}{mtk.69}
\bookmark[dest={mtk.69},level=2]{อาฬวกภิกฺขุวตฺถุ}
\proseitem{341}{เตน โข ปน สมเยน อญฺญตโร ปุริโส อญฺญตรํ ภิกฺขุํ อาณาเปสิ\csromansharedfootnote{662544f7ef3efa8d}{อาณาเปติ (สฺยา, ก)} “คจฺฉ ภนฺเต อิตฺถนฺนามํ อิตฺถิํ วีมํสา”ติ. โส คนฺตฺวา มนุสฺเส ปุจฺฉิ “กหํ อิตฺถนฺนามา”ติ. สุตฺตา ภนฺเตติ. ตสฺส \csromanfolio{218}กุกฺกุจฺจํ อโหสิ “ภควตา สิกฺขาปทํ ปญฺญตฺตํ, กจฺจิ นุ โข อหํ สํฆาทิเสสํ อาปตฺติํ อาปนฺโน”ติ. ภควโต เอตมตฺถํ อาโรเจสิ. อนาปตฺติ ภิกฺขุ สํฆาทิเสสสฺส, อาปตฺติ ทุกฺกฏสฺสาติ. (1)}

\prose{เตน โข ปน สมเยน อญฺญตโร ปุริโส อญฺญตรํ ภิกฺขุํ อาณาเปสิ “คจฺฉ ภนฺเต อิตฺถนฺนามํ อิตฺถิํ วีมํสา”ติ. โส คนฺตฺวา มนุสฺเส ปุจฺฉิ “กหํ อิตฺถนฺนามา”ติ. มตา ภนฺเตติ ฯเปฯ. นิกฺขนฺตา ภนฺเตติ. อนิตฺถี ภนฺเตติ. อิตฺถิปณฺฑกา ภนฺเตติ. ตสฺส กุกฺกุจฺจํ อโหสิ ฯเปฯ. อนาปตฺติ ภิกฺขุ สํฆาทิเสสสฺส, อาปตฺติ ทุกฺกฏสฺสาติ. \mbox{(2-5)}}

\prose{เตน โข ปน สมเยน อญฺญตรา อิตฺถี สามิเกน สห ภณฺฑิตฺวา มาตุฆรํ อคมาสิ. กุลูปโก ภิกฺขุ สมฺโมทนียํ อกาสิ. ตสฺส กุกฺกุจฺจํ อโหสิ ฯเปฯ. อลํวจนียา ภิกฺขูติ. นาลํวจนียา ภควาติ. อนาปตฺติ ภิกฺขุ นาลํ\-วจนียายาติ. (6)}

\prosewithcloserruled{เตน โข ปน สมเยน อญฺญตโร ภิกฺขุ ปณฺฑเก สญฺจริตฺตํ สมาปชฺชิ. ตสฺส กุกฺกุจฺจํ อโหสิ ฯเปฯ. อนาปตฺติ ภิกฺขุ สํฆาทิเสสสฺส, อาปตฺติ ถุลฺลจฺจยสฺสาติ. (7)}{สญฺจริตฺต\-สิกฺขาปทํ นิฏฺฐิตํ ปญฺจมํ.}
\csromanheader{1}{6. \textbf{กุฏิการ}\-\textbf{สิกฺขาปท}}
\proseitem{342}{เตน สมเยน พุทฺโธ ภควา ราชคเห วิหรติ เวฬุวเน กลนฺทก\-นิวาเป. เตน โข ปน สมเยน อาฬวกา ภิกฺขู สญฺญาจิกาโย กุฏิโย การาเปนฺติ อสฺสามิกาโย อตฺตุทฺเทสิกาโย อปฺปมาณิกาโย, ตาโย น นิฏฺฐานํ คจฺฉนฺติ, เต ยาจนพหุลา วิญฺญตฺติพหุลา วิหรนฺติ “ปุริสํ เทถ ปุริส\-ตฺถกรํ เทถ โคณํ เทถ สกฏํ เทถ วาสิํ เทถ ปรสุํ เทถ กุฐาริํ เทถ กุทาลํ เทถ นิขาทนํ เทถ วลฺลิํ เทถ เวฬุํ เทถ มุญฺชํ เทถ ปพฺพชํ เทถ ติณํ เทถ มตฺติกํ เทถา”ติ. มนุสฺสา อุปทฺทุตา ยาจนาย อุปทฺทุตา วิญฺญตฺติยา ภิกฺขู ทิสฺวา อุพฺพิชฺชนฺติปิ อุตฺตสนฺติปิ ปลายนฺติปิ อญฺเญนปิ คจฺฉนฺติ อญฺเญนปิ มุขํ กโรนฺติ ทฺวารํปิ ถเกนฺติ คาวิมฺปิ ทิสฺวา ปลายนฺติ ภิกฺขูติ มญฺญมานา.}

\prose{\csromanfolio{219}อถ โข อายสฺมา มหากสฺสโป ราชคเห วสฺสํวุฏฺโฐ เยน อาฬวี เตน ปกฺกามิ. อนุปุพฺเพน เยน อาฬวี ตทวสริ. ตตฺร สุทํ อายสฺมา มหากสฺสโป อาฬวิยํ วิหรติ อคฺคาฬเว เจติเย. อถ โข อายสฺมา มหากสฺสโป ปุพฺพณฺหสมยํ นิวาเสตฺวา ปตฺต\-จีวรมา\-ทาย อาฬวิํ ปิณฺฑาย ปาวิสิ. มนุสฺสา อายสฺมนฺตํ มหากสฺสปํ ปสฺสิตฺวา อุพฺพิชฺชนฺติปิ อุตฺตสนฺติปิ ปลายนฺติปิ อญฺเญนปิ คจฺฉนฺติ อญฺเญนปิ มุขํ กโรนฺติ ทฺวารํปิ ถเกนฺติ. อถ โข อายสฺมา มหากสฺสโป อาฬวิยํ ปิณฺฑาย จริตฺวา ปจฺฉาภตฺตํ ปิณฺฑปาต\-ปฺปฏิกฺกนฺโต ภิกฺขู อามนฺเตสิ “ปุพฺพายํ อาวุโส อาฬวี สุภิกฺขา อโหสิ สุลภปิณฺฑา สุกรา อุญฺเฉน ปคฺคเหน ยาเปตุํ, เอตรหิ ปนายํ อาฬวี ทุพฺภิกฺขา ทุลฺลภปิณฺฑา น สุกรา อุญฺเฉน ปคฺคเหน ยาเปตุํ, โก นุ โข อาวุโส เหตุ โก ปจฺจโย, เยนายํ อาฬวี ทุพฺภิกฺขา ทุลฺลภปิณฺฑา น สุกรา อุญฺเฉน ปคฺคเหน ยาเปตุนฺ”ติ. อถ โข เต ภิกฺขู อายสฺมโต มหากสฺสปสฺส เอตมตฺถํ อาโรเจสุํ.}

\csromanmatikaanchor{mtk.70}
\csromantocmarkref{subsection}{มณิกณฺฐนคราชวตฺถุ}{mtk.70}
\bookmark[dest={mtk.70},level=2]{มณิกณฺฐนคราชวตฺถุ}
\proseitem{343}{อถ โข ภควา ราชคเห ยถาภิรนฺตํ วิหริตฺวา เยน อาฬวี เตน จาริกํ ปกฺกามิ, อนุปุพฺเพน จาริกํ จรมาโน เยน อาฬวี ตทวสริ. ตตฺรสุทํ ภควา อาฬวิยํ วิหรติ อคฺคาฬเว เจติเย. อถ โข อายสฺมา มหากสฺสโป เยน ภควา เตนุปสงฺกมิ, อุปสงฺกมิตฺวา ภควนฺตํ อภิวาเทตฺวา เอกมนฺตํ นิสีทิ, เอกมนฺตํ นิสินฺโน โข อายสฺมา มหากสฺสโป ภควโต เอตมตฺถํ อาโรเจสิ. อถ โข ภควา เอตสฺมิํ นิทาเน เอตสฺมิํ ปกรเณ ภิกฺขุสํฆํ สนฺนิปาตาเปตฺวา อาฬวเก ภิกฺขู ปฏิปุจฺฉิ “สจฺจํ กิร ตุเมฺห ภิกฺขเว สญฺญาจิกาโย กุฏิโย การาเปถ อสฺสามิกาโย อตฺตุทฺเทสิกาโย อปฺปมาณิกาโย, ตาโย น นิฏฺฐานํ คจฺฉนฺติ, เต ตุเมฺห ยาจนพหุลา วิญฺญตฺติพหุลา วิหรถ ‘ปุริสํ เทถ ปุริส\-ตฺถกรํ เทถ ฯเปฯ ติณํ เทถ มตฺติกํ เทถา’ติ, มนุสฺสา อุปทฺทุตา ยาจนาย อุปทฺทุตา วิญฺญตฺติยา ภิกฺขู ทิสฺวา อุพฺพิชฺชนฺติปิ อุตฺตสนฺติปิ ปลายนฺติปิ อญฺเญนปิ คจฺฉนฺติ อญฺเญนปิ มุขํ กโรนฺติ ทฺวารํปิ ถเกนฺติ คาวิมฺปิ ทิสฺวา ปลายนฺติ ภิกฺขูติ มญฺญมานา”ติ. สจฺจํ ภควาติ. วิครหิ พุทฺโธ ภควา ฯเปฯ. กถํ หิ นาม ตุเมฺห โมฆปุริสา สํยาจิกาโย กุฏิโย การาเปสฺสถ อสฺสามิกาโย อตฺตุทฺเทสิกาโย อปฺปมาณิกาโย, ตาโย \csromanfolio{220}น นิฏฺฐานํ คจฺฉนฺติ, เต ตุเมฺห ยาจนพหุลา วิญฺญตฺติพหุลา วิหริสฺสถ ปุริสํ เทถ ปุริส\-ตฺถกรํ เทถ ฯเปฯ ติณํ เทถ มตฺติกํ เทถาติ, เนตํ โมฆปุริสา อปฺปสนฺนานํ วา ปสาทาย ฯเปฯ วิครหิตฺวา ธมฺมิํ กถํ กตฺวา ภิกฺขู อามนฺเตสิ\csromandash{}}

\proseitem{344}{ภูตปุพฺพํ ภิกฺขเว เทฺว ภาตโร อิสโย คงฺคํ นทิํ อุปนิสฺสาย วิหริํสุ. อถ โข ภิกฺขเว มณิกณฺโฐ นาคราชา คงฺคํ นทิํ อุตฺตริตฺวา เยน กนิฏฺโฐ อิสิ เตนุปสงฺกมิ, อุปสงฺกมิตฺวา กนิฏฺฐํ อิสิํ สตฺตกฺขตฺตุํ โภเคหิ ปริกฺขิปิตฺวา อุปริมุทฺธนิ มหนฺตํ ผณํ กริตฺวา อฏฺฐาสิ. อถ โข ภิกฺขเว กนิฏฺโฐ อิสิ ตสฺส นาคสฺส ภยา กิโส อโหสิ ลูโข ทุพฺพณฺโณ อุปฺปณฺฑุ\-ปฺปณฺฑุก\-ชาโต ธมนิสนฺถต\-คตฺโต. อทฺทส โข ภิกฺขเว เชฏฺโฐ อิสิ กนิฏฺฐํ อิสิํ กิสํ ลูขํ ทุพฺพณฺณํ อุปฺปณฺฑุ\-ปฺปณฺฑุก\-ชาตํ ธมนิสนฺถต\-คตฺตํ. ทิสฺวาน กนิฏฺฐํ อิสิํ เอตทโวจ “กิสฺส ตฺวํ โภ กิโส ลูโข ทุพฺพณฺโณ อุปฺปณฺฑุ\-ปฺปณฺฑุก\-ชาโต ธมนิสนฺถต\-คตฺโต”ติ. อิธ โภ มณิกณฺโฐ นาคราชา คงฺคํ นทิํ อุตฺตริตฺวา เยนาหํ เตนุปสงฺกมิ, อุปสงฺกมิตฺวา มํ สตฺตกฺขตฺตุํ โภเคหิ ปริกฺขิปิตฺวา อุปริมุทฺธนิ มหนฺตํ ผณํ กริตฺวา อฏฺฐาสิ, ตสฺสาหํ โภ นาคสฺส ภยา\csromansharedfootnote{1792d7be297d59db}{ภยามฺหิ (สี)} กิโส ลูโข ทุพฺพณฺโณ อุปฺปณฺฑุ\-ปฺปณฺฑุก\-ชาโต ธมนิสนฺถต\-คตฺโตติ. อิจฺฉสิ ปน ตฺวํ โภ ตสฺส นาคสฺส อนาคมนนฺติ. อิจฺฉามหํ โภ ตสฺส นาคสฺส อนาคมนนฺติ. อปิ ปน ตฺวํ โภ ตสฺส นาคสฺส กิญฺจิ ปสฺสสีติ. ปสฺสามหํ โภ มณิมสฺส\csromansharedfootnote{350b251964d1bdb0}{มณิสฺส (สี, ก)} กณฺเฐ ปิลนฺธนนฺติ. เตน หิ ตฺวํ โภ ตํ นาคํ มณิํ ยาจ “มณิํ เม โภ เทหิ, มณินา เม อตฺโถ”ติ.}

\prose{อถ โข ภิกฺขเว มณิกณฺโฐ นาคราชา คงฺคํ นทิํ อุตฺตริตฺวา เยน กนิฏฺโฐ อิสิ เตนุปสงฺกมิ, อุปสงฺกมิตฺวา เอกมนฺตํ อฏฺฐาสิ, เอกมนฺตํ ฐิตํ โข ภิกฺขเว มณิกณฺฐํ นาคราชานํ กนิฏฺโฐ อิสิ เอตทโวจ “มณิํ เม โภ เทหิ, มณินา เม อตฺโถ”ติ. อถ โข ภิกฺขเว มณิกณฺโฐ นาคราชา “ภิกฺขุ มณิํ ยาจติ, ภิกฺขุสฺส มณินา อตฺโถ”ติ ขิปฺปญฺเญว อคมาสิ. ทุติยมฺปิ โข ภิกฺขเว มณิกณฺโฐ นาคราชา คงฺคํ นทิํ อุตฺตริตฺวา เยน กนิฏฺโฐ อิสิ เตนุปสงฺกมิ. \csromanfolio{221}อทฺทส โข ภิกฺขเว กนิฏฺโฐ อิสิ มณิกณฺฐํ นาคราชานํ ทูรโตว อาคจฺฉนฺตํ, ทิสฺวาน มณิกณฺฐํ นาคราชานํ เอตทโวจ “มณิํ เม โภ เทหิ, มณินา เม อตฺโถ”ติ. อถ โข ภิกฺขเว มณิกณฺโฐ นาคราชา “ภิกฺขุ มณิํ ยาจติ, ภิกฺขุสฺส มณินา อตฺโถ”ติ ตโตว ปฏินิวตฺติ. ตติยมฺปิ โข ภิกฺขเว มณิกณฺโฐ นาคราชา คงฺคํ นทิํ อุตฺตรติ. อทฺทส โข ภิกฺขเว กนิฏฺโฐ อิสิ มณิกณฺฐํ นาคราชานํ คงฺคํ นทิํ อุตฺตรนฺตํ, ทิสฺวาน มณิกณฺฐํ นาคราชานํ เอตทโวจ “มณิํ เม โภ เทหิ, มณินา เม อตฺโถ”ติ. อถ โข ภิกฺขเว มณิกณฺโฐ นาคราชา กนิฏฺฐํ อิสิํ คาถาหิ อชฺฌภาสิ\csromandash{}}

\setlength{\csromangathapagewidth}{0pt}
\setlength{\csromangathawakwidth}{0pt}
\csromangathameasuregroup{}{\textsuperscript{*}“มมนฺนปานํ วิปุลํ อุฬารํ, \\ อุปฺปชฺชตีมสฺส มณิสฺส เหตุ. \\ ตํ เต น ทสฺสํ อติยาจเก, \\ น จาปิ เต อสฺสมมาคมิสฺสํ. \\ \textsuperscript{*}สุสู ยถา สกฺขรโธตปาณี, \\ ตาเสสฺสิ มํ เสลมายาจมาโน. \\ ตํ เต น ทสฺสํ อติยาจโกสิ, \\ น จาปิ เต อสฺสมมาคมิสฺสนฺ”ติ.}
\csromangathameasurewak{\textsuperscript{*}“มมนฺนปานํ วิปุลํ อุฬารํ, \\ อุปฺปชฺชตีมสฺส มณิสฺส เหตุ. \\ ตํ เต น ทสฺสํ อติยาจเก, \\ น จาปิ เต อสฺสมมาคมิสฺสํ. \\ \textsuperscript{*}สุสู ยถา สกฺขรโธตปาณี, \\ ตาเสสฺสิ มํ เสลมายาจมาโน. \\ ตํ เต น ทสฺสํ อติยาจโกสิ, \\ น จาปิ เต อสฺสมมาคมิสฺสนฺ”ติ.}
\setlength{\csromangathaleftcol}{0pt}
\csromangathagroup{\csromangathastanza{\csromangathawak{\csromansymbolfootnote{*}{ขุ 5. 71-72 ปิฏฺเฐสุ มณิกณฺฐชาตเกปิ.}“มมนฺนปานํ วิปุลํ อุฬารํ, \\ อุปฺปชฺชตีมสฺส มณิสฺส เหตุ. \\ ตํ เต น ทสฺสํ อติยาจเก, \\ น จาปิ เต อสฺสมมาคมิสฺสํ.}}\csromangathastanza{\csromangathawak{\csromansymbolmark{*}สุสู ยถา สกฺขร\-โธต\-ปาณี, \\ ตาเสสฺสิ มํ เสลมายาจมาโน. \\ ตํ เต น ทสฺสํ อติยาจโกสิ, \\ น จาปิ เต อสฺสมมาคมิสฺสนฺ”ติ.}}}

\prose{อถ โข ภิกฺขเว มณิกณฺโฐ นาคราชา “ภิกฺขุ มณิํ ยาจติ ภิกฺขุสฺส มณินา อตฺโถ”ติ ปกฺกามิ, ตถา ปกฺกนฺโตว\csromansharedfootnote{13d5a98903ae8188}{ตทา ปกฺกนฺโตว (ก)} อโหสิ, น ปุน ปจฺจาคญฺฉิ. อถ โข ภิกฺขเว กนิฏฺโฐ อิสิ ตสฺส นาคสฺส ทสฺสนียสฺส อทสฺสเนน ภิยฺโยโส\-มตฺตาย กิโส อโหสิ ลูโข ทุพฺพณฺโณ อุปฺปณฺฑุ\-ปฺปณฺฑุก\-ชาโต ธมนิสนฺถต\-คตฺโต. อทฺทส โข ภิกฺขเว เชฏฺโฐ อิสิ กนิฏฺฐํ อิสิํ ภิยฺโยโส\-มตฺตาย กิสํ ลูขํ ทุพฺพณฺณํ อุปฺปณฺฑุ\-ปฺปณฺฑุก\-ชาตํ ธมนิสนฺถต\-คตฺตํ, ทิสฺวาน กนิฏฺฐํ อิสิํ เอตทโวจ “กิสฺส ตฺวํ โภ ภิยฺโยโส\-มตฺตาย กิโส ลูโข ทุพฺพณฺโณ อุปฺปณฺฑุ\-ปฺปณฺฑุก\-ชาโต ธมนิสนฺถต\-คตฺโต”ติ. ตสฺสาหํ โภ นาคสฺส ทสฺสนียสฺส อทสฺสเนน ภิยฺโยโส\-มตฺตาย กิโส ลูโข ทุพฺพณฺโณ อุปฺปณฺฑุ\-ปฺปณฺฑุก\-ชาโต ธมนิสนฺถต\-คตฺโตติ. อถ โข ภิกฺขเว เชฏฺโฐ อิสิ กนิฏฺฐํ อิสิํ คาถาย อชฺฌาภาสิ\csromandash{}}

\setlength{\csromangathapagewidth}{0pt}
\setlength{\csromangathawakwidth}{0pt}
\csromangathameasuregroup{}{\textsuperscript{*}“น ตํ ยาเจ ยสฺส ปิยํ ชิคีเส, \\ วิเทสฺโส\textsuperscript{1} โหติ อติยาจนาย. \\ นาโค มณิํ ยาจิโต พฺราหฺมเณน, \\ อทสฺสนญฺเญว ตทชฺฌคมา”ติ.}
\csromangathameasurewak{\textsuperscript{*}“น ตํ ยาเจ ยสฺส ปิยํ ชิคีเส, \\ วิเทสฺโส\textsuperscript{1} โหติ อติยาจนาย. \\ นาโค มณิํ ยาจิโต พฺราหฺมเณน, \\ อทสฺสนญฺเญว ตทชฺฌคมา”ติ.}
\setlength{\csromangathaleftcol}{0pt}
\csromangathagroup{\csromangathastanza{\csromangathawak{\csromanfolio{222}\csromansymbolfootnote{*}{ขุ 5. 72 ปิฏฺเฐ มณิกณฺฐชาตเกปิ.}“น ตํ ยาเจ ยสฺส ปิยํ ชิคีเส, \\ วิเทสฺโส\csromansharedfootnote{02178f58a5f53f0f}{เทสฺโส (สี), เทสฺโส จ (สฺยา)} โหติ อติยาจนาย. \\ นาโค มณิํ ยาจิโต พฺราหฺมเณน, \\ อทสฺสนญฺเญว ตทชฺฌคมา”ติ.}}}

\prose{เตสํ หิ นาม ภิกฺขเว ติรจฺฉาน\-คตานํ ปาณานํ อมนาปา ภวิสฺสติ ยาจนา อมนาปา วิญฺญตฺติ, กิมงฺคํ\csromansharedfootnote{000f45a8a6c5a7bf}{กิมงฺค (สี)} ปน มนุสฺส\-ภูตานํ.}

\proseitem{345}{ภูตปุพฺพํ ภิกฺขเว อญฺญตโร ภิกฺขุ หิมวนฺตปสฺเส วิหรติ อญฺญตรสฺมิํ วนสณฺเฑ. ตสฺส โข ภิกฺขเว วนสณฺฑสฺส อวิทูเร มหนฺตํ นินฺนํ ปลฺลลํ, อถ โข ภิกฺขเว มหาสกุณสํโฆ ตสฺมิํ ปลฺลเล ทิวสํ โคจรํ จริตฺวา สายํ ตํ วนสณฺฑํ วาสาย อุปคจฺฉติ. อถ โข ภิกฺขเว โส ภิกฺขุ ตสฺส สกุณ\-สํฆสฺส สทฺเทน อุพฺพาโฬฺห เยนาหํ เตนุปสงฺกมิ, อุปสงฺกมิตฺวา มํ อภิวาเทตฺวา เอกมนฺตํ นิสีทิ, เอกมนฺตํ นิสินฺนํ โข อหํ ภิกฺขเว ตํ ภิกฺขุํ เอตทโวจ “กจฺจิ ภิกฺขุ ขมนียํ, กจฺจิ ยาปนียํ, กจฺจิสิ อปฺป\-กิลมเถน อทฺธานํ อาคโต, กุโต จ ตฺวํ ภิกฺขุ อาคจฺฉสี”ติ. ขมนียํ ภควา, ยาปนียํ ภควา, อปฺป\-กิลมเถน จาหํ ภนฺเต อทฺธานํ อาคโต, อตฺถิ ภนฺเต หิมวนฺตปสฺเส มหาวนสณฺโฑ, ตสฺส โข ปน ภนฺเต วนสณฺฑสฺส อวิทูเร มหนฺตํ นินฺนํ ปลฺลลํ, อถ โข ภนฺเต มหาสกุณสํโฆ ตสฺมิํ ปลฺลเล ทิวสํ โคจรํ จริตฺวา สายํ ตํ วนสณฺฑํ วาสาย อุปคจฺฉติ, ตโต อหํ ภควา อาคจฺฉามิ ตสฺส สกุณ\-สํฆสฺส สทฺเทน อุพฺพาโฬฺหติ. อิจฺฉสิ ปน ตฺวํ ภิกฺขุ ตสฺส สกุณ\-สํฆสฺส อนาคมนนฺติ. อิจฺฉามหํ ภควา ตสฺส สกุณ\-สํฆสฺส อนาคมนนฺติ. เตน หิ ตฺวํ ภิกฺขุ ตตฺถ คนฺตฺวา ตํ วนสณฺฑํ อชฺโฌคาเหตฺวา รตฺติยา ปฐมํ ยามํ ติกฺขตฺตุํ สทฺทมนุสฺสาเวหิ “สุณนฺตุ เม โภนฺโต สกุณา ยาวติกา อิมสฺมิํ วนสณฺเฑ วาสํ อุปคตา, ปตฺเตน เม อตฺโถ, เอเกกํ เม โภนฺโต ปตฺตํ ททนฺตู”ติ. รตฺติยา มชฺฌิมํ ยามํ. รตฺติยา ปจฺฉิมํ ยามํ ติกฺขตฺตุํ สทฺทมนุสฺสาเวหิ “สุณนฺตุ เม โภนฺโต สกุณา ยาวติกา อิมสฺมิํ วนสณฺเฑ วาสํ อุปคตา, ปตฺเตน เม อตฺโถ, เอเกกํ เม โภนฺโต ปตฺตํ ททนฺตู”ติ.}

\csromanmatikaanchor{mtk.71}
\csromantocmarkref{subsection}{ปญฺญตฺติ}{mtk.71}
\bookmark[dest={mtk.71},level=2]{ปญฺญตฺติ}
\prose{\csromanfolio{223}อถ โข ภิกฺขเว โส ภิกฺขุ ตตฺถ คนฺตฺวา ตํ วนสณฺฑํ อชฺโฌคาเหตฺวา รตฺติยา ปฐมํ ยามํ ติกฺขตฺตุํ สทฺทม\-นุสฺสาเวสิ “สุณนฺตุ เม โภนฺโต สกุณา ยาวติกา อิมสฺมิํ วนสณฺเฑ วาสํ อุปคตา, ปตฺเตน เม อตฺโถ, เอเกกํ เม โภนฺโต ปตฺตํ ททนฺตู”ติ. รตฺติยา มชฺฌิมํ ยามํ. รตฺติยา ปจฺฉิมํ ยามํ ติกฺขตฺตุํ สทฺทม\-นุสฺสาเวสิ “สุณนฺตุ เม โภนฺโต สกุณา ยาวติกา อิมสฺมิํ วนสณฺเฑ วาสํ อุปคตา, ปตฺเตน เม อตฺโถ, เอเกกํ เม โภนฺโต ปตฺตํ ททนฺตู”ติ. อถ โข ภิกฺขเว โส สกุณสํโฆ “ภิกฺขุ ปตฺตํ ยาจติ ภิกฺขุสฺส ปตฺเตน อตฺโถ”ติ ตมฺหา วนสณฺฑา ปกฺกามิ, ตถา ปกฺกนฺโตว อโหสิ, น ปุน ปจฺจาคญฺฉิ. เตสํ หิ นาม ภิกฺขเว ติรจฺฉาน\-คตานํ ปาณานํ อมนาปา ภวิสฺสติ ยาจนา อมนาปา วิญฺญตฺติ, กิมงฺคํ ปน มนุสฺส\-ภูตานํ.}

\proseitem{346}{ภูตปุพฺพํ ภิกฺขเว รฏฺฐปาลสฺส กุลปุตฺตสฺส ปิตา รฏฺฐปาลํ กุลปุตฺตํ คาถาย อชฺฌาภาสิ\csromandash{}}

\setlength{\csromangathapagewidth}{0pt}
\setlength{\csromangathawakwidth}{0pt}
\csromangathameasuregroup{“อปาหํ เต น ชนามิ, \\ \textsuperscript{*}เต มํ สํคมฺม ยาจนฺติ, \\ \textsuperscript{*}ยาจโก อปฺปิโย โหติ, \\ ตสฺมาหํ ตํ น ยาจามิ,}{\csromangathabat{“อปาหํ เต น ชนามิ,}{รฏฺฐปาล พหู ชนา.} \\ \csromangathabat{\textsuperscript{*}เต มํ สํคมฺม ยาจนฺติ,}{กสฺมา มํ ตฺวํ น ยาจสีติ.} \\ \csromangathabat{\textsuperscript{*}ยาจโก อปฺปิโย โหติ,}{ยาจํ อททมปฺปิโย.} \\ \csromangathabat{ตสฺมาหํ ตํ น ยาจามิ,}{มา เม วิเทสฺสนา อหู”ติ.}}
\csromangathasetleft{“อปาหํ เต น ชนามิ, \\ \textsuperscript{*}เต มํ สํคมฺม ยาจนฺติ, \\ \textsuperscript{*}ยาจโก อปฺปิโย โหติ, \\ ตสฺมาหํ ตํ น ยาจามิ,}
\csromangathagroup{\csromangathastanza{\csromangathabat{“อปาหํ เต น ชนามิ,}{รฏฺฐปาล พหู ชนา.} \\ \csromangathabat{\csromansymbolfootnote{*}{ขุ 5. 161 ปิฏฺเฐปิ.}เต มํ สํคมฺม ยาจนฺติ,}{กสฺมา มํ ตฺวํ น ยาจสีติ.}}\csromangathastanza{\csromangathabat{\csromansymbolmark{*}ยาจโก อปฺปิโย โหติ,}{ยาจํ อททมปฺปิโย.} \\ \csromangathabat{ตสฺมาหํ ตํ น ยาจามิ,}{มา เม วิเทสฺสนา อหู”ติ.}}}

\prose{โส หิ นาม ภิกฺขเว รฏฺฐปาโล กุลปุตฺโต สกํ ปิตรํ เอวํ วกฺขติ, กิมงฺคํ ปน ชโน ชนํ.}

\proseitem{347}{คิหีนํ ภิกฺขเว ทุสฺสํหรานิ โภคานิ สมฺภตานิปิ ทุรกฺขิยานิ. ตตฺถ นาม ตุเมฺห โมฆปุริสา เอวํ ทุสฺสํหเรสุ โภเคสุ สมฺภเตสุปิ ทุรกฺขิเยสุ ยาจนพหุลา วิญฺญตฺติพหุลา วิหริสฺสถ “ปุริสํ เทถ ปุริส\-ตฺถกรํ เทถ โคณํ เทถ สกฏํ เทถ วาสิํ เทถ ปรสุํ เทถ กุฐาริํ เทถ กุทาลํ เทถ นิขาทนํ เทถ วลฺลิํ เทถ เวฬุํ เทถ มุญฺชํ เทถ ปพฺพชํ เทถ ติณํ เทถ มตฺติกํ เทถา”ติ, เนตํ โมฆปุริสา อปฺปสนฺนานํ วา ปสาทาย ฯเปฯ. เอวญฺจ ปน ภิกฺขเว อิมํ สิกฺขาปทํ อุทฺทิเสยฺยาถ\csromandash{}}

\csromanmatikaanchor{mtk.72}
\csromantocmarkref{subsection}{สิกฺขาปทวิภงฺค, ปทภาชนีย}{mtk.72}
\bookmark[dest={mtk.72},level=2]{สิกฺขาปทวิภงฺค, ปทภาชนีย}
\proseitem{348}{\textbf{“สญฺญาจิกาย ปน ภิกฺขุนา กุฏิํ การยมาเนน อสฺสามิกํ อตฺตุทฺเทสํ ปมาณิกา กาเรตพฺพา, ตตฺริทํ ปมาณํ, ทีฆโส} \csromanfolio{224}\textbf{ทฺวาทส วิทตฺถิโย สุคต}\-\textbf{วิทตฺถิยา, ติริยํ สตฺตนฺตรา, ภิกฺขู อภิเนตพฺพา วตฺถุเทสนาย, เตหิ ภิกฺขูหิ วตฺถุ เทเสตพฺพํ อนารมฺภํ}\csromansharedfootnote{4579198a4ceb92b8}{อนารพฺภํ (ก)} \textbf{สปริกฺกมนํ, สารมฺเภ}\csromansharedfootnote{b2cc0084ea80d109}{สารพฺเภ (ก)} \textbf{เจ ภิกฺขุ วตฺถุสฺมิํ อปริกฺกมเน สญฺญาจิกาย กุฏิํ กาเรยฺย, ภิกฺขู วา อนภิเนยฺย วตฺถุเทสนาย, ปมาณํ วา อติกฺกาเมยฺย, สํฆาทิเสโส”}ติ.}

\proseitem{349}{\textbf{สญฺญาจิกา} นาม สยํ ยาจิตฺวา ปุริสํปิ ปุริส\-ตฺถกรํปิ โคณํปิ สกฏํปิ วาสิํปิ ปรสุํปิ กุฐาริํปิ กุทาลํปิ นิขาทนํปิ ปลฺลิํปิ เวฬุํปิ มุญฺชํปิ ปพฺพชํปิ ติณํปิ มตฺติกํปิ.}

\prose{\textbf{กุฏิ} นาม อุลฺลิตฺตา วา โหติ อวลิตฺตา วา อุลฺลิตฺตาวลิตฺตา วา.}

\prose{\textbf{การยมาเนนา}ติ กโรนฺโต วา การาเปนฺโต วา.}

\prose{\textbf{อสฺสามิกนฺ}ติ น อญฺโญ โกจิ สามิโก โหติ อิตฺถี วา ปุริโส วา คหฏฺโฐ วา ปพฺพชิโต วา.}

\prose{\textbf{อตฺตุทฺเทสนฺ}ติ อตฺตโน อตฺถาย.}

\prose{\textbf{ปมาณิกา กาเรตพฺพา, ตตฺริทํ ปมาณํ, ทีฆโส ทฺวาทส วิทตฺถิโย สุคต}\-\textbf{วิทตฺถิยา}ติ พาหิริเมน มาเนน.}

\prose{\textbf{ติริยํ สตฺตนฺตรา}ติ อพฺภนฺตริเมน มาเนน.}

\prose{\textbf{ภิกฺขู อภิเนตพฺพา วตฺถุเทสนายา}ติ เตน กุฏิการเกน ภิกฺขุนา กุฏิวตฺถุํ โสเธตฺวา สํฆํ อุปสงฺกมิตฺวา เอกํสํ อุตฺตราสงฺคํ กริตฺวา วุฑฺฒานํ ภิกฺขูนํ ปาเท วนฺทิตฺวา อุกฺกุฏิกํ นิสีทิตฺวา อญฺชลิํ ปคฺคเหตฺวา เอวมสฺส วจนีโย “อหํ ภนฺเต สญฺญาจิกาย กุฏิํ กตฺตุกาโม อสฺสามิกํ อตฺตุทฺเทสํ, โสหํ ภนฺเต สํฆํ กุฏิวตฺถุ\-โอโลกนํ ยาจามี”ติ. ทุติยมฺปิ ยาจิตพฺพา. ตติยมฺปิ ยาจิตพฺพา. สเจ สพฺโพ สํโฆ อุสฺสหติ กุฏิวตฺถุํ โอโลเกตุํ, สพฺเพน สํเฆน โอโลเกตพฺพํ, โน เจ สพฺโพ สํโฆ อุสฺสหติ กุฏิวตฺถุํ โอโลเกตุํ, เย ตตฺถ โหนฺติ ภิกฺขู พฺยตฺตา ปฏิพลา สารมฺภํ อนารมฺภํ สปริกฺกมนํ อปริกฺกมนํ ชานิตุํ, เต ยาจิตฺวา สมฺมนฺนิตพฺพา, เอวญฺจ ปน ภิกฺขเว สมฺมนฺนิตพฺพา. พฺยตฺเตน ภิกฺขุนา ปฏิพเลน สํโฆ ญาเปตพฺโพ\csromandash{}}

\proseitem{350}{\csromanfolio{225}“สุณาตุ เม ภนฺเต สํโฆ, อยํ อิตฺตนฺนาโม ภิกฺขุ สญฺญาจิกาย กุฏิํ กตฺตุกาโม อสฺสามิกํ อตฺตุทฺเทสํ, โส สํฆํ กุฏิวตฺถุ\-โอโลกนํ ยาจติ, ยทิ สํฆสฺส ปตฺตกลฺลํ, สํโฆ อิตฺถนฺนามญฺจ อิตฺถนฺนามญฺจ ภิกฺขู สมฺมนฺเนยฺย อิตฺถนฺนามสฺส ภิกฺขุโน กุฏิวตฺถุํ โอโลเกตุํ, เอสา ญตฺติ.}

\prose{สุณาตุ เม ภนฺเต สํโฆ, อยํ อิตฺถนฺนาโม ภิกฺขุ สญฺญาจิกาย กุฏิํ กตฺตุกาโม อสฺสามิกํ อตฺตุทฺเทสํ, โส สํฆํ กุฏิวตฺถุ\-โอโลกนํ ยาจติ, สํโฆ อิตฺถนฺนามญฺจ อิตฺถนฺนามญฺจ ภิกฺขู สมฺมนฺนติ อิตฺถนฺนามสฺส ภิกฺขุโน กุฏิวตฺถุํ โอโลเกตุํ, ยสฺสายสฺมโต ขมติ อิตฺถนฺนามสฺส จ อิตฺถนฺนามสฺส จ ภิกฺขูนํ สมฺมุติ\csromansharedfootnote{5c6f5c0884ad928c}{สมฺมติ (สฺยา)} อิตฺถนฺนามสฺส ภิกฺขุโน กุฏิวตฺถุํ โอโลเกตุํ, โส ตุณฺหสฺส, ยสฺส นกฺขมติ, โส ภาเสยฺย.}

\prose{สมฺมตา สํเฆน อิตฺถนฺนาโม จ อิตฺถนฺนาโม จ ภิกฺขู อิตฺถนฺนามสฺส ภิกฺขุโน กุฏิวตฺถุํ โอโลเกตุํ, ขมติ สํฆสฺส, ตสฺมา ตุณฺหี, เอวเมตํ ธารยามี”ติ.}

\proseitem{351}{เตหิ สมฺมเตหิ ภิกฺขูหิ ตตฺถ คนฺตฺวา กุฏิวตฺถุ โอโลเกตพฺพํ, สารมฺภํ อนารมฺภํ สปริกฺกมนํ อปริกฺกมนํ ชานิตพฺพํ. สเจ สารมฺภํ โหติ อปริกฺกมนํ, “มา อิธ กรี”ติ วตฺตพฺโพ. สเจ อนารมฺภํ โหติ สปริกฺกมนํ, สํฆสฺส อาโรเจตพฺพํ “อนารมฺภํ สปริกฺกมนนฺ”ติ. เตน กุฏิการเกน ภิกฺขุนา สํฆํ อุปสงฺกมิตฺวา เอกํสํ อุตฺตราสงฺคํ กริตฺวา วุฑฺฒานํ ภิกฺขูนํ ปาเท วนฺทิตฺวา อุกฺกุฏิกํ นิสีทิตฺวา อญฺชลิํ ปคฺคเหตฺวา เอวมสฺส วจนีโย “อหํ ภนฺเต สญฺญาจิกาย กุฏิํ กตฺตุกาโม อสฺสามิกํ อตฺตุทฺเทสํ, โสหํ ภนฺเต สํฆํ กุฏิวตฺถุ\-เทสนํ ยาจามี”ติ. ทุติยมฺปิ ยาจิตพฺพา. ตติยมฺปิ ยาจิตพฺพา. พฺยตฺเตน ภิกฺขุนา ปฏิพเลน สํโฆ ญาเปตพฺโพ\csromandash{}}

\proseitem{352}{“สุณาตุ เม ภนฺเต สํโฆ, อยํ อิตฺถนฺนาโม ภิกฺขุ สญฺญาจิกาย กุฏิํ กตฺตุกาโม อสฺสามิกํ อตฺตุทฺเทสํ, โส สํฆํ กุฏิวตฺถุ\-เทสนํ ยาจติ, ยทิ สํฆสฺส ปตฺตกลฺลํ, สํโฆ อิตฺถนฺนามสฺส ภิกฺขุโน กุฏิวตฺถุํ เทเสยฺย, เอสา ญตฺติ.}

\prose{\csromanfolio{226}สุณาตุ เม ภนฺเต สํโฆ, อยํ อิตฺถนฺนาโม ภิกฺขุ สญฺญาจิกาย กุฏิํ กตฺตุกาโม อสฺสามิกํ อตฺตุทฺเทสํ, โส สํฆํ กุฏิวตฺถุ\-เทสนํ ยาจติ, สํโฆ อิตฺถนฺนามสฺส ภิกฺขุโน กุฏิวตฺถุํ เทเสติ, ยสฺสายสฺมโต ขมติ อิตฺถนฺนามสฺส ภิกฺขุโน กุฏิวตฺถุสฺส เทสนา, โส ตุณฺหสฺส, ยสฺส นกฺขมติ, โส ภาเสยฺย.}

\prose{เทสิตํ สํเฆน อิตฺถนฺนามสฺส ภิกฺขุโน กุฏิวตฺถุ, ขมติ สํฆสฺส, ตสฺมา ตุณฺหี, เอวเมตํ ธารยามี”ติ.}

\proseitem{353}{\textbf{สารมฺภํ} นาม กิปิลฺลิกานํ วา อาสโย โหติ, อุปจิกานํ วา อาสโย โหติ, อุนฺทูรานํ วา อาสโย โหติ, อหีนํ วา อาสโย โหติ, วิจฺฉิกานํ วา อาสโย โหติ, สตปทีนํ วา อาสโย โหติ, หตฺถีนํ วา อาสโย โหติ, อสฺสานํ วา อาสโย โหติ, สีหานํ วา อาสโย โหติ, พฺยคฺฆานํ วา อาสโย โหติ, ทีปีนํ วา อาสโย โหติ, อจฺฉานํ วา อาสโย โหติ, ตรจฺฉานํ วา อาสโย โหติ, เยสํ เกสญฺจิ ติรจฺฉาน\-คตานํ ปาณานํ อาสโย โหติ, ปุพฺพณฺณ\-นิสฺสิตํ วา โหติ, อปรณฺณ\-นิสฺสิตํ วา โหติ, อพฺภาฆาต\-นิสฺสิตํ วา โหติ, อาฆาตน\-นิสฺสิตํ วา โหติ, สุสานนิสฺสิตํ วา โหติ, อุยฺยานนิสฺสิตํ วา โหติ, ราชวตฺถุ\-นิสฺสิตํ วา โหติ, หตฺถิสาลา\-นิสฺสิตํ วา โหติ, อสฺสสาลา\-นิสฺสิตํ วา โหติ, พนฺธนาคาร\-นิสฺสิตํ วา โหติ, ปานาคาร\-นิสฺสิตํ วา โหติ, สูนนิสฺสิตํ วา โหติ, รจฺฉานิสฺสิตํ วา โหติ, จจฺจรนิสฺสิต`ํ วา โหติ, สภานิสฺสิตํ วา โหติ, สํสรณ\-นิสฺสิตํ วา\csromansharedfootnote{f881278652a770a7}{สญฺจรณนิสฺสิตํ วา (ก)} โหติ, เอตํ สารมฺภํ นาม.}

\prose{\textbf{อปริกฺกมนํ} นาม น สกฺกา โหติ ยถายุตฺเตน สกเฏน อนุปริคนฺตุํ, สมนฺตา นิสฺเสณิยา อนุปริคนฺตุํ, เอตํ อปริกฺกมนํ นาม.}

\prose{\textbf{อนารมฺภํ} นาม น กิปิลฺลิกานํ วา อาสโย โหติ, น อุปจิกานํ วา อาสโย โหติ, น อุนฺทูรานํ วา อาสโย โหติ, น อหีนํ วา อาสโย โหติ, น วิจฺฉิกานํ วา อาสโย โหติ, น สตปทีนํ วา อาสโย โหติ ฯเปฯ น สํสรณ\-นิสฺสิตํ วา โหติ, เอตํ อนารมฺภํ นาม.}

\prose{\csromanfolio{227}\textbf{สปริกฺกมนํ} นาม สกฺกา โหติ ยถายุตฺเตน สกเฏน อนุปริคนฺตุํ, สมนฺตา นิสฺเสณิยา อนุปริคนฺตุํ, เอตํ สปริกฺกมนํ นาม.}

\prose{\textbf{สญฺญาจิกา} นาม สยํ ยาจิตฺวา ปุริสํปิ ปุริส\-ตฺถกรํปิ ฯเปฯ มตฺติกํปิ.}

\prose{\textbf{กุฏิ} นาม อุลฺลิตฺตา วา โหติ อวลิตฺตา วา อุลฺลิตฺตาวลิตฺตา วา.}

\prose{\textbf{กาเรยฺยา}ติ กโรติ วา การาเปติ วา.}

\prose{\textbf{ภิกฺขู วา อนภิเนยฺย วตฺถุเทสนาย ปมาณํ วา อติกฺกาเมยฺยา}ติ ญตฺติทุติเยน กมฺเมน กุฏิวตฺถุํ น เทสาเปตฺวา อายามโต วา วิตฺถารโต วา อนฺตมโส เกสคฺค\-มตฺตํปิ อติกฺกาเมตฺวา กโรติ วา การาเปติ วา, ปโยเค ทุกฺกฏํ. เอกํ ปิณฺฑํ อนาคเต อาปตฺติ ถุลฺลจฺจยสฺส. ตสฺมิํ ปิณฺเฑ อาคเต อาปตฺติ สํฆาทิเสสสฺส.}

\prose{\textbf{สํฆาทิเสโส}ติ ฯเปฯ เตนปิ วุจฺจติ “สํฆาทิเสโส”ติ.}

\proseitem{354}{ภิกฺขุ กุฏิํ กโรติ อเทสิต\-วตฺถุกํ สารมฺภํ อปริกฺกมนํ, อาปตฺติ สํฆาทิเสเสน ทฺวินฺนํ ทุกฺกฏานํ. ภิกฺขุ กุฏิํ กโรติ อเทสิต\-วตฺถุกํ สารมฺภํ สปริกฺกมนํ, อาปตฺติ สํฆาทิเสเสน ทุกฺกฏสฺส. ภิกฺขุ กุฏิํ กโรติ อเทสิต\-วตฺถุกํ อนารมฺภํ อปริกฺกมนํ, อาปตฺติ สํฆาทิเสเสน ทุกฺกฏสฺส. ภิกฺขุ กุฏิํ กโรติ อเทสิต\-วตฺถุกํ อนารมฺภํ สปริกฺกมนํ, อาปตฺติ สํฆาทิเสสสฺส.}

\proseitem{355}{ภิกฺขุ กุฏิํ กโรติ เทสิต\-วตฺถุกํ สารมฺภํ อปริกฺกมนํ, อาปตฺติ ทฺวินฺนํ ทุกฺกฏานํ. ภิกฺขุ กุฏิํ กโรติ เทสิต\-วตฺถุกํ สารมฺภํ สปริกฺกมนํ, อาปตฺติ ทุกฺกฏสฺส. ภิกฺขุ กุฏิํ กโรติ เทสิต\-วตฺถุกํ อนารมฺภํ อปริกฺกมนํ, อาปตฺติ ทุกฺกฏสฺส. ภิกฺขุ กุฏิํ กโรติ เทสิต\-วตฺถุกํ อนารมฺภํ สปริกฺกมนํ, อนาปตฺติ.}

\prose{ภิกฺขุ กุฏิํ กโรติ ปมาณา\-ติกฺกนฺตํ สารมฺภํ อปริกฺกมนํ, อาปตฺติ สํฆาทิเสเสน ทฺวินฺนํ ทุกฺกฏานํ. ภิกฺขุ กุฏิํ กโรติ ปมาณา\-ติกฺกนฺตํ สารมฺภํ สปริกฺกมนํ, อาปตฺติ สํฆาทิเสเสน ทุกฺกฏสฺส. ภิกฺขุ กุฏิํ กโรติ ปมาณา\-ติกฺกนฺตํ อนารมฺภํ อปริกฺกมนํ, อาปตฺติ สํฆาทิเสเสน ทุกฺกฏสฺส. ภิกฺขุ กุฏิํ กโรติ ปมาณา\-ติกฺกนฺตํ อนารมฺภํ สปริกฺกมนํ, อาปตฺติ สํฆาทิเสสสฺส.}

\prose{\csromanfolio{228}ภิกฺขุ กุฏิํ กโรติ ปมาณิกํ สารมฺภํ อปริกฺกมนํ, อาปตฺติ ทฺวินฺนํ ทุกฺกฏานํ. ภิกฺขุ กุฏิํ กโรติ ปมาณิกํ สารมฺภํ สปริกฺกมนํ, อาปตฺติ ทุกฺกฏสฺส. ภิกฺขุ กุฏิํ กโรติ ปมาณิกํ อนารมฺภํ อปริกฺกมนํ, อาปตฺติ ทุกฺกฏสฺส. ภิกฺขุ กุฏิํ กโรติ ปมาณิกํ อนารมฺภํ สปริกฺกมนํ, อนาปตฺติ.}

\prose{ภิกฺขุ กุฏิํ กโรติ อเทสิต\-วตฺถุกํ ปมาณา\-ติกฺกนฺตํ สารมฺภํ อปริกฺกมนํ, อาปตฺติ ทฺวินฺนํ สํฆาทิเสเสน ทฺวินฺนํ ทุกฺกฏานํ. ภิกฺขุ กุฏิํ กโรติ อเทสิต\-วตฺถุกํ ปมาณา\-ติกฺกนฺตํ สารมฺภํ สปริกฺกมนํ, อาปตฺติ ทฺวินฺนํ สํฆาทิเสเสน ทุกฺกฏสฺส. ภิกฺขุ กุฏิํ กโรติ อเทสิต\-วตฺถุกํ ปมาณา\-ติกฺกนฺตํ อนารมฺภํ อปริกฺกมนํ, อาปตฺติ ทฺวินฺนํ สํฆาเทเสเสน ทุกฺกฏสฺส. ภิกฺขุ กุฏิํ กโรติ อเทสิต\-วตฺถุกํ ปมาณา\-ติกฺกนฺตํ อนารมฺภํ สปริกฺกมนํ, อาปตฺติ ทฺวินฺนํ สํฆา\-ทิ\-เสสานํ.}

\prose{ภิกฺขุ กุฏิํ กโรติ เทสิต\-วตฺถุกํ ปมาณิกํ สารมฺภํ อปริกฺกมนํ, อาปตฺติ ทฺวินฺนํ ทุกฺกฏานํ. ภิกฺขุ กุฏิํ กโรติ เทสิต\-วตฺถุกํ ปมาณิกํ สารมฺภํ สปริกฺกมนํ, อาปตฺติ ทุกฺกฏสฺส. ภิกฺขุ กุฏิํ กโรติ เทสิต\-วตฺถุกํ ปมาณิกํ อนารมฺภํ อปริกฺกมนํ, อาปตฺติ ทุกฺกฏสฺส. ภิกฺขุ กุฏิํ กโรติ เทสิต\-วตฺถุกํ ปมาณิกํ อนารมฺภํ สปริกฺกมนํ, อนาปตฺติ.}

\proseitem{356}{ภิกฺขุ สมาทิสติ “กุฏิํ เม กโรถา”ติ. ตสฺส กุฏิํ กโรนฺติ อเทสิต\-วตฺถุกํ สารมฺภํ อปริกฺกมนํ, อาปตฺติ สํฆาทิเสเสน ทฺวินฺนํ ทุกฺกฏานํ ฯเปฯ สารมฺภํ สปริกฺกมนํ, อาปตฺติ สํฆาทิเสเสน กุกฺกฏสฺส ฯเปฯ อนารมฺภํ อปริกฺกมนํ, อาปตฺติ สํฆาทิเสเสน ทุกฺกฏสฺส ฯเปฯ อนารมฺภํ สปริกฺกมนํ, อาปตฺติ สํฆาทิเสสสฺส.}

\prose{ภิกฺขุ สมาทิสติ “กุฏิํ เม กโรถา”ติ. ตสฺส กุฏิํ กโรนฺติ เทสิต\-วตฺถุกํ สารมฺภํ อปริกฺกมนํ, อาปตฺติ ทฺวินฺนํ ทุกฺกฏานํ ฯเปฯ สารมฺภํ สปริกฺกมนํ, อาปตฺติ ทุกฺกฏสฺส ฯเปฯ อนารมฺภํ อปริกฺกมนํ, อาปตฺติ ทุกฺกฏสฺส ฯเปฯ อนารมฺภํ สปริกฺกมนํ, อนาปตฺติ.}

\prose{ภิกฺขุ สมาทิสติ “กุฏิํ เม กโรถา”ติ. ตสฺส กุฏิํ กโรนฺติ ปมาณา\-ติกฺกนฺตํ สารมฺภํ อปริกฺกมนํ, อาปตฺติ สํฆาทิเสเสน ทฺวินฺนํ ทุกฺกฏานํ ฯเปฯ สารมฺภํ สปริกฺกมนํ, อาปตฺติ สํฆาทิเสเสน ทุกฺกฏสฺส ฯเปฯ \csromanfolio{229}อนารมฺภํ อปริกฺกมนํ, อาปตฺติ สํฆาทิเสเสน ทุกฺกฏสฺส ฯเปฯ อนารมฺภํ สปริกฺกมนํ, อาปตฺติ สํฆาทิเสสสฺส.}

\prose{ภิกฺขุ สมาทิสติ “กุฏิํ เม กโรถา”ติ. ตสฺส กุฏิํ กโรนฺติ ปมาณิกํ สารมฺภํ อปริกฺกมนํ, อาปตฺติ ทฺวินฺนํ ทุกฺกฏานํ ฯเปฯ สารมฺภํ สปริกฺกมนํ, อาปตฺติ ทุกฺกฏสฺส ฯเปฯ อนารมฺภํ อปริกฺกมนํ, อาปตฺติ ทุกฺกฏสฺส ฯเปฯ อนารมฺภํ สปริกฺกมนํ, อนาปตฺติ.}

\prose{ภิกฺขุ สมาทิสติ “กุฏิํ เม กโรถา”ติ. ตสฺส กุฏิํ กโรนฺติ อเทสิต\-วตฺถุกํ ปมาณา\-ติกฺกนฺตํ สารมฺภํ อปริกฺกมนํ, อาปตฺติ ทฺวินฺนํ สํฆาทิเสเสน ทฺวินฺนํ ทุกฺกฏานํ ฯเปฯ สารมฺภํ สปริกฺกมนํ, อาปตฺติ ทฺวินฺนํ สํฆาทิเสเสน ทุกฺกฏสฺส ฯเปฯ อนารมฺภํ อปริกฺกมนํ, อาปตฺติ ทฺวินฺนํ สํฆาทิเสเสน ทุกฺกฏสฺส ฯเปฯ อนารมฺภํ สปริกฺกมนํ, อาปตฺติ ทฺวินฺนํ สํฆาทิเสเสนํ.}

\prose{ภิกฺขุ สมาทิสติ “กุฏิํ เม กโรถา”ติ. ตสฺส กุฏิํ กโรนฺติ เทสิต\-วตฺถุกํ ปมาณิกํ สารมฺภํ อปริกฺกมนํ, อาปตฺติ ทฺวินฺนํ ทุกฺกฏานํ ฯเปฯ สารมฺภํ สปริกฺกมนํ, อาปตฺติ ทุกฺกฏสฺส ฯเปฯ อนารมฺภํ อปริกฺกมนํ, อาปตฺติ ทุกฺกฏสฺส ฯเปฯ อนารมฺภํ สมริกฺกมนํ, อนาปตฺติ.}

\proseitem{357}{ภิกฺขุ สมาทิสิตฺวา ปกฺกมติ “กุฏิํ เม กโรถา”ติ, น จ สมาทิสติ “เทสิตวตฺถุกา จ โหตุ อนารมฺภา จ สปริกฺกมนา จา”ติ. ตสฺส กุฏิํ กโรนฺติ อเทสิต\-วตฺถุกํ สารมฺภํ อปริกฺกมนํ, อาปตฺติ สํฆาทิเสเสน ทฺวินฺนํ ทุกฺกฏานํ ฯเปฯ สารมฺภํ สปริกฺกมนํ, อาปตฺติ สํฆาทิเสเสน ทุกฺกฏสฺส ฯเปฯ อนารมฺภํ อปริกฺกมนํ, อาปตฺติ สํฆาทิเสเสน ทุกฺกฏสฺส ฯเปฯ อนารมฺภํ สปริกฺกมนํ, อาปตฺติ สํฆาทิเสสสฺส.}

\prose{ภิกฺขุ สมาทิสิตฺวา ปกฺกมติ “กุฏิํ เม กโรถา”ติ, น จ สมาทิสติ “เทสิตวตฺถุกา จ โหตุ อนารมฺภา จ สปริกฺกมนา จา”ติ. ตสฺส กุฏิํ กโรนฺติ เทสิต\-วตฺถุกํ สารมฺภํ อปริกฺกมนํ, อาปตฺติ ทฺวินฺนํ ทุกฺกฏานํ ฯเปฯ สารมฺภํ สปริกฺกมนํ, อาปตฺติ ทุกฺกฏสฺส ฯเปฯ อนารมฺภํ อปริกฺกมนํ, อาปตฺติ ทุกฺกฏสฺส ฯเปฯ อนารมฺภํ สปริกฺกมนํ, อนาปตฺติ.}

\prose{\csromanfolio{230}ภิกฺขุ สมาทิสิตฺวา ปกฺกมติ “กุฏิํ เม กโรถา”ติ, น จ สมาทิสติ “ปมาณิกา จ โหตุ อนารมฺภา จ สปริกฺกมนา จา”ติ. ตสฺส กุฏิํ กโรนฺติ ปมาณา\-ติกฺกนฺตํ สารมฺภํ อปริกฺกมนํ, อาปตฺติ สํฆาทิเสเสน ทฺวินฺนํ ทุกฺกฏานํ ฯเปฯ สารมฺ ภํ สปริกฺกมนํ, อาปตฺติ สํฆาทิเสเสน ทุกฺกฏสฺส ฯเปฯ อนารมฺภํ อปริกฺกมนํ, อาปตฺติ สํฆาทิเสเสน ทุกฺกฏสฺส ฯเปฯ อนารมฺภํ สปริกฺกมนํ, อาปตฺติ สํฆาทิเสสสฺส.}

\prose{ภิกฺขุ สมาทิสิตฺวา ปกฺกมติ “กุฏิํ เม กโรถา”ติ, น จ สมาทิสติ “ปมาณิกา จ โหตุ อนารมฺภา จ สปริกฺกมนา จา”ติ. ตสฺส กุฏิํ กโรนฺติ ปมาณิกํ สารมฺภํ อปริกฺกมนํ, อาปตฺติ ทฺวินฺนํ ทุกฺกฏานํ ฯเปฯ สารมฺภํ สปริกฺกมนํ, อาปตฺติ ทุกฺกฏสฺส ฯเปฯ อนารมฺภํ อปริกฺกมนํ, อาปตฺติ ทุกฺกฏสฺส ฯเปฯ อนารมฺภํ สปริกฺกมนํ, อนาปตฺติ.}

\prose{ภิกฺขุ สมาทิสิตฺวา ปกฺกมติ “กุฏิํ เม กโรถา”ติ, น จ สมาทิสติ “เทสิตวตฺถุกา จ โหตุ ปมาณิกา จ อนารมฺภา จ สปริกฺกมนา จา”ติ. ตสฺส กุฏิํ กโรนฺติ อเทสิต\-วตฺถุกํ ปมาณา\-ติกฺกนฺตํ สารมฺภํ อปริกฺกมนํ, อาปตฺติ ทฺวินฺนํ สํฆาทิเสเสน ทฺวินฺนํ ทุกฺกฏานํ ฯเปฯ สารมฺภํ สปริกฺกมนํ, อาปตฺติ ทฺวินฺนํ สํฆาทิเสเสน ทุกฺกฏสฺส ฯเปฯ อนารมฺภํ อปริกฺกมนํ, อาปตฺติ ทฺวินฺนํ สํฆาทิเสเสน ทุกฺกฏสฺส ฯเปฯ อนารมฺภํ สปริกฺกมนํ, อาปตฺติ ทฺวินฺนํ สํฆา\-ทิ\-เสสานํ.}

\prose{ภิกฺขุ สมาทิสิตฺวา ปกฺกมติ “กุฏิํ เม กโรถา”ติ, น จ สมาทิสติ “เทสิตวตฺถุกา จ โหตุ ปมาณิกา จ อนารมฺภา จ สปริกฺกมนา จา”ติ. ตสฺส กุฏิํ กโรนฺติ เทสิต\-วตฺถุกํ ปมาณิกํ สารมฺภํ อปริกฺกมนํ, อาปตฺติ ทฺวินฺนํ ทุกฺกฏานํ ฯเปฯ สารมฺภํ สปริกฺกมนํ, อาปตฺติ ทุกฺกฏสฺส ฯเปฯ อนารมฺภํ อปริกฺกมนํ, อาปตฺติ ทุกฺกฏสฺส ฯเปฯ อนารมฺภํ สปริกฺกมนํ, อนาปตฺติ.}

\proseitem{358}{ภิกฺขุ สมาทิสิตฺวา ปกฺกมติ “กุฏิํ เม กโรถา”ติ, สมาทิสติ จ “เทสิตวตฺถุกา จ โหตุ อนารมฺภา จ สปริกฺกมนา จา”ติ. ตสฺส กุฏิํ กโรนฺติ อเทสิต\-วตฺถุกํ สารมฺภํ อปริกฺกมนํ. โส สุณาติ “กุฏิ กิร เม กยิรติ อเทสิต\-วตฺถุกา สารมฺภา อปริกฺกมนา”ติ. เตน ภิกฺขุนา สามํ วา คนฺตพฺพํ ทูโต วา ปาเหตพฺโพ “เทสิตวตฺถุกา จ โหตุ อนารมฺภา จ สปริกฺกมนา จา”ติ. โน เจ สามํ วา คจฺเฉยฺย ทูตํ วา ปหิเณยฺย, อาปตฺติ ทุกฺกฏสฺส.}

\prose{\csromanfolio{231}ภิกฺขุ สมาทิสิตฺวา ปกฺกมติ “กุฏิํ เม กโรถา”ติ, สมาทิสติ จ “เทสิตวตฺถุกา จ โหตุ อนารมฺภา จ สปริกฺกมนา จา”ติ. ตสฺส กุฏิํ กโรนฺติ อเทสิต\-วตฺถุกํ สารมฺภํ สปริกฺกมนํ. โส สุณาติ “กุฏิ กิร เม กยิรติ อเทสิต\-วตฺถุกา สารมฺภา สปริกฺกมนา”ติ. เตน ภิกฺขุนา สามํ วา คนฺตพฺพํ ทูโต วา ปาเหตพฺโพ “เทสิตวตฺถุกา จ โหตุ อนารมฺภา จา”ติ. โน เจ สามํ วา คจฺเฉยฺย ทูตํ วา ปหิเณยฺย, อาปตฺติ ทุกฺกฏสฺส.}

\prose{ภิกฺขุ สมาทิสิตฺวา ปกฺกมติ “กุฏิํ เม กโรถา”ติ, สมาทิสติ จ “เทสิตวตฺถุกา จ โหตุ อนารมฺภา จ สปริกฺกมนา จา”ติ. ตสฺส กุฏิํ กโรนฺติ อเทสิต\-วตฺถุกํ อนารมฺภํ อปริกฺกมนํ. โส สุณาติ “กุฏิ กิร เม กยิรติ อเทสิต\-วตฺถุกา อนารมฺภา อปริกฺกมนา”ติ. เตน ภิกฺขุนา สามํ วา คนฺตพฺพํ ทูโต วา ปาเหตพฺโพ “เทสิตวตฺถุกา จ โหตุ สปริกฺกมนา จา”ติ. โน เจ สามํ วา คจฺเฉยฺย ทูตํ วา ปหิเณยฺย, อาปตฺติ ทุกฺกฏสฺส.}

\prose{ภิกฺขุ สมาทิสิตฺวา ปกฺกมติ “กุฏิํ เม กโรถา”ติ, สมาทิสติ จ “เทสิตวตฺถุกา จ โหตุ อนารมฺภา จ สปริกฺกมนา จา”ติ. ตสฺส กุฏิํ กโรนฺติ อเทสิต\-วตฺถุกํ อนารมฺภํ สปริกฺกมนํ. โส สุณาติ “กุฏิ กิร เม กยิรติ อเทสิต\-วตฺถุกา อนารมฺภา สปริกฺกมนา”ติ. เตน ภิกฺขุนา สามํ วา คนฺตพฺพํ ทูโต วา ปาเหตพฺโพ “เทสิตวตฺถุกา โหตู”ติ. โน เจ สามํ วา คจฺเฉยฺย ทูตํ วา ปหิเณยฺย, อาปตฺติ ทุกฺกฏสฺส.}

\prose{ภิกฺขุ สมาทิสิตฺวา ปกฺกมติ “กุฏิํ เม กโรถา”ติ, สมาทิสติ จ “เทสิตวตฺถุกา จ โหตุ อนารมฺภา จ สปริกฺกมนา จา”ติ. ตสฺส กุฏิํ กโรนฺติ เทสิต\-วตฺถุกํ สารมฺภํ อปริกฺกมนํ. โส สุณาติ “กุฏิ กิร เม กยิรติ เทสิตวตฺถุกา สารมฺภา อปริกฺกมนา”ติ. เตน ภิกฺขุนา สามํ วา คนฺตพฺพํ ทูโต วา ปาเหตพฺโพ “อนารมฺภา จ โหตุ สปริกฺกมนา จา”ติ. โน เจ สามํ วา คจฺเฉยฺย ทูตํ วา ปหิเณยฺย, อาปตฺติ ทุกฺกฏสฺส.}

\prose{ภิกฺขุ สมาทิสิตฺวา ปกฺกมติ “กุฏิํ เม กโรถา”ติ, สมาทิสติ จ “เทสิตวตฺถุกา จ โหตุ อนารมฺภา จ สปริกฺกมนา จา”ติ. ตสฺส กุฏิํ กโรนฺติ เทสิต\-วตฺถุกํ สารมฺภํ สปริกฺกมนํ. โส สุณาติ \csromanfolio{232}“กุฏิ กิร เม กยิรติ เทสิตวตฺถุกา สารมฺภา สปริกฺกมนา”ติ. เตน ภิกฺขุนา สามํ วา คนฺตพฺพํ ทูโต วา ปาเหตพฺโพ “อนารมฺภา โหตู”ติ. โน เจ สามํ วา คจฺเฉยฺย ทูตํ วา ปหิเณยฺย, อาปตฺติ ทุกฺกฏสฺส.}

\prose{ภิกฺขุ สมาทิสิตฺวา ปกฺกมติ “กุฏิํ เม กโรถา”ติ, สมาทิสติ จ “เทสิตวตฺถุกา จ โหตุ อนารมฺภา จ สปริกฺกมนา จา”ติ. ตสฺส กุฏิํ กโรนฺติ เทสิต\-วตฺถุกํ อนารมฺภํ อปริกฺกมนํ. โส สุณาติ “กุฏิ กิร เม กยิรติ เทสิตวตฺถุกา อนารมฺภา อปริกฺกมนาติ”. เตน ภิกฺขุนา สามํ วา คนฺตพฺพํ ทูโต วา ปาเหตพฺโพ “สปริกฺกมนา โหตู”ติ. โน เจ สามํ วา คจฺเฉยฺย ทูตํ วา ปหิเณยฺย, อาปตฺติ ทุกฺกฏสฺส.}

\prose{ภิกฺขุ สมาทิสิตฺวา ปกฺกมติ “กุฏิํ เม กโรถา”ติ, สมาทิสติ จ “เทสิตวตฺถุกา จ โหตุ อนารมฺภา จ สปริกฺกมนา จา”ติ. ตสฺส กุฏิํ กโรนฺติ เทสิต\-วตฺถุกํ อนารมฺภํ สปริกฺกมนํ, อนาปตฺติ.}

\proseitem{359}{ภิกฺขุ สมาทิสิตฺวา ปกฺกมติ “กุฏิํ เม กโรถา”ติ, สมาทิสติ จ “ปมาณิกา จ โหตุ อนารมฺภา จ สปริกฺกมนา จา”ติ. ตสฺส กุฏิํ กโรนฺติ ปมาณา\-ติกฺกนฺตํ สารมฺภํ อปริกฺกมนํ. โส สุณาติ “กุฏิ กิร เม กยิรติ ปมาณาติกฺกนฺตา สารมฺภา อปริกฺกมนา”ติ. เตน ภิกฺขุนา สามํ วา คนฺตพฺพํ ทูโต วา ปาเหตพฺโพ “ปมาณิกา จ โหตุ อนารมฺภา จ สปริกฺกมนา จา”ติ ฯเปฯ “ปมาณิกา จ โหตุ อนารมฺภา จา”ติ ฯเปฯ “ปมาณิกา จ โหตุ สปริกฺกมนา จา”ติ ฯเปฯ “ปมาณิกา โหตู”ติ. โน เจ สามํ วา คจฺเฉยฺย ทูตํ วา ปหิเณยฺย, อาปตฺติ ทุกฺกฏสฺส.}

\prose{ภิกฺขุ สมาทิสิตฺวา ปกฺกมติ “กุฏิํ เม กโรถา”ติ, สมาทิสติ จ “ปมาณิกา จ โหตุ อนารมฺภา จ สปริกฺกมนา จา”ติ. ตสฺส กุฏิํ กโรนฺติ ปมาณิกํ สารมฺภํ อปริกฺกมนํ. โส สุณาติ “กุฏิ กิร เม กยิรติ ปมาณิกา สารมฺภา อปริกฺกมนา”ติ. เตน ภิกฺขุนา สามํ วา คนฺตพฺพํ ทูโต วา ปาเหตพฺโพ “อนารมฺภา จ โหตุ สปริกฺกมนา จา”ติ ฯเปฯ “อนารมฺภา โหตู”ติ ฯเปฯ “สปริกฺกมนา โหตู”ติ ฯเปฯ. อนาปตฺติ.}

\proseitem{360}{\csromanfolio{233}ภิกฺขุ สมาทิสิตฺวา ปกฺกมติ “กุฏิํ เม กโรถา”ติ, สมาทิสติ จ “เทสิตวตฺถุกา จ โหตุ ปมาณิกา จ อนารมฺภา จ สปริกฺกมนา จา”ติ. ตสฺส กุฏิํ กโรนฺติ อเทสิต\-วตฺถุกํ ปมาณา\-ติกฺกนฺตํ สารมฺภํ อปริกฺกมนํ. โส สุณาติ “กุฏิ กิร เม กยิรติ อเทสิต\-วตฺถุกา ปมาณาติกฺกนฺตา สารมฺภา อปริกฺกมนา”ติ. เตน ภิกฺขุนา สามํ วา คนฺตพฺพํ ทูโต วา ปาเหตพฺโพ “เทสิตวตฺถุกา จ โหตุ ปมาณิกา จ อนารมฺภา จ สปริกฺกมนา จา”ติ ฯเปฯ “เทสิตวตฺถุกา จ โหตุ ปมาณิกา จ อนารมฺภา จา”ติ ฯเปฯ “เทสิตวตฺถุกา จ โหตุ ปมาณิกา จ สปริกฺกมนา จา”ติ ฯเปฯ “เทสิตวตฺถุกา จ โหตุ ปมาณิกา จา”ติ. โน เจ สามํ วา คจฺเฉยฺย ทูตํ วา ปหิเณยฺย, อาปตฺติ ทุกฺกฏสฺส.}

\prose{ภิกฺขุ สมาทิสิตฺวา ปกฺกมติ “กุฏิํ เม กโรถา”ติ, สมาทิสติ จ “เทสิตวตฺถุกา จ โหตุ ปมาณิกา จ อนารมฺภา จ สปริกฺกมนา จา”ติ. ตสฺส กุฏิํ กโรนฺติ เทสิต\-วตฺถุกํ ปมาณิกํ สารมฺภํ อปริกฺกมนํ. โส สุณาติ “กุฏิ กิร เม กยิรติ เทสิตวตฺถุกา ปมาณิกา สารมฺภา อปริกฺกมนา”ติ. เตน ภิกฺขุนา สามํ วา คนฺตพฺพํ ทูโต วา ปาเหตพฺโพ “อนารมฺภา จ โหตุ สปริกฺกมนา จา”ติ ฯเปฯ “อนารมฺภา โหตู”ติ ฯเปฯ “สปริกฺกมนา โหตู”ติ ฯเปฯ. อนาปตฺติ.}

\prose{ภิกฺขุ สมาทิสิตฺวา ปกฺกมติ “กุฏิํ เม กโรถา”ติ, สมาทิสติ จ “เทสิตวตฺถุกา จ โหตุ อนารมฺภา จ สปริกฺกมนา จา”ติ. ตสฺส กุฏิํ กโรนฺติ อเทสิต\-วตฺถุกํ สารมฺภํ อปริกฺกมนํ, อาปตฺติ การุกานํ ติณฺณํ ทุกฺกฏานํ ฯเปฯ สารมฺภํ สปริกฺกมนํ, อาปตฺติ การุกานํ ทฺวินฺนํ ทุกฺกฏานํ ฯเปฯ อนารมฺภํ อปริกฺกมนํ, อาปตฺติ การุกานํ ทฺวินฺนํ ทุกฺกฏานํ ฯเปฯ อนารมฺภํ สปริกฺกมนํ, อาปตฺติ การุกานํ ทุกฺกฏสฺส.}

\prose{ภิกฺขุ สมาทิสิตฺวา ปกฺกมติ “กุฏิํ เม กโรถา”ติ, สมาทิสติ จ “เทสิตวตฺถุกา จ โหตุ อนารมฺภา จ สปริกฺกมนา จา”ติ. ตสฺส กุฏิํ กโรนฺติ เทสิต\-วตฺถุกํ สารมฺภํ อปริกฺกมนํ, อาปตฺติ การุกานํ ทฺวินฺนํ ทุกฺกฏานํ ฯเปฯ สารมฺภํ สปริกฺกมนํ, อาปตฺติ การุกานํ ทุกฺกฏสฺส ฯเปฯ อนารมฺภํ อปริกฺกมนํ, อาปตฺติ การุกานํ ทุกฺกฏสฺส ฯเปฯ อนารมฺภํ สปริกฺกมนํ, อนาปตฺติ.}

\prose{\csromanfolio{234}ภิกฺขุ สมาทิสิตฺวา ปกฺกมติ “กุฏิํ เม กโรถา”ติ, สมาทิสติ จ “ปมาณิกา จ โหตุ อนารมฺภา จ สปริกฺกมนา จา”ติ. ตสฺส กุฏิํ กโรนฺติ ปมาณา\-ติกฺกนฺตํ สารมฺภํ อปริกฺกมนํ, อาปตฺติ การุกานํ ติณฺณํ ทุกฺกฏานํ ฯเปฯ สารมฺภํ สปริกฺกมนํ, อาปตฺติ การุกานํ ทฺวินฺนํ ทุกฺกฏานํ ฯเปฯ อนารมฺภํ อปริกฺกมนํ, อาปตฺติ การุกานํ ทฺวินฺนํ ทุกฺกฏานํ ฯเปฯ อนารมฺภํ สปริกฺกมนํ, อาปตฺติ การุกานํ ทุกฺกฏสฺส.}

\prose{ภิกฺขุ สมาทิสิตฺวา ปกฺกมติ “กุฏิํ เม กโรถา”ติ, สมาทิสติ จ “ปมาณิกา จ โหตุ อนรมฺภา จ สปริกฺกมนา จา”ติ. ตสฺส กุฏิํ กโรนฺติ ปมาณิกํ สารมฺภํ อปริกฺกมนํ, อาปตฺติ การุกานํ ทฺวินฺนํ ทุกฺกฏานํ ฯเปฯ สารมฺภํ สปริกฺกมนํ, อาปตฺติ การุกานํ ทุกฺกฏสฺส ฯเปฯ อนารมฺภํ อปริกฺกมนํ, อาปตฺติ การุกานํ ทุกฺกฏสฺส ฯเปฯ อนารมฺภํ สปริกฺกมนํ, อนาปตฺติ.}

\prose{ภิกฺขุ สมาทิสิตฺวา ปกฺกมติ “กุฏิํ เม กโรถา”ติ, สมาทิสติ จ “เทสิตวตฺถุกา จ โหตุ ปมาณิกา จ อนารมฺภา จ สปริกฺกมนา จา”ติ. ตสฺส กุฏิํ กโรนฺติ อเทสิต\-วตฺถุกํ ปมาณา\-ติกฺกนฺตํ สารมฺภํ อปริกฺกมนํ, อาปตฺติ การุกานํ จตุนฺนํ ทุกฺกฏานํ ฯเปฯ สารมฺภํ สปริกฺกมนํ, อาปตฺติ การุกานํ ติณฺณํ ทุกฺกฏานํ ฯเปฯ อนารมฺภํ อปริกฺกมนํ, อาปตฺติ การุกานํ ติณฺณํ ทุกฺกฏานํ ฯเปฯ อนารมฺภํ สปริกฺกมนํ, อาปตฺติ การุกานํ ทฺวินฺนํ ทุกฺกฏานํ.}

\prose{ภิกฺขุ สมาทิสิตฺวา ปกฺกมติ “กุฏิํ เม กโรถา”ติ, สมาทิสติ จ “เทสิตวตฺถุกา จ โหตุ ปมาณิกา จ อนารมฺภา จ สปริกฺกมนา จา”ติ. ตสฺส กุฏิํ กโรนฺติ เทสิต\-วตฺถุกํ ปมาณิกํ สารมฺภํ อปริกฺกมนํ, อาปตฺติ การุกานํ ทฺวินฺนํ ทุกฺกฏานํ ฯเปฯ สารมฺภํ สปริกฺกมนํ, อาปตฺติ การุกานํ ทุกฺกฏสฺส ฯเปฯ อนารมฺภํ อปริกฺกมนํ, อาปตฺติ การุกานํ ทุกฺกฏสฺส ฯเปฯ อนารมฺภํ สปริกฺกมนํ, อนาปตฺติ.}

\proseitem{361}{ภิกฺขุ สมาทิสิตฺวา ปกฺกมติ “กุฏิํ เม กโรถา”ติ. ตสฺส กุฏิํ กโรนฺติ อเทสิต\-วตฺถุกํ สารมฺภํ อปริกฺกมนํ. โส เจ วิปฺปกเต อาคจฺฉติ, เตน ภิกฺขุนา สา กุฏิ อญฺญสฺส วา ทาตพฺพา ภินฺทิตฺวา วา ปุน กาตพฺพา, โน เจ อญฺญสฺส วา ทเทยฺย ภินฺทิตฺวา วา ปุน กาเรยฺย, อาปตฺติ สํฆาทิเสเสน ทฺวินฺนํ ทุกฺกฏานํ.}

\prose{\csromanfolio{235}ภิกฺขุ สมาทิสิตฺวา ปกฺกมติ “กุฏิํ เม กโรถา”ติ. ตสฺส กุฏิํ กโรนฺติ อเทสิต\-วตฺถุกํ สารมฺภํ สปริกฺกมนํ. โส เจ วิปฺปกเต อาคจฺฉติ, เตน ภิกฺขุนา สา กุฏิ อญฺญสฺส วา ทาตพฺพา ภินฺทิตฺวา วา ปุน กาตพฺพา, โน เจ อญฺญสฺส วา ทเทยฺย ภินฺทิตฺวา วา ปุน กาเรยฺย, อาปตฺติ สํฆาทิเสเสน ทุกฺกฏสฺส ฯเปฯ อนารมฺภํ อปริกฺกมนํ ฯเปฯ อาปตฺติ สํฆาทิเสเสน ทุกฺกฏสฺส ฯเปฯ อนารมฺภํ สปริกฺกมนํ ฯเปฯ อาปตฺติ สํฆาทิเสสสฺส.}

\prose{ภิกฺขุ สมาทิสิตฺวา ปกฺกมติ “กุฏิํ เม กโรถา”ติ. ตสฺส กุฏิํ กโรนฺติ เทสิต\-วตฺถุกํ สารมฺภํ อปริกฺกมนํ. โส เจ วิปฺปกเต อาคจฺฉติ, เตน ภิกฺขุนา สา กุฏิ อญฺญสฺส วา ทาตพฺพา ภินฺทิตฺวา วา ปุน กาตพฺพา, โน เจ อญฺญสฺส วา ทเทยฺย ภินฺทิตฺวา วา ปุน กาเรยฺย, อาปตฺติ ทฺวินฺนํ ทุกฺกฏานํ ฯเปฯ สารมฺภํ สปริกฺกมนํ. อาปตฺติ ทุกฺกฏสฺส ฯเปฯ อนารมฺภํ อปริกฺกมนํ. อาปตฺติ ทุกฺกฏสฺส ฯเปฯ อนารมฺภํ สปริกฺกมนํ. อนาปตฺติ.}

\proseitem{362}{ภิกฺขุ สมาทิสิตฺวา ปกฺกมติ “กุฏิํ เม กโรถา”ติ. ตสฺส กุฏิํ กโรนฺติ ปมาณา\-ติกฺกนฺตํ สารมฺภํ อปริกฺกมนํ. โส เจ วิปฺปกเต อาคจฺฉติ, เตน ภิกฺขุนา สา กุฏิ อญฺญสฺส วา ทาตพฺพา ภินฺทิตฺวา วา ปุน กาตพฺพา, โน เจ อญฺญสฺส วา ทเทยฺย ภินฺทิตฺวา วา ปุน กาเรยฺย, อาปตฺติ สํฆาทิเสเสน ทฺวินฺนํ ทุกฺกฏานํ ฯเปฯ สารมฺภํ สปริกฺกมนํ. อาปตฺติ สํฆาทิเสเสน ทุกฺกฏสฺส ฯเปฯ อนารมฺภํ อปริกฺกมนํ. อาปตฺติ สํฆาทิเสเสน ทุกฺกฏสฺส ฯเปฯ อนารมฺภํ สปริกฺกมนํ, อาปตฺติ สํฆาทิเสสสฺส.}

\prose{ภิกฺขุ สมาทิสิตฺวา ปกฺกมติ “กุฏิํ เม กโรถา”ติ. ตสฺส กุฏิํ กโรนฺติ ปมาณิกํ สารมฺภํ อปริกฺกมนํ. โส เจ วิปฺปกเต อาคจฺฉติ, เตน ภิกฺขุนา สา กุฏิ อญฺญสฺส วา ทาตพฺพา ภินฺทิตฺวา วา ปุน กาตพฺพา, โน เจ อญฺญสฺส วา ทเทยฺย ภินฺทิตฺวา วา ปุน กาเรยฺย, อาปตฺติ ทฺวินฺนํ ทุกฺกฏานํ ฯเปฯ สารมฺภํ สปริกฺกมนํ. อาปตฺติ ทุกฺกฏสฺส ฯเปฯ อนารมฺภํ อปริกฺกมนํ. อาปตฺติ ทุกฺกฏสฺส ฯเปฯ อนารมฺภํ สปริกฺกมนํ, อนาปตฺติ.}

\prose{ภิกฺขุ สมาทิสิตฺวา ปกฺกมติ “กุฏิํ เม กโรถา”ติ. ตสฺส กุฏิํ กโรนฺติ อเทสิต\-วตฺถุกํ ปมาณา\-ติกฺกนฺตํ สารมฺภํ อปริกฺกมนํ. โส \csromanfolio{236}เจ วิปฺปกเต อาคจฺฉติ, เตน ภิกฺขุนา สา กุฏิ อญฺญสฺส วา ทาตพฺพา ภินฺทิตฺวา วา ปุน กาตพฺพา, โน เจ อญฺญสฺส วา ทเทยฺย ภินฺทิตฺวา วา ปุน กาเรยฺย, อาปตฺติ ทฺวินฺนํ สํฆาทิเสเสน ทฺวินฺนํ ทุกฺกฏานํ ฯเปฯ สารมฺภํ สปริกฺกมนํ. อาปตฺติ ทฺวินฺนํ สํฆาเทเสเสน ทุกฺกฏสฺส ฯเปฯ อนารมฺภํ อปริกฺกมนํ. อาปตฺติ ทฺวินฺนํ สํฆาทิเสเสน ทุกฺกฏสฺส ฯเปฯ อนารมฺภํ สปริกฺกมนํ, อาปตฺติ ทฺวินฺนํ สํฆา\-ทิ\-เสสานํ.}

\prose{ภิกฺขุ สมาทิสิตฺวา ปกฺกมติ “กุฏิํ เม กโรถา”ติ. ตสฺส กุฏิํ กโรนฺติ เทสิต\-วตฺถุกํ ปมาณิกํ สารมฺภํ อปริกฺกมนํ. โส เจ วิปฺปกเต อาคจฺฉติ, เตน ภิกฺขุนา สา กุฏิ อญฺญสฺส วา ทาตพฺพา ภินฺทิตฺวา วา ปุน กาตพฺพา, โน เจ อญฺญสฺส วา ทเทยฺย ภินฺทิตฺวา วา ปุน กาเรยฺย, อาปตฺติ ทฺวินฺนํ ทุกฺกฏานํ ฯเปฯ สารมฺภํ สปริกฺกมนํ, อาปตฺติ ทุกฺกฏสฺส ฯเปฯ อนารมฺภํ อปริกฺกมนํ, อาปตฺติ ทุกฺกฏสฺส.}

\prose{ภิกฺขุ สมาทิสิตฺวา ปกฺกมติ “กุฏิํ เม กโรถา”ติ. ตสฺส กุฏิํ กโรนฺติ เทสิต\-วตฺถุกํ ปมาณิกํ อนารมฺภํ สปริกฺกมนํ, อนาปตฺติ.}

\proseitem{363}{อตฺตนา วิปฺปกตํ อตฺตนา ปริโยสาเปติ, อาปตฺติ สํฆาทิเสสสฺส.}

\prose{อตฺตนา วิปฺปกตํ ปเรหิ ปริโยสาเปติ\csromansharedfootnote{b33a6a7871dc62e6}{ปริโยสาวาเปติ (ก)}, อาปตฺติ สํฆาทิเสสสฺส.}

\prose{ปเรหิ วิปฺปกตํ อตฺตนา ปริโยสาเปติ, อาปตฺติ สํฆาทิเสสสฺส.}

\prose{ปเรหิ วิปฺปกตํ ปเรหิ ปริโยสาเปติ\csromansharedfootnote{b33a6a7871dc62e6}{ปริโยสาวาเปติ (ก)} อาปตฺติ สํฆาทิเสสสฺส.}

\proseitem{364}{อนาปตฺติ เลเณ คุหาย ติณกุฏิกาย อญฺญสฺสตฺถาย วาสาคารํ ฐเปตฺวา สพฺพตฺถ, อนาปตฺติ อุมฺมตฺตกสฺส อาทิกมฺมิกสฺสาติ.}

\vspace{\tipitakaparskip}
\csromancenter{กุฏิการ\-สิกฺขาปทํ นิฏฺฐิตํ ฉฏฺฐํ.}
\csromanmatikaanchor{mtk.73}
\csromantocmarknopage{section}{7. วิหารการสิกฺขาปท}{mtk.73}
\bookmark[dest={mtk.73},level=1]{7. วิหารการสิกฺขาปท}
\csromanheader{1}{7. \textbf{วิหารการ}\-\textbf{สิกฺขาปท}}
\csromanmatikaanchor{mtk.74}
\csromanmatikaanchor{mtk.75}
\csromanmatikaanchor{mtk.76}
\csromantocmarkref{subsection}{ฉนฺนภิกฺขุวตฺถุ}{mtk.74}
\bookmark[dest={mtk.74},level=2]{ฉนฺนภิกฺขุวตฺถุ}
\csromantocmarkref{subsection}{ปญฺญตฺติ}{mtk.75}
\bookmark[dest={mtk.75},level=2]{ปญฺญตฺติ}
\csromantocmarkref{subsection}{สิกฺขาปทวิภงฺค, ปทภาชนีย}{mtk.76}
\bookmark[dest={mtk.76},level=2]{สิกฺขาปทวิภงฺค, ปทภาชนีย}
\proseitem{365}{\csromanfolio{237}เตน สมเยน พุทฺโธ ภควา โกสมฺพิยํ วิหรติ โฆสิตาราเม. เตน โข ปน สมเยน อายสฺมโต ฉนฺนสฺส อุปฏฺฐาโก คหปติ อายสฺมนฺตํ ฉนฺนํ เอตทโวจ “วิหารวตฺถุํ ภนฺเต ชานาหิ, อยฺยสฺส วิหารํ การาเปสฺสามี”ติ. อถ โข อายสฺมา ฉนฺโน วิหารวตฺถุํ โสเธนฺโต อญฺญตรํ เจติยรุกฺขํ เฉทาเปสิ คามปูชิตํ นิคมปูชิตํ นครปูชิตํ ชนปท\-ปูชิตํ รฏฺฐปูชิตํ. มนุสฺสา อุชฺฌายนฺติ ขิยฺยนฺติ วิปาเจนฺติ “กถํ หิ นาม สมณา สกฺยปุตฺติยา เจติยรุกฺขํ เฉทาเปสฺสนฺติ คามปูชิตํ นิคมปูชิตํ นครปูชิตํ ชนปท\-ปูชิตํ รฏฺฐปูชิตํ, เอกินฺทฺริยํ สมณา สกฺยปุตฺติยา ชีวํ วิเหเฐนฺตี\csromansharedfootnote{b172837ba8dff7d3}{วิเหเฐสฺสนฺติ (กตฺถจิ)} ติ. อสฺโสสุํ โข ภิกฺขู เตสํ มนุสฺสานํ อุชฺฌายนฺตานํ ขิยฺยนฺตานํ วิปาเจนฺตานํ, เย เต ภิกฺขู อปฺปิจฺฉา ฯเปฯ เต อุชฺฌายนฺติ ขิยฺยนฺติ วิปาเจนฺติ “กถํ หิ นาม อายสฺมา ฉนฺโน เจติยรุกฺขํ เฉทาเปสฺสติ คามปูชิตํ ฯเปฯ รฏฺฐปูชิตนฺ”ติ. อถ โข เต ภิกฺขู อายสฺมนฺตํ ฉนฺนํ อเนก\-ปริยาเยน วิครหิตฺวา ภควโต เอตมตฺถํ อาโรเจสุํ ฯเปฯ “สจฺจํ กิร ตฺวํ ฉนฺน เจติยรุกฺขํ เฉทาเปสิ คามปูชิตํ ฯเปฯ รฏฺฐปูชิตนฺ”ติ. สจฺจํ ภควาติ. วิครหิ พุทฺโธ ภควา ฯเปฯ. กถํ หิ นาม ตฺวํ โมฆปุริส เจติยรุกฺขํ เฉทาเปสฺสสิ คามปูชิตํ นิคมปูชิตํ นครปูชิตํ ชนปท\-ปูชิตํ รฏฺฐปูชิตํ, ชีวสญฺญิโน หิ โมฆปุริส มนุสฺสา รุกฺขสฺมิํ, เนตํ โมฆปุริส อปฺปสนฺนานํ วา ปสาทาย ฯเปฯ. เอวญฺจ ปน ภิกฺขเว อิมํ สิกฺขาปทํ อุทฺทิเสยฺยาถ\csromandash{}}

\proseitem{366}{\textbf{“มหลฺลกํ ปน ภิกฺขุนา วิหารํ การยมาเนน สสฺสามิกํ อตฺตุทฺเทสํ ภิกฺขู อภิเนตพฺพา วตฺถุเทสนาย, เตหิ ภิกฺขูหิ วตฺถุ เทเสตพฺพํ อนรมฺภํ สปริกฺกมนํ. สารมฺเภ เจ ภิกฺขุ วตฺถุสฺมิํ อปริกฺกมเน มหลฺลกํ วิหารํ กาเรยฺย ภิกฺขู วา อนภิเนยฺย วตฺถุเทสนาย, สํฆาทิเสโส”}ติ.}

\proseitem{367}{\textbf{มหลฺลโก} นาม วิหาโร สสฺสามิโก วุจฺจติ.}

\prose{\textbf{วิหาโร} นาม อุลฺลิตฺโต วา โหติ อวลิตฺโต วา อุลฺลิตฺตาวลิตฺโต วา.}

\prose{\csromanfolio{238}\textbf{การยมาเนนา}ติ กโรนฺโต วา การาเปนฺโต วา.}

\prose{\textbf{สสฺสามิกนฺ}ติ อญฺโญ โกจิ สามิโก โหติ อิตฺถี วา ปุริโส วา คหฏฺโฐ วา ปพฺพชิโต วา.}

\prose{\textbf{อตฺตุ ทฺเทสนฺ}ติ อตฺตโน อตฺถาย.}

\prose{\textbf{ภิกฺขู อภิเนตพฺพา วตฺถุเทสนายา}ติ เตน วิหารการเกน ภิกฺขุนา วิหารวตฺถุํ โสเธตฺวา สํฆํ อุปสงฺกมิตฺวา เอกํสํ อุตฺตราสงฺคํ กริตฺวา วุฑฺฒานํ ภิกฺขูนํ ปาเท วนฺทิตฺวา อุกฺกุฏิกํ นิสีทิตฺวา อญฺชลิํ ปคฺคเหตฺวา เอวมสฺส วจนีโย “อหํ ภนฺเต มหลฺลกํ วิหารํ กตฺตุกาโม สสฺสามิกํ อตฺตุทฺเทสํ, โสหํ ภนฺเต สํฆํ วิหารวตฺถุ- โอโลกนํ ยาจามี”ติ. ทุติยมฺปิ ยาจิตพฺพา. ตติยมฺปิ ยาจิตพฺพา. สเจ สพฺโพ สํโฆ อุสฺสหติ วิหารวตฺถุํ โอโลเกตุํ, สพฺเพน สํเฆน โอโลเกตพฺพํ. โน เจ สพฺโพ สํโฆ อุสฺสหติ วิหารวตฺถุํ โอโลเกตุํ, เย ตตฺถ โหนฺติ ภิกฺขู พฺยตฺตา ปฏิพลา สารมฺภํ อนารมฺภํ สปริกฺกมนํ อปริกฺกมนํ ชานิตุํ, เต ยาจิตฺวา สมฺมนฺนิตพฺพา. เอวญฺจ ปน ภิกฺขเว สมฺมนฺนิตพฺพา, พฺยตฺเตน ภิกฺขุนา ปฏิพเลน สํโฆ ญาเปตพฺโพ\csromandash{}}

\proseitem{368}{“สุณาตุ เม ภนฺเต สํโฆ, อยํ อิตฺถนฺนาโม ภิกฺขุ มหลฺลกํ วิหารํ กตฺตุกาโม สสฺสามิกํ อตฺตุทฺเทสํ, โส สํฆํ วิหารวตฺถุ- โอโลกนํ ยาจติ, ยทิ สํฆสฺส ปตฺตกลฺลํ, สํโฆ อิตฺถนฺนามญฺจ อิตฺถนฺนามญฺจ ภิกฺขู สมฺมนฺเนยฺย อิตฺถนฺนามสฺส ภิกฺขุโน วิหารวตฺถุํ โอโลเกตุํ, เอสา ญตฺติ.}

\prose{สุณาตุ เม ภนฺเต สํโฆ, อยํ อิตฺถนฺนาโม ภิกฺขุ มหลฺลกํ วิหารํ กตฺตุกาโม สสฺสามิกํ อตฺตุทฺเทสํ, โส สํฆํ วิหารวตฺถุ- โอโลกนํ ยาจติ, สํโฆ อิตฺถนฺนามญฺจ อิตฺถนฺนามญฺจ ภิกฺขู สมฺมนฺนาติ อิตฺถนฺนามสฺส ภิกฺขุโน วิหารวตฺถุํ โอโลเกตุํ, ยสฺสายสฺมโต ขมติ อิตฺถนฺนามสฺส จ อิตฺถนฺนามสฺส จ ภิกฺขูนํ สมฺมุติ อิตฺถนฺนามสฺส ภิกฺขุโน วิหารวตฺถุํ โอโลเกตุํ, โส ตุณฺหสฺส, ยสฺส นกฺขมติ, โส ภาเสยฺย.}

\prose{สมฺมตา สํเฆน อิตฺถนฺนาโม จ อิตฺถนฺนาโม จ ภิกฺขู อิตฺถนฺนามสฺส ภิกฺขุโน วิหารวตฺถุํ โอโลเกตุํ, ขมติ สํฆสฺส, ตสฺมา ตุณฺหี, เอวเมตํ ธารยามี”ติ.}

\proseitem{369}{\csromanfolio{239}เตหิ สมฺมเตหิ ภิกฺขูหิ ตตฺถ คนฺตฺวา วิหารวตฺถุ โอโลเกตพฺพํ, สารมฺภํ อนารมฺภํ สปริกฺกมนํ อปริกฺกมนํ ชานิตพฺพํ. สเจ สารมฺภํ โหติ อปริกฺกมนํ, “มายิธ กรี”ติ วตฺตพฺโพ. สเจ อนารมฺภํ โหติ สปริกฺกมนํ, สํฆสฺส อาโรเจตพฺพํ “อนารมฺภํ สปริกฺกมนนฺ”ติ. เตน วิหารการเกน ภิกฺขุนา สํฆํ อุปสงฺกมิตฺวา เอกํสํ อุตฺตราสงฺคํ กริตฺวา วุฑฺฒานํ ภิกฺขูนํ ปาเท วนฺทิตฺวา อุกฺกุฏิกํ นิสีทิตฺวา อญฺชลิํ ปคฺคเหตฺวา เอวมสฺส วจนีโย “อหํ ภนฺเต มหลฺลกํ วิหารํ กตฺตุกาโม สสฺสามิกํ อตฺตุทฺเทสํ, โสหํ ภนฺเต สํฆํ วิหารวตฺถุ\-เทสนํ ยาจามี”ติ. ทุติยมฺปิ ยาจิตพฺพา. ตติยมฺปิ ยาจิตพฺพา. พฺยตฺเตน ภิกฺขุนา ปฏิพเลน สํโฆ ญาเปตพฺโพ\csromandash{}}

\proseitem{370}{“สุณาตุ เม ภนฺเต สํโฆ, อยํ อิตฺถนฺนาโม ภิกฺขุ มหลฺลกํ วิหารํ กตฺตุกาโม สสฺสามิกํ อตฺตุทฺเทสํ, โส สํฆํ วิหารวตฺถุ เทสนํ ยาจติ, ยทิ สํฆสฺส ปตฺตกลฺลํ, สํโฆ อิตฺถนฺนามสฺส ภิกฺขุโน วิหารวตฺถุํ เทเสยฺย, เอสา ญตฺติ.}

\prose{สุณาตุ เม ภนฺเต สํโฆ, อยํ อิตฺถนฺนาโม ภิกฺขุ มหลฺลกํ วิหารํ กตฺตุกาโม สสฺสามิกํ อตฺตุทฺเทสํ, โส สํฆํ วิหารวตฺถุ\-เทสนํ ยาจติ, สํโฆ อิตฺถนฺนามสฺส ภิกฺขุโน วิหารวตฺถุํ เทเสติ, ยสฺสายสฺมโต ขมติ อิตฺถนฺนามสฺส ภิกฺขุโน วิหาร\-วตฺถุสฺส เทสนา, โส ตุณฺหาสฺส, ยสฺส นกฺขมติ, โส ภาเสยฺย.}

\prose{เทสิตํ สํเฆน อิตฺถนฺนามสฺส ภิกฺขุโน วิหารวตฺถุ, ขมติ สํฆสฺส, ตสฺมา ตุณฺหี, เอวเมตํ ธารยามี”ติ.}

\proseitem{371}{\textbf{สารมฺภํ} นาม กิปิลฺลิกานํ วา อาสโย โหติ, อุปจิกานํ วา อาสโย โหติ, อุนฺทูรานํ วา ฯเปฯ อหีนํ วา. วิจฺฉิกานํ วา. สตปทีนํ วา. หตฺถีนํ วา. อสฺสานํ วา. สีหานํ วา. พฺยคฺฆานํ วา. ทีปีนํ วา. อจฺฉานํ วา. ตรจฺฉานํ วา อาสโย โหติ, เยสํ เกสญฺจิ ติรจฺฉาน\-คตานํ ปาณานํ อาสโย โหติ, ปุพฺพณฺณ\-นิสฺสิตํ วา โหติ, อปรณฺณ\-นิสฺสิตํ วา โหติ, อพฺภาฆาต\-นิสฺสิตํ วา โหติ, อาฆาตน\-นิสฺสิตํ วา โหติ, สุสานนิสฺสิตํ วา โหติ, อุยฺยานนิสฺสิตํ วา โหติ, ราชวตฺถุ\-นิสฺสิตํ วา โหติ, หตฺถิสาลา\-นิสฺสิตํ วา โหติ, อสฺสสาลา\-นิสฺสิตํ วา โหติ, พนฺธนาคาร\-นิสฺสิตํ วา โหติ, \csromanfolio{240}ปานาคาร\-นิสฺสิตํ วา โหติ, สูนนิสฺสิตํ วา โหติ, รจฺฉานิสฺสิตํ วา โหติ, จจฺจร\-นิสฺสิตํ วา โหติ, สภานิสฺสิตํ วา โหติ, สํสรณ\-นิสฺสิตํ วา โหติ, เอตํ สารมฺภํ นาม.}

\prose{\textbf{อปริกฺกมนํ} นาม น สกฺกา โหติ ยถายุตฺเตน สกเฏน อนุปริคนฺตุํ, สมนฺตา นิสฺเสณิยา อนุปริคนฺตุํ, เอตํ อปริกฺกมนํ นาม.}

\prose{\textbf{อนารมฺภํ} นาม น กิปิลฺลิกานํ วา อาสโย โหติ ฯเปฯ น สํสรณ\-นิสฺสิตํ วา โหติ, เอตํ อนารมฺภํ นาม.}

\prose{\textbf{สปริกฺกมนํ} นาม สกฺกา โหติ ยถายุตฺเตน สกเฏน อนุปริคนฺตุํ, สมนฺตา นิสฺเสณิยา อนุปริคนฺตุํ, เอตํ สปริกฺกมนํ นาม.}

\prose{\textbf{มหลฺลโก} นาม วิหาโร สสฺสามิโก วุจฺจติ.}

\prose{\textbf{วิหาโร} นาม อุลฺลิตฺโต วา โหติ อวลิตฺโต วา อุลฺลิตฺตาวลิตฺโต วา.}

\prose{\textbf{กาเรยฺยา}ติ กโรติ วา การาเปติ วา.}

\prose{\textbf{ภิกฺขู วา อนภิเนยฺย วตฺถุเทสนายา}ติ ญตฺติทุติเยน กมฺเมน วิหารวตฺถุํ น เทสาเปตฺวา กโรติ วา การาเปติ วา, ปโยเค ทุกฺกฏํ. เอกํ ปิณฺฑํ อนาคเต อาปตฺติ ถุลฺลจฺจยสฺส. ตสฺมิํ ปิณฺเฑ อาคเต อาปตฺติ สํฆาเทเสสสฺส.}

\prose{\textbf{สํฆาทิเสโส}ติ ฯเปฯ เตนปิ วุจฺจติ “สํฆาทิเสโส”ติ.}

\proseitem{372}{ภิกฺขุ วิหารํ กโรติ อเทสิต\-วตฺถุกํ สารมฺภํ อปริกฺกมนํ, อาปตฺติ สํฆาทิเสเสน ทฺวินฺนํ ทุกฺกฏานํ. ภิกฺขุ วิหารํ กโรติ อเทสิต\-วตฺถุกํ สารมฺภํ สปริกฺกมนํ, อาปตฺติ สํฆาทิเสเสน ทุกฺกฏสฺส. ภิกฺขุ วิหารํ กโรติ อเทสิต\-วตฺถุกํ อนรมฺภํ อปริกฺกมนํ, อาปตฺติ สํฆาทิเสเสน ทุกฺกฏสฺส. ภิกฺขุ วิหารํ กโรติ อเทสิต\-วตฺถุกํ อนารมฺภํ สปริกฺกมนํ, อาปตฺติ สํฆาทิเสสสฺส.}

\prose{ภิกฺขุ วิหารํ กโรติ เทสิต\-วตฺถุกํ สารมฺภํ อปริกฺกมนํ, อาปตฺติ ทฺวินฺนํ ทุกฺกฏานํ. ภิกฺขุ วิหารํ กโรติ เทสิต\-วตฺถุกํ สารมฺภํ สปริกฺกมนํ, อาปตฺติ ทุกฺกฏสฺส. ภิกฺขุ วิหารํ กโรติ เทสิต\-วตฺถุกํ อนารมฺภํ อปริกฺกมนํ, อาปตฺติ ทุกฺกฏสฺส. ภิกฺขุ วิหารํ กโรติ เทสิต\-วตฺถุกํ อนารมฺภํ สปริกฺกมนํ, อนาปตฺติ.}

\proseitem{373}{\csromanfolio{241}ภิกฺขุ สมาทิสติ “วิหารํ เม กโรถา”ติ. ตสฺส วิหารํ กโรนฺติ อเทสิต\-วตฺถุกํ สารมฺภํ อปริกฺกมนํ, อาปตฺติ สํฆาทิเสเสน ทฺวินฺนํ ทุกฺกฏานํ ฯเปฯ สารมฺภํ สปริกฺกมนํ, อาปตฺติ สํฆาทิเสเสน ทุกฺกฏสฺส ฯเปฯ อนารมฺภํ อปริกฺกมนํ, อาปตฺติ สํฆาทิเสเสน ทุกฺกฏสฺส ฯเปฯ อนารมฺภํ สปริกฺกมนํ, อาปตฺติ สํฆาทิเสสสฺส.}

\prose{ภิกฺขุ สมาทิสติ “วิหารํ เม กโรถา”ติ. ตสฺส วิหารํ กโรนฺติ เทสิต\-วตฺถุกํ สารมฺภํ อปริกฺกมนํ, อาปตฺติ ทฺวินฺนํ ทุกฺกฏานํ ฯเปฯ สารมฺภํ สปริกฺกมนํ, อาปตฺติ ทุกฺกฏสฺส ฯเปฯ อนารมฺภํ อปริกฺกมนํ, อาปตฺติ ทุกฺกฏสฺส ฯเปฯ อนารมฺภํ สปริกฺกมนํ, อนาปตฺติ.}

\proseitem{374}{ภิกฺขุ สมาทิสิตฺวา ปกฺกมติ “วิหารํ เม กโรถา”ติ, น จ สมาทิสติ “เทสิตวตฺถุโก จ โหตุ อนารมฺโภ จ สปริกฺกมโน จา”ติ. ตสฺส วิหารํ กโรนฺติ อเทสิต\-วตฺถุกํ สารมฺภํ อปริกฺกมนํ, อาปตฺติ สํฆาทิเสเสน ทฺวินฺนํ ทุกฺกฏานํ ฯเปฯ สารมฺภํ สปริกฺกมนํ อาปตฺติ สํฆาทิเสเสน ทุกฺกฏสฺส ฯเปฯ อนารมฺภํ อปริกฺกมนํ, อาปตฺติ สํฆาทิเสเสน ทุกฺกฏสฺส ฯเปฯ อนารมฺภํ สปริกฺกมนํ, อาปตฺติ สํฆาทิเสสสฺส.}

\prose{ภิกฺขุ สมาทิสิตฺวา ปกฺกมติ “วิหารํ เม กโรถา”ติ, น จ สมาทิสติ “เทสิตวตฺถุโก จ โหตุ อนารมฺโภ จ สปริกฺกมโน จา”ติ. ตสฺส วิหารํ กโรนฺติ เทสิต\-วตฺถุกํ สารมฺภํ อปริกฺกมนํ, อาปตฺติ ทฺวินฺนํ ทุกฺกฏานํ ฯเปฯ สารมฺภํ สปริกฺกมนํ, อาปตฺติ ทุกฺกฏสฺส ฯเปฯ อนารมฺภํ อปริกฺกมนํ, อาปตฺติ ทุกฺกฏสฺส ฯเปฯ อนารมฺภํ สปริกฺกมนํ, อนาปตฺติ.}

\proseitem{375}{ภิกฺขุ สมาทิสิตฺวา ปกฺกมติ “วิหารํ เม กโรถา”ติ, สมาทิสติ จ “เทสิตวตฺถุโก จ โหตุ อนารมฺโภ จ สปริกฺกมโน จา”ติ. ตสฺส วิหารํ กโรนฺติ อเทสิต\-วตฺถุกํ สารมฺภํ อปริกฺกมนํ. โส สุณาติ “วิหาโร กิร เม กยิรติ อเทสิต\-วตฺถุโก สารมฺโภ อปริกฺกมโน”ติ. เตน ภิกฺขุนา สามํ วา คนฺตพฺพํ ทูโต วา ปาเหตพฺโพ “เทสิตวตฺถุโก จ โหตุ อนารมฺโภ จ สปริกฺกมโน จา”ติ ฯเปฯ “เทสิตวตฺถุโก จ โหตุ อนารมฺโภ จา”ติ ฯเปฯ “เทสิตวตฺถุโก จ โหตุ สปริกฺกมโน \csromanfolio{242}จา”ติ ฯเปฯ “เทสิตวตฺถุโก โหตู”ติ. โน เจ สามํ วา คจฺเฉยฺย ทูตํ วา ปหิเณยฺย, อาปตฺติ ทุกฺกฏสฺส.}

\prose{ภิกฺขุ สมาทิสิตฺวา ปกฺกมติ “วิหารํ เม กโรถา”ติ, สมาทิสติ จ “เทสิตวตฺถุโก จ โหตุ อนารมฺโภ จ สปริกฺกมโน จา”ติ. ตสฺส วิหารํ กโรนฺติ เทสิต\-วตฺถุกํ สารมฺภํ อปริกฺกมนํ. โส สุณาติ “วิหาโร กิร เม กยิรติ เทสิตวตฺถุโก สารมฺโภ อปริกฺกมโน”ติ. เตน ภิกฺขุนา สามํ วา คนฺตพฺพํ ทูโต วา ปาเหตพฺโพ “อนารมฺโภ จ โหตุ สปริกฺกมโน จา”ติ ฯเปฯ “อนารมฺโภ โหตู”ติ. “สปริกฺกมโน โหตู”ติ. อนาปตฺติ.}

\proseitem{376}{ภิกฺขุ สมาทิสิตฺวา ปกฺกมติ “วิหารํ เม กโรถา”ติ, สมาทิสติ จ “เทสิตวตฺถุโก จ โหตุ อนารมฺโภ จ สปริกฺกมโน จา”ติ. ตสฺส วิหารํ กโรนฺติ อเทสิต\-วตฺถุกํ สารมฺภํ อปริกฺกมนํ, อาปตฺติ การุกานํ ติณฺณํ ทุกฺกฏานํ ฯเปฯ สารมฺภํ สปริกฺกมนํ, อาปตฺติ การุกานํ ทฺวินฺนํ ทุกฺกฏานํ ฯเปฯ อนารมฺภํ อปริกฺกมนํ, อาปตฺติ การุกานํ ทฺวินฺนํ ทุกฺกฏานํ ฯเปฯ อนารมฺภํ สปริกฺกมนํ, อาปตฺติ การุกานํ ทุกฺกฏสฺส.}

\prose{ภิกฺขุ สมาทิสิตฺวา ปกฺกมติ “วิหารํ เม กโรถา”ติ, สมาทิสติ จ “เทสิตวตฺถุโก จ โหตุ อนารมฺโภ จ สปริกฺกมโน จา”ติ. ตสฺส วิหารํ กโรนฺติ เทสิต\-วตฺถุกํ สารมฺภํ อปริกฺกมนํ, อาปตฺติ การุกานํ ทฺวินฺนํ ทุกฺกฏานํ ฯเปฯ สารมฺภํ สปริกฺกมนํ, อาปตฺติ การุกานํ ทุกฺกฏสฺส ฯเปฯ อนารมฺภํ อปริกฺกมนํ, อาปตฺติ การุกานํ ทุกฺกฏสฺส ฯเปฯ อนารมฺภํ สปริกฺกมนํ, อนาปตฺติ.}

\proseitem{377}{ภิกฺขุ สมาทิสิตฺวา ปกฺกมติ “วิหารํ เม กโรถา”ติ. ตสฺส วิหารํ กโรนฺติ อเทสิต\-วตฺถุกํ สารมฺภํ อปริกฺกมนํ. โส เจ วิปฺปกเต อาคจฺฉติ, เตน ภิกฺขุนา โส วิหาโร อญฺญสฺส วา ทาตพฺโพ ภินฺทิตฺวา วา ปุน กาตพฺโพ, โน เจ อญฺญสฺส วา ทเทยฺย ภินฺทิตฺวา วา ปุน กาเรยฺย, อาปตฺติ สํฆาเทเสเสน ทฺวินฺนํ ทุกฺกฏานํ ฯเปฯ สารมฺภํ สปริกฺกมนํ, อาปตฺติ สํฆาทิเสเสน ทุกฺกฏสฺส ฯเปฯ อนารมฺภํ อปริกฺกมนํ, อาปตฺติ สํฆาทิเสเสน ทุกฺกฏสฺส ฯเปฯ อนารมฺภํ สปริกฺกมนํ, อาปตฺติ สํฆาทิเสสสฺส.}

\csromanmatikaanchor{mtk.77}
\csromanmatikaanchor{mtk.78}
\csromantocmarknopage{section}{8. ทุฏฺฐโทสสิกฺขาปท}{mtk.77}
\bookmark[dest={mtk.77},level=1]{8. ทุฏฺฐโทสสิกฺขาปท}
\csromantocmarkref{subsection}{ทพฺพมลฺลปุตฺตวตฺถุ}{mtk.78}
\bookmark[dest={mtk.78},level=2]{ทพฺพมลฺลปุตฺตวตฺถุ}
\prose{\csromanfolio{243}ภิกฺขุ สมาทิสิตฺวา ปกฺกมติ “วิหารํ เม กโรถา”ติ. ตสฺส วิหารํ กโรนฺติ เทสิต\-วตฺถุกํ สารมฺภํ อปริกฺกมนํ. โส เจ วิปฺปกเต อาคจฺฉติ, เตน ภิกฺขุนา โส วิหาโร อญฺญสฺส วา ทาตพฺโพ ภินฺทิตฺวา วา ปุน กาตพฺโพ, โน เจ อญฺญสฺส วา ทเทยฺย ภินฺทิตฺวา วา ปุน กาเรยฺย, อาปตฺติ ทฺวินฺนํ ทุกฺกฏานํ ฯเปฯ สารมฺภํ สปริกฺกมนํ, อาปตฺติ ทุกฺกฏสฺส ฯเปฯ อนารมฺภํ อปริกฺกมนํ, อาปตฺติ ทุกฺกฏสฺส ฯเปฯ อนารมฺภํ สปริกฺกมนํ, อนาปตฺติ.}

\proseitem{378}{อตฺตนา วิปฺปกตํ อตฺตนา ปริโยสาเปติ, อาปตฺติ สํฆาเทเสสสฺส.}

\prose{อตฺตนา วิปฺปกตํ ปเรหิ ปริโยสาเปติ, อาปตฺติ สํฆาทิเสสสฺส.}

\prose{ปเรหิ วิปฺปกตํ อตฺตนา ปริโยสาเปติ, อาปตฺติ สํฆาทิเสสสฺส.}

\prose{ปเรหิ วิปฺปกตํ ปเรหิ ปริโยสาเปติ, อาปตฺติ สํฆาทิเสสสฺส.}

\proseitemwithcloserruled{379}{อนาปตฺติ เลเณ คุหาย ติณกุฏิกาย อญฺญสฺสตฺถาย วาสาคารํ ฐเปตฺวา สพฺพตฺถ, อนาปตฺติ อุมฺมตฺตกสฺส อาทิกมฺมิกสฺสาติ.}{วิหารการ\-สิกฺขาปทํ นิฏฺฐิตํ สตฺตมํ.}
\csromanheader{1}{8. \textbf{ทุฏฺฐ}\-\textbf{โทส}\-\textbf{สิกฺขาปท}}
\proseitem{380}{\csromansymbolfootnote{*}{อิทํ วตฺถุ วิ 4. 191 ปิฏฺฐาทีสุปิ อาคตํ.}เตน สมเยน พุทฺโธ ภควา ราชคเห วิหรติ เวฬุวเน กลนฺทก\-นิวาเป. เตน โข ปน สมเยน อายสฺมตา ทพฺเพน มลฺลปุตฺเตน ชาติยา สตฺตวสฺเสน อรหตฺตํ สจฺฉิกตํ โหติ, ยํ กิญฺจิ\csromansharedfootnote{3dd8a0fe7d537136}{ยญฺจ กิญฺจิ (สี, ก)} สาวเกน ปตฺตพฺพํ, สพฺพํ เตน อนุปฺปตฺตํ โหติ, นตฺถิ จสฺส กิญฺจิ อุตฺตริ กรณียํ, กตสฺส วา ปติจโย. อถ โข อายสฺมโต ทพฺพสฺส มลฺลปุตฺตสฺส รโหคตสฺส ปฏิสลฺลีนสฺส เอวํ เจตโส ปริวิตกฺโก อุทปาทิ “มยา โข ชาติยา สตฺตวสฺเสน อรหตฺตํ สจฺฉิกตํ, ยํ กิญฺจิ สาวเกน ปตฺตพฺพํ, สพฺพํ มยา อนุปฺปตฺตํ, นตฺถิ จ เม กิญฺจิ อุตฺตริ กรณียํ, กตสฺส วา ปติจโย, กิํ นุ โข อหํ สํฆสฺส เวยฺยาวจฺจํ กเรยฺยนฺ”ติ.}

\prose{อถ โข อายสฺมโต ทพฺพสฺส มลฺลปุตฺตสฺส เอตทโหสิ “ยํ นูนาหํ สํฆสฺส เสนาสนญฺจ ปญฺญเปยฺยํ ภตฺตานิ จ อุทฺทิเสยฺยนฺ”ติ. อถ โข \csromanfolio{244}อายสฺมา ทพฺโพ มลฺลปุตฺโต สายนฺหสมยํ ปฏิสลฺลานา วุฏฺฐิโต เยน ภควา เตนุปสงฺกมิ, อุปสงฺกมิตฺวา ภควนฺตํ อภิวาเทตฺวา เอกมนฺตํ นิสีทิ, เอกมนฺตํ นิสินฺโน โข อายสฺมา ทพฺโพ มลฺลปุตฺโต ภควนฺตํ เอตทโวจ “อิธ มยฺหํ ภนฺเต รโหคตสฺส ปฏิสลฺลีนสฺส เอวํ เจตโส ปริวิตกฺโก อุทปาทิ มยา โข ชาติยา สตฺตวสฺเสน อรหตฺตํ สจฺฉิกตํ, ยํ กิญฺจิ สาวเกน ปตฺตพฺพํ, สพฺพํ มยา อนุปฺปตฺตํ, นตฺถิ จ เม กิญฺจิ อุตฺตริ กรณียํ, กตสฺส วา ปติจโย, กิํ นุ โข อหํ สํฆสฺส เวยฺยาวจฺจํ กเรยฺยนฺ”ติ. ตสฺส มยฺหํ ภนฺเต เอตทโหสิ ‘ยํนูนาหํ สํฆสฺส เสนาสนญฺจ ปญฺญเปยฺยํ, ภตฺตานิ จ อุทฺทิเสยฺยนฺ’ติ, อิจฺฉามหํ ภนฺเต สํฆสฺส เสนาสนญฺจ ปญฺญเปตุํ ภตฺตานิ จ อุทฺทิสิตุนฺ”ติ. สาธุ สาธุ ทพฺพ, เตน หิ ตฺวํ ทพฺพ สํฆสฺส เสนาสนญฺจ ปญฺญเปหิ ภตฺตานิ จ อุทฺทิสาติ. “เอวํ ภนฺเต”ติ โข อายสฺมา ทพฺโพ มลฺลปุตฺโต ภควโต ปจฺจสฺโสสิ. อถ โข ภควา เอตสฺมิํ นิทาเน เอตสฺมิํ ปกรเณ ธมฺมิํ กถํ กตฺวา ภิกฺขู อามนฺเตสิ “เตน หิ ภิกฺขเว สํโฆ ทพฺพํ มลฺลปุตฺตํ เสนาสนปญฺญาปกญฺจ ภตฺตุทฺเทสกญฺจ สมฺมนฺนตุ, เอวญฺจ ปน ภิกฺขเว สมฺมนฺนิตพฺโพ, ปฐมํ ทพฺโพ มลฺลปุตฺโต ยาจิตพฺโพ. ยาจิตฺวา พฺยตฺเตน ภิกฺขุนา ปฏิพเลน สํโฆ ญาเปตพฺโพ\csromandash{}}

\proseitem{381}{“สุณาตุ เม ภนฺเต สํโฆ, ยทิ สํฆสฺส ปตฺตกลฺลํ, สํโฆ อายสฺมนฺตํ ทพฺพํ มลฺลปุตฺตํ เสนาสนปญฺญาปกญฺจ ภตฺตุทฺเทสกญฺจ สมฺมนฺเนยฺย, เอสา ญตฺติ.}

\prose{สุณาตุ เม ภนฺเต สํโฆ, สํโฆ อายสฺมนฺตํ ทพฺพํ มลฺลปุตฺตํ เสนาสนปญฺญาปกญฺจ ภตฺตุทฺเทสกญฺจ สมฺมนฺนติ, ยสฺสายสฺมโต ขมติ อายสฺมโต ทพฺพสฺส มลฺลปุตฺตสฺส เสนาสน\-ปญฺญาปกสฺส จ ภตฺตุ\-ทฺเทสกสฺส จ สมฺมุติ, โส ตุณฺหสฺส, ยสฺส นกฺขมติ, โส ภาเสยฺย.}

\prose{สมฺมโต สํเฆน อายสฺมา ทพฺโพ มลฺลปุตฺโต เสนาสน\-ปญฺญาปโก จ ภตฺตุทฺเทสโก จ, ขมติ สํฆสฺส, ตสฺมา ตุณฺหี, เอวเมตํ ธารยามี”ติ.}

\proseitem{382}{สมฺมโต จ ปนายสฺมา ทพฺโพ มลฺลปุตฺโต สภาคานํ ภิกฺขูนํ เอกชฺฌํ เสนาสนํ ปญฺญเปติ, เย เต ภิกฺขู สุตฺตนฺติกา, เตสํ เอกชฺฌํ \csromanfolio{245}เสนาสนํ ปญฺญเปติ “เต อญฺญมญฺญํ สุตฺตนฺตํ สงฺคายิสฺสนฺตี”ติ. เย เต ภิกฺขู วินยธรา, เตสํ เอกชฺฌํ เสนาสนํ ปญฺญเปติ “เต อญฺญมญฺญํ วินยํ วินิจฺฉินิ\-สฺสนฺตี”ติ\csromansharedfootnote{39e0eb03fca631f6}{ปินิจฺฉิสฺสนฺตีติ (ก)}. เย เต ภิกฺขู ธมฺมกถิกา, เตสํ เอกชฺฌํ เสนาสนํ ปญฺญเปติ “เต อญฺญมญฺญํ ธมฺมํ สากจฺฉิสฺสนฺตี”ติ. เย เต ภิกฺขู ฌายิโน, เตสํ เอกชฺฌํ เสนาสนํ ปญฺญเปติ “เต อญฺญมญฺญํ น พฺยาพาธิสฺสนฺตี”ติ\csromansharedfootnote{355212f63f4faf69}{น พฺยาพาหิสฺสนฺตีติ (ก)}. เย เต ภิกฺขู ติรจฺฉาน\-กถิกา กายทฬฺหิ\-พหุลา\csromansharedfootnote{3fb9caade9574027}{กายทฑฺฒิพหุลา (สี)} วิหรนฺติ, เตสํปิ เอกชฺฌํ เสนาสนํ ปญฺญเปติ “อิมายปิเม อายสฺมนฺโต รติยา อจฺฉิสฺสนฺตี”ติ. เยปิ เต ภิกฺขู วิกาเล อาคจฺฉนฺติ, เตสํปิ เตโชธาตุํ สมาปชฺชิตฺวา เตเนว อาโลเกน เสนาสนํ ปญฺญเปติ, อปิสุ ภิกฺขู สญฺจิจฺจ วิกาเล อาคจฺฉนฺติ “มยํ อายสฺมโต ทพฺพสฺส มลฺลปุตฺตสฺส อิทฺธิ\-ปาฏิหาริยํ ปสฺสิสฺสามา”ติ.}

\prose{เต อายสฺมนฺตํ ทพฺพํ มลฺลปุตฺตํ อุปสงฺกมิตฺวา เอวํ วทนฺติ “อมฺหากํ อาวุโส ทพฺพ เสนาสนํ ปญฺญเปหี”ติ. เต อายสฺมา ทพฺโพ มลฺลปุตฺโต เอวํ วเทติ “กตฺถายสฺมนฺตา อิจฺฉนฺติ กตฺถ ปญฺญเปมี”ติ. เต สญฺจิจฺจ ทูเร อปทิสนฺติ “อมฺหากํ อาวุโส ทพฺพ คิชฺฌกูเฏ ปพฺพเต เสนาสนํ ปญฺญเปหิ, อมฺหากํ อาวุโส โจรปปาเต เสนาสนํ ปญฺญเปหิ, อมฺหากํ อาวุโส อิสิคิลิปสฺเส กาฬสิลายํ เสนาสนํ ปญฺญเปหิ, อมฺหากํ อาวุโส เวภารปสฺเส สตฺตปณฺณิ\-คุหายํ เสนาสนํ ปญฺญเปหิ, อมฺหากํ อาวุโส สีตวเน สปฺป\-โสณฺฑิก\-ปพฺภาเร เสนาสนํ ปญฺญเปหิ, อมฺหากํ อาวุโส โคตมก\-กนฺทรายํ เสนาสนํ ปญฺญเปหิ, อมฺหากํ อาวุโส ตินฺทุก\-กนฺทรายํ เสนาสนํ ปญฺญเปหิ, อมฺหากํ อาวุโส ตโปท\-กนฺทรายํ เสนาสนํ ปญฺญเปหิ, อมฺหากํ อาวุโส ตโปทาราเม เสนาสนํ ปญฺญเปหิ, อมฺหากํ อาวุโส ชีวกมฺพวเน เสนาสนํ ปญฺญเปหิ, อมฺหากํ อาวุโส มทฺท\-กุจฺฉิสฺมิํ มิคทาเย เสนาสนํ ปญฺญเปหี”ติ.}

\prose{เตสํ อายสฺมา ทพฺโพ มลฺลปุตฺโต เตโชธาตุํ สมาปชฺชิตฺวา องฺคุลิยา ชลมานาย ปุรโต ปุรโต คจฺฉติ. เตปิ เตเนว อาโลเกน อายสฺมโต ทพฺพสฺส มลฺลปุตฺตสฺส ปิฏฺฐิโต ปิฏฺฐิโต คจฺฉนฺติ. เตสํ อายสฺมา ทพฺโพ มลฺลปุตฺโต เอวํ เสนาสนํ ปญฺญเปติ “อยํ มญฺโจ, อิทํ ปีฐํ, อยํ ภิสิ, อิทํ พิพฺโพหนํ, อิทํ วจฺจฏฺฐานํ, อิทํ \csromanfolio{246}ปสฺสวฏฺฐานํ, อิทํ ปานียํ, อิทํ ปริโภชนียํ, อยํ กตฺตรทณฺโฑ, อิทํ สํฆสฺส กติก\-สณฺฐานํ, อิมํ กาลํ ปวิสิตพฺพํ, อิมํ กาลํ นิกฺขมิตพฺพนฺ”ติ. เตสํ อายสฺมา ทพฺโพ มลฺลปุตฺโต เอวํ เสนาสนํ ปญฺญเปตฺวา ปุนเทว เวฬุวนํ ปจฺจาคจฺฉติ.}

\proseitem{383}{เตน โข ปน สมเยน เมตฺติย\-ภูมชกา\csromansharedfootnote{ee3e7b05283a693a}{ภุมฺมชกา (สี, สฺยา)} ภิกฺขู นวกา เจว โหนฺติ อปฺปปุญฺญา จ, ยานิ สํฆสฺส ลามกานิ เสนาสนานิ, ตานิ เตสํ ปาปุณนฺติ ลามกานิ จ ภตฺตานิ. เตน โข ปน สมเยน ราชคเห มนุสฺสา อิจฺฉนฺติ เถรานํ ภิกฺขูนํ อภิสงฺขาริกํ ปิณฺฑปาตํ ทาตุํ สปฺปิมฺปิ เตลมฺปิ อุตฺตริ\-ภงฺคมฺปิ. เมตฺติย\-ภูมชกานํ ปน ภิกฺขูนํ ปากติกํ เทนฺติ ยถารนฺธํ กณาชกํ พิลงฺคทุติยํ. เต ปจฺฉาภตฺตํ ปิณฺฑปาต\-ปฺปฏิกฺกนฺตา เถเร ภิกฺขู ปุจฺฉนฺติ “ตุมฺหากํ อาวุโส ภตฺตคฺเค กิํ อโหสิ, ตุมฺหากํ อาวุโส ภตฺตคฺเค กิํ อโหสี”ติ. เอกจฺเจ เถรา เอวํ วทนฺติ “อมฺหากํ อาวุโส สปฺปิ อโหสิ, เตลํ อโหสิ, อุตฺตริภงฺคํ อโหสี”ติ. เมตฺติย\-ภูมชกา ปน ภิกฺขู เอวํ วทนฺติ “อมฺหากํ อาวุโส น กิญฺจิ อโหสิ, ปากติกํ ยถารนฺธํ กณาชกํ พิลงฺคทุติยนฺ”ติ.}

\prose{เตน โข ปน สมเยน กลฺยาณภตฺติโก คหปติ สํฆสฺส จตุกฺกภตฺตํ เทติ นิจฺจภตฺตํ, โส ภตฺตคฺเค สปุตฺตทาโร อุปติฏฺฐิตฺวา ปริวิสติ, อญฺเญ โอทเนน ปุจฺฉนฺติ, อญฺเญ สูเปน ปุจฺฉนฺติ, อญฺเญ เตเลน ปุจฺฉนฺติ, อญฺเญ อุตฺตริภงฺเคน ปุจฺฉนฺติ. เตน โข ปน สมเยน กลฺยาณภตฺติกสฺส คหปติโน ภตฺตํ สฺวาตนาย เมตฺติย\-ภูมชกานํ ภิกฺขูนํ อุทฺทิฏฺฐํ โหติ. อถ โข กลฺยาณภตฺติโก คหปติ อารามํ อคมาสิ เกนจิเทว กรณีเยน. โส เยนายสฺมา ทพฺโพ มลฺลปุตฺโต เตนุปสงฺกมิ, อุปสงฺกมิตฺวา อายสฺมนฺตํ ทพฺพํ มลฺลปุตฺตํ อภิวาเทตฺวา เอกมนฺตํ นิสีทิ, เอกมนฺตํ นิสินฺนํ โข กลฺยาณภตฺติกํ คหปติํ อายสฺมา ทพฺโพ มลฺลปุตฺโต ธมฺมิยา กถาย สนฺทสฺเสสิ สมาทเปสิ สมุตฺเตเชสิ สมฺปหํเสสิ. อถ โข กลฺยาณภตฺติโก คหปติ อายสฺมตา ทพฺเพน มลฺลปุตฺเตน ธมฺมิยา กถาย สนฺทสฺสิโต สมาทปิโต สมุตฺเตชิโต สมฺปหํสิโต อายสฺมนฺตํ ทพฺพํ มลฺลปุตฺตํ เอตทโวจ “กสฺส ภนฺเต อมฺหากํ ฆเร สฺวาตนาย ภตฺตํ อุทฺทิฏฺฐนฺ”ติ. เมตฺติย\-ภูมชกานํ โข คหปติ ภิกฺขูนํ \csromanfolio{247}ตุมฺหากํ ฆเร สฺวาตนาย ภตฺตํ อุทฺทิฏฺฐนฺติ. อถ โข กลฺยาณภตฺติโก คหปติ อนตฺตมโน อโหสิ “กถํ หิ นาม ปาปภิกฺขู อมฺหากํ ฆเร ภุญฺชิสฺสนฺตี”ติ ฆรํ คนฺตฺวา ทาสิํ อาณาเปสิ “เย เช เสฺว ภตฺติกา อาคจฺฉนฺติ, เต โกฏฺฐเก อาสนํ ปญฺญเปตฺวา กณาชเกน พิลงฺค\-ทุติเยน ปริวิสา”ติ. เอวํ อยฺยาติ โข สา ทาสี กลฺยาณภตฺติกสฺส คหปติโน ปจฺจสฺโสสิ.}

\prose{อถ โข เมตฺติย\-ภูมชกา ภิกฺขู “หิยฺโย โข อาวุโส อมฺหากํ กลฺยาณภตฺติกสฺส คหปติโน ภตฺตํ อุทฺทิฏฺฐํ, เสฺว อเมฺห กลฺยาณภตฺติโก คหปติ สปุตฺตทาโร อุปติฏฺฐิตฺวา ปริวิสิสฺสติ, อญฺเญ โอทเนน ปุจฺฉิสฺสนฺติ, อญฺเญ สูเปน ปุจฺฉิสฺสนฺติ, อญฺเญ เตเลน ปุจฺฉิสฺสนฺติ, อญฺเญ อุตฺตริภงฺเคน ปุจฺฉิสฺสนฺตี”ติ เต เตเนว โสมนสฺเสน น จิตฺตรูปํ รตฺติยา สุปิํสุ. อถ โข เมตฺติย\-ภูมชกา ภิกฺขู ปุพฺพณฺหสมยํ นิวาเสตฺวา ปตฺตจีวรํ อาทาย เยน กลฺยาณภตฺติกสฺส คหปติโน นิเวสนํ เตนุ\-ปสงฺกมิํสุ. อทฺทสา โข สา ทาสี เมตฺติย\-ภูมชเก ภิกฺขู ทูรโตว อาคจฺฉนฺเต, ทิสฺวาน โกฏฺฐเก อาสนํ ปญฺญเปตฺวา เมตฺติย\-ภูมชเก ภิกฺขู เอตทโวจ “นิสีทถ ภนฺเต”ติ. อถ โข เมตฺติย\-ภูมชกานํ ภิกฺขูนํ เอตทโหสิ “นิสฺสํสยํ โข น ตาว ภตฺตํ สิทฺธํ ภวิสฺสติ, ยถา\csromansharedfootnote{1d781ac9c8c095c8}{ยาว (สี)} มยํ โกฏฺฐเก นิสีเทยฺยามา”ติ\csromansharedfootnote{0d8f6512085c892a}{นิสีทาปีเยยฺยามาติ (สี), นิสีทาปียามาติ (สฺยา)}. อถ โข สา ทาสี กณาชเกน พิลงฺค\-ทุติเยน อุปคจฺฉิ “ภุญฺชถ ภนฺเต”ติ. มยํ โข ภคินิ นิจฺจภตฺติกาติ. ชานามิ อยฺยา นิจฺจภตฺติกาติ, อปิจาหํ หิยฺโยว คหปตินา อาณตฺตา “เย เช เสฺว ภตฺติกา อาคจฺฉนฺติ, เต โกฏฺฐเก อาสนํ ปญฺญเปตฺวา กณาชเกน พิลงฺค\-ทุติเยน ปริวิสา”ติ, ภุญฺชถ ภนฺเตติ. อถ โข เมตฺติย\-ภูมชกา ภิกฺขู “หิยฺโย โข อาวุโส กลฺยาณภตฺติโก คหปติ อารามํ อคมาสิ ทพฺพสฺส มลฺลปุตฺตสฺส สนฺติเก, นิสฺสํสยํ โข มยํ ทพฺเพน มลฺลปุตฺเตน คหปติโน อนฺตเร\csromansharedfootnote{fd8df2b8b430675c}{สนฺติเก (สฺยา, กํ, ก)} ปริภินฺนา”ติ เต เตเนว โทมนสฺเสน น จิตฺตรูปํ ภุญฺชิํสุ. อถ โข เมตฺติย\-ภูมชกา ภิกฺขู ปจฺฉาภตฺตํ ปิณฺฑปาต\-ปฺปฏิกฺกนฺตา อารามํ คนฺตฺวา ปตฺตจีวรํ ปฏิสาเมตฺวา พหารามโกฏฺเฐเก สงฺฆาฏิ\-ปลฺลตฺถิกาย นิสีทิํสุ ตุณฺหีภูตา มงฺกุภูตา ปตฺตกฺขนฺธา อโธมุขา ปชฺฌายนฺตา อปฺปฏิภานา.}

\prose{\csromanfolio{248}อถ โข เมตฺติยา ภิกฺขุนี เยน เมตฺติย\-ภูมชกา ภิกฺขู เตนุปสงฺกมิ, อุปสงฺกมิตฺวา เมตฺติย\-ภูมชเก ภิกฺขู เอตทโวจ “วนฺทามิ อยฺยา”ติ. เอวํ วุตฺเต เมตฺติย\-ภูมชกา ภิกฺขู นาลปิํสุ. ทุติยมฺปิ โข ฯเปฯ. ตติยมฺปิ โข เมตฺติยา ภิกฺขุนี เมตฺติย\-ภูมชเก ภิกฺขู เอตทโวจ “วนฺทามิ อยฺยา”ติ. ตติยมฺปิ โข เมตฺติย\-ภูมชกา ภิกฺขู นาลปิํสุ. กฺยาหํ อยฺยานํ อปรชฺฌามิ, กิสฺส มํ อยฺยา นาลปนฺตีติ. ตถา หิ ปน ตฺวํ ภคินิ อเมฺห ทพฺเพน มลฺลปุตฺเตน วิเหฐียมาเน อชฺฌุเปกฺขสีติ. กฺยาหํ อยฺยา กโรมีติ. สเจ โข ตฺวํ ภคินิ อิจฺเฉยฺยาสิ, อชฺเชว ภควา ทพฺพํ มลฺลปุตฺตํ นาสาเปยฺยาติ. กฺยาหํ อยฺยา กโรมิ, กิํ มยา สกฺกา กาตุนฺติ. เอหิ ตฺวํ ภคินิ เยน ภควา เตนุปสงฺกมิ, อุปสงฺกมิตฺวา ภควนฺตํ เอวํ วเทหิ “อิทํ ภนฺเต นจฺฉนฺนํ นปฺปติรูปํ, ยายํ ภนฺเต ทิสา อภยา อนีติกา อนุปทฺทวา, สายํ ทิสา สภยา สอีติกา สอุปทฺทวา, ยโต นิวาตํ ตโต สเวตํ\csromansharedfootnote{c5626aa331030b5a}{ปวาตํ (สี, สฺยา)}, อุทกํ มญฺเญ อาทิตฺตํ, อยฺเยนมฺหิ ทพฺเพน มลฺลปุตฺเตน ทูสิตา”ติ. เอวํ อยฺยาติ โข เมตฺติยา ภิกฺขุนี เมตฺติย\-ภูมชกานํ ภิกฺขูนํ ปฏิสฺสุตฺวา เยน ภควา เตนุปสงฺกมิ, อุปสงฺกมิตฺวา ภควนฺตํ อภิวาเทตฺวา เอกมนฺตํ อฏฺฐาสิ, เอกมนฺตํ ฐิตา โข สา เมตฺติยา ภิกฺขุนี ภควนฺตํ เอตทโวจ “อิทํ ภนฺเต นจฺฉนฺนํ นปฺปติรูปํ, ยายํ ภนฺเต ทิสา อภยา อนีติกา อนุปทฺทวา, สายํ ทิสา สภย สอีติกา สอุปทฺทวา, ยโต นิวาตํ ตโต สวาตํ, อุทกํ มญฺเญ อาทิตฺตํ, อยฺเยนมฺหิ ทพฺเพน มลฺลปุตฺเตน ทูสิตา”ติ.}

\csromanmatikaanchor{mtk.79}
\csromanmatikaanchor{mtk.80}
\csromantocmarkref{subsection}{ปญฺญตฺติ}{mtk.79}
\bookmark[dest={mtk.79},level=2]{ปญฺญตฺติ}
\csromantocmarkref{subsection}{สิกฺขาปทวิภงฺค, ปทภาชนีย}{mtk.80}
\bookmark[dest={mtk.80},level=2]{สิกฺขาปทวิภงฺค, ปทภาชนีย}
\proseitem{384}{อถ โข ภควา เอตสฺมิํ นิทาเน เอตสฺมิํ ปกรเณ ภิกฺขุสํฆํ สนฺนิปาตาเปตฺวา อายสฺมนฺตํ ทพฺพํ มลฺลปุตฺตํ ปฏิปุจฺฉิ “สรสิ ตฺวํ ทพฺพ เอวรูปํ กตฺตา ยถายํ ภิกฺขุนี อาหา”ติ. ยถา มํ ภนฺเต ภควา ชานาตีติ. ทุติยมฺปิ โข ภควา ฯเปฯ. ตติยมฺปิ โข ภควา อายสฺมนฺตํ ทพฺพํ มลฺลปุตฺตํ เอตทโวจ “สรสิ ตฺวํ ทพฺพ เอวรูปํ กตฺตา ยถายํ ภิกฺขุนี อาหา”ติ. ยถา มํ ภนฺเต ภควา ชานาตีติ. น โข ทพฺพ ทพฺพา เอวํ นิพฺเพเฐนฺติ, สเจ ตยา กตํ กตนฺติ วเทหิ, สเจ ตยา อกตํ อกตนฺติ วเทหีติ. ยโต อหํ ภนฺเต ชาโต นาภิชานามิ สุปินนฺเตนปิ เมถุนํ ธมฺมํ ปฏิเสวิตา, ปเคว ชาคโรติ. \csromanfolio{249}อถ โข ภควา ภิกฺขู อามนฺเตสิ “เตน หิ ภิกฺขเว เมตฺติยํ ภิกฺขุนิํ นาเสถ, อิเม จ ภิกฺขู อนุยุญฺชถา”ติ. อิทํ วตฺวา ภควา \textbf{อุฏฺฐายาสนา วิหารํ ปาวิสิ.}}

\prose{อถ โข เต ภิกฺขู เมตฺติยํ ภิกฺขุนิํ นาเสสุํ. อถ โข เมตฺติย\-ภูมชกา ภิกฺขู เต ภิกฺขู เอตทโวจุํ “มาวุโส เมตฺติยํ ภิกฺขุนิํ นาเสถ, น สา กิญฺจิ อปรชฺฌติ, อเมฺหหิ สา อุสฺสาหิตา กุปิเตหิ อนตฺตมเนหิ จาวนาธิปฺปาเยหี”ติ. กิํ ปน ตุเมฺห อาวุโส อายสฺมนฺตํ ทพฺพํ มลฺลปุตฺตํ อมูลเกน ปาราชิเกน ธมฺเมน อนุทฺธํเส\-ถาติ. เอวมาวุโสติ. เย เต ภิกฺขู อปฺปิจฺฉา ฯเปฯ. เต อุชฺฌายนฺติ ขิยฺยนฺติ วิปาเจนฺติ “กถํ หิ นาม เมตฺติย\-ภูมชกา ภิกฺขู อายสฺมนฺตํ ทพฺพํ มลฺลปุตฺตํ อมูลเกน ปาราชิเกน ธมฺเมน อนุทฺธํเสสฺสนฺตี”ติ. อถ โข เต ภิกฺขู เมตฺติย\-ภูมชเก ภิกฺขู อเนก\-ปริยาเยน วิครหิตฺวา ภควโต เอตมตฺถํ อาโรเจสุํ ฯเปฯ “สจฺจํ กิร ตุเมฺห ภิกฺขเว ทพฺพํ มลฺลปุตฺตํ อมูลเกน ปาราชิเกน ธมฺเมน อนุทฺธํเสถา”ติ. สจฺจํ ภควาติ. วิครหิ พุทฺโธ ภควา ฯเปฯ. กถํ หิ นาม ตุเมฺห โมฆปุริสา ทพฺพํ มลฺลปุตฺตํ อมูลเกน ปาราชิเกน ธมฺเมน อนุทฺธํเสสฺสถ, เนตํ โมฆปุริสา อปฺปสนฺนานํ วา ปสาทาย ฯเปฯ. เอวญฺจ ปน ภิกฺขเว อิมํ สิกฺขาปทํ อุทฺทิเสยฺยถ\csromandash{}}

\proseitem{385}{\textbf{“โย ปน ภิกฺขุ ภิกฺขุํ ทุฏฺโฐ โทโส อปฺปตีโต อมูลเกน ปาราชิเกน ธมฺเมน อนุทฺธํเสยฺย ‘อปฺเปว นาม นํ อิมมฺหา พฺรหฺมจริยา จาเวยฺยนฺ’ติ, ตโต อปเรน สมเยน สมนุ}\-\textbf{คฺคาหีย}\-\textbf{มาโน วา อสมนุ}\-\textbf{คฺคาหีย}\-\textbf{มาโน วา อมูลกญฺเจว ตํ อธิกรณํ โหติ, ภิกฺขุ จ โทสํ ปติฏฺฐาติ, สํฆาทิเสโส”}ติ.}

\proseitem{386}{\textbf{โย ปนา}ติ โย ยาทิโส ฯเปฯ. \textbf{ภิกฺขู}ติ ฯเปฯ อยํ อิมสฺมิํ อตฺเถ อธิปฺเปโต “ภิกฺขู”ติ.}

\prose{\textbf{ภิกฺขุนฺ}ติ อญฺญํ ภิกฺขุํ.}

\prose{\textbf{ทุฏฺโฐ โทโส}ติ กุปิโต อนตฺตมโน อนภิรทฺโธ อาหตจิตฺโต ขิลชาโต.}

\prose{\textbf{อปฺปตีโต}ติ เตน จ โกเปน เตน จ โทเสน ตาย จ อนตฺตมนตาย ตาย จ อนภิรทฺธิยา อปฺปตีโต โหติ.}

\prose{\csromanfolio{250}\textbf{อมูลกํ} นาม อทิฏฺฐํ อสุตํ อปริสงฺกิตํ.}

\prose{\textbf{ปาราชิเกน ธมฺเมนา}ติ จตุนฺนํ อญฺญตเรน.}

\prose{\textbf{อนุทฺธํเสยฺยา}ติ โจเทติ วา โจทาเปติ วา.}

\prose{\textbf{อปฺเปว นาม นํ อิมมฺหา พฺรหฺมจริยา จาเวยฺยนฺ}ติ ภิกฺขุภาวา จาเวยฺยํ, สมณธมฺมา จาเวยฺยํ, สีลกฺขนฺธา จาเวยฺยํ, ตโปคุณา จาเวยฺยํ.}

\prose{\textbf{ตโต อปเรน สมเรนา}ติ ยสฺมิํ ขเณ อนุทฺธํสิโต โหติ, ตํ ขณํ ตํ ลยํ ตํ มุหุตฺตํ วีติวตฺเต.}

\prose{\textbf{สมนุ}\-\textbf{คฺคาหีย}\-\textbf{มาโน}ติ เยน วตฺถุนา อนุทฺธํสิโต โหติ, ตสฺมิํ วตฺถุสฺมิํ สมนุ\-คฺคาหีย\-มาโน.}

\prose{\textbf{อสมนุคฺคาหียมาโน}ติ น เกนจิ วุจฺจมาโน.}

\prose{\textbf{อธิกรณํ} นาม จตฺตาริ อธิกรณานิ วิวาทา\-ธิกรณํ อนุวาทา\-ธิกรณํ อาปตฺตา\-ธิกรณํ กิจฺจา\-ธิกรณํ.}

\prose{\textbf{ภิกฺขุ จ โทสํ ปติฏฺฐาตี}ติ ตุจฺฉกํ มยา ภณิตํ, มุสา มยา ภณิตํ, อภูตํ มยา ภณิตํ, อชานนฺเตน มยา ภณิตํ.}

\prose{\textbf{สํฆาทิเสโส}ติ ฯเปฯ เตนปิ วุจฺจติ “สํฆาทิเสโส”ติ.}

\proseitem{387}{อทิฏฺฐสฺส โหติ ปาราชิกํ ธมฺมํ อชฺฌาปชฺชนฺโต, ตญฺเจ โจเทติ “ทิฏฺโฐ มยา ปาราชิกํ ธมฺมํ อชฺฌาปนฺโนสิ อสฺสมโณสิ อสกฺยปุตฺติโยสิ, นตฺถิ ตยา สทฺธิํ อุโปสโถ วา ปวารณา วา สํฆกมฺมํ วา”ติ, อาปตฺติ วาจาย วาจาย สํฆาทิเสสสฺส.}

\prose{อสุตสฺส โหติ “ปาราชิกํ ธมฺมํ อชฺฌาปนฺโน”ติ, ตญฺเจ โจเทติ “สุโต มยา ปาราชิกํ ธมฺมํ อชฺฌาปนฺโนสิ อสฺสมโณสิ อสกฺยปุตฺติโยสิ, นตฺถิ ตยา สทฺธิํ อุโปสโถ วา ปวารณา วา สํฆกมฺมํ วา”ติ, อาปตฺติ วาจาย วาจาย สํฆาทิเสสสฺส.}

\prose{อปริสงฺกิตสฺส โหติ “ปาราชิกํ ธมฺมํ อชฺฌาปนฺโน”ติ, ตญฺเจ โจเทติ “ปริสงฺกิโต มยา ปาราชิกํ ธมฺมํ อชฺฌาปนฺโนสิ อสฺสมโณสิ อสกฺยปุตฺติโยสิ, นตฺถิ ตยา สทฺธิํ อุโปสโถ วา ปวารณา วา สํฆกมฺมํ วา”ติ, อาปตฺติ วาจาย วาจาย สํฆาทิเสสสฺส.}

\prose{\csromanfolio{251}อทิฏฺฐสฺส โหติ ปาราชิกํ ธมฺมํ อชฺฌาปชฺชนฺโต, ตญฺเจ โจเทติ “ทิฏฺโฐ มยา สุโต จ ปาราชิกํ ธมฺมํ อชฺฌาปนฺโนสิ อสฺสมโณสิ ฯเปฯ อาปตฺติ วาจาย วาจาย สํฆาทิเสสสฺส.}

\prose{อทิฏฺฐสฺส โหติ ปาราชิกํ ธมฺมํ อชฺฌาปชฺชนฺโต, ตญฺเจ โจเทติ “ทิฏฺโฐ มยา ปริสงฺกิโต จ ปาราชิกํ ธมฺมํ อชฺฌาปนฺโนสิ อสฺสมโณสิ ฯเปฯ อาปตฺติ วาจาย วาจาย สํฆาทิเสสสฺส.}

\prose{อทิฏฺฐสฺส โหติ ปาราชิกํ ธมฺมํ อชฺฌาปชฺชนฺโต, ตญฺเจ โจเทติ “ทิฏฺโฐ มยา สุโต จ ปริสงฺกิโต จ ปาราชิกํ ธมฺมํ อชฺฌาปนฺโนสิ อสฺสมโณสิ ฯเปฯ อาปตฺติ วาจาย วาจาย สํฆาทิเสสสฺส.}

\prose{อสุตสฺส โหติ “ปาราชิกํ ธมฺมํ อชฺฌาปนฺโน”ติ, ตญฺเจ โจเทติ “สุโต มยา ปริสงฺกิโต จ ฯเปฯ. สุโต มยา ทิฏฺโฐ จ ฯเปฯ. สุโต มยา ปริสงฺกิโต จ ทิฏฺโฐ จ ปาราชิกํ ธมฺมํ อชฺฌาปนฺโนสิ อสฺสมโณสิ ฯเปฯ อาปตฺติ วาจาย วาจาย สํฆาทิเสสสฺส.}

\prose{อปริสงฺกิตสฺส โหติ “ปาราชิกํ ธมฺมํ อชฺฌาปนฺโน”ติ, ตญฺเจ โจเทติ ปริสงฺกิโต มยา ทิฏฺโฐ จ ฯเปฯ ปริสงฺกิโต มยา สุโต จ ฯเปฯ ปริสงฺกิโต มยา ทิฏฺโฐ จ สุโต จ ปาราชิกํ ธมฺมํ อชฺฌาปนฺโนสิ อสฺสมโณสิ ฯเปฯ อาปตฺติ วาจาย วาจาย สํฆาทิเสสสฺส.}

\prose{ทิฏฺฐสฺส โหติ ปาราชิกํ ธมฺมํ อชฺฌาปชฺชนฺโต, ตญฺเจ โจเทติ “สุโต มยา ปาราชิกํ ธมฺมํ อชฺฌาปนฺโนสิ อสฺสมโณสิ ฯเปฯ อาปตฺติ วาจาย วาจาย สํฆาทิเสสสฺส.}

\prose{ทิฏฺฐสฺส โหติ ปาราชิกํ ธมฺมํ อชฺฌาปชฺชนฺโต, ตญฺเจ โจเทติ “ปริสงฺกิโต มยา ปาราชิกํ ธมฺมํ อชฺฌาปนฺโนสิ ฯเปฯ “สุโต มยา ปริสงฺกิโต จ ปาราชิกํ ธมฺมํ อชฺฌาปนฺโนสิ อสฺสมโณสิ ฯเปฯ อาปตฺติ วาจาย วาจาย สํฆาทิเสสสฺส.}

\prose{สุตสฺส โหติ “ปาราชิกํ ธมฺมํ อชฺฌาปนฺโน”ติ, ตญฺเจ โจเทติ “ปริสงฺกิโต มยา ปาราชิกํ ธมฺมํ อชฺฌาปนฺโนสิ ฯเปฯ “ทิฏฺโฐ มยา ปาราชิกํ ธมฺมํ อชฺฌาปนฺโนสิ ฯเปฯ ปริสงฺกิโต มยา ทิฏฺโฐ จ ปาราชิกํ ธมฺมํ อชฺฌาปนฺโนสิ อสฺสมโณสิ ฯเปฯ อาปตฺติ วาจาย วาจาย สํฆาทิเสสสฺส.}

\prose{\csromanfolio{252}ปริสงฺกิตสฺส โหติ “ปาราชิกํ ธมฺมํ อชฺฌาปนฺโน”ติ, ตญฺเจ โจเทติ “ทิฏฺโฐ มยา ปาราชิกํ ธมฺมํ อชฺฌาปนฺโนสิ ฯเปฯ “สุโต มยา ปาราชิกํ ธมฺมํ อชฺฌาปนฺโนสิ ฯเปฯ “ทิฏฺโฐ มยา สุโต จ ปาราชิกํ ธมฺมํ อชฺฌาปนฺโนสิ อสฺสมโณสิ อสกฺยปุตฺติโยสิ ฯเปฯ อาปตฺติ วาจาย วาจาย สํฆาทิเสสสฺส.}

\prose{ทิฏฺฐสฺส โหติ ปาราชิกํ ธมฺมํ อชฺฌาปชฺชนฺโต, ทิฏฺเฐ เวมติโก ทิฏฺฐํ โน กปฺเปติ ทิฏฺฐํ นสฺสรติ ทิฏฺฐํ ปมุฏฺโฐ โหติ ฯเปฯ สุเต เวมติโก สุตํ โน กปฺเปติ สุตํ นสฺสรติ สุตํ ปมุฏฺโฐ โหติ ฯเปฯ ปริสงฺกิเต เวมติโก ปริสงฺกิตํ โน กปฺเปติ ปริสงฺกิตํ นสฺสรติ ปริสงฺกิตํ ปมุฏฺโฐ โหติ. ตญฺเจ โจเทติ “ปริสงฺกิโต มยา ทิฏฺโฐ จ ฯเปฯ “ปริสงฺกิโต มยา สุโต จ ฯเปฯ “ปริสงฺกิโต มยา ทิฏฺโฐ จ สุโต จ ปาราชิกํ ธมฺมํ อชฺฌาปนฺโนสิ อสฺสมโณสิ อสกฺยปุตฺติโยสิ, นตฺถิ ตยา สทฺธิํ อุโปสโถ วา ปวารณา วา สํฆกมฺมํ วา”ติ, อาปตฺติ วาจาย วาจาย สํฆาทิเสสสฺส.}

\proseitem{388}{อทิฏฺฐสฺส โหติ ปาราชิกํ ธมฺมํ อชฺฌาปชฺชนฺโต, ตญฺเจ โจทาเปติ “ทิฏฺโฐสิ ปาราชิกํ ธมฺมํ อชฺฌาปนฺโนสิ อสฺสมโณสิ อสกฺยปุตฺติโยสิ, นตฺถิ ตยา สทฺธิํ อุโปสโถ วา ปวารณา วา สํฆกมฺมํ วา”ติ, อาปตฺติ วาจาย วาจาย สํฆาทิเสสสฺส.}

\prose{อสุตสฺส โหติ “ปาราชิกํ ธมฺมํ อชฺฌาปนฺโน”ติ ฯเปฯ. อปริสงฺกิตสฺส โหติ “ปาราชิกํ ธมฺมํ อชฺฌาปนฺโน”ติ, ตญฺเจ โจทาเปติ ปริสงฺกิโตสิ ปาราชิกํ ธมฺมํ อชฺฌาปนฺโนสิ ฯเปฯ อาปตฺติ วาจาย วาจาย สํฆาทิเสสสฺส.}

\prose{อทิฏฺฐสฺส โหติ ปาราชิกํ ธมฺมํ อชฺฌาปชฺชนฺโต, ตญฺเจ โจทาเปติ ทิฏฺโฐสิ สุโตสิ ฯเปฯ ทิฏฺโฐสิ ปริสงฺกิโตสิ ฯเปฯ ทิฏฺโฐสิ สุโตสิ ปริสงฺกิโตสิ ปาราชิกํ ธมฺมํ อชฺฌาปนฺโนสิ ฯเปฯ. อสุตสฺส โหติ “ปาราชิกํ ธมฺมํ อชฺฌาปนฺโน”ติ ฯเปฯ. อปริสงฺกิตสฺส โหติ “ปาราชิกํ ธมฺมํ อชฺฌาปนฺโน”ติ, ตญฺเจ โจทาเปติ ปริสงฺกิโตสิ ทิฏฺโฐสิ ฯเปฯ ปริสงฺกิโตสิ สุโตสิ ฯเปฯ ปริสงฺกิโตสิ ทิฏฺโฐสิ สุโตสิ ปาราชิกํ ธมฺมํ อชฺฌาปนฺโนสิ อสฺสมโณสิ ฯเปฯ อาปตฺติ วาจาย วาจาย สํฆาทิเสสสฺส.}

\prose{\csromanfolio{253}ทิฏฺฐสฺส โหติ ปาราชิกํ ธมฺมํ อชฺฌาปชฺชนฺโต, ตญฺเจ โจทาเปติ สุโตสิ ฯเปฯ ตญฺเจ โจทาเปติ ปริสงฺกิโตสิ ฯเปฯ ตญฺเจ โจทาเปติ สุโตสิ ปริสงฺกิโตสิ ปาราชิกํ ธมฺมํ อชฺฌาปนฺโนสิ อสฺสมโณสิ ฯเปฯ อาปตฺติ วาจาย วาจาย สํฆาทิเสสสฺส.}

\prose{สุตสฺส โหติ “ปาราชิกํ ธมฺมํ อชฺฌาปนฺโน”ติ ฯเปฯ. ปริสงฺกิตสฺส โหติ “ปาราชิกํ ธมฺมํ อชฺฌาปนฺโน”ติ, ตญฺเจโจทาเปติ ทิฏฺโฐสิ ฯเปฯ ตญฺเจ โจทาเปติ สุโตสิ ฯเปฯ ตญฺเจ โจทาเปติ ทิฏฺโฐสิ สุโตสิ ปาราชิกํ ธมฺมํ อชฺฌาปนฺโนสิ อสฺสมโณสิ ฯเปฯ อาปตฺติ วาจาย วาจาย สํฆาทิเสสสฺส.}

\prose{ทิฏฺฐสฺส โหติ ปาราชิกํ ธมฺมํ อชฺฌาปชฺชนฺโต, ทิฏฺเฐ เวมติโก ทิฏฺฐํ โน กปฺเปติ ทิฏฺฐํ นสฺสรติ ทิฏฺฐํ ปมุฏฺโฐ โหติ ฯเปฯ สุเต เวมติโก สุตํ โน กปฺเปติ สุตํ นสฺสรติ สุตํ ปมุฏฺโฐ โหติ ฯเปฯ ปริสงฺกิเต เวมติโก ปริสงฺกิตํ โน กปฺเปติ ปริสงฺกิตํ นสฺสรติ ปริสงฺกิตํ ปมุฏฺโฐ โหติ, ตญฺเจ โจทาเปติ “ปริสงฺกิโตสิ ทิฏฺโฐสิ ฯเปฯ ปริสงฺกิตํ ปมุฏฺโฐ โหติ, ตญฺเจ โจทาเปติ ปริสงฺกิโตสิ สุโตสิ ฯเปฯ ปริสงฺกิตํ ปมุฏฺโฐ โหติ, ตญฺเจ โจทาเปติ ปริสงฺกิโตสิ ทิฏฺโฐสิ สุโตสิ ปาราชิกํ ธมฺมํ อชฺฌาปนฺโนสิ อสฺสมโณสิ อสกฺยปุตฺติโยสิ, นตฺถิ ตยา สทฺธิํ อุโปสโถ วา ปวารณา วา สํฆกมฺมํ วา”ติ, อาปตฺติ วาจาย วาจาย สํฆาทิเสสสฺส.}

\proseitem{389}{อสุทฺเธ สุทฺธทิฏฺฐิ, สุทฺเธ อสุทฺธทิฏฺฐิ, อสุทฺเธ อสุทฺธทิฏฺฐิ, สุทฺเธ สุทฺธทิฏฺฐิ.}

\prose{อสุทฺโธ โหติ ปุคฺคโล อญฺญตรํ ปาราชิกํ ธมฺมํ อชฺฌาปนฺโน, ตญฺเจ สุทฺธทิฏฺฐิ สมาโน อโนกาสํ การาเปตฺวา จาวนาธิปฺปาโย วเทติ, อาปตฺติ สํฆาทิเสเสน ทุกฺกฏสฺส.}

\prose{อสุทฺโธ โหติ ปุคฺคโล อญฺญตรํ ปาราชิกํ ธมฺมํ อชฺฌาปนฺโน, ตญฺเจ สุทฺธทิฏฺฐิ สมาโน โอกาสํ การาเปตฺวา จาวนาธิปฺปาโย วเทติ, อาปตฺติ สํฆาทิเสสสฺส.}

\prose{\csromanfolio{254}อสุทฺโธ โหติ ปุคฺคโล อญฺญตรํ ปาราชิกํ ธมฺมํ อชฺฌาปนฺโน, ตญฺเจ สุทฺธทิฏฺฐิ สมาโน อโนกาสํ การาเปตฺวา อกฺโกสาธิปฺปาโย วเทติ, อาปตฺติ โอมสวาเทน ทุกฺกฏสฺส.}

\prose{อสุทฺโธ โหติ ปุคฺคโล อญฺญตรํ ปาราชิกํ ธมฺมํ อชฺฌาปนฺโน, ตญฺเจ สุทฺธทิฏฺฐิ สมาโน โอกาสํ การาเปตฺวา อกฺโกสาธิปฺปาโย วเทติ, อาปตฺติ โอมสวาทสฺส.}

\prose{สุทฺโธ โหติ ปุคฺคโล อญฺญตรํ ปาราชิกํ ธมฺมํ อนชฺฌาปนฺโน, ตญฺเจ อสุทฺธทิฏฺฐิ สมาโน อโนกาสํ การาเปตฺวา จาวนาธิปฺปาโย วเทติ, อาปตฺติ ทุกฺกฏสฺส.}

\prose{สุทฺโธ โหติ ปุคฺคโล อญฺญตรํ ปาราชิกํ ธมฺมํ อนชฺฌาปนฺโน, ตญฺเจ อสุทฺธทิฏฺฐิ สมาโน โอกาสํ การาเปตฺวา จาวนาธิปฺปาโย วเทติ, อนาปตฺติ.}

\prose{สุทฺโธ โหติ ปุคฺคโล อญฺญตรํ ปาราชิกํ ธมฺมํ อนชฺฌาปนฺโน, ตญฺเจ อสุทฺธทิฏฺฐิ สมาโน อโนกาสํ การาเปตฺวา อกฺโกสาธิปฺปาโย วเทติ, อาปตฺติ โอมสวาเทน ทุกฺกฏสฺส.}

\prose{สุทฺโธ โหติ ปุคฺคโล อญฺญตรํ ปาราชิกํ ธมฺมํ อนชฺฌาปนฺโน, ตญฺเจ อสุทฺธทิฏฺฐิ สมาโน โอกาสํ การาเปตฺวา อกฺโกสาธิปฺปาโย วเทติ, อาปตฺติ โอมสวาทสฺส.}

\prose{อสุทฺโธ โหติ ปุคฺคโล อญฺญตรํ ปาราชิกํ ธมฺมํ อชฺฌาปนฺโน, ตญฺเจ อสุทฺธทิฏฺฐิ สมาโน อโนกาสํ การาเปตฺวา จาวนาธิปฺปาโย วเทติ, อาปตฺติ ทุกฺกฏสฺส.}

\prose{อสุทฺโธ โหติ ปุคฺคโล อญฺญตรํ ปาราชิกํ ธมฺมํ อชฺฌาปนฺโน, ตญฺเจ อสุทฺธทิฏฺฐิ สมาโน โอกาสํ การาเปตฺวา จาวนาธิปฺปาโย วเทติ, อนาปตฺติ.}

\prose{อสุทฺโธ โหติ ปุคฺคโล อญฺญตรํ ปาราชิกํ ธมฺมํ อชฺฌาปนฺโน, ตญฺเจ อสุทฺธทิฏฺฐิ สมาโน อโนกาสํ การาเปตฺวา อกฺโกสาธิปฺปาโย วเทติ, อาปตฺติ โอมสวาเทน ทุกฺกฏสฺส.}

\prose{อสุทฺโธ โหติ ปุคฺคโล อญฺญตรํ ปาราชิกํ ธมฺมํ อชฺฌาปนฺโน, ตญฺเจ อสุทฺธทิฏฺฐิ สมาโน โอกาสํ การาเปตฺวา อกฺโกสาธิปฺปาโย วเทติ, อาปตฺติ โอมสวาทสฺส.}

\csromanmatikaanchor{mtk.81}
\csromanmatikaanchor{mtk.82}
\csromantocmarknopage{section}{9. ทุติยทุฏฺฐโทสสิกฺขาปท}{mtk.81}
\bookmark[dest={mtk.81},level=1]{9. ทุติยทุฏฺฐโทสสิกฺขาปท}
\csromantocmarkref{subsection}{เมตฺติยภูมชกภิกฺขุวตฺถุ}{mtk.82}
\bookmark[dest={mtk.82},level=2]{เมตฺติยภูมชกภิกฺขุวตฺถุ}
\prose{\csromanfolio{255}สุทฺโธ โหติ ปุคฺคโล อญฺญตรํ ปาราชิกํ ธมฺมํ อนชฺฌาปนฺโน, ตญฺเจ สุทฺธทิฏฺฐิ สมาโน อโนกาสํ การาเปตฺวา จาวนาธิปฺปาโย วเทติ, อาปตฺติ สํฆาทิเสเสน ทุกฺกฏสฺส.}

\prose{สุทฺโธ โหติ ปุคฺคโล อญฺญตรํ ปาราชิกํ ธมฺมํ อนชฺฌาปนฺโน, ตญฺเจ สุทฺธทิฏฺฐิ สมาโน โอกาสํ การาเปตฺวา จาวนาธิปฺปาโย วเทติ, อาปตฺติ สํฆาทิเสสสฺส.}

\prose{สุทฺโธ โหติ ปุคฺคโล อญฺญตรํ ปาราชิกํ ธมฺมํ อนชฺฌาปนฺโน, ตญฺเจ สุทฺธทิฏฺฐิ สมาโน อโนกาสํ การาเปตฺวา อกฺโกสาธิปฺปาโย วเทติ, อาปตฺติ โอมสวาเทน ทุกฺกฏสฺส.}

\prose{สุทฺโธ โหติ ปุคฺคโล อญฺญตรํ ปาราชิกํ ธมฺมํ อนชฺฌาปนฺโน, ตญฺเจ สุทฺธทิฏฺฐิ สมาโน โอกาสํ การาเปตฺวา อกฺโกสาธิปฺปาโย วเทติ, อาปตฺติ โอมสวาทสฺส.}

\proseitemwithcloserruled{390}{อนาปตฺติ สุทฺเธ อสุทฺธ\-ทิฏฺฐิสฺส อสุทฺเธ อสุทฺธ\-ทิฏฺฐิสฺส อุมฺมตฺตกสฺส อาทิกมฺมิกสฺสาติ.}{อมูลก\-สิกฺขาปทํ นิฏฺฐิตํ อฏฺฐมํ.}
\csromanheader{1}{9. \textbf{ทุติย}\-\textbf{ทุฏฺฐ}\-\textbf{โทส}\-\textbf{สิกฺขาปท}}
\proseitem{391}{เตน สมเยน พุทฺโธ ภควา ราชคเห วิหรติ เวฬุวเน กลนฺทก\-นิวาเป. เตน โข ปน สมเยน เมตฺติย\-ภูมชกา ภิกฺขู คิชฺฌกูฏา ปพฺพตา โอโรหนฺตา อทฺทสํสุ ฉคลกํ\csromansharedfootnote{211a34502b25db36}{ฉกลกํ (สฺยา)} อชิกาย วิปฺปฏิปชฺชนฺตํ, ทิสฺวาน เอวมาหํสุ “หนฺท มยํ อาวุโส อิมํ ฉคลกํ ทพฺพํ มลฺลปุตฺตํ นาม กโรม, อิมํ อชิกํ เมตฺติยํ นาม ภิกฺขุนิํ กโรม, เอวํ มยํ โวหริสฺสาม ปุพฺเพ มยํ อาวุโส ทพฺพํ มลฺลปุตฺตํ สุเตน อโวจุมฺหา, อิทานิ ปน อเมฺหหิ สามํ ทิฏฺโฐ เมตฺติยาย ภิกฺขุนิยา วิปฺปฏิปชฺชนฺโต”ติ. เต ตํ ฉคลกํ ทพฺพํ มลฺลปุตฺตํ นาม อกํสุ, ตํ อชิกํ เมตฺติยํ นาม ภิกฺขุนิํ อกํสุ. เต ภิกฺขูนํ อาโรเจสุํ “ปุพฺเพ มยํ อาวุโส ทพฺพํ มลฺลปุตฺตํ สุเตน อโวจุมฺหา, อิทานิ ปน อเมฺหหิ สามํ ทิฏฺโฐ เมตฺติยาย ภิกฺขุนิยา วิปฺปฏิปชฺชนฺโต”ติ. ภิกฺขู เอวมาหํสุ “มาวุโส เอวํ อวจุตฺถ, นายสฺมา ทพฺโพ มลฺลปุตฺโต เอวํ กริสฺสตี”ติ.}

\csromanmatikaanchor{mtk.83}
\csromantocmarkref{subsection}{ปญฺญตฺติ}{mtk.83}
\bookmark[dest={mtk.83},level=2]{ปญฺญตฺติ}
\prose{\csromanfolio{256}อถ โข เต ภิกฺขู ภควโต เอตมตฺถํ อาโรเจสุํ. อถ โข ภควา เอตสฺมิํ นิทาเน เอตสฺมิํ ปกรเณ ภิกฺขุสํฆํ สนฺนิปาตาเปตฺวา อายสฺมนฺตํ ทพฺพํ มลฺลปุตฺตํ ปฏิปุจฺฉิ “สรสิ ตฺวํ ทพฺพ เอวรูปํ กตฺตา ยถยิเม ภิกฺขู อาหํสู”ติ. ยถา มํ ภนฺเต ภควา ชานาตีติ. ทุติยมฺปิ โข ภควา ฯเปฯ. ตติยมฺปิ โข ภควา อายสฺมนฺตํ ทพฺพํ มลฺลปุตฺตํ เอตทโวจ “สรสิ ตฺวํ ทพฺพ เอวรูปํ กตฺตา ยถยิเม ภิกฺขู อาหํสู”ติ. ยถา มํ ภนฺเต ภควา ชานาตีติ. น โข ทพฺพ ทพฺพา เอวํ นิพฺเพเฐนฺติ, สเจ ตยา กตํ กตนฺติ วเทหิ, สเจ ตยา อกตํ อกตนฺติ วเทหีติ. ยโต อหํ ภนฺเต ชาโต นาภิชานามิ สุปินนฺเตนปิ เมถุนธมฺมํ ปฏิเสวิตา, ปเคว ชาคโรติ. อถ โข ภควา ภิกฺขู อามนฺเตสิ “เตน หิ ภิกฺขเว อิเม ภิกฺขู อนุยุญฺชถา”ติ. อิทํ วตฺวา ภควา อุฏฺฐายาสนา วิหารํ ปาวิสิ.}

\prose{อถ โข เต ภิกฺขู เมตฺติย\-ภูมชเก ภิกฺขู อนุยุญฺชิํสุ. เต ภิกฺขูหิ อนุยุญฺชียมานา ภิกฺขูนํ เอตมตฺถํ อาโรเจสุํ. กิํ ปน ตุเมฺห อาวุโส อายสฺมนฺตํ ทพฺพํ มลฺลปุตฺตํ อญฺญภาคิยสฺส อธิกรณสฺส กิญฺจิเทสํ เลสมตฺตํ อุปาทาย ปาราชิเกน ธมฺเมน อนุทฺธํเส\-ถาติ. เอวมาวุโสติ. เย เต ภิกฺขู อปฺปิจฺฉา ฯเปฯ เต อุชฺฌายนฺติ ขิยฺยนฺติ วิปาเจนฺติ “กถํ หิ นาม เมตฺติย\-ภูมชกา ภิกฺขู อายสฺมนฺตํ ทพฺพํ มลฺลปุตฺตํ อญฺญภาคิยสฺส อธิกรณสฺส กิญฺจิ เทสํ เลสมตฺตํ อุปาทาย ปาราชิเกน ธมฺเมน อนุทฺธํเสสฺสนฺตี”ติ. อถ โข เต ภิกฺขู เมตฺติย\-ภูมชเก ภิกฺขู อเนก\-ปริยาเยน วิครหิตฺวา ภควโต เอตมตฺถํ อาโรเจสุํ ฯเปฯ “สจฺจํ กิร ตุเมฺห ภิกฺขเว ทพฺพํ มลฺลปุตํ อญฺญภาคิยสฺส อธิกรณสฺส กิญฺจิ เทสํ เลสมตฺตํ อุปาทาย ปาราชิเกน ธมฺเมน อนุทฺธํเสถา”ติ. สจฺจํ ภควาติ. วิครหิ พุทฺโธ ภควา ฯเปฯ. กถํ หิ นาม ตุเมฺห โมฆปุริสา ทพฺพํ มลฺลปุตฺตํ อญฺญภาคิยสฺส อธิกรณสฺส กิญฺจิ เทสํ เลสมตฺตํ อุปาทาย ปาราชิเกน ธมฺเมน อนุทฺธํเสสฺสถ, เนตํ โมฆปุริสา อปฺปสนฺนานํ วา ปสาทาย ฯเปฯ. เอวญฺจ ปน ภิกฺขเว อิมํ สิกฺขาปทํ อุทฺทิเสยฺยาถ\csromandash{}}

\csromanmatikaanchor{mtk.84}
\csromantocmarkref{subsection}{สิกฺขาปทวิภงฺค, ปทภาชนีย}{mtk.84}
\bookmark[dest={mtk.84},level=2]{สิกฺขาปทวิภงฺค, ปทภาชนีย}
\proseitem{392}{\textbf{“โย ปน ภิกฺขุ ภิกฺขุํ ทุฏฺโฐ โทโส อปฺปตีโต อญฺญภาคิยสฺส อธิกรณสฺส กิญฺจิเทสํ เลสมตฺตํ อุปาทาย ปาราชิเกน ธมฺเมน อนุทฺธํเสยฺย ‘อปฺเปว นาม นํ อิมมฺหา พฺรหฺมจริยา จาเวยฺยนฺ’ติ,} \csromanfolio{257}\textbf{ตโต อปเรน สมเยน สมนุ}\-\textbf{คฺคาหิย}\-\textbf{มาโน วา อสมนุคฺคาหียมาโน วา อญฺญภาคิยญฺเจว ตํ อธิกรณํ โหติ โกจิเทโส เลสมตฺโต อุปาทินฺโน, ภิกฺขุ จ โทสํ ปติฏฺฐาติ, สํฆาทิเสโส”}ติ.}

\proseitem{393}{\textbf{โย ปนา}ติ โย ยาทิโส ฯเปฯ. พฺ\textbf{หิกฺขู}ติ ฯเปฯ อยํ อิมสฺมิํ อตฺเถ อธิปฺเปโต “ภิกฺขู”ติ.}

\prose{\textbf{ภิกฺขุนฺ}ติ อญฺญํ ภิกฺขุํ.}

\prose{\textbf{ทุฏฺโฐ โทโส}ติ กุปิโต อนตฺตมโน อนภิรทฺโธ อาหตจิตฺโต ขิลชาโต.}

\prose{\textbf{อปฺปตีโต}ติ เตน จ โกเปน เตน จ โทเสน ตาย จ อนตฺตมนตาย ตาย จ อนภิรทฺธิยา อปฺปตีโต โหติ.}

\prose{\textbf{อญฺญภาคิยสฺส อธิกรณสฺสา}ติ อาปตฺต\-ญฺญภาคิยํ วา โหติ อธิกรณ\-ญฺญภาคิยํ วา. กถํ อธิกรณํ อธิกรณสฺส อญฺญภาคิยํ. วิวาทา\-ธิกรณํ อนุวาทา\-ธิกรณสฺส อาปตฺตา\-ธิกรณสฺส กิจฺจา\-ธิกรณสฺส อญฺญภาคิยํ. อนุวาทา\-ธิกรณํ อาปตฺตา\-ธิกรณสฺส กิจฺจา\-ธิกรณสฺส วิวาทา\-ธิกรณสฺส อญฺญภาคิยํ. อาปตฺตา\-ธิกรณํ กิจฺจา\-ธิกรณสฺส วิวาทา\-ธิกรณสฺส อนุวาทา\-ธิกรณสฺส อญฺญภาคิยํ. กิจฺจา\-ธิกรณํ วิวาทา\-ธิกรณสฺส อนุวาทา\-ธิกรณสฺส อาปตฺตา\-ธิกรณสฺส อญฺญภาคิยํ. เอวํ อธิกรณํ อธิกรณสฺส อญฺญภาคิยํ.}

\prose{กถํ อธิกรณํ อธิกรณสฺส ตพฺภาคิยํ. วิวาทา\-ธิกรณํ วิวาทา\-ธิกรณสฺส ตพฺภาคิยํ. อนุวาทา\-ธิกรณํ อนุวาทา\-ธิกรณสฺส ตพฺภาคิยํ. อาปตฺตา\-ธิกรณํ อาปตฺติธิกรณสฺส สิยา ตพฺภาคิยํ. สิยา อญฺญภาคิยํ.}

\prose{กถํ อาปตฺตา\-ธิกรณํ อาปตฺตา\-ธิกรณสฺส อญฺญภาคิยํ. เมถุนธมฺม\-ปาราชิกาปตฺติ อทินฺนาทาน\-ปาราชิกาปตฺติยา มนุสฺสวิคฺคห\-ปาราชิกาปตฺติยา อุตฺตริมนุสฺสธมฺม\-ปาราชิกาปตฺติยา อญฺญภาคิยา. อทินฺนาทาน\-ปาราชิกาปตฺติ มนุสฺสวิคฺคห\-ปาราชิกาปตฺติยา อุตฺตริมนุสฺสธมฺม\-ปาราชิกาปตฺติยา เมถุนธมฺม\-ปาราชิกาปตฺติยา อญฺญภาคิยา. มนุสฺสวิคฺคห\-ปาราชิกาปตฺติ อุตฺตริมนุสฺสธมฺม\-ปาราชิกาปตฺติยา เมถุนธมฺม\-ปาราชิกาปตฺติยา \csromanfolio{258}อทินฺนาทาน\-ปาราชิกาปตฺติยา อญฺญภาคิยา. อุตฺตริมนุสฺสธมฺม\-ปาราชิกาปตฺติ เมถุนธมฺม\-ปาราชิกาปตฺติยา อทินฺนาทานปาราชิกปตฺติยา มนุสฺสวิคฺคห\-ปาราชิกาปตฺติยา อญฺญภาคิยา. เอวํ อาปตฺตา\-ธิกรณํ อาปตฺตา\-ธิกรณสฺส อญฺญภาคิยํ.}

\prose{กถํ อาปตฺตธิกรณํ อาปตฺตา\-ธิกรณสฺส ตพฺภาคิยํ. เมถุนธมฺม\-ปาราชิกาปตฺติ เมถุนธมฺม\-ปาราชิกาปตฺติยา ตพฺภาคิยา. อทินฺนาทาน\-ปาราชิกาปตฺติ อทินฺนาทาน\-ปาราชิกาปตฺติยา ตพฺภาคิยา. มนุสฺสวิคฺคห\-ปาราชิกาปตฺติ มนุสฺสวิคฺคห\-ปาราชิกาปตฺติยา ตพฺภาคิยา. อุตฺตริมนุสฺสธมฺม\-ปาราชิกาปตฺติ อุตฺตริมนุสฺสธมฺม\-ปาราชิกาปตฺติยา ตพฺภาคิยา. เอวํ อาปตฺตา\-ธิกรณํ อาปตฺตา\-ธิกรณสฺส ตพฺภาคิยํ.}

\prose{กิจฺจา\-ธิกรณํ กิจฺจา\-ธิกรณสฺส ตพฺภาคิยํ. เอวํ อธิกรณํ อธิกรณสฺส ตพฺภาคิยํ.}

\proseitem{394}{\textbf{กิญฺจิ เทสํ เลสมตฺตํ อุปาทายา}ติ เลโส นาม ทส เลสา ชาติเลโส นามเลโส โคตฺตเลโส ลิงฺคเลโส อาปตฺติเลโส ปตฺตเลโส จีวรเลโส อุปชฺฌายเลโส อาจริยเลโส เสนาสนเลโส.}

\proseitem{395}{\textbf{ชาติเลโส} นาม ขตฺติโย ทิฏฺโฐ โหติ ปาราชิกํ ธมฺมํ อชฺฌาปชฺชนฺโต, อญฺญํ ขตฺติยํ ปสฺสิตฺวา โจเทติ “ขตฺติโย มยา ทิฏฺโฐ\csromansharedfootnote{1579eddaaba8a580}{เอตฺถ “ปาราชิกํ ธมฺมํ อชฺฌาปชฺชนฺโต”ติ ปาโฐ อูโน มญฺเญ, อฏฺฐกถายํ หิ “ขตฺติโย มยา ทิฏฺโฐ ปาราชิกํ ธมฺมํ อชฺฌาปชฺชนฺโต”ติ ทิสฺสติ.}, ปาราชิกํ ธมฺมํ อชฺฌาปนฺโนสิ อสฺสมโณสิ อสกฺยปุตฺติโยสิ, นตฺถิ ตยา สทฺธิํ อุโปสโถ วา ปวารณา วา สํฆกมฺมํ วา”ติ, อาปตฺติ วาจาย วาจาย สํฆาทิเสสสฺส.}

\prose{พฺราหฺมโณ ทิฏฺโฐ โหติ ฯเปฯ. เวสฺโส ทิฏฺโฐ โหติ ฯเปฯ. สุทฺโท ทิฏฺโฐ โหติ ปาราชิกํ ธมฺมํ อชฺฌาปชฺชนฺโต, อญฺญํ สุทฺทํ ปสฺสิตฺวา โจเทติ “สุทฺโท มยา ทิฏฺโฐ, ปาราชิกํ ธมฺมํ อชฺฌาปนฺโนสิ อสฺสมโณสิ อสกฺยปุตฺติโยสิ ฯเปฯ อาปตฺติ วาจาย วาจาย สํฆาทิเสสสฺส.}

\proseitem{396}{\textbf{นามเลโส} นาม พุทฺธรกฺขิโต ทิฏฺโฐ โหติ ฯเปฯ. ธมฺมรกฺขิโต ทิฏฺโฐ โหติ ฯเปฯ. สํฆรกฺขิโต ทิฏฺโฐ โหติ ปาราชิกํ ธมฺมํ \csromanfolio{259}อชฺฌาปชฺชนฺโต, อญฺญํ สํฆรกฺขิตํ ปสฺสิตฺวา โจเทติ “สํฆรกฺขิโต มยา ทิฏฺโฐ, ปาราชิกํ ธมฺมํ อชฺฌาปนฺโนสิ อสฺสมโณสิ อสกฺยปุตฺติโยสิ ฯเปฯ อาปตฺติ วาจาย วาจาย สํฆาทิเสสสฺส.}

\proseitem{397}{\textbf{โคตฺตเลโส} นาม โคตโม ทิฏฺโฐ โหติ ฯเปฯ. โมคฺคลฺลาโน ทิฏฺโฐ โหติ ฯเปฯ. กจฺจายโน ทิฏฺโฐ โหติ ฯเปฯ. วาสิฏฺโฐ ทิฏฺโฐ โหติ ปาราชิกํ ธมฺมํ อชฺฌาปชฺชนฺโต, อญฺญํ วาสิฏฺฐํ ปสฺสิตฺวา โจเทติ “วาสิฏฺโฐ มยา ทิฏฺโฐ, ปาราชิกํ ธมฺมํ อชฺฌาปนฺโนสิ อสฺสมโณสิ อสกฺยปุตฺติโยสิ ฯเปฯ อาปตฺติ วาจาย วาจาย สํฆาทิเสสสฺส.}

\proseitem{398}{\textbf{ลิงฺคเลโส} นาม ทีโฆ ทิฏฺโฐ โหติ ฯเปฯ. รสฺโส ทิฏฺโฐ โหติ ฯเปฯ .กโณฺห ทิฏฺโฐ โหติ ฯเปฯ. โอทาโต ทิฏฺโฐ โหติ ปาราชิกํ ธมฺมํ อชฺฌาปชฺชนฺโต, อญฺญํ โอทาตํ ปสฺสิตฺวา โจเทติ “โอทาโต มยา ทิฏฺโฐ ปาราชิกํ ธมฺมํ อชฺฌาปนฺโนสิ อสฺสมโณสิ อสกฺยปุตฺติโยสิ ฯเปฯ อาปตฺติ วาจาย วาจาย สํฆาทิเสสสฺส.}

\proseitem{399}{\textbf{อาปตฺติเลโส} นาม ลหุกํ อาปตฺติํ อาปชฺชนฺโต ทิฏฺโฐ โหติ. ตญฺเจ ปาราชิเกน โจเทติ “อสฺสมโณสิ อสกฺยปุตฺติโยสิ ฯเปฯ อาปตฺติ วาจาย วาจาย สํฆาทิเสสสฺส.}

\proseitem{400}{\textbf{ปตฺตเลโส} นาม โลหปตฺตธโร ทิฏฺโฐ โหติ ฯเปฯ. สาฏก\-ปตฺต\-ธโร ทิฏฺโฐ โหติ ฯเปฯ. สุมฺภก\-ปตฺต\-ธโร ทิฏฺโฐ โหติ ปาราชิกํ ธมฺมํ อชฺฌาปชฺชนฺโต, อญฺญํ สุมฺภก\-ปตฺตธรํ ปสฺสิตฺวา โจเทติ “สุมฺภก\-ปตฺต\-ธโร มยา ทิฏฺโฐ, ปาราชิกํ ธมฺมํ อชฺฌาปนฺโนสิ อสฺสมโณสิ อสกฺยปุตฺติโยสิ ฯเปฯ อาปตฺติ วาจาย วาจาย สํฆาทิเสสสฺส.}

\proseitem{401}{\textbf{จีวรเลโส} นาม ปํสุกูลิโก ทิฏฺโฐ โหติ ฯเปฯ. คหปติจีวร\-ธโร ทิฏฺโฐ โหติ ปาราชิกํ ธมฺมํ อชฺฌาปชฺชนฺโต, อญฺญํ คหปติจีวรธรํ ปสฺสิตฺวา โจเทติ “คหปติจีวร\-ธโร มยา ทิฏฺโฐ, ปาราชิกํ ธมฺมํ อชฺฌาปนฺโนสิ อสฺสมโณสิ อสกฺยปุตฺติโยสิ ฯเปฯ อาปตฺติ วาจาย วาจาย สํฆาทิเสสสฺส.}

\proseitem{402}{\csromanfolio{260}\textbf{อุปชฺฌายเลโส} นาม อิตฺถนฺนามสฺส สทฺธิวิหาริโก ทิฏฺโฐ โหติ ปาราชิกํ ธมฺมํ อชฺฌาปชฺชนฺโต, อญฺญํ อิตฺถนฺนามสฺส สทฺธิ\-วิหาริกํ ปสฺสิตฺวา โจเทติ “อิตฺถนฺนามสฺส สทฺธิวิหาริโก มยา ทิฏฺโฐ, ปาราชิกํ ธมฺมํ อชฺฌาปนฺโนสิ อสฺสมโณสิ อสกฺยปุตฺติโยสิ ฯเปฯ อาปตฺติ วาจาย วาจาย สํฆาทิเสสสฺส.}

\proseitem{403}{\textbf{อาจริยเลโส} นาม อิตฺถนฺนามสฺส อนฺเตวาสิโก ทิฏฺโฐ โหติ ปาราชิกํ ธมฺมํ อชฺฌาปชฺชนฺโต, อญฺญํ อิตฺถนฺนามสฺส อนฺเตวาสิกํ ปสฺสิตฺวา โจเทติ “อิตฺถนฺนามสฺส อนฺเตวาสิโก มยา ทิฏฺโฐ ปาราชิกํ ธมฺมํ อชฺฌาปนฺโนสิ อสฺสมโณสิ อสกฺยปุตฺติโยสิ ฯเปฯ อาปตฺติ วาจาย วาจาย สํฆาทิเสสสฺส.}

\proseitem{404}{\textbf{เสนาสนเลโส} นาม อิตฺถนฺนาม\-เสนาสน\-วาสิโก ทิฏฺโฐ โหติ ปาราชิกํ ธมฺมํ อชฺฌาปชฺชนฺโต, อญฺญํ อิตฺถนฺนาม\-เสนาสน\-วาสิกํ ปสฺสิตฺวา โจเทติ “อิตฺถนฺนาม\-เสนาสน\-วาสิโก มยา ทิฏฺโฐ, ปาราชิกํ ธมฺมํ อชฺฌาปนฺโนสิ อสฺสมโณสิ อสกฺยปุตฺติโยสิ, นตฺถิ ตยา สทฺธิํ อุโปสโถ วา ปวารณา วา สํฆกมฺมํ วา”ติ, อาปตฺติ วาจาย วาจาย สํฆาทิเสสสฺส.}

\proseitem{405}{\textbf{ปาราชิเกน ธมฺเมนา}ติ จตุนฺนํ อญฺญตเรน.}

\prose{\textbf{อนุทฺธํเสยฺยา}ติ โจเทติ วา โจทาเปติ วา.}

\prose{\textbf{อปฺเปว นาม นํ อิมมฺหา พฺรหฺมจริยา จาเวยฺยนฺ}ติ ภิกฺขุภาวา จาเวยฺยํ, สมณธมฺมา จาเวยฺยํ, สีลกฺขนฺธา จาเวยฺยํ, ตโปคุณา จาเวยฺยํ.}

\prose{\textbf{ตโต อปเรน สมเยนา}ติ ยสฺมิํ ขเณ อนุทฺธํสิโต โหติ, ตํ ขณํ ตํ ลยํ ตํ มุหุตฺตํ วีติวตฺเต.}

\prose{\textbf{สมนุ}\-\textbf{คฺคาหีย}\-\textbf{มาโน}ติ เยน วตฺถุนา อนุทฺธํสิโต โหติ, ตสฺมิํ วตฺถุสฺมิํ สมนุ\-คฺคาหีย\-มาโน.}

\prose{\textbf{อสมนุคฺคาหียมาโน}ติ น เกนจิ วุจฺจมาโน.}

\prose{\textbf{อธิกรณํ} นาม จตฺตาริ อธิกรณานิ วิวาทา\-ธิกรณํ อนุวาทา\-ธิกรณํ อาปตฺตา\-ธิกรณํ กิจฺจา\-ธิกรณํ.}

\prose{\csromanfolio{261}\textbf{โกจิ เทโส เลสมตฺโต อุปาทินฺโน}ติ เตสํ เลสานํ อญฺญตโร เลโส อุปาทินฺโน โหติ.}

\prose{\textbf{ภิกฺขุ จ โทสํ ปติฏฺฐาตี}ติ ตุจฺฉกํ มยา ภณิตํ, มุสา มยา ภณิตํ, อภูตํ มยา ภณิตํ, อชานนฺเตน มยา ภณิตํ.}

\prose{\textbf{สํฆาทิเสโส}ติ ฯเปฯ เตนปิ วุจฺจติ “สํฆาทิเสโส”ติ.}

\proseitem{406}{ภิกฺขุ สํฆาทิเสสํ อชฺฌาปชฺชนฺโต ทิฏฺโฐ โหติ, สํฆาทิเสเส สํฆาทิเสสทิฏฺฐิ โหติ, ตญฺเจ ปาราชิเกน โจเทติ “อสฺสมโณสิ อสกฺยปุตฺติโยสิ, นตฺถิ ตยา สทฺธิํ อุโปสโถ วา ปวารณา วา สํฆกมฺมํ วา”ติ, เอวํปิ อาปตฺต\-ญฺญภาคิยํ โหติ เลโส จ อุปาทินฺโน, อาปตฺติ วาจาย วาจาย สํฆาทิเสสสฺส.}

\prose{ภิกฺขุ สํฆาทิเสสํ อชฺฌาปชฺชนฺโต ทิฏฺโฐ โหติ, สํฆาทิเสเส ถุลฺลจฺจย\-ทิฏฺฐิ โหติ ฯเปฯ ปาจิตฺติย\-ทิฏฺฐิ โหติ. ปาฏิเทสนีย\-ทิฏฺฐิ โหติ. ทุกฺกฏทิฏฺฐิ โหติ. ทุพฺภาสิต\-ทิฏฺฐิ โหติ, ตญฺเจ ปาราชิเกน โจเทติ “อสฺสมโณสิ ฯเปฯ เอวํปิ อาปตฺต\-ญฺญภาคิยํ โหติ เลโส จ อุปาทินฺโน, อาปตฺติ วาจาย วาจาย สํฆาทิเสสสฺส.}

\prose{ภิกฺขุ ถุลฺลจฺจยํ อชฺฌาปชฺชนฺโต ทิฏฺโฐ โหติ, ถุลฺลจฺจเย ถุลฺลจฺจย\-ทิฏฺฐิ โหติ ฯเปฯ ถุลฺลจฺจเย ปาจิตฺติย\-ทิฏฺฐิ โหติ. ปาฏิเทสนีย\-ทิฏฺฐิ โหติ. ทุกฺกฏทิฏฺฐิ โหติ. ทุพฺภาสิต\-ทิฏฺฐิ โหติ. สํฆาทิเสสทิฏฺฐิ โหติ, ตญฺเจ ปาราชิเกน โจเทติ “อสฺสมโณสิ ฯเปฯ เอวํปิ อาปตฺต\-ญฺญภาคิยํ โหติ เลโส จ อุปาทินฺโน. อาปตฺติ วาจาย วาจาย สํฆาทิเสสสฺส.}

\prose{ภิกฺขุ ปาจิตฺติยํ ฯเปฯ ปาฏิเทสนียํ. ทุกฺกฏํ. ทุพฺภาสิตํ อชฺฌาปชฺชนฺโต ทิฏฺโฐ โหติ, ทุพฺภาสิเต ทุพฺภาสิต\-ทิฏฺฐิ โหติ ฯเปฯ ทุพฺภาสิเต สํฆาทิเสสทิฏฺฐิ โหติ. ถุลฺลจฺจย\-ทิฏฺฐิ โหติ. ปาจิตฺติย\-ทิฏฺฐิ โหติ. ปาฏิเทสนีย\-ทิฏฺฐิ โหติ. ทุกฺกฏทิฏฺฐิ โหติ, ตญฺเจ ปาราชิเกน โจเทติ “อสฺสมโณสิ อสกฺยปุตฺติโยสิ, นตฺถิ ตยา สทฺธิํ อุโปสโถ วา ปวารณา วา สํฆกมฺมํ วา”ติ, เอวํปิ อาปตฺต\-ญฺญภาคิยํ โหติ เลโส จ อุปาทินฺโน, อาปตฺติ วาจาย วาจาย สํฆาทิเสสสฺส.}

\vspace{\tipitakaparskip}
\csromancenter{เอเกกํ มูลํ กาตุน จกฺกํ พนฺธิตพฺพํ.}
\proseitem{407}{\csromanfolio{262}ภิกฺขุ สํฆาทิเสสํ อชฺฌาปชฺชนฺโต ทิฏฺโฐ โหติ, สํฆาทิเสเส สํฆาทิเสสทิฏฺฐิ โหติ, ตญฺเจ ปาราชิเกน โจทาเปติ “อสฺสมโณสิ ฯเปฯ เอวํปิ อาปตฺต\-ญฺญภาคิยํ โหติ เลโส จ อุปาทินฺโน, อาปตฺติ วาจาย วาจาย สํฆาทิเสสสฺส.}

\prose{ภิกฺขุ สํฆาทิเสสํ อชฺฌาปชฺชนฺโต ทิฏฺโฐ โหติ, สํฆาทิเสเส ถุลฺลจฺจย\-ทิฏฺฐิ โหติ ฯเปฯ ปาจิตฺติย\-ทิฏฺฐิ โหติ. ปาฏิเทสนีย\-ทิฏฺฐิ โหติ. ทุกฺกฏทิฏฺฐิ โหติ. ทุพฺภาสิต\-ทิฏฺฐิ โหติ, ตญฺเจ ปาราชิเกน โจทาเปติ “อสฺสมโณสิ ฯเปฯ เอวํปิ อาปตฺต\-ญฺญภาคิยํ โหติ เลโส จ อุปาทินฺโน, อาปตฺติ วาจาย วาจาย สํฆาทิเสสสฺส.}

\prose{ภิกฺขุ ถุลฺลจฺจยํ อชฺฌาปชฺชนฺโต ทิฏฺโฐ โหติ, ถุลฺลจฺจเย ถุลฺลจฺจย\-ทิฏฺฐิ โหติ ฯเปฯ ถุลฺลจฺจเย ปาจิตฺติย\-ทิฏฺฐิ โหติ. ปาฏิเทสนีย\-ทิฏฺฐิ โหติ. ทุกฺกฏทิฏฺฐิ โหติ. ทุพฺภาสิต\-ทิฏฺฐิ โหติ. สํฆาทิเสสทิฏฺฐิ โหติ, ตญฺเจ ปาราชิเกน โจทาเปติ “อสฺสมโณสิ ฯเปฯ เอวํปิ อาปตฺต\-ญฺญภาคิยํ โหติ เลโส จ อุปาทินฺโน, อาปตฺติ วาจาย วาจาย สํฆาทิเสสสฺส.}

\prose{ภิกฺขุ ปาจิตฺติยํ ฯเปฯ ปาฏิเทสนียํ. ทุกฺกฏํ. ทุพฺภาสิตํ อชฺฌาปชฺชนฺโต ทิฏฺโฐ โหติ, ทุพฺภาสิเต ทุพฺภาสิต\-ทิฏฺฐิ โหติ ฯเปฯ ทุพฺภาสิเต สํฆาทิเสสทิฏฺฐิ โหติ. ถุลฺลจฺจย\-ทิฏฺฐิ โหติ. ปาจิตฺติย\-ทิฏฺฐิ โหติ. ปาฏิเทสนีย\-ทิฏฺฐิ โหติ. ทุกฺกฏทิฏฺฐิ โหติ, ตญฺเจ ปาราชิเกน โจทาเปติ “อสฺสมโณสิ อสกฺยปุตฺติโยสิ, นตฺถิ ตยา สทฺธิํ อุโปสโถ วา ปวารณา วา สํฆกมฺมํ วา”ติ, เอวํปิ อาปตฺต\-ญฺญภาคิยํ โหติ เลโส จ อุปาทินฺโน, อาปตฺติ วาจาย วาจาย สํฆาทิเสสสฺส.}

\proseitem{408}{อนาปตฺติ ตถาสญฺญี โจเทติ วา โจทาเปติ วา อุมฺมตฺตกสฺส อาทิกมฺมิกสฺสาติ. (อญฺญาภาคิย) กิญฺจิ\-เลส\-สิกฺขาปทํ นิฏฺฐิตํ นวมํ.}

\csromanmatikaanchor{mtk.85}
\csromantocmarknopage{section}{10. สํฆเภทสิกฺขาปท}{mtk.85}
\bookmark[dest={mtk.85},level=1]{10. สํฆเภทสิกฺขาปท}
\csromanheader{1}{10. \textbf{สํฆเภทสิกฺขาปท}}
\csromanmatikaanchor{mtk.86}
\csromantocmarkref{subsection}{เทวทตฺตภิกฺขุวตฺถุ}{mtk.86}
\bookmark[dest={mtk.86},level=2]{เทวทตฺตภิกฺขุวตฺถุ}
\proseitem{409}{\csromanfolio{263}เตน สมเยน พุทฺโธ ภควา ราชคเห วิหรติ เวฬุวเน กลนฺทก\-นิวาเป. อถ โข เทวทตฺโต เยน โกกาลิโก กฏโมทกติสฺสโก\csromansharedfootnote{cdf23afe6ca59bf2}{กฏโมรกติสฺสโก (สี, สฺยา)} ขณฺฑเทวิยา ปุตฺโต สมุทฺททตฺโต เตนุปสงฺกมิ, อุปสงฺกมิตฺวา โกกาลิกํ กฏโมทกติสฺสกํ ขณฺฑเทวิยา ปุตฺตํ สมุทฺททตฺตํ เอตทโวจ “เอถ มยํ อาวุโส สมณสฺส โคตมสฺส สํฆเภทํ กริสฺสาม จกฺกเภทนฺ”ติ. เอวํ วุตฺเต โกกาลิโก เทวทตฺตํ เอตทโวจ “สมโณ โข อาวุโส โคตโม มหิทฺธิโก มหานุภาโว, กถํ มยํ สมณสฺส โคตมสฺส สํฆเภทํ กริสฺสาม จกฺกเภทนฺ”ติ. เอถ มยํ อาวุโส สมณํ โคตมํ อุปสงฺกมิตฺวา ปญฺจ วตฺถูนิ ยาจิสฺสาม “ภควา ภนฺเต อเนก\-ปริยาเยน อปฺปิจฺฉสฺส สนฺตุฏฺฐสฺส สลฺเลขสฺส ธุตสฺส ปาสาทิกสฺส อปจยสฺส วีริยารพฺภสฺส\csromansharedfootnote{9267ffdd082a3e79}{วิริยารมฺภสฺส (สี, สฺยา)} วณฺณวาที, อิมานิ ภนฺเต ปญฺจ วตฺถูนิ อเนก\-ปริยาเยน อปฺปิจฺฉตาย สนฺตุฏฺฐิยา สลฺเลขาย ธุตตาย ปาสาทิกตาย อปจยาย วีริยารพฺภาย สํวตฺตนฺติ, สาธุ ภนฺเต ภิกฺขู ยาวชีวํ อารญฺญิกา อสฺสุ, โย คามนฺตํ โอสเรยฺย, วชฺชํ นํ ผุเสยฺย. ยาวชีวํ ปิณฺฑปาติกา อสฺสุ, โย นิมนฺตนํ สาทิเยยฺย, วชฺชํ นํ ผุเสยฺย. ยาวชีวํ ปํสุกูลิกา อสฺสุ, โย คหปติ\-จีวรํ สาทิเยยฺย, วชฺชํ นํ ผุเสยฺย. ยาวชีวํ รุกฺขมูลิกา อสฺสุ, โย ฉนฺนํ อุปคจฺเฉยฺย, วชฺชํ นํ ผุเสยฺย. ยาวชีวํ มจฺฉมํสํ น ขาเทยฺยุํ, โย มจฺฉมํสํ ขาเทยฺย, วชฺชํ นํ ผุเสยฺยา”ติ. อิมานิ สมโณ โคตโม นานุชานิสฺสติ, เต มยํ อิเมหิ ปญฺจหิ วตฺถูหิ ชนํ สญฺญาเปสฺสามาติ. สกฺกา โข อาวุโส อิเมหิ ปญฺจหิ วตฺถูหิ สมณสฺส โคตมสฺส สํฆเภโท กาตุํ จกฺกเภโท, ลูขปสนฺนา หิ อาวุโส มนุสฺสาติ.}

\prose{อถ โข เทวทตฺโต สปริโส เยน ภควา เตนุปสงฺกมิ, อุปสงฺกมิตฺวา ภควนฺตํ อภิวาเทตฺวา เอกมนฺตํ นิสีทิ, เอกมนฺตํ นิสินฺโน โข เทวทตฺโต ภควนฺตํ เอตทโวจ “ภควา ภนฺเต อเนก\-ปริยาเยน อปฺปิจฺฉสฺส ฯเปฯ วีริยารพฺภสฺส วณฺณวาที, อิมานิ ภนฺเต ปญฺจ วตฺถูนิ อเนก\-ปริยาเยน อปฺปิจฺฉตาย ฯเปฯ วีริยารพฺภาย สํวตฺตนฺติ, สาธุ ภนฺเต \csromanfolio{264}ภิกฺขู ยาวชีวํ อารญฺญิกา อสฺสุ, โย คามนฺตํ โอสเรยฺย, วชฺชํ นํ ผุเสยฺย ฯเปฯ. ยาวชีวํ มจฺฉมํสํ น ขาเทยฺยุํ, โย มจฺฉมํสํ ขาเทยฺย, วชฺชํ นํ ผุเสยฺยา”ติ. อลํ เทวทตฺต โย อิจฺฉติ อารญฺญิโก โหตุ, โย อิจฺฉติ คามนฺเต วิหรตุ, โย อิจฺฉติ ปิณฺฑปาติโก โหตุ, โย อิจฺฉติ นิมนฺตนํ สาทิยตุ, โย อิจฺฉติ ปํสุกูลิโก โหตุ, โย อิจฺฉติ คหปติ\-จีวรํ สาทิยตุ, อฏฺฐมาเส โข มยา เทวทตฺต รุกฺขมูล\-เสนาสนํ อนุญฺญาตํ, ติโกฏิ\-ปริสุทฺธํ มจฺฉมํสํ อทิฏฺฐํ อสุตํ อปริสงฺกิตนฺติ. อถ โข เทวทตฺโต “น ภควา อิมานิ \textbf{ปญฺจ วตฺถูนิ อนุชานาตี”ติ หฏฺโฐ อุทคฺโค สปริโส อุฏฺฐายาสนา ภควนฺตํ} \textbf{อภิวาเทตฺวา ปทกฺขิณํ กตฺวา ปกฺกามิ.}}

\csromanmatikaanchor{mtk.87}
\csromanmatikaanchor{mtk.88}
\csromantocmarkref{subsection}{ปญฺญตฺติ}{mtk.87}
\bookmark[dest={mtk.87},level=2]{ปญฺญตฺติ}
\csromantocmarkref{subsection}{สิกฺขาปทวิภงฺค, ปทภาชนีย}{mtk.88}
\bookmark[dest={mtk.88},level=2]{สิกฺขาปทวิภงฺค, ปทภาชนีย}
\proseitem{410}{อถ โข เทวทตฺโต สปริโส ราชคหํ ปวิสิตฺวา ปญฺจหิ วตฺถูหิ ชนํ สญฺญาเปสิ “มยํ อาวุโส สมณํ โคตมํ อุปสงฺกมิตฺวา ปญฺจ วตฺถูนิ ยาจิมฺหา ‘ภควา ภนฺเต อเนก\-ปริยาเยน อปฺปิจฺฉสฺส สนฺตุฏฺฐสฺส สลฺเลขสฺส ธุตสฺส ปาสาทิกสฺส อปจยสฺส วีริยา\-รมฺภสฺส วณฺณวาที, อิมานิ ภนฺเต ปญฺจ วตฺถูนิ อเนก\-ปริยาเยน อปฺปิจฺฉตาย สนฺตุฏฺฐิยา สลฺเลขาย ธุตตาย ปาสาทิกตาย อปจยาย วีริยารมฺภาย สํวตฺตนฺติ, สาธุ ภนฺเต ภิกฺขู ยาวชีวํ อารญฺญิกา อสฺสุ, โย คามนฺตํ โอสเรยฺย, วชฺชํ นํ ผุเสยฺย. ยาวชีวํ ปิณฺฑปาติกา อสฺสุ, โย นิมนฺตนํ สาทิเยยฺย, วชฺชํ นํ ผุเสยฺย. ยาวชีวํ ปํสุกูลิยา อสฺสุ, โย คหปติ\-จีวรํ สาทิเยยฺย, วชฺชํ นํ ผุเสยฺย. ยาวชีวํ รุกฺขมูลิกา อสฺสุ, โย ฉนฺนํ อุปคจฺเฉยฺย, วชฺชํ นํ ผุเสยฺย. ยาวชีวํ มจฺฉมํสํ น ขาเทยฺยุํ, โย มจฺฉมํสํ ขาเทยฺย, วชฺชํ นํ ผุเสยฺยา’ติ, อิมานิ สมโณ โคตโม นานุชานาติ, เต มยํ อิเมหิ ปญฺชหิ วตฺถูหิ สมาทาย วตฺตามา”ติ. ตตฺถ เย เต มนุสฺสา อสฺสทฺธา อปฺปสนฺนา ทุพฺพุทฺธิโน, เต เอวมาหํสุ “อิเม โข สมาณา สกฺยปุตฺติยา ธุตา สลฺเลข\-วุตฺติโน, สมโณ ปน โคตโม พาหุลฺลิโก พาหุลฺลาย เจเตตี”ติ. เย ปน เต มนุสฺสา สทฺธา ปสนฺนา ปณฺฑิตา พฺยตฺตา พุทฺธิมนฺโต, เต อุชฺฌายนฺติ ขิยฺยนฺติ วิปาเจนฺติ “กถํ หิ นาม เทวทตฺโต ภควโต สํฆเภทาย ปรกฺกมิสฺสติ จกฺกเภทายา”ติ. อสฺโสสุํ โข ภิกฺขู เตสํ มนุสฺสานํ อุชฺฌายนฺตานํ ขิยฺยนฺตานํ วิปาเจนฺตานํ. เย เต ภิกฺขู อปฺปิจฺฉา ฯเปฯ เต อุชฺฌายนฺติ ขิยฺยนฺติ \csromanfolio{265}วิปาเจนฺติ “กถํ หิ นาม เทวทตฺโต สํฆเภทาย ปรกฺกมิสฺสติ จกฺกเภทายา”ติ. อถ โข เต ภิกฺขู เทวทตฺตํ อเนก\-ปริยาเยน วิครหิตฺวา ภควโต เอตมตฺถํ อาโรเจสุํ ฯเปฯ “สจฺจํ กิร ตฺวํ เทวทตฺต สํฆเภทาย ปรกฺกมสิ จกฺกเภทายา”ติ. สจฺจํ ภควาติ. วิครหิ พุทฺโธ ภควา ฯเปฯ. กถํ หิ นาม ตฺวํ โมฆปุริส สํฆเภทาย ปรกฺกมิสฺสสิ จกฺกเภทาย, เนตํ โมฆปุริส อปฺปสนฺนานํ วา ปสาทาย ฯเปฯ. เอวญฺจ ปน ภิกฺขเว อิมํ สิกฺขาปทํ อุทฺทิเสยฺยาถ\csromandash{}}

\proseitem{411}{\textbf{“โย ปน ภิกฺขุ สมคฺคสฺส สํฆสฺส เภทาย ปรกฺกเมยฺย, เภทน}\-\textbf{สํวตฺตนิกํ วา อธิกรณํ สมาทาย ปคฺคยฺห ติฏฺเฐยฺย, โส ภิกฺขุ ภิกฺขูหิ เอวมสฺส วจนีโย ‘มายสฺมา สมคฺคสฺส สํฆสฺส เภทาย ปรกฺกมิ, เภทน}\-\textbf{สํวตฺตนิกํ วา อธิกรณํ สมาทาย ปคฺคยฺห อฏฺฐาสิ, สเมตายสฺมา สํเฆน, สมคฺโค หิ สํโฆ สมฺโมทมาโน อวิวทมาโน เอกุทฺเทโส ผาสุ วิหรตี”ติ, เอวญฺจ โส ภิกฺขุ ภิกฺขูหิ วุจฺจมาโน ตเถว ปคฺคเณฺหยฺย, โส ภิกฺขุ ภิกฺขูหิ ยาวตติยํ สมนุภาสิตพฺโพ ตสฺส ปฏินิสฺสคฺคาย, ยาวตติยญฺเจ สมนุ}\-\textbf{ภาสิย}\-\textbf{มาโน ตํ ปฏินิสฺสชฺเชยฺย, อิจฺเจตํ กุสลํ, โน เจ ปฏินิสฺสชฺเชยฺย, สํฆาทิเสโส”}ติ.}

\proseitem{412}{\textbf{โย ปนา}ติ โย ยาทิโส ฯเปฯ. \textbf{ภิกฺขู}ติ ฯเปฯ อยํ อิมสฺมิํ อตฺเถ อธิปฺเปโต “ภิกฺขู”ติ.}

\prose{\textbf{สมคฺโค} นาม \textbf{สํโฆ} สมาน\-สํวาสโก สมานสีมายํ ฐิโต.}

\prose{\textbf{เภทาย ปรกฺกเมยฺยา}ติ กถํ อิเม นานา อสฺสุ วินา อสฺสุ วคฺคา อสฺสูติ ปกฺขํ ปริเยสติ คณํ พนฺธติ.}

\prose{\textbf{เภทน}\-\textbf{สํวตฺตนิกํ วา อธิกรณนฺ}ติ อฏฺฐารส\-เภทกร\-วตฺถูนิ.}

\prose{\textbf{สมาทายา}ติ อาทาย.}

\prose{\textbf{ปคฺคยฺหา}ติ ทีเปยฺย.}

\prose{\textbf{ติฏฺเฐยฺยา}ติ น ปฏินิสฺสชฺเชยฺย.}

\prose{\textbf{โส ภิกฺขู}ติ โย โส สํฆเภทโก ภิกฺขุ.}

\prose{\textbf{ภิกฺขูหี}ติ อญฺเญหิ ภิกฺขูหิ.}

\prose{\csromanfolio{266}เย ปสฺสนฺติ, เย สุณนฺติ, เตหิ วตฺตพฺโพ “มายสฺมา สมคฺคสฺส สํฆสฺส เภทาย ปรกฺกมิ, เภทน\-สํวตฺตนิกํ วา อธิกรณํ สมาทาย ปคฺคยฺห อฏฺฐาสิ. สเมตายสฺมา สํเฆน, สมคฺโค หิ สํโฆ สมฺโมทมาโน อวิวทมาโน เอกุทฺเทโส ผาสุ วิหรตี”ติ. ทุติยมฺปิ วตฺตพฺโพ. ตติยมฺปิ วตฺตพฺโพ. สเจ ปฏินิสฺสชฺชติ, อิจฺเจตํ กุสลํ. โน เจ ปฏินิสฺสชฺชติ, อาปตฺติ ทุกฺกฏสฺส. สุตฺวา น วทนฺติ, อาปตฺติ ทุกฺกฏสฺส. โส ภิกฺขุ สํฆมชฺฌมฺปิ อากฑฺฒิตฺวา วตฺตพฺโพ “มายสฺมา สมคฺคสฺส สํฆสฺส เภทาย ปรกฺกมิ, เภทน\-สํวตฺตนิกํ วา อธิกรณํ สมาทาย ปคฺคยฺห อฏฺฐาสิ, สเมตายสฺมา สํเฆน, สมคฺโค หิ สํโฆ สมฺโมทมาโน อวิวทมาโน เอกุทฺเทโส ผาสุ วิหรตี”ติ. ทุติยมฺปิ วตฺตพฺโพ. ตติยมฺปิ วตฺตพฺโพ. สเจ ปฏินิสฺสชฺชติ, อิจฺเจตํ กุสลํ. โน เจ ปฏินิสฺสชฺชติ, อาปตฺติ ทุกฺกฏสฺส. โส ภิกฺขุ สมนุภาสิตพฺโพ, เอวญฺจ ปน ภิกฺขเว สมนุภาสิตพฺโพ. พฺยตฺเตน ภิกฺขุนา ปฏิพเลน สํโฆ ญาเปตพฺโพ.}

\proseitem{413}{“สุณตุ เม ภนฺเต สํโฆ, อยํ อิตฺถนฺนาโม ภิกฺขุ สมคฺคสฺส สํฆสฺส เภทาย ปรกฺกมติ, โส ตํ วตฺถุํ น ปฏินิสฺสชฺชติ, ยทิ สํฆสฺส ปตฺตกลฺลํ, สํโฆ อิตฺถนฺนามํ ภิกฺขุํ สมนุภาเสยฺย ตสฺส วตฺถุสฺส ปฏินิสฺสคฺคาย, เอสา ญตฺติ.}

\prose{สุณาตุ เม ภนฺเต สํโฆ, อยํ อิตฺถนฺนาโม ภิกฺขุ สมคฺคสฺส สํฆสฺส เภทาย ปรกฺกมติ, โส ตํ วตฺถุํ น ปฏินิสฺสชฺชติ, สํโฆ อิตฺถนฺนามํ ภิกฺขุํ สมนุภาสติ ตสฺส วตฺถุสฺส ปฏินิสฺสคฺคาย, ยสฺสายสฺมโต ขมติ อิตฺถนฺนามสฺส ภิกฺขุโน สมนุภาสนา ตสฺส วตฺถุสฺส ปฏินิสฺสคฺคาย, โส ตุณฺหสฺส, ยสฺส นกฺขมติ, โส ภาเสยฺย.}

\prose{ทุติยมฺปิ เอตมตฺถํ วทามิ ฯเปฯ. ตติยมฺปิ เอตมตฺถํ วาทามิ. สุณาตุ เม ภนฺเต สํโฆ, อยํ อิตฺถนฺนาโม ภิกฺขุ สมคฺคสฺส สํฆสฺส เภทาย ปรกฺกมติ, โส ตํ วตฺถุํ น ปฏินิสฺสชฺชติ, สํโฆ อิตฺถนฺนามํ ภิกฺขุํ สมนุภาสติ ตสฺส วตฺถุสฺส ปฏินิสฺสคฺคาย, ยสฺสายสฺมโต ขมติ อิตฺถนฺนามสฺส ภิกฺขุโน สมนุภาสนา ตสฺส วตฺถุสฺส ปฏินิสฺสคฺคาย, โส ตุณฺหสฺส, ยสฺส นกฺขมติ, โส ภาเสยฺย.}

\prose{สมนุภฏฺโฐ สํเฆน อิตฺถนฺนาโม ภิกฺขุ ตสฺส วตฺถุสฺส ปฏินิสฺสคฺคาย, ขมติ สํฆสฺส, ตสฺมา ตุณฺหี, เอวเมตํ ธารยามี”ติ.}

\csromanmatikaanchor{mtk.89}
\csromanmatikaanchor{mtk.90}
\csromantocmarknopage{section}{11. เภทานุวตฺตกสิกฺขาปท}{mtk.89}
\bookmark[dest={mtk.89},level=1]{11. เภทานุวตฺตกสิกฺขาปท}
\csromantocmarkref{subsection}{เภทานุวตฺตกภิกฺขุวตฺถุ}{mtk.90}
\bookmark[dest={mtk.90},level=2]{เภทานุวตฺตกภิกฺขุวตฺถุ}
\proseitem{414}{\csromanfolio{267}ญตฺติยา ทุกฺกฏํ, ทฺวีหิ กมฺมวาจาหิ ถุลฺลจฺจยา, กมฺมวาจา\-ปริโยสาเน อาปตฺติ สํฆาทิเสสสฺส. สํฆาทิเสสํ อชฺฌาปชฺชนฺตสฺส ญตฺติยา ทุกฺกฏํ ทฺวีหิ กมฺมวาจาหิ ถุลฺลจฺจยา ปฏิปฺปสฺสมฺภนฺติ.}

\prose{\textbf{สํฆาทิเสโส}ติ ฯเปฯ เตนปิ วุจฺจติ “สํฆาทิเสโส”ติ.}

\proseitem{415}{ธมฺมกมฺเม ธมฺมกมฺม\-สญฺญี น ปฏินิสฺสชฺชติ, อาปตฺติ สํฆาทิเสสสฺส.}

\prose{ธมฺมกมฺเม เวมติโก น ปฏินิสฺสชฺชติ, อาปตฺติ สํฆาทิเสสสฺส.}

\prose{ธมฺมกมฺเม อธมฺมกมฺม\-สญฺญี น ปฏินิสฺสชฺชติ, อาปตฺติ สํฆาทิเสสสฺส.}

\prose{อธมฺมกมฺเม ธมฺมกมฺม\-สญฺญี, อาปตฺติ ทุกฺกฏสฺส.}

\prose{อธมฺมกมฺเม เวมติโก, อาปตฺติ ทุกฺกฏสฺส.}

\prose{อธมฺมกมฺเม อธมฺมกมฺม\-สญฺญี, อาปตฺติ ทุกฺกฏสฺส.}

\proseitemwithcloserruled{416}{อนาปตฺติ อสมนุภาสนฺตสฺส ปฏินิสฺสชฺชนฺตสฺส อุมฺมตฺตกสฺส ขิตฺตจิตฺตสฺส เวทนาฏฺฏสฺส อาทิกมฺมิกสฺสาติ.}{สํฆเภทสิกฺขาปทํ นิฏฺฐิตํ ทสมํ.}
\csromanheader{1}{11. \textbf{เภทานุวตฺตก}\-\textbf{สิกฺขาปท}}
\csromanmatikaanchor{mtk.91}
\csromanmatikaanchor{mtk.92}
\csromantocmarkref{subsection}{ปญฺญตฺติ}{mtk.91}
\bookmark[dest={mtk.91},level=2]{ปญฺญตฺติ}
\csromantocmarkref{subsection}{สิกฺขาปทวิภงฺค, ปทภาชนีย}{mtk.92}
\bookmark[dest={mtk.92},level=2]{สิกฺขาปทวิภงฺค, ปทภาชนีย}
\proseitem{417}{เตน สมเยน พุทฺโธ ภควา ราชคเห วิหรติ เวฬุวเน กลนฺทก\-นิวาเป. เตน โข ปน สมเยน เทวทตฺโต สํฆเภทาย ปรกฺกมติ จกฺกเภทาย. ภิกฺขู เอวมาหํสุ “อธมฺมวาที เทวทตฺโต, อวินยวาที เทวทตฺโต, กถํ หิ นาม เทวทตฺโต สํฆเภทาย ปรกฺกมิสฺสติ จกฺกเภทายา”ติ. เอวํ วุตฺเต โกกาลิโก กฏโมทกติสฺสโก ขณฺฑเทวิยา ปุตฺโต สมุทฺททตฺโต เต ภิกฺขู เอตทโวจุํ “มายสฺมนฺโต เอวํ อวจุตฺถ, ธมฺมวาที เทวทตฺโต, วินยวาที เทวทตฺโต, อมฺหากญฺจ เทวทตฺโต ฉนฺทญฺจ รุจิญฺจ อาทาย โวหรติ ชานาติ โน ภาสติ อมฺหากํเปตํ ขมตี”ติ. เย เต ภิกฺขู อปฺปิจฺฉา ฯเปฯ เต อุชฺฌายนฺติ ขิยฺยนฺติ วิปาเจนฺติ “กถํ หิ นาม ภิกฺขู \csromanfolio{268}เทวทตฺตสฺส สํฆเภทาย ปรกฺกม\-นฺตสฺส อนุวตฺตกา ภวิสฺสนฺติ วคฺควาทกา”ติ. อถ โข เต ภิกฺขู เต อนุวตฺเตเก ภิกฺขู อเนก\-ปริยาเยน วิครหิตฺวา ภควโต เอตมตฺถํ อาโรเจสุํ ฯเปฯ “สจฺจํ กิร ภิกฺขเว ภิกฺขู เทวทตฺตสฺส สํฆเภทาย ปรกฺกม\-นฺตสฺส อนุวตฺตากา โหนฺติ วคฺควาทกา”ติ. สจฺจํ ภควาติ. วิครหิ พุทฺโธ ภควา ฯเปฯ. กถํ หิ นาม เต ภิกฺขเว โมฆปุริสา เทวทตฺตสฺส สํฆเภทาย ปรกฺกม\-นฺตสฺส อนุวตฺตกา ภวิสฺสนฺติ วคฺควาทกา, เนตํ ภิกฺขเว อปฺปสนฺนานํ วา ปสาทาย ฯเปฯ. เอวญฺจ ปน ภิกฺขเว อิมํ สิกฺขาปทํ อุทฺทิเสยฺยาถ\csromandash{}}

\proseitem{418}{\textbf{“ตสฺเสว โข ปน ภิกฺขุสฺส ภิกฺขู โหนฺติ อนุวตฺตกา วคฺควาทกา เอโก วา เทฺว วา ตโย วา, เต เอวํ วเทยฺยุํ “มายสฺมนฺโต เอตํ ภิกฺขุํ กิญฺจิ อวจุตฺถ, ธมฺมวาที เจโส ภิกฺขุ, วินยวาที เจโส ภิกฺขุ, อมฺหากญฺเจโส ภิกฺขุ ฉนฺทญฺจ รุจิญฺจ อาทาย โวหรติ ชานาติ โน ภาสติ อมฺหากํเปตํ ขมตี’ติ, เต ภิกฺขู ภิกฺขูหิ เอวมสฺสุวจนียา ‘มายสฺมนฺโต เอวํ อวจุตฺถ, น เจโส ภิกฺขุ ธมฺมวาที, น เจโส ภิกฺขุ วินยวาที, มายสฺมนฺตานํปิ สํฆเภโท รุจฺจิตฺถ, สเมตา}\-\textbf{ยสฺมนฺตานํ สํเฆน, สมคฺโค หิ สํโฆ สมฺโมทมาโน อวิวทมาโน เอกุทฺเทโส ผาสุ วิหรตี’ติ, เอวญฺจ เต ภิกฺขู ภิกฺขูหิ วุจฺจมานา ตเถว ปคฺคเณฺหยฺยุํ, เต ภิกฺขู ภิกฺขูหิ ยาวตติยํ สมนุภาสิตพฺพา ตสฺส ปฏินิสฺสคฺคาย, ยาวตติยญฺเจ สมนุ}\-\textbf{ภาสิย}\-\textbf{มานา ตํ ปฏินิสฺสชฺเชยฺยุํ, อิจฺเจตํ กุสลํ, โน เจ ปฏินิสฺสชฺเชยฺยุํ, สํฆาทิเสโส”}ติ.}

\proseitem{419}{\textbf{ตสฺเสว โข ปนา}ติ ตสฺส สํฆเภทกสฺส ภิกฺขุโน.}

\prose{\textbf{ภิกฺขู โหนฺตี}ติ อญฺเญ ภิกฺขู โหนฺติ.}

\prose{\textbf{อนุวตฺตกา}ติ ยํทิฏฺฐิโก โส โหติ ยํขนฺติโก ยํรุจิโก, เตปิ ตํทิฏฺฐิกา โหนฺติ ตํขนฺติกา ตํรุจิกา.}

\prose{\textbf{วคฺควาทกา}ติ ตสฺส วณฺณาย ปกฺขาย ฐิตา โหนฺติ.}

\prose{\textbf{เอโก วา เทฺว วา ตโย วา}ติ เอโก วา โหติ เทฺว วา ตโย วา. เต เอวํ วเทยฺยุํ “มายสฺมนฺโต เอตํ ภิกฺขุํ กิญฺจิ อวจุตฺถ, ธมฺมวาที \csromanfolio{269}เจโส ภิกฺขุ, วินยวาที เจโส ภิกฺขุ, อมฺหากญฺเจโส ภิกฺขุ ฉนฺทญฺจ รุจิญฺจ อาทาย โวหรติ ชานาติ โน ภาสติ อมฺหากํเปตํ ขมตี”ติ.}

\prose{\textbf{เต ภิกฺขู}ติ เย เต อนุวตฺตกา ภิกฺขู.}

\prose{\textbf{ภิกฺขูหี}ติ อญฺเญหิ ภิกฺขูหิ.}

\prose{เย ปสฺสนฺติ, เย สุณนฺติ, เตหิ วตฺตพฺพา “มายสฺมนฺโต เอวํ อวจุตฺถ, น เจโส ภิกฺขุ ธมฺมวาที, น เจโส ภิกฺขุ วินยวาที, มายสฺมนฺตานํปิ สํฆเภโท รุจฺจิตฺถ, สเมตา\-ยสฺมนฺตานํ สํเฆน, สมคฺโค หิ สํโฆ สมฺโมทมาโน อวิวทมาโน เอกุทฺเทโส ผาสุ วิหรตี”ติ. ทุติยมฺปิ วตฺตพฺพา. ตติยมฺปิ วตฺตพฺพา. สเจ ปฏินิสฺสชฺชนฺติ, อิจฺเจตํ กุสลํ. โน เจ ปฏินิสฺสชฺชนฺติ, อาปตฺติ ทุกฏสฺส. สุตฺวา น วทนฺติ, อาปตฺติ ทุกฺกฏสฺส. เต ภิกฺขู สํฆมชฺฌมฺปิ อากฑฺฒิตฺวา วตฺตพฺพา “มายสฺมนฺโต เอวํ อวจุตฺถ, น เจโส ภิกฺขุ ธมฺมวาที, น เจโส ภิกฺขุ วินยวาที, มายสฺมนฺตานํปิ สํฆเภโท รุจฺจิตฺถ, สเมตา\-ยสฺมนฺตานํ สํเฆน, สมคฺโค หิ สํโฆ สมฺโมทมาโน อวิวทมาโน เอกุทฺเทโส ผาสุ วิหรตี”ติ. ทุติยมฺปิ วตฺตพฺพา. ตติยมฺปิ วตฺตพฺพา. สเจ ปฏินิสฺสชฺชนฺติ, อิจฺเจตํ กุสลํ. โน เจ ปฏิสฺสชฺชนฺติ, อาปตฺติ ทุกฺกฏสฺส. เต ภิกฺขู สมนุภาสิตพฺพา, เอวญฺจ ปน ภิกฺขเว สมนุภาสิตพฺพา. พฺยตฺเตน ภิกฺขุนา ปฏิพเลน สํโฆ ญาเปตพฺโพ\csromandash{}}

\proseitem{420}{“สุณาตุ เม ภนฺเต สํโฆ, อิตฺถนฺนาโม จ อิตฺถนฺนาโม จ ภิกฺขู อิตฺถนฺนามสฺส ภิกฺขุโน สํฆเภทาย ปรกฺกม\-นฺตสฺส อนุวตฺตกา วคฺควาทกา, เต ตํ วตฺถุํ น ปฏินิสฺสชฺชนฺติ, ยทิ สํฆสฺส ปตฺตกลฺลํ, สํโฆ อิตฺถนฺนามญฺจ อิตฺถนฺนามญฺจ ภิกฺขู สมนุภาเสยฺย ตสฺส วตฺถุสฺส ปฏินิสฺสคฺคาย, เอสา ญตฺติ.}

\prose{สุณาตุ เม ภนฺเต สํโฆ, อิตฺถนฺนาโม จ อิตฺถนฺนาโม จ ภิกฺขู อิตฺถนฺนามสฺส ภิกฺขุโน สํฆเภทาย ปรกฺกม\-นฺตสฺส อนุวตฺตกา วคฺควาทกา, เต ตํ วตฺถุํ น ปฏินิสฺสชฺชนฺติ, สํโฆ อิตฺถนฺนามญฺจ อิตฺถนฺนามญฺจ ภิกฺขู สมนุภาสติ ตสฺส วตฺถุสฺส ปฏินิสฺสคฺคาย, ยสฺสายสฺมโต ขมติ อิตฺถนฺธามสฺส จ อิตฺถนฺนามสฺส จ ภิกฺขุนํ สมนุภาสนา ตสฺส \csromanfolio{270}วตฺถุสฺส ปฏินิสฺสคฺคาย, โส ตุณฺหสฺส, ยสฺส นกฺขมติ, โส ภาเสยฺย.}

\prose{ทุติยมฺปิ เอตมตฺถํ วทามิ ฯเปฯ. ตติยมฺปิ เอตมตฺถํ วทามิ. สุณาตุ เม ภนฺเต สํโฆ, อิตฺถนฺนาโม จ อิตฺถนฺนาโม จ ภิกฺขู อิตฺถนฺนามสฺส ภิกฺขุโน สํฆเภทาย ปรกฺกม\-นฺตสฺส อนุวตฺตกา วคฺควาทกา, เต ตํ วตฺถุํ น ปฏินิสฺสชฺชนฺติ, สํโฆ อิตฺถนฺนามญฺจ อิตฺถนฺนามญฺจ ภิกฺขู สมนุภาสติ ตสฺส วตฺถุสฺส ปฏินิสฺสคฺคาย, ยสฺสายสฺมโต ขมติ อิตฺถนฺนามสฺส จ อิตฺถนฺนามสฺส จ ภิกฺขูนํ สมนุภาสนา ตสฺส วตฺถุสฺส ปฏินิสฺสคฺคาย, โส ตุณฺหาสฺส, ยสฺส นกฺขมติ, โส ภาเสยฺย.}

\prose{สมนุภฏฺฐา สํเฆน อิตฺถนฺนาโม จ อิตฺถนฺนาโม จ ภิกฺขู ตสฺส วตฺถุสฺส ปฏินิสฺสคฺคาย, ขมติ สํฆสฺส, ตสฺมา ตุณฺหี, เอวเมตํ ธารยามี”ติ.}

\proseitem{421}{ญตฺติยา ทุกฺกฏํ, ทฺวีหิ กมฺมวาจาหิ ถุลฺลจฺจยา, กมฺมวาจา\-ปริโยสาเน อาปตฺติ สํฆาทิเสสสฺส. สํฆาทิเสสํ อชฺฌาปชฺชนฺตานํ ญตฺติยา ทุกฺกฏํ ทฺวีหิ กมฺมวาจาหิ ถุลฺลจฺจยา ปฏิปฺปสฺสมฺภนฺติ. เทฺว ตโย เอกโต สมนุภาสิตพฺพา, ตตุตฺตริ น สมนุภาสิตพฺพา.}

\prose{\textbf{สํฆาทิเสโส}ติ ฯเปฯ เตนปิ วุจฺจติ “สํฆาทิเสโส”ติ.}

\proseitem{422}{ธมฺมกมฺเม ธมฺมกมฺม\-สญฺญี น ปฏินิสฺสชฺชนฺติ อาปตฺติ สํฆาทิเสสสฺส.}

\prose{ธมฺมกมฺเม เวมติกา น ปฏินิสฺสชฺชนฺติ, อาปตฺติ สํฆาทิเสสสฺส.}

\prose{ธมฺมกมฺเม อธมฺมกมฺม\-สญฺญี น ปฏินิสฺสชฺชนฺติ, อาปตฺติ สํฆาทิเสสสฺส.}

\prose{อธมฺมกมฺเม ธมฺมกมฺม\-สญฺญิ, อาปตฺติ ทุกฺกฏสฺส.}

\prose{อธมฺมกมฺเม เวมติกา, อาปตฺติ ทุกฺกฏสฺส.}

\prose{อธมฺมกมฺม อธมฺมกมฺม\-สญฺญี, อาปตฺติ ทุกฺกฏสฺส.}

\proseitem{423}{อนาปตฺติ อสมนุภาสนฺตานํ ปฏินิสฺสชฺชนฺตานํ อุมฺมตฺตกานํ ขิตฺตจิตฺตานํ เวทนาฏฺฏานํ อาทิ\-กมฺมิกานนฺติ.}

\vspace{\tipitakaparskip}
\csromancenter{เภทานุวตฺตก\-สิกฺขาปทํ นิฏฺฐิตํ เอกาทสมํ.}
\csromanmatikaanchor{mtk.93}
\csromantocmarknopage{section}{12. ทุพฺพจสิกฺขาปท}{mtk.93}
\bookmark[dest={mtk.93},level=1]{12. ทุพฺพจสิกฺขาปท}
\csromanheader{1}{12. \textbf{ทุพฺพจ}\-\textbf{สิกฺขาปท}}
\csromanmatikaanchor{mtk.94}
\csromanmatikaanchor{mtk.95}
\csromantocmarkref{subsection}{ฉนฺนภิกฺขุวตฺถุ}{mtk.94}
\bookmark[dest={mtk.94},level=2]{ฉนฺนภิกฺขุวตฺถุ}
\csromantocmarkref{subsection}{ปญฺญตฺติ}{mtk.95}
\bookmark[dest={mtk.95},level=2]{ปญฺญตฺติ}
\proseitem{424}{\csromanfolio{271}เตน สมเยน พุทฺโธ ภควา โกสมฺพิยํ วิหรติ โฆสิตาราเม. เตน โข ปน สมเยน อายสฺมา ฉนฺโน อนาจารํ อาจรติ. ภิกฺขู เอวมาหํสุ “มาวุโส ฉนฺน เอวรูปํ อกาสิ, เนตํ กปฺปตี”ติ. โส เอวํ วเทติ “กิํ นุ โข นาม ตุเมฺห อาวุโส มํ วตฺตพฺพํ มญฺญถ, อหํ โข นาม ตุเมฺห วเทยฺยํ, อมฺหากํ พุทฺโธ อมฺหากํ ธมฺโม อมฺหากํ อยฺยปุตฺเตน ธมฺโม อภิสมิโต, เสยฺยถาปิ นาม มหาวาโต วายนฺโต ติณกฏฺฐ\-ปณฺณสฏํ\csromansharedfootnote{9aa451e6d32f5490}{ติณกฏฺฐปณฺณกสฏํ (ก)} เอกโต อุสฺสาเรยฺย, เสยฺยถา วา ปน นที ปพฺพเตยฺยา สงฺข\-เสวาล\-ปณกํ เอกโต อุสฺสาเรยฺย, เอวเมว ตุเมฺห นานานามา นานาโคตฺตา นานาชจฺจา นานากุลา ปพฺพชิตา เอกโต อุสฺสาริตา, กิํ นุ โข นาม ตุเมฺห อาวุโส มํ วตฺตพฺพํ มญฺญถ, อหํ โข นาม ตุเมฺห วเทยฺยํ, อมฺหากํ พุทฺโธ อมฺหากํ ธมฺโม อมฺหากํ อยฺยปุตฺเตน ธมฺโม อภิสมิโต”ติ. เย เต ภิกฺขู อปฺปิจฺฉา ฯเปฯ เต อุชฺฌายนฺติ ขิยฺยนฺติ วิปาเจนฺติ “กถํ หิ นาม อายสฺมา ฉนฺโน ภิกฺขูหิ สหธมฺมิกํ วุจฺจมาโน อตฺตานํ อวจนียํ กริสฺสตี”ติ. อถ โข เต ภิกฺขู อายสฺมนฺตํ ฉนฺนํ อเนก\-ปริยาเยน วิครหิตฺวา ภควโต เอตมตฺถํ อาโรเจสุํ ฯเปฯ “สจฺจํ กิร ตฺวํ ฉนฺน ภิกฺขูหิ สหธมฺมิกํ วุจฺจมาโน อตฺตานํ อวจนียํ กโรสี”ติ. สจฺจํ ภควาติ. วิครหิ พุทฺโธ ภควา ฯเปฯ. กถํ หิ นาม ตฺวํ โมฆปุริส ภิกฺขูหิ สหธมฺมิกํ วุจฺจมาโน อตฺตานํ อวจนียํ กริสฺสสิ, เนตํ โมฆปุริส อปฺปสนฺนานํ วา ปสาทาย ฯเปฯ. เอวญฺจ ปน ภิกฺขเว อิมํ สิกฺขาปทํ อุทฺทิเสยฺยาถ\csromandash{}}

\csromanmatikaanchor{mtk.96}
\csromantocmarkref{subsection}{สิกฺขาปทวิภงฺค, ปทภาชนีย}{mtk.96}
\bookmark[dest={mtk.96},level=2]{สิกฺขาปทวิภงฺค, ปทภาชนีย}
\proseitem{425}{\textbf{“ภิกฺขุ ปเนว ทุพฺพจชาติโก โหติ อุทฺเทส}\-\textbf{ปริยาปนฺเนสุ สิกฺขาปเทสุ ภิกฺขูหิ สหธมฺมิกํ วุจฺจมาโน อตฺตานํ อวจนียํ กโรติ ‘มา มํ อายสฺมนฺโต กิญฺจิ อวจุตฺถ กลฺยาณํ วา ปาปกํ วา, อหํปายสฺมนฺเต น กิญฺจิ วกฺขามิ กลฺยาณํ วา ปาปกํ วา, วิรมถา}\-\textbf{ยสฺมนฺโต มม วจนายา’ติ, โส ภิกฺขุ ภิกฺขูหิ เอวมสฺส วจนีโย ‘มายสฺมา อตฺตานํ อวจนียํ อกาสิ, วจนียเม}\-\textbf{วา}\-\textbf{ยสฺมา อตฺตานํ กโรตุ, อายสฺมาปิ ภิกฺขู วเทตุ สหธมฺเมน, ภิกฺขูปิ อายสฺมนฺตํ วกฺขนฺติ สหธมฺเมน, เอวํ สํวทฺธา หิ ตสฺส ภควโต ปริสา ยทิทํ} \csromanfolio{272}\textbf{อญฺญม}\-\textbf{ญฺญ}\-\textbf{วจเนน อญฺญมญฺญวุฏฺฐาปเนนา’ติ, เอวญฺจ โส ภิกฺขุ ภิกฺขูหิ วุจฺจมาโน ตเถว ปคฺคเณฺหยฺย, โส ภิกฺขุ ภิกฺขูหิ ยาวตติยํ สมนุภาสิตพฺโพ ตสฺส ปฏินิสฺสคฺคาย, ยาวตติยญฺเจ สมนุ}\-\textbf{ภาสีย}\-\textbf{มาโน ตํ ปฏินิสฺสชฺเชยฺย, อิจฺเจตํ กุสลํ, โน เจ ปฏินิสฺสชฺเชยฺย, สํฆาทิเสโส”}ติ.}

\proseitem{426}{\textbf{ภิกฺขุ ปเนว ทุพฺพจชาติโก โหตี}ติ ทุพฺพโจ โหติ โทวจสฺส\-กรเณหิ ธมฺเมหิ สมนฺนาคโต อกฺขโม อปฺปทกฺขิณคฺคาหี อนุสาสนิํ.}

\prose{\textbf{อุทฺเทส}\-\textbf{ปริยาปนฺเนสุ สิกฺขาปเทสู}ติ ปาติโมกฺข\-ปริยาปนฺเนสุ สิกฺขาปเทสุ.}

\prose{\textbf{ภิกฺขูหี}ติ อญฺเญหิ ภิกฺขูหิ.}

\prose{\textbf{สหธมฺมิกํ} นาม ยํ ภาควตา ปญฺญตฺตํ สิกฺขาปทํ, เอตํ สหธมฺมิกํ นาม.}

\prose{เตน วุจฺจมาโน อตฺตานํ อวจนียํ กโรติ “มา มํ อายสฺมนฺโต กิญฺจิ อวจุตฺถ กลฺยาณํ วา ปาปกํ วา, อหํปายสฺมนฺเต น กิญฺจิ วกฺขามิ กลฺยาณํ วา ปาปกํ วา, วิรมถา\-ยสฺมนฺโต มม วจนายา”ติ.}

\prose{\textbf{โส ภิกฺขู}ติ โย โส ทุพฺพจชาติโก ภิกฺขุ.}

\prose{\textbf{ภิกฺขูหี}ติ อญฺเญหิ ภิกฺขูหิ. เย ปสฺสนฺติ เย สุณนฺติ, เตหิ วตฺตพฺโพ “มายสฺมา อตฺตานํ อวจนียํ อกาสิ, วจนียเมว อายสฺมา อตฺตานํ กโรตุ, อายสฺมาปิ ภิกฺขู วเทตุ สหธมฺเมน, ภิกฺขูปิ อายสฺมนฺตํ วกฺขนฺติ สหธมฺเมน, เอวํ สํวทฺธา หิ ตสฺส ภควโต ปริสา ยทิทํ อญฺญม\-ญฺญ\-วจเนน อญฺญมญฺญวุฏฺฐาปเนนา”ติ. ทุติยมฺปิ วตฺตพฺโพ. ตติยมฺปิ วตฺตพฺโพ. สเจ ปฏินิสฺสชฺชติ, อิจฺเจตํ กุสลํ. โน เจ ปฏินิสฺสชฺชติ, อาปตฺติ ทุกฺกฏสฺส. สุตฺวา น วทนฺติ, อาปตฺติ ทุกฺกฏสฺส. โส ภิกฺขุ สํฆมชฺฌมฺปิ อากฑฺฒิตฺวา วตฺตพฺโพ “มายสฺมา อตฺตานํ อวจนียํ อกาสิ ฯเปฯ อญฺญมญฺญวุฏฺฐาปเนนา”ติ. ทุติยมฺปิ วตฺตพฺโพ. ตติยมฺปิ วตฺตพฺโพ. สเจ ปฏินิสฺสชฺชติ, อิจฺเจตํ กุสลํ. โน เจ ปฏินิสฺสชฺชติ, อาปตฺติ ทุกฺกฏสฺส. โส ภิกฺขุ สมนุภาสิตพฺโพ, เอวญฺจ ปน ภิกฺขเว สมนุภาสิตพฺโพ, พฺยตฺเตน ภิกฺขุนา ปฏิพเลน สํโฆ ญาเปตพฺโพ\csromandash{}}

\proseitem{427}{\csromanfolio{273}“สุณาตุ เม ภนฺเต สํโฆ, อยํ อิตฺถนฺนาโม ภิกฺขุ ภิกฺขูหิ สหธมฺมิกํ วุจฺจมาโน อตฺตานํ อวจนียํ กโรติ, โส ตํ วตฺถุํ น ปฏินิสฺสชฺชติ, ยทิ สํฆสฺส ปตฺตกลฺลํ, สํโฆ อิตฺถนฺนามํ ภิกฺขุํ สมนุภาเสยฺย ตสฺส วตฺถุสฺส ปฏินิสฺสคฺคาย, เอสา ญตฺติ.}

\prose{สุณาตุ เม ภนฺเต สํโฆ, อยํ อิตฺถนฺนาโม ภิกฺขุ ภิกฺขูหิ สหธมฺมิกํ วุจฺจมาโน อตฺตานํ อวจนียํ กโรติ, โส ตํ วตฺถุํ ปฏินิสฺสชฺชติ, สํโฆ อิตฺถนฺนามํ ภิกฺขุํ สมนุภาสติ ตสฺส วตฺถุสฺส ปฏินิสฺสคฺคาย, ยสฺสายสฺมโต ขมติ อิตฺถนฺนามสฺส ภิกฺขุโน สมนุภาสนา ตสฺส วตฺถุสฺส ปฏินิสฺสคฺคาย, โส ตุณฺหสฺส, ยสฺส นกฺขมติ, โส ภาเสยฺย.}

\prose{ทุติยมฺปิ เอตมตฺถํ วทามิ ฯเปฯ. ตติยมฺปิ เอตมตฺถํ วทามิ. สุณาตุ เม ภนฺเต สํโฆ, อยํ อิตฺถนฺนาโม ภิกฺขุ ภิกฺขูหิ สหธมฺมิกํ วุจฺจมาโน อเตนํ อวจนียํ กโรติ, โส ตํ วตฺถุํ น ปฏินิสฺสชฺชติ, สํโฆ อิตฺถนฺนามํ ภิกฺขุํ สมนุภาสติ ตสฺส วตฺถุสฺส ปฏินิสฺสคฺคาย, ยสฺสายสฺมโต ขมติ อิตฺถนฺนามสฺส ภิกฺขุโน สมนุภาสนา ตสฺส วตฺถุสฺส ปฏินิสฺสคฺคาย, โส ตุณฺหสฺส, ยสฺส นกฺขมติ, โส ภาเสยฺย.}

\prose{สมนุภฏฺโฐ สํเฆน อิตฺถนฺนาโม ภิกฺขุ ตสฺส วตฺถุสฺส ปฏินิสฺสคฺคาย, ขมติ สํฆสฺส, ตสฺมา ตุณฺหี, เอวมาตํ ธารยามี”ติ.}

\proseitem{428}{ญตฺติยา ทุกฺกฏํ, ทฺวีหิ กมฺมวาจาหิ ถุลฺลจฺจยา, กมฺมวาจา\-ปริโยสาเน อาปตฺติ สํฆาทิเสสสฺส. สํฆาทิเสสํ อชฺฌาปชฺชนฺตสฺส ญตฺติยา ทุกฺกฏํ ทฺวีหิ กมฺมวาจาหิ ถุลฺลจฺจยา ปฏิปฺปสฺสมฺภนฺติ.}

\prose{\textbf{สํฆาทิเสโส}ติ ฯเปฯ เตนปิ วุจฺจติ “สํฆาทิเสโส”ติ.}

\proseitem{429}{ธมฺมกมฺเม ธมฺมกมฺม\-สญฺญี น ปฏินิสฺสชฺชติ, อาปตฺติ สํฆาทิเสสสฺส.}

\prose{ธมฺมกมฺเม เวมติโก น ปฏินิสฺสชฺชติ, อาปตฺติ สํฆาทิเสสสฺส.}

\prose{ธมฺมกมฺเม อธมฺมกมฺม\-สญฺญี น ปฏินิสฺสชฺชติ, อาปตฺติ สํฆาทิเสสสฺส.}

\prose{อธมฺมกมฺเม ธมฺมกมฺม\-สญฺญี, อาปตฺติ ทุกฺกฏสฺส.}

\prose{อธมฺมกมฺเม เวมติโก, อาปตฺติ ทุกฺกฏสฺส.}

\prose{อธมฺมกมฺเม อธมฺมกมฺม\-สญฺญี, อาปตฺติ ทุกฺกฏสฺส.}

\csromanmatikaanchor{mtk.97}
\csromanmatikaanchor{mtk.98}
\csromantocmarknopage{section}{13. กุลทูสกสิกฺขาปท}{mtk.97}
\bookmark[dest={mtk.97},level=1]{13. กุลทูสกสิกฺขาปท}
\csromantocmarkref{subsection}{อสฺสชิปุนพฺพสุกภิกฺขุวตฺถุ}{mtk.98}
\bookmark[dest={mtk.98},level=2]{อสฺสชิปุนพฺพสุกภิกฺขุวตฺถุ}
\proseitemwithcloserruled{430}{\csromanfolio{274}อนาปตฺติ อสมนุภาสนฺตสฺส ปฏินิสฺสชฺชนฺตสฺส อุมฺมตฺตกสฺส อาทิกมฺมิกสฺสาติ.}{ทุพฺพจ\-สิกฺขาปทํ นิฏฺฐิตํ ทฺวาทสมํ.}
\csromanheader{1}{13. \textbf{กุลทูสก}\-\textbf{สิกฺขาปท}}
\proseitem{431}{\csromansymbolfootnote{*}{อิทํ วตฺถุ วิ 4. 22 ปิฏฺฐาทีสุปิ อาคตํ.}เตน สมเยน พุทฺโธ ภควา สาวตฺถิยํ วิหรติ เชตวเน อนาถ\-ปิณฺฑิกสฺส อาราเม. เตน โข ปน สมเยน อสฺสชิ\-ปุนพฺพสุกา นาม\csromansharedfootnote{e177e442d5af50ad}{นาม ภิกฺขู (ก)} กีฏาคิริสฺมิํ อาวาสิกา โหติ อลชฺชิโน ปาปภิกฺขู. เต\csromansymbolfootnote{+}{วิ 4. 284 ปิฏฺฐาทีสุปิ.}เอวรูปํ อนาจารํ อาจรนฺติ, มาลาวจฺฉํ โรเปนฺติปิ โรปาเปนฺติปิ, สิญฺจนฺติปิ สิญฺจาเปนฺติปิ, โอจินนฺติปิ โอจินาเปนฺติปิ, คนฺเถนฺติปิ คนฺถาเปนฺติปิ, เอกโตวณฺฏิกมาลํ กโรนฺติปิ การาเปนฺติปิ, อุภโตวณฺฏิกมาลํ กโรนฺติปิ การาเปนฺติปิ, มญฺชริกํ กโรนฺติปิ การาเปนฺติปิ, วิธูติกํ กโรนฺติปิ การาเปนฺติปิ, วฏํสกํ กโรนฺติปิ การาเปนฺติปิ, อาเวฬํ กโรนฺติปิ กราเปนฺติปิ, อุรจฺฉทํ กโรนฺติปิ การาเปนฺติปิ. เต กุลิตฺถีนํ กุลธีตานํ กุลกุมารีนํ กุลสุณฺหานํ กุลทาสีนํ เอกโตวณฺฏิกมาลํ หรนฺติปิ หราเปนฺติปิ, อุภโตวณฺฏิกมาลํ หรนฺติปิ หราเปนฺติปิ, มญฺชริกํ หราเปนฺติปิ, วิธูติกํ หรนฺติปิ หราเปนฺติปิ, วฏํสกํ หรนฺติปิ หราเปนฺติปิ, อาเวฬํ หรนฺติปิ หราเปนฺติปิ, อุรจฺฉทํ หรนฺติปิ หราเปนฺติปิ. เต กุลิตฺถีหิ กุลธีตาหิ กุลกุมารีหิ กุลสุณฺหาหิ กุลทาสีหิ สทฺธิํ เอกภาชเนปิ ภุญฺชนฺติ, เอกถาลเกปิ ปิวนฺติ, เอกาสเนปิ นิสีทนฺติ, เอกมญฺเจปิ ตุวฏฺเฏนฺติ, เอกตฺถรณาปิ ตุวฏฺเฏนฺติ, เอกวาวุรณาปิ ตุวฏฺเฏนฺติ, เอก\-ตฺถรณ\-ปาวุรณาปิ ตุวฏฺเฏนฺติ, วิกาเลปิ ภุญฺชนฺติ, มชฺชมฺปิ ปิวนฺติ, มาลา\-คนฺธ\-วิเลปนมฺปิ ธาเรนฺติ, นจฺจนฺติปิ คายนฺติปิ วาเทนฺติปิ ลาเสนฺติปิ, นจฺจนฺติยาปิ นจฺจนฺติ นจฺจนฺติยาปิ คายนฺติ นจฺจนฺติยาปิ วาเทนฺติ นจฺจนฺติยาปิ ลาเสนฺติ, คายนฺติยาปิ นจฺจนฺติ คายนฺติยาปิ คายนฺติ คายนฺติยาปิ วาเทนฺติ คายนฺติยาปิ ลาเสนฺติ, วาเทนฺติยาปิ นจฺจนฺติ วาเทนฺติยาปิ คายนฺติ วาเทนฺติยาปิ วาเทนฺติ วาเทนฺติยาปิ ลาเสนฺติ, ลาเสนฺติยาปิ นจฺจนฺติ ลาเสนฺติยาปิ คายนฺติ ลาเสนฺติยาปิ วาเทนฺติ \csromanfolio{275}ลาเสนฺติยาปิ ลาเสนฺติ, อฏฺฐปเทปิ กีฬนฺติ, ทสปเทปิ กีฬนฺติ, อากาเสปิ กีฬนฺติ, ปริหาร\-ปเถปิ กีฬนฺติ, สนฺติกายปิ กีฬนฺติ, ขลิกายปิ กีฬนฺติ, ฆฏิกายปิ กีฬนฺติ, สลากหตฺเถนปิ กีฬนฺติ, อกฺเขนปิ กีฬนฺติ, ปงฺคจีเรนปิ กีฬนฺติ, วงฺกเกนปิ กีฬนฺติ, โมกฺขจิกายปิ กีฬนฺติ, จิงฺคุลเกนปิ กีฬนฺติ, ปตฺตาฬฺหเกนปิ กีฬนฺติ, รถเกนปิ กีฬนฺติ, ธนุเกนปิ กีฬนฺติ, อกฺขริกายปิ กีฬนฺติ, มเนสิกายปิ กีฬนฺติ, ยถาวชฺเชนปี กีฬนฺติ, หตฺถิสฺมิํปิ สิกฺขนฺติ, อสฺสสฺมิํปิ สิกฺขนฺติ, รถสฺมิํปิ สิกฺขนฺติ, ธนุสฺมิํปิ สิกฺขนฺติ, ถรุสฺมิํปิ สิกฺขานฺติ, หตฺถิสฺสปิ ปุรโต ธาวนฺติ, อสฺสสฺสปิ ปุรโต ธาวนฺติ, รถสฺสปิ ปุรโต\csromansharedfootnote{3c8b67f5c4ccb131}{ปุรโต ธาวนฺติ (สฺยา)} ธาวนฺติปิ อาธาวนฺติปิ, อุสฺเสเฬนฺติปิ, อปฺโผเฏนฺติปิ, นิพฺพุชฺฌนฺติปิ, มุฏฺฐีหิปิ ยุชฺฌนฺติ, รงฺคมชฺเฌปิ สงฺฆาฏิํ ปตฺถริตฺวา นจฺจกิํ\csromansharedfootnote{f605615cc9978f89}{นจฺจนฺติํ (สี, สฺยา)} เอวํ วทนฺติ “อิธ ภคินิ นจฺจสฺสู”ติ, นลาฏิกํปิ เทนฺติ, วิวิธมฺปิ อนาจารํ อาจรนฺติ.}

\proseitem{432}{เตน โข ปน สมเยน อญฺญตโร ภิกฺขุ กาสีสุ วสฺสํวุฏฺโฐ สาวตฺถิํ คจฺฉนฺโต ภควนฺตํ ทสฺสนาย เยน กีฏาคิริ ตทวสริ. อถ โข โส ภิกฺขุ ปุพฺพณฺหสมยํ นิวาเสตฺวา ปตฺต\-จีวรมา\-ทาย กีฏาคิริํ ปิณฺฑาย ปาวิสิ ปาสาทิเกน อภิกฺกนฺเตน ปฏิกฺกนฺเตน อาโลกิเตน วิโลกิเตน สมิญฺชิเตน ปสาริเตน โอกฺขิตฺตจกฺขุ อิริยาปถ\-สมฺปนฺโน. มนุสฺสา ตํ ภิกฺขุํ ปสฺสิตฺวา เอวมาหํสุ “กฺวายํ อพลพโล วิย มนฺทมนฺโท วิย ภากุฏิก\-ภากุฏิโก วิย, โก อิมสฺส อุปคตสฺส ปิณฺฑกํ ทสฺสติ, อมฺหากํ ปน อยฺยา อสฺสชิ\-ปุนพฺพสุกา สณฺหา สขิลา สุขสมฺภาสา มิหิต\-ปุพฺพงฺคมา เอหิ\-สฺวาคต\-วาทิโน อพฺภากุฏิกา อุตฺตานมุขา ปุพฺพภาสิโน, เตสํ โข นาม ปิณฺโฑ ทาตพฺโพ”ติ.}

\prose{อทฺทสา โข อญฺญตโร อุปาสโก ตํ ภิกฺขุํ กีฏาคิริสฺมิํ ปิณฺฑาย จรนฺตํ, ทิสฺวาน เยน โส ภิกฺขุ เตนุปสงฺกมิ, อุปสงฺกมิตฺวา ตํ ภิกฺขุํ อภิวาเทตฺวา เอตทโวจ “อปิ ภนฺเต ปิณฺโฑ ลพฺภตี”ติ. น โข อาวุโส ปิณฺโฑ ลพฺภตีติ. เอหิ ภนฺเต ฆรํ คมิสฺสามาติ. อถ โข โส อุปาสโก ตํ ภิกฺขุํ ฆรํ เนตฺวา โภเชตฺวา เอตทโวจ “กหํ ภนฺเต อยฺโย คมิสฺสตี”ติ. สาวตฺถิํ โข อหํ อาวุโส คมิสฺสามิ ภควนฺตํ ทสฺสนายาติ. เตน หิ ภนฺเต มม วจเนน \csromanfolio{276}ภควโต ปาเท สิรสา วนฺท, เอวญฺจ วเทหิ “ทุฏฺโฐ ภนฺเต กีฏาคิริสฺมิํ อาวาโส, อสฺสชิ\-ปุนพฺพสุกา นาม กีฏาคิริสฺมิํ อาวาสิกา อลชฺชิโน ปาปภิกฺขู, เต เอวรูปํ อนาจารํ อาจรนฺติ, มาลาวจฺฉํ โรเปนฺติปิ โรปาเปนฺติปิ, สิญฺจนฺติปิ สิญฺจาเปนฺติปิ ฯเปฯ วิวิธมฺปิ อนาจารํ อาจรนฺติ, เยปิ เต ภนฺเต มนุสฺสา ปุพฺเพ สทฺธา อเหสุํ ปสนฺนา, เตปิ เอตรหิ อสฺสทฺธา อปฺปสนฺนา, ยานิปิ ตานิ สํฆสฺส ปุพฺเพ ทานปถานิ, ตานิปิ เอตรหิ อุปจฺฉินฺนานิ, ริญฺจนฺติ เปสลา ภิกฺขู, นิวสนฺติ ปาปภิกฺขู, สาธุ ภนฺเต ภควา กีฏาคิริํ ภิกฺขู ปหิเณยฺย, ยถายํ กีฏาคิริสฺมิํ อาวาโส สณฺฐเหยฺยา”ติ.}

\prose{“เอวมาวุโส”ติ โข โส ภิกฺขุ ตสฺส อุปาสกสฺส ปฏิสฺสุตฺวา เยน สาวตฺถิ เตน ปกฺกามิ, อนุปุพฺเพน เยน สาวตฺถิ เชตวนํ อนาถ\-ปิณฺฑิกสฺส อาราโม, เยน ภควา เตนุปสงฺกมิ, อุปสงฺกมิตฺวา ภควนฺตํ อภิวาเทตฺวา เอกมนฺตํ นิสีทิ, อาจิณฺณํ โข ปเนตํ พุทฺธานํ ภควนฺตานํ อาคนฺตุเกหิ ภิกฺขูหิ สทฺธิํ ปฏิสมฺโมทิตุํ. อถ โข ภควา ตํ ภิกฺขุํ เอตทโวจ “กจฺจิ ภิกฺขุ ขมนียํ, กจฺจิ ยาปนียํ, กจฺจิสิ อปฺป\-กิลมเถน อทฺธานํ อาคโต, กุโต จ ตฺวํ ภิกฺขุ อาคจฺฉสี”ติ. ขมนียํ ภควา, ยาปนียํ ภควา, อปฺป\-กิลมเถน จาหํ ภนฺเต อทฺธานํ อาคโต, อิธาหํ ภนฺเต กาสีสุ วสฺสํวุฏฺโฐ สาวตฺถิํ อาคจฺฉนฺโต ภควนฺตํ ทสฺสนาย เยน กีฏาคิริ ตทวสริํ. อถ ขฺวาหํ ภนฺเต ปุพฺพณฺหสมยํ นิวาเสตฺวา ปตฺต\-จีวรมา\-ทาย กีฏาคิริํ ปิณฺฑย ปาวิสิํ, อทฺทสา โข มํ ภนฺเต อญฺญตโร อุปาสโก กีฏาคิริสฺมิํ ปิณฺฑาย จรนฺตํ, ทิสฺวาน เยนาหํ เตนุปสงฺกมิ, อุปสงฺกมิตฺวา มํ อภิวาเทตฺวา เอตทโวจ “อปิ ภนฺเต ปิณฺโฑ ลพฺภตี”ติ, น โข อาวุโส ปิณฺโฑ ลพฺภตีติ, เอหิ ภนฺเต ฆรํ คมิสฺสมาติ, อถ โข ภนฺเต โส อุปาสโก มํ ฆรํ เนตฺวา โภเชตฺวา เอตทโวจ “กหํ ภนฺเต อยฺโย คมิสฺสตี”ติ, สาวตฺถิํ โข อหํ อาวุโส คมิสฺสามิ ภควนฺตํ ทสฺสนายาติ, เตน หิ ภนฺเต มม วจเนน ภควโต ปาเท สิรสา วนฺท, เอวญฺจ วเทหิ “ทุฏฺโฐ ภนฺเต กีฏาคิริสฺมิํ อาวาโส, อสฺสชิ\-ปุนพฺพสุกา นาม กีฏาคิริสฺมิํ อาวาสิกา อลชฺชิโน ปาปภิกฺขู, เต เอวรูปํ อนาจารํ อาจรนฺติ, มาลาวจฺฉํ โรเปนฺติปิ โรปาเปนฺติปิ, สิญฺจนฺติปิ สิญฺจาเปนฺติปิ ฯเปฯ วิวิธมฺปิ อนาจารํ อาจรนฺติ, เยปิ เต ภนฺเต มนุสฺสา ปุพฺเพ สทฺธา \csromanfolio{277}อเหสุํ ปสนฺนา, เตปิ เอตรหิ อสฺสทฺธา อปฺปสนฺนา, ยานิปิ ตานิ สํฆสฺส ปุพฺเพ ทานปถานิ, ตานิปิ เอตรหิ อุปจฺฉินฺนานิ, ริญฺจนฺติ เปสลา ภิกฺขู, นิวสนฺติ ปาปภิกฺขู, สาธุ ภนฺเต ภควา กีฏาคิริํ ภิกฺขู ปหิเณยฺย, ยถายํ กีฏาคิริสฺมิํ อาวาโส สณฺฐเหยฺยา”ติ, ตโต อหํ ภควา อาคจฺฉามีติ.}

\proseitem{433}{อถ โข ภควา เอตสฺมิํ นิทาเน เอตสฺมิํ ปกรเณ ภิกฺขุสํฆํ สนฺนิปาตาเปตฺวา ภิกฺขู ปฏิปุจฺฉิ “สจฺจํ กิร ภิกฺขเว อสฺสชิ\-ปุนพฺพสุกา นาม กีฏาคิริสฺมิํ อาวาสิกา อลชฺชิโน ปาปภิกฺขู, เต เอวรูปํ อนาจารํ อาจรนฺติ, มาลาวจฺฉํ โรเปนฺติปิ โรปาเปนฺติปิ, สิญฺจนฺติปิ สิญฺจาเปนฺติปิ ฯเปฯ วิวิธมฺปิ อนาจารํ อาจรนฺติ, เยปิ เต ภิกฺขเว มนุสฺสา ปุพฺเพ สทฺธา อเหสุํ ปสนฺนา, เตปิ เอตรหิ อสฺสทฺธา อปฺปสนฺนา, ยานิปิ ตานิ สํฆสฺส ปุพฺเพ ทานปถานิ, ตานิปิ เอตรหิ อุปจฺฉินฺนานิ, ริญฺจนฺติ เปสลา ภิกฺขู, นิวสนฺติ ปาปภิกฺขู”ติ. สจฺจํ ภควาติ. วิครหิ พุทฺโธ ภควา ฯเปฯ. กถํ หิ นาม เต ภิกฺขเว โมฆปุริสา เอวรูปํ อนาจารํ อาจริสฺสนฺติ, มาลาวจฺฉํ โรเปสฺสนฺติปิ โรปาเปสฺสนฺติปิ, สิญฺจิสฺสนฺติปิ สิญฺจาเปสฺสนฺติปิ, โอจินิสฺสนฺติปิ โอจินาเปสฺสนฺติปิ, คนฺถิสฺสนฺติปิ คนฺถาเปสฺสนฺติปิ, เอกโตวณฺฏิกมาลํ กริสฺสนฺติปิ การาเปสฺสนฺติปิ, อุภโตวณฺฏิกมาลํ กริสฺสนฺติปิ การาเปสฺสนฺติปิ, มญฺชริกํ กริสฺสนฺติปิ การาเปสฺสนฺติปิ, วิธูติกํ กริสฺสนฺติปิ การาเปสฺสนฺติปิ, วฏํสกํ กริสฺสนฺติปิ การาเปสฺสนฺติปิ, อาเวฬํ กริสฺสนฺติปิ การาเปสฺสนฺติปิ, อุรจฺฉทํ กริสฺสนฺติปิ การาเปสฺสนฺติปิ. เต กุลิตฺถีนํ กุลธีตานํ กุลกุมารีนํ กุลสุณฺหานํ กุลทาสีนํ เอกโตวณฺฏิกมาลํ หริสฺสนฺติปิ หาราเปสฺสนฺติปิ, อุภโตวณฺฏิกมาลํ หริสฺสนฺติปิ หาราเปสฺสนฺติปิ, มญฺชริกํ หริสฺสนฺติปิ หาราเปสฺสนฺติปิ, วิธูติกํ หริสฺสนฺติปิ หาราเปสฺสนฺติปิ, วฏํสกํ หริสฺสนฺติปิ หาราเปสฺสนฺติปิ, อาเวฬํ หริสฺสนฺติปิ หาราเปสฺสนฺติปิ, อุรจฺฉทํ หริสฺสนฺติปิ หาราเปสฺสนฺติปิ. เต กุลิตฺถีหิ กุลธีตาหิ กุลกุมารีหิ กุลสุณฺหาหิ กุลทาสีหิ สทฺธิํ เอกภาชเนปิ ภุญฺชิสฺสนฺติ, เอกถาลเกปิ ปิวิสฺสนฺติ, เอกาสเนปิ นิสีทิสฺสนฺติ, เอกมญฺเจปิ ตุวฏฺฏิสฺสนฺติ, เอกตฺถรณาปิ ตุวฏฺฏิสฺสนฺติ, เอกปาวุรณาปิ ตุวฏฺฏิสฺสนฺติ, เอกตฺถรณาปาวุรณาปิ ตุวฏฺฏิสฺสนฺติ, วิกาเลปิ ภุญฺชิสฺสนฺติ, มชฺชมฺปิ ปิวิสฺสนฺติ, มาลา\-คนฺธ\-วิเลปนมฺปิ ธาริสฺสนฺติ, นจฺจิสฺสนฺติปิ \csromanfolio{278}คายิสฺสนฺติปิ วาเทสฺสนฺติปิ ลาเสสฺสนฺติปิ, นจฺจนฺติยาปิ นจฺจิสฺสนฺติ นจฺจนฺติยาปิ คายิสฺสนฺติ นจฺจนฺติยาปิ วาเทสฺสนฺติ นจฺจนฺติยาปิ ลาเสสฺสนฺติ, คายนฺติยาปิ นจฺจิสฺสนฺติ คายนฺติยาปิ คายิสฺสนฺติ คายนฺติยาปิ วาเทสฺสนฺติ คายนฺติยาปิ ลาเสสฺสนฺติ, วาเทนฺติยาปิ นจฺจิสฺสนฺติ วาเทนฺติยาปิ คายิสฺสนฺติ วาเทนฺติยาปิ วาเทสฺสนฺติ วาเทนฺติยาปิ ลาเสสฺสนฺติ, ลาเสนฺติยาปิ นจฺจิสฺสนฺติ ลาเสนฺติยาปิ คายิสฺสนฺติ ลาเสนฺติยาปิ วาเทสฺสนฺติ ลาเสนฺติยาปิ ลาเสสฺสนฺติ, อฏฺฐปเทปิ กีฬิสฺสนฺติ, ทสปเทปิ กีฬิสฺสนฺติ, อากาเสปิ กีฬิสฺสนฺติ, ปริหาร\-ปเถปิ กีฬิสฺสนฺติ, สนฺติกายปิ กีฬิสฺสนฺติ, ขลิกายปิ กีฬิสฺสนฺติ, ฆฏิกายปิ กีฬิสฺสนฺติ, สลากหตฺเถนปิ กีฬิสฺสนฺติ, อกฺเขนปิ กีฬิสฺสนฺติ, ปงฺคจีเรนปิ กีฬิสฺสนฺติ, วงฺกเกนปิ กีฬิสฺสนฺติ, โมกฺขจิกายปิ กีฬิสฺสนฺติ, จิงฺคุลเกนปิ กีฬิสฺสนฺติ, ปตฺตาฬฺหเกนปิ กีฬิสฺสนฺติ, รถเกนปิ กีฬิสฺสนฺติ, ธนุเกนปิ กีฬิสฺสนฺติ, อกฺขริกายปิ กีฬิสฺสนฺติ, มเนสิกายปิ กีฬิสฺสนฺติ, ยถาวชฺเชนปิ กีฬิสฺสนฺติ, หตฺถิสฺมิํปิ สิกฺขิสฺสนฺติ, อสฺสสฺมิํปิ สิกฺขิสฺสนฺติ, รถสฺมิํปิ สิกฺขิสฺสนฺติ, ธนุสฺมิํปิ สิกฺขิสฺสนฺติ, ถรุสฺมิํปิ สิกฺขิสฺสนฺติ, หตฺถิสฺสปิ ปุรโต ธาวิสฺสนฺติ, อสฺสสฺสปิ ปุรโต ธาวิสฺสนฺติ, รถสฺสปิ ปุรโต\csromansharedfootnote{5947dce77d2f62c8}{ปุรโต ธาวิสฺสนฺติ (สฺยา)} ธาวิสฺสนฺติปิ อาธาวิสฺสนฺติปิ, อุสฺเสเฬสฺสนฺติปิ, อปฺโผเฏสฺสนฺติปิ, นิพฺพุชฺฌิสฺสนฺติปิ, มุฏฺฐีหิปิ ยุชฺฌิสฺสนฺติ, รงฺคมชฺเฌปิ สงฺฆาฏิํ ปตฺถริตฺวา นจฺจกิํ เอวํ วกฺขนฺติ “อิธ ภคินิ นจฺจสฺสู”ติ, นลาฏิกมฺปิ ทสฺสนฺติ, วิวิธมฺปิ อนาจารํ อาจริสฺสนฺติ, เนตํ ภิกฺขเว อปฺปสนฺนานํ วา ปสาทาย ฯเปฯ วิครหิตฺวา ธมฺมิํ กถํ กตฺวา สาริปุตฺต\-โมคฺคลฺลาเน อามนฺเตสิ “คจฺฉถ ตุเมฺห สาริปุตฺตา กีฏาคิริํ คนฺตฺวา อสฺสชิ\-ปุนพฺพสุกานํ ภิกฺขูนํ กีฏาคิริสฺมา ปพฺพาชนีย\-กมฺมํ กโรถ ตุมฺหากํ เอเต สทฺธิวิหาริกา”ติ.}

\prose{กถํ มยํ ภนฺเต อสฺสชิ\-ปุนพฺพสุกานํ ภิกฺขูนํ กีฏาคิริสฺมา ปพฺพาชนีย\-กมฺมํ กโรม จณฺฑา เต ภิกฺขู ผรุสาติ. เตน หิ ตุเมฺห สาริปุตฺตา พหุเกหิ ภิกฺขูหิ สจฺฉิํ คจฺฉถาติ. เอวํ ภนฺเตติ โข สาริปุตฺต\-โมคฺคลฺลานา ภควโต ปจฺจสฺโสสุํ. เอวญฺจ ปน ภิกฺขเว กาตพฺพํ, ปฐมํ อสฺสชิ\-ปุนพฺพสุกา ภิกฺขู โจเทตพฺพา, โจเทตฺวา สาเรตพฺพา, สาเรตฺวา อาปตฺติํ โรเปตพฺพา, อาปตฺติํ โรเปตฺวา พฺยตฺเตน ภิกฺขุนา ปฏิพเลน สํโฆ ญาเปตพฺโพ\csromandash{}}

\proseitem{434}{\csromanfolio{279}“สุณาตุ เม ภนฺเต สํโฆ, อิเม อสฺสชิ\-ปุนพฺพสุกา ภิกฺขู กุลทูสกา ปาปสมาจารา, อิเมสํ ปาปกา สมาจารา ทิสฺสนฺติ เจว สุยฺยนฺติ จ, กุลานิ จ อิเมหิ ทุฏฺฐานิ ทิสฺสนฺติ เจว สุยฺยนฺติ จ, ยทิ สํฆสฺส ปตฺตกลฺลํ, สํโฆ อสฺสชิ\-ปุนพฺพสุกานํ ภิกฺขูนํ กีฏาคิริสฺมา ปพฺพาชนีย\-กมฺมํ กเรยฺย น อสฺสชิปุนพฺพสุเกหิ ภิกฺขูหิ กีฏคิริสฺมิํ วตฺถพฺพนฺติ, เอสา ญตฺติ.}

\prose{สุณาตุ เม ภนฺเต สํโฆ, อิเม อสฺสชิ\-ปุนพฺพสุกา ภิกฺขู กุลทูสกา ปาปสมาจารา, อิเมสํ ปาปกา สมาจารา ทิสฺสนฺติ เจว สุยฺยนฺติ จ, กุลานิ จ อิเมหิ ทุฏฺฐานิ ทิสฺสนฺติ เจว สุยฺยนฺติ จ, สํโฆ อสฺสชิ\-ปุนพฺพสุกานํ ภิกฺขูนํ กีฏาคิริสฺมา ปพฺพาชนีย\-กมฺมํ กโรติ น อสฺสชิปุนพฺพสุเกหิ ภิกฺขูหิ กีฏาคิริสฺมิํ วตฺถพฺพนฺติ, ยสฺสายสฺมตฺโต ขมติ อสฺสชิ\-ปุนพฺพสุกานํ ภิกฺขูนํ กีฏาคิริสฺมา ปพฺพาชนีย\-กมฺมสฺส กรณํ น อสฺสชิปุนพฺพสุเกหิ ภิกฺขูหิ กีฏาคิริสฺมิํ วตฺถพฺพนฺติ, โส ตุณฺหสฺส, ยสฺส นกฺขมติ, โส ภาเสยฺย.}

\prose{ทุติยมฺปิ เอตมตฺถํ วทามิ ฯเปฯ. ตติยมฺปิ เอตมตฺถํ วทามิ. สุณาตุ เม ภนฺเต สํโฆ ฯเปฯ โส ภาเสยฺย.}

\prose{กตํ สํเฆน อสฺสชิ\-ปุนพฺพสุกานํ ภิกฺขูนํ กีฏาคิริสฺมา ปพฺพาชนีย\-กมฺมํ ‘น อสฺสชิปุนพฺพสุเกหิ ภิกฺขูหิ กีฏาคิริสฺมิํ วตฺถพฺพนฺ’ติ, ขมติ สํฆสฺส, ตสฺมา ตุณฺหี, เอวเมตํ ธารยามี”ติ.}

\csromanmatikaanchor{mtk.99}
\csromanmatikaanchor{mtk.100}
\csromantocmarkref{subsection}{ปญฺญตฺติ}{mtk.99}
\bookmark[dest={mtk.99},level=2]{ปญฺญตฺติ}
\csromantocmarkref{subsection}{สิกฺขาปทวิภงฺค, ปทภาชนีย}{mtk.100}
\bookmark[dest={mtk.100},level=2]{สิกฺขาปทวิภงฺค, ปทภาชนีย}
\proseitem{435}{อถ โข สาริปุตฺต\-โมคฺคลฺลาน\-ปฺปมุโข ภิกฺขุสํโฆ กีฏาคิริํ คนฺตฺวา อสฺสชิ\-ปุนพฺพสุกานํ ภิกฺขูนํ กีฏาคิริสฺมา ปพฺพาชนีย\-กมฺมํ อกาสิ ‘น อสฺสชิปุนพฺพสุเกหิ ภิกฺขูหิ กีฏาคิริสฺมิํ วตฺถพฺพนฺ’ติ. เต สํเฆน ปพฺพาชนีย\-กมฺมกตา น สมฺมา วตฺตนฺติ, น โลมํ ปาเตนฺติ, น เนตฺถารํ วตฺตนฺติ, น ภิกฺขู ขมาเปนฺติ, อกฺโกสนฺติ, ปริภาสนฺติ, ฉนฺทคามิตา โทสคามิตา โมหคามิตา ภยคามิตา ปาเปนฺติ, ปกฺกมนฺติปิ วิพฺภมนฺติปิ. เย เต ภิกฺขู อปฺปิจฺฉา ฯเปฯ เต อุชฺฌายนฺติ ขิยฺยนฺติ วิปาเจนฺติ “กถํ หิ นาม อสฺสชิ\-ปุนพฺพสุกา ภิกฺขู สํเฆน ปพฺพาชนีย\-กมฺมกตา น สมฺมา วตฺติสฺสนฺติ, น โลมํ ปาเตสฺสนฺติ, น เนตฺถารํ วตฺติสฺสนฺติ, น ภิกฺขู ขมาเปสฺสนฺติ, อกฺโกสิสฺสนฺติ, ปริภาสิสฺสนฺติ, ฉนฺทคามิตา โทสคามิตา โมหคามิตา ภยคามิตา ปาเปสฺสนฺติ, ปกฺกมิสฺสนฺติปิ \csromanfolio{280}วิพฺภมิสฺสนฺติปี”ติ. อถ โข เต ภิกฺขู อสฺสชิปุนพฺพสุเก ภิกฺขู อเนก\-ปริยาเยน วิครหิตฺวา ภควโต เอตมตฺถํ อาโรเจสุํ ฯเปฯ “สจฺจํ กิร ภิกฺขเว อสฺสชิ\-ปุนพฺพสุกา ภิกฺขู สํเฆน ปพฺพาชนีย\-กมฺมกตา น สมฺมา วตฺตนฺติ ฯเปฯ วิพฺภมนฺติปี”ติ. สจฺจํ ภควาติ. วิครหิ พุทฺโธ ภควา ฯเปฯ. เอวญฺจ ปน ภิกฺขเว อิมํ สิกฺขาปทํ อุทฺทิเสยฺยาถ\csromandash{}}

\proseitem{436}{\textbf{“ภิกฺขุ ปเนว อญฺญตรํ คามํ วา นิคมํ วา อุปนิสฺสาย วิหรติ กุลทูสโก ปาปสมาจาโร, ตสฺส โข ปาปกา สมาจารา ทิสฺสนฺติ เจว สุยฺยนฺติ จ, กุลานิ จ เตน ทุฏฺฐานิ ทิสฺสนฺติ เจว สุยฺยนฺติ จ, โส ภิกฺขุ ภิกฺขูหิ เอวมสฺส วจนีโย ‘อายสฺมา โข กุลทูสโก ปาปสมาจาโร, อายสฺมโต โข ปาปกา สมาจารา ทิสฺสนฺติ เจว สุยฺยนฺติ จ, กุลานิ จายสฺมตา ทุฏฺฐานิ ทิสฺสนฺติ เจว สุยฺยนฺติ จ, ปกฺกมตายสฺมา อิมมฺหา อาวาสา, อลํ เต อิธ วาเสนา’ติ. เอวญฺจ โส ภิกฺขุ ภิกฺขูหิ วุจฺจมาโน เต ภิกฺขู เอวํ วเทยฺย ‘ฉนฺทคามิโน จ ภิกฺขู โทสคามิโน จ ภิกฺขู โมหคามิโน จ ภิกฺขู ภยคามิโน จ ภิกฺขู ตาทิสิกาย อาปตฺติยา เอกจฺจํ ปพฺพาเชนฺติ เอกจฺจํ น ปพฺพาเชนฺตี’ติ. โส ภิกฺขุ ภิกฺขูหิ เอวมสฺส วจนีโย ‘มายสฺมา เอวํ อวจ, น จ ภิกฺขู ฉนฺทคามิโน น จ ภิกฺขู โทสคามิโน น จ ภิกฺขู โมหคามิโน น จ ภิกฺขู ภยคามิโน, อายสฺมา โข กุลทูสโก ปาปสมาจาโร, อายสฺมโต โข ปาปกา สมาจารา ทิสฺสนฺติ เจว สุยฺยนฺติจ, กุลานิ จายสฺมตา ทุฏฺฐานิ ทิสฺสนฺติ เจว สุยฺยนฺติจ, ปกฺกมตายสฺมา อิมมฺหา อาวาสา, อลํ เต อิธ วาเสนา’ติ. เอวญฺจ โส ภิกฺขุ ภิกฺขูหิ วุจฺจมาโน ตเถว ปคฺคเณฺหยฺย, โส ภิกฺขุ ภิกฺขูหิ ยาวตติยํ สมนุภาสิตพฺโพ ตสฺส ปฏินิสฺสคฺคาย, ยาวตติยญฺเจ สมนุ}\-\textbf{ภาสีย}\-\textbf{มาโน ตํ ปฏินิสฺสชฺเชยฺย, อิจฺเจตํ กุสลํ, โน เจ ปฏินิสฺสชฺเชยฺย, สํฆาทิเสโส”}ติ.}

\proseitem{437}{\textbf{ภิกฺขุ ปเนว อญฺญตรํ คามํ วา นิคมํ} \textbf{วา}ติ คาโมปิ นิคโมปิ นครํปิ คาโม เจว นิคโม จ.}

\prose{\textbf{อุปนิสฺสาย วิหรตี}ติ ตตฺถ ปฏิพทฺธา จีรปิณฺฑปาต\-เสนาสน\-คิลานปฺปจฺจย\-เภสชฺช\-ปริกฺขารา\csromansharedfootnote{cf2a2be673868a02}{จีร- ปาโฐ, จีวร- \csromandash{} ม.พ.ป.}.}

\prose{\csromanfolio{281}\textbf{กุลํ} นาม จตฺตาริ กุลานิ ขตฺติยกุลํ พฺราหฺมณกุลํ เวสฺสกุลํ สุทฺทกุลํ.}

\prose{\textbf{กุลทูสโก}ติ กุลานิ ทูเสติ ปุปฺเผน วา ผเลน วา จุณฺเณน วา มตฺติกาย วา ทนฺตกฏฺเฐน วา เวฬุยา วา\csromansharedfootnote{59198a94e207834d}{เวฬุนา วา (สฺยา)} เวชฺชิกาย วา ชงฺฆ\-เปสนิ\-เกน วา.}

\prose{\textbf{ปาปสมาจาโร}ติ มาลาวจฺฉํ โรเปติปิ, โรปาเปติปิ, สิญฺจติปิ, สิญฺจาเปติปิ, โอจินาติปิ, โอจินาเปติปิ, คนฺเถติปิ, คนฺถาเปติปิ.}

\prose{\textbf{ทิสฺสนฺติ เจว สุยฺยนฺติ จา}ติ เย สํมุขา เต ปสฺสนฺติ, เย ติโรกฺขา เต สุณนฺติ.}

\prose{\textbf{กุลานิ จ เตน ทุฏฺฐานี}ติ ปุพฺเพ สทฺธา หุตฺวา ตํ อาคมฺม อสฺสทฺธา โหนฺติ, ปสนฺนา หุตฺวา อปฺปสนฺนา โหนฺติ.}

\prose{\textbf{ทิสฺสนฺติ เจว สุยฺยนฺติ จา}ติ เย สํมุขา เต ปสฺสนฺติ, เย ติโรกฺขา เต สุณนฺติ.}

\prose{\textbf{โส ภิกฺขู}ติ โย โส กุลทูลโก ภิกฺขุ.}

\prose{\textbf{ภิกฺขูหี}ติ อญฺเญหิ ภิกฺขูหิ. เย ปสฺสนฺติ, เย สุณนฺติ, เตหิ วตฺตพฺโพ “อายสฺมา โข กุลทูสโก ปาปสมาจาโร, อายสฺมโต โข ปาปกา สมาจารา ทิสฺสนฺติ เจว สุยฺยนฺติ จ, กุลานิ จายสฺมตา ทุฏฺฐานิ ทิสฺสนฺติ เจว สุยฺยนฺติ จ, ปกฺกมตายสฺมา อิมมฺหา อาวาสา, อลํ เต อิธ วาเสนา”ติ.}

\prose{เอวญฺจ โส ภิกฺขุ ภิกฺขูหิ วุจฺจมาโน เต ภิกฺขู เอวํ วเทยฺย “ฉนฺทคามิโน จ ภิกฺขู โทสคามิโน จ ภิกฺขู โมหคามิโน จ ภิกฺขู ภยคามิโน จ ภิกฺขู ตาทิสิกาย อาปตฺติยา เอกจฺจํ ปพฺพาเชนฺติ เอกจฺจํ น ปพฺพาเชนฺตี”ติ.}

\prose{\textbf{โส ภิกฺขู}ติ โย โส กมฺมกโต ภิกฺขุ.}

\prose{\textbf{ภิกฺขูหี}ติ อญฺเญหิ ภิกฺขูหิ. เย ปสฺสนฺติ, เย สุณนฺติ, เตหิ วตฺตพฺโพ “มายสฺมา เอวํ อวจ, น จ ภิกฺขู ฉนฺทคามิโน น จ ภิกฺขู โทสคามิโน น จ ภิกฺขู โมหคามิโน น จ ภิกฺขู ภยคามิโน, \csromanfolio{282}อายสฺมา โข กุลทูสโก ปาปสมาจาโร, อายสฺมโต โข ปาปกา สมาจารา ทิสฺสนฺติ เจว สุยฺยนฺติ จ, กุลานิ จายสฺมตา ทุฏฺฐานิ ทิสฺสนฺติ เจว สุยฺยนฺติ จ, ปกฺกมตายสฺมา อิมมฺหา อาวาสา, อลํ เต อิธ วาเสนา”ติ. ทุติยมฺปิ วตฺตพฺโพ. ตติยมฺปิ วตฺตพฺโพ.}

\prose{สเจ ปฏินิสฺสชฺชติ, อิจฺเจตํ กุสลํ. โน เจ ปฏินิสฺสชฺชติ, อาปตฺติ ทุกฺกฏสฺส. สุตฺวา น วทนฺติ, อาปตฺติ ทุกฺกฏสฺส. โส ภิกฺขุ สํฆมชฺฌมฺปิ อากฑฺฒิตฺวา วตฺตพฺโพ “มายสฺมา เอวํ อวจ, น จ ภิกฺขู ฉนฺทคามิโน น จ ภิกฺขู โทสคามิโน น จ ภิกฺขู โมหคามิโน น จ ภิกฺขู ภยคามิโน, อายสฺมา โข กุลทูสโก ปาปสมาจาโร, อายสฺมโต โข ปาปกา สมาจารา ทิสฺสนฺติ เจว สุยฺยนฺติ จ, กุลานิ จายสฺมตา ทุฏฺฐานิ ทิสฺสนฺติ เจว สุยฺยนฺติ จ, ปกฺกมตายสฺมา อิมมฺหา อาวาสา, อลํ เต อิธ วาเสนา”ติ. ทุติยมฺปิ วตฺตพฺโพ. ตติยมฺปิ วตฺตพฺโพ. สเจ ปฏินิสฺสชฺชติ, อิจฺเจตํ กุสลํ. โน เจ ปฏินิสฺสชฺชติ, อาปตฺติ ทุกฺกฏสฺส. โส ภิกฺขุ สมนุภาสิตพฺโพ, เอวญฺจ ปน ภิกฺขเว สมนุภาสิตพฺโพ, พฺยตฺเตน ภิกฺขุนา ปฏิพเลน สํโฆ ญาเปตพฺโพ\csromandash{}}

\proseitem{438}{“สุณาตุ เม ภนฺเต สํโฆ, อยํ อิตฺถนฺนาโม ภิกฺขุ สํเฆน ปพฺพาชนีย\-กมฺมกโต ภิกฺขู ฉนฺทคามิตา โทสคามิตา โมหคามิตา ภยคามิตา ปาเปติ, โส ตํ วตฺถุํ น ปฏินิสฺสชฺชติ, ยทิ สํฆสฺส ปตฺตกลฺลํ, สํโฆ อิตฺถนฺนามํ ภิกฺขุํ สมนุภาเสยฺย ตสฺส วตฺถุสฺส ปฏินิสฺสคฺคาย, เอสา ญตฺติ.}

\prose{สุณาตุ เม ภนฺเต สํโฆ, อยํ อิตฺถนฺนาโม ภิกฺขุ สํเฆน ปพฺพาชนีย\-กมฺมกโต ภิกฺขู ฉนฺทคามิตา โทสคามิตา โมหคามิตา ภยคามิตา ปาเปติ, โส ตํ วตฺถุํ น ปฏินิสฺสชฺชติ, สํโฆ อิตฺถนฺนามํ ภิกฺขุํ สมนุภาสติ ตสฺส วตฺถุสฺส ปฏินิสฺสคฺคาย, ยสฺสายสฺมโต ขมติ อิตฺถนฺนามสฺส ภิกฺขโน สมนุภาสนา ตสฺส วตฺถุสฺส ปฏินิสฺสคฺคาย, โส ตุณฺหสฺส, ยสฺส นกฺขมติ, โส ภาเสยฺย.}

\prose{ทุติยมฺปิ เอตมตฺถํ วทามิ ฯเปฯ. ตติยมฺปิ เอตมตฺถํ วทามิ ฯเปฯ.}

\prose{สมนุภฏฺโฐ สํเฆน อิตฺถนฺนาโม ภิกฺขุ ตสฺส วตฺถุสฺส ปฏินิสฺสคฺคาย, ขมติ สํฆสฺส, ตสฺมา ตุณฺหี, เอวเมตํ ธารยามี”ติ.}

\proseitem{439}{\csromanfolio{283}ญตฺติยา ทุกฺกฏํ, ทฺวีหิ กมฺมวาจาหิ ถุลฺลจฺจยา, กมฺมวาจา\-ปริโยสาเน อาปตฺติ สํฆาทิเสสสฺส. สํฆาทิเสสํ อชฺฌาปชฺชนฺตสฺส ญตฺติยา ทุกฺกฏํ ทฺวีหิ กมฺมวาจาหิ ถุลฺลจฺจยา ปฏิปฺปสฺสมฺภนฺติ.}

\prose{\textbf{สํฆาทิเสโส}ติ สํโฆว ตสฺสา อาปตฺติยา ปริวาสํ เทติ มูลาย ปฏิกสฺสติ มานตฺตํ เทติ อพฺเภติ, น สมฺพหุลา, น เอกปุคฺคโล, เตน วุจฺจติ “สํฆาทิเสโส”ติ. ตสฺเสว อาปตฺติ\-นิกายสฺส นามกมฺมํ อธิวจนํ, เตนปิ วุจฺจติ “สํฆาทิเสโส”ติ.}

\proseitem{440}{ธมฺมกมฺเม ธมฺมกมฺม\-สญฺญี น ปฏินิสฺสชฺชติ, อาปตฺติ สํฆาทิเสสสฺส.}

\prose{ธมฺมกมฺเม เวมติโก น ปฏินิสฺสชฺชติ, อาปตฺติ สํฆาทิเสสสฺส.}

\prose{ธมฺมกมฺเม อธมฺมกมฺม\-สญฺญี น ปฏินิสฺสชฺชติ, อาปตฺติ สํฆาทิเสสสฺส.}

\prose{อธมฺมกมฺเม ธมฺมกมฺม\-สญฺญี, อาปตฺติ ทุกฺกฏสฺส.}

\prose{อธมฺมกมฺเม เวมติโก, อาปตฺติ ทุกฺกฏสฺส.}

\prose{อธมฺมกมฺเม อธมฺมกมฺม\-สญฺญี, อาปตฺติ ทุกฺกฏสฺส.}

\proseitem{441}{อนาปตฺติ อสมนุภาสนฺตสฺส ปฏินิสฺสชฺชนฺตสฺส อุมฺมตฺตกสฺส อาทิกมฺมิกสฺสาติ.}

\vspace{\tipitakaparskip}
\csromancenter{กุลทูสก\-สิกฺขาปทํ นิฏฺฐิตํ เตรสมํ.}
\proseitemwithcloser{442}{อุทฺทิฏฺฐา โข อายสฺมนฺโต เตรส สํฆาทิเสสา ธมฺมา นว ปฐมาปตฺติกา จตฺตาโร ยาวตติยกา, เยสํ ภิกฺขุ อญฺญตรํ วา อญฺญตรํ วา อาปชฺชิตฺวา ยาวตีหํ ชานํ ปฏิจฺฉาเทติ, ตาวตีหํ เตน ภิกฺขุนา อกามา ปริวตฺถพฺพํ. ปริวุตฺถ\-ปริวาเสน ภิกฺขุนา อุตฺตริ ฉารตฺตํ ภิกฺขุ\-มานตฺตาย ปฏิปชฺชิตพฺพํ. จิณฺณมานตฺโต ภิกฺขุ ยตฺถ สิยา วีสติคโณ ภิกฺขุสํโฆ, ตตฺถ โส ภิกฺขุ อพฺเภตพฺโพ. เอเกนปิ เจ อูโน วีสติคโณ ภิกฺขุสํโฆ ตํ ภิกฺขุํ อพฺเภยฺย, โส จ ภิกฺขุ อนพฺภิโต, เต จ ภิกฺขู คารยฺหา, อยํ ตตฺถ สามีจิ, ตตฺถายสฺมนฺเต ปุจฺฉามิ กจฺจิตฺถ ปริสุทฺธา, ทุติยมฺปิ ปุจฺฉามิ กจฺจิตฺถ ปริสุทฺธา, ตติยมฺปิ ปุจฺฉามิ กจฺจิตฺถ ปริสุทฺธา, ปริสุทฺเธ\-ตฺถา\-ยสฺมนฺโต, ตสฺมา ตุณฺหี, เอวเมตํ ธารยามีติ.}{เตรสกํ นิฏฺฐิตํ.}
\csromancenter{ตสฺสุทฺทานํ}
\csromanmatikaanchor{mtk.101}
\csromantocmarkref{section}{อุทฺทานคาถา}{mtk.101}
\bookmark[dest={mtk.101},level=1]{อุทฺทานคาถา}
\setlength{\csromangathapagewidth}{0pt}
\setlength{\csromangathawakwidth}{0pt}
\csromangathameasuregroup{วิสฺสฏฺฐิ กายสํสคฺคํ, \\ สญฺจริตฺตํ กุฏี เจว, \\ กิญฺจิเลสญฺจ เภโท จ, \\ ทุพฺพจํ กุลทูสญฺจ,}{\csromangathabat{วิสฺสฏฺฐิ กายสํสคฺคํ,}{ทุฏฺฐุลฺลํ อตฺตกามญฺจ.} \\ \csromangathabat{สญฺจริตฺตํ กุฏี เจว,}{วิหาโร จ อมูลกํ.} \\ \csromangathabat{กิญฺจิเลสญฺจ เภโท จ,}{ตสฺเสว อนุวตฺตกา.} \\ \csromangathabat{ทุพฺพจํ กุลทูสญฺจ,}{สํฆาทิเสสา เตรสาติ.}}
\csromangathasetleft{วิสฺสฏฺฐิ กายสํสคฺคํ, \\ สญฺจริตฺตํ กุฏี เจว, \\ กิญฺจิเลสญฺจ เภโท จ, \\ ทุพฺพจํ กุลทูสญฺจ,}
\csromangathagroupwithclosermajor{\csromangathastanza{\csromanfolio{284}\csromangathabat{วิสฺสฏฺฐิ กายสํสคฺคํ,}{ทุฏฺฐุลฺลํ อตฺตกามญฺจ.} \\ \csromangathabat{สญฺจริตฺตํ กุฏี เจว,}{วิหาโร จ อมูลกํ.}}\csromangathastanza{\csromangathabat{กิญฺจิเลสญฺจ เภโท จ,}{ตสฺเสว อนุวตฺตกา.} \\ \csromangathabat{ทุพฺพจํ กุลทูสญฺจ,}{สํฆาทิเสสา เตรสาติ.}}}{\textbf{สํฆาทิเสสกณฺฑํ นิฏฺฐิตํ.}}
\csromanmatikaanchor{mtk.102}
\csromantocmarknopage{chapter}{3. อนิยตกณฺฑ}{mtk.102}
\bookmark[dest={mtk.102},level=0]{3. อนิยตกณฺฑ}
\chapterheadpageread{3. \textbf{อนิยตกณฺฑ}}
\csromanheader{1}{1. ปฐม\-อนิยต\-สิกฺขาปท}
\vspace{\tipitakaparskip}
\csromancenter{อิเม โข ปนายสฺมนฺโต เทฺว อนิยตา ธมฺมา \\ อุทฺเทสํ อาคจฺฉนฺติ.}
\csromanmatikaanchor{mtk.103}
\csromanmatikaanchor{mtk.104}
\csromantocmarknopage{section}{1. ปฐมอนิยตสิกฺขาปท}{mtk.103}
\bookmark[dest={mtk.103},level=1]{1. ปฐมอนิยตสิกฺขาปท}
\csromantocmarkref{subsection}{อุทายิภิกฺขุวตฺถุ}{mtk.104}
\bookmark[dest={mtk.104},level=2]{อุทายิภิกฺขุวตฺถุ}
\proseitem{443}{\csromanfolio{285}เตน สมเยน พุทฺโธ ภควา สาวตฺถิยํ วิหรติ เชตวเน อนาถ\-ปิณฺฑิกสฺส อาราเม. เตน โข ปน สมเยน อายสฺมา อุทายี สาวตฺถิยํ กุลูปโก โหติ, พหุกานิ กุลานิ อุปสงฺกมติ. เตน โข ปน สมเยน อายสฺมโต อุทายิสฺส อุปฏฺฐาก\-กุลสฺส กุมาริกา อญฺญตรสฺส กุลสฺส กุมารกสฺส ทินฺนา โหติ. อถ โข อายสฺมา อุทายี ปุพฺพณฺหสมยํ นิวาเสตฺวา ปตฺต\-จีวรมา\-ทาย เยน ตํ กุลํ เตนุปสงฺกมิ, อุปสงฺกมิตฺวา มนุสฺเส ปุจฺฉิ “กหํ อิตฺถนฺนามา”ติ. เต เอวมาหํสุ “ทินฺนา ภนฺเต อมุกสฺส กุลสฺส กุมารกสฺสา”ติ. ตํปิ โข กุลํ อายสฺมโต อุทายิสฺส อุปฏฺฐากํ โหติ. อถ โข อายสฺมา อุทายี เยน ตํ กุลํ เตนุปสงฺกมิ, อุปสงฺกมิตฺวา มนุสฺเส ปุจฺฉิ “กหํ อิตฺถนฺนามา”ติ. เต เอวมาหํสุ “เอสายฺย โอวรเก นิสินฺนา”ติ. อถ โข อายสฺมา อุทายี เยน สา กุมาริกา เตนุปสงฺกมิ, อุปสงฺกมิตฺวา ตสฺสา กุมาริกาย สทฺธิํ เอโก เอกาย รโห ปฏิจฺฉนฺเน อาสเน อลํกมฺมนิเย นิสชฺชํ กปฺเปสิ กาลยุตฺตํ สมุลฺลปนฺโต กาลยุตฺตํ ธมฺมํ ภณนฺโต.}

\csromanmatikaanchor{mtk.105}
\csromanmatikaanchor{mtk.106}
\csromantocmarkref{subsection}{ปญฺญตฺติ}{mtk.105}
\bookmark[dest={mtk.105},level=2]{ปญฺญตฺติ}
\csromantocmarkref{subsection}{สิกฺขาปทวิภงฺค, ปทภาชนีย}{mtk.106}
\bookmark[dest={mtk.106},level=2]{สิกฺขาปทวิภงฺค, ปทภาชนีย}
\prose{เตน โข ปน สมเยน วิสาขา มิคารมาตา พหุปุตฺตา โหติ พหุนตฺตา อโรคปุตฺตา อโรคนตฺตา อภิมงฺคลสมฺมตา, มนุสฺสา ยญฺเญสุ ฉเณสุ อุสฺสเวสุ วิสาขํ มิคารมาตรํ ปฐมํ โภเชนฺติ. อถ โข วิสาขา มิคารมาตา นิมนฺติตา ตํ กุลํ อคมาสิ. อทฺทสา โข วิสาขา มิคารมาตา อายสฺมนฺตํ อุทายิํ ตสฺสา กุมาริกาย สทฺธิํ เอกํ เอกาย รโห ปฏิจฺฉนฺเน อาสเน อลํกมฺมนิเย นิสินฺนํ, ทิสฺวาน อายสฺมนฺตํ อุทายิํ เอตทโวจ “อิทํ ภนฺเต นจฺฉนฺนํ นปฺปติรูปํ, ยํ อยฺโย มาตุคาเมน สทฺธิํ เอโก เอกาย รโห ปฏิจฺฉนฺเน \csromanfolio{286}อาสเน อลํกมฺมนิเย นิสชฺชํ กปฺเปติ, กิญฺจาปิ ภนฺเต อยฺโย อนตฺถิโก เตน ธมฺเมน, อปิ จ ทุสฺสทฺธาปยา อปฺปสนฺนา มนุสฺสา”ติ. เอวํปิ โข อายสฺมา อุทายี วิสาขาย มิคารมาตุยา วุจฺจมาโน นาทิยิ. อถ โข วิสาขา มิคารมาตา นิกฺขมิตฺวา ภิกฺขูนํ เอตมตฺถํ อาโรเจสิ. เย เต ภิกฺขู อปฺปิจฺฉา ฯเปฯ เต อุชฺฌายนฺติ ขิยฺยนฺติ วิปาเจนฺติ “กถํ หิ นาม อายสฺมา อุทายี มาตุคาเมน สทฺธิํ เอโก เอกาย รโห ปฏิจฺฉนฺเน อาสเน อลํกมฺมนิเย นิสชฺชํ กปฺเปสฺสตี”ติ. อถ โข เต ภิกฺขู อายสฺมนฺตํ อุทายิํ อเนก\-ปริยาเยน วิครหิตฺวา ภควโต เอตมตฺถํ อาโรเจสุํ ฯเปฯ “สจฺจํ กิร ตฺวํ อุทายิ มาตุคาเมน สทฺธิํ เอโก เอกาย รโห ปฏิจฺฉนฺเน อาสเน อลํกมฺมนิเย นิสชฺชํ กปฺเปสี”ติ. สจฺจํ ภควาติ. วิครหิ พุทฺโธ ภควา ฯเปฯ. กถํ หิ นาม ตฺวํ โมฆปุริส มาตุคาเมน สทฺธิํ เอโก เอกาย รโห ปฏิจฺฉนฺเน อาสเน อลํกมฺมนิเย นิสชฺชํ กปฺเปสฺสสิ, เนตํ โมฆปุริส อปฺปสนฺนานํ วา ปสาทาย ฯเปฯ. เอวญฺจ ปน ภิกฺกเว อิมํ สิกฺขาปทํ อุทฺทิเสยฺยาถ\csromandash{}}

\proseitem{444}{\textbf{“โย ปน ภิกฺขุ มาตุคาเมน สทฺธิํ เอโก เอกาย รโห ปฏิจฺฉนฺเน อาสเน อลํกมฺมนิเย นิสชฺชํ กปฺเปยฺย, ตเมนํ สทฺเธยฺย}\-\textbf{วจสา อุปาสิกา ทิสฺวา ติณฺณํ ธมฺมานํ อญฺญตเรน วเทยฺย ปาราชิเกน วา สํฆาทิเสเสน วา ปาจิตฺติเยน วา, นิสชฺชํ ภิกฺขุ ปฏิชานมาโน ติณฺณํ ธมฺมานํ อญฺญตเรน กาเรตพฺโพ ปาราชิเกน วา สํฆาทิเสเสน วา ปาจิตฺติเยน วา, เยน วา สา สทฺเธยฺย}\-\textbf{วจสา อุปาสิกา วเทยฺย, เตน โส ภิกฺขุ กาเรตพฺโพ, อยํ ธมฺโม อนิยโต”}ติ.}

\proseitem{445}{\textbf{โย ปนา}ติ โย ยาทิโส ฯเปฯ. \textbf{ภิกฺขู}ติ ฯเปฯ อยํ อิมสฺมิํ อตฺเถ อธิปฺเปโต “ภิกฺขู”ติ.}

\prose{\textbf{มาตุคาโม} นาม มนุสฺสิตฺถี, น ยกฺขี, น เปตี, น ติรจฺฉานคตา, อนฺตมโส ตทหุชาตาปิ ทาริกา, ปเคว มหตฺตรี.}

\prose{\textbf{สทฺธินฺ}ติ เอกโต.}

\prose{\textbf{เอโก เอกายา}ติ ภิกฺขุ เจว โหติ มาตุคาโม จ.}

\prose{\csromanfolio{287}\textbf{รโห} นาม จกฺขุสฺส รโห โสตสฺส รโห. \textbf{จกฺขุสฺส รโห} นาม น สกฺกา โหติ อกฺขิํ วา นิขณียมาเน ภมุกํ วา อุกฺขิปิยมาเน สีสํ วา อุกฺขิปียมาเน ปสฺสิตุํ. \textbf{โสตสฺส รโห} นาม น สกฺกา โหติ ปกติกถา โสตุํ.}

\prose{\textbf{ปฏิจฺฉนฺนํ} นาม อาสนํ กุฏฺเฏน วา\csromansharedfootnote{d8cbedaa538e3c21}{กุฑฺเฑน วา (สี, สฺยา)} กวาเฏน วา กิลญฺเชน วา สาณิปากาเรน วา รุกฺเขน วา ถมฺเภน วา โกตฺถฬิยา วา เยน เกนจิ ปฏิจฺฉนฺนํ โหติ.}

\prose{\textbf{อลํกมฺมนิเย}ติ สกฺกา โหติ เมถุนํ ธมฺมํ ปฏิเสวิตุํ.}

\prose{\textbf{นิสชฺชํ กปฺเปยฺยา}ติ มาตุคาเม นิสินฺเน ภิกฺขุ อุปนิสินฺโน วา โหติ อุปนิปนฺโน วา, ภิกฺขุ นิสินฺเน มาตุคาโม อุปนิสินฺโน วา โหติ อุปนิปนฺโน วา, อุโภ วา นิสินฺนา โหนฺติ อุโภ วา นิปนฺนา.}

\prose{\textbf{สทฺเธยฺย}\-\textbf{วจสา} นาม อาคตผลา อภิสเมตาวินี วิญฺญาตสาสนา.}

\prose{\textbf{อุปาสิกา} นาม พุทฺธํ สรณํ คตา, ธมฺมํ สรณํ คตา, สํฆํ สรณํ คตา.}

\prose{\textbf{ทิสฺวา}ติ ปสฺสิตฺวา.}

\prose{ติณฺณํ ธมฺมานํ อญฺญตเรน วเทยฺย ปาราชิเกน วา สํฆาทิเสเสน วา ปาจิตฺติเยน วา, นิสชฺชํ ภิกฺขุ ปฏิชานมาโน ติณฺณํ ธมฺมานํ อญฺญตเรน กาเรตพฺโพ ปาราชิเกน วา สํฆาเทเสเสน วา ปาจิตฺติเยน วา, เยน วา สา สทฺเธยฺย\-วจสา อุปาสิกา วเทยฺย, เตน โส ภิกฺขุ กาเรตพฺโพ.}

\proseitem{446}{สา เจ เอวํ วเทยฺย “อยฺโย มยา ทิฏฺโฐ นิสินฺโน มาตุคามสฺส เมถุนํ ธมฺมํ ปฏิเสวนฺโต”ติ. โส จ ตํ ปฏิชานาติ, อาปตฺติยา กาเรตพฺโพ. สา เจ เอวํ วเทยฺย “อยฺโย มยา ทิฏฺโฐ นิสินฺโน มาตุคามสฺส เมถุนํ ธมฺมํ ปฏิเสวนฺโต”ติ. โส เจ เอวํ วเทยฺย “สจฺจาหํ นิสินฺโน โน จ โข เมถุนํ ธมฺมํ ปฏิเสวินฺ”ติ, นิสชฺชาย กาเรตพฺโพ. สา เจ เอวํ วเทยฺย “อยฺโย มยา ทิฏฺโฐ \csromanfolio{288}นิสินฺโน มาตุคามสฺส เมถุนํ ธมฺมํ ปฏิเสวนฺโต”ติ. โส เจ เอวํ วเทยฺย “นาหํ นิสินฺโน อปิ จ โข นิปนฺโน”ติ, นิปชฺชาย กาเรตพฺโพ. สา เจ เอวํ วเทยฺย “อยฺโย มยา ทิฏฺโฐ นิสินฺโน มาตุคามสฺส เมถุนํ ธมฺมํ ปฏิเสวนฺโต”ติ. โส เจ เอวํ วเทยฺย “นาหํ นิสินฺโน อปิ จ โข ฐิโต”ติ, น กาเรตพฺโพ.}

\proseitem{447}{สา เจ เอวํ วเทยฺย “อยฺโย มยา ทิฏฺโฐ นิปนฺโน มาตุคามสฺส เมถุนํ ธมฺมํ ปฏิเสวนฺโต”ติ. โส จ ตํ ปฏิชานาติ, อาปตฺติยา กาเรตพฺโพ. สา เจ เอวํ วเทยฺย “อยฺโย มยา ทิฏฺโฐ นิปนฺโน มาตุคามสฺส เมถุนํ ธมฺมํ ปฏิเสวนฺโต”ติ. โส เจ เอวํ วเทยฺย “สจฺจาหํ นิปนฺโน โน จ โข เมถุนํ ธมฺมํ ปฏิเสวินฺ”ติ, นิปชฺชาย กาเรตพฺโพ. สา เจ เอวํ วเทยฺย “อยฺโย มยา ทิฏฺโฐ นิปนฺโน มาตุคามสฺส เมถุนํ ธมฺมํ ปฏิเสวนฺโต”ติ. โส เจ เอวํ วเทยฺย “นาหํ นิปนฺโน อปิ จ โข นิสินฺโน”ติ, นิสชฺชาย กาเรตพฺโพ. สา เจ เอวํ วเทยฺย “อยฺโย มยา ทิฏฺโฐ นิปนฺโน มาตุคามสฺส เมถุนํ ธมฺมํ ปฏิเสวนฺโต”ติ. โส เจ เอวํ วเทยฺย “นาหํ นิปนฺโน อปิ จ โข ฐิโต”ติ, น กาเรตพฺโพ.}

\proseitem{448}{สา เจ เอวํ วเทยฺย “อยฺโย มยา ทิฏฺโฐ นิสินฺโน มาตุคาเมน สทฺธิํ กายสํสคฺคํ สมาปชฺชนฺโต”ติ. โส จ ตํ ปฏิชานาติ, อาปตฺติยา กาเรตพฺโพ ฯเปฯ “สจฺจาหํ นิสินฺโน โน จ โข กายสํสคฺคํ สมาปชฺชินฺ”ติ, นิสชฺชาย กาเรตพฺโพ ฯเปฯ “นาหํ นิสินฺโน อปิ จ โข นิปนฺโน”ติ, นิปชฺชาย กาเรตพฺโพ ฯเปฯ “นาหํ นิสินฺโน อปิ จ โข ฐิโต”ติ, น กาเรตพฺโพ.}

\prose{สา เจ เอวํ วเทยฺย “อยฺโย มยา ทิฏฺโฐ นิปนฺโน มาตุคาเมน สทฺธิํ กายสํสคฺคํ สมาปชฺชนฺโต”ติ. โส จ ตํ ปฏิชานาติ, อาปตฺติยา กาเรตพฺโพ ฯเปฯ “สจฺจาหํ นิปนฺโน โน จ โข กายสํสคฺคํ สมาปชฺชินฺ”ติ, นิปชฺชาย กาเรตพฺโพ ฯเปฯ “นาหํ นิปนฺโน อปิ จ โข นิสินฺโน”ติ, นิสชฺชาย กาเรตพฺโพ ฯเปฯ “นาหํ นิปนฺโน อปิ จ โข ฐิโต”ติ, น กาเรตพฺโพ.}

\csromanmatikaanchor{mtk.107}
\csromanmatikaanchor{mtk.108}
\csromantocmarknopage{section}{2. ทุติยอนิยตสิกฺขาปท}{mtk.107}
\bookmark[dest={mtk.107},level=1]{2. ทุติยอนิยตสิกฺขาปท}
\csromantocmarkref{subsection}{อุทายิภิกฺขุวตฺถุ}{mtk.108}
\bookmark[dest={mtk.108},level=2]{อุทายิภิกฺขุวตฺถุ}
\proseitem{449}{สา เจ เอวํ วเทยฺย “อยฺโย มยา ทิฏฺโฐ มาตุคาเมน สทฺธิํ เอโก เอกาย รโห ปฏิจฺฉนฺเน อาสเน อลํกมฺมนิเย \csromanfolio{289}นิสินฺโน”ติ, โส จ ตํ ปฏิชานาติ, นิสชฺชาย กาเรตพฺโพ ฯเปฯ “นาหํ นิสินฺโน อปิ จ โข นิปนฺโน”ติ, นิปชฺชาย กาเรตพฺโพ ฯเปฯ “นาหํ นิสินฺโน อปิ จ โข ฐิโต”ติ, น กาเรตพฺโพ.}

\proseitem{450}{สา เจ เอวํ วเทยฺย “อยฺโย มยา ทิฏฺโฐ มาตุคาเมน สทฺธิํ เอโก เอกาย รโห ปฏิจฺฉนฺเน อาสเน อลํกมฺมนิเย นิปนฺโน”ติ, โส จ ตํ ปฏิชานาติ, นิปชฺชาย กาเรตพฺโพ ฯเปฯ “นาหํ นิปนฺโน อปิ จ โข นิสินฺโน”ติ, นิสชฺชาย กาเรตพฺโพ. “นาหํ นิปนฺโน อปิ จ โข ฐิโต”ติ, น กาเรตพฺโพ.}

\prose{\textbf{อนิยโต}ติ น นิยโต ปาราชิกํ วา สํฆาทิเสโส วา ปาจิตฺติยํ วา.}

\proseitem{451}{คมนํ ปฏิชานาติ, นิสชฺชํ ปฏิชานาติ, อาปตฺติํ ปฏิชานาติ, อาปตฺติยา กาเรตพฺโพ. คมนํ ปฏิชานาติ, นิสชฺชํ น ปฏิชานาติ, อาปตฺติํ ปฏิชานาติ, อาปตฺติยา กาเรตพฺโพ. คมนํ ปฏิชานาติ, นิสชฺชํ ปฏิชานาติ, อาปตฺติํ น ปฏิชานาติ, นิสชฺชาย กาเรตพฺโพ. คมนํ ปฏิชานาติ, นิสชฺชํ น ปฏิชานาติ, อาปตฺติํ น ปฏิชานาติ, น กาเรตพฺโพ.}

\prosewithcloserruled{คมนํ น ปฏิชานาติ, นิสชฺชํ ปฏิชานาติ, อาปตฺติํ ปฏิชานาติ, อาปตฺติยา กาเรตพฺโพ. คมนํ น ปฏิชานาติ, นิสชฺชํ น ปฏิชานาติ, อาปตฺติํ ปฏิชานาติ, อาปตฺติยา กาเรตพฺโพ. คมนํ น ปฏิชานาติ, นิสชฺชํ ปฏิชานาติ, อาปตฺติํ น ปฏิชานาติ, นิสชฺชาย กาเรตพฺโพ. คมนํ น ปฏิชานาติ, นิสชฺชํ น ปฏิชานาติ, อาปตฺติํ น ปฏิชานาติ, น กาเรตพฺโพติ.}{ปฐโม อนิยโต นิฏฺฐิโต.}
\csromanheader{1}{2. ทุติย\-อนิยต\-สิกฺขาปท}
\csromanmatikaanchor{mtk.109}
\csromanmatikaanchor{mtk.110}
\csromantocmarkref{subsection}{ปญฺญตฺติ}{mtk.109}
\bookmark[dest={mtk.109},level=2]{ปญฺญตฺติ}
\csromantocmarkref{subsection}{สิกฺขาปทวิภงฺค, ปทภาชนีย}{mtk.110}
\bookmark[dest={mtk.110},level=2]{สิกฺขาปทวิภงฺค, ปทภาชนีย}
\proseitem{452}{เตน สมเยน พุทฺโธ ภควา สาวตฺถิยํ วิหรติ เชตวเน อนาถ\-ปิณฺฑิกสฺส อาราเม. เตน โข ปน สมเยน อายสฺมา อุทายี ภควตา ปฏิกฺขิตฺตํ “มาตุคาเมน สทฺธิํ เอโก เอกาย รโห ปฏิจฺฉนฺเน อาสเน อลํกมฺมนิเย นิสชฺชํ กปฺเปตุนฺ”ติ ตสฺสา เยว \csromanfolio{290}กุมาริกาย สทฺทิํ เอโก เอกาย รโห นิสชฺชํ กปฺเปสิ กาลยุตฺตํ \textbf{สมุลฺลปนฺโต กาลยุตฺตํ ธมฺมํ ภณนฺโต. ทุติยมฺปิ โข วิสาขา} \textbf{มิคารมาตา นิมนฺติตา ตํ กุลํ อคมาสิ. อทฺทสา โข วิสาขา มิคารมาตา} นิมนฺติตา ตํ กุลํ อคมาสิ. อทฺทสา โข วิสาขา มิคารมาตา อายสฺมนฺตํ อุทายิํ ตสฺสาเยว กุมาริกาย สทฺธิํ เอกํ เอกาย รโห นิสินฺนํ, ทิสฺวาน อายสฺมนฺตํ อุทายิํ เอตทโวจ “อิทํ ภนฺเต นจฺฉนฺนํ นปฺปติรูปํ, ยํ อยฺโย มาตุคาเมน สทฺธิํ เอโก เอกาย รโห นิสชฺชํ กปฺเปติ, กิญฺจาปิ ภนฺเต อยฺโย อนตฺถิโก เตน ธมฺเมน, อปิ จ ทุสฺสทฺธาปยา อปฺปสนฺนา มนุสฺสา”ติ. เอวํปิ โข อายสฺมา อุทายี วิสาขาย มิคารมาตุยา วุจฺจมาโน นาทิยิ. อถ โข วิสาขา มิคารมาตา นิกฺขมิตฺวา ภิกฺขูนํ เอตมตฺถํ อาโรเจสิ. เย เต ภิกฺขู อปฺปิจฺฉา ฯเปฯ เต อุชฺฌายนฺติ ขิยฺยนฺติ วิปาเจนฺติ “กถํ หิ นาม อายสฺมา อุทายี มาตุคาเมน สทฺธิํ เอโก เอกาย รโห นิสชฺชํ กปฺเปสฺสตี”ติ. อถ โข เต ภิกฺขู อายสฺมนฺตํ อุทายิํ อเนก\-ปริยาเยน วิครหิตฺวา ภควโต เอตมตฺถํ อาโรเจสุํ ฯเปฯ “สจฺจํ กิร ตฺวํ อุทายิ มาตุคาเมน สทฺธิํ เอโก เอกาย รโห นิสชฺชํ กปฺเปสี”ติ. สจฺจํ ภควาติ. วิครหิ พุทฺโธ ภควา ฯเปฯ. กถํ หิ นาม ตฺวํ โมฆปุริส มาตุคาเมน สทฺธิํ เอโก เอกาย รโห นิสชฺชํ กปฺเปสฺสสิ, เนตํ โมฆปุริส อปฺปสนฺนานํ วา ปสาทาย ฯเปฯ. เอวญฺจ ปน ภิกฺขเว อิมํ สิกฺขาปทํ อุทฺทิเสยฺยาถ\csromandash{}}

\proseitem{453}{\textbf{“นเหว โข ปน ปฏิจฺฉนฺนํ อาสนํ โหติ นาลํ}\-\textbf{กมฺมนิยํ, อลญฺจ โข โหติ มาตุคามํ ทุฏฺฐุลฺลาหิ วา จาหิ โอภาสิตุํ, โย ปน ภิกฺขุ ตถารูเป อาสเน มาตุคาเมน สทฺธิํ เอโก เอกาย รโห นิสชฺชํ กปฺเปยฺย, ตเมนํ สทฺเธยฺย}\-\textbf{วจสา อุปาสิกา ทิสฺวา ทฺวินฺนํ ธมฺมานํ อญฺญตเรน วเทยฺย สํฆาทิเสเสน วา ปาจิตฺติเยน วา, นิสชฺชํ ภิกฺขุ ปฏิชานมาโน ทฺวินฺนํ ธมฺมานํ อญฺญตเรน กาเรตพฺโพ สํฆาทิเสเสน วา ปาจิตฺติเยน วา, เยน วา สา สทฺเธยฺย}\-\textbf{วจสา อุปาสิกา วเทยฺย, เตน โส ภิกฺขุ กาเรตพฺโพ, อยํปิ ธมฺโม อนิยโต”}ติ.}

\proseitem{454}{\textbf{น เหว โข ปน ปฏิจฺฉนฺนํ อาสนํ โหตี}ติ อปฺปฏิจฺฉนํ โหติ กุฏฺเฏน วา กวาเฏน วา กิลญฺเชน วา สาณิปากาเรน วา รุกฺเขน วา ถมฺเภน วา โกตฺถฬิยา วา เยน เกนจิ อปฺปฏิจฺฉนฺนํ โหติ.}

\prose{\csromanfolio{291}\textbf{นาลํ กมฺมนิยนฺ}ติ น สกฺกา โหติ เมถุนํ ธมฺมํ ปฏิเสวิตุํ.}

\prose{\textbf{อลญฺจ โข โหติ มาตุคามํ ทุฏฺฐุลฺลาหิ วาจาหิ โอภาสิตุนฺ}ติ สกฺกา โหติ มาตุคามํ ทุฏฺฐุลฺลาหิ วาจาหิ โอภาสิตุํ.}

\prose{\textbf{โย ปนา}ติ โย ยาทิโส ฯเปฯ. \textbf{ภิกฺขู}ติ ฯเปฯ อยํ อิมสฺมิํ อตฺเถ อธิปฺเปโต “ภิกฺขู”ติ.}

\prose{\textbf{ตถารูเป อาสเน}ติ เอวรูเป อาสเน.}

\prose{\textbf{มาตุคาโม} นาม มนุสฺสิตฺถี, น ยกฺขี, น เปตี, น ติรจฺฉานคตา, วิญฺญู ปฏิพลา สุภาสิต\-ทุพฺภาสิตํ ทุฏฺฐุลฺลา\-ทุฏฺฐุลฺลํ อาชานิตุํ.}

\prose{\textbf{สทฺธินฺ}ติ เอกโต.}

\prose{\textbf{เอโก เอกายา}ติ ภิกฺขุ เจว โหติ มาตุคาโม จ.}

\prose{\textbf{รโห} นาม จกฺขุสฺส รโห โสตสฺส รโห. \textbf{จกฺขุสฺส รโห} นาม น สกฺกา โหติ อกฺขิํ วา นิขณียมาเน ภมุกํ วา อุกฺขิปียมาเน สีสํ วา อุกฺขิปียมาเน ปสฺสิตุํ. \textbf{โสตสฺส รโห} นาม น สกฺกา โหติ ปกติกถา โสตุํ.}

\prose{\textbf{นิสชฺชํ กปฺเปยฺยา}ติ มาตุคาเม นิสินฺเน ภิกฺขุ อุปนิสินฺโน วา โหติ อุปนิปนฺโน วา, ภิกฺขุ นิสินฺเน มาตุคาโม อุปนิสินฺโน วา โหติ อุปนิปนฺโน วา, อุโภ วา นิสินฺนา โหนฺติ อุโภ วา นิปนฺนา.}

\prose{\textbf{สทฺเธยฺย}\-\textbf{วจสา} นาม อาคตผลา อภิสเมตาวินี วิญฺญาตสาสนา.}

\prose{\textbf{อุปาสิกา} นาม พุทฺธํ สรณํ คตา, ธมฺมํ สรณํ คตา, สํฆํ สรณํ คตา.}

\prose{\textbf{ทิสฺวา}ติ ปสฺสิตฺวา.}

\prose{ทฺวินฺนํ ธมฺมานํ อญฺญตเรน วเทยฺย สํฆาทิเสเสน วา ปาจิตฺติเยน วา, นิสชฺชํ ภิกฺขุ ปฏิชานมาโน ทฺวินฺนํ ธมฺมานํ อญฺญตเรน กาเรตพฺโพ สํฆาทิเสเสน วา ปาจิตฺติเยน วา เยน วา สา สทฺเธยฺย\-วจสา อุปาสิกา วเทยฺย, เตน โส ภิกฺขุ กาเรตพฺโพ.}

\proseitem{455}{สา เจ เอวํ วเทยฺย “อยฺโย มยา ทิฏฺโฐ นิสินฺโน มาตุคาเมน สทฺธิํ กายสํสคฺคํ สมาปชฺชนฺโต”ติ, โส จ ตํ ปฏิชานาติ, อาปตฺติยา กาเรตพฺโพ. สา เจ เอวํ วเทยฺย “อยฺโย \csromanfolio{292}มยา ทิฏฺโฐ นิสินฺโน มาตุคาเมน สทฺธิํ กายสํสคฺคํ สมาปชฺชนฺโต”ติ. โส เจ เอวํ วเทยฺย “สจฺจาหํ นิสินฺโน โน จ โข กายสํสคฺคํ สมาปชฺชินฺ”ติ, นิสชฺชาย กาเรตพฺโพ ฯเปฯ “นาหํ นิสินฺโน อปิ จ โข นิปนฺโน”ติ, นิปชฺชาย กาเรตพฺโพ ฯเปฯ “นาหํ นิสินฺโน อปิ จ โข ฐิโต”ติ, น \textbf{กาเรตพฺโพ.}}

\prose{สา เจ เอวํ วเทยฺย “อยฺโย มยา ทิฏฺโฐ นิปนฺโน มาตุคาเมน สทฺธิํ กายสํสคฺคํ สมาปชฺชนฺโต”ติ. โส จ ตํ ปฏิชานาติ, อาปตฺติยา กาเรตพฺโพ ฯเปฯ “สจฺจาหํ นิปนฺโน โน จ โข กายสํสคฺคํ สมาปชฺชินฺ”ติ, นิปชฺชาย กาเรตพฺโพ ฯเปฯ “นาหํ นิปนฺโน อปิ จ โข นิสินฺโน”ติ, นิสชฺชาย กาเรตพฺโพ ฯเปฯ “นาหํ นิปนฺโน อปิ จ โข ฐิโต”ติ, น กาเรตพฺโพ.}

\prose{สา เจ เอวํ วเทยฺย “อยฺยสฺส มยา สุตํ นิสินฺนสฺส มาตุคามํ ทุฏฺฐุลฺลาหิ วาจาหิ โอภาสนฺตสฺสา”ติ. โส จ ตํ ปฏิชานาติ, อาปตฺติยา กาเรตพฺโพ. สา เจ เอวํ วเทยฺย “อยฺยสฺส มยา สุตํ นิสินฺนสฺส มาตุคามํ ทุฏฺฐุลฺลาหิ วาจาหิ โอภาสนฺตสฺสา”ติ. โส เจ เอวํ วเทยฺย “สจฺจาหํ นิสินฺโน โน จ โข ทุฏฺฐุลฺลาหิ วาจาหิ โอภาสินฺ”ติ, นิสชฺชาย กาเรตพฺโพ ฯเปฯ “นาหํ นิสินฺโน อปิ จ โข นิปนฺโน”ติ, นิปชฺชาย กาเรตพฺโพ ฯเปฯ “นาหํ นิสินฺโน อปิ จ โข ฐิโต”ติ, น กาเรตพฺโพ.}

\prose{สา เจ เอวํ วเทยฺย “อยฺยสฺส มยา สุตํ นิปนฺนสฺส มาตุคามํ ทุฏฺฐุลฺลาหิ วาจาหิ โอภาสนฺตสฺสา”ติ, โส จ ตํ ปฏิชานาติ, อาปตฺติยา กาเรตพฺโพ ฯเปฯ “สจฺจาหํ นิปนฺโน โน จ โข ทุฏฺฐุลฺลาหิ วาจาหิ โอภาสินฺ”ติ, นิปชฺชาย กาเรตพฺโพ ฯเปฯ “นาหํ นิปนฺโน อปิ จ โข นิสินฺโน”ติ, นิสชฺชาย กาเรตพฺโพ ฯเปฯ “นาหํ นิปนฺโน อปิ จ โข ฐิโต”ติ, น กาเรตพฺโพ.}

\proseitem{456}{สา เจ เอวํ วเทยฺย “อยฺโย มยา ทิฏฺโฐ มาตุคาเมน สทฺธิํ เอโก เอกาย รโห นิสินฺโน”ติ. โส จ ตํ ปฏิชานาติ, นิสชฺชาย กาเรตพฺโพ ฯเปฯ “นาหํ นิสินฺโน อปิ จ โข นิปนฺโน”ติ, นิปชฺชาย กาเรตพฺโพ ฯเปฯ “นาหํ นิสินฺโน อปิ จ โข ฐิโต”ติ, น กาเรตพฺโพ.}

\csromanmatikaanchor{mtk.111}
\csromantocmarkref{section}{อุทฺทานคาถา}{mtk.111}
\bookmark[dest={mtk.111},level=1]{อุทฺทานคาถา}
\prose{\csromanfolio{293}สา เจ เอวํ วเทยฺย “อยฺโย มยา ทิฏฺโฐ มาตุคาเมน สทฺธิํ เอโก เอกาย รโห นิปนฺโน”ติ. โส จ ตํ ปฏิชานาติ, นิปชฺชาย กาเรตพฺโพ ฯเปฯ “นาหํ นิปนฺโน อปิ จ โข นิสินฺโน”ติ, นิสชฺชาย กาเรตพฺโพ ฯเปฯ “นาหํ นิปนฺโน อปิ จ โข ฐิโต”ติ, น กาเรตพฺโพ.}

\prose{\textbf{อยมฺปิ}ติ ปุริมํ อุปาทาย วุจฺจติ.}

\prose{\textbf{อนิยโต}ติ น นิยโต สํฆาทิเสโส วา ปาจิตฺติยํ วา.}

\proseitem{457}{คมนํ ปฏิชานาติ, นิสชฺชํ ปฏิชานาติ, อาปตฺติํ ปฏิชานาติ, อาปตฺติยา กาเรตพฺโพ. คมนํ ปฏิชานาติ, นิสชฺชํ น ปฏิชานาติ, อาปตฺติํ ปฏิชานาติ, อาปตฺติยา กาเรตพฺโพ. คมนํ ปฏิชานาติ, นิสชฺชํ ปฏิชานาติ, อาปตฺติํ น ปฏิชานาติ, นิสชฺชาย กาเรตพฺโพ. คมนํ ปฏิชานาติ, นิสชฺชํ น ปฏิชานาติ, อาปตฺติํ น ปฏิชานาติ, น กาเรตพฺโพ.}

\prosewithcloser{คมนํ น ปฏิชานาติ, นิสชฺชํ ปฏิชานาติ, อาปตฺติํ ปฏิชานาติ, อาปตฺติยา กาเรตพฺโพ. คมนํ น ปฏิชานาติ, นิสชฺชํ น ปฏิชานาติ, อาปตฺติํ ปฏิชานาติ, อาปตฺติยา กาเรตพฺโพ. คมนํ น ปฏิชานาติ, นิสชฺชํ ปฏิชานาติ, อาปตฺติํ น ปฏิชานาติ, นิสชฺชาย กาเรตพฺโพ. คมนํ น ปฏิชานาติ, นิสชฺชํ น ปฏิชานาติ, อาปตฺติํ น ปฏิชานาติ, น กาเรตพฺโพติ.}{ทุติโย อนิยโต นิฏฺฐิโต.}
\proseitem{458}{อุทฺทิฏฺฐา โข อายสฺมนฺโต เทฺว อนิยตา ธมฺมา. ตตฺถายสฺมนฺเต ปุจฺฉามิ กจฺจิตฺถ ปริสุทฺธา, ทุติยมฺปิ ปุจฺฉามิ กจฺจิตฺถ ปริสุทฺธา, ตติยมฺปิ ปุจฺฉามิ กจฺจิตฺถ ปริสุทฺธา, ปริสุทฺเธ\-ตฺถา\-ยสฺมนฺโต, ตสฺมา ตุณฺหี, เอวเมตํ ธารยามีติ.}

\vspace{\tipitakaparskip}
\csromancenter{ตสฺสุทฺทานํ}
\setlength{\csromangathapagewidth}{0pt}
\setlength{\csromangathawakwidth}{0pt}
\csromangathameasuregroup{อลํ กมฺมนิยญฺเจว, \\ อนิยตา สุปญฺญตฺตา,}{\csromangathabat{อลํ กมฺมนิยญฺเจว,}{ตเถว จ นเหว โข.} \\ \csromangathabat{อนิยตา สุปญฺญตฺตา,}{พุทฺธเสฏฺเฐน ตาทินาติ.}}
\csromangathasetleft{อลํ กมฺมนิยญฺเจว, \\ อนิยตา สุปญฺญตฺตา,}
\csromangathagroupwithclosermajor{\csromangathastanza{\csromangathabat{อลํ กมฺมนิยญฺเจว,}{ตเถว จ นเหว โข.} \\ \csromangathabat{อนิยตา สุปญฺญตฺตา,}{พุทฺธเสฏฺเฐน ตาทินาติ.}}}{\textbf{อนิยตกณฺฑํ นิฏฺฐิตํ.}}
\csromanmatikaanchor{mtk.112}
\csromantocmarknopage{chapter}{4. นิสฺสคฺคิยกณฺฑ}{mtk.112}
\bookmark[dest={mtk.112},level=0]{4. นิสฺสคฺคิยกณฺฑ}
\chapterheadpageread{4. \textbf{นิสฺสคฺคิย}\-\textbf{กณฺฑ}}
\csromanmatikaanchor{mtk.113}
\csromanmatikaanchor{mtk.114}
\csromantocmarknopage{section}{1. จีวรวคฺค}{mtk.113}
\bookmark[dest={mtk.113},level=1]{1. จีวรวคฺค}
\csromantocmarknopage{subsection}{1. ปฐมกถินสิกฺขาปท}{mtk.114}
\bookmark[dest={mtk.114},level=2]{1. ปฐมกถินสิกฺขาปท}
\csromanmarkh{1. จีวรวคฺค}\csromanmarkhii{1. ปฐม\-กถิน\-สิกฺขาปท}\csromanheadernomark{2}{1. \textbf{จีวรวคฺค 1. ปฐม}\-\textbf{กถิน}\-\textbf{สิกฺขาปท}}
\vspace{\tipitakaparskip}
\csromancenter{อิเม โข ปนายสฺมนฺโต ติํส นิสฺสคฺคิยา ปาจิตฺติยา \\ ธมฺมา อุทฺเทสํ อาคจฺฉนฺติ.}
\csromanmatikaanchor{mtk.115}
\csromanmatikaanchor{mtk.116}
\csromantocmarkref{subsubsection}{จพฺพคฺคิยภิกฺขุวตฺถุ}{mtk.115}
\bookmark[dest={mtk.115},level=3]{จพฺพคฺคิยภิกฺขุวตฺถุ}
\csromantocmarkref{subsubsection}{ปฐมปญฺญตฺติ}{mtk.116}
\bookmark[dest={mtk.116},level=3]{ปฐมปญฺญตฺติ}
\proseitem{459}{\csromanfolio{294}เตน สเมเยน พุทฺโธ ภควา เวสาลิยํ วิหรติ โคตมเก เจติเย. เตน โข ปน สมเยน ภควตา ภิกฺขูนํ ติจีวรํ อนุญฺญาตํ โหติ. ฉพฺพคฺคิยา ภิกฺขู ภควตา ติจีวรํ อนุญฺญาตนฺติ อญฺเญเนว ติจีวเรน คามํ ปวิสนฺติ, อญฺเญน ติจีวเรน อาราเม อจฺฉนฺติ, อญฺเญน ติจีวเรน นหานํ โอตรนฺติ. เย เต ภิกฺขู อปฺปิจฺฉา ฯเปฯ เต อุชฺฌายนฺติ ขิยฺยนฺติ วิปาเจนฺติ “กถํ หิ นาม ฉพฺพคฺคิยา ภิกฺขู อติเรกจีวรํ ธาเรสฺสนฺตี”ติ. อถ โข เต ภิกฺขู ฉพฺพคฺคิเย ภิกฺขู อเนก\-ปริยาเยน วิครหิตฺวา ภควโต เอตมตฺถํ อาโรเจสุํ ฯเปฯ “สจฺจํ กิร ตุเมฺห ภิกฺขเว อติเรกจีวรํ ธาเรถา”ติ. สจฺจํ ภควาติ. วิครหิ พุทฺโธ ภควา ฯเปฯ. กถํ หิ นาม ตุเมฺห โมฆปุริสา อติเรกจีวรํ ธาเรสฺสถ, เนตํ โมฆปุริสา อปฺปสนฺนานํ วา ปสาทาย ฯเปฯ. เอวญฺจ ปน ภิกฺขเว อิมํ สิกฺขาปทํ อุทฺทิเสยฺยาถ\csromandash{}}

\proseitem{460}{\textbf{“โย ปน ภิกฺขุ อติเรกจีวรํ ธาเรยฺย, นิสฺสคฺคิยํ ปาจิตฺติยนฺ”}ติ.}

\prose{เอวญฺจิทํ ภควตา ภิกฺขูนํ สิกฺขาปทํ ปญฺญตฺตํ โหติ.}

\csromanmatikaanchor{mtk.117}
\csromanmatikaanchor{mtk.118}
\csromantocmarkref{subsubsection}{อนุปญฺญตฺติ}{mtk.117}
\bookmark[dest={mtk.117},level=3]{อนุปญฺญตฺติ}
\csromantocmarkref{subsubsection}{สิกฺขาปทวิภงฺค, ปทภาชนีย}{mtk.118}
\bookmark[dest={mtk.118},level=3]{สิกฺขาปทวิภงฺค, ปทภาชนีย}
\proseitem{461}{\csromansymbolfootnote{*}{วิ 3 . 402 ปิฏฺเฐปิ.}เตน โข ปน สมเยน อายสฺมโต อานนฺทสฺส อติเรกจีวรํ อุปฺปนฺนํ โหติ. อายสฺมา จ อานนฺโท ตํ จีวรํ อายสฺมโต สาริปุตฺตสฺส ทาตุกาโม โหติ. อายสฺมา จ สาริปุตฺโต สาเกเต วิหรติ. อถ โข อายสฺมโต อานนฺทสฺส เอตทโหสิ “ภควตา สิกฺขาปทํ ปญฺญตฺตํ ‘น อติเรกจีวรํ ธาเรตพฺพนฺ’ติ, อิทญฺจ เม อติเรกจีวรํ อุปฺปนฺนํ, อหญฺจิมํ จีวรํ อายสฺมโต สาริปุตฺตสฺส ทาตุกาโม, อายสฺมา จ สาริปุตฺโต สาเกเต วิหรติ, กถํ นุ โข มยา ปฏิปชฺชิตพฺพนฺ”ติ. อถ โข อายสฺมา อานนฺโท ภควโต \csromanfolio{295}เอตมตฺถํ อาโรเจสิ. กีวจิรํ ปนานนฺท สาริปุตฺโต อาคจฺฉิสฺสตีติ. นวมํ วา ภควา ทิวสํ ทสมํ วาติ. อถ โข ภควา เอตสฺมิํ นิทาเน เอตสฺมิํ ปกรเณ ธมฺมิํ กถํ กตฺวา ภิกฺขู อามนฺเตสิ อนุชานามิ ภิกฺขเว ทสาหปรมํ อตฺติเรกจีวรํ ธาเรตุํ. เอวญฺจ ปน ภิกฺขเว อิมํ สิกฺขาปทํ อุทฺทิเสยฺยาถ\csromandash{}}

\proseitem{462}{\textbf{“นิฏฺฐิต}\-\textbf{จีวรสฺมิํ ภิกฺขุนา อุพฺภตสฺมิํ กถิเน ทสาหปรมํ อติเรกจีวรํ ธาเรตพฺพํ, ตํ อติกฺกามยโต นิสฺสคฺคิยํ ปาจิตฺติยนฺ”}ติ.}

\proseitem{463}{\textbf{นิฏฺฐิต}\-\textbf{จีวร}\-\textbf{สฺมินฺ}ติ ภิกฺขุโน จีวรํ กตํ วา โหติ นฏฺฐํ วา วินฏฺฐํ วา ทฑฺฒํ วา จีวราสา วา อุปจฺฉินฺนา.}

\prose{\textbf{อุพฺภตสฺมิํ กถิเน}ติ อฏฺฐนฺนํ มาติกานํ อญฺญตราย มาติกาย อุพฺภตํ โหติ, สํเฆน วา อนฺตรา อุพฺภตํ โหติ.}

\prose{\textbf{ทสา}\-\textbf{ห}\-\textbf{ปรมนฺ}ติ ทสาหปรมตา ธาเรตพฺพํ.}

\prose{\textbf{อติเรกจีวรํ} นาม อนธิฏฺฐิตํ อวิกปฺปิตํ.}

\prose{\textbf{จิวรํ} นาม ฉนฺนํ จีวรานํ อญฺญตรํ จีวรํ วิกปฺปนุปคํ ปจฺฉิมํ.}

\prose{\textbf{ตํ อติกฺกามยโต นิสฺสคฺคิยํ โหตี}\csromansharedfootnote{87c5bd1ef08fb0e5}{โหตีติ อิทํ ปทํ สพฺพโปตฺถเกสุ อตฺถิ, สิกฺขาปเท ปน นตฺถิ, เอวมุปริปิ.} ติ เอกาทเส อรุณุคฺคมเน นิสฺสคฺคิยํ โหติ. นิสฺสชฺชิตพฺพํ สํฆสฺส วา คณสฺส วา ปุคฺคลสฺส วา. เอวญฺจ ปน ภิกฺขเว นิสฺสชฺชิตพฺพํ. เตน ภิกฺขุนา สํฆํ อุปสงฺกมิตฺวา เอกํสํ อุตฺตราสงฺคํ กริตฺวา วุฑฺฒานํ ภิกฺขูนํ ปาเท วนฺทิตฺวา อุกฺกุฏิกํ นิสีทิตฺวา อญฺชลิํ ปคฺคเหตฺวา เอวมสฺส วจนีโย “อิทํ เม ภนฺเต จีวรํ ทสาหา\-ติกฺกนฺตํ นิสฺสคฺคิยํ, อิมาหํ สํฆสฺส นิสฺสชฺชามี”ติ. นิสฺสชฺชิตฺวา อาปตฺติ เทเสตพฺพา, พฺยตฺเตน ภิกฺขุนา ปฏิพเลน อาปตฺติ ปฏิคฺคเหตพฺพา, นิสฺสฏฺฐ\-จีวรํ ทาตพฺพํ\csromandash{}}

\proseitem{464}{“สุณาตุ เม ภนฺเต สํโฆ, อิทํ จีวรํ อิตฺถนฺนามสฺส ภิกฺขุโน นิสฺสคฺคิยํ สํฆสฺส นิสฺสฏฺฐํ, ยทิ สํฆสฺส ปตฺตกลฺลํ, สํโฆ อิมํ จีวรํ อิตฺถนฺนามสฺส ภิกฺขุโน ทเทยฺยา”ติ.}

\proseitem{465}{\csromanfolio{296}เตน ภิกฺขุนา สมฺพหุเล ภิกฺขู อุปสงฺกมิตฺวา เอกํสํ อุตฺตราสงฺคํ กริตฺวา วุฑฺฒานํ ภิกฺขูนํ ปาเท วนฺทิตฺวา อุกฺกุฏิกํ นิสีทิตฺวา อญฺชลิํ ปคฺคเหตฺวา เอวมสฺสุ วจนียา “อิทํ เม ภนฺเต จีวรํ ทสาหา\-ติกฺกนฺตํ นิสฺสคฺคิยํ, อิมาหํ อายสฺมนฺตานํ นิสฺสชฺชามี”ติ. นิสฺสชฺชิตฺวา อาปตฺติ เทเสตพฺพา, พฺยตฺเตน ภิกฺขุนา ปฏิพเลน อาปตฺติ ปฏิคฺคเหตพฺพา, นิสฺสฏฺฐ\-จีวรํ ทาตพฺพํ\csromandash{}}

\proseitem{466}{“สุณนฺตุ เม อายสฺมนฺตา, อิทํ จีวรํ อิตฺถนฺนามสฺส ภิกฺขุโน นิสฺสคฺคิยํ อายสฺมนฺตานํ นิสฺสฏฺฐํ, ยทา\-ยสฺมนฺตานํ ปตฺตกลฺลํ, อายสฺมนฺตา อิมํ จีวรํ อิตฺถนฺนามสฺส ภิกฺขุโน ทเทยฺยุนฺ”ติ.}

\proseitem{467}{เตน ภิกฺขุนา เอกํ ภิกฺขุํ อุปสงฺกมิตฺวา เอกํสํ อุตฺตราสงฺคํ กริตฺวา อุกฺกุฏิกํ นิสีทิตฺวา อญฺชลิํ ปคฺคเหตฺวา เอวมสฺส วจนีโย “อิทํ เม อาวุโส จีวรํ ทสาหา\-ติกฺกนฺตํ นิสฺสคฺคิยํ, อิมาหํ อายสฺมโต นิสฺสชฺชามี”ติ. นิสฺสชฺชิตฺวา อาปตฺติ เทเสตพฺพา, เตน ภิกฺขุนา อาปตฺติ ปฏิคฺคเหตพฺพา, นิสฺสฏฺฐ\-จีวรํ ทาตพฺพํ “อิมํ จีวรํ อายสฺมโต ทมฺมี”ติ.}

\proseitem{468}{ทสาหาติกฺกนฺเต อติกฺกนฺตสญฺญี, นิสฺสคฺคิยํ ปาจิตฺติยํ. ทสาหาติกฺกนฺเต เวมติโก, นิสฺสคฺคิยํ ปาจิตฺติยํ. ทสาหาติกฺกนฺเต อนติกฺกนฺต\-สญฺญี, นิสฺสคฺคิยํ ปาจิตฺติยํ. อนธิฏฺฐิเต อธิฏฺฐิต\-สญฺญี, นิสฺสคฺคิยํ ปาจิตฺติยํ. อวิกปฺปิเต วิกปฺปิตสญฺญี, นิสฺสคฺคิยํ ปาจิตฺติยํ. อวิสฺสชฺชิเต วิสฺสชฺชิต\-สญฺญี, นิสฺสคฺคิยํ ปาจิตฺติยํ. อนฏฺเฐ นฏฺฐสญฺญี, นิสฺสคฺคิยํ ปาจิตฺติยํ. อวินฏฺเฐ วินฏฺฐสญฺญี, นิสฺสคฺคิยํ ปาจิตฺติยํ. อทฑฺเฒ ทฑฺฒสญฺญี, นิสฺสคฺคิยํ ปาจิตฺติยํ. อวิลุตฺเต วิลุตฺตสญฺญี, นิสฺสคฺคิยํ ปาจิตฺติยํ.}

\prose{นิสฺสคฺคิยํ จีวรํ อนิสฺสชฺชิตฺวา ปริภุญฺชติ, อาปตฺติ ทุกฺกฏสฺส. ทสาหา\-นติกฺกนฺเต อติกฺกนฺตสญฺญี, อาปตฺติ ทุกฺกฏสฺส. ทสาหา\-นติกฺกนฺเต เวมติโก, อาปตฺติ ทุกฺกฏสฺส. ทสาหา\-นติกฺกนฺเต อนติกฺกนฺต\-สญฺญี, อนาปตฺติ.}

\proseitem{469}{อนาปตฺติ อนฺโตทสาหํ อธิฏฺเฐติ, วิกปฺเปติ, วิสฺสชฺเชติ, นสฺสติ, วินสฺสติ, ฑยฺหติ, อจฺฉินฺทิตฺวา คณฺหนฺติ, วิสฺสาสํ คณฺหนฺติ, อุมฺมตฺตกสฺส อาทิกมฺมิกสฺสาติ.}

\csromanmatikaanchor{mtk.119}
\csromanmatikaanchor{mtk.120}
\csromanmatikaanchor{mtk.121}
\csromantocmarknopage{subsection}{2. อุโทสิตสิกฺขาปท}{mtk.119}
\bookmark[dest={mtk.119},level=2]{2. อุโทสิตสิกฺขาปท}
\csromantocmarkref{subsubsection}{สาวตฺถิภิกฺขุวตฺถุ}{mtk.120}
\bookmark[dest={mtk.120},level=3]{สาวตฺถิภิกฺขุวตฺถุ}
\csromantocmarkref{subsubsection}{ปฐมปญฺญตฺติ}{mtk.121}
\bookmark[dest={mtk.121},level=3]{ปฐมปญฺญตฺติ}
\proseitemwithcloserruled{470}{\csromanfolio{297}เตน โข ปน สมเยน ฉพฺพคฺคิยา ภิกฺขู นิสฺสฏฺฐ\-จีวรํ น เทนฺติ. ภควโต เอตมตฺถํ อาโรเจสุํ. น ภิกฺขเว นิสฺสฏฺฐ\-จีวรํ น ทาตพฺพํ, โย น ทเทยฺย, อาปตฺติ ทุกฺกฏสฺสาติ.}{กถิน\-สิกฺขาปทํ นิฏฺฐิตํ ปฐมํ.}
\csromanmarkh{1. จีวรวคฺค}\csromanmarkhii{2. อุโทสิต\-สิกฺขาปท}\csromanheadernomark{2}{1. \textbf{จีวรวคฺค 2. อุโทสิต}\-\textbf{สิกฺขาปท}}
\proseitem{471}{เตน สมเยน พุทฺโธ ภควา สาวตฺถิยํ วิหรติ เชตวเน อนาถ\-ปิณฺฑิกสฺส อาราเม. เตน โข ปน สมเยน ภิกฺขู ภิกฺขูนํ หตฺเถ จีวรํ นิกฺขิปิตฺวา สนฺตรุตฺตเรน ชนปท\-จาริกํ ปกฺกมนฺติ, ตานิ จีวรานิ จิรํ นิกฺขิตฺตานิ กณฺณกิตานิ โหนฺติ, ตานิ ภิกฺขู โอตาเปนฺติ. อทฺทสา โข อายสฺมา อานนฺโท เสนาสนจาริกํ อาหิณฺฑนฺโต เต ภิกฺขู ตานิ จีวรานิ โอตาเปนฺเต, ทิสฺวาน เยน เต ภิกฺขู เตนุปสงฺกมิ, อุปสงฺกมิตฺวา เต ภิกฺขู เอตทโวจ “กสฺสิมานิ อาวุโส จีวรานิ กณฺณกิตานี”ติ. อถ โข เต ภิกฺขู อายสฺมโต อานนฺทสฺส เอตมตฺถํ อาโรเจสุํ. อายสฺมา อานนฺโท อุชฺฌายติ ขิยฺยติ วิปาเจติ “กถํ หิ นาม ภิกฺขู ภิกฺขูนํ หตฺเถ จีวรํ นิกฺขิปิตฺวา สนฺตรุตฺตเรน ชนปท\-จาริกํ ปกฺกมิสฺสนฺตี”ติ. อถ โข อายสฺมา อานนฺโท เต ภิกฺขู อเนก\-ปริยาเยน วิครหิตฺวา ภควโต เอตมตฺถํ อาโรเจสิ ฯเปฯ “สจฺจํ กิร ภิกฺขเว ภิกฺขู ภิกฺขูนํ หตฺเถ จีวรํ นิกฺขิปิตฺวา สนฺตรุตฺตเรน ชนปท\-จาริกํ ปกฺกมนฺตี”ติ. สจฺจํ ภควาติ. วิครหิ พุทฺโธ ภควา ฯเปฯ. กถํ หิ นาม เต ภิกฺขเว โมฆปุริสา ภิกฺขูนํ หตฺเถ จีวรํ นิกฺขิปิตฺวา สนฺตรุตฺตเรน ชนปท\-จาริกํ ปกฺกมิสฺสนฺติ, เนตํ ภิกฺขเว อปฺปสนฺนานํ วา ปสาทาย ฯเปฯ. เอวญฺจ ปน ภิกฺขเว อิมํ ทิกฺขาปทํ อุทฺทิเสยฺยาถ\csromandash{}}

\proseitem{472}{\textbf{“นิฏฺฐิต}\-\textbf{จีวรสฺมิํ ภิกฺขุนา อุพฺภตสฺมิํ กถิเน เอกรตฺตํปิ เจ ภิกฺขุ ติจีวเรน วิปฺปวเสยฺย, นิสฺสคฺคิยํ ปาจิตฺติยนฺ”}ติ.}

\prose{\textbf{เอวญฺจิทํ ภควตา ภิกฺขูนํ สิกฺขาปทํ ปญฺญตฺตํ โหติ.}}

\csromanmatikaanchor{mtk.122}
\csromanmatikaanchor{mtk.123}
\csromantocmarkref{subsubsection}{อนุปญฺญตฺติ}{mtk.122}
\bookmark[dest={mtk.122},level=3]{อนุปญฺญตฺติ}
\csromantocmarkref{subsubsection}{สิกฺขาปทวิภงฺค, ปทภาชนีย}{mtk.123}
\bookmark[dest={mtk.123},level=3]{สิกฺขาปทวิภงฺค, ปทภาชนีย}
\proseitem{473}{เตน โข ปน สมเยน อญฺญตโร ภิกฺขุ โกสมฺพิยํ คิลาโน โหติ, ญาตกา ตสฺส ภิกฺขุโน สนฺติเก ทูตํ ปาเหสุํ “อาคจฺฉตุ ภทนฺโต มยํ อุปฏฺฐหิสฺสามา”ติ. ภิกฺขูปิ เอวมาหํสุ \csromanfolio{298}“คจฺฉาวุโส ญาตกา ตํ อุปฏฺฐหิสฺสนฺตี”ติ. โส เอวมาห “ภควตาวุโส สิกฺขาปทํ ปญฺญตฺตํ ‘น ติจีวเรน วิปฺปวสิตพฺพนฺ’ติ, อหญฺจมฺหิ คิลาโน \textbf{น สกฺโกมิ ติจีวรํ อาทาย ปกฺกมิตุํ, นาหํ คมิสฺสามี”ติ. ภควโต} \textbf{เอตมตฺถํ อาโรเจสุํ. อถ โข ภควา เอตสฺมิํ นิทาเน เอตสฺมิํ ปกเรเณ} ธมฺมิํ กถํ กตฺวา ภิกฺขู อามนฺเตสิ “อนุชานามิ ภิกฺขเว คิลานสฺส ภิกฺขุโน ติจีวเรน อวิปฺปวาส\-สมฺมุติํ ทาตุํ, เอวญฺจ ปน ภิกฺขเว ทาตพฺพา, เตน คิลาเนน ภิกฺขุนา สํฆํ อุปสงฺกมิตฺวา เอกํสํ อุตฺตราสงฺคํ กริตฺวา วุฑฺฒานํ ภิกฺขูนํ ปาเท วนฺทิตฺวา อุกฺกุฏิกํ นิสีทิตฺวา อญฺชลิํ ปคฺคเหตฺวา เอวมสฺส วจนีโย “อหํ ภนฺเต คิลาโน น สกฺโกมิ ติจีวรํ อาทาย ปกฺกมิตุํ, โสหํ ภนฺเต สํฆํ ติจีวเรน อวิปฺปวาส\-สมฺมุติํ ยาจามี”ติ. ทุติยมฺปิ ยาจิตพฺพา. ตติยมฺปิ ยาจิตพฺพา. พฺยตฺเตน ภิกฺขุนา ปฏิพเลน สํโฆ ญาเปตพฺโพ\csromandash{}}

\proseitem{474}{“สุณาตุ เม ภนฺเต สํโฆ, อยํ อิตฺถนฺนาโม ภิกฺขุ คิลาโน น สกฺโกติ ติจีวรํ อาทาย ปกฺกมิตุํ, โส สํฆํ ติจีวเรน อวิปฺปวาส\-สมฺมุติํ ยาจติ, ยทิ สํฆสฺส ปตฺตกลฺลํ, สํโฆ อิตฺถนฺนามสฺส ภิกฺขุโน ติจีวเรน อวิปฺปวาส\-สมฺมุติํ ทเทยฺย, เอสา ญตฺติ.}

\prose{สุณาตุ เม ภนฺเต สํโฆ อยํ อิตฺถนฺนาโม ภิกฺขุ คิลาโน น สกฺโกติ ติจีวรํ อาทาย ปกฺกมิตุํ, โส สํฆํ ติจีวเรน อวิปฺปวาส\-สมฺมุติํ ยาจติ, สํโฆ อิตฺถนฺนามสฺส ภิกฺขุโน ติจีวเรน อวิปฺปวาส\-สมฺมุติํ เทติ, ยสฺสายสฺมโต ขมติ อิตฺถนฺนามสฺส ภิกฺขุโน ติจีวเรน อวิปฺปวาส\-สมฺมุติยา ทานํ, โส ตุณฺหสฺส, ยสฺส นกฺขมติ, โส ภาเสยฺย.}

\prose{ทินฺนา สํเฆน อิตฺถนฺนามสฺส ภิกฺขุโน ติจีวเรน อวิปฺปวาส\-สมฺมุติ, ขมติ สํฆสฺส, ตสฺมา ตุณฺหี, เอวเมตํ ธารยามี”ติ.}

\prose{เอวญฺจ ปน ภิกฺขเว อิมํ สิกฺขาปทํ อุทฺทิเสยฺยาถ\csromandash{}}

\proseitem{475}{\textbf{“นิฏฺฐิต}\-\textbf{จีวรสฺมิํ ภิกฺขุนา อุพฺภตสฺมิํ กถิเน เอกรตฺตํปิ เจ ภิกฺขุ ติจีวเรน วิปฺปวเสยฺย, อญฺญตฺร ภิกฺขุ}\-\textbf{สมฺมุติยา นิสฺสคฺคิยํ ปาจิตฺติยนฺ”}ติ.}

\proseitem{476}{\textbf{นิฏฺฐิต}\-\textbf{จีวร}\-\textbf{สฺมินฺ}ติ ภิกฺขุโน จีวรํ กตํ วา โหติ นฏฺฐํ วา วินฏฺฐํ วา ทฑฺฒํ วา จีวราสา วา อุปจฺฉินฺนา.}

\prose{\csromanfolio{299}\textbf{อุพฺภตสฺมิํ กถิเน}ติ อฏฺฐนฺนํ มาติกานํ อญฺญตราย มาติกาย อุพฺภตํ โหติ, สํเฆน วา อนฺตรา อุพฺภตํ โหติ.}

\prose{\textbf{เอกรตฺตํปิ เจ ภิกฺขุ ติจีวเรน วิปฺปวเสยฺยา}ติ สํฆาฏิยา วา อุตฺตราสงฺเคน วา อนฺตรวาสเกน วา.}

\prose{\textbf{อญฺญตฺร ภิกฺขุ}\-\textbf{สมฺมุติยา}ติ ฐเปตฺวา ภิกฺขุ\-สมฺมุติํ.}

\prose{\textbf{นิสฺสคฺคิยํ โหตี}ติ สห อรุณุคฺคมนา\csromansharedfootnote{852db0cfd3a93aa3}{อรุณุคฺคมเนน (สี, สฺยา)} นิสฺสคฺคิยํ โหติ. นิสฺสชฺชิตพฺพํ สํฆสฺส วา คณสฺส วา ปุคฺคลสฺส วา. เอวญฺจ ปน ภิกฺขเว นิสฺสชฺชิตพฺพํ ฯเปฯ “อิทํ เม ภนฺเต จีวรํ รตฺติ\-วิปฺปวุตฺถํ\csromansharedfootnote{c9e125122a668deb}{รตฺติํ วิปฺปวุตฺถํ (สี)} อญฺญตฺร ภิกฺขุ\-สมฺมุติยา นิสฺสคฺคิยํ, อิมาหํ สํฆสฺส นิสฺสชฺชามี”ติ ฯเปฯ ทเทยฺยาติ ฯเปฯ ทเทยฺยุนฺติ ฯเปฯ อายสฺมโต ทมฺมีติ.}

\proseitem{477}{คาโม เอกูปจาโร นานูปจาโร. นิเวสนํ เอกูปจารํ นานูปจารํ. อุโทสิโต เอกูปจาโร นานูปจาโร. อฏฺโฏ เอกูปจาโร นานูปจาโร. มาโฬ เอกูปจาโร นานูปจาโร. ปาสาโท เอกูปจาโร นานูปจาโร. หมฺมิยํ เอกูปจารํ นานูปจารํ. นาวา เอกูปจารา นานูปจารา. สตฺโถ เอกูปจาโร นานูปจาโร. เขตฺตํ เอกูปจารํ นานูปจารํ. ธญฺญกรณํ เอกูปจารํ นานูปจารํ. อาราโม เอกูปจาโร นานูปจาโร. วิหาโร เอกูปจาโร นานูปจาโร. รุกฺขมูลํ เอกูปจารํ นานูปจารํ. อชฺโฌกาโส เอกูปจาโร นานูปเจโร.}

\proseitem{478}{\textbf{คาโม เอกูปจาโร} นาม เอกกุลสฺส คาโม โหติ ปริกฺขิตฺโต จ. อนฺโตคาเม จีวรํ นิกฺขิปิตฺวา อนฺโตคาเม วตฺถพฺพํ. อปริกฺขิตฺโต โหติ, ยสฺมิํ ฆเร จีวรํ นิกฺขิตฺตํ โหติ, ตสฺมิํ ฆเร วตฺถพฺพํ, หตฺถปาสา วา น วิชหิตพฺพํ.}

\proseitem{479}{นานากุลสฺส คาโม โหติ ปริกฺขิตฺโต จ, ยสฺมิํ ฆเร จีวรํ นิกฺขิตฺตํ โหติ, ตสฺมิํ ฆเร วตฺถพฺพํ สภาเย วา ทฺวารมูเล วา, หตฺถปาสา วา น วิชหิตพฺพํ. สภายํ คจฺฉนฺเตน หตฺถปาเส จีวรํ นิกฺขิปิตฺวา สภาเย วา วตฺถพฺพํ ทฺวารมูเล วา, หตฺถปาสา วา น วิชหิตพฺพํ. สภาเย จีวรํ นิกฺขิปิตฺวา สภาเย วา วตฺถพฺพํ ทฺวารมูเล วา, หตฺถปาสา วา น วิชหิตพฺพํ. อปริกฺขิตฺโต โหติ, ยสฺมิํ ฆเร จีวรํ นิกฺขิตฺตํ โหติ, ตสฺมิํ ฆเร วตฺถพฺพํ, หตฺถปาสา วา น วิชหิตพฺพํ.}

\proseitem{480}{\csromanfolio{300}เอกกุลสฺส นิเวสนํ โหติ ปริกฺขิตฺตญฺจ, นานาคพฺภา นานา- โอวรกา. อนฺโตนิเวสเน จีวรํ นิกฺขิปิตฺวา อนฺโตนิเวสเน วตฺถพฺพํ. อปริกฺขิตฺตํ โหติ, ยสฺมิํ คพฺเภ จีวรํ นิกฺขิตฺตํ โหติ, ตสฺมิํ คพฺเภ วตฺถพฺพํ, หตฺถปาสา วา น วิชหิตพฺพํ.}

\proseitem{481}{นานากุลสฺส นิเวสนํ โหติ ปริกฺขิตฺตญฺจ, นานาคพฺภา นานา- โอวรกา. ยสฺมิํ คพฺเภ จีวรํ นิกฺขิตฺตํ โหติ, ตสฺมิํ คพฺเภ วตฺถพฺพํ ทฺวารมูเล วา, หตฺถปาสา วา น วิชหิตพฺพํ. อปริกฺขิตฺตํ โหติ, ยสฺมิํ คพฺเภ จีวรํ นิกฺขิตฺตํ โหติ, ตสฺมิํ คพฺเภ วตฺถพฺพํ, หตฺถปาสา วา น วิชหิตพฺพํ.}

\proseitem{482}{เอกกุลสฺส อุโทสิโต โหติ ปริกฺขิตฺโต จ, นานาคพฺภา นานาโอวรกา. อนฺโตอุโทสิเต จีวรํ นิกฺขิปิตฺวา อนฺโตอุโทสิเต วตฺถพฺพํ. อปริกฺขิตฺโต โหติ, ยสฺมิํ คพฺเภ จีวรํ นิกฺขิตฺตํ โหติ, ตสฺมิํ คพฺเภ วตฺถพฺพํ, หตฺถปาสา วา น วิชหิตพฺพํ.}

\proseitem{483}{นานากุลสฺส อุโทสิโต โหติ ปริกฺขิตฺโต จ, นานาคพฺภา นานา- โอวรกา. ยสฺมิํ คพฺเภ จีวรํ นิกฺขิตฺตํ โหติ, ตสฺมิํ คพฺเภ วตฺถพฺพํ ทฺวารมูเล วา, หตฺถปาสา วา น วิชหิตพฺพํ. อปริกฺขิตฺโต โหติ, ยสฺมิํ คพฺเภ จีวรํ นิกฺขิตฺตํ โหติ, ตสฺมิํ คพฺเภ วตฺถพฺพํ, หตฺถปาสา วา น วิชหิตพฺพํ.}

\proseitem{484}{เอกกุลสฺส อฏฺโฏ โหติ, อนฺโตอฏฺเฏ จีวรํ นิกฺขิปิตฺวา อนฺโตอฏฺเฏ วตฺถพฺพํ. นานากุลสฺส อฏฺโฏ โหติ, นานาคพฺภา นานาโอวรกา. ยสฺมิํ คพฺเภ จีวรํ นิกฺขิตฺตํ โหติ, ตสฺมิํ คพฺเภ วตฺถพฺพํ ทฺวารมูเล วา, หตฺถปาสา วา น วิชหิตพฺพํ.}

\proseitem{485}{เอกกุลสฺส มาโฬ โหติ, อนฺโตมาเฬ จีวรํ นิกฺขิปิตฺวา อนฺโตมาเฬ วตฺถพฺพํ. นานากุลสฺส มาโฬ โหติ, นานาคพฺภา นานาโอวรกา. ยสฺมิํ คพฺเภ จีวรํ นิกฺขิตฺตํ โหติ, ตสฺมิํ คพฺเภ วตฺถพฺพํ ทฺวารมูเล วา, หตฺถปาสา วา น วิชหิตพฺพํ.}

\proseitem{486}{เอกกุลสฺส ปาสาโท โหติ, อนฺโตปาสาเท จีวรํ นิกฺขิปิตฺวา อนฺโตปาสาเท วตฺตพฺพํ. นานากุลสฺส ปาสาโท โหติ, \csromanfolio{301}นานาคพฺภา นานาโอวรกา. ยสฺมิํ คพฺเภ จีวรํ นิกฺขิตฺตํ โหติ, ตสฺมิํ คพฺเภ วตฺถพฺพํ ทฺวารมูเล วา, หตฺถปาสา วา น วิชหิตพฺพํ.}

\proseitem{487}{เอกกุลสฺส หมฺมิยํ โหติ, อนฺโตหมฺมิเย จีวรํ นิกฺขิปิตฺวา อนฺโตหมฺมิเย วตฺถพฺพํ. นานากุลสฺส หมฺมิยํ โหติ, นานาคพฺภา นานา- โอวรกา. ยสฺมิํ คพฺเภ จีวรํ นิกฺขิตฺตํ โหติ, ตสฺมิํ คพฺเภ วตฺถพฺพํ ทฺวารมูเล วา, หตฺถปาสา วา น วิชหิตพฺพํ.}

\proseitem{488}{เอกกุลสฺส นาวา โหติ, อนฺโตนาวาย จีวรํ นิกฺขิปิตฺวา อนฺโตนาวาย วตฺถพฺพํ. นานากุลสฺส นาวา โหติ, นานาคพฺภา นานาโอวรกา. ยสฺมิํ โอวรเก จีวรํ นิกฺขิตฺตํ โหติ, ตสฺมิํ โอวรเก วตฺถพฺพํ, หตฺถปาสา วา น วิชหิตพฺพํ.}

\proseitem{489}{เอกกุลสฺส สตฺโถ โหติ, สตฺเถ จีวรํ นิกฺขิปิตฺวา ปุรโต วา ปจฺฉโต วา สตฺตพฺภนฺตรา น วิชหิตพฺพา, ปสฺสโต อพฺภนฺตรํ น วิชหิตพฺพํ. นานากุลสฺส สตฺโถ โหติ, สตฺเถ จีวรํ นิกฺขิปิตฺวา หตฺถปาสา น วิชหิตพฺพํ.}

\proseitem{490}{เอกกุลสฺส เขตฺตํ โหติ ปริกฺขิตฺตญฺจ, อนฺโตเขตฺเต จีวรํ นิกฺขิปิตฺวา อนฺโตเขตฺเต วตฺถพฺพํ. อปริกฺขิตฺตํ โหติ, หตฺถปาสา น วิชหิตพฺพํ. นานากุลสฺส เขตฺตํ โหติ ปริกฺขิตฺตญฺจ, อนฺโตเขตฺเต จีวรํ นิกฺขิปิตฺวา ทฺวารมูเล วา วตฺถพฺพํ, หตฺถปาสา วา น วิชหิตพฺพํ. อปริกฺขิตฺตํ โหติ, หตฺถปาสา น วิชหิตพฺพํ.}

\proseitem{491}{เอกกุลสฺส ธญฺญกรณํ โหติ ปริกฺขิตฺตญฺจ, อนฺโต\-ธญฺญกรเณ จีวรํ นิกฺขิปิตฺวา อนฺโต\-ธญฺญกรเณ วตฺถพฺพํ. อปริกฺขิตฺตํ โหติ, หตฺถปาสา น วิชหิตพฺพํ. นานากุลสฺส ธญฺญกรณํ โหติ ปริกฺขิตฺตญฺจ, อนฺโต\-ธญฺญกรเณ จีวรํ นิกฺขิปิตฺวา ทฺวารมูเล วา วตฺถพฺพํ, หตฺถปาสา วา น วิชหิตพฺพํ. อปริกฺขิตฺตํ โหติ, หตฺถวาสา น วิชหิตพฺพํ.}

\proseitem{492}{เอกกุลสฺส อาราโม โหติ ปริกฺขิตฺโต จ, อนฺโตอาราเม จีวรํ นิกฺขิปิตฺวา อนฺโตอาราเม วตฺถพฺพํ. อปริกฺขิตฺโต โหติ, หตฺถปาสา น วิชหิตพฺพํ. นานากุลสฺส อาราโม โหติ ปริกฺขิตฺโต จ, อนฺโตอาราเม จีวรํ นิกฺขิปิตฺวา ทฺวารมูเล วา วตฺถพฺพํ, หตฺถปาสา วา น วิชหิตพฺพํ. อปริกฺขิตฺโต โหติ, หตฺถปาสา น วิชหิตพฺพํ.}

\proseitem{493}{\csromanfolio{302}เอกกุลสฺส วิหาโร โหติ ปริกฺขิตฺโต จ, อนฺโตวิหาเร จีวรํ นิกฺขิปิตฺวา อนฺโตวิหาเร วตฺถพฺพํ. อปริกฺขิตฺโต โหติ, ยสฺมิํ วิหาเร จีวรํ นิกฺขิตฺตํ โหติ, ตสฺมิํ วิหาเร วตฺถพฺพํ, หตฺถปาสา วา น วิชหิตพฺพํ. นานากุลสฺส วิหาโร โหติ ปริกฺขิตฺโต จ, ยสฺมิํ วิหาเร จีวรํ นิกฺขิตฺตํ โหติ, ตสฺมิํ วิหาเร วตฺถพฺพํ ทฺวารมูเล วา, หตฺถปาสา วา น วิชหิตพฺพํ. อปริกฺขิตฺโต โหติ, ยสฺมิํ วิหาเร จีวรํ นิกฺขิตฺตํ โหติ, ตสฺมิํ วิหาเร วตฺถพฺพํ, หตฺถปาสา วา น วิชหิตพฺพํ.}

\proseitem{494}{เอกกุลสฺส รุกฺขมูลํ โหติ, ยํ มชฺฌนฺหิเก กาเล สมนฺตา ฉายา ผรติ, อนฺโตฉายาย จีวรํ นิกฺขิปิตฺวา อนฺโตฉายาย วตฺถพฺพํ. นานากุลสฺส รุกฺขมูลํ โหติ, หตฺถปาสา น วิชหิตพฺพํ.}

\prose{อชฺโฌกาโส เอกูปจาโร นาม อคามเก อรญฺเญ สมนฺตา สตฺตพฺภนฺตรา เอกูปจาโร. ตโต ปรํ นานูปจาโร.}

\proseitem{495}{วิปฺปวุตฺเถ วิปฺปวุตฺถ\-สญฺญี อญฺญตฺร ภิกฺขุ\-สมฺมุติยา, นิสฺสคฺคิยํ ปาจิตฺติยํ. วิปฺปวุตฺเถ เวมติโก อญฺญตฺร ภิกฺขุ\-สมฺมุติยา, นิสฺสคฺคิยํ ปาจิตฺติยํ. วิปฺปวุตฺเถ อวิปฺปวุตฺถ\-สญฺญี อญฺญตฺร ภิกฺขุ\-สมฺมุติยา, นิสฺสคฺคิยํ ปาจิตฺติยํ. อปฺปจฺจุทฺธเฏ ปจฺจุทฺธฏ\-สญฺญี ฯเปฯ. อวิสฺสชฺชิเต วิสฺสชฺชิต\-สญฺญี. อนฏฺเฐ นฏฺฐสญฺญี. อวินฏฺเฐ วินฏฺฐสญฺญี. อทฑฺเฒ ทฑฺฒสญฺญี. อวิลุตฺเต วิลุตฺตสญฺญี อญฺญตฺร ภิกฺขุ\-สมฺมุติยา, นิสฺสคิยํ ปาจิตฺติยํ.}

\prose{นิสฺสคฺคิยํ จีวรํ อนิสฺสชฺชิตฺวา ปริภุญฺชติ, อาปตฺติ ทุกฺกฏสฺส. อวิปฺปวุตฺเถ วิปฺปวุตฺถ\-สญฺญี, อาปตฺติ ทุกฺกฏสฺส. อวิปฺปวุตฺเถ เวมติโก, อาปตฺติ ทุกฺกฏสฺส. อวิปฺปวุตฺเถ อวุปฺปวุตฺถสญฺญี, อนาปตฺติ.}

\proseitem{496}{อนาปตฺติ อนฺโตอรุเณ ปจฺจุทฺธรติ, วิสฺสชฺเชติ, นสฺสติ, วินสฺสติ, ฑยฺหติ, อจฺฉินฺทิตฺวา คณฺหนฺติ, วิสฺสาสํ คณฺหนฺติ, ภิกฺขุ\-สมฺมุติยา, อุมฺมตฺตกสฺส อาทิกมฺมิกสฺสาติ.}

\vspace{\tipitakaparskip}
\csromancenter{อุโทสิต\-สิกฺขาปทํ นิฏฺฐิตํ ทุติยํ.}
\csromanmatikaanchor{mtk.124}
\csromantocmarknopage{subsection}{3. ตติยกถินสิกฺขาปท}{mtk.124}
\bookmark[dest={mtk.124},level=2]{3. ตติยกถินสิกฺขาปท}
\csromanmarkh{1. จีวรวคฺค}\csromanmarkhii{3. ตติย\-กถิน\-สิกฺขาปท}\csromanheadernomark{2}{1. \textbf{จีวรวคฺค 3. ตติย}\-\textbf{กถิน}\-\textbf{สิกฺขาปท}}
\csromanmatikaanchor{mtk.125}
\csromantocmarkref{subsubsection}{อญฺญตรภิกฺขุวตฺถุ}{mtk.125}
\bookmark[dest={mtk.125},level=3]{อญฺญตรภิกฺขุวตฺถุ}
\proseitem{497}{\csromanfolio{303}เตน สมเยน พุทฺโธ ภควา สาวตฺถิยํ วิหรติ เชตวเน อนาถ\-ปิณฺฑิกสฺส อาราเม. เตน โข ปน สมเยน อญฺญตรสฺส ภิกฺขุโน อกาลจีวรํ อุปฺปนฺนํ โหติ, ตสฺส ตํ จีวรํ กยิรมานํ นปฺปโหติ. อถ โข โส ภิกฺขุ ตํ จีวรํ อุสฺสาเปตฺวา ปุนปฺปุนํ วิมชฺชติ. อทฺทสา โข ภควา เสนาสนจาริกํ อาหิณฺฑนฺโต ตํ ภิกฺขุํ ตํ จีวรํ อุสฺสาเปตฺวา ปุนปฺปุนํ วิมชฺชนฺตํ, ทิสฺวาน เยน โส ภิกฺขุ เตนุปสงฺกมิ, อุปสงฺกมิตฺวา ตํ ภิกฺขุํ เอตทโวจ “กิสฺส ตฺวํ ภิกฺขุ อิมํ จีวรํ อุสฺสาเปตฺวา ปุนปฺปุนํ วิมชฺชสี”ติ. อิทํ เม ภนฺเต อกาลจีวรํ อุปฺปนฺนํ กยิรมานํ นปฺปโหติ, เตนาหํ อิมํ จีวรํ อุสฺสาเปตฺวา ปุนปฺปุนํ วิมชฺชามีติ. อตฺถิ ปน เต ภิกฺขุ จีวรปจฺจาสาติ. อตฺถิ ภควาติ. อถ โข ภควา เอตสฺมิํ นิทาเน เอตสฺมิํ ปกรเณ ธมฺมิํ กถํ กตฺวา ภิกฺขู อามนฺเตสิ “อนุชานามิ ภิกฺขเว อกาลจีวรํ ปฏิคฺคเหตฺวา จีวรปจฺจาสา นิกฺขิปิตุนฺ”ติ.}

\proseitem{498}{เตน โข ปน สมเยน ภิกฺขู ภควตา อนุญฺญตํ “อกาลจีวรํ ปฏิคฺคเหตฺวา จีวรปจฺจาสา นิกฺขิปิตุนฺ”ติ อกาลจีวรานิ ปฏิคฺคเหตฺวา อติเรกมาสํ นิกฺขิปนฺติ, ตานิ จีวรานิ จีวรวํเส ภณฺฑิกา\-พทฺธานิ ติฏฺฐนฺติ. อทฺทสา โข อายสฺมา อานนฺโท เสนาสนจาริกํ อาหิณฺฑนฺโต ตานิ จีวรานิ จีวรวํเส ภณฺฑิกา\-พทฺธานิ ติฏฺฐนฺเต. ทิสฺวาน ภิกฺขู อามนฺเตสิ “กสฺสิมานิ อาวุโส จีวรานิ จีวรวํเส ภณฺฑิกา\-พทฺธานิ ติฏฺฐนฺตี”ติ. อมฺหากํ อาวุโส อกาลจีวรานิ จีวรปจฺจาสา นิกฺขิตฺตานีติ. กีวจิรํ ปนาวุโส อิมานิ จีวรานิ นิกฺขิตฺตานีติ. อติเรกมาสํ อาวุโสติ. อายสฺมา อานนฺโท อุชฺฌายติ ขิยฺยติ วิปาเจติ “กถํ หิ นาม ภิกฺขู อกาลจีวรํ ปฏิคฺคเหตฺวา อติเรกมาสํ นิกฺขิปิสฺสนฺตี”ติ. อถ โข อายสฺมา อานนฺโท เต ภิกฺขู อเนกปริยเยน วิครหิตฺวา ภควโต เอตมตฺถํ อาโรเจสิ ฯเปฯ “สจฺจํ กิร ภิกฺขเว ภิกฺขู อกาลจีวรํ ปฏิคฺคเหตฺวา อติเรกมาสํ นิกฺขิปนฺตี”ติ. สจฺจํ ภควาติ. วิครหิ พุทฺโธ ภควา ฯเปฯ. “กถํ หิ นาม เต ภิกฺขเว โมฆปุริสา อกาลจีวรํ ปฏิคฺคเหตฺวา อติเรกมาสํ นิกฺขิปิสฺสนฺติ, เนตํ ภิกฺขเว อปฺปสนฺนานํ วา ปสาทาย ฯเปฯ. เอวญฺจ ปน ภิกฺขเว อิมํ สิกฺขาปทํ อุทฺทิเสยฺยาถ\csromandash{}}

\csromanmatikaanchor{mtk.126}
\csromanmatikaanchor{mtk.127}
\csromantocmarkref{subsubsection}{ปญฺญตฺติ}{mtk.126}
\bookmark[dest={mtk.126},level=3]{ปญฺญตฺติ}
\csromantocmarkref{subsubsection}{สิกฺขาปทวิภงฺค, ปทภาชนีย}{mtk.127}
\bookmark[dest={mtk.127},level=3]{สิกฺขาปทวิภงฺค, ปทภาชนีย}
\proseitem{499}{\csromanfolio{304}\textbf{“นิฏฺฐิต}\-\textbf{จีวรสฺมิํ ภิกฺขุนา อุพฺภตสฺมิํ กถิเน ภิกฺขุโน ปเนว อกาลจีวรํ อุปฺปชฺเชยฺย, อากงฺขมาเนน ภิกฺขุนา ปฏิคฺคเหตพฺพํ, ปฏิคฺคเหตฺวา ขิปฺปเมว กาเรตพฺพํ, โนจสฺส ปาริปูริ, มาสปรมํ เตน ภิกฺขุนา ตํ จีวรํ นิกฺขิปิตพฺพํ อูนสฺส ปาริปูริยา สติยา ปจฺจาสาย. ตโต เจ อุตฺตริ}\csromansharedfootnote{721fc33e8dfeba2a}{อุตฺตริํ (สี, สฺยา)} \textbf{นิกฺขิเปยฺย สติยาปิ ปจฺจาสาย, นิสฺสคฺคิยํ ปาจิตฺติยนฺ”}ติ.}

\proseitem{500}{\textbf{นิฏฺฐิต}\-\textbf{จีวร}\-\textbf{สฺมินฺ}ติ ภิกฺขุโน จีวรํ กตํ วา โหติ นฏฺฐํ วา วินฏฺฐํ วา ทฑฺฒํ วา จีวราสา วา อุปจฺฉินฺนา.}

\prose{\textbf{อุพฺภตสฺมิํ กถิเน}ติ อฏฺฐนฺนํ มาติกานํ อญฺญตราย มาติกาย อุพฺภตํ โหติ, สํเฆน วา อนฺตรา อุพฺภตํ โหติ.}

\prose{\textbf{อกาลจีวรํ} นาม อนตฺถเต กถิเน เอกาทสมาเส อุปฺปนฺนํ, อตฺถเต กถิเน สตฺตมาเส อุปฺปนฺนํ, กาเลปิ อาทิสฺส ทินฺนํ, เอตํ อกาลจีวรํ นาม.}

\prose{\textbf{อุปฺปชฺเชยฺยา}ติ อุปฺปชฺเชยฺย สํฆโต วา คณโต วา ญาติโต วา มิตฺตโต วา ปํสุกูลํ วา อตฺตโน วา ธเนน.}

\prose{\textbf{อากงฺข}\-\textbf{มาเนนา}ติ อิจฺฉมาเนน ปฏิคฺคเหตพฺพํ.}

\prose{\textbf{ปฏิคฺคเหตฺวา ขิปฺปเมว กาเรตพฺพนฺ}ติ ทสาหา กาเรตพฺพํ.}

\prose{\textbf{โน จสฺส ปาริปูรี}ติ กยิรมานํ นปฺปโหติ.}

\prose{\textbf{มาสปรมํ เตน ภิกฺขุนา ตํ จีวรํ นิกฺขิปิตพฺพนฺ}ติ มาสปรมตา นิกฺขิปิตพฺพํ.}

\prose{อู\textbf{นสฺส ปาริปูริยา}ติ อูนสฺส ปาริปูรตฺถาย.}

\prose{\textbf{สติยา ปจฺจาสายา}ติ ปจฺจาสา โหติ สํฆโต วา คณโต วา ญาติโต วา มิตฺตโต วา ปํสุกูลํ วา อตฺตโน วา ธเนน.}

\prose{\textbf{ตโต เจ อุตฺตริ นิกฺขิเปยฺย สติยาปิ ปจฺจาสายา}ติ ตทหุปฺปนฺเน มูลจีวเร ปจฺจาสาจีวรํ อุปฺปชฺชติ, ทสาหา กาเรตพฺพํ. ทฺวีหุปฺปนฺเน มูลจีวเร ปจฺจาสาจีวรํ อุปฺปชฺชติ, ทสาหา กาเรตพฺพํ. ตีหุปฺปนฺเน มูลจีวเร ฯเปฯ. จตูหุปฺปนฺเน. ปญฺจาหุปฺปนฺเน. ฉาหุปฺปนฺเน. สตฺตาหุปฺปนฺเน. \csromanfolio{305}\textbf{อฏฺฐาหุปฺปนฺเน. นวาหุปฺปนฺเน. ทสาหุปฺปนฺเน มูลจีวเร ปจฺจาสาจีวรํ} \textbf{อุปฺปชฺชติ, ทสาหา กาเรตพฺพํ. เอกาทเส อุปฺปนฺเน ฯเปฯ. ทฺวาทเส อุปฺปนฺเน.} เตรเส อุปฺปนฺเน. จุทฺทเส อุปฺปนฺเน. ปนฺนรเส อุปฺปนฺเน. โสฬเส อุปฺปนฺเน. สตฺตรเส อุปฺปนฺเน. อฏฺฐารเส อุปฺปนฺเน. เอกูนวีเส อุปฺปนฺเน. วีเส อุปฺปนฺเน มูลจีวเร ปจฺจาสาจีวรํ อุปฺปชฺชติ, ทสาหา กาเรตพฺพํ. เอกวีเส อุปฺปนฺเน มูลจีวเร ปจฺจาสาจีวรํ อุปฺปชฺชติ, นวาหา กาเรตพฺพํ. ทฺวาวีเส อุปฺปนฺเน ฯเปฯ. เตวีเส อุปฺปนฺเน. จตุวีเส อุปฺปนฺเน. ปญฺจวีเส อุปฺปนฺเน. ฉพฺพีเส อุปฺปนฺเน. สตฺตวีเส อุปฺปนฺเน. อฏฺฐวีเส อุปฺปนฺเน. เอกูนติํเส อุปฺปนฺเน มูลจีวเร ปจฺจาสาจีวรํ อุปฺปชฺชติ, เอกาหํ กาเรตพฺพํ. ติํเส อุปฺปนฺเน มูลจีวเร ปจฺจาสาจีวรํ อุปฺปชฺชติ, ตทเหว อธิฏฺฐาตพฺพํ วิกปฺเปตพฺพํ วิสฺสชฺเชตพฺพํ. โน เจ อธิฏฺเฐยฺย วา วิกปฺเปยฺย วา วิสฺสชฺเชยฺย วา, เอกติํเส อรุณุคฺคมเน นิสฺสคฺคิยํ โหติ. นิสฺสชฺชิตพฺพํ สํฆสฺส วา คณสฺส วา ปุคฺคลสฺส วา. เอวญฺจ ปน ภิกฺขเว นิสฺสชฺชิตพฺพํ ฯเปฯ “อิทํ เม ภนฺเต อกาลจีวรํ มาสาติกฺกนฺตํ นิสฺสคฺคิยํ, อิมาหํ สํฆสฺส นิสฺสชฺชามี”ติ ฯเปฯ ทเทยฺยาติ ฯเปฯ ทเทยฺยุนฺติ ฯเปฯ อายสฺมโต ทมฺมีติ.}

\prose{วิสภาเค อุปฺปนฺเน มูลจีวเร ปจฺจาสาจีวรํ อุปฺปชฺชติ, รตฺติโย จ เสสา โหนฺติ, น อกามา กาเรตพฺพํ.}

\proseitem{501}{มาสาติกฺกนฺเต อติกฺกนฺตสญฺญี, นิสฺสคฺคิยํ ปาจิตฺติยํ. มาสาติกฺกนฺเต เวมติโก, นิสฺสคฺคิยํ ปาจิตฺติยํ. มาสาติกฺกนฺเต อนติกฺกนฺต\-สญฺญี, นิสฺสคฺคิยํ ปาจิตฺติยํ. อนธิฏฺฐิเต อธิฏฺฐิต\-สญฺญี ฯเปฯ. อวิกปฺปิเต วิกปฺปิตสญฺญี. อวิสฺสชฺชิเต วิสฺสชฺชิต\-สญฺญี. อนฏฺเฐ น ฏฺฐสญฺญี. อวินฏฺเฐ วินฏฺฐสญฺญี. อทฑฺเฒ ทฑฺฒสญฺญี. อวิลุตฺเต วิลุตฺตสญฺญี, นิสฺสคฺคิยํ ปาจิตฺติยํ.}

\prose{นิสฺสคฺคิยํ จีวรํ อนิสฺสชฺชิตฺวา ปริภุญฺชติ, อาปตฺติ ทุกฺกฏสฺส. มาสานติกฺกนฺเต อติกฺกนฺตสญฺญี, อาปตฺติ ทุกฺกฏสฺส. มาสานติกฺกนฺเต เวมติโก, อาปตฺติ ทุกฺกฏสฺส. มาสานติกฺกนฺเต อนติกฺกนฺต\-สญฺญี, อนาปตฺติ.}

\proseitemwithcloser{502}{อนาปตฺติ อนฺโตมาเส อธิฏฺเฐติ, วิกปฺเปติ, วิสฺสชฺเชติ, นสฺสติ, วินสฺสติ, ฑยฺหติ, อจฺฉินฺทิตฺวา คณฺหนฺติ, วิสฺสาสํ คณฺหนฺติ, อุมฺมตฺตกสฺส อาทิกมฺมิกสฺสาติ.}{ตติย\-กถิน\-สิกฺขาปทํ นิฏฺฐิตํ.}
\csromanmatikaanchor{mtk.128}
\csromantocmarknopage{subsection}{4. ปุราณจีวรสิกฺขาปท}{mtk.128}
\bookmark[dest={mtk.128},level=2]{4. ปุราณจีวรสิกฺขาปท}
\csromanmarkh{1. จีวรวคฺค}\csromanmarkhii{4. ปุราณจีวร\-สิกฺขาปท}\csromanheadernomark{2}{1. \textbf{จีวรวคฺค 4. ปุราณจีวร}\-\textbf{สิกฺขาปท}}
\csromanmatikaanchor{mtk.129}
\csromantocmarkref{subsubsection}{อุทายิภิกฺขุวตฺถุ}{mtk.129}
\bookmark[dest={mtk.129},level=3]{อุทายิภิกฺขุวตฺถุ}
\proseitem{503}{\csromanfolio{306}เตน สมเยน พุทฺโธ ภควา สาวตฺถิยํ วิหรติ เชตวเน อนาถ\-ปิณฺฑิกสฺส อาราเม. เตน โข ปน สมเยน อายสฺมโต อุทายิสฺส ปุราณทุติยิกา ภิกฺขุนีสุ ปพฺพชิตา โหติ, สา อายสฺมโต อุทายิสฺส สนฺติเก อภิกฺขณํ อาคจฺฉติ, อายสฺมาปิ อุทายี ตสฺสา ภิกฺขุนิยา สนฺติเก อภิกฺขณํ คจฺฉติ. เตน โข ปน สมเยน อายสฺมา อุทายี ตสฺสา ภิกฺขุนิยา สนฺติเก ภตฺตวิสฺสคฺคํ กโรติ, อถ โข อายสฺมา อุทายี ปุพฺพณฺหสมยํ นิวาเสตฺวา ปตฺต\-จีวรมา\-ทาย เยน สา ภิกฺขุนี เตนุปสงฺกมิ, อุปสงฺกมิตฺวา ตสฺสา ภิกฺขุนิยา ปุรโต องฺคชาตํ วิวริตฺวา อาสเน นิสีทิ, สาปิ โข ภิกฺขุนี อายสฺมโต อุทายิสฺส ปุรโต องฺคชาตํ วิวริตฺวา อาสเน นิสีทิ. อถ โข อายสฺมา อุทายี สารตฺโต ตสฺสา ภิกฺขุนิยา องฺคชาตํ อุปนิชฺฌายิ, ตสฺส อสุจิ มุจฺจิ. อถ โข อายสฺมา อุทายี ตํ ภิกฺขุนิํ เอตทโวจ “คจฺฉ ภคินิ อุทกํ อาหร, อนฺตรวาสกํ โธวิสฺสามี”ติ. “อาหรยฺย อหเมว โธวิสฺสามี”ติ ตํ อสุจิํ เอกเทสํ มุเขน อคฺคเหสิ, เอกเทสํ องฺคชาเต ปกฺขิปิ, สา เตน คพฺภํ คณฺหิ. ภิกฺขุนิโย เอวมาหํสุ “อพฺรหฺมจารินี อยํ ภิกฺขุนี คพฺภินี”ติ. “นาหํ อยฺเย อพฺรหฺมจารินี”ติ ภิกฺขุนีนํ เอตมตฺถํ อาโรเจสิ. ภิกฺขุนิโย อุชฺฌายนฺติ ขิยฺยนฺติ วิปาเจนฺติ “กถํ หิ นาม อยฺโย อุทายี ภิกฺขุนิยา ปุราณจีวรํ โธวาเปสฺสตี”ติ. อถ โข ตา ภิกฺขุนิโย ภิกฺขูนํ เอตมตฺถํ อาโรเจสุํ. เย เต ภิกฺขู อปฺปิจฺฉา ฯเปฯ เต อุชฺฌายนฺติ ขิยฺยนฺติ วิปาเจนฺติ “กถํ หิ นาม อายสฺมา อุทายี ภิกฺขุนิยา ปุราณจีวรํ โธวาเปสฺสตี”ติ. อถ โข เต ภิกฺขู อายสฺมนฺตํ อุทายิํ อเนก\-ปริยาเยน วิครหิตฺวา ภควโต เอตมตฺถํ อาโรเจสุํ ฯเปฯ “สจฺจํ กิร ตฺวํ อุทายิ ภิกฺขุนิยา ปุราณจีวรํ โธวาเปสี”ติ. สจฺจํ ภควาติ. ญาติกา เต อุทายิ อญฺญาติกาติ. อญฺญาติกา ภควาติ. อญฺญาตโก โมฆปุริส อญฺญาติกาย น ชานาติ ปติรูปํ วา อปฺปติรูปํ วา ปาสาทิกํ วา อปาสาทิกํ วา, ตตฺถ นาม ตฺวํ โมฆปุริส อญฺญาติกาย ภิกฺขุนิยา ปุราณจีวรํ โธวาเปสฺสสิ, เนตํ โมฆปุริส อปฺปสนฺนานํ วา ปสาทาย ฯเปฯ. เอวญฺจ ปน ภิกฺขเว อิมํ สิกฺขาปทํ อุทฺทิเสยฺยาถ\csromandash{}}

\csromanmatikaanchor{mtk.130}
\csromanmatikaanchor{mtk.131}
\csromantocmarkref{subsubsection}{ปญฺญตฺติ}{mtk.130}
\bookmark[dest={mtk.130},level=3]{ปญฺญตฺติ}
\csromantocmarkref{subsubsection}{สิกฺขาปทวิภงฺค, ปทภาชนีย}{mtk.131}
\bookmark[dest={mtk.131},level=3]{สิกฺขาปทวิภงฺค, ปทภาชนีย}
\proseitem{504}{\csromanfolio{307}\textbf{“โย ปน ภิกฺขุ อญฺญาติกาย ภิกฺขุนิยา ปุราณจีวรํ โธวาเปยฺย วา รชาเปยฺย วา อาโกฏาเปยฺย วา, นิสฺสคฺคิยํ ปาจิตฺติยนฺ”}ติ.}

\proseitem{505}{\textbf{โย ปนา}ติ โย ยาทิโส ฯเปฯ. พฺ\textbf{หิกฺขู}ติ ฯเปฯ อยํ อิมสฺมิํ อตฺเถ อธิปฺเปโต “ภิกฺขู”ติ.}

\prose{\textbf{อญฺญาติกา} นาม มาติโต วา ปิติโต วา ยาว สตฺตมา ปิตามหยุคา อสมฺพทฺธา.}

\prose{\textbf{ภิกฺขุนี} นาม อุภโตสํเฆ อุปสมฺปนฺนา.}

\prose{\textbf{ปุราณจีวรํ} นาม สกิํ นิวตฺถํปิ สกิํ ปารุตํปิ.}

\prose{โธวาติ อาณาเปติ, อาปตฺติ ทุกฺกฏสฺส. โธตํ นิสฺสคฺคิยํ โหติ. รชาติ อาณาเปติ, อาปตฺติ ทุกฺกฏสฺส. รตฺตํ นิสฺสคฺคิยํ โหติ. อาโกเฏหีติ อาณาเปติ, อาปตฺติ ทุกฺกฏสฺส. สกิํ ปาณิปฺปหารํ วา มุคฺคร\-ปฺปหารํ วา ทินฺเน นิสฺสคฺคิยํ โหติ. นิสฺสชฺชิตพฺพํ สํฆสฺส วา คณสฺส วา ปุคฺคลสฺส วา. เอวญฺจ ปน ภิกฺขเว นิสฺสชฺชิตพฺพํ ฯเปฯ “อิทํ เม ภนฺเต ปุราณจีวรํ อญฺญาติกาย ภิกฺขุนิยา โธวาปิตํ นิสฺสคฺคิยํ, อิมาหํ สํฆสฺส นิสฺสชฺชามี”ติ ฯเปฯ ทเทยฺยาติ ฯเปฯ ทเทยฺยุนฺติ ฯเปฯ อายสฺมโต ทมฺมีติ.}

\proseitem{506}{อญฺญาติกาย อญฺญาติกสญฺญี ปุราณจีวรํ โธวาเปติ, นิสฺสคฺคิยํ ปาวิตฺติยํ. อญฺญาติกาย อญฺญาติกสญฺญี ปุราณจีวรํ โธวาเปติ รชาเปติ, นิสฺสคฺคิเยน อาปตฺติ ทุกฺกฏสฺส. อญฺญาติกาย อญฺญาติกสญฺญี ปุราณจีวรํ โธวาเปติ อาโกฏาเปติ, นิสฺสคฺคิเยน อาปตฺติ ทุกฺกฏสฺส. อญฺญาติกาย อญฺญาติกสญฺญี ปุราณจีวรํ โธวาเปติ รชาเปติ อาโกฏาเปติ, นิสฺสคฺคิเยน อาปตฺติ ทฺวินฺนํ ทุกฺกฏานํ.}

\prose{อญฺญาติกาย อญฺญาติกสญฺญี ปุราณจีวรํ รชาเปติ, นิสฺสคฺคิยํ ปาจิตฺติยํ. อญฺญาติกาย อญฺญาติกสญฺญี ปุราณจีวรํ รชาเปติ อาโกฏาเปติ, นิสฺสคฺคิเยน อาปตฺติ ทุกฺกฏสฺส. อญฺญาติกาย อญฺญาติกสญฺญี ปุราณจีวรํ รชาเปติ โธวาเปติ, นิสฺสคฺคิเยน อาปตฺติ ทุกฺกฏสฺส. อญฺญาติกาย อญฺญาติกสญฺญี ปุราณจีวรํ รชาเปติ อาโกฏาเปติ โธวาเปติ, นิสฺสคฺคิเยน อาปตฺติ ทฺวินฺนํ ทุกฺกฏานํ.}

\csromanmatikaanchor{mtk.132}
\csromanmatikaanchor{mtk.133}
\csromantocmarknopage{subsection}{5. จีวรปฏิคฺคหณสิกฺขาปท}{mtk.132}
\bookmark[dest={mtk.132},level=2]{5. จีวรปฏิคฺคหณสิกฺขาปท}
\csromantocmarkref{subsubsection}{อุปฺปลวณฺณาภิกฺกุนีวตฺถุ}{mtk.133}
\bookmark[dest={mtk.133},level=3]{อุปฺปลวณฺณาภิกฺกุนีวตฺถุ}
\prose{\csromanfolio{308}อญฺญาติกาย อญฺญาติกสญฺญี ปุราณจีวรํ อาโกฏาเปติ, นิสฺสคฺคิยํ ปาจิตฺติยํ. อญฺญาติกาย อญฺญาติกสญฺญี ปุราณจีวรํ อาโกฏาเปติ โธวาเปติ, นิสฺสคฺคิเยน อาปตฺติ ทุกฺกฏสฺส. อญฺญาติกาย อญฺญาติกสญฺญี ปุราณจีวรํ อาโกฏาเปติ รชาเปติ, นิสฺสคฺคิเยน อาปตฺติ ทุกฺกฏสฺส. อญฺญาติกาย อญฺญาติกสญฺญี ปุราณจีวรํ อาโกฏาเปติ โธวาเปติ รชาเปติ, นิสฺสคฺคิเยน อาปตฺติ ทฺวินฺนํ ทุกฺกฏานํ.}

\prose{อญฺญาติกาย เวมติโก ฯเปฯ. อญฺญาติกาย ญาติกสญฺญี ฯเปฯ อญฺญสฺส ปุราณจีวรํ โธวาเปติ, อาปตฺติ ทุกฺกฏสฺส. นิสีทน\-ปจฺจตฺถรณํ โธวาเปติ, อาปตฺติ ทุกฺกฏสฺส. เอกโต\-อุปสมฺปนฺนาย โธวาเปติ, อาปตฺติ ทุกฺกฏสฺส. ญาติกาย อญฺญาติกสญฺญี, อาปตฺติ ทุกฺกฏสฺส. ญาติกาย เวมติโก, อาปตฺติ ทุกฺกฏสฺส. ญาติกาย ญาติกสญฺญี, อนาปตฺติ.}

\proseitemwithcloserruled{507}{อนาปตฺติ ญาติกาย โธวนฺติยา อญฺญาติกา ทุติยา โหติ, อวุตฺตา โธวติ, อปริภุตฺตํ โธวาเปติ, จีวรํ ฐเปตฺวา อญฺญํ ปริกฺขารํ โธวาเปติ, สิกฺขมานาย สามเณริยา อุมฺมตฺตกสฺส อาทิกมฺมิกสฺสาติ.}{ปุราณจีวร\-สิกฺขาปทํ นิฏฺฐิตํ จตุตฺถํ.}
\csromanmatikaanchor{mtk.134}
\csromanmatikaanchor{mtk.219}
\csromantocmarkref{subsubsection}{ปญฺญตฺติ}{mtk.134}
\bookmark[dest={mtk.134},level=3]{ปญฺญตฺติ}
\csromanmarkh{1. จีวรวคฺค}\csromanmarkhii{5. จีวร\-ปฏิคฺคหณ\-สิกฺขาปท}\csromanheadernomark{2}{1. \textbf{จีวรวคฺค 5. จีวร}\-\textbf{ปฏิคฺคหณ}\-\textbf{สิกฺขาปท}}
\proseitem{508}{เตน สมเยน พุทฺโธ ภควา ราชคเห วิหรติ เวฬุวเน กลนฺทก\-นิวาเป. เตน โข ปน สมเยน อุปฺปลวณฺณา ภิกฺขุนี สาวตฺถิยํ วิหรติ. อถ โข อุปฺปลวณฺณา ภิกฺขุนี ปุพฺพณฺหสมยํ นิวาเสตฺวา ปตฺต\-จีวรมา\-ทาย สาวตฺถิํ ปิณฺฑาย ปาวิสิ, สาวตฺถิยํ ปิณฺฑย จริตฺวา ปจฺฉาภตฺตํ ปิณฺฑปาต\-ปฺปฏิกฺกนฺตา เยน อนฺธวนํ เตนุปสงฺกมิ ทิวาวิหาราย, อนฺธวนํ อชฺโฌคาเหตฺวา อญฺญตรสฺมิํ รุกฺขมูเล ทิวาวิหารํ นิสีทิ. เตน โข ปน สมเยน โจรา กตกมฺมา คาวิํ วธิตฺวา มํสํ คเหตฺวา อนฺธวนํ ปวิสิํสุ. อทฺทสา โข โจรคามณิโก อุปฺปลวณฺณํ ภิกฺขุนิํ อญฺญตรสฺมิํ รุกฺขมูเล \csromanfolio{309}ทิวาวิหารํ นิสินฺนํ, ทิสฺวานสฺส เอตทโหสิ “สเจ เม ปุตฺตภาตุกา ปสฺสิสฺสนฺติ วิเหฐิสฺสนฺติ อิมํ ภิกฺขุนินฺ”ติ อญฺเญน มคฺเคน อคมาสิ. อถ โข โส โจรคามณิโก มํเส ปกฺเก วรมํสานิ คเหตฺวา ปณฺณปุฏํ\csromansharedfootnote{9ade05b3dbe18484}{ปณฺเณน ปุฏํ (สฺยา)} พนฺธิตฺวา อุปฺปลวณฺณาย ภิกฺขุนิยา อวิทูเร รุกฺเข อาลคฺเคตฺวา “โย ปสฺสติ สมโณ วา พฺราหฺมโณ วา, ทินฺนํ เยว หรตู”ติ วตฺวา ปกฺกามิ. อสฺโสสิ โข อุปฺปลวณฺณา ภิกฺขุนี สมาธิมฺหา วุฏฺฐหิตฺวา ตสฺส โจร\-คามณิกสฺส อิมํ วาจํ ภาสมานสฺส. อถ โข อุปฺปลวณฺณา ภิกฺขุนี ตํ มํสํ คเหตฺวา อุปสฺสยํ อคมาสิ. อถ โข อุปฺปลวณฺณา ภิกฺขุนี ตสฺสา รตฺติยา อจฺจเยน ตํ มํสํ สมฺปาเทตฺวา อุตฺตราสงฺเคน ภณฺฑิกํ พนฺธิตฺวา \textbf{เวหาสํ อพฺภุคฺคนฺตฺวา เวฬุวเน ปจฺจุฏฺฐาสิ}\csromansharedfootnote{035dc7a5419cb609}{ปจฺจุปฏฺฐาสิ (?)}\textbf{.}}

\csromanmatikaanchor{mtk.135}
\csromantocmarkref{subsubsection}{สิกฺขาปทวิภงฺค, ปทภาชนีย}{mtk.135}
\bookmark[dest={mtk.135},level=3]{สิกฺขาปทวิภงฺค, ปทภาชนีย}
\prose{เตน โข ปน สมเยน ภควา คามํ ปิณฺฑาย ปวิฏฺโฐ โหติ. อายสฺมา อุทายี โอหิยฺยโก โหติ วิหารปาโล. อถ โข อุปฺปลวณฺณา ภิกฺขุนี เยนายสฺมา อุทายี เตนุปสงฺกมิ, อุปสงฺกมิตฺวา อายสฺมนฺตํ อุทายิํ เอตทโวจ “กหํ ภนฺเต ภควา”ติ. ปวิฏฺโฐ ภคินิ ภควา คามํ ปิณฺฑายาติ. อิมํ ภนฺเต มํสํ ภควโต เทหีติ. สนฺตปฺปิโต ตยา ภคินิ ภควา มํเสน, สเจ เม ตฺวํ อนฺตรวาสกํ ทเทยฺยาสิ, เอวํ อหมฺปิ สนฺตปฺปิโต ภเวยฺยํ อนฺตรวาสเกนาติ. มยํ โข ภนฺเต มาตุคามา นาม กิจฺฉลาภา, อิทญฺจ เม อนฺติมํ ปญฺจมํ จีวรํ, นาหํ ทสฺสามีติ. เสยฺยถาปิ ภคินิ ปุริโส หตฺถิํ ทตฺวา กจฺเฉ สชฺเชยฺย\csromansharedfootnote{b5a4f98cc7aae59d}{วิสฺสชฺเชยฺย (สฺยา)}, เอวเมว โข ตฺวํ ภคินิ ภควโต มํสํ ทตฺวา มยิ อนฺตรวาสเก สชฺชสีติ\csromansharedfootnote{9a46fe49a50e519b}{อนฺตรวาสกํ น สชฺชสีติ (ก), มยฺหํ อนฺตรวาสกํ วิสฺสชฺเชหีติ (สฺยา)}. อถ โข อุปฺปลวณฺณา ภิกฺขุนี อายสฺมตา อุทายินา นิปฺปีฬิยมานา อนฺตรวาสกํ ทตฺวา อุปสฺสยํ อคมาสิ. ภิกฺขุนิโย อุปฺปลวณฺณาย ภิกฺขุนิยา ปตฺตจีวรํ ปฏิคฺคณฺหนฺติโย อุปฺปลวณฺณํ ภิกฺขุนิํ เอตทโวจุํ “กหํ เต อยฺเย อนฺตรวาสโก”ติ. อุปฺปลวณฺณา ภิกฺขุนี ภิกฺขุนีนํ เอตมตฺถํ อาโรเจสิ. ภิกฺขุนิโย อุชฺฌายนฺติ ขิยฺยนฺติ วิปาเจนฺติ “กถํ หิ นาม อยฺโย อุทายี ภิกฺขุนิยา จีวรํ ปฏิคฺคเหสฺสติ, กิจฺฉลาโภ มาตุคาโม”ติ. อถ โข ตา ภิกฺขุนิโย ภิกฺขูนํ เอตมตฺถํ อาโรเจสุํ. เย เต ภิกฺขู อปฺปิจฺฉา ฯเปฯ เต อุชฺฌายนฺติ ขิยฺยนฺติ วิปาเจนฺติ “กถํ หิ นาม อายสฺมา อุทายี ภิกฺขุนิยา จีวรํ ปฏิคฺคเหสฺสตี”ติ. \csromanfolio{310}อถ โข เต ภิกฺขู อายสฺมนฺตํ อุทายิํ อเนก\-ปริยาเยน วิครหิตฺวา ภควโต เอตมตฺถํ อาโรเจสุํ ฯเปฯ “สจฺจํ กิร ตฺวํ อุทายิ ภิกฺขุนิยา จีวรํ ปฏิคฺคเหสี”ติ. สจฺจํ ภควาติ. ญาติกา เต อุทายิ อญฺญาติกาติ. อญฺญาติกา ภควาติ. อญฺญาตโก โมฆปุริส อญฺญาติกาย น ชานาติ ปติรูปํ วา อปฺปติรูปํ วา สนฺตํ วา อสนฺตํ วา, ตตฺถ นาม ตฺวํ โมฆปุริส อญฺญาติกาย ภิกฺขุนิยา หตฺถโต จีวรํ ปฏิคฺคเหสฺสสิ, เนตํ โมฆปุริส อปฺปสนฺนานํ วา ปสาทาย ฯเปฯ. เอวญฺจ ปน ภิกฺขเว อิมํ สิกฺขาปทํ อุทฺทิเสยฺยาถ\csromandash{}}

\proseitem{509}{\textbf{“โย ปน ภิกฺขุ อญฺญาติกาย ภิกฺขุนิยา หตฺถโต จีวรํ ปฏิคฺคเณฺหยฺย, นิสฺสคฺคิยํ ปาจิตฺติยนฺ”}ติ.}

\prose{เอวญฺจิทํ ภควตา ภิกฺขูนํ สิกฺขาปทํ ปญฺญตฺตํ โหติ.}

\proseitem{510}{เตน โข ปน สมเยน ภิกฺขู กุกฺกุจฺจายนฺตา ภิกฺขุนีนํ ปาริวตฺตก\-จีวรํ น ปฏิคฺคณฺหนฺติ. ภิกฺขุนิโย อุชฺฌยนฺติ ขิยฺยนฺติ วิปาเจนฺติ “กถํ หิ นาม อยฺยา อมฺหากํ ปาริวตฺตก\-จีวรํ น ปฏิคฺคเหสฺสนฺตี”ติ. อสฺโสสุํ โข ภิกฺขู ตาสํ ภิกฺขุนีนํ อุชฺฌายนฺตีนํ ขิยฺยนฺตีนํ วิปาเจนฺตีนํ. อถ โข เต ภิกฺขู ภควโต เอตมตฺถํ อาโรเจสุํ. อถ โข ภควา เอตสฺมิํ นิทาเน เอตสฺมิํ ปกรเณ ธมฺมิํ กถํ กตฺวา ภิกฺขู อามนฺเตสิ “อนุชานามิ ภิกฺขเว ปญฺจนฺนํ ปาริวตฺตกํ ปฏิคฺคเหตุํ ภิกฺขุสฺส ภิกฺขุนิยา สิกฺขมานาย สามเณรสฺส สามเณริยา, อนุชานามิ ภิกฺขเว อิเมสํ ปญฺจนฺนํ ปาริวตฺตกํ ปฏิคฺคเหตุํ. เอวญฺจ ปน ภิกฺขเว อิมํ สิกฺขาปทํ อุทฺทิเสยฺยาถ\csromandash{}}

\proseitem{511}{“\textbf{โย ปน ภิกฺขุ อญฺญาติกาย ภิกฺขุนิยา หตฺถโต จีวรํ} ปฏิคฺคเณฺหยฺย อญฺญตฺร ปาริวตฺตกา, นิสฺสคฺคิยํ ปาจิตฺติยนฺ”ติ.}

\proseitem{512}{\textbf{โย ปนา}ติ โย ยาทิโส ฯเปฯ. \textbf{ภิกฺขู}ติ ฯเปฯ อยํ อิมสฺมิํ อตฺเถ อธิปฺเปโต “ภิกฺขู”ติ.}

\prose{\textbf{อญฺญาติกา} นาม มาติโต วา ปิติโต วา ยาว สตฺตมา ปิตามหยุคา อสมฺพทฺธา.}

\prose{\textbf{ภิกฺขุนี} นาม อุภโตสํเฆ อุปสมฺปนฺนา.}

\csromanmatikaanchor{mtk.136}
\csromanmatikaanchor{mtk.137}
\csromantocmarknopage{subsection}{6. อญฺญาตกวิญฺญตฺติสิกฺขาปท}{mtk.136}
\bookmark[dest={mtk.136},level=2]{6. อญฺญาตกวิญฺญตฺติสิกฺขาปท}
\csromantocmarkref{subsubsection}{อุปนนฺทภิกฺขุวตฺถุ}{mtk.137}
\bookmark[dest={mtk.137},level=3]{อุปนนฺทภิกฺขุวตฺถุ}
\prose{\csromanfolio{311}\textbf{จีวรํ} นาม ฉนฺนํ จีวรานํ อญฺญตรํ จีวรํ วิกปฺปนุปคํ ปจฺฉิมํ.}

\prose{\textbf{อญฺญตฺร ปาริวตฺตกา}ติ ฐเปตฺวา ปาริวตฺตกํ.}

\prose{ปฏิคฺคณฺหาติ ปโยเค ทุกฺกฏํ. ปฏิลาเภน นิสฺสคฺคิยํ โหติ. นิสฺสชฺชิตพฺพํ สํฆสฺส วา คณสฺส วา ปุคฺคลสฺส วา, เอวญฺจ ปน ภิกฺขเว นิสฺสชฺชิตพฺพํ ฯเปฯ “อิทํ เม ภนฺเต จีวรํ อญฺญาติกาย ภิกฺขุนิยา หตฺถโต ปฏิคฺคหิตํ อญฺญตฺร ปาริวตฺตกา นิสฺสคฺคิยํ, อิมาหํ สํฆสฺส นิสฺสชฺชามี”ติ ฯเปฯ ทเทยฺยาติ ฯเปฯ ทเทยฺยุนฺติ ฯเปฯ อายสฺมโต ทมฺมีติ.}

\proseitem{513}{อญฺญาติกาย อญฺญาติกสญฺญี จีวรํ ปฏิคฺคณฺหาติ อญฺญตฺร ปาริวตฺตกา, นิสฺสคฺคิยํ ปาจิตฺติยํ. อญฺญาติกาย เวมติโก จีวรํ ปฏิคฺคณฺหาติ อญฺญตฺร ปาริวตฺตกา, นิสฺสคฺคิยํ ปาจิตฺติยํ. อญฺญาติกาย ญาติกสญฺญี จีวรํ ปฏิคฺคณฺหาติ อญฺญตฺร ปาริวตฺตกา, นิสฺสคฺคิยํ ปาจิตฺติยํ.}

\prose{เอกโต\-อุปสมฺปนฺนาย หตฺถโต จีวรํ ปฏิคฺคณฺหาติ อญฺญตฺร ปาริวตฺตกา, อาปตฺติ ทุกฺกฏสฺส. ญาติกาย อญฺญาติกสญฺญี, อาปตฺติ ทุกฺกฏสฺส. ญาติกาย เวมติโก, อาปตฺติ ทุกฺกฏสฺส. ญาติกาย ญาติกสญฺญี, อนาปตฺติ.}

\proseitemwithcloserruled{514}{อนาปตฺติ ญาติกาย, ปาริวตฺตกํ ปริตฺเตน วา วิปุลํ, วิปุเลน วา ปริตฺตํ, ภิกฺขุ วิสฺสาสํ คณฺหาติ, ตาวกาลิกํ คณฺหาติ, จีวรํ ฐเปตฺวา อญฺญํ ปริกฺขารํ คณฺหาติ, สิกฺขมานาย สามเณริยา อุมฺมตฺตกสฺส อาทิกมฺมิกสฺสาติ.}{จีวร\-ปฏิคฺคหณ\-สิกฺขาปทํ นิฏฺฐิตํ ปญฺจมํ.}
\csromanmarkh{1. จีวรวคฺค}\csromanmarkhii{6. อญฺญาตก\-วิญฺญตฺติ\-สิกฺขาปท}\csromanheadernomark{2}{1. \textbf{จีวรวคฺค 6. อญฺญาตก}\-\textbf{วิญฺญตฺติ}\-\textbf{สิกฺขาปท}}
\proseitem{515}{เตน สมเยน พุทฺโธ ภควา สาวตฺถิยํ วิหรติ เชตวเน อนาถ\-ปิณฺฑิกสฺส อาราเม. เตน โข ปน สมเยน อายสฺมา อุปนนฺโท สกฺยปุตฺโต ปฏฺโฏ\csromansharedfootnote{b387556fb62cee3c}{ปฏฺโฐ (สฺยา, ก)} โหติ ธมฺมิํ กถํ กาตุํ. อถ โข อญฺญตโร เสฏฺฐิปุตฺโต เยนายสฺมา อุปนนฺโท สกฺยปุตฺโต \csromanfolio{312}เตนุปสงฺกมิ, อุปสงฺกมิตฺวา อายสฺมนฺตํ อุปนนฺทํ สกฺยปุตฺตํ อภิวาเทตฺวา เอกมนฺตํ นิสีทิ, เอกมนฺตํ นิสินฺนํ โข ตํ เสฏฺฐิปุตฺตํ อายสฺมา อุปนนฺโท สกฺยปุตฺโต ธมฺมิยา กถาย สนฺทสฺเสสิ สมาทเปสิ สมุตฺเตเชสิ สมฺปหํเสสิ. อถ โข โส เสฏฺฐิปุตฺโต อายสฺมตา อุปนนฺเทน สกฺยปุตฺเตน ธมฺมิยา กถาย สนฺทสฺสิโต สมาทปิโต สมุตฺเตชิโต สมฺปหํสิโต อายสฺมนฺตํ อุปนนฺทํ สกฺยปุตฺตํ เอตทโวจ “วเทยฺยาถ ภนฺเต เยน อตฺโถ ปฏิพลา มยํ อยฺยสฺส ทาตุํ, ยทิทํ จีวร\-ปิณฺฑปาต\-เสนาสน\-คิลานปฺปจฺจย\-เภสชฺช\-ปริกฺขารนฺ”ติ. สเจ เม ตฺวํ อาวุโส ทาตุกาโมสิ, อิโต เอกํ สาฏกํ เทหีติ. อมฺหากํ โข ภนฺเต กุลปุตฺตานํ กิสฺมิํ วิย เอกสาฏกํ คนฺตุํ, อาคเมหิ ภนฺเต ยาว ฆรํ คจฺฉามิ, ฆรํ คโต อิโต วา เอกํ สาฏกํ ปหิณิสฺสามิ อิโต วา \textbf{สุนฺทรตรนฺติ. ทุติยมฺปิ โข อายสฺมา อุปนนฺโท สกฺยปุตฺโต ตํ} \textbf{เสฏฺฐิปุตฺตํ เอตทโวจ “สเจ เม ตฺวํ อาวุโส ทาตุกาโมสิ, อิโต เอกํ สาฏกํ} \textbf{เทหี”ติ. อมฺหากํ โข ภนฺเต กุลปุตฺตานํ กิสฺมิํ วิย เอกสาฏกํ} คนฺตุํ, อาคเมหิ ภนฺเต ยาว ฆรํ คจฺฉามิ, ฆรํ คโต อิโต วา เอกํ สาฏกํ ปหิณิสฺสามิ อิโต วา สุนฺทรตรนฺติ. ตติยมฺปิ โข อายสฺมา อุปนนฺโท สกฺยปุตฺโต ตํ เสฏฺฐิปุตฺตํ เอตทโวจ “สเจ เม ตฺวํ อาวุโส ทาตุกาโมสิ, อิโต เอกํ สาฏกํ เทหี”ติ. อมฺหากํ โข ภนฺเต กุลปุตฺตานํ กิสฺมิํ วิย เอกสาฏกํ คนฺตุํ, อาคเมหิ ภนฺเต ยาว ฆรํ คจฺฉามิ, ฆรํ คโต อิโต วา เอกํ สาฏกํ ปหิณิสฺสามิ อิโต วา สุนฺทรตรนฺติ. กิํ ปน ตยา อาวุโส อทาตุกาเมน ปวาริเตน ยํ ตฺวํ ปวาเรตฺวา น เทสีติ.}

\csromanmatikaanchor{mtk.138}
\csromantocmarkref{subsubsection}{ปฐมปญฺญตฺติ}{mtk.138}
\bookmark[dest={mtk.138},level=3]{ปฐมปญฺญตฺติ}
\prose{อถ โข โส เสฏฺฐิปุตฺโต อายสฺมตา อุปนนฺเทน สกฺยปุตฺเตน นิปฺปีฬิยมาโน เอกํ สาฏกํ ทตฺวา อคมาสิ. มนุสฺสา ตํ เสฏฺฐิปุตฺตํ เอตทโวจุํ “กิสฺส ตฺวํ อยฺโย เอกสาฏโก อาคจฺฉสี”ติ. อถ โข โส เสฏฺฐิปุตฺโต เตสํ มนุสฺสานํ เอตมตฺถํ อาโรเจสิ. มนุสฺสา อุชฺฌายนฺติ ขิยฺยนฺติ วิปาเจนฺติ “มหิจฺฉา อิเม สมณา สกฺยปุตฺติยา อสนฺตุฏฺฐา, นยิเมสํ สุกรา ธมฺมนิมนฺตนาปิ กาตุํ\csromansharedfootnote{59e57c05da804e33}{นยิเม สุกรา ธมฺมนิมนฺตนายปิ กาตุํ (สฺยา)}, กถํ หิ นาม เสฏฺฐิปุตฺเตน ธมฺม\-นิมนฺตนาย กยิรมานาย สาฏกํ คเหสฺสนฺตี”ติ. \csromanfolio{313}อสฺโสสุํ โข ภิกฺขู เตสํ มนุสฺสานํ อุชฺฌยนฺตานํ ขิยฺยนฺตานํ วิปาเจนฺตานํ. เย เต ภิกฺขู อปฺปิจฺฉา ฯเปฯ เต อุชฺฌายนฺติ ขิยฺยนฺติ วิปาเจนฺติ “กถํ หิ นาม อายสฺมา อุปนนฺโท สกฺยปุตฺโต เสฏฺฐิปุตฺตํ จีวรํ วิญฺญาเปสฺสตี”ติ. อถ โข เต ภิกฺขู อายสฺมนฺตํ อุปนนฺทํ สกฺยปุตฺตํ อเนก\-ปริยาเยน วิครหิตฺวา ภควโต เอตมตฺถํ อาโรเจสุํ ฯเปฯ “สจฺจํ กิร ตฺวํ อุปนนฺท เสฏฺฐิปุตฺตํ จีวรํ วิญฺญาเปสี”ติ. สจฺจํ ภควาติ. ญาตโก เต อุปนนฺท อญฺญาตโกติ. อญฺญาตโก ภควาติ. อญฺญาตโก โมฆปุริส อญฺญาตกสฺส น ชานาติ ปติรูปํ วา อปฺปติรูปํ วา สนฺตํ วา อสนฺตํ วา, ตตฺถ นาม ตฺวํ โมฆปุริส อญฺญาตกํ เสฏฺฐิปุตฺตํ จีวรํ \textbf{วิญฺญาเปสฺสสิ, เนตํ โมฆปุริส อปฺปสนฺนานํ วา ปสาทาย ฯเปฯ. เอวญฺจ} \textbf{ปน ภิกฺขเว อิมํ สิกฺขาปทํ อุทฺทิเสยฺยาถ\csromandash{}}}

\proseitem{516}{\textbf{“โย ปน ภิกฺขุ อญฺญาตกํ คหปติํ วา คหปตานิํ วา จีวรํ วิญฺญาเปยฺย, นิสฺสคฺคิยํ ปาจิตฺติยนฺ”}ติ.}

\prose{เอวญฺจิทํ ภควตา ภิกฺขูนํ สิกฺขาปทํ ปญฺญตฺตํ โหติ.}

\csromanmatikaanchor{mtk.139}
\csromanmatikaanchor{mtk.140}
\csromantocmarkref{subsubsection}{อนุปญฺญตฺติ}{mtk.139}
\bookmark[dest={mtk.139},level=3]{อนุปญฺญตฺติ}
\csromantocmarkref{subsubsection}{สิกฺขาปทวิภงฺค, ปทภาชนีย}{mtk.140}
\bookmark[dest={mtk.140},level=3]{สิกฺขาปทวิภงฺค, ปทภาชนีย}
\proseitem{517}{เตน โข ปน สมเยน สมฺพหุลา ภิกฺขู สาเกตา สาวตฺถิํ อทฺธาน\-มคฺค\-ปฺปฏิปนฺนา โหนฺติ. อนฺตรามคฺเค โจรา นิกฺขมิตฺวา เต ภิกฺขู อจฺฉินฺทิํสุ. อถ โข เต ภิกฺขู ภควตา ปฏิกฺขิตฺตํ “อญฺญาตกํ คหปติํ วา คหปตานิํ วา จีวรํ วิญฺญาเปตุนฺ”ติ กุกฺกุจฺจายนฺตา น วิญฺญาเปสุํ, ยถานคฺคาว สาวตฺถิํ คนฺตฺวา ภิกฺขู อภิวาเทนฺติ. ภิกฺขู เอวมาหํสุ “สุนฺทรา โข อิเม อาวุโส อาชีวกา, เย อิเม ภิกฺขูสุ อภิวาเทนฺตี”ติ. เต เอวมาหํสุ “น มยํ อาวุโส อาชีวกา, ภิกฺขู มยนฺ”ติ. ภิกฺขู อายสฺมนฺตํ อุปาลิํ เอตทโวจุํ “อิงฺฆาวุโส อุปาลิ อิเม อนุยุญฺชาหี”ติ. อถ โข อายสฺมตา อุปาลินา อนุยุญฺชิยมานา เต ภิกฺขู เอตมตฺถํ อาโรเจสุํ. อถ โข อายสฺมา อุปาลิ เต ภิกฺขู อนุยุญฺชิตฺวา ภิกฺขู เอตทโวจ “ภิกฺขู อิเม อาวุโส เทถ เนสํ จีวรานี”ติ. เย เต ภิกฺขู อปฺปิจฺฉา ฯเปฯ เต อุชฺฌายนฺติ ขิยฺยนฺติ วิปาเจนฺติ “กถํ หิ นาม ภิกฺขู นคฺคา อาคจฺฉิสฺสนฺติ, นนุ นาม ติเณน วา ปณฺเณน วา ปฏิจฺฉาเทตฺวา อาคนฺตพฺพนฺ”ติ. อถ โข เต ภิกฺขู เต อเนก\-ปริยาเยน วิครหิตฺวา ภควโต เอตมตฺถํ อาโรเจสุํ. อถ โข ภควา เอตสฺมิํ นิทาเน เอตสฺมิํ ปกรเณ ธมฺมิํ กถํ กตฺวา \csromanfolio{314}ภิกฺขู อามนฺเตสิ “อนุชานามิ ภิกฺขเว อจฺฉินฺน\-จีวรสฺส วา นฏฺฐ\-จีวรสฺส วา อญฺญาตกํ คหปติํ วา คหปตานิํ วา จีวรํ วิญฺญาเปตุํ, ยํ อาวาสํ ปฐมํ อุปคจฺฉติ, สเจ ตตฺถ โหติ สํฆสฺส วิหารจีวรํ วา อุตฺตร\-ตฺถรณํ วา ภูมตฺถรณํ วา ภิสิจฺฉวิ วา, ตํ คเหตฺวา ปารุปิตุํ ‘ลภิตฺวา โอทหิสฺสามี’ติ, โน เจ โหติ สํฆสฺส วิหารจีวรํ วา อุตฺตร\-ตฺถรณํ วา ภูมตฺถรณํ วา ภิสิจฺฉวิ วา, ติเณน วา ปณฺเณน วา ปฏิจฺฉาเทตฺวา อาคนฺตพฺพํ, น เตฺวว นคฺเคน อาคนฺตพฺพํ, โย อาคจฺเฉยฺย, อาปตฺติ ทุกฺกฏสฺส. เอวญฺจ ปน ภิกฺขเว อิมํ สิกฺขาปทํ อุทฺทิเสยฺยาถ\csromandash{}}

\proseitem{518}{\textbf{“โย ปน ภิกฺขุ อญฺญาตกํ คหปติํ วา คหปตานิํ วา จีวรํ วิญฺญาเปยฺย อญฺญตฺร สมยา, นิสฺสคฺคิยํ ปาจิตฺติยํ. ตตฺถายํ สมโย, อจฺฉินฺนจีวโร วา โหติ ภิกฺขุ นฏฺฐจีวโร วา, อยํ ตตฺถ สมโย”}ติ.}

\proseitem{519}{\textbf{โย ปนา}ติ โย ยาทิโส ฯเปฯ. \textbf{ภิกฺขู}ติ ฯเปฯ อยํ อิมสฺมิํ อตฺเถ อธิปฺเปโต “ภิกฺขู”ติ.}

\prose{\textbf{อญฺญาตโก} นาม มาติโต วา ปิติโต วา ยาว สตฺตมา ปิตามหยุคา อสมฺพทฺโธ.}

\prose{\textbf{คหปติ} นาม โย โกจิ อคารํ อชฺฌาวสติ.}

\prose{\textbf{คหปตานี} นาม ยากาจิ อคารํ อชฺฌาวสติ.}

\prose{\textbf{จีวรํ} นาม ฉนฺนํ จีวรานํ อญฺญตรํ จีวรํ วิกปฺปนุปคํ ปจฺฉิมํ.}

\prose{\textbf{อญฺญตฺร สมยา}ติ ฐเปตฺวา สมยํ.}

\prose{\textbf{อจฺฉินฺนจีวโร} นาม ภิกฺขุสฺส จีวรํ อจฺฉินฺนํ โหติ ราชูหิ วา โจเรหิ วา ธุตฺเตหิ วา เยหิ เกหิจิ วา อจฺฉินฺนํ โหติ.}

\prose{\textbf{นฏฺฐจีวโร} นาม ภิกฺขุสฺส จิวรํ อคฺคินา วา ทฑฺฒํ โหติ, อุทเกน วา วูฬฺหํ โหติ, อุนฺทูเรหิ วา อุปจิกาหิ วา ขายิตํ โหติ, ปริโภค\-ชิณฺณํ วา โหติ.}

\csromanmatikaanchor{mtk.141}
\csromanmatikaanchor{mtk.142}
\csromantocmarknopage{subsection}{7. ตตุตฺตริสิกฺขาปท}{mtk.141}
\bookmark[dest={mtk.141},level=2]{7. ตตุตฺตริสิกฺขาปท}
\csromantocmarkref{subsubsection}{ฉพฺพคฺคิยภิกฺขุวตฺถุ}{mtk.142}
\bookmark[dest={mtk.142},level=3]{ฉพฺพคฺคิยภิกฺขุวตฺถุ}
\prose{อญฺญตฺร สมยา วิญฺญาเปติ, ปโยเค ทุกฺกฏํ. ปฏิลาเภน นิสฺสคฺคิยํ โหติ. นิสฺสชฺชิตพฺพํ สํฆสฺส วา คณสฺส วา ปุคฺคลสฺส วา. เอวญฺจ ปน \csromanfolio{315}ภิกฺขเว นิสฺสชฺชิตพฺพํ ฯเปฯ “อิทํ เม ภนฺเต จีวรํ อญฺญาตกํ คหปติกํ อญฺญตฺร สมยา วิญฺญาปิตํ นิสฺสคฺคิยํ, อิมาหํ สํฆสฺส นิสฺสชฺชามี”ติ ฯเปฯ ทเทยฺยาติ ฯเปฯ ทเทยฺยุนฺติ ฯเปฯ อายสฺมโต ทมฺมีติ.}

\proseitem{520}{อญฺญตเก อญฺญาตกสญฺญี อญฺญตฺร สมยา จีวรํ วิญฺญาเปติ, นิสฺสคฺคิยํ ปาจิตฺติยํ. อญฺญาตเก เวมติโก อญฺญตฺร สมยา จีวรํ วิญฺญาเปติ, นิสฺสคฺคิยํ ปาจิตฺติยํ. อญฺญาตเก ญาตกสญฺญี อญฺญตฺร สมยา จีวรํ วิญฺญาเปติ, นิสฺสคฺคิยํ ปาจิตฺติยํ.}

\prose{ญาตเก อญฺญตกสญฺญี, อาปตฺติ ทุกฺกฏสฺส. ญาตเก เวมติโก, อาปตฺติ ทุกฺกฏสฺส. ญาตเก ญาตกสญฺญี, อนาปตฺติ.}

\proseitemwithcloserruled{521}{อนาปตฺติ สมเย ญาติกานํ ปวาริตานํ อญฺญสฺสตฺถาย อตฺตโน ธเนน อุมฺมตฺตกสฺส อาทิกมฺมิกสฺสาติ.}{อญฺญาตก\-วิญฺญตฺติ\-สิกฺขาปทํ นิฏฺฐิตํ ฉฏฺฐํ.}
\csromanmarkh{1. จีวรวคฺค}\csromanmarkhii{7. ตตุตฺตริ\-สิกฺขาปท}\csromanheadernomark{2}{1. \textbf{จีวรวคฺค 7. ตตุตฺตริ}\-\textbf{สิกฺขาปท}}
\proseitem{522}{เตน สมเยน พุทฺโธ ภควา สาวตฺถิยํ วิหรติ เชตวเน อนาถ\-ปิณฺฑิกสฺส อาราเม. เตน โข ปน สมเยน ฉพฺพคฺคิยา ภิกฺขู อจฺฉินฺน\-จีวรเก ภิกฺขู อุปสงฺกมิตฺวา เอวํ วทนฺติ “ภควตา อาวุโส อนุญฺญาตํ อจฺฉินฺน\-จีวรสฺส วา นฏฺฐ\-จีวรสฺส วา อญฺญาตกํ คหปติํ วา คหปตานิํ วา จีวรํ วิญฺญาเปตุํ, วิญฺญาเปถ อาวุโส จีวรนฺ”ติ. อลํ อาวุโส ลทฺธํ อเมฺหหิ จีวรนฺติ. มยํ อายสฺมนฺตานํ วิญฺญาเปมาติ. วิญฺญาเปถ อาวุโสติ. อถ โข ฉพฺพคฺคิยา ภิกฺขู คหปติเก อุปสงฺกมิตฺวา เอตทโวจุํ “อจฺฉินฺน\-จีวรกา อาวุโส ภิกฺขู อาคตา, เทถ เนสํ จีวรานี”ติ พหุํ จีวรํ วิญฺญาเปสุํ.}

\csromanmatikaanchor{mtk.143}
\csromanmatikaanchor{mtk.144}
\csromantocmarkref{subsubsection}{ปญฺญตฺติ}{mtk.143}
\bookmark[dest={mtk.143},level=3]{ปญฺญตฺติ}
\csromantocmarkref{subsubsection}{สิกฺขาปทวิภงฺค, ปทภาชนีย}{mtk.144}
\bookmark[dest={mtk.144},level=3]{สิกฺขาปทวิภงฺค, ปทภาชนีย}
\prose{เตน โข ปน สมเยน อญฺญตโร ปุริโส สภายํ นิสินฺโน อญฺญตรํ ปุริสํ เอตทโวจ “อจฺฉินฺน\-จีวรกา อยฺโย ภิกฺขู อาคตา, เตสํ มยา จีวรํ ทินฺนนฺ”ติ. โสปิ เอวมาห “มยาปิ ทินฺนนฺ”ติ. อปโรปิ เอวมาห “มยาปิ ทินฺนนฺ”ติ. เต อุชฺฌายนฺติ ขิยฺยนฺติ วิปาเจนฺติ “กถํ หิ นาม สมณา สกฺยปุตฺติยา น มตฺตํ ชานิตฺวา พหุํ จีวรํ วิญฺญาเปสฺสนฺติ, \csromanfolio{316}ทุสฺสวาณิชฺชํ วา สมณา สกฺยปุตฺติยา กริสฺสนฺติ, ปคฺคาหิกสาลํ\csromansharedfootnote{15f3e34bb08ff691}{ปฏคฺคาหิกสาลํ (?)} วา ปสาเรสฺสนฺตี”ติ. อสฺโสสุํ โข ภิกฺขู เตสํ มนุสฺสานํ อุชฺฌายนฺตานํ ขิยฺยนฺตานํ วิปาเจนฺตานํ. เย เต ภิกฺขู อปฺปิจฺฉา ฯเปฯ เต อุชฺฌายนฺติ ขิยฺยนฺติ วิปาเจนฺติ “กถํ หิ นาม ฉพฺพคฺคิยา ภิกฺขู น มตฺถํ ชานิตฺวา พหุํ จีวรํ วิญฺญาเปสฺสนฺตี”ติ. อถ โข เต ภิกฺขู ฉพฺพคฺคิเย ภิกฺขู อเนก\-ปริยาเยน วิครหิตฺวา ภควโต เอตมตฺถํ อาโรเจสุํ ฯเปฯ “สจฺจํ กิร ตุเมฺห ภิกฺขเว น มตฺตํ ชานิตฺวา พหุํ จีวรํ วิญฺญาเปถา”ติ. สจฺจํ ภควาติ. วิครหิ พุทฺโธ ภควา ฯเปฯ. กถํ หิ นาม ตุเมฺห โมฆปุริสา น มตฺตํ ชานิตฺวา พหุํ จีวรํ วิญฺญาเปสฺสถ, เนตํ โมฆปุริสา อปฺปสนฺนานํ วา ปสาทาย ฯเปฯ. เอวญฺจ ปน ภิกฺขเว อิมํ สิกฺขาปทํ อุทฺทิเสยฺยาถ\csromandash{}}

\proseitem{523}{\textbf{“ตญฺเจ อญฺญาตโก คหปติ วา คหปตานี วา พหูหิ จีวเรหิ อภิหฏฺฐุํ ปเวเรยฺย, สนฺตรุ}\-\textbf{ตฺตร}\-\textbf{ปรมํ เตน ภิกฺขุนา ตโต จีวรํ สาทิตพฺพํ, ตโต เจ อุตฺตริ สาทิเยยฺย, นิสฺสคฺคิยํ ปาจิตฺติยนฺ”}ติ.}

\proseitem{524}{\textbf{ตญฺเจ}ติ อจฺฉินฺน\-จีวรกํ ภิกฺขุํ.}

\prose{\textbf{อญฺญาตโก} นาม มาติโต วา ปิติโต วา ยาว สตฺตมา ปิตามหยุคา อสมฺพทฺโธ.}

\prose{\textbf{คหปติ} นาม โย โกจิ อคารํ อชฺฌาวสติ.}

\prose{\textbf{คหปตานี} นาม ยา กาจิ อคารํ อชฺฌาวสติ.}

\prose{\textbf{พหูหิ จีวเรหี}ติ พหุเกหิ จีวเรหิ.}

\prose{\textbf{อภิหฏฺฐุํ ปวาเรยฺยา}ติ ยาวตกํ อิจฺฉสิ, ตาวตกํ คณฺหาหีติ.}

\prose{\textbf{สนฺตรุ}\-\textbf{ตฺตร}\-\textbf{ปรมํ เตน ภิกฺขุนา ตโต จีวรํ สาทิตพฺพนฺ}ติ สเจ ตีณิ นฏฺฐานิ โหนฺติ, เทฺว สาทิตพฺพานิ. เทฺว นฏฺฐานิ, เอกํ สาทิตพฺพํ. เอกํ นฏฺฐํ, น กิญฺจิ สาทิตพฺพํ.}

\csromanmatikaanchor{mtk.145}
\csromanmatikaanchor{mtk.146}
\csromantocmarknopage{subsection}{8. อุปกฺขฏสิกฺขาปท}{mtk.145}
\bookmark[dest={mtk.145},level=2]{8. อุปกฺขฏสิกฺขาปท}
\csromantocmarkref{subsubsection}{อุปนนฺทภิกฺขุวตฺถุ}{mtk.146}
\bookmark[dest={mtk.146},level=3]{อุปนนฺทภิกฺขุวตฺถุ}
\prose{\textbf{ตโต เจ อุตฺตริ สาทิเยยฺยา}ติ ตตุตฺตริ วิญฺญาเปติ, ปโยเค ทุกฺกฏํ. ปฏิลาเภน นิสฺสคฺคิยํ โหติ. นิสฺสชฺชิตพฺพํ สํฆสฺส วา คณสฺส วา ปุคฺคลสฺส วา. เอวญฺจ ปน ภิกฺขเว นิสฺสชฺชิตพฺพํ ฯเปฯ “อิทํ เม ภนฺเต จีวรํ \csromanfolio{317}อญฺญาตกํ คหปติกํ อุปสงฺกมิตฺวา ตตุตฺตริ วิญฺญาปิตํ นิสฺสคฺคิยํ, อิมาหํ สํฆสฺส นิสฺสชฺชามี”ติ ฯเปฯ ทเทยฺยาติ ฯเปฯ ทเทยฺยุนฺติ ฯเปฯ อายสฺมโต ทมฺมีติ.}

\proseitem{525}{อญฺญาตเก อญฺญาตกสญฺญี ตตุตฺตริ จีวรํ วิญฺญาเปติ, นิสฺสคฺคิยํ ปาจิตฺติยํ. อญฺญาตเก เวมติโก ตตุตฺตริ จีวรํ วิญฺญาเปติ, นิสฺสคฺคิยํ ปาจิตฺติยํ. อญฺญตเก ญาตกสญฺญี ตตุตฺตริ จีวรํ วิญฺญาเปติ, นิสฺสคฺคิยํ ปาจิตฺติยํ.}

\prose{ญาตเก อญฺญาตกสญฺญี, อาปตฺติ ทุกฺกฏสฺส. ญาตเก เวมติโก, อาปตฺติ ทุกฺกฏสฺส. ญาตเก ญาตกสญฺญี, อนาปตฺติ.}

\proseitemwithcloserruled{526}{อนาปตฺติ เสสกํ อาหริสฺสามีติ หรนฺโต คจฺฉติ, เสสกํ ตุเยฺหว โหตูติ เทนฺติ, น อจฺฉินฺนการณา เทนฺติ, น นฏฺฐาการณา เทนฺติ, ญาตกานํ ปวาริตานํ อตฺตโน ธเนน อุมฺมตฺตกสฺส อาทิกมฺมิกสฺสาติ.}{ตตุตฺตริ\-สิกฺขาปทํ นิฏฺฐิตํ สตฺตมํ.}
\csromanmarkh{1. จีวรวคฺค}\csromanmarkhii{8. อุปกฺขฏ\-สิกฺขาปท}\csromanheadernomark{2}{1. \textbf{จีวรวคฺค 8. อุปกฺขฏ}\-\textbf{สิกฺขาปท}}
\proseitem{527}{เตน สมเยน พุทฺโธ ภควา สาวตฺถิยํ วิหรติ เชตวเน อนาถ\-ปิณฺฑิกสฺส อาราเม. เตน โข ปน สมเยน อญฺญตโร ปุริโส ปชาปติํ เอตทโวจ “อยฺยํ อุปนนฺทํ จีวเรน อจฺฉาเทสฺสามี”ติ. อสฺโสสิ โข อญฺญตโร ปิณฺฑจาริโก ภิกฺขุ ตสฺส ปุริสสฺส อิมํ วาจํ ภาสมานสฺส. อถ โข โส ภิกฺขุ เยนายสฺมา อุปนนฺโท สกฺยปุตฺโต เตนุปสงฺกมิ, อุปสงฺกมิตฺวา อายสฺมนฺตํ อุปนนฺทํ สกฺยปุตฺตํ เอตทโวจ “มหาปุญฺโญสิ ตฺวํ อาวุโส อุปนนฺท, อมุกสฺมิํ โอกาเส อญฺญตโร ปุริโส ปชาปติํ เอตทโวจ อยฺยํ อุปนนฺทํ จีวเรน อจฺฉาเทสฺสามี”ติ. อตฺถาวุโส มํ โส อุปฏฺฐาโกติ. อถ โข อายสฺมา อุปนนฺโท สกฺยปุตฺโต เยน โส ปุริโส เตนุปสงฺกมิ, อุปสงฺกมิตฺวา ตํ ปุริสํ เอตทโวจ “สจฺจํ กิร มํ ตฺวํ อาวุโส จีวเรน อจฺฉาเทตุกาโมสี”ติ. อปิ เมยฺย เอวํ โหติ อยฺยํ อุปนนฺทํ จีวเรน อจฺฉาเทสฺสามีติ. สเจ โข มํ ตฺวํ อาวุโส จีวเรน อจฺฉาเทตุกาโมสิ, เอวรูเปน จีวเรน อจฺฉาเทหิ, กฺยาหํ เตน อจฺฉนฺโนปิ กริสฺสามิ, ยาหํ น ปริภุญฺชิสฺสามีติ.}

\csromanmatikaanchor{mtk.147}
\csromanmatikaanchor{mtk.148}
\csromantocmarkref{subsubsection}{ปญฺญตฺติ}{mtk.147}
\bookmark[dest={mtk.147},level=3]{ปญฺญตฺติ}
\csromantocmarkref{subsubsection}{สิกฺขาปทวิภงฺค, ปทภาชนีย}{mtk.148}
\bookmark[dest={mtk.148},level=3]{สิกฺขาปทวิภงฺค, ปทภาชนีย}
\prose{\csromanfolio{318}อถ โข โส ปุริโส อุชฺฌายติ ขิยฺยติ วิปาเจติ “มหิจฺฉา อิเม สมณา สกฺยปุตฺติยา อสนฺตุฏฺฐา, นยิเม สุกรา จีวเรน อจฺฉาเทตุํ, กถํ หิ นาม อยฺโย อุปนนฺโท มยา ปุพฺเพ อปฺปวาริโต มํ อุปสงฺกมิตฺวา จีวเร วิกปฺปํ อาปชฺชิสฺสตี”ติ. อสฺโสสุํ โข ภิกฺขู ตสฺส ปุริสสฺส อุชฺฌายนฺตสฺส ขิยฺยนฺตสฺส วิปาเจนฺตสฺส. เย เต ภิกฺขู อปฺปิจฺฉา ฯเปฯ เต อุชฺฌายฺยนฺติ ขิยฺยนฺติ วิปาเจนฺติ “กถํ หิ นาม อายสฺมา อุปนนฺโท สกฺยปุตฺโต ปุพฺเพ อปฺปวาริโต คหปติกํ อุปงฺกมิตฺวา จีวเร วิกปฺปํ อาปชฺชิสฺสตี”ติ. อถ โข เต ภิกฺขู อายสฺมนฺตํ อุปนนฺทํ สกฺยปุตฺตํ อเนก\-ปริยาเยน วิครหิตฺวา ภควโต เอตมตฺตํ อาโรเจสุํ ฯเปฯ “สจฺจํ กิร ตฺวํ อุปนนฺท ปุพฺเพ อปฺปาวาริโต คหปติกํ อุปสงฺกมิตฺวา จีวเร วิกปฺปํ อาปชฺชสี”ติ. สจฺจํ ภควาติ. ญาตโก เต อุปนนฺท อญฺญาตโกติ. อญฺญาตโก ภควาติ. อญฺญตโก โมฆปุริส อญฺญาตกสฺส น ชานาติ ปติรูปํ วา อปฺปติรูปํ วา สนฺตํ วา อสนฺตํ วา, ตตฺถ นาม ตฺวํ โมฆปุริส ปุพฺเพ อปฺปวาริโต อญฺญาตกํ คหปติกํ อุปสงฺกมิตฺวา จีวเร วิกปฺปํ อาปชฺชิสฺสสิ, เนตํ โมฆปุริส อปฺปสนฺนานํ วา ปสาทาย ฯเปฯ. เอวญฺจ ปน ภิกฺขเว อิมํ สิกฺขาปทํ อุทฺทิเสยฺยาถ\csromandash{}}

\proseitem{528}{\textbf{“ภิกฺขุํ ปเนว อุทฺทิสฺส อญฺญาตกสฺส คหปติสฺส วา คหปตานิยา วา จีวร}\-\textbf{เจตาปนฺนํ}\csromansharedfootnote{ebe070b961164e90}{จีวรเจตาปนํ (สฺยา)} \textbf{อุปกฺขฏํ โหติ ‘อิมินา จีวร}\-\textbf{เจตาปนฺเนน จีวรํ เจตาเปตฺวา อิตฺถนฺนามํ ภิกฺขุํ จีวเรน อจฺฉาเทสฺสามี’ติ, ตตฺร เจ โส ภิกฺขุ ปุพฺเพ อปฺปวาริโต อุปสงฺกมิตฺวา จีวเร วิกปฺปํ อาปชฺเชยฺย ‘สาธุ วต มํ อายสฺมา อิมินา จีวร}\-\textbf{เจตาปนฺเนน เอวรูปํ วา เอวรูปํ วา จีวรํ เจตาเปตฺวา อจฺฉาเทหี’ติ กลฺยาณ}\-\textbf{กมฺยตํ อุปาทาย, นิสฺสคฺคิยํ ปาจิตฺติยนฺ”}ติ.}

\proseitem{529}{\textbf{ภิกฺขุํ ปเนว อุทฺทิสฺสา}ติ ภิกฺขุสฺส\-ตฺถาย ภิกฺขุํ อารมฺมณํ กริตฺวา ภิกฺขุํ อจฺฉาเทตุกาโม.}

\prose{\textbf{อญฺญาตโก} นาม มาติโต วา ปิติโต วา ยาว สตฺตมา ปิตามหยุคา อสมฺพทฺโธ.}

\prose{\csromanfolio{319}\textbf{คหปติ} นาม โย โกจิ อคารํ อชฺฌาวสติ.}

\prose{\textbf{คหคตานี} นาม ยากาจิ อคารํ อชฺฌาวสติ.}

\prose{\textbf{จีวร}\-\textbf{เจตาปนฺนํ} นาม หิรญฺญํ วา สุวณฺณํ วา มุตฺตา วา มณิ วา ปวาโฬ วา ผลิโก วา ปฏโก วา สุตฺตํ วา กปฺปาโส วา.}

\prose{\textbf{อิมินา จีวรเจตาปนฺเนนา}ติ ปจฺจุปฏฺฐิเตน.}

\prose{\textbf{เจตาเปตฺวา}ติ ปริวตฺเตตฺวา.}

\prose{\textbf{อจฺฉาเทสฺสามี}ติ ทสฺสามิ.}

\prose{\textbf{ตตฺร เจ โส ภิกฺขู}ติ ยํ ภิกฺขุํ อุทฺทิสฺส จีวร\-เจตาปนฺนํ อุปกฺขฏํ โหติ, โส ภิกฺขุ. \textbf{ปุพฺเพ อปวาริโต}ติ ปุพฺเพ อวุตฺโต โหติ “กีทิเสน เต ภนฺเต จีวเรน อตฺโถ, กีทิสํ เต จีวรํ เจตาเปมี”ติ.}

\prose{\textbf{อุปสงฺกมิตฺวา}ติ ฆรํ คนฺตฺวา ยตฺถ กตฺถจิ อุปสงฺกมิตฺวา.}

\prose{\textbf{จีวเร วิกปฺปํ อาปชฺเชยฺยา}ติ อายตํ วา โหตุ วิตฺถตํ วา อปฺปิตํ วา สณฺหํ วา.}

\prose{\textbf{อิมินา จีวรเจตาปนฺเนนา}ติ ปจฺจุปฏฺฐิเตน.}

\prose{\textbf{เอวรูปํ วา เอวรูปํ วา}ติ อายตํ วา วิตฺถตํ วา อปฺปิตํ วา สณฺหํ วา.}

\prose{\textbf{เจตาเปตฺวา}ติ ปริวตฺเตตฺวา.}

\prose{\textbf{อจฺฉาเทหี}ติ ทชฺเชหิ.}

\prose{\textbf{กลฺยาณ}\-\textbf{กมฺยตํ อุปาทายา}ติ สาธตฺถิโก\csromansharedfootnote{432493805b5becc7}{สาธุตฺถิโก (สฺยา)} มหคฺฆตฺถิโก. ตสฺส วจเนน อายตํ วา วิตฺถตํ วา อปฺปิตํ วา สณฺหํ วา เจตาเปติ, ปโยเค ทุกฺกฏํ. ปฏิลาเภน นิสฺสคฺคิยํ โหติ. นิสฺสชฺชิตพฺพํ สํฆสฺส วา คณสฺสวา ปุคฺคลสฺสวา. เอวญฺจ ปน ภิกฺขเว นิสฺสชฺชิตพฺพํ ฯเปฯ “อิทํ เม ภนฺเต จีวรํ ปุพฺเพ อปฺปวาริโต อญฺญากกํ คหปติกํ อุปสงฺกมิตฺวา จีวเร วิกปฺปํ อาปนฺนํ นิสฺสคฺคิยํ, อิมาหํ สํฆสฺส นิสฺสชฺชามี”ติ ฯเปฯ ทเทยฺยาติ ฯเปฯ ทเทยฺยุนฺติ ฯเปฯ อายสฺมโต ทมฺมีติ.}

\csromanmatikaanchor{mtk.149}
\csromanmatikaanchor{mtk.150}
\csromantocmarknopage{subsection}{9. ทุติยอุปกฺขฏสิกฺขาปท}{mtk.149}
\bookmark[dest={mtk.149},level=2]{9. ทุติยอุปกฺขฏสิกฺขาปท}
\csromantocmarkref{subsubsection}{อุปนนฺทภิกฺขุวตฺถุ}{mtk.150}
\bookmark[dest={mtk.150},level=3]{อุปนนฺทภิกฺขุวตฺถุ}
\proseitem{530}{\csromanfolio{320}อญฺญาตเก อญฺญาตกสญฺญี ปุพฺเพ อปฺปวาริโต คหปติกํ อุปสงฺกมิตฺวา จีวเร วิกปฺปํ อาปชฺชติ, นิสฺสคฺคิยํ ปาจิตฺติยํ. อญฺญาตเก เวมติโก ปุพฺเพ อปฺปวาริโต คหปติกํ อุปสงฺกมิตฺวา จีวเร วิกปฺปํ อาปชฺชติ, นิสฺสคฺคิยํ ปาจิตฺติยํ. อญฺญาตเก ญาตกสญฺญี ปุพฺเพ อปฺปวาริโต คหปติกํ อุปสงฺกมิตฺวา จีวเร วิกปฺปํ อาปชฺชติ, นิสฺสคฺคิยํ ปาจิตฺติยํ.}

\prose{ญาตเก อญฺญาตกสญฺญี, อาปตฺติ ทุกฺกฏสฺส. ญาตเก เวมติโก, อาปตฺติ ทุกฺกฏสฺส. ญาตเก ญาตกสญฺญี. อนาปตฺติ.}

\proseitemwithcloserruled{531}{อนาปตฺติ ญาตกานํ ปวาริตานํ อญฺญสฺสตฺถาย อตฺตโน ธเนน มหคฺฆํ เจตาเปตุ\-กามสฺส อปฺปคฺฆํ เจตาเปติ อุมฺมตฺตกสฺส อาทิกมฺมิกสฺสาติ.}{อุปกฺขฏ\-สิกฺขาปทํ นิฏฺฐิตํ อฏฺฐมํ.}
\csromanmarkh{1. จีวรวคฺค}\csromanmarkhii{9. ทุติยอุปกฺขฏสิกฺขาปท}\csromanheadernomark{2}{1. จีวรวคฺค 9. ทุติยอุปกฺขฏสิกฺขาปท}
\proseitem{532}{เตน สมเยน พุทฺโธ ภควา สาวตฺถิยํ วิหรติ เชตวเน อนาถ\-ปิณฺฑิกสฺส อาราเม. เตน โข ปน สมเยน อญฺญตโร ปุริโส อญฺญตรํ ปุริสํ เอตทโวจ “อยฺยํ อุปนนฺทํ จีวเรน อจฺฉาเทสฺสามี”ติ. โสปิ เอวมาห “อหํปิ อยฺยํ อุปนนฺทํ จีวเรน อจฺฉาเทสฺสามี”ติ. อสฺโสสิ โข อญฺญตโร ปิณฺฑจาริโก ภิกฺขุ เตสํ ปุริสานํ อิมํ กถสลฺลาปํ. อถ โข โส ภิกฺขุ เยนายสฺมา อุปนนฺโท สกฺยปุตฺโต เตนุปสงฺกมิ, อุปสงฺกมิตฺวา อายสฺมนฺตํ อุปนนฺทํ สกฺยปุตฺตํ เอตทโวจ “มหาปุญฺโญสิ ตฺวํ อาวุโส อุปนนฺท, อมุกสฺมิํ โอกาเส อญฺญตโร ปุริโส อญฺญตรํ ปุริสํ เอตทโวจ “อยฺยํ อุปนนฺทํ จีวเรน อจฺฉาเทสฺสามี”ติ. โสปิ เอวมาห “อหํปิ อยฺยํ อุปนนฺทํ จีวเรน อจฺฉาเทสฺสามี”ติ. อตฺถาวุโส มํ เต อุปฏฺฐากาติ.}

\csromanmatikaanchor{mtk.151}
\csromanmatikaanchor{mtk.152}
\csromantocmarkref{subsubsection}{ปญฺญตฺติ}{mtk.151}
\bookmark[dest={mtk.151},level=3]{ปญฺญตฺติ}
\csromantocmarkref{subsubsection}{สิกฺขาปทวิภงฺค, ปทภาชนีย}{mtk.152}
\bookmark[dest={mtk.152},level=3]{สิกฺขาปทวิภงฺค, ปทภาชนีย}
\prose{อถ โข อายสฺมา อุปนนฺโท สกฺยปุตฺโต เยน เต ปุริสา เตนุปสงฺกมิ, อุปสงฺกมิตฺวา เต ปุริเส เอตทโวจ “สจฺจํ กิร มํ ตุเมฺห อาวุโส จีวเรหิ อจฺฉาเทตุกามาตฺถา”ติ. อปิ นยฺย เอวํ โหติ “อยฺยํ อุปนนฺทํ จีวเรหิ อจฺฉาเทสฺสามา”ติ. สเจ โข มํ \csromanfolio{321}ตุเมฺห อาวุโส จีวเรหิ อจฺฉาเทตุกามาตฺถ, เอวรูเปน จีวเรน อจฺฉาเทถ, กฺยาหํ เตหิ อจฺฉนฺโนปิ กริสฺสามิ, ยานาหํ น ปริภุญฺชิสฺสามีติ. อถ โข เต ปุริสา อุชฺฌายนฺติ ขิยฺยนฺติ วิปาเจนฺติ “มหิจฺฉา อิเม สมณา สกฺยปุตฺติยา อสนฺตุฏฺฐา, นยิเม สุกรา จีวเรหิ อจฺฉาเทตุํ, กถํ หิ นาม อยฺโย อุปนนฺโท อเมฺหหิ ปุพฺเพ อปฺปวาริโต อุปสงฺกมิตฺวา จีวเร วิกปฺปํ \textbf{อาปชฺชิสฺสตี”ติ.}}

\prose{อสฺโสสุํ โข ภิกฺขู เตสํ ปุริสานํ อุชฺฌายนฺตานํ ขิยฺยนฺตานํ วิปาเจนฺตานํ. เย เต ภิกฺขู อปฺปิจฺฉา ฯเปฯ เต อุชฺฌายนฺติ ขิยฺยนฺติ วิปาเจนฺติ “กถํ หิ นาม อายสฺมา อุปนนฺโท สกฺยปุตฺโต ปุพฺเพ อปฺปวาริโต คหปติเก อุปสงฺกมิตฺวา จีวเร วิกปฺปํ อาปชฺชิสฺสตี”ติ. อถ โข เต ภิกฺขู อายสฺมนฺตํ อุปนนฺทํ สกฺยปุตฺตํ อเนก\-ปริยาเยน วิครหิตฺวา ภควโต เอตมตฺถํ อาโรเจสุํ ฯเปฯ “สจฺจํ กิร ตฺวํ อุปนนฺท ปุพฺเพ อปฺปวาริโต คหปติเก อุปสงฺกมิตฺวา จีวเร วิกปฺปํ อาปชฺชสี”ติ. สจฺจํ ภควาติ. ญาตกา เต อุปนนฺท อญฺญาตกาติ. อญฺญาตกา ภควาติ. อญฺญาตโก โมฆปุริส อญฺญาตกานํ น ชานาติ ปติรูปํ วา อปฺปติรูปํ วา สนฺตํ วา อสนฺตํ วา, ตตฺถ นาม ตฺวํ โมฆปุริส ปุพฺเพ อปฺปวาริโต อญฺญาตเก คหปติเก อุปสงฺกิมิตฺวา จีวเร วิกปฺปํ อาปชฺชิสฺสสิ, เนตํ โมฆปุริส อปฺปสนฺนานํ วา ปสาทาย ฯเปฯ. เอวญฺจ ปน ภิกฺขเว อิมํ สิกฺขาปทํ อุทฺทิเสยฺยาถ\csromandash{}}

\proseitem{533}{\textbf{“ภิกฺขุํ ปเนว อุทฺทิสฺส อุภินฺนํ อญฺญาตกานํ คหปตีนํ วา คหปตานีนํ วา ปจฺเจก}\-\textbf{จีวร}\-\textbf{เจตาปนฺนานิ อุปกฺขฏานิ โหนฺติ ‘อิเมหิ มยํ ปจฺเจก}\-\textbf{จีวร}\-\textbf{เจตาปนฺเนหิ ปจฺเจก}\-\textbf{จีวรานิ เจตาเปตฺวา อิตฺถนฺนามํ ภิกฺขุํ จีวเรหิ อจฺฉาเทสฺสามา’ติ, ตตฺร เจ โส ภิกฺขุ ปุพฺเพ อปฺปวาริโต อุปสงฺกมิตฺวา จีวเร วิกปฺปํ อาปชฺเชยฺย ‘สาธุ วต มํ อายสฺมนฺโต อิเมหิ ปจฺเจก}\-\textbf{จีวร}\-\textbf{เจตาปนฺเนหิ เอวรูปํ วา เอวรูปํ วา จีวรํ เจตาเปตฺวา อจฺฉาเทถ อุโภว สนฺตา เอเกนา’ติ กลฺยาณ}\-\textbf{กมฺยตํ อุปาทาย, นิสฺสคฺคิยํ ปาจิตฺติยนฺ”}ติ.}

\proseitem{534}{\textbf{ภิกฺขุํ ปเนว อุทฺทิสฺสา}ติ ภิกฺขุสฺส\-ตฺถาย ภิกฺขุํ อารมฺมณํ กริตฺวา ภิกฺขุํ อจฺฉาเทตุกามา.}

\prose{\csromanfolio{322}\textbf{อุภินฺนนฺ}ติ ทฺวินฺนํ.}

\prose{\textbf{อญฺญาตกา} นาม มาติโต วา ปิติโต วา ยาว สตฺตมา ปิตามหยุคา อสมฺพทฺธา.}

\prose{\textbf{คหปตี} นาม เย เกจิ อคารํ อชฺฌาวสนฺติ.}

\prose{\textbf{คหปตานิโย} นาม ยา กาจิ อคารํ อชฺฌาวสนฺติ.}

\prose{\textbf{จีวร}\-\textbf{เจตาปนฺนานิ} นาม หิรญฺญา วา สุวณฺณา วา มุตฺตา วา มณี วา ปวาฬา วา ผลิกา วา ปฏกา วา สุตฺตา วา กปฺปาสา วา.}

\prose{\textbf{อิเมหิ ปจฺเจก}\-\textbf{จีวร}\-\textbf{เจตาปนฺเนหี}ติ ปจฺจุปฏฺฐิเตหิ.}

\prose{\textbf{เจตาเปตฺวา}ติ ปริวตฺเตตฺวา.}

\prose{\textbf{อจฺฉาเทสฺสามา}ติ ทสฺสาม.}

\prose{\textbf{ตตฺร เจ โส ภิกฺขู}ติ ยํ ภิกฺขุํ อุทฺทิสฺส จีวร\-เจตาปนฺนานิ อุปกฺขฏานิ โหนฺติ, โส ภิกฺขุ.}

\prose{\textbf{ปุพฺเพ อปฺปาวาริโต}ติ ปุพฺเพ อวุตฺโต โหติ “กีทิเสน เต ภนฺเต จีวเรน อตฺโถ กีทิสํ เต จีวรํ เจตาเปมา”ติ.}

\prose{\textbf{อุปสงฺกมิตฺวา}ติ ฆรํ คนฺตฺวา ยตฺถ กตฺถจิ อุปสงฺกมิตฺวา.}

\prose{\textbf{จีวเร วิกปฺปํ อาปชฺเชยฺยา}ติ อายตํ วา โหตุ วิตฺถตํ วา อปฺปิตํ วา สณฺหํ วา.}

\prose{\textbf{อิเมหิ ปจฺเจก}\-\textbf{จีวร}\-\textbf{เจตาปนฺเนหี}ติ ปจฺจุปฏฺฐิเตหิ.}

\prose{\textbf{เอวรูปํ วา เอวรูปํ วา}ติ อายตํ วา วิตฺถตํ วา อปฺปิตํ วา สณฺหํ วา.}

\prose{\textbf{เจตาเปตฺวา}ติ ปริวตฺเตตฺวา.}

\prose{\textbf{อจฺฉาเทถา}ติ ทชฺเชถ.}

\prose{\textbf{อุโภว สนฺตา เอเกนา}ติ เทฺวปิ ชนา เอเกน.}

\prose{\textbf{กลฺยาณ}\-\textbf{กมฺยตํ อุปาทายา}ติ สาธตฺถิโก มหคฺฆตฺถิโก.}

\csromanmatikaanchor{mtk.153}
\csromanmatikaanchor{mtk.154}
\csromantocmarknopage{subsection}{10. ราชสิกฺขาปท}{mtk.153}
\bookmark[dest={mtk.153},level=2]{10. ราชสิกฺขาปท}
\csromantocmarkref{subsubsection}{อุปนนฺทภิกฺขุวตฺถุ}{mtk.154}
\bookmark[dest={mtk.154},level=3]{อุปนนฺทภิกฺขุวตฺถุ}
\prose{ตสฺส วจเนน อายตํ วา วิตฺถตํ วา อปฺปิตํ วา สณฺหํ วา เจตาเปนฺติ, ปโยเค ทุกฺกฏํ. ปฏิลาเภน นิสฺสคฺคิยํ โหติ. นิสฺสชฺชิตพฺพํ สํฆสฺส วา คณสฺส วา ปุคฺคลสฺส วา, เอวญฺจ ปน ภิกฺขเว นิสฺสชฺชิตพฺพํ ฯเปฯ “อิทํ เม ภนฺเต จีวรํ ปุพฺเพ อปฺปวาริโต อญฺญาตเก คหปติเก อุปสงฺกมิตฺวา จีวเร วิกปฺปํ อาปนฺนํ นิสฺสคฺคิยํ, อิมาหํ \csromanfolio{323}สํฆสฺส นิสฺสชฺชามี”ติ ฯเปฯ ทเทยฺยาติ ฯเปฯ ทเทยฺยุนฺติ ฯเปฯ อายสฺมโต ทมฺมีติ.}

\proseitem{535}{อญฺญาตเก อญฺญาตกสญฺญี ปุพฺเพ อปฺปวาริโต คหปติเก อุปสงฺกมิตฺวา จีวเร วิกปฺปํ อาปชฺชติ, นิสฺสคฺคิยํ ปาจิตฺติยํ. อญฺญาตเก เวมติโก ปุพฺเพ อปฺปวาริโต คหปติเก อุปสงฺกมิตฺวา จีวเร วิกปฺปํ อาปชฺชติ, นิสฺสคฺคิยํ ปาจิตฺติยํ. อญฺญาตเก ญาตกสญฺญี ปุพฺเพ อปฺปวาริโต คหปติเก อุปสงฺกมิตฺวา จีวเร วิกปฺปํ อาปชฺชติ, นิสฺสคฺคิยํ วาจิตฺติยํ.}

\prose{ญาตเก อญฺญาตกสญฺญี, อาปตฺติ ทุกฺกฏสฺส. ญาตเก เวมติโก, อาปตฺติ ทุกฺกฏสฺส. ญาตเก ญาตกสญฺญี, อนาปตฺติ.}

\proseitemwithcloserruled{536}{อนาปตฺติ ญาตกานํ, ปวาริตานํ, อญฺญสฺสตฺถาย, อตฺตโน ธเนน, มหคฺฆํ เจตาเปตุ\-กามานํ อปฺปคฺฆํ เจตาเปติ, อุมฺมตฺตกสฺส อาทิกมฺมิกสฺสาติ.}{ทุติยอุปกฺขฏสิกฺขาปทํ นิฏฺฐิตํ นวมํ.}
\csromanmarkh{1. จีวรวคฺค}\csromanmarkhii{10. ราชสิกฺขาปท}\csromanheadernomark{2}{1. \textbf{จีวรวคฺค} 10. ราชสิกฺขาปท}
\proseitem{537}{เตน สมเยน พุทฺโธ ภควา สาวตฺถิยํ วิหรติ เชตวเน อนาถ\-ปิณฺฑิกสฺส อาราเม. เตน โข ปน สมเยน อายสฺมโต อุปนนฺทสฺส สกฺยปุตฺตสฺส อุปฏฺฐาโก มหามตฺโต อายสฺมโต อุปนนฺทสฺส สกฺยปุตฺตสฺส ทูเตน จีวร\-เจตาปนฺนํ ปาเหสิ “อิมินา จีวร\-เจตาปนฺเนน จีวรํ เจตาเปตฺวา อยฺยํ อุปนนฺทํ จีวเรน อจฺฉาเทหี”ติ. อถ โข โส ทูโต เยนายสฺมา อุปนนฺโท สกฺยปุตฺโต เตนุปสงฺกมิ, อุปสงฺกมิตฺวา อายสฺมนฺตํ อุปนนฺทํ สกฺยปุตฺตํ เอตทโวจ “อิทํ โข ภนฺเต อายสฺมนฺตํ อุทฺทิสฺส จีวร\-เจตาปนฺนํ อาภตํ, ปฏิคฺคณฺหาตุ อายสฺมา จีวรเจตาปนฺนนฺ”ติ. เอวํ วุตฺเต อายสฺมา อุปนนฺโท สกฺยปุตฺโต ตํ ทูตํ เอตทโวจ “น โข มยํ อาวุโส จีวร\-เจตาปนฺนํ ปฏิคฺคณฺหาม, จีวรญฺจ โข มยํ ปฏิคฺคณฺหาม กาเลน กปฺปิยนฺ”ติ. เอวํ วุตฺเต โส ทูโต อายสฺมนฺตํ อุปนนฺทํ สกฺยปุตฺตํ เอตทโวจ “อตฺถิ ปนายสฺมโต โกจิ \csromanfolio{324}เวยฺยาวจฺจกโร”ติ. เตน โข ปน สมเยน อญฺญตโร อุปาสโก อารามํ อคมาสิ เกนจิเทว กรณีเยน. อถ โข อายสฺมา อุปนนฺโท สกฺยปุตฺโต ตํ ทูตํ เอตทโวจ “เอโส โข อาวุโส อุปาสโก ภิกฺขูนํ เวยฺยาวจฺจกโร”ติ. อถ โข โส ทูโต ตํ อุปาสกํ สญฺญาเปตฺวา เยนายสฺมา อุปนนฺโท สกฺยปุตฺโต เตนุปสงฺกมิ, อุปสงฺกมิตฺวา อายสฺมนฺตํ อุปนนฺทํ สกฺยปุตฺตํ เอตทโวจ “ยํ โข ภนฺเต อายสฺมา เวยฺยาวจฺจกรํ นิทฺทิสิ, สญฺญตฺโต โส มยา, อุปสงฺกมตุ อายสฺมา กาเลน, จีวเรน ตํ อจฺฉาเทสฺสตี”ติ.}

\prose{เตน โข ปน สมเยน โส มหามตฺโต อายสฺมโต อุปนนฺทสฺส สกฺยปุตฺตสฺส สนฺติเก ทูตํ ปาเหสิ “ปริภุญฺชตุ อยฺโย ตํ จีวรํ, อิจฺฉาม มยํ อยฺเยน ตํ จีวรํ ปริภุตฺตนฺ”ติ, อถ โข อายสฺมา อุปนนฺโท สกฺยปุตฺโต ตํ อุปาสกํ น กิญฺจิ อวจาสิ. ทุติยมฺปิ โข โส มหามตฺโต อายสฺมโต อุปนนฺทสฺส สกฺยปุตฺตสฺส สนฺติเก ทูตํ ปาเหสิ “ปริภุญฺชตุ อยฺโย ตํ จีวรํ, อิจฺฉาม มยํ อยฺเยน ตํ จีวรํ ปริภุตฺตนฺ”ติ. ทุติยมฺปิ โข อายสฺมา อุปนนฺโท สกฺยปุตฺโต ตํ อุปาสกํ น กิญฺจิ อวจาสิ. ตติยมฺปิ โข โส มหามตฺโต อายสฺมโต อุปนนฺทสฺส สกฺยปุตฺตสฺส สนฺติเก ทูตํ ปาเหสิ “ปริภุญฺชตุ อยฺโย ตํ จีวรํ, อิจฺฉาม มยํ อยฺเยน ตํ จีวรํ ปริภุตฺตนฺ”ติ.}

\prose{เตน โข ปน สมเยน เนคมสฺส สมโย โหติ, เนคเมน จ กติกา กตา โหติ “โย ปจฺฉา อาคจฺฉติ ปญฺญาสํ พทฺโธ”ติ. อถ โข อายสฺมา อุปนนฺโท สกฺยปุตฺโต เยน โส อุปาสโก เตนุปสงฺกมิ, อุปสงฺกมิตฺวา ตํ อุปาสกํ เอตทโวจ “อตฺโถ เม อาวุโส จีวเรนา”ติ. อชฺชโณฺห ภนฺเต อาคเมหิ, อชฺช เนคมสฺส สมโย, เนคเมน จ กติกา กตา โหติ “โย ปจฺฉา อาคจฺฉติ ปญฺญาสํ พทฺโธ”ติ. “อชฺเชว เม อาวุโส จีวรํ เทหี”ติ โอวฏฺฏิกาย ปรามสิ. อถ โข โส อุปาสโก อายสฺมตา อุปนนฺเทน สกฺยปุตฺเตน นิปฺปีฬิยมาโน อายสฺมโต อุปนนฺทสฺส สกฺยปุตฺตสฺส จีวรํ เจตาเปตฺวา ปจฺฉา อคมาสิ. มนุสฺสา ตํ อุปาสกํ เอตทโวจุํ “กิสฺส ตฺวํ อยฺโย ปจฺฉา อาคโต ปญฺญาสํ ชีโนสี”ติ.}

\csromanmatikaanchor{mtk.155}
\csromantocmarkref{subsubsection}{ปญฺญตฺติ}{mtk.155}
\bookmark[dest={mtk.155},level=3]{ปญฺญตฺติ}
\prose{อถ โข โส อุปาสโก เตสํ มนุสฺสานํ เอตมตฺถํ อาโรเจสิ. มนุสฺสา อุชฺฌยนฺติ ขิยฺยนฺติ วิปาเจนฺติ “มหิจฺฉา อิเม สมณา สกฺยปุตฺติยา \csromanfolio{325}\textbf{อสนฺตุฏฺฐา, นยิเมสํ สุกรํ เวยฺยาวจฺจํปิ กาตุํ, กถํ หิ นาม} \textbf{อายสฺมา อุปนนฺโท อุปาสเกน ‘อชฺชโณฺห ภนฺเต อาคเมหี’ติ วุจฺจมาโน} \textbf{นาคเมสฺสตี”ติ. อสฺโสสุํ โข ภิกฺขู เตสํ มนุสฺสานํ อุชฺฌายนฺตานํ} \textbf{ขิยฺยนฺตานํ วิปาเจนฺตานํ. เย เต ภิกฺขู อปฺปิจฺฉา, เต อุชฺฌายนฺติ} \textbf{ขิยฺยนฺติ วิปาเจนฺติ “กถํ หิ นาม อายสฺมา อุปนนฺโท สกฺยปุตฺโต} \textbf{อุปาสเกน ‘อชฺชโณฺห ภนฺเต อาคเมหี’ติ วุจฺจมาโน นาคเมสฺสตี”ติ. อถ โข} \textbf{เต ภิกฺขู อายสฺมนฺตํ อุปนนฺทํ สกฺยปุตฺตํ อเนก}\-\textbf{ปริยาเยน} \textbf{วิครหิตฺวา ภควโต เอตมตฺถํ อาโรเจสุํ ฯเปฯ “สจฺจํ กิร ตฺวํ} \textbf{อุปนนฺท อุปาสเกน ‘อชฺชโณฺห ภนฺเต อาคเมหี’ติ วุจฺจมาโน นาคเมสี”ติ.} สจฺจํ ภควาติ. วิครหิ พุทฺโธ ภควา ฯเปฯ. กถํ หิ นาม ตฺวํ โมฆปุริส อุปาสเกน ‘อชฺชโณฺห ภนฺเต อาคเมหี’ติ วุจฺจมาโน นาคเมสฺสสิ, เนตํ โมฆปุริส อปฺปสนฺนานํ วา ปสาทาย ฯเปฯ. เอวญฺจ ปน ภิกฺขเว อิมํ สิกฺขาปทํ อุทฺทิเสยฺยาถ\csromandash{}}

\csromanmatikaanchor{mtk.156}
\csromantocmarkref{subsubsection}{สิกฺขาปทวิภงฺค, ปทภาชนีย}{mtk.156}
\bookmark[dest={mtk.156},level=3]{สิกฺขาปทวิภงฺค, ปทภาชนีย}
\proseitem{538}{\textbf{“ภิกฺขุํ ปเนว อุทฺทิสฺส ราชา วา ราชโภคฺโค วา พฺราหฺมโณ วา คหปติโก วา ทูเตน จีวร}\-\textbf{เจตาปนฺนํ ปหิเณยฺย ‘อิมินา จีวร}\-\textbf{เจตาปนฺเนน จีวรํ เจตาเปตฺวา อิตฺถนฺนามํ ภิกฺขุํ จีวเรน อจฺฉาเทหี’ติ. โส เจ ทูโต ตํ ภิกฺขุํ อุปสงฺกมิตฺวา เอวํ วเทยฺย ‘อิทํ โข ภนฺเต อายสฺมนฺตํ อุทฺทิสฺส จีวร}\-\textbf{เจตาปนฺนํ อาภตํ, ปฏิกฺกณฺหาตุ อายสฺมา จีวรเจตาปนฺนนฺ’ติ. เตน ภิกฺขุนา โส ทูโต เอวมสฺส วจนีโย ‘น โข มยํ อาวุโส จีวร}\-\textbf{เจตาปนฺนํ ปฏิคฺคณฺหาม, จีวรญฺจ โข มยํ ปฏิคฺคณฺหาม กาเลน กปฺปิยนฺ’ติ. โส เจ ทูโต ตํ ภิกฺขุํ เอวํ วเทยฺย ‘อตฺถิ ปนายสฺมโต โกจิ เวยฺยาวจฺจกโร’ติ. จีวรตฺถิเกน ภิกฺขเว ภิกฺขุนา เวยฺยาวจฺจกโร นิทฺทิสิตพฺโพ อารามิโก วา อุปาสโก วา ‘เอโส โข อาวุโส ภิกฺขูนํ เวยฺยาวจฺจกโร’ติ. โส เจ ทูโต ตํ เวยฺยาวจฺจกรํ สญฺญาเปตฺวา ตํ ภิกฺขุํ อุปสงฺกมิตฺวา เอวํ วเทยฺย ‘ยํ โข ภนฺเต อายสฺมา เวยฺยาวจฺจกรํ นิทฺทิสิ, สญฺญตฺโต โส มยา, อุปสงฺกมตุ อายสฺมา กาเลน, จีวเรน ตํ อจฺฉาเทสฺสตี’ติ. จีวรตฺถิเกน ภิกฺขเว ภิกฺขุนา เวยฺยาวจฺจกโร อุปสงฺกมิตฺวา ทฺวตฺติกฺขตฺตุํ โจเทตพฺโพ สาเรตพฺโพ ‘อตฺโถ เม อาวุโส จีวเรนา’ติ, ทฺวตฺติกฺขตฺตุํ โจทยมาโน สารยมาโน ตํ จีวรํ อภินิปฺผาเทยฺย, อิจฺเจตํ กุสลํ, โน เจ อภินิปฺผาเทยฺย,} \csromanfolio{326}\textbf{จตุกฺขตฺตุํ ปญฺจกฺขตฺตุํ ฉกฺขตฺตุปรมํ ตุณฺหีภูเตน อุทฺทิสฺส ฐาตพฺพํ, จตุกฺขตฺตุํ ปญฺจกฺขตฺตุํ ฉกฺขตฺตุปรมํ ตุณฺหีภูโต อุทฺทิสฺส ติฏฺฐามาโน ตํ จีวรํ อภินิปฺผาเทยฺย, อิจฺเจตํ กุสลํ, ตโต เจ อุตฺตริ วายมมาโน ตํ จีวรํ อภินิปฺผาเทยฺย, นิสฺสคฺคิยํ ปาจิตฺติยํ. โน เจ อภินิปฺผาเทยฺย, ยตสฺส จีวร}\-\textbf{เจตาปนฺนํ อาภตํ, ตตฺถ สามํ วา คนฺตพฺพํ ทูโต วา ปาเหตพฺโพ ‘ยํ โข ตุเมฺห อายสฺมนฺโต ภิกฺขุํ อุทฺทิสฺส จีวร}\-\textbf{เจตาปนฺนํ ปหิณิตฺถ, น ตํ ตสฺส ภิกฺขุโน กิญฺจิ อตฺถํ อนุโภติ, ยุญฺชนฺตา}\-\textbf{ยสฺมนฺโต สกํ, มา โว สกํ วินสฺสา’ติ, อยํ ตตฺถ สามีจี”}ติ.}

\proseitem{539}{\textbf{ภิกฺขุํ ปเนว อุทฺทิสฺสา}ติ ภิกฺขุสฺส\-ตฺถาย ภิกฺขุํ อารมฺมณํ กริตฺวา ภิกฺขุํ อจฺฉาเทตุกาโม.}

\prose{\textbf{ราชา} นาม โย โกจิ รชฺชํ กาเรติ.}

\prose{\textbf{ราชโภคฺโค} นาม โย โกจิ รญฺโญ ภตฺต\-เวตนา\-หาโร.}

\prose{\textbf{พฺราหฺมโณ} นาม ชาติยา พฺราหฺมโณ.}

\prose{\textbf{คหปติโก} นาม ฐาเปตฺวา ราชํ ราชโภคฺคํ พฺราหฺมณํ อวเสโส คหปติโก นาม.}

\prose{\textbf{จีวร}\-\textbf{เจตาปนฺนํ} นาม หิรญฺญํ วา สุวณฺณํ วา มุตฺตา วา มณิ วา.}

\prose{\textbf{อิมินา จีวรเจตาปนฺเนนา}ติ ปจฺจุปฏฺฐิเตน.}

\prose{\textbf{เจตาเปตฺวา}ติ ปริวตฺเตตฺวา.}

\prose{\textbf{อจฺฉาเทหี}ติ ทชฺเชหิ.}

\prose{โส เจ ทูโต ตํ ภิกฺขุํ อุปสงฺกมิตฺวา เอวํ วเทยฺย “อิทํ โข ภนฺเต อายสฺมนฺตํ อุทฺทิสฺส จีวร\-เจตาปนฺนํ อาภตํ, ปฏิคฺคณฺหาตุ อายสฺมา จีวรเจตาปนฺนนฺ”ติ. เตน ภิกฺขุนา โส ทูโต เอวมสฺส วจนีโย “น โข มยํ อาวุโส จีวร\-เจตาปนฺนํ ปฏิคฺคณฺหาม, จีวรญฺจ โข มยํ ปฏิคฺคณฺหาม กาเลน กปฺปิยนฺ”ติ. โส เจ ทูโต ตํ ภิกฺขุํ เอวํ วเทยฺย “อตฺถิ ปนายสฺมโต โกจิ เวยฺยาวจฺจกโร”ติ, จีวรตฺถิเกน ภิกฺขเว ภิกฺขุนา เวยฺยาวจฺจกโร นิทฺทิสิตพฺโพ อารามิโก วา อุปาสโก วา “เอโส โข อาวุโส ภิกฺขูนํ เวยฺยาวจฺจกโร”ติ. น วตฺตพฺโพ “ตสฺส เทหีติ วา, โส วา นิกฺขิปิสฺสติ, โส วา ปริวตฺเตสฺสติ, โส วา เจตาเปสฺสตี”ติ.}

\prose{\csromanfolio{327}โส เจ ทูโต ตํ เวยฺยาวจฺจกรํ สญฺญาเปตฺวา ตํ ภิกฺขุํ อุปสงฺกมิตฺวา เอวํ วเทยฺย “ยํ โข ภนฺเต อายสฺมา เวยฺยาวจฺจกรํ นิทฺทิสิ, สญฺญตฺโต โส มยา, อุปสงฺกมตุ อายสฺมา กาเลน, จีวเรน ตํ อจฺฉาเทสฺสตี”ติ. จีวรตฺถิเกน ภิกฺขเว ภิกฺขุนา เวยฺยาวจฺจกโร อุปสงฺกมิตฺวา ทฺวตฺติกฺขตฺตุํ โจเทตพฺโพ สาเรตพฺโพ “อตฺโถ เม อาวุโส จีวเรนา”ติ. น วตฺตพฺโพ “เทหิ เม จีวรํ, อาหร เม จีวรํ, ปริวตฺเตหิ เม จีวรํ, เจตาเปหิ เม จีวรนฺ”ติ. ทุติยมฺปิ วตฺตพฺโพ. ตติยมฺปิ วตฺตพฺโพ. สเจ อภินิปฺผาเทติ, อิจฺเจตํ กุสลํ, โน เจ อภินิปฺผาเทติ, ตตฺถ คนฺตฺวา ตุณฺหีภูเตน อุทฺทิสฺส ฐาตพฺพํ, น อาสเน นิสีทิตพฺพํ, น อามิสํ ปฏิคฺคเหตพฺพํ, น ธมฺโม ภาสิตพฺโพ, “กิํการณา อาคโตสี”ติ ปุจฺฉิยมาโน “ชานาหิ อาวุโส”ติ วตฺตพฺโพ. สเจ อาสเน วา นิสีทติ อามิสํ วา ปฏิคฺคณฺหติ ธมฺมํ วา ภาสติ, ฐานํ ภญฺชติ. ทุติยมฺปิ ฐาตพฺพํ. ตติยมฺปิ ฐาตพฺพํ. จตุกฺขตฺตุํ โจเทตฺวา จตุกฺขตฺตุํ ฐาตพฺพํ, ปญฺจกฺขตฺตุํ โจเทตฺวา ทฺวิกฺขตฺตุํ ฐาตพฺพํ, ฉกฺขตฺตุํ โจเทตฺวา น ฐาตพฺพํ, ตโต เจ อุตฺตริ วายมมาโน ตํ จีวรํ อภินิปฺผาเทติ, ปโยเค ทุกฺกฏํ. ปฏิลาเภน นิสฺสคฺคิยํ โหติ. นิสฺสชฺชิตพฺพํ สํฆสฺส วา คณสฺส วา ปุคฺคลสฺส วา. เอวญฺจ ปน ภิกฺขเว นิสฺสชฺชิตพฺพํ ฯเปฯ “อิทํ เม ภนฺเต จีวรํ อติเรก\-ติ\-กฺขตฺตุํ โจทนาย, อติเรก\-ฉ\-กฺขตฺตุํ ฐาเนน อภินิ\-ปฺผาทิตํ นิสฺสคฺคิยํ. อิมาหํ สํฆสฺส นิสฺสชฺชามี”ติ ฯเปฯ ทเทยฺยาติ ฯเปฯ ทเทยฺยุนฺติ ฯเปฯ อายสฺมโต ทมฺมีติ.}

\prose{โน เจ อภินิปฺผาเทยฺย, ยตสฺส จีวร\-เจตาปนฺนํ อาภตํ, ตตฺถ สามํ วา คนฺตพฺพํ, ทูโต วา ปาเหตพฺโพ “ยํ โข ตุเมฺห อายสฺมนฺโต ภิกฺขุํ อุทฺทิสฺส จีวร\-เจตาปนฺนํ ปหิณิตฺถ, น ตํ ตสฺส ภิกฺขุโน กิญฺจิ อตฺถํ อนุโภติ, ยุญฺชนฺตา\-ยสฺมนฺโต สกํ, มา โว สกํ วินสฺสา”ติ.}

\prose{\textbf{อยํ ตตฺถ สามีจี}ติ อยํ ตตฺถ อนุธมฺมตา.}

\proseitem{540}{อภิเรกติกฺขตฺตุํ โจทนาย อติเรก\-ฉ\-กฺขตฺตุํ ฐาเนน อติเรกสญฺญี อภินิปฺผาเทติ, นิสฺสคฺคิยํ ปาจิตฺติยํ. อติเรก\-ติ\-กฺขตฺตุํ โจทนาย อติเรก\-ฉ\-กฺขตฺตุํ ฐาเนน เวมติโก อภินิปฺผาเทติ, นิสฺสคฺคิยํ ปาจิตฺติยํ. อติเรก\-ติ\-กฺขตฺตุํ โจทนาย อติเรก\-ฉ\-กฺขตฺตุํ ฐาเนน อูนกสญฺญี อภินิปฺผาเทติ, นิสฺสคฺคิยํ ปาจิตฺติยํ.}

\csromanmatikaanchor{mtk.157}
\csromanmatikaanchor{mtk.160}
\csromantocmarkref{subsection}{อุทฺทานคาถา}{mtk.157}
\bookmark[dest={mtk.157},level=2]{อุทฺทานคาถา}
\prose{\csromanfolio{328}อูนก\-ติ\-กฺขตฺตุํ โจทนาย อูนก\-ฉ\-กฺขตฺตุํ ฐาเนน อติเรกสญฺญี, อาปตฺติ ทุกฺกฏสฺส. อูนก\-ติ\-กฺขตฺตุํ โจทนาย อูนก\-ฉ\-กฺขตฺตุํ ฐาเนน เวมติโก, อาปตฺติ ทุกฺกฏสฺส. อูนก\-ติ\-กฺขตฺตุํ โจทนาย อูนก\-ฉ\-กฺขตฺตุํ ฐาเนน อูนกสญฺญี, อนาปตฺติ.}

\proseitem{541}{อนาปตฺติ ติกฺขตฺตุํ โจทนาย ฉกฺขตฺตุํ ฐาเนน, อูนก\-ติ\-กฺขตฺตุํ โจทนาย อูนก\-ฉ\-กฺขตฺตุํ ฐาเนน, อโจทิยมาโน เทติ, สามิกา โจเทตฺวา เทนฺติ, อุมฺมตฺตกสฺส อาทิกมฺมิกสฺสาติ.}

\vspace{\tipitakaparskip}
\csromancenter{ราช\-สิกฺขาปทํ นิฏฺฐิตํ ทสมํ.}
\nitthitammidruled{กถินวคฺโค ปฐโม.}
\csromancenter{ตสฺสุทฺทานํ}
\setlength{\csromangathapagewidth}{0pt}
\setlength{\csromangathawakwidth}{0pt}
\csromangathameasuregroup{อุพฺภตํ กถินํ ตีณิ, \\ อญฺญาตกานิ ตีเณว,}{\csromangathabat{อุพฺภตํ กถินํ ตีณิ,}{โธวนญฺจ ปฏิคฺคโห.} \\ \csromangathabat{อญฺญาตกานิ ตีเณว,}{อุภินฺนํ ทูตเกน จาติ.}}
\csromangathasetleft{อุพฺภตํ กถินํ ตีณิ, \\ อญฺญาตกานิ ตีเณว,}
\csromangathagroup{\csromangathastanza{\csromangathabat{อุพฺภตํ กถินํ ตีณิ,}{โธวนญฺจ ปฏิคฺคโห.} \\ \csromangathabat{อญฺญาตกานิ ตีเณว,}{อุภินฺนํ ทูตเกน จาติ.}}}

\csromanmatikaanchor{mtk.158}
\csromanmatikaanchor{mtk.159}
\csromantocmarknopage{section}{2. โกสิยวคฺค}{mtk.158}
\bookmark[dest={mtk.158},level=1]{2. โกสิยวคฺค}
\csromantocmarknopage{subsection}{1. โกสิยสิกฺขาปท}{mtk.159}
\bookmark[dest={mtk.159},level=2]{1. โกสิยสิกฺขาปท}
\csromantocmarkref{subsubsection}{ฉพฺพคฺคิยภิกฺขุวตฺถุ}{mtk.160}
\bookmark[dest={mtk.160},level=3]{ฉพฺพคฺคิยภิกฺขุวตฺถุ}
\csromanmarkh{2. โกสิยวคฺค}\csromanmarkhii{1. โกสิย\-สิกฺขาปท}\csromanheadernomark{2}{2. \textbf{โกสิยวคฺค 1. โกสิย}\-\textbf{สิกฺขาปท}}
\csromanmatikaanchor{mtk.161}
\csromanmatikaanchor{mtk.162}
\csromantocmarkref{subsubsection}{ปญฺญตฺติ}{mtk.161}
\bookmark[dest={mtk.161},level=3]{ปญฺญตฺติ}
\csromantocmarkref{subsubsection}{สิกฺขาปทวิภงฺค, ปทภาชนีย}{mtk.162}
\bookmark[dest={mtk.162},level=3]{สิกฺขาปทวิภงฺค, ปทภาชนีย}
\proseitem{542}{เตน สมเยน พุทฺโธ ภควา อาฬวิยํ วิหรติ อคฺคาฬเว เจติเย. เตน โข ปน สมเยน ฉพฺพคฺคิยา ภิกฺขู โกสิยการเก อุปสงฺกมิตฺวา เอวํ วทนฺติ “พหู อาวุโส โกสการเก ปจถ, อมฺหากํปิ ทสฺสถ, มยมฺปิ อิจฺฉาม โกสิยมิสฺสกํ สนฺถตํ กาตุนฺ”ติ. เต อุชฺฌายนฺติ ขิยฺยนฺติ วิปาเจนฺติ “กถํ หิ นาม สมณา สกฺยปุตฺติยา อเมฺห อุปสงฺกมิตฺวา เอวํ วกฺขนฺติ ‘พหู อาวุโส โกสการเก ปจถ, อมฺหากํปิ ทสฺสถ, มยํปิ อิจฺฉาม โกสิยมิสฺสกํ สนฺถตํ กาตุนฺ’ติ, อมฺหากํปิ อลาภา, อมฺหากํปิ ทุลฺลทฺธํ, เย มยํ อาชีวสฺส เหตุ ปุตฺตทารสฺส การณา พหู ขุทฺทเก ปาเณ สงฺฆาตํ อาปาเทมา”ติ. อสฺโสสุํ โข ภิกฺขู เตสํ มนุสฺสานํ อุชฺฌายนฺตานํ ขิยฺยนฺตานํ วิปาเจนฺตานํ. เย เต ภิกฺขู อปฺปิจฺฉา ฯเปฯ เต อุชฺฌายนฺติ ขิยฺยนฺติ วิปาเจนฺติ “กถํ หิ นาม ฉพฺพคฺคิยา ภิกฺขู โกสิยการเก อุปสงฺกมิตฺวา เอวํ วกฺขนฺติ “พหู อาวุโส โกสการเก ปจถ, อมฺหากํปิ ทสฺสถ, มยํปิ อิจฺฉาม โกสิยมิสฺสกํ \csromanfolio{329}สนฺถตํ กาตุนฺ”ติ. อถ โข เต ภิกฺขู ฉพฺพคฺคิเย ภิกฺขู อเนก\-ปริยาเยน วิครหิตฺวา ภควโต เอตมตฺถํ อาโรเจสุํ ฯเปฯ “สจฺจํ กิร ตุเมฺห ภิกฺขเว โกสิยการเก อุปสงฺกมิตฺวา เอวํ วเทถ พหู อาวุโส โกสการเก ปจถ, อมฺหากํปิ ทสฺสถ, มยํปิ อิจฺฉาม โกสิยมิสฺสกํ สนฺถตํ กาตุนฺ”ติ. สจฺจํ ภควาติ. วิครหิ พุทฺโธ ภควา ฯเปฯ. กถํ หิ นาม ตุเมฺห โมฆปุริสา โกสิยการเก อุปสงฺกมิตฺวา เอวํ วกฺขถ “พหู อาวุโส โกสการเก ปจถ, อมฺหากํปิ ทสฺสถ, มยํปิ อิจฺฉาม โกสิยมิสฺสกํ สนฺถตํ กาตุนฺ”ติ. เนตํ โมฆปุริสา อปฺปสนฺนานํ วา ปสาทาย ฯเปฯ. เอวญฺจ ปน ภิกฺขเว อิมํ สิกฺขาปทํ อุทฺทิเสยฺยาถ\csromandash{}}

\proseitem{543}{\textbf{“โย ปน ภิกฺขุ โกสิยมิสฺสกํ สนฺถตํ การาเปยฺย, นิสฺสคฺคิยํ ปาจิตฺติยนฺ”}ติ.}

\proseitem{544}{\textbf{โย ปนา}ติ โย ยาทิโส ฯเปฯ. \textbf{ภิกฺขู}ติ ฯเปฯ อยํ อิมสฺมิํ อตฺเถ อธิปฺเปโต “ภิกฺขู”ติ.}

\prose{\textbf{สนฺถตํ} นาม สนฺถริตฺวา กตํ โหติ อวายิมํ.}

\prose{\textbf{การาเปยฺยา}ติ เอเกนปิ โกสิยํสุนา มิสฺสิตฺวา กโรติ วา การาเปติ วา, ปโยเค ทุกฺกฏํ. ปฏิลาเภน นิสฺสคฺคิยํ โหติ. นิสฺสชฺชิตพฺพํ สํฆสฺส วา คณสฺส วา ปุคฺคลสฺส วา, เอวญฺจ ปน ภิกฺขเว นิสฺสชฺชิตพฺพํ ฯเปฯ “อิทํ เม ภนฺเต โกสิยมิสฺสกํ สนฺถตํ การาปิตํ นิสฺสคฺคิยํ, อิมาหํ สํฆสฺส นิสฺสชฺชามี”ติ ฯเปฯ ทเทยฺยาติ ฯเปฯ ทเทยฺยุนฺติ ฯเปฯ อายสฺมโต ทมฺมีติ.}

\proseitem{545}{อตฺตนา วิปฺปกตํ อตฺตนา ปริโยสาเปติ, นิสฺสคฺคิยํ ปาจิตฺติยํ. อตฺตนา วิปฺปกตํ ปเรหิ ปริโยสาเปติ, นิสฺสคฺคิยํ ปาจิตฺติยํ. ปเรหิ วิปฺปกตํ อตฺตนา ปริโยสาเปติ, นิสฺสคฺคิยํ ปาจิตฺติยํ. ปเรหิ วิปฺปกตํ ปเรหิ ปริโยสาเปติ, นิสฺสคฺคิยํ ปาจิตฺติยํ.}

\prose{อญฺญสฺสตฺถาย กโรติ วา การาเปติ วา, อาปตฺติ ทุกฺกฏสฺส. อญฺเญน กตํ ปฏิลภิตฺวา ปริภุญฺชติ, อาปตฺติ ทุกฺกฏสฺส.}

\csromanmatikaanchor{mtk.163}
\csromanmatikaanchor{mtk.164}
\csromanmatikaanchor{mtk.165}
\csromanmatikaanchor{mtk.166}
\csromantocmarknopage{subsection}{2. สุทฺธกาฬกสิกฺขาปท}{mtk.163}
\bookmark[dest={mtk.163},level=2]{2. สุทฺธกาฬกสิกฺขาปท}
\csromantocmarkref{subsubsection}{ฉพฺพคฺคิยภิกฺขุวตฺถุ}{mtk.164}
\bookmark[dest={mtk.164},level=3]{ฉพฺพคฺคิยภิกฺขุวตฺถุ}
\csromantocmarkref{subsubsection}{ปญฺญตฺติ}{mtk.165}
\bookmark[dest={mtk.165},level=3]{ปญฺญตฺติ}
\csromantocmarkref{subsubsection}{สิกฺขาปทวิภงฺค, ปทภาชนีย}{mtk.166}
\bookmark[dest={mtk.166},level=3]{สิกฺขาปทวิภงฺค, ปทภาชนีย}
\proseitemwithcloserruled{546}{\csromanfolio{330}อนาปตฺติ วิตานํ วา ภูมตฺถรณํ วา สาณิปาการํ วา ภิสิํ วา พิพฺโพหนํ วา กโรติ, อุมฺมตฺตกสฺส อาทิกมฺมิกสฺสาติ.}{โกสิย\-สิกฺขาปทํ นิฏฺฐิตํ ปฐมํ.}
\csromanmarkh{2. โกสิยวคฺค}\csromanmarkhii{2. สุทฺธกาฬก\-สิกฺขาปท}\csromanheadernomark{2}{2. \textbf{โกสิยวคฺค 2. สุทฺธกาฬก}\-\textbf{สิกฺขาปท}}
\proseitem{547}{เตน สมเยน พุทฺโธ ภควา เวสาลิยํ วิหรติ มหาวเน กูฏาคาร\-สาลายํ. เตน โข ปน สมเยน ฉพฺพคฺคิยา ภิกฺขู สุทฺธ\-กาฬกานํ เอฬกโลมานํ สนฺถตํ การาเปนฺติ. มนุสฺสา วิหารจาริกํ อาหิณฺฑนฺตา ปสฺสิตฺวา อุชฺฌายนฺติ ขิยฺยนฺติ วิปาเจนฺติ “กถํ หิ นาม สมณา สกฺยปุตฺติยา สุทฺธ\-กาฬกานํ เอฬกโลมานํ สนฺถตํ การาเปสฺสนฺติ เสยฺยถาปิ คิหี กามโภคิโน”ติ. อสฺโสสุํ โข ภิกฺขู เตสํ มนุสฺสานํ อุชฺฌายนฺตานํ ขิยฺยนฺตานํ วิปาเจนฺตานํ. เย เต ภิกฺขู อปฺปิจฺฉา ฯเปฯ เต อุชฺฌายนฺติ ขิยนฺติ วิปาเจนฺติ “กถํ หิ นาม ฉพฺพคฺคิยา ภิกฺขู สุทฺธ\-กาฬกานํ เอฬกโลมานํ สนฺถตํ การาเปสฺสนฺตี”ติ. อถ โข เต ภิกฺขู ฉพฺพคฺคิเย ภิกฺขู อเนก\-ปริยาเยน วิครหิตฺวา ภควโต เอตมตฺถํ อาโรเจสุํ ฯเปฯ “สจฺจํ กิร ตุเมฺห ภิกฺขเว สุทฺธ\-กาฬกานํ เอฬกโลมานํ สนฺถตํ การาเปถา”ติ. สจฺจํ ภควาติ. วิครหิ พุทฺโธ ภควา ฯเปฯ. กถํ หิ นาม ตุเมฺห โมฆปุริสา สุทฺธ\-กาฬกานํ เอฬกโลมานํ สนฺถตํ การาเปสฺสถ, เนตํ โมฆปุริสา อปฺปสนฺนานํ วา ปสาทาย ฯเปฯ. เอวญฺจ ปน ภิกฺขเว อิมํ สิกฺขาปทํ อุทฺทิเสยฺยาถ\csromandash{}}

\proseitem{548}{\textbf{“โย ปน ภิกฺขุ สุทฺธ}\-\textbf{กาฬกานํ เอฬกโลมานํ สนฺถตํ การาเปยฺย, นิสฺสคฺคิยํ ปาจิตฺติยนฺ”}ติ.}

\proseitem{549}{\textbf{โย ปนา}ติ โย ยาทิโส ฯเปฯ. \textbf{ภิกฺขู}ติ ฯเปฯ อยํ อิมสฺมิํ อตฺเถ อธิปฺเปโต “ภิกฺขู”ติ.}

\prose{\textbf{กาฬกํ} นาม เทฺว กาฬกานิ ชาติยา กาฬกํ วา รชนกาฬกํ วา.}

\prose{\textbf{สนฺถตํ} นาม สนฺถริตฺวา กตํ โหติ อวายิมํ.}

\csromanmatikaanchor{mtk.167}
\csromanmatikaanchor{mtk.168}
\csromantocmarknopage{subsection}{3. เทฺวภาคสิกฺขาปท}{mtk.167}
\bookmark[dest={mtk.167},level=2]{3. เทฺวภาคสิกฺขาปท}
\csromantocmarkref{subsubsection}{ฉพฺพคฺคิยภิกฺขุวตฺถุ}{mtk.168}
\bookmark[dest={mtk.168},level=3]{ฉพฺพคฺคิยภิกฺขุวตฺถุ}
\prose{\textbf{การาเปยฺยา}ติ กโรติ วา การาเปติ วา, ปโยเค ทุกฺกฏํ. ปฏิลาเภน นิสฺสคฺคิยํ โหติ. นิสฺสชฺชิตพฺพํ สํฆสฺส วา คณสฺส วา \csromanfolio{331}ปุคฺคลสฺส วา. เอวญฺจ ปน ภิกฺขเว นิสฺสชฺชิตพฺพํ ฯเปฯ “อิทํ เม ภนฺเต สุทฺธ\-กาฬกานํ เอฬกโลมานํ สนฺถตํ การาปิตํ นิสฺสคฺคิยํ, อิมาหํ สํฆสฺส นิสฺสชฺชามี”ติ ฯเปฯ ทเทยฺยาติ ฯเปฯ ทเทยฺยุนฺติ ฯเปฯ อายสฺมโต \textbf{ทมฺมีติ.}}

\proseitem{550}{อตฺตนา วิปฺปกตํ อตฺตนา ปริโยสาเปติ, นิสฺสคฺคิยํ ปาจิตฺติยํ. อตฺตนา วิปฺปกตํ ปเรหิ ปริโยสาเปติ, นิสฺสคฺคิยํ ปาจิตฺติยํ. ปเรหิ วิปฺปกตํ อตฺตนา ปริโยสาเปติ, นิสฺสคฺคิยํ ปาจิตฺติยํ. ปเรหิ วิปฺปกตํ ปเรหิ ปริโยสาเปติ, นิสฺสคฺคิยํ ปาจิตฺติยํ.}

\prose{อญฺญสฺสตฺถาย กโรติ วา การาเปติ วา, อาปตฺติ ทุกฺกฏสฺส. อญฺเญน กตํ ปฏิลภิตฺวา ปริภุญฺชติ, อาปตฺติ ทุกฺกฏสฺส.}

\proseitemwithcloserruled{551}{อนาปตฺติ วิตานํ วา ภูมตฺถรณํ วา สาณิปาการํ วา ภิสิํ วา พิพฺโพหนํ วา กโรติ, อุมฺมตฺตกสฺส อาทิกมฺมิกสฺสาติ.}{สุทฺธกาฬก\-สิกฺขาปทํ นิฏฺฐิตํ ทุติยํ.}
\csromanmarkh{2. โกสิยวคฺค}\csromanmarkhii{3. เทฺวภาค\-สิกฺขาปท}\csromanheadernomark{2}{2. \textbf{โกสิยวคฺค 3. เทฺวภาค}\-\textbf{สิกฺขาปท}}
\csromanmatikaanchor{mtk.169}
\csromanmatikaanchor{mtk.170}
\csromantocmarkref{subsubsection}{ปญฺญตฺติ}{mtk.169}
\bookmark[dest={mtk.169},level=3]{ปญฺญตฺติ}
\csromantocmarkref{subsubsection}{สิกฺขาปทวิภงฺค, ปทภาชนีย}{mtk.170}
\bookmark[dest={mtk.170},level=3]{สิกฺขาปทวิภงฺค, ปทภาชนีย}
\proseitem{552}{เตน สมเยน พุทฺโธ ภควา สาวตฺถิยํ วิหรติ เชตวเน อนาถ\-ปิณฺฑิกสฺส อาราเม. เตน โข ปน สมเยน ฉพฺพคฺคิยา ภิกฺขู “ภควตา ปฏิกฺขิตฺตํ สุทฺธ\-กาฬกานํ เอฬกโลมานํ สนฺถตํ การาเปตุนฺ”ติ เต โถกํ เยว โอทาตํ อนฺเต อาทิยิตฺวา ตเถว สุทฺธ\-กาฬกานํ เอฬกโลมานํ สนฺถตํ การาเปนฺติ. เย เต ภิกฺขู อปฺปิจฺฉา, เต อุชฺฌายนฺติ ขิยฺยนฺติ วิปาเจนฺติ “กถํ หิ นาม ฉพฺพคฺคิยา ภิกฺขู โถกํ เยว โอทาตํ อนฺเต อาทิยิตฺวา ตเถว สุทฺธ\-กาฬกานํ เอฬกโลมานํ สนฺถตํ การาเปสฺสนฺตี”ติ. อถ โข เต ภิกฺขู ฉพฺพคฺคิเย ภิกฺขู อเนก\-ปริยาเยน วิครหิตฺวา ภควโต เอตมตฺถํ อาโรเจสุํ ฯเปฯ “สจฺจํ กิร ตุเมฺห ภิกฺขเว โถกํ เยว โอทาตํ อนฺเต อาทิยิตฺวา ตเถว สุทฺธ\-กาฬกานํ เอฬกโลมานํ สนฺถตํ การาเปถา”ติ. สจฺจํ ภควาติ. วิครหิ พุทฺโธ ภควา ฯเปฯ. กถํ หิ นาม ตุเมฺห โมฆปุริสา \csromanfolio{332}โถกํ เยว โอทาตํ อนฺเต อาทิยิตฺวา ตเถว สุทฺธ\-กาฬกานํ เอฬกโลมานํ สนฺถตํ การาเปสฺสถ, เนตํ โมฆปุริสา อปฺปสนฺนานํ วา ปสาทาย ฯเปฯ. เอวญฺจ ปน ภิกฺขเว อิมํ สิกฺขาปทํ อุทฺทิเสยฺยาถ\csromandash{}}

\proseitem{553}{\textbf{“นวํ ปน ภิกฺขุนา สนฺถตํ การยมาเนน เทฺว ภาคา สุทฺธ}\-\textbf{กาฬกานํ เอฬกโลมานํ อาทาตพฺพา ตติยํ โอทาตานํ จตุตฺถํ โคจริยานํ. อนาทา เจ ภิกฺขุ เทฺว ภาเค สุทฺธ}\-\textbf{กาฬกานํ เอฬกโลมานํ ตติยํ โอทาตานํ จตุตฺถํ โคจริยานํ นวํ สนฺถตํ การาเปยฺย, นิสฺสคฺคิยํ ปาจิตฺติยนฺ”}ติ.}

\proseitem{554}{\textbf{นวํ} นาม กรณํ อุปาทาย วุจฺจติ.}

\prose{\textbf{สนฺถตํ} นาม สนฺถริตฺวา กตํ โหติ อวายิมํ.}

\prose{\textbf{การยมาเนนา}ติ กโรนฺโต วา การาเปนฺโต วา.}

\prose{\textbf{เทฺว ภาคา สุทฺธ}\-\textbf{กาฬกานํ เอฬกโลมานํ อาทาตพฺพา}ติ ธารยิตฺวา เทฺว ตุลา อาทาตพฺพา.}

\prose{\textbf{ตติยํ โอทาตานนฺ}ติ ตุลํ โอทาตานํ.}

\prose{\textbf{จตุตฺถํ โคจริยานนฺ}ติ ตุลํ โคจริยานํ.}

\prose{\textbf{อนาทา เจ ภิกฺขุ เทฺว ภาเค สุทฺธ}\-\textbf{กาฬกานํ เอฬกโลมานํ ตติยํ โอทาตานํ จตุตฺถํ โคจริยานนฺ}ติ อนาทิยิตฺวา เทฺว ตุเล สุทฺธ\-กาฬกานํ เอฬกโลมานํ ตุลํ โอทาตานํ ตุลํ โคจริยานํ นวํ สนฺถตํ กโรติ วา การาเปติ วา, ปโยเค ทุกฺกฏํ. ปฏิลาเภน นิสฺสคฺคิยํ โหติ. นิสฺสชฺชิตพฺพํ สํฆสฺส วา คณสฺส วา ปุคฺคลสฺส วา. เอวญฺจ ปน ภิกฺขเว นิสฺสชฺชิตพฺพํ ฯเปฯ “อิทํ เม ภนฺเต สนฺถตํ อนาทิยิตฺวา เทฺว ตุเล สุทฺธ\-กาฬกานํ เอฬกโลมานํ ตุลํ โอทาตานํ ตุลํ โคจริยานํ การาปิตํ นิสฺสคฺคิยํ, อิมาหํ สํฆสฺส นิสฺสชฺชามี”ติ ฯเปฯ ทเทยฺยาติ ฯเปฯ ทเทยฺยุนฺติ ฯเปฯ อายสฺมโต ทมฺมีติ.}

\proseitem{555}{อตฺตนา วิปฺปกตํ อตฺตนา ปริโยสาเปติ, นิสฺสคฺคิยํ ปาจิตฺติยํ. อตฺตนา วิปฺปกตํ ปเรหิ ปริโยสาเปติ, นิสฺสคฺคิยํ ปาจิตฺติยํ. ปเรหิ วิปฺปกตํ อตฺตนา ปริโยสาเปติ, นิสฺสคฺคิยํ ปาจิตฺติยํ. ปเรหิ วิปฺปกตํ ปเรหิ ปริโยสาเปติ, นิสฺสคฺคิยํ ปาจิตฺติยํ.}

\csromanmatikaanchor{mtk.171}
\csromanmatikaanchor{mtk.172}
\csromantocmarknopage{subsection}{4. ฉพฺพสฺสสิกฺขาปท}{mtk.171}
\bookmark[dest={mtk.171},level=2]{4. ฉพฺพสฺสสิกฺขาปท}
\csromantocmarkref{subsubsection}{สาวตฺถิภิกฺขุวตฺถุ}{mtk.172}
\bookmark[dest={mtk.172},level=3]{สาวตฺถิภิกฺขุวตฺถุ}
\prose{\csromanfolio{333}อญฺญสฺสตฺถาย กโรติ วา การาเปติ วา, อาปตฺติ ทุกฺกฏสฺส. อญฺเญน กตํ ปฏิลภิตฺวา ปริภุญฺชติ, อาปตฺติ ทุกฺกฏสฺส.}

\proseitemwithcloserruled{556}{อนาปตฺติ ตุลํ โอทาตานํ ตุลํ โคจริยานํ อาทิยิตฺวา กโรติ, พหุตรํ โอทาตานํ พหุตรํ โคจริยานํ อาทิยิตฺวา กโรติ, สุทฺธํ โอทาตานํ สุทฺธํ โคจริยานํ อาทิยิตฺวา กโรติ, วิตานํ วา ภูมตฺถรณํ วา สาณิปาการํ วา ภิสิํ วา พิพฺโพหนํ วา กโรติ, อุมฺมตฺตกสฺส อาทิกมฺมิกสฺสาติ.}{เทฺวภาค\-สิกฺขาปทํ นิฏฺฐิตํ ตติยํ.}
\csromanmarkh{2. โกสิยวคฺค}\csromanmarkhii{4. ฉพฺพสฺส\-สิกฺขาปท}\csromanheadernomark{2}{2. \textbf{โกสิยวคฺค 4. ฉพฺพสฺส}\-\textbf{สิกฺขาปท}}
\proseitem{557}{เตน สมเยน พุทฺโธ ภควา สาวตฺถิยํ วิหรติ เชตวเน อนาถ\-ปิณฺฑิกสฺส อาราเม. เตน โข ปน สมเยน ภิกฺขู อนุวสฺสํ สนฺถตํ การาเปนฺติ. เต ยาจนพหุลา วิญฺญตฺติพหุลา วิหรนฺติ “เอฬกโลมานิ เทถ, เอฬกโลเมหิ อตฺโถ”ติ. มนุสฺสา อุชฺฌายนฺติ ขิยฺยนฺติ วิปาเจนฺติ “กถํ หิ นาม สมณา สกฺยปุตฺติยา อนุวสฺสํ สนฺถตํ การาเปสฺสนฺติ, ยาจนพหุลา วิญฺญตฺติพหุลา วิหริสฺสนฺติ ‘เอฬกโลมานิ เทถ, เอฬกโลเมหิ อตฺโถ’ติ. อมฺหากํ ปน สกิํ กตานิ สนฺถตานิ ปญฺจปิ ฉปิ วสฺสานิ โหนฺติ, เยสํ โน ทารกา อุหทนฺติปิ อุมฺมิหนฺติปิ อุนฺทูเรหิปิ ขชฺชนฺติ, อิเม ปน สมณา สกฺยปุตฺติยา อนุวสฺสํ สนฺถตํ การาเปนฺติ, ยาจนพหุลา วิญฺญตฺติพหุลา วิหรนฺติ ‘เอฬกโลมานิ เทถ, เอฬกโลเมหิ อตฺโถ”ติ.}

\csromanmatikaanchor{mtk.173}
\csromantocmarkref{subsubsection}{ปฐมปญฺญตฺติ}{mtk.173}
\bookmark[dest={mtk.173},level=3]{ปฐมปญฺญตฺติ}
\prose{อสฺโสสุํ โข ภิกฺขู เตสํ มนุสฺสานํ อุชฺฌายนฺตานํ ขิยฺยนฺตานํ วิปาเจนฺตานํ. เย เต ภิกฺขู อปฺปิจฺฉา ฯเปฯ เต อุชฺฌายนฺติ ขิยฺยนฺติ วิปาเจนฺติ “กถํ หิ นาม ภิกฺขู อนุวสฺสํ สนฺถตํ การาเปสฺสนฺติ, ยาจนพหุลา วิญฺญตฺติพหุลา วิหริสฺสนฺติ เอฬกโลมานิ เทถ, เอฬกโลเมหิ อตฺโถ”ติ. อถ โข เต ภิกฺขู เต อเนก\-ปริยาเยน วิครหิตฺวา ภควโต เอตมตฺถํ อาโรเจสุํ ฯเปฯ “สจฺจํ กิร ภิกฺขเว ภิกฺขู อนุวสฺสํ สนฺถตํ การาเปนฺติ, ยาจนพหุลา วิญฺญตฺติพหุลา วิหรนฺติ เอฬกโลมานิ เทถ, เอฬกโลเมหิ อตฺโถ”ติ. สจฺจํ ภควาติ. วิครหิ \csromanfolio{334}พุทฺโธ ภควา ฯเปฯ. กถํ หิ นาม เต ภิกฺขเว โมฆปุริสา อนุวสฺสํ สนฺถตํ การาเปสฺสนฺติ, ยาจนพหุลา วิญฺญตฺติพหุลา วิหริสฺสนฺติ “เอฬกโลมานิ เทถ, เอฬกโลเมหิ อตฺโถ”ติ, เนตํ ภิกฺขเว อปฺปสนฺนานํ วา ปสาทาย ฯเปฯ. เอวญฺจ ปน ภิกฺขเว อิมํ สิกฺขาปทํ อุทฺทิเสยฺยาถ\csromandash{}}

\proseitem{558}{\textbf{“นวํ ปน ภิกฺขุนา สนฺถตํ การาเปตฺวา ฉพฺพสฺสานิ ธาเรตพฺพํ, โอเรน เจ ฉนฺนํ วสฺสานํ ตํ สนฺถตํ วิสฺสชฺเชตฺวา วา อวิสฺสชฺเชตฺวา วา อญฺญํ นวํ สนฺถตํ การาเปยฺย, นิสฺสคฺคิยํ ปาจิตฺติยนฺ”}ติ.}

\prose{เอวญฺจิทํ ภควตา ภิกฺขูนํ สิกฺขาปทํ ปญฺญตฺตํ โหติ.}

\proseitem{559}{เตน โข ปน สมเยน อญฺญตโร ภิกฺขุ โกสมฺพิยํ คิลาโน โหติ. ญาตกา ตสฺส ภิกฺขุโน สนฺติเก ทูตํ ปาเหสุํ “อาคจฺฉตุ ภทนฺโต มยํ อุปฏฺฐหิสฺสามา”ติ. ภิกฺขูปิ เอวมาหํสุ “คจฺฉาวุโส ญาตกา ตํ อุปฏฺฐหิสฺสนฺตี”ติ. โส เอวมาห “ภควตา อาวุโส สิกฺขาปทํ ปญฺญตฺตํ ‘นวํ ปน ภิกฺขุนา สนฺถตํ การาเปตฺวา ฉพฺพสฺสานิ ธาเรตพฺพนฺ’ติ, อหญฺจมฺหิ คิลาโน, น สกฺโกมิ สนฺถตํ อาทาย ปกฺกมิตุํ, มยฺหญฺจ วินา สนฺถตา น ผาสุ โหติ, นาหํ คมิสฺสามี”ติ. ภควโต เอตมตฺถํ อาโรเจสุํ. อถ โข ภควา เอตสฺมิํ นิทาเน เอตสฺมิํ ปกรเณ ธมฺมิํ กถํ กตฺวา ภิกฺขู อามนฺเตสิ อนุชานามิ ภิกฺขเว คิลานสฺส ภิกฺขุโน สนฺถต\-สมฺมุติํ ทาตุํ, เอวญฺจ ปน ภิกฺขเว ทาตพฺพา\csromandash{}เตน คิลาเนน ภิกฺขุนา สํฆํ อุปสงฺกมิตฺวา เอกํสํ อุตฺตราสงฺคํ กริตฺวา วุฑฺฒานํ ภิกฺขูนํ ปาเท วนฺทิตฺวา อุกฺกุฏิกํ นิสีทิตฺวา อญฺชลิํ ปคฺคเหตฺวา เอวมสฺส วจนีโย “อหํ ภนฺเต คิลาโน น สกฺโกมิ สนฺถตํ อาทาย ปกฺกมิตุํ, โสหํ ภนฺเต สํฆํ สนฺถต\-สมฺมุติํ ยาจามี”ติ. ทุติยมฺปิ ยาจิตพฺพา. ตติยมฺปิ ยาจิตพฺพา. พฺยตฺเตน ภิกฺขุนา ปฏิพเลน สํโฆ ญาเปตพฺโพ\csromandash{}}

\proseitem{560}{“สุณาตุ เม ภนฺเต สํโฆ, อยํ อิตฺถนฺนาโม ภิกฺขุ คิลาโน น สกฺโกติ สนฺถตํ อาทาย ปกฺกมิตุํ, โส สํฆํ สนฺถต\-สมฺมุติํ ยาจติ, ยทิ สํฆสฺส ปตฺตกลฺลํ, สํโฆ อิตฺถนฺนามสฺส ภิกฺขุโน สนฺถต\-สมฺมุติํ ทเทยฺย, เอสา ญตฺติ.}

\csromanmatikaanchor{mtk.174}
\csromanmatikaanchor{mtk.175}
\csromantocmarkref{subsubsection}{อนุปญฺญตฺติ}{mtk.174}
\bookmark[dest={mtk.174},level=3]{อนุปญฺญตฺติ}
\csromantocmarkref{subsubsection}{สิกฺขาปทวิภงฺค, ปทภาชนีย}{mtk.175}
\bookmark[dest={mtk.175},level=3]{สิกฺขาปทวิภงฺค, ปทภาชนีย}
\prose{\csromanfolio{335}สุณาตุ เม ภนฺเต สํโฆ, อยํ อิตฺถนฺนาโม ภิกฺขุ คิลาโน น สกฺโกติ สนฺถตํ อาทาย ปกฺกมิตุํ, โส สํฆํ สนฺถต\-สมฺมุติํ ยาจติ, สํโฆ อิตฺถนฺนามสฺส ภิกฺขุโน สนฺถต\-สมฺมุติํ เทติ, ยสฺสายสฺมโต ขมติ อิตฺถนฺนามสฺส ภิกฺขุโน สนฺถต\-สมฺมุติยา ทานํ, โส ตุณฺหสฺส, ยสฺส นกฺขมติ, โส ภาเสยฺย.}

\prose{ทินฺนา สํเฆน อิตฺถนฺนามสฺส ภิกฺขุโน สนฺถต\-สมฺมุติ, ขมติ สํฆสฺส, ตสฺมา ตุณฺหี, เอวเมตํ ธารยามี”ติ.}

\prose{เอวญฺจ ปน ภิกฺขเว อิมํ สิกฺขาปทํ อุทฺทิเสยฺยาถ\csromandash{}}

\proseitem{561}{\textbf{“นวํ ปน ภิกฺขุนา สนฺถตํ การาเปตฺวา ฉพฺพสฺสานิ ธาเรตพฺพํ, โอเรน เจ ฉนฺนํ วสฺสานํ ตํ สนฺถตํ วิสฺสชฺเชตฺวา วา อวิสฺสชฺเชตฺวา วา อญฺญํ นวํ สนฺถตํ การาเปยฺย อญฺญตฺร ภิกฺขุ}\-\textbf{สมฺมุติยา, นิสฺสคฺคิยํ ปาจิตฺติยนฺ”}ติ.}

\proseitem{562}{\textbf{นวํ} นาม กรณํ อุปาทาย วุจฺจติ.}

\prose{\textbf{สนฺถตํ} นาม สนฺถริตฺวา กตํ โหติ อวายิมํ.}

\prose{\textbf{การาเปตฺวา}ติ กริตฺวา วา การาเปตฺวา วา.}

\prose{\textbf{ฉพฺพสฺสานิ ธาเรตพฺพนฺ}ติ ฉพฺพสฺส\-ปรมตา ธาเรตพฺพํ.}

\prose{\textbf{โอเรน เจ ฉนฺนํ วสฺสานนฺ}ติ อูนก\-ฉพฺพสฺสานิ.}

\prose{\textbf{ตํ สนฺถตํ วิสฺสชฺเชตฺวา}ติ อญฺเญสํ ทตฺวา.}

\prose{\textbf{อวิสฺสชฺเชตฺวา}ติ น กสฺสจิ ทตฺวา.}

\prose{\textbf{อญฺญตฺร ภิกฺขุ}\-\textbf{สมฺมุติยา}ติ ฐเปตฺวา ภิกฺขุ\-สมฺมุติํ, อญฺญํ นวํ สนฺถตํ กโรติ วา การาเปติ วา, ปโยเค ทุกฺกฏํ. ปฏิลาเภน นิสฺสคฺคิยํ โหติ. นิสฺสชฺชิตพฺพํ สํฆสฺส วา คณสฺส วา ปุคฺคลสฺส วา. เอวญฺจ ปน ภิกฺขเว นิสฺสชฺชิตพฺพํ ฯเปฯ “อิทํ เม ภนฺเต สนฺถตํ อูนก\-ฉพฺพสฺสานิ การาปิตํ อญฺญตฺร ภิกฺขุ\-สมฺมุติยา นิสฺสคฺคิยํ, อิมาหํ สํฆสฺส นิสฺสชฺชามี”ติ ฯเปฯ ทเทยฺยาติ ฯเปฯ ทเทยฺยุนฺติ ฯเปฯ อายสฺมโต ทมฺมีติ.}

\csromanmatikaanchor{mtk.176}
\csromanmatikaanchor{mtk.177}
\csromantocmarknopage{subsection}{5. นิสีทนสนฺถตสิกฺขาปท}{mtk.176}
\bookmark[dest={mtk.176},level=2]{5. นิสีทนสนฺถตสิกฺขาปท}
\csromantocmarkref{subsubsection}{อุปเสนภิกฺขุวตฺถุ}{mtk.177}
\bookmark[dest={mtk.177},level=3]{อุปเสนภิกฺขุวตฺถุ}
\proseitem{563}{อตฺตนา วิปฺปกตํ อตฺตนา ปริโยสาเปติ, นิสฺสคฺคิยํ ปาจิตฺติยํ. อตฺตนา วิปฺปกตํ ปเรหิ ปริโยสาเปติ, นิสฺสคฺคิยํ \csromanfolio{336}ปาจิตฺติยํ. ปเรหิ วิปฺปกตํ อตฺตนา ปริโยสาเปติ, นิสฺสคฺคิยํ ปาจิตฺติยํ. ปเรหิ วิปฺปกตํ ปเรหิ ปริโยสาเปติ, นิสฺสคฺคิยํ ปาจิตฺติยํ.}

\proseitemwithcloserruled{564}{อนาปตฺติ ฉพฺพสฺสานิ กโรติ, อติเรก\-ฉพฺพสฺสานิ กโรติ, อญฺญสฺสตฺถาย กโรติ วา การาเปติ วา, อญฺเญน กตํ ปฏิลภิตฺวา ปริภุญฺชติ, วิตานํ วา ภูมตฺถรณํ วา สาณิปาการํ วา ภิสิํ วา พิพฺโพหนํ วา กโรติ, ภิกฺขุ\-สมฺมุติยา, อุมฺมตฺตกสฺส อาทิกมฺมิกสฺสาติ.}{ฉพฺพสฺส\-สิกฺขาปทํ นิฏฺฐิตํ จตุตฺถํ.}
\csromanmarkh{2. โกสิยวคฺค}\csromanmarkhii{5. นิสีทนสนฺถต\-สิกฺขาปท}\csromanheadernomark{2}{2. \textbf{โกสิยวคฺค 5. นิสีทนสนฺถต}\-\textbf{สิกฺขาปท}}
\proseitem{565}{เตน สมเยน พุทฺโธ ภควา สาวตฺถิยํ วิหรติ เชตวเน อนาถ\-ปิณฺฑิกสฺส อาราเม. อถ โข ภควา ภิกฺขู อามนฺเตสิ “อิจฺฉามหํ ภิกฺขเว เตมาสํ ปฏิสลฺลียิตุํ, นมฺหิ เกนจิ อุปสงฺกมิตพฺโพ อญฺญตฺร เอเกน ปิณฺฑปาต\-นีหารเกนา”ติ. เอวํ ภนฺเตติ โข เต ภิกฺขู ภควโต ปฏิสฺสุณิตฺวา นาสฺสุธ โกจิ ภควนฺตํ อุปสงฺกมติ อญฺญตฺร เอเกน ปิณฺฑปาต\-นีหารเกน. เตน โข ปน สมเยน สาวตฺถิยา สํเฆน กติกา กตา โหติ “อิจฺฉาตาวุโส ภควา เตมาสํ ปฏิสฺสลฺลียิตุํ, น ภควา เกนจิ อุปสงฺกมิตพฺโพ อญฺญตฺร เอเกน ปิณฺฑปาต\-นีหารเกน, โย ภควนฺตํ อุปสงฺกมติ, โส ปาจิตฺติยํ เทสาเปตพฺโพ”ติ. อถ โข อายสฺมา อุปเสโน วงฺคนฺตปุตฺโต สปริโส เยน ภควา เตนุปสงฺกมิ, อุปสงฺกมิตฺวา ภควนฺตํ อภิวาเทตฺวา เอกมนฺตํ นิสีทิ. อาจิณฺณํ โข ปเนตํ พุทฺธานํ ภควนฺตานํ อาคนฺตุเกหิ ภิกฺขูหิ สทฺธิํ ปฏิสมฺโมทิตุํ. อถ โข ภควา อายสฺมนฺตํ อุปเสนํ วงฺคนฺตปุตฺตํ เอตทโวจ “กจฺจิ โว อุปเสน ขมนียํ, กจฺจิ ยาปนียํ, กจฺจิตฺถ อปฺป\-กิลมเถน อทฺธานํ อาคตา”ติ. ขมนียํ ภควา, ยาปนียํ ภควา, อปฺป\-กิลมเถน จ มยํ ภนฺเต อทฺธานํ อาคตาติ.}

\prose{เตน โข ปน สมเยน อายสฺมโต อุปเสนสฺส วงฺคนฺต\-ปุตฺตสฺส สทฺธิวิหาริโก ภิกฺขุ ภควโต อวิทูเร นิสินฺโน โหติ. อถ โข \csromanfolio{337}ภควา ตํ ภิกฺขุํ เอตทโวจ “มนาปานิ เต ภิกฺขุ ปํสุกูลานี”ติ. น โข เม ภนฺเต มนาปานิ ปํสุกูลานีติ. กิสฺส ปน ตฺวํ ภิกฺขุ ปํสุกูลิโกติ. อุปชฺฌาโย เม ภนฺเต ปํสุกูลิโก, เอวํ อหํปิ ปํสุกูลิโกติ. อถ โข ภควา อายสฺมนฺตํ อุปเสนํ วงฺคนฺตปุตฺตํ เอตทโวจ “ปาสาทิกา โข ตฺยายํ อุปเสน ปริสา, กถํ ตฺวํ อุปเสน ปริสํ วิเนสี”ติ. โย มํ ภนฺเต อุปสมฺปทํ ยาจติ, ตมหํ\csromansharedfootnote{8deef69c5a52540f}{ตาหํ (ก)} เอวํ วทามิ “อหํ โข อาวุโส อารญฺญิโก ปิณฺฑปาติโก ปํสุกูลิโก, สเจ ตฺวํปิ อารญฺญิโก ภวิสฺสสิ ปิณฺฑปาติโก ปํสุกูลิโก, เอวาหํ ตํ อุปสมฺปาเทสฺสามี”ติ. สเจ เม ปฏิสฺสุณาติ อุปสมฺปาเทมิ, โน เจ เม ปฏิสฺสุณาติ น อุปสมฺปาเทมิ. โย มํ นิสฺสยํ ยาจติ, ตมหํ\csromansharedfootnote{8deef69c5a52540f}{ตาหํ (ก)} เอวํ วทามิ “อหํ โข อาวุโส อารญฺญิโก ปิณฺฑปาติโก ปํสุกูลิโก, สเจ ตฺวํปิ อารญฺญิโก ภวิสฺสสิ ปิณฺฑปาติโก ปํสุกูลิโก, เอวาหํ เต นิสฺสยํ ทสฺสามี”ติ. สเจ เม ปฏิสฺสุณาติ นิสฺสยํ เทมิ, โน เจ เม ปฏิสฺสุณาติ น นิสฺสยํ เทมิ, เอวํ โข อหํ ภนฺเต ปริสํ วิเนมีติ.}

\prose{สาธุ สาธุ อุปเสน, สาธุ โข ตฺวํ อุปเสน ปริสํ วิเนสิ, ชานาสิ ปน ตฺวํ อุปเสน สาวตฺถิยา สํฆสฺส กติกนฺติ. น โข อหํ ภนฺเต ชานามิ สาวตฺถิยา สํฆสฺส กติกนฺติ. สาวตฺถิยา โข อุปเสน สํเฆน กติกา กตา “อิจฺฉาตาวุโส ภควา เตมาสํ ปฏิสลฺลียิตุํ, น ภควา เกนจิ อุปสงฺกมิตพฺโพ อญฺญตฺร เอเกน ปิณฺฑปาต\-นีหารเกน, โย ภควนฺตํ อุปสงฺกมติ, โส ปาจิตฺติยํ เทสาเปตพฺโพ”ติ. ปญฺญายิสฺสติ ภนฺเต สาวตฺถิยา สํโฆ สกาย กติกาย, น มยํ อปญฺญตฺตํ ปญฺญเปสฺสาม, ปญฺญตฺตํ วา น สมุจฺฉินฺทิสฺสาม, ยถา\-ปญฺญตฺเตสุ สิกฺขาปเทสุ สมาทาย วตฺติสฺสามาติ. สาธุ สาธุ อุปเสน, อปญฺญตฺตํ น ปญฺญเปตพฺพํ, ปญฺญตฺตํ วา น สมุจฺฉินฺทิตพฺพํ, ยถา\-ปญฺญตฺเตสุ สิกฺขาปเทสุ สมาทาย วตฺติตพฺพํ. อนุชานามิ อุปเสน “เย เต ภิกฺขู อารญฺญิกา ปิณฺฑปาติกา ปํสุกูลิกา ยถาสุขํ มํ ทสฺสนาย อุปสงฺกมนฺตู”ติ.}

\proseitem{566}{เตน โข ปน สมเยน สมฺพหุลา ภิกฺขู พหิทฺวาร\-โกฏฺฐเก ฐิตา โหนฺติ “มยํ อายสฺมนฺตํ อุปเสนํ วงฺคนฺตปุตฺตํ ปาจิตฺติยํ \csromanfolio{338}\textbf{เทสาเปสฺสามา”ติ. อถ โข อายสฺมา อุปเสโน วงฺคนฺตปุตฺโต สปริโส} \textbf{อุฏฺฐายาสนา ภควนฺตํ อภิวาเทตฺวา ปทกฺขิณํ กตฺวา ปกฺกามิ. อถ} \textbf{โข เต ภิกฺขู อายสฺมนฺตํ อุปเสนํ วงฺคนฺตปุตฺตํ เอตทโวจุํ “ชานาสิ} \textbf{ตฺวํ อาวุโส อุปเสน สาวตฺถิยา สํฆสฺส กติกนฺ”ติ. ภควาปิ มํ} อาวุโส เอวมาห “ชานาสิ ปน ตฺวํ อุปเสน สาวตฺถิยา สํฆสฺส กติกนฺ”ติ. น โข อหํ ภนฺเต ชานามิ สาวตฺถิยา สํฆสฺส กติกนฺติ. สาวตฺถิยา โข อุปเสน สํเฆน กติกา กตา “อิจฺฉตาวุโส ภควา เตมาสํ ปฏิสลฺลียิตุํ, น ภควา เกนจิ อุปสงฺกมิตพฺโพ อญฺญตฺร เอเกน ปิณฺฑปาต\-นีหารเกน, โย ภควนฺตํ อุปสงฺกมติ, โส ปาจิตฺติยํ \textbf{เทสาเปตพฺโพ”ติ. ปญฺญยิสฺสติ ภนฺเต สาวตฺถิยา สํโฆ สกาย กติกาย, น} \textbf{มยํ อปญฺญตฺตํ ปญฺญเปสฺสาม, ปญฺญตฺตํ วา น สมุจฺฉินฺทิสฺสาม,} ยถา\-ปญฺญตฺเตสุ สิกฺขาปเทสุ สมาทาย วตฺติสฺสามาติ. อนุญฺญาตาวุโส \textbf{ภควตา “เย เต ภิกฺขู อารญฺญิกา ปิณฺฑปาติกา ปํสุกูลิกา ยถาสุขํ} \textbf{มํ ทสฺสนาย อุปสงฺกมนฺตู”ติ.}}

\prose{อถ โข เต ภิกฺขู สจฺจํ โข อายสฺมา อุปเสโน อาห “น อปญฺญตฺตํ ปญฺญเปตพฺพํ, ปญฺญตฺตํ วา น สมุจฺฉินฺทิตพฺพํ, ยถา\-ปญฺญตฺเตสุ สิกฺขาปเทสุ สมาทาย วตฺติตพฺพนฺ”ติ. อสฺโสสุํ โข ภิกฺขู อนุญฺญาตา กิร ภควตา “เย เต ภิกฺขู อารญฺญิกา ปิณฺฑปาติกา ปํสุกูลิกา ยถาสุขํ มํ ทสฺสนาย อุปสงฺกมนฺตู”ติ. เต ภควนฺตํ ทสฺสนํ ปิเหนฺตา\csromansharedfootnote{88a0c800db9d6a63}{ทสฺสนาย ปิหยมานา (สฺยา)} สนฺถตานิ อุชฺฌิตฺวา อารญฺญิกงฺคํ ปิณฺฑปาติ\-กงฺคํ ปํสุกูลิ\-กงฺคํ สมาทิยิํสุ. อถ โข ภควา สมฺพหุเลหิ ภิกฺขูหิ สทฺธิํ เสนาสนจาริกํ อาหิณฺฑนฺโต อทฺทส สนฺถตานิ ตหํ ตหํ อุชฺฌิตานิ, ปสฺสิตฺวา ภิกฺขู อามนฺเตสิ “กสฺสิมานิ ภิกฺขเว สนฺถตานิ ตหํ ตหํ อุชฺฌิตานี”ติ. อถ โข เต ภิกฺขู ภควโต เอตมตฺถํ อาโรเจสุํ. อถ โข ภควา เอตสฺมิํ นิทาเน เอตสฺมิํ ปกรเณ ธมฺมิํ กถํ กตฺวา ภิกฺขู อามนฺเตสิ เตน หิ ภิกฺขเว ภิกฺขูนํ สิกฺขาปทํ ปญฺญเปสฺสามิ ทส อตฺถวเส ปฏิจฺจ สํฆสุฏฺฐุตาย สํฆผาสุตาย ฯเปฯ. เอวญฺจ ปน ภิกฺขเว อิมํ สิกฺขาปทํ อุทฺทิเสยฺยาถ\csromandash{}}

\csromanmatikaanchor{mtk.178}
\csromanmatikaanchor{mtk.179}
\csromantocmarkref{subsubsection}{ปญฺญตฺติ}{mtk.178}
\bookmark[dest={mtk.178},level=3]{ปญฺญตฺติ}
\csromantocmarkref{subsubsection}{สิกฺขาปทวิภงฺค, ปทภาชนีย}{mtk.179}
\bookmark[dest={mtk.179},level=3]{สิกฺขาปทวิภงฺค, ปทภาชนีย}
\proseitem{567}{\csromanfolio{339}\textbf{“นิสีทน}\-\textbf{สนฺถตํ ปน ภิกฺขุนา การยมาเนน ปุราณ}\-\textbf{สนฺถตสฺส สามนฺตา สุคตวิทตฺถิ อาทาตพฺพา ทุพฺพณฺณ}\-\textbf{กรณาย, อนาทา เจ ภิกฺขุ ปุราณ}\-\textbf{สนฺถตสฺส สามนฺตา สุคต}\-\textbf{วิทตฺถิํ นวํ นิสีทน}\-\textbf{สนฺถตํ การาเปยฺย, นิสฺสคฺคิยํ ปาจิตฺติยนฺ”}ติ.}

\proseitem{568}{\textbf{นิสีทนํ} นาม สทสํ วุจฺจติ.}

\prose{\textbf{สนฺถตํ} นาม สนฺถริตฺวา กตํ โหติ อวายิมํ.}

\prose{\textbf{การยมาเนนา}ติ กโรนฺโต วา การาเปนฺโต วา.}

\prose{\textbf{ปุราณ}\-\textbf{สนฺถตํ} นาม สกิํ นิวตฺถํปิ สกิํ ปารุตํปิ.}

\prose{\textbf{สามนฺตา สุคตวิทตฺถิ อาทาตพฺพา ทุพฺพณฺณกรณายา}ติ ถิรภาวาย วฏฺฏํ วา จตุรสฺสํ วา ฉินฺทิตฺวา เอกเทเส วา สนฺถริตพฺพํ, วิชเฏตฺวา วา สนฺถริตพฺพํ.}

\prose{\textbf{อนาทา เจ ภิกฺขุ ปุราณ}\-\textbf{สนฺถตสฺส สามนฺตา สุคต}\-\textbf{วิทตฺถินฺ}ติ อนาทิยิตฺวา ปุราณ\-สนฺถตสฺส สมนฺตา สุคต\-วิทตฺถิํ. นวํ นิสีทน\-สนฺถตํ กโรติ วา การาเปติ วา, ปโยเค ทุกฺกฏํ. ปฏิลาเภน นิสฺสคฺคิยํ โหติ. นิสฺสชฺชิตพฺพํ สํฆสฺส วา คณสฺส วา ปุคฺคลสฺส วา. เอวญฺจ ปน ภิกฺขเว นิสฺสชฺชิตพฺพํ ฯเปฯ “อิทํ เม ภนฺเต นิสีทน\-สนฺถตํ อนาทิยิตฺวา ปุราณ\-สนฺถตสฺส สามนฺตา สุคต\-วิทตฺถิํ การาปิตํ นิสฺสคฺคิยํ, อิมาหํ สํฆสฺส นิสฺสชฺชามี”ติ ฯเปฯ ทเทยฺยาติ ฯเปฯ ทเทยฺยุนฺติ ฯเปฯ อายสฺมโต ทมฺมีติ.}

\proseitem{569}{อตฺตานา วิปฺปกตํ อตฺตนา ปริโยสาเปติ, นิสฺสคฺคิยํ ปาจิตฺติยํ. อตฺตนา วิปฺปกตํ ปเรหิ ปริโยสาเปติ, นิสฺสคฺคิยํ ปาจิตฺติยํ. ปเรหิ วิปฺปกตํ อตฺตานา ปริโยสาเปติ, นิสฺสคฺคิยํ ปาจิตฺติยํ. ปเรหิ วิปฺปกตํ ปเรหิ ปริโยสาเปติ, นิสฺสคฺคิยํ ปาจิตฺติยํ.}

\prose{อญฺญสฺสตฺถาย กโรติ วา การาเปติ วา, อาปตฺติ ทุกฺกฏสฺส.}

\proseitem{570}{อนาปตฺติ ปุราณ\-สนฺถตสฺส สามนฺตา สุคต\-วิทตฺถิํ อาทิยิตฺวา กโรติ, อลภนฺโต โถกตรํ อาทิยิตฺวา กโรติ, อลภนฺโต อนาทิยิตฺวา กโรติ, อญฺเญน กตํ ปฏิลภิตฺวา ปริภุญฺชติ, วิตานํ วา ภูมตฺถรณํ วา สาณิปาการํ วา ภิสิํ วา พิพฺโพหนํ วา กโรติ, อุมฺมตฺตกสฺส อาทิกมฺมิกสฺสาติ.}

\vspace{\tipitakaparskip}
\csromancenter{นิสีทนสนฺถต\-สิกฺขาปทํ นิฏฺฐิตํ ปญฺจมํ.}
\csromanmatikaanchor{mtk.180}
\csromantocmarknopage{subsection}{6. เอฬกโลมสิกฺขาปท}{mtk.180}
\bookmark[dest={mtk.180},level=2]{6. เอฬกโลมสิกฺขาปท}
\csromanmarkh{2. โกสิยวคฺค}\csromanmarkhii{6. เอฬกโลม\-สิกฺขาปท}\csromanheadernomark{2}{2. \textbf{โกสิยวคฺค 6. เอฬกโลม}\-\textbf{สิกฺขาปท}}
\csromanmatikaanchor{mtk.181}
\csromanmatikaanchor{mtk.182}
\csromanmatikaanchor{mtk.183}
\csromantocmarkref{subsubsection}{อญฺญตรภิกฺขุวตฺถุ}{mtk.181}
\bookmark[dest={mtk.181},level=3]{อญฺญตรภิกฺขุวตฺถุ}
\csromantocmarkref{subsubsection}{ปญฺญตฺติ}{mtk.182}
\bookmark[dest={mtk.182},level=3]{ปญฺญตฺติ}
\csromantocmarkref{subsubsection}{สิกฺขาปทวิภงฺค, ปทภาชนีย}{mtk.183}
\bookmark[dest={mtk.183},level=3]{สิกฺขาปทวิภงฺค, ปทภาชนีย}
\proseitem{571}{\csromanfolio{340}เตน สมเยน พุทฺโธ ภควา สาวตฺถิยํ วิหรติ เชตวเน อนาถ\-ปิณฺฑิกสฺส อาราเม. เตน โข ปน สมเยน อญฺญตรสฺส ภิกฺขุโน โกสเลสุ ชนปเท สาวตฺถิํ คจฺฉนฺตสฺส อนฺตรามคฺเค เอฬกโลมานิ อุปฺปชฺชิํสุ. อถ โข โส ภิกฺขู ตานิ เอฬกโลมานิ อุตฺตราสงฺเคน ภณฺฑิกํ พนฺธิตฺวา อคมาสิ. มนุสฺสา ตํ ภิกฺขุํ ปสฺสิตฺวา อุปฺปณฺเฑสุํ “กิตฺตเกน เต ภนฺเต กีตานิ, กิตฺตโก อุทโย ภวิสฺสตี”ติ. โส ภิกฺขุ เตหิ มนุสฺเสหิ อุปฺปณฺฑิยมาโน มงฺกุ อโหสิ. อถ โข โส ภิกฺขุ สาวตฺถิํ คนฺตฺวา ตานิ เอฬกโลมานิ ฐิตโกว อาสุมฺภิ. ภิกฺขู ตํ ภิกฺขุํ เอตทโวจุํ “กิสฺส ตฺวํ อาวุโส อิมานิ เอฬกโลมานิ ฐิตโกว อาสุมฺภสี”ติ. ตถา หิ ปนาหํ อาวุโส อิเมสํ เอฬกโลมานํ การณา มนุสฺเสหิ อุปฺปณฺฑิโตติ. กีว ทูรโต ปน ตฺวํ อาวุโส อิมานิ เอฬกโลมานิ อาหรีติ. อติเรก\-ติโยชนํ อาวุโสติ. เย เต ภิกฺขู อปฺปิจฺฉา ฯเปฯ เต อุชฺฌายนฺติ ขิยฺยนฺติ วิปาเจนฺติ “กถํ หิ นาม ภิกฺขุ อติเรก\-ติโยชนํ เอฬกโลมานิ อาหริสฺสตี”ติ. อถ โข เต ภิกฺขู ตํ ภิกฺขุํ อเนก\-ปริยาเยน วิครหิตฺวา ภควโต เอตมตฺถํ อาโรเจสุํ ฯเปฯ “สจฺจํ กิร ตฺวํ ภิกฺขุ อติเรก\-ติโยชนํ เอฬกโลมานิ อาหรี”ติ. สจฺจํ ภควาติ. วิครหิ พุทฺโธ ภควา ฯเปฯ. กถํ หิ นาม ตฺวํ โมฆปุริส อติเรก\-ติโยชนํ เอฬกโลมานิ อาหริสฺสสิ, เนตํ โมฆปุริส อปฺปสนฺนานํ วา ปสาทาย ฯเปฯ. เอวญฺจ ปน ภิกฺขเว อิมํ สิกฺขาปทํ อุทฺทิเสยฺยาถ\csromandash{}}

\proseitem{572}{\textbf{“ภิกฺขุโน ปเนว อทฺธาน}\-\textbf{มคฺค}\-\textbf{ปฺปฏิปนฺนสฺส เอฬกโลมานิ อุปฺปชฺเชยฺยุํ, อากงฺขมาเนน ภิกฺขุนา ปฏิคฺคเหตพฺพานิ, ปฏิคฺคเหตฺวา ติโยชนปรมํ สหตฺถา หริตพฺพานิ}\csromansharedfootnote{8afd19e8c4fe01d4}{หาเรตพฺพานิ (สี, สฺยา, ก)} \textbf{อสนฺเต หารเก, ตโต เจ อุตฺตริ หเรยฺย อสนฺเตปิ หารเก, นิสฺสคฺคิยํ ปาจิตฺติยนฺ”}ติ.}

\proseitem{573}{\textbf{ภิกฺขุโน ปเนว อทฺธาน}\-\textbf{มคฺค}\-\textbf{ปฺปฏิปนฺนสฺสา}ติ ปนฺถํ คจฺฉนฺตสฺส.}

\prose{\textbf{เอฬกโลมานิ อุปฺปชฺเชยฺยุนฺ}ติ อุปฺปชฺเชยฺยุํ สํฆโต วา คณโต วา ญาติโต วา มิตฺตโต วา ปํสุกูลํ วา อตฺตโน วา ธเนน.}

\prose{\csromanfolio{341}\textbf{อากงฺข}\-\textbf{มาเนนา}ติ อิจฺฉมาเนน ปฏิคฺคเหตพฺพานิ.}

\prose{\textbf{ปฏิคฺคเหตฺวา ติโยชนปรมํ สหตฺถา หริตพฺพานี}ติ ติโยชน\-ปรมตา สหตฺถา หริตพฺพานิ.}

\prose{\textbf{อสนฺเต หารเก}ติ นาญฺโญ โกจิ หารโก โหติ อิตฺถี วา ปุริโส วา คหฏฺโฐ วา ปพฺพชิโต วา. \textbf{ตโต เจ อุตฺตริ หเรยฺย อสนฺเตปิ หารเก}ติ ปฐมํ ปาทํ ติโยชนํ อติกฺกาเมติ, อาปตฺติ ทุกฺกฏสฺส. ทุติยํ ปาทํ อติกฺกาเมติ, นิสฺสคฺคิยํ ปาจิตฺติยํ\csromansharedfootnote{a6ce1fdcc94e5017}{นิสฺสคฺคิยานิ โหนฺติ (สฺยา)}. อนฺโตติโยชเน ฐิโต พหิติโยชนํ ปาเตติ, นิสฺสคฺคิยํ ปาจิตฺติยํ\csromansharedfootnote{a6ce1fdcc94e5017}{นิสฺสคฺคิยานิ โหนฺติ (สฺยา)}. อญฺญสฺส ยาเน วา ภณฺเฑ วา อชานนฺตสฺส ปกฺขิปิตฺวา ติโยชนํ อติกฺกาเมติ, นิสฺสคฺคิยานิ โหนฺติ. นิสฺสชฺชิตพฺพานิ สํฆสฺส วา คณสฺส วา ปุคฺคลสฺส วา. เอวญฺจ ปน ภิกฺขเว นิสฺสชฺชิตพฺพานิ ฯเปฯ “อิมานิ เม ภนฺเต เอฬกโลมานิ ติโยชนํ อติกฺกามิตานิ นิสฺสคฺคิยานิ, อิมานาหํ สํฆสฺส นิสฺสชฺชามี”ติ ฯเปฯ ทเทยฺยาติ ฯเปฯ ทเทยฺยุนฺติ ฯเปฯ อายสฺมโต ทมฺมีติ.}

\proseitem{574}{อติเรก\-ติโยชเน อติเรกสญฺญี อติกฺกาเมติ\csromansharedfootnote{e6feb8c186cd9736}{ติโยชนํ อติกฺกาเมติ (สฺยา)}, นิสฺสคฺคิยํ ปาจิตฺติยํ. อติเรก\-ติโยชเน เวมติโก อติกฺกาเมติ\csromansharedfootnote{e6feb8c186cd9736}{ติโยชนํ อติกฺกาเมติ (สฺยา)}, นิสฺสคฺคิยํ ปาจิตฺติยํ. อติเรก\-ติโยชเน อูนกสญฺญี อติกฺกาเมติ\csromansharedfootnote{e6feb8c186cd9736}{ติโยชนํ อติกฺกาเมติ (สฺยา)}, นิสฺสคฺคิยํ ปาจิตฺติยํ.}

\prose{อูนกติโยชเน อติเรกสญฺญี, อาปตฺติ ทุกฺกฏสฺส. อูนกติโยชเน เวมติโก, อาปตฺติ ทุกฺกฏสฺส. อูนกติโยชเน อูนกสญฺญี, อนาปตฺติ.}

\proseitem{575}{อนาปตฺติ ติโยชนํ หรติ, อูนกติโยชนํ หรติ, ติโยชนํ หรติปิ ปจฺจาหรติปิ, ติโยชนํ วาสาธิปฺปาโย คนฺตฺวา ตโต ปรํ หรติ, อจฺฉินฺนํ ปฏิลภิตฺวา หรติ, นิสฺสฏฺฐํ ปฏิลภิตฺวา หรติ, อญฺญํ หราเปติ, กตภณฺฑํ, อุมฺมตฺตกสฺส อาทิกมฺมิกสฺสาติ.}

\vspace{\tipitakaparskip}
\csromancenter{เอฬกโลม\-สิกฺขาปทํ นิฏฺฐิตํ ฉฏฺฐํ.}
\csromanmatikaanchor{mtk.184}
\csromantocmarknopage{subsection}{7. เอฬกโลมโธวาปนสิกฺขาปท}{mtk.184}
\bookmark[dest={mtk.184},level=2]{7. เอฬกโลมโธวาปนสิกฺขาปท}
\csromanmarkh{2. โกสิยวคฺค}\csromanmarkhii{7. เอฬกโลม\-โธวาปน\-สิกฺขาปท}\csromanheadernomark{2}{2. \textbf{โกสิยวคฺค 7. เอฬกโลม}\-\textbf{โธวาปน}\-\textbf{สิกฺขาปท}}
\csromanmatikaanchor{mtk.185}
\csromanmatikaanchor{mtk.186}
\csromantocmarkref{subsubsection}{ฉพฺพคฺคิยภิกฺขุวตฺถุ}{mtk.185}
\bookmark[dest={mtk.185},level=3]{ฉพฺพคฺคิยภิกฺขุวตฺถุ}
\csromantocmarkref{subsubsection}{ปญฺญตฺติ}{mtk.186}
\bookmark[dest={mtk.186},level=3]{ปญฺญตฺติ}
\proseitem{576}{\csromanfolio{342}เตน สมเยน พุทฺโธ ภควา สกฺเกสุ วิหรติ กปิล\-วตฺถุสฺมิํ นิโคฺรธาราเม. เตน โข ปน สมเยน ฉพฺพคฺคิยา ภิกฺขู ภิกฺขุนีหิ เอฬกโลมานิ โธวาเปนฺติปิ รชาเปนฺติปิ วิชฏาเปนฺติปิ. ภิกฺขุนิโย เอฬกโลมานิ โธวนฺติโย รชนฺติโย วิชเฏนฺติโย ริญฺจนฺติ อุทฺเทสํ ปริปุจฺฉํ อธิสีลํ อธิจิตฺตํ อธิปญฺญํ. อถ โข มหาปชาปติ โคตมี เยน ภควา เตนุปสงฺกมิ, อุปสงฺกมิตฺวา ภควนฺตํ อภิวาเทตฺวา เอกมนฺตํ อฏฺฐาสิ, เอกมนฺตํ ฐิตํ โข มหาปชาปติํ โคตมิํ ภควา เอตทโวจ “กจฺจิ โคตมิ ภิกฺขุนิโย อปฺปมตฺตา อาตาปินิโย ปหิตตฺตา วิหรนฺตี”ติ. กุโต ภนฺเต ภิกฺขุนีนํ อปฺปมาโท, อยฺยา ฉพฺพคฺคิยา ภิกฺขุนีหิ เอฬกโลมานิ โธวาเปนฺติปิ รชาเปนฺติปิ วิชฏาเปนฺติปิ, ภิกฺขุนิโย เอฬกโลมานิ โธวนฺติโย รชนฺติโย วิชเฏนฺติโย ริญฺจนฺติ อุทฺเทสํ ปริปุจฺฉํ อธิสีลํ อธิจิตฺตํ อธิปญฺญนฺติ.}

\prose{อถ โข ภควา มหาปชาปติํ โคตมิํ ธมฺมิยา กถาย สนฺทสฺเสสิ สมาทเปสิ สมุตฺเตเชสิ สมฺปหํเสสิ. อถ โข มหาปชาปติ โคตมี ภควตา ธมฺมิยา กถาย สนฺทสฺสิตา สมาทปิตา สมุตฺเตชิตา สมฺปหํสิตา ภควนฺตํ อภิวาเทตฺวา ปทกฺขิณํ กตฺวา ปกฺกามิ. อถ โข ภควา เอตสฺมิํ นิทาเน เอตสฺมิํ ปกรเณ ภิกฺขุสํฆํ สนฺนิปาตาเปตฺวา ฉพฺพคฺคิเย ภิกฺขู ปฏิปุจฺฉิ “สจฺจํ กิร ตุเมฺห ภิกฺขเว ภิกฺขุนีหิ เอฬกโลมานิ โธวาเปถปิ รชาเปถปิ วิชฏาเปถปี”ติ. สจฺจํ ภควาติ. ญาติกาโย ตุมฺหากํ ภิกฺขเว อญฺญาติกาโยติ. อญฺญาติกาโย ภควาติ. อญฺญาตกา โมฆปุริสา อญฺญาติกานํ น ชานนฺติ ปติรูปํ วา อปฺปติรูปํ วา ปาสาทิกํ วา อปาสาทิกํ วา, ตตฺถ นาม ตุเมฺห โมฆปุริสา อญฺญาติกาหิ ภิกฺขุนีหิ เอฬกโลมานิ โธวาเปสฺสถปิ รชาเปสฺสถปิ วิชฏาเปสฺสถปิ. เนตํ โมฆปุริสา อปฺปสนฺนานํ วา ปสาทาย ฯเปฯ. เอวญฺจ ปน ภิกฺขเว อิมํ สิกฺขาปทํ อุทฺทิเสยฺยาถ\csromandash{}}

\proseitem{577}{\textbf{“โย ปน ภิกฺขุ อญฺญาติกาย ภิกฺขุนิยา เอฬกโลมานิ โธวาเปยฺย วา รชาเปยฺย วา วิชฏาเปยฺย วา, นิสฺสคฺคิยํ ปาจิตฺติยนฺ”}ติ.}

\csromanmatikaanchor{mtk.187}
\csromantocmarkref{subsubsection}{สิกฺขาปทวิภงฺค, ปทภาชนีย}{mtk.187}
\bookmark[dest={mtk.187},level=3]{สิกฺขาปทวิภงฺค, ปทภาชนีย}
\proseitem{578}{\csromanfolio{343}\textbf{โย ปนา}ติ โย ยาทิโส ฯเปฯ. \textbf{ภิกฺขู}ติ ฯเปฯ อยํ อิมสฺมิํ อตฺเถ อธิปฺเปโต “ภิกฺขู”ติ.}

\prose{\textbf{อญฺญาติกา} นาม มาติโต วา ปิติโต วา ยาว สตฺตมา ปิตามหยุคา อสมฺพทฺธา.}

\prose{\textbf{ภิกฺขุนี} นาม อุภโตสํเฆ อุปสมฺปนฺนา.}

\prose{โธวาติ อาณาเปติ, อาปตฺติ ทุกฺกฏสฺส. โธตานิ นิสฺสคฺคิยานิ โหนฺติ. รชาติ อาณาเปติ, อาปตฺติ ทุกฺกฏสฺส. รตฺตานิ นิสฺสคฺคิยานิ โหนฺติ. วิชเฏหีติ อาณาเปติ, อาปตฺติ ทุกฺกฏสฺส. วิชฏิตานิ นิสฺสคฺคิยานิ โหนฺติ. นิสฺสชฺชิตพฺพานิ สํฆสฺส วา คณสฺส วา ปุคฺคลสฺส วา. เอวญฺจ ปน ภิกฺขเว นิสฺสชฺชิตพฺพานิ ฯเปฯ “อิมานิ เม ภนฺเต เอฬกโลมานิ อญฺญาติกาย ภิกฺขุนิยา โธวาปิตานิ นิสฺสคฺคิยานิ, อิมานาหํ สํฆสฺส นิสฺสชฺชามี”ติ ฯเปฯ ทเทยฺยาติ ฯเปฯ ทเทยฺยุนฺติ ฯเปฯ อายสฺมโต ทมฺมีติ.}

\proseitem{579}{อญฺญาติกาย อญฺญาติกสญฺญี เอฬกโลมานิ โธวาเปติ, นิสฺสคฺคิยํ ปาจิตฺติยํ. อญฺญาติกาย อญฺญาติกสญฺญี เอฬกโลมานิ โธวาเปติ รชาเปติ, นิสฺสคฺคิเยน อาปตฺติ ทุกฺกฏสฺส. อญฺญาติกาย อญฺญาติกสญฺญี เอฬกโลมานิ โธวาเปติ วิชฏาเปติ, นิสฺสคฺคิเยน อาปตฺติ ทุกฺกฏสฺส. อญฺญาติกาย อญฺญาติกสญฺญี เอฬกโลมานิ โธวาเปติ รชาเปติ วิชฏาเปติ, นิสฺสคฺคิเยน อาปตฺติ ทฺวินฺนํ ทุกฺกฏานํ.}

\prose{อญฺญาติกาย อญฺญาติกสญฺญี เอฬกโลมานิ รชาเปติ, นิสฺสคฺคิยํ ปาจิตฺติยํ. อญฺญาติกาย อญฺญาติกสญฺญี เอฬกโลมานิ รชาเปติ วิชฏาเปติ, นิสฺสคฺคิเยน อาปตฺติ ทุกฺกฏสฺส. อญฺญาติกาย อญฺญาติกสญฺญี เอฬกโลมานิ รชาเปติ โธวาเปติ, นิสฺสคฺคิเยน อาปตฺติ ทุกฺกฏสฺส. อญฺญาติกาย อญฺญาติกสญฺญี เอฬกโลมานิ รชาเปติ วิชฏาเปติ โธวาเปติ, นิสฺสคฺคิเยน อาปตฺติ ทฺวินฺนํ ทุกฺกฏานํ.}

\csromanmatikaanchor{mtk.188}
\csromanmatikaanchor{mtk.189}
\csromantocmarknopage{subsection}{8. รูปิยสิกฺขาปท}{mtk.188}
\bookmark[dest={mtk.188},level=2]{8. รูปิยสิกฺขาปท}
\csromantocmarkref{subsubsection}{อุปนนฺทภิกฺขุวตฺถุ}{mtk.189}
\bookmark[dest={mtk.189},level=3]{อุปนนฺทภิกฺขุวตฺถุ}
\prose{อญฺญาติกาย อญฺญาติกสญฺญี เอฬกโลมานิ วิชฏาเปติ, นิสฺสคฺคิยํ ปาจิตฺติยํ. อญฺญาติกาย อญฺญาติกสญฺญี เอฬกโลมานิ วิชฏาเปติ โธวาเปติ, นิสฺสคฺคิเยน อาปตฺติ ทุกฺกฏสฺส. อญฺญาติกาย อญฺญาติกสญฺญี เอฬกโลมานิ วิชฏาเปติ รชาเปติ, นิสฺสคฺคิเยน อาปตฺติ \csromanfolio{344}ทุกฺกฏสฺส. อญฺญาติกาย อญฺญาติกสญฺญี เอฬกโลมานิ วิชฏาเปติ โธวาเปติ รชาเปติ, นิสฺสคฺคิเยน อาปตฺติ ทฺวินฺนํ ทุกฺกฏานํ.}

\proseitem{580}{อญฺญาติกาย เวมติโก ฯเปฯ. อญฺญาติกาย ญาติกสญฺญี ฯเปฯ อญฺญสฺส เอฬกโลมานิ โธวาเปติ, อาปตฺติ ทุกฺกฏสฺส. เอกโต อุปสมฺปนฺนาย โธวาเปติ, อาปตฺติ ทุกฺกฏสฺส. ญาติกาย อญฺญาติกสญฺญี, อาปตฺติ ทุกฺกฏสฺส. ญาติกาย เวมติโก, อาปตฺติ ทุกฺกฏสฺส. ญาติกาย ญาติกสญฺญี, อนาปตฺติ.}

\proseitemwithcloserruled{581}{อนาปตฺติ ญาติกาย โธวนฺติยา อญฺญาติกา ทุติยา โหติ, อวุตฺตา โธวติ, อปริภุตฺตํ กตภณฺฑํ โธวาเปติ, สิกฺขมานาย สามเณริยา อุมฺมตฺตกสฺส อาทิกมฺมิกสฺสาติ.}{เอฬกโลม\-โธวาปน\-สิกฺขาปทํ นิฏฺฐิตํ สตฺตมํ.}
\csromanmarkh{2. โกสิยวคฺค}\csromanmarkhii{8. รูปิย\-สิกฺขาปท}\csromanheadernomark{2}{2. \textbf{โกสิยวคฺค 8. รูปิย}\-\textbf{สิกฺขาปท}}
\proseitem{582}{เตน สมเยน พุทฺโธ ภควา ราชคเห วิหรติ เวฬุวเน กลนฺทก\-นิวาเป. เตน โข ปน สมเยน อายสฺมา อุปนนฺโท สกฺยปุตฺโต อญฺญตรสฺส กุลสฺส กุลูปโก โหติ นิจฺจภตฺติโก. ยํ ตสฺมิํ กุเล อุปฺปชฺชติ ขาทนียํ วา โภชนียํ วา, ตโต อายสฺมโต อุปนนฺทสฺส สกฺยปุตฺตสฺส ปฏิวิโส ฐปิยฺยติ. เตน โข ปน สมเยน สายํ ตสฺมิํ กุเล มํสํ อุปฺปนฺนํ โหติ, ตโต อายสฺมโต อุปนนฺทสฺส สกฺยปุตฺตสฺส ปฏิวิโส ฐปิโต โหติ. ตสฺส กุลสฺส ทารโก รตฺติยา ปจฺจูสสมยํ ปจฺจุฏฺฐาย โรทติ “มํสํ เม เทถา”ติ. อถ โข โส ปุริโส ปชาปติํ เอตทโวจ “อยฺยสฺส ปฏิวิสํ ทารกสฺส เทหิ, อญฺญํ เจตาเปตฺวา อยฺยสฺส ทสฺสามา”ติ.}

\csromanmatikaanchor{mtk.190}
\csromanmatikaanchor{mtk.191}
\csromantocmarkref{subsubsection}{ปญฺญตฺติ}{mtk.190}
\bookmark[dest={mtk.190},level=3]{ปญฺญตฺติ}
\csromantocmarkref{subsubsection}{สิกฺขาปทวิภงฺค, ปทภาชนีย}{mtk.191}
\bookmark[dest={mtk.191},level=3]{สิกฺขาปทวิภงฺค, ปทภาชนีย}
\prose{อถ โข อายสฺมา อุปนนฺโท สกฺยปุตฺโต ปุพฺพณฺหสมยํ นิวาเสตฺวา ปตฺตจีวรํ อาทาย เยน ตํ กุลํ เตนุปสงฺกมิ, อุปสงฺกมิตฺวา ปญฺญตฺเต อาสเน นิสีทิ. อถ โข โส ปุริโส เยนายสฺมา อุปนนฺโท สกฺยปุตฺโต เตนุปสงฺกมิ, อุปสงฺกมิตฺวา อายสฺมนฺตํ อุปนนฺทํ สกฺยปุตฺตํ อภิวาเทตฺวา เอกมนฺตํ นิสีทิ. เอกมนฺตํ นิสินฺโน โข โส ปุริโส \csromanfolio{345}อายสฺมนฺตํ อุปนนฺทํ สกฺยปุตฺตํ เอตทโวจ “หิยฺโย โข ภนฺเต สายํ มํสํ อุปฺปนฺนํ อโหสิ, ตโต อยฺยสฺส ปฏิวิโส ฐปิโต, อยํ ภนฺเต ทารโก รตฺติยา ปจฺจูสสมยํ ปจฺจุฏฺฐาย โรทติ ‘มํสํ เม เทถา’ติ, อยฺยสฺส ปฏิวิโส ทารกสฺส ทินฺโน, กหาปเณน ภนฺเต กิํ อาหริยฺยตู”ติ. ปริจฺจตฺโต เม อาวุโส กหาปโณติ. อาม ภนฺเต ปริจฺจตฺโตติ. ตญฺเญว เม อาวุโส กหาปณํ เทหีติ.}

\prose{อถ โข โส ปุริโส อายสฺมโต อุปนนฺทสฺส สกฺยปุตฺตสฺส กหาปณํ ทตฺวา อุชฺฌายติ ขิยฺยติ วิปาเจติ “ยเถว มยํ รูปิยํ ปฏิคฺคณฺหาม, เอวเมวิเม สมณา สกฺยปุตฺติยา รูปิยํ ปฏิคฺคณฺหนฺตี”ติ. อสฺโสสุํ โข ภิกฺขู ตสฺส ปุริสสฺส อุชฺฌายนฺตสฺส ขิยฺยนฺตสฺส วิปาเจนฺตสฺส. เย เต ภิกฺขู อปฺปิจฺฉา, เต อุชฺฌายนฺติ ขิยฺยนฺติ วิปาเจนฺติ “กถํ หิ นาม อายสฺมา อุปนนฺโท สกฺยปุตฺโต รูปิยํ ปฏิคฺคเหสฺสตี”ติ. อถ โข เต ภิกฺขู อายสฺมนฺตํ อุปนนฺทํ สกฺยปุตฺตํ อเนก\-ปริยาเยน วิครหิตฺวา ภควโต เอตมตฺถํ อาโรเจสุํ ฯเปฯ “สจฺจํ กิร ตฺวํ อุปนนฺท รูปิยํ ปฏิคฺคเหสี”ติ\csromansharedfootnote{562d106575aa7469}{ปฏิคฺคณฺหาสีติ (สฺยา)}. สจฺจํ ภควาติ. วิครหิ พุทฺโธ ภควา ฯเปฯ. กถํ หิ นาม ตฺวํ โมฆปุริส รูปิยํ ปฏิคฺคเหสฺสสิ, เนตํ โมฆปุริส อปฺปสนฺนานํ วา ปสาทาย ฯเปฯ. เอวญฺจ ปน ภิกฺขเว อิมํ สิกฺขาปทํ อุทฺทิเสยฺยาถ\csromandash{}}

\proseitem{583}{\textbf{“โย ปน ภิกฺขุ ชาตรูป}\-\textbf{รชตํ อุคฺคเณฺหยฺย วา อุคฺคณฺหาเปยฺย วา อุปนิกฺขิตฺตํ วา สาทิเยยฺย, นิสฺสคฺคิยํ ปาจิตฺติยนฺ”}ติ.}

\proseitem{584}{\textbf{โย ปนา}ติ โย ยาทิโส ฯเปฯ. \textbf{ภิกฺขู}ติ ฯเปฯ อยํ อิมสฺมิํ อตฺเถ อธิปฺเปโต “ภิกฺขู”ติ.}

\prose{\textbf{ชาตรูปํ} นาม สตฺถุวณฺโณ วุจฺจติ.}

\prose{\textbf{รชตํ} นาม กหาปโณ โลหมาสโก ทารุมาสโก ชตุมาสโก, เย โวหารํ คจฺฉนฺติ.}

\prose{\textbf{อุคฺคเณฺหยฺยา}ติ สยํ คณฺหาติ, นิสฺสคฺคิยํ ปาจิตฺติยํ\csromansharedfootnote{7174dcf91c017a08}{นิสฺสคฺคิยํ โหติ (สฺยา)}.}

\prose{\textbf{อุคฺคณฺหาเปยฺยา}ติ อญฺญํ คาหาเปติ, นิสฺสคฺคิยํ ปาจิตฺติยํ\csromansharedfootnote{7174dcf91c017a08}{นิสฺสคฺคิยํ โหติ (สฺยา)}.}

\prose{\csromanfolio{346}\textbf{อุปนิกฺขิตฺตํ วา สาทิเยยฺยา}ติ อิทํ อยฺยสฺส โหตูติ อุปนิกฺขิตฺตํ สาทิยติ, นิสฺสคฺคิยํ โหติ. สํฆมชฺเฌ นิสฺสชฺชิตพฺพํ. เอวญฺจ ปน ภิกฺขเว นิสฺสชฺชิตพฺพํ, เตน ภิกฺขุนา สํฆํ อุปสงฺกมิตฺวา เอกํสํ อุตฺตราสงฺคํ กริตฺวา วุฑฺฒานํ ภิกฺขูนํ ปาเท วนฺทิตฺวา อุกฺกุฏิกํ นิสีทิตฺวา อญฺชลิํ ปคฺคเหตฺวา เอวมสฺส วจนีโย “อหํ ภนฺเต รูปิยํ ปฏิคฺคเหสิํ, อิทํ เม นิสฺสคฺคิยํ, อิมาหํ สํฆสฺส นิสฺสชฺชามี”ติ. นิสฺสชฺชิตฺวา อาปตฺติ เทเสตพฺพา. พฺยตฺเตน ภิกฺขุนา ปฏิพเลน อาปตฺติ ปฏิคฺคเหตพฺพา. สเจ ตตฺถ อาคจฺฉติ อารามิโก วา อุปาสโก วา, โส วตฺตพฺโพ “อาวุโส อิมํ ชานาหี”ติ. สเจ โส ภณติ “อิมินา กิํ อาหริยฺยตู”ติ. น วตฺตพฺโพ “อิมํ วา อิมํ วา อาหรา”ติ, กปฺปิยํ อาจิกฺขิตพฺพํ สปฺปิ วา เตลํ วา มธุ วา ผาณิตํ วา. สเจ โส เตน ปริวตฺเตตฺวา กปฺปิยํ อาหรติ, รูปิย\-ปฺปฏิคฺคาหกํ ฐเปตฺวา สพฺเพเหว ปริภุญฺชิ\-ตพฺพํ. เอวญฺเจตํ ลเภถ, อิจฺเจตํ กุสลํ. โน เจ ลเภถ, โส วตฺตพฺโพ “อาวุโส อิมํ ฉฑฺเฑหี”ติ. สเจ โส ฉฑฺเฑติ, อิจฺเจตํ กุสลํ. โน เจ ฉฑฺเฑติ, ปญฺจหงฺเคหิ สมนฺนาคโต ภิกฺขุ รูปิยฉฑฺฑโก สมฺมนฺนิตพฺโพ, โย น ฉนฺทาคติํ คจฺฉยฺย, น โทสาคติํ คจฺเฉยฺย, น โมหาคติํ คจฺเฉยฺย, น ภยาคติํ คจฺเฉยฺย, ฉฑฺฑิตาฉฑฺฑิตญฺจ ชาเนยฺย. เอวญฺจ ปน ภิกฺขเว สมฺมนฺนิตพฺโพ, ปฐมํ ภิกฺขุ ยาจิตพฺโพ, ยาจิตฺวา พฺยตฺเตน ภิกฺขุนา ปฏิพเลน สํโฆ ญาเปตพฺโพ\csromandash{}}

\proseitem{585}{“สุณาตุ เม ภนฺเต สํโฆ, ยทิ สํฆสฺส ปตฺตกลฺลํ, สํโฆ อิตฺถนฺนามํ ภิกฺขุํ รูปิย\-ฉฑฺฑกํ สมฺมนฺเนยฺย, เอสา ญตฺติ.}

\prose{สุณาตุ เม ภนฺเต สํโฆ, สํโฆ อิตฺถนฺนามํ ภิกฺขุํ รูปิย\-ฉฑฺฑกํ สมฺมนฺนติ, ยสฺสายสฺมโต ขมติ อิตฺถนฺนามสฺส ภิกฺขุโน รูปิย\-ฉฑฺฑกสฺส สมฺมุติ, โส ตุณฺหสฺส, ยสฺส นกฺขมติ, โส ภาเสยฺย.}

\prose{สมฺมโต สํเฆน อิตฺถนฺนาโม ภิกฺขุ รูปิยฉฑฺฑโก, ขมติ สํฆสฺส, ตสฺมา ตุณฺหี, เอวเมตํ ธารยามี”ติ.}

\prose{เตน สมฺมเตน ภิกฺขุนา อนิมิตฺตํ กตฺวา ปาเตตพฺพํ. สเจ นิมิตฺตํ กตฺวา ปาเตติ, อาปตฺติ ทุกฺกฏสฺส.}

\proseitem{586}{รูปิเย รูปิยสญฺญี รูปิยํ ปฏิคฺคณฺหาติ, นิสฺสคฺคิยํ ปาจิตฺติยํ. รูปิเย เวมติโก รูปิยํ ปฏิคฺคณฺหาติ, นิสฺสคฺคิยํ ปาจิตฺติยํ. รูปิเย อรูปิยสญฺญี รูปิยํ ปฏิคฺคณฺหาติ, นิสฺสคฺคิยํ ปาจิตฺติยํ.}

\csromanmatikaanchor{mtk.192}
\csromanmatikaanchor{mtk.193}
\csromanmatikaanchor{mtk.194}
\csromanmatikaanchor{mtk.195}
\csromantocmarknopage{subsection}{9. รูปิยสํโวหารสิกฺขาปท}{mtk.192}
\bookmark[dest={mtk.192},level=2]{9. รูปิยสํโวหารสิกฺขาปท}
\csromantocmarkref{subsubsection}{ฉพฺพคฺคิยภิกฺขุวตฺถุ}{mtk.193}
\bookmark[dest={mtk.193},level=3]{ฉพฺพคฺคิยภิกฺขุวตฺถุ}
\csromantocmarkref{subsubsection}{ปญฺญตฺติ}{mtk.194}
\bookmark[dest={mtk.194},level=3]{ปญฺญตฺติ}
\csromantocmarkref{subsubsection}{สิกฺขาปทวิภงฺค, ปทภาชนีย}{mtk.195}
\bookmark[dest={mtk.195},level=3]{สิกฺขาปทวิภงฺค, ปทภาชนีย}
\prose{\csromanfolio{347}อรูปิเย รูปิยสญฺญี, อาปตฺติ ทุกฺกฏสฺส. อรูปิเย เวมติโก, อาปตฺติ ทุกฺกฏสฺส. อรูปิเย อรูปิยสญฺญี, อนาปตฺติ.}

\prosewithcloserruled{อนาปตฺติ อชฺฌาราเม วา อชฺฌาวสเถ วา อุคฺคเหตฺวา วา อุคฺคหาเปตฺวา วา นิกฺขิปติ “ยสฺส ภวิสฺสติ โส หริสฺสตี”ติ, อุมฺมตฺตกสฺส อาทิกมฺมิกสฺสติ.}{รูปิย\-สิกฺขาปทํ นิฏฺฐิตํ อฏฺฐมํ.}
\csromanmarkh{2. โกสิยวคฺค}\csromanmarkhii{9. รูปิยสํโวหาร\-สิกฺขาปท}\csromanheadernomark{2}{2. \textbf{โกสิยวคฺค 9. รูปิยสํโวหาร}\-\textbf{สิกฺขาปท}}
\proseitem{587}{เตน สมเยน พุทฺโธ ภควา สาวตฺถิยํ วิหรติ เชตวเน อนาถ\-ปิณฺฑิกสฺส อาราเม. เตน โข ปน สมเยน ฉพฺพคฺคิยา ภิกฺขู นานปฺปการกํ รูปิย\-สํโวหารํ สมาปชฺชนฺติ. มนุสฺสา อุชฺฌายนฺติ ขิยฺยนฺติ วิปาเจนฺติ “กถํ หิ นาม สมณา สกฺยปุตฺติยา นานปฺปการกํ รูปิย\-สํโวหารํ สมาปชฺชิสฺสนฺติ เสยฺยถาปิ คิหี กามโภคิโน”ติ. อสฺโสสุํ โข ภิกฺขู เตสํ มนุสฺสานํ อุชฺฌยนฺตานํ ขิยฺยนฺตานํ วิปาเจนฺตานํ. เย เต ภิกฺขู อปฺปิจฺฉา ฯเปฯ เต อุชฺฌายนฺติ ขิยฺยนฺติ วิปาเจนฺติ “กถํ หิ นาม ฉพฺพคฺคิยา ภิกฺขู นานปฺปการกํ รูปิย\-สํโวหารํ สมาปชฺชิสฺสนฺตี”ติ. อถ โข เต ภิกฺขู ฉพฺพคฺคิเย ภิกฺขู อเนก\-ปริยาเยน วิครหิตฺวา ภควโต เอตมตฺถํ อาโรเจสุํ ฯเปฯ “สจฺจํ กิร ตุเมฺห ภิกฺขเว นานปฺปการกํ รูปิย\-สํโวหารํ สมาปชฺชถา”ติ. สจฺจํ ภควาติ. วิครหิ พุทฺโธ ภควา ฯเปฯ. กถํ หิ นาม ตุเมฺห โมฆปุริสา นานปฺปการกํ รูปิย\-สํโวหารํ สมาปชฺชิสฺสถ, เนตํ โมฆปุริสา อปฺปสนฺนานํ วา ปสาทาย ฯเปฯ. เอวญฺจ ปน ภิกฺขเว อิมํ สิกฺขาปทํ อุทฺทิเสยฺยาถ\csromandash{}}

\proseitem{588}{\textbf{“โย ปน ภิกฺขุ นานปฺปการกํ รูปิย}\-\textbf{สํโวหารํ สมาปชฺเชยฺย, นิสฺสคฺคิยํ ปาจิตฺติยนฺ”}ติ.}

\proseitem{589}{\textbf{โย ปนา}ติ โย ยาทิโส ฯเปฯ. \textbf{ภิกฺขู}ติ ฯเปฯ อยํ อิมสฺมิํ อตฺเถ อธิปฺเปโต “ภิกฺขู”ติ.}

\prose{\csromanfolio{348}\textbf{นานปฺปการกํ} นาม กตํปิ อกตํปิ กตากตํปิ. \textbf{กตํ} นาม สีสูปคํ คีวูปคํ หตฺถูปคํ ปาทูปคํ กฏูปคํ. \textbf{อกตํ} นาม ฆนกตํ วุจฺจติ. กตากตํ นาม ตทุภยํ.}

\prose{\textbf{รูปิยํ} นาม สตฺถุวณฺโณ กหาปโณ โลหมาสโก ทารุมาสโก ชตุมาสโก, เย โวหารํ คจฺฉนฺติ.}

\prose{\textbf{สมาปชฺเชยฺยา}ติ กเตน กตํ เจตาเปติ, นิสฺสคฺคิยํ ปาจิตฺติยํ\csromansharedfootnote{7174dcf91c017a08}{นิสฺสคฺคิยํ โหติ (สฺยา)}. กเตน อกตํ เจตาเปติ, นิสฺสคฺคิยํ ปาจิตฺติยํ. กเตน กตากตํ เจตาเปติ, นิสฺสคฺคิยํ ปาจิตฺติยํ. อกเตน กตํ เจตาเปติ, นิสฺสคฺคิยํ ปาจิตฺติยํ. อกเตน อกตํ เจตาเปติ, นิสฺสคฺคิยํ ปาจิตฺติยํ. อกเตน กตากตํ เจตาเปติ, นิสฺสคฺคิยํ ปาจิตฺติยํ. กตากเตน กตํ เจตาเปติ, นิสฺสคฺคิยํ ปาจิตฺติยํ. กตากเตน อกตํ เจตาเปติ, นิสฺสคฺคิยํ ปาจิตฺติยํ. กตากเตน กตากตํ เจตาเปติ, นิสฺสคฺคิยํ ปาจิตฺติยํ. สํฆมชฺเฌ นิสฺสชฺชิตพฺพํ. เอวญฺจ ปน ภิกฺขเว นิสฺสชฺชิตพฺพํ. เตน ภิกฺขุนา สํฆํ อุปสงฺกมิตฺวา เอกํสํ อุตฺตราสงฺคํ กริตฺวา วุฑฺฒานํ ภิกฺขูนํ ปาเท วนฺทิตฺวา อุกฺกุฏิกํ นิสีทิตฺวา อญฺชลิํ ปคฺคเหตฺวา เอวมสฺส วจนีโย “อหํ ภนฺเต นานปฺปการกํ รูปิย\-สํโวหารํ สมาปชฺชิํ, อิทํ เม นิสฺสคฺคิยํ, อิมาหํ สํฆสฺส นิสฺสชฺชามี”ติ. นิสฺสชฺชิตฺวา อาปตฺติ เทเสตพฺพา, พฺยตฺเตน ภิกฺขุนา ปฏิพเลน อาปตฺติ ปฏิคฺคเหตพฺพา, สเจ ตตฺถ อาคจฺฉติ อารามิโก วา อุปาสโก วา, โส วตฺตพฺโพ “อาวุโส อิมํ ชานาหี”ติ. สเจ โส ภณติ “อิมินา กิํ อาหริยฺยตู”ติ. น วตฺตพฺโพ “อิมํ วา อิมํ วา อาหรา”ติ. กปฺปิยํ อาจิกฺขิตพฺพํ สปฺปิ วา เตลํ วา มธุ วา ผาณิตํ วา. สเจ โส เตน ปริวตฺเตตฺวา กปฺปิยํ อาหรติ, รูปิย\-เจตาปกํ ฐเปตฺวา สพฺเพเหว ปริภุญฺชิ\-ตพฺพํ. เอวญฺเจตํ ลเภถ, อิจฺเจตํ กุสลํ. โน เจ ลเภถ, โส วตฺตพฺโพ “อาวุโส อิมํ ฉฑฺเฑหี”ติ. สเจ โส ฉฑฺเฑติ, อิจฺเจตํ กุสลํ. โน เจ ฉฑฺเฑติ, ปญฺจหงฺเคหิ สมนฺนาคโต ภิกฺขุ รูปิยฉฑฺฑโก สมฺมนฺนิตพฺโพ, โย น ฉนฺทาคติํ คจฺเฉยฺย, น โทสาคติํ คจฺเฉยฺย, น โมหาคติํ คจฺเฉยฺย, น ภยาคติํ คจฺเฉยฺย, ฉฑฺฑิตาฉฑฺฑิตญฺจ ชาเนยฺย. เอวญฺจ ปน ภิกฺขเว สมฺมนฺนิตพฺโพ, ปฐมํ ภิกฺขุ ยาจิตพฺโพ, ยาจิตฺวา พฺยตฺเตน ภิกฺขุนา ปฏิพเลน สํโฆ ญาเปตพฺโพ\csromandash{}}

\csromanmatikaanchor{mtk.196}
\csromanmatikaanchor{mtk.197}
\csromantocmarknopage{subsection}{10. กยวิกฺกยสิกฺขาปท}{mtk.196}
\bookmark[dest={mtk.196},level=2]{10. กยวิกฺกยสิกฺขาปท}
\csromantocmarkref{subsubsection}{อุปนนฺทภิกฺขุวตฺถุ}{mtk.197}
\bookmark[dest={mtk.197},level=3]{อุปนนฺทภิกฺขุวตฺถุ}
\proseitem{590}{\csromanfolio{349}“สุณาตุ เม ภนฺเต สํโฆ, ยทิ สํฆสฺส ปตฺตกลฺลํ, สํโฆ อิตฺถนฺนามํ ภิกฺขุํ รูปิย\-ฉฑฺฑกํ สมฺมนฺเนยฺย, เอสา ญตฺติ.}

\prose{สุณาตุ เม ภนฺเต สํโฆ, สํโฆ อิตฺถนฺนามํ ภิกฺขุํ รูปิย\-ฉฑฺฑกํ สมฺมนฺนติ, ยสฺสายสฺมโต ขมติ อิตฺถนฺนามสฺส ภิกฺขุโน รูปิย\-ฉฑฺฑกสฺส สมฺมุติ, โส ตุณฺหสฺส, ยสฺส นกฺขมติ, โส ภาเสยฺย.}

\prose{สมฺมโต สํเฆน อิตฺถนฺนาโม ภิกฺขุ รูปิยฉฑฺฑโก, ขมติ สํฆสฺส, ตสฺมา ตุณฺหี, เอวเมตํ ธารยามี”ติ.}

\prose{เตน สมฺมเตน ภิกฺขุนา อนิมิตฺตํ กตฺวา ปาเตตพฺพํ. สเจ นิมิตฺตํ กตฺวา ปาเตติ, อาปตฺติ ทุกฺกฏสฺส.}

\proseitem{591}{รูปิเย รูปิยสญฺญี รูปิยํ เจตาเปติ, นิสฺสคฺคิยํ ปาจิตฺติยํ. รูปิเย เวมติโก รูปิยํ เจตาเปติ, นิสฺสคฺคิยํ ปาจิตฺติยํ. รูปิเย อรูปิยสญฺญี รูปิยํ เจตาเปติ, นิสฺสคฺคิยํ ปาจิตฺติยํ. อรูปิเย รูปิยสญฺญี รูปิยํ เจตาเปติ, นิสฺสคฺคิยํ ปาจิตฺติยํ. อรูปิเย เวมติโก รูปิยํ เจตาเปติ, นิสฺสคฺคิยํ ปาจิตฺติยํ. อรูปิเย อรูปิยสญฺญี รูปิยํ เจตาเปติ, นิสฺสคฺคิยํ ปาจิตฺติยํ.}

\prose{อรูปิเย รูปิยสญฺญี, อาปตฺติ ทุกฺกฏสฺส. อรูปิเย เวมติโก, อาปตฺติ ทุกฺกฏสฺส. อรูปิเย อรูปิยสญฺญี อนาปตฺติ.}

\proseitemwithcloserruled{392}{อนาปตฺติ อุมฺมตฺตกสฺส อาทิกมฺมิกสฺสาติ.}{รูปิยสํโวหาร\-สิกฺขาปทํ นิฏฺฐิตํ นวมํ.}
\csromanmarkh{2. โกสิยวคฺค}\csromanmarkhii{10. กยวิกฺกย\-สิกฺขาปท}\csromanheadernomark{2}{2. โกสิยวคฺค 10. กยวิกฺกย\-สิกฺขาปท}
\csromanmatikaanchor{mtk.198}
\csromanmatikaanchor{mtk.199}
\csromantocmarkref{subsubsection}{ปญฺญตฺติ}{mtk.198}
\bookmark[dest={mtk.198},level=3]{ปญฺญตฺติ}
\csromantocmarkref{subsubsection}{สิกฺขาปทวิภงฺค, ปทภาชนีย}{mtk.199}
\bookmark[dest={mtk.199},level=3]{สิกฺขาปทวิภงฺค, ปทภาชนีย}
\proseitem{593}{เตน สมเยน พุทฺโธ ภควา สาวตฺถิยํ วิหรติ เชตวเน อนาถ\-ปิณฺฑิกสฺส อาราเม. เตน โข ปน สมเยน อายสฺมา อุปนนฺโท สกฺยปุตฺโต ปฏฺโฏ โหติ จีวรกมฺมํ กาตุํ. โส ปฏปิโลติกานํ สงฺฆาฏิํ กริตฺวา สุรตฺตํ สุปริกมฺมกตํ กตฺวา ปารุปิ. อถ โข อญฺญตโร ปริพฺพาชโก มหคฺฆํ ปฏํ ปารุปิตฺวา เยนายสฺมา อุปนนฺโท สกฺยปุตฺโต เตนุปสงฺกมิ, อุปสงฺกมิตฺวา อายสฺมนฺตํ อุปนนฺทํ \csromanfolio{350}สกฺยปุตฺตํ เอตทโวจ “สุนฺทรา โข ตฺยายํ อาวุโส สงฺฆาฏิ, เทหิ เม ปเฏนา”ติ. ชานาหิ อาวุโสติ. อามาวุโส ชานามีติ. หนฺทาวุโสติ อทาสิ. อถ โข โส ปริพฺพาชโก ตํ สงฺฆาฏิํ ปารุปิตฺวา ปริพฺพาชกา\-รามํ อคมาสิ. ปริพฺพาชกา ตํ ปริพฺพาชกํ เอตทโวจุํ “สุนฺทรา โข ตฺยายํ อาวุโส สงฺฆาฏิ, กุโต ตยา ลทฺธา”ติ. เตน เม อาวุโส ปเฏน ปริวตฺติตาติ. กติหิปิ ตฺยายํ อาวุโส สงฺฆาฏิ ภวิสฺสติ, โส เยว เต ปโฏ วโรติ.}

\prose{อถ โข โส ปริพฺพาชโก “สจฺจํ โข ปริพฺพาชกา อาหํสุ, กติหิปิ มฺยายํ สงฺฆาฏิ ภวิสฺสติ, โส เยว เม ปโฏ วโร”ติ เยนายสฺมา อุปนนฺโท สกฺยปุตฺโต เตนุปสงฺกมิ , อุปสงฺกมิตฺวา อายสฺมนฺตํ อุปนนฺทํ สกฺยปุตฺตํ เอตทโวจ “หนฺท เต อาวุโส สงฺฆาฏิ\csromansharedfootnote{9d3a468491b6d026}{สงฺฆาฏิํ (สฺยา, ก)}, เทหิ เม ปฏนฺ”ติ. นนุ ตฺวํ อาวุโส มยา วุตฺโต “ชานาหิ อาวุโส”ติ, นาหํ ทสฺสามีติ. อถ โข โส ปริพฺพาชโก อุชฺฌายติ ขิยฺยติ วิปาเจติ “คิหีปิ นํ คิหิสฺส วิปฺปฏิสาริสฺส เทนฺติ, กิํ ปน ปพฺพชิโต ปพฺพชิตสฺส น ทสฺสตี”ติ. อสฺโสสุํ โข ภิกฺขู ตสฺส ปริพฺพาชกสฺส อุชฺฌายนฺตสฺส ขิยฺยนฺตสฺส วิปาเจนฺตสฺส. เย เต ภิกฺขู อปฺปิจฺฉา ฯเปฯ เต อุชฺฌายนฺติ ขิยฺยนฺติ วิปาเจนฺติ “กถํ หิ นาม อายสฺมา อุปนนฺโท สกฺยปุตฺโต ปริพฺพาชเกน สทฺธิํ กยวิกฺกยํ สมาปชฺชิสฺสตี”ติ. อถ โข เต ภิกฺขู อายสฺมนฺตํ อุปนนฺทํ สกฺยปุตฺตํ อเนก\-ปริยาเยน วิครหิตฺวา ภควโต เอตมตฺถํ อาโรเจสุํ ฯเปฯ “สจฺจํ กิร ตฺวํ อุปนนฺท ปริพฺพาชเกน สทฺธิํ กยวิกฺกยํ สมาปชฺชสี”ติ. สจฺจํ ภควาติ. วิครหิ พุทฺโธ ภควา ฯเปฯ. กถํ หิ นาม ตฺวํ โมฆปุริส ปริพฺพาชเกน สทฺธิํ กยวิกฺกยํ สมาปชฺชิสฺสสิ, เนตํ โมฆปุริส อปฺปสนฺนานํ วา ปสาทาย ฯเปฯ. เอวญฺจ ปน ภิกฺขเว อิมํ สิกฺขาปทํ อุทฺทิเสยฺยาถ\csromandash{}}

\proseitem{594}{\textbf{“โย ปน ภิกฺขุ นานปฺปการกํ กยวิกฺกยํ สมาปชฺเชยฺย, นิสฺสคฺคิยํ ปาจิตฺติยนฺ”ติ}.}

\proseitem{595}{\textbf{โย ปนา}ติ โย ยาทิโส ฯเปฯ. \textbf{ภิกฺขู}ติ ฯเปฯ อยํ อิมสฺมิํ อตฺเถ อธิปฺเปโต “ภิกฺขู”ติ.}

\csromanmatikaanchor{mtk.200}
\csromantocmarkref{subsection}{อุทฺทานคาถา}{mtk.200}
\bookmark[dest={mtk.200},level=2]{อุทฺทานคาถา}
\prose{\csromanfolio{351}\textbf{นานปฺปการกํ} นาม จีวร\-ปิณฺฑปาต\-เสนาสน\-คิลานปฺปจฺจย\-เภสชฺช\-ปริกฺขารา, อนฺตมโส จุณฺณปิณฺโฑปิ ทนฺตกฏฺฐํปิ ทสิกสุตฺตํปิ.}

\prose{\textbf{กยวิกฺกยํ สมาปชฺเชยฺยา}ติ “อิมินา อิมํ เทหิ, อิมินา อิมํ อาหร, อิมินา อิมํ ปริวตฺเตหิ, อิมินา อิมํ เจตาเปหี”ติ อชฺฌาจรติ, อาปตฺติ ทุกฺกฏสฺส. ยโต กยิตญฺจ โหติ วิกฺกยิตญฺจ, อตฺตโน ภณฺฑํ ปรหตฺถคตํ ปรภณฺฑํ อตฺตโน หตฺถคตํ, นิสฺสคฺคิยํ โหติ. นิสฺสชฺชิตพฺพํ สํฆสฺส วา คณสฺส วา ปุคฺคลสฺส วา. เอวญฺจ ปน ภิกฺขเว นิสฺสชฺชิตพฺพํ ฯเปฯ “อหํ ภนฺเต นานปฺปการกํ กยวิกฺกยํ สมาปชฺชิํ, อิทํ เม นิสฺสคฺคิยํ, อิมาหํ สํฆสฺส นิสฺสชฺชามี”ติ ฯเปฯ ทเทยฺยาติ ฯเปฯ ทเทยฺยุนฺติ ฯเปฯ อายสฺมโต ทมฺมีติ.}

\proseitem{596}{กยวิกฺกเย กยวิกฺกย\-สญฺญี, นิสฺสคฺคิยํ ปาจิตฺติยํ. กยวิกฺกเย เวมติโก, นิสฺสคฺคิยํ ปาจิตฺติยํ. กยวิกฺกเย นกยวิกฺกย\-สญฺญี, นิสฺสคฺคิยํ ปาจิตฺติยํ.}

\prose{นกยวิกฺกเย กยวิกฺกย\-สญฺญี, อาปตฺติ ทุกฺกฏสฺส. นกยวิกฺกเย เวมติโก, อาปตฺติ ทุกฺกฏสฺส. นกยวิกฺกเย นกยวิกฺกย\-สญฺญี, อนาปตฺติ.}

\proseitem{597}{อนาปตฺติ อคฺฆํ ปุจฺฉติ, กปฺปิย\-การกสฺส อาจิกฺขติ, “อิทํ อมฺหากํ อตฺถิ อมฺหากญฺจ อิมินา จ อิมินา จ อตฺโถ”ติ ภณติ, อุมฺมตฺตกสฺส อาทิกมฺมิกสฺสาติ.}

\vspace{\tipitakaparskip}
\csromancenter{กยวิกฺกสิกฺขาปทํ นิฏฺฐิตํ ทสมํ.}
\nitthitammidruled{โกสิยวคฺโค ทุติโย.}
\csromancenter{ตสฺสุทฺทานํ}
\setlength{\csromangathapagewidth}{0pt}
\setlength{\csromangathawakwidth}{0pt}
\csromangathameasuregroup{โกสิยา สุทฺธเทฺวภาคา, \\ เทฺว จ โลมานิ อุคฺคเณฺห,}{\csromangathabat{โกสิยา สุทฺธเทฺวภาคา,}{ฉพฺพสฺสานิ นิสีทนํ.} \\ \csromangathabat{เทฺว จ โลมานิ อุคฺคเณฺห,}{อุโภ นานปฺปการกาติ.}}
\csromangathasetleft{โกสิยา สุทฺธเทฺวภาคา, \\ เทฺว จ โลมานิ อุคฺคเณฺห,}
\csromangathagroup{\csromangathastanza{\csromangathabat{โกสิยา สุทฺธเทฺวภาคา,}{ฉพฺพสฺสานิ นิสีทนํ.} \\ \csromangathabat{เทฺว จ โลมานิ อุคฺคเณฺห,}{อุโภ นานปฺปการกาติ.}}}

\csromanmatikaanchor{mtk.201}
\csromanmatikaanchor{mtk.202}
\csromantocmarknopage{section}{3. ปตฺตวคฺค}{mtk.201}
\bookmark[dest={mtk.201},level=1]{3. ปตฺตวคฺค}
\csromantocmarknopage{subsection}{1. ปตฺตสิกฺขาปท}{mtk.202}
\bookmark[dest={mtk.202},level=2]{1. ปตฺตสิกฺขาปท}
\csromanmarkh{3. ปตฺตวคฺค}\csromanmarkhii{1. ปตฺต\-สิกฺขาปท}\csromanheadernomark{2}{3. \textbf{ปตฺตวคฺค 1. ปตฺต}\-\textbf{สิกฺขาปท}}
\csromanmatikaanchor{mtk.203}
\csromanmatikaanchor{mtk.204}
\csromantocmarkref{subsubsection}{ฉพฺพคฺคิยภิกฺขุวตฺถุ}{mtk.203}
\bookmark[dest={mtk.203},level=3]{ฉพฺพคฺคิยภิกฺขุวตฺถุ}
\csromantocmarkref{subsubsection}{ปฐมปญฺญตฺติ}{mtk.204}
\bookmark[dest={mtk.204},level=3]{ปฐมปญฺญตฺติ}
\proseitem{598}{\csromanfolio{352}เตน สมเยน พุทฺโธ ภควา สาวตฺถิยํ วิหรติ เชตวเน อนาถ\-ปิณฺฑิกสฺส อาราเม. เตน โข ปน สมเยน ฉพฺพคฺคิยา ภิกฺขู พหู ปตฺเต สนฺนิจยํ กโรนฺติ. มนุสฺสา วิหารจาริกํ อาหิณฺฑนฺตา ปสฺสิตฺวา อุชฺฌายนฺติ ขิยฺยนฺติ วิปาเจนฺติ “กถํ หิ นาม สมณา สกฺยปุตฺติยา พหู ปตฺเต สนฺนิจยํ กริสฺสนฺติ, ปตฺตวาณิชฺชํ วา สมณา สกฺยปุตฺติยา กริสฺสนฺติ, อามตฺติกาปณํ วา ปสาเรสฺสนฺตี”ติ. อสฺโสสุํ โข ภิกฺขู เตสํ มนุสฺสานํ อุชฺฌายนฺตานํ ขิยฺยนฺตานํ วิปาเจนฺตานํ. เย เต ภิกฺขู อปฺปิจฺฉา ฯเปฯ เต อุชฺฌายนฺติ ขิยฺยนฺติ วิปาเจนฺติ “กถํ หิ นาม ฉพฺพคฺคิยา ภิกฺขู อติเรกปตฺตํ ธาเรสฺสนฺตี”ติ. อถ โข เต ภิกฺขู ฉพฺพคฺคิเย ภิกฺขู อเนก\-ปริยาเยน วิครหิตฺวา ภควโต เอตมตฺถํ อาโรเจสุํ ฯเปฯ “สจฺจํ กิร ตุเมฺห ภิกฺขเว อติเรกปตฺตํ ธาเรถา”ติ. สจฺจํ ภควาติ. วิครหิ พุทฺโธ ภควา ฯเปฯ. กถํ หิ นาม ตุเมฺห โมฆปุริสา อติเรกปตฺตํ ธาเรสฺสถ, เนตํ โมฆปุริสา อปฺปสนฺนานํ วา ปสาทาย ฯเปฯ. เอวญฺจ ปน ภิกฺขเว อิมํ สิกฺขาปทํ อุทฺทิเสยฺยาถ\csromandash{}}

\proseitem{599}{\textbf{“โย ปน ภิกฺขเว อติเรกปตฺตํ ธาเรยฺย, นิสฺสคฺคิยํ ปาจิตฺติยนฺ”}ติ.}

\prose{เอวญฺจิทํ ภควตา ภิกฺขูนํ สิกฺขาปทํ ปญฺญตฺตํ โหติ.}

\csromanmatikaanchor{mtk.205}
\csromanmatikaanchor{mtk.206}
\csromantocmarkref{subsubsection}{อนุปญฺญตฺติ}{mtk.205}
\bookmark[dest={mtk.205},level=3]{อนุปญฺญตฺติ}
\csromantocmarkref{subsubsection}{สิกฺขาปทวิภงฺค, ปทภาชนีย}{mtk.206}
\bookmark[dest={mtk.206},level=3]{สิกฺขาปทวิภงฺค, ปทภาชนีย}
\proseitem{600}{เตน โข ปน สมเยน อายสฺมโต อานนฺทสฺส อติเรกปตฺโต อุปฺปนฺโน โหติ, อายสฺมา จ อานนฺโท ตํ ปตฺตํ อายสฺมโต สาริปุตฺตสฺส ทาตุกาโม โหติ, อายสฺมา จ สาริปุตฺโต สาเกเต วิหรติ. อถ โข อายสฺมโต อานนฺทสฺส เอตทโหสิ “ภควตา สิกฺขาปทํ ปญฺญตฺตํ ‘น อติเรกปตฺโต ธาเรตพฺโพ’ติ, อยญฺจ เม อติเรกปตฺโต อุปฺปนฺโน, อหญฺจิมํ ปตฺตํ อายสฺมโต สาริปุตฺตสฺส ทาตุกาโม, อายสฺมา จ สาริปุตฺโต สาเกเต วิหรติ, กถํ นุ โข มยา ปฏิปชฺชิตพฺพนฺ”ติ. ภควโต เอตมตฺถํ อาโรเจสิ. กีวจิรํ ปนานนฺท สาริปุตฺโต อาคจฺฉิสฺสตีติ. นวมํ ภควา ทิวสํ ทสมํ วาติ. อถ โข ภควา เอตสฺมิํ นิทาเน เอตสฺมิํ ปกรเณ ธมฺมิํ กถํ กตฺวา ภิกฺขู อามนฺเตสิ อนุชานามิ \csromanfolio{353}ภิกฺขเว ทสาหปรมํ อติเรกปตฺตํ ธาเรตุํ. เอวญฺจ ปน ภิกฺขเว อิมํ สิกฺขาปทํ อุทฺทิเสยฺยาถ\csromandash{}}

\proseitem{601}{\textbf{“ทสาหปรมํ อติเรกปตฺโต ธาเรตพฺโพ, ตํ อติกฺกามยโต นิสฺสคฺคิยํ ปาจิตฺติยนฺ”}ติ.}

\proseitem{602}{\textbf{ทสา}\-\textbf{ห}\-\textbf{ปรมนฺ}ติ ทสาหปรมตา ธาเรตพฺโพ.}

\prose{\textbf{อติเรกปตฺโต} นาม อนธิฏฺฐิโต อวิกปฺปิโต.}

\prose{\textbf{ปตฺโต} นาม เทฺว ปตฺตา อโยปตฺโต มตฺติกาปตฺโตติ.}

\prose{ตโย ปตฺตสฺส วณฺณา อุกฺกฏฺโฐ ปตฺโต มชฺฌิโม ปตฺโต โอมโก ปตฺโต. \textbf{อุกฺกฏฺโฐ} นาม \textbf{ปตฺโต} อฑฺฒา\-ฬฺหโกทนํ คณฺหาติ จตุภาคํ ขาทนํ ตทุปิยํ พฺยญฺชนํ. \textbf{มชฺฌิโม} นาม, ปตฺโต นาฬิโกทนํ คณฺหาติ จตุภาคํ ขาทนํ ตทุปิยํ พฺยญฺชนํ. \textbf{โอมโก} นาม \textbf{ปตฺโต} ปตฺโถทนํ คณฺหาติ จตุภาคํ ขาทนํ ตทุปิยํ พฺยญฺชนํ. ตโต อุกฺกฏฺโฐ อปตฺโต โอมโก อปตฺโต.}

\prose{\textbf{ตํ อติกฺกามยโต นิสฺสคฺคิโย โหตี}ติ เอกาทเส อรุณุคฺคมเน นิสฺสคฺคิโย โหติ. นิสฺสชฺชิตพฺโพ สํฆสฺส วา คณสฺส วา ปุคฺคลสฺส วา. เอวญฺจ ปน ภิกฺขเว นิสฺสชฺชิตพฺโพ. เตน ภิกฺขุนา สํฆํ อุปสงฺกมิตฺวา เอกํสํ อุตฺตราสงฺคํ กริตฺวา วุฑฺฒานํ ภิกฺขูนํ ปาเท วนฺทิตฺวา อุกฺกุฏิกํ นิสีเทตฺวา อญฺชลิํ ปคฺคเหตฺวา เอวมสฺส วจนีโย “อยํ เม ภนฺเต ปตฺโต ทสาหาติกฺกนฺโต นิสฺสคฺคิโย, อิมาหํ สํฆสฺส นิสฺสชฺชามี”ติ. นิสฺสชฺชิตฺวา อาปตฺติ เทเสตพฺพา, พฺยตฺเตน ภิกฺขุนา ปฏิพเลน อาปตฺติ ปฏิคฺคเหตพฺพา, นิสฺสฏฺฐปตฺโต ทาตพฺโพ\csromandash{}}

\proseitem{603}{“สุณาตุ เม ภนฺเต สํโฆ, อยํ ปตฺโต อิตฺถนฺนามสฺส ภิกฺขุโน นิสฺสคฺคิโย สํฆสฺส นิสฺสฏฺโฐ, ยทิ สํฆสฺส ปตฺตกลฺลํ, สํโฆ อิมํ ปตฺตํ อิตฺถนฺนามสฺส ภิกฺขุโน ทเทยฺยา”ติ.}

\proseitem{604}{เตน ภิกฺขุนา สมฺพหุเล ภิกฺขู อุปสงฺกมิตฺวา เอกํสํ อุตฺตราสงฺคํ กริตฺวา วุฑฺฒานํ ภิกฺขูนํ ปาเท วนฺทิตฺวา อุกฺกุฏิกํ นิสีทิตฺวา อญฺชลิํ ปคฺคเหตฺวา เอวมสฺสุ วจนียา “อยํ เม ภนฺเต ปตฺโต ทสาหาติกฺกนฺโต นิสฺสคฺคิโย, อิมาหํ อายสฺมนฺตานํ นิสฺสชฺชามี”ติ. \csromanfolio{354}นิสฺสชฺชิตฺวา อาปตฺติ เทเสตพฺพา, พฺยตฺเตน ภิกฺขุนา ปฏิพเลน อาปตฺติ ปฏิคฺคเหตพฺพา, นิสฺสฏฺฐปตฺโต ทาตพฺโพ\csromandash{}}

\proseitem{605}{“สุณนฺตุ เม อายสฺมนฺตา, อยํ ปตฺโต อิตฺถนฺนามสฺส ภิกฺขุโน นิสฺสคฺคิโย อายสฺมนฺตานํ นิสฺสฏฺโฐ, ยทา\-ยสฺมนฺตานํ ปตฺตกลฺลํ, อายสฺมนฺตา อิมํ ปตฺตํ อิตฺถนฺนามสฺส ภิกฺขุโน ทเทยฺยุนฺ”ติ.}

\proseitem{606}{เตน ภิกฺขุนา เอกํ ภิกฺขุํ อุปสงฺกมิตฺวา เอกํสํ อุตฺตราสงฺคํ กริตฺวา อุกฺกุฏิกํ นิสีทิตฺวา อญฺชลิํ ปคฺคเหตฺวา เอวมสฺส วจนีโย “อยํ เม อาวุโส ปตฺโต ทสาหาติกฺกนฺโต นิสฺสคฺคิโย, อิมาหํ อายสฺมโต นิสฺสชฺชามี”ติ. นิสฺสชฺชิตฺวา อาปตฺติ เทเสตพฺพา, เตน ภิกฺขุนา อาปตฺติ ปฏิคฺคเหตพฺพา, นิสฺสฏฺฐปตฺโต ทาตพฺโพ “อิมํ ปตฺตํ อายสฺมโต ทมฺมี”ติ.}

\proseitem{607}{ทสาหาติกฺกนฺเต อติกฺกนฺตสญฺญี, นิสฺสคฺคิยํ ปาจิตฺติยํ. ทสาหาติกฺกนฺเต เวมติโก, นิสฺสคฺคิยํ ปาจิตฺติยํ. ทสาหาติกฺกนฺเต อนติกฺกนฺต\-สญฺญี, นิสฺสคฺคิยํ ปาจิตฺติยํ. อนธิฏฺฐิเต อธิฏฺฐิต\-สญฺญี, นิสฺสคฺคิยํ ปาจิตฺติยํ. อวิกปฺปิเต วิกปฺปิตสญฺญี, นิสฺสคฺคิยํ ปาจิตฺติยํ. อวิสฺสชฺชิเต วิสฺสชฺชิต\-สญฺญี, นิสฺสคฺคิยํ ปาจิตฺติยํ. อนฏฺเฐ นฏฺฐสญฺญี, นิสฺสคฺคิยํ ปาจิตฺติยํ. อวินฏฺเฐ วินฏฺฐสญฺญี, นิสฺสคฺคิยํ ปาจิตฺติยํ. อภินฺเน ภินฺนสญฺญี, นิสฺสคฺคิยํ ปาจิตฺติยํ. อวิลุตฺเต วิลุตฺตสญฺญี, นิสฺสคฺคิยํ ปาจิตฺติยํ.}

\prose{นิสฺสคฺคิยํ ปตฺตํ อนิสฺสชฺชิตฺวา ปริภุญฺชติ, อาปตฺติ ทุกฺกฏสฺส. ทสาหา\-นติกฺกนฺเต อติกฺกนฺตสญฺญี อาปตฺติ ทุกฺกฏสฺส. ทสาหา\-นติกฺกนฺเต เวมติโก, อาปตฺติ ทุกฺกฏสฺส. ทสาหา\-นติกฺกนฺเต อนติกฺกนฺต\-สญฺญี, อนาปตฺติ.}

\proseitem{608}{อนาปตฺติ อนฺโตทสาหํ อธิฏฺเฐติ, วิกปฺเปติ, วิสฺสชฺเชติ, นสฺสติ, วินสฺสติ, ภิชฺชติ, อจฺฉินฺทิตฺวา คณฺหนฺติ, วิสฺสาสํ คณฺหนฺติ, อุมฺมตฺตกสฺส อาทิกมฺมิกสฺสาติ.}

\prose{เตน โข ปน สมเยน ฉพฺพคฺคิยา ภิกฺขู นิสฺสฏฺฐ\-ปตฺตํ น เทนฺติ. ภควโต เอตมตฺถํ อาโรเจสุํ. น ภิกฺขเว นิสฺสฏฺฐปตฺโต น ทาตพฺโพ, โย น ทเทยฺย, อาปตฺติ ทุกฺกฏสฺสาติ.}

\vspace{\tipitakaparskip}
\csromancenter{ปตฺต\-สิกฺขาปทํ นิฏฺฐิตํ ปฐมํ.}
\csromanmatikaanchor{mtk.207}
\csromantocmarknopage{subsection}{2. อูนปญฺจพนฺธนสิกฺขาปท}{mtk.207}
\bookmark[dest={mtk.207},level=2]{2. อูนปญฺจพนฺธนสิกฺขาปท}
\csromanmarkh{3. ปตฺตวคฺค}\csromanmarkhii{2. อูนปญฺจพนฺธน\-สิกฺขาปท}\csromanheadernomark{2}{3. ปตฺตวคฺค 2. อู\textbf{นปญฺจพนฺธน}\-\textbf{สิกฺขาปท}}
\csromanmatikaanchor{mtk.208}
\csromantocmarkref{subsubsection}{กปิลวตฺถุภิกฺขุวตฺถุ}{mtk.208}
\bookmark[dest={mtk.208},level=3]{กปิลวตฺถุภิกฺขุวตฺถุ}
\proseitem{609}{\csromanfolio{355}เตน สมเยน พุทฺโธ ภควา สกฺเกสุ วิหรติ กปิล\-วตฺถุสฺมิํ นิโคฺรธาราเม. เตน โข ปน สมเยน อญฺญตเรน กุมฺภกาเรน ภิกฺขู ปวาริตา โหนฺติ “เยสํ อยฺยานํ ปตฺเตน อตฺโถ อหํ ปตฺเตนา”ติ. เตน โข ปน สมเยน ภิกฺขู น มตฺตํ ชานิตฺวา พหู ปตฺเต วิญฺญาเปนฺติ. เยสํ ขุทฺทกา ปตฺตา, เต มหนฺเต ปตฺเต วิญฺญาเปนฺติ. เยสํ มหนฺตา ปตฺตา, เต ขุทฺทเก ปตฺเต วิญฺญาเปนฺติ. อถ โข โส กุมฺภกาโร ภิกฺขูนํ พหู ปตฺเต กโรนฺโต น สกฺโกติ อญฺญํ วิกฺกายิกํ ภณฺฑํ กาตุํ. อตฺตนาปิ น ยาเปติ, ปุตฺตทาราปิสฺส กิลมนฺติ. มนุสฺสา อุชฺฌายนฺติ ขิยฺยนฺติ วิปาเจนฺติ “กถํ หิ นาม สมณา สกฺยปุตฺติยา น มตฺตํ ชานิตฺวา พหู ปตฺเต วิญฺญาเปสฺสนฺติ, อยํ อิเมสํ พหู ปตฺเต กโรนฺโต น สกฺโกติ อญฺญํ วิกฺกายิกํ ภณฺฑํ กาตุํ, อตฺตนาปิ น ยาเปติ ปุตฺตทาราปิสฺส กิลมนฺตี”ติ.}

\prose{อสฺโสสุํ โข ภิกฺขู เตสํ มนุสฺสานํ อุชฺฌายนฺตานํ ขิยฺยนฺตานํ วิปาเจนฺตานํ. เย เต ภิกฺขู อปฺปิจฺฉา ฯเปฯ เต อุชฺฌายนฺติ ขิยฺยนฺติ วิปาเจนฺติ “กถํ หิ นาม ภิกฺขู น มตฺตํ ชานิตฺวา พหู ปตฺเต วิญฺญาเปสฺสนฺตี”ติ. อถ โข เต ภิกฺขู เต อเนก\-ปริยาเยน วิครหิตฺวา ภควโต เอตมตฺถํ อาโรเจสุํ ฯเปฯ “สจฺจํ กิร ภิกฺขเว ภิกฺขู น มตฺตํ ชานิตฺวา พหู ปตฺเต วิญฺญาเปนฺตี”ติ. สจฺจํ ภควาติ. วิครหิ พุทฺโธ ภควา ฯเปฯ. กถํ หิ นาม เต ภิกฺขเว โมฆปุริสา น มตฺตํ ชานิตฺวา พหู ปตฺเต วิญฺญาเปสฺสนฺติ, เนตํ ภิกฺขเว อปฺปสนฺนานํ วา ปสาทาย ฯเปฯ วิครหิตฺวา ธมฺมิํ กถํ กตฺวา ภิกฺขู อามนฺเตสิ “น ภิกฺขเว ปตฺโต วิญฺญาเปตพฺโพ, โย วิญฺญาเปยฺย, อาปตฺติ ทุกฺกฏสฺสา”ติ.}

\csromanmatikaanchor{mtk.209}
\csromanmatikaanchor{mtk.210}
\csromantocmarkref{subsubsection}{ปญฺญตฺติ}{mtk.209}
\bookmark[dest={mtk.209},level=3]{ปญฺญตฺติ}
\csromantocmarkref{subsubsection}{สิกฺขาปทวิภงฺค, ปทภาชนีย}{mtk.210}
\bookmark[dest={mtk.210},level=3]{สิกฺขาปทวิภงฺค, ปทภาชนีย}
\proseitem{610}{เตน โข ปน สมเยน อญฺญตรสฺส ภิกฺขุโน ปตฺโต ภินฺโน โหติ. อถ โข โส ภิกฺขุ ภควตา ปฏิกฺขิตฺตํ ปตฺตํ วิญฺญาเปตุนฺติ กุกฺกุจฺจายนฺโต น วิญฺญาเปติ, หตฺเถสุ ปิณฺฑาย จรติ. มนุสฺสา อุชฺฌายนฺติ ขิยฺยนฺติ วิปาเจนฺติ “กถํ หิ นาม สมณา สกฺยปุตฺติยา หตฺเถสุ ปิณฺฑาย จริสฺสนฺติ เสยฺยถาปิ ติตฺถิยา”ติ. อสฺโสสุํ โข ภิกฺขู เตสํ มนุสฺสานํ อุชฺฌายนฺตานํ ขิยฺยนฺตานํ วิปาเจนฺตานํ. อถ โข เต ภิกฺขู ภควโต เอตมตฺถํ อาโรเจสุํ. อถ โข ภควา เอตสฺมิํ \csromanfolio{356}นิทาเน เอตสฺมิํ ปกรเณ ธมฺมิํ กถํ กตฺวา ภิกฺขู อามนฺเตสิ “อนุชานามิ ภิกฺขเว นฏฺฐปตฺตสฺส วา ภินฺนปตฺตสฺส วา ปตฺตํ วิญฺญาเปตุนฺ”ติ.}

\proseitem{611}{เตน โข ปน สมเยน ฉพฺพคฺคิยา ภิกฺขู ภควตา อนุญฺญาตํ “นฏฺฐปตฺตสฺส วา ภินฺนปตฺตสฺส วา ปตฺตํ วิญฺญาเปตุนฺ”ติ อปฺป\-มตฺต\-เกนปิ ภินฺเนน อปฺป\-มตฺต\-เกนปิ ขณฺเฑน วิลิขิตมตฺเตนปิ พหู ปตฺเต วิญฺญาเปนฺติ. อถ โข โส กุมฺภกาโร ภิกฺขูนํ ตเถว พหู ปตฺเต กโรนฺโต น สกฺโกติ อญฺญํ วิกฺกายิกํ ภณฺฑํ กาตุํ, อตฺตนาปิ น ยาเปติ, ปุตฺตทาราปิสฺส กิลมนฺติ. มนุสฺสา ตเถว อุชฺฌายนฺติ ขิยฺยนฺติ วิปาเจนฺติ “กถํ หิ นาม สมณา สกฺยปุตฺติยา น มตฺตํ ชานิตฺวา พหู ปตฺเต วิญฺญาเปสฺสนฺติ, อยํ อิเมสํ พหู ปตฺเต กโรนฺโต น สกฺโกติ อญฺญํ วิกฺกายิกํ ภณฺฑํ กาตุํ, อตฺตนาปิ น ยาเปติ, ปุตฺตทาราปิสฺส กิลมนฺตี”ติ.}

\prose{อสฺโสสุํ โข ภิกฺขู เตสํ มนุสฺสานํ อุชฺฌายนฺตานํ ขิยฺยนฺตานํ วิปาเจนฺตานํ. เย เต ภิกฺขู อปฺปิจฺฉา, เต อุชฺฌายนฺติ ขิยฺยนฺติ วิปาเจนฺติ “กถํ หิ นาม ฉพฺพคฺคิยา ภิกฺขู อปฺป\-มตฺต\-เกนปิ ภินฺเนน อปฺป\-มตฺต\-เกนปิ ขณฺเฑน วิลิขิตมตฺเตนปิ พหู ปตฺเต วิญฺญาเปสฺสนฺตี”ติ. อถ โข เต ภิกฺขู ฉพฺพคฺคิเย ภิกฺขู อเนก\-ปริยาเยน วิครหิตฺวา ภควโต เอตมตฺถํ อาโรเจสุํ ฯเปฯ “สจฺจํ กิร ตุเมฺห ภิกฺขเว อปฺป\-มตฺต\-เกนปิ ภินฺเนน อปฺป\-มตฺต\-เกนปิ ขณฺเฑน วิลิขิตมตฺเตนปิ พหู ปตฺเต วิญฺญาเปถา”ติ. สจฺจํ ภควาติ. วิครหิ พุทฺโธ ภควา ฯเปฯ. กถํ หิ นาม ตุเมฺห โมฆปุริสา อปฺป\-มตฺต\-เกนปิ ภินฺเนน อปฺป\-มตฺต\-เกนปิ ขณฺเฑน วิลิขิตมตฺเตนปิ พหู ปตฺเต วิญฺญาเปสฺสถ, เนตํ โมฆปุริสา อปฺปสนฺนานํ วา ปสาทาย ฯเปฯ. เอวญฺจ ปน ภิกฺขเว อิมํ สิกฺขาปทํ อุทฺทิเสยฺยาถ\csromandash{}}

\proseitem{612}{\textbf{“โย ปน ภิกฺขุ อูนปญฺจ}\-\textbf{พนฺธเนน ปตฺเตน อญฺญํ นวํ ปตฺตํ เจตาเปยฺย, นิสฺสคฺคิยํ ปาจิตฺติยํ. เตน ภิกฺขุนา โส ปตฺโต ภิกฺขุ}\-\textbf{ปริสาย นิสฺสชฺชิตพฺโพ, โย จ ตสฺสา ภิกฺขุ}\-\textbf{ปริสาย ปตฺตปริยนฺโต, โส ตสฺส ภิกฺขุโน ปทาตพฺโพ ‘อยํ เต ภิกฺขุ ปตฺโต ยาว เภทนาย ธาเรตพฺโพ’ติ, อยํ ตตฺถ สามีจี”}ติ.}

\proseitem{613}{\textbf{โย ปนา}ติ โย ยาทิโส ฯเปฯ. \textbf{ภิกฺขู}ติ ฯเปฯ อยํ อิมสฺมิํ อตฺเถ อธิปฺเปโต “ภิกฺขู”ติ.}

\prose{\csromanfolio{357}อู\textbf{นปญฺจ}\-\textbf{พนฺธโน} นาม ปตฺโต อพนฺธโน วา เอกพนฺธโน วา ทฺวิพนฺธโน วา ติพนฺธโน วา จตุพนฺธโน วา. อพนฺธโนกาโส นาม ปตฺโต ยสฺส ทฺวงฺคุลา ราชิ น โหติ. พนฺธโนกาโส นาม ปตฺโต ยสฺส ทฺวงฺคุลา ราชิ โหติ. \textbf{นโว} นาม ปตฺโต วิญฺญตฺติํ อุปาทาย วุจฺจติ.}

\prose{\textbf{เจตาเปยฺยา}ติ วิญฺญาเปติ, ปโยเค ทุกฺกฏํ. ปฏิลาเภน นิสฺสคฺคิโย โหติ. สํฆมชฺเฌ นิสฺสชฺชิตพฺโพ. สพฺเพเหว อธิฏฺฐิต\-ปตฺตํ คเหตฺวา สนฺนิปติตพฺพํ. น ลามโก ปตฺโต อธิฏฺฐาตพฺโพ “มหคฺฆํ ปตฺตํ คเหสฺสามี”ติ. สเจ ลามกํ ปตฺตํ อธิฏฺเฐติ “มหคฺฆํ ปตฺตํ คเหสฺสามี”ติ, อาปตฺติ ทุกฺกฏสฺส. เอวญฺจ ปน ภิกฺขเว นิสฺสชฺชิตพฺโพ. เตน ภิกฺขุนา สํฆํ อุปสงฺกมิตฺวา เอกํสํ อุตฺตราสงฺคํ กริตฺวา วุฑฺฒานํ ภิกฺขูนํ ปาเท วนฺทิตฺวา อุกฺกุฏิกํ นิสีทิตฺวา อญฺชลิํ ปคฺคเหตฺวา เอวมสฺส วจนีโย “อยํ เม ภนฺเต ปตฺโต อูนปญฺจ\-พนฺธเนน ปตฺเตน เจตาปิโต นิสฺสคฺคิโย, อิมาหํ สํฆสฺส นิสฺสชฺชามี”ติ. นิสฺสชฺชิตฺวา อาปตฺติ เทเสตพฺพา, พฺยตฺเตน ภิกฺขุนา ปฏิพเลน อาปตฺติ ปฏิคฺคเหตพฺพา. ปญฺจหงฺเคหิ สมนฺนาคโต ภิกฺขุ ปตฺตคฺคาหาปโก สมฺมนฺนิตพฺโพ, โย น ฉนฺทาคติํ คจฺเฉยฺย, น โทสาคติํ คจฺเฉยฺย, น โมหาคติํ คจฺเฉยฺย, น ภยาคติํ คจฺเฉยฺย, คาหิตาคาหิตญฺจ ชาเนยฺย. เอวญฺจ ปน ภิกฺขเว สมฺมนฺนิตพฺโพ. ปฐมํ ภิกฺขุ ยาจิตพฺโพ, ยาจิตฺวา พฺยตฺเตน ภิกฺขุนา ปฏิพเลน สํโฆ ญาเปตพฺโพ\csromandash{}}

\proseitem{614}{“สุณาตุ เม ภนฺเต สํโฆ, ยทิ สํฆสฺส ปตฺตกลฺลํ, สํโฆ อิตฺถนฺนามํ ภิกฺขุํ ปตฺต\-คฺคาหาปกํ สมฺมนฺเนยฺย, เอสา ญตฺติ.}

\prose{สุณาตุ เม ภนฺเต สํโฆ, สํโฆ อิตฺถนฺนามํ ภิกฺขุํ ปตฺต\-คฺคาหาปกํ สมฺมนฺนติ, ยสฺสายสฺมโต ขมติ อิตฺถนฺนามสฺส ภิกฺขุโน ปตฺต\-คฺคาหาปกสฺส สมฺมุติ, โส ตุณฺหสฺส, ยสฺส นกฺขมติ, โส ภาเสยฺย.}

\prose{สมฺมโต สํเฆน อิตฺถนฺนาโม ภิกฺขุ ปตฺตคฺคาหาปโก, ขมติ สํฆสฺส, ตสฺมา ตุณฺหี, เอวเมตํ ธารยามี”ติ.}

\proseitem{615}{เตน สมฺมเตน ภิกฺขุนา ปตฺโต คาเหตพฺโพ. เถโร วตฺตพฺโพ “คณฺหาตุ ภนฺเต เถโร ปตฺตนฺ”ติ. สเจ เถโร คณฺหาติ, \csromanfolio{358}เถรสฺส ปตฺโต ทุติยสฺส คาเหตพฺโพ. น จ ตสฺส อนุทฺทยตาย น คเหตพฺโพ, โย น คเณฺหยฺย, อาปตฺติ ทุกฺกฏสฺส. อปตฺตกสฺส น คาเหตพฺโพ. เอเตเนว อุปาเยน ยาว สํฆนวกา คาเหตพฺโพ. โย จ ตสฺสา ภิกฺขุ\-ปริสาย ปตฺตปริยนฺโต, โส ตสฺส ภิกฺขุโน ปทาตพฺโพ “อยํ เต ภิกฺขุ ปตฺโต ยาว เภทนาย ธาเรตพฺโพ”ติ.}

\prose{เตน ภิกฺขุนา โส ปตฺโต น อเทเส นิกฺขิปิตพฺโพ, น อโภเคน ภุญฺชิตพฺโพ, น วิสฺสชฺเชตพฺโพ “กถายํ ปตฺโต นสฺเสยฺย วา วินสฺเสยฺย วา ภิชฺเชยฺยวา”ติ. สเจ อเทเส วา นิกฺขิปติ อโภเคน วา ภุญฺชติ วิสฺสชฺเชติ วา, อาปตฺติ ทุกฺกฏสฺส.}

\prose{\textbf{อยํ ตตฺถ สามีจี}ติ อยํ ตตฺถ อนุธมฺมตา.}

\proseitem{616}{อพนฺธเนน ปตฺเตน อพนฺธนํ ปตฺตํ เจตาเปติ, นิสฺสคฺคิยํ ปาจิตฺติยํ. อพนฺธเนน ปตฺเตน เอกพนฺธนํ ปตฺตํ เจตาเปติ, นิสฺสคฺคิยํ ปาจิตฺติยํ. อพนฺธเนน ปตฺเตน ทฺวิพนฺธนํ ปตฺตํ เจตาเปติ, นิสฺสคฺคิยํ ปาจิติยํ. อพนฺธเนน ปตฺเตน ติพนฺธนํ ปตฺตํ เจตาเปติ, นิสฺสคฺคิยํ ปาจิตฺติยํ. อพนฺธเนน ปตฺเตน จตุพนฺธนํ ปตฺตํ เจตาเปติ, นิสฺสคฺคิยํ ปาจิตฺติยํ.}

\prose{เอกพนฺธเนน ปตฺเตน อพนฺธนํ ปตฺตํ เจตาเปติ, นิสฺสคฺคิยํ ปาจิตฺติยํ. เอกพนฺธเนน ปตฺเตน เอกพนฺธนํ ปตฺตํ เจตาเปติ, นิสฺสคฺคิยํ ปาจิตฺติยํ. เอกพนฺธเนน ปตฺเตน ทฺวิพนฺธนํ ปตฺตํ เจตาเปติ, นิสฺสคฺคิยํ ปาจิตฺติยํ. เอกพนฺธเนน ปตฺเตน ติพนฺธนํ ปตฺตํ เจตาเปติ, นิสฺสคฺคิยํ ปาจิตฺติยํ. เอกพนฺธเนน ปตฺเตน จตุพนฺธนํ ปตฺตํ เจตาเปติ, นิสฺสคฺคิยํ ปาจิตฺติยํ.}

\prose{ทฺวิพนฺธเนน ปตฺเตน อพนฺธนํ ปตฺตํ เจตาเปติ, นิสฺสคฺคิยํ ปาจิตฺติยํ. ทฺวิพนฺธเนน ปตฺเตน เอกพนฺธนํ ปตฺตํ ฯเปฯ ทฺวิพนฺธนํ ปตฺตํ. ติพนฺธนํ ปตฺตํ. จตุพนฺธนํ ปตฺตํ เจตาเปติ, นิสฺสคฺคิยํ ปาจิตฺติยํ.}

\prose{ติพนฺธเนน ปตฺเตน อพนฺธนํ ปตฺตํ ฯเปฯ เอกพนฺธนํ ปตฺตํ. ทฺวิพนฺธนํ ปตฺตํ. ติพนฺธนํ ปตฺตํ. จตุพนฺธนํ ปตฺตํ เจตาเปติ, นิสฺสคฺคิยํ ปาจิตฺติยํ.}

\prose{จตุพนฺธเนน ปตฺเตน อพนฺธนํ ปตฺตํ ฯเปฯ เอกพนฺธนํ ปตฺตํ. ทฺวิพนฺธนํ ปตฺตํ. ติพนฺธนํ ปตฺตํ. จตุพนฺธนํ ปตฺตํ เจตาเปติ, นิสฺสคฺคิยํ ปาจิตฺติยํ.}

\prose{อพนฺธเนน ปตฺเตน อพนฺธโนกาสํ ปตฺตํ เจตาเปติ, นิสฺสคฺคิยํ ปาจิตฺติยํ. อพนฺธเนน ปตฺเตน เอก\-พนฺธโนกาสํ ปตฺตํ ฯเปฯ ทฺวิพนฺธโนกาสํ \csromanfolio{359}\textbf{ปตฺตํ. ติพนฺธโนกาสํ ปตฺตํ. จตุพนฺธโนกาสํ ปตฺตํ เจตาเปติ,} นิสฺสคฺคิยํ ปาจิตฺติยํ.}

\prose{เอกพนฺธเนน ปตฺเตน อพนฺธโนกาสํ ปตฺตํ ฯเปฯ เอก\-พนฺธโนกาสํ ปตฺตํ. ทฺวิพนฺธโนกาสํ ปตฺตํ. ติพนฺธโนกาสํ ปตฺตํ. จตุพนฺธโนกาสํ ปตฺตํ เจตาเปติ, นิสฺสคฺคิยํ ปาจิตฺติยํ.}

\prose{ทฺวิพนฺธเนน ปตฺเตน อพนฺธโนกาสํ ปตฺตํ ฯเปฯ จตุพนฺธโนกาสํ ปตฺตํ เจตาเปติ, นิสฺสคฺคิยํ ปาจิตฺติยํ.}

\prose{ติพนฺธเนน ปตฺเตน อพนฺธโนกาสํ ปตฺตํ ฯเปฯ จตุพนฺธโนกาสํ ปตฺตํ เจตาเปติ, นิสฺสคฺคิยํ ปาจิตฺติยํ.}

\prose{จตุพนฺธเนน ปตฺเตน อพนฺธโนกาสํ ปตฺตํ ฯเปฯ เอก\-พนฺธโนกาสํ ปตฺตํ. ทฺวิพนฺธโนกาสํ ปตฺตํ. ติพนฺธโนกาสํ ปตฺตํ. จตุพนฺธโนกาสํ ปตฺตํ เจตาเปติ, นิสฺสคฺคิยํ ปาจิตฺติยํ.}

\prose{อพนฺธโนกาเสน ปตฺเตน อพนฺธนํ ปตฺตํ เจตาเปติ, นิสฺสคฺคิยํ ปาจิตฺติยํ. อพนฺธโนกาเสน ปตฺเตน เอกพนฺธนํ ปตฺตํ ฯเปฯ ทฺวิพนฺธนํ ปตฺตํ. ติพนฺธนํ ปตฺตํ. จตุพนฺธนํ ปตฺตํ เจตาเปติ, นิสฺสคฺคิยํ ปาจิตฺติยํ.}

\prose{จตุพนฺธโนกาเสน ปตฺเตน อพนฺธนํ ปตฺตํ เจตาเปติ, นิสฺสคฺคิยํ ปาจิตฺติยํ. จตุพนฺธโนกาเสน ปตฺเตน เอกพนฺธนํ ปตฺตํ ฯเปฯ ทฺวิพนฺธนํ ปตฺตํ. ติพนฺธนํ ปตฺตํ. จตุพนฺธนํ ปตฺตํ เจตาเปติ, นิสฺสคฺคิยํ ปาจิตฺติยํ.}

\prose{อพนฺธโนกาเสน ปตฺเตน อพนฺธโนกาสํ ปตฺตํ ฯเปฯ เอก\-พนฺธโนกาสํ ปตฺตํ. ทฺวิพนฺธโนกาสํ ปตฺตํ. ติพนฺธโนกาสํ ปตฺตํ. จตุพนฺธโนกาสํ ปตฺตํ เจตาเปติ, นิสฺสคฺคิยํ ปาจิตฺติยํ.}

\prose{จตุพนฺธโนกาเสน ปตฺเตน อพนฺธโนกาสํ ปตฺตํ ฯเปฯ เอก\-พนฺธโนกาสํ ปตฺตํ. ทฺวิพนฺธโนกาสํ ปตฺตํ. ติพนฺธโนกาสํ ปตฺตํ. จตุพนฺธโนกาสํ ปตฺตํ เจตาเปติ, นิสฺสคฺคิยํ ปาจิตฺติยํ.}

\proseitem{617}{อนาปตฺติ นฏฺฐปตฺตสฺส ภินฺนปตฺตสฺส ญาตกานํ ปวาริตานํ อญฺญสฺสตฺถาย อตฺตโน ธเนน อุมฺมตฺตกสฺส อาทิกมฺมิกสฺสาติ.}

\vspace{\tipitakaparskip}
\csromancenter{อูนปญฺจพนฺธน\-สิกฺขาปทํ นิฏฺฐิตํ ทุติยํ.}
\csromanmatikaanchor{mtk.211}
\csromantocmarknopage{subsection}{3. เภสชฺชสิกฺขาปท}{mtk.211}
\bookmark[dest={mtk.211},level=2]{3. เภสชฺชสิกฺขาปท}
\csromanmarkh{3. ปตฺตวคฺค}\csromanmarkhii{3. เภสชฺช\-สิกฺขาปท}\csromanheadernomark{2}{3. \textbf{ปตฺตวคฺค 3. เภสชฺช}\-\textbf{สิกฺขาปท}}
\csromanmatikaanchor{mtk.212}
\csromantocmarkref{subsubsection}{ปิลินฺทวจฺฉภิกฺขุวตฺถุ}{mtk.212}
\bookmark[dest={mtk.212},level=3]{ปิลินฺทวจฺฉภิกฺขุวตฺถุ}
\proseitem{618}{\csromanfolio{360}\csromansymbolfootnote{*}{อิทํ วตฺถุ วิ 3. 299 ปิฏฺฐาทีสุปิ อาคตํ.}เตน สมเยน พุทฺโธ ภควา สาวตฺถิยํ วิหรติ เชตวเน อนาถ\-ปิณฺฑิกสฺส อาราเม. เตน โข ปน สมเยน อายสฺมา ปิลินฺทวจฺโฉ ราชคเห ปพฺภารํ โส ธาเปติ เลณํ กตฺตุกาโม. อถ โข ราชา มาคโธ เสนิโย พิมฺพิสาโร เยนายสฺมา ปิลินฺทวจฺโฉ เตนุปสงฺกมิ, อุปสงฺกมิตฺวา อายสฺมนฺตํ ปิลินฺทวจฺฉํ อภิวาเทตฺวา เอกมนฺตํ นิสีทิ. เอกมนฺตํ นิสินฺโน โข ราชา มาคโธ เสนิโย พิมฺพิสาโร อายสฺมนฺตํ ปิลินฺทวจฺฉํ เอตทโวจ “กิํ ภนฺเต เถโร การาเปตี”ติ. ปพฺภารํ มหาราช โสธาเปมิ เลณํ กตฺถุกาโมติ. อตฺโถ ภนฺเต อยฺยสฺส อารามิเกนาติ. น โข มหาราช ภควตา อารามิโก อนุญฺญาโตติ. เตน หิ ภนฺเต ภควนฺตํ ปฏิปุจฺฉิตฺวา มม อาโรเจยฺยาถาติ. “เอวํ มหาราชา”ติ โข อายสฺมา ปิลินฺทวจฺโฉ รญฺโญ มาคธสฺส เสนิยสฺส พิมฺพิสารสฺส ปจฺจสฺโสสิ. อถ โข อายสฺมา ปิลินฺทวจฺโฉ ราชานํ มาคธํ เสนิยํ พิมฺภิสารํ ธมฺมิยา กถาย สนฺทสฺเสสิ สมาทเปสิ สมุตฺเตเชสิ สมฺปหํเสสิ. อถ โข ราชา มาคโธ เสนิโย พิมฺพิสาโร อายสฺมตา ปิลินฺท\-วจฺเฉน ธมฺมิยา กถาย สนฺทสฺสิโต สมาทปิโต สมุตฺเตชิโต สมฺปหํสิโต อุฏฺฐายาสนา อายสฺมนฺตํ ปิลินฺทวจฺฉํ อภิวาเทตฺวา ปทกฺขิณํ กตฺวา ปกฺกามิ.}

\proseitem{619}{อถ โข อายสฺมา ปิลินฺทวจฺโฉ ภควโต สนฺติเก ทูตํ ปาเหสิ “ราชา ภนฺเต มาคโธ เสนิโย พิมฺพิสาโร อารามิกํ ทาตุกาโม, กถํ นุ โข ภนฺเต มยา ปฏิปชฺชิตพฺพนฺ”ติ. อถ โข ภควา เอตสฺมิํ นิทาเน เอตสฺมิํ ปกรเณ ธมฺมิํ กถํ กตฺวา ภิกฺขู อามนฺเตสิ “อนุชานามิ ภิกฺขเว อารามิกนฺ”ติ. ทุติยมฺปิ โข ราชา มาคโธ เสนิโย พิมฺพิสาโร เยนายสฺมา ปิลินฺทวจฺโฉ เตนุปสงฺกมิ, อุปสงฺกมิตฺวา อายสฺมนฺตํ ปิลินฺทวจฺฉํ อภิวาเทตฺวา เอกมนฺตํ นิสีทิ. เอกมนฺตํ นิสินฺโน โข ราชา มาคโธ เสนิโย พิมฺพิสาโร อายสฺมนฺตํ ปิลินฺทวจฺฉํ เอตทโวจ “อนุญฺญาโต ภนฺเต ภควตา อารามิโก”ติ. เอวํ มหาราชาติ. เตน หิ ภนฺเต อยฺยสฺส อารามิกํ ทมฺมีติ. อถ โข ราชา มาคโธ เสนิโย พิมฺพิสาโร อายสฺมโต ปิลินฺท\-วจฺฉสฺส \csromanfolio{361}อารามิกํ ปฏิสฺสุณิตฺวา วิสฺสริตฺวา จิเรน สติํ ปฏิลภิตฺวา อญฺญตรํ สพฺพตฺถกํ มหามตฺตํ อามนฺเตสิ “โย มยา ภเณ อยฺยสฺส อารามิโก ปฏิสฺสุโต, ทินฺโน โส อารามิโก”ติ. น โข เทว อยฺยสฺส อารามิโก ทินฺโนติ. กีวจิรํ นุ โข ภเณ อิโต หิ ตํ โหตีติ. อถ โข โส มหามตฺโต รตฺติโย คเณตฺวา\csromansharedfootnote{344bed4c41f70cef}{วิคเณตฺวา (ก)} ราชานํ มาคธํ เสนิยํ พิมฺพิสารํ เอตทโวจ “ปญฺจ เทว รตฺติสตานี”ติ. เตน หิ ภเณ อยฺยสฺส ปญฺจ อารามิกสตานิ เทหีติ. “เอวํ เทวา”ติ โข โส มหามตฺโต รญฺโญ มาคธสฺส เสนิยสฺส พิมฺพิสารสฺส ปฏิสฺสุณิตฺวา อายสฺมโต ปิลินฺท\-วจฺฉสฺส ปญฺจ อารามิกสนานิ ปาทาสิ\csromansharedfootnote{5b0c729655da8e57}{อทาสิ (สฺยา)}. ปาฏิเยกฺโก คาโม นิวิสิ. อารามิกคามโกติปิ นํ อาหํสุ, ปิลินฺทคามโกติปิ นํ อาหํสุ.}

\proseitem{620}{เตน โข ปน สมเยน อายสฺมา ปิลินฺทวจฺโฉ ตสฺมิํ คามเก กุลูปโก โหติ. อถ โข อายสฺมา ปิลินฺทวจฺโฉ ปุพฺพณฺหสมยํ นิวาเสตฺวา ปตฺตจีวรํ อาทาย ปิลินฺทคามกํ ปิณฺฑาย ปาวิสิ. เตน โข ปน สมเยน ตสฺมิํ คามเก อุสฺสโว โหติ, ทารกา อลงฺกตา มาลากิตา กีฬนฺติ. อถ โข อายสฺมา ปิลินฺทวจฺโฉ ปิลินฺทคามเก สปทานํ ปิณฺฑาย จรมาโน เยน อญฺญตรสฺส อารามิกสฺส นิเวสนํ เตนุปสงฺกมิ, อุปสงฺกมิตฺวา ปญฺญตฺเต อาสเน นิสีทิ. เตน โข ปน สมเยน ตสฺสา อารามิกินิยา ธีตา อญฺเญ ทารเก อลงฺกเต มาลากิเต ปสฺสิตฺวา โรทติ “มาลํ เม เทถ, อลงฺการํ เม เทถา”ติ. อถ โข อายสฺมา ปิลินฺทวจฺโฉ ตํ อารามิกินิํ เอตทโวจ “กิสฺสายํ ทาริกา โรทตี”ติ. อยํ ภนฺเต ทาริกา อญฺเญ ทารเก อลงฺกเต มาลากิเต ปสฺสิตฺวา โรทติ “มาลํ เม เทถ, อลงฺการํ เม เทถา”ติ. กุโต อมฺหากํ ทุคฺคตานํ มาลา, กุโต อลงฺกาโรติ. อถ โข อายสฺมา ปิลินฺทวจฺโฉ อญฺญตรํ ติณณฺฑุปกํ คเหตฺวา ตํ อารามิกินิํ เอตทโวจ “หนฺทิมํ ติณณฺฑุปกํ ตสฺสา ทาริกาย สีเส ปฏิมุญฺจา”ติ. อถ โข สา อารามิกินี ตํ ติณณฺฑุปกํ คเหตฺวา ตสฺสา ทาริกาย สีเส ปฏิมุญฺจิ. สา อโหสิ สุวณฺณมาลา อภิรูปา ทสฺสนียา ปาสาทิกา. นตฺถิ ตาทิสา รญฺโญปิ อนฺเตปุเร สุวณฺณมาลา. มนุสฺสา รญฺโญ มาคธสฺส เสนิยสฺส พิมฺพิสารสฺส \csromanfolio{362}อาโรเจสุํ “อมุกสฺส เทว อารามิกสฺส ฆเร สุวณฺณมาลา อภิรูปา ทสฺสนียา ปาสาทิกา, นตฺถิ ตาทิสา เทวสฺสาปิ อนฺเตปุเร สุวณฺณมาลา, กุโต ตสฺส ทุคฺคตสฺส, นิสฺสํสยํ โจริกาย อาภตา”ติ. อถ โข ราชา มาคโธ เสนิโย พิมฺภิสาโร ตํ อารามิกกุลํ พนฺธาเปสิ. ทุติยมฺปิ โข อายสฺมา ปิลินฺทวจฺโฉ ปุพฺพณฺหสมยํ นิวาเสตฺวา ปตฺตจีวรํ อาทาย ปิลินฺทคามกํ ปิณฺฑาย ปาวิสิ. ปิลินฺทคามเก สปทานํ ปิณฺฑาย จรมาโน เยน ตสฺส อารามิกสฺส นิเวสนํ เตนุปสงฺกมิ, อุปสงฺกมิตฺวา ปฏิวิสฺสเก ปุจฺฉิ “กหํ อิมํ อารามิกกุลํ คตนฺ”ติ. เอติสฺสา ภนฺเต สุวณฺณมาลาย การณา รญฺญา พนฺธาปิตนฺติ.}

\csromanmatikaanchor{mtk.213}
\csromanmatikaanchor{mtk.214}
\csromantocmarkref{subsubsection}{ปญฺญตฺติ}{mtk.213}
\bookmark[dest={mtk.213},level=3]{ปญฺญตฺติ}
\csromantocmarkref{subsubsection}{สิกฺขาปทวิภงฺค, ปทภาชนีย}{mtk.214}
\bookmark[dest={mtk.214},level=3]{สิกฺขาปทวิภงฺค, ปทภาชนีย}
\proseitem{621}{อถ โข อายสฺมา ปิลินฺทวจฺโฉ เยน รญฺโญ มาคธสฺส เสนิยสฺส พิมฺพิสารสฺส นิเวสนํ เตนุปสงฺกมิ, อุปสงฺกมิตฺวา ปญฺญตฺเต อาสเน นิสีทิ. อถ โข ราชา มาคโธ เสนิโย พิมฺพิสาโร เยนายสฺมา ปิลินฺทวจฺโฉ เตนุปสงฺกมิ, อุปสงฺกมิตฺวา อายสฺมนํ ปิลินฺทวจฺฉํ อภิวาเทตฺวา เอกมนฺตํ นิสีทิ, เอกมนฺตํ นิสินฺนํ โข ราชานํ มาคธํ เสนิยํ พิมฺพิสารํ อายสฺมา ปิลินฺทวจฺโฉ เอตทโวจ “กิสฺส มหาราช อารามิกกุลํ พนฺธาปิตนฺ”ติ. ตสฺส ภนฺเต อารามิกสฺส ฆเร สุวณฺณมาลา อภิรูปา ทสฺสนียา ปาสาทิกา, นตฺถิ ตาทิสา อมฺหากํปิ อนฺเตปุเร สุวณฺณมาลา, กุโต ตสฺส ทุคฺคตสฺส, นิสฺสํสยํ โจริกาย อาภตาติ. อถ โข อายสฺมา ปิลินฺทวจฺโฉ รญฺโญ มาคธสฺส เสนิยสฺส พิมฺพิสารสฺส ปาสาทํ สุวณฺณนฺติ อธิมุจฺจิ. โส อโหสิ สพฺพ\-โสวณฺณมโย. อิทํ ปน เต มหาราช ตาว พหุํ สุวณฺณํ กุโตติ. อญฺญาตํ ภนฺเต, อยฺยสฺเสเวโส อิทฺธานุภาโวติ ตํ อารามิกกุลํ มุญฺจาเปสิ. มนุสฺสา อยฺเยน กิร ปิลินฺท\-วจฺเฉน สราชิกาย ปริสาย อุตฺตริ\-มนุสฺสธมฺมํ อิทฺธิ\-ปาฏิหาริยํ ทสฺสิตนฺติ อตฺตมนา อภิปฺปสนฺนา อายสฺมโต ปิลินฺท\-วจฺฉสฺส ปญฺจ เภสชฺชานิ อภิหริํสุ. เสยฺยถิทํ, สปฺปิ นวนีตํ เตลํ มธุ ผาณิตํ. ปกติยาปิ จ อายสฺมา ปิลินฺทวจฺโฉ ลาภี โหติ ปญฺจนฺนํ เภสชฺชานํ, ลทฺธํ ลทฺธํ ปริสาย วิสฺสชฺเชติ. ปริสา จสฺส โหติ พาหุลฺลิกา ลทฺธํ ลทฺธํ โกลมฺเพปิ ฆเฏปิ ปูเรตฺวา ปฏิสาเมติ. ปริสฺสาว\-นานิปิ ถวิกาโยปิ ปูเรตฺวา วาตปาเนสุ ลคฺเคติ. ตานิ โอลีนวิลีนานิ ติฏฺฐนฺติ. \csromanfolio{363}อุนฺทูเรหิปิ วิหารา โอกิณฺณวิกิณฺณา โหนฺติ. มนุสฺสา วิหารจาริกํ อาหิณฺฑนฺตา ปสฺสิตฺวา อุชฺฌายนฺติ ขิยฺยนฺติ วิปาเจนฺติ “อนฺโตโกฏฺฐาคาริกา อิเม สมณา สกฺยปุตฺติยา เสยฺยถาปิ ราชา มาคโธ เสนิโย พิมฺพิสาโร”ติ. อสฺโสสุํ โข ภิกฺขู เตสํ มนุสฺสานํ อุชฺฌายนฺตานํ ขิยฺยนฺตาณํ วิปาเจนฺตานํ. เย เต ภิกฺขู อปฺปิจฺฉา ฯเปฯ เต อุชฺฌายนฺติ ขิยฺยนฺติ วิปาเจนฺติ “กถํ หิ นาม ภิกฺขู เอวรูปาย พาหุลฺลาย เจเตสฺสนฺตี”ติ. อถ โข เต ภิกฺขู เต อเนก\-ปริยาเยน วิครหิตฺวา ภควโต เอตมตฺถํ อาโรเจสุํ ฯเปฯ “สจฺจํ กิร ภิกฺขเว ภิกฺขู เอวรูปาย พาหุลฺลาย เจเตนฺตี”ติ. สจฺจํ ภควาติ. วิครหิ พุทฺโธ ภควา ฯเปฯ. กถํ หิ นาม เต ภิกฺขเว โมฆปุริสา เอวรูปาย พาหุลฺลาย เจเตสฺสนฺติ, เนตํ ภิกฺขเว อปฺปสนฺนานํ วา ปสาทาย ฯเปฯ. เอวญฺจ ปน ภิกฺขเว อิมํ สิกฺขาปทํ อุทฺทิเสยฺยาถ\csromandash{}}

\proseitem{622}{\textbf{“ยานิ โข ปน ตานิ คิลานานํ ภิกฺขูนํ ปฏิสายนียานิ เภสชฺชานิ, เสยฺยถิทํ, สปฺปิ นวนีตํ เตลํ มธุ ผาณิตํ, ตานิ ปฏิคฺคเหตฺวา อตฺตาหปรมํ สนฺนิธิการกํ ปริภุญฺชิ}\-\textbf{ตพฺพานิ, ตํ อติกฺกามยโต นิสฺสคฺคิยํ ปาจิตฺติยนฺ”}ติ.}

\proseitem{623}{\textbf{ยานิ โข ปน ตานิ คิลานานํ ภิกฺขูนํ ปฏิสายนียานิ} เภสชฺชานีติ สปฺปิ นาม โคสปฺปิ วา อชิกาสปฺปิ วา มหิํสสปฺปิ\csromansharedfootnote{696f515e7e16fb01}{มหิสสปฺปิ (สี, สฺยา)} วา เยสํ มํสํ กปฺปติ เตสํ สปฺปิ. นวนีตํ นาม เตสํ เยว\csromansharedfootnote{9a3af51fcf41c94d}{เตสํเยว นิคฺคหีตสนฺธิ \csromandash{} ม.พ.ป.} นวนีตํ. เตลํ นาม ติลเตลํ สาสปเตลํ มธุกเตลํ เอรณฺฑเตลํ วสาเตลํ. มธุ นาม มกฺขิกามธุ. ผาณิตํ นาม อุจฺฉุมฺหา นิพฺพตฺตํ.}

\prose{\textbf{ตานิ ปฏิคฺคเหตฺวา สตฺตาหปรมํ สนฺนิธิการกํ ปริภุญฺชิ}\-\textbf{ตพฺพานีติ} สตฺตาห\-ปรมตา ปริภุญฺชิ\-ตพฺพานิ.}

\prose{\textbf{ตํ อติกฺกามยโต นิสฺสคฺคิยํ โหตี}ติ อฏฺฐเม อรุณุคฺคมเน นิสฺสคฺคิยํ โหติ. นิสฺสชฺชิตพฺพํ สํฆสฺส วา คณสฺส วา ปุคฺคลสฺส วา. เอวญฺจ ปน ภิกฺกเว นิสฺสชฺชิตพฺพํ ฯเปฯ “อิทํ เม ภนฺเต เภสชฺชํ สตฺตาหา\-ติกฺกนฺตํ นิสฺสคฺคิยํ, อิมาหํ สํฆสฺส นิสฺสชฺชามี”ติ ฯเปฯ ทเทยฺยาติ ฯเปฯ ทเทยฺยุนฺติ ฯเปฯ อายสฺมโต ทมฺมีติ.}

\csromanmatikaanchor{mtk.215}
\csromanmatikaanchor{mtk.216}
\csromantocmarknopage{subsection}{4. วสฺสิกสาฏิกสิกฺขาปท}{mtk.215}
\bookmark[dest={mtk.215},level=2]{4. วสฺสิกสาฏิกสิกฺขาปท}
\csromantocmarkref{subsubsection}{ฉพฺพคฺคิยภิกฺขุวตฺถุ}{mtk.216}
\bookmark[dest={mtk.216},level=3]{ฉพฺพคฺคิยภิกฺขุวตฺถุ}
\proseitem{624}{\csromanfolio{364}สตฺตาหา\-ติกฺกนฺเต อติกฺกนฺตสญฺญี, นิสฺสคฺคิยํ ปาจิตฺติยํ. สตฺตาหา\-ติกฺกนฺเต เวมติโก, นิสฺสคฺคิยํ ปาจิตฺติยํ. สตฺตาหาติกฺกานฺเต อนติกฺกนฺต\-สญฺญี, นิสฺสคฺคิยํ ปาจิตฺติยํ. อนธิฏฺฐิเต อธิฏฺฐิต\-สญฺญี, นิสฺสคฺคิยํ ปาจิตฺติยํ. อวิสฺสชฺชิเต วิสฺสชฺชิต\-สญฺญี, นิสฺสคฺคิยํ ปาจิตฺติยํ. อนฏฺเฐ นฏฺฐสญฺญี, นิสฺสคฺคิยํ ปาจิตฺติยํ. อวินฏฺเฐ วินฏฺฐสญฺญี, นิสฺสคฺคิยํ ปาจิตฺติยํ. อทฑฺเฒ ทฑฺฒสญฺญี, นิสฺสคฺคิยํ ปาจิตฺติยํ. อวิลุตฺเต วิลุตฺตสญฺญี, นิสฺสคฺคิยํ ปาจิตฺติยํ.}

\prose{นิสฺสฏฺฐํ ปฏิลภิตฺวา น กายิเกน ปริโภเคน ปริภุญฺชิ\-ตพฺพํ, น อชฺโฌหริตพฺพํ, ปทีเป วา กาฬวณฺเณ วา อุปเนตพฺพํ, อญฺเญน ภิกฺขุนา กายิเกน ปริโภเคน ปริภุญฺชิ\-ตพฺพํ, น อชฺโฌหริตพฺพํ.}

\prose{สตฺตาหา\-นติกฺกนฺเต อติกฺกนฺเตสญฺญี, อาปตฺติ ทุกฺกฏสฺส. สตฺตาหา\-นติกฺกนฺเต เวมติโก, อาปตฺติ ทุกฺกฏสฺส. สตฺตาหา\-นติกฺกนฺเต อนติกฺกนฺต\-สญฺญี, อนาปตฺติ.}

\proseitemwithcloserruled{625}{อนาปตฺติ อนฺโตสตฺตาหํ อธิฏฺเฐติ, วิสฺสชฺเชติ, นสฺสติ, วินสฺสติ, ฑยฺหติ, อจฺฉินฺทิตฺวา คณฺหนฺติ, วิสฺสาสํ คณฺหนฺติ, อนุปสมฺปนฺนสฺส จตฺเตน วนฺเตน มุตฺเตน อนเปกฺโข ทตฺวา ปฏิลภิตฺวา ปริภุญฺชติ, อุมฺมตฺตกสฺส อาทิกมฺมิกสฺสาติ.}{เภสชฺช\-สิกฺขาปทํ นิฏฺฐิตํ ตติยํ.}
\csromanmarkh{3. ปตฺตวคฺค}\csromanmarkhii{4. วสฺสิกสาฏิก\-สิกฺขาปท}\csromanheadernomark{2}{3. \textbf{ปตฺตวคฺค 4. วสฺสิกสาฏิก}\-\textbf{สิกฺขาปท}}
\csromanmatikaanchor{mtk.217}
\csromanmatikaanchor{mtk.218}
\csromantocmarkref{subsubsection}{ปญฺญตฺติ}{mtk.217}
\bookmark[dest={mtk.217},level=3]{ปญฺญตฺติ}
\csromantocmarkref{subsubsection}{สิกฺขาปทวิภงฺค, ปทภาชนีย}{mtk.218}
\bookmark[dest={mtk.218},level=3]{สิกฺขาปทวิภงฺค, ปทภาชนีย}
\csromantocmarknopage{subsection}{5. จีวรอจฺฉินฺทนสิกฺขาปท}{mtk.219}
\bookmark[dest={mtk.219},level=2]{5. จีวรอจฺฉินฺทนสิกฺขาปท}
\proseitem{626}{เตน สมเยน พุทฺโธ ภควา สาวตฺถิยํ วิหรติ เชตวเน อนาถ\-ปิณฺฑิกสฺส อาราเม. เตน โข ปน สมเยน ภควตา ภิกฺขูนํ วสฺสิกสาฏิกา อนุญฺญาตา โหติ. ฉพฺพคฺคิยา ภิกฺขู ภควตา วสฺสิกสาฏิกา อนุญฺญาตาติ ปฏิกจฺเจว\csromansharedfootnote{07508a4902455451}{ปฏิคจฺเจว (สี)} วสฺสิก\-สาฏิก\-จีวรํ ปริเยสนฺติ, ปฏิกจฺเจว กตฺวา นิวาเสนฺติ, ชิณฺณาย วสฺสิก\-สาฏิกาย นคฺคา กายํ โอวสฺสาเปนฺติ. เย เต ภิกฺขู อปฺปิจฺฉา ฯเปฯ เต อุชฺฌายนฺติ ขิยฺยนฺติ วิปาเจนฺติ “กถํ หิ นาม ฉพฺพคฺคิยา ภิกฺขู ปฏิกจฺเจว วสฺสิก\-สาฏิก\-จีวรํ ปริเยสิสฺสนฺติ, ปฏิกจฺเจว กตฺวา นิวาเสสฺสนฺติ, \csromanfolio{365}ชิณฺณาย วสฺสิก\-สาฏิกาย นคฺคา กายํ โอวสฺสาเปสฺสนฺตี”ติ. อถ โข เต ภิกฺขู ฉพฺพคฺคิเย ภิกฺขู อเนก\-ปริยาเยน วิครหิตฺวา ภควโต เอตมตฺถํ อาโรเจสุํ ฯเปฯ “สจฺจํ กิร ตุเมฺห ภิกฺขเว ปฏิกจฺเจว วสฺสิก\-สาฏิก\-จีวรํ ปริเยสถ, ปฏิกจฺเจว กตฺวา นิวาเสถ, ชิณฺณาย วสฺสิก\-สาฏิกาย นคฺคา กายํ โอวสฺสาเปถา”ติ. สจฺจํ ภควาติ. วิครหิ พุทฺโธ ภควา ฯเปฯ. กถํ หิ นาม ตุเมฺห โมฆปุริสา ปฏิกจฺเจว วสฺสิก\-สาฏิก\-จีวรํ ปริเยสิสฺสถ, ปฏิกจฺเจว กตฺวา นิวาเสสฺสถ, ชิณฺณาย วสฺสิก\-สาฏิกาย นคฺคา กายํ โอวสฺสาเปสฺสถ, เนตํ โมฆปุริสา อปฺปสนฺนานํ วา ปสาทาย ฯเปฯ. เอวญฺจ ปน ภิกฺขเว อิมํ สิกฺขาปทํ อุทฺทิเสยฺยาถ\csromandash{}}

\proseitem{627}{\textbf{“มาโส เสโส คิมฺหานนฺติ ภิกฺขุนา วสฺสิก}\-\textbf{สาฏิก}\-\textbf{จีวรํ ปริเยสิตพฺพํ, อทฺธมาโส เสโส คิมฺหานนฺติ กตฺวา นิวาเสตพฺพํ, โอเรน เจ มาโส เสโส คิมฺหานนฺติ วสฺสิก}\-\textbf{สาฏิก}\-\textbf{จีวรํ ปริเยเสยฺย, โอเรนทฺธมาโส เสโส คิมฺหานนฺติ กตฺวา นิวาเสยฺย, นิสฺสคฺคิยํ ปาจิตฺติยนฺ”}ติ.}

\proseitem{628}{\textbf{มาโส เสโส คิมฺหานนฺติ ภิกฺขุนา วสฺสิก}\-\textbf{สาฏิก}\-\textbf{จีวรํ} ปริเยสิตพฺพนฺติ เย มนุสฺสา ปุพฺเพปิ วสฺสิก\-สาฏิก\-จีวรํ เทนฺติ, เต อุปสงฺกมิตฺวา เอวมสฺสุ วจนียา “กาโล วสฺสิก\-สาฏิกาย สมโย วสฺสิก\-สาฏิกาย, อญฺเญปิ มนุสฺสา วสฺสิก\-สาฏิก\-จีวรํ เทนฺตี”ติ, น วตฺตพฺพา “เทถ เม วสฺสิก\-สาฏิก\-จีวรํ, อาหรถ เม วสฺสิก\-สาฏิก\-จีวรํ, ปริวตฺเตถ เม วสฺสิก\-สาฏิก\-จีวรํ, เจตาเปถ เม วสฺสิกสาฏิกจีวรนฺ”ติ.}

\prose{\textbf{อทฺธมาโส เสโส คิมฺหานนฺติ กตฺวา นิวาเสตพฺพนฺ}ติ อทฺธมาเส เสเส คิมฺหาเน กตฺวา นิวาเสตพฺพํ.}

\prose{\textbf{โอเรน เจ มาโส เสโส คิมฺหานนฺ}ติ อติเรกมาเส เสเส คิมฺหาเน วสฺสิก\-สาฏิก\-จีวรํ ปริเยสติ, นิสฺสคฺคิยํ ปาจิตฺติยํ.}

\prose{\textbf{โอเรนทฺธมาโส เสโส คิมฺหานนฺ}ติ อติเรก\-ทฺธมาเส เสเส คิมฺหาเน กตฺวา นิวาเสติ, นิสฺสคฺคิยํ โหติ. นิสฺสชฺชิตพฺพํ สํฆสฺส วา คณสฺส วา ปุคฺคลสฺส วา. เอวญฺจ ปน ภิกฺขเว นิสฺสชฺชิตพฺพํ ฯเปฯ “อิทํ เม \csromanfolio{366}\textbf{ภนฺเต วสฺสิกสาฏิกจิวรํ อติเรกมาเส เสเส คิมฺหาเน ปริยิฏฺฐํ} \textbf{อติเรก}\-\textbf{ทฺธมาเส เสเส คิมฺหาเน กตฺวา ปริทหิตํ นิสฺสคฺคิยํ, อิมาหํ} สํฆสฺส นิสฺสชฺชามี”ติ ฯเปฯ ทเทยฺยาติ ฯเปฯ ทเทยฺยุนฺติ ฯเปฯ อายสฺมโต ทมฺมีติ.}

\proseitem{629}{อติเรกมาเส เสเส คิมฺหาเน อติเรกสญฺญี วสฺสิก\-สาฏิก\-จีวรํ ปริเยสติ, นิสฺสคฺคิยํ ปาจิตฺติยํ. อติเรกมาเส เสเส คิมฺหาเน เวมติโก วสฺสิก\-สาฏิก\-จีวรํ ปริเยสติ, นิสฺสคฺคิยํ ปาจิตฺติยํ. อติเรกมาเส เสเส คิมฺหาเน อูนกสญฺญี วสฺสิก\-สาฏิก\-จีวรํ ปริเยสติ, นิสฺสคฺคิยํ ปาจิตฺติยํ.}

\prose{อติเรก\-ทฺธมาเส เสเส คิมฺหาเน อติเรกสญฺญี กตฺวา นิวาเสติ, นิสฺสคฺคิยํ ปาจิตฺติยํ. อติเรก\-ทฺธมาเส เสเส คิมฺหาเน เวมติโก กตฺวา นิวาเสติ, นิสฺสคฺคิยํ ปาจิตฺติยํ. อติเรก\-ทฺธมาเส เสเส คิมฺหาเน อูนกสญฺญี กตฺวา นิวาเสติ, นิสฺสคฺคิยํ ปาจิตฺติยํ.}

\prose{สติยา วสฺสิก\-สาฏิกาย นคฺโค กายํ โอวสฺสาเปติ, อาปตฺติ ทุกฺกฏสฺส. อูนกมาเส เสเส คิมฺหาเน อติเรกสญฺญี, อาปตฺติ ทุกฺกฏสฺส. อูนกมาเส เสเส คิมฺหาเน เวมติโก, อาปตฺติ ทุกฺกฏสฺส. อูนกมาเส เสเส คิมฺหาเน อูนกสญฺญี, อนาปตฺติ.}

\prose{อูนกทฺธมาเส เสเส คิมฺหาเน อติเรกสญฺญี, อาปตฺติ ทุกฺกฏสฺส. อูนกทฺธมาเส เสเส คิมฺหาเน เวมติโก, อาปตฺติ ทุกฺกฏสฺส. อูนกทฺธมาเส เสเส คิมฺหาเน อูนกสญฺญี, อนาปตฺติ.}

\proseitem{630}{อนาปตฺติ มาโส เสโส คิมฺหานนฺติ วสฺสิก\-สาฏิก\-จีวรํ ปริเยสติ, อทฺธมาโส เสโส คิมฺหานนฺติ กตฺวา นิวาเสติ. อูนกมาโส เสโส คิมฺหานนฺติ วสฺสิก\-สาฏิก\-จีวรํ ปริเยสติ, อูนกทฺธมาโส เสโส คิมฺหานนฺติ กตฺวา นิวาเสติ. ปริยิฏฺฐาย วสฺสิก\-สาฏิกาย วสฺสํ อุกฺกฑฺฒิยฺยติ, นิวตฺถาย วสฺสิก\-สาฏิกาย วสฺสํ อุกฺกฑฺฒิยฺยติ, โธวิตฺวา นิกฺขิปิตพฺพํ, สมเย นิวาเสตพฺพํ. อจฺฉินฺน\-จีวรสฺส นฏฺฐ\-จีวรสฺส อาปทาสุ อุมฺมตฺตกสฺส อาทิกมฺมิกสฺสาติ.}

\vspace{\tipitakaparskip}
\csromancenter{วสฺสิกสาฏิก\-สิกฺขาปทํ นิฏฺฐิตํ จตุตฺถํ.}
\csromanmarkh{3. ปตฺตวคฺค}\csromanmarkhii{5. จีวร\-อจฺฉินฺทน\-สิกฺขาปท}\csromanheadernomark{1}{3. ปตฺตวคฺค 5. จีวร\-อจฺฉินฺทน\-สิกฺขาปท}
\csromanmatikaanchor{mtk.220}
\csromanmatikaanchor{mtk.221}
\csromanmatikaanchor{mtk.222}
\csromantocmarkref{subsubsection}{อุปนนฺทภิกฺขุวตฺถุ}{mtk.220}
\bookmark[dest={mtk.220},level=3]{อุปนนฺทภิกฺขุวตฺถุ}
\csromantocmarkref{subsubsection}{ปญฺญตฺติ}{mtk.221}
\bookmark[dest={mtk.221},level=3]{ปญฺญตฺติ}
\csromantocmarkref{subsubsection}{สิกฺขาปทวิภงฺค, ปทภาชนีย}{mtk.222}
\bookmark[dest={mtk.222},level=3]{สิกฺขาปทวิภงฺค, ปทภาชนีย}
\proseitem{631}{\csromanfolio{367}เตน สมเยน พุทฺโธ ภควา สาวตฺถิยํ วิหรติ เชตวเน อนาถ\-ปิณฺฑิกสฺส อาราเม. เตน โข ปน สมเยน อายสฺมา อุปนนฺโท สกฺยปุตฺโต ภาตุโน สทฺธิ\-วิหาริกํ ภิกฺขุํ เอตทโวจ “เอหาวุโส ชนปท\-จาริกํ ปกฺกมิสฺสามา”ติ. นาหํ ภนฺเต คมิสฺสามิ ทุพฺพลจีวโรมฺหีติ. เอหาวุโส อหํ เต จีวรํ ทสฺสามีติ ตสฺส จีวรํ อทาสิ. อสฺโสสิ โข โส ภิกฺขุ “ภควา กิร ชนปท\-จาริกํ ปกฺกมิสฺสตี”ติ. อถ โข ตสฺส ภิกฺขุโน เอตทโหสิ “น ทานาหํ อายสฺมตา อุปนนฺเทน สกฺยปุตฺเตน สทฺธิํ ชนปท\-จาริกํ ปกฺกมิสฺสามิ, ภควตา สทฺธิํ ชนปท\-จาริกํ ปกฺกมิสฺสามี”ติ. อถ โข อายสฺมา อุปนนฺโท สกฺยปุตฺโต ตํ ภิกฺขุํ เอตทโวจ “เอหิ ทานิ อาวุโส ชนปท\-จาริกํ ปกฺกมิสฺสามา”ติ. นาหํ ภนฺเต ตยา สทฺธิํ ชนปท\-จาริกํ ปกฺกมิสฺสามิ, ภควตา สทฺธิํ ชนปท\-จาริกํ ปกฺกมิสฺสามีติ. ยํปิ ตฺยาหํ อาวุโส จีวรํ อทาสิํ “มยา สทฺธิํ ชนปท\-จาริกํ ปกฺกมิสฺสตี”ติ กุปิโต อนตฺตมโน อจฺฉินฺทิ.}

\prose{อถ โข โส ภิกฺขุ ภิกฺขูนํ เอตมตฺถํ อาโรเจสิ. เย เต ภิกฺขู อปฺปิจฺฉา ฯเปฯ เต อุชฺฌายนฺติ ขิยฺยนฺติ วิปาเจนฺติ “กถํ หิ นาม อายสฺมา อุปนนฺโท สกฺยปุตฺโต ภิกฺขุสฺส สามํ จีวรํ ทตฺวา กุปิโต อนตฺตมโน อจฺฉินฺทิสฺสตี”ติ. อถ โข เต ภิกฺขู อายสฺมนฺตํ อุปนนฺทํ สกฺยปุตฺตํ อเนก\-ปริยาเยน วิครหิตฺวา ภควโต เอตมตฺถํ อาโรเจสุํ ฯเปฯ “สจฺจํ กิร ตฺวํ อุปนนฺท ภิกฺขุสฺส สามํ จีวรํ ทตฺวา กุปิโต อนตฺตมโน อจฺฉินฺที”ติ. สจฺจํ ภควาติ. วิครหิ พุทฺโธ ภควา ฯเปฯ. กถํ หิ นาม ตฺวํ โมฆปุริส ภิกฺขุสฺส สามํ จีวรํ ทตฺวา กุปิโต อนตฺตมโน อจฺฉินฺทิสฺสสิ, เนตํ โมฆปุริส อปฺปสนฺนานํ วา ปสาทาย ฯเปฯ. เอวญฺจ ปน ภิกฺขเว อิมํ สิกฺขาปทํ อุทฺทิเสยฺยาถ\csromandash{}}

\proseitem{632}{\textbf{“โย ปน ภิกฺขุ ภิกฺขุสฺส สามํ จีวรํ ทตฺวา กุปิโต อนตฺตมโน อจฺฉินฺเทยฺย วา อจฺฉินฺทาเปยฺย วา, นิสฺสคฺคิยํ ปาจิตฺติยนฺ”}ติ.}

\proseitem{633}{\textbf{โย ปนา}ติ โย ยาทิโส ฯเปฯ. \textbf{ภิกฺขู}ติ ฯเปฯ อยํ อิมสฺมิํ อตฺเถ อธิปฺเปโต “ภิกฺขู”ติ.}

\csromanmatikaanchor{mtk.223}
\csromanmatikaanchor{mtk.224}
\csromantocmarknopage{subsection}{6. สุตฺตวิญฺญตฺติสิกฺขาปท}{mtk.223}
\bookmark[dest={mtk.223},level=2]{6. สุตฺตวิญฺญตฺติสิกฺขาปท}
\csromantocmarkref{subsubsection}{ฉพฺพคฺคิยภิกฺขุวตฺถุ}{mtk.224}
\bookmark[dest={mtk.224},level=3]{ฉพฺพคฺคิยภิกฺขุวตฺถุ}
\prose{\csromanfolio{368}\textbf{ภิกฺขุสฺสา}ติ อญฺญสฺส ภิกฺขุสฺส.}

\prose{\textbf{สามนฺ}ติ สยํ ทตฺวา.}

\prose{\textbf{จีวรํ} นาม ฉนฺนํ จีวรานํ อญฺญตรํ จีวรํ วิกปฺปนุปคํ ปจฺฉิมํ.}

\prose{\textbf{กุปิโต อนตฺตมโน}ติ อนภิรทฺโธ อาหตจิตฺโต ขิลชาโต.}

\prose{\textbf{อจฺฉินฺเทยฺยา}ติ สยํ อจฺฉินฺทติ, นิสฺสคฺคิยํ ปาจิตฺติยํ\csromansharedfootnote{7174dcf91c017a08}{นิสฺสคฺคิยํ โหติ (สฺยา)}.}

\prose{\textbf{อจฺฉินฺทาเปยฺยา}ติ อญฺญํ อาณาเปติ, อาปตฺติ ทุกฺกฏสฺส. สกิํ อาณตฺโต พหุกํปิ อจฺฉินฺทติ, นิสฺสคฺคิยํ โหติ. นิสฺสชฺชิตพฺพํ สํฆสฺส วา คณสฺส วา ปุคฺคลสฺส วา. เอวญฺจ ปน ภิกฺขเว นิสฺสชฺชิตพฺพํ ฯเปฯ “อิทํ เม ภนฺเต จีวรํ ภิกฺขุสฺส สามํ ทตฺวา อจฺฉินฺนํ นิสฺสคฺคิยํ, อิมาหํ สํฆสฺส นิสฺสชฺชามี”ติ ฯเปฯ ทเทยฺยาติ ฯเปฯ ทเทยฺยุนฺติ ฯเปฯ อายสฺมโต ทมฺมีติ.}

\proseitem{634}{อุปสมฺปนฺเน อุปสมฺปนฺน\-สญฺญี จีวรํ ทตฺวา กุปิโต อนตฺตมโน อจฺฉินฺทติ วา อจฺฉินฺทาเปติ วา, นิสฺสคฺคิยํ ปาจิตฺติยํ. อุปสมฺปนฺเน เวมติโก จีวรํ ทตฺวา กุปิโต อนตฺตมโน อจฺฉินฺทติ วา อจฺฉินฺทาเปติ วา, นิสฺสคฺคิยํ ปาจิตฺติยํ. อุปสมฺปนฺเน อนุปสมฺปนฺน\-สญฺญี จีวรํ ทตฺวา กุปิโต อนตฺตมโน อจฺฉินฺทติ วา อจฺฉินฺทาเปติ วา, นิสฺสคฺคิยํ ปาจิตฺติยํ.}

\prose{อญฺญํ ปริกฺขารํ ทตฺวา กุปิโต อนตฺตมโน อจฺฉินฺทติ วา อจฺฉินฺทาเปติ วา, อาปตฺติ ทุกฺกฏสฺส. อนุปสมฺปนฺนสฺส จีวรํ วา อญฺญํ วา ปริกฺขารํ ทตฺวา กุปิโต อนตฺตมโน อจฺฉินฺทติ วา อจฺฉินฺทาเปติ วา, อาปตฺติ ทุกฺกฏสฺส. อนุปสมฺปนฺเน อุปสมฺปนฺน\-สญฺญี, อาปตฺติ ทุกฺกฏสฺส. อนุปสมฺปนฺเน เวมติโก, อาปตฺติ ทุกฺกฏสฺส. อนุปสมฺปนฺเน อนุปสมฺปนฺน\-สญฺญี, อาปตฺติ ทุกฺกฏสฺส.}

\proseitemwithcloserruled{635}{อนาปตฺติ โส วา เทติ, ตสฺส วา วิสฺสสนฺโต คณฺหาติ, อุมฺมตฺตกสฺส อาทิกมฺมิกสฺสาติ.}{จีวร\-อจฺฉินฺทน\-สิกฺขาปทํ นิฏฺฐิตํ ปญฺจมํ.}
\csromanmarkh{3. ปตฺตวคฺค}\csromanmarkhii{6. สุตฺต\-วิญฺญตฺติ\-สิกฺขาปท}\csromanheadernomark{2}{3. \textbf{ปตฺตวคฺค 6. สุตฺต}\-\textbf{วิญฺญตฺติ}\-\textbf{สิกฺขาปท}}
\csromanmatikaanchor{mtk.225}
\csromanmatikaanchor{mtk.226}
\csromantocmarkref{subsubsection}{ปญฺญตฺติ}{mtk.225}
\bookmark[dest={mtk.225},level=3]{ปญฺญตฺติ}
\csromantocmarkref{subsubsection}{สิกฺขาปทวิภงฺค, ปทภาชนีย}{mtk.226}
\bookmark[dest={mtk.226},level=3]{สิกฺขาปทวิภงฺค, ปทภาชนีย}
\proseitem{636}{เตน สมเยน พุทฺโธ ภควา ราชคเห วิหรติ เวฬุวเน กลนฺทก\-นิวาเป. เตน โข ปน สมเยน ฉพฺพคฺคิยา ภิกฺขู จีวร\-การ\-สมเย \csromanfolio{369}พหุํ สุตฺตํ วิญฺญาเปสุํ. กเตปิ จีวเร พหุํ สุตฺตํ อวสิฏฺฐํ โหติ. อถ โข ฉพฺพคฺคิยานํ ภิกฺขูนํ เอตทโหสิ “หนฺท มยํ อาวุโส อญฺญํปิ สุตฺตํ วิญฺญาเปตฺวา ตนฺตวาเยหิ จีวรํ วายาเปมา”ติ. อถ โข ฉพฺพคฺคิยา ภิกฺขู อญฺญํปิ สุตฺตํ วิญฺญาเปตฺวา ตนฺตวาเยหิ จีวรํ วายาเปสุํ. วีเตปิ จีวเร พหุํ สุตฺตํ อวสิฏฺฐํ โหติ. ทุติยมฺปิ โข ฉพฺพคฺคิยา ภิกฺขู อญฺญํปิ สุตฺตํ วิญฺญาเปตฺวา ตนฺตวาเยหิ จีวรํ วายาเปสุํ. วีเตปิ จีวเร พหุํ สุตฺตํ อวสิฏฺฐํ โหติ. ตติยมฺปิ โข ฉพฺพคฺคิยา ภิกฺขู อญฺญํปิ สุตฺตํ วิญฺญาเปตฺวา ตนฺตวาเยหิ จีวรํ วายาเปสุํ. มนุสฺสา อุชฺฌายนฺติ ขิยฺยนฺติ วิปาเจนฺติ “กถํ หิ นาม สมณา สกฺยปุตฺติยา สามํ สุตฺตํ วิญฺญาเปตฺวา ตนฺตวาเยหิ จีวรํ วายาเปสฺสนฺตี”ติ.}

\prose{อสฺโสสุํ โข ภิกฺขู เตสํ มนุสฺสานํ อุชฺฌายนฺตานํ ขิยฺยนฺตานํ วิปาเจนฺตานํ. เย เต ภิกฺขู อปฺปิจฺฉา ฯเปฯ เต อุชฺฌายนฺติ ขิยฺยนฺติ วิปาเจนฺติ “กถํ หิ นาม ฉพฺพคฺคิยา ภิกฺขู สามํ สุตฺตํ วิญฺญาเปตฺวา ตนฺตวาเยหิ จีวรํ วายาเปสฺสนฺตี”ติ. อถ โข เต ภิกฺขู ฉพฺพคฺคิเย ภิกฺขู อเนก\-ปริยาเยน วิครหิตฺวา ภควโต เอตมตฺถํ อาโรเจสุํ ฯเปฯ “สจฺจํ กิร ตุเมฺห ภิกฺขเว สามํ สุตฺตํ วิญฺญาเปตฺวา ตนฺตวาเยหิ จีวรํ วายาเปถา”ติ. สจฺจํ ภควาติ. วิครหิ พุทฺโธ ภควา ฯเปฯ. กถํ หิ นาม ตุเมฺห โมฆปุริสา สามํ สุตฺตํ วิญฺญาเปตฺวา ตนฺตวาเยหิ จีวรํ วายาเปสฺสถ, เนตํ โมฆปุริสา อปฺปสนฺนานํ วา ปสาทาย ฯเปฯ. เอวญฺจ ปน ภิกฺขเว อิมํ สิกฺขาปทํ อุทฺทิเสยฺยาถ\csromandash{}}

\proseitem{637}{\textbf{“โย ปน ภิกฺขุ สามํ สุตฺตํ วิญฺญาเปตฺวา ตนฺตวาเยหิ จีวรํ วายาเปยฺย, นิสฺสคฺคิยํ ปาจิตฺติยนฺ”}ติ.}

\proseitem{638}{\textbf{โย ปนา}ติ โย ยาทิโส ฯเปฯ. \textbf{ภิกฺขู}ติ ฯเปฯ อยํ อิมสฺมิํ อตฺเถ อธิปฺเปโต “ภิกฺขู”ติ.}

\prose{\textbf{สามนฺ}ติ สยํ วิญฺญาเปตฺวา.}

\prose{\textbf{สุตฺตํ} นาม ฉ สุตฺตานิ โขมํ กปฺปาสิกํ โกเสยฺยํ กมฺพลํ สาณํ ภงฺคํ.}

\csromanmatikaanchor{mtk.227}
\csromanmatikaanchor{mtk.228}
\csromantocmarknopage{subsection}{7. มหาเปสการสิกฺขาปท}{mtk.227}
\bookmark[dest={mtk.227},level=2]{7. มหาเปสการสิกฺขาปท}
\csromantocmarkref{subsubsection}{อุปนนฺทภิกฺขุวตฺถุ}{mtk.228}
\bookmark[dest={mtk.228},level=3]{อุปนนฺทภิกฺขุวตฺถุ}
\prose{\csromanfolio{370}\textbf{ตนฺตวาเยหี}ติ เปสกาเรหิ วายาเปติ, ปโยเค ปโยเค ทุกฺกฏํ\csromansharedfootnote{e030b675c2fdc3ba}{วายาเปติ, ปโยเค ทุกฺกฏํ (สฺยา)}. ปฏิลาเภน นิสฺสคฺคิยํ โหติ. นิสฺสชฺชิตพฺพํ สํฆสฺส วา คณสฺส วา ปุคฺคลสฺส วา. เอวญฺจ ปน ภิกฺขเว นิสฺสชฺชิตพฺพํ ฯเปฯ “อิทํ เม ภนฺเต จีวรํ สามํ สุตฺตํ วิญฺญาเปตฺวา ตนฺตวาเยหิ วายาปิตํ นิสฺสคฺคิยํ, อิมาหํ สํฆสฺส นิสฺสชฺชามี”ติ ฯเปฯ ทเทยฺยาติ ฯเปฯ ทเทยฺยุนฺติ ฯเปฯ อายสฺมโต ทมฺมิติ.}

\proseitem{639}{วายาปิเต วายาปิตสญฺญี, นิสฺสคฺคิยํ ปาจิตฺติยํ. วายาปิเต เวมติโก, นิสฺสคฺคิยํ ปาจิตฺติยํ. วายาปิเต อวายาปิตสญฺญี, นิสฺสคฺคิยํ ปาจิตฺติยํ.}

\prose{อวายาปิเต วายาปิตสญฺญี, อาปตฺติ ทุกฺกฏสฺส. อวายาปิเต เวมติโก, อาปตฺติ ทุกฺกฏสฺส. อวายาปิเต อวายาปิตสญฺญี, อนาปตฺติ.}

\proseitemwithcloserruled{640}{อนาปตฺติ จีวรํ สิพฺเพตุํ, อาโยเค, กายพนฺธเน, อํสพทฺธเก\csromansharedfootnote{4633f88d2081f312}{อํสวฏฺฏเก (สี)}, ปตฺต\-ตฺถวิกาย, ปริสฺสาวเน, ญาตกานํ, ปวาริตานํ, อญฺญสฺสตฺถาย, อตฺตโน ธเนน, อุมฺมตฺตกสฺส อาทิกมฺมิกสฺสาติ.}{สุตฺต\-วิญฺญตฺติ\-สิกฺขาปทํ นิฏฺฐิตํ ฉฏฺฐํ.}
\csromanmarkh{3. ปตฺตวคฺค}\csromanmarkhii{7. มหา\-เปสการ\-สิกฺขาปท}\csromanheadernomark{2}{3. \textbf{ปตฺตวคฺค 7. มหา}\-\textbf{เปสการ}\-\textbf{สิกฺขาปท}}
\proseitem{641}{เตน สมเยน พุทฺโธ ภควา สาวตฺถิยํ วิหรติ เชตวเน อนาถ\-ปิณฺฑิกสฺส อาราเม. เตน โข ปน สมเยน อญฺญตโร ปุริโส ปวาสํ คจฺฉนฺโต ปชาปติํ เอตทโวจ “สุตฺตํ ธารยิตฺวา อมุกสฺส ตนฺตวายสฺส เทหิ, จีวรํ วายาเปตฺวา นิกฺขิป, อาคโต อยฺยํ อุปนนฺทํ จีวเรน อจฺฉาเทสฺสามี”ติ. อสฺโสสิ โข อญฺญตโร ปิณฺฑจาริโก ภิกฺขุ ตสฺส ปุริสสฺส อิมํ วาจํ ภาสมานสฺส. อถ โข โส ภิกฺขุ เยนายสฺมา อุปนนฺโท สกฺยปุตฺโต เตนุปสงฺกมิ, อุปสงฺกมิตฺวา อายสฺมนฺตํ อุปนนฺทํ สกฺยปุตฺตํ เอตทโวจ “มหาปุญฺโญสิ ตฺวํ อาวุโส อุปนนฺท, อมุกสฺมิํ โอกาเส อญฺญตโร ปุริโส ปวาสํ \csromanfolio{371}คจฺฉนฺโต ปชาปติํ เอตทโวจ “สุตฺตํ ธารยิตฺวา อมุกสฺส ตนฺตวายสฺส เทหิ, จีวรํ วายาเปตฺวา นิกฺขิป, อาคโต อยฺยํ อุปนนฺทํ จีวเรน อจฺฉาเทสฺสามี”ติ. อตฺถาวุโส มํ โส อุปฏฺฐาโกติ. โสปิ โข ตนฺตวาโย อายสฺมโต อุปนนฺทสฺส สกฺยปุตฺตสฺส อุปฏฺฐาโก โหติ. อถ โข อายสฺมา อุปนนฺโท สกฺยปุตฺโต เยน โส ตนฺตวาโย เตนุปสงฺกมิ, อุปสงฺกมิตฺวา ตํ ตนฺตวายํ เอตทโวจ “อิทํ โข อาวุโส จีวรํ มํ อุทฺทิสฺส วิยฺยติ, อายตญฺจ กโรหิ วิตฺถตญฺจ อปฺปิตญฺจ สุวีตญฺจ สุปฺปวายิตญฺจ สุวิเลขิตญฺจ สุวิตจฺฉิตญฺจ กโรหี”ติ. เอเต โข เม ภนฺเต สุตฺตํ ธารยิตฺวา อทํสุ “อิมินา สุตฺเตน จีวรํ วินาหี”ติ, น ภนฺเต สกฺกา อายตํ วา วิตฺถตํ วา อปฺปิตํ ว กาตุํ, สกฺกา จ โข ภนฺเต สุวีตญฺจ สุปฺปวายิตญฺจ สุวิเลขิตญฺจ สุวิตจฺฉิตญฺจ กาตุนฺติ. อิงฺฆ ตฺวํ อาวุโส อายตญฺจ กโรหิ วิตฺถตญฺจ อปฺปิตญฺจ, น เตน สุตฺเตน ปฏิพทฺธํ ภวิสฺสตีติ.}

\csromanmatikaanchor{mtk.229}
\csromanmatikaanchor{mtk.230}
\csromantocmarkref{subsubsection}{ปญฺญตฺติ}{mtk.229}
\bookmark[dest={mtk.229},level=3]{ปญฺญตฺติ}
\csromantocmarkref{subsubsection}{สิกฺขาปทวิภงฺค, ปทภาชนีย}{mtk.230}
\bookmark[dest={mtk.230},level=3]{สิกฺขาปทวิภงฺค, ปทภาชนีย}
\prose{อถ โข โส ตนฺตวาโย ยถาภตํ สุตฺตํ ตนฺเต อุปเนตฺวา เยน สา อิตฺถี เตนุปสงฺกมิ, อุปสงฺกมิตฺวา ตํ อิตฺถิํ เอตทโวจ “สุตฺเตน อยฺเย อตฺโถ”ติ. นนุ ตฺวํ อยฺโย\csromansharedfootnote{3fe1aebe00dfa21d}{อยฺย (สฺยา)} มยา วุตฺโต “อิมินา สุตฺเตน จีวรํ วินาหี”ติ. สจฺจาหํ อยฺเย ตยา วุตฺโต “อิมินา สุตฺเตน จีวรํ วินาหี”ติ. อปิ จ มํ อยฺโย อุปนนฺโท เอวมาห “อิงฺฆ ตฺวํ อาวุโส อายตญฺจ กโรหิ วิตฺถตญฺจ อปฺปิตญฺจ, น เตน สุตฺเตน ปฏิพทฺธํ ภวิสฺสตี”ติ. อถ โข สา อิตฺถี ยตฺตกํ เยว สุตฺตํ ปฐมํ อทาสิ, ตตฺตกํ ปจฺฉา อทาสิ. อสฺโสสิ โข อายสฺมา อุปนนฺโท สกฺยปุตฺโต “โส กิร ปุริโส ปวาสโต อาคโต”ติ. อถ โข อายสฺมา อุปนนฺโท สกฺยปุตฺโต เยน ตสฺส ปุริสสฺส นิเวสนํ เตนุปสงฺกมิ, อุปสงฺกมิตฺวา ปญฺญตฺเต อาสเน นิสีทิ. อถ โข โส ปุริโส เยนายสฺมา อุปนนฺโท สกฺยปุตฺโต เตนุปสงฺกมิ, อุปสงฺกมิตฺวา อายสฺมนฺตํ อุปนนฺทํ สกฺยปุตฺตํ อภิวาเทตฺวา เอกมนฺตํ นิสีทิ, เอกมนฺตํ นิสินฺโน โข โส ปุริโส ปชาปติํ เอตทโวจ “วีตํ ตํ จีวรนฺ”ติ. อามายฺย วีตํ ตํ จีวรนฺติ. อาหร อยฺยํ อุปนนฺทํ จีวเรน อจฺฉาเทสฺสามีติ. อถ โข สา อิตฺถี ตํ จีวรํ นีหริตฺวา สามิกสฺส ทตฺวา เอตมตฺถํ อาโรเจสิ. อถ โข โส ปุริโส อายสฺมโต อุปนนฺทสฺส \csromanfolio{372}สกฺยปุตฺตสฺส จีวรํ ทตฺวา อุชฺฌายติ ขิยฺยติ วิปาเจติ “มหิจฺฉา อิเม สมณา สกฺยปุตฺติยา อสนฺตุฏฺฐา, นยิเม สุกรา จีวเรน อจฺฉาเทตุํ, กถํ หิ นาม อยฺโย อุปนนฺโท มยา ปุพฺเพ อปฺปวาริโต ตนฺตวาเย\csromansharedfootnote{712b90db12530bf8}{คหปติกสฺส ตนฺตวาเย (ก)} อุปสงฺกมิตฺวา จีวเร วิกปฺปํ อาปชฺชิสฺสตี”ติ.}

\prose{อสฺโสสุํ โข ภิกฺขู ตสฺส ปุริสสฺส อุชฺฌายนฺตสฺส ขิยฺยนฺตสฺส วิปาเจนฺตสฺส. เย เต ภิกฺขู อปฺปิจฺฉา ฯเปฯ เต อุชฺฌายนฺติ ขิยฺยนฺติ วิปาเจนฺติ “กถํ หิ นาม อายสฺมา อุปนนฺโท สกฺยปุตฺโต ปุพฺเพ อปฺปวาริโต คหปติกสฺส ตนฺตวาเย อุปสงฺกมิตฺวา จีวเร วิกปฺปํ อาปชฺชิสฺสตี”ติ. อถ โข เต ภิกฺขู อายสฺมนฺตํ อุปนนฺทํ สกฺยปุตฺตํ อเนก\-ปริยาเยน วิครหิตฺวา ภควโต เอตมตฺถํ อาโรเจสุํ ฯเปฯ “สจฺจํ กิร ตฺวํ อุปนนฺท ปุพฺเพ อปฺปวาริโต คหปติกสฺส ตนฺตวาเย อุปสงฺกมิตฺวา จีวเร วิกปฺปํ อาปชฺชี”ติ. สจฺจํ ภควาติ. ญาตโก เต อุปนนฺท อญฺญาตโกติ. อญฺญาตโก ภควาติ. อญฺญาตโก โมฆปุริส อญฺญาตกสฺส น ชานาติ ปติรูปํ วา อปฺปติรูปํ วา สนฺตํ วา อสนฺตํ วา, ตตฺถ นาม ตฺวํ โมฆปุริส ปุพฺเพ อปฺปวาริโต อญฺญาตกสฺส คหปติกสฺส ตนฺตวาเย อุปสงฺกมิตฺวา จีวเร วิกปฺปํ อาปชฺชิสฺสสิ, เนตํ โมฆปุริส อปฺปสนฺนานํ วา ปสาทาย ฯเปฯ. เอวญฺจ ปน ภิกฺขเว อิมํ สิกฺขาปทํ อุทฺทิเสยฺยาถ\csromandash{}}

\proseitem{642}{\textbf{“ภิกฺขุํ ปเนว อุทฺทิสฺส อญฺญาตโก คหปติ วา คหปตานี วา ตนฺตวาเยหิ จีวรํ วายาเปยฺย, ตตฺร เจ โส ภิกฺขุ ปุพฺเพ อปฺปวาริโต ตนฺตวาเย อุปสงฺกมิตฺวา จีวเร วิกปฺปํ อาปชฺเชยฺย ‘อิทํ โข อาวุโส จีวรํ มํ อุทฺทิสฺส วิยฺยติ, อายตญฺจ กโรถ, วิตฺถตญฺจ อปฺปิตญฺจ สุวีตญฺจ สุปฺปวายิตญฺจ สุวิเลขิตญฺจ สุวิตจฺฉิตญฺจ กโรถ, อปฺเปว นาม มยํปิ อายสฺมนฺตานํ กิญฺจิมตฺตํ อนุปทชฺเชยฺยามา’ติ. เอวญฺจ โส ภิกฺขุ วตฺวา กิญฺจิมตฺตํ อนุปทชฺเชยฺย อนฺตมโส ปิณฺฑปาตมตฺตํปิ, นิสฺสคฺคิยํ ปาจิตฺติยนฺ”}ติ.}

\proseitem{643}{\textbf{ภิกฺขิํ ปเนว อุทฺทิสฺสา}ติ ภิกฺขุสฺส\-ตฺถาย ภิกฺขุํ อารมฺมณํ กริตฺวา ภิกฺขุํ อจฺฉาเทตุกาโม.}

\prose{\textbf{อญฺญาตโก} นาม มาติโต วา ปิติโต วา ยาว สตฺตมา ปิตามหยุคา อสมฺพทฺโธ.}

\prose{\csromanfolio{373}\textbf{คหปติ} นาม โย โกจิ อคารํ อชฺฌาวสติ.}

\prose{\textbf{คหปตานี} นาม ยา กาจิ อคารํ อชฺฌาวสติ.}

\prose{\textbf{ตนฺต วาเยหี}ติ เปสกาเรหิ.}

\prose{\textbf{จีวรํ} นาม ฉนฺนํ จีวรานํ อญฺญตรํ จีวรํ วิกปฺปนุปคํ ปจฺฉิมํ.}

\prose{\textbf{วายาเปยฺยา}ติ วินาเปติ.}

\prose{\textbf{ตตฺร เจ โส ภิกฺขู}ติ ยํ ภิกฺขุํ อุทฺทิสฺส จีวรํ วิยฺยติ โส ภิกฺขุ.}

\prose{\textbf{ปุพฺเพ อปฺปวาริโต}ติ ปุพฺเพ อวุตฺโต โหติ “กีทิเสน เต ภนฺเต จีวเรน อตฺโถ, กีทิสํ เต จีวรํ วายาเปมี”ติ.}

\prose{\textbf{ตนฺตวาเย อุปสงฺกมิตฺวา}ติ ฆรํ คนฺตฺวา ยตฺถ กตฺถจิ อุปสงฺกมิตฺวา.}

\prose{\textbf{จีวเร วิกปฺปํ อาปชฺเชยฺยา}ติ “อิทํ โข อาวุโส จีวรํ มํ อุทฺทิสฺส วิยฺยติ, อายตญฺจ กโรถ, วิตฺถตญฺจ อปฺปิตญฺจ สุวีตญฺจ สุปฺปวายิตญฺจ สุวิเลขิตญฺจ สุวิตจฺฉิตญฺจ กโรถ, อปฺเปว นาม มยํปิ อายสฺมนฺตานํ กิญฺจิมตฺตํ อนุปทชฺเชยฺยามา”ติ.}

\prose{\textbf{เอวญฺจ โส ภิกฺขุ วตฺวา กิญฺจิมตฺตํ อนุปทชฺเชยฺย อนฺตมโส ปิณฺฑปาตมตฺตํปี}ติ \textbf{ปิณฺฑปาโต} นาม ยาคุปิ ภตฺตํปิ ขาทนียํปิ จุณฺณปิณฺโฑปิ ทนฺตกฏฺฐํปิ ทสิกสุตฺตํปิ, อนฺตมโส ธมฺมํปิ ภณติ.}

\prose{ตสฺส วจเนน อายตํ วา วิตฺถตํ วา อปฺปิตํ วา กโรติ, ปโยเค ทุกฺกฏํ. ปฏิลาเภน นิสฺสคฺคิยํ โหติ. นิสฺสชฺชิตพฺพํ สํฆสฺส วา คณสฺส วา ปุคฺคลสฺส วา. เอวญฺจ ปน ภิกฺขเว นิสฺสชฺชิตพฺพํ ฯเปฯ “อิทํ เม ภนฺเต จีวรํ ปุพฺเพ อปฺปวาริโต อญฺญาตกสฺส คหปติกสฺส ตนฺตวาเย อุปสงฺกมิตฺวา จีวเร วิกปฺปํ อาปนฺนํ นิสฺสคฺคิยํ, อิมาหํ สํฆสฺส นิสฺสชฺชามี”ติ ฯเปฯ ทเทยฺยาติ ฯเปฯ ทเทยฺยุนฺติ ฯเปฯ อายสฺมโต ทมฺมีติ.}

\proseitem{644}{อญฺญาตเก อญฺญาตกสญฺญี ปุพฺเพ อปฺปวาริโต คหปติกสฺส ตนฺตวาเย อุปสงฺกมิตฺวา จีวเร วิกปฺปํ อาปชฺชติ, นิสฺสคฺคิยํ ปาจิตฺติยํ. อญฺญาตเก เวมติโก ปุพฺเพ อปฺปวาริโต คหปติกสฺส ตนฺตวาเย อุปสงฺกมิตฺวา จีวเร วิกปฺปํ อาปชฺชติ, นิสฺสคฺคิยํ ปาจิตฺติยํ. อญฺญาตเก ญาตกสญฺญี ปุพฺเพ อปฺปวาริโต คหปติกสฺส ตนฺตวาเย อุปสงฺกมิตฺวา จีวเร วิกปฺปํ อาปชฺชติ, นิสฺสคฺคิยํ ปาจิตฺติยํ.}

\csromanmatikaanchor{mtk.231}
\csromanmatikaanchor{mtk.232}
\csromantocmarknopage{subsection}{8. อจฺเจกจีวรสิกฺขาปท}{mtk.231}
\bookmark[dest={mtk.231},level=2]{8. อจฺเจกจีวรสิกฺขาปท}
\csromantocmarkref{subsubsection}{สาวตฺถิภิกฺขุวตฺถุ}{mtk.232}
\bookmark[dest={mtk.232},level=3]{สาวตฺถิภิกฺขุวตฺถุ}
\prose{\csromanfolio{374}ญาตเก อญฺญาตกสญฺญี, อาปตฺติ ทุกฺกฏสฺส. ญาตเก เวมติโก, อาปตฺติ ทุกฺกฏสฺส. ญาตเก ญาตกสญฺญี, อนาปตฺติ.}

\proseitemwithcloserruled{645}{อนาปตฺติ ญาตกานํ, ปวาริตานํ, อญฺญสฺสตฺถาย, อตฺตโน ธเนน, มหคฺฆํ วายาเปตุ\-กามสฺส อปฺปคฺฆํ วายาเปติ, อุมฺมตฺตกสฺส อาทิกมฺมิกสฺสาติ.}{มหา\-เปสการ\-สิกฺขาปทํ นิฏฺฐิตํ สตฺตมํ.}
\csromanmarkh{3. ปตฺตวคฺค}\csromanmarkhii{8. อจฺเจกจีวร\-สิกฺขาปท}\csromanheadernomark{2}{3. \textbf{ปตฺตวคฺค 8. อจฺเจกจีวร}\-\textbf{สิกฺขาปท}}
\proseitem{646}{เตน สมเยน พุทฺโธ ภควา สาวตฺถิยํ วิหรติ เชตวเน อนาถ\-ปิณฺฑิกสฺส อาราเม. เตน โข ปน สมเยน อญฺญตโร มหามตฺโต ปวาสํ คจฺฉนฺโต ภิกฺขูนํ สนฺติเก ทูตํ ปาเหสิ “อาคจฺฉนฺตุ ภทนฺตา วสฺสาวาสิกํ ทสฺสามี”ติ. ภิกฺขู “วสฺสํ\-วุฏฺฐานํ ภควตา วสฺสาวาสิกํ อนุญฺญาตนฺ”ติ กุกฺกุจฺจายนฺตา นาคมํสุ. อถ โข โส มหามตฺโต อุชฺฌายติ ขิยฺยติ วิปาเจติ “กถํ หิ นาม ภทนฺตา มยา ทูเต ปหิเต นาคจฺฉิสฺสนฺติ, อหํ หิ เสนาย คจฺฉามิ, ทุชฺชานํ ชีวิตํ, ทุชฺชานํ มรณนฺ”ติ. อสฺโสสุํ โข ภิกฺขู ตสฺส มหามตฺตสฺส อุชฺฌายนฺตสฺส ขิยฺยนฺตสฺส วิปาเจนฺตสฺส. อถ โข เต ภิกฺขู ภควโต เอตมตฺถํ อาโรเจสุํ. อถ โข ภควา เอตสฺมิํ นิทาเน เอตสฺมิํ ปกรเณ ธมฺมิํ กถํ กตฺวา ภิกฺขู อามนฺเตสิ “อนุชานามิ ภิกฺขเว อจฺเจกจีวรํ ปฏิคฺคเหตฺวา นิกฺขิปิตุนฺ”ติ.}

\csromanmatikaanchor{mtk.233}
\csromanmatikaanchor{mtk.234}
\csromantocmarkref{subsubsection}{ปญฺญตฺติ}{mtk.233}
\bookmark[dest={mtk.233},level=3]{ปญฺญตฺติ}
\csromantocmarkref{subsubsection}{สิกฺขาปทวิภงฺค, ปทภาชนีย}{mtk.234}
\bookmark[dest={mtk.234},level=3]{สิกฺขาปทวิภงฺค, ปทภาชนีย}
\proseitem{647}{เตน โข ปน สมเยน ภิกฺขู “ภควตา อนุญฺญาตํ อจฺเจกจีวรํ ปฏิคฺคเหตฺวา นิกฺขิปิตุนฺ”ติ อจฺเจกจีวรานิ ปฏิคฺคเหตฺวา จีวร\-กาล\-สมยํ อติกฺกาเมนฺติ, ตานิ จีวรานิ จีวรวํเส ภณฺฑิกา\-พทฺธานิ ติฏฺฐนฺติ. อทฺทส โข อายสฺมา อานนฺโท เสนาสนจาริกํ อาหิณฺฑนฺโต ตานิ จีวรานิ จีวรวํเส ภณฺฑิกา\-พทฺธานิ ติฏฺฐนฺเต, ทิสฺวา ภิกฺขู เอตทโวจ “กสฺสิมานิ อาวุโส จีวรานิ จีวรวํเส ภณฺฑิกา\-พทฺธานิ ติฏฺฐนฺตี”ติ. อมฺหากํ อาวุโส อจฺเจกจีวรานีติ. กีวจิรํ ปนาวุโส อิมานิ จีวรานิ นิกฺขิตฺตานีติ. อถ โข เต ภิกฺขู \csromanfolio{375}อายสฺมโต อานนฺทสฺส ยถา\-นิกฺขิตฺตํ อาโรเจสุํ. อายสฺมา อานนฺโท อุชฺฌายติ ขิยฺยติ วิปาเจติ “กถํ หิ นาม ภิกฺขู อจฺเจกจีวรํ ปฏิคฺคเหตฺวา จีวร\-กาล\-สมยํ อติกฺกาเมสฺสนฺตี”ติ. อถ โข อายสฺมา อานนฺโท เต ภิกฺขู อเนก\-ปริยาเยน วิครหิตฺวา ภควโต เอตมตฺถํ อาโรเจสิ ฯเปฯ “สจฺจํ กิร ภิกฺขเว ภิกฺขู อจฺเจกจีวรํ ปฏิคฺคเหตฺวา จีวร\-กาล\-สมยํ อติกฺกาเมนฺตี”ติ. สจฺจํ ภควาติ. วิครหิ พุทฺโธ ภควา ฯเปฯ. กถํ หิ นาม เต ภิกฺขเว โมฆปุริสา อจฺเจกจีวรํ ปฏิคฺคเหตฺวา จีวร\-กาล\-สมยํ อติกฺกาเมสฺสนฺติ, เนตํ ภิกฺขเว อปฺปสนฺนานํ วา ปสาทาย ฯเปฯ. เอวญฺจ ปน ภิกฺขเว อิมํ สิกฺขาปทํ อุทฺทิเสยฺยาถ\csromandash{}}

\proseitem{648}{\textbf{“ทสาหานาคตํ กตฺติกเตมาสิก ปุณฺณมํ ภิกฺขุโน ปเนว อจฺเจกจีวรํ อุปฺปชฺเชยฺย, อจฺเจกํ มญฺญมาเนน ภิกฺขุนา ปฏิคฺคเหตพฺพํ, ปฏิคฺคเหตฺวา ยาว จีวร}\-\textbf{กาล}\-\textbf{สมยํ นิกฺขิปิตพฺพํ, ตโต เจ อุตฺตริ นิกฺขิเปยฺย, นิสฺสคฺคิยํ ปาจิตฺติยนฺ”}ติ.}

\proseitem{649}{\textbf{ทสาหา}\-\textbf{นาคตนฺ}ติ ทสาหานาคตาย ปวารณาย.}

\prose{\textbf{กตฺติก}\-\textbf{เตมาสิก}\-\textbf{ปุณฺณมนฺ}ติ ปวารณา กตฺติกา วุจฺจติ.}

\prose{\textbf{อจฺเจกจีวรํ} นาม เสนาย วา คนฺตุกาโม โหติ, ปวาสํ วา คนฺตุกาโม โหติ, คิลาโน วา โหติ, คพฺภินี วา โหติ, อสฺสทฺธสฺส วา สทฺธา อุปฺปนฺนา โหติ, อปฺปสนฺนสฺส วา ปสาโท อุปฺปนฺโน โหติ, โส เจ ภิกฺขูนํ สนฺติเก ทูตํ ปหิเณยฺย “อาคจฺฉนฺตุ ภทนฺตา วสฺสาวาสิกํ ทสฺสามี”ติ, เอตํ อจฺเจกจีวรํ นาม.}

\prose{\textbf{อจฺเจกํ มญฺญมาเนน ภิกฺขุนา ปฏิคฺคเหตพฺพํ, ปฏิคฺคเหตฺวา ยาว จีวร}\-\textbf{กาล}\-\textbf{สมยํ นิกฺขิปิตพฺพนฺ}ติ สญฺญาณํ กตฺวา นิกฺขิปิตพฺพํ “อิทํ อจฺเจกจีวรนฺ”ติ.}

\prose{\textbf{จีวร}\-\textbf{กาล}\-\textbf{สมโย} นาม อนตฺถเต กถิเน วสฺสานสฺส ปจฺฉิโม มาโส, อตฺถเต กถิเน ปญฺจมาสา.}

\csromanmatikaanchor{mtk.235}
\csromanmatikaanchor{mtk.236}
\csromantocmarknopage{subsection}{9. สาสงฺกสิกฺขาปท}{mtk.235}
\bookmark[dest={mtk.235},level=2]{9. สาสงฺกสิกฺขาปท}
\csromantocmarkref{subsubsection}{สาวตฺถิภิกฺขุวตฺถุ}{mtk.236}
\bookmark[dest={mtk.236},level=3]{สาวตฺถิภิกฺขุวตฺถุ}
\prose{\textbf{ตโต เจ อุตฺตริ นิกฺขิเปยฺยา}ติ อนตฺถเต กถิเน วสฺสานสฺส ปจฺฉิมํ ทิวสํ อติกฺกาเมติ, นิสฺสคฺคิยํ ปาจิตฺติยํ\csromansharedfootnote{7174dcf91c017a08}{นิสฺสคฺคิยํ โหติ (สฺยา)}. อตฺถเต กถิเน \csromanfolio{376}กถินุทฺธาร\-ทิวสํ อติกฺกาเมติ, นิสฺสคฺคิยํ โหติ. นิสฺสชฺชิตพฺพํ สํฆสฺส วา คณสฺส วา ปุคฺคลสฺส วา. เอวญฺจ ปน ภิกฺขเว นิสฺสชฺชิตพฺพํ ฯเปฯ “อิทํ เม ภนฺเต อจฺเจกจีวรํ จีวร\-กาล\-สมยํ อติกฺกามิตํ นิสฺสคฺคิยํ, อิมาหํ สํฆสฺส นิสฺสชฺชามี”ติ ฯเปฯ ทเทยฺยติ ฯเปฯ ทเทยฺยุนฺติ ฯเปฯ อายสฺมโต ทมฺมีติ.}

\proseitem{650}{อจฺเจกจีวเร อจฺเจกจีวร\-สญฺญี จีวร\-กาล\-สมยํ อติกฺกาเมติ, นิสฺสคฺคิยํ ปาจิตฺติยํ. อจฺเจกจีวเร เวมติโก จีวร\-กาล\-สมยํ อติกฺกาเมติ, นิสฺสคฺคิยํ ปาจิตฺติยํ. อจฺเจกจีวเร อนจฺเจกจีวร\-สญฺญี จีวร\-กาล\-สมยํ อติกฺกาเมติ, นิสฺสคฺคิยํ ปาจิตฺติยํ. อนธิฏฺฐิเต อธิฏฺฐิต\-สญฺญี ฯเปฯ. อวิกปฺปิเต วิกปฺปิตสญฺญี. อวิสฺสชฺชิเต วิสฺสชฺชิต\-สญฺญี. อนฏฺเฐ นฏฺฐสญฺญี. อวินฏฺเฐ วินฏฺฐสญฺญี. อทฑฺเฒ ทฑฺฒสญฺญี. อวิลุตฺเต วิลุตฺตสญฺญี, จีวร\-กาล\-สมยํ อติกฺกาเมติ, นิสฺสคฺคิยํ ปาจิตฺติยํ.}

\prose{นิสฺสคฺคิยํ จีวรํ อนิสฺสชฺชิตฺวา ปริภุญฺชติ, อาปตฺติ ทุกฺกฏสฺส. อนจฺเจกจีวเร อจฺเจกจีวร\-สญฺญี, อาปตฺติ ทุกฺกฏสฺส. อนจฺเจกจีวเร เวมติโก, อาปตฺติ ทุกฺกฏสฺส. อนจฺเจกจีวเร อนจฺเจกจีวร\-สญฺญี, อนาปตฺติ.}

\proseitemwithcloserruled{651}{อนาปตฺติ อนฺโตสมเย อธิฏฺเฐติ, วิกปฺเปติ, วิสฺสชฺเชติ, นสฺสติ, วินสฺสติ, ฑยฺหติ, อจฺฉินฺทิตฺวา คณฺหนฺติ, วิสฺสาสํ คณฺหนฺติ, อุมฺมตฺตกสฺส อาทิกมฺมิกสฺสาติ.}{อจฺเจกจีวร\-สิกฺขาปทํ นิฏฺฐิตํ อฏฺฐมํ.}
\csromanmarkh{3. ปตฺตวคฺค}\csromanmarkhii{9. สาสงฺก\-สิกฺขาปท}\csromanheadernomark{2}{3. \textbf{ปตฺตวคฺค 9. สาสงฺก}\-\textbf{สิกฺขาปท}}
\proseitem{652}{เตน สมเยน พุทฺโธ ภควา สาวตฺถิยํ วิหรติ เชตวเน อ นาถปิณฺฑิกสฺส อาราเม. เตน โข ปน สมเยน ภิกฺขู วุฏฺฐวสฺสา อรญฺญเกสุ เสนาสเนสุ วิหรนฺติ. กตฺติกโจรกา “ภิกฺขู ลทฺธลาภา”ติ ปริปาเตนฺติ. ภควโต เอตมตฺถํ อาโรเจสุํ. อถ โข ภควา เอตสฺมิํ นิทาเน เอตสฺมิํ ปกรเณ ธมฺมิํ กถํ กตฺวา ภิกฺขู อามนฺเตสิ “อนุชานามิ ภิกฺขเว อารญฺญเกสุ เสนาสเนสุ วิหรนฺเตน ติณฺณํ จีวรานํ อญฺญตรํ จีวรํ อนฺตรฆเร นิกฺขิปิตุนฺ”ติ.}

\csromanmatikaanchor{mtk.237}
\csromanmatikaanchor{mtk.238}
\csromantocmarkref{subsubsection}{ปญฺญตฺติ}{mtk.237}
\bookmark[dest={mtk.237},level=3]{ปญฺญตฺติ}
\csromantocmarkref{subsubsection}{สิกฺขาปทวิภงฺค, ปทภาชนีย}{mtk.238}
\bookmark[dest={mtk.238},level=3]{สิกฺขาปทวิภงฺค, ปทภาชนีย}
\prose{\csromanfolio{377}เตน โข ปน สมเยน ภิกฺขู ภควตา อนุญฺญาตํ “อารญฺญเกสุ เสนาสเนสุ วิหรนฺเตน ติณฺณํ จีวรานํ อญฺญตรํ จีวรํ อนฺตรฆเร นิกฺขิปิตุนฺ”ติ ติณฺณํ จีวรานํ อญฺญตรํ จีวรํ อนฺตรฆเร นิกฺขิปิตฺวา อติเรกฉารตฺตํ วิปฺปวสนฺติ, ตานิ จีวรานิ นสฺสนฺติปิ วินสฺสนฺติปิ ฑยฺหนฺติปิ อุนฺทูเรหิปิ ขชฺชนฺติ. ภิกฺขู ทุจฺโจฬา โหนฺติ ลูขจีวรา. ภิกฺขู เอวมาหํสุ “กิสฺส ตุเมฺห อาวุโส ทุจฺโจฬา ลูขจีวรา”ติ. อถ โข เต ภิกฺขู ภิกฺขูนํ เอตมตฺถํ อาโรเจสุํ. เย เต ภิกฺขู อปฺปิจฺฉา, เต อุชฺฌายนฺติ ขิยฺยนฺติ วิปาเจนฺติ “กถํ หิ นาม ภิกฺขู ติณฺณํ จีวรานํ อญฺญตรํ จีวรํ อนฺตรฆเร นิกฺขิปิตฺวา อติเรกฉารตฺตํ วิปฺปวสิสฺสนฺตี”ติ. อถ โข เต ภิกฺขู เต อเนก\-ปริยาเยน วิครหิตฺวา ภควโต เอตมตฺถํ อาโรเจสุํ ฯเปฯ “สจฺจํ กิร ภิกฺขเว ภิกฺขู ติณฺณํ จีวรานํ อญฺญตรํ จีวรํ อนฺตรฆเร นิกฺขิปิตฺวา อติเรกฉารตฺตํ วิปฺปวสนฺตี”ติ. สจฺจํ ภควาติ. วิครหิ พุทฺโธ ภควา ฯเปฯ. กถํ หิ นาม เต ภิกฺขเว โมฆปุริสา ติณฺณํ จีวรานํ อญฺญตรํ จีวรํ อนฺตรฆเร นิกฺขิปิตฺวา อติเรกฉารตฺตํ วิปฺปวสิสฺสนฺติ, เนตํ ภิกฺขเว อปฺปสนฺนานํ วา ปสาทาย ฯเปฯ. เอวญฺจ ปน ภิกฺขเว อิมํ สิกฺขาปทํ อุทฺทิเสยฺยาถ\csromandash{}}

\proseitem{653}{\textbf{“อุปวสฺสํ โข ปน กตฺติก}\-\textbf{ปุณฺณมํ ยานิ โข ปน ตานิ อารญฺญกานิ เสนาสนานิ สาสงฺก}\-\textbf{สมฺมตานิ สปฺปฏิภยานิ, ตถารูเปสุ ภิกฺขุ เสนาสเนสุ วิหรนฺโต อากงฺขมาโน ติณฺณํ จีวรานํ อญฺญตรํ จีวรํ อนฺตรฆเร นิกฺขิเปยฺย, สิยา จ ตสฺส ภิกฺขุโน โกจิเทว ปจฺจโย เตน จีวเรน วิปฺปวาสาย, ฉารตฺต}\-\textbf{ปรมํ เตน ภิกฺขุนา เตน จีวเรน วิปฺปวสิตพฺพํ, ตโต เจ อุตฺตริ วิปฺปวเสยฺย อญฺญตฺร ภิกฺขุ}\-\textbf{สมฺมุติยา, นิสฺสคฺคิยํ ปาจิตฺติยนฺ”ติ}.}

\proseitem{654}{\textbf{อุปวสฺสํ โข ปนา}ติ วุฏฺฐวสฺสานํ.}

\prose{\textbf{กตฺติกปุณฺณามนฺ}ติ กตฺติก\-จาตุมาสินี วุจฺจติ.}

\prose{\csromansymbolfootnote{*}{วิ 2. 236 ปิฏฺเฐปิ.}\textbf{ยานิ โข ปน ตานิ อารญฺญกานิ เสนาสนานี}ติ \textbf{อารญฺญกํ} นาม เสนาสนํ ปญฺจ\-ธนุ\-สติกํ ปจฺฉิมํ.}

\prose{\csromansymbolmark{*}\textbf{สาสงฺกํ} นาม อาราเม อารามูปจาเร โจรานํ นิวิฏฺโฐกาโส ทิสฺสติ, ภุตฺโตกาโส ทิสฺสติ, ฐิโตกาโส ทิสฺสติ, นิสินฺโนกาโส ทิสฺสติ, นิปนฺโนกาโส ทิสฺสติ.}

\prose{\csromanfolio{378}\csromansymbolfootnote{*}{วิ 2. 236 ปิฏฺเฐปิ.}\textbf{สปฺปฏิภยํ} นาม อาราเม อารามูปจาเร โจเรหิ มนุสฺสา หตา ทิสฺสนฺติ, วิลุตฺตา ทิสฺสนฺติ, อาโกฏิตา ทิสฺสนฺติ.}

\prose{\csromansymbolmark{*}\textbf{ตถารูเปสุ ภิกฺขุ เสนาสเนสุ วิหรนฺโต}ติ เอวรูเปสุ ภิกฺขุ เสนาสเนสุ วิหรนฺโต.}

\prose{\textbf{อากงฺขมาโน}ติ อิจฺฉามาโน.}

\prose{\textbf{ติณฺณํ จีวรานํ อญฺญตรํ จีวรนฺ}ติ สํฆาฏิํ วา อุตฺตราสงฺคํ วา อนฺตรวาสกํ วา.}

\prose{\textbf{อนฺตรฆเร นิกฺขิเปยฺยา}ติ สมนฺตา โคจรคาเม นิกฺขิเปยฺย.}

\prose{\textbf{สิยา จ ตสฺส ภิกฺขุโน โกจิเทว ปจฺจโย เตน จีวเรน วิปฺปวาสายา}ติ สิยา ปจฺจโย, สิยา กรณียํ.}

\prose{\textbf{ฉารตฺต}\-\textbf{ปรมํ เตน ภิกฺขุนา เตน จีวเรน วิปฺปวสิตพฺพนฺ}ติ ฉารตฺต\-ปรมตา วิปฺปวสิตพฺพํ.}

\prose{\textbf{อญฺญตฺร ภิกฺขุ}\-\textbf{สมฺมุติยา}ติ ฐเปตฺวา ภิกฺขุ\-สมฺมุติํ.}

\prose{\textbf{ตโต เจ อุตฺตริ วิปฺปวเสยฺยา}ติ สตฺตเม อรุณุคฺคมเน นิสฺสคฺคิยํ โหติ. นิสฺสชฺชิตพฺพํ สํฆสฺส วา คณสฺส วา ปุคฺคลสฺส วา. เอวญฺจ ปน ภิกฺขเว นิสฺสชฺชิตพฺพํ ฯเปฯ “อิทํ เม ภนฺเต จีวรํ อติเรกฉารตฺตํ วิปฺปวุฏฺฐํ อญฺญตฺร ภิกฺขุ\-สมฺมุติยา นิสฺสคฺคิยํ, อิมาหํ สํฆสฺส นิสฺสชฺชามี”ติ ฯเปฯ ทเทยฺยาติ ฯเปฯ ทเทยฺยุนฺติ ฯเปฯ อายสฺมโต ทมฺมีติ.}

\proseitem{655}{อติเรก\-ฉารตฺเต อติเรกสญฺญี วิปฺปวสติ อญฺญตฺร ภิกฺขุ\-สมฺมุติยา, นิสฺสคฺคิยํ ปาจิตฺติยํ. อติเรก\-ฉารตฺเต เวมติโก วิปฺปวสติ อญฺญตฺร ภิกฺขุ\-สมฺมุติยา, นิสฺสคฺคิยํ ปาจิตฺติยํ. อติเรก\-ฉารตฺเต อูนกสญฺญี วิปฺปวสติ อญฺญตฺร ภิกฺขุ\-สมฺมุติยา, นิสฺสคฺคิยํ ปาจิตฺติยํ. อปฺปจฺจุทฺธเฏ ปจฺจุทฺธฏ\-สญฺญี ฯเปฯ. อวิสฺสชฺชิเต วิสฺสชฺชิต\-สญฺญี. อนฏฺเฐ นฏฺฐสญฺญี. อวินฏฺเฐ วินฏฺฐสญฺญี. อทฑฺเฒ ทฑฺฒสญฺญี. อวิลุตฺเต วิปุตฺตสญฺญี วิปฺปวสติ อญฺญตฺร ภิกฺขุ\-สมฺมุติยา, นิสฺสคฺคิยํ ปาจิตฺติยํ.}

\prose{นิสฺสคฺคิยํ จีวรํ อนิสฺสชฺชิตฺวา ปริภุญฺชติ, อาปตฺติ ทุกฺกฏสฺส. อูนกฉารตฺเต อติเรกสญฺญี, อาปตฺติ ทุกฺกฏสฺส. อูนกฉารตฺเต เวมติโก, อาปตฺติ ทุกฺกฏสฺส. อูนกฉารตฺเต อูนกสญฺญี, อนาปตฺติ.}

\csromanmatikaanchor{mtk.239}
\csromanmatikaanchor{mtk.240}
\csromantocmarknopage{subsection}{10. ปริณตสิกฺขาปท}{mtk.239}
\bookmark[dest={mtk.239},level=2]{10. ปริณตสิกฺขาปท}
\csromantocmarkref{subsubsection}{ฉพฺพคฺคิยภิกฺขุวตฺถุ}{mtk.240}
\bookmark[dest={mtk.240},level=3]{ฉพฺพคฺคิยภิกฺขุวตฺถุ}
\proseitemwithcloserruled{656}{\csromanfolio{379}อนาปตฺติ ฉารตฺตํ วิปฺปวสติ, อูนกฉารตฺตํ วิปฺปวสติ, ฉารตฺตํ วิปฺปวสิตฺวา ปุน คามสีมํ โอกฺกมิตฺวา วสิตฺวา ปกฺกมติ, อนฺโต ฉารตฺตํ ปจฺจุทฺธรติ, วิสฺสชฺเชติ, นสฺสติ, วินสฺสติ, ฑยฺหติ, อจฺฉินฺทิตฺวา คณฺหนฺติ, วิสฺสาสํ คณฺหนฺติ, ภิกฺขุ\-สมฺมุติยา, อุมฺมตฺตกสฺส อาทิกมฺมิกสฺสาติ.}{สาสงฺก\-สิกฺขาปทํ นิฏฺฐิตํ นวมํ.}
\csromanmarkh{3. ปตฺตวคฺค}\csromanmarkhii{10. ปริณต\-สิกฺขาปท}\csromanheadernomark{2}{3. \textbf{ปตฺตวคฺค 1}0. ปริณต\-สิกฺขาปท}
\csromanmatikaanchor{mtk.241}
\csromanmatikaanchor{mtk.242}
\csromantocmarkref{subsubsection}{ปญฺญตฺติ}{mtk.241}
\bookmark[dest={mtk.241},level=3]{ปญฺญตฺติ}
\csromantocmarkref{subsubsection}{สิกฺขาปทวิภงฺค, ปทภาชนีย}{mtk.242}
\bookmark[dest={mtk.242},level=3]{สิกฺขาปทวิภงฺค, ปทภาชนีย}
\proseitem{657}{เตน สมเยน พุทฺโธ ภควา สาวตฺถิยํ วิหรติ เชตวเน อนาถ\-ปิณฺฑิกสฺส อาราเม. เตน โข ปน สมเยน สาวตฺถิยํ อญฺญตรสฺส ปูคสฺส สํฆสฺส สจีวร\-ภตฺตํ ปฏิยตฺตํ โหติ “โภเชตฺวา จีวเรน อจฺฉาเทสฺสามา”ติ. อถ โข ฉพฺพคฺคิยา ภิกฺขู เยน โส ปูโค เตนุ\-ปสงฺกมิํสุ, อุปสงฺกมิตฺวา ตํ ปูคํ เอตทโวจุํ “เทถาวุโส อมฺหากํ อิมานิ จีวรานี”ติ. น มยํ ภนฺเต ทสฺสาม, อมฺหากํ สํฆสฺส อนุวสฺสํ สจีวร\-ภิกฺขา ปญฺญตฺตาติ. พหู อาวุโส สํฆสฺส ทายกา พหู สํฆสฺส ภตฺตา\csromansharedfootnote{8fc5c5cfa51835f2}{ภทฺทา (ก)}, มยํ ตุเมฺห นิสฺสาย ตุเมฺห สมฺปสฺสนฺตา อิธ วิหราม, ตุเมฺห เจ อมฺหากํ น ทสฺสถ, อถ โก จรหิ อมฺหากํ ทสฺสติ, เทถาวุโส อมฺหากํ อิมานิ จีวรานีติ. อถ โข โส ปูโค ฉพฺพคฺคิเยหิ ภิกฺขูหิ นิปฺปีฬิยมาโน ยถา\-ปฏิยตฺตํ จีวรํ ฉพฺพคฺคิยานํ ภิกฺขูนํ ทตฺวา สํฆํ ภตฺเตน ปริวิสิ. เย เต ภิกฺขู ชานนฺติ สํฆสฺส สจีวร\-ภตฺตํ ปฏิยตฺตํ, น จ ชานนฺติ “ฉพฺพคฺคิยานํ ภิกฺขูนํ ทินฺนนฺ”ติ, เต เอวมาหํสุ “โอโณเชถาวุโส สํฆสฺส จีวรนฺ”ติ. นตฺถิ ภนฺเต, ยถา\-ปฏิยตฺตํ จีวรํ อยฺยา ฉพฺพคฺคิยา อตฺตโน ปริณาเมสุนฺติ. เย เต ภิกฺขู อปฺปิจฺฉา ฯเปฯ เต อุชฺฌายนฺติ ขิยฺยนฺติ วิปาเจนฺติ “กถํ หิ นาม ฉพฺพคฺคิยา ภิกฺขู ชานํ สํฆิกํ ลาภํ ปริณตํ อตฺตโน ปริณาเมสฺสนฺตี”ติ. อถ โข เต ภิกฺขู ฉพฺพคฺคิเย ภิกฺขู อเนก\-ปริยาเยน วิครหิตฺวา ภควโต เอตมตฺถํ อาโรเจสุํ ฯเปฯ “สจฺจํ กิร ตุเมฺห ภิกฺขเว ชานํ สํฆิกํ ลาภํ ปริณตํ อตฺตโน ปริณาเมถา”ติ. สจฺจํ ภควาติ. วิครหิ พุทฺโธ ภควา ฯเปฯ. กถํ หิ นาม ตุเมฺห โมฆปุริสา ชานํ สํฆิกํ ลาภํ ปริณตํ \csromanfolio{380}อตฺตโน ปริณาเมสฺสถ, เนตํ โมฆปุริสา อปฺปสนฺนานํ วา ปสาทาย ฯเปฯ. เอวญฺจ ปน ภิกฺขเว อิมํ สิกฺขาปทํ อุทฺทิเสยฺยาถ\csromandash{}}

\proseitem{658}{\textbf{“โย ปน ภิกฺขุ ชานํ สํฆิกํ ลาภํ ปริณตํ อตฺตโน ปริณาเมยฺย, นิสฺสคฺคิยํ ปาจิตฺติยนฺ”}ติ.}

\proseitem{659}{\textbf{โย ปนา}ติ โย ยาทิโส ฯเปฯ. \textbf{ภิกฺขู}ติ ฯเปฯ อยํ อิมสฺมิํ อตฺเถ อธิปฺเปโต “ภิกฺขู”ติ.}

\prose{\csromansymbolfootnote{*}{วิ 2. 204 ปิฏฺเฐปิ.}\textbf{ชานาติ} นาม สามํ วา ชานาติ, อญฺเญ วา ตสฺส อาโรเจนฺติ, โส วา อาโรเจติ.}

\prose{\csromansymbolmark{*}\textbf{สํฆิกํ} นาม สํฆสฺส ทินฺนํ โหติ ปริจฺจตฺตํ.}

\prose{\csromansymbolmark{*}\textbf{ลาโภ} นาม จีวรปิณฺฑปาตเสนาสน\-คิลานปฺปจฺจย\-เภสชฺช- ปริกฺขารา, อนฺตมโส จุณฺณปิณฺโฑปิ ทนฺตกฏฺฐํปิ ทสิกสุตฺตํปิ.}

\prose{\csromansymbolmark{*}\textbf{ปริณตํ} นาม “ทสฺสาม กริสฺสามา”ติ วาจา ภินฺนา โหติ.}

\prose{อตฺตโน ปริณาเมติ, ปโยเค ทุกฺกฏํ. ปฏิลาเภน นิสฺสคฺคิยํ โหติ. นิสฺสชฺชิตพฺพํ สํฆสฺส วา คณสฺส วา ปุคฺคลสฺส วา. เอวญฺจ ปน ภิกฺขเว นิสฺสชฺชิตพฺพํ ฯเปฯ “อิทํ เม ภนฺเต ชานํ สํฆิกํ ลาภํ ปริณตํ อตฺตโน ปริณามิตํ นิสฺสคฺคิยํ, อิมาหํ สํฆสฺส นิสฺสชฺชามี”ติ ฯเปฯ ทเทยฺยาติ ฯเปฯ ทเทยฺยุนฺติ ฯเปฯ อายสฺมโต ทมฺมีติ.}

\proseitem{660}{ปริณเต ปริณตสญฺญี อตฺตโน ปริณาเมติ, นิสฺสคฺคิยํ ปาจิตฺติยํ.}

\prose{ปริณเต เวมติโก อตฺตโน ปริณาเมติ, อาปตฺติ ทุกฺกฏสฺส. ปริณเต อปริณตสญฺญี อตฺตโน ปริณาเมติ, อนาปตฺติ. สํฆสฺส ปริณตํ อญฺญสํฆสฺส วา เจติยสฺส วา ปริณาเมติ, อาปตฺติ ทุกฺกฏสฺส. เจติยสฺส ปริณตํ อญฺญเจติยสฺส วา สํฆสฺส วา ปุคฺคลสฺส วา ปริณาเมติ, อาปตฺติ ทุกฺกฏสฺส. ปุคฺคลสฺส ปริณตํ อญฺญปุคฺคลสฺส วา สํฆสฺส วา เจติยสฺส วา ปริณาเมติ, อาปตฺติ ทุกฺกฏสฺส. อปริณเต ปริณตสญฺญี, อาปตฺติ ทุกฺกฏสฺส. อปริณเต เวมติโก อาปตฺติ ทุกฺกฏสฺส. อปริณเต อปริณตสญฺญี, อนาปตฺติ.}

\csromanmatikaanchor{mtk.243}
\csromantocmarkref{subsection}{อุทฺทานคาถา}{mtk.243}
\bookmark[dest={mtk.243},level=2]{อุทฺทานคาถา}
\proseitem{661}{\csromanfolio{381}อนาปตฺติ “กตฺถ เทมา”ติ ปุจฺฉียมาโน “ยตฺถ ตุมฺหากํ เทยฺยธมฺโม ปริโภคํ วา ลเภยฺย, ปฏิสงฺขารํ วา ลเภยฺย, จิรฏฺฐิติโก วา อสฺส, ยตฺถ วา ปน ตุมฺหากํ จิตฺตํ ปสีทติ, ตตฺถ เทถา”ติ ภณติ, อุมฺมตฺตกสฺส อาทิกมฺมิกสฺสาติ.}

\vspace{\tipitakaparskip}
\csromancenter{ปริณต\-สิกฺขาปทํ นิฏฺฐิตํ ทสมํ.}
\nitthitammidruled{ปตฺตวคฺโค ตติโย.}
\csromancenter{ตสฺสุทฺทานํ}
\setlength{\csromangathapagewidth}{0pt}
\setlength{\csromangathawakwidth}{0pt}
\csromangathameasuregroup{เทฺว จ ปตฺตานิ เภสชฺชํ, \\ สามํ วายาปนจฺเจโก,}{\csromangathabat{เทฺว จ ปตฺตานิ เภสชฺชํ,}{วสฺสิกา ทานปญฺจมํ.} \\ \csromangathabat{สามํ วายาปนจฺเจโก,}{สาสงฺกํ สํฆิเกน จาติ.}}
\csromangathasetleft{เทฺว จ ปตฺตานิ เภสชฺชํ, \\ สามํ วายาปนจฺเจโก,}
\csromangathagroup{\csromangathastanza{\csromangathabat{เทฺว จ ปตฺตานิ เภสชฺชํ,}{วสฺสิกา ทานปญฺจมํ.} \\ \csromangathabat{สามํ วายาปนจฺเจโก,}{สาสงฺกํ สํฆิเกน จาติ.}}}

\proseitemwithclosermajor{662}{อุทฺทิฏฺฐา โข อายสฺมนฺโต ติํส นิสฺสคฺคิยา ปาจิตฺติยา ธมฺมา, ตตฺถายสฺมนฺเต ปุจฺฉามิ กจฺจิตฺถ ปริสุทฺธา, ทุติยมฺปิ ปุจฺฉามิ กจฺจิตฺถ ปริสุทฺธา, ตติยมฺปิ ปุจฺฉามิ กจฺจิตฺถ ปริสุทฺธา, ปริสุทฺเธ\-ตฺถา\-ยสฺมนฺโต, ตสฺมา ตุณฺหี, เอวเมตํ ธารยามีติ.}{นิสฺสคฺคิย\-กณฺฑํ นิฏฺฐิตํ.}
\nitthitammajor{\textbf{ปาราชิกปาฬิ นิฏฺฐิตา}}
